\documentclass[10pt,a4paper]{article}
\usepackage[margin=18mm]{geometry}
\usepackage[T1]{fontenc}
\usepackage[utf8]{inputenc}
\usepackage{lmodern,amsmath,amssymb,amsthm,mathtools,mathrsfs}
\usepackage{longtable,booktabs,array,calc}
\usepackage{graphicx,xcolor,pdflscape}
\usepackage[unicode,colorlinks=true,linkcolor=blue!45!black,urlcolor=blue!45!black,bookmarksdepth=3]{hyperref}
\hypersetup{pdftitle={The ordinary QRH half-plane: a compact working proof -- Cycle092 edition},pdfauthor={Hailey Collet},pdfsubject={Full public proof with dated formalization supplement}}
\setcounter{secnumdepth}{-1}
\setcounter{tocdepth}{2}
\sloppy
\setlength{\emergencystretch}{0.4\textwidth}
\setlength{\parskip}{3pt}
\setlength{\parindent}{0pt}
\providecommand{\tightlist}{\setlength{\itemsep}{0pt}\setlength{\parskip}{0pt}}
\providecommand{\passthrough}[1]{#1}
\renewcommand{\labelitemi}{\ensuremath{\bullet}}
\renewcommand{\labelitemii}{\ensuremath{\circ}}
\DeclareMathSizes{14.4}{10}{7}{5}
\DeclareMathSizes{12}{10}{7}{5}
\DeclareFontShape{OT1}{cmr}{m}{n}{<->ssub*lmr/m/n}{}
\newsavebox{\displaybox}
\newcommand{\fitmath}[1]{\sbox{\displaybox}{$\displaystyle #1$}\ifdim\wd\displaybox>0.90\linewidth\resizebox{0.90\linewidth}{!}{\usebox{\displaybox}}\else\usebox{\displaybox}\fi}
\newcommand{\fitinline}[1]{\sbox{\displaybox}{$#1$}\ifdim\wd\displaybox>0.90\linewidth\resizebox{0.90\linewidth}{!}{\usebox{\displaybox}}\else\usebox{\displaybox}\fi}
\allowdisplaybreaks
\begin{document}
\hypersetup{pageanchor=false}
\begin{titlepage}
\centering
\vspace*{22mm}
{\Large\bfseries The ordinary QRH half-plane:\\[3mm] a compact working proof\par}
\vspace{12mm}
{\normalsize Hailey Collet\par}
\vspace{3mm}
\href{mailto:HaileyCollet@gmail.com}{HaileyCollet@gmail.com}\par
\vspace{14mm}
{\normalsize Cycle092 public proof edition\par}
\vspace{4mm}
Written submission: 19 November 2026\\
Formalization supplement: 20 November 2026 (UTC)\par
\vspace{12mm}
\(c=63333603/500000000=.126667206\)\\[2mm]
\(\theta=436666397/500000000=.873332794\)\par
\vfill
\begin{minipage}{0.88\textwidth}\small\raggedright
This edition typesets the complete public Markdown proof at commit
\texttt{1affb4664d550a568e32be35238445e3831ed237}.
The proof text, its mathematics and its inherited attribution are retained.
The dated supplement records the later local formalization status separately.
The unchanged Markdown, finite certificates and replay program accompany this edition in \texttt{public-source/}.

Copyright 2026 Hailey Collet. New material is licensed under Apache-2.0.
Inherited work retains its authorship, licences and notices; see the manuscript and the accompanying \texttt{NOTICE.md} and \texttt{LICENSE}.
Hailey Collet is the sole human author of this contribution. Research, drafting and code development were substantially AI-assisted; this does not constitute independent verification.
\end{minipage}
\end{titlepage}
\hypersetup{pageanchor=true}
\section*{Formalization supplement: 20 November 2026 (UTC)}
\addcontentsline{toc}{section}{Formalization supplement: 20 November 2026 (UTC)}
This supplement records the status after the written Cycle092 submission. The complete proof that follows is the public text submitted on 19 November 2026; its historical statements about adopted inputs and verification retain their original dates and scopes.

The final local formalization establishes the advertised strict half-plane at the exact value
\[
 \theta=\frac{436666397}{500000000},\qquad
 c=\frac{63333603}{500000000}.
\]
Its declaration \texttt{QRH.Cycle092.dirichlet\_nonvanishing} has the ordinary all-Dirichlet signature: every positive presented modulus, every complex Dirichlet character (including imprimitive presentations), and every complex argument with \(\Re s>\theta\), excluding \(\chi=1,s=1\). The companion \texttt{QRH.Cycle092.hecke\_and\_dirichlet\_record} also gives the Hecke zero-supremum and ordinary target conclusions. No boundary-line assertion or uniformity in growing moduli is added.

The implementation retains the original amplitude-phase cells and their full class means when choosing the finite return ticket. One same-zero bank is used through the response descriptors and original error masks. The resulting packet decomposition uses the actual original coefficients and prime marks; the large-coefficient contribution is paid once on the full phase cell for each original mask. This realizes the physical alignment required in Appendices P--Q rather than replacing an original factor by an unrelated moment.

For the physical denominator the implementation uses the direct reciprocal route of K.1: the buffered estimate for the actual \(1/L(s)\) is paid separately from the numerator packet. The inverse detector at the same zero supplies its own count estimate; its length is not identified with a physical reciprocal annulus. The positive reserve \(\rho=1/50000000000\) in P.2 is retained, together with positive inner decrements and meshes chosen in dependency order (in particular \(0<\delta\le\rho/10000\)). The fixed base and its positive collar are retained in the marked source and restoration steps; they are not removed by setting a reserve to zero.

The source-bound local receipt is dated 20 November 2026, 04:06:57 UTC. Lean 4.34.1 compiled the unconditional final root from source with exit status zero against its development dependency objects. A separate process checked the exact all-Dirichlet signature and audited both declarations; their reported axioms are \texttt{Classical.choice}, \texttt{Quot.sound} and \texttt{propext}. The final-root source SHA-256 is
\begin{quote}\small\ttfamily\nolinkurl{005e7ae5a54853586d9d591ce45814ec4ca023191ca5fafc935eedd0509711cd}\end{quote}
This is an unconditional local development check. A complete clean submission build, canonical independent kernel replay and protected tracker acceptance are separate checks and are not asserted by this edition. The written proof's originally adopted H6/SN9 input remains identified in its own historical text; the local final-root statement has no extra analytic premise.

The public source is pinned at
\href{https://github.com/HaileyStorm/quasi-rh-hillclimb/blob/1affb4664d550a568e32be35238445e3831ed237/submissions/hailey-collet-cycle092/packet/CLOSED_PROOF.md}{the complete Cycle092 Markdown proof}; its SHA-256 is
\begin{quote}\small\ttfamily\nolinkurl{c1876d05ad7376e6b9ecf91005b6be496d061b304cfd598adb54ac2603f87247}\end{quote}
The accompanying unchanged \texttt{public-source/} directory contains that Markdown, its source manifest, finite certificates and the original replay entry point. The public written submission is \href{https://github.com/rohanarun/quasi-rh-hillclimb/pull/45}{PR45}. Its submission and this dated local receipt do not constitute hosted acceptance.

\clearpage
\tableofcontents
\clearpage
\hypertarget{the-ordinary-qrh-half-plane-a-compact-working-proof}{%
\section{The ordinary QRH half-plane: a compact working proof}\label{the-ordinary-qrh-half-plane-a-compact-working-proof}}

\hypertarget{the-ordinary-theorem}{%
\subsection{1. The ordinary theorem}\label{the-ordinary-theorem}}

\textbf{Theorem Q (Local Codex Cycle 090).} Put \[\fitmath{
  c_{90}=\frac{63333603}{500000000}=.126667206,\qquad
  \theta_{90}=1-c_{90}=\frac{436666397}{500000000}=.873332794.
}\] At every fixed positive presented modulus \(\fitinline{q}\), for every ordinary Dirichlet character \(\fitinline{\chi\bmod q}\), \(\fitinline{L(s,\chi)\ne0}\) throughout \(\fitinline{\Re s>\theta_{90}}\), excluding \(\fitinline{s=1}\) for a principal character. This includes imprimitive presentations and all heights. No statement is asserted on the boundary or uniformly in growing moduli.

The gain over the previous local endpoint \(\fitinline{c=.126667205}\) is exactly \(\fitinline{\frac{1}{1000000000}=.000000001}\). \protect\hyperlink{appendix-q}{Appendix Q} gives the proof. The same symmetric ordinary source now clears the full physical return after tightening the witness allowances and refining the finite cutoff bank. The separate signed research toward .12668 and .12725 is not an input. The retained natural selected-cell input SN9 and the other explicitly stated foundations keep their original scopes.

\hypertarget{the-retained-cycle-080-theorem}{%
\subsubsection{The retained Cycle 080 theorem}\label{the-retained-cycle-080-theorem}}

\textbf{Theorem H (Local Codex Cycle 080).} Put \[\fitmath{
 \theta_{80}=\frac{17466917}{20000000}=.87334585,\qquad
 c_{80}=1-\theta_{80}=\frac{2533083}{20000000}=.12665415.
}\] For every fixed positive integer \(\fitinline{q}\), every ordinary Dirichlet character presented modulo \(\fitinline{q}\), and every complex \(\fitinline{s}\) with \(\fitinline{\Re s>\theta_{80}}\), \(\fitinline{L(s,\chi)\ne0}\), excluding \(\fitinline{s=1}\) for a principal character. This includes imprimitive presentations and all heights. There is no assertion on the boundary or of uniform constants in growing moduli.

The exact increase over the attained PR143 theorem (T2) is \(\fitinline{\theta_{\rm rec}-\theta_{80}=.00000000610671857680118479606507\ldots}\). The proof first obtains (T2). \protect\hyperlink{appendix-h}{Appendix H} then uses that previous bound on the \emph{actual} zero supremum, proves the improvement in the two signed row estimates, and returns through the original physical denominator and its independent numerator. Its gap to \(\fitinline{c=.127}\) was \(\fitinline{6917/20000000=.00034585}\).

\hypertarget{previous-attained-theorem-used-in-the-proof}{%
\subsubsection{1.1. Previous attained theorem used in the proof}\label{previous-attained-theorem-used-in-the-proof}}

Throughout the retained argument and Appendix H, \(\fitinline{\theta_{\rm rec}}\) and \(\fitinline{c_{\rm rec}}\) denote the PR143 contact in (T1). The boundary and coefficient of Theorem H are \(\fitinline{\theta_{80}}\) and \(\fitinline{c_{80}}\).

Let \(\fitinline{\theta_{\rm rec}}\) be the unique root in \(\fitinline{(13/15,7/8)}\) of \[\fitmath{
3888\,19^5(15t-13)^4(1-t)^5-7(7-6t)^4(9t-7)^5=0,
\qquad c_{\rm rec}=1-\theta_{\rm rec}.                    
}\tag{T1}\] \textbf{Theorem T.} For every positive integer \(\fitinline{q}\), every ordinary complex Dirichlet character \(\fitinline{\chi}\) presented modulo \(\fitinline{q}\), and every \(\fitinline{s\in\mathbf C}\), \[\fitmath{
\Re s>\theta_{\rm rec},\qquad
\neg(\chi\text{ principal and }s=1)
\quad\Longrightarrow\quad L(s,\chi)\ne0.                 
}\tag{T2}\] Here \[\fitmath{
c_{\rm rec}=.12665414\allowbreak{}38932814\allowbreak{}23198815\allowbreak{}20393492\allowbreak{}76432295\allowbreak{}58060156\allowbreak{}0188104\ldots.
}\] Each analytic estimate fixes the presented modulus. The conclusion is at all heights. It asserts neither nonvanishing on the boundary line nor constants uniform in a growing modulus.

The main argument and its mathematical appendices A--Q are contained in this standalone manuscript. Appendices A--C retain the original foundations; \protect\hyperlink{appendix-d}{Appendix D} proves marked-anchor induction; \protect\hyperlink{appendix-e}{Appendix E} proves the continuum supplier, averaged floor and complete rational return; \protect\hyperlink{appendix-f}{Appendix F} supplies its finite strict source construction. \protect\hyperlink{appendix-g}{Appendix G} parameterizes the actual source and physical return by arbitrary positive costs and proves (T2) by a final order argument. Every strict target uses its own finite source, mesh, arithmetic data, height and constants. No limit of analytic estimates or their implicit constants is taken. The finite mathematical data named by these appendices remain in the packet. Original source copies and SOURCE\_USE records are provenance; a declaration or finite arithmetic check supplies no missing analytic argument. The adopted natural selected-cell foundation in \S{}3 keeps its stated scope.

For the retained initialization and alternative suppliers write \[\fitmath{
c_*=(75+\sqrt{1081})/852,\qquad
\theta_*=(777-\sqrt{1081})/852.                            
}\tag{T3}\] Sections 2--10 and Appendices A--C retain that earlier theorem and its more general source/count/return interfaces. Section 11 and Appendix E retain the proved PR141 rational endpoint \(\fitinline{R_{141}=N/D}\), with the exact integers displayed there. Appendix G improves the ordinary coefficient by exactly \(\fitinline{R_{141}-\theta_{\rm rec}>0}\), approximately \(\fitinline{3.2768490251484275310018257340\times10^{-254}}\). The ranked rational witness is unchanged; it does not describe the full exact theorem.

The compensated probe has two exact representations. Completed theta energy pays its signed physical form; Poisson gives a nonzero reciprocal signal and the complete nonprincipal numerator. Same-zero witnesses count its rows, marked moments bound their actual amplitudes, and the averaged floor pays its complete sector. The full return improves the actual zero supremum, without assuming a zero attains that supremum. Appendix G then bounds the supremum by the exact contact.

\begin{longtable}[]{@{}
  >{\raggedright\arraybackslash}p{(\columnwidth - 2\tabcolsep) * \real{0.5000}}
  >{\raggedright\arraybackslash}p{(\columnwidth - 2\tabcolsep) * \real{0.5000}}@{}}
\toprule\noalign{}
\begin{minipage}[b]{\linewidth}\raggedright
Part
\end{minipage} & \begin{minipage}[b]{\linewidth}\raggedright
Role in the current theorem
\end{minipage} \\
\midrule\noalign{}
\endhead
\bottomrule\noalign{}
\endlastfoot
\S{}\S{}2--3, Appendices A--B & Arithmetic, selected norm, completed energy, actual probe, witnesses and original initialization. \\
\S{}\S{}4--6, Appendix C & Literal source restoration and the strict family \(\fitinline{G(B,L,c)}\). \\
\S{}\S{}7--10 & Count/physical interfaces, complementary sectors, feedback and retained algebraic baseline. \\
Appendix D & Coupled marked-anchor induction beyond the older source cone. \\
Appendices E--F, \S{}11 & Continuum sources, averaged floor and complete positive-cost PR141 rational theorem. \\
Appendix G & Variable positive-cost analytic return and previous ordinary contact (T2). \\
Appendix H & Two signed side bounds, uncentered root restoration, original physical integral and strict ordinary return at the new rational target. \\
Appendix I & Strict ordinary return at c=.126655; scope and changed estimates are given there. \\
Appendix J & Strict ordinary return at c=.126657; scope and changed estimates are given there. \\
Appendix K & Strict ordinary return at c=.126658; scope and changed estimates are given there. \\
Appendix L & Strict ordinary return at c=.12666414; scope and changed estimates are given there. \\
Appendix M & Strict ordinary return at c=.126666; scope and changed estimates are given there. \\
Appendix N & Strict ordinary return at c=.12666645; scope and changed estimates are given there. \\
Appendix O & Strict ordinary return at c=.126667; scope and changed estimates are given there. \\
Appendix P & Strict ordinary return at c=.126667205; scope and changed estimates are given there. \\
Appendix Q & Strict ordinary return at c=.126667206; scope and changed estimates are given there. \\
\end{longtable}

\textbf{Provenance.} The earlier algebraic endpoint came from supplied GPT6-Pro C074; admitted Tracker PR129 formalizes that same contact. Sil's admitted PR137 supplies marked-anchor induction. Joseph Bebel's admitted PR141 supplies the continuum extension and rational theorem; his admitted PR143 supplies the complete variable-cost return and exact-contact theorem, with AI assistance. PR142 adds the isolated-polynomial-root verification grammar; it supplies no new analytic theorem. OpenAI, Hailey Collet, ArviFoxy, Joe Richey, Jacob Sussman and all other inherited contributors retain their attribution, licenses and notices. Local Codex Cycle077 reconstructs and integrates this public advance; it claims no independent new endpoint or local Lean replay. The near-record continuation in Cycles 078--080 uses SovereignSteak's draft PR43 at commit c838a750\allowbreak{}079f7f86\allowbreak{}a573df9c\allowbreak{}9f729b8b\allowbreak{}c92d125e for the two-ray and physical reductions. The source retains its author, assistance credits, license and notice. Appendix H gives the arguments it uses and the further steps from this local investigation. Hailey Collet directed the work. The new endpoint is a locally reviewed written result on the stated foundations; no new Lean run, human peer review or public acceptance is claimed. The separate angular \texttt{.\allowbreak{}1\allowbreak{}2\allowbreak{}5\allowbreak{}2\allowbreak{}2\allowbreak{}7\allowbreak{}0\allowbreak{}1\allowbreak{}3} theorem retains its fixed-family scope. The complete \texttt{.\allowbreak{}1\allowbreak{}2\allowbreak{}7} return remains conditional in the research guide.

\hypertarget{shared-arithmetic-and-analytic-contract}{%
\subsection{2. Shared arithmetic and analytic contract}\label{shared-arithmetic-and-analytic-contract}}

Use \(\fitinline{K=\mathbf Q(\sqrt{-3})}\), \(\fitinline{\mathcal O=\mathbf Z[\omega]}\), \(\fitinline{\omega=e^{2\pi i/3}}\), \(\fitinline{\lambda=1+2\omega}\), and \(\fitinline{q_a=Na}\). An ideal sum uses the fixed primary generator and counts each ideal once; an element-row sum retains all six units. The fixed bad set contains the target conductor primes, primes above six, and the subsequently fixed small-prime cutoff. It is chosen before the asymptotic scale. The Fourier character, measure, Gauss phases and classical inputs are stated in Appendix A.1--A.2 and used without a change of convention.

Every integer power of a residue character, including powers zero and six and negative powers, is zero at a nonunit. Only the absent modulus-one factor is identically one. Removing a phase therefore never removes its zero set. A presented character keeps primitive-conductor zeros, redundant presentation zeros, fixed bad-prime zeros and all displayed moving zeros. A common puncture may depend on frozen outer labels but is fixed across the live row sum. Moving arithmetic support is charged as width; it is not a smooth-profile seminorm.

Let \(\fitinline{\beta}\) be the supremum of \(\fitinline{1/2}\) and the real parts of the positive-real-part zeros in the finite-order, trivial-infinity Eisenstein Hecke family, omitting the principal pole. Primitive and presented characters have the same positive-half-plane zeros, because the extra Euler factors have no zero there. Appendix B proves \(\fitinline{\beta\le7/8}\) from the original inputs. Write \(\fitinline{\kappa=2\beta-1}\). Every use of the source family is in the standing case \(\fitinline{\beta\ge51/100}\); if that condition fails, Theorem T already follows by the ordinary norm pullback.

\hypertarget{the-literal-global-source}{%
\subsubsection{2.1 The literal global source}\label{the-literal-global-source}}

Fix a nonzero ray modulus \(\fitinline{\mathfrak M}\), a subgroup containing global units, fixed anchor characters \(\fitinline{\eta,\eta_0}\), and \(\fitinline{Q\subseteq\mathfrak M}\). Put \(\fitinline{Q_0=Q\cap\mathfrak q_{\eta_0}}\), and require \(\fitinline{Q_0\ne0,\mathcal O}\) and \(\fitinline{Q_0\subseteq(72)}\). For each slot, \(\fitinline{\vartheta_i}\) is a character of this fixed ray quotient and \(\fitinline{\nu_i=\vartheta_i\overline{\eta_0}}\), with both zero sets retained. Its pool \(\fitinline{\mathcal P_i^{\rm can}}\) is the original identity ray class of good primes at scale \(\fitinline{P_i=Z^{w_i}}\). The common smooth \(\fitinline{W}\) is supported in \(\fitinline{[1,2]}\). The normalized slot is \[\fitmath{
Q_i(z)=P_i^{-1/2}\sum_{\mathfrak p\in\mathcal P_i^{\rm can}}
 \nu_i(\mathfrak p)c_z(\mathfrak p)
 W(N\mathfrak p/P_i)
 (N\mathfrak p/P_i)^{-33/50+iv_i},
}\tag{S1}\] where \(\fitinline{c_z(I)=\eta(I)\chi_I(m_{\rm bad}^{6}z)}\), with \(\fitinline{m_{\rm bad}}\) a fixed nonzero element containing the prescribed bad-prime support. The sixth power retains exactly its zero mask and does not alter a nonzero sextic value. A common whole-product norm phase \(\fitinline{(NI)^{it_0}}\) may multiply this coefficient in every occurrence; its real height \(\fitinline{t_0}\) is a separate common parameter and is included in the bound \(\fitinline{V}\) and the polynomial height control. The two ordinary factors have independent norm phases: \[\fitmath{
A_j(z;X_j,p_j,t_j)=X_j^{-1/2}
 \sum_{I\ne0}c_z(I)(NI)^{it_j}p_j(NI/X_j),\qquad j=1,2.
}\tag{S2}\] The profiles have the original fixed annular supports inside \(\fitinline{[1/4,9/4]}\). They remain independent. Both factors have the same whole row-feature orientation at this source inlet. A physical reflection must restore that orientation before applying the source.

A row is \emph{fixed inducing} when the primitive character inducing its actual natural presentation has conductor dividing \(\fitinline{Q_0}\). Define \[\fitmath{
\mathrm{initialKeep}_{\eta,Q_0}(z)
 =\{z\ne0:\text{the natural presentation }c_z
       \text{ is not fixed inducing}\}.
}\tag{S3}\] Additional common Euler deletions preserve this primitive inducing predicate. The source energy is the literal positive sum \[\fitmath{
\mathcal E_Z
=\sum_{z:\mathrm{initialKeep}_{\eta,Q_0}(z)}
 |A_1(z)A_2(z)|^2\prod_{i\in T}|Q_i(z)|^2\Re\Phi(Nz/Z),
}\tag{S4}\] including all roots, units and natural zeros. Here \(\fitinline{\Phi}\) is Schwartz, \(\fitinline{\Re\Phi\ge0}\), and \(\fitinline{\operatorname{supp}\Phi\subset(-\infty,b_\Phi]}\) for a fixed \(\fitinline{b_\Phi\ge2}\). Finite row subsets are allowed by positivity.

Write \(\fitinline{\ell_Z(X)=\log_Z\max(1,X)}\). The assertion \(\fitinline{G(B,L,c)}\) means that for every positive loss \(\fitinline{\epsilon}\), \[\fitmath{
\mathcal E_Z\le C\,\mathcal D(\Phi)
 \mathcal S_{\mathcal F}(p_1,p_2)^2
 (1+|t_1|+|t_2|+V)^J Z^{1+\epsilon}
}\tag{S5}\] throughout \[\fitmath{
\ell_Z(X_1)+\ell_Z(X_2)+L\sum_{i\in T}w_i\le B,
\qquad \sum_{i\in T}w_i\le c,
\qquad0\le w_i\le\mathrm{mesh},\quad |v_i|\le V,\quad |t_0|\le V.
}\tag{S6}\] The profile control is the product of the two finite Schwartz-seminorm controls; its square gives the quartic homogeneity of (S4). One may take \(\fitinline{\mathcal D(\Phi)=16(1+\pi)^2\max_{0\le a\le2}\sup_x|x|^a|\Phi(x)|}\). The allowed real powers and all further ordinary shortenings retain the same bounded-range contract. The source is not an arbitrary coefficient or arbitrary prime-pool theorem.

The order of choices is part of (S5). Requested loss and bounded geometric ranges determine the mesh \textbf{before} the finite slot type. After the finite type and common marked profile are fixed, take the finite union of profile orders and the maximum height degree. Fixed compatible characters, ray data and support caps determine constants and onset. Those precede the live slot subset, widths, heights, ordinary profiles and lengths, radial test and retained row family. All statements below respect this order; neither a live growing number of slots nor boundary-uniform source constants are asserted.

For a packed sublattice row \(\fitinline{Vz}\) at radial scale \(\fitinline{K}\), positive enlargement uses \(\fitinline{U=N(V)K=Z^M}\). The actual keep and all masks are packed first. There is no omitted lattice-density factor. The canonical energy at \(\fitinline{U}\) bounds the packed energy; later complete restoration retains its original outside weights.

\hypertarget{common-analytic-operations}{%
\subsubsection{2.2 Common analytic operations}\label{common-analytic-operations}}

Compact joint logarithmic profiles admit Fourier expansions with a finite weighted \(\fitinline{L^1}\) norm. Minkowski, Cauchy and the common-profile moment therefore permit independent unit-modulus norm powers and rowwise smooth parameters at a finite seminorm cost. The proof and required derivatives are in Appendix A.1 and the moment proofs. A moving arithmetic mask may not be hidden in this control. Every tail is estimated on the original expression with a fixed polynomial envelope before selecting the decay order.

The actual row decomposition is \(\fitinline{z=u e^6}\), with \(\fitinline{u}\) sixth-power-free. No coprimality between \(\fitinline{u,e}\) is required. Its character is the original core character with the exact additional mask \(\fitinline{\mathbf1_{(I,e)=1}}\). All root sums below include this mask. Bounded nonempty ordinary scales retain the unit; if an annulus is empty it contributes zero. In particular a nonempty profile supported below \(\fitinline{9X/4}\) has \(\fitinline{X\ge4/9}\), so subunit central normalization is bounded.

The natural unmarked paired moment is available for all bounded ordinary lengths, with independent profiles/heights and polynomial-size common puncture fixed across the row sum. It excludes principal inducing rows, and includes fixed nonprincipal characters. Its full proof is Appendix A.5; its marked extension and the inducing-predicate treatment are Appendix A.7. These scopes are not interchangeable with an arbitrary selected operator norm.

\hypertarget{the-accepted-natural-selected-cell-input}{%
\subsection{3. The accepted natural selected-cell input}\label{the-accepted-natural-selected-cell-input}}

\hypertarget{the-foundational-boundary}{%
\subsubsection{The foundational boundary}\label{the-foundational-boundary}}

The following theorem is the PR95 natural selected-cell norm, at the exact evidence scope already adopted by the supplied proof and used in its half-price source construction. It is an explicit mathematical input to this manuscript, alongside the original OpenAI analytic inputs listed in the foundation ledger. Its statement here is not a claim that this consolidation independently proves the upstream PR95 analytic theorem, and a declaration name or successful code check is not used as a substitute for its statement. All applications, coefficient adapters, root sums, source constructions and endpoint returns made from this input are proved locally below.

The original source is Joseph Bebel's supplied PR95 head \texttt{a\allowbreak{}0\allowbreak{}6\allowbreak{}1\allowbreak{}f\allowbreak{}9\allowbreak{}9\allowbreak{}c\allowbreak{}a\allowbreak{}b\allowbreak{}6\allowbreak{}d\allowbreak{}d\allowbreak{}d\allowbreak{}6\allowbreak{}e\allowbreak{}d\allowbreak{}f\allowbreak{}b\allowbreak{}c\allowbreak{}7\allowbreak{}9\allowbreak{}f\allowbreak{}b\allowbreak{}a\allowbreak{}a\allowbreak{}8\allowbreak{}d\allowbreak{}a\allowbreak{}7\allowbreak{}4\allowbreak{}4\allowbreak{}7\allowbreak{}a\allowbreak{}9\allowbreak{}1\allowbreak{}c\allowbreak{}9\allowbreak{}c\allowbreak{}c}. The precise interface is \texttt{J\allowbreak{}o\allowbreak{}i\allowbreak{}n\allowbreak{}t\allowbreak{}D\allowbreak{}e\allowbreak{}t\allowbreak{}e\allowbreak{}c\allowbreak{}t\allowbreak{}o\allowbreak{}r\allowbreak{}/\allowbreak{}N\allowbreak{}a\allowbreak{}t\allowbreak{}u\allowbreak{}r\allowbreak{}a\allowbreak{}l\allowbreak{}S\allowbreak{}e\allowbreak{}l\allowbreak{}e\allowbreak{}c\allowbreak{}t\allowbreak{}e\allowbreak{}d\allowbreak{}N\allowbreak{}o\allowbreak{}r\allowbreak{}m\allowbreak{}C\allowbreak{}o\allowbreak{}n\allowbreak{}t\allowbreak{}r\allowbreak{}a\allowbreak{}c\allowbreak{}t\allowbreak{}.\allowbreak{}l\allowbreak{}e\allowbreak{}a\allowbreak{}n}, definition \texttt{N\allowbreak{}a\allowbreak{}t\allowbreak{}u\allowbreak{}r\allowbreak{}a\allowbreak{}l\allowbreak{}S\allowbreak{}e\allowbreak{}l\allowbreak{}e\allowbreak{}c\allowbreak{}t\allowbreak{}e\allowbreak{}d\allowbreak{}C\allowbreak{}e\allowbreak{}l\allowbreak{}l\allowbreak{}N\allowbreak{}o\allowbreak{}r\allowbreak{}m}; its supplying argument ends at \texttt{N\allowbreak{}a\allowbreak{}t\allowbreak{}u\allowbreak{}r\allowbreak{}a\allowbreak{}l\allowbreak{}U\allowbreak{}n\allowbreak{}i\allowbreak{}f\allowbreak{}o\allowbreak{}r\allowbreak{}m\allowbreak{}N\allowbreak{}e\allowbreak{}a\allowbreak{}r\allowbreak{}.\allowbreak{}l\allowbreak{}e\allowbreak{}a\allowbreak{}n}, theorem \texttt{n\allowbreak{}a\allowbreak{}t\allowbreak{}u\allowbreak{}r\allowbreak{}a\allowbreak{}l\allowbreak{}\_\allowbreak{}s\allowbreak{}e\allowbreak{}l\allowbreak{}e\allowbreak{}c\allowbreak{}t\allowbreak{}e\allowbreak{}d\allowbreak{}\_\allowbreak{}c\allowbreak{}e\allowbreak{}l\allowbreak{}l\allowbreak{}\_\allowbreak{}n\allowbreak{}o\allowbreak{}r\allowbreak{}m}. The supplied adoption and scope are recorded in \texttt{r\allowbreak{}e\allowbreak{}s\allowbreak{}e\allowbreak{}a\allowbreak{}r\allowbreak{}c\allowbreak{}h\allowbreak{}/\allowbreak{}c\allowbreak{}y\allowbreak{}c\allowbreak{}l\allowbreak{}e\allowbreak{}0\allowbreak{}6\allowbreak{}7\allowbreak{}/\allowbreak{}a\allowbreak{}n\allowbreak{}a\allowbreak{}l\allowbreak{}y\allowbreak{}t\allowbreak{}i\allowbreak{}c\allowbreak{}/\allowbreak{}P\allowbreak{}R\allowbreak{}9\allowbreak{}5\allowbreak{}\_\allowbreak{}A\allowbreak{}N\allowbreak{}A\allowbreak{}L\allowbreak{}Y\allowbreak{}T\allowbreak{}I\allowbreak{}C\allowbreak{}\_\allowbreak{}I\allowbreak{}N\allowbreak{}T\allowbreak{}A\allowbreak{}K\allowbreak{}E\allowbreak{}.\allowbreak{}m\allowbreak{}d}, and its use with an arbitrary fixed positive gap is made explicit in the C071/C072 source arguments. These are provenance pointers for the unchanged input, not additional assumptions and not substitutes for the mathematical contract below. The supplied original sources and their copyright notices remain recoverable in this package.

\hypertarget{original-feature-and-finite-operator}{%
\subsubsection{3.1. Original feature and finite operator}\label{original-feature-and-finite-operator}}

Let \(\fitinline{\mathcal O}\) be the Eisenstein integers and \(\fitinline{NI}\) the norm of a nonzero integral ideal. We give the ordinary row feature explicitly. For a prime ideal \(\fitinline{\mathfrak p\ne(\lambda)}\), let \(\fitinline{c_{\mathfrak p}}\) be the cubic residue character: on a nonzero residue \(\fitinline{z}\), its value is the unique complex cube root of unity whose Eisenstein reduction is \(\fitinline{z^{(N\mathfrak p-1)/3}}\). Put its value at zero equal to zero. Let \(\fitinline{q_{\mathfrak p}}\) be the quadratic character of \(\fitinline{\mathcal O/\mathfrak p}\), zero at zero; in characteristic two it is one on every nonzero residue. Set \[\fitmath{
 \chi_{\mathfrak p}(z)=c_{\mathfrak p}(z)^2q_{\mathfrak p}(z),
 \qquad \chi_{(\lambda)}(z)=0.
}\tag{SN0a}\] At primes outside six this is the usual sextic residue character of Appendix A.1: squaring gives the cubic character and cubing the quadratic character. At the dyadic prime \(\fitinline{(2)}\), of norm four, it is the square of the cubic character; that raw source coordinate is not silently deleted. The fixed physical bad-prime mask later deletes dyadic columns where required by the source contract.

Define, including all multiplicities, \[\fitmath{
 \mathsf X(z,0)=0,\qquad \mathsf X(z,\mathcal O)=1,
 \qquad
 \mathsf X(z,I)=\prod_{\mathfrak p^a\parallel I}
                    \chi_{\mathfrak p}(z)^a\quad(I\ne0).
}\tag{SN0b}\] This is the original \(\fitinline{\chi_I(z)}\) on the supported ordinary ideals. On all ideals it is the exact natural feature denoted \texttt{i\allowbreak{}d\allowbreak{}e\allowbreak{}a\allowbreak{}l\allowbreak{}R\allowbreak{}o\allowbreak{}w\allowbreak{}H\allowbreak{}o\allowbreak{}m} in the supplied source interface. It contains no Gauss factor or completion. The source definition and bad-prime dictionary are recorded, with pinned original bytes, in \texttt{c\allowbreak{}o\allowbreak{}n\allowbreak{}s\allowbreak{}o\allowbreak{}l\allowbreak{}i\allowbreak{}d\allowbreak{}a\allowbreak{}t\allowbreak{}i\allowbreak{}o\allowbreak{}n\allowbreak{}/\allowbreak{}f\allowbreak{}e\allowbreak{}a\allowbreak{}t\allowbreak{}u\allowbreak{}r\allowbreak{}e\allowbreak{}\_\allowbreak{}s\allowbreak{}o\allowbreak{}u\allowbreak{}r\allowbreak{}c\allowbreak{}e\allowbreak{}s\allowbreak{}/\allowbreak{}S\allowbreak{}O\allowbreak{}U\allowbreak{}R\allowbreak{}C\allowbreak{}E\allowbreak{}\_\allowbreak{}B\allowbreak{}I\allowbreak{}N\allowbreak{}D\allowbreak{}I\allowbreak{}N\allowbreak{}G\allowbreak{}S\allowbreak{}.\allowbreak{}j\allowbreak{}s\allowbreak{}o\allowbreak{}n}; they are provenance for (SN0a)--(SN0b), whose definitions are local here.

Every prime factor present in \(\fitinline{I}\) retains its zero at \(\fitinline{z\in\mathfrak p}\), even for exponents divisible by six. The unit ideal is the absent product and has value one. A column divisible by \(\fitinline{(\lambda)}\) is zero. The feature is multiplicative in each variable and has modulus at most one; all fixed ray phases and further presentation masks stay in the actual coefficients. In particular, multiplying ordinary polynomials merges their ideals and never replaces this feature by a power.

A \textbf{natural free row} is one of the original nonzero element rows whose principal ideal is sixth-power-free. Units, every allowed residue class, bad-prime states, and all valuations from zero through five are included. No squarefree or supported-row restriction is imposed on this natural row set. For a finite family \(\fitinline{\mathcal F}\) of such rows and a finite family \(\fitinline{\mathcal S}\) of ordinary integral ideals, define

\[\fitmath{
 \mathcal N(\mathcal F,\mathcal S)
 =\left\|\big(\mathsf X(z,I)\big)_{I\in\mathcal S,z\in\mathcal F}\right\|_{\ell^2(\mathcal F)\to\ell^2(\mathcal S)}^2.
}\tag{SN1}\]

Thus, for every complex row vector \(\fitinline{(a_z)}\),

\[\fitmath{
 \sum_{I\in\mathcal S}
  \left|\sum_{z\in\mathcal F}a_z\mathsf X(z,I)\right|^2
 \le\mathcal N(\mathcal F,\mathcal S)\sum_{z\in\mathcal F}|a_z|^2.
}\tag{SN2}\]

The adjoint inequality gives the corresponding ordinary-source estimate; conjugating all entries preserves the operator norm. The source family is arbitrary: it may contain the unit and nonsquarefree ideals, and ideals in different factors may share primes. It need not avoid the fixed bad primes. A zero-ideal index may be retained as its zero source coordinate or erased exactly; it contributes nothing for these nonzero rows.

\hypertarget{literal-selected-slot-descriptor}{%
\subsubsection{3.2. Literal selected-slot descriptor}\label{literal-selected-slot-descriptor}}

For a fixed finite number \(\fitinline{n}\) of slots, a descriptor consists of real widths \(\fitinline{w_i}\), real energy labels \(\fitinline{e_i}\), finite sets \(\fitinline{\mathcal S_i}\) of nonzero ordinary ideals, and common coefficients \(\fitinline{b_i(I)}\). At scale \(\fitinline{Z\ge1}\), put

\[\fitmath{
 Q_i(z)=\sum_{I\in\mathcal S_i}b_i(I)\mathsf X(z,I).
}\tag{SN3}\]

For a mesh \(\fitinline{\omega\ge0}\) and coefficient bound \(\fitinline{C_{\rm slot}\ge0}\), the literal bounds are

\[\fitmath{
 \begin{gathered}
 0\le w_i\le\omega,\qquad
 NI\le2Z^{w_i}\quad(I\in\mathcal S_i),\\
 \sum_{I\in\mathcal S_i}|b_i(I)|^2\le C_{\rm slot}.
 \end{gathered}
}\tag{SN4}\]

The finite row family belongs to one common individual-amplitude cell if

\[\fitmath{
 Z^{e_i}\le|Q_i(z)|^2
 \qquad(z\in\mathcal F,\ 1\le i\le n).
}\tag{SN5}\]

Its total width and slope hypotheses are

\[\fitmath{
 0<w_{\min}\le\sum_iw_i,
 \qquad v\sum_iw_i\le\sum_ie_i.
}\tag{SN6}\]

The energy labels need not be nonnegative. Slot width zero is permitted, and all slots remain in the descriptor. The ordinary source variable \(\fitinline{I}\) in (SN1) does not enter the choice of the amplitude cell or its row family.

\hypertarget{natural-selected-cell-theorem-with-full-order-of-choices}{%
\subsubsection{3.3. Natural selected-cell theorem, with full order of choices}\label{natural-selected-cell-theorem-with-full-order-of-choices}}

For every fixed

\[\fitmath{
 A>0,\qquad P\ge3A+1,\qquad
 0<g<1/2,\qquad\varepsilon>0,
}\]

there exists an \textbf{early} \(\fitinline{\omega_0>0}\) with the following property. For every later finite slot number \(\fitinline{n}\), every \(\fitinline{C_{\rm slot}\ge0}\), and every fixed \(\fitinline{w_{\min}>0}\), there are constants \(\fitinline{C>0}\) and \(\fitinline{Z_0}\) such that, simultaneously for every

\[\fitmath{
 Z\ge\max(1,Z_0),\quad0\le\omega\le\omega_0,
 \quad\frac12+g\le v\le\frac34,
}\tag{SN7}\]

for every literal descriptor satisfying (SN4)--(SN6), every finite selected natural row family \(\fitinline{\mathcal F}\), every finite ordinary source family \(\fitinline{\mathcal S}\), and every pair of real scales

\[\fitmath{
 \begin{gathered}
 1\le Y\le Z^A,\qquad N(z)\le Y\quad(z\in\mathcal F),\\
 1\le X\le Z^P,\qquad NI\le X\quad(I\in\mathcal S),
 \end{gathered}
}\tag{SN8}\]

one has

\[\fitmath{
 \boxed{\displaystyle
 \mathcal N(\mathcal F,\mathcal S)
 \le C Z^{\varepsilon}\bigl(X+Y^{1-v}X^v\bigr).}
}\tag{SN9}\]

The constants are uniform in every descriptor, common coefficient mask, admissible \(\fitinline{v,Y,X}\), finite row selection and finite ordinary source family quantified after \(\fitinline{C,Z_0}\). Dependence on \(\fitinline{n,C_{\rm slot},w_{\min},A,P,g,\varepsilon}\) is permitted; dependence of the early mesh on the later slot number, later coefficient bounds, later positive total-width floor, row labels or source ideals is not permitted.

This theorem has a fixed positive gap. It asserts neither \(\fitinline{v=1/2}\), a gap tending to zero with \(\fitinline{Z}\), nor \(\fitinline{v<1/2}\). It supplies the original sextic feature and its complex conjugate. It does not license separate natural-row bounds for its square, cube or fourth power.

\hypertarget{exact-coefficient-and-quotient-adapters}{%
\subsubsection{3.4. Exact coefficient and quotient adapters}\label{exact-coefficient-and-quotient-adapters}}

The following consequences are elementary and are proved here so that applications of the input retain the original source.

\textbf{Common masks.} If \(\fitinline{|m(I)|\le1}\), replacing \(\fitinline{b_i(I)}\) by \(\fitinline{b_i(I)m(I)}\) leaves the support and widths unchanged and decreases its squared coefficient mass. The response of the new descriptor must still be checked on the selected rows; it is not inferred from the old response for an arbitrary mask. Fixed anchor characters, common punctures and presentation-zero factors are handled as these actual common coefficients. Independent real heights belong to their corresponding coefficient functions as unit-modulus norm phases. They do not change the squared coefficient mass: \(\fitinline{C_{\rm slot}}\) is fixed from the original profile and bounded real-power interval, independently of all live heights. Polynomial height dependence enters only the subsequent finite smooth-profile and Sobolev controls, after the selected norm. Allowing the norm constant to depend on a height-dependent coefficient bound would not establish the stated polynomial height uniformity and is not used.

\textbf{Literal quotient descendants.} Fix an element \(\fitinline{a}\). Put

\[\fitmath{
 b_i^{(a)}(I)=b_i(I)\mathsf X(a,I),\qquad
 Q_i^{(a)}(z)=\sum_Ib_i^{(a)}(I)\mathsf X(z,I).
}\tag{SN10}\]

Since \(\fitinline{|\mathsf X(a,I)|\le1}\), the descriptor bounds are preserved. Multiplicativity in the element argument gives the exact identity

\[\fitmath{
 Q_i^{(a)}(z)=Q_i(az).
}\tag{SN11}\]

Therefore any actual descendant family \(\fitinline{\mathcal F_a}\) for which \(\fitinline{az\in\mathcal F}\) inherits the same individual energy labels and total slope under this quotient descriptor. One must still check that the descendants are in the natural row class and their norm bound is the claimed \(\fitinline{Y}\). This is a transfer along an exact product-row identity, not transportation of a zero label or replacement by arbitrary powers of a character. Shared-prime zeros are retained in \(\fitinline{\mathsf X(a,I)}\).

\textbf{Merged ordinary ideals.} For finite polynomials in the original feature,

\[\fitmath{
 \left(\sum_I b(I)\mathsf X(z,I)\right)
 \left(\sum_J c(J)\mathsf X(z,J)\right)
 =\sum_H\left(\sum_{IJ=H}b(I)c(J)\right)\mathsf X(z,H).
}\tag{SN12}\]

This follows term by term from multiplicativity in the ideal variable. Every product fibre is merged as written, including units and shared primes. No squarefree projection, coprimality deletion or feature power appears. Cauchy within each fibre, followed by the ordinary ideal-divisor bound, gives for fixed \(\fitinline{\eta>0}\) and \(\fitinline{NH\le T}\)

\[\fitmath{
 \sum_H\left|\sum_{IJ=H}b(I)c(J)\right|^2
 \le C_\eta T^\eta
       \left(\sum_I|b(I)|^2\right)
       \left(\sum_J|c(J)|^2\right).
}\tag{SN13}\]

Indeed there are at most \(\fitinline{d(H)\le C_\eta(NH)^\eta}\) ideal factorizations \(\fitinline{H=IJ}\); expanding the resulting nonnegative sum gives the product on the right. Repeating the argument a fixed finite number of times yields the corresponding fixed-word coefficient mass, with a late constant and an arbitrarily assigned small power. The original detector coefficients, including the exact M\"{o}bius sign for an inverse polynomial or coefficient one for an ordinary polynomial, remain inside the convolution.

\textbf{Restrictions and conjugation.} Restricting \(\fitinline{\mathcal F}\) or \(\fitinline{\mathcal S}\) contracts the corresponding coordinate projections and therefore cannot increase (SN1). Thus original fixed-inducing predicates, initial keep, rays and finite source cuts may be retained. Both the feature and all coefficients may be conjugated: complex conjugation is an isometry, so the same operator bound holds. This proves the orientation adapter without identifying independently extracted heights or profiles.

An aggregate product-response condition is not literally (SN5). Its passage to common individual cells, including inactive-slot padding and a positive retained active width, is a separate finite histogram argument in the half-price proof. That argument must be supplied before invoking (SN9); the product condition is not substituted silently for the selected-cell hypothesis.

\hypertarget{native-initiation-and-the-complete-half-deletion-source}{%
\subsection{4. Native initiation and the complete half-deletion source}\label{native-initiation-and-the-complete-half-deletion-source}}

The source construction has two independent parts. Appendix C proves the complete arithmetic assembly: a priced packed energy, with its geometric and reference reserves, produces the original global source (S5). This section supplies that packed energy. The argument uses actual amplitudes on the packed rows; it does not transport an earlier row-selection mask through Poisson.

\hypertarget{the-generic-assembly-interface}{%
\subsubsection{4.1 The generic assembly interface}\label{the-generic-assembly-interface}}

For a strictly preceding parent \(\fitinline{G(B_0,L_0,c_0)}\), define \[\fitmath{
D(P,S,M)=\max\left(0,\min\left(S,
       \max\left(S-c_0M,{P+L_0S-B_0M\over L_0}\right)\right)\right).
}\tag{HP1}\] For nonnegative \(\fitinline{P,S}\), this is the width which must be deleted to enter the parent, clipped at deletion of all slots. The full parent is used only when \(\fitinline{P\le B_0M}\); if even the empty-slot point is outside the parent, the unrestricted paired ordinary moment supplies the empty product. This fallback is essential.

A \emph{price} \(\fitinline{\Pi(P,S,M)}\) bounds the literal packed energy by \(\fitinline{\mathcal A Z^{M+\Pi+\epsilon}}\), where \(\fitinline{\mathcal A}\) is exactly the radial, profile and polynomial-height control of (S5). In Appendix C's source geometry, the retained ordinary width is \(\fitinline{P}\), the retained slot width is \(\fitinline{S}\), and \(\fitinline{M=1+t-6u}\) is the actual packed height. The selected width \(\fitinline{u}\) is paid when \(\fitinline{\Pi(P,S,M)\le u+\epsilon/16}\), on the full perturbed geometric domain. A monotone locally Lipschitz price permits the small Fourier and whole-slot errors already charged in that domain.

The four reference reserves for a child \(\fitinline{G(B,L,c)}\), at root singleton width \(\fitinline{W}\) and selected width \(\fitinline{u_0(W)}\), are \[\fitmath{
\begin{aligned}
1-B+(L-2)W-\tfrac92u_0,&\qquad
\tfrac1{20}+\tfrac32(1-B)+(\tfrac32L-\tfrac52)W-\tfrac92u_0,\\
1-B+(L-1)W-4u_0,&\qquad
\tfrac13-2W-3u_0.
\end{aligned}
}\tag{HP2}\] Their positive uniform reserves pay, respectively, the complete reference interval, second reference, diagonal budget and positive reference upper endpoint. Appendix C supplies the actual common-divisor normalization, the masked packing identity, both restoring directions, LOW's constant \(\fitinline{2/9}\) comparison, HIGH's reference choice and all zero, fixed-inducing and tail contributions. Consequently this interface is an implication for the original source, not merely its nonexceptional positive leaf.

\hypertarget{a-native-source-for-every-strict-intercept}{%
\subsubsection{4.2 A native source for every strict intercept}\label{a-native-source-for-every-strict-intercept}}

\textbf{Lemma HP-N.} In the standing case \(\fitinline{\beta\ge51/100}\), for every \(\fitinline{0<B<1}\), \[\fitmath{
G(B,4,1/10)
}\tag{HP3}\] holds, including zero total marked width.

\textbf{Proof.} The natural canonical moment in Appendix A.7, after the already proved initialization \(\fitinline{\beta\le7/8}\), supplies the native price \(\fitinline{\Pi(P,S,M)=\tfrac34S}\) for all admitted bounded ordinary lengths. One can view the parent as \(\fitinline{G(B,5,-1)}\): it is vacuous, because a nonnegative width cannot be at most a negative multiple of the positive row height. Thus no prior improved endpoint or positive-cap source is required.

Use child slope \(\fitinline{4}\), cap \(\fitinline{1/10}\), crossing zero and \(\fitinline{u_0(W)=11W/25}\). In the actual perturbed geometry, if \(\fitinline{f}\) is the frozen width, the selection satisfies \[\fitmath{
{11\over25}(W-f)\le u\le {11\over25}(W-f)+m,
\quad 0\le S\le W-f-u,\quad0\le W\le1/10.
}\] Hence \(\fitinline{\tfrac34S\le u}\), since \(\fitinline{\tfrac34(W-f-u)-u\le-\tfrac1{50}(W-f)}\). Use the floor \(\fitinline{\epsilon/16}\), mesh at most \(\fitinline{1/1000}\), and geometric tolerance at most \(\fitinline{1/100}\). The height satisfies \[\fitmath{
1+t-6u\ge1-1/100-6(11/250+1/1000)>7/10,
\qquad 1+t-6u\le3.
}\] The upper bound is the actual defect bound \(\fitinline{t\le1+6u+1/100}\). The four expressions (HP2) are \[\fitmath{
1-B+W/50,\quad
1/20+3(1-B)/2+38W/25,\quad
1-B+31W/25,\quad
1/3-83W/25.
}\] They have the common positive lower bound \(\fitinline{\min(1-B,1/750)}\). All choices precede the live slots. The complete assembly of Appendix C now gives (HP3). Its zero-width branch uses the ordinary paired moment directly, so a strict positive width has not been assumed. This is a mathematical transcription of Joseph Bebel's supplied native initiation; no formal declaration is used in place of the argument. \ensuremath{\blacksquare}

\hypertarget{the-aggregate-selected-norm-at-an-arbitrary-fixed-gap}{%
\subsubsection{4.3 The aggregate selected norm at an arbitrary fixed gap}\label{the-aggregate-selected-norm-at-an-arbitrary-fixed-gap}}

Use the natural selected-cell norm of \S{}3 with fixed caps \(\fitinline{A>0}\), \(\fitinline{P_{\rm norm}\ge3A+1}\). For \(\fitinline{0<g\le1/8}\) and requested loss \(\fitinline{\lambda>0}\), it implies the aggregate estimate \[\fitmath{
\mathcal N(\mathcal R_{\rm high},\mathcal I_X)
 \ll Z^\lambda\{X+Y^{1-v}X^v\},
\quad 1/2+g\le v\le3/4,
}\tag{HP4}\] whenever the actual product of the original selected slot responses is at least \(\fitinline{Z^{vS}}\). The early mesh is independent of the subsequent finite slot number. All ordinary source ideals, units, common root masks and natural zeros remain.

Here is the finite passage from cells to the aggregate, including its choice order. Put \[\fitmath{
b=\min\{g/4,\lambda/[8(A+1)]\}.
}\] Choose the cell mesh for gap \(\fitinline{g/2}\) and loss \(\fitinline{\lambda/4}\), and reduce it to at most one. Only afterwards fix the slot number \(\fitinline{N}\), a coefficient bound \(\fitinline{C_{\rm slot}\ge1}\), and lower aggregate width \(\fitinline{s_{\min}>0}\). Choose \(\fitinline{t=\min(bs_{\min}/(N+1),1/2)}\) and \(\fitinline{n=\lceil2/t\rceil}\). The actual slot bound is \(\fitinline{|Q_i|^2\le256C_{\rm slot}Z^{w_i}}\). At the late onset \(\fitinline{256C_{\rm slot}\le Z^t}\), it is at most \(\fitinline{Z^{w_i+t}}\), with \(\fitinline{w_i\le1}\) and \(\fitinline{nt\ge2}\).

Define the label of a response to be the largest \(\fitinline{j\in\{0,\ldots,n\}}\) with \(\fitinline{Z^{jt}\le|Q_i|^2}\), or zero if that set is empty. For label zero replace that slot, on this entire cell, by the unit polynomial \(\fitinline{1}\), width zero and energy zero. Positive-label slots retain their original polynomial, width and coefficient mass. Thus the cell has one common legal descriptor, independent of the ordinary source ideal; the replacement is used to invoke the norm and does not assert equality of physical integrands.

Put \(\fitinline{E=\sum_i j_i t}\), and let \(\fitinline{S'}\) be the sum of widths of positive-label slots. The label bounds and the original aggregate response imply \[\fitmath{
 vS\le E+Nt,\qquad E\le S'+Nt,\qquad
 Nt\le b s_{\min}\le bS.
}\] It follows that \[\fitmath{
 S'\ge(v-2b)S\ge s_{\min}/4,
 \qquad E\ge(v-b)S\ge(v-b)S'.
}\] Consequently the padded cell satisfies the selected-cell slope condition at \(\fitinline{v-b\ge1/2+g/2}\), with lower width \(\fitinline{s_{\min}/4}\). There are at most \(\fitinline{(n+1)^N}\) label vectors. Apply the cell theorem to each actual restricted row family and sum its squared row norms. This argument covers arbitrarily small nonzero responses; it does not attempt to place unbounded negative logarithms in a finite histogram.

The envelope changes by at most \(\fitinline{Y^b\le Z^{Ab}}\), since \[\fitmath{
X+Y^{1-v+b}X^{v-b}
 \le Z^{Ab}\{X+Y^{1-v}X^v\}\qquad(X,Y\ge1).
}\] The exponent is \(\fitinline{\lambda/4+Ab\le\lambda}\). The finite label count multiplies the constant. It is chosen after the dictionary and before live data. This proves (HP4), including bounded row radii, with no live growing slot family and no norm theorem at \(\fitinline{v=1/2}\).

\hypertarget{low-delete-the-smallest-actual-responses}{%
\subsubsection{4.4 LOW: delete the smallest actual responses}\label{low-delete-the-smallest-actual-responses}}

Fix the original packed state, positive row profile and keep. Put \(\fitinline{U=Z^M}\), \(\fitinline{P=\ell_Z(X_1)+\ell_Z(X_2)}\), and \(\fitinline{S=\sum_iw_i}\), with \(\fitinline{0<M_{\min}\le M\le M_{\max}}\). There is one fixed \(\fitinline{C_Q\ge1}\), independent of rows, roots and heights, such that \(\fitinline{|Q_i(z)|^2\le C_Q^2Z^{w_i}}\). Set \(\fitinline{B_i=|Q_i|^2/C_Q^2}\), and choose \(\fitinline{q=1/2+g}\). Partition this very same packed energy according to \(\fitinline{\prod_iB_i\le Z^{qS}}\) or \(\fitinline{>Z^{qS}}\).

On LOW a zero factor makes the row zero. Otherwise delete all zero-width slots and order the positive-width slots by \(\fitinline{\log B_i/(w_i\log Z)}\). Delete the smallest prefix whose width \(\fitinline{e}\) reaches \(\fitinline{D=D(P,S,M)}\). Then \(\fitinline{D\le e\le D+\mathrm{mesh}}\), or every slot is deleted and \(\fitinline{e=S}\). The deleted normalized mass is at most \(\fitinline{Z^{qe}}\). Indeed, if \(\fitinline{l_0\le0}\) is the total zero-width logarithm, \(\fitinline{A}\) is the positive-width logarithm and \(\fitinline{A_e}\) its sorted-prefix sum, weighted ordering gives \(\fitinline{A_e\le eA/S}\). Thus \[\fitmath{
l_0+A_e\le(e/S)(l_0+A)+(1-e/S)l_0\le qe\log Z.
}\tag{HP5}\] When \(\fitinline{S=0}\), the assertion is immediate from \(\fitinline{B_i\le1}\). No lower bound on a nonzero response or on its logarithm is used.

Partition LOW by the deleted subset. On each of these at most \(\fitinline{2^N}\) parts, the surviving subset has common original coefficients, keep and profiles. Apply the complete parent source if \(\fitinline{P\le B_0M}\); otherwise \(\fitinline{D=S}\) and apply the unmarked ordinary pair. Positivity permits each row restriction. Assign the whole-slot overshoot and parent loss before the finite slot type. This proves \[\fitmath{
\mathcal E_{\rm LOW}\ll\mathcal A Z^{M+qD+\epsilon}.
}\tag{HP6}\] The factors \(\fitinline{2^N,C_Q^{2N}}\) are fixed constants, not scale powers.

\hypertarget{high-convolve-ordinary-ideals-and-sum-every-root}{%
\subsubsection{4.5 HIGH: convolve ordinary ideals and sum every root}\label{high-convolve-ordinary-ideals-and-sum-every-root}}

Both ordinary factors use the same row feature. Convolving their actual coefficients gives \[\fitmath{
A_1(u)A_2(u)=\sum_I c(I)\chi_I(u),\quad
\operatorname{supp}c\subseteq\{NI\le(81/16)Z^P\},\quad
\sum_I|c(I)|^2\ll Z^\epsilon\mathcal P(p_1,p_2)^2.
}\tag{HP7}\] For each factor the squared coefficient mass follows from ideal counting and central normalization. Cauchy on the complete multiplication fibre costs the fixed-factor ideal-divisor bound. This proves the last estimate; it retains units, shared primes and all ordinary ideals. At a fixed sixth root \(\fitinline{e}\), replace \(\fitinline{c(I)}\) by \(\fitinline{\mathbf1_{(I,e)=1}c(I)}\). This is exactly the original root mask, because \(\fitinline{(I_1I_2,e)=1}\) is equivalent to both factors being coprime to \(\fitinline{e}\). Its coefficient square mass contracts.

The actual root factory gives the same masked prime polynomials on the core \(\fitinline{u}\), even when \(\fitinline{(u,e)\ne1}\). With \(\fitinline{Y\ll Z^M/(Ne)^6}\), (HP4) at \(\fitinline{v=q}\) and (HP7) yield \[\fitmath{
\mathcal E_e\ll\mathcal A Z^{S+\epsilon}
 \{Z^P+Y^{1-q}Z^{qP}\}.
}\tag{HP8}\] There is a second bound on the identical root subset: \[\fitmath{
\mathcal E_e\ll\mathcal A Z^{S+\epsilon}Y.
}\tag{HP9}\] To prove it, use the trivial original prime-product bound and enlarge the positive row set to the nonprincipal core rows with its fixed-across-core puncture. Apply Appendix A.5's full paired moment. The initial keep continues to exclude principal/fixed-inducing rows; none is added by this enlargement. The ambient scale is \(\fitinline{Z}\), so bounded \(\fitinline{Y}\) and ordinary lengths as large as fixed powers of \(\fitinline{Z}\) are included. Fixed support factors such as \(\fitinline{81/16}\) are put into the admitted scale before applying the norm; an empty core family contributes zero.

Combine (HP8) and (HP9) \textbf{before} summing roots. For \(\fitinline{0\le P\le M}\), ideal counting and a split at \(\fitinline{Ne=Z^{(M-P)/6}}\) give \[\fitmath{
\sum_e\min\{Z^P,Z^M/(Ne)^6\}
 \ll Z^{(5P+M)/6}.
}\tag{HP10}\] The other norm term sums to \[\fitmath{
Z^{(1-q)M+qP}\sum_e(Ne)^{-6(1-q)}
 \ll Z^{(1-q)M+qP},
}\tag{HP11}\] because \(\fitinline{q\le5/8<5/6}\). The exponent in (HP10) is at most that in (HP11), since their difference is \(\fitinline{(5/6-q)(M-P)\ge0}\). If \(\fitinline{P\ge M}\), the full unmarked pair bounds the energy by \(\fitinline{\mathcal A Z^{M+S+\epsilon}}\). Therefore HIGH always has price at most \[\fitmath{
H_{1/2}(P,S,M)=S+(P-M)/2.
}\tag{HP12}\] For \(\fitinline{P\le M}\), the actual price \(\fitinline{S+q(P-M)}\) is no larger than (HP12); for \(\fitinline{P\ge M}\), the preceding unmarked price \(\fitinline{S}\) is no larger. No first norm term, root or cross term has been discarded.

The aggregate estimate was used for \(\fitinline{S\ge s_{\min}>0}\). If \(\fitinline{S<s_{\min}}\), the trivial prime product and unmarked paired moment cost \(\fitinline{Z^{M+s_{\min}+\epsilon}}\). Choose \(\fitinline{s_{\min}}\) from the final loss before slots, and use the nonnegative price below. This pays all small and empty slot sets.

\hypertarget{the-exact-half-price-with-unbounded-total-slot-width}{%
\subsubsection{4.6 The exact half price, with unbounded total slot width}\label{the-exact-half-price-with-unbounded-total-slot-width}}

\textbf{Theorem HP.} A complete preceding parent with \(\fitinline{0<B_0<1}\), \(\fitinline{L_0>1}\), \(\fitinline{0<c_0<1/6}\) supplies \[\fitmath{
\boxed{\Pi_{1/2}(P,S,M)=
 \max\{D(P,S,M)/2,\ S+(P-M)/2\}.}
}\tag{HP13}\] It is a full priced packed-source estimate with the shared contract of \S{}2, not a norm at the boundary slope.

\textbf{Proof.} Put \(\fitinline{T=\Pi_{1/2}}\). The HIGH estimate is (HP12). For LOW, \[\fitmath{
qD\le T+2gM_{\max}.
}\tag{HP14}\] If \(\fitinline{S\le2M_{\max}}\), use \(\fitinline{D\le S}\) to bound \(\fitinline{qD=D/2+gD\le T+2gM_{\max}}\). If \(\fitinline{S\ge2M_{\max}}\), then \[\fitmath{
qD\le3S/4\le S-M_{\max}/2\le S+(P-M)/2\le T.
}\] Choose a fixed \(\fitinline{g>0}\) with \(\fitinline{g\le1/8}\) and \(\fitinline{2gM_{\max}}\) below its assigned share of the final loss. This choice precedes the selected-cell mesh, the parent mesh, the dictionary and every live parameter. Thus (HP6), the complete HIGH estimate and the small-slot payment prove (HP13). No bound on total slot width was silently chosen before the dictionary.

The deletion function is monotone in \(\fitinline{P,S}\), and its finite Lipschitz constants follow from maxima, minima and the affine functions in (HP1). The HIGH term has coordinate constants \(\fitinline{1/2,1,1/2}\). Taking the maximum gives the generic price needed by Appendix C; a specialized adapter with a smaller hardcoded slot constant is not used. Independent ordinary heights enter only through common coefficients and the already fixed profile/height controls. Original sublattice packing is restored by its positive enlargement before the split. All roots, shortenings, principal/reference payments and coefficient-box operations are those proved in Appendix C. \ensuremath{\blacksquare}

\hypertarget{complete-global-sources-on-the-open-parameter-domain}{%
\subsection{5. Complete global sources on the open parameter domain}\label{complete-global-sources-on-the-open-parameter-domain}}

\hypertarget{statement-and-analytic-inputs}{%
\subsubsection{5.1. Statement and analytic inputs}\label{statement-and-analytic-inputs}}

Under the standing case \(\fitinline{\beta\ge51/100}\), the literal global source contract holds for

\[\fitmath{
\boxed{G(B,L,c)\quad\text{whenever}\quad
0<B<1,\qquad L>\max(1,2B),\qquad0<c<1/6.}
}\tag{UF1}\]

Every source retains both original ordinary factors, their literal ideal coefficients and real annular weights, all selected original primes, every natural sixth root, common root masks, natural and presentation zeros, the initial keep, independent ordinary profiles/heights, shortenings, reflection/Fourier operations and coefficient boxes. Its early mesh precedes the finite slot dictionary. Each strict triple has a finite proof. No uniform constant is asserted as \(\fitinline{B\uparrow1}\), \(\fitinline{L\downarrow2B}\), or \(\fitinline{c\uparrow1/6}\), and no boundary source is used by this theorem.

The locally proved native source gives

\[\fitmath{
B<1,\quad\beta\ge51/100\quad\Longrightarrow\quad G(B,4,1/10),
}\tag{UF4}\]

including zero marked width. The full packed half-price theorem gives, from any strictly preceding parent \(\fitinline{G(B,L_0,c_0)}\),

\[\fitmath{
\Pi(P,s,M)=\max\left\{
\tfrac12\operatorname{deletionNeed}(B,L_0,c_0;P,s,M),
s+\tfrac12(P-M)\right\}.
}\tag{UF3}\]

The price is derived using a fixed positive selected-norm gap chosen from the requested final loss; no selected operator norm at the one-half boundary is assumed. The following finite induction applies the complete source-step assembly to this price. Its scalar checks are proved uniformly in the stage index and in the full closed marked-width interval.

\hypertarget{a-finite-source-path-with-bounds-uniform-in-its-index}{%
\subsubsection{5.2. A finite source path with bounds uniform in its index}\label{a-finite-source-path-with-bounds-uniform-in-its-index}}

First suppose \[\fitmath{
            \max(1,2B)<L<4,\qquad0<c<1/6.
}\] Put \[\fitmath{
\begin{aligned}
 c_{\rm seed}&=\min(c,1/10),&d_L&=4-L,&d_c&=c-c_{\rm seed},\\
 D_*&=\tfrac32(L-2B),&
 r_1&=(1-B)(1-2c),&
 r_4&=\tfrac13-2c .
\end{aligned}                                                     
}\tag{UF5}\] Here \(\fitinline{c_{\rm seed},d_L,D_*,r_1,r_4}\) are positive and \(\fitinline{d_c\ge0}\). Restricting the native source (UF4) to cap \(\fitinline{c_{\rm seed}}\) gives the complete initial parent.

Choose the explicit constants \[\fitmath{
\begin{aligned}
 K&=1+\frac{d_L}{D_*}
       +\frac{d_c}{(1-6c)c_{\rm seed}},\\
 a&=\min\left\{\frac12,\frac1{120c},
                    \frac{r_1}{9c},\frac{r_4}{6c}\right\},\\
 N&=\left\lfloor K/a\right\rfloor+1,\qquad A=K/N .
\end{aligned}                                                     
}\tag{UF6}\] Then \(\fitinline{N\ge1}\) is a finite integer and \(\fitinline{0<A<a}\). Set \[\fitmath{
 d=d_L/N,\qquad e=d_c/N,\qquad
 L_i=4-id,\qquad c_i=c_{\rm seed}+ie
       \quad(0\le i\le N).                                       
}\tag{UF7}\] The intercept is \(\fitinline{B}\) at every stage. The selector is \(\fitinline{u_0=AW}\), the crossing is zero, and the last source is exactly \(\fitinline{G(B,L,c)}\). Every stage has \[\fitmath{
 L\le L_i<L_{i-1}\le4,\quad
 c_{\rm seed}\le c_{i-1}\le c_i\le c<1/6,\quad
 D_i:=\tfrac32(L_{i-1}-2B)\ge D_*>0.                              
}\tag{UF8}\]

The number \(\fitinline{N}\) can be very large near a strict boundary. No iteration over that number is needed to verify the algebra: the following bounds hold for every integer \(\fitinline{1\le i\le N}\) and every real \(\fitinline{0\le W\le c_i}\). They give a genuine finite induction, not an infinite limiting argument. Constants and onset may depend on this fixed finite \(\fitinline{N}\), as the literal source contract allows.

\hypertarget{low-payment-and-the-actual-retained-parent-cap}{%
\subsubsection{5.3. LOW payment and the actual retained parent cap}\label{low-payment-and-the-actual-retained-parent-cap}}

At stage \(\fitinline{i}\), write \(\fitinline{L_0=L_{i-1}}\), \(\fitinline{c_0=c_{i-1}}\), and \[\fitmath{
 P=B-L_iW,\qquad s=W-u_0,\qquad M=1-6u_0,\qquad u_0=AW.
}\] The half-price LOW overshoot denominator is \(\fitinline{D_i>0}\). The coefficient of \(\fitinline{W}\) in the LOW payment reserve is \[\fitmath{
\begin{aligned}
 L_0A-\tfrac12\{d+(6B-L_0)A\}
 &=D_iA-\tfrac12d\\
 &\ge\frac{D_*K-d_L/2}{N}>0 .                                   
\end{aligned}
}\tag{UF9}\] The final strict inequality follows directly from \(\fitinline{K=1+d_L/D_*+d_c/((1-6c)c_{\rm seed})}\). The payment itself is nonnegative at \(\fitinline{W=0}\), as required; a strictly positive payment at zero width is not asserted or needed.

For the parent cap, the exact reserve is \[\fitmath{
 R_C(W)=c_0(1-6AW)-(W-AW)
       =c_0+\{-1+(1-6c_0)A\}W.                                  
}\tag{UF10}\] Its slope is negative because \(\fitinline{0<c_0<1/6}\) and \(\fitinline{A<1/2}\). At the child endpoint \(\fitinline{W=c_i=c_0+e}\), \[\fitmath{
\begin{aligned}
 R_C(c_i)
   &=-e+(1-6c_0)Ac_i\\
   &\ge\frac{(1-6c)c_{\rm seed}K-d_c}{N}>0 .
\end{aligned}                                                     
}\tag{UF11}\] Every factor used for the lower bound is nonnegative. In particular, the cap increase has been paid by the actual retained parent cap, not by an artificial requirement that caps decrease. Define the positive uniform cap reserve \[\fitmath{
 r_C=\min\left(c_{\rm seed},
        \frac{(1-6c)c_{\rm seed}K-d_c}{N}\right)>0 .                        
}\tag{UF12}\] Then \(\fitinline{R_C(W)\ge r_C}\) throughout every stage and full width interval.

\hypertarget{reference-height-and-high-gates}{%
\subsubsection{5.4. Reference, height and HIGH gates}\label{reference-height-and-high-gates}}

Let \(\fitinline{u_0=AW}\), and write the four literal reference reserves as \[\fitmath{
\begin{aligned}
 R_1&=1-B+(L_i-2)W-\tfrac92u_0,\\
 R_2&=\tfrac1{20}+\tfrac32(1-B)
                         +(\tfrac32L_i-\tfrac52)W-\tfrac92u_0,\\
 R_3&=1-B+(L_i-1)W-4u_0,\\
 R_4&=\tfrac13-2W-3u_0 .
\end{aligned}                                                     
}\tag{UF13}\] The unselected first reserve satisfies the exact decomposition \[\fitmath{
 1-B+(L_i-2)W
 =(1-B)(1-2W)+(L_i-2B)W\ge r_1.                                 
}\tag{UF14}\] Here \(\fitinline{L_i>2B}\), \(\fitinline{W\ge0}\), and \(\fitinline{W\le c<1/6}\) justify every sign. Therefore the choice \(\fitinline{A<r_1/(9c)}\) gives \[\fitmath{
 R_1\ge r_1-\tfrac92Ac>\tfrac12r_1>0 .
}\] The remaining reserves obey \[\fitmath{
\begin{aligned}
 R_2&=\tfrac32R_1+\tfrac1{20}+\tfrac12W+\tfrac94u_0>0,\\
 R_3&=R_1+W+\tfrac12u_0>0,\\
 R_4&\ge r_4-3Ac>\tfrac12r_4>0 .
\end{aligned}                                                     
}\tag{UF15}\] Thus a single positive reference tolerance, for example the minimum of \(\fitinline{r_1/2}\) and \(\fitinline{r_4/2}\), works uniformly in all indices and widths.

The nominal height has the positive reserve \[\fitmath{
 r_H=\tfrac1{10}-6Ac>\tfrac1{20},\qquad
              M=1-6u_0\ge9/10+r_H .                             
}\tag{UF16}\] This uses \(\fitinline{A<1/(120c)}\). The half HIGH and complete root-cone identities are \[\fitmath{
\begin{aligned}
 u_0-\{s+\tfrac12(P-M)\}
        &=\tfrac12R_1+\tfrac54u_0
          \ge\tfrac12(r_1-\tfrac92Ac)=:r_J>0,\\
 M-P-2s&=R_1+\tfrac12u_0>0 .
\end{aligned}                                                     
}\tag{UF17}\] The first reference condition therefore pays HIGH; no additional unproved source-length estimate is inserted.

These are bounds over the whole closed real interval \(\fitinline{0\le W\le c_i}\). They do not infer an interval statement from sampled contacts.

\hypertarget{the-full-perturbed-ledger}{%
\subsubsection{5.5. The full perturbed ledger}\label{the-full-perturbed-ledger}}

Use the full half-price theorem (HP13). Its HIGH coordinate Lipschitz constants are \(\fitinline{1/2,1,1/2}\); the deletion-need Lipschitz constant, enlarged to at least one, controls the complete price. The generic source assembly uses this full price directly.

At any stage and any requested child error \(\fitinline{\varepsilon>0}\), put \[\fitmath{
\begin{aligned}
 m&=\min(r_H,r_J),&
 \mathrm{floor}&=\varepsilon/16,\\
 \mathrm{meshCap}&=m/12,&
 \mathrm{tailCap}
 &=\min\{m/4,r_C,1,L_0\varepsilon/24,\varepsilon/32\}.
\end{aligned}                                                     
}\tag{UF18}\] All are positive and precede the actual geometry and finite slots. The actual source-side geometry is \[\fitmath{
\begin{gathered}
 u_0\le u\le u_0+\mathrm{mesh},\quad s\le W-u,\quad
 P_C\le B-L_iW+\mathrm{tl},\\
 M=1+t-6u,\qquad
 -\mathrm{tl}\le t\le1+6u+\mathrm{tl}.
\end{gathered}
}\] The chosen budgets give \(\fitinline{9/10\le M\le3}\). More explicitly, the lower height loses at most \(\fitinline{6\mathrm{mesh}+\mathrm{tl}\le3m/4\le3r_H/4}\); the upper height follows from \(\fitinline{M\le2+\mathrm{tl}\le3}\).

The perturbed parent-cap reserve is still complete: \[\fitmath{
 s-c_0M
 \le-r_C+(6c_0-1)(u-u_0)+c_0\mathrm{tl}
 \le-r_C+c_0\mathrm{tl}\le0 .                                   
}\tag{UF19}\] The overshoot coefficient is negative, and \(\fitinline{c_0<1}\), \(\fitinline{\mathrm{tl}\le r_C}\).

Let \(\fitinline{F=P_C+L_0s-BM}\). The inherited defect calculation, now with zero crossing and constant intercept, gives \[\fitmath{
 F\le dW+(6B-L_0)u+(B+1)\mathrm{tl}.
}\] Equations (UF9) and (UF8) imply, first at \(\fitinline{u_0=AW}\) and then at \(\fitinline{u\ge u_0}\), \[\fitmath{
 \frac{F}{2L_0}
 \le u+\frac{(B+1)\mathrm{tl}}{2L_0}
 \le u+\varepsilon/16 .
}\] Since (UF19) holds, the complete inherited inequality \(\fitinline{\operatorname{deletionNeed}\le\max(0,F/L_0)}\) applies. The right side \(\fitinline{u+\varepsilon/16}\) is nonnegative, so the same LOW payment holds when \(\fitinline{F<0}\). This includes width zero and the entire small-width floor.

For HIGH, the exact perturbed bound is \[\fitmath{
 s+\tfrac12(P_C-M)-u
       \le-r_J+\mathrm{mesh}+\mathrm{tl}\le0 .
}\] Indeed the two perturbations total at most \(\fitinline{m/3\le r_J/3}\). The unchanged more conservative ledger allowance \(\fitinline{5\mathrm{mesh}+2\mathrm{tl}\le11m/12\le r_J}\) is also valid. Finally \(\fitinline{\varepsilon/16\le\varepsilon/(16(3/4))}\), so the generic floor field remains paid. These estimates prove the full \texttt{L\allowbreak{}e\allowbreak{}d\allowbreak{}g\allowbreak{}e\allowbreak{}r\allowbreak{}A\allowbreak{}t}, not merely its nominal counterpart.

\hypertarget{complete-induction-and-quantifier-order}{%
\subsubsection{5.6. Complete induction and quantifier order}\label{complete-induction-and-quantifier-order}}

At each stage take the generic data with constant intercept \(\fitinline{B}\), parent \(\fitinline{(L_{i-1},c_{i-1})}\), child \(\fitinline{(L_i,c_i)}\), crossing zero, \(\fitinline{H_{\min}=9/10}\), \(\fitinline{H_{\max}=3}\), \(\fitinline{L_{\mathrm{plain}}=4}\), \(\fitinline{\kappa=3/4}\), selector \(\fitinline{AW}\), and price (UF3). The face, crossing, envelope, selector, and reference fields follow from (UF8) and (UF13)--(UF17). The selector is monotone on \(\fitinline{[0,c_i]}\) and lies between zero and \(\fitinline{W}\) because \(\fitinline{0<A<1}\). The generic source-step ledger is fully paid by (UF18)--(UF19).

Assume only that the strictly preceding parent is complete. The full half-price theorem supplies the packed positive-energy estimate at price (UF3), with row exponent between \(\fitinline{9/10}\) and \(\fitinline{4}\) and ordinary source range bounded by \(\fitinline{4}\). Apply the locally proved generic cube-free source assembly: its contrary-side second-divisor provider, diagonal supplier, principal/reference suppliers and zero-width complement have exactly these inputs. The constant-intercept, zero-crossing geometry is permitted because its crossing obligation is \(\fitinline{0\le s_0}\), \(\fitinline{(L_0-L)s_0\le B_0-B}\), with a parent-cap condition only when \(\fitinline{s_0>0}\); here both sides are zero. The source-step theorem therefore gives the complete child from the complete parent.

This composition preserves all reference, diagonal, second-divisor, principal, floor, profile, small-height and large-height sectors. Both original ordinary factors and every selected original prime retain their coefficients, masks, roots, and natural zeros. The complete source is global; restricting it to an original detector band uses positivity. There is no Poisson transport of an original zero label.

For each requested child error, request the preceding parent at the smaller error prescribed by the complete source assembly, for example \(\fitinline{\varepsilon/(4\max(1,4))=\varepsilon/16}\). On packing to base \(\fitinline{Z^M}\), \(\fitinline{M\le4}\) costs at most \(\fitinline{\varepsilon/4}\). Choose the fixed positive norm gap for the half-price inlet from this requested final error. Choose the child mesh below \(\fitinline{9/10}\) times the parent mesh, below that inlet's early mesh, and below all geometric and reference tolerances. Only then choose the finite slot type, fixed ray data, characters, profiles and radial caps.

The supplier and assembly then choose finite seminorm controls, constants, polynomial height degrees and onset before the live widths, shortenings and heights. A finite union and maximum combine these controls; fixed height degrees precede the later small height exponent, and the external tail order is chosen last. Each stage is a theorem for every requested error. There is no accumulation of one fixed error through \(\fitinline{N}\) stages.

Finite induction from the complete native source (UF4), restricted to \(\fitinline{c_{\rm seed}}\), proves the final source \(\fitinline{G(B,L,c)}\). If \(\fitinline{L\ge4}\), first apply the proved case at slope \(\fitinline{3}\), which exceeds \(\fitinline{\max(1,2B)}\). Then \(\fitinline{P+LS\le B}\) and \(\fitinline{S\ge0}\) imply \(\fitinline{P+3S\le B}\), so cone inclusion gives the desired larger-slope source. This proves all of (UF1).

\hypertarget{a-convenient-common-count-source}{%
\subsubsection{5.7. A convenient common count source}\label{a-convenient-common-count-source}}

A fixed specialization used by the main count bank is

\[\fitmath{
G\left(\frac{999999999999}{10^{12}},
       \frac{20000000001}{10^{10}},\frac4{25}\right).
}\tag{UF2}\]

For this example, the explicit construction above has

\[\fitmath{
K=\frac{2000000002348}{153},\quad
 a=\frac{17}{36000000000000},\quad
 N=27681660932152249134949,
}\] \[\fitmath{
 A=\frac{2000000002348}{4235294122619294117647197},\quad
 D_*=\frac{153}{10^{12}},\quad
 r_1=\frac{17}{25000000000000},\quad r_4=\frac1{75}.
}\]

The stage-uniform proof, rather than enumeration of that integer, establishes every step. It asserts a finite mathematical induction, not a new kernel run. No endpoint is a premise of the construction beyond the fixed original standing case. The same theorem supplies each nearby strict source used in the algebraic endpoint return.

\hypertarget{one-source-for-the-full-inherited-count-bank}{%
\subsection{6. One source for the full inherited count bank}\label{one-source-for-the-full-inherited-count-bank}}

This is an equivalent main-path replacement made in the one-time consolidation, not an endpoint improvement. It removes numerical source-generation ancestry from the attained endpoint's dependency chain while preserving every actual count branch and admission window. The old source constructions remain available as research alternatives; their statements and historical evidence are not altered.

\hypertarget{the-domain-inclusion-lemma}{%
\subsubsection{The domain-inclusion lemma}\label{the-domain-inclusion-lemma}}

Write a source face as

\[\fitmath{
 F=(B,L,a,b,c_-,c_+),
}\]

meaning that at row exponent \(\fitinline{M\ge0}\), its two ordinary exponents and total marked width satisfy

\[\fitmath{
 aM\le m_1,m_2\le bM,\qquad c_-M\le s\le c_+M,
 \qquad m_1+m_2+Ls\le BM.
}\]

The complete global source \(\fitinline{G(B,L,c)}\) has face \(\fitinline{(B,L,0,B,0,c)}\); the counting inlet uses the broad face \(\fitinline{(B,L,0,1/2,0,c)}\). All other data in these contracts are literal and identical: ideal coefficients and real annular powers, original ray-prime pools, fixed masks and holes, inducing predicates, natural/presentation zeros and initial keep, both ordinary profiles and their independent real heights, root masks, reflection/Fourier operations, shortenings, and coefficient boxes.

\textbf{Lemma D1 (exact face inclusion).} Let \(\fitinline{G(B_u,L_u,c_u)}\) be a complete global source with \(\fitinline{B_u\ge1/2}\). A broad face with parameters \(\fitinline{(B_b,L_b,c_b)}\) is supplied by this source if

\[\fitmath{
 0\le c_b\le c_u,\qquad
 B_b+\max(L_u-L_b,0)c_b\le B_u. 
}\tag{D1}\]

\textbf{Proof.} The ordinary bounds \(\fitinline{0\le m_i\le M/2}\) lie in \(\fitinline{[0,B_uM]}\), and \(\fitinline{0\le s\le c_bM\le c_uM}\). Moreover

\[\fitmath{
\begin{aligned}
 m_1+m_2+L_us
 &=(m_1+m_2+L_bs)+(L_u-L_b)s\\
 &\le \{B_b+\max(L_u-L_b,0)c_b\}M\le B_uM.
\end{aligned}
}\]

Thus every admissible input to the broad face is an admissible input to the same complete global estimate, with the same coefficients and data. Restriction of a quantified estimate proves the claim. Its early mesh, seminorm orders, constants, height degrees and onset retain their original order and uniformity. No new row selection, coefficient class, moment, or lower-width bound is introduced. \ensuremath{\blacksquare}

The same argument gives global-to-global inclusion if additionally the weaker ordinary upper bound is at most \(\fitinline{B_u}\). For present purposes only the broad face is required. In particular a small slope \(\fitinline{L_b<2B_b}\) does not preclude inclusion on a bounded marked cap: the correct test is (D1), not a comparison of slopes alone.

\hypertarget{a-single-fixed-source-for-all-24-enhanced-branches}{%
\subsubsection{A single fixed source for all 24 enhanced branches}\label{a-single-fixed-source-for-all-24-enhanced-branches}}

The complete open source theorem supplies

\[\fitmath{
 B_u=\frac{999999999999}{10^{12}},\qquad
 L_u=\frac{20000000001}{10^{10}},\qquad
 c_u=\frac4{25},\qquad G_u=G(B_u,L_u,c_u). 
}\tag{D2}\]

Indeed \(\fitinline{0<B_u<1}\), \(\fitinline{L_u>\max(1,2B_u)}\), and \(\fitinline{0<c_u<1/6}\). This source uses only the standing case \(\fitinline{\beta\ge51/100}\), which holds in any contrary-supremum argument at the attained endpoint. It does not use an intermediate improved endpoint as a premise.

In the table, \(\fitinline{c_b}\) is the \emph{declared count cap}. The actual count width is \(\fitinline{z_b=c_b-10^{-7}}\), so checking the larger declared cap also checks the actual one. Affine and capped indicate the actual common-witness formula, not a modification of the source. The old natural count is separate and remains available everywhere.

\begin{landscape}\small\renewcommand{\arraystretch}{1.8}
\begin{longtable}[]{@{}
  >{\raggedright\arraybackslash}p{(\columnwidth - 8\tabcolsep) * \real{0.1000}}
  >{\raggedright\arraybackslash}p{(\columnwidth - 8\tabcolsep) * \real{0.1000}}
  >{\raggedleft\arraybackslash}p{(\columnwidth - 8\tabcolsep) * \real{0.3500}}
  >{\raggedleft\arraybackslash}p{(\columnwidth - 8\tabcolsep) * \real{0.3500}}
  >{\raggedleft\arraybackslash}p{(\columnwidth - 8\tabcolsep) * \real{0.1000}}@{}}
\toprule\noalign{}
\begin{minipage}[b]{\linewidth}\raggedright
Branch
\end{minipage} & \begin{minipage}[b]{\linewidth}\raggedright
Mode
\end{minipage} & \begin{minipage}[b]{\linewidth}\raggedleft
\(\fitinline{B_b}\)
\end{minipage} & \begin{minipage}[b]{\linewidth}\raggedleft
\(\fitinline{L_b}\)
\end{minipage} & \begin{minipage}[b]{\linewidth}\raggedleft
\(\fitinline{c_b}\)
\end{minipage} \\
\midrule\noalign{}
\endhead
\bottomrule\noalign{}
\endlastfoot
Cnew & affine & \(\fitinline{\dfrac{755433010311}{800000000000}}\) & \(\fitinline{\dfrac{499999854189139951396678196647534309}{291621720097206643606704931382000000}}\) & \(\fitinline{\dfrac{10049}{100000}}\) \\
Pnew & capped & \(\fitinline{\dfrac{132365500000000000000000000000000000}{134691681706757286391944389662954587}}\) & \(\fitinline{\dfrac{1000000000000000000000000000000}{457834420624410799685731440459}}\) & \(\fitinline{\dfrac{31}{500}}\) \\
Cphys & affine & \(\fitinline{\dfrac{132365500000000000000000000000000000}{134691681706757286391944389662954587}}\) & \(\fitinline{\dfrac{1000000000000000000000000000000}{457834420624410799685731440459}}\) & \(\fitinline{\dfrac{9999}{100000}}\) \\
Cold & affine & \(\fitinline{\dfrac{19452098647453}{20000000000000}}\) & \(\fitinline{\dfrac{4230401}{2000000}}\) & \(\fitinline{\dfrac{9999}{100000}}\) \\
Pold & capped & \(\fitinline{\dfrac{3945996453113}{4000000000000}}\) & \(\fitinline{\dfrac{4603201}{2000000}}\) & \(\fitinline{\dfrac{31}{500}}\) \\
Cphysold & affine & \(\fitinline{\dfrac{3945996453113}{4000000000000}}\) & \(\fitinline{\dfrac{4603201}{2000000}}\) & \(\fitinline{\dfrac{9999}{100000}}\) \\
Pwide & capped & \(\fitinline{\dfrac{132365500000000000000000000000000000}{134691681706757286391944389662954587}}\) & \(\fitinline{\dfrac{1000000000000000000000000000000}{457834420624410799685731440459}}\) & \(\fitinline{\dfrac{9}{100}}\) \\
Pcount & capped & \(\fitinline{\dfrac{9442912138791}{10000000000000}}\) & \(\fitinline{\dfrac{249999854189139951396678196647534309}{145810860048603321803352465691000000}}\) & \(\fitinline{\dfrac{10099}{100000}}\) \\
Pcountmid & capped & \(\fitinline{\dfrac{9442912138791}{10000000000000}}\) & \(\fitinline{\dfrac{249999854189139951396678196647534309}{145810860048603321803352465691000000}}\) & \(\fitinline{\dfrac{10049}{100000}}\) \\
Cbridge0 & affine & \(\fitinline{\dfrac{1868587268906691992741201750000000000000000000000000000}{1899396507298621824297773961734367699086839004747994531}}\) & \(\fitinline{\dfrac{500000000000000000000000000000}{227175886267711795724080853911}}\) & \(\fitinline{\dfrac{9999}{100000}}\) \\
Cbridge1 & affine & \(\fitinline{\dfrac{21644250000000000000000000000000000}{22150759854695188032554610654516749}}\) & \(\fitinline{\dfrac{500000000000000000000000000000}{238015564731513913664253158079}}\) & \(\fitinline{\dfrac{9999}{100000}}\) \\
Cbridge2 & affine & \(\fitinline{\dfrac{190423250000000000000000000000000000}{196413576346749754754208100812485763}}\) & \(\fitinline{\dfrac{500000000000000000000000000000}{250183199839442597183478478997}}\) & \(\fitinline{\dfrac{9999}{100000}}\) \\
Cbridge3 & affine & \(\fitinline{\dfrac{745505000000000500000000000000000000000000000}{775501918972004435116107964306693560032938203}}\) & \(\fitinline{\dfrac{1000000000000000000000000000000}{526467158830439879510010979401}}\) & \(\fitinline{\dfrac{9999}{100000}}\) \\
Cbridge4 & affine & \(\fitinline{\dfrac{45625312500000031250000000000000000000000000}{47901013384894191411174859411308084300058581}}\) & \(\fitinline{\dfrac{62500000000000000000000000000}{34710124696487820758100019527}}\) & \(\fitinline{\dfrac{9999}{100000}}\) \\
Cbridge5 & affine & \(\fitinline{\dfrac{368752500000000250000000000000000000000000000}{385367862672968203250537142863633598483352623}}\) & \(\fitinline{\dfrac{500000000000000000000000000000}{270427894621844727456161117541}}\) & \(\fitinline{\dfrac{9999}{100000}}\) \\
Cbridge6 & affine & \(\fitinline{\dfrac{722505000000000500000000000000000000000000000}{762245343028001044640210045205416176788051141}}\) & \(\fitinline{\dfrac{1000000000000000000000000000000}{570957463898189053255596017047}}\) & \(\fitinline{\dfrac{9999}{100000}}\) \\
Cprevious & affine & \(\fitinline{\dfrac{9442912138791}{10000000000000}}\) & \(\fitinline{\dfrac{249999854189139951396678196647534309}{145810860048603321803352465691000000}}\) & \(\fitinline{\dfrac{10099}{100000}}\) \\
C072 & affine & \(\fitinline{\dfrac{999999}{1000000}}\) & \(\fitinline{\dfrac{51}{25}}\) & \(\fitinline{\dfrac{1399}{10000}}\) \\
P072i & capped & \(\fitinline{\dfrac{999999}{1000000}}\) & \(\fitinline{\dfrac{51}{25}}\) & \(\fitinline{\dfrac{1}{25}}\) \\
P072j & capped & \(\fitinline{\dfrac{999999}{1000000}}\) & \(\fitinline{\dfrac{51}{25}}\) & \(\fitinline{\dfrac{1}{20}}\) \\
P072a & capped & \(\fitinline{\dfrac{999999}{1000000}}\) & \(\fitinline{\dfrac{51}{25}}\) & \(\fitinline{\dfrac{21}{200}}\) \\
P072b & capped & \(\fitinline{\dfrac{999999}{1000000}}\) & \(\fitinline{\dfrac{51}{25}}\) & \(\fitinline{\dfrac{11}{100}}\) \\
P072c & capped & \(\fitinline{\dfrac{999999}{1000000}}\) & \(\fitinline{\dfrac{51}{25}}\) & \(\fitinline{\dfrac{23}{200}}\) \\
C073 & affine & \(\fitinline{\dfrac{999999999999}{1000000000000}}\) & \(\fitinline{\dfrac{20000000001}{10000000000}}\) & \(\fitinline{\dfrac{1599}{10000}}\) \\
\end{longtable}
\end{landscape}

For each of the seventeen PR95 faces, the budget reserve in (D1) is strictly greater than \(\fitinline{1/100}\). For each of the six C072 faces it is exactly \(\fitinline{999999/10^{12}>0}\). For the C073 face it is zero because its \(\fitinline{B,L}\) agree with (D2). Every cap reserve is positive; the smallest is \(\fitinline{1/10000}\), for C073. The attached exact rational receipt records each individual reserve. These are direct comparisons of the displayed fractions, not fitted or sampled inequalities.

\textbf{Corollary D2.} Every enhanced source used by the 813 inherited count windows is supplied by \(\fitinline{G_u}\) at precisely its original face. Consequently the existing same-zero count proof and its 15 admission guards remain valid without dependence on the older numerical source chains. The 284 PR95, 219 C072 and 310 C073 full upward windows, their actual coefficient classes, lower witness floors, clipped cap-to-affine continuations, upper natural joins, derivative guards, finite loss coefficients and height controls remain unchanged. The old branch plus these windows still defines the same conductor-independent antitone count minimum.

The nearby strict physical source \(\fitinline{G(1-t^2,2+t^2,4/25)}\) used in the algebraic endpoint argument is a separate source choice. It supplies the physical interval and the local contact count. The fixed source \(\fitinline{G_u}\) supplies the inherited bank, independently of \(\fitinline{t}\). Thus no limit of source constants is taken and no \(\fitinline{t}\)-dependent weakening of an inherited count is required.

\hypertarget{consequences-for-dependency-closure-and-retained-generality}{%
\subsubsection{Consequences for dependency closure and retained generality}\label{consequences-for-dependency-closure-and-retained-generality}}

The retained algebraic endpoint needs the original foundations; the native global source; the complete selected natural-row norm and its original-slot/full-root half-price adapter; the generic source-step assembly; the complete open source theorem; the common-witness count theorem; the explicit count/window and return certificates; and the generic physical/complement/all-height return. It no longer needs the numerical PR95 repeated-mark source curves or the C063--C073 sequence of intermediate endpoint proofs to establish its count sources. The preliminary \(\fitinline{\beta\le7/8}\) bound already suffices for the generic feedback range.

This replacement does not declare the older methods obsolete. Their rate formulas, admissible coefficient classes, selected-count cutoffs, source-curve continuations, LOW-only bootstrap, and potential beyond the common-source cone remain useful research material. The exact older statements and their proof/evidence scope remain preserved and are described in the research guide. A failed construction of one newer source would not retroactively invalidate an earlier supplier.

\hypertarget{scope-of-the-new-finite-check}{%
\subsubsection{Scope of the new finite check}\label{scope-of-the-new-finite-check}}

\texttt{c\allowbreak{}o\allowbreak{}n\allowbreak{}s\allowbreak{}o\allowbreak{}l\allowbreak{}i\allowbreak{}d\allowbreak{}a\allowbreak{}t\allowbreak{}i\allowbreak{}o\allowbreak{}n\allowbreak{}/\allowbreak{}C\allowbreak{}O\allowbreak{}U\allowbreak{}N\allowbreak{}T\allowbreak{}\_\allowbreak{}F\allowbreak{}A\allowbreak{}C\allowbreak{}E\allowbreak{}\_\allowbreak{}D\allowbreak{}O\allowbreak{}M\allowbreak{}I\allowbreak{}N\allowbreak{}A\allowbreak{}T\allowbreak{}I\allowbreak{}O\allowbreak{}N\allowbreak{}.\allowbreak{}j\allowbreak{}s\allowbreak{}o\allowbreak{}n} binds the three immutable parameter inputs, all 24 exact rows and all exact reserves. The check uses Python \texttt{F\allowbreak{}r\allowbreak{}a\allowbreak{}c\allowbreak{}t\allowbreak{}i\allowbreak{}o\allowbreak{}n} arithmetic; it verifies (D1), cap inclusion and the common ordinary upper bound for the entire faces, so there is no grid or remaining real-variable interpolation. It does not certify the analytic source theorem or rerun the inherited count/return tower. The mathematical proof of the replacement is Lemma D1 above.

\hypertarget{the-complete-same-zero-count}{%
\subsection{7. The complete same-zero count}\label{the-complete-same-zero-count}}

The count is a statement about the original zero/amplitude fibres. Its inverse and plain witnesses belong to the same zero, character presentation and height. No zero label is transported by Poisson, and the underlying moment estimates retain their original coefficients, masks, natural and presentation zeros, all roots, independent profiles/heights, shortenings and zero-width alternatives.

\hypertarget{the-scalar-curves-and-exact-crossings}{%
\subsubsection{7.1. The scalar curves and exact crossings}\label{the-scalar-curves-and-exact-crossings}}

For \(\fitinline{0\le\delta\le3/4}\), \(\fitinline{0\le p\le1/2}\), put \(\fitinline{V=5/6-\delta>0}\). The natural branch, available on the whole closed rectangle, is

\[\fitmath{
\begin{gathered}
D_O=3-p-\frac{2p}{3\kappa},\qquad
A_O=2(1-p)\left(1-\frac p{3\kappa}\right),\\
J_O=VD_O+\delta A_O,\qquad
N_O=(1-\delta)J_O+V\delta A_O/2,\qquad R_O=N_O/J_O,
\end{gathered}
}\tag{C1}\]

where \(\fitinline{7/10\le\kappa\le3/4}\). Its denominator is positive throughout the domain. The complete natural inverse/plain argument establishes this bound, with the same actual witness and all small-width fallbacks.

Let a complete broad source face have intercept \(\fitinline{B}\), slope \(\fitinline{L>1}\) and declared count cap \(\fitinline{\bar z}\). Set

\[\fitmath{
\eta=10^{-7},\qquad z=\bar z-\eta,\qquad
s_-=3/100,\qquad m_H=173/400.
}\tag{C2}\]

The source's actual width at a plain witness length \(\fitinline{m}\) is always

\[\fitmath{
z(m)=\min\{z,(B-2m)/L\}.
}\tag{C3}\]

The two count modes are represented by one set of formulas:

\begin{longtable}[]{@{}llll@{}}
\toprule\noalign{}
Mode & \(\fitinline{D}\) & \(\fitinline{A}\) & \(\fitinline{C_0}\) \\
\midrule\noalign{}
\endhead
\bottomrule\noalign{}
\endlastfoot
Affine & \(\fitinline{3-p-4p/L}\) & \(\fitinline{2-4p/L}\) & \(\fitinline{1-2pB/L}\) \\
Capped & \(\fitinline{3-p}\) & \(\fitinline{2}\) & \(\fitinline{1-2pz}\) \\
\end{longtable}

In either mode define

\[\fitmath{
\begin{gathered}
Z_0=1-C_0,\quad U=(1-p)C_0,\quad
J=VD+\delta(1-p)A,\quad T=J+\delta U,\\
N=(1-\delta)J+V\delta U,\quad
R_N=AT+(Z_0-p)J,\quad D_J=DJ,\quad M_N=DT-R_N.
\end{gathered}
}\tag{C4}\]

The candidate count is \(\fitinline{R=N/J}\). With \(\fitinline{t_*=T/J}\), \(\fitinline{r_*=R_N/D_J}\), and \(\fitinline{m_*=M_N/D_J}\), direct substitution proves

\[\fitmath{
m_*=t_*-r_*,\qquad
1-\delta[p+(1-p)r_*]
=1-\delta A m_*-\delta Z_0
=1-\delta+V(t_*-1)=R.
}\tag{C5}\]

The middle expression is the literal affine or capped plain curve in its designated segment. The clipped continuation (C3), rather than a constant cap extended past its cone, is used outside that segment.

\hypertarget{complete-admission-and-free-nonnegative-witness-floor}{%
\subsubsection{7.2. Complete admission and free nonnegative witness floor}\label{complete-admission-and-free-nonnegative-witness-floor}}

Choose a fixed \(\fitinline{\mu\ge0}\). For an affine branch set

\[\fitmath{
 m_L=\max\{\mu+\eta,(B-Lz)/2+\eta\},
\qquad
 m_U=\min\{11/25-\eta,m_H-\eta,
                 (B-L(s_-+\eta))/2-\eta\}.
}\tag{C6a}\]

For a capped branch set

\[\fitmath{
 m_L=\mu+\eta,\qquad
 m_U=\min\{11/25-\eta,m_H-\eta,(B-Lz)/2-\eta\}.
}\tag{C6b}\]

On an upward closed window \(\fitinline{Q=[\delta_0,\delta_1]\times[p_0,1/2]}\) inside the full rectangle, require all fifteen of the following polynomials to be strictly positive:

\[\fitmath{
\begin{gathered}
D,\ A,\ C_0,\ J,\ T-J,\ (12499/10000)J-T,\\
R_N-(23/37)D_J,\ D_J-R_N,\ M_N-m_LD_J,\ m_UD_J-M_N,\\
N-(1-2\delta m_H-2\delta p s_-)J,\ UJ_p-U_pJ,\\
11/25-m_H-\eta,\ B-2m_H-L(s_-+\eta)-2\eta,\ \Xi,
\end{gathered}
}\tag{C7}\]

where \(\fitinline{\Xi=1}\) in affine mode and \(\fitinline{\Xi=(z-s_--\eta)D_J}\) in capped mode. Require also \(\fitinline{\eta<z<7/37}\), actual source-cap inclusion, and the literal source/inducing/height hypotheses. The numerical condition \(\fitinline{B>22/23}\) appearing in some earlier convenient wrappers is not part of the general source-face theorem: the full PR95 bank has faces below it. Here the exact general face and all guards (C7) are used. Lemma D1 supplies those actual broad faces without a wrapper restriction.

\textbf{Theorem C (complete count).} Under these hypotheses, the original detector count on \(\fitinline{Q}\) is bounded with exponent \(\fitinline{R=N/J}\), arbitrary requested small loss and the original finite polynomial height factor. The lower witness floor may be any fixed \(\fitinline{\mu\ge0}\). It does not change the physical bookkeeping floor \(\fitinline{2/5}\).

\textbf{Proof.} Set \(\fitinline{m_{\rm Lo}=m_L-\eta}\) and \(\fitinline{m_{\rm Hi}=m_H}\). These satisfy \(\fitinline{m_{\rm Lo}\ge0}\) and \(\fitinline{2/5\le m_{\rm Hi}<1/2}\). Positive denominators and the cutoff/inverse guards give

\[\fitmath{
1<t_*<12499/10000<3/2,\qquad 23/37<r_*<1,
\qquad m_L<m_*<m_U.
}\tag{C8}\]

For \(\fitinline{0\le\epsilon\le\eta}\), therefore \(\fitinline{m_{\rm Lo}+\epsilon\le m_L<m_*}\). The actual same-zero witness relation is

\[\fitmath{
t_*-r-\epsilon\le m,\qquad r\le t_*+\epsilon.
}\tag{C9}\]

Before splitting the witnesses, check that the marked widths are actual whole-slot selections in the original fibre. At row scale \(\fitinline{U=Z^d}\), write the physical prime widths as \(\fitinline{z_i=\ell_i/d}\), their total as \(\fitinline{L_{\rm supply}=\ell/d}\), and their binned squared amplitudes as

\[\fitmath{
 |Q_{u,i}|^2\ge U^{2z_i g_i-o(1)},\qquad
 0\le g_i\le\delta/2,\qquad
 \sum_i z_i g_i=\delta pL_{\rm supply}.
}\tag{C9a}\]

Arrange the common cell labels \(\fitinline{g_i}\) in decreasing order. The average on any initial subinterval of the resulting step function is at least its average on the whole interval. For each desired \(\fitinline{0\le s\le L_{\rm supply}}\), take this initial interval of length \(\fitinline{s}\) and discard its last fractional slot. The resulting genuine subset \(\fitinline{I}\) has total width in \(\fitinline{[s-\mathrm{mesh},s]}\) and

\[\fitmath{
 \prod_{i\in I}|Q_{u,i}|^2
 \ge U^{2\delta p s-\delta\,\mathrm{mesh}-o(1)}.
}\tag{C9b}\]

The selection depends only on the fixed amplitude cell and widths; it precedes the ordinary summation variables. Whole-polynomial conjugation and multiplication by the actual fixed character put these same polynomials in the witness orientation without changing their moduli. Natural and presentation zeros, the fixed inducing exclusion, punctures and common masks remain in the resulting source.

The return geometry supplies \(\fitinline{\ell/d>1/5}\) throughout the conductor range. For every \(\fitinline{r\in[r_*,1]}\), the desired inverse width is less than

\[\fitmath{
 (1-r_*)/2<(1-23/37)/2=7/37<1/5.
}\tag{C9c}\]

If \(\fitinline{(1-r)/2>\nu>0}\), choose \(\fitinline{s=(1-r)/2-\nu}\). Both literal inverse gates hold, including after discarding the final fractional slot:

\[\fitmath{
 r+2s=1-2\nu<1,\qquad
 2r+8s=4-2r-8\nu<3,
}\tag{C9d}\]

because \(\fitinline{r>23/37>1/2}\). If \(\fitinline{(1-r)/2\le\nu}\), select no mark and use the unmarked inverse estimate; the omitted marked gain is at most \(\fitinline{2\delta p\nu}\). In particular \(\fitinline{r=1}\) is covered without pretending a positive marked width exists.

For the plain witness use \(\fitinline{s=z(m)-\nu}\), if positive. Its supply follows from \(\fitinline{z(m)\le z<7/37<1/5}\). The source face is applicable to two copies of the actual ordinary polynomial at length \(\fitinline{U^m}\): \(\fitinline{2m+Ls\le B-L\nu<B}\), and \(\fitinline{s\le z<\bar z}\). The guard at \(\fitinline{m_H}\) in (C7) makes \(\fitinline{z(m)>\eta}\) for all \(\fitinline{m\le m_H}\), so this width is positive whenever \(\fitinline{\nu\le\eta}\). The natural plain continuation uses \(\fitinline{s=(1-2m)/(6\kappa_{\rm plain})-\nu}\) when positive; for \(\fitinline{m\ge m_H}\) its undecremented value is at most \(\fitinline{(1-2m_H)/(6\cdot7/10)<1/5}\). Otherwise its unmarked fallback is available.

The same-zero witness lemma gives \(\fitinline{|M_u(U^r)|^2\ge U^{\delta r-O(\epsilon)}}\) and \(\fitinline{|S_u(U^m)|^4\ge U^{2\delta m-O(\epsilon)}}\). Combining these lower bounds with (C9b) and the corresponding positive moment sums, then dividing, gives the actual cardinality exponents

\[\fitmath{
 \begin{aligned}
 I(r)&=1-\delta[p+(1-p)r],\\
 P(m)&=1-2\delta m-2\delta pz(m).
 \end{aligned}
}\tag{C9e}\]

The errors are precisely the witness, moment, amplitude-cell, final-slot and \(\fitinline{\nu}\) decrements kept in (C11). For \(\fitinline{r\ge1}\), the full unmarked inverse moment has exponent \(\fitinline{(1+5r)/6}\). Dividing by its same-zero spike gives

\[\fitmath{
 1-\delta+(5/6-\delta)(r-1)
 \le1-\delta+V(t_*-1)+O(\epsilon),
}\tag{C9f}\]

using \(\fitinline{r\le t_*+\epsilon}\). Thus these are counts of the original fibre with its actual same-zero witnesses, not independent optimizations over unrelated events.

Every witness belongs to one of the following exhaustive cases.

If \(\fitinline{r\ge1}\), the complete long unmarked inverse moment gives exponent \(\fitinline{1-\delta+V(t_*-1)+O(\epsilon)=R+O(\epsilon)}\). If \(\fitinline{r_*\le r<1}\), the marked short-inverse curve \(\fitinline{1-\delta[p+(1-p)r]}\) decreases in \(\fitinline{r}\) and is at most its crossing value \(\fitinline{R}\). When the requested marked width is at most the small-width parameter \(\fitinline{\nu}\), its original unmarked inverse fallback supplies the same estimate with the allowed loss; otherwise use the marked inverse estimate. Width zero is included in that fallback.

If \(\fitinline{r<r_*}\), (C9) implies \(\fitinline{m>m_*-\epsilon\ge m_{\rm Lo}}\). For \(\fitinline{m\le m_H}\), the broad source and (C3) apply. Its two plain-curve slopes are \(\fitinline{-2\delta}\) and \(\fitinline{-2\delta(1-2p/L)}\), both nonpositive because \(\fitinline{L>1}\). Both have magnitude at most \(\fitinline{2}\), hence the clipped curve satisfies

\[\fitmath{
P(m-g)\le P(m)+2g\qquad(g\ge0).
}\tag{C10}\]

Together with (C5), this bounds the plain witness by \(\fitinline{R}\) plus the already budgeted witness displacement. There is no gap where the cap branch passes to the affine segment, even if one segment is empty. The source-width guard in (C7) implies \(\fitinline{(B-2m)/L>\eta}\) for every \(\fitinline{m\le m_H}\), so all \(\fitinline{\nu\le\eta}\) are below the clipped width over the enlarged witness interval.

For \(\fitinline{m_H<m<1/2}\), use the natural width \(\fitinline{(1-2m)/(6\kappa_{\rm plain})}\), with \(\fitinline{7/10\le\kappa_{\rm plain}\le3/4}\). Its value at \(\fitinline{m_H}\) is at least \(\fitinline{s_-}\). The natural-tail polynomial in (C7) pays the count at that join. Thereafter the natural plain curve has derivative \(\fitinline{-2\delta(1-p/(3\kappa_{\rm plain}))\le0}\). If its width is at most \(\fitinline{\nu}\), use its unmarked plain fallback. At \(\fitinline{m\ge1/2}\), the same unmarked moment pays the endpoint and the original bounded witness extension. Thus no inverse, plain, natural-tail or zero-width alternative is lost.

The same-zero and height bookkeeping, together with the displayed moment-to-count conversion, retains the exact aggregate loss

\[\fitmath{
159\epsilon+\epsilon_{\rm moment}+\mathrm{mesh}+7\nu.
}\tag{C11}\]

Its fibre constant and one finite polynomial height degree are unchanged. The early mesh precedes the finite slot type and row labels; the two ordinary profiles and real heights remain independent. The value of \(\fitinline{\mu}\) enters only the admissible interval, so changing it cannot change those loss coefficients. This proves the count.

Finally

\[\fitmath{
\partial_pR=-V\delta\,(UJ_p-U_pJ)/J^2\le0.
}\tag{C12}\]

The old branch is also antitone. Since each enhanced window is upward in \(\fitinline{p}\), its bound remains available as \(\fitinline{p}\) increases, and the finite minimum of all available bounds is antitone. It is independent of the conductor exponent. \ensuremath{\blacksquare}

\hypertarget{the-one-maintained-finite-bank}{%
\subsubsection{7.3. The one maintained finite bank}\label{the-one-maintained-finite-bank}}

Use the exact 284 PR95 windows in \texttt{r\allowbreak{}e\allowbreak{}s\allowbreak{}e\allowbreak{}a\allowbreak{}r\allowbreak{}c\allowbreak{}h\allowbreak{}/\allowbreak{}c\allowbreak{}y\allowbreak{}c\allowbreak{}l\allowbreak{}e\allowbreak{}0\allowbreak{}6\allowbreak{}7\allowbreak{}/\allowbreak{}r\allowbreak{}e\allowbreak{}t\allowbreak{}u\allowbreak{}r\allowbreak{}n\allowbreak{}/\allowbreak{}P\allowbreak{}R\allowbreak{}9\allowbreak{}5\allowbreak{}\_\allowbreak{}R\allowbreak{}E\allowbreak{}T\allowbreak{}U\allowbreak{}R\allowbreak{}N\allowbreak{}\_\allowbreak{}D\allowbreak{}A\allowbreak{}T\allowbreak{}A\allowbreak{}.\allowbreak{}j\allowbreak{}s\allowbreak{}o\allowbreak{}n}, the 219 added windows in \texttt{r\allowbreak{}e\allowbreak{}s\allowbreak{}e\allowbreak{}a\allowbreak{}r\allowbreak{}c\allowbreak{}h\allowbreak{}/\allowbreak{}c\allowbreak{}y\allowbreak{}c\allowbreak{}l\allowbreak{}e\allowbreak{}0\allowbreak{}7\allowbreak{}2\allowbreak{}/\allowbreak{}r\allowbreak{}e\allowbreak{}t\allowbreak{}u\allowbreak{}r\allowbreak{}n\allowbreak{}/\allowbreak{}R\allowbreak{}E\allowbreak{}T\allowbreak{}U\allowbreak{}R\allowbreak{}N\allowbreak{}\_\allowbreak{}C\allowbreak{}E\allowbreak{}R\allowbreak{}T\allowbreak{}I\allowbreak{}F\allowbreak{}I\allowbreak{}C\allowbreak{}A\allowbreak{}T\allowbreak{}E\allowbreak{}.\allowbreak{}j\allowbreak{}s\allowbreak{}o\allowbreak{}n}, and the 310 added windows in \texttt{r\allowbreak{}e\allowbreak{}s\allowbreak{}e\allowbreak{}a\allowbreak{}r\allowbreak{}c\allowbreak{}h\allowbreak{}/\allowbreak{}c\allowbreak{}y\allowbreak{}c\allowbreak{}l\allowbreak{}e\allowbreak{}0\allowbreak{}7\allowbreak{}3\allowbreak{}/\allowbreak{}r\allowbreak{}e\allowbreak{}t\allowbreak{}u\allowbreak{}r\allowbreak{}n\allowbreak{}/\allowbreak{}R\allowbreak{}E\allowbreak{}T\allowbreak{}U\allowbreak{}R\allowbreak{}N\allowbreak{}\_\allowbreak{}C\allowbreak{}E\allowbreak{}R\allowbreak{}T\allowbreak{}I\allowbreak{}F\allowbreak{}I\allowbreak{}C\allowbreak{}A\allowbreak{}T\allowbreak{}E\allowbreak{}.\allowbreak{}j\allowbreak{}s\allowbreak{}o\allowbreak{}n}. Their source parameters are exactly the 24 rows in Lemma D1's table. The first 23 branches use \(\fitinline{\mu=377/1000}\); C073 uses \(\fitinline{\mu=3/10}\). Every window is upward, all of (C7) are verified on its complete closed rectangle, and all previous stated margins remain unchanged. Their original receipts bind these exact immutable data; no arithmetic receipt is being presented as the proof of the underlying analytic moments.

The active count \(\fitinline{R_*}\) is the minimum of (C1) and every available enhanced branch. Its definition is given once by these formulas and the finite window data. The full bank, including branches not selected by the final return cover, remains reusable. This gives one maintained parameter bank rather than a second narrative catalogue.

\hypertarget{natural-count-monotonicity-and-supremum-feedback}{%
\subsubsection{7.4 Natural-count monotonicity and supremum feedback}\label{natural-count-monotonicity-and-supremum-feedback}}

This calculation uses the complete natural count already obtained from the same-zero inverse/plain moments. It changes no source, amplitude cell or admission window. Throughout

\[\fitmath{
0\le\delta\le3/4,\qquad0\le p\le1/2,\qquad
7/10\le a\le\kappa\le3/4,\qquad V=5/6-\delta>0.
}\]

Write

\[\fitmath{
D_\kappa=3\kappa(3-p)-2p,\quad
C_\kappa=\frac{2(1-p)(3\kappa-p)}{D_\kappa},\quad
S_\kappa=V+\delta C_\kappa.
}\tag{FB1}\]

The denominators in the natural count are positive on the entire closed rectangle, including its boundary. Indeed

\[\fitmath{
 D_\kappa\ge3(7/10)(5/2)-1=17/4>0,\qquad
 1-p\ge1/2,\qquad3\kappa-p\ge8/5,
}\]

so \(\fitinline{C_\kappa>0}\) and \(\fitinline{S_\kappa\ge V\ge1/12>0}\). In the notation of (C1),

\[\fitmath{
 D_O=D_\kappa/(3\kappa)>0,\qquad
 A_O=D_OC_\kappa>0,\qquad J_O=D_OS_\kappa>0.
}\tag{FB1a}\]

Thus the natural count (C1), after division of numerator and denominator by \(\fitinline{D_O}\), is exactly

\[\fitmath{
R_O(\kappa)=1-\delta+
 \frac{V\delta C_\kappa}{2S_\kappa}.
}\tag{FB2}\]

Its antitonicity in the amplitude parameter \(\fitinline{p}\) follows from an explicit sign calculation. Put

\[\fitmath{
 \begin{aligned}
 H_\kappa(p)
  &=(3\kappa+2)p^2-18\kappa p+18\kappa^2+3\kappa\\
  &=(3\kappa+2)p^2+18\kappa(1/2-p)
                       +6\kappa(3\kappa-1)>0.
 \end{aligned}
}\tag{FB2a}\]

Every term in the last line is nonnegative, and its last term is strictly positive because \(\fitinline{\kappa\ge7/10>1/3}\). Differentiating the displayed rational functions, with \(\fitinline{\delta,\kappa}\) fixed, gives

\[\fitmath{
 \partial_p C_\kappa=-\frac{2H_\kappa(p)}{D_\kappa^2}<0,
 \qquad
 \partial_p R_O
 =\frac{\delta V^2}{2S_\kappa^2}\,\partial_p C_\kappa\le0.
}\tag{FB2b}\]

The displayed derivative of \(\fitinline{R_O}\) is zero when \(\fitinline{\delta=0}\) and strictly negative when \(\fitinline{\delta>0}\). Hence the natural branch is available everywhere and is antitone on the full closed amplitude interval. This establishes the positivity and monotonicity used by the complete count without an admission restriction.

For the feedback in \(\fitinline{\kappa}\), subtracting the two rational functions gives the exact identity

\[\fitmath{
C_\kappa-C_a=
 \frac{6p(1-p)^2(\kappa-a)}{D_\kappa D_a}.
}\tag{FB3}\]

On \(\fitinline{[0,1/2]}\), the maximum of \(\fitinline{p(1-p)^2}\) is \(\fitinline{4/27}\): its derivative is \(\fitinline{(1-p)(1-3p)}\), so the only interior maximum is \(\fitinline{p=1/3}\). Consequently

\[\fitmath{
0\le C_\kappa-C_a\le
 \frac{128}{2601}(\kappa-a)<\frac{\kappa-a}{18}.
}\tag{FB4}\]

A second exact subtraction, using \(\fitinline{S_\kappa=V+\delta C_\kappa}\), yields

\[\fitmath{
R_O(\kappa)-R_O(a)
=\frac{\delta V^2}{2S_\kappa S_a}(C_\kappa-C_a)
\le\frac5{216}(\kappa-a).
}\tag{FB5}\]

Indeed \(\fitinline{V^2/(S_\kappa S_a)\le1}\) and \(\fitinline{\delta/2\le3/8<5/12}\), so (FB4) gives the displayed inherited constant (and even a slightly smaller scalar constant on this rectangle). The proof uses only the stated rectangle; it assumes no intermediate endpoint ceiling.

All enhanced count candidates have their actual source and upward admission windows fixed independently of \(\fitinline{\kappa}\). Let \(\fitinline{R_*(\kappa)}\) be the minimum of the natural candidate and those available enhanced candidates at the fixed point \(\fitinline{(\delta,p)}\). Choose a candidate attaining this finite minimum at \(\fitinline{a}\). If it is enhanced, its value does not change with \(\fitinline{\kappa}\); if it is natural, (FB5) applies. Since the minimum at \(\fitinline{\kappa}\) is no greater than that same candidate's value,

\[\fitmath{
R_*(\kappa)-R_*(a)\le\frac5{216}(\kappa-a).
}\tag{FB6}\]

The new local affine count used in the algebraic endpoint construction is also fixed independently of the unknown actual supremum, so the same argument includes it. No differentiation of a changing minimizing branch is required.

In the physical return functional, the count is multiplied by \(\fitinline{3d/4}\). Its feedback cost is therefore at most

\[\fitmath{
 \frac{3d}{4}\frac5{216}(\kappa-a)
 =\frac{5d}{288}(\kappa-a)
 <\frac1{50}(\kappa-a)\qquad(0\le d<1).
}\tag{FB7}\]

The strict old-source capacity cost is separate. Its exact change is

\[\fitmath{
 \frac{\delta p\min(d,2y-d)}{12}
      (a^{-1}-\kappa^{-1})
 \le\frac{25}{784}(\kappa-a),
}\tag{FB8}\]

because \(\fitinline{\delta p\le3/8}\), \(\fitinline{\min(d,2y-d)\le y<1/2}\), and \(\fitinline{a\kappa\ge(7/10)^2}\). Thus the total old-row feedback coefficient is smaller than

\[\fitmath{
 \frac1{50}+\frac{25}{784}<\frac1{10}.
}\tag{FB9}\]

Enhanced physical capacities are fixed and pay only (FB7). These are exactly the bounds needed by the all-height return; they do not assert a new endpoint or boundary conclusion.

\hypertarget{the-complete-physical-return}{%
\subsection{8. The complete physical return}\label{the-complete-physical-return}}

Fix theta with 87/100\textless=theta\textless7/8 and put

\[\fitmath{
 \ell=11/3-4\theta,\qquad x=1+\ell-h,\qquad
 y=h-2\ell,\qquad k=(4+y-x)/6,\qquad\kappa_*=2\theta-1.
}\tag{GT1}\]

Thus x+y=1-ell and

\[\fitmath{
 \theta-k=x/2+(y-x)/12.
}\tag{GT2}\]

Retain the conservative scalar lower edge d\_0=1/200. Choose a fixed zeta\textgreater0 and set d\_+=h+zeta. The actual small-row/central-band split remains the original one at Z\^{}(1/100); its dyadic scalar extension uses d\_0. Require

\[\fitmath{
 0<\ell<x\le y<1/2,\quad \ell<1/5,\quad 2x-y-\ell>0,
 \quad d_0<y,\quad h>0,\quad d_+<\min(1,2y).
}\tag{GT3}\]

These imply 2(x+y)-1-3ell=1-5ell\textgreater0, h\textless2y, and 4ell\textless h. They do not require ell\textless.171, h\textless.806, y\textgreater.458, or any other old decimal box endpoint.

Choose a finite partition of {[}d\_0,d\_+{]}. On each closed partition interval use either the old natural physical capacity or one specified \textbf{already proved} complete physical source G(B\_j,L\_j,c\_j). A source interval {[}a\_j,b\_j{]} must have a\_j\textgreater y. For that source independently choose an actual physical marked cap sigma\_j and a declared physical cap barsigma\_j with

\[\fitmath{
 0<\sigma_j<\bar\sigma_j<c_j,\qquad
 \bar\sigma_j<1/5<\ell/d_+,
}\tag{GT4}\] \[\fitmath{
 0<B_j\le1,\qquad 2\le L_j\le9/2,\qquad
 B_j/2+(9/2-L_j/2)\bar\sigma_j<1,
}\tag{GT5}\] \[\fitmath{
 2y-(2-B_j)b_j>0.
}\tag{GT6}\]

The strict barsigma\_j buffer pays whole-slot rounding and permits a common mesh below barsigma\_j-sigma\_j. The fallback guard is stated at barsigma\_j, which is stronger than necessary for an exact selection of width at most sigma\_j. Neither physical cap is a cap of a counting window.

The capacities on the assigned intervals are

\[\fitmath{
 C_O(\kappa,d)=\frac{\min(d,2y-d)}{6\kappa},\qquad
 C_j(d)=\min\left(\sigma_jd,
                    \frac{2y-(2-B_j)d}{L_j}\right).
}\tag{GT7}\]

Retain a single common, conductor-independent count envelope R\_*(kappa,delta,p): the minimum of the natural branch (C1) and finitely many complete enhanced branches (C4), each only on its own full admission windows (C7). Its windows are upward in p and its derivative guard makes it antitone. Every count source, declared count cap, witness endpoint and admission buffer is the actual one used in that branch. The entire closed domain is 0\textless=delta\textless=3/4, 0\textless=p\textless=1/2. The old branch is available everywhere.

For some fixed r\textgreater r\_0\textgreater0, require on the entire domain at \textbf{both endpoints of every assigned conductor interval}

\[\fitmath{
 E_{\kappa_*}(d,C;\delta,p)\le\theta-r,
}\tag{GT8}\] where the literal complete exponent is \[\fitmath{
 E_\kappa=
 \frac12+\frac{13h}{75}-\frac y2-\frac\ell2-\frac{9d}{100}
 +\frac\delta2+\delta p(\ell-C/2)+\frac{3d}{4}R_*.
}\tag{GT9}\]

Each endpoint may be proved by a finite closed rectangle cover, choosing any admitted branch on each rectangle and clearing its positive denominator as in the finite-certification section. Such different choices still bound the same R\_*.

Finally define the actual complementary margins

\[\fitmath{
 \Delta_f=\theta-k-F_{\rm fl}(d_+)>0,\qquad
 F_{\rm fl}(d)=4x/25-3/20-51y/100+67d/100-3\ell/20,
}\tag{GT10}\] \[\fitmath{
 \Delta_s=S_{\rm small}-(7/8-\theta)>0,\qquad
 S_{\rm small}=y/2-13h/75-1/50,
}\tag{GT11}\] and, retaining the conservative principal contours and error allowances, require \[\fitmath{
 g_w=y/20-(1+h)/1000-1/4000>0,\qquad
 g_z=h/600-1/1000-1/4000>0.
}\tag{GT12}\]

\textbf{Generic conclusion.} Under the complete source/count hypotheses and (GT1)--(GT12), every ordinary character at every fixed positive presented modulus is zero-free for Re(s)\textgreater theta at all heights, excluding the actual principal pole. No extremal zero or attained supremum is assumed.

The analytic proof follows. Its only numerical obligation left to an application is to certify its chosen source/count data and the stated finite inequalities. The new theorem includes every complementary sector.

\hypertarget{independent-caps-every-shift-and-shortening}{%
\subsubsection{8.1 Independent caps, every shift and shortening}\label{independent-caps-every-shift-and-shortening}}

We prove the physical restriction with its independently chosen cap sigma\_j. The proof uses the full source (S5), not a numerical wrapper with a fixed counting cap.

Fix one source interval and suppress j. Put m=y/d.~By a\_j\textgreater y and b\_j\textless2y, 1/2\textless m\textless1. For t\textgreater=0 define

\[\fitmath{
 n(t)=1-m+t,\quad A(t)=(B-2n(t))/L,\quad
 C(t)=\min(\sigma,A(t)),\quad Z(t)=\max(0,C(t)).
}\tag{GT13}\]

If a nonempty selected marked family has total width 0\textless z\textless=Z(t), then z\textless=sigma\textless barsigma\textless c\_j and 2n(t)+Lz\textless=B. Also n(t)\textgreater=0 and n(t)\textless=B/2\textless=1/2. Every nonnegative shortened pair n\_1,n\_2\textless=n(t), and every remaining subset of width z'\textless=z, therefore lies in the same literal G source. This is the complete global source restricted to its actual ordinary and marked domain; no lower marked-width window and no upper shift window is needed.

For a nonpositive ordinary exponent retain the original bounded-scale treatment. Its natural budget has kappa\_nat=3/4, supplied by the proved 7/8 bound, and hence price 6*kappa\_nat=9/2. If n\_1\textless=0, n\_2\textless=n(t), then

\[\fitmath{
 n_1+n_2+(9/2)z'
 \le B/2+(9/2-L/2)z\le1.
}\tag{GT14}\]

This pays both orientations by (GT5). The target kappa\_* is not substituted for this natural source parameter.

If Z(t)=0 there is exactly no positive-width nonempty selection. Use the retained empty-slot reflection/clipping branch. Its actual coefficient bins retain 0\textless=n(t)+ann\textless=2 and X\_i\textless=U\textsuperscript{(n(t)+ann)\textless=U}2, supplied by the original bin condition 2xi+ann\textless=1/2. No moment for arbitrary unbounded ordinary lengths is inferred from an empty product.

The clipping maps are monotone and 1-Lipschitz. With v=2t,

\[\fitmath{
 C(0)-Z(v/2)\le v/L.
}\] For 0\textless=a\textless=delta/2 and every v\textgreater=0, \[\fitmath{
 \frac a2 C(0)+\frac{\delta v}{1000}
 \le \frac{\delta v}{4}+\frac a2 Z(v/2).
}\tag{GT15}\] Indeed the capacity loss is at most delta\emph{v/(4L)\textless=delta}v/8, and 1/8+1/1000\textless1/4. All shifted tails keep this damping. The source calls at positive capacity, their natural fallbacks, and the empty-selection branch are precisely those just proved; (GT15) is not used to assert a moment outside those domains.

The actual prime supply is unchanged: sigma\textless barsigma\textless1/5\textless ell/d\_+\textless=ell/d. The physical slot at Re(z)=17/50 is exactly the original inverse-character annular coefficient of real power 33/50, with conjugated profile and opposite imaginary part. This equality holds on every ideal, including every natural zero. It is substituted into the finite pair before taking positive energy. The original relative characters, fixed modulus and Euler factors stay fixed.

The physical free rows satisfy U\^{}(1/100)\textless=N(u)\textless=U. After fixing the source and its inducing data, the fixed-inducing exceptions have bounded norm and are absent for sufficiently large U. Hence the original physical keep implies the required source initial keep on these rows. The threshold is allowed to depend on those fixed data. This is an eventual predicate inclusion, not an identification of different keeps.

The two ordinary factors and both Fourier heights remain independent. The same nonnegative radial majorant bounds the finite physical energy. The source provides its original seminorms, polynomial height factors and common profile controls. A physical source mesh is taken separately from every count-source mesh, and their finite minimum is chosen before the slot type. If original widths are at most mesh/200 and d\textgreater=1/200, the rescaled widths are at most mesh; (Z\textsuperscript{d)}(w/d)=Z\^{}w preserves the same prime pools exactly.

Consequently (GT7)'s new capacity is the full physical consumer, for every actual shift, shortening and remaining subset. The declared cap 31/500 of the retained capP \textbf{count} branch is absent from this proof. Each new application separately chooses a physical cap below its available global source cap; the counting cap and all counting windows remain their literal independent inputs. The algebraic application below makes its physical choices separately.

\hypertarget{edge-reduction-and-actual-supremum-feedback}{%
\subsubsection{8.2 Edge reduction and actual-supremum feedback}\label{edge-reduction-and-actual-supremum-feedback}}

For fixed delta,p, the count envelope is independent of d.~Each capacity in (GT7) is a minimum of two affine functions of d.~Because delta\emph{p\textgreater=0, -delta}p*C/2 is convex in d; all other d terms in (GT9) are affine. Therefore (GT8) bounds the whole assigned conductor interval. Different intervals may use different already-proved physical suppliers. No equality of capacities at a join is needed, and both estimates remain valid at any shared endpoint. The finite partition covers every conductor exponent from d\_0 to d\_+.

Suppose the actual retained primitive-family zero supremum beta exceeds theta. Set mu=beta-theta\textgreater0 and kappa\_beta=2beta-1. The preliminary bound gives kappa\_\emph{\textless=kappa\_beta\textless=3/4, while theta\textgreater=.87 implies kappa\_}\textgreater=.74\textgreater7/10. All enhanced sources, windows, physical intervals and capacities are fixed independently of kappa\_beta. Only the old count and natural capacity vary.

The natural-count feedback calculation in the count section, followed by the finite minimum, gives \[\fitmath{
 R_*(\kappa_\beta)-R_*(\kappa_*)
 \le\frac5{216}(\kappa_\beta-\kappa_*).
}\tag{GT16}\] Its count cost in (GT9) is at most 5d\_+(kappa\_beta-kappa\_\emph{)/288 \textless{} (kappa\_beta-kappa\_})/50, since d\_+\textless1. Every new physical capacity is fixed. On old intervals the additional capacity cost is at most \[\fitmath{
 \frac{\delta p\min(d,2y-d)}{12}
 (\kappa_*^{-1}-\kappa_\beta^{-1})
 \le\frac{25}{784}(\kappa_\beta-\kappa_*),
}\tag{GT17}\] using y\textless1/2, delta\emph{p\textless=3/8 and kappa\_beta}kappa\_*\textgreater=(7/10)\^{}2. Because 5/288+25/784\textless1/8, the retained reserves are \[\fitmath{
 E_{\kappa_\beta}-\beta\le-r_0-24\mu/25
 \quad\hbox{on new intervals},
}\] \[\fitmath{
 E_{\kappa_\beta}-\beta\le-r_0-3\mu/4
 \quad\hbox{on old intervals}.
}\tag{GT18}\]

Thus a finite physical partition and additional independent count branches do not alter the feedback. No actual zero is assumed to attain beta.

\hypertarget{every-complementary-term}{%
\subsubsection{8.3 Every complementary term}\label{every-complementary-term}}

\hypertarget{floor}{%
\paragraph{Floor}\label{floor}}

Equation (234) is a complete positive floor-row bound, including cube/ray subfamilies, natural zeros and the original per-slot prime bound. Its coefficient of d is 67/100\textgreater0, so d\_+ bounds every floor row. Its subsequent loss is exactly bounded by \[\fitmath{
 26e+(N+8)\epsilon_{\rm tr}+t+\rho\ell+\nu+\sigma_0.
}\tag{GT19}\] The old proof chose this below 1/200. That number is an error allowance, not an arithmetic term of F\_fl.~Here choose (GT19) below Delta\_f/2; (GT10) then pays the whole floor. No floor summand is omitted.

Using (GT1) and d\_+=h+zeta, direct algebra gives \[\fitmath{
 F_{\rm fl}(d_+)=1/100+103\ell/100+67\zeta/100,
}\] \[\fitmath{
 \Delta_f=61/150-h/3-39\ell/50-67\zeta/100.
}\tag{GT20}\] The intrinsic floor condition is consequently \[\fitmath{
 h<61/50-117\ell/50-201\zeta/100.
}\tag{GT21}\] The old guard subtracting 1/200 is stronger by exactly that chosen allowance. It is not retained as a mathematical obstruction.

\hypertarget{small-rows-and-nonprincipal-units}{%
\paragraph{Small rows and nonprincipal units}\label{small-rows-and-nonprincipal-units}}

The original small-row estimate and calibrated dyadic payment retain the complete exponent \[\fitmath{
 7/8-k-S_{\rm small}+8e+\nu.
}\tag{GT22}\] The constant 1/50 in S\_small pays the original small-tail bookkeeping, including the actual Z\^{}(1/100) cutoff; it is retained. All five nonprincipal units remain in this estimate. By (GT11), imposing sigma\_0+8e+nu\textless Delta\_s/2 makes (GT22) strictly below beta-k-sigma\_0. There is no need for the extra coarse pair S\_small\textgreater7/100 and theta\textgreater.873. The original contour estimate (235) is unchanged.

\hypertarget{large-rows}{%
\paragraph{Large rows}\label{large-rows}}

For N(u)\textgreater Z\^{}(h+zeta), the absolute contour (236) gives \[\fitmath{
 O\!\left(Z^{x/2+1+y+(h+\zeta)(1+\varepsilon)-\zeta v+\varepsilon}\right).
}\tag{GT23}\] After fixed zeta\textgreater0 choose v\textgreater2 sufficiently large for arbitrary saving, including dyadic summation and its fixed polynomial factors. The central estimate extends to d\_+=h+zeta, so no interval is lost at the transition.

\hypertarget{principal-contours-and-finite-slot-correction}{%
\paragraph{Principal contours and finite-slot correction}\label{principal-contours-and-finite-slot-correction}}

Keep the original exact two-pole calculation, with the same w=1 and z=1/6 residues. The conservative contour choice w=19/20, z=33/200 lies inside Re(s)\textgreater=theta\textgreater=.87, Re(w)\textgreater=.9, Re(z)\textgreater=.16. The original local defect bounds remain summable there. Their proof uses theta\textgreater=.87 and these contour coordinates, not any bound ell\textless.171. A fixed deletion cutoff supplies the same holomorphic box.

The scale saves y/20 and h/600 from the double residue. For e\textless=1/1000, (GT12) leaves the retained 1/4000 reserve. After dividing by the original nonzero logarithmic prime normalizer, the remainder powers are \[\fitmath{
 Z^{\beta-k-1/4000+\nu},\qquad
 Z^{\beta-k+e-\theta\ell_{\min}}.
}\tag{GT24}\] Choose sigma\_0+nu\textless1/4000 and sigma\_0+e\textless theta*ell\_min. The principal marked ratio is used only on the residue locus; the quotient-free product is kept on other contours. The finite Z-dependent correction (202) remains in this paid error and is not inserted into the fixed Euler numerator.

The retained absolute residual power Z\^{}(-1+nu) is also harmless: beta-k\textgreater= theta-k=x/2+(y-x)/12\textgreater0, and the small positive nu,sigma\_0 can be chosen with nu+sigma\_0\textless1. Fixed-inducing terms, zero frequencies, units, finite-slot corrections, signed compensators and original profile tails retain their literal suppliers. No count gain replaces one of these sector payments.

\hypertarget{one-common-choice-of-mesh-losses-and-heights}{%
\subsubsection{8.4 One common choice of mesh, losses and heights}\label{one-common-choice-of-mesh-losses-and-heights}}

Choose a\textgreater0 below 1/1000 times the minimum of 1, Delta\_f, Delta\_s, g\_w, g\_z, every strict geometric/source-cap margin, and all finite source/admission/ row reserves. Under mu\textgreater0 put \[\fitmath{
 b_0=\min(a,\mu/10000,r_0/8,10^{-8}),\qquad
 \epsilon_P=\min(a,\mu/10000).
}\tag{GT25}\]

For every physical source obtain its mesh at epsilon\_P, and for every count error t\textgreater0 obtain the finite common count mesh at t. Include all natural, geometric, declared-cap and shortening meshes. Their finite positive minimum with t and a is m(t). Define the actual mesh function rho(t)=min(commonMesh(t),m(t)) \textbf{before} constructing any finite slots. The retained joint construction gives t,epsilon\_D and one finite slot system with \[\fitmath{
 159\epsilon_D+609t\le b_0,\qquad
 0<\ell_i\le\rho(t)/200,\qquad\sum_i\ell_i=\ell.
}\tag{GT26}\] Set rho=rho(t)\textless=t. The slots already satisfy the actual mesh; there is no post-slot shrinking of the mesh. Only now fix N, ell\_min, phases, and choose \[\fitmath{
 e\le\min(e_D,t/10000,\phi/(32t),\theta\ell_{\min}/4),
 \quad\phi=\ell_{\min}\rho/2,
}\] \[\fitmath{
 \epsilon_{\rm tr}=t/[100(N+8)],\qquad\nu=t/8,\qquad
 0<\sigma_0\le\min(\theta\ell_{\min}/4,t/8).
}\tag{GT27}\]

The fine-partition step uses only ell\textgreater0 and rho(t)\textgreater0. After t and its actual mesh are fixed, take N sufficiently large and choose positive widths close to ell/N, summing to ell and each below rho(t)/200. The retained finite distinctness conditions are achieved by a small perturbation inside that open simplex; their excluded equalities are finitely many proper hyperplanes. All restrictions involving N or ell\_min are imposed afterwards in (GT27). Consequently this joint construction does not import an old upper decimal bound on ell, h or y through its geometric mesh.

Because ell\textless1/5, the full floor loss (GT19) is \textless3t\textless Delta\_f/2. Also sigma\_0+8e+nu\textless t\textless Delta\_s/2. The existing phase inequalities hold, e\textless1/1000, sigma\_0+nu\textless1/4000, and sigma\_0+e\textless theta*ell\_min. Thus all newly used intrinsic margins are compatible with the same joint construction.

The complete physical loss vector is \[\fitmath{
 L_f=3t,\ L_\omega=17t/25,\ L_{\rm slack}=t,\
 L_s=2\epsilon_{\rm tr},\ L_m=t,\ L_c=b_0,
}\] \[\fitmath{
 L_M=2t+\epsilon_P,\ L_{\rm sel}=0,\
 L_{\rm mesh}=\eta_{\rm nat}(t)+4t,\quad\eta_{\rm nat}(t)\le t.
}\tag{GT28}\] Using d\textless=d\_+\textless1, delta\textless=3/4 and ell\textless1/5 gives exactly the retained bound \[\fitmath{
 16e+\mathcal L+t
 <9t+\mu/10000<\mu/4.
}\tag{GT29}\] Together with (GT18), this pays the nonfloor numerator on every assigned conductor interval, including all main/error choices and the full local mass (222). Error slots have zero amplitude and remain controlled by the same antitone count envelope.

Fix the source system, then the target, its finite character dictionary and all internal profile/Sobolev/moment degrees. Choose a positive target-dependent height exponent below their finitely many caps and the original requirements tau\textless d\_0/2, 4tau\textless d\_0*cost, 2tau\textless=t, tau(2+4epsilon\_tr)\textless t, and 2tau(1+J)\textless=t. The three external height coordinates Im(s+z), Im(z), Im(w) remain independent. The Gaussian and Mellin weights control their tails. Fix the raw polynomial envelope first and choose the external tail order last; it does not feed back into the mesh, slot widths or internal degrees.

The exact weighted-numerator identity in Appendix B, with (GT9), has main exponent is E\_kappa-k. Every factor, including delta\emph{p}(ell-C/2), remains. The full probe therefore satisfies \[\fitmath{
 J_\eta(Z)-f_\eta(Z)\ll_\eta Z^{\beta-k-\sigma_1}
 \quad\hbox{for a common }\sigma_1>0.
}\tag{GT30}\]

\hypertarget{signed-physical-estimate-and-all-height-conclusion}{%
\subsubsection{8.5 Signed physical estimate and all-height conclusion}\label{signed-physical-estimate-and-all-height-conclusion}}

The complete signed Gram/marked-energy estimate (198)--(199) requires 0\textless ell\textless x\textless=y, x+y+ell\textless=1 and y-2x+ell\textless0. These are exactly (GT1),(GT3). Its excess 1+3ell-2(x+y)=5ell-1 is negative by ell\textless1/5. All subset terms are retained by the original proof, whose subset-width slopes are -1 and -7/8. The same complete probe and nonzero logarithmic normalizer thus obey \[\fitmath{
 J_\eta(Z)\ll_\eta Z^{x/2+(y-x)/12+\omega}
 =Z^{\theta-k+\omega}\qquad(\omega>0).
}\tag{GT31}\]

Its signal is still the exact fixed-numerator integral (203), \[\fitmath{
 f_\eta(Z)=\frac1{2\pi i}\int_{(2)}
 Z^{s-k}e^{(s-5/6)^2}\frac{H_\eta(s)}{L_K^S(s,\eta)}\,ds.
}\tag{GT32}\] The Euler numerator (200)--(201) is fixed after the original data and is nonzero on Re(s)\textgreater=.87. No changed geometric parameter creates a moving numerator. The finite correction (202) has already been paid in (GT24).

Choose omega\textless mu/4 and 0\textless m\textless min(mu-omega,sigma\_1). Then (GT30)--(GT31) give f\_eta(Z)\textless\textless Z\^{}(beta-k-m) at infinity, while its absolute defining line gives arbitrary power decay at zero. Mellin inversion continues the displayed quotient holomorphically to Re(s)\textgreater beta-m. Since beta-m\textgreater theta\textgreater=.87, the numerator cannot cancel an original zero there. The primitive-supremum argument supplies a zero beyond beta-m/2 and gives the contradiction without assuming supremum attainment. The original ordinary norm pullback and restoration of fixed Euler factors complete the strict all-height ordinary conclusion, excluding only the actual principal pole.

\hypertarget{the-retained-algebraic-endpoint}{%
\subsection{9. The retained algebraic endpoint}\label{the-retained-algebraic-endpoint}}

With the original fixed presented modulus, ordinary character and principal-pole convention, the retained baseline theorem is

\[\fitmath{
\boxed{L(s,\chi)\ne0\quad\text{for all}\quad
\Re s>\theta_*:=\frac{777-\sqrt{1081}}{852},\quad s\ne1
\text{ when the actual principal pole is present}.}
}\tag{AE1}\]

Equivalently \(\fitinline{c_*=(75+\sqrt{1081})/852}\). The theorem covers every height and every ordinary character at each fixed positive presented modulus; it asserts neither the boundary nor uniform constants for growing moduli. The separate angular theorem retains its original fixed-family scope.

The proof uses the complete open source theorem, the exact source-face inclusion D1, the complete common-witness count theorem and the generic physical/complement/all-height return. The fixed source (D2) supplies all 24 inherited enhanced count faces, so every one of the 813 inherited windows is available without reproducing the older numerical source chains. The old natural branch remains available everywhere. All 1,826 original return rows are retained as the finite covering skeleton.

The preliminary theorem \(\fitinline{\beta\le7/8}\) suffices. Under a contrary supremum above any target used below, \(\fitinline{\beta\ge51/100}\) holds and all sources are available independently of the desired endpoint. Count and physical sources have separate roles: (D2) supplies the fixed inherited bank, while a nearby strict source selected below supplies the physical interval and the new local contact branch.

We first certify a nominal algebraic cover. Its contact is then paid by a quantitatively controlled family of strict source and return parameters. Taking the union of the resulting open half-planes gives (AE1); no boundary source or limiting uniform estimate is used.

\hypertarget{exact-nominal-parameters-and-the-complete-count-near-the-contact}{%
\subsubsection{9.1. Exact nominal parameters and the complete count near the contact}\label{exact-nominal-parameters-and-the-complete-count-near-the-contact}}

Write \(\fitinline{r=\sqrt{1081}}\), and define

\[\fitmath{
 v_*=(61-r)/66,\qquad \ell_*=(4+r)/213,
 \quad h_* =\frac{32}{71}+\frac{838r}{76751}.                  
}\tag{AE3}\]

Thus \(\fitinline{v_*=.4260823568\ldots}\), \(\fitinline{h_*=.8096863494\ldots}\), and

\[\fitmath{
 h_* =\frac{16(1/2+\ell_*)}{11+25/(5-4v_*)^2},\qquad
 \ell_*=11/3-4\theta_* .                                    
}\tag{AE4}\]

Set \(\fitinline{x_*=1+\ell_*-h_*}\), \(\fitinline{y_*=h_*-2\ell_*}\), and \(\fitinline{\kappa_*=2\theta_*-1}\). The conductor edges are initially \(\fitinline{d_0=1/200}\), \(\fitinline{d_m=19/25}\), and \(\fitinline{h_*}\). At this nominal stage only, the upper extension is zero. We subsequently replace it by a strictly positive extension; no large-row argument uses zero extension.

For the new local branch use the nominal source parameters \(\fitinline{B=1,L=2}\), source cap \(\fitinline{4/25}\), declared count cap \(\fitinline{1599/10000}\), actual width \(\fitinline{z=1598999/10000000}\), witness floor \(\fitinline{\mu=3/10}\), \(\fitinline{\eta=10^{-7}}\), \(\fitinline{m_H=173/400}\), and \(\fitinline{s_-=3/100}\). Its true endpoints are

\[\fitmath{
 m_L=\frac{3401002}{10000000},\qquad
 m_U=\frac{4324999}{10000000}.                              
}\tag{AE5}\]

On the closed upward rectangle

\[\fitmath{
 Q=[21/50,11/25]\times[49/100,1/2]
       \quad\text{in }(\delta,p),                           
}\tag{AE6}\]

put \(\fitinline{V=5/6-\delta}\) and \(\fitinline{H=5/2-\delta(1+2p)}\). The nominal branch is

\[\fitmath{
 R(\delta,p)=1-\delta+
       \frac{\delta V(1-p)}{H},\qquad J=(1-p)H>0.            
}\tag{AE7}\]

This is the literal affine formula, with its actual clipped width \(\fitinline{\min\{z,(B-2m)/L\}}\), not a constant-width continuation beyond the source cone. All fifteen (C7) polynomials, using (AE5), have positive exact Bernstein coefficients on (AE6). The nominal certificate has 75 such coefficients. \S{}9.4 supplies a uniform strict-source version of these admissions.

\hypertarget{the-exact-local-contact-inequality}{%
\subsubsection{9.2. The exact local contact inequality}\label{the-exact-local-contact-inequality}}

At the nominal upper edge \(\fitinline{d=h_*}\), the physical capacity is

\[\fitmath{
 C_* =\min\{7h_*/50,\ y_*-h_*/2\}=h_*/2-2\ell_* .          
}\tag{AE8}\]

The equality and a positive cap reserve are checked exactly. With the complete exponent GT9, substitution of (AE7) at \(\fitinline{p=1/2}\) gives

\[\fitmath{
 E_*(\delta,1/2)=\frac12+\frac{\ell_*}{2}
 +\delta\left(\frac12+\ell_*\right)
 +h_*\left(\frac13-\frac{11\delta}{16}
       -\frac{5\delta}{16(5-4\delta)}\right).                
}\tag{AE9}\]

Define

\[\fitmath{
 q_* =2+4\ell_*-11h_*/4>0 .                                
}\tag{AE10}\]

The following is a polynomial identity over \(\fitinline{\mathbb Q(r)}\):

\[\fitmath{
 \boxed{\theta_*-E_*(\delta,1/2)
      =\frac{q_*(\delta-v_*)^2}{5-4\delta}.}                 
}\tag{AE11}\]

The denominator is positive throughout \(\fitinline{0\le\delta\le3/4}\). In particular (AE11) is nonnegative on the entire density interval; its sole zero is \(\fitinline{v_*}\). The identity follows either by expanding the two quadratic numerators, or by using \(\fitinline{33v_*^2-61v_*+20=0}\), (AE4), and equality at \(\fitinline{v_*}\). It is a direct polynomial identity; the finite coefficient records also contain this equality.

It remains to control the full local rectangle, including \(\fitinline{p<1/2}\). From (AE7),

\[\fitmath{
 \partial_p R=-\frac{3\delta V^2}{H^2},\qquad
 \partial_p E_*=\delta\left(2\ell_*-h_*/4
          -\frac{9h_*V^2}{4H^2}\right).                     
}\tag{AE12}\]

For fixed \(\fitinline{\delta}\), \(\fitinline{H\ge5/2-2\delta}\) when \(\fitinline{p\le1/2}\). Moreover

\[\fitmath{
 \frac{d}{d\delta}\frac{5/6-\delta}{5/2-2\delta}
        =-\frac{5/6}{(5/2-2\delta)^2}<0.
}\]

The bracket in (AE12) is therefore bounded below on \(\fitinline{Q}\) by

\[\fitmath{
 b_*:=2\ell_*-h_*/4
   -\frac{9h_*}{4}
      \left(\frac{5/6-21/50}{5/2-42/50}\right)^2>0.          
}\tag{AE13}\]

Its exact value in \(\fitinline{\mathbb Q(r)}\) is recorded in the certificate; it exceeds .0309. Thus \(\fitinline{E_*(\delta,p)\le E_*(\delta,1/2)}\) on \(\fitinline{Q}\). Equations (AE11)--(AE13) prove \(\fitinline{E_*\le\theta_*}\) on every contact cell, with equality only at \(\fitinline{(v_*,1/2)}\). No sample or mesh limit proves this assertion.

\hypertarget{the-complete-nominal-return-away-from-contact}{%
\subsubsection{9.3. The complete nominal return away from contact}\label{the-complete-nominal-return-away-from-contact}}

\texttt{r\allowbreak{}e\allowbreak{}s\allowbreak{}e\allowbreak{}a\allowbreak{}r\allowbreak{}c\allowbreak{}h\allowbreak{}/\allowbreak{}c\allowbreak{}y\allowbreak{}c\allowbreak{}l\allowbreak{}e\allowbreak{}0\allowbreak{}7\allowbreak{}4\allowbreak{}/\allowbreak{}m\allowbreak{}a\allowbreak{}t\allowbreak{}h\allowbreak{}/\allowbreak{}B\allowbreak{}O\allowbreak{}U\allowbreak{}N\allowbreak{}D\allowbreak{}A\allowbreak{}R\allowbreak{}Y\allowbreak{}\_\allowbreak{}C\allowbreak{}E\allowbreak{}R\allowbreak{}T\allowbreak{}I\allowbreak{}F\allowbreak{}I\allowbreak{}C\allowbreak{}A\allowbreak{}T\allowbreak{}E\allowbreak{}.\allowbreak{}j\allowbreak{}s\allowbreak{}o\allowbreak{}n} retains all 1,826 original rows from \texttt{r\allowbreak{}e\allowbreak{}s\allowbreak{}e\allowbreak{}a\allowbreak{}r\allowbreak{}c\allowbreak{}h\allowbreak{}/\allowbreak{}c\allowbreak{}y\allowbreak{}c\allowbreak{}l\allowbreak{}e\allowbreak{}0\allowbreak{}6\allowbreak{}7\allowbreak{}/\allowbreak{}r\allowbreak{}e\allowbreak{}t\allowbreak{}u\allowbreak{}r\allowbreak{}n\allowbreak{}/\allowbreak{}P\allowbreak{}R\allowbreak{}9\allowbreak{}5\allowbreak{}\_\allowbreak{}R\allowbreak{}E\allowbreak{}T\allowbreak{}U\allowbreak{}R\allowbreak{}N\allowbreak{}\_\allowbreak{}D\allowbreak{}A\allowbreak{}T\allowbreak{}A\allowbreak{}.\allowbreak{}j\allowbreak{}s\allowbreak{}o\allowbreak{}n}. It subdivides them into 1,920 closed cells:

\begin{longtable}[]{@{}lllr@{}}
\toprule\noalign{}
Edge & Conductor & Capacity & Cells \\
\midrule\noalign{}
\endhead
\bottomrule\noalign{}
\endlastfoot
oldSmall & \(\fitinline{1/200}\) & natural & 448 \\
oldSwitch & \(\fitinline{19/25}\) & natural & 448 \\
newSwitch & \(\fitinline{19/25}\) & nominal \(\fitinline{B=1,L=2}\) & 448 \\
newTop & \(\fitinline{h_*}\) & nominal \(\fitinline{B=1,L=2}\) & 576 \\
\end{longtable}

There are 1,844 strictly positive polynomial cells and 76 contact cells. Each contact cell lies wholly inside (AE6) and uses \S{}9.2. Every other cell uses an inherited admitted branch, only within a complete retained upward admission window containing that cell. The old branch is available on the whole rectangle. The four covers equal the full closed domain \(\fitinline{[0,3/4]\times[0,1/2]}\), with exact original-row subdivision and no gap.

For an assigned branch \(\fitinline{R_b=N_b/J_b}\), the strict cells certify

\[\fitmath{
 J_b(\theta_*-E_{\mathrm{base},*})-(3d/4)N_b>0,
       \qquad J_b>0                                         
}\tag{AE14}\]

by positive tensor Bernstein coefficients. All coefficients are exact pairs \(\fitinline{[a,b]}\) representing \(\fitinline{a+b\sqrt{1081}}\). The certificate has 19,650 positive return coefficients and 9,412 positive denominator coefficients. \S{}10 specifies the exact field arithmetic, polynomial conversion, domain coverage and the implication of these finite inequalities. The unchanged historical independent receipt is evidence for these exact data; its scope is finite arithmetic.

For reference, a uniform lower bound for the actual strict-cell energy margin is

\[\fitmath{
 \mu_*=
 \frac{82148464380569}{461721298125000}
 -\frac{43214656680291589}{7985931572370000000}\sqrt{1081}
 >0,
 \quad \mu_*=.0000004965320729\ldots .                        
}\tag{AE15}\]

This bound is the minimum, over the strict cells, of the least numerator Bernstein coefficient divided by the greatest denominator Bernstein coefficient. Exact ordered-field comparisons establish it. It is not a claimed minimum of the actual energy margin.

The nominal finite cover alone is not the endpoint proof: \(\fitinline{d_+=h_*}\) lacks the necessary positive large-row extension, and the contact has no positive return reserve. We now repair both by an explicit family of valid applications of the generic theorem.

\hypertarget{an-explicit-strict-family-and-its-complete-sourcecount-admissions}{%
\subsubsection{9.4. An explicit strict family and its complete source/count admissions}\label{an-explicit-strict-family-and-its-complete-sourcecount-admissions}}

Fix an arbitrary real parameter

\[\fitmath{
 0<t\le T:=10^{-10},\qquad \rho=t^2,\qquad\zeta=t^2.          
}\tag{AE16}\]

Here \(\fitinline{t}\) is an endpoint-approximation parameter; it is distinct from the later coefficient shift and the later mesh/loss variables in GT13 and GT25--GT29. Choose

\[\fitmath{
\begin{gathered}
 \theta_t=\theta_*+t,\quad c_t=c_*-t,\quad h_t=h_*,\quad
 \ell_t=\ell_*-4t,\\
 x_t=x_*-4t,\quad y_t=y_*+8t,\quad
 \kappa_t=\kappa_*+2t,\quad d_+=h_*+\zeta,\\
 B_\rho=1-\rho,\quad L_\rho=2+\rho,\quad c_s=4/25,
 \quad\sigma=7/50,\quad\bar\sigma=141/1000.                  
\end{gathered}
}\tag{AE17}\]

The complete open source theorem UF1 applies because \(\fitinline{0<B_\rho<1}\), \(\fitinline{L_\rho>1}\), \(\fitinline{L_\rho-2B_\rho=3\rho>0}\), and \(\fitinline{0<c_s<1/6}\). It supplies one fixed complete source

\[\fitmath{
                 G(1-t^2,2+t^2,4/25)                         
}\tag{AE18}\]

for the entire new conductor interval and every density band. The source is fixed for the chosen \(\fitinline{t}\), before any subsequent requested analytic loss, mesh, slots, profiles or live heights. For each requested loss its UF1 finite induction chooses its positive norm gap and meshes in the required order. The stage count can grow with \(\fitinline{t}\); each application is a fixed finite induction with the universal-stage inequalities of UF1. No uniform constant as \(\fitinline{t\downarrow0}\), enumerated stage execution, selected norm at one half, or boundary physical-source consumer is invoked.

Use this same strict source for the local affine branch, with the fixed \(\fitinline{z,\bar z,\mu,\eta,m_H,s_-}\) of \S{}9.1. Its actual lower endpoint is

\[\fitmath{
 m_L(\rho)=(B_\rho-L_\rho z)/2+\eta,
 \quad m_U=m_H-\eta.                                       
}\tag{AE19}\]

The maximum and minimum in (C6a) and the literal upper endpoint really do select these terms: the exact interval checker verifies their strict separations. In particular the source cone determines the lower endpoint, and the unchanged natural join determines the upper one.

\texttt{r\allowbreak{}e\allowbreak{}s\allowbreak{}e\allowbreak{}a\allowbreak{}r\allowbreak{}c\allowbreak{}h\allowbreak{}/\allowbreak{}c\allowbreak{}y\allowbreak{}c\allowbreak{}l\allowbreak{}e\allowbreak{}0\allowbreak{}7\allowbreak{}4\allowbreak{}/\allowbreak{}m\allowbreak{}a\allowbreak{}t\allowbreak{}h\allowbreak{}/\allowbreak{}P\allowbreak{}E\allowbreak{}R\allowbreak{}T\allowbreak{}U\allowbreak{}R\allowbreak{}B\allowbreak{}A\allowbreak{}T\allowbreak{}I\allowbreak{}O\allowbreak{}N\allowbreak{}\_\allowbreak{}C\allowbreak{}E\allowbreak{}R\allowbreak{}T\allowbreak{}I\allowbreak{}F\allowbreak{}I\allowbreak{}C\allowbreak{}A\allowbreak{}T\allowbreak{}E\allowbreak{}.\allowbreak{}j\allowbreak{}s\allowbreak{}o\allowbreak{}n} proves all fifteen (C7) inequalities on the entire product

\[\fitmath{
 (\delta,p)\in Q,\qquad 0\le\rho\le10^{-12}.                
}\tag{AE20}\]

It forms the literal rational coefficient expressions, encloses each coefficient by rational interval arithmetic in \(\fitinline{\rho}\), and applies the tensor Bernstein conversion on the full \(\fitinline{(\delta,p)}\) rectangle. All 79 resulting interval coefficient lower bounds are positive; the least is

\[\fitmath{
                 \frac{389644444334009}{66666666666700000}>0.
}\]

Interval arithmetic retains the whole continuous parameter range. It does not sample parameter values. The slightly higher coefficient count than in the nominal certificate comes from retaining enclosing intervals for coefficients that cancel exactly at \(\fitinline{\rho=0}\). Since \(\fitinline{\rho=t^2\le10^{-20}}\), (AE20) covers every actual application.

Theorem C now supplies the complete strict-source count on the full upward window \(\fitinline{Q}\), with every original same-zero witness alternative, the source continuation \(\fitinline{\min\{z,(B_\rho-2m)/L_\rho\}}\), natural tail, small and zero widths, original loss and finite polynomial height degree. Its derivative guard makes it antitone in \(\fitinline{p}\). Adding this window to the unchanged 813 admitted windows preserves the single conductor-independent antitone count minimum.

\hypertarget{quantitative-transfer-of-every-nominal-cell}{%
\subsubsection{9.5. Quantitative transfer of every nominal cell}\label{quantitative-transfer-of-every-nominal-cell}}

We compare the actual complete exponent for (AE17) with the nominal one on the same closed cells. These estimates preserve the whole GT9 exponent, including its marked-capacity coefficient and conductor factor.

\hypertarget{capacities}{%
\subsubsection{Capacities}\label{capacities}}

On a natural edge the conductor is unchanged. With \(\fitinline{n_t=\min(d,2y_t-d)}\), one has \(\fitinline{n_t\ge n_0}\), and hence

\[\fitmath{
 C_{O,t}-C_{O,0}
 \ge-\frac{n_0t}{3\kappa_*\kappa_t}\ge-\frac{100}{147}t>-t.
}\tag{AE21}\]

Here \(\fitinline{d<1}\), \(\fitinline{0\le n_0\le d}\), and both kappas exceed \(\fitinline{7/10}\). On a new edge put \(\fitinline{d_t=d+\Delta d}\), where \(\fitinline{\Delta d=0}\) at the switch and \(\fitinline{\Delta d=\rho}\) at the top. The cap term \(\fitinline{\sigma d_t}\) increases. For the affine term, subtracting the nominal value after multiplication by \(\fitinline{2+\rho}\) gives

\[\fitmath{
 16t-\Delta d-\rho d-\rho\Delta d
      -\rho(y_*-d/2)\ge16t-3\rho-\rho^2>0.                  
}\tag{AE22}\]

The last inequality follows from \(\fitinline{\rho=t^2}\) and (AE16). All quantities bounded by one in this estimate are positive and less than one by the geometric checks. Both terms of the minimum increase, so \(\fitinline{C_{G,t}\ge C_{G,0}}\). Thus (AE21)'s weaker lower bound \(\fitinline{C_t-C_0\ge-t}\) is valid for every edge.

\hypertarget{fixed-inherited-counts}{%
\subsubsection{Fixed inherited counts}\label{fixed-inherited-counts}}

Every imported nonold branch is independent of \(\fitinline{\kappa_t}\) and is unchanged. For the old branch, write

\[\fitmath{
 D=3-p-2p/(3\kappa),\quad
 A_0=2(1-p)(1-p/(3\kappa)),\quad J=VD+\delta A_0,
}\] \[\fitmath{
 R_{\rm old}=1-\delta+\frac{\delta A_0V}{2J}.
}\]

Direct differentiation yields

\[\fitmath{
 \partial_\kappa R_{\rm old}
 =\frac{\delta V^2p(1-p)^2}{3\kappa^2J^2}.                  
}\tag{AE23}\]

For \(\fitinline{7/10\le\kappa\le3/4}\), \(\fitinline{0\le\delta\le3/4}\), \(\fitinline{0\le p\le1/2}\), one has \(\fitinline{D\ge85/42>2}\), \(\fitinline{A_0\ge16/21>3/4}\), and \(\fitinline{J\ge35/48>7/10}\). Consequently

\[\fitmath{
 0\le\partial_\kappa R_{\rm old}
       \le15625/43218<1,
 \qquad R_{{\rm old},t}-R_{{\rm old},0}\le2t.                
}\tag{AE24}\]

Every assigned branch obeys \(\fitinline{0\le R_b\le1}\). For a nonold branch, \(\fitinline{R_b=1-\delta+\delta VU/J}\), and its positive admission gives \(\fitinline{U\le D}\) and \(\fitinline{J\ge VD}\); in the affine formula \(\fitinline{D-U=2-4p/L+2pB(1-p)/L\ge2-2/L>0}\), and the capped formula has \(\fitinline{D\ge5/2}\), \(\fitinline{U\le1}\). For the old branch \(\fitinline{A_0/2\le D}\). These establish the claimed upper bounds, including the zero-density endpoint. Positivity follows from \(\fitinline{1-\delta\ge1/4}\) and the nonnegative correction. This explicitly pays the top-edge change in the factor \(\fitinline{dR_b}\).

On every strictly certified imported cell, use (AE21), (AE24), \(\fitinline{\delta p/2\le3/16}\), \(\fitinline{d_t<1}\), and \(\fitinline{\Delta d\le\rho}\). The changes \(\fitinline{y_t-y_*=8t}\), \(\fitinline{\ell_t-\ell_*=-4t}\) give

\[\fitmath{
 E_t-E_0\le-2t+\frac{3t}{16}
                   +\frac34(2t+\rho)
             =-\frac{5t}{16}+\frac{3\rho}{4}.                
}\tag{AE25}\]

The omitted terms \(\fitinline{-4\delta pt}\) and \(\fitinline{-9\Delta d/100}\) are nonpositive, so this is a valid upper bound on the full exponent. Therefore

\[\fitmath{
 \theta_t-E_t\ge\mu_*+\frac{21t}{16}-\frac{3t^2}{4}.
}\tag{AE26}\]

\hypertarget{the-local-count-and-contact-cells}{%
\subsubsection{The local count and contact cells}\label{the-local-count-and-contact-cells}}

For the actual local branch let

\[\fitmath{
\begin{aligned}
 D_\rho&=3-p-4p/(2+\rho),&
 A_\rho&=2-4p/(2+\rho),\\
 C_\rho&=1-2p(1-\rho)/(2+\rho),&
 U_\rho&=(1-p)C_\rho,\\
 J_\rho&=VD_\rho+\delta(1-p)A_\rho, &
 R_\rho&=1-\delta+V\delta U_\rho/J_\rho.
\end{aligned}                                                   
}\tag{AE27}\]

On (AE20), \(\fitinline{D_\rho,A_\rho}\) increase, and \(\fitinline{J_\rho\ge J_0\ge81/100>4/5}\). Also \(\fitinline{J_\rho\le5/2}\), \(\fitinline{0\le U_\rho\le1}\), \(\fitinline{0\le U'_\rho\le3/4}\), and \(\fitinline{0\le J'_\rho\le5/12}\). Using \(\fitinline{V\delta\le25/144}\), we obtain

\[\fitmath{
 |\partial_\rho R_\rho|
 \le\frac{25}{144}\frac{(3/4)(5/2)+5/12}{(4/5)^2}
 =\frac{34375}{55296}<1.
 \quad |R_\rho-R_0|\le\rho.                                
}\tag{AE28}\]

Every contact cell is at the new top edge, where the physical capacity increases by (AE22). Both the change in its count and the change in its conductor therefore cost at most \(\fitinline{3\rho/2}\). The contact identity and capacity comparison give

\[\fitmath{
 \boxed{\theta_t-E_t\ge3t-\frac{3t^2}{2}>2t
                   \quad\text{on every contact cell}.}     
}\tag{AE29}\]

Thus every cell now has a positive reserve, including the former equality point. A concrete choice is

\[\fitmath{
       r_t=t,\qquad \mathrm{rowFee}=t/100,
                    \qquad r_{0,t}=t/10 .                    
}\tag{AE30}\]

Both (AE26) and (AE29) exceed \(\fitinline{101t/100}\) on (AE16), so the row fee is paid in addition to the full reserve. These are explicit inequalities on an entire parameter interval, not a sequence of successful numerical runs.

\hypertarget{complementary-sectors-and-the-all-height-composition}{%
\subsubsection{9.6. Complementary sectors and the all-height composition}\label{complementary-sectors-and-the-all-height-composition}}

The perturbation certificate encloses \(\fitinline{c_*}\) and \(\fitinline{h_*}\) between explicit rational endpoints, checked by exact ordered-field arithmetic. It then verifies 36 positive rational interval guards for all \(\fitinline{0\le t\le10^{-10}}\), \(\fitinline{0\le\rho,\zeta\le10^{-20}}\). The actual relation \(\fitinline{\rho=\zeta=t^2}\) lies inside this enclosing box. The two strict source fields \(\fitinline{1-B=\rho>0}\) and \(\fitinline{L-2=\rho>0}\), the UF field \(\fitinline{L-2B=3\rho>0}\), and the large-row extension \(\fitinline{\zeta>0}\) are paid directly by (AE16), not by claiming positivity at the nominal boundary.

The checked guards include all GT1--GT12 geometry, the independent actual/declared/source cap buffers, prime supply, the declared-cap natural fallback, positive affine capacity, the intrinsic floor after its \(\fitinline{1/1000}\) allowance, the complete small-conductor estimate, both principal contours, \(\fitinline{7/10<\kappa_t<3/4}\), and \(\fitinline{.87<\theta_t<436691/500000}\). The exact identities GT1--GT2 hold by (AE17). The source interval \(\fitinline{[19/25,h_*+t^2]}\) starts above \(\fitinline{y_t}\); the natural interval \(\fitinline{[1/200,19/25]}\) covers the rest of the central conductors.

For clarity, the principal and complementary margins checked are the actual GT expressions or the proved conservative lower bounds:

\[\fitmath{
\begin{aligned}
 \Delta_f&=61/150-h_*/3-39\ell_t/50-67t^2/100,\\
 \Delta_s&=49h_*/150-5\ell_t/4+13/600,\\
 g_w&=49h_*/1000-\ell_t/10-1/800,\\
 g_z&=h_*/600-1/800>0.
\end{aligned}                                                   
}\tag{AE31}\]

The conservative \(\fitinline{(6\ell_t-1)/2400}\) lower bound for the last margin is also positive. The bound itself is not confused with equality to the actual contour margin.

For every coefficient shift \(\fitinline{s\ge0}\), the actual physical selection is

\[\fitmath{
 \max\left(0,\min\left(7/50,
 \frac{B_\rho-2+2y_t/d-2s}{L_\rho}\right)\right).             
}\tag{AE32}\]

The \(\fitinline{-2s}\) shift is retained without a finite cutoff. The actual declared cap pays whole-slot rounding, and the strict source cap pays all positive selections, shortenings and remaining subsets. GT13--GT15 provide the nonpositive-ordinary natural fallback, bounded empty-selection case, full shift damping and prime supply, with unchanged original characters, masks, roots and coefficients. The natural fallback uses \(\fitinline{\kappa_{\rm nat}=3/4}\), never the new desired endpoint.

On each fixed conductor interval the actual capacity is a minimum of affine functions of \(\fitinline{d}\). Its coefficient in GT9 is nonpositive; the other conductor dependence is affine and the count envelope is conductor-independent. Convexity therefore extends the complete endpoint bounds (AE26), (AE29) over the entire intervals. Both valid suppliers hold at the common edge. Every central conductor, density, amplitude and coefficient shift is covered.

The floor, all nonprincipal units, small conductors, large rows, principal residues, fixed numerator, finite-slot correction, fixed-inducing terms, zero rows and frequencies, signed compensators, prime powers and profile tails are the same terms as GT19--GT24. Their numerical guards hold uniformly in (AE16). For the large rows, first fix \(\fitinline{\zeta=t^2>0}\), then choose the large external contour parameter in GT23; it is not chosen uniformly as \(\fitinline{t}\) tends to zero.

Suppose the actual primitive-family supremum satisfies \(\fitinline{\beta>\theta_t}\). The original preliminary theorem gives \(\fitinline{\beta\le7/8}\), and the standing source condition \(\fitinline{\beta\ge51/100}\) follows. Thus \(\fitinline{\kappa_t\le2\beta-1\le3/4}\), exactly the range needed for the generic feedback. No intermediate numerical endpoint is a premise. All enhanced sources and count windows for this \(\fitinline{t}\) have already been fixed and are independent of \(\fitinline{\beta-\theta_t}\). Only the old count and natural capacity vary with \(\fitinline{\kappa_\beta=2\beta-1}\). GT16--GT18 consequently give, with \(\fitinline{\nu=\beta-\theta_t>0}\),

\[\fitmath{
 E_{\kappa_\beta}-\beta\le-r_{0,t}-24\nu/25
       \quad\text{on the new interval},
}\] \[\fitmath{
 E_{\kappa_\beta}-\beta\le-r_{0,t}-3\nu/4
       \quad\text{on the natural interval}.                 
}\tag{AE33}\]

Take the positive minimum of the finitely many actual source, admission, cap, return and complementary reserves for this fixed \(\fitinline{t}\). GT25--GT29 choose the common actual mesh before the finite slot dictionary. The source and count theorems are invoked at their prescribed smaller losses. The two ordinary profiles and all Fourier heights stay independent. Finite seminorm orders and polynomial height degrees are fixed before the later small physical-height exponent, and the external tail order is chosen last. No source-length or root term is discarded in these choices.

The complete compensated probe and its signed physical estimate now satisfy GT30--GT31 at \(\fitinline{\theta_t}\). The unchanged nonvanishing Mellin numerator in GT32, on \(\fitinline{\Re s\ge.87}\), contradicts \(\fitinline{\beta>\theta_t}\) using zeros approaching the actual supremum. No maximal zero is assumed. This proves, for every fixed \(\fitinline{0<t\le10^{-10}}\), the ordinary all-height theorem

\[\fitmath{
                L(s,\chi)\ne0\qquad(\Re s>\theta_*+t),      
}\tag{AE34}\]

at each fixed positive presented modulus, excluding the actual principal pole.

Finally let \(\fitinline{s}\) satisfy \(\fitinline{\Re s>\theta_*}\) and be different from that pole. Choose \(\fitinline{0<t<\min(10^{-10},\Re s-\theta_*)}\). The proved theorem (AE34) applies to this one fixed \(\fitinline{t}\) and excludes a zero at \(\fitinline{s}\). This proves (AE1). The last step is a union of already proved open half-planes, not a passage to a uniform bound or a claim on the boundary.

\hypertarget{exact-finite-certification}{%
\subsection{10. Exact finite certification}\label{exact-finite-certification}}

The finite data are part of this packet. Their role is to specify exact parameters, admission rectangles and closed return cells. The analytic source, count and physical-return proofs are given in the manuscript. Neither a catalogue entry nor a program's success message supplies an analytic estimate.

\hypertarget{exact-data-and-their-mathematical-interpretation}{%
\subsubsection{10.1. Exact data and their mathematical interpretation}\label{exact-data-and-their-mathematical-interpretation}}

\begin{longtable}[]{@{}
  >{\raggedright\arraybackslash}p{(\columnwidth - 2\tabcolsep) * \real{0.5000}}
  >{\raggedright\arraybackslash}p{(\columnwidth - 2\tabcolsep) * \real{0.5000}}@{}}
\toprule\noalign{}
\begin{minipage}[b]{\linewidth}\raggedright
Immutable packet file
\end{minipage} & \begin{minipage}[b]{\linewidth}\raggedright
Data used in the proof
\end{minipage} \\
\midrule\noalign{}
\endhead
\bottomrule\noalign{}
\endlastfoot
\texttt{r\allowbreak{}e\allowbreak{}s\allowbreak{}e\allowbreak{}a\allowbreak{}r\allowbreak{}c\allowbreak{}h\allowbreak{}/\allowbreak{}c\allowbreak{}y\allowbreak{}c\allowbreak{}l\allowbreak{}e\allowbreak{}0\allowbreak{}6\allowbreak{}7\allowbreak{}/\allowbreak{}r\allowbreak{}e\allowbreak{}t\allowbreak{}u\allowbreak{}r\allowbreak{}n\allowbreak{}/\allowbreak{}P\allowbreak{}R\allowbreak{}9\allowbreak{}5\allowbreak{}\_\allowbreak{}R\allowbreak{}E\allowbreak{}T\allowbreak{}U\allowbreak{}R\allowbreak{}N\allowbreak{}\_\allowbreak{}D\allowbreak{}A\allowbreak{}T\allowbreak{}A\allowbreak{}.\allowbreak{}j\allowbreak{}s\allowbreak{}o\allowbreak{}n} & Eighteen literal branch definitions, including the old branch; 284 full upward windows; the 1,826 original closed return rows with their edge and rectangle. \\
\texttt{r\allowbreak{}e\allowbreak{}s\allowbreak{}e\allowbreak{}a\allowbreak{}r\allowbreak{}c\allowbreak{}h\allowbreak{}/\allowbreak{}c\allowbreak{}y\allowbreak{}c\allowbreak{}l\allowbreak{}e\allowbreak{}0\allowbreak{}7\allowbreak{}2\allowbreak{}/\allowbreak{}r\allowbreak{}e\allowbreak{}t\allowbreak{}u\allowbreak{}r\allowbreak{}n\allowbreak{}/\allowbreak{}R\allowbreak{}E\allowbreak{}T\allowbreak{}U\allowbreak{}R\allowbreak{}N\allowbreak{}\_\allowbreak{}C\allowbreak{}E\allowbreak{}R\allowbreak{}T\allowbreak{}I\allowbreak{}F\allowbreak{}I\allowbreak{}C\allowbreak{}A\allowbreak{}T\allowbreak{}E\allowbreak{}.\allowbreak{}j\allowbreak{}s\allowbreak{}o\allowbreak{}n} & Six additional branch definitions and 219 full upward admission windows. Only those data and their bound admission coefficients are needed here; its older endpoint is not a premise. \\
\texttt{r\allowbreak{}e\allowbreak{}s\allowbreak{}e\allowbreak{}a\allowbreak{}r\allowbreak{}c\allowbreak{}h\allowbreak{}/\allowbreak{}c\allowbreak{}y\allowbreak{}c\allowbreak{}l\allowbreak{}e\allowbreak{}0\allowbreak{}7\allowbreak{}3\allowbreak{}/\allowbreak{}r\allowbreak{}e\allowbreak{}t\allowbreak{}u\allowbreak{}r\allowbreak{}n\allowbreak{}/\allowbreak{}R\allowbreak{}E\allowbreak{}T\allowbreak{}U\allowbreak{}R\allowbreak{}N\allowbreak{}\_\allowbreak{}C\allowbreak{}E\allowbreak{}R\allowbreak{}T\allowbreak{}I\allowbreak{}F\allowbreak{}I\allowbreak{}C\allowbreak{}A\allowbreak{}T\allowbreak{}E\allowbreak{}.\allowbreak{}j\allowbreak{}s\allowbreak{}o\allowbreak{}n} & The common source's affine branch and 310 full upward admission windows, with lower floor \(\fitinline{3/10}\). Its older endpoint is not a premise. \\
\texttt{r\allowbreak{}e\allowbreak{}s\allowbreak{}e\allowbreak{}a\allowbreak{}r\allowbreak{}c\allowbreak{}h\allowbreak{}/\allowbreak{}c\allowbreak{}y\allowbreak{}c\allowbreak{}l\allowbreak{}e\allowbreak{}0\allowbreak{}7\allowbreak{}4\allowbreak{}/\allowbreak{}m\allowbreak{}a\allowbreak{}t\allowbreak{}h\allowbreak{}/\allowbreak{}B\allowbreak{}O\allowbreak{}U\allowbreak{}N\allowbreak{}D\allowbreak{}A\allowbreak{}R\allowbreak{}Y\allowbreak{}\_\allowbreak{}C\allowbreak{}E\allowbreak{}R\allowbreak{}T\allowbreak{}I\allowbreak{}F\allowbreak{}I\allowbreak{}C\allowbreak{}A\allowbreak{}T\allowbreak{}E\allowbreak{}.\allowbreak{}j\allowbreak{}s\allowbreak{}o\allowbreak{}n} & The exact \(\fitinline{\mathbb Q(\sqrt{1081})}\) parameters, local contact box/admission, and 1,920 return cells. Each records its original row, edge, complete box, assigned branch and proof method; strict cells also record all numerator and denominator Bernstein coefficients. \\
\texttt{r\allowbreak{}e\allowbreak{}s\allowbreak{}e\allowbreak{}a\allowbreak{}r\allowbreak{}c\allowbreak{}h\allowbreak{}/\allowbreak{}c\allowbreak{}y\allowbreak{}c\allowbreak{}l\allowbreak{}e\allowbreak{}0\allowbreak{}7\allowbreak{}4\allowbreak{}/\allowbreak{}m\allowbreak{}a\allowbreak{}t\allowbreak{}h\allowbreak{}/\allowbreak{}P\allowbreak{}E\allowbreak{}R\allowbreak{}T\allowbreak{}U\allowbreak{}R\allowbreak{}B\allowbreak{}A\allowbreak{}T\allowbreak{}I\allowbreak{}O\allowbreak{}N\allowbreak{}\_\allowbreak{}C\allowbreak{}E\allowbreak{}R\allowbreak{}T\allowbreak{}I\allowbreak{}F\allowbreak{}I\allowbreak{}C\allowbreak{}A\allowbreak{}T\allowbreak{}E\allowbreak{}.\allowbreak{}j\allowbreak{}s\allowbreak{}o\allowbreak{}n} & Rational interval coefficients for every local admission over \(\fitinline{0\le\rho\le10^{-12}}\); exact constant enclosures; the 36 geometric/complementary interval guards; derivative and perturbation bounds. \\
\end{longtable}

The branch data mean exactly (C1)--(C7). For the first seventeen enhanced PR95 branches and six C072 branches, use \(\fitinline{\mu=377/1000}\); for C073 use \(\fitinline{\mu=3/10}\). The local branch has the parameters in \S{}9.1 and \S{}9.4. All declared caps are converted to actual width by subtracting \(\fitinline{\eta=10^{-7}}\). No field named a physical cap is used as a count cap or conversely.

The three window lists form a fixed list of 813 complete upward windows. The common source (D2) supplies their actual broad faces by D1, including the faces whose slopes are below \(\fitinline{2B}\) and intercepts below an older wrapper's numerical box. Their original finite admission receipts are bound to their exact unchanged data. The new source-domain replacement changes no polynomial or count conclusion.

\hypertarget{bernstein-positivity-on-a-complete-closed-rectangle}{%
\subsubsection{10.2. Bernstein positivity on a complete closed rectangle}\label{bernstein-positivity-on-a-complete-closed-rectangle}}

Let a polynomial \(\fitinline{P(\delta,p)}\) have degree at most \(\fitinline{n}\) in \(\fitinline{\delta}\) and at most \(\fitinline{m}\) in \(\fitinline{p}\). On \(\fitinline{[a,b]\times[c,d]}\), make the exact affine substitution

\[\fitmath{
\delta=a+(b-a)u,\qquad p=c+(d-c)v,
}\]

and write the result as \(\fitinline{\sum_{k=0}^n\sum_{l=0}^m a_{kl}u^kv^l}\). Its tensor Bernstein coefficients are

\[\fitmath{
b_{ij}=\sum_{k\le i}\sum_{l\le j}
 a_{kl}\frac{\binom{i}{k}}{\binom{n}{k}}
         \frac{\binom{j}{l}}{\binom{m}{l}}
\quad(0\le i\le n,\ 0\le j\le m).
}\tag{K1}\]

Indeed \(\fitinline{u^k=\sum_{i=k}^n\binom{i}{k}\binom{n}{k}^{-1} \binom ni u^i(1-u)^{n-i}}\), and the analogous formula holds for \(\fitinline{v}\). Expanding proves (K1), including all closed boundaries. The tensor basis functions are nonnegative and sum to one on the unit square. Therefore

\[\fitmath{
\min_{i,j} b_{ij}\le P(\delta,p)\le\max_{i,j}b_{ij}.
}\tag{K2}\]

Strict positivity of every coefficient is a sufficient proof of positivity on the entire rectangle; no sampling or mesh limit enters. The method is not asserted to be necessary or optimal.

For a branch \(\fitinline{R_b=N_b/J_b}\), the strict nominal cell polynomial is

\[\fitmath{
P_b=J_b(\theta_*-E_{\mathrm{base},*})-(3d/4)N_b.
}\tag{K3}\]

The certificate provides the tensor coefficients of both \(\fitinline{P_b}\) and \(\fitinline{J_b}\). Their positivity gives \(\fitinline{J_b>0}\), \(\fitinline{E_b<\theta_*}\), and the quantitative lower bound

\[\fitmath{
\theta_*-E_b\ge
\frac{\min\operatorname{Bern}(P_b)}{
      \max\operatorname{Bern}(J_b)}>0.
}\tag{K4}\]

Since the actual common count is the minimum of all available branch counts, \(\fitinline{R_*\le R_b}\), and its coefficient in the physical exponent is positive, the same bound holds for the actual exponent. This explains why choosing different admitted branches on different cells is valid.

For the 76 contact cells, the certificate uses the exact contact factorization and positive \(\fitinline{p}\)-derivative in \S{}9.2. It does not claim positive Bernstein coefficients at the equality point. The 1,844 other cells have 19,650 positive numerator and 9,412 positive denominator coefficients. The common lower bound in (AE15) is the least of the ratios (K4), so it is an explicit lower bound, not a claimed minimum of the underlying continuous energy margin.

\hypertarget{ordered-quadratic-field-and-interval-arithmetic}{%
\subsubsection{10.3. Ordered quadratic-field and interval arithmetic}\label{ordered-quadratic-field-and-interval-arithmetic}}

A pair of rational strings \(\fitinline{[a,b]}\) represents \(\fitinline{a+b\sqrt{1081}}\). Addition and multiplication are

\[\fitmath{
(a,b)+(c,d)=(a+c,b+d),\quad
(a,b)(c,d)=(ac+1081bd,ad+bc).
}\tag{K5}\]

Division uses conjugation and the nonzero rational denominator \(\fitinline{a^2-1081b^2}\). The integer 1081 is not a rational square. Signs are determined exactly: equal-sign terms have that sign; for opposing signs compare \(\fitinline{a^2}\) with \(\fitinline{1081b^2}\), retaining the sign of the larger-magnitude term. Alternatively one may use rational square-root enclosures whose squared endpoints bracket 1081. No floating-point assertion proves a sign in this argument.

For the strict perturbation, rational intervals enclose every coefficient before conversion (K1). Addition is endpointwise, multiplication takes the minimum and maximum of the four endpoint products, and inversion is used only on an interval wholly of one nonzero sign. These operations preserve enclosure. Positive lower endpoints of the Bernstein intervals therefore prove all guards for the entire parameter interval simultaneously, even when dependency among repeated parameter occurrences widens an enclosure.

The interval proof uses

\[\fitmath{
\begin{aligned}
\frac{126618033386792}{10^{15}}&<c_*<
\frac{126618033386793}{10^{15}},\\
\frac{809686349433473}{10^{15}}&<h_*<
\frac{809686349433475}{10^{15}}.
\end{aligned}
}\tag{K6}\]

These are exact quadratic-field comparisons. The 79 local admission coefficient lower endpoints are positive on \(\fitinline{0\le\rho\le10^{-12}}\), with least value

\[\fitmath{
\frac{389644444334009}{66666666666700000}>0.
}\]

All 36 listed geometric/complementary guards have positive interval lower endpoints on the enclosing box

\[\fitmath{
0\le t\le10^{-10},\qquad
0\le\rho,\zeta\le10^{-20}.
}\tag{K7}\]

The actual relations \(\fitinline{\rho=\zeta=t^2}\) lie in this box. The strict fields \(\fitinline{\rho>0}\), \(\fitinline{1-B=\rho}\), \(\fitinline{L-2=\rho}\), \(\fitinline{L-2B=3\rho}\), and \(\fitinline{\zeta>0}\) are proved analytically from \(\fitinline{t>0}\), not asserted at the nominal value zero. The extra historical guard \(\fitinline{\theta_t<436691/500000}\) is valid but no longer needed as an inherited endpoint premise; the preliminary \(\fitinline{7/8}\) bound supplies the full feedback range.

\hypertarget{exact-domain-coverage}{%
\subsubsection{10.4. Exact domain coverage}\label{exact-domain-coverage}}

Every certificate cell is contained in its named original closed row rectangle. If the cell uses an enhanced branch, its whole box is contained in the named complete upward admission window. Each original row is tiled by its assigned cells: their interiors are disjoint and their exact rational areas sum to the original area. The same exact comparison for each edge gives the whole closed density/amplitude domain \(\fitinline{[0,3/4]\times[0,1/2]}\). A finite union of closed rectangles with full area and no overlap of interiors cannot omit a relative-open piece of the domain; any such omitted piece would have positive area. The finite union is closed, so it cannot omit an isolated boundary point either. Thus the area and containment statements certify the full closed cover, not only almost-everywhere coverage.

The edge counts are 448 for oldSmall, 448 for oldSwitch, 448 for newSwitch and 576 for newTop, totaling 1,920 cells. They retain all 1,826 original rows. On every conductor interval the capacity is the minimum of two affine functions and the count is independent of conductor. The physical exponent is consequently convex in conductor. The two endpoint inequalities therefore cover the full interval, as proved in the generic return.

\hypertarget{scope-of-the-evidence}{%
\subsubsection{10.5. Scope of the evidence}\label{scope-of-the-evidence}}

The preserved receipts establish the indicated finite arithmetic properties at the hashes of these unchanged data. The consolidation adds a direct exact check of the changed source-domain inclusion and comparison review of the rewritten arguments/interfaces. It does not replay the inherited validation tower or make old receipts certify new prose. The analytic theorem follows from the manuscript's source, count, physical, complement, feedback and Mellin arguments composed with the finite inequalities just specified.

\hypertarget{the-marked-anchor-rational-endpoint}{%
\subsection{11. The marked-anchor rational endpoint}\label{the-marked-anchor-rational-endpoint}}

This section retains the proved PR141 rational theorem. Its target is R141, while the current exact-contact theorem is completed in Appendix G.

Use the arithmetic conventions, natural zeros and actual Hecke zero supremum \(\fitinline{\beta}\) of \S{}2. Appendices A--B supply \(\fitinline{\beta\le7/8}\). If \(\fitinline{\beta<51/100}\), (T2) follows immediately by the original norm pullback. All further source calls are in the standing case \(\fitinline{\beta\ge51/100}\).

Appendix D proves the complete simultaneous inverse and paired-ordinary anchor contract. Its inverse polynomial retains M\"{o}bius coefficients and all nonzero rows, including principal and fixed-inducing rows. Its marks are whole canonical pools, with one common relative ray twist, independent bounded heights and possibly overlapping pools. The ordinary half retains the original initialKeep predicate, independent profiles and its allowed per-slot twists. Neither statement is a moment for an arbitrary inverse-times-ordinary product.

Appendix F extends the paired analytic step with a variable positive tolerance and no fixed per-stage loss. Appendix E uses a finite induction for each strict curve to obtain the open anchor region \[\fitmath{
P\le r-d,\qquad S\le(1-r)/2+I(r)-d,\qquad
I(r)=11r/16-7/24-\frac5{16}\{7r^4(r-2/3)\}^{1/5},
}\] for \(\fitinline{3/4\le r<7/9}\) and every fixed \(\fitinline{d>0}\), with the full coefficient and analytic contracts of Appendices D--E. Its exact finite lower chain supplies the remaining required interval. This is not a passage of analytic constants through a limit. The early mesh precedes an arbitrary fixed finite slot type; late constants and derivative/height orders may depend on that type.

The tangent count branch and every native continuation are admitted on their entire stated windows. The source/count and physical capacities remain independent. Appendix E's averaged floor is proved on the actual reflected numerator, by the retained empty-selection fourth moment, positive coefficient mass, complete mask sum, original contours and nonzero normalizer. It replaces the stronger pointwise floor premise only where that proof actually applies.

Let \(\fitinline{t=R_{141}=N/D}\), with the exact integers of Appendix E, and let \(\fitinline{\alpha}\) be its degree-nine contact. Its exact source and return data prove the nonnegative scalar return at \(\fitinline{\alpha}\), including the zero-margin critical rectangle. They then retain the fixed positive source cost, conductor extension and row fee \(\fitinline{s_0=10^{-258}}\). The exact rational separation \(\fitinline{t-\alpha>25600s_0}\) pays all three and every complementary term. All conductor intervals, capacity switches and strict source admissions are included, rather than inferred from the contact scalar alone.

For clarity, the closing feedback uses the actual supremum. Suppose \(\fitinline{\beta>t}\), put \(\fitinline{g=\beta-t>0}\) and \(\fitinline{\kappa_\beta=2\beta-1}\). At the chosen target parameter \(\fitinline{K_t=2t-1}\), the old and new complete row exponents are at most \(\fitinline{t-s_0}\). The quantitative count and capacity transport of \S{}\S{}7.4--8.2, with Appendix E's literal branches, raises the new row by at most \(\fitinline{(\kappa_\beta-K_t)/50=g/25}\), and the old row by at most \(\fitinline{(\kappa_\beta-K_t)/8=g/4}\). Hence \[\fitmath{
E_{\rm new}(\kappa_\beta)-\beta\le-s_0-24g/25<0,
\qquad E_{\rm old}(\kappa_\beta)-\beta\le-s_0-3g/4<0.
}\] The actual natural cap \(\fitinline{\kappa_\beta\le3/4}\) and the chosen target parameter \(\fitinline{K_t}\) have separate roles. They have not been equated. Appendix E also pays the actual supported floor with a fixed positive margin. Choose the remaining analytic losses inside these strict budgets. The finite profile and height degrees are fixed before the late positive height exponent. The principal signal, every common-divisor/root sector, Fourier tails and complementary conductors are restored by the retained full probe argument.

The resulting error is smaller by a fixed power than the exact nonzero reciprocal signal. The original all-height Mellin argument therefore improves the actual zero supremum strictly, contradicting its definition. Consequently \(\fitinline{\beta\le t}\). Appendix B.10's norm pullback then proves the retained rational corollary for every positive presented ordinary modulus, with the actual principal pole excluded. This fixed-cost argument asserts no nonvanishing at \(\fitinline{\Re s=t}\). Appendix G supplies the separate variable-cost argument needed for the strict half-plane above the algebraic contact \(\fitinline{\alpha}\). The earlier full source families and alternative stronger suppliers remain available at their original scopes.

\hypertarget{appendix-a.-arithmetic-completed-reflection-and-the-original-moments}{%
\subsection{Appendix A. Arithmetic, completed reflection and the original moments}\label{appendix-a.-arithmetic-completed-reflection-and-the-original-moments}}

This appendix proves the original analytic inputs used above. Sections A.1--A.8 retain the coefficients and full domains of the supplied proofs, including the more general moving-support and puncture families. Their original equation labels and stable anchors are kept: e.g.~(MF3.40) belongs to A.7, and (67) belongs to A.5. These labels are mathematical cross-references inside this file, not requests to read a historical manuscript. No previous optimized endpoint is assumed. The preliminary seven-eighths bound is proved in Appendix B only after these moments have been established.

\hypertarget{a.1-arithmetic-reflection-and-polynomial-moments}{%
\subsubsection{A.1 Arithmetic, reflection and polynomial moments}\label{a.1-arithmetic-reflection-and-polynomial-moments}}

Lattice Poisson summation transforms a sextic row into a finite Gauss coefficient; cubic-theta automorphy reflects its completed cubic coefficient. After fixing phases and zeros, we prove completed energy, inverse and mixed plain moments, and the probe's exact Poisson factorization, without assuming a zero-free boundary.

These specialized arguments reconstruct Hailey Collet's \emph{A zero-free half-plane from a single compensating prime}, revision \texttt{0\allowbreak{}b\allowbreak{}9\allowbreak{}7\allowbreak{}b\allowbreak{}f\allowbreak{}d\allowbreak{}a\allowbreak{}e\allowbreak{}9\allowbreak{}f\allowbreak{}0\allowbreak{}2\allowbreak{}6\allowbreak{}d\allowbreak{}9\allowbreak{}4\allowbreak{}4\allowbreak{}6\allowbreak{}d\allowbreak{}0\allowbreak{}7\allowbreak{}7\allowbreak{}9\allowbreak{}5\allowbreak{}2\allowbreak{}b\allowbreak{}2\allowbreak{}b\allowbreak{}4\allowbreak{}9\allowbreak{}8\allowbreak{}0\allowbreak{}7\allowbreak{}0\allowbreak{}0\allowbreak{}e\allowbreak{}7\allowbreak{}4\allowbreak{}2}, arithmetic conventions and Appendices A--D (with the source's attribution to OpenAI's GPT-6 Pro). The probe and moment framework originates in OpenAI's manuscripts of 30 September and 5 October 2026; their terminal zero-free claims are not inputs.

\hypertarget{cf:classical-inputs}{}

\hypertarget{a.1.1-classical-inputs-and-their-scope}{%
\paragraph{A.1.1 Classical inputs and their scope}\label{a.1.1-classical-inputs-and-their-scope}}

Classical sieve inputs are Goldmakher--Louvel, \emph{A quadratic large sieve inequality over number fields}, Math. Proc. Cambridge Philos. Soc. 154 (2013), 193--212, arXiv:1112.1642v2, Theorem 1.1 and Corollary 1.2, specialized and verified in \protect\hyperlink{cf:eq:B.11}{(27)}, and Heath-Brown, \emph{Kummer's conjecture for cubic Gauss sums}, Israel J. Math. 120 (2000), 97--124, Theorem 2, DOI:10.1007/s11856-000-1273-y, specialized in \protect\hyperlink{cf:eq:B.19}{(36)}. Both require squarefree \textbf{ideals}, not squarefree rational norms. The quadratic sieve requires a fixed ray modulus and row-independent coefficients.

The cubic theta function, Kubota multiplier and cusp expansions are the unconditional formulas in Dunn--Radziwi\l{}\l{}, \emph{Bias in cubic Gauss sums: Patterson's conjecture}, Ann. of Math. 200 (2024), 967--1057, arXiv:2109.07463v3, \S{}5.1, (5.1)--(5.17) and cusp tables. Their needed coefficients, three cusp sequences and archimedean factors appear in the reflection proof. We use neither its conditional prime asymptotics nor conditional improved sieve.

For a primitive finite-order Hecke character \(\fitinline{\chi}\) of \(\fitinline{K=\mathbf Q(\sqrt{-3})}\), of conductor norm \(\fitinline{Q}\), Hecke's theorem continues \(\fitinline{L_K(s,\chi)}\) (entire unless principal) and gives

\hypertarget{cf:hecke-classical}{}

\[\fitmath{
\Lambda(s,\chi)=C_\chi^{s/2}\Gamma(s)L_K(s,\chi),\qquad
\Lambda(s,\chi)=\epsilon_\chi\Lambda(1-s,\overline\chi),\qquad
C_\chi=\frac{3Q}{(2\pi)^2},\quad |\epsilon_\chi|=1.
\qquad\text{(1)}
}\]

For principal \(\fitinline{\chi}\), \(\fitinline{\Lambda}\) has only simple poles at \(\fitinline{0,1}\). Hecke's theorem also gives finite order, polynomial growth on closed vertical strips after removing poles, and absolute Euler convergence for \(\fitinline{\Re s>1}\). The normalization is Goldmakher--Louvel, \S{}3, (3.1)--(3.2), with \(\fitinline{(r_1,r_2,|d_K|)=(0,1,3)}\) and conductor norm absorbed into \(\fitinline{C_\chi}\). Other classical inputs are cubic reciprocity and supplementary laws, finite-field Gauss--Jacobi identities, Fourier inversion, self-dual Poisson summation, Mellin inversion, Stirling and Phragm\'{e}n--Lindel\"{o}f. Their uses, including moving masks, are explicit below and assume no sought zero-free boundary.

\hypertarget{cf:prime-supply-and-boundary}{}

\textbf{Additional classical inputs used below.} The fixed-ray prime ideal theorem is used with the field and ray modulus fixed before the asymptotic scale. For every integral ray modulus \(\fitinline{\mathfrak m}\) of \(\fitinline{K}\), every class \(\fitinline{C}\) of its finite ray-class group, and every fixed finite excluded prime set \(\fitinline{S}\), \[\fitmath{
\#\{\mathfrak p\notin S:(\mathfrak p,\mathfrak m)=1,
 [\mathfrak p]=C,\ N\mathfrak p\le X\}
\sim \operatorname{Li}(X)/h_{\mathfrak m}.
}\] A fixed union \(\fitinline{B}\) of classes has density \(\fitinline{|B|/h_{\mathfrak m}}\). No uniformity in a moving modulus is asserted. For fixed nonnegative, nonzero \(\fitinline{w\in C_c^\infty((1,2))}\), partial summation consequently gives \[\fitmath{
\sum_{\mathfrak p\in B} w(N\mathfrak p/P)(N\mathfrak p)^{-5/6}
\sim \frac{|B|}{h_{\mathfrak m}}\frac{P^{1/6}}{\log P}
 \int_0^\infty w(t)t^{-5/6}\,dt.
}\] Indeed Stieltjes partial summation against the stated counting asymptotic, followed by \(\fitinline{x=Pt}\), gives this formula uniformly on the fixed compact support of \(\fitinline{w}\). It supplies the fixed-ray amplifier pool in Appendix A.7 and the nonzero principal normalizer in Appendix B.3; the latter's all-good-prime pool has density one. Finite excluded primes do not change the asymptotic.

We also use classical Dirichlet nonvanishing on the line \(\fitinline{\Re s=1}\): for every modulus \(\fitinline{q\ge1}\) and Dirichlet character \(\fitinline{\chi\bmod q}\), \(\fitinline{L(s,\chi)\ne0}\) when \(\fitinline{\Re s=1}\), excluding the principal pole at \(\fitinline{s=1}\). Together with absolute Euler convergence for \(\fitinline{\Re s>1}\), this is the boundary input in Appendix B.10. An imprimitive character's deleted factors \(\fitinline{1-\chi^*(p)p^{-s}}\) are nonzero on \(\fitinline{\Re s\ge1}\) because \(\fitinline{|\chi^*(p)p^{-s}|\le p^{-1}<1}\). Boundary nonvanishing is a separate classical theorem; it is not deduced here from the functional equation.

\hypertarget{cf:sec:1.1}{}

\hypertarget{a.1.2-arithmetic-conventions}{%
\paragraph{A.1.2 Arithmetic conventions}\label{a.1.2-arithmetic-conventions}}

Put \(\fitinline{\mathcal O=\mathbf Z[\omega]}\), \(\fitinline{\omega=e^{2\pi i/3}}\), \(\fitinline{\lambda=1+2\omega}\), \(\fitinline{q_a=Na}\). This norm-Euclidean ring has six units; every ideal prime to \(\fitinline{3}\) has a unique primary generator \(\fitinline{\equiv1\pmod{3\mathcal O}}\). Ideal sums count it once; element sums retain units. Ideal and lattice counts through norm \(\fitinline{X\ge1}\) are \(\fitinline{O(X)}\).

For each target character fix, before varying scales or rows, a finite \(\fitinline{S}\) containing conductor primes, primes above \(\fitinline{6}\), and a small-prime cutoff. Columns avoid \(\fitinline{S}\) unless specified; primes outside \(\fitinline{S}\) are good. Superscript \(\fitinline{S}\) deletes these Euler factors, a star means squarefree ideals, and \(\fitinline{\sum^{(6)}}\) means element rows with valuations \(\fitinline{0,\ldots,5}\).

At a good prime \(\fitinline{\mathfrak p}\), \(\fitinline{\chi_{\mathfrak p}(v)}\) is the sextic residue symbol, the sixth root of unity \(\fitinline{\equiv v^{(q_{\mathfrak p}-1)/6}}\) at units, extended multiplicatively in the denominator. \textbf{All integer character powers vanish at nonunits}, including zero and negative powers; only an absent modulus-one factor is identically one. Write

\hypertarget{cf:eq:1.1}{}

\[\fitmath{
\psi_u(n)=\chi_n(u),\qquad \alpha(a)=a/|a|,\qquad
 e(z)=\exp\!\left(2\pi i\operatorname{Tr}(z/\lambda)\right),\qquad
 d\mu(z)=\frac2{\sqrt3}\,dx\,dy.
\qquad\text{(2)}
}\]

Here \(\fitinline{\operatorname{Tr}z=z+\bar z}\); \(\fitinline{e(zy)}\) is self-dual on \(\fitinline{\mathcal O}\), with Fourier kernel \(\fitinline{e(-zy)}\). For a zero-extended character \(\fitinline{\chi}\), put

\hypertarget{cf:eq:1.2}{}

\[\fitmath{
g_\chi(c,h)=\sum_{d\bmod c}\chi(d)e(hd/c),\qquad
 \gamma_j(c)=q_c^{-1/2}g_{\chi_c^j}(c,1)\quad(c\text{ squarefree}).
\qquad\text{(3)}
}\]

The inverse-character sum \(\fitinline{\gamma_{-1}}\) does not mean conjugation. For odd \(\fitinline{a+b\omega}\), define

\hypertarget{cf:eq:1.3}{}

\[\fitmath{
\Gamma_{\rm q}(a+b\omega)=\tfrac12(1+i^{-b}+i^a+i^{b-a}),\quad
 G(c)=\overline{\chi_c(4)}\Gamma_{\rm q}(c),\quad
 R(v,w)=\frac{\Gamma_{\rm q}(vw)}{\Gamma_{\rm q}(v)\Gamma_{\rm q}(w)}.
\qquad\text{(4)}
}\]

Lemma \protect\hyperlink{cf:finite-phases}{A.2.1} proves that \(\fitinline{G}\) is a finite-ray unit phase, \(\fitinline{R}\) a symmetric sign-valued bicharacter, and \(\fitinline{\chi_b(a)=R(a,b)\chi_a(b)}\) for coprime primary indices. Expand finite-ray functions into genuine characters before using multiplicativity.

The ideal character \(\fitinline{\psi_u}\) is the arithmetic-Frobenius Kummer character of \(\fitinline{K(u^{1/6})/K}\), retaining even redundant zeros at primes of \(\fitinline{u}\). At a good prime with valuation \(\fitinline{j}\), tame inertia has exponent \(\fitinline{\pm j\bmod6}\) and conductor exponent one unless \(\fitinline{6\mid j}\). At \(\fitinline{S}\), Hensel makes deep principal units sixth powers. Thus fixed finite-ray twists of \(\fitinline{\psi_u^{\pm1}}\) have conductor norm \(\fitinline{O_S(q_u)}\); every nonunit sixth-power-free row prime to \(\fitinline{S}\) is nonprincipal.

\hypertarget{cf:sec:1.2}{}

\hypertarget{a.1.3-smooth-weights-and-uniformity}{%
\paragraph{A.1.3 Smooth weights and uniformity}\label{a.1.3-smooth-weights-and-uniformity}}

An annular weight is smooth and compactly supported in \(\fitinline{(0,\infty)}\), with \(\fitinline{\widehat W(s)=\int_0^\infty W(y)y^s\,dy/y}\). Constants may depend on fixed characters, supports and positive losses, never moving arithmetic labels. Real exponents lie in fixed bounded sets. All estimates use finitely many smooth seminorms; multiplying a test by \(\fitinline{y^{b+it}}\), for bounded \(\fitinline{b}\), costs \(\fitinline{(1+|t|)^{O(1)}}\).

Profiles are common; arithmetic coefficients are row-independent. Integration by parts gives a compact joint logarithmic profile a Fourier expansion with finite weighted \(\fitinline{L^1}\) norm. Minkowski and Cauchy reduce its row norm to separated unit-modulus norm powers. Thus if \(\fitinline{C_j(W)=\sum_nc_{j,n}W(x_{j,n})}\) satisfies \(\fitinline{\sum_j|C_j(W)|^2\le M\|W\|_{C^m(I')}^2}\) for every common test in \(\fitinline{I'}\), then kernels supported in fixed \(\fitinline{I^2\Subset(I')^2}\) satisfy

\hypertarget{cf:eq:1.4}{}

\[\fitmath{
\left|\sum_jb_j\sum_{n_1,n_2}c_{j,n_1}\overline{c_{j,n_2}}
 K_j(x_{j,n_1},x_{j,n_2})\right|
 \ll M\sup_j\|K_j\|_{C^{2m+4}},\qquad |b_j|\le1.
\qquad\text{(5)}
}\]

Fourier decay \(\fitinline{(1+|t_1|)^{-m-2}(1+|t_2|)^{-m-2}}\) and the common-test bounds leave integrable powers \(\fitinline{-2}\). Sobolev on unit intervals in log scale, real power and height similarly permits rowwise smooth parameters after summing common-profile bounds. Choose all internal smooth orders before the external height cutoff.

\hypertarget{cf:finite-arithmetic}{}

\hypertarget{a.2-finite-arithmetic-for-both-transforms}{%
\subsubsection{A.2 Finite arithmetic for both transforms}\label{a.2-finite-arithmetic-for-both-transforms}}

\hypertarget{cf:sec:A.1}{}

\hypertarget{a.2.1-gauss-phases}{%
\paragraph{A.2.1 Gauss phases}\label{a.2.1-gauss-phases}}

\textbf{Lemma A.2.1} (Finite Gauss phases). \emph{\hypertarget{cf:finite-phases}{} For squarefree \(\fitinline{n}\), and coprime squarefree \(\fitinline{a,b}\), the conventions of Section~\protect\hyperlink{cf:sec:1.1}{A.1.2} give}

\hypertarget{cf:eq:A.1}{}

\[\fitmath{
\begin{gathered}
\chi_b(a)=R(a,b)\chi_a(b),\qquad R(a,b)=R(b,a)\in\{\pm1\},\quad (a,b)=1,\\
\gamma_2(n)^3=\mu(n)\alpha(n),\qquad
\gamma_1(n)\gamma_2(n)=\mu(n)\alpha(n)G(n),\\
\mu(n)\gamma_{-1}(n)=\chi_n(-1)G(n)^{-1}\overline{\alpha(n)}\gamma_2(n),\\
a_\nu(ab)=a_\nu(a)a_\nu(b)\chi_b(a)^4,
\qquad a_\nu(n)=\overline{\alpha(n)}\gamma_2(n)\nu(n),\quad (a,b)=1.
\end{gathered}
\qquad\text{(6)}
}\]

\emph{Proof.} At a prime \(\fitinline{p}\), put \(\fitinline{J(A,B)=\sum_xA(x)B(1-x)}\) with zero extensions. Substitution \(\fitinline{u=2x-1}\) gives \(\fitinline{\chi_p(4)J(\chi_p,\chi_p)=J(\chi_p,\chi_p^3)}\); Gauss--Jacobi and \(\fitinline{\gamma_2(p)\gamma_4(p)=1}\) give

\[\fitmath{
\gamma_1(p)\gamma_2(p)
=\overline{\chi_p(4)}\gamma_3(p)\gamma_2(p)^3.
}\]

The Jacobi sum \(\fitinline{J(\chi_p^2,\chi_p^2)}\) has norm \(\fitinline{q_p}\) and is divisible by \(\fitinline{p}\), by expanding \(\fitinline{x^{(q_p-1)/3}(1-x)^{(q_p-1)/3}}\) modulo \(\fitinline{p}\) and summing powers. Its unit follows modulo three: \(\fitinline{\prod_{x\ne0,1}x(1-x)=1}\) implies \(\fitinline{\sum j_x\equiv0\pmod3}\) for character values \(\fitinline{\omega^{j_x}}\), so the sum is \(\fitinline{-1\pmod3}\). Since \(\fitinline{p\equiv1\pmod3}\), it equals \(\fitinline{-p}\), giving \(\fitinline{\gamma_2(p)^3=-\alpha(p)}\).

For odd \(\fitinline{c}\), Gaussian-regularized Poisson summation gives

\[\fitmath{
|c|^{-1}\sum_{x\bmod c}e(x^2/c)
=\tfrac12\sum_{y\bmod2\mathcal O}e(-cy^2/4)=\Gamma_{\rm q}(c).
}\]

For regulator \(\fitinline{e^{-\pi\eta|z|^2}}\), the transformed factor has \(\fitinline{r_\eta=1+3q_c\eta^2/16}\). Replacing \(\fitinline{r_\eta}\) by one costs \(\fitinline{O_c(1)}\); divide by Gaussian mass \(\fitinline{\asymp\eta^{-1}}\) and let \(\fitinline{\eta\downarrow0}\). The square classes \(\fitinline{1,-1,\lambda,-\lambda}\) have values \(\fitinline{1,1,i,-i}\), whose quotient is the symmetric bicharacter

\[\fitmath{
R((-1)^e\lambda^f,(-1)^g\lambda^h)=(-1)^{eh+fg+fh}.
}\]

Grouping by \(\fitinline{y=x^2}\) gives \(\fitinline{\Gamma_{\rm q}(p)=\gamma_3(p)}\); CRT and cubic reciprocity give sextic reciprocity with \(\fitinline{R}\). Also \(\fitinline{\chi_c(4)=(-2/c)_3=(c/(-2))_3}\) is a fixed character modulo two. Hence \(\fitinline{G(c)=\overline{\chi_c(4)}\Gamma_{\rm q}(c)}\) is a fixed ray phase, and

\hypertarget{cf:eq:A.2}{}

\[\fitmath{
\begin{gathered}
G(ab)=G(a)G(b)R(a,b),\qquad R(a,a)=\chi_a(-1),\\
G(ba^{-1})=\chi_a(-1)\overline{G(a)}G(b)R(a,b).
\end{gathered}
\qquad\text{(7)}
}\]

CRT extends the prime identities to squarefree indices; \(\fitinline{\gamma_2(ab)=\gamma_2(a)\gamma_2(b)\chi_a(b)^2\chi_b(a)^2}\) gives the \(\fitinline{a_\nu}\) identity. Finally \(\fitinline{\gamma_1(n)\gamma_{-1}(n)=\chi_n(-1)}\) gives inverse conversion in \protect\hyperlink{cf:eq:A.1}{(6)}. All equalities retain nonunit zeros. \ensuremath{\square}

\hypertarget{cf:sec:A.4}{}

\hypertarget{a.2.2-finite-fourier-correlations}{%
\paragraph{A.2.2 Finite Fourier correlations}\label{a.2.2-finite-fourier-correlations}}

For primary moduli distinguish the normalized Gauss sum from \(\fitinline{G}\):

\[\fitmath{
\mathfrak g(a,h)=q_a^{-1/2}\sum_{x\bmod a}\chi_a(x)e(hx/a),\qquad \mathfrak g(1,h)=1.
}\]

For a good prime \(\fitinline{p}\), \(\fitinline{P=q_p}\), \(\fitinline{j\ge1}\), lifting \(\fitinline{x=x_0+py}\) gives

\[\fitmath{
|\mathfrak g(p^j,h)|=P^{(j-1)/2}\mathbf1_{v_p(h)=j-1}\quad(6\nmid j),
}\]

\hypertarget{cf:eq:A.12}{}

\[\fitmath{
\mathfrak g(p^j,h)=P^{j/2}\mathbf1_{p^j\mid h}
 -P^{j/2-1}\mathbf1_{p^{j-1}\mid h}\quad(6\mid j).
\qquad\text{(8)}
}\]

Here \(\fitinline{v_p(0)=\infty}\); a sixth-power character retains its coprimality mask.

Define the full correlation

\hypertarget{cf:eq:A.13}{}

\[\fitmath{
F(u,v;j)=\sum_{\substack{x\bmod u,\ y\bmod v\\vx-uy\equiv j\pmod{uv}}}
 \chi_u(x)\overline{\chi_v(y)}.
\qquad\text{(9)}
}\]

Additive orthogonality gives, with negative Fourier sign,

\hypertarget{cf:eq:A.14}{}

\[\fitmath{
F(u,v;j)=\frac1{\sqrt{q_uq_v}}
 \sum_{h\bmod uv}\mathfrak g(u,h)\overline{\mathfrak g(v,h)}e(-jh/(uv)).
\qquad\text{(10)}
}\]

At \(\fitinline{j=0}\), both nonzero residues are units, forcing \(\fitinline{u\mid v}\) and \(\fitinline{v\mid u}\). Hence

\hypertarget{cf:eq:A.15}{}

\[\fitmath{
F(u,v;0)=\mathbf1_{u=v}\phi(u).
\qquad\text{(11)}
}\]

For complete common-support extraction, let

\[\fitmath{
u=Da,\quad v=Eb,\quad (a,b)=1,\quad (ab,DE)=1,
}\]

where every prime of \(\fitinline{D}\) also divides \(\fitinline{E}\), and conversely. Then

\hypertarget{cf:eq:A.16}{}

\[\fitmath{
F(Da,Eb;j)=F(D,E;j)R(a,E)R(b,D)R(a,b)
 \chi_a(j)\overline{\chi_b(-j)}.
\qquad\text{(12)}
}\]

At primes of \(\fitinline{a,b}\), reduction gives \(\fitinline{\chi_a(j)\overline{\chi_a(Eb)}}\) and \(\fitinline{\overline{\chi_b(-j)}\chi_b(Da)}\). At common primes substitute \(\fitinline{x'=bx}\), \(\fitinline{y'=ay}\), leaving \(\fitinline{\overline{\chi_D(b)}\chi_E(a)}\). The three unit quotients give the three \(\fitinline{R}\) factors. All inversions are at units; nonunit \(\fitinline{j}\) gives the same zero on both sides.

From \protect\hyperlink{cf:eq:A.13}{(9)}, at common multiplicities \(\fitinline{i,j}\) the correlation vanishes unless \(\fitinline{p^{\min(i,j)}\mid h}\). For divided frequency \(\fitinline{k}\) and \(\fitinline{i=j}\), its local value is

\hypertarget{cf:eq:A.17}{}

\[\fitmath{
P^{i-1}\begin{cases}
P-1,&p\mid k,\\
-1,&p\nmid k,\ 6\nmid i,\\
P-2,&p\nmid k,\ 6\mid i.
\end{cases}
\qquad\text{(13)}
}\]

For \(\fitinline{i>j}\), it vanishes unless \(\fitinline{6\mid j}\) and \(\fitinline{p\nmid k}\), in which case it is

\hypertarget{cf:eq:A.18}{}

\[\fitmath{
P^{j-1}(P-1)\chi_p(k)^{i-j}.
\qquad\text{(14)}
}\]

For \(\fitinline{j>i}\), conjugate and negate the frequency. Indeed dividing the common power fixes one variable when multiplicities differ; equal multiplicities give a character sum over \(\fitinline{\mathbb F_P^\times\setminus\{1\}}\).

\textbf{Artificial extension rule.} \protect\hyperlink{cf:eq:A.16}{(12)} need not hold on noncoprime residual pairs. Extend its right side to all individually allowed \(\fitinline{a,b}\) and first insert

\hypertarget{cf:eq:A.19}{}

\[\fitmath{
\sum_{(a,b)=1}c(a,b)F(Da,Eb;j)
=\sum_t\mu(t)\sum_{t\mid a,b}c(a,b)\widetilde F_{D,E}(a,b;j).
\qquad\text{(15)}
}\]

The complete M\"{o}bius sum removes artificial pairs before factorwise estimates. Its extra primes are outside genuine common support and inherit none of its valuation restrictions.

\hypertarget{cf:completed-plan}{}

\hypertarget{a.2.3-the-completed-coefficient-and-the-two-transform-plan}{%
\paragraph{A.2.3 The completed coefficient and the two-transform plan}\label{a.2.3-the-completed-coefficient-and-the-two-transform-plan}}

For a fixed finite-ray character \(\fitinline{\nu}\) set

\hypertarget{cf:completed-definition}{}

\[\fitmath{
\mathcal C_V(N;k,\nu)=\sum_{c\ {\rm sf},b}
\frac{\overline{\alpha(cb^3)}\gamma_2(c)\nu(cb^3)\chi_{cb^3}(k)}{\sqrt{q_c}q_b}
V(q_cq_b^3/N).
\qquad\text{(16)}
}\]

The ideals \(\fitinline{c,b}\) may share primes. Exact cube inversion removes the completing \(\fitinline{b}\)-sum for the inverse moment. The plain moment instead closes a two-factor family with fixed displayed moving radical: its width includes that entire radical, even at exponent zero. Later marked coefficients must use the appropriate closure.

\hypertarget{cf:sec:B}{}

\hypertarget{a.3-cubic-theta-reflection-and-completed-energy}{%
\subsubsection{A.3 Cubic-theta reflection and completed energy}\label{a.3-cubic-theta-reflection-and-completed-energy}}

\hypertarget{cf:sec:B.1}{}

\hypertarget{a.3.1-cubic-reflection}{%
\paragraph{A.3.1 Cubic reflection}\label{a.3.1-cubic-reflection}}

For \(\fitinline{k=u_0k_Sk_{\rm good}}\), the unit, \(\fitinline{S}\)-valuations modulo six and good row's ray class give finitely many characters \(\fitinline{\nu(A)\chi_A(u_0k_S)R(A,k_{\rm good})}\). Remaining dependence is in local good-prime exponents; the reflection's fixed modulus contains no moving prime.

Let \(\fitinline{\nu_0}\) lie in a fixed finite set of ray characters with prescribed zeros. At finitely many good primes choose \(\fitinline{j_{\mathfrak p}\in\{0,\ldots,5\}}\) and set \(\fitinline{\Psi(A)=\nu_0(A)\prod_{\mathfrak p}\chi_{\mathfrak p}(A)^{j_{\mathfrak p}}}\). Define

\hypertarget{cf:eq:B.1}{}

\[\fitmath{
\mathcal C_V(X;\Psi)=
\sum_{n\ \mathrm{sf},b}
\frac{\gamma_2(n)\overline{\alpha(nb^3)}\Psi(nb^3)}{\sqrt{q_n}\,q_b}
V(q_nq_b^3/X).
\qquad\text{(17)}
}\]

\textbf{Proposition A.3.1} (Cubic reflection). \emph{\hypertarget{cf:reflection}{} For annular \(\fitinline{V}\), or the Gaussian used in the probe, this sum is a divisor-bounded finite sum of branches \hypertarget{cf:eq:B.2}{}}

\[\fitmath{
\zeta\sum_{\mu\ne0}\frac{d(\mu)\alpha(\mu)\vartheta(\lambda^4\mu)}{\sqrt{q_\mu}}
\prod_{\mathfrak p\in A}B_{\mathfrak p,j_{\mathfrak p}}(\lambda^4\mu)
V^\sharp\!\left(\frac{q_\mu X}{q_{c_F}^{2}q_r^{2}}\right),
\quad r=\prod_{\mathfrak p\in A}\mathfrak p,
\qquad\text{(18)}
}\]

\emph{where \(\fitinline{c_F}\) has finitely many fixed possibilities and the local factors are \hypertarget{cf:eq:B.3}{}}

\[\fitmath{
B_{\mathfrak p,j}(x)=
\begin{cases}
\chi_{\mathfrak p}(x)^{-j-2},&j\notin\{0,4\},\\
q_{\mathfrak p}^{-1/2}(-1+q_{\mathfrak p}1_{\mathfrak p\mid x}),&j=4,\\
q_{\mathfrak p}^{-1/2}\chi_{\mathfrak p}(x)^{-2},&j=0\text{ active}.
\end{cases}
\qquad\text{(19)}
}\]

\emph{Every nonzero exponent is active. An exponent-zero prime additionally has an inactive branch, scalar \(\fitinline{1-q_{\mathfrak p}^{-1}}\), omitting it from \(\fitinline{r}\).}

\emph{The branch scalar, needed when the prime varies, is \hypertarget{cf:eq:B.4}{}}

\[\fitmath{
\zeta=-\frac{i}{81}\overline{\alpha(c)}^{\,2}\widehat\phi(h_0)
\overline{\kappa_F}
\prod_{\mathfrak p\in I}(1-q_{\mathfrak p}^{-1})
\prod_{\mathfrak p\in A}\chi_{\mathfrak p}(\sigma_{\mathfrak p})^{-2}
\omega_{\mathfrak p,j_{\mathfrak p}},\qquad c=c_Fr,
\qquad\text{(20)}
}\]

\emph{Here \(\fitinline{I}\) is the inactive set, \(\fitinline{\widehat\phi}\) the normalized fixed-modulus Fourier coefficient below, and}

\[\fitmath{
\begin{gathered}
\sigma_{\mathfrak p}=\lambda^2c/\mathfrak p,\qquad
\epsilon_{\mathfrak p}=-\bigl((\lambda^3c/\mathfrak p)\sigma_{\mathfrak p}\bigr)^{-1}
\pmod{\mathfrak p},\\
\omega_{\mathfrak p,j}=
\begin{cases}
\chi_{\mathfrak p}(-1)^j\gamma_j(\mathfrak p)\gamma_{j+2}(\mathfrak p)
\chi_{\mathfrak p}(\epsilon_{\mathfrak p})^{-j-2},&j\notin\{0,4\},\\
\gamma_4(\mathfrak p),&j=4,\\
-\gamma_2(\mathfrak p)\chi_{\mathfrak p}(\epsilon_{\mathfrak p})^{-2},&j=0.
\end{cases}
\end{gathered}
}\]

\emph{The inverse is in the residue field at a unit. Each \(\fitinline{|\omega_{\mathfrak p,j}|=1}\), hence \(\fitinline{|\zeta|\le1/81}\). The three fixed cusp sequences satisfy \hypertarget{cf:eq:B.5}{}}

\[\fitmath{
\mu=u\lambda^mnb^3,\qquad m\ge-4,
\qquad |d(u\lambda^mnb^3)|\le27\,3^{m/6}\sqrt{q_b}.
\qquad\text{(21)}
}\]

\emph{The dual primary ideals \(\fitinline{n,b}\) avoid \(\fitinline{\lambda}\); \(\fitinline{n}\) is squarefree. They may share primes and meet \(\fitinline{S}\), without the primal \(\fitinline{S}\)-exclusion. Fixing the active radical modulo one fixed ray modulus makes \(\fitinline{d,\vartheta,c_F}\) independent of its moving primes; only a bounded scalar and the displayed local characters vary.}

\emph{With \(\fitinline{\widehat V(t)=\int_0^\infty V(x)x^t\,dx/x}\), the exact transform is \hypertarget{cf:eq:B.6}{}}

\[\fitmath{
V^\sharp(y)=\frac1{2\pi i}\int_{(0)}\widehat V(-t)
\prod_{\pm}\frac{\Gamma(1+t\pm1/6)}{\Gamma(1-t\pm1/6)}
\left(\frac{(2\pi)^4y}{27}\right)^{-t}\,dt.
\qquad\text{(22)}
}\]

\emph{For \(\fitinline{V(y)=y^{1/2}W(y)}\), with annular \(\fitinline{W}\), every fixed \(\fitinline{A,j}\) has a finite \(\fitinline{J}\) such that \hypertarget{cf:eq:B.7}{}}

\[\fitmath{
|(y\partial_y)^jV^\sharp(y)|
\ll_{A,j}\|W\|_{C^J}\min(y^{1/4},y^{-A}).
\qquad\text{(23)}
}\]

\emph{The same bound holds for \(\fitinline{V_G}\), with a fixed constant in place of the annular seminorm.}

\emph{Proof.} Use the classical cubic-theta automorphy and cusp expansions of \S{}\protect\hyperlink{cf:classical-inputs}{A.1.1}. Put \(\fitinline{\check e(z)=e(\lambda z)}\) and \(\fitinline{L(t)=\left(\begin{smallmatrix}1&0\\t&1\end{smallmatrix}\right)}\). On \(\fitinline{\Gamma_1(3)=\{g\in\operatorname{SL}_2(\mathcal O):g\equiv I\pmod3\}}\), the theta multiplier is \(\fitinline{\kappa(\left(\begin{smallmatrix}a&b\\c&d\end{smallmatrix}\right))=(c/a)_3}\) for \(\fitinline{c\ne0}\), otherwise one; extend to \(\fitinline{\operatorname{SL}_2(\mathbf Z)\Gamma_1(3)}\) with value one on \(\fitinline{\operatorname{SL}_2(\mathbf Z)}\). The conjugate theta coefficient at \(\fitinline{\lambda^{-3}nb^3}\) is \(\fitinline{3^{5/2}|b|\chi_n(\lambda)^{-2}\gamma_2(n)}\), its constant term \(\fitinline{3^{5/2}v^{2/3}/2}\). At \(\fitinline{I,L(\lambda),L(-\lambda)}\) the expansions are

\hypertarget{cf:eq:B.8}{}

\[\fitmath{
\sigma_Hv^{2/3}+\sum_{\mu\ne0}d_H(\mu)vK_{1/3}(4\pi|\mu|v)\check e(\mu z).
\qquad\text{(24)}
}\]

These cusp expansions give \protect\hyperlink{cf:eq:B.5}{(21)}; their twisted reflection below uses no prime asymptotic.

Fix \(\fitinline{L}\) containing primary and supplementary conditions and fixed conductors. Set \(\fitinline{\phi(x)=\chi_x(\lambda)^2\nu_0(x)}\) with prescribed primary and nonunit zeros, and \(\fitinline{\widehat\phi(h_0)=q_L^{-1}\sum_{x\bmod L}\phi(x)e(-h_0x/L)}\). Expand modulo \(\fitinline{L}\). At a moving prime,

\hypertarget{cf:eq:B.9}{}

\[\fitmath{
q^{-1}\sum_{x\bmod\mathfrak p}\chi_{\mathfrak p}(x)^j e(-hx/\mathfrak p)
=\begin{cases}
q^{-1/2}\tau^-_{\mathfrak p,j}\chi_{\mathfrak p}(h)^{-j},&j\ne0, h\ne0,\\
0,&j\ne0, h=0,\\
-q^{-1},&j=0, h\ne0,\\
1-q^{-1},&j=0, h=0,
\end{cases}
\qquad\text{(25)}
}\]

Here \(\fitinline{|\tau^-_{\mathfrak p,j}|=1}\). An active branch translates by \(\fitinline{z_h=\lambda^2(h_0/L+\sum_{\mathfrak p\mid r}h_{\mathfrak p}/\mathfrak p)}\), all \(\fitinline{h_{\mathfrak p}\ne0}\).

Reduce \(\fitinline{\lambda^2h_0/L=a_F/c_F}\), unit-normalized with \(\fitinline{a_F\equiv1\pmod3}\) if \(\fitinline{\lambda\mid c_F}\), otherwise \(\fitinline{c_F\equiv1\pmod3}\). Then

\[\fitmath{
c=c_Fr,\qquad a=a_Fr+\lambda^2c_F\sum_{\mathfrak p\mid r}(r/\mathfrak p)h_{\mathfrak p},
\qquad z_h=a/c.
}\]

Fix \(\fitinline{r\bmod M^2}\), \(\fitinline{M=\lambda^{12}L^4}\), and lift \(\fitinline{h_{\mathfrak p}\equiv0\pmod{M^2}}\). Complete \(\fitinline{(a,c)}\) to \(\fitinline{g=\left(\begin{smallmatrix}a&b_0\\c&\delta\end{smallmatrix}\right)}\) with \(\fitinline{\delta\equiv a^{-1}}\) to the required depth at primes of \(\fitinline{c_F}\), \(\fitinline{\delta\equiv0}\) at primes of \(\fitinline{M}\) outside \(\fitinline{c_F}\), and \(\fitinline{\delta\equiv a^{-1}\pmod{\mathfrak p}}\) for \(\fitinline{\mathfrak p\mid r}\). Choose

\[\fitmath{
H=\begin{cases}
I,&3\mid c,\\
L(u_0),\quad u_0\in\{\lambda,-\lambda\},\ u_0\equiv c\pmod3,
&v_\lambda(c)=1,\\
EL(-u_0),\quad u_0\equiv a\pmod3,&(c,\lambda)=1,
\end{cases}
\qquad E=\begin{pmatrix}0&-1\\1&0\end{pmatrix}.
}\]

In the last case \(\fitinline{u_0}\) ranges over fixed representatives of \(\fitinline{\mathcal O/3\mathcal O}\), not merely \(\fitinline{\{\pm\lambda\}}\). These congruences give \(\fitinline{gH^{-1}\equiv I\pmod3}\); trivial multipliers on \(\fitinline{\operatorname{SL}_2(\mathbf Z)}\) and translations by \(\fitinline{\mathbf Z+3\mathcal O}\) reduce to the same three cusps.

Multiplier reciprocity gives \(\fitinline{\kappa(gH^{-1})=\kappa_F\prod_{\mathfrak p\mid r}\chi_{\mathfrak p}(a)^2}\). The fixed-sector factors are respectively \(\fitinline{(c_F/a)_3}\), \(\fitinline{(-u_0/(a-u_0b_0))_3((c_F/u_0)/a)_3}\), \(\fitinline{(a/c_F)_3}\), independent of the residual column.

The coordinate formula

\hypertarget{cf:eq:B.10}{}

\[\fitmath{
g^{-1}(a/c+z,v)=
\left(-\frac\delta c-\frac{\bar z}{c^2(|z|^2+v^2)},
\frac v{q_c(|z|^2+v^2)}\right)
\qquad\text{(26)}
}\]

shows that \(\fitinline{\partial_{\bar z}}\) at zero kills the constant term and equals \(\fitinline{-(cv)^{-2}\partial_{z'}}\). For \(\fitinline{x=\lambda^4\mu}\), \(\fitinline{f=\lambda^3c_F}\), additive CRT gives

\[\fitmath{
\check e(-\delta\mu/c)=\vartheta(x)
\prod_{\mathfrak p\mid r}e(\epsilon_{\mathfrak p}h_{\mathfrak p}^{-1}x/\mathfrak p),
\quad \vartheta(x)=e(-\delta_Fr^{-1}x/f),
\quad\epsilon_{\mathfrak p}\equiv-(\lambda^5c_F^2(r/\mathfrak p)^2)^{-1}\pmod{\mathfrak p}.
}\]

The conjugate multiplier contributes \(\fitinline{\chi_{\mathfrak p}(h_{\mathfrak p})^{-2}}\); substituting \(\fitinline{y=h_{\mathfrak p}^{-1}}\) gives exponent \(\fitinline{j+2}\) and exactly the Gauss/Ramanujan factors \protect\hyperlink{cf:eq:B.3}{(19)}, with their zeros. Their normalized scalars are \(\fitinline{\omega_{\mathfrak p,j}}\) and \protect\hyperlink{cf:eq:B.4}{(20)}, recording the entire active conductor's dependence.

To check the archimedean normalization, define

\[\fitmath{
\mathcal D_\Psi(s)=\sum_{n\ {\rm sf},b}\gamma_2(n)\overline{\alpha(nb^3)}\Psi(nb^3)q_n^{-s}q_b^{1/2-3s}.
}\]

Multiply the conjugate theta \(\fitinline{\mu}\)-coefficient by \(\fitinline{\phi(\lambda^3\mu)\prod_{\mathfrak p}\chi_{\mathfrak p}(\lambda^3\mu)^{j_{\mathfrak p}}}\), obtaining \(\fitinline{\Theta_\Psi}\). Differentiation supplies \(\fitinline{2\pi i\overline\mu}\), and

\[\fitmath{
\int_0^\infty v^{2s}K_{1/3}(4\pi|\mu|v)\,dv
=\frac{2^{2s-1}\Gamma(s+1/3)\Gamma(s+2/3)}{(4\pi|\mu|)^{2s+1}}.
}\]

Consequently, initially for \(\fitinline{\Re s>1}\),

\hypertarget{cf:direct-mellin-identity}{}

\[\fitmath{
\int_0^\infty\partial_{\bar z}\Theta_\Psi(0,v)v^{2s-1}\,dv
=\frac{3^{5/2}}4\left(\frac{27}{(2\pi)^2}\right)^s
\Gamma(s+1/3)\Gamma(s+2/3)\mathcal D_\Psi(s).
}\]

Apply \protect\hyperlink{cf:eq:B.10}{(26)} and change height to \(\fitinline{1/(q_cv)}\), contributing \(\fitinline{-c^{-2}q_c^{1-2t}}\), with \(\fitinline{|c^{-2}q_c|=1}\). Division by the direct Mellin factor and \(\fitinline{t\mapsto-t}\) give \protect\hyperlink{cf:eq:B.6}{(22)} and scalar \(\fitinline{1/81}\). Differentiated cusp expansions decay at both endpoints, so their Mellin integrals are entire. Absolute convergence on the direct/reflected half-planes, Stirling and Phragm\'{e}n--Lindel\"{o}f give intervening polynomial vertical bounds. Rapid test decay permits shifting to \(\fitinline{\Re t=-1/4}\) without gamma poles, giving small-\(\fitinline{y}\) decay; arbitrary positive lines give rapid large-\(\fitinline{y}\) decay. This proves \protect\hyperlink{cf:eq:B.2}{(18)}--\protect\hyperlink{cf:eq:B.7}{(23)}, including finite seminorm control. \ensuremath{\square}

\hypertarget{cf:sec:B.2}{}

\hypertarget{a.3.2-unmarked-energy-including-a-frozen-label}{%
\paragraph{A.3.2 Unmarked energy, including a frozen label}\label{a.3.2-unmarked-energy-including-a-frozen-label}}

We use the quadratic large sieve (Goldmakher--Louvel, Theorem 1.1 and Corollary 1.2) in the form

\hypertarget{cf:eq:B.11}{}

\[\fitmath{
\sum_{q_R\le K}^{*}\left|\sum_{q_n\le N}^{*}c_n\chi_R(n)^3\right|^2
\ll_\varepsilon(KN)^\varepsilon(K+N)\sum_n|c_n|^2.
\qquad\text{(27)}
}\]

Coefficients are \(\fitinline{R}\)-independent. The primary quadratic symbol has primitive conductor \(\fitinline{R\lambda^{e(R)}}\), \(\fitinline{e(R)=(q_R-1)/2\pmod2}\); primary adjustment contributes that power of the character modulo \(\fitinline{\lambda}\). A primary square class modulo four fixes \(\fitinline{e(R)}\). Coprime indices in it have primitive product conductor their product and reciprocity factor \(\fitinline{R(\cdot,\cdot)}\), verifying the sieve's Hecke-family hypotheses. Finite ray partition and transpose duality give \protect\hyperlink{cf:eq:B.11}{(27)}. Row restrictions decrease its positive sum; column restrictions are zero coefficients.

For a squarefree ideal \(\fitinline{f}\), set

\[\fitmath{
\Psi_k(a)=\nu(a)\chi_a(k)\chi_a(f)^4,
}\]

and, with \(\fitinline{V(y)=y^{1/2}W(y)}\), define

\hypertarget{cf:eq:B.12}{}

\[\fitmath{
T(X;k,f)=\sum_{n\ \mathrm{sf},b}
\frac{\overline{\alpha(n)}\gamma_2(n)\Psi_k(n)
\overline{\alpha(b)}^{3}\Psi_k(b)^3}{\sqrt{q_n}\,q_b}
V(q_nq_b^3/X).
\qquad\text{(28)}
}\]

\textbf{Proposition A.3.2} (Completed energy with a frozen label). \emph{\hypertarget{cf:completed-energy}{} For every \(\fitinline{\varepsilon>0}\), \hypertarget{cf:eq:B.13}{}}

\[\fitmath{
\sum_{0<q_k\le H}|T(X;k,f)|^2
\ll(HXq_f)^\varepsilon\|W\|_{C^J}^2
\left(H+\frac{H^2q_f}{X}\right),\qquad H,X\ge1.
\qquad\text{(29)}
}\]

\emph{The constants and smooth order are uniform in \(\fitinline{f}\).}

\emph{Proof.} Allow an arbitrary ideal \(\fitinline{g}\) in place of \(\fitinline{f}\) and squarefree rows \(\fitinline{u_0s}\). Write \(\fitinline{s=tk_0}\), \(\fitinline{t\mid gS}\), \(\fitinline{H_0=H/q_t}\), \(\fitinline{(k_0,gS)=1}\). Residual exponent one turns \protect\hyperlink{cf:eq:B.3}{(19)} into \(\fitinline{\chi_{\mathfrak p}(x)^3}\). Frozen exponents are \(\fitinline{j_{\mathfrak p}\equiv1_{\mathfrak p\mid t}+4v_{\mathfrak p}(g)\pmod6}\), including every collision; absorb fixed primes in the finite sector.

After fixing unit, \(\fitinline{\lambda}\)-valuation and a finite ray class, a reflected branch is

\hypertarget{cf:eq:B.14}{}

\[\fitmath{
S_m[A,Y](k_0)=3^{-m/3}\sum_{n\ \mathrm{sf},b}
\frac{A(n,b)\chi_{k_0}(nb)^3}{\sqrt{q_n}\,q_b}
V^\sharp\!\left(\frac{3^mq_nq_b^3H_0^2}{Yq_{k_0}^2}\right),
\qquad\text{(30)}
}\]

times a bounded row scalar, with \(\fitinline{k_0}\)-independent \(\fitinline{A}\), initially \(\fitinline{|A|\le a=1}\), \(\fitinline{Y\asymp H_0^2/X}\). For \(\fitinline{j=4}\) set \(\fitinline{A_n=1_{\mathfrak p\nmid n}A(\mathfrak pn,b)}\), \(\fitinline{A_b=1_{\mathfrak p\nmid n}A(n,\mathfrak pb)}\). Splitting \(\fitinline{1_{\mathfrak p\mid nb^3}=1_{\mathfrak p\mid n}+1_{\mathfrak p\nmid n,\mathfrak p\mid b}}\) gives, for \(\fitinline{q=q_{\mathfrak p}}\),

\hypertarget{cf:eq:B.15}{}

\[\fitmath{
S_m[AB_{\mathfrak p,4},q^2Y]
=-q^{-1/2}S_m[A,q^2Y]
+\chi_{k_0}(\mathfrak p)^3S_m[A_n,qY]
+q^{-1/2}\chi_{k_0}(\mathfrak p)^3S_m[A_b,Y/q].
\qquad\text{(31)}
}\]

This preserves squarefreeness and shared primes. The branch updates of \(\fitinline{(a^2,Y)}\) are \(\fitinline{(a^2,q^2Y)}\) for \(\fitinline{j\notin\{0,4\}}\); \(\fitinline{(a^2/q,q^2Y)}\) for active \(\fitinline{j=0}\); \(\fitinline{(a^2/q,q^2Y),(a^2,qY),(a^2/q,Y/q)}\) for \(\fitinline{j=4}\). Inactive primes increase neither entry. Consequently,

\hypertarget{cf:eq:B.16}{}

\[\fitmath{
a^2\le1,\qquad a^2Y\ll H_0^2q_t^2q_g/X.
\qquad\text{(32)}
}\]

A prime of \(\fitinline{t}\) costs at most \(\fitinline{q^2}\); a valuation-one prime of \(\fitinline{g}\) outside \(\fitinline{t}\) has \(\fitinline{j=4}\), costing \(\fitinline{q}\); higher valuations cost at most \(\fitinline{q^2\le q^{v_{\mathfrak p}(g)}}\).

Dual stars impose squarefreeness without \(\fitinline{S}\)-exclusion. Split the finitely many squarefree \(\fitinline{S}\)-parts of \(\fitinline{n}\) into bounded row factors before \protect\hyperlink{cf:eq:B.11}{(27)}. Retain the actual row cutoff in \S{}\protect\hyperlink{cf:sec:1.2}{A.1.3}'s joint profile; on \(\fitinline{q_n\asymp U,q_b\asymp B}\), separate into

\[\fitmath{
G_{b,\tau}(k_0)=\sum_{q_n\asymp U}^{*}
\frac{A(n,b)}{\sqrt{q_n}}(q_n/U)^{i\tau}\chi_{k_0}(n)^3.
}\]

Their Fourier densities are \(\fitinline{\ll(1+3^mUB^3/Y)^{-A}(1+|\tau|)^{-2}}\). Cauchy and \protect\hyperlink{cf:eq:B.11}{(27)} give

\[\fitmath{
\sum_{q_{k_0}\le H_0}^{*}|G_{b,\tau}(k_0)|^2
\ll a^2(H_0U)^\varepsilon(H_0+U).
}\]

Weighted Cauchy--Schwarz in \(\fitinline{b}\), using \(\fitinline{\sum_{q_b\asymp B}q_b^{-1}\ll1}\), now yields

\[\fitmath{
\|S_{m;U,B}\|_2
\ll a\mathcal D^\varepsilon3^{-m/3}(\sqrt{H_0}+\sqrt U)
(1+3^mUB^3/Y)^{-2},
}\]

For \(\fitinline{\mathcal D\ge2+H+X+q_g}\), use \(\fitinline{\sum_{U,B}(1+UB^3/Z)^{-2}\ll\log^2(2+Z)}\) and \(\fitinline{\sum_{U,B}\sqrt U(1+UB^3/Z)^{-2}\ll\sqrt Z}\), then sum \(\fitinline{m\ge-4}\), obtaining \(\fitinline{\|\mathrm{branch}\|_2^2\ll a^2\mathcal D^\varepsilon(H_0+Y)}\). Equation \protect\hyperlink{cf:eq:B.16}{(32)} and divisor-many branches give

\[\fitmath{
\sum_{q_s\le H}^{*}|T(X;u_0s,g)|^2
\ll(HXq_g)^\varepsilon\|W\|_{C^J}^2
\left(H+H^2q_g/X\right).
}\]

Finally write uniquely \(\fitinline{k=u_0sv^2}\), \(\fitinline{s}\) squarefree, allowing \(\fitinline{(s,v)>1}\). Multiplicativity with zeros and \(\fitinline{8\equiv2\pmod6}\) give \(\fitinline{T(X;u_0sv^2,f)=T(X;u_0s,fv^2)}\). The last bound at \(\fitinline{(H/q_v^2,g=fv^2)}\) has main expression \(\fitinline{q_v^{-2}(H+H^2q_f/X)}\). Sum \(\fitinline{\sum_vq_v^{-2}<\infty}\) and six units to obtain \protect\hyperlink{cf:eq:B.13}{(29)}. \ensuremath{\square}

Taking \(\fitinline{f=1}\), with \(\fitinline{V(y)=y^{1/2}W(y)}\), gives

\hypertarget{cf:eq:4.4}{}

\[\fitmath{
\sum_{0<q_k\le H}|\mathcal C_V(N;k,\nu)|^2\ll(HN)^\varepsilon(H+H^2/N)\|W\|_{C^J}^2.\qquad\text{(33)}
}\]

A log-annular partition extends this to \(\fitinline{V_G}\), whose seminorms are summable against every fixed exponential; the row-norm triangle sum converges.

\hypertarget{cf:sec:B.3}{}

\hypertarget{a.3.3-a-quadraticcubic-norm-for-one-prime}{%
\paragraph{A.3.3 A quadratic--cubic norm for one prime}\label{a.3.3-a-quadraticcubic-norm-for-one-prime}}

Let \(\fitinline{k,n}\) be squarefree, \(\fitinline{q_k\ll K}\), \(\fitinline{q_n\asymp N}\), \(\fitinline{q_b\asymp B}\) with arbitrary \(\fitinline{b}\). Rows avoid \(\fitinline{S}\); dual indices may meet \(\fitinline{S\setminus\{\lambda\}}\). For \(\fitinline{k}\)-independent \(\fitinline{\beta(n,b)}\) write

\[\fitmath{
b=gt^2,\quad c=(n,g),\quad n=cm,\quad g=ch,
\qquad q_c\asymp C,\quad q_g\asymp G,\quad q_t\asymp T,
\quad B\asymp GT^2.
}\]

Here \(\fitinline{g}\) is squarefree, \(\fitinline{c,m,h}\) are pairwise coprime squarefree, and \(\fitinline{g,t}\) may share primes. Define the dyadic quadratic row norm

\[\fitmath{
\mathcal Q_{C,G,T}(\beta)=
\sum_{\substack{k\ {\rm sf},\ (k,S)=1\\q_k\ll K}}
\left|\sum_{\substack{n\ {\rm sf},\ q_n\asymp N,\ q_b\asymp B\\
 b=gt^2,\ g\ {\rm sf},\ q_g\asymp G,\ q_t\asymp T\\
 c=(n,g),\ q_c\asymp C}}
\beta(n,b)\chi_k(nb)^3\right|^2.
}\]

All column restrictions are \(\fitinline{k}\)-independent. Then

\hypertarget{cf:eq:B.17}{}

\[\fitmath{
\mathcal Q_{C,G,T}(\beta)
\ll (KNB)^\varepsilon T
\sum_{q_t\asymp T}\sum_{\substack{r\ {\rm sf},\ (r,S)=1\\q_r\ll\min(K,C)}}
\left(\frac K{q_r}+\frac{NG}{C^2}\right)\frac C{q_r}
\sum_{c',m,h}|\beta(rc'm,rc'ht^2)|^2.
\qquad\text{(34)}
}\]

The last sum retains the support, \(\fitinline{c=rc'}\). Fixed row and coefficient restrictions are allowed.

\emph{Proof.} The zero-preserving identity is

\[\fitmath{
\chi_k(nb)^3=\chi_k(mh)^3\mathbf1_{(k,ct)=1}.
}\]

Cauchy over \(\fitinline{O(T)}\) choices of \(\fitinline{t}\) supplies \(\fitinline{T}\). Fixing \(\fitinline{t}\), keep \(\fitinline{(k,t)=1}\) as a row restriction and expand \(\fitinline{1_{(k,c)=1}=\sum_{r\mid(k,c)}\mu(r)}\). Divisor Cauchy costs \(\fitinline{(KNB)^\varepsilon}\). With \(\fitinline{k=rk'}\), \(\fitinline{c=rc'}\), squarefreeness imposes \(\fitinline{(k',r)=(c',r)=1}\); no \(\fitinline{(k',c')}\)-mask is reinserted termwise.

Absorb \(\fitinline{\chi_r(mh)^3}\) into column coefficients; split finite squarefree \(\fitinline{S}\)-parts of \(\fitinline{m,h}\) into bounded row factors. Their residual product is squarefree and good, of norm \(\fitinline{O(NG/C^2)}\). Apply \protect\hyperlink{cf:eq:B.11}{(27)} at \(\fitinline{K/q_r}\), group by product at divisor cost, and use Cauchy over \(\fitinline{O(C/q_r)}\) choices of \(\fitinline{c'}\). This proves \protect\hyperlink{cf:eq:B.17}{(34)}; restore \(\fitinline{S}\)-parts inside its positive sum. Empty quotients give zero; in nonempty bounded ranges replace scales by \(\fitinline{\max(1,\text{scale})}\) at constant cost. \ensuremath{\square}

For \(\fitinline{|\beta|\le1}\), ideal counting and dyadic summation in \protect\hyperlink{cf:eq:B.17}{(34)} give

\hypertarget{cf:eq:B.18}{}

\[\fitmath{
\sum_k\left|\sum_{n,b}\beta(n,b)\chi_k(nb)^3\right|^2
\ll(KNB)^\varepsilon(K+NB)NB.
\qquad\text{(35)}
}\]

Indeed for \(\fitinline{q_r\asymp r_0}\), coefficient count times \(\fitinline{C/q_r}\) is \(\fitinline{O(NG/r_0^2)}\); \(\fitinline{O(r_0)}\) choices of \(\fitinline{r}\) and both \(\fitinline{T}\)-factors give \(\fitinline{NB/r_0}\).

The cubic sieve (Heath-Brown, Theorem 2) for squarefree primary Eisenstein ideals and fixed coefficients is

\hypertarget{cf:eq:B.19}{}

\[\fitmath{
\sum_{q_a\le U}^{*}\left|\sum_{q_b\le V}^{*}c_b\chi_a(b)^2\right|^2
 \ll(UV)^\varepsilon\{U+V+(UV)^{2/3}\}\sum_b|c_b|^2.
\qquad\text{(36)}
}\]

Rational norms need not be squarefree; reciprocity and conjugation allow either cubic orientation.

\textbf{Lemma A.3.3} (One-prime hybrid norm). \emph{\hypertarget{cf:hybrid-norm}{} Let \(\fitinline{p}\) run over good primes with \(\fitinline{q_p\asymp L}\). Suppose \(\fitinline{|a(p)|\le1}\) is independent of \(\fitinline{k,n,b}\), and \(\fitinline{|\beta(n,b)|\le1}\) is independent of \(\fitinline{k,p}\). With the other variables as above, \hypertarget{cf:eq:B.20}{}}

\[\fitmath{
\begin{aligned}
 &\sum_k\left|\sum_p\frac{a(p)}{\sqrt{q_p}}\mathbf1_{p\nmid k}
 \sum_{n,b}\beta(n,b)\chi_k(nb)^3\chi_p(n)^{-2}\mathbf1_{p\nmid b}\right|^2\\
 &\hspace{15mm}\ll(KNBL)^\varepsilon(K+NB)B\{N+L+(NL)^{2/3}\}.
\end{aligned}
\qquad\text{(37)}
}\]

\emph{Separate fixed restrictions on the row, prime and dual supports, including ray sectors, are allowed.}

\emph{Proof.} First omit \(\fitinline{1_{p\nmid k}}\). Insert the prime sum into \protect\hyperlink{cf:eq:B.17}{(34)}'s column coefficient. Its decomposition gives

\[\fitmath{
\chi_p(rc'm)^{-2}\mathbf1_{p\nmid rc'ht^2}
 =\chi_p(r)^{-2}\chi_p(c'm)^{-2}\mathbf1_{p\nmid ht}.
}\]

The character retains zeros at \(\fitinline{r,c',m}\). Fix \(\fitinline{r,h,t}\) and absorb \(\fitinline{\chi_p(r)^{-2}1_{p\nmid ht}}\) into \(\fitinline{a(p)}\), whose central squared mass is \(\fitinline{\sum_{q_p\asymp L}q_p^{-1}\ll1}\). In the positive sum discard \(\fitinline{|\beta|^2\le1}\), group by squarefree \(\fitinline{j=c'm=n/r}\) at divisor cost, and apply \protect\hyperlink{cf:eq:B.19}{(36)}, giving

\[\fitmath{
(KNBL)^\varepsilon\left\{\frac N{q_r}+L+
                 \left(\frac{NL}{q_r}\right)^{2/3}\right\}.
}\]

The \(\fitinline{O(G/C)}\) choices of \(\fitinline{h}\), both \(\fitinline{T}\)-factors and \(\fitinline{C/q_r}\) give \(\fitinline{B/q_r}\); the dyadic \(\fitinline{r}\)-count cancels its denominator. Since \(\fitinline{NG/C^2\ll NB}\), summing gives \protect\hyperlink{cf:eq:B.20}{(37)} without the mask.

For collisions put \(\fitinline{k=pk'}\). Divisor Cauchy and \protect\hyperlink{cf:eq:B.18}{(35)}, fixing \(\fitinline{p}\) and absorbing its characters into \(\fitinline{\beta}\), give

\[\fitmath{
\ll(KNBL)^\varepsilon\sum_{q_p\asymp L}\frac1{q_p}
       \left(\frac K{q_p}+NB\right)NB
 \ll(KNBL)^\varepsilon(K/L+NB)NB.
}\]

This is bounded by \protect\hyperlink{cf:eq:B.20}{(37)}, whose brace is \(\fitinline{\ge N}\). Triangle restores the mask. An absent prime uses \protect\hyperlink{cf:eq:B.18}{(35)}, equivalent to \(\fitinline{L=1}\). \ensuremath{\square}

\hypertarget{cf:sec:B.4}{}

\hypertarget{a.3.4-whole-index-compensation-and-the-required-energy}{%
\paragraph{A.3.4 Whole-index compensation and the required energy}\label{a.3.4-whole-index-compensation-and-the-required-energy}}

Insert \(\fitinline{1_{p\mid nb^3}}\) in \protect\hyperlink{cf:eq:B.1}{(17)} for a selected good prime \(\fitinline{p}\). If \(\fitinline{p}\) is already a base local prime, its original zero kills the sum. Otherwise write

\[\fitmath{
\mathbf1_{p\mid nb^3}=1-\chi_p(nb^3)^0
}\]

and reflect both terms. Inactive coefficients combine to \(\fitinline{1-(1-q_p^{-1})=q_p^{-1}}\), with \(\fitinline{p}\) absent from the conductor. The active dual factor is

\hypertarget{cf:eq:B.21}{}

\[\fitmath{
q_p^{-1/2}\chi_p(\lambda^4\mu)^{-2},
\qquad\text{(38)}
}\]

with scalar \protect\hyperlink{cf:eq:B.4}{(20)} replacing \(\fitinline{\omega_{p,0}}\) by \(\fitinline{-\omega_{p,0}}\). Identical fixed Fourier term and active radical give identical cusp data, so subtract before estimating, even when primal squarefree and cube indices share \(\fitinline{p}\).

For distinct active \(\fitinline{r,p}\), scalar \protect\hyperlink{cf:eq:B.4}{(20)} has pair factor

\[\fitmath{
\chi_r(p)^{2j_r+2}\chi_p(r)^{2j_p+2}
=(p/r)_3^{j_r+j_p+2}.
}\]

For a residual row prime \(\fitinline{j_r=1}\) and marked prime \(\fitinline{j_p=0}\), this equals one. Freeze other active primes, including every \(\fitinline{j=4}\) prime. Separately fixing row product \(\fitinline{R}\) and \(\fitinline{P=p}\) modulo \S{}\protect\hyperlink{cf:sec:B.1}{A.3.1}'s full \(\fitinline{M^2}\) fixes cusp and additive data without pair restrictions. The remaining scalar separates into bounded row and prime factors.

For dual index a fixed factor times \(\fitinline{nb^3}\), the moving characters are

\[\fitmath{
\chi_R(nb^3)^3=\chi_R(nb)^3,\qquad
\chi_P(nb^3)^{-2}=\chi_P(n)^{-2}\mathbf1_{(P,b)=1}.
}\]

Primal row--prime collisions retain \(\fitinline{1_{(P,R)=1}}\); frozen collisions give separate row, prime or dual restrictions. Thus Lemma \protect\hyperlink{cf:hybrid-norm}{A.3.3}'s two independence hypotheses hold. Logarithmic Fourier separation handles joint completed-norm/prime-norm weights with finite weighted seminorms.

For a bounded row-independent coefficient \(\fitinline{a(\mathfrak p)}\), define

\hypertarget{cf:eq:4.5}{}

\[\fitmath{
\mathcal T(m)=\sum_{q_{\mathfrak p}\asymp Z^{5/33}}a(\mathfrak p)
 \sum_{c\ {\rm sf},n}\frac{\overline{\alpha(cn^3)}\gamma_2(c)\nu(cn^3)\chi_{cn^3}(m)}
 {\sqrt{q_c}\,q_n}\mathbf1_{\mathfrak p\mid cn^3}
 V\!\left(\frac{q_cq_n^3}{Zq_{\mathfrak p}}\right).
\qquad\text{(39)}
}\]

Including a common joint norm profile, the estimate is

\hypertarget{cf:eq:4.6}{}

\[\fitmath{
\sum_{0<q_m\ll Z^{28/33}}|\mathcal T(m)|^2\ll Z^{28/33+\varepsilon}.
\qquad\text{(40)}
}\]

The prime stays inside the norm throughout the proof.

\textbf{Proposition A.3.4} (Whole-index one-prime energy). \emph{\hypertarget{cf:one-prime-energy}{} The sum \protect\hyperlink{cf:eq:4.5}{(39)} satisfies \protect\hyperlink{cf:eq:4.6}{(40)}.}

\emph{Proof.} Locally put \(\fitinline{M=28/33}\), \(\fitinline{z=5/33}\), \(\fitinline{N_*=1+z=38/33}\). Freeze the row's maximal powerful part, \(\fitinline{S}\)-supported valuation-one part and unit. Their norm dyad has exponent \(\fitinline{O}\), the remaining squarefree row exponent \(\fitinline{H\le M-O+O(1/\log Z)}\). Apply \S{}\protect\hyperlink{cf:sec:B.1}{A.3.1}'s finite sector reduction.

For active frozen good primes, let \(\fitinline{A_0}\) be radical exponent; \(\fitinline{S_0}\) count small Ramanujan and active zero-mask factors; \(\fitinline{N_0,B_0}\) count forced squarefree/cube factors. Splitting as in \protect\hyperlink{cf:eq:B.15}{(31)} leaves shared cube occurrences in \(\fitinline{b}\) on the squarefree branch, or extracts one cube occurrence with factor \(\fitinline{q_p^{-1/2}}\). Powerful support gives

\[\fitmath{
2A_0\le O,\qquad 0\le N_0\le A_0,\qquad S_0,B_0\ge0.
}\]

Inactive marked primes have mass \(\fitinline{\sum_{q_p\asymp Z^z}q_p^{-1}\ll Z^\varepsilon}\); remove by triangle and set \(\fitinline{z_a=0}\). Active primes stay inside the norm with \(\fitinline{z_a=z}\).

For residual dual squarefree, cube and ramified exponents \(\fitinline{v,b,e_\lambda}\), the central coefficient and dual length are

\hypertarget{cf:eq:B.22}{}

\[\fitmath{
Z^{-v/2-b-e_\lambda/3-(S_0+B_0+z_a)/2},\qquad
 T_d=2H+2A_0+2z_a-N_*-N_0-3B_0.
\qquad\text{(41)}
}\]

The actual kernel argument is a fixed factor times \(\fitinline{Z^{v+3b+e_\lambda-T_d}y_ny_b^3y_R^{-2}y_P^{-2}}\), each \(\fitinline{y}\) normalized by its dyadic center, omitting \(\fitinline{y_P}\) when inactive. By \protect\hyperlink{cf:eq:B.7}{(23)}, extract \(\fitinline{Z^{-(T_d-v-3b-e_\lambda)_+/4}}\) from the \textbf{linear} vector before Fourier separation; the normalized profile has bounded seminorms. Then enlarge the positive row norm to \(\fitinline{q_R\ll Z^H}\).

Apply \protect\hyperlink{cf:eq:B.20}{(37)} at \(\fitinline{(K,N,B,L)=(Z^H,Z^v,Z^b,Z^{z_a})}\), or \protect\hyperlink{cf:eq:B.18}{(35)} when inactive. Its prime inverse square root is already included. Multiply the remaining squared coefficient and count powerful rows; with \(\fitinline{u_h=\min\{v,z_a,(v+z_a)/3\}}\) the exponent is

\hypertarget{cf:eq:B.23}{}

\[\fitmath{
\begin{aligned}
 E={}&O/2+\max(H,v+b)-S_0-B_0+z_a-u_h-b-2e_\lambda/3\\
    &-\tfrac12(T_d-v-3b-e_\lambda)_+.
\end{aligned}
\qquad\text{(42)}
}\]

Powerful ideals number \(\fitinline{O(Z^{O/2})}\): write \(\fitinline{x^2y^3}\), \(\fitinline{y}\) squarefree, and sum \(\fitinline{Z^{O/2}q_y^{-3/2}}\). Fixed-supported parts and allocations have divisor-many choices. Their Cauchy cost is the single \(\fitinline{O/2+\varepsilon}\) already in \protect\hyperlink{cf:eq:B.23}{(42)}.

Retain \(\fitinline{v+3b+e_\lambda\le T_d}\), suppressing tiny buffers. Since

\[\fitmath{
T_d-H\le M+2z-N_*=0,
}\]

we have \(\fitinline{\max(H,v+b)=H}\). If \(\fitinline{u_h=z_a}\), then \(\fitinline{E\le M}\); otherwise \(\fitinline{u_h\ge v/2}\). Put \(\fitinline{y_d=v+3b+e_\lambda}\); then

\[\fitmath{
u_h+b+2e_\lambda/3+\tfrac12(T_d-y_d)_+
 \ge y_d/4+\tfrac12(T_d-y_d)_+\ge T_d/4.
}\]

Substitution in \protect\hyperlink{cf:eq:B.23}{(42)}, with \(\fitinline{\Delta_H=M-O-H\ge0}\), gives

\hypertarget{cf:eq:B.24}{}

\[\fitmath{
4(E-M)\le N_*+2z_a-2M-2\Delta_H-(2A_0-N_0+B_0+4S_0)
 \le N_*+2z-2M=-\frac8{33}.
\qquad\text{(43)}
}\]

Hence all retained blocks have \(\fitinline{E\le M}\).

Discard whole tail dyads outside a small retention buffer before separation; their kernel argument exceeds a positive power of \(\fitinline{Z}\). For \(\fitinline{|K(y)|\le C_Ay^{-A}}\), \(\fitinline{A>1}\), lattice shells give

\hypertarget{cf:eq:B.25}{}

\[\fitmath{
\sum_{\substack{0\ne\mu\in\lambda^{-4}\mathcal O\\a q_\mu>U}}
 |K(aq_\mu)|\ll_A C_A(1+a^{-1})U^{1-A},\qquad U\ge1.
\qquad\text{(44)}
}\]

Indeed each dyadic shell has \(\fitinline{O(1+2^jU/a)}\) points. Raw coefficients, rows, primes and branches have fixed polynomial cost; choose \(\fitinline{A}\) afterwards for an arbitrarily small tail. The lower bound \(\fitinline{e_\lambda\ge-4\log_Z3}\) and fixed support shifts cost \(\fitinline{O(1/\log Z)}\). Buffers, dyads and finite profile orders fit every requested loss, proving \protect\hyperlink{cf:eq:4.6}{(40)}. \ensuremath{\square}

\hypertarget{cf:sec:C}{}

\hypertarget{a.4-the-inverse-second-moment}{%
\subsubsection{A.4 The inverse second moment}\label{a.4-the-inverse-second-moment}}

The two Poisson transforms recover the essential M\"{o}bius coefficients through the cubic Gauss identities.

\textbf{Theorem A.4.1} (Inverse second moment). \emph{\hypertarget{cf:inverse-moment}{} Fix \(\fitinline{C,\kappa,\varepsilon>0}\). For \(\fitinline{Z\ge2}\), \(\fitinline{1\le D,H\le Z^C}\), and \(\fitinline{H\ge DZ^\kappa}\), \hypertarget{cf:eq:C.1}{}}

\[\fitmath{
\sum_{0<q_u\le H}\left|D^{-1/2}
\sum_n\mu(n)\nu(n)\chi_n(u)^{\pm1}W(q_n/D)\right|^2
\ll HZ^\varepsilon\|W\|_{C^J}^2.
\qquad\text{(45)}
}\]

\emph{The fixed \(\fitinline{\nu}\) includes its presentation's Euler deletions. The order depends only on \(\fitinline{C,\kappa,\varepsilon}\), support and fixed arithmetic data; test norm twists cost a fixed height power.}

\emph{Proof.} We prove it through the positive energy

\hypertarget{cf:eq:C.2}{}

\[\fitmath{
E(H,X,F;\xi,W)=\frac1{XF}
\sum_{F\le q_f<2F}^{*}\sum_{0<q_k\le H}
\left|\sum_n^{*}a_\xi(n)\chi_n(k)\chi_n(f)^4W(q_n/X)\right|^2.
\qquad\text{(46)}
}\]

The label \(\fitinline{f}\) is averaged. Write \(\fitinline{E^{(r)}}\) for an added inner mask \(\fitinline{1_{(n,r)=1}}\); \(\fitinline{T}\) is \protect\hyperlink{cf:eq:B.12}{(28)}.

\hypertarget{cf:sec:C.1}{}

\hypertarget{a.4.1-poisson-smooth-separation-and-punctures}{%
\paragraph{A.4.1 Poisson, smooth separation, and punctures}\label{a.4.1-poisson-smooth-separation-and-punctures}}

For a primitive character \(\fitinline{\psi}\) modulo \(\fitinline{m}\), masked Poisson summation is

\hypertarget{cf:eq:C.3}{}

\[\fitmath{
\sum_{k\in\mathcal O}\psi(k)1_{(k,r)=1}\Phi(q_k/H)
=\frac{H\gamma(\psi)}{\sqrt{q_m}}
\sum_{d\mid\operatorname{rad}r}\frac{\mu(d)\psi(d)}{q_d}
\sum_h\overline{\psi(h)}\widehat\Phi\!\left(\frac{Hq_h}{q_dq_m}\right).
\qquad\text{(47)}
}\]

Here \(\fitinline{\Phi}\) is radial Schwartz and \(\fitinline{\gamma(\psi)=q_m^{-1/2}\sum_{x\bmod m}\psi(x)e(x/m)}\). Expand the mask as \(\fitinline{\sum_{d\mid(r,k)}\mu(d)}\), substitute \(\fitinline{k=d\ell}\), and Poisson-sum modulo \(\fitinline{m}\). The finite transform is \(\fitinline{\sqrt{q_m}\gamma(\psi)\overline{\psi(h)}}\), with nonunit zeros. Primitive principal modulus one has value one at \(\fitinline{h=0}\), contributing \(\fitinline{H\widehat\Phi(0)\prod_{\mathfrak p\mid r}(1-q_{\mathfrak p}^{-1})}\). Shared primes of \(\fitinline{m,r}\) vanish through \(\fitinline{\psi(d)}\).

Smooth quadratic kernels are separated by \protect\hyperlink{cf:eq:1.4}{(5)}.

M\"{o}bius inversion and \protect\hyperlink{cf:eq:A.1}{(6)} remove a puncture without changing \(\fitinline{XF}\):

\hypertarget{cf:eq:C.4}{}

\[\fitmath{
E^{(r)}(H,X,F;\xi,W)
\le\tau(\operatorname{rad}r)
\sum_{d\mid\operatorname{rad}r}E(H,X/q_d,Fq_d;\xi,W).
\qquad\text{(48)}
}\]

Set \(\fitinline{n=dm}\), changing \(\fitinline{f}\) to \(\fitinline{df}\); \(\fitinline{\chi_m(d)^4}\) retains \(\fitinline{(m,d)=1}\), and \(\fitinline{(d,f)>1}\) vanishes. Cauchy in \(\fitinline{d}\) and positive enlargement of the injective image \(\fitinline{f\mapsto df}\) give \protect\hyperlink{cf:eq:C.4}{(48)}, preserving \(\fitinline{XF}\).

\hypertarget{cf:sec:C.2}{}

\hypertarget{a.4.2-canonical-estimate-and-cube-completion}{%
\paragraph{A.4.2 Canonical estimate and cube completion}\label{a.4.2-canonical-estimate-and-cube-completion}}

We claim that, for every fixed \(\fitinline{C_0\ge1,\kappa>0}\),

\hypertarget{cf:eq:C.5}{}

\[\fitmath{
E(H,X,F;\xi,W)\ll Z^\varepsilon\Sigma\|W\|_{C^J}^2,
\quad \Sigma=XF\le Z^{C_0},\quad
H,X,F\ge1,\quad H\le\Sigma Z^{-\kappa}.
\qquad\text{(49)}
}\]

The derivative order is independent of the location of the fixed support interval.

For \(\fitinline{\Psi_{k,f}(n)=\xi(n)\chi_n(k)\chi_n(f)^4}\), cube inversion gives the exact identity

\hypertarget{cf:eq:C.6}{}

\[\fitmath{
X^{-1/2}\sum_n^{*}a_\xi(n)\chi_n(k)\chi_n(f)^4W(q_n/X)
=\sum_h\frac{\mu(h)\overline{\alpha(h)}^{3}\Psi_{k,f}(h)^3}{q_h}
T(X/q_h^3;k,f).
\qquad\text{(50)}
}\]

Opening \(\fitinline{T}\), each total cube \(\fitinline{b}\) has divisor factor \(\fitinline{\sum_{h\mid b}\mu(h)=1_{b=1}}\), by complete multiplicativity with zeros.

Put \(\fitinline{H_c^3=\min(X,X^2/H^2)}\), \(\fitinline{L_b=X/q_b^3}\). Weighted Cauchy, \protect\hyperlink{cf:eq:B.13}{(29)} and the \(\fitinline{f}\)-average bound \(\fitinline{q_h\le H_c}\) by

\[\fitmath{
Z^\varepsilon\|W\|_{C^J}^2\left(H+\frac{H^2F H_c^3}{X}\right)
\ll Z^\varepsilon\Sigma\|W\|_{C^J}^2\quad(H\le\Sigma).
}\]

Reopen the remaining \(\fitinline{T}\) terms. A total cube \(\fitinline{b}\) has coefficient \(\fitinline{\overline{\alpha(b)}^3\xi(b)^3 1_{(b,f)=1}\sum_{h\mid b,q_h>H_c}\mu(h)}\), bounded by \(\fitinline{\tau(b)}\), and bounded row multiplier \(\fitinline{\chi_b(k)^3}\). A second weighted Cauchy uses \(\fitinline{\sum_{q_b\le R}\tau(b)/q_b\ll\log^2(2R)}\) to give

\hypertarget{cf:eq:C.7}{}

\[\fitmath{
E(H,X,F;\xi,W)
\ll Z^\varepsilon\left\{\Sigma\|W\|_{C^J}^2+
\sup_{q_b>H_c,\ L_b>1}E(H,L_b,F;\xi,W)\right\}.
\qquad\text{(51)}
}\]

For \(\fitinline{L_b\le1}\), nonempty annular support has boundedly many columns, costing \(\fitinline{O(H)}\). The supremum is empty when \(\fitinline{H^2\le X}\).

\hypertarget{cf:sec:C.3}{}

\hypertarget{a.4.3-two-poisson-transformations-and-the-local-projector}{%
\paragraph{A.4.3 Two Poisson transformations and the local projector}\label{a.4.3-two-poisson-transformations-and-the-local-projector}}

For \(\fitinline{1\le H,L,F,\Sigma\le Z^{C_0}}\) and \(\fitinline{\max(H,LF)\le\Sigma}\), we prove

\hypertarget{cf:eq:C.8}{}

\[\fitmath{
E(H,L,F;\xi,W)
\ll Z^\varepsilon\Sigma\|W\|_{C^{4m+12}}^2
\left(1+\sup\frac{E(H',X',F';\xi',U)}{\Sigma'}\right),
\qquad\text{(52)}
}\]

where \(\fitinline{\Sigma'=X'F'}\), \(\fitinline{H',X',F'\ge1}\), \(\fitinline{\|U\|_{C^m}\le1}\) on a fixed enlargement of the support, and

\hypertarget{cf:eq:C.9}{}

\[\fitmath{
H'\le\frac{HL}{\Sigma F},\qquad
\frac{H'}{\Sigma'}\le\frac H\Sigma,\qquad \Sigma'\le L.
\qquad\text{(53)}
}\]

Square a rescaled real Schwartz function of compact Fourier support to obtain radial \(\fitinline{\Phi\ge1}\) on \(\fitinline{[0,1]}\), \(\fitinline{\Phi\ge0}\), with \(\fitinline{\operatorname{supp}\widehat\Phi\subset[0,C_\Phi]}\). For \(\fitinline{\operatorname{supp}W\subset[u,v]}\), fix \(\fitinline{c_I\ge\max(1,4C_\Phi v^2)}\) and enlarged support \(\fitinline{[u/2,2v]}\).

Apply \protect\hyperlink{cf:eq:C.3}{(47)} to the square. Write \(\fitinline{n_i=Cu_i}\), \(\fitinline{C=(n_1,n_2)}\); the primitive row conductor is \(\fitinline{u_1u_2}\), residual mask \(\fitinline{C}\), divisor \(\fitinline{d\mid C}\). Zero frequency costs \(\fitinline{O(H\|W\|_\infty^2)}\). Use \(\fitinline{a_\xi(n)\gamma_1(n)=\mu(n)G(n)\xi(n)}\), expand its fixed ray factor, and remove \(\fitinline{(u_1,u_2)=1}\) by complete M\"{o}bius inversion, \(\fitinline{u_i=tx_i}\). A fixed ray component is

\hypertarget{cf:eq:C.10}{}

\[\fitmath{
P_{C,d,t}(y;U)=\sum_n^{*}\mu(n)\xi_1(n)1_{(n,t)=1}
\overline{\chi_n(y)}\chi_n(C)^4\chi_n(d)U(q_n/\ell),
\qquad \ell=\frac L{q_Cq_t}.
\qquad\text{(54)}
}\]

Here \(\fitinline{C,t}\) are coprime squarefree, \(\fitinline{d\mid C}\), \(\fitinline{q_Cq_t\le2vL}\). Combine frequency and label as \(\fitinline{y=hf^2}\), retaining \(\fitinline{(f,Ct)=(h,t)=1}\). This map has divisor-bounded fibres; enlarge its nonzero image to \(\fitinline{q_y\le Y_{C,d}}\), where

\[\fitmath{
w_{C,d}=\frac{Hq_C}{q_dL^2F},\qquad
Y_{C,d}=\frac{c_I\Sigma LFq_d}{Hq_C^2}.
}\]

Since \(\fitinline{\Sigma\ge LF}\), Cauchy and smooth separation give the positive form

\hypertarget{cf:eq:C.11}{}

\[\fitmath{
Q_{\xi_1}(U)=
\sum_{\substack{C,t\ \mathrm{sf}\\(C,t)=1\\q_Cq_t\le2vL}}
\sum_{d\mid C}w_{C,d}\sum_y\Phi(q_y/Y_{C,d})|P_{C,d,t}(y;U)|^2.
\qquad\text{(55)}
}\]

If \(\fitinline{Q_{\xi_1}(U)\le M\|U\|_{C^j}^2}\) for all common \(\fitinline{U}\), then \protect\hyperlink{cf:eq:1.4}{(5)} gives \(\fitinline{E\ll Z^\varepsilon(\Sigma+M)\|W\|_{C^{2j+4}}^2}\). Kernels combine annular weights, inverse roots and \(\fitinline{\widehat\Phi(A/(x_1x_2))}\) on bounded support, with uniform derivatives.

Apply \protect\hyperlink{cf:eq:C.3}{(47)} to \(\fitinline{y}\), write columns \(\fitinline{gz_i}\), take \(\fitinline{e\mid g}\), and remove \(\fitinline{(z_1,z_2)=1}\) by \(\fitinline{z_i=wn_i}\). Nonzero terms have

\[\fitmath{
(g,Ct)=(w,gCth)=1,\qquad e\mid g,\qquad h\ne0,
}\]

and sign \(\fitinline{\mu(e)\mu(w)}\). Inverse conversion \protect\hyperlink{cf:eq:A.1}{(6)} restores cubic coefficients. The ray phase is \(\fitinline{\xi_1(z_1z_2^{-1})G(z_2z_1^{-1})}\); expanding \(\fitinline{\xi_1(z)G(z^{-1})=\sum_{\xi'}c_{\xi'}\xi'(z)}\) gives \(\fitinline{a_{\xi'}(n)}\) independent of reconstruction choices.

Reconstruct the puncture, label and row by

\hypertarget{cf:eq:C.12}{}

\[\fitmath{
r=tg/e,\qquad f'=Cew,\qquad k'=deh,
\qquad a^\sharp(n)=a_{\xi'}(n)1_{(n,r)=1}\chi_n(k')\chi_n(f')^4.
\qquad\text{(56)}
}\]

The ideals \(\fitinline{r,f'}\) are coprime squarefree. Locally,

\hypertarget{cf:eq:C.13}{}

\[\fitmath{
\chi_n(dh)\overline{\chi_n(e)}\chi_n(C)^4
=\chi_n(deh)\chi_n(Ce)^4;
\qquad\text{(57)}
}\]

Both sides have unit exponent five at \(\fitinline{e}\) and vanish at nonunits; \(\fitinline{w}\) contributes \protect\hyperlink{cf:eq:A.1}{(6)}'s multiplicative correction.

For \(\fitinline{U_0(x)=x^{-1/2}U(x)}\), the scale, scalar and kernel are

\[\fitmath{
X'_0=\frac\ell{q_gq_w}=\frac L{q_rq_{f'}},
\qquad w_{C,d}\frac{Y_{C,d}q_g}{q_e\ell}=\frac{c_I\Sigma q_r}{L^2},
}\]

\hypertarget{cf:eq:C.14}{}

\[\fitmath{
K_{r,f',k'}(x_1,x_2)=U_0(x_1)\overline{U_0(x_2)}
\widehat\Phi\!\left(\frac{c_I\Sigma Fq_{k'}q_r^2}{HLx_1x_2}\right).
\qquad\text{(58)}
}\]

These depend only on \(\fitinline{(r,f',k')}\). Support forces \(\fitinline{q_rq_{f'}\le2vL}\) and \(\fitinline{0<q_{k'}\le HL/(\Sigma Fq_r^2)}\).

At fixed \(\fitinline{(r,f',k')}\), all preimages are

\[\fitmath{
t\mid r,\quad f'=Cew,\quad e\mid k',\quad(w,k')=1,
\quad d\mid(C,k'),\quad g=e(r/t),\quad h=k'/(de),
}\]

with disjoint \(\fitinline{f'=Cew}\). For \(\fitinline{\mathfrak p\mid f'}\), write \(\fitinline{I_{\mathfrak p}=1_{\mathfrak p\mid k'}}\); allocation to \(\fitinline{C,e,w}\) gives

\hypertarget{cf:eq:C.15}{}

\[\fitmath{
\underbrace{(1+I_{\mathfrak p})}_{\mathfrak p\mid C}
+\underbrace{(-I_{\mathfrak p})}_{\mathfrak p\mid e}
+\underbrace{(-(1-I_{\mathfrak p}))}_{\mathfrak p\mid w}
=I_{\mathfrak p}.
\qquad\text{(59)}
}\]

The \(\fitinline{C}\) terms distinguish \(\fitinline{\mathfrak p\nmid d}\) and \(\fitinline{\mathfrak p\mid d}\). Shared kernel, scalar and character permit this signed cancellation before absolute values: it forces \(\fitinline{f'\mid k'}\), while \(\fitinline{t\mid r}\) contributes \(\fitinline{\tau(r)}\).

Writing \(\fitinline{Z_Q}\) for the second zero frequency, we have

\hypertarget{cf:eq:C.16}{}

\[\fitmath{
Q_{\xi_1}(U)-Z_Q=
\frac{c_I\Sigma}{L^2}\sum_{\xi'}c_{\xi'}
\sum_{\substack{r,f'\ \mathrm{sf}\\(r,f')=1}}q_r\tau(r)
\sum_{\substack{k'\ne0\\f'\mid k'}}\sum_{n_1,n_2}^{*}
a^\sharp(n_1)\overline{a^\sharp(n_2)}
K_{r,f',k'}(q_{n_1}/X'_0,q_{n_2}/X'_0).
\qquad\text{(60)}
}\]

By \protect\hyperlink{cf:eq:C.3}{(47)}, \(\fitinline{|Z_Q|\ll Z^\varepsilon\Sigma\|U\|_\infty^2}\), since its absolute divisor sum is \(\fitinline{\ll\Sigma\sum_C\tau(C)q_C^{-2}\sum_{q_t\ll L}q_t^{-1}}\).

On dyads \(\fitinline{q_r\asymp R,q_{f'}\asymp F'}\), set

\hypertarget{cf:eq:C.17}{}

\[\fitmath{
X'=\frac L{RF'},\qquad \Sigma'=\frac LR,
\qquad H'=\frac{HL}{\Sigma FR^2}.
\qquad\text{(61)}
}\]

Nonempty support and the projector give \(\fitinline{F'\le q_{k'}\le HL/(\Sigma Fq_r^2)}\), hence \(\fitinline{H'\ge1}\). For \(\fitinline{j\mid r}\),

\hypertarget{cf:eq:C.18}{}

\[\fitmath{
\frac{X'}{q_j}\ge\frac{\Sigma Fq_r^2}{HRq_j}
\ge\frac{\Sigma F}{H}\ge1.
\qquad\text{(62)}
}\]

Puncture removal therefore preserves admissibility and \(\fitinline{\Sigma'}\). Since \(\fitinline{1\le X'/X'_0<4}\), rescaled kernels have fixed enlarged support and uniform \(\fitinline{C^{2m+4}}\) norm \(\fitinline{\ll\|U\|_{C^{2m+4}}^2}\).

Writing \(\fitinline{S_m}\) for \protect\hyperlink{cf:eq:C.8}{(52)}'s supremum, remove punctures and positively enlarge \(\fitinline{(f',k')}\) to bound each retained \(\fitinline{r}\)'s common-test sum by \(\fitinline{Z^\varepsilon(\Sigma')^2S_m\|V\|_{C^m}^2}\). Apply \protect\hyperlink{cf:eq:1.4}{(5)} and count \(\fitinline{O(R)}\) values of \(\fitinline{r}\); the scalar is

\[\fitmath{
\frac{\Sigma R}{L^2}\,R(\Sigma')^2=\Sigma.
}\]

Absorb divisor and dyadic counts in \(\fitinline{Z^\varepsilon}\). This proves \protect\hyperlink{cf:eq:C.8}{(52)}; \protect\hyperlink{cf:eq:C.17}{(61)} gives all inequalities \protect\hyperlink{cf:eq:C.9}{(53)}.

\hypertarget{cf:sec:C.4}{}

\hypertarget{a.4.4-closing-the-finite-depth-induction}{%
\paragraph{A.4.4 Closing the finite-depth induction}\label{a.4.4-closing-the-finite-depth-induction}}

Inductively prove \protect\hyperlink{cf:eq:C.5}{(49)} for \(\fitinline{H\le Z^{j\kappa}}\). At \(\fitinline{j=0}\), counting gives \(\fitinline{E\ll HX\|W\|_\infty^2\le\Sigma\|W\|_\infty^2}\). Next use \protect\hyperlink{cf:eq:C.7}{(51)}. A nonempty supremum has \(\fitinline{H^2>X}\), \(\fitinline{H_c^3=X^2/H^2}\), \(\fitinline{q_b>H_c}\). Apply \protect\hyperlink{cf:eq:C.8}{(52)} at \(\fitinline{L=L_b}\), keeping \(\fitinline{\Sigma=XF}\); every child satisfies

\[\fitmath{
H'\le\frac{HL_b}{\Sigma F}=\frac H{q_b^3F^2}
<H\left(\frac H\Sigma\right)^2
\le Z^{-2\kappa}H\le Z^{j\kappa},
\quad H'/\Sigma'\le Z^{-\kappa},\quad \Sigma'\le\Sigma.
}\]

Induction applies. Distribute losses among completion, transfer and child, propagating orders backwards by \(\fitinline{m\mapsto4m+12}\) and retaining the completed-energy order. Depth \(\fitinline{\le\lceil C_0/\kappa\rceil}\) gives finite orders and support enlargements, fixed before external heights. This proves \protect\hyperlink{cf:eq:C.5}{(49)}.

\hypertarget{cf:sec:C.5}{}

\hypertarget{a.4.5-return-to-the-inverse-polynomial}{%
\paragraph{A.4.5 Return to the inverse polynomial}\label{a.4.5-return-to-the-inverse-polynomial}}

Majorize the \(\fitinline{u}\)-ball by \(\fitinline{\Phi(q_u/H)}\), expand and apply \protect\hyperlink{cf:eq:C.3}{(47)}. Write columns \(\fitinline{gz_i}\), mask divisor \(\fitinline{e\mid g}\), frequency \(\fitinline{h\ne0}\); zero frequency costs \(\fitinline{O(H\|W\|_\infty^2)}\). Inverse conversion \protect\hyperlink{cf:eq:A.1}{(6)} and fixed ray expansion restore cubic coefficients. Remove \(\fitinline{(z_1,z_2)=1}\) by \(\fitinline{z_i=vn_i}\). Retaining zeros, the bijection is

\hypertarget{cf:eq:C.19}{}

\[\fitmath{
b=g/e,\qquad f=ev,\qquad k=eh;
\qquad e=(f,k),\quad v=f/e,\quad g=be,\quad h=k/e.
\qquad\text{(63)}
}\]

Here \(\fitinline{b,f}\) are coprime squarefree. The identities \(\fitinline{gv=bf}\), \(\fitinline{he^5v^4=kf^4}\), \(\fitinline{\mu(e)\mu(v)=\mu(f)}\) and \(\fitinline{q_g/q_e=q_b}\) give each nonzero ray component

\hypertarget{cf:eq:C.20}{}

\[\fitmath{
\begin{aligned}
\frac H{D^2}\sum_{\substack{b,f\ \mathrm{sf}\\(b,f)=1}}\mu(f)q_b
\sum_{k\ne0}\sum_{n_1,n_2}^{*}
&a_\xi(n_1)\overline{a_\xi(n_2)}
\chi_{n_1}(kf^4)\overline{\chi_{n_2}(kf^4)}1_{(n_1n_2,b)=1}\\
&\times W_0(q_{bfn_1}/D)\overline{W_0(q_{bfn_2}/D)}
\widehat\Phi\!\left(\frac{Hq_k}{q_f^2q_{n_1}q_{n_2}}\right),
\end{aligned}
\qquad\text{(64)}
}\]

with \(\fitinline{W_0(x)=x^{-1/2}W(x)}\) and support \(\fitinline{q_bq_f\ll D}\), \(\fitinline{q_k\ll D^2/(Hq_b^2)}\). On \(\fitinline{q_b\asymp B,q_f\asymp F}\) put

\[\fitmath{
X=\frac D{BF},\qquad \mathcal H=\frac{C D^2}{HB^2},\qquad \Sigma=XF=\frac DB,
}\]

Fixed \(\fitinline{C}\) covers supports. Since \(\fitinline{\mathcal H/\Sigma=CD/(HB)\ll Z^{-\kappa}}\), apply \protect\hyperlink{cf:eq:C.5}{(49)} with gap \(\fitinline{\kappa/2}\), then \protect\hyperlink{cf:eq:1.4}{(5)}, \protect\hyperlink{cf:eq:C.4}{(48)}. Kernels have uniform derivatives, giving dyad bound

\[\fitmath{
Z^\varepsilon\frac{HB}{D^2}\,B(XF)^2\|W\|_{C^J}^2
=HZ^\varepsilon\|W\|_{C^J}^2.
}\]

Subunit column scales are empty or have boundedly many columns; count them. Subunit row scales are empty. Sum dyads and fixed ray components, including zero frequency, to prove \protect\hyperlink{cf:eq:C.1}{(45)}. Conjugation with \(\fitinline{\nu\mapsto\bar\nu}\) gives the other orientation. Finite seminorms of \(\fitinline{W(y)y^{it}}\) cost \(\fitinline{(1+|t|)^{O(1)}}\).

For \protect\hyperlink{cf:eq:3.2}{(95)}, take \(\fitinline{Z=\max(2,D,H)}\). If \(\fitinline{D\ge2}\) and \(\fitinline{H\ge D^{1+\vartheta}}\), choose \(\fitinline{\kappa=\vartheta/(1+\vartheta)}\); rename the loss to get \(\fitinline{H(HD)^\varepsilon}\). Count bounded columns. All element rows, both orientations and common-profile control follow. \ensuremath{\square}

\hypertarget{cf:sec:D}{}

\hypertarget{a.5-the-plain-moment-with-moving-radical}{%
\subsubsection{A.5 The plain moment with moving radical}\label{a.5-the-plain-moment-with-moving-radical}}

\hypertarget{cf:sec:D.1}{}

\hypertarget{a.5.1-equal-product-rectangle-cancellation}{%
\paragraph{A.5.1 Equal-product rectangle cancellation}\label{a.5.1-equal-product-rectangle-cancellation}}

Two Poisson transforms decrease width except for fixed-character rows; equal-product rectangle subtraction cancels their principal volume.

Fix a finite group \(\fitinline{\Theta}\) of \(\fitinline{S}\)-supported ray characters closed under reciprocity, conjugation and \(\fitinline{n\mapsto\chi_n(-1)}\). For \(\fitinline{\vartheta\in\Theta}\), squarefree \(\fitinline{q_{R_*}\le Z^B}\) and fixed-annular \(\fitinline{W_1,W_2}\), put

\[\fitmath{
L_i(X,t)=\sum_n\vartheta(n)\mathbf1_{(n,R_*)=1}q_n^{it}W_i(q_n/X).
}\]

Poisson summation gives, uniformly for \(\fitinline{X>0}\),

\hypertarget{cf:eq:D.1}{}

\[\fitmath{
L_i(X,t)=c_{\vartheta,R_*}X^{1+it}I_i(t)
 +O\bigl(Z^\epsilon(1+|t|)^J\bigr),\qquad |L_i(X,t)|\ll X,
\qquad\text{(65)}
}\]

Here \(\fitinline{I_i(t)}\) is a fixed lattice-measure multiple of \(\fitinline{\int W_i(y)y^{it}\,dy}\), and \(\fitinline{c_{\vartheta,R_*}}\) is \(\fitinline{X,t}\)-independent.

Without the mask, Poisson on fixed lattice cosets for \(\fitinline{X^{it}f_t(z/\sqrt X)}\) gives the zero-frequency main term. Integration by parts bounds other frequencies by \(\fitinline{(1+|t|)^J(1+|\xi|)^{-J}}\); \(\fitinline{J>4}\) handles \(\fitinline{X\ge1}\), counting \(\fitinline{X<1}\). Insert

\[\fitmath{
L_i(X,t)=\sum_{d\mid R_*}\mu(d)\vartheta(d)q_d^{it}L_i^0(X/q_d,t).
}\]

The main coefficient is

\[\fitmath{
c_{\vartheta,R_*}=c_\vartheta\sum_{d\mid R_*}\frac{\mu(d)\vartheta(d)}{q_d}.
}\]

Norm powers cancel in this coefficient; divisor-many errors cost \(\fitinline{Z^\epsilon}\). Counting gives \(\fitinline{|L_i|\ll X}\).

Consequently, if

\[\fitmath{
U_1U_2=V_1V_2=T,\qquad U_i,V_i\ge Z^r\quad(i=1,2),\quad r\ge0,
}\]

then

\hypertarget{cf:eq:D.2}{}

\[\fitmath{
T^{-1/2}\left|L_1(U_1,t)L_2(U_2,t)-L_1(V_1,t)L_2(V_2,t)\right|
\ll Z^\epsilon(1+|t|)^{2J}T^{1/2}Z^{-r}.
\qquad\text{(66)}
}\]

Main terms equal \(\fitinline{c_{\vartheta,R_*}^2T^{1+it}I_1(t)I_2(t)}\). Error scales are \(\fitinline{\le TZ^{-r}}\); the error product also fits since \(\fitinline{T\ge Z^{2r}}\). Without a positive lower scale, use volume \(\fitinline{O(T^{1/2})}\).

The character, mask and norm power must agree across both factors and rectangles.

\hypertarget{cf:sec:D.2}{}

\hypertarget{a.5.2-the-family-closed-under-both-transforms}{%
\paragraph{A.5.2 The family closed under both transforms}\label{a.5.2-the-family-closed-under-both-transforms}}

Moving twists and a common exclusion close the family under both transforms; no selected prime or zero-free hypothesis is used.

Fix bounded nonnegative real ranges for \(\fitinline{m,q,n_1,n_2,B}\). A fixed twist within a row sum has the form

\[\fitmath{
\tau(n)=\nu(n)\prod_{p\mid A}\chi_n(p)^{e_p},\qquad \nu\in\Theta,\quad q_{\operatorname{rad}A}\le Z^q,
}\]

Retain all displayed zeros, including exponents \(\fitinline{0\bmod6}\), and count their moving support in \(\fitinline{q}\). Reciprocity reduces other residue products to this form. A squarefree puncture \(\fitinline{q_R\le Z^B}\) is fixed across the row sum and all plain factors. Define

\[\fitmath{
\psi_k(n)=\tau(n)\chi_n(k)\mathbf1_{(n,R)=1},\qquad
S_k(N;W)=N^{-1/2}\sum_n\psi_k(n)W(q_n/N).
}\]

\textbf{Theorem A.5.1} (Moving-radical mixed plain moment). \emph{\hypertarget{cf:mixed-plain-moment}{} For every fixed \(\fitinline{\epsilon>0}\),}

\hypertarget{cf:eq:D.3}{}

\[\fitmath{
\sum_{\substack{0<q_k\le Z^m\\\psi_k^*\ne1}}
 |S_k(Z^{n_1};W_1)S_k(Z^{n_2};W_2)|^2
\ll Z^{m+q+\epsilon}.
\qquad\text{(67)}
}\]

\emph{No conductor lower bound or condition on \(\fitinline{n_1+n_2}\) is imposed. The bound carries \(\fitinline{\|W_1\|_{C^J}^2\|W_2\|_{C^J}^2}\), uniformly over moving supports, punctures and bounded lengths, with \S{}\protect\hyperlink{cf:sec:1.2}{A.1.3}'s real-power and height control. Increase \(\fitinline{J}\) for any fixed list of Euler/profile derivatives. With no moving twist and two copies of one polynomial,}

\hypertarget{cf:eq:D.4}{}

\[\fitmath{
\sum_{\substack{0<q_k\le H\\(\nu\chi_\bullet(k))^*\ne1}}
 |S_k(N;W)|^4\ll H^{1+\epsilon}
\qquad\text{(68)}
}\]

\emph{for bounded polynomial ranges of \(\fitinline{N}\) relative to \(\fitinline{H}\).}

\emph{Proof.} Descend in effective width \(\fitinline{M=m+q}\); fixed-conductor twists alone do not close under the first transform.

\hypertarget{cf:sec:D.3}{}

\hypertarget{a.5.3-reflection-and-the-centered-comparison}{%
\paragraph{A.5.3 Reflection and the centered comparison}\label{a.5.3-reflection-and-the-centered-comparison}}

A natural presentation has exactly fixed, displayed moving and row zeros. Remove extra punctures by

\hypertarget{cf:eq:D.5}{}

\[\fitmath{
S_{k,R}(N;W)=\sum_{d\mid R}\mu(d)\psi_k^0(d)q_d^{-1/2}S_k^0(N/q_d;W).
\qquad\text{(69)}
}\]

Its coefficient mass is \(\fitinline{Z^\epsilon}\). At each width prove all natural lengths before removing extra punctures; recursive children have smaller width.

Let \(\fitinline{\psi_k^*}\) have conductor norm \(\fitinline{Q_k}\), and put \(\fitinline{C_k=3Q_k/(2\pi)^2}\). If \(\fitinline{R_{0,k}}\) is the redundant natural radical, then

\hypertarget{cf:eq:D.6}{}

\[\fitmath{
C_kq_{R_{0,k}}\ll Z^M.
\qquad\text{(70)}
}\]

Each good prime is counted once in primitive conductor or redundant radical, whose union lies in row radical and displayed moving support. Fixed primes cost a constant.

The primitive functional equation gives exactly

\hypertarget{cf:eq:D.7}{}

\[\fitmath{
S_{\psi_k^*}(N;W)=\epsilon_k S_{\overline{\psi_k^*}}(C_k/N;W^\sharp),\qquad
\mathcal MW^\sharp(s)=\mathcal MW(1-s)\frac{\Gamma(s)}{\Gamma(1-s)}.
\qquad\text{(71)}
}\]

Central normalization adds no \(\fitinline{C_k,N}\) power. Inside an absolute product, conjugate the entire reflected factor and profile to restore orientation. Delete redundant Euler factors before reflection, restore geometrically afterwards: scales are \(\fitinline{C_kq_d/(Nq_h)}\), \(\fitinline{d\mid R_{0,k}}\), \(\fitinline{\operatorname{rad}h\mid R_{0,k}}\), with coefficients \(\fitinline{\epsilon_k\mu(d)\psi_k^*(d)\overline{\psi_k^*(h)}/\sqrt{q_dq_h}}\). Their mass is \(\fitinline{Z^\epsilon}\) and largest scale \(\fitinline{\ll Z^M/N}\).

The gamma quotient gives bounded Euler derivatives at zero and rapid decay at infinity, hence a centrally summable annular partition with polynomial height cost. With padding \(\fitinline{\xi>0}\), retain dual lengths \(\fitinline{\le M-\log_ZN+\xi}\) and arbitrarily small tails. Apply \S{}\protect\hyperlink{cf:sec:1.2}{A.1.3}'s Sobolev common-profile bounds to rowwise scales.

Reflect any original factor longer than \(\fitinline{M/2}\). It remains to prove a padded core with

\hypertarget{cf:eq:D.8}{}

\[\fitmath{
n_i\le M/2+\xi,\qquad A:=n_1+n_2\le M+2\xi.
\qquad\text{(72)}
}\]

At each width first prove the uncentered range \(\fitinline{A\le5M/6}\). Otherwise put \(\fitinline{L=M/4}\). If one factor is shorter than \(\fitinline{L}\), reflect the other: total becomes \(\fitinline{\le3M/2-A+\xi<2M/3+\xi}\), already covered. If both are \(\fitinline{\ge L}\), compare

\[\fitmath{
S_k(X_1;W_1)S_k(X_2;W_2)
\quad\hbox{with}\quad
S_k(Y_1;W_1)S_k(Y_2;W_2),
}\]

where \(\fitinline{Y_1=Z^L}\), \(\fitinline{Y_2=X_1X_2/Y_1}\). Reflecting the latter's second factor again gives total \(\fitinline{\le3M/2-A+\xi}\), so the comparison is known at this width. Retain the difference as one coefficient; its positive squared norm may include all rows. Subtraction controls added principal rows.

For fixed extracted ideals \(\fitinline{b_1,b_2}\), the coefficient to preserve is

\hypertarget{cf:eq:D.9}{}

\[\fitmath{
\mathcal D_b(l_1,l_2)=
W_1(q_{b_1}q_{l_1}/X_1)W_2(q_{b_2}q_{l_2}/X_2)
-W_1(q_{b_1}q_{l_1}/Y_1)W_2(q_{b_2}q_{l_2}/Y_2).
\qquad\text{(73)}
}\]

All four formal scales are \(\fitinline{\ge Z^L}\) and extracted products stay equal. Retain empty rectangles until cancellation.

\hypertarget{cf:sec:D.4}{}

\hypertarget{a.5.4-first-poisson-transform}{%
\paragraph{A.5.4 First Poisson transform}\label{a.5.4-first-poisson-transform}}

Set \(\fitinline{H=Z^m}\), \(\fitinline{X=Z^A}\) and majorize rows by fixed nonnegative radial Schwartz \(\fitinline{\Phi_1}\); central squared normalization is \(\fitinline{X^{-1}}\). Zero frequency requires every prime's net sextic exponent to vanish, hence no prime occurs once. The combined ideal is powerful, norm \(\fitinline{O(X^2)}\): \(\fitinline{O(X^{1+\epsilon})}\) possibilities with divisor-many four-factor allocations, costing \(\fitinline{O(HZ^\epsilon)\le Z^{M+\epsilon}}\).

For nonzero frequency extract full common support as \(\fitinline{Ca,Db}\), \(\fitinline{(a,b)=(ab,CD)=1}\), \(\fitinline{\operatorname{rad}C=\operatorname{rad}D}\). Put

\[\fitmath{
c=\log_Zq_C,\ d=\log_Zq_D,\ p=\log_Zq_{\operatorname{rad}C}.
}\]

Let \(\fitinline{r}\) contain common primes of nonzero net exponent modulo six, \(\fitinline{R=\log_Zq_r}\), with primitive common character \(\fitinline{\xi_r}\). Expand the remaining common mask by \(\fitinline{e\mid\operatorname{rad}C/r}\), \(\fitinline{E=\log_Zq_e}\); thus \(\fitinline{R+E\le p}\).

For the allocated convolution coefficients \(\fitinline{A_C(a),A_D(b)}\), residue-class Poisson gives, up to a frozen scalar of modulus at most one,

\hypertarget{cf:eq:D.10}{}

\[\fitmath{
\frac{H}{Xq_e\sqrt{q_r}}
\sum_{h\ne0}\mathfrak g_{\xi_r}(r,h)
\sum_{(a,b)=1}\frac{A_C(a)\overline{A_D(b)}}{\sqrt{q_aq_b}}
\tau_C(a)\overline{\tau_D(b)}R(a,b)
\mathfrak g(a,h)\overline{\mathfrak g(b,-h)}
\widehat\Phi_1\!\left(\frac{Hq_h}{q_eq_rq_aq_b}\right).
\qquad\text{(74)}
}\]

Here \(\fitinline{\mathfrak g_{\xi_r}}\) is the normalized primitive Gauss sum, of modulus at most one, and one consistent orientation is

\[\fitmath{
\tau_C(a)=\tau(a)\chi_a(er)\xi_r(a),\qquad
\tau_D(b)=\tau(b)\chi_b(er)\overline{\xi_r(b)}.
}\]

Indeed \(\fitinline{k=ek'}\) and Poisson modulo \(\fitinline{rab}\) give \(\fitinline{H/(q_eq_rq_aq_b)}\); three Gauss sums restore \(\fitinline{\sqrt{q_rq_aq_b}}\) and central normalization supplies \(\fitinline{X^{-1}}\). Reciprocity gives \(\fitinline{R(a,b)}\). Full displayed moving length becomes

\hypertarget{cf:eq:D.11}{}

\[\fitmath{
\widetilde q=q+R+E.
\qquad\text{(75)}
}\]

Old moving zeros puncture both products, so discard colliding allocations first; new common support is disjoint even when exponents cancel.

The natural first dual scale has logarithmic length

\hypertarget{cf:eq:D.12}{}

\[\fitmath{
K_0=2A-c-d+R+E-m.
\qquad\text{(76)}
}\]

Retain full smooth dyads \(\fitinline{q_h\asymp Z^K}\), \(\fitinline{K\le K_0+\xi}\). Separate only row and two \textbf{whole-product} norms with one Fourier measure; expand \(\fitinline{R(a,b)}\) on its ray group. Before Cauchy insert \protect\hyperlink{cf:eq:A.19}{(15)}'s full \(\fitinline{\sum_{s\mid a,b}\mu(s)}\), and put \(\fitinline{s_0=\log_Zq_s}\).

The two resulting positive Gauss norms have the form

\hypertarget{cf:eq:D.13}{}

\[\fitmath{
\mathcal N_C=
\sum_h\omega_K(q_h/Z^K)
\left|Z^{-(A-c)/2}\sum_{s\mid a}A_C(a)\tau'_C(a)q_a^{it}\mathfrak g(a,h)\right|^2,
\qquad\text{(77)}
}\]

and similarly for \(\fitinline{D}\). Within each norm the profiles and masks are common; row dependence is a bounded scalar. Coefficient \protect\hyperlink{cf:eq:D.9}{(73)} retains a common whole-product character and norm power.

A sufficient estimate is

\hypertarget{cf:eq:D.14}{}

\[\fitmath{
\mathcal N_C\ll Z^{A-c+\widetilde q+B_c-s_0+\epsilon},\qquad
B_c=\frac{(3c-5d-R)_+}{6}.
\qquad\text{(78)}
}\]

The complete exponent balance is

\[\fitmath{
\underbrace{m-A-R/2-E}_{\text{Poisson prefactor}}
+\underbrace{p+s_0}_{\text{common-support and divisor counts}}
+\frac12\{A-c+\widetilde q+B_c-s_0+A-d+\widetilde q+B_d-s_0\}
}\]

\hypertarget{cf:eq:D.15}{}

\[\fitmath{
=M+\frac12\{B_c+B_d-(c+d-2p-R)\}\le M.
\qquad\text{(79)}
}\]

Locally at multiplicities \(\fitinline{i\ge j\ge1}\), put \(\fitinline{r_p=1_{6\nmid i-j}}\). The contribution is \(\fitinline{\le(3i-5j-r_p)_+/6}\), compared with nonnegative \(\fitinline{i+j-2-r_p}\). When the former is positive, six times their difference is \(\fitinline{3i+11j-12-5r_p\ge0}\). Subadditivity of positive parts proves \protect\hyperlink{cf:eq:D.15}{(79)}. For fixed radical, Rankin and divisor bounds give \(\fitinline{Z^\epsilon}\) allowed powers and allocations, verifying all label counts.

\hypertarget{cf:sec:D.5}{}

\hypertarget{a.5.5-second-transform-and-width-decrease}{%
\paragraph{A.5.5 Second transform and width decrease}\label{a.5.5-second-transform-and-width-decrease}}

Set \(\fitinline{a_0=A-c}\), \(\fitinline{D_0=d+(K_0-K)}\), and

\hypertarget{cf:eq:D.16}{}

\[\fitmath{
J=D_0-c,\qquad g=J_++\sigma,
\qquad\text{(80)}
}\]

Choose small width step \(\fitinline{\sigma>0}\). Enlarge the positive Gauss norm to a smooth ball of length \(\fitinline{K+g}\), adding its zero row. By \protect\hyperlink{cf:eq:A.12}{(8)}, only sixth-power columns contribute: \(\fitinline{O(Z^{a_0/6+\epsilon})}\) columns, Gauss magnitude \(\fitinline{\le Z^{a_0/2}}\), hence centrally normalized square \(\fitinline{Z^{a_0/3+\epsilon}}\). Nonempty \(\fitinline{s\mid a}\) forces \(\fitinline{s_0\le a_0/6+O(1/\log Z)}\), so \protect\hyperlink{cf:eq:D.14}{(78)} covers it.

For the remaining norm, \protect\hyperlink{cf:eq:A.14}{(10)} and lattice Poisson give the exact expansion

\hypertarget{cf:eq:D.17}{}

\[\fitmath{
Z^{K+g-a_0}\sum_{u,v}
\frac{B(u)\overline{B(v)}\tau_1(u)\overline{\tau_1(v)}q_u^{it}q_v^{-it}}{\sqrt{q_uq_v}}
\sum_jF(u,v;j)\widehat\Phi_2\!\left(\frac{Z^{K+g}q_j}{q_uq_v}\right),
\qquad\text{(81)}
}\]

Here \(\fitinline{B}\) retains the whole rectangle difference and \(\fitinline{s\mid u,v}\). At \(\fitinline{j=0}\), \protect\hyperlink{cf:eq:A.15}{(11)} forces \(\fitinline{u=v}\), giving exponent \(\fitinline{K+g-s_0}\) and excess over \protect\hyperlink{cf:eq:D.14}{(78)}

\hypertarget{cf:eq:D.18}{}

\[\fitmath{
(A-M)-D_0-B_c+g\le (A-M)+\sigma+O(\xi).
\qquad\text{(82)}
}\]

Since \(\fitinline{D_0\ge-\xi}\), \(\fitinline{-D_0+(D_0-c)_+\le O(\xi)}\); thus padded \(\fitinline{A\le M+2\xi}\) suffices.

For \(\fitinline{j\ne0}\), retain full-kernel dyads through length \(\fitinline{2a_0-K-g+\xi}\), then extract common support \(\fitinline{u=D_2a,v=E_2b}\). Put

\[\fitmath{
c_2=\log_Zq_{D_2},\ d_2=\log_Zq_{E_2},\ b_2=(c_2+d_2)/2,
}\]

\[\fitmath{
p_2=\log_Zq_{\operatorname{rad}D_2},\quad g_2=\log_Zq_{(D_2,E_2)}.
}\]

At equal non-six-divisible multiplicities, split divided frequencies into units and nonunits, whose radical lengths are \(\fitinline{t_2,V}\). By \protect\hyperlink{cf:eq:A.17}{(13)}--\protect\hyperlink{cf:eq:A.18}{(14)},

\[\fitmath{
j=(D_2,E_2)V_{\rm id}h',\qquad
|F(D_2,E_2;j)|\le Z^{g_2-t_2},
}\]

where \(\fitinline{V_{\rm id}}\) is the product of the nonunit primes. Set the scalar to zero off its partition before enlarging rows. New row and displayed moving lengths are

\[\fitmath{
m'=2a_0-K-g-g_2-V,\quad q'=\widetilde q+t_2+V,
}\]

\hypertarget{cf:eq:D.19}{}

\[\fitmath{
M'=m'+q'=M+J-g-g_2+t_2\le M-\sigma.
\qquad\text{(83)}
}\]

Here \(\fitinline{g_2\ge t_2}\). Sixth-divisible fixed factors give a common puncture with retained ray phase, without a new pair-dependent mask.

Since \(\fitinline{s\mid u,v}\) lies in the second common radical, its count is \(\fitinline{Z^{p_2-s_0+\epsilon}}\). The ledger is

\hypertarget{cf:eq:D.20}{}

\[\fitmath{
\underbrace{K+g-a_0}_{\text{Poisson and roots}}
+\underbrace{p_2-s_0}_{\text{common supports}}
+\underbrace{g_2-t_2}_{\text{correlation}}
-\underbrace{b_2}_{\text{remaining central conversion}},
\qquad\text{(84)}
}\]

so the inner residual products are allowed exponent

\hypertarget{cf:eq:D.21}{}

\[\fitmath{
M'+\Delta_{\rm child},\qquad \Delta_{\rm child}=b_2-p_2+B_c\ge0.
\qquad\text{(85)}
}\]

Separate \protect\hyperlink{cf:eq:D.17}{(81)}'s kernel and inverse roots only in row and whole-product norms. Remove residual coprimality by \protect\hyperlink{cf:eq:A.19}{(15)}; for each prime of its divisor \(\fitinline{t}\), use

\hypertarget{cf:eq:D.22}{}

\[\fitmath{
\mathbf1_{p\mid l_1l_2}=\mathbf1_{p\mid l_1}+\mathbf1_{p\mid l_2}
-\mathbf1_{p\mid l_1}\mathbf1_{p\mid l_2}.
\qquad\text{(86)}
}\]

In each selected factor put \(\fitinline{l_i=pl_i'}\) without \(\fitinline{(l_i',p)=1}\). The common character supplies a scalar, dividing corresponding scales in both rectangles by \(\fitinline{P}\); mask and norm power remain common. On \(\fitinline{q_t\asymp Z^{t_-}}\), each side extracts total length \(\fitinline{r_j\ge t_-+O(1/\log Z)}\). Central factors \(\fitinline{Z^{-r_j/2}}\) cancel the \(\fitinline{Z^{t_-}}\) count; allocations cost divisors.

Both children are the original whole-product twist times \(\fitinline{\chi_n((D_2,E_2)V_{\rm id}h')}\), differing by \(\fitinline{\Theta}\), so either both induce a member of \(\fitinline{\Theta}\) or neither. In the latter case apply Cauchy, then split centered children, discard empty annuli and rescale bounded subunit annuli. The completed theorem at smaller width \(\fitinline{M'+O(\xi)}\) gives \(\fitinline{Z^{M'+\epsilon}}\), within \protect\hyperlink{cf:eq:D.21}{(85)}, with common polynomial masks. It remains to treat fixed-character children.

\hypertarget{cf:sec:D.6}{}

\hypertarget{a.5.6-exceptional-children}{%
\paragraph{A.5.6 Exceptional children}\label{a.5.6-exceptional-children}}

Outside fixed support, exceptionality prescribes \(\fitinline{v_p(h')\bmod6}\), so \(\fitinline{h'=h_0v^6}\) for fixed sixth-power-free \(\fitinline{h_0}\). Let \(\fitinline{v_1}\) be radical length of nonunit common multiplicity-one primes, \(\fitinline{f=2v_1}\). Their fixed row valuation is two, forcing \(\fitinline{v_p(h')\equiv4\pmod6}\). Old moving support is disjoint by its zeros, hence \(\fitinline{q_{h_0}\ge Z^{4v_1}}\). The weaker sufficient count is

\hypertarget{cf:eq:D.23}{}

\[\fitmath{
\#\{h'\text{ exceptional}\}\ll Z^{(m'-f)/6+\epsilon}.
\qquad\text{(87)}
}\]

This covers all \(\fitinline{\Theta}\) and positively added rows; a forced \(\fitinline{h_0}\) that does not fit gives an empty range.

After \protect\hyperlink{cf:eq:D.22}{(86)}, paired normalized convolution volume is \(\fitinline{\ll Z^{a_0-b_2+\epsilon}}\). Its excess, with \protect\hyperlink{cf:eq:D.23}{(87)}, over \protect\hyperlink{cf:eq:D.21}{(85)} is

\hypertarget{cf:eq:D.24}{}

\[\fitmath{
\frac{m'-f}{6}+a_0-b_2-(M'+\Delta_{\rm child})
=A-\frac56M-F_1-F_2,
\qquad\text{(88)}
}\]

where

\[\fitmath{
F_1=\frac c6+\frac{5D_0}{6}+\frac{\widetilde q}{6}+B_c-\frac{5g}{6},
}\]

\hypertarget{cf:eq:D.25}{}

\[\fitmath{
F_2=2b_2-\frac56g_2-p_2+t_2+\frac V6+\frac f6.
\qquad\text{(89)}
}\]

There are the two local lower bounds

\hypertarget{cf:eq:D.26}{}

\[\fitmath{
F_1\ge\frac23c-\frac56\sigma-O(\xi),\qquad F_2\ge\frac23b_2.
\qquad\text{(90)}
}\]

If \(\fitinline{D_0\ge c}\), substitution gives \(\fitinline{F_1=c+\widetilde q/6+B_c-5\sigma/6}\). Otherwise \(\fitinline{\widetilde q\ge R}\) implies \(\fitinline{\widetilde q+6B_c\ge(3c-5D_0)_+-O(\xi)}\); use \(\fitinline{c+5D_0+(3c-5D_0)_+\ge4c}\).

For \(\fitinline{F_2}\), the local contributions, in units of \(\fitinline{\log_ZP}\), are:

\[\fitmath{
\begin{array}{c|c|c}
\text{Common multiplicities; divided frequency}&F_2&b_2\\\hline
i=i,\ 6\nmid i,\ \text{unit}&7i/6&i\\
i=i,\ 6\nmid i,\ \text{nonunit}&(7i-5)/6+\tfrac13\mathbf1_{i=1}&i\\
i=i,\ 6\mid i&7i/6-1&i\\
i>j,\ 6\mid j,\ \text{unit}&i+j/6-1&(i+j)/2
\end{array}
}\]

The middle column is always \(\fitinline{\ge2/3}\) of the last, with equality at nonunit multiplicity one; these exhaust \protect\hyperlink{cf:eq:A.17}{(13)}--\protect\hyperlink{cf:eq:A.18}{(14)}.

For \(\fitinline{A\le5M/6}\), \protect\hyperlink{cf:eq:D.24}{(88)}--\protect\hyperlink{cf:eq:D.26}{(90)} finish the uncentered stage, with terminal loss \(\fitinline{5\sigma/6+O(\xi)}\).

For the centered stage, the common-character calculation gives, on an exceptional row, a coefficient of the form

\hypertarget{cf:eq:D.27}{}

\[\fitmath{
\vartheta(l_1)\vartheta(l_2)\mathbf1_{(l_1l_2,R_*)=1}q_{l_1l_2}^{it}\mathcal D_b(l_1,l_2).
\qquad\text{(91)}
}\]

The common mask contains all redundant natural zeros. Before final extraction, one paired side has all four formal scales at least

\[\fitmath{
Z^{L-c-\min(c_2,d_2)}.
}\]

Set \(\fitinline{r=(L-c-\min(c_2,d_2))_+}\). If that side extracts \(\fitinline{r_1<r}\), \protect\hyperlink{cf:eq:D.2}{(66)} saves \(\fitinline{Z^{-(r-r_1)}}\); otherwise use volume. Divisor counts satisfy

\hypertarget{cf:eq:D.28}{}

\[\fitmath{
t_--r_1-r_2-(r-r_1)_+\le-r
\qquad\text{(92)}
}\]

so all allocated paired products cost \(\fitinline{Z^{a_0-b_2-r+\epsilon}}\). Triangle before the full divisor sum loses this saving.

With \(\fitinline{v=c+\min(c_2,d_2)}\) and \(\fitinline{b_2\ge\min(c_2,d_2)}\), the remaining deficit is at most

\hypertarget{cf:eq:D.29}{}

\[\fitmath{
A-\frac56M-\frac23v-(L-v)_++\frac56\sigma+O(\xi)
\le A-M+\frac56\sigma+O(\xi).
\qquad\text{(93)}
}\]

The maximum over \(\fitinline{v\ge0}\) occurs at \(\fitinline{v=L=M/4}\); padded \(\fitinline{A\le M+2\xi}\) leaves only a small terminal loss.

\hypertarget{cf:sec:D.7}{}

\hypertarget{a.5.7-closing-the-finite-descent}{%
\paragraph{A.5.7 Closing the finite descent}\label{a.5.7-closing-the-finite-descent}}

Choose small floor \(\fitinline{\rho>0}\), step \(\fitinline{\sigma>0}\), then padding \(\fitinline{\xi>0}\). At \(\fitinline{M\le\rho}\), reflection and counting in \protect\hyperlink{cf:eq:D.8}{(72)} lose at most \(\fitinline{\rho+2\xi}\). Above it, \protect\hyperlink{cf:eq:D.19}{(83)} decreases width by \(\fitinline{\sigma}\) before enclosure; \(\fitinline{\xi\ll\sigma}\) leaves decrease \(\fitinline{\ge\sigma/2}\). Depth is at most \(\fitinline{D=2+\lceil2M_{\max}/\sigma\rceil}\).

Within each band prove: uncentered range, centered core, reflection to all lengths, then extra punctures by \protect\hyperlink{cf:eq:D.5}{(69)}. Comparison uses the already proved uncentered range with positive margin; other children use strictly smaller completed bands.

Make every terminal loss (floor, diagonal, exceptional) \(\fitinline{<\epsilon/4}\) by choosing \(\fitinline{\rho,\sigma}\). Choose padding and local divisor/Schwartz losses \(\fitinline{\ll\epsilon/D}\). Triangle takes the maximum of child exponents after counting mass; Cauchy averages them. Each branch incurs one terminal loss plus edge losses.

Measure powers before dyadic rounding. Use joint log boxes, selecting whole frequency dyads from the outer box's scale supremum with padding \(\fitinline{Z^{\xi/2}}\), independently of live columns. Discarded genuine Poisson arguments are \(\fitinline{\ge Z^{\xi/4}}\). Bounds \(\fitinline{|\mathfrak g(a,h)|\le\sqrt{q_a}}\), \(\fitinline{|F(u,v;j)|\le q_uq_v}\) and polynomial counts make tails arbitrarily small by Schwartz decay \textbf{before} separation.

Retained \(\fitinline{\widehat\Phi(R\exp L(x))}\) has uniform Euler derivatives for \(\fitinline{R>0}\) and a common Fourier measure of finite weighted \(\fitinline{L^1}\) norm. Remove central powers before inverse-root separation. Finite-depth support enlargements cost \(\fitinline{O(1/\log Z)}\); rescale subunit annuli only in positive child norms, preserving formal centered scales. Formulas \protect\hyperlink{cf:eq:D.10}{(74)}, \protect\hyperlink{cf:eq:A.16}{(12)}, \protect\hyperlink{cf:eq:D.22}{(86)} preserve character/common-mask closure with polynomial norms. Choose seminorm and height orders backwards, then thresholds and external tail orders. This proves \protect\hyperlink{cf:eq:D.3}{(67)}. \ensuremath{\square}

No moving twist and two identical profiles give \protect\hyperlink{cf:eq:3.3}{(96)}; restricting the positive row sum to sixth-power-free rows decreases it. Conjugation gives the inverse orientation.

\hypertarget{cf:moment-corollaries}{}

\hypertarget{a.5.8-coefficient-specific-corollaries-and-sixth-power-amplification}{%
\paragraph{A.5.8 Coefficient-specific corollaries and sixth-power amplification}\label{a.5.8-coefficient-specific-corollaries-and-sixth-power-amplification}}

For a fixed finite-ray character \(\fitinline{\nu}\), let \(\fitinline{\varphi_u=\nu\psi_u^{\pm1}}\), with its original zeros, and define

\hypertarget{cf:eq:3.1}{}

\[\fitmath{
M_u(D;W)=D^{-1/2}\sum_n\mu(n)\varphi_u(n)W(q_n/D),\qquad
 S_u(N;W)=N^{-1/2}\sum_n\varphi_u(n)W(q_n/N).
\qquad\text{(94)}
}\]

The preceding theorems imply the following two coefficient-specific estimates:

\hypertarget{cf:eq:3.2}{}

\[\fitmath{
\sum_{0<q_u\le H}|M_u(D;W)|^2\ll H(HD)^\varepsilon\|W\|_{C^J}^2,
 \qquad H\ge D^{1+\vartheta},\quad \vartheta>0;
\qquad\text{(95)}
}\]

\hypertarget{cf:eq:3.3}{}

\[\fitmath{
\sum_{\substack{q_u\asymp U\text{ sixth-power-free}\\\varphi_u^*\ne1}}
 |S_u(U^m;W)|^4\ll U^{1+\varepsilon}\|W\|_{C^J}^4.
\qquad\text{(96)}
}\]

Here \(\fitinline{m\ge0}\) is bounded and \(\fitinline{\varphi_u^*}\) is the primitive inducing character. Equations \protect\hyperlink{cf:eq:3.2}{(95)}, \protect\hyperlink{cf:eq:3.3}{(96)} follow from Theorems 4.1 and 5.1, with \S{}\protect\hyperlink{cf:sec:1.2}{A.1.3}'s uniformity and no zero-free assumption; \protect\hyperlink{cf:eq:3.2}{(95)} includes sixth-power rows.

Sixth-power amplification extends \protect\hyperlink{cf:eq:3.2}{(95)} to longer inverse polynomials. The exact reconstruction identity is

\hypertarget{cf:eq:3.4}{}

\[\fitmath{
M_u(D;W)=\sum_{d\mid\operatorname{rad}(a)}
 \mu(d)\varphi_u(d)q_d^{-1/2}M_{ua^6}(D/q_d;W).
\qquad\text{(97)}
}\]

Separate squarefree column primes dividing \(\fitinline{a}\) to prove \protect\hyperlink{cf:eq:3.4}{(97)}, retaining shared \(\fitinline{u,a}\) zeros. Its coefficient mass is \(\fitinline{\ll q_a^\varepsilon}\). Log-scale differentiation sends \(\fitinline{W}\) to \(\fitinline{-W/2-yW'}\), so Sobolev and \protect\hyperlink{cf:eq:3.2}{(95)} give

\hypertarget{cf:eq:3.5}{}

\[\fitmath{
\sum_{0<q_v\ll H}\sup_{0<E\le D}|M_v(E;W)|^2\ll H(HD)^\varepsilon,
 \qquad H\ge D^{1+\vartheta}.
\qquad\text{(98)}
}\]

Only \(\fitinline{O(\log(2+D))}\) log intervals occur; subunit scales have boundedly many columns and eventually none.

For sixth-power-free \(\fitinline{q_u\asymp U}\), put \(\fitinline{H=\max(2U,D^{1+\vartheta})}\) and average \protect\hyperlink{cf:eq:3.4}{(97)} over the \(\fitinline{\gg_S(H/U)^{1/6}}\) primary ideals \(\fitinline{q_a\le(H/U)^{1/6}}\) prime to \(\fitinline{S}\). The map \(\fitinline{(u,a)\mapsto ua^6}\) is injective: valuations modulo six recover \(\fitinline{u}\), primary generators its unit. Equation \protect\hyperlink{cf:eq:3.5}{(98)} gives \(\fitinline{\sum^{(6)}|M_u(D)|^2\ll U^{1/6}H^{5/6}(UD)^\varepsilon}\). For bounded \(\fitinline{r}\), take \(\fitinline{D=U^r}\) and small \(\fitinline{\vartheta}\):

\hypertarget{cf:eq:3.6}{}

\[\fitmath{
\sum_{q_u\asymp U}^{(6)}|M_u(U^r;W)|^2
 \ll U^{e(r)+\varepsilon},\qquad e(r)=\max\{1,(1+5r)/6\}.
\qquad\text{(99)}
}\]

\hypertarget{cf:probe-arithmetic}{}

\hypertarget{a.6-the-exact-poisson-series-of-the-common-probe}{%
\subsubsection{A.6 The exact Poisson series of the common probe}\label{a.6-the-exact-poisson-series-of-the-common-probe}}

Here \(\fitinline{X,Y,Z>0}\) are independent. This unmarked probe supplies the exact coefficient interface; compensated or affine probes must derive their combinations from it.

Fix primitive \(\fitinline{\eta}\) and let \(\fitinline{b_*}\) generate the product of \(\fitinline{S}\). Choose primitive \(\fitinline{\xi\bmod b_*}\), quadratic at odd residue characteristic and cubic at norm four; thus \(\fitinline{\xi^6=1}\) on units. Put

\[\fitmath{
\tau=q_{b_*}^{-1/2}g_\xi(b_*,1),\qquad
 \Xi(A)=\xi(A)\chi_A(b_*),\qquad |\tau|=1.
}\]

Choose nonnegative nonzero annular weights \(\fitinline{W_0,W_1,w}\), and set

\hypertarget{cf:eq:2.2}{}

\[\fitmath{
V_G(r)=\frac{e^{-(\log r)^2/4}}{2\sqrt\pi},\qquad \widehat V_G(t)=e^{t^2},
\qquad\text{(100)}
}\]

The ideals \(\fitinline{c,n}\) may share primes. With \(\fitinline{A=cn^3}\), define

\hypertarget{cf:eq:2.3}{}

\[\fitmath{
T_{m,a,\eta}(t)=\sum_{c\ {\rm sf},n}
 \gamma_2(c)\overline{\alpha(A)G(A)}\eta(A)\Xi(A)^{-1}R(A,a)
 \chi_A(m)q_c^{-1/2-t}q_n^{-1-3t},
\qquad\text{(101)}
}\]

\hypertarget{cf:eq:2.4}{}

\[\fitmath{
\begin{aligned}
 I_\eta(X,Y,Z)=\frac1Y\sum_a
 &\frac{W_1(q_a/Y)\chi_a(b_*)}{\tau\xi(a)\sqrt{q_{b_*}q_aX}}
 \sum_m\xi(m)q_a^{-1/2}g_{\chi_a}(a,-m)
 W_0\!\left(\frac{q_m}{q_{b_*}q_aX}\right)\\
 &\times\frac1{2\pi i}\int_{(4)}Z^te^{t^2}T_{m,a,\eta}(t)\,dt.
\end{aligned}
\qquad\text{(102)}
}\]

Annular weights make \(\fitinline{a,m}\) finite, \(\fitinline{m\ne0}\); remaining sums/integrals converge absolutely. Expand \(\fitinline{\overline G}\) on its ray group and partition \(\fitinline{a}\) into classes \(\fitinline{\sigma}\) for \(\fitinline{R(A,a)}\). Then \(\fitinline{B_{m,\sigma}(Z)}\) is a fixed finite combination of \protect\hyperlink{cf:completed-definition}{(16)} with \(\fitinline{V_G}\). Put \(\fitinline{Q_0=q_{b_*}XY}\) and fix \(\fitinline{\Omega=1}\) on \(\fitinline{\operatorname{supp}W_0\operatorname{supp}W_1}\). Mellin inversion gives a finite sum of

\hypertarget{cf:eq:2.6}{}

\[\fitmath{
\frac{Q_0^{-1/2}}{2\pi}\int_{\mathbf R}\widehat W_0(iv)
 \sum_{m\ne0}\Omega(q_m/Q_0)\xi(m)(q_m/Q_0)^{-iv}
 A_{m,\sigma,v}(Y)B_{m,\sigma}(Z)\,dv,
\qquad\text{(103)}
}\]

where

\hypertarget{cf:eq:2.7}{}

\[\fitmath{
A_{m,\sigma,v}(Y)=Y^{-3/2}\sum_{a\in\sigma}
 \frac{\chi_a(b_*)}{\tau\xi(a)}W_1(q_a/Y)(q_a/Y)^{-1+iv}g_{\chi_a}(a,-m).
\qquad\text{(104)}
}\]

The normalization follows from \(\fitinline{q_m/(q_{b_*}q_aX)=(q_m/Q_0)/(q_a/Y)}\).

Let \(\fitinline{\widetilde W_0(|y|^2)}\) be the planar Fourier transform of \(\fitinline{W_0(|z|^2)}\), and \(\fitinline{\mathcal M(z)=\int_0^\infty\widetilde W_0(r)r^{z-1}\,dr}\), holomorphic for \(\fitinline{\Re z>0}\) and rapidly decreasing vertically on compact substrips.

\textbf{Proposition A.6.1} (Exact common probe series). \hypertarget{cf:exact-probe-series}{}

\hypertarget{cf:eq:2.8}{}

\[\fitmath{
\begin{aligned}
 I_\eta(X,Y,Z)=\frac1{(2\pi i)^3}\int_{(3)}\int_{(3)}\int_{(2)}
 &X^{1/2-z}Z^{s+z-1}Y^{w-1}e^{(s+z-1)^2}
 \mathcal M(z)\widehat W_1(w)\\
 &\times\sum_u^{(6)}q_u^{-z}\overline{\xi(u)}
 \frac{\zeta_K^S(6z)L_K^S(w,\psi_u)}{L_K^S(s,\eta\overline{\psi_u})}
 H_{\eta,u}(s,w,z)\,dz\,dw\,ds.
\end{aligned}
\qquad\text{(105)}
}\]

\emph{Here \(\fitinline{u\ne0}\) is prime to \(\fitinline{S}\), including all units.}

Section \protect\hyperlink{cf:sec:A.3}{A.6.2} derives the local factors \(\fitinline{H_{\mathfrak p}}\) and full/prime-divisible series \(\fitinline{P_{\mathfrak p},P_{\mathfrak p}^*}\).

\emph{Proof.}

\hypertarget{cf:sec:A.2}{}

\hypertarget{a.6.1-the-finite-transform-and-phase-cancellation}{%
\paragraph{A.6.1 The finite transform and phase cancellation}\label{a.6.1-the-finite-transform-and-phase-cancellation}}

Expanding the Gauss sum in the physical row gives the finite coefficient

\hypertarget{cf:eq:A.3}{}

\[\fitmath{
\mathcal F(a,A,h)=
\sum_{\substack{d\bmod a\\h\equiv b_*Ad\pmod a}}
\chi_a(d)\,g_{\xi\chi_A}\!\left(b_*A,\frac{h-b_*Ad}{a}\right).
\qquad\text{(106)}
}\]

The finite transform of \(\fitinline{\xi(m)\chi_A(m)q_a^{-1/2}g_{\chi_a}(a,-m)}\) modulo full \(\fitinline{ab_*A}\) equals \(\fitinline{q_a^{1/2}\mathcal F(a,A,h)}\): summing lifts modulo \(\fitinline{b_*A}\) imposes the congruence, including common primes of \(\fitinline{a,A}\).

Poisson at scale \(\fitinline{q_{b_*}q_aX}\) gives factor \(\fitinline{X/q_A}\) and argument \(\fitinline{Xq_h/q_A}\). Combining \(\fitinline{q_a^{1/2}}\) with \protect\hyperlink{cf:eq:2.4}{(102)}'s square root leaves \(\fitinline{X^{1/2}/(\sqrt{q_{b_*}}q_A)}\) and the character/\(\fitinline{Y^{-1}}\) factors. Primitivity at \(\fitinline{b_*}\) kills \(\fitinline{(h,S)>1}\), including \(\fitinline{h=0}\). Write uniquely \(\fitinline{h=ur^6}\), \(\fitinline{u}\) sixth-power-free, \(\fitinline{r}\) primary.

Mellin inversion of the two radial weights, followed by \(\fitinline{s=t+1-z}\), gives the row Dirichlet series

\hypertarget{cf:eq:A.4}{}

\[\fitmath{
\mathcal F_{\eta,u}(s,w,z)=
\sum_{c\ {\rm sf},n,a,r} K_\eta(c,n,a,r;u)
q_c^{-1/2-s}q_n^{-1-3s}q_a^{-w}q_r^{-6z},
\qquad\text{(107)}
}\]

where

\hypertarget{cf:eq:A.5}{}

\[\fitmath{
K_\eta=
\frac{\mathcal F(a,A,ur^6)\chi_a(b_*)}
{\sqrt{q_{b_*}}\tau\xi(a)\overline{\xi(u)}}
\gamma_2(c)\overline{\alpha(A)G(A)}\eta(A)\Xi(A)^{-1}R(A,a).
\qquad\text{(108)}
}\]

The remaining row factor is \(\fitinline{q_u^{-z}\overline{\xi(u)}}\), and the global weight is

\[\fitmath{
X^{1/2-z}Z^{s+z-1}Y^{w-1}e^{(s+z-1)^2}\mathcal M(z)\widehat W_1(w).
}\]

Interchanges are absolute at \(\fitinline{(\Re s,\Re w,\Re z)=(3,3,2)}\) by \(\fitinline{|\mathcal F(a,A,h)|\le q_aq_{b_*A}}\), counting and Mellin decay. We factor \protect\hyperlink{cf:eq:A.4}{(107)}.

At good \(\fitinline{p}\), put \(\fitinline{Q=q_p}\), \(\fitinline{e_0=v_p(c)\in\{0,1\}}\), \(\fitinline{l=v_p(n)}\), \(\fitinline{k=v_p(a)}\), \(\fitinline{m=v_p(r)}\), \(\fitinline{t=e_0+3l}\), \(\fitinline{j=v_p(u)}\), \(\fitinline{j'=j+6m}\), \(\fitinline{\epsilon_p=\chi_p(-1)}\). The local convolution is

\hypertarget{cf:eq:A.6}{}

\[\fitmath{
C_{\mathfrak p}(t,k,j')=
\sum_{\substack{d\bmod\mathfrak p^k\\\mathfrak p^{j'}\equiv\mathfrak p^td\pmod{\mathfrak p^k}}}
\chi_{\mathfrak p}(d)^k
g_{\chi_{\mathfrak p}^t}\!\left(\mathfrak p^t,
\frac{\mathfrak p^{j'}-\mathfrak p^td}{\mathfrak p^k}\right).
\qquad\text{(109)}
}\]

At \(\fitinline{k=0}\) the absent outer character is one; at \(\fitinline{t=0}\) the Gauss factor is one. Residue lifting gives

\hypertarget{cf:eq:A.7}{}

\[\fitmath{
C_p(t,k,j')=
\begin{cases}
1,&t=k=0,\\
\mathbf1_{j'=0},&t=0,\ k>0,\\
g_{\chi_p^t}(p^t,p^{j'}),&t>0,\ k=0,\\
\epsilon_p^{-1}g_{\chi_p}(p,1)g_{\chi_p^{t-1}}(p^t,p^{j'-1}),
&t>0,\ k=1,\ j'\ge1,\\
0,&\text{otherwise}.
\end{cases}
\qquad\text{(110)}
}\]

For case four, expanding the Gauss sum and summing \(\fitinline{d\bmod p}\) changes exponent \(\fitinline{t}\) to \(\fitinline{t-1}\). For \(\fitinline{k\ge2}\), the final \(\fitinline{d}\)-lift preserves its character but has nontrivial additive phase, hence cancels. Principal local characters retain their nonunit zeros.

For unit parts \(\fitinline{h_p=h/p^{j'}}\), \(\fitinline{a_p^\circ=a/p^k}\), \(\fitinline{A_p^\circ=A/p^t}\), CRT and \(\fitinline{d=h_p(b_*A_p^\circ)^{-1}d'}\) give

\[\fitmath{
\chi_p(h_p)^{k-t}\chi_p(a_p^\circ)^t
\chi_p(A_p^\circ)^{t-k}\chi_p(b_*)^{t-k}.
}\]

Put \(\fitinline{\rho=\chi_p(u/p^j)}\). The \(\fitinline{b_*}\) Gauss factor is \(\fitinline{\sqrt{q_{b_*}}\tau\xi(A)\overline{\xi(h/a)}}\), with \(\fitinline{\overline{\xi(h/a)}=\overline{\xi(u)}\xi(a)}\) since \(\fitinline{\xi^6=1}\). Thus \protect\hyperlink{cf:eq:A.5}{(108)}'s fixed \(\fitinline{\xi}\) factors cancel; the preceding \(\fitinline{b_*}\) factors cancel \(\fitinline{\chi_a(b_*)\chi_A(b_*)^{-1}}\).

For distinct good primes, reciprocity gives

\[\fitmath{
R(A,a)\prod_p\chi_p(a_p^\circ)^{t_p}\chi_p(A_p^\circ)^{-k_p}
=\prod_pR(p,p)^{t_pk_p}.
}\]

Use \(\fitinline{\gamma_2(c)=\mu(c)\alpha(c)G(c)/\gamma_1(c)}\) and CRT for \(\fitinline{\gamma_1(c)}\). At \(\fitinline{p<q}\), the remaining exponent is \(\fitinline{t_pt_q-e_pe_q=3(e_pl_q+e_ql_p+3l_pl_q)}\); its sign product is

\[\fitmath{
\prod_{p<q}R(p,q)^{e_pl_q+e_ql_p+l_pl_q}
=R(c,n)G(n^3)\prod_p\overline{G(p^3)}^{\,l_p}
R(p,p)^{-e_pl_p-l_p(l_p-1)/2}.
}\]

Expand \protect\hyperlink{cf:eq:A.2}{(7)} to verify this. Then \(\fitinline{G(c)R(c,n)G(n^3)=G(cn^3)}\) cancels the conjugate phase; \(\fitinline{G(p^3)=\gamma_3(p)}\) and \(\fitinline{R(p,p)=\epsilon_p}\). The coefficient therefore factors at every prime without ray restriction.

\hypertarget{cf:sec:A.3}{}

\hypertarget{a.6.2-local-euler-factors}{%
\paragraph{A.6.2 Local Euler factors}\label{a.6.2-local-euler-factors}}

At \(\fitinline{\mathfrak p\notin S}\), put

\[\fitmath{
Q=q_{\mathfrak p},\quad j=v_{\mathfrak p}(u)\in\{0,\ldots,5\},\quad
\rho=\chi_{\mathfrak p}(u/\mathfrak p^j),
}\]

\[\fitmath{
a_{\mathfrak p}=\overline{\alpha(\mathfrak p)}^3\eta(\mathfrak p)^3,
\quad b_{\mathfrak p}=\overline{\alpha(\mathfrak p)}^2\eta(\mathfrak p)^2
\overline{\chi_{\mathfrak p}(4)},\quad b_{\mathfrak p}^3=a_{\mathfrak p}^2,
}\]

\hypertarget{cf:eq:A.8}{}

\[\fitmath{
V=Q^{-6z},\quad R=a_{\mathfrak p}^2Q^{4-6s-6z},\quad
W^\circ=\rho Q^{-w},\quad
W=\chi_{\mathfrak p}(u)Q^{-w},\quad
D=\eta(\mathfrak p)\overline{\chi_{\mathfrak p}(u)}Q^{-s}.
\qquad\text{(111)}
}\]

For \(\fitinline{j>0}\), \(\fitinline{W=D=0}\) while \(\fitinline{W^\circ}\) retains its unit phase. The marked and full series are

\[\fitmath{
P_{\mathfrak p}^*=\frac1{1-R}
\left\{\frac{R(1-Q^{-1})-\eta(\mathfrak p)(Q-1)Q^{-s-w}V^{\mathbf1_{j\le1}}}{1-V}+J_j\right\},
}\]

\hypertarget{cf:eq:A.9}{}

\[\fitmath{
P_{\mathfrak p}=\frac1{1-V}+\mathbf1_{j=0}\frac{W^\circ}{1-W^\circ}+P_{\mathfrak p}^*,
\qquad
H_{\mathfrak p}=P_{\mathfrak p}\frac{(1-V)(1-W)}{1-D},
\quad H_{\eta,u}=\prod_{\mathfrak p\notin S}H_{\mathfrak p},
\qquad\text{(112)}
}\]

with

\[\fitmath{
\begin{array}{c|l}
j&J_j\\\hline
0&-\eta(\mathfrak p)\rho^{-1}Q^{-s}+W^\circ R\\
1&\eta(\mathfrak p)Q^{-s-w}\\
2&a_{\mathfrak p}\rho^{-3}Q^{3/2-3s}\\
3&-\eta(\mathfrak p)b_{\mathfrak p}\rho^{-2}Q^{2-3s-w}+b_{\mathfrak p}^2\rho^{-4}Q^{2-4s}\\
4&-\eta(\mathfrak p)a_{\mathfrak p}\rho^{-3}Q^{5/2-4s-w}\\
5&-a_{\mathfrak p}^2Q^{3-6s}
\end{array}
}\]

The indicator exponent gives \(\fitinline{V}\) for \(\fitinline{j=0,1}\), otherwise one.

Equivalently, using \protect\hyperlink{cf:eq:A.8}{(111)},

\hypertarget{cf:eq:A.10}{}

\[\fitmath{
\begin{aligned}
P_p=\sum_{\substack{e_0=0,1\\l,k,m\ge0}}
&(-1)^{e_0}\gamma_1(p)^{-e_0}\eta(p)^{e_0}
\bigl(a_p\overline{\gamma_3(p)}\bigr)^l
\rho^{k-t}\epsilon_p^{tk-e_0l-l(l-1)/2}\\
&\times C_p(t,k,j+6m)
Q^{-(s+1/2)e_0-(1+3s)l-wk-6zm},\qquad t=e_0+3l.
\end{aligned}
\qquad\text{(113)}
}\]

Restricting to \(\fitinline{t>0}\) gives the marked series, with families

\hypertarget{cf:eq:A.11}{}

\[\fitmath{
\begin{array}{c|cc}
 t&k=0&k=1\\\hline
 6r+6&j'\ge t-1&j'=t\\
 6r+1&j'=t-1&j'\ge t\\
 6r+3&j'=t-1&j'=t\\
 6r+4&j'=t-1&j'=t
\end{array}\qquad r\ge0.
\qquad\text{(114)}
}\]

Residue lifting supports a nonprincipal Gauss sum exactly at valuation \(\fitinline{t-1}\); a principal one has its negative boundary there and value \(\fitinline{Q^t-Q^{t-1}}\) above. In \protect\hyperlink{cf:eq:A.7}{(110)}, use exponents \(\fitinline{t}\) at \(\fitinline{k=0}\) and \(\fitinline{t-1}\) at \(\fitinline{k=1}\).

Insert \(\fitinline{j'=j+6m}\). Increasing \(\fitinline{r}\) and first allowed \(\fitinline{m}\) by one multiplies by \(\fitinline{R}\); further \(\fitinline{m}\) increments multiply by \(\fitinline{V}\). The strict tails are

\[\fitmath{
\frac{R(1-Q^{-1})}{(1-R)(1-V)},\qquad
 -\frac{\eta(p)(Q-1)Q^{-s-w}V^{\mathbf1_{j\le1}}}{(1-R)(1-V)}.
}\]

The second starts at \(\fitinline{m=r+1}\) for \(\fitinline{j\le1}\), otherwise \(\fitinline{m=r}\). Equality cases give \(\fitinline{J_j}\) by \protect\hyperlink{cf:eq:A.1}{(6)}: e.g.~\(\fitinline{j=2}\) from \(\fitinline{(t,k,m)=(3,0,0)}\) and \(\fitinline{j=3}\) from \(\fitinline{(3,1,0),(4,0,0)}\). The restored cube phase is \(\fitinline{\overline{\gamma_3(p)}}\), retaining parity through \(\fitinline{\gamma_3(p)^2=\chi_p(-1)}\).

For \(\fitinline{t=0}\), \(\fitinline{k=0}\) gives \(\fitinline{(1-V)^{-1}}\); \(\fitinline{k>0}\) requires \(\fitinline{j=m=0}\) and gives \(\fitinline{W^\circ/(1-W^\circ)}\). Thus \protect\hyperlink{cf:eq:A.9}{(112)} holds in every valuation, including collisions. Extracting \(\fitinline{(1-V)^{-1},(1-W)^{-1},1-D}\) gives

\[\fitmath{
\mathcal F_{\eta,u}(s,w,z)=
 \frac{\zeta_K^S(6z)L_K^S(w,\psi_u)}{L_K^S(s,\eta\overline{\psi_u})}
 H_{\eta,u}(s,w,z).
}\]

Together with Section~\protect\hyperlink{cf:sec:A.2}{A.6.1} this proves \protect\hyperlink{cf:eq:2.8}{(105)}.~\ensuremath{\square}

\hypertarget{cf:interfaces}{}

\hypertarget{a.6.3-interfaces-retained-for-later-chapters}{%
\paragraph{A.6.3 Interfaces retained for later chapters}\label{a.6.3-interfaces-retained-for-later-chapters}}

The retained contracts are: \protect\hyperlink{cf:eq:B.13}{(29)}, arbitrary squarefree frozen \(\fitinline{f}\) at cost \(\fitinline{q_f}\); \protect\hyperlink{cf:eq:C.1}{(45)}, exact M\"{o}bius coefficients, common profile and power gap; \protect\hyperlink{cf:eq:D.3}{(67)}, two factors with common character, natural zeros, common puncture and row-fixed moving radical paid in \(\fitinline{m+q}\). Identities \protect\hyperlink{cf:eq:A.12}{(8)}--\protect\hyperlink{cf:eq:A.19}{(15)}, \protect\hyperlink{cf:eq:C.3}{(47)} preserve collision masks; \protect\hyperlink{cf:eq:2.8}{(105)} retains its quotient and all valuation factors.

Selected-prime, merged inverse--plain, relative-character and incumbent-dependent coefficients require their own contracts. In particular, \protect\hyperlink{cf:eq:4.6}{(40)} is only the stated-scale theta energy, not an inverse-times-prime or plain-fourth-times-prime moment.

\hypertarget{mf:chapter}{}

\hypertarget{a.7-marked-completed-energy-and-the-incumbent-moments}{%
\subsubsection{A.7 Marked completed energy and the incumbent moments}\label{a.7-marked-completed-energy-and-the-incumbent-moments}}

Using \S{}\protect\hyperlink{cf:reflection}{A.3.1}'s reflection and the finite identities \protect\hyperlink{cf:eq:A.12}{(8)}--\protect\hyperlink{cf:eq:A.19}{(15)}, \protect\hyperlink{cf:eq:C.3}{(47)}, we prove multi-slot completed energy, inverse-times-prime and natural marked plain moments. Original MF equation labels are retained; identical foundational formulas are aliases. Here \(\fitinline{\beta}\) is the actual Hecke zero supremum, and every prime estimate assumes \(\fitinline{2\beta-1\le\kappa}\).

The original and refined bootstrap ranges remain distinct. Every \(\fitinline{Z^\epsilon}\) means a loss chosen before the scale. Rows are Eisenstein elements with all six units and arbitrary valuations.

\hypertarget{the-completed-reflection-with-product-marks}{}

\hypertarget{a.7.1-the-completed-reflection-with-product-marks}{%
\paragraph{A.7.1 The completed reflection with product marks}\label{a.7.1-the-completed-reflection-with-product-marks}}

Whole-index marks change conductors at vanishing local frequency. Retain compatible branches until row and mark phases separate, then apply the quadratic--cubic sieve.

\hypertarget{the-completed-family-and-fixed-conventions}{}

\hypertarget{a.7.1.1-the-completed-family-and-fixed-conventions}{%
\subparagraph{A.7.1.1 The completed family and fixed conventions}\label{a.7.1.1-the-completed-family-and-fixed-conventions}}

Use \S{}\protect\hyperlink{cf:sec:1.1}{A.1.2}'s arithmetic conventions, \(\fitinline{q_a=|a|^2}\), \(\fitinline{Z\ge2}\), and fixed \(\fitinline{S}\) containing primes above \(\fitinline{6}\). Primary generators represent ideals prime to \(\fitinline{\lambda}\). In particular,

\hypertarget{mf:eq:B2.1}{}

\[\fitmath{
\chi_p(A)^0=1_{p\nmid A}.
\qquad\text{(MF2.1)}
}\]

Exponent reduction never deletes this mask.

Fix finitely many finite-ray \(\fitinline{\nu}\), zero outside primary \(\fitinline{S}\)-units. A squarefree good label \(\fitinline{f}\) and squarefree puncture \(\fitinline{r_\rho}\) may vary between invocations but are fixed within a row norm:

\hypertarget{mf:eq:B2.2}{}

\[\fitmath{
\rho(A)=1_{(A,r_\rho)=1},\qquad
\Psi_k(A)=\nu(A)\rho(A)\chi_A(k)\chi_A(f)^4.
\qquad\text{(MF2.2)}
}\]

The puncture may depend on frozen outer labels, but neither the current row nor a subsequently averaged \(\fitinline{f}\).

For a fixed number of pairwise disjoint prime lists \(\fitinline{\mathcal P_i}\), with \(\fitinline{q_p\asymp Z^{z_i}}\) on the \(\fitinline{i}\)-th list and \(\fitinline{|a_i(p)|\le1}\), put

\hypertarget{mf:eq:B2.3}{}

\[\fitmath{
d(A)=\sum_{(p_i)\in\prod_i\mathcal P_i}
       \prod_i a_i(p_i)\,1_{\prod_i p_i\mid A},
\qquad z=\sum_i z_i.
\qquad\text{(MF2.3)}
}\]

The completed sum at length \(\fitinline{X}\) is

\hypertarget{mf:eq:B2.4}{}

\[\fitmath{
\mathcal T_d(X;k;V)=
\sum_{\substack{n\ {\rm sf},\,b\\ (nb,S)=1}}
\frac{\overline{\alpha(n)}\gamma_2(n)\Psi_k(n)\,
      \overline{\alpha(b)}^{\,3}\Psi_k(b)^3}
     {\sqrt{q_n}\,q_b}\,
d(nb^3)V(q_nq_b^3/X).
\qquad\text{(MF2.4)}
}\]

Primal \(\fitinline{n,b}\) may share primes. Tests are annular smooth or rapidly decreasing with all Euler derivatives at both endpoints; include the central conversion factor \(\fitinline{y^{1/2}}\) in \(\fitinline{V}\).

A common smooth profile in completed norm and individual prime norms, on fixed logarithmic prime boxes, is allowed inside the tuple sum. Logarithmic Fourier inversion reduces it to an integrable common family of \(\fitinline{V}\); bounds use finitely many weighted Euler seminorms, uniformly in lists, bounded coefficients, \(\fitinline{k,f,r_\rho}\) and bounded logarithmic lengths. Heights cost fixed polynomials. Product coefficients are required during assignment to extracted factors; afterwards surviving tuple coefficients may be arbitrary bounded functions independent of row and dual indices.

Classical theta and sieve inputs are exactly \S{}\protect\hyperlink{cf:classical-inputs}{A.1.1}'s unconditional results; moving-conductor and marked reductions follow below.

\hypertarget{exact-reflection-with-the-local-zeros-retained}{}

\hypertarget{a.7.1.2-exact-reflection-with-the-local-zeros-retained}{%
\subparagraph{A.7.1.2 Exact reflection, with the local zeros retained}\label{a.7.1.2-exact-reflection-with-the-local-zeros-retained}}

\hypertarget{the-finite-theta-source}{}

The finite theta source

With \(\fitinline{\check e(z)=\exp(2\pi i(z+\bar z))}\) and \(\fitinline{e(z)=\check e(z/\lambda)}\), the theta input {[}DR, (5.4), (5.7), (5.9), (5.12)--(5.17), Appendix A{]} gives

\hypertarget{mf:eq:B2.5}{}

\[\fitmath{
\Gamma_2=\langle {\rm SL}_2(\mathbb Z),\Gamma(3)\rangle,\qquad
\theta(gw)=\kappa(g)\theta(w),\qquad
\kappa\begin{pmatrix}a&b\\c&d\end{pmatrix}
=\left(\frac ca\right)_3\quad(g\in\Gamma(3),\ c\ne0).
\qquad\text{(MF2.5)}
}\]

The multiplier is one on \(\fitinline{{\rm SL}_2(\mathbb Z)}\) and on the \(\fitinline{c=0}\) elements of \(\fitinline{\Gamma(3)}\). Let

\[\fitmath{
L(t)=\begin{pmatrix}1&0\\t&1\end{pmatrix},\qquad
\Theta_j(w)=\theta(L(t_j)w),\qquad
(t_0,t_1,t_2)=(0,\omega^2,\omega).
}\]

Translations by \(\fitinline{\mathbb Z+3\mathcal O}\) have multiplier one, so \(\fitinline{L(\pm\lambda)}\) give \(\fitinline{\Theta_1,\Theta_2}\). For \(\fitinline{u_0=s+t\omega}\), \(\fitinline{0\le s,t<3}\), \(\fitinline{\theta(\left(\begin{smallmatrix}u_0&-1\\1&0\end{smallmatrix}\right)w)=\Theta_t(w)}\).

For each of these three sources define \(\fitinline{d_H}\) by

\hypertarget{mf:eq:B2.6}{}

\[\fitmath{
\overline{\theta(H(z,v))}
=C_H(v)+\sum_{\mu\ne0}d_H(\mu)\,vK_{1/3}(4\pi|\mu|v)\check e(\mu z).
\qquad\text{(MF2.6)}
}\]

Their support and coefficient bound are

\hypertarget{mf:eq:B2.7}{}

\textbf{(MF2.7)} is the cusp support and bound \protect\hyperlink{cf:eq:B.5}{(21)}, with \(\fitinline{m=h}\), \(\fitinline{d=d_H}\), and \(\fitinline{\sqrt{q_b}=|b|}\).

The dual unit \(\fitinline{u}\) and primary ideals \(\fitinline{n,b}\) retain the full cusp support: \(\fitinline{n,b}\) may share primes and meet \(\fitinline{S\setminus\{\lambda\}}\), without primal \(\fitinline{S}\)-exclusion.

For the precise source check, the ordinary cubic sums

\[\fitmath{
\widetilde g_3(n)=q_n^{-1/2}\sum_{v\bmod n}\left(\frac vn\right)_3\check e(v/n)
}\]

have modulus one for primary squarefree \(\fitinline{n}\) by finite Gauss and CRT, including the norm-four prime \(\fitinline{-2}\) using only its cubic symbol. The three source magnitudes are \(\fitinline{3^{h/6+8/3}|b|}\), \(\fitinline{3^{h/6+3}|b|}\), and \(\fitinline{9|b|}\) at \(\fitinline{h=-4}\), each within \protect\hyperlink{mf:eq:B2.7}{(MF2.7)}. Valuations uniquely determine \(\fitinline{n,b}\). At infinity, \(\fitinline{\lambda^3\mu\equiv1\pmod3}\) selects

\hypertarget{mf:eq:B2.8}{}

\[\fitmath{
d_I(\lambda^{-3}nb^3)=3^{5/2}|b|\widetilde g_3(n),\qquad
\widetilde g_3(n)=\chi_n(\lambda)^{-2}\gamma_2(n)\quad((n,S)=1).
\qquad\text{(MF2.8)}
}\]

Shared \(\fitinline{n,b}\) primes remain allowed.

\hypertarget{all-nonzero-element-rows-use-one-fixed-modulus}{}

All nonzero element rows use one fixed modulus

Write uniquely \(\fitinline{k=uk_Sk_{\rm good}}\) with fixed generators at \(\fitinline{S}\) and primary generators elsewhere. On primary \(\fitinline{S}\)-units, reciprocity gives

\hypertarget{mf:eq:B2.9}{}

\[\fitmath{
\chi_A(k)=\chi_A(uk_S)R(A,k_{\rm good})
          \prod_{p\mid k_{\rm good}}\chi_p(A)^{v_p(k)}.
\qquad\text{(MF2.9)}
}\]

At collisions both sides vanish, including \(\fitinline{6\mid v_p(k)}\). Fixing \(\fitinline{u}\), \(\fitinline{S}\)-valuations modulo six and the reciprocity ray class of \(\fitinline{k_{\rm good}}\) puts the first two factors times \(\fitinline{\nu}\) in a fixed finite \(\fitinline{S}\)-supported family \(\fitinline{\Psi_0}\). Supplementary characters have only these finitely many exponents. The fourth-power label adds no sign because \(\fitinline{R(A,f)^4=1}\).

Fix \(\fitinline{E}\) with support exactly \(\fitinline{S}\), divisible by these conductors. Choose \(\fitinline{L}\) on the same support, \(\fitinline{18\mid L}\), \(\fitinline{E\mid L}\), so

\[\fitmath{
\varphi(x)=
\begin{cases}
\chi_x(\lambda)^2\Psi_0(x),&x\equiv1\pmod3,\ (x,E)=1,\\
0,&\text{otherwise}
\end{cases}
}\]

is periodic: \(\fitinline{(\lambda/x)_3}\) has modulus \(\fitinline{9}\) by the supplementary law; all other characters and zeros have fixed moduli. Evaluate the zero branch first. Set

\hypertarget{mf:eq:B2.10}{}

\[\fitmath{
\widehat\varphi(h_0)=q_L^{-1}\sum_{x\bmod L}\varphi(x)e(-h_0x/L),
\qquad \mathfrak M=\lambda^{12}L^4.
\qquad\text{(MF2.10)}
}\]

Keep every good prime of \(\fitinline{k,f,r_\rho}\) in a moving local set \(\fitinline{\mathcal P}\), combining exponents modulo six; none enters \(\fitinline{L}\). Reflect \(\fitinline{\Psi(A)=\Psi_0(A)\prod_{p\in\mathcal P}\chi_p(A)^{j_p}}\), \(\fitinline{0\le j_p<6}\).

\hypertarget{the-reflection-identity}{}

The reflection identity

For \(\fitinline{m\not\equiv0\pmod6}\), let

\[\fitmath{
\tau_{p,m}^{\pm}=q_p^{-1/2}\sum_{y\bmod p}\chi_p(y)^m e(\pm y/p),
\qquad |\tau_{p,m}^{\pm}|=1.
}\]

A branch chooses \(\fitinline{h_0\bmod L}\) and an active subset containing all \(\fitinline{j_p\ne0}\). Let \(\fitinline{r}\) be its product. Reduce \(\fitinline{\lambda^2h_0/L=a_F/c_F}\), making the numerator primary if \(\fitinline{\lambda\mid c_F}\), otherwise the denominator. Put \(\fitinline{c=c_Fr}\) and, modulo active \(\fitinline{p}\), define

\[\fitmath{
\sigma_p=\lambda^2c/p,\qquad
\epsilon_p=-\bigl((\lambda^3c/p)\sigma_p\bigr)^{-1}
          =-\lambda^{-5}(c/p)^{-2},
}\]

\hypertarget{mf:eq:B2.11}{}

\[\fitmath{
B_{p,j}(x)=
\begin{cases}
\chi_p(x)^{-j-2},&j\ne0,4,\\
q_p^{-1/2}(-1+q_p1_{p\mid x}),&j=4,\\
q_p^{-1/2}\chi_p(x)^{-2},&j=0,
\end{cases}
\quad
\omega_{p,j}=
\begin{cases}
\tau_{p,j}^-\tau_{p,j+2}^+\chi_p(\epsilon_p)^{-j-2},&j\ne0,4,\\
\tau_{p,4}^-,&j=4,\\
-\tau_{p,2}^+\chi_p(\epsilon_p)^{-2},&j=0.
\end{cases}
\qquad\text{(MF2.11)}
}\]

Indices on \(\fitinline{\tau}\) are reduced modulo six.

Define

\hypertarget{mf:eq:B2.12}{}

\[\fitmath{
K_*=\frac{(2\pi)^4}{27},\qquad
\mathcal R(t)=\prod_\pm\frac{\Gamma(1+t\pm1/6)}{\Gamma(1-t\pm1/6)},\qquad
V^\sharp(x)=\frac1{2\pi i}\int_{(\sigma)}
\widehat V(-t)\mathcal R(t)(K_*x)^{-t}\,dt,\quad \sigma\ge0.
\qquad\text{(MF2.12)}
}\]

The integral is independent of \(\fitinline{\sigma\ge0}\). Let \(\fitinline{\mathcal T_1(X;\Psi;V)}\) be \protect\hyperlink{mf:eq:B2.4}{(MF2.4)} with no mark and character \(\fitinline{\Psi}\). Each branch's source \(\fitinline{d}\), additive character \(\fitinline{\vartheta}\) and unit \(\fitinline{\kappa_F}\) depend only on \(\fitinline{h_0,r\bmod\mathfrak M^2}\).

\hypertarget{mf:eq:B2.13}{}

\textbf{(MF2.13)} is \protect\hyperlink{cf:eq:B.2}{(18)} with scalar \protect\hyperlink{cf:eq:B.4}{(20)}, under \(\fitinline{\mathcal T_1=\mathcal C_V(X;\Psi)}\), \(\fitinline{\widehat\phi=\widehat\varphi}\), \(\fitinline{c=c_Fr}\), and the local factors \protect\hyperlink{mf:eq:B2.11}{(MF2.11)}. The scalars agree since \(\fitinline{\tau_{p,j}^-=\chi_p(-1)^j\gamma_j(p)}\) and \(\fitinline{\tau_{p,j}^+=\gamma_j(p)}\).

Dual sums converge absolutely, \(\fitinline{|\zeta|\le1/81}\), with \(\fitinline{O_L(2^{|\mathcal P|})}\) branches, or \(\fitinline{O_L(4^{|\mathcal P|})}\) after splitting \(\fitinline{j=4}\). Choices are simultaneous: adding an inactive \(\fitinline{j=0}\) prime multiplies by \(\fitinline{1-q_p^{-1}}\) while preserving conductor, source, additive character, other phases and kernel.

\textbf{Conductor and phase proof.} The fixed-modulus construction of \S{}\protect\hyperlink{cf:reflection}{A.3.1} applies; we retain its extra explicit sector checks here. Its local Fourier identity is

\hypertarget{mf:eq:B2.14}{}

\textbf{(MF2.14)} is \protect\hyperlink{cf:eq:B.9}{(25)}, with \(\fitinline{q=q_p}\), \(\fitinline{j=j_p}\) and \(\fitinline{\tau^-_{\mathfrak p,j}=\tau^-_{p,j}}\).

Hence the masked theta function is

\hypertarget{mf:eq:B2.15}{}

\[\fitmath{
F(z,v)=\sum_{h_0,(h_p)}
\widehat\varphi(h_0)\prod_pC_{p,j_p}(h_p)\,
\overline{\theta(z+z_h,v)},\qquad
z_h=\lambda^2\left(h_0/L+\sum_ph_p/p\right).
\qquad\text{(MF2.15)}
}\]

At infinity \(\fitinline{\lambda^3\mu=nb^3}\), the mask first deletes \(\fitinline{(nb,E)>1}\), leaving multiplier \(\fitinline{\chi_n(\lambda)^2\Psi(n)\Psi(b)^3}\), also at shared moving primes. Thus \protect\hyperlink{mf:eq:B2.8}{(MF2.8)} identifies the completed direct series.

Fix \(\fitinline{h_0}\) and the active frequencies. Then

\hypertarget{mf:eq:B2.16}{}

\[\fitmath{
z_h=a/c,\qquad c=c_Fr,\qquad
a=a_Fr+\lambda^2c_F\sum_{p\mid r}(r/p)h_p.
\qquad\text{(MF2.16)}
}\]

The numerator is a unit at \(\fitinline{c_F}\) and at \(\fitinline{p\mid r}\) (there \(\fitinline{a\equiv\sigma_ph_p}\)), so the fraction is reduced even for nonunit \(\fitinline{h_0}\). Lift \(\fitinline{h_p\equiv0\pmod{\mathfrak M^2}}\) while retaining its \(\fitinline{p}\)-class and fix \(\fitinline{r\bmod\mathfrak M^2}\); then \(\fitinline{a\bmod\mathfrak Mc_F}\) is fixed. Choose by CRT

\[\fitmath{
\begin{aligned}
\delta'&\equiv a^{-1}\pmod{\mathfrak Mc_F}\text{ at primes of }c_F,\\
\delta'&\equiv0\pmod{\mathfrak M}\text{ at other primes of }\mathfrak M,\\
\delta'&\equiv a^{-1}\pmod p\ (p\mid r).
\end{aligned}
}\]

Put \(\fitinline{b=(a\delta'-1)/c}\), \(\fitinline{g=\left(\begin{smallmatrix}a&b\\c&\delta'\end{smallmatrix}\right)}\). At \(\fitinline{c_F}\), \(\fitinline{b\equiv0\pmod{\mathfrak M}}\); elsewhere in \(\fitinline{\mathfrak M}\), \(\fitinline{b\equiv-c^{-1}}\). Thus fixed-prime matrix residues are frequency-independent. Periodic lifts and CRT-modulus shifts of \(\fitinline{\delta'}\) avoid zero entries if needed for symbols.

If \(\fitinline{v_\lambda(c)\ge2}\), take \(\fitinline{H=I}\): \(\fitinline{g\in\Gamma(3)}\) and \(\fitinline{\kappa(g)=(c_F/a)_3(a/r)_3}\). If \(\fitinline{v_\lambda(c)=1}\), take \(\fitinline{H=L(u_0)}\), \(\fitinline{u_0\in\{\pm\lambda\}}\), \(\fitinline{u_0\equiv c\pmod3}\), and \(\fitinline{A=a-u_0b}\). Then \(\fitinline{G=gH^{-1}\in\Gamma(3)}\), \(\fitinline{a,A}\) are coprime primary, and

\[\fitmath{
a(c-u_0\delta')=cA-u_0,\qquad bc\equiv-1\pmod a.
}\]

Cubic reciprocity gives

\[\fitmath{
\kappa(G)=
\left(\frac{-u_0}{A}\right)_3\left(\frac aA\right)_3^{-1}
=\left(\frac{-u_0}{A}\right)_3
 \left(\frac{c_F/u_0}{a}\right)_3\left(\frac ar\right)_3.
}\]

If \(\fitinline{\lambda\nmid c}\), choose \(\fitinline{u_0=s+t\omega\equiv a\pmod3}\), \(\fitinline{0\le s,t<3}\), and \(\fitinline{H=\bigl(\begin{smallmatrix}u_0&-1\\1&0\end{smallmatrix}\bigr)}\). Then

\[\fitmath{
G=gH^{-1}=
\begin{pmatrix}-b&a+bu_0\\-\delta'&c+\delta'u_0\end{pmatrix}\in\Gamma(3).
}\]

Both \(\fitinline{-b,c}\) are primary and \(\fitinline{-bc=1-a\delta'\equiv1\pmod9}\). For \(\fitinline{\delta'=u\lambda^h\delta_0}\), supplementary laws handle unit/\(\fitinline{\lambda}\) factors modulo \(\fitinline{9}\); reciprocity gives numerator \(\fitinline{1\bmod\delta_0}\), so \(\fitinline{(\delta'/(-bc))_3=1}\). Hence \(\fitinline{\kappa(G)=(a/c_F)_3(a/r)_3}\). In all cases,

\hypertarget{mf:eq:B2.17}{}

\[\fitmath{
\kappa(G)=\kappa_F\prod_{p\mid r}\chi_p(a)^2,
\qquad\text{(MF2.17)}
}\]

The sector fixes \(\fitinline{\kappa_F,d_H}\), without requiring unit-frequency support for \(\fitinline{\widehat\varphi}\).

For \(\fitinline{x=\lambda^4\mu}\) and \(\fitinline{D_F=\lambda^3c_F}\), additive CRT gives

\hypertarget{mf:eq:B2.18}{}

\[\fitmath{
\check e(-\delta'\mu/c)
=\vartheta(x)\prod_{p\mid r}e(\epsilon_ph_p^{-1}x/p),\qquad
\vartheta(x)=e(-\delta_F'r^{-1}x/D_F).
\qquad\text{(MF2.18)}
}\]

Fixed congruences determine \(\fitinline{\delta_F'\bmod D_F}\), and the sector fixes \(\fitinline{r^{-1}}\). Conjugating \protect\hyperlink{mf:eq:B2.17}{(MF2.17)} contributes \(\fitinline{\overline{\kappa_F}\prod_p\chi_p(\sigma_p)^{-2}\chi_p(h_p)^{-2}}\). Multiply \protect\hyperlink{mf:eq:B2.14}{(MF2.14)} and set \(\fitinline{y=h_p^{-1}}\): Gauss exponent \(\fitinline{j+2}\), or \(\fitinline{-1+q_p1_{p\mid x}}\) at \(\fitinline{j=4}\), gives exactly \(\fitinline{\omega_{p,j}B_{p,j}(x)}\) with zeros. Inactive and absent zero-exponent primes have identical active radical, proving simultaneous compatibility.

\textbf{Proof of the Mellin normalization and continuation.} Put

\[\fitmath{
D(s,\Psi)=\sum_{n\ {\rm sf}}\overline{\alpha(n)}\gamma_2(n)\Psi(n)q_n^{-s},
\quad L_0(s,\Psi)=\sum_b\overline{\alpha(b)}^{\,3}\Psi(b)^3q_b^{-3s+1/2},
\quad T=L_0D.
}\]

These primary \(\fitinline{S}\)-unit sums converge for \(\fitinline{\Re s>1}\). Put \(\fitinline{J(s)=\int_0^\infty\partial_{\bar z}F(0,v)v^{2s-1}\,dv}\); the Bessel integral is

\[\fitmath{
\int_0^\infty K_{1/3}(y)y^{w-1}\,dy
=2^{w-2}\Gamma((w-1/3)/2)\Gamma((w+1/3)/2).
}\]

Using \(\fitinline{\partial_{\bar z}\check e(\mu z)=2\pi i\bar\mu\check e(\mu z)}\) and \(\fitinline{\overline{\alpha(\lambda^{-3})}=-i}\),

\hypertarget{mf:eq:B2.19}{}

\textbf{(MF2.19)} is the \protect\hyperlink{cf:direct-mellin-identity}{direct Mellin identity}, with \(\fitinline{\Theta_\Psi=F}\), \(\fitinline{\mathcal D_\Psi(s)=T(s,\Psi)}\) and left side \(\fitinline{J(s)}\).

The inverse cusp action is \protect\hyperlink{cf:eq:B.10}{(26)}, with \(\fitinline{\delta=\delta'}\).

At \(\fitinline{z=0}\) the chain rule gives \(\fitinline{-\overline{\kappa(G)}(cv)^{-2}\partial_z}\) at \(\fitinline{(-\delta'/c,1/(q_cv))}\), deleting every cusp constant. Substitute \(\fitinline{v\mapsto1/(q_cv)}\) and the Bessel integral: for \(\fitinline{\Re s<0}\), \(\fitinline{J(s)}\) equals

\hypertarget{mf:eq:B2.20}{}

\[\fitmath{
-\frac i4(2\pi)^{2s-2}\overline{\alpha(c)}^{\,2}q_c^{1-2s}
\overline{\kappa(G)}
\prod_\pm\Gamma(3/2-s\pm1/6)
\sum_{\mu\ne0}d_H(\mu)\alpha(\mu)\check e(-\delta'\mu/c)q_\mu^{-(1-s)}.
\qquad\text{(MF2.20)}
}\]

By \protect\hyperlink{mf:eq:B2.7}{(MF2.7)}, this converges absolutely. Direct and inverse-cusp expansions give exponential endpoint decay; log-height integration by parts makes \(\fitinline{J}\) entire and rapidly decreasing on fixed strips. Equation \protect\hyperlink{mf:eq:B2.19}{(MF2.19)} makes \(\fitinline{T}\) entire, bounded right of \(\fitinline{1}\) and polynomial left of \(\fitinline{0}\) via \protect\hyperlink{mf:eq:B2.20}{(MF2.20)}. Its preliminary intervening bound is \(\fitinline{(1+|s|)^C e^{\pi|\Im s|}}\). Apply the maximum principle to \(\fitinline{T(s)(B+s)^{-C'}e^{\epsilon s^2}}\), \(\fitinline{-B}\) outside the strip, then let \(\fitinline{\epsilon\downarrow0}\), obtaining polynomial strip growth.

Mellin inversion gives \(\fitinline{(2\pi i)^{-1}\int_{(a)}\widehat V(s-1/2)X^{s-1/2}T(s,\Psi)\,ds}\), \(\fitinline{a>1}\). Shift left of zero, use \protect\hyperlink{mf:eq:B2.20}{(MF2.20)}/\protect\hyperlink{mf:eq:B2.19}{(MF2.19)} and set \(\fitinline{t=1/2-s}\). The gamma quotient is \(\fitinline{\mathcal R(t)}\); the remaining factor is

\[\fitmath{
-\frac i{81}\overline{\alpha(c)}^{\,2}q_\mu^{-1/2}
(K_*q_\mu X/q_c^2)^{-t}.
}\]

Absolute convergence and Mellin decay justify interchanges and shifts. The finite factors give \protect\hyperlink{mf:eq:B2.13}{(MF2.13)}, also for the rapidly decreasing tests admitted here.

\hypertarget{common-profiles-and-whole-series-tails}{}

\hypertarget{a.7.1.3-common-profiles-and-whole-series-tails}{%
\subparagraph{A.7.1.3 Common profiles and whole-series tails}\label{a.7.1.3-common-profiles-and-whole-series-tails}}

For \(\fitinline{A,B\ge0}\), with \(\fitinline{B}\) an integer, set

\[\fitmath{
\mathcal M_{A,B}(V)=
\sup_{-A\le\eta\le1/4}\int_{\mathbb R}
(1+|t|)^B|\widehat V(\eta+it)|\,dt.
}\]

The first negative numerator pole is \(\fitinline{-5/6}\) and \(\fitinline{|\mathcal R(\sigma+it)|\ll_A(1+|t|)^{4\sigma}}\) for \(\fitinline{-1/4\le\sigma\le A}\). Shifts to \(\fitinline{-1/4,0,A}\) give

\hypertarget{mf:eq:B2.21}{}

\[\fitmath{
|(x\partial_x)^jV^\sharp(x)|
\ll_{A,j}\min\{x^{1/4},(1+x)^{-A}\}
\mathcal M_{A,\lceil4A\rceil+j+2}(V).
\qquad\text{(MF2.21)}
}\]

These are finite seminorms: \(\fitinline{B+2}\) integrations by parts bound \(\fitinline{\mathcal M_{A,B}}\) by

\[\fitmath{
\ll_{A,B}\int_0^\infty
(1+x^{-A}+x^{1/4})\bigl(|V(x)|+|(x\partial_x)^{B+2}V(x)|\bigr)\frac{dx}{x}.
}\]

Replacing \(\fitinline{V}\) by \(\fitinline{x^{it_0}V}\) costs at most \(\fitinline{(1+|t_0|)^B}\).

On a fixed compact log box let \(\fitinline{F_Y(\mathbf y)=W(\mathbf y)V^\sharp(Y\prod_i y_i^{a_i})}\), \(\fitinline{Y}\) fixed per block. Equation \protect\hyperlink{mf:eq:B2.21}{(MF2.21)} bounds each fixed derivative by

\[\fitmath{
m_A(Y)=\min\{Y^{1/4},(1+Y)^{-A}\}
}\]

times fixed seminorms. Logarithmic Fourier integration by parts gives the same weighted \(\fitinline{L^1}\) scalar. For a linear row vector, set \(\fitinline{F_Y=m_A(Y)\widetilde F_Y}\) and apply one common Fourier inversion and Minkowski:

\hypertarget{mf:eq:B2.22}{}

\[\fitmath{
\|U(F_Y)\|_{\ell^2(I)}
\le m_A(Y)\int |\widehat{\widetilde F_Y}(\mathbf t)|
\|B(\mathbf t)\|_{\ell^2(I)}\,\frac{d\mathbf t}{(2\pi)^d}.
\qquad\text{(MF2.22)}
}\]

Then enlarge the positive row set using the separated definition of \(\fitinline{B}\) on added rows, without asserting kernel bounds there. Squaring gives \(\fitinline{m_A(Y)^2}\), in particular small-argument energy gain \(\fitinline{Y^{1/2}}\). Row-dependent Fourier measures would invalidate this step.

For \(\fitinline{\Lambda=\mathcal O}\) or \(\fitinline{\lambda^{-4}\mathcal O}\), \(\fitinline{|K(x)|\le C_Ax^{-A}}\) (\(\fitinline{x\ge1}\), \(\fitinline{A>1}\)) gives

\hypertarget{mf:eq:B2.23}{}

\textbf{(MF2.23)} is the shell bound \protect\hyperlink{cf:eq:B.25}{(44)}, valid for both lattices with \(\fitinline{a>0}\), \(\fitinline{U\ge1}\).

Its \(\fitinline{O(1+2^jU/a)}\) shell count controls the whole discarded series independently of retained-index bounds.

\hypertarget{the-quadraticcubic-norm-including-both-collision-masks}{}

\hypertarget{a.7.1.4-the-quadraticcubic-norm-including-both-collision-masks}{%
\subparagraph{A.7.1.4 The quadratic--cubic norm, including both collision masks}\label{a.7.1.4-the-quadraticcubic-norm-including-both-collision-masks}}

\hypertarget{imported-sieves-and-their-domains}{}

Imported sieves and their domains

For squarefree primary good ideals, use the quadratic sieve

\hypertarget{mf:eq:B2.24}{}

\textbf{(MF2.24)} is \protect\hyperlink{cf:eq:B.11}{(27)} under \(\fitinline{(R,n,K,N)=(a,b,U,V)}\), with the same row-independent coefficients.

For clarity, the primitive ideal character in \S{}\protect\hyperlink{cf:sec:B.2}{A.3.2}'s family check is

\[\fitmath{
\vartheta_a((x))=\epsilon_\lambda(x)^{e(a)}(x/a)_2,\qquad
e(a)=(q_a-1)/2\bmod2,
}\]

Here \(\fitinline{\epsilon_\lambda}\) is the nontrivial character modulo \(\fitinline{\lambda}\). Primary adjustment uses square powers of \(\fitinline{\omega}\); its sign contributes \(\fitinline{\epsilon_\lambda^{e(a)}}\). CRT gives conductor \(\fitinline{a\lambda^{e(a)}}\), trivial infinite type, omitting \(\fitinline{\lambda}\) when \(\fitinline{e(a)=0}\). Fixed primary square classes modulo four make \(\fitinline{e(a)}\) equal; coprime product conductors are exactly products, and reciprocity is a fixed-ray bicharacter. These are {[}GL, Definition 1{]}'s three conditions; apply its theorem and transpose duality.

Heath-Brown's cubic sieve gives

\hypertarget{mf:eq:B2.25}{}

\textbf{(MF2.25)} is \protect\hyperlink{cf:eq:B.19}{(36)}, since \(\fitinline{(b/a)_3=\chi_a(b)^2}\), with squarefree primary ideals, lengths \(\fitinline{U,V}\) and row-independent \(\fitinline{c_b}\).

Neither sieve requires squarefree rational norms or exclusion of rational prime factors. Conjugation/reciprocity changes cubic orientation. Fixed row restrictions decrease the positive norm; fixed column restrictions are zero coefficients. Resolve mixed collision masks first.

\hypertarget{quadratic-reduction-for-a-squarefree-index-times-a-cube}{}

Quadratic reduction for a squarefree index times a cube

Use \S{}\protect\hyperlink{cf:sec:B.3}{A.3.3}'s variables: squarefree good \(\fitinline{q_k\ll K}\), squarefree primary \(\fitinline{q_n\asymp N}\), primary \(\fitinline{q_b\asymp B}\), allowing dual \(\fitinline{S\setminus\{\lambda\}}\) primes and arbitrary \(\fitinline{k}\)-independent \(\fitinline{\beta(n,b)}\) with fixed supports.

Write uniquely

\[\fitmath{
b=gt^2,\qquad c=(n,g),\qquad n=cm,\qquad g=ch,
}\]

Here \(\fitinline{g}\) is squarefree; \(\fitinline{c,m,h}\) are pairwise coprime squarefree, allowing \(\fitinline{(g,t)>1}\). On \(\fitinline{q_c\asymp C,q_g\asymp G,q_t\asymp T}\), \(\fitinline{B\asymp GT^2}\), and

\hypertarget{mf:eq:B2.26}{}

\[\fitmath{
\chi_k(nb)^3=\chi_k(mh)^3\,1_{(k,ct)=1}.
\qquad\text{(MF2.26)}
}\]

Apply the quadratic reduction proved in \S{}\protect\hyperlink{cf:sec:B.3}{A.3.3}: Cauchy in \(\fitinline{t}\), retain \(\fitinline{(k,t)=1}\), expand \(\fitinline{1_{(k,c)=1}}\) completely, and set \(\fitinline{k=rk'}\), \(\fitinline{c=rc'}\). Keep \(\fitinline{(k',r)=(c',r)=1}\); \(\fitinline{(r,t)>1}\) vanishes. No \(\fitinline{(k',c')}\)-mask is restored termwise.

The same fixed-\(\fitinline{S}\) split, squarefree product grouping and Cauchy over \(\fitinline{O(C/q_r)}\) choices of \(\fitinline{c'}\) give, with \(\fitinline{\mathcal D_r(t)}\) retaining the original support,

\hypertarget{mf:eq:B2.27}{}

\textbf{(MF2.27)} is the positive quadratic bound \protect\hyperlink{cf:eq:B.17}{(34)}, with \(\fitinline{Q_{C,G,T}=\mathcal Q_{C,G,T}}\) and final sum restricted to \(\fitinline{\mathcal D_r(t)}\). All its row/column independence and common-prime conventions are unchanged.

This positive form preserves tuple independence. For \(\fitinline{|\beta|\le1}\), the triple count is \(\fitinline{O(NG/(q_rC))}\); on \(\fitinline{q_r\asymp r_0}\), \(\fitinline{O(r_0)}\) choices of \(\fitinline{r}\) and \(\fitinline{O(T)}\) choices of \(\fitinline{t}\) give

\[\fitmath{
\ll (KNB)^\epsilon\frac{NB}{r_0}
\left(\frac K{r_0}+\frac{NG}{C^2}\right)
\ll(KNB)^\epsilon NB(K+NB).
}\]

Triangle over \(\fitinline{C,G,T}\) dyads yields

\hypertarget{mf:eq:B2.28}{}

\textbf{(MF2.28)} is \protect\hyperlink{cf:eq:B.18}{(35)}, with the same \(\fitinline{K,N,B}\), supports and \(\fitinline{|\beta|\le1}\).

Subunit quotients are empty; bounded ranges, including the unit, cost constants.

\hypertarget{hybrid-norm}{}

Hybrid norm

Let \(\fitinline{P}\) be squarefree good with \(\fitinline{q_P\asymp L}\). Suppose \(\fitinline{|a(P)|,|\beta(n,b)|\le1}\), with \(\fitinline{a}\) independent of \(\fitinline{k,n,b}\) and \(\fitinline{\beta}\) independent of \(\fitinline{k,P}\). Then

\hypertarget{mf:eq:B2.29}{}

\[\fitmath{
\begin{aligned}
\mathcal H&=
\sum_k\left|
\sum_P\frac{a(P)}{\sqrt{q_P}}1_{(P,k)=1}
\sum_{n,b}\beta(n,b)\chi_k(nb)^3\chi_P(n)^{-2}1_{(P,b)=1}
\right|^2\\
&\ll_\epsilon(KNBL)^\epsilon
(K+NB)B\{N+L+(NL)^{2/3}\}.
\end{aligned}
\qquad\text{(MF2.29)}
}\]

Fixed supports and ray sectors may be imposed while preserving both independence conditions.

Expand \(\fitinline{1_{(P,k)=1}}\) completely over \(\fitinline{D\mid(P,k)}\) and apply divisor Cauchy. Put \(\fitinline{P=DP'}\), \(\fitinline{k=Dk'}\) and

\[\fitmath{
a_D(P')=a(DP')1_{(D,P')=1},\qquad
\beta_D(n,b)=\beta(n,b)\chi_D(nb)^3\chi_D(n)^{-2}1_{(D,b)=1}.
}\]

Energy gains \(\fitinline{q_D^{-1}}\), with lengths \(\fitinline{K/q_D,L/q_D}\). Do not restore \(\fitinline{(P',k')=1}\). Apply \protect\hyperlink{mf:eq:B2.27}{(MF2.27)} to \(\fitinline{k'}\) with column coefficient

\[\fitmath{
\beta_D(n,b)\sum_{P'}a_D(P')q_{P'}^{-1/2}
\chi_{P'}(n)^{-2}1_{(P',b)=1}
}\]

as its row-independent coefficient. In its notation the positive output for one \(\fitinline{C,G,T}\)-block is bounded by

\hypertarget{mf:eq:B2.30}{}

\[\fitmath{
\begin{aligned}
(KNBL)^\epsilon
\sum_D\frac1{q_D}T\sum_{q_t\asymp T}\sum_r
\left(\frac K{q_Dq_r}+\frac{NG}{C^2}\right)\frac C{q_r}
\sum_{c',m,h}|\beta_D(rc'm,rc'ht^2)|^2\\
\hspace{15mm}\cdot
\left|\sum_{P'}\frac{a_{D,r}(P')}{\sqrt{q_{P'}}}
\chi_{P'}(c'm)^{-2}1_{(P',ht)=1}\right|^2,
\qquad a_{D,r}(P')=a_D(P')\chi_{P'}(r)^{-2}.
\end{aligned}
\qquad\text{(MF2.30)}
}\]

The identity used to obtain it, with zeros included, is

\[\fitmath{
\chi_{P'}(rc'm)^{-2}1_{(P',rc'ht^2)=1}
=\chi_{P'}(r)^{-2}\chi_{P'}(c'm)^{-2}1_{(P',ht)=1}.
}\]

The character supplies the zeros at \(\fitinline{r,c',m}\).

Only in this positive sum discard \(\fitinline{|\beta_D|^2\le1}\). For fixed \(\fitinline{h,t,D,r}\), coefficients are independent of squarefree primary \(\fitinline{j=c'm=n/r}\), with mass \(\fitinline{\sum_{P'}q_{P'}^{-1}|a_{D,r}(P')|^2 1_{(P',ht)=1}\ll(KNBL)^\epsilon}\). Group by \(\fitinline{j}\), retaining its \(\fitinline{S}\)-parts and norm \(\fitinline{O(N/q_r)}\); \protect\hyperlink{mf:eq:B2.25}{(MF2.25)} and \(\fitinline{O(G/C)}\) values of \(\fitinline{h}\) give on \(\fitinline{q_D\asymp D_0,q_r\asymp r_0}\)

\[\fitmath{
\ll(KNBL)^\epsilon\frac GC
\left\{\frac N{r_0}+\frac L{D_0}+
\left(\frac{NL}{r_0D_0}\right)^{2/3}\right\}.
}\]

The \(\fitinline{t}\)-sum and \(\fitinline{C/q_r}\) give \(\fitinline{T^2(C/r_0)(G/C)\asymp B/r_0}\); \(\fitinline{O(r_0)}\) values of \(\fitinline{r}\) cancel its denominator and \(\fitinline{\sum_{q_D\asymp D_0}q_D^{-1}\ll1}\). Replace \(\fitinline{D_0,r_0}\) by one in increasing sieve factors, use \(\fitinline{NG/C^2\ll NB}\) and sum dyads. This proves \protect\hyperlink{mf:eq:B2.29}{(MF2.29)} with all original masks.

\hypertarget{whole-index-marks-and-separation-of-their-phases}{}

\hypertarget{a.7.1.5-whole-index-marks-and-separation-of-their-phases}{%
\subparagraph{A.7.1.5 Whole-index marks and separation of their phases}\label{a.7.1.5-whole-index-marks-and-separation-of-their-phases}}

For marks \(\fitinline{\mathcal Q}\) outside the base local set, expand \(\fitinline{\prod_q1_{q\mid A}=\prod_q(1-\chi_q(A)^0)}\) and apply compatible \protect\hyperlink{mf:eq:B2.13}{(MF2.13)}. Each inactive contribution is

\[\fitmath{
1-(1-q_q^{-1})=q_q^{-1}.
}\]

The subtracted active branch contributes \(\fitinline{q_q^{-1/2}\chi_q(\lambda^4\mu)^{-2}}\), scalar \(\fitinline{-\omega_{q,0}}\). Thus

\hypertarget{mf:eq:B2.31}{}

\[\fitmath{
\begin{aligned}
\mathcal T_{\mathcal Q}(X;\Psi;V)
&=\sum_{\mathcal B}\zeta_{\mathcal Q}
\sum_{\mu\ne0}\frac{d(\mu)\alpha(\mu)\vartheta(\lambda^4\mu)}{\sqrt{q_\mu}}
\prod_{p\in\mathcal A_0}B_{p,j_p}(\lambda^4\mu)
\prod_{q\in A_{\mathcal Q}}q_q^{-1/2}\chi_q(\lambda^4\mu)^{-2}
V^\sharp(q_\mu X/q_c^2),\\
\zeta_{\mathcal Q}
&=-\frac i{81}\overline{\alpha(c)}^{\,2}
\widehat\varphi(h_0)\overline{\kappa_F}
\prod_{p\in\mathcal I_0}(1-q_p^{-1})\prod_{q\in I_{\mathcal Q}}q_q^{-1}
\prod_{p\in\mathcal A_0}\chi_p(\sigma_p)^{-2}\omega_{p,j_p}
\prod_{q\in A_{\mathcal Q}}\chi_q(\sigma_q)^{-2}(-\omega_{q,0}).
\end{aligned}
\qquad\text{(MF2.31)}
}\]

Here base sets are \(\fitinline{\mathcal A_0,\mathcal I_0}\), marks split as \(\fitinline{\mathcal Q=A_{\mathcal Q}\sqcup I_{\mathcal Q}}\), and \(\fitinline{c=c_F\prod_{p\in\mathcal A_0\cup A_{\mathcal Q}}p}\). Inactive marks are absent from conductor and kernel. A mark already in the base set conflicts with its primal zero, so delete that term first.

The phase at an active prime has the exact form

\hypertarget{mf:eq:B2.32}{}

\[\fitmath{
\chi_p(\sigma_p)^{-2}\omega_{p,j}
=\xi_{p,j}\chi_p(c/p)^{2j+2},
\qquad |\xi_{p,j}|=1,
\qquad\text{(MF2.32)}
}\]

Here \(\fitinline{\xi_{p,j}}\) depends only on \(\fitinline{p,j}\). Substitute \(\fitinline{\epsilon_p=-\lambda^{-5}(c/p)^{-2}}\) in \protect\hyperlink{mf:eq:B2.11}{(MF2.11)}, using \(\fitinline{2j+2\equiv-2\pmod6}\) at \(\fitinline{j=4}\). Distinct active primes therefore give

\hypertarget{mf:eq:B2.33}{}

\[\fitmath{
\chi_p(q)^{2j_p+2}\chi_q(p)^{2j_q+2}
=\left(\frac qp\right)_3^{j_p+j_q+2}.
\qquad\text{(MF2.33)}
}\]

Residual row exponent one and mark exponent zero give phase one. Exponent four would retain a cubic row interaction, so freeze every nonresidual base prime first.

Let \(\fitinline{R}\) be residual row, \(\fitinline{P}\) active mark product, \(\fitinline{F_{\rm act}}\) frozen active product. Separately fix \(\fitinline{R,P\bmod\mathfrak M^2}\), determining \(\fitinline{d,\vartheta,\kappa_F}\) without pair restrictions. Equation \protect\hyperlink{mf:eq:B2.33}{(MF2.33)} assigns row--row/frozen-row phases to a row scalar and mark--mark/frozen-mark phases to a tuple scalar; \(\fitinline{\overline{\alpha(c)}^2}\) separates. Primal collisions retain \(\fitinline{1_{(P,R)=1}}\).

After fixing the dual unit, ramified power, and extracted frozen factors, the remaining dual part is \(\fitinline{nb^3}\). Its moving columns are exactly

\hypertarget{mf:eq:B2.34}{}

\[\fitmath{
\chi_R(nb^3)^3=\chi_R(nb)^3,\qquad
\chi_P(nb^3)^{-2}=\chi_P(n)^{-2}1_{(P,b)=1}.
\qquad\text{(MF2.34)}
}\]

Source, additive character and frozen columns depend only on the full extracted dual index. Frozen collisions are fixed row, tuple or dual restrictions. Common separation therefore gives tuple coefficients independent of \(\fitinline{R,n,b}\), dual coefficients independent of \(\fitinline{R,P}\), and exactly \protect\hyperlink{mf:eq:B2.29}{(MF2.29)}'s mixed masks.

\hypertarget{reflected-energy-and-its-two-consequences}{}

\hypertarget{a.7.1.6-reflected-energy-and-its-two-consequences}{%
\subparagraph{A.7.1.6 Reflected energy and its two consequences}\label{a.7.1.6-reflected-energy-and-its-two-consequences}}

\hypertarget{block-parameters-and-estimate}{}

Block parameters and estimate

Fix \(\fitinline{q_f\asymp Z^V}\), \(\fitinline{Q=\log_Zq_{r_\rho}}\), \(\fitinline{X=Z^{N_*}}\), and rows \(\fitinline{0<q_k\ll Z^M}\). Split the ideal of \(\fitinline{k}\) as

\[\fitmath{
(k)=k_{\rm pow}k_{\rm sup}R.
}\]

The maximal powerful part contains valuations \(\fitinline{\ge2}\); \(\fitinline{k_{\rm sup}}\) contains valuation-one primes in \(\fitinline{S\cup\operatorname{supp}(fr_\rho)}\); residual \(\fitinline{R}\) is squarefree, with disjoint supports. Freeze both actual parts, unit and sectors. On dyads \(\fitinline{q_{k_{\rm pow}}\asymp Z^O}\), \(\fitinline{q_R\asymp Z^H}\), retain the actual row cutoff:

\hypertarget{mf:eq:B2.35}{}

\[\fitmath{
0\le H\le M-O+O(1/\log Z).
\qquad\text{(MF2.35)}
}\]

Equality need not hold.

Choose the active zero-exponent base primes and split every \(\fitinline{j=4}\) factor. For a large \(\fitinline{j=4}\) term use the disjoint assignment

\[\fitmath{
1_{p\mid nb^3}=1_{p\mid n}+1_{p\nmid n}1_{p\mid b}.
}\]

Let \(\fitinline{A_0}\) be active non-slot radical length outside \(\fitinline{R}\); \(\fitinline{S_0}\) counts small \(\fitinline{j=4}\) and active base-zero factors; \(\fitinline{N_0,B_0}\) count large factors assigned to squarefree/cube indices. Let \(\fitinline{z_a=\sum_{i\,\mathrm{active}}z_i\le z}\), so \(\fitinline{q_P\asymp Z^{z_a}}\). All are nonnegative.

For actual frozen products \(\fitinline{F_{\rm act},N_F,B_F}\) of logarithmic norms \(\fitinline{A_0,N_0,B_0}\), write a dual index as

\[\fitmath{
\mu=u\lambda^hN_Fn(B_Fb)^3,\quad h\ge-4,\quad
q_n\asymp Z^v,\quad q_b\asymp Z^{\ell_b},\quad
e_\lambda=h\log_Z3,
}\]

Here \(\fitinline{v,\ell_b\ge0}\); retain source restrictions and allowed \(\fitinline{S}\)-parts. Define

\hypertarget{mf:eq:B2.36}{}

\[\fitmath{
T_d=2H+2A_0+2z_a-N_*-N_0-3B_0,\qquad
u_h=\min\{v,z_a,(v+z_a)/3\}.
\qquad\text{(MF2.36)}
}\]

For any fixed \(\fitinline{\tau_{\rm ref}>0}\), all dyads outside

\hypertarget{mf:eq:B2.37}{}

\[\fitmath{
v+3\ell_b+e_\lambda\le T_d+\tau_{\rm ref}
\qquad\text{(MF2.37)}
}\]

have arbitrarily small total size. A retained dyad, including the sum over active tuples, satisfies

\hypertarget{mf:eq:B2.38}{}

\[\fitmath{
\begin{aligned}
E_{\rm ref}
={}&\frac O2+\max(H,v+\ell_b)-S_0-B_0+z_a-u_h-\ell_b-\frac23e_\lambda\\
&-\frac12(T_d-v-3\ell_b-e_\lambda)_+.
\end{aligned}
\qquad\text{(MF2.38)}
}\]

For each frozen pair, energy over \(\fitinline{R}\) is \(\fitinline{\ll Z^{E_{\rm ref}-O/2+\epsilon}}\); summing the \(\fitinline{O}\)-dyad and local choices gives \(\fitinline{Z^{E_{\rm ref}+\epsilon}}\). Uniformity is as in Appendix A.7.1.1. Row restrictions are independent of active tuples/dual indices except \(\fitinline{1_{(P,R)=1}}\); character, puncture and original product mark are independent of subsequently averaged \(\fitinline{f}\). Thresholds depend only on bounded log ranges, profiles, fixed moduli, \(\fitinline{\epsilon,\tau_{\rm ref}}\); negative \(\fitinline{e_\lambda=O(1/\log Z)}\).

\textbf{Proof.} Each inactive slot has mass \(\fitinline{\sum_{p\in\mathcal P_i}q_p^{-1}\ll Z^\epsilon}\), including bounded lists. Freeze it by row-space triangle, making collisions fixed row restrictions; active tuples stay inside the norm.

Equation \protect\hyperlink{mf:eq:B2.7}{(MF2.7)}, before the local factors, gives

\[\fitmath{
\frac{|d(\mu)|}{\sqrt{q_\mu}}
\le27\,q_{\lambda^h}^{-1/3}q_{N_Fn}^{-1/2}q_{B_Fb}^{-1}.
}\]

Small \(\fitinline{j=4}\) and active zero factors contribute \(\fitinline{q_p^{-1/2}}\). Large \(\fitinline{j=4}\) contributes \(\fitinline{q_p^{1/2}}\): squarefree assignment cancels its root denominator, cube assignment leaves \(\fitinline{q_p^{-1/2}}\). Shared \(\fitinline{N_F,b}\) cube occurrences stay in \(\fitinline{b}\). Hence the common coefficient per tuple has exponent

\hypertarget{mf:eq:B2.39}{}

\[\fitmath{
-v/2-\ell_b-e_\lambda/3-(S_0+B_0+z_a)/2.
\qquad\text{(MF2.39)}
}\]

The denominator is \(\fitinline{c=c_FF_{\rm act}RP}\). With actual annular coordinates

\[\fitmath{
y_R=q_R/Z^H,\quad y_n=q_n/Z^v,\quad
y_b=q_b/Z^{\ell_b},\quad y_i=q_{p_i}/Z^{z_i},
}\]

the kernel argument is exactly

\hypertarget{mf:eq:B2.40}{}

\[\fitmath{
\frac{q_\mu X}{q_c^2}
=Y\,y_ny_b^3y_R^{-2}\prod_{i\ {\rm active}}y_i^{-2},
\qquad
Y=q_{c_F}^{-2}Z^{v+3\ell_b+e_\lambda-T_d}.
\qquad\text{(MF2.40)}
}\]

This remains exact at shared primes. The annular monomial is bounded above/below, so \protect\hyperlink{mf:eq:B2.22}{(MF2.22)} extracts

\[\fitmath{
m_{\rm ker}=Z^{-(T_d-v-3\ell_b-e_\lambda)_+/4}
}\]

to be factored from the whole linear row vector, with bounded normalized seminorms.

Before separation remove tails: violating \protect\hyperlink{mf:eq:B2.37}{(MF2.37)} forces argument \(\fitinline{>Z^{\tau_{\rm ref}/2}}\) for large \(\fitinline{Z}\). Use \(\fitinline{|d(\mu)|/\sqrt{q_\mu}\le27\,3^{4/3}}\) and local factors \(\fitinline{\le q_p^{1/2}}\). Rows, tuples, branches and \(\fitinline{1+q_c^2/X}\) cost \(\fitinline{Z^{B_{\rm tail}}}\), fixed before kernel order. Apply \protect\hyperlink{mf:eq:B2.23}{(MF2.23)} at \(\fitinline{a=X/q_c^2}\), \(\fitinline{U=Z^{\tau_{\rm ref}/2}}\) with

\[\fitmath{
A>1+2(B_{\rm tail}+D)/\tau_{\rm ref}.
}\]

This bounds all discarded \(\fitinline{\mu}\) by \(\fitinline{O(Z^{-D})}\) times a finite seminorm and fixed height power.

Separate each retained \protect\hyperlink{mf:eq:B2.40}{(MF2.40)} once, then enlarge the positive row norm to \(\fitinline{q_R\ll Z^H}\). Phase separation gives \protect\hyperlink{mf:eq:B2.29}{(MF2.29)} with

\[\fitmath{
K=Z^H,\qquad N=Z^v,\qquad B=Z^{\ell_b},\qquad L=Z^{z_a}.
}\]

Leave \(\fitinline{q_P^{-1/2}}\) inside that norm. The outside squared coefficient is \(\fitinline{Z^{-v-2\ell_b-2e_\lambda/3-S_0-B_0+\epsilon}}\). Since

\[\fitmath{
\max\{v,z_a,2(v+z_a)/3\}=v+z_a-u_h,
}\]

the hybrid norm gives \protect\hyperlink{mf:eq:B2.38}{(MF2.38)}'s first line without \(\fitinline{O/2}\); \(\fitinline{m_{\rm ker}^2}\) supplies its last term.

The powerful count proved after \protect\hyperlink{cf:eq:B.23}{(42)} is \(\fitinline{O(Z^{O/2})}\); the supported part divides \(\fitinline{\operatorname{rad}(Sfr_\rho)}\), giving \(\fitinline{Z^\epsilon}\) choices, as do local branches. Divisor Cauchy and the positive frozen-part sum add \(\fitinline{O/2+\epsilon}\) once; active tuples are already included. This proves \protect\hyperlink{mf:eq:B2.38}{(MF2.38)}.

\hypertarget{general-label-and-puncture-width}{}

General label and puncture width

At an active non-slot prime, the coefficient of its log-norm in \(\fitinline{2A_0-N_0-S_0-4B_0}\) is

\begin{longtable}[]{@{}ll@{}}
\toprule\noalign{}
Local choice & Coefficient \\
\midrule\noalign{}
\endhead
\bottomrule\noalign{}
\endlastfoot
\(\fitinline{j}\) odd or \(\fitinline{j=2}\) & \(\fitinline{2}\) \\
\(\fitinline{j=0}\), active & \(\fitinline{1}\) \\
\(\fitinline{j=4}\), small or assigned to \(\fitinline{n}\) & \(\fitinline{1}\) \\
\(\fitinline{j=4}\), assigned to \(\fitinline{b}\) & \(\fitinline{-2}\) \\
\(\fitinline{j=0}\), inactive & \(\fitinline{0}\) \\
\end{longtable}

Residual primes contribute two. Primewise these coefficients are bounded by valuations of \(\fitinline{(q_k^2/q_{k_{\rm pow}})q_fq_{r_\rho}}\): powerful valuation \(\fitinline{e\ge2}\) supplies \(\fitinline{e}\), simple rows two, label/puncture-only primes one; overlaps increase availability. Thus

\[\fitmath{
q_R^2 Z^{2A_0-N_0-S_0-4B_0}
\le C_S\,\frac{q_k^2}{q_{k_{\rm pow}}}\,q_fq_{r_\rho},
}\]

and therefore

\hypertarget{mf:eq:B2.41}{}

\[\fitmath{
T_d-S_0-B_0
\le2M-O+Q+V-N_*+2z_a+O(1/\log Z).
\qquad\text{(MF2.41)}
}\]

This remains uniform in actual \(\fitinline{f,r_\rho}\) when completion gives \(\fitinline{N_*=N-3\ell_h}\), without replacing \(\fitinline{H}\) by \(\fitinline{M-O}\).

With no slots, \(\fitinline{z_a=u_h=0}\). On a retained dyad, \protect\hyperlink{mf:eq:B2.38}{(MF2.38)} gives

\[\fitmath{
E_{\rm ref}\le O/2+\max\{H,T_d-S_0-B_0\}
+O(1/\log Z)+\tau_{\rm ref}.
}\]

Drop nonpositive terms for the row branch; for the column branch use \(\fitinline{v\le T_d-3\ell_b-e_\lambda+\tau_{\rm ref}}\). Apply \protect\hyperlink{mf:eq:B2.35}{(MF2.35)}, \protect\hyperlink{mf:eq:B2.41}{(MF2.41)}, choose small buffer/losses and sum dyads:

\hypertarget{mf:eq:B2.42}{}

\[\fitmath{
\sum_{\substack{0<q_k\ll Z^M\\q_{k_{\rm pow}}\asymp Z^O}}
|\mathcal T_1(Z^{N_*};k;V)|^2
\ll_\epsilon
Z^{\,O/2+\max\{M-O,\;2M-O+Q+V-N_*\}+\epsilon}.
\qquad\text{(MF2.42)}
}\]

In particular, with \(\fitinline{f=r_\rho=1}\), summing \(\fitinline{O}\) gives the unmarked completed bound

\hypertarget{mf:eq:B2.43}{}

\textbf{(MF2.43)} is the unmarked bound \protect\hyperlink{cf:eq:4.4}{(33)}, with \(\fitinline{(H,N)=(Z^M,Z^{N_*})}\) and equivalent power notation \(\fitinline{Z^{\max(M,2M-N_*)+\epsilon}}\), for the tests admitted in Appendix A.7.1.1. The preceding argument extends its annular/Gaussian statement to that test class.

This includes the balanced energy \(\fitinline{Z^{1+\epsilon}}\) at \(\fitinline{M=N_*=1}\).

\hypertarget{short-marked-energy-without-label-or-puncture}{}

Short marked energy without label or puncture

When \(\fitinline{f=r_\rho=1}\), every active non-slot prime outside \(\fitinline{R}\) divides the powerful part. Thus

\[\fitmath{
2A_0\le O+O(1/\log Z),\qquad 0\le N_0\le A_0.
}\]

\textbf{Theorem A.7.1} (Short marked completed energy). \emph{\hypertarget{mf:marked-energy}{}If \(\fitinline{M+2z\le N_*}\), then for every \(\fitinline{\epsilon>0}\), \hypertarget{mf:eq:B2.44}{}}

\[\fitmath{
\sum_{0<q_k\ll Z^M}|\mathcal T_d(Z^{N_*};k;V)|^2
\ll_\epsilon
Z^{\,M+(N_*+2z-2M)_+/4+\epsilon}.
\qquad\text{(MF2.44)}
}\]

\emph{Proof.} Use \protect\hyperlink{mf:eq:B2.38}{(MF2.38)}, suppressing \(\fitinline{O(1/\log Z)}\) and tiny \(\fitinline{\tau_{\rm ref}}\). Since

\[\fitmath{
T_d-H\le M+2z-N_*\le0,
}\]

retention gives \(\fitinline{\max(H,v+\ell_b)=H+O(\tau_{\rm ref}+1/\log Z)}\). If \(\fitinline{u_h=z_a}\), then \(\fitinline{E_{\rm ref}\le M}\) up to this error; otherwise \(\fitinline{u_h\ge v/2}\). Put \(\fitinline{y_d=v+3\ell_b+e_\lambda}\); \(\fitinline{\ell_b\ge0}\), \(\fitinline{e_\lambda\ge-O(1/\log Z)}\) imply

\[\fitmath{
u_h+\ell_b+\frac23e_\lambda
+\frac12(T_d-y_d)_+
\ge \frac14y_d+\frac12(T_d-y_d)_+-O(1/\log Z)
\ge\frac14T_d-O(1/\log Z).
}\]

With \(\fitinline{\Delta_H=M-O-H\ge-O(1/\log Z)}\), substitution into \protect\hyperlink{mf:eq:B2.38}{(MF2.38)} gives

\[\fitmath{
\begin{aligned}
4(E_{\rm ref}-M)
&\le N_*+2z_a-2M-2\Delta_H
-(2A_0-N_0+B_0+4S_0)+O(\tau_{\rm ref}+1/\log Z)\\
&\le N_*+2z-2M+O(\tau_{\rm ref}+1/\log Z).
\end{aligned}
}\]

Restore the row branch, tails and losses to obtain \protect\hyperlink{mf:eq:B2.44}{(MF2.44)}, independently of moments or zero-free conclusions. \ensuremath{\square}

\hypertarget{the-inverse-and-plain-moment-inductions}{}

\hypertarget{a.7.2-the-inverse-and-plain-moment-inductions}{%
\paragraph{A.7.2 The inverse and plain moment inductions}\label{a.7.2-the-inverse-and-plain-moment-inductions}}

We prove inverse moments and unconditional zero-slot/canonical common-profile positive-slot plain moments from the finite identities and reflected energy above. Inverse descent terminates at reflection; plain descent is independent of the inverse moment.

\hypertarget{common-conventions-and-analytic-bookkeeping}{}

\hypertarget{a.7.2.1-common-conventions-and-analytic-bookkeeping}{%
\subparagraph{A.7.2.1 Common conventions and analytic bookkeeping}\label{a.7.2.1-common-conventions-and-analytic-bookkeeping}}

Retain Appendix A.7.1's conventions: primary ideal columns avoid \(\fitinline{S}\); all element rows and displayed zeros remain. Fixed ray moduli lie in \(\fitinline{S}\), moving punctures are separate. Put

\[\fitmath{
a_0(n)=\overline{\alpha(n)}\gamma_2(n)\qquad(n\text{ squarefree}).
}\]

For disjoint prime lists \(\fitinline{\mathcal P_i}\), fixed throughout the current row and label averages, define a \textbf{product mark}

\hypertarget{mf:eq:B3.1}{}

\textbf{(MF3.1)} is the product mark \protect\hyperlink{mf:eq:B2.3}{(MF2.3)} on slot subset \(\fitinline{I}\), with lists and coefficients fixed throughout the current row and label averages.

Slot count is fixed; coefficients need not be multiplicative and include individual smooth weights. Product form, beyond a divisor bound, allows

\hypertarget{mf:eq:B3.2}{}

\[\fitmath{
\mathbf1_{p\mid D_1\cdots D_h n}
=\sum_{j=1}^h\mathbf1_{p\mid D_j}\prod_{a<j}\mathbf1_{p\nmid D_a}
 +\mathbf1_{p\mid n}\prod_{a\le h}\mathbf1_{p\nmid D_a}
\qquad\text{(MF3.2)}
}\]

to assign each slot first to one extracted factor or retain it on the residual column. Assigned coefficients are divisor-bounded; survivors retain product form. Allocate the two copies of a squared mark independently.

\textbf{Common-profile lemma.} For finitely many compact normalized-norm profiles, radial Schwartz kernels and fixed-ray Fourier expansions, keep nonsmooth arithmetic restrictions in outer weights. Multiply a coupled \(\fitinline{H(y_1,\ldots,y_d)}\) by larger individual cutoffs equal to one on current supports and use

\hypertarget{mf:eq:B3.3}{}

\[\fitmath{
H(y)=(2\pi)^{-d}\int_{\mathbf R^d}\widehat H(t)\prod_jy_j^{it_j}\,dt,
\qquad
\int|\widehat H(t)|\langle t\rangle^Jdt
 \ll_{J,d}\max_{|\beta|\le2r}\|D^\beta H\|_\infty,
\quad 2r>J+d,
\qquad\text{(MF3.3)}
}\]

Here \(\fitinline{D_j=y_j\partial_{y_j}}\), \(\fitinline{\langle t\rangle=1+|t|}\). Define the measure on the ambient box before choosing row or live label; norm relations restrict evaluation points only. Include inverse roots once in the coupled profile; individual column tests get only cutoffs and Fourier powers.

Logarithmic integration by parts with \(\fitinline{(1-\Delta)^r}\) proves the bound. For \(\fitinline{K(A\prod y_j^{e_j})}\), Euler derivatives are fixed multiples of \(\fitinline{(x\partial_x)^rK(x)}\) uniformly in \(\fitinline{A>0}\), so fixed derivative orders add no \(\fitinline{Z}\) power; norm heights cost polynomials. Also,

\hypertarget{mf:eq:B3.4}{}

\textbf{(MF3.4)} is \protect\hyperlink{cf:eq:B.25}{(44)} on \(\fitinline{\mathcal O}\), with \(\fitinline{(U,A)=(Y,B)}\) and the kernel seminorm absorbed in the constant. It holds for \(\fitinline{a>0}\), \(\fitinline{Y\ge1}\), \(\fitinline{B>1}\).

Discard tails on the genuine expression before coprime extension or absolute Fourier integration. Choose a sector-fixed outer ball whose complement forces kernel argument \(\fitinline{>Z^\tau}\), retaining the full kernel inside. If raw factors and \(\fitinline{1+a^{-1}}\) cost \(\fitinline{Z^{B_0}}\), choose \(\fitinline{B>1+(B_0+T)/\tau}\) in \protect\hyperlink{mf:eq:B3.4}{(MF3.4)} for \(\fitinline{O(Z^{-T})}\). Polynomial scale bounds fix \(\fitinline{B_0}\) before \(\fitinline{B}\).

For uniformity choose terminal profile/height orders, then propagate backwards through the bounded-depth acyclic graph using \protect\hyperlink{mf:eq:B3.3}{(MF3.3)}. Normalize discrete counting measures after extracting their \(\fitinline{Z^b}\) mass once. Cauchy averages child exponents; Minkowski takes their maximum after coefficient mass, so numerical loss grows only with path length. Fix quotient/separation cutoffs and bounded clipping families before the final \(\fitinline{Z}\) threshold. A nonempty nominal negative-length window has \(\fitinline{1\le Z^{-n}\le B}\) (fixed upper endpoint), hence clips to length zero by bounded dilation.

\textbf{Masked primitive Poisson.} For primitive zero-extended \(\fitinline{\psi\bmod\mathfrak m}\) and radial Schwartz \(\fitinline{\Phi}\), use

\hypertarget{mf:eq:B3.5}{}

\textbf{(MF3.5)} is \protect\hyperlink{cf:eq:C.3}{(47)}, under \(\fitinline{(H,m,r)=(K,\mathfrak m,R)}\), \(\fitinline{\gamma(\psi)=\gamma(\psi;\mathfrak m)}\) and \(\fitinline{\widehat\Phi=\mathcal F\Phi}\).

The proof is \S{}\protect\hyperlink{cf:sec:C.1}{A.4.1}'s complete mask expansion and primitive finite transform; Parseval gives \(\fitinline{|\gamma|=1}\). For the primitive principal character use modulus one, \(\fitinline{\gamma=1}\), including \(\fitinline{h=0}\). Its nonzero frequencies cost \(\fitinline{O_\Phi(d(\operatorname{rad}R))}\) since \(\fitinline{A\sum_{h\ne0}|\mathcal F\Phi(Aq_h)|\ll1}\); restore them with the original outer mask at diagonal cost. Expand fixed rays into characters first, keeping deleted conductor zeros in \(\fitinline{R}\).

\hypertarget{a-canonical-marked-energy}{}

\hypertarget{a.7.2.2-a-canonical-marked-energy}{%
\subparagraph{A.7.2.2 A canonical marked energy}\label{a.7.2.2-a-canonical-marked-energy}}

Fix squarefree \(\fitinline{\rho(n)=1_{(n,r_\rho)=1}}\) independent of row \(\fitinline{k}\) and label \(\fitinline{f}\), with \(\fitinline{Q=\log_Zq_{r_\rho}}\). Let \(\fitinline{q_f\asymp Z^V}\) and label-only \(\fitinline{0\le w(f)\le Cd(f)^C}\). For \(\fitinline{F=N+V}\), define

\hypertarget{mf:eq:B3.6}{}

\[\fitmath{
\mathcal E(N,V,M)=Z^{-V}\sum_f w(f)\sum_{0<q_k\ll Z^M}
 \left|Z^{-N/2}\sum_{n\ {\rm sf}}a_0(n)\nu(n)\rho(n)
       \chi_n(k)\chi_n(f)^4d_I(n)W(q_n/Z^N)\right|^2.
\qquad\text{(MF3.6)}
}\]

Here \(\fitinline{\nu}\) is fixed finite-ray, lengths are bounded, slot sum is \(\fitinline{\le z_0}\), and no additional \(\fitinline{n}\)-coefficient is allowed.

\textbf{Proposition A.7.2} (Canonical marked energy). \emph{\hypertarget{mf:canonical-energy}{} For any fixed \(\fitinline{c_*>0}\), \hypertarget{mf:eq:B3.7}{}}

\[\fitmath{
F-M-Q-z_0\ge c_*,\qquad4F-3M-6z_0\ge c_*
\quad\Longrightarrow\quad
\mathcal E(N,V,M)\ll Z^{F+\epsilon}.
\qquad\text{(MF3.7)}
}\]

\emph{Uniformity includes moving punctures, the stated label weights and every surviving slot subset, always with the original cap \(\fitinline{z_0}\).}

\emph{Proof.} \textbf{Terminal completion.} Choose \(\fitinline{0<d\le c_*/200}\). Set \(\fitinline{\Psi_k(n)=\nu(n)\rho(n)\chi_n(k)\chi_n(f)^4}\) multiplicatively on all ideals. For \(\fitinline{V_*(y)=y^{1/2}W(y)}\) define

\[\fitmath{
T_{d_*}(X;k)=\sum_{n\ {\rm sf},b}
 \frac{a_0(n)\Psi_k(n)\overline{\alpha(b)}^3\Psi_k(b)^3}
      {\sqrt{q_n}\,q_b}\,d_*(nb^3)V_*(q_nq_b^3/X).
}\]

The exact cube inversion is

\hypertarget{mf:eq:B3.8}{}

\[\fitmath{
Z^{-N/2}\sum_{n\ {\rm sf}}a_0(n)\Psi_k(n)d_I(n)W(q_n/Z^N)
=\sum_h\frac{\mu(h)\overline{\alpha(h)}^3\Psi_k(h)^3}{q_h}
 T_{d_h}(Z^N/q_h^3;k),\quad d_h(A)=d_I(h^3A).
\qquad\text{(MF3.8)}
}\]

Opening the completion and setting \(\fitinline{c=hb}\) leaves mark \(\fitinline{d_I(nc^3)}\) independent of \(\fitinline{h}\); \(\fitinline{\sum_{h\mid c}\mu(h)=1_{c=1}}\) verifies the identity, normalization and masks.

Use fixed dyadic windows with unit center zero. If \(\fitinline{V<d}\) and the \(\fitinline{h}\)-center \(\fitinline{\ell_1<d}\), freeze \(\fitinline{h}\), assign slots by \protect\hyperlink{mf:eq:B3.2}{(MF3.2)} and apply reflected energy (RE) at \(\fitinline{N_*=N-3\log_Zq_h}\), surviving cap \(\fitinline{z_a\le z_0}\). The factor \(\fitinline{\Psi_k(h)^3}\) contracts the joint \(\fitinline{(f,k)}\) norm and \(\fitinline{\sum_{q_h\asymp Z^{\ell_1}}q_h^{-1}\ll1}\). If \(\fitinline{c}\) is the current margin in \protect\hyperlink{mf:eq:B3.7}{(MF3.7)} and \(\fitinline{\eta}\) bounds window errors, RE gives

\hypertarget{mf:eq:B3.9}{}

\[\fitmath{
H\le M-O+\eta,\qquad
T_d-S_0-B_0\le M-O+2z_a-z_0-c+5d+4\eta.
\qquad\text{(MF3.9)}
}\]

For the second, apply \protect\hyperlink{mf:eq:B2.41}{(MF2.41)}, the first margin, and \(\fitinline{2V+3\log_Zq_h<5d+3\eta}\).

Retention gives \(\fitinline{v+3b+e_\lambda\le T_d+\tau_{\rm ref}}\), \(\fitinline{e_\lambda\ge-\eta}\). The row branch yields \(\fitinline{E\le M+z_a+5\eta/3\le F-c+5\eta/3}\). The column branch gives \(\fitinline{E\le O/2+T_d-S_0-B_0+5\eta/3+\tau_{\rm ref}}\) if \(\fitinline{u=z_a}\); otherwise \(\fitinline{u\ge v/2}\) gives \(\fitinline{E\le O/2+(T_d-S_0-B_0)/2+z_a+7\eta/6+\tau_{\rm ref}/2}\). Substitution and the respective margins give

\hypertarget{mf:eq:B3.10}{}

\[\fitmath{
E\le F-2c+5d+17\eta/3+\tau_{\rm ref},\qquad
E\le F-M/4-3c/4+5d/2+19\eta/6+\tau_{\rm ref}/2.
\qquad\text{(MF3.10)}
}\]

These are \(\fitinline{<F}\) for \(\fitinline{c\ge c_*/2}\) and \(\fitinline{\eta,\tau_{\rm ref}}\) and reflected divisor loss \(\fitinline{\le c_*/1000}\). Summing labels costs \(\fitinline{Z^{V+o(1)}}\), canceled by \(\fitinline{Z^{-V}}\).

\textbf{Nonterminal block.} If \(\fitinline{V<d}\), \(\fitinline{\ell_1\ge d}\), reopen \protect\hyperlink{mf:eq:B3.8}{(MF3.8)}, combine its cube with \(\fitinline{h}\) into \(\fitinline{b}\), center \(\fitinline{\ell}\), and put \(\fitinline{r=N-3\ell}\). If \(\fitinline{V\ge d}\), take \(\fitinline{b=1}\), \(\fitinline{\ell=0}\), \(\fitinline{r=N}\). Joint Hilbert-space triangle and common-profile separation reduce to

\hypertarget{mf:eq:B3.11}{}

\[\fitmath{
Z^{-r-2\ell-V}\sum_f w(f)\sum_k\Phi(q_k/Z^M)
 \left|\sum_{q_b\asymp Z^\ell}\beta(b,f)\chi_b(k)^3
  \sum_{n\ {\rm sf}}a_0(n)\nu(n)\rho(n)\chi_n(k)\chi_n(f)^4
        d_I(nb^3)W(q_n/Z^r)\right|^2,
\qquad\text{(MF3.11)}
}\]

Here \(\fitinline{V+\ell\ge d}\), \(\fitinline{\beta}\) is divisor-bounded, independent of \(\fitinline{k,n}\), and includes \(\fitinline{(b,f)=1}\); \(\fitinline{b,n}\) may share primes. We return this to \protect\hyperlink{mf:eq:B3.6}{(MF3.6)} with shorter rows.

\textbf{First Poisson transformation.} In the square write

\[\fitmath{
\mathfrak B=\operatorname{rad}(b_1b_2),\qquad
n_i=A_iCu_i,\quad A_i=(n_i,\mathfrak B),\quad
C=(n_1/A_1,n_2/A_2).
}\]

Then \(\fitinline{u_1,u_2}\) are coprime squarefree and avoid \(\fitinline{C\mathfrak B}\). At \(\fitinline{p\mid\mathfrak B}\), put \(\fitinline{a_i=1_{p\mid A_i}}\), \(\fitinline{\pi=v_p(b_1b_2)\bmod2}\), \(\fitinline{t_p=a_1-a_2+3\pi\bmod6}\), and \(\fitinline{R_1=\prod_{t_p\ne0}p}\). Primitive character and residual mask are

\hypertarget{mf:eq:B3.12}{}

\[\fitmath{
\chi_{u_1}\overline{\chi_{u_2}}\prod_{t_p\ne0}\chi_p^{t_p},
\qquad\mathfrak m_1=u_1u_2R_1,
\qquad R_{\rm mask}=C\mathfrak B/R_1.
\qquad\text{(MF3.12)}
}\]

Apply \protect\hyperlink{mf:eq:B3.5}{(MF3.5)}, with \(\fitinline{d_k\mid C\mathfrak B/R_1}\) and frequency \(\fitinline{h}\). Also denote the centers of \(\fitinline{A_i,C,R_1,d_k}\) by \(\fitinline{A_i,B,R,\delta}\). Actual conductor length is \(\fitinline{\log_Z(q_{n_1}q_{n_2}/(q_{A_1}q_{A_2}q_C^2))+\log_Zq_{R_1}}\). The principal case forces \(\fitinline{u_1=u_2=1}\), \(\fitinline{\pi=0}\), \(\fitinline{a_1=a_2}\) everywhere, hence \(\fitinline{n_1=n_2}\) and square \(\fitinline{b_1b_2}\). Counting even all cube pairs gives normalized \(\fitinline{O(Z^{M+o(1)})}\), also covering restored principal nonzero frequencies with the same mask.

CRT, \(\fitinline{\gamma_1(u)\gamma_2(u)=\mu(u)\alpha(u)G(u)}\) and reciprocity give inverse residual coefficients. Expand their joint fixed-ray factor \(\fitinline{G([u_1u_2^{-1}])R([u_1u_2^{-1}],[\prod p^{t_p}])}\) before coprime extension. Set \(\fitinline{E=\prod_{\pi=1}p^{a_1+a_2}}\), \(\fitinline{y=hf^2E}\). The residual exponents \(\fitinline{4a_i+1_{t_p\ne0}(1\pm t_p)+\pi(a_1+a_2)\bmod6}\) agree on both sides by this exhaustive table:

\begin{longtable}[]{@{}lllll@{}}
\toprule\noalign{}
\(\fitinline{\pi}\) & \(\fitinline{a_1}\) & \(\fitinline{a_2}\) & \(\fitinline{t_p}\) & Residual exponent \\
\midrule\noalign{}
\endhead
\bottomrule\noalign{}
\endlastfoot
0 & 0 & 0 & 0 & 0 \\
0 & 1 & 0 & 1 & 0 \\
0 & 0 & 1 & 5 & 0 \\
0 & 1 & 1 & 0 & 4 \\
1 & 0 & 0 & 3 & 4 \\
1 & 1 & 0 & 4 & 4 \\
1 & 0 & 1 & 2 & 4 \\
1 & 1 & 1 & 3 & 4 \\
\end{longtable}

An exponent zero still means a puncture. Write \(\fitinline{b_1b_2=q^2s}\), \(\fitinline{s}\) squarefree, and put

\[\fitmath{
J_2=\prod_{\pi=0,\ a_1=a_2=1}p,
\quad J=sJ_2,\quad q=J_2q_0.
}\]

Then \(\fitinline{J}\) is squarefree and the residual factor is \(\fitinline{\chi_n(J)^4 1_{(n,\operatorname{rad}q_0)=1}}\). For center \(\fitinline{j}\) of \(\fitinline{J}\), the identities

\hypertarget{mf:eq:B3.13}{}

\[\fitmath{
\log_Zq_{q_0}=\tfrac12\log_Z(q_{b_1}q_{b_2})
 -\tfrac12\log_Zq_s-\log_Zq_{J_2}\ge0,
\qquad
\log_Z\frac{q_{R_1}q_E}{q_{A_1}q_{A_2}}
=\log_Zq_s-2\log_Zq_{J_2}
\qquad\text{(MF3.13)}
}\]

give the retained row enclosure

\hypertarget{mf:eq:B3.14}{}

\[\fitmath{
H_c=2r-2B-M+4\ell+2V+\delta-j,
\qquad H_+=H_c+12\eta+\tau,
\qquad 0<q_y\le Z^{H_+}.
\qquad\text{(MF3.14)}
}\]

The actual argument is \(\fitinline{Z^Mq_h/(q_{d_k}q_{\mathfrak m_1})}\); if \(\fitinline{\le Z^\tau}\), \protect\hyperlink{mf:eq:B3.13}{(MF3.13)} gives \(\fitinline{\log_Zq_y\le H_c+12\eta+\tau}\). Discard only this outer ball's complement by \protect\hyperlink{mf:eq:B3.4}{(MF3.4)}; \(\fitinline{H_+<0}\) leaves no nonzero rows.

Extend the quotient-free expression to individually allowed \(\fitinline{u_1,u_2}\), retaining full kernel, formal product conductor and outer ball. Insert \(\fitinline{1_{(u_1,u_2)=1}=\sum_{t'\mid(u_1,u_2)}\mu(t')}\) and write \(\fitinline{u_i=t'x_i}\), center \(\fitinline{t}\). The complete M\"{o}bius sum removes artificial noncoprime conductors; squarefreeness retains \(\fitinline{(x_i,t')=1}\). With

\hypertarget{mf:eq:B3.15}{}

\[\fitmath{
s_i=r-A_i-B-t,\qquad
\varkappa_i=M-2r-2\ell-V-\delta+A_i+B-R/2,
\qquad\text{(MF3.15)}
}\]

the separated inverse polynomial on side \(\fitinline{i}\) is

\hypertarget{mf:eq:B3.16}{}

\[\fitmath{
P_i(y)=\sum_{x\ {\rm sf}}\mu(x)\nu_i(x)\rho(x)
 \overline{\chi_x(y)}\chi_x(CJ)^4\chi_x(d_k)
 \mathbf1_{(x,\operatorname{rad}(q_0t'))=1}d_{I_i}(x)W_i(q_x/Z^{s_i}).
\qquad\text{(MF3.16)}
}\]

The ray character \(\fitinline{\nu_i}\) is fixed. The \(\fitinline{f}\) factor became \(\fitinline{\overline{\chi_x(f^2)}}\), explaining row \(\fitinline{y}\) and conjugation. Priority allocation to \(\fitinline{(b_i,A_i,C,t';x_i)}\) leaves original slot subsets; \(\fitinline{CJ,q_0,t'}\) supply all earlier survival zeros.

The first common profile uses independent normalized norms

\[\fitmath{
(A_1,A_2,C,d_k,R_1,t',h,x_1,x_2).
}\]

Apart from its larger cutoffs, its real root and kernel are

\hypertarget{mf:eq:B3.17}{}

\[\fitmath{
y_d^{-1}y_R^{-1/2}y_t^{-1}(y_1y_2)^{-1/2}
 W(y_{A_1}y_Cy_ty_1)\overline{W(y_{A_2}y_Cy_ty_2)}
 K_1\!\left(A_1^{\rm ker}\frac{y_h}{y_dy_Ry_t^2y_1y_2}\right),
\quad A_1^{\rm ker}=Z^{M+h_1-\delta-R-2t-s_1-s_2},
\qquad\text{(MF3.17)}
}\]

Here \(\fitinline{h_1}\) centers bare frequency. Multiplying the profile by \(\fitinline{Z^{-9\eta/2}}\) leaves outside side factors \(\fitinline{Z^{(\varkappa_i+9\eta/2)/2}}\), exactly separating prefactor and kernel. Fresh supports satisfy \(\fitinline{|\log_Zq_{x_i}-s_i|\le4\eta}\); profile bounds are uniform for \(\fitinline{A_1^{\rm ker}>0}\).

Apply weighted Cauchy to the \textbf{whole} outer tuple before bounding weights. At fixed reconstruction and \(\fitinline{y\ne0}\),

\hypertarget{mf:eq:B3.18}{}

\[\fitmath{
\sum_{f^2E\mid(y)}w(f)\ll d((y))^{C+1}\ll Z^{o(1)}.
\qquad\text{(MF3.18)}
}\]

Thus \(\fitinline{f,h}\) disappear together through \(\fitinline{h=y/(f^2E)}\), retaining \(\fitinline{-V}\) in \(\fitinline{\varkappa_i}\). Only now smoothly majorize the positive row ball at \(\fitinline{Z^{H_+}}\). Average \(\fitinline{J}\), freeze \(\fitinline{q_0}\). If \(\fitinline{r_{\rm diff}}\) is the radical length of even-parity primes with \(\fitinline{a_1\ne a_2}\), the table gives \(\fitinline{\log_Zq_{R_1}=\log_Zq_s+r_{\rm diff}}\); hence

\hypertarget{mf:eq:B3.19}{}

\[\fitmath{
\log_Zq_{q_0}\le\ell+R/2-j+5\eta/2.
\qquad\text{(MF3.19)}
}\]

Counting \(\fitinline{q_0}\) and \(\fitinline{J}\) costs \(\fitinline{Z^{\ell+R/2+7\eta/2+o(1)}}\) once.

\textbf{Second Poisson.} On one first-Cauchy side expand \(\fitinline{|P_i(y)|^2}\), write \(\fitinline{x_a=g'z_a}\), \(\fitinline{g'=(x_1,x_2)}\) with center \(\fitinline{g}\), and apply \protect\hyperlink{mf:eq:B3.5}{(MF3.5)}: character \(\fitinline{\overline{\chi_{z_1}}\chi_{z_2}}\), modulus \(\fitinline{z_1z_2}\), mask \(\fitinline{g'}\), divisor \(\fitinline{d_2\mid g'}\) with center \(\fitinline{\theta}\), frequency \(\fitinline{k''}\). Put

\hypertarget{mf:eq:B3.20}{}

\[\fitmath{
\lambda_c=H_c-\theta-(s_i-g),\qquad
M_c=2(s_i-g)-H_c+\delta+2\theta
=M-4\ell-2A_i-2t-2g+2\theta-2V+j,
\quad M'=M_c+\eta.
\qquad\text{(MF3.20)}
}\]

Retain \(\fitinline{0<q_{d_kd_2k''}\le Z^{M'}}\). Conductor error is \(\fitinline{10\eta}\); two \(\fitinline{d_2}\) and one \(\fitinline{d_k}\) contribute \(\fitinline{3\eta}\). Subtracting \(\fitinline{H_+=H_c+12\eta+\tau}\) cancels the same \(\fitinline{\tau}\), making the complement an actual-ratio tail under \protect\hyperlink{mf:eq:B3.4}{(MF3.4)}.

The principal case \(\fitinline{z_1=z_2=1}\) has nominal total exponent

\hypertarget{mf:eq:B3.21}{}

\[\fitmath{
\varkappa_i+H_c+s_i+B+t+\ell+R/2=F-B-j.
\qquad\text{(MF3.21)}
}\]

Prefactor, enclosure, diagonal window and outer counts add \(\fitinline{\le26\eta+\tau+o(1)}\), within target, including restored principal nonzero frequencies with the same mask.

Use \(\fitinline{\mu(z)\gamma_{-1}(z)=a_0(z)\chi_z(-1)\overline{G(z)}}\). The paired ray factor \(\fitinline{G([z_2z_1^{-1}])}\) survives common-factor removal. Extend separate columns/full kernel off coprimality, insert the full M\"{o}bius sum \(\fitinline{z_a=v'n_a}\), \(\fitinline{v'}\) squarefree with center \(\fitinline{v}\). Retain \(\fitinline{(v',g'CJ)=(g',CJ)=(v',k'')=1}\). Set

\hypertarget{mf:eq:B3.22}{}

\[\fitmath{
\gamma=(q_0,t',r_g),\quad r_g=g'/d_2,\qquad
k_{\rm new}=d_kd_2k'',\qquad f_{\rm new}=JCd_2v',
\qquad\rho_\gamma(n)=\rho(n)\mathbf1_{(n,\operatorname{rad}(q_0t'r_g))=1}.
\qquad\text{(MF3.22)}
}\]

Zeros make \(\fitinline{J,C,d_2,v'}\) pairwise coprime squarefree, hence \(\fitinline{f_{\rm new}}\) squarefree, and

\hypertarget{mf:eq:B3.23}{}

\[\fitmath{
\chi_n(d_k)\chi_n(k'')\overline{\chi_n(d_2)}
 \chi_n(CJ)^4\chi_n(v')^4
=\chi_n(k_{\rm new})\chi_n(f_{\rm new})^4.
\qquad\text{(MF3.23)}
}\]

Use \(\fitinline{\overline{\chi_n(d_2)}=\chi_n(d_2)^5}\) with zeros and \(\fitinline{a_0(v'n)=a_0(v')a_0(n)\chi_n(v')^4}\). Expand \(\fitinline{G([n_2n_1^{-1}])}\) to get the same fixed \(\fitinline{\nu'}\) on both sides. Priority allocation leaves original slot subsets, with survival zeros supplied by \(\fitinline{f_{\rm new},\rho_\gamma}\). Freeze only \(\fitinline{\gamma}\); average all four label factors.

The formal child lengths and their actual enclosing windows are

\hypertarget{mf:eq:B3.24}{}

\[\fitmath{
\begin{aligned}
N_c&=r-A_i-B-t-g-v,& V_c&=B+\theta+v+j,\\
F_c&=N_c+V_c,& N'&=\max(0,N_c),\\
F'&=N'+V_c,&0\le\delta_N&=N'-N_c\le6\eta.
\end{aligned}
\qquad\text{(MF3.24)}
}\]

Fresh column and label errors are \(\fitinline{6\eta,4\eta}\). Omit empty windows; other clipping is bounded dilation. Put \(\fitinline{D_c=\ell+A_i+t+g-\theta+V}\); then

\hypertarget{mf:eq:B3.25}{}

\[\fitmath{
\begin{aligned}
F_c&=F-D_c-2\ell+j,&M_c&=M-2D_c-2\ell+j,\\
F_c-M_c&=F-M+D_c,&
4F_c-3M_c&=4F-3M+2(A_i+t+g-\theta+V)+j.
\end{aligned}\qquad\text{(MF3.25)}
}\]

Project fixed triples from the genuine reconstruction set \textbf{before} positive enlargement. Each has a source witness, giving \(\fitinline{g-\theta\ge-2\eta}\), \(\fitinline{j\le2\ell+3\eta}\), \(\fitinline{\log_Zq_{q_0}\le\ell+\eta}\), hence \(\fitinline{Q'\le Q+\ell+t+g-\theta+4\eta\le Q+D_c+4\eta}\). Therefore

\hypertarget{mf:eq:B3.26}{}

\[\fitmath{
F'-M'-Q'-z_0\ge c-5\eta,\quad
4F'-3M'-6z_0\ge c-7\eta,\quad
M-M'\ge2(V+\ell)-8\eta,\quad F'\le F+11\eta.
\qquad\text{(MF3.26)}
}\]

These hold on structural triples, not every point of a containing ball.

The second common profile uses \(\fitinline{(g',d_2,v',k'',n_1,n_2)}\), root \(\fitinline{y_d^{-1}y_v^{-1}(y_1y_2)^{-1/2}}\), windows \(\fitinline{W_i(y_gy_vy_1)\overline{W_i(y_gy_vy_2)}}\), and kernel \(\fitinline{K_2(A_2^{\rm ker}y_k/(y_dy_v^2y_1y_2))}\), \(\fitinline{A_2^{\rm ker}=Z^{H_++h_2-\theta-2v-2N_c}}\). Multiplying by \(\fitinline{Z^{-6\eta}}\) leaves exact outside factor \(\fitinline{Z^{\lambda_c+18\eta+\tau}}\). Fresh tests \(\fitinline{\omega_a(Z^{\delta_N}y)y^{\pm it_a}}\) and density are independent of \(\fitinline{k_{\rm new},f_{\rm new},\gamma}\). Retain both extracted squared magnitudes, priority masks, \(\fitinline{\mu(d_2)\mu(v')}\) and frequency mask in signed outer weights. Fourier inversion restores terms exactly; then weighted Cauchy over all outer tuples permits bounding common zeros by one.

At fixed \(\fitinline{\gamma,f_{\rm new},k_{\rm new}}\), with at most \(\fitinline{K_0}\) original slots, split \(\fitinline{f_{\rm new}=JCd_2v'}\), \(\fitinline{J=sJ_2}\), reconstruct \(\fitinline{b_1b_2=q_0^2J_2^2s}\), and choose \(\fitinline{A_i,d_k}\). Multiplicity is at most

\hypertarget{mf:eq:B3.27}{}

\[\fitmath{
D(f_{\rm new})C(\gamma),\qquad
D(f)=d(f)^{9+4K_0},\quad
C(\gamma)=d(q_0)^{5+2K_0}d(t')^{2K_0}d(r_g)^{2K_0}.
\qquad\text{(MF3.27)}
}\]

Indeed \(\fitinline{d_k\mid f_{\rm new}\operatorname{rad}q_0}\), and extracted slot primes divide reconstructed factors, giving those divisor powers. The frequency is then determined; old \(\fitinline{f,h}\) are absent. By \protect\hyperlink{mf:eq:B3.19}{(MF3.19)}, distinct triples have weighted count

\hypertarget{mf:eq:B3.28}{}

\[\fitmath{
\ll Z^{C_c+11\eta/2+o(1)},\qquad
C_c=\ell+R/2-j+t+g-\theta.
\qquad\text{(MF3.28)}
}\]

Extract that count once and normalize its measure. The unnormalized squared child is \(\fitinline{Z^{F'}\mathcal E}\); a child bound \(\fitinline{Z^{F'+\epsilon_{\rm child}}}\) therefore contributes \(\fitinline{Z^{2F'+\epsilon_{\rm child}}}\). The exact balance

\hypertarget{mf:eq:B3.29}{}

\[\fitmath{
\varkappa_i+\lambda_c+C_c+2F_c=F
\qquad\text{(MF3.29)}
}\]

gives exponent \(\fitinline{F+40\eta+\tau+\pi+\epsilon_{\rm child}}\): profiles and triple count contribute \(\fitinline{9\eta/2+18\eta+11\eta/2=28\eta}\), clipping \(\fitinline{2\delta_N\le12\eta}\). Here \(\fitinline{\pi}\) combines freely chosen divisor, ray, dyadic and coefficient losses. Cauchy takes geometric means; neither exponent nor triple mass doubles.

\textbf{Closure.} Put \(\fitinline{D=\lceil(M_{\max}+2)/d\rceil+1}\). At depth \(\fitinline{h}\) use margin \(\fitinline{c_h=c_*-7h\eta}\) and cap \(\fitinline{M\le M_{\max}-hd}\). Choose

\hypertarget{mf:eq:B3.30}{}

\[\fitmath{
\eta\le\min\{d/16,c_*/(14D),\epsilon/(160D),c_*/1000\},
\qquad\tau,\pi\le\epsilon/(4D),\qquad
\tau\le c_*/1000,
\qquad\text{(MF3.30)}
}\]

with the reflected losses above. By \protect\hyperlink{mf:eq:B3.26}{(MF3.26)}, rows decrease by \(\fitinline{\ge d}\), margins stay \(\fitinline{\ge c_*/2}\), and \(\fitinline{F}\) remains bounded through depth \(\fitinline{D}\); nonnegative row lengths preclude a strict call there. Allocate losses within \(\fitinline{\pi}\), fix raw tails by \protect\hyperlink{mf:eq:B3.4}{(MF3.4)}, propagate finite smooth/height orders backwards, then take \(\fitinline{Z}\) large for window errors. Path loss is \(\fitinline{<\epsilon}\). Completion, principal counts and strict descent prove \protect\hyperlink{mf:eq:B3.7}{(MF3.7)}. \ensuremath{\square}

\hypertarget{the-marked-inverse-estimate-and-inverse-amplification}{}

\hypertarget{a.7.2.3-the-marked-inverse-estimate-and-inverse-amplification}{%
\subparagraph{A.7.2.3 The marked inverse estimate and inverse amplification}\label{a.7.2.3-the-marked-inverse-estimate-and-inverse-amplification}}

Let \(\fitinline{\varphi_u(n)=\nu(n)\chi_n(u)^{\varepsilon_\chi}}\), with one common sign \(\fitinline{\varepsilon_\chi\in\{\pm1\}}\), and put

\[\fitmath{
M_u(Z^r)=Z^{-r/2}\sum_n\mu(n)\varphi_u(n)W(q_n/Z^r),\qquad
Q_u=\prod_i\left(Z^{-z_i/2}\sum_{p\in\mathcal P_i}a_i(p)\varphi_u(p)\right).
}\]

Lists are fixed disjoint, \(\fitinline{|a_i(p)|\le1}\), \(\fitinline{q_p\asymp Z^{z_i}}\); individual annular weights are absorbed at fixed supremum cost. Retain natural zeros. Length ranges and slot count are bounded/fixed; \(\fitinline{c_1,c_2>0}\) are fixed.

\textbf{Theorem A.7.3} (Marked inverse moment). \emph{\hypertarget{mf:inverse}{}The full-width inverse estimate is \hypertarget{mf:eq:B3.31}{}}

\[\fitmath{
r+2z\le m-c_1,\qquad2r+8z\le3m-c_2
\quad\Longrightarrow\quad
\sum_{0<q_u\ll Z^m}|M_u(Z^r)Q_u|^2\ll Z^{m+\epsilon}.
\qquad\text{(MF3.31)}
}\]

\emph{The prime coefficients are arbitrary bounded row-independent coefficients, and no mesh is required.}

\emph{Proof.} Set \(\fitinline{P=\prod_i p_i}\), \(\fitinline{j=(n,P)}\) and fix its assigned slots. Write \(\fitinline{n=jn_0}\), \(\fitinline{P=jP_0}\), \(\fitinline{c=n_0P_0}\); \(\fitinline{c}\) is squarefree, \(\fitinline{(c,j)=1}\). The identities \(\fitinline{\mu(n_0)=\mu(c)\mu(P_0)}\) and \(\fitinline{\varphi_u(n)\varphi_u(P)=\varphi_u(j)^2\varphi_u(c)}\) retain zeros. For assigned length \(\fitinline{G}\), put

\hypertarget{mf:eq:B3.32}{}

\[\fitmath{
z_0=z-G,\qquad D_0=r+z-2G,
\qquad Z^{-(r+z)/2}=Z^{-G}Z^{-D_0/2}.
\qquad\text{(MF3.32)}
}\]

The \(\fitinline{O(Z^{G+o(1)})}\) assigned tuples have triangle mass canceled by \(\fitinline{Z^{-G}}\); \(\fitinline{\varphi_u(j)^2}\) contracts. For fixed \(\fitinline{j}\), the profile argument is \(\fitinline{(q_j/Z^G)(q_c/Z^{D_0})/\prod_{i\,\mathrm{surviving}}(q_{p_i}/Z^{z_i})}\). Separate it commonly before choosing survivors, absorb \(\fitinline{\mu(p_i)}\) into coefficients, and conjugate the whole polynomial if needed. It remains to bound

\hypertarget{mf:eq:B3.33}{}

\[\fitmath{
R_0(u)=Z^{-D_0/2}\sum_{c\ {\rm sf}}\mu(c)\nu_0(c)
 \mathbf1_{(c,j)=1}\overline{\chi_c(u)}d_0(c)W_0(q_c/Z^{D_0}).
\qquad\text{(MF3.33)}
}\]

The Fourier measure is uniform over the fixed \(\fitinline{q_j/Z^G}\) interval.

Square in a smooth row ball, write \(\fitinline{c_i=Cz_i}\), \(\fitinline{(z_1,z_2)=1}\), and apply \protect\hyperlink{mf:eq:B3.5}{(MF3.5)} with \(\fitinline{\overline{\chi_{z_1}}\chi_{z_2}}\), \(\fitinline{d'\mid C}\), frequency \(\fitinline{h'}\). For centers \(\fitinline{B,\theta}\) of \(\fitinline{C,d'}\) put \(\fitinline{P_1=B-\theta\ge-2\eta}\) and retain

\hypertarget{mf:eq:B3.34}{}

\[\fitmath{
M_c=2D_0-m-2P_1,\qquad M_{\rm init}=M_c+6\eta+\tau_{\rm init},
\qquad0<q_{d'h'}\le Z^{M_{\rm init}}.
\qquad\text{(MF3.34)}
}\]

Conductor error is \(\fitinline{4\eta}\) and two \(\fitinline{d'}\) appearances add \(\fitinline{2\eta}\), so the complement is an actual-ratio tail. Principal diagonal costs \(\fitinline{O(Z^{m+o(1)})}\), including bounded nonempty negative-\(\fitinline{D_0}\) windows.

As in the second canonical transform, Gauss conversion gives \(\fitinline{a_0(z)\nu_*(z)\overline{\chi_z(d')}\chi_z(h')}\), with common \(\fitinline{\nu_*}\) after expanding \(\fitinline{G([z_2z_1^{-1}])}\). Insert full M\"{o}bius \(\fitinline{z_a=sn_a}\), retaining \(\fitinline{(C,j)=(s,Cj)=(s,h')=1}\). Set

\hypertarget{mf:eq:B3.35}{}

\[\fitmath{
t=C/d',\qquad k=d'h',\qquad f=d's,\qquad
\rho_{t,j}(n)=\mathbf1_{(n,\operatorname{rad}(tj))=1}.
\qquad\text{(MF3.35)}
}\]

The zero-preserving identity

\hypertarget{mf:eq:B3.36}{}

\[\fitmath{
a_0(sn)\nu_*(sn)\overline{\chi_{sn}(d')}\chi_{sn}(h')
=a_0(s)\nu_*(s)\overline{\chi_s(d')}\chi_s(h')
 a_0(n)\nu_*(n)\chi_n(k)\chi_n(f)^4
\qquad\text{(MF3.36)}
}\]

holds for \(\fitinline{(s,n)=1}\) by \(\fitinline{\overline{\chi_n(d')}=\chi_n(d')^5}\), including zeros. The squarefree \(\fitinline{f}\) and priority allocation to \(\fitinline{(t,d',s;n)}\) preserve original slot subsets, with survival zeros supplied by \(\fitinline{\rho_{t,j}}\) and \(\fitinline{\chi_n(f)^4}\).

If \(\fitinline{v}\) centers \(\fitinline{s}\), the formal child parameters are

\hypertarget{mf:eq:B3.37}{}

\[\fitmath{
N_c=D_0-B-v,\quad V_c=\theta+v,\quad F_c=D_0-P_1,
\quad N'=\max(0,N_c),\quad0\le\delta_N=N'-N_c\le3\eta,
\quad F'=F_c+\delta_N.
\qquad\text{(MF3.37)}
}\]

For normalized \(\fitinline{(C,d',s,h',n_1,n_2)}\) use root \(\fitinline{y_d^{-1}y_s^{-1}(y_1y_2)^{-1/2}}\), windows \(\fitinline{W_0(y_Cy_sy_1)\overline{W_0(y_Cy_sy_2)}}\), kernel \(\fitinline{K_{\rm init}(A^{\rm ker}y_h/(y_dy_s^2y_1y_2))}\), \(\fitinline{A^{\rm ker}=Z^{m+H-\theta-2(D_0-B)}}\) (\(\fitinline{H}\) the bare-frequency center). Multiplying the profile by \(\fitinline{Z^{-3\eta}}\) leaves \(\fitinline{Z^{\varkappa_c+3\eta}}\), \(\fitinline{\varkappa_c=m-2D_0+P_1}\). Fresh tests are bounded clipped cutoffs times Fourier powers. Inversion restores the full kernel before Cauchy, retaining \protect\hyperlink{mf:eq:B3.34}{(MF3.34)}'s mask and all three zeros outside it.

At fixed \(\fitinline{t,f,k}\), choose \(\fitinline{d'\mid f}\); then \(\fitinline{s=f/d'}\), \(\fitinline{C=td'}\), \(\fitinline{h'=k/d'}\). With at most \(\fitinline{K_0}\) slots, multiplicity is \(\fitinline{\le d(f)^{1+2K_0}d(t)^{2K_0}}\). Project distinct \(\fitinline{t}\) once from the genuine outer set: \(\fitinline{\log_Zq_t\le P_1+2\eta}\), giving weighted count \(\fitinline{Z^{P_1+2\eta+o(1)}}\); extract it and normalize. The actual puncture has \(\fitinline{Q'\le P_1+G+3\eta}\). Then

\hypertarget{mf:eq:B3.38}{}

\[\fitmath{
\begin{aligned}
F'-M_{\rm init}-Q'-z_0&\ge c_1-9\eta-\tau_{\rm init},\\
4F'-3M_{\rm init}-6z_0&\ge c_2-22\eta-3\tau_{\rm init},\\
\varkappa_c+P_1+2F_c&=m.
\end{aligned}\qquad\text{(MF3.38)}
}\]

The formal differences are \(\fitinline{m-r-2z+2G}\) and \(\fitinline{3m-2r-8z+10G+2P_1}\). Choose \(\fitinline{\eta,\tau_{\rm init}}\) small and apply canonical margin \(\fitinline{\min(c_1,c_2)/2}\). The last identity and \(\fitinline{2\delta_N\le6\eta}\) give exponent \(\fitinline{m+11\eta+\pi_{\rm init}+\epsilon_{\rm canonical}}\). Allocate these within \(\fitinline{\epsilon}\), then choose tails and backward seminorm orders as above. This proves \protect\hyperlink{mf:eq:B3.31}{(MF3.31)}, hence MI at \(\fitinline{m=1}\). \ensuremath{\square}

\textbf{Unmarked consequence (UM).} With no slots, \(\fitinline{Z=H}\), \(\fitinline{m=1}\), \(\fitinline{r=\log_HD}\), the condition \(\fitinline{H\ge D^{1+\vartheta}}\) gives \(\fitinline{1-r\ge\vartheta/(1+\vartheta)}\) and \(\fitinline{3-2r\ge1}\); count bounded scales. Thus

\hypertarget{mf:eq:B3.39}{}

\textbf{(MF3.39)} is \protect\hyperlink{cf:eq:3.2}{(95)}, with \(\fitinline{\varphi_u=\nu\chi_\bullet(u)^{\varepsilon_\chi}}\), both signs and every element row, under \(\fitinline{H\ge D^{1+\vartheta}}\).

All original zeros, finite-seminorm and polynomial-height control remain.

\hypertarget{the-enlarged-plain-moment-family}{}

\hypertarget{a.7.2.4-the-enlarged-plain-moment-family}{%
\subparagraph{A.7.2.4 The enlarged plain-moment family}\label{a.7.2.4-the-enlarged-plain-moment-family}}

Use \(\fitinline{\psi_k(n)=\tau(n)\chi_n(k)}\) on \(\fitinline{0<q_k\ll Z^m}\), where \(\fitinline{\tau}\) combines a fixed ray character and displayed moving symbols. Their union \textbf{before} cancellations/exponent reduction has norm \(\fitinline{\le Z^q}\); width is \(\fitinline{M=m+q}\). A squarefree common mask \(\fitinline{q_R\le Z^B}\) may depend on frozen labels but is fixed over the row sum and applies to every plain and prime factor. Canceled row-character primes retain natural zeros; they do not become freely chosen row masks.

For positive slots fix nonzero integral ray modulus \(\fitinline{\mathfrak M}\), a subgroup containing global units, common annular profile \(\fitinline{W}\), anchor \(\fitinline{\eta_0}\), and ideal \(\fitinline{Q\subseteq\mathfrak M}\). Require \(\fitinline{Q_0=Q\cap\mathfrak q_{\eta_0}}\) to be nonzero, proper and contained in \(\fitinline{(72)}\). Each slot has \(\fitinline{\nu_i=\theta_i\overline{\eta_0}}\) for a fixed-ray-quotient character \(\fitinline{\theta_i}\), bounded real power \(\fitinline{\sigma_i}\) and height \(\fitinline{v_i}\):

\[\fitmath{
Q_{k,i}=Z^{-z_i/2}\sum_{p\in\mathcal P_i^{\rm can}}
 \psi_k(p)\nu_i(p)W(q_p/Z^{z_i})
 (q_p/Z^{z_i})^{-\sigma_i+iv_i}.
}\]

Here \(\fitinline{\mathcal P_i^{\rm can}}\) is the genuine canonical good-prime pool at \(\fitinline{Z^{z_i}}\), retaining natural zeros and common puncture. MP uses this pool and common \(\fitinline{W}\); it asserts neither arbitrary subsets, separate slot profiles nor a finite-character-expansion extension. MI permits arbitrary subsets. Physical scales are separated annuli.

Put \(\fitinline{z=\sum_i z_i}\), \(\fitinline{A=n_1+n_2+z}\). The two centrally normalized plains have independent profiles and heights with polynomial costs.

\textbf{Theorem A.7.4} (Natural marked plain moment). \emph{\hypertarget{mf:plain}{} The zero-slot assertion is the unconditional Theorem \protect\hyperlink{cf:mixed-plain-moment}{A.5.1}. For positive slots require \(\fitinline{0<z_i\le\eta_{\rm mesh}}\), with mesh chosen below before the slot system. Absorb zero-width bounded factors into fixed arithmetic data. Assume}

\[\fitmath{
7/10\le\kappa\le1,\qquad \beta\ge51/100,\qquad
2\beta-1\le\kappa.
}\]

\emph{The centrally normalized plain sums are}

\[\fitmath{
S_{k,i}(X)=X^{-1/2}\sum_{n}\psi_k(n)\mathbf1_{(n,R)=1}
 W_i(q_n/X),
}\]

\emph{with all integral ideals \(\fitinline{n}\) and the fixed bad-prime zeros understood. Then \hypertarget{mf:eq:B3.40}{}}

\[\fitmath{
\begin{array}{ll}
z=0:&\displaystyle
 \sum_{\psi_k^*\ne1}|S_k(Z^{n_1})S_k(Z^{n_2})|^2\ll Z^{M+\epsilon}
 \quad\text{for all bounded }n_1,n_2\ge0;\\[2mm]
z>0:&\displaystyle
 n_1+n_2+6\kappa z\le M
 \quad\Longrightarrow\quad
 \sum_{k\ {\rm eligible}}|S_k(Z^{n_1})S_k(Z^{n_2})\prod_iQ_{k,i}|^2
 \ll Z^{M+\epsilon}.
\end{array}\qquad\text{(MF3.40)}
}\]

\emph{A natural state has a whole-product presentation, complete natural zeros, fixed nonzero puncture and nonnegative smooth radial measure. Its width counts row-ball exponent plus displayed moving modulus before cancellation. Keep exactly nonzero rows not induced from \(\fitinline{\mathfrak Q_0=Q\cap\mathfrak q_{\eta_0}}\): fixed-inducing means a primitive finite Hecke character of conductor dividing \(\fitinline{\mathfrak Q_0}\) induces the actual row function with natural zeros. Deleting an extra puncture preserves this predicate, since additional zero Euler factors do not change the primitive inducing character. Initially,}

\[\fitmath{
\operatorname{initialKeep}(\eta,\mathfrak Q_0)
=\{k\in\mathcal O:k\ne0,\ \text{the natural presentation of }n\mapsto\eta(n)\chi_n(k)
\text{ is not induced from }\mathfrak Q_0\}.
}\]

\emph{The presentation retains fixed bad-prime zeros, primitive-conductor zeros, redundant moving zeros and row radical. Eligibility excludes every fixed-inducing presentation, not just principal characters or a chosen list.}

\emph{Quantifier order: for a requested loss and bounded width/puncture/length ranges choose one mesh, uniformly for \(\fitinline{\kappa\in[7/10,1]}\), before arithmetic family, slot count, profiles or target height. After fixing a finite slot index type and common \(\fitinline{W}\), choose finite seminorms and height order before fixed character/ideal; their constant and eventual threshold precede live subset, slot characters, lengths, powers, heights, radial state and plain profiles. Positive slots allow a common whole-product norm power; independent plain twists are controlled by seminorms. Zero slots retain independent heights, all bounded lengths and every nonprincipal inducing character, including fixed-conductor rows.}

\emph{Absorb fixed conductors in arbitrarily small width padding charged to the requested loss before the eventual threshold. Moving support remains width, never a profile seminorm.}

\emph{The physical initial state has scale \(\fitinline{U}\), puncture one, fixed source \(\fitinline{\eta}\) and width \(\fitinline{1+\delta}\) once \(\fitinline{q_{\mathfrak q_\eta}\le U^\delta}\). Its measure is \(\fitinline{1_{\operatorname{initialKeep}(\eta,\mathfrak Q_0)}(k)\Phi(q_k/U)}\), \(\fitinline{\Re\Phi\ge0}\), with the literal natural plains and slots above, and relative character}

\[\fitmath{
\nu_i(n)=\theta_i(n)\overline{\eta_0(n)},
}\]

\emph{Lift \(\fitinline{\theta_i}\) from \(\fitinline{(\mathcal O/\mathfrak M)^\times/H}\); both factors retain their zeros. The prime pool is the identity ray class, cut at \(\fitinline{\max(1,b)U^{z_i}}\); \(\fitinline{W}\) selects its exact annulus. Physically \(\fitinline{W=\overline{F.W}}\), \(\fitinline{\theta_i=1}\), real power \(\fitinline{33/50}\), height \(\fitinline{-\Im z_i}\). Fixed-inducing sixth-power-free rows have prescribed good-prime valuations modulo six, hence bounded norm and leave a growing lower annulus.}

\emph{The zero-slot endpoint restores fixed nonprincipal characters by masked lattice cancellation below. Width counts full moving support but excludes the puncture. For exceptional children fix a finite conjugation-stable \(\fitinline{S}\)-supported group \(\fitinline{\Theta}\) containing fixed phases, inverses and \(\fitinline{n\mapsto\chi_n(-1)}\); it does not replace positive-slot eligibility. Taking \(\fitinline{m=1,q=0,n_1=n_2}\) gives MP in the positive family and MP0 for all nonprincipal zero-slot rows.}

\emph{Proof.} \textbf{Extra masks.} A natural presentation has only primitive, redundant displayed, row and \(\fitinline{S}\) zeros. Delete its additional fixed \(\fitinline{R}\) mask to obtain \(\fitinline{\psi_k^0}\). Multiplicativity gives

\hypertarget{mf:eq:B3.41}{}

\[\fitmath{
S_{k,R}(X)=\sum_{d\mid R}\frac{\mu(d)\psi_k^0(d)}{\sqrt{q_d}}
 S_k^0(X/q_d),\qquad
Q_{k,i,R}=Q_{k,i}^0-Z^{-z_i/2}\sum_{p\mid R}\psi_k^0(p)\nu_i(p)W(q_p/Z^{z_i})(q_p/Z^{z_i})^{-\sigma_i+iv_i}.
\qquad\text{(MF3.41)}
}\]

Total coefficient mass is \(\fitinline{Z^{o(1)}}\): \(\fitinline{\prod_{p\mid R}(1-q_p^{-1/2})^{-1}\ll_\epsilon Z^\epsilon}\), and a frozen slot has \(\fitinline{Z^{-z_i/2}\ll q_p^{-1/2}}\). Apply both identities and Minkowski. The affine length \(\fitinline{A+(6\kappa-1)z}\) decreases; if slots vanish the zero-slot row set contains the original eligible rows. Nonempty subunit plains clip by bounded dilation without exceeding old nonnegative lengths. Prove natural states uniformly, retaining this deletion at every child.

\textbf{Whole-factor reflection.} Use \S{}\protect\hyperlink{cf:sec:D.3}{A.5.3}'s natural reflection, with primitive \(\fitinline{\psi_k^*}\), conductor norm \(\fitinline{Q_k}\), \(\fitinline{C_k=3Q_k/(2\pi)^2}\) and redundant radical \(\fitinline{R_{0,k}}\).

\hypertarget{mf:eq:B3.42}{}

\textbf{(MF3.42)} is \protect\hyperlink{cf:eq:D.6}{(70)}, \(\fitinline{C_kq_{R_{0,k}}\ll Z^M}\), with the same full-support width.

Unit and \(\fitinline{S}\)-valuations modulo six give finite ray sectors. For nonprincipal inducing characters, the central primitive reflection has scale \(\fitinline{C_k/X}\), unit root number and profile

\hypertarget{mf:eq:B3.43}{}

\textbf{(MF3.43)} is \protect\hyperlink{cf:eq:D.7}{(71)}'s Mellin profile formula, \(\fitinline{\widehat{W^\sharp}(s)=\widehat W(1-s)\Gamma(s)/\Gamma(1-s)}\).

The proof in \S{}\protect\hyperlink{cf:sec:D.3}{A.5.3} substitutes the functional equation, sets \(\fitinline{s\mapsto1-s}\) and shifts without an \(\fitinline{L}\)-pole. It gives bounded Euler derivatives at zero, rapid decay at infinity, and finite-seminorm/polynomial-height control by Stirling and \(\fitinline{\widehat{Wy^{it}}(s)=\widehat W(s+it)}\).

Deleting redundant factors before reflection and restoring geometrically gives exactly

\hypertarget{mf:eq:B3.44}{}

\[\fitmath{
S_k^0(X;W)=\varepsilon_k
 \sum_{d_0\mid R_{0,k}}\sum_{\operatorname{rad}(h_0)\mid R_{0,k}}
 \frac{\mu(d_0)\psi_k^*(d_0)\overline{\psi_k^*(h_0)}}{\sqrt{q_{d_0}q_{h_0}}}
 S_{\overline{\psi_k^0}}\!\left(\frac{C_kq_{d_0}}{Xq_{h_0}};W^\sharp\right).
\qquad\text{(MF3.44)}
}\]

Its mass is \(\fitinline{Z^{o(1)}}\). Lower annuli of \(\fitinline{W^\sharp}\) have central coefficients \(\fitinline{O(2^{-j/2})}\); upper annuli decay arbitrarily rapidly. Discard above relative scale \(\fitinline{Z^{\xi/2}}\). A main scale \(\fitinline{Y<Z^{-\xi}}\) contributes \(\fitinline{O_B(Y^{B-1/2})}\) by counting and decay, for arbitrary fixed \(\fitinline{B}\). Handle row-dependent coefficients by mass, then scale suprema: integrate function and log-scale derivative over \(\fitinline{O(\log Z)}\) unit boxes, in two variables for two scales. The derivative is \(\fitinline{-W/2-yW'}\), so common-profile bounds apply. Conjugate the entire reflected factor to restore orientation. By \protect\hyperlink{mf:eq:B3.42}{(MF3.42)}, retained scales satisfy

\hypertarget{mf:eq:B3.45}{}

\[\fitmath{
n_{\rm ref}\le M-n+\xi.
\qquad\text{(MF3.45)}
}\]

Absorb fixed endpoints and boxes in the final threshold; negative right side gives an empty range. Reflecting plains longer than \(\fitinline{M/2}\) reduces zero slots to

\hypertarget{mf:eq:B3.46}{}

\textbf{(MF3.46)} is the padded core \protect\hyperlink{cf:eq:D.8}{(72)}, with \(\fitinline{2\xi\le\delta}\): \(\fitinline{n_i\le M/2+\xi}\) and \(\fitinline{A\le M+2\xi\le M+\delta}\).

\textbf{Prime cost.} On eligible rows each slot remains nonprincipal after expanding the canonical identity-ray projector as the average of all fixed-quotient characters. Its product with \(\fitinline{\nu_i=\theta_i\overline{\eta_0}}\) has conductor dividing \(\fitinline{Q_0}\); a principal resulting row would therefore be induced from \(\fitinline{Q_0}\). Multiplication by such a fixed character preserves eligibility by inverse multiplication. The finite average changes only constants; no arbitrary-pool/profile extension follows. Under \(\fitinline{2\beta-1\le\kappa}\), move the smooth von Mangoldt contour to \(\fitinline{s_\kappa+e}\), \(\fitinline{s_\kappa=(1+\kappa)/2}\). There \(\fitinline{|L'/L|\ll_e\log(Q(3+|t|)^2)}\): on the zero-free disk centered at \(\fitinline{2+it}\) of radius \(\fitinline{2-s_\kappa-e/4}\), the Euler branch \(\fitinline{g=\log L}\) has \(\fitinline{g(2+it)=O(1)}\) and \(\fitinline{\Re g\ll\log(Q(3+|t|)^2)}\) by functional-equation strip growth. Borel--Carath\'{e}odory and then Cauchy's derivative bound on successively smaller disks containing \(\fitinline{s_\kappa+e+it}\) prove the claim; all gaps are fixed multiples of \(\fitinline{e}\). Prime powers of degree \(\fitinline{\ge2}\) cost \(\fitinline{O(P^{1/2+\epsilon})}\). Divide the annular weight by \(\fitinline{\log(Py)}\); redundant zeros delete only \(\fitinline{O(\log Z)}\) primes. Central normalization over the fixed slots gives

\hypertarget{mf:eq:B3.47}{}

\[\fitmath{
\left|\prod_iQ_{k,i}\right|^2
 \ll Z^{\kappa z+\epsilon_1}\,p_J(W)^b\langle\mathbf t\rangle^h.
\qquad\text{(MF3.47)}
}\]

Contour displacement costs \(\fitinline{2ez}\), depending on total length, not slot count. Heights translate rapidly decreasing Mellin transforms; required derivative/height orders do not affect the exponent.

\hypertarget{induction-order-and-the-centered-comparison}{}

\hypertarget{a.7.2.5-induction-order-and-the-centered-comparison}{%
\subparagraph{A.7.2.5 Induction order and the centered comparison}\label{a.7.2.5-induction-order-and-the-centered-comparison}}

Choose floor \(\fitinline{\rho>0}\), step \(\fitinline{\sigma>0}\), padding \(\fitinline{\delta\ll\min(\rho,\sigma)}\) and bands of length \(\fitinline{\sigma/4}\). Use the zero-slot theorem first, then prove positive bands increasingly; within each band prove \(\fitinline{A\le5M/6}\) before centering. The zero-slot proof in Appendix A.5 supplies all lengths by reflection and padded-floor excess \(\fitinline{\le\rho+\delta}\). At \(\fitinline{M\le\rho}\), positive slots follow from it and \protect\hyperlink{mf:eq:B3.47}{(MF3.47)}, since

\hypertarget{mf:eq:B3.48}{}

\[\fitmath{
z\le M/(6\kappa),\qquad\kappa z\le M/6.
\qquad\text{(MF3.48)}
}\]

Its additional terminal excess is at most \(\fitinline{\rho/6}\).

For \(\fitinline{A>5M/6}\), put \(\fitinline{L=M/4}\). If one plain is shorter than \(\fitinline{L}\), reflect the other; otherwise compare lengths \(\fitinline{L,A-z-L}\) with the same profiles, reflecting the larger comparison factor only to estimate the comparison. Since \(\fitinline{q_\kappa=6\kappa-1>2}\) and \(\fitinline{z<(M/6)/q_\kappa<M/12}\), both comparison lengths are \(\fitinline{\ge L}\). Reflection gives

\hypertarget{mf:eq:B3.49}{}

\[\fitmath{
A_{\rm comp}\le3M/2-A+2z+\xi,
\quad
\min\{5M/6-A_{\rm comp},\,M-A_{\rm comp}-q_\kappa z\}
\ge \mathfrak m(7/10)M-\xi,
\quad
\mathfrak m(\kappa)=\frac{2\kappa-1}{2(6\kappa-1)}.
\qquad\text{(MF3.49)}
}\]

Indeed \(\fitinline{t=M-A<M/6}\), \(\fitinline{z\le t/q_\kappa}\) give margins \(\fitinline{\ge M/3-t(1+2/q_\kappa)-\xi}\) and \(\fitinline{M/2-t(2+2/q_\kappa)-\xi}\). Zero-slot comparison margin is \(\fitinline{\ge M/6-\xi}\). Choose \(\fitinline{\xi<\rho\mathfrak m(7/10)/2}\); each comparison then uses its already proved uncentered stage.

For both plains \(\fitinline{\ge L}\), put \(\fitinline{X_i=Z^{n_i}}\), \(\fitinline{Y_1=Z^L}\), \(\fitinline{Y_2=X_1X_2/Y_1}\). Subtract and retain the equal-product rectangles; after extracting \(\fitinline{b_i}\), the coefficient is

\hypertarget{mf:eq:B3.50}{}

\textbf{(MF3.50)} is the exact rectangle difference \protect\hyperlink{cf:eq:D.9}{(73)}, with \(\fitinline{D_b=\mathcal D_b}\) and \(\fitinline{X_1X_2=Y_1Y_2}\).

Use the same \(\fitinline{b_i}\) even when one rectangle vanishes. The comparison is already bounded on eligible rows; enlarge the positive norm of the \textbf{whole difference} to a smooth ball of all rows. In the uncentered range enlarge the original squared norm. Never Poisson-sum an eligibility indicator.

The invariant coefficient carried through the transformations is

\hypertarget{mf:eq:B3.51}{}

\[\fitmath{
\prod_{i\in I'}\nu_i(p_i)W(q_{p_i}/Z^{z_i})(q_{p_i}/Z^{z_i})^{-\sigma_i+iv_i}\,
D_b(l_1,l_2)\,\tau_1(u)q_u^{it}\,\mathbf1_{(u,R)=1}\mathbf1_{s\mid u},
\qquad u=\Bigl(\prod_{i\in I'}p_i\Bigr)l_1l_2,
\qquad\text{(MF3.51)}
}\]

Here \(\fitinline{s}\) is fixed squarefree; uncentered inputs omit the second rectangle. The character, mask and norm power act on the entire product, preserving collisions by zero extension.

Use one aggregate tolerance. For fixed slot count \(\fitinline{N}\), let \(\fitinline{H_N}\) bound sums of absolute logarithmic endpoints of every slot/plain window and all finite later enlargements; set \(\fitinline{\theta_N=H_N/\log Z}\). Every slot subset's total ratio error is \(\fitinline{\le\theta_N}\). Extract plain powers using exact norms. Stage errors are \(\fitinline{O(\theta_N)}\) with numerical coefficient independent of \(\fitinline{N}\), not a loss per slot.

\hypertarget{first-poisson-from-plain-products-to-gauss-polynomials}{}

\hypertarget{a.7.2.6-first-poisson-from-plain-products-to-gauss-polynomials}{%
\subparagraph{A.7.2.6 First Poisson: from plain products to Gauss polynomials}\label{a.7.2.6-first-poisson-from-plain-products-to-gauss-polynomials}}

Set \(\fitinline{H=Z^m}\), \(\fitinline{X=Z^A}\). The zero-frequency argument of \S{}\protect\hyperlink{cf:sec:D.4}{A.5.4} applies unchanged to full products: zero net sextic exponents exclude multiplicity-one primes, giving a powerful ideal of norm \(\fitinline{O(X^2)}\), \(\fitinline{O(X^{1+\epsilon})}\) choices and divisor-bounded slot/plain allocations. Multiplying central \(\fitinline{X^{-1}}\) and row mean \(\fitinline{\le H}\) gives \(\fitinline{O(Z^{m+o(1)})}\).

For nonzero frequency write full products \(\fitinline{Ca,Db}\) with complete common support, \(\fitinline{(a,b)=(ab,CD)=1}\). Fix each common prime's full plain valuations and slot occurrences, extract them, freeze supplying slots and puncture every residual factor. Apply the same allocation to both rectangles; retain any scalar zero before enlargement.

Let \(\fitinline{c,d,p}\) be lengths of \(\fitinline{C,D,\operatorname{rad}(CD)}\). Let \(\fitinline{r}\) contain nonzero net exponents modulo six, length \(\fitinline{R}\), primitive character \(\fitinline{\xi_r}\); remaining common-mask divisor \(\fitinline{e}\) has \(\fitinline{E\le p-R}\). For allocated coefficients without the whole-product twist, set \(\fitinline{\tau_C(a)=\tau(a)\chi_a(er)\overline{\xi_r(a)}}\) and similarly for \(\fitinline{\tau_D}\). Poisson/CRT give, up to bounded frozen scalar,

\hypertarget{mf:eq:B3.52}{}

\[\fitmath{
\frac{H}{Xq_e\sqrt{q_r}}\sum_{h\ne0}\mathcal G_{\xi_r}(r,h)
\sum_{(a,b)=1}\frac{A_C(a)\overline{A_D(b)}}{\sqrt{q_aq_b}}
 \tau_C(a)\overline{\tau_D(b)}R(a,b)
 \mathcal G(a,h)\overline{\mathcal G(b,-h)}
 K_1\!\left(\frac{Hq_h}{q_eq_rq_aq_b}\right).
\qquad\text{(MF3.52)}
}\]

Here \(\fitinline{\mathcal G}\) is the normalized sum with local values \protect\hyperlink{cf:eq:A.12}{(8)}, \(\fitinline{|\mathcal G_{\xi_r}|\le1}\). As in \S{}\protect\hyperlink{cf:sec:D.4}{A.5.4}, Poisson after \(\fitinline{k=ek'}\) contributes \(\fitinline{H/(q_eq_rq_aq_b)}\); Gauss sums restore \(\fitinline{\sqrt{q_rq_aq_b}}\). CRT gives \(\fitinline{\chi_a(br)\overline{\chi_b(ar)}\overline{\xi_r(ab)}}\), whose \(\fitinline{a,b}\) ratio is \(\fitinline{R(a,b)}\) by reciprocity. New moving primes puncture every residual factor even if exponents cancel; their support and frequency lengths are

\hypertarget{mf:eq:B3.53}{}

\textbf{(MF3.53)} combines \protect\hyperlink{cf:eq:D.11}{(75)} and \protect\hyperlink{cf:eq:D.12}{(76)}: \(\fitinline{\widetilde q=q+R+E}\), \(\fitinline{K_0=2A-c-d+R+E-m}\), with this chapter's full-product length \(\fitinline{A=n_1+n_2+z}\).

The first kernel's nominal frequency scale is \(\fitinline{Z^{K_0}}\).

For both Poisson transforms use one nonnegative norm-dyadic partition. In an outer sector retain whole weights meeting \(\fitinline{q_h\le T_{\rm sec}Z^{\xi/2}}\), \(\fitinline{T_{\rm sec}}\) the maximum nominal scale. For large \(\fitinline{Z}\), upper length is \(\fitinline{K\le K_0+\delta_1}\), \(\fitinline{0\le\delta_1<\xi}\); discarded genuine arguments are \(\fitinline{\ge Z^{\xi/4}}\), since full-product support costs \(\fitinline{e^{2H_N}}\). Apply \protect\hyperlink{mf:eq:B3.4}{(MF3.4)} with raw bounds \(\fitinline{|\mathcal G|\le\sqrt{q_a}}\), \(\fitinline{|F|\le q_uq_v}\) and polynomial column/count cost, fixing that cost before decay order. Retained weights/full kernels are common to rectangles and independent of live columns.

On each dyad extend the separated Gauss bridge off coprimality using formal norms, original masks and bicharacter. Insert the complete \(\fitinline{\sum_{s\mid(a,b)}\mu(s)}\) before absolute values; let \(\fitinline{s_0=\log_Zq_s}\). Separate kernel/inverse roots in row and whole-product norms, adding an aggregate outer-ratio coordinate if needed. Each column receives one common norm power across its plains and rectangles. Then Cauchy and \(\fitinline{|\mathcal G_{\xi_r}|\le1}\) give

\hypertarget{mf:eq:B3.54}{}

\textbf{(MF3.54)} is \protect\hyperlink{cf:eq:D.13}{(77)} with \(\fitinline{\mathfrak g=\mathcal G}\), this chapter's \(\fitinline{A_C,\tau_C'}\) and dyadic \(\fitinline{h\ne0}\) support. Its complete coefficient is \protect\hyperlink{mf:eq:B3.51}{(MF3.51)}.

The exchanged \(\fitinline{D}\) norm has the same form. Outside exponent is \(\fitinline{m-A-R/2-E}\); common-support and M\"{o}bius counts are \(\fitinline{Z^{p+s_0+o(1)}}\). Count radicals first, then bounded powers via \(\fitinline{\#\{v:\operatorname{rad}v\mid c,q_v\le Z^T\}\le Z^{aT}\prod_{p\mid c}(1-q_p^{-a})^{-1}}\) for small \(\fitinline{a>0}\); allocations cost divisors.

It suffices to prove

\hypertarget{mf:eq:B3.55}{}

\textbf{(MF3.55)} is the target bound \protect\hyperlink{cf:eq:D.14}{(78)}, here with loss \(\fitinline{\epsilon_G}\): \(\fitinline{\mathcal N_C\ll Z^{A-c+\widetilde q+B_c-s_0+\epsilon_G}}\), \(\fitinline{B_c=(3c-5d-R)_+/6}\).

Its exchanged counterpart and outside factors have exponent \(\fitinline{M+\{B_c+B_d-(c+d-2p-R)\}/2}\). The local comparison is

\hypertarget{mf:eq:B3.56}{}

\[\fitmath{
B_c+B_d\le c+d-2p-R.
\qquad\text{(MF3.56)}
}\]

This is the primewise inequality proved after \protect\hyperlink{cf:eq:D.15}{(79)}, with identical multiplicities and \(\fitinline{r_p=1_{6\nmid i-j}}\); slot occurrences are already included in those multiplicities.

\hypertarget{enlarging-a-gauss-row-norm}{}

\hypertarget{a.7.2.7-enlarging-a-gauss-row-norm}{%
\subparagraph{A.7.2.7 Enlarging a Gauss-row norm}\label{a.7.2.7-enlarging-a-gauss-row-norm}}

Extract length \(\fitinline{w}\) from \protect\hyperlink{mf:eq:B3.54}{(MF3.54)}, adding moving radical \(\fitinline{w_o\le w}\) with squared scalar \(\fitinline{Z^{-w_o}}\). Put \(\fitinline{a_0=A-c-w}\); \protect\hyperlink{mf:eq:B3.55}{(MF3.55)} permits the unweighted norm exponent

\hypertarget{mf:eq:B3.57}{}

\[\fitmath{
\Lambda_c=a_0+\widetilde q+w_o+w+B_c-s_0+\epsilon_G.
\qquad\text{(MF3.57)}
}\]

Define \(\fitinline{J=d-c+(K_0-K)-2w+w_o}\), \(\fitinline{J_+=\max(J,0)}\). With zero slots take \(\fitinline{w=w_o=0}\), \(\fitinline{\ell=0}\), and enlarge the row norm to a smooth ball of length \(\fitinline{K+g}\), \(\fitinline{g=J_++\sigma}\).

For positive slots choose \(\fitinline{\ell_* =\sigma/3}\), mesh \(\fitinline{<\sigma/6}\), and pool \(\fitinline{\mathcal P=\{p:Z^{\ell_*}/2<q_p\le Z^{\ell_*},p\notin S\}}\). The fixed-field prime ideal theorem gives \(\fitinline{|\mathcal P|=Z^{\ell_*+o(1)}}\). For large \(\fitinline{Z}\) it avoids all live slots. For each \(\fitinline{h\ne0}\), delete the \(\fitinline{O(\log Z)}\) primes dividing \(\fitinline{h,s}\) or frozen supports.

For \(\fitinline{p\nmid uh}\), the prime-power Gauss identities \protect\hyperlink{cf:eq:A.12}{(8)} and CRT give

\hypertarget{mf:eq:B3.58}{}

\[\fitmath{
\mathcal G(p^iu,h)=R(p,u)^i\chi_p(u)^{2i}\mathcal G(p^i,h)\mathcal G(u,h),
\qquad\mathcal G(u,hp^6)=\mathcal G(u,h).
\qquad\text{(MF3.58)}
}\]

Replacing \(\fitinline{h}\) by \(\fitinline{hp^6}\) changes only valuations \(\fitinline{i=1,6,7}\). Assign powers to the two plains (pool and slots are disjoint), puncturing the residual product and updating both rectangles. Jensen over eligible primes and these finite allocations bounds the original square by amplified square and extracted errors. Each output \(\fitinline{hp^6}\) has boundedly many representations, since pool length is fixed positive and output length bounded; density is \(\fitinline{Z^{-\ell_*+o(1)}}\). Take \(\fitinline{\ell=\ell_*}\), \(\fitinline{g=J_++2\sigma}\), enclosing amplified rows since \(\fitinline{6\ell_*=2\sigma}\).

The extracted alternatives are exactly:

\begin{longtable}[]{@{}
  >{\raggedright\arraybackslash}p{(\columnwidth - 6\tabcolsep) * \real{0.2933}}
  >{\raggedright\arraybackslash}p{(\columnwidth - 6\tabcolsep) * \real{0.4267}}
  >{\raggedright\arraybackslash}p{(\columnwidth - 6\tabcolsep) * \real{0.1467}}
  >{\raggedright\arraybackslash}p{(\columnwidth - 6\tabcolsep) * \real{0.1333}}@{}}
\toprule\noalign{}
\begin{minipage}[b]{\linewidth}\raggedright
Column valuation \(\fitinline{i}\)
\end{minipage} & \begin{minipage}[b]{\linewidth}\raggedright
Squared extraction coefficient
\end{minipage} & \begin{minipage}[b]{\linewidth}\raggedright
\(\fitinline{w}\)
\end{minipage} & \begin{minipage}[b]{\linewidth}\raggedright
\(\fitinline{w_o}\)
\end{minipage} \\
\midrule\noalign{}
\endhead
\bottomrule\noalign{}
\endlastfoot
1 & \(\fitinline{q_p^{-1}}\) & \(\fitinline{\ell_p}\) & \(\fitinline{\ell_p}\) \\
6 & \(\fitinline{(1-q_p^{-1})^2}\) & \(\fitinline{6\ell_p}\) & 0 \\
7 & \(\fitinline{q_p^{-1}}\) & \(\fitinline{7\ell_p}\) & \(\fitinline{\ell_p}\) \\
\end{longtable}

Here \(\fitinline{\ell_p=\log_Zq_p\in[\sigma/6,\sigma/3]}\). At valuation six, \(\fitinline{\chi_p(u)^{12}}\) is the existing common puncture; retain \(\fitinline{R(p,u)^6}\). Valuations one/seven retain a moving character counted in \(\fitinline{w_o}\). Freeze each error's \(\fitinline{p}\), take \(\fitinline{\ell=0}\), \(\fitinline{g=J_++\sigma}\), and transform without re-amplifying. The normalized prime average has no extra count. Always

\hypertarget{mf:eq:B3.59}{}

\[\fitmath{
0\le w_o\le w\le7\sigma/3.
\qquad\text{(MF3.59)}
}\]

Coefficient \protect\hyperlink{mf:eq:B3.51}{(MF3.51)}, multiplied by \(\fitinline{\mathcal G(u,h)}\), retains its difference, puncture and \(\fitinline{s}\)-divisibility.

Add zero only at the final smooth ball, never amplify it. By \protect\hyperlink{cf:eq:A.12}{(8)}, \(\fitinline{\mathcal G(u,0)}\) vanishes unless \(\fitinline{u}\) is sixth-power, and is \(\fitinline{\le\sqrt{q_u}}\). For nonempty upper bound \(\fitinline{Y=e^{H_N}Z^{a_0}}\) and \(\fitinline{s\mid u}\), \(\fitinline{q_s^6\le Y}\). Counting \(\fitinline{O(Y^{1/6})}\) columns gives normalized square \(\fitinline{Z^{-a_0}Y^{4/3+o(1)}\le Z^{a_0-s_0+2\theta_N+o(1)}}\), within \protect\hyperlink{mf:eq:B3.57}{(MF3.57)}.

\hypertarget{second-transform-and-the-actual-children}{}

\hypertarget{a.7.2.8-second-transform-and-the-actual-children}{%
\subparagraph{A.7.2.8 Second transform and the actual children}\label{a.7.2.8-second-transform-and-the-actual-children}}

Restore live slot sums before squaring. If \(\fitinline{B(u)}\) is the complete convolution coefficient, with \(\fitinline{s\mid u}\), \protect\hyperlink{cf:eq:A.14}{(10)} gives

\hypertarget{mf:eq:B3.60}{}

\[\fitmath{
\begin{aligned}
&Z^{-\ell}\sum_h\Phi_2(q_h/Z^{K+g})
 \left|Z^{-a_0/2}\sum_u B(u)\tau_1(u)q_u^{it}\mathcal G(u,h)\right|^2
\\
&\quad=Z^{K+g-\ell-a_0}\sum_{u,v}
 \frac{B(u)\overline{B(v)}\tau_1(u)\overline{\tau_1(v)}q_u^{it}q_v^{-it}}
      {\sqrt{q_uq_v}}
\\
&\qquad{}\times\sum_jF(u,v;j)K_2\!\left(\frac{Z^{K+g}q_j}{q_uq_v}\right).
\end{aligned}
\qquad\text{(MF3.60)}
}\]

At \(\fitinline{j=0}\), \protect\hyperlink{cf:eq:A.15}{(11)} forces \(\fitinline{u=v}\) and \(\fitinline{F(u,u;0)=\phi(u)\le q_u}\). There are \(\fitinline{O(Z^{a_0-s_0+\theta_N})}\) supported multiples of \(\fitinline{s}\), giving exponent \(\fitinline{K+g-\ell-s_0+\theta_N+o(1)}\) and nominal excess

\hypertarget{mf:eq:B3.61}{}

\[\fitmath{
(A-M)-d-(K_0-K)-w_o-B_c+g-\ell
\le A-M+5\sigma/3+\delta_1.
\qquad\text{(MF3.61)}
}\]

If \(\fitinline{J\ge0}\), \(\fitinline{-d-(K_0-K)-w_o+J_+=-c-2w}\); otherwise it is \(\fitinline{\le\delta_1}\). The largest increment is \(\fitinline{2\sigma-\ell_*=5\sigma/3}\), a terminal loss, plus \(\fitinline{\delta}\) in the zero-slot padded core.

Localize genuine nonzero \(\fitinline{j}\) at nominal \(\fitinline{Z^{2a_0-K-g}}\) by the whole-dyad rule before common extraction \(\fitinline{u=D_2a,v=E_2b}\). Put

\[\fitmath{
c_2=\log_Zq_{D_2},\quad d_2=\log_Zq_{E_2},\quad b_2=(c_2+d_2)/2,
\quad p_2=\log_Zq_{\operatorname{rad}(D_2E_2)},\quad G_c=(D_2,E_2),\quad g_2=\log_Zq_{G_c}.
}\]

Apply the complete-support factorization

\hypertarget{mf:eq:B3.62}{}

\textbf{(MF3.62)} is \protect\hyperlink{cf:eq:A.16}{(12)}, with \(\fitinline{(D,E)=(D_2,E_2)}\) and the same residual \(\fitinline{a,b,j}\).

It holds only for \(\fitinline{(a,b)=(ab,D_2E_2)=1}\), allowing arbitrary powers and slot/plain coincidences. Extracted primes puncture residual factors; freeze supplying slots. A plain \(\fitinline{p^5}\) and slot \(\fitinline{p}\) therefore retain \(\fitinline{\chi_p(j)^6}\), not one.

Equations \protect\hyperlink{cf:eq:A.17}{(13)}--\protect\hyperlink{cf:eq:A.18}{(14)} force \(\fitinline{G_c\mid j}\). At equal non-six-divisible multiplicities split divided frequency into units/nonunits, radical lengths \(\fitinline{t_2,V}\) and nonunit radical \(\fitinline{V_{\rm id}}\). Set \(\fitinline{j=G_cV_{\rm id}h'}\). Extend \(\fitinline{1_{\rm part}F(D_2,E_2;G_cV_{\rm id}h')/Z^{g_2-t_2}}\) by zero off its partition; it has modulus \(\fitinline{\le1}\) by

\begin{longtable}[]{@{}ll@{}}
\toprule\noalign{}
Common multiplicities and divided frequency & Bound \\
\midrule\noalign{}
\endhead
\bottomrule\noalign{}
\endlastfoot
\(\fitinline{i=i,\ 6\nmid i}\), unit & \(\fitinline{P^{i-1}}\) \\
\(\fitinline{i=i,\ 6\nmid i}\), nonunit & \(\fitinline{P^i}\) \\
\(\fitinline{i=i,\ 6\mid i}\), either & \(\fitinline{P^i}\) \\
\(\fitinline{i>j_0,\ 6\mid j_0}\), unit & \(\fitinline{P^{j_0}}\) \\
\end{longtable}

Other unequal cases vanish. In \(\fitinline{v=G_cV_{\rm id}}\), reduce valuations modulo six: nonzero residues retain moving characters, zeros the existing puncture. Active primes lie in \(\fitinline{t_2+V}\). Factor only \(\fitinline{\chi_n(v)}\), preserving varying-row zeros of \(\fitinline{\chi_n(h')}\). Thus

\hypertarget{mf:eq:B3.63}{}

\[\fitmath{
m'=2a_0-K-g-g_2-V,\quad q'=\widetilde q+w_o+t_2+V,
\qquad M'=m'+q'=M+J-g-g_2+t_2\le M-\sigma.
\qquad\text{(MF3.63)}
}\]

The pulled-back retained weight has \(\fitinline{q_{h'}\le Z^{m'+\delta_2}}\), \(\fitinline{0\le\delta_2<\xi}\). Nonempty ranges have \(\fitinline{m'+\delta_2\ge0}\); declare \(\fitinline{m'_{\rm act}=m'+\delta_2}\), \(\fitinline{M'_{\rm act}=M'+\delta_2}\). Extend weight and partition scalar by zero across this ball, keeping them signed.

The condition \(\fitinline{s\mid u,v}\) forces \(\fitinline{s}\) into the common radical, so labels cost \(\fitinline{Z^{p_2-s_0+o(1)}}\). The ledger is

\begin{longtable}[]{@{}ll@{}}
\toprule\noalign{}
Factor & Exponent \\
\midrule\noalign{}
\endhead
\bottomrule\noalign{}
\endlastfoot
Original squared central normalization & \(\fitinline{-a_0}\) \\
Poisson factor and inverse roots & \(\fitinline{K+g-a_0}\) \\
Amplifier density & \(\fitinline{-\ell}\) \\
Common-support count with fixed \(\fitinline{s}\) & \(\fitinline{p_2-s_0}\) \\
Common correlation & \(\fitinline{g_2-t_2}\) \\
Conversion to normalized residual products & \(\fitinline{a_0-b_2}\) \\
\end{longtable}

Subtracting this sum from \protect\hyperlink{mf:eq:B3.57}{(MF3.57)} leaves the inner allowance

\hypertarget{mf:eq:B3.64}{}

\[\fitmath{
M'+\Delta,\qquad \Delta=b_2-p_2+w+B_c+\ell\ge w+\ell.
\qquad\text{(MF3.64)}
}\]

The last inequality uses \(\fitinline{b_2\ge p_2}\).

Extend \protect\hyperlink{mf:eq:B3.62}{(MF3.62)}'s right side to individually allowed pairs, retaining original common-support masks, formal full norms, roots, kernel and pulled-back row weight. Insert full \(\fitinline{\sum_{t\mid(a,b)}\mu(t)}\); only this sum restores the genuine correlation. Its artificial common primes do not enter \(\fitinline{D_2,E_2}\) or frequency restrictions. Separate in row/whole-product norms, then allocate primes of \(\fitinline{t}\) by

\hypertarget{mf:eq:B3.65}{}

\[\fitmath{
\mathbf1_{p\mid\prod_i n_i}
=\sum_{\varnothing\ne J}(-1)^{|J|+1}\prod_{i\in J}\mathbf1_{p\mid n_i},
\qquad\text{(MF3.65)}
}\]

Factors are the two plains and live slots. Selected plains become \(\fitinline{n_i=pl_i}\) without residual puncture, dividing corresponding scales in both rectangles by \(\fitinline{q_p}\). Selected slots freeze; selecting distinct primes in one prime slot vanishes. Quotients may retain \(\fitinline{p}\).

Thus each child retains a whole-product character, common mask and norm power across plains/rectangles; live slots keep original \(\fitinline{\nu_i}\) without Gauss factors. Sides may have different whole-product norm powers, but their fixed-inducing predicates agree: conjugation and multiplication by fixed-modulus characters preserve conductor divisibility, including unchanged redundant zeros. At each fixed-modulus stage enlarge \(\fitinline{\Theta}\) to all primitive characters dividing that modulus plus the finitely many fixed ray twists; it remains finite. On an exceptional row inducing \(\fitinline{\vartheta}\), freeze live slots to obtain

\hypertarget{mf:eq:B3.66}{}

\textbf{(MF3.66)} is the exceptional rectangle coefficient \protect\hyperlink{cf:eq:D.27}{(91)}, with \(\fitinline{D_b=\mathcal D_b}\), this child's \(\fitinline{\vartheta,t}\), and its complete squarefree mask \(\fitinline{R_*}\).

The polynomial-size \(\fitinline{R_*}\) is common to both variables/rectangles. Retain extracted row scalars, including zeros, until positive or absolute estimation.

\textbf{Extraction normalization.} Put \(\fitinline{\alpha_1=a_0-c_2}\), \(\fitinline{\alpha_2=a_0-d_2}\). Let \(\fitinline{F_j}\) be genuinely frozen entrance slots and \(\fitinline{I_j}\) the live slots before final allocation. Their pre-allocation plain product \(\fitinline{T_j}\) satisfies \(\fitinline{T_j\prod_{i\in I_j}Z^{z_i}=Z^{\alpha_j+e_j}}\), \(\fitinline{e_j=\sum_{i\in F_j}\log_Z(q_{p_i}/Z^{z_i})}\). Let \(\fitinline{d_{j,i}}\) be final products selected in plain \(\fitinline{i}\), \(\fitinline{J_j\subset I_j}\) newly frozen slots, and put

\[\fitmath{
a_{j,i}=\log_Zq_{d_{j,i}},\quad
r_j=a_{j,1}+a_{j,2}+\sum_{i\in J_j}z_i,\quad
\omega_j=\sum_{i\in J_j}\log_Z(q_{p_i}/Z^{z_i}).
}\]

Use lower-endpoint dyads \(\fitinline{Z^{t_-}\le q_t<2Z^{t_-}}\). Every selected product is divisible by \(\fitinline{t}\), so

\hypertarget{mf:eq:B3.67}{}

\[\fitmath{
r_j+\omega_j\ge t_-,\qquad |\omega_j|\le\theta_N,
\qquad |e_j+\omega_j|\le\theta_N.
\qquad\text{(MF3.67)}
}\]

The last bound uses disjoint subsets of current entrance slots, without recounting ancestors. Raw child normalizer is \(\fitinline{Z^{-(\alpha_j-r_j)/2}}\), extraction scalar \(\fitinline{Z^{-r_j/2}}\), formal scale product \(\fitinline{Z^{\alpha_j+e_j-r_j}}\). Signed raw rectangles may have negative individual lengths.

\textbf{Child eligibility.} Split by the same fixed-inducing predicate with complete source masks. On kept rows, \protect\hyperlink{mf:eq:B3.62}{(MF3.62)}--\protect\hyperlink{mf:eq:B3.65}{(MF3.65)} give exactly the whole-product character in nonzero \(\fitinline{h'}\) with pulled-back radial weight. Frozen zeros belong to its natural presentation; the extra common mask is row-fixed. Delete it by \protect\hyperlink{mf:eq:B3.41}{(MF3.41)}, obtaining a natural state of width \(\fitinline{M'_{\rm act}}\) with precisely the non-fixed-inducing keep set. Each plain and surviving slot has this same character/mask, and Fourier density precedes live labels/rows. Hence the smaller-width theorem applies to the actual positive norm; complementary fixed-inducing rows remain for separate evaluation.

\textbf{Nonexceptional children.} Use Cauchy, then split centered children into rectangles; only now bound common weights/scalars by one and enlarge the ball. For pre-allocation plain lengths \(\fitinline{x_i}\) and post-lengths \(\fitinline{y_i=x_i-a_{j,i}}\), put \(\fitinline{\pi_0=\sum(-x_i)_+}\), \(\fitinline{\pi=\sum(-y_i)_+}\). Omit rectangles whose plain upper support is below the first ideal; otherwise \(\fitinline{0\le\pi_0\le\pi\le\theta_N}\). The clipped reduction and total are

\hypertarget{mf:eq:B3.68}{}

\[\fitmath{
r_{{\rm clip},j}=r_j-(\pi-\pi_0)\ge0,\qquad
A_{{\rm clip},j}=\alpha_j+e_j-r_j+\pi,
\qquad\text{(MF3.68)}
}\]

and its exact scalar relative to the pre-normalizer is \(\fitinline{Z^{(-r_{{\rm clip},j}+e_j+\pi_0)/2}}\). For a paired divisor dyad, the count and both scalars satisfy

\hypertarget{mf:eq:B3.69}{}

\[\fitmath{
t_--\frac{r_{{\rm clip},1}+r_{{\rm clip},2}}2
 +\frac{e_1+\pi_{0,1}+e_2+\pi_{0,2}}2\le2\theta_N.
\qquad\text{(MF3.69)}
}\]

Substitute \protect\hyperlink{mf:eq:B3.68}{(MF3.68)} and \protect\hyperlink{mf:eq:B3.67}{(MF3.67)}; there are at most four rectangle pairs. Do not assign a common clipped reduction to a signed difference.

With no surviving slot, \protect\hyperlink{mf:eq:B3.69}{(MF3.69)} and the smaller-width zero-slot bound fit \protect\hyperlink{mf:eq:B3.64}{(MF3.64)}, \(\fitinline{\Delta\ge0}\). Otherwise side one's pre-allocation total is \(\fitinline{A'=A-c-w-c_2}\), and

\hypertarget{mf:eq:B3.70}{}

\[\fitmath{
(A'-M')-(A-M)
=g+w-d-c_2-(K_0-K)-w_o+g_2-t_2
\le6(w+\ell)+\delta_1.
\qquad\text{(MF3.70)}
}\]

Use \(\fitinline{g_2-t_2\le c_2}\). In the amplified norm, \(\fitinline{g\le d+(K_0-K)+6\ell+\delta_1}\); in errors, \(\fitinline{g\le d+(K_0-K)+\sigma+\delta_1}\) and the extraction table gives \(\fitinline{w-w_o+\sigma\le6w}\). Exchange \(\fitinline{c_2,d_2}\) for the other side.

Delete extra masks. For each actual natural child with live slots put \(\fitinline{F_{\rm act}=(A_{\rm act}-M'_{\rm act}+(6\kappa-1)z_{\rm act})_+}\). Equations \protect\hyperlink{mf:eq:B3.68}{(MF3.68)}--\protect\hyperlink{mf:eq:B3.70}{(MF3.70)}, parent capacity and aggregate endpoint bound give

\hypertarget{mf:eq:B3.71}{}

\[\fitmath{
F_{\rm act}\le6(w+\ell)+\delta_1+2\theta_N.
\qquad\text{(MF3.71)}
}\]

Use \protect\hyperlink{mf:eq:B3.47}{(MF3.47)} to remove whole slots until defect vanishes or no slot remains. A length-\(\fitinline{d}\) slot releases \(\fitinline{6\kappa d}\) capacity and costs \(\fitinline{\kappa d}\) energy. A greedy prefix overshoots by at most one mesh, so

\hypertarget{mf:eq:B3.72}{}

\[\fitmath{
\kappa d_z\le F_{\rm act}/6+\kappa\eta_{\rm mesh}
\le\Delta+\delta_1/6+\theta_N/3+\eta_{\rm mesh}.
\qquad\text{(MF3.72)}
}\]

If slots remain use the smaller-width positive theorem; otherwise use unrestricted zero-slot bounds. Equations \protect\hyperlink{mf:eq:B3.69}{(MF3.69)}, \protect\hyperlink{mf:eq:B3.72}{(MF3.72)} and \(\fitinline{M'_{\rm act}-M'=\delta_2}\) give edge cost \(\fitinline{\le\delta_2+\delta_1/6+7\theta_N/3+\eta_{\rm mesh}}\) plus its small-power share. Cauchy averages costs, including the single mesh.

\hypertarget{fixed-character-rows-and-rectangle-cancellation}{}

\hypertarget{a.7.2.9-fixed-character-rows-and-rectangle-cancellation}{%
\subparagraph{A.7.2.9 Fixed-character rows and rectangle cancellation}\label{a.7.2.9-fixed-character-rows-and-rectangle-cancellation}}

\textbf{Exceptional count.} Genuine second allocations avoid old moving support by original zeros, so no old character cancels \(\fitinline{G_cV_{\rm id}}\) there. For nonunit common multiplicity-one primes of radical length \(\fitinline{v_1}\), put \(\fitinline{f=2v_1}\). Their fixed row valuation is two; unramified fixed-modulus induction forces \(\fitinline{v_p(h')\equiv4\pmod6}\). More generally each good prime has a prescribed valuation class modulo six, zero outside frozen support. Units and \(\fitinline{S}\)-valuations give finitely many sectors, with sixth powers free. Thus \(\fitinline{(h')=h_0v^6}\), \(\fitinline{h_0}\) fixed sixth-power-free, \(\fitinline{q_{h_0}\ge Z^{4v_1}=Z^{2f}}\). Counting gives \(\fitinline{(m'_{\rm act}-2f)/6}\); we use only

\hypertarget{mf:eq:B3.73}{}

\[\fitmath{
\#\{h'\text{ exceptional}:q_{h'}\ll Z^{m'_{\rm act}}\}
\ll Z^{(m'-f)/6+\delta_2/6+o(1)}.
\qquad\text{(MF3.73)}
}\]

Below the forced initial norm the range is empty up to fixed endpoints. This counts principal and all other fixed-inducing rows, including positively added partition rows; no \(\fitinline{t_2}\) saving is used.

After full M\"{o}bius allocation, side \(\fitinline{j}\)'s normalized absolute volume is \(\fitinline{O(Z^{(\alpha_j-r_j)/2+e_j})}\). Both extraction scalars and divisor count give paired volume \(\fitinline{\le Z^{a_0-b_2+2\theta_N+o(1)}}\) by \protect\hyperlink{mf:eq:B3.67}{(MF3.67)}, since \(\fitinline{t_--r_1-r_2\le\omega_2}\), \(\fitinline{r_1\ge0}\). Including the exceptional count, its nominal excess is

\hypertarget{mf:eq:B3.74}{}

\[\fitmath{
\frac{m'-f}{6}+a_0-b_2-(M'+\Delta)
=A-\frac56M-F_1-F_2,
\qquad\text{(MF3.74)}
}\]

where

\hypertarget{mf:eq:B3.75}{}

\[\fitmath{
\begin{aligned}
F_1&=c/6+5\{d+(K_0-K)\}/6+w/3+\widetilde q/6+w_o+B_c-5g/6+\ell,\\
F_2&=2b_2-5g_2/6-p_2+t_2+V/6+f/6.
\end{aligned}\qquad\text{(MF3.75)}
}\]

These expressions obey

\hypertarget{mf:eq:B3.76}{}

\[\fitmath{
F_1\ge\frac23(c+w)-3\sigma-5\delta_1/6,
\qquad F_2\ge\frac23b_2.
\qquad\text{(MF3.76)}
}\]

Initially take \(\fitinline{g=J_+}\), \(\fitinline{\ell=0}\), \(\fitinline{D_0=d+K_0-K}\). For \(\fitinline{J\ge0}\), \(\fitinline{F_1=c+2w+\widetilde q/6+w_o/6+B_c}\). Otherwise \(\fitinline{\widetilde q\ge R}\) gives \(\fitinline{B_c+\widetilde q/6\ge(3c-5D_0)_+/6-5\delta_1/6}\), hence \(\fitinline{F_1\ge2(c+w)/3-w/3-5\delta_1/6}\). Restoring \(\fitinline{g\le J_++2\sigma}\), \(\fitinline{\ell\ge0}\), \(\fitinline{w\le7\sigma/3}\) costs \(\fitinline{<3\sigma}\). For \(\fitinline{F_2}\) use the exact local table of \S{}\protect\hyperlink{cf:sec:D.6}{A.5.6}, with \(\fitinline{j=j_0}\):

The four pairs \(\fitinline{(F_2,b_2)}\) are \(\fitinline{(7i/6,i)}\), \(\fitinline{((7i-5)/6+1_{i=1}/3,i)}\), \(\fitinline{(7i/6-1,i)}\), and \(\fitinline{(i+j_0/6-1,(i+j_0)/2)}\) for, respectively, equal non-six-divisible/unit, equal non-six-divisible/nonunit, equal six-divisible, and unequal \(\fitinline{i>j_0}\), \(\fitinline{6\mid j_0}\)/unit cases.

Each first entry is \(\fitinline{\ge2/3}\) of the second; at nonunit \(\fitinline{i=1}\), \(\fitinline{f/6}\) gives equality. Thus \protect\hyperlink{mf:eq:B3.76}{(MF3.76)} holds and uncentered \(\fitinline{A\le5M/6}\) costs at most \(\fitinline{3\sigma+O(\xi+\theta_N)}\), validating comparisons.

\textbf{Masked lattice cancellation.} Apply \S{}\protect\hyperlink{cf:sec:D.1}{A.5.1} to fixed \(\fitinline{\vartheta\in\Theta}\), polynomial squarefree \(\fitinline{R_*}\), annular \(\fitinline{W_i}\) and common real \(\fitinline{t}\).

\hypertarget{mf:eq:B3.77}{}

\textbf{(MF3.77)} is \protect\hyperlink{cf:eq:D.1}{(65)}, with the same \(\fitinline{L_i(X,t)}\), \(\fitinline{c_{\vartheta,R_*}}\) and \(\fitinline{I_i(t)}\), loss \(\fitinline{\epsilon_1}\) and \(\fitinline{\langle t\rangle=1+|t|}\). It holds uniformly for \(\fitinline{X>0}\).

Its proof there is fixed-lattice Poisson followed by complete mask inversion; in particular \(\fitinline{c_{\vartheta,R_*}=c_\vartheta\sum_{d\mid R_*}\mu(d)\vartheta(d)/q_d}\). Primary/\(\fitinline{S}\) restrictions are fixed lattice classes, and already excluded primes can be omitted. Thus it applies to every retained exceptional mask.

For \(\fitinline{U_i=X_i/q_{b_i}}\), \(\fitinline{V_i=Y_i/q_{b_i}}\), \(\fitinline{T=U_1U_2=V_1V_2}\), the sum of \protect\hyperlink{mf:eq:B3.66}{(MF3.66)} is the difference of the two \(\fitinline{L_1L_2}\) products. If all four scales are \(\fitinline{\ge Z^r}\), \(\fitinline{r\ge0}\), then

\hypertarget{mf:eq:B3.78}{}

\textbf{(MF3.78)} is the normalized rectangle bound \protect\hyperlink{cf:eq:D.2}{(66)}, with these \(\fitinline{U_i,V_i,T}\), loss \(\fitinline{\epsilon_1}\) and height factor \(\fitinline{\langle t\rangle^{2J}}\).

The same row, mask and norm power make main terms equal. Without a nonnegative lower scale use the all-scale volume \(\fitinline{O(T^{1/2})}\).

For centered exceptional children, pre-allocation lower lengths are \(\fitinline{L-c-w-c_2}\) and \(\fitinline{L-c-w-d_2}\). Put \(\fitinline{r=(L-c-w-\min(c_2,d_2))_+}\) and call the saving side one. Final extraction removes at most \(\fitinline{r_1}\) from either plain: use \protect\hyperlink{mf:eq:B3.78}{(MF3.78)} if \(\fitinline{r_1<r}\), otherwise volume. Preserve raw equal-product scales without clipping. Exactly,

\hypertarget{mf:eq:B3.79}{}

\[\fitmath{
t_--r_1-r_2-(r-r_1)_+
=(t_--r_2)-\max(r,r_1)\le\omega_2-r.
\qquad\text{(MF3.79)}
}\]

Paired volume is therefore \(\fitinline{Z^{a_0-b_2-r+2\theta_N+o(1)}}\), uniformly in live slots. Their labels never entered Fourier density; sum them now by volume with \protect\hyperlink{mf:eq:B3.67}{(MF3.67)}'s normalization.

Put \(\fitinline{v=c+w+\min(c_2,d_2)}\). From \protect\hyperlink{mf:eq:B3.76}{(MF3.76)}, \(\fitinline{F_1+F_2\ge2v/3-3\sigma-5\delta_1/6}\). The remaining centered exceptional deficit is bounded by

\hypertarget{mf:eq:B3.80}{}

\[\fitmath{
A-5M/6-2v/3-(L-v)_+\le(A-M)_+,
\qquad L=M/4.
\qquad\text{(MF3.80)}
}\]

As in \S{}\protect\hyperlink{cf:sec:D.6}{A.5.6}, the maximum occurs at \(\fitinline{v=L}\) and equals \(\fitinline{A-M}\). Positive slots have \(\fitinline{A\le M}\), zero-slot padding \(\fitinline{A\le M+\delta}\); exceptional terminal cost is \(\fitinline{\le3\sigma+\delta}\), with reserved support/frequency errors.

\hypertarget{closure-and-uniform-common-mesh}{}

\hypertarget{a.7.2.10-closure-and-uniform-common-mesh}{%
\subparagraph{A.7.2.10 Closure and uniform common mesh}\label{a.7.2.10-closure-and-uniform-common-mesh}}

Strict nominal widths decrease by \(\fitinline{\sigma}\). Actual enclosure/support corrections cost \(\fitinline{C_*\xi}\) for numerical \(\fitinline{C_*}\) independent of slot count: they involve one aggregate \(\fitinline{\theta_N}\), two frequency enclosures, at most two plain clippings and finitely many aggregate lengths in \protect\hyperlink{mf:eq:B3.53}{(MF3.53)}--\protect\hyperlink{mf:eq:B3.80}{(MF3.80)}. Ledgers are affine or positive-part affine; positive capacity varies by at most \(\fitinline{|\Delta M|+|\Delta A|+5|\Delta z|}\), since \(\fitinline{6\kappa-1\le5}\). Greedy removal charges one mesh.

Choose \(\fitinline{\rho,\sigma,\delta}\) from the bounded real ranges and final \(\fitinline{\epsilon}\), with \(\fitinline{\delta\ll\min(\rho,\sigma)}\), so that

\[\fitmath{
T_{\rm term}:=\rho+\delta+\rho/6+3\sigma+5\sigma/3<\epsilon/4.
}\]

Put \(\fitinline{D=2+\lceil2M_{\max}/\sigma\rceil}\), and, still before specifying the fixed slot system, choose

\hypertarget{mf:eq:B3.81}{}

\[\fitmath{
\xi<\min\{\delta/2,\rho\mathfrak m(7/10)/2,\sigma/(4C_*),
                   \epsilon/(16C_*D)\},\qquad
\eta_{\rm mesh}<\min\{\sigma/6,\epsilon/(16C_*D)\},\qquad
\epsilon_0<\epsilon/(16C_*D).
\qquad\text{(MF3.81)}
}\]

These preserve comparison margins, padding and width decrease \(\fitinline{\ge\sigma/2}\). Each strict child leaves its band; depth is \(\fitinline{\le D-2}\). Same-band comparisons use earlier uncentered stages, each plain reflects at most once before centering, and errors are not re-amplified. Thus the induction is acyclic.

Now fix arithmetic data, profiles and finite slot count, then \(\fitinline{H_N}\) and boxes through depth \(\fitinline{D}\). Eventually \(\fitinline{\theta_N\ll\xi}\), tails satisfy whole-dyad inclusions and pool avoids slots. Allocate stage losses within \(\fitinline{\epsilon_0}\): e.g.~at most \(\fitinline{N+2}\) mask masses get \(\fitinline{\epsilon_0/(10(N+2))}\) each; use global divisor bounds \(\fitinline{C_N^{\omega(v)},d_{N+2}(v)^C\ll_{N,C,a}q_v^a}\) for one small \(\fitinline{a}\). Prime displacement costs \(\fitinline{2ez}\), independent of slot count. These affect constants, not the already chosen mesh.

With \(\fitinline{d}\) strict calls remaining allow \(\fitinline{E_d=T_{\rm term}+(d+1)C_*(\eta_{\rm mesh}+\xi+\epsilon_0)}\). Floor, diagonal and exception fit \(\fitinline{E_0}\); \protect\hyperlink{mf:eq:B3.69}{(MF3.69)}--\protect\hyperlink{mf:eq:B3.72}{(MF3.72)} bound a strict edge by its child plus one allowance, including comparisons/reflection masses. Cauchy averages and triangle takes maxima, so terminal loss occurs once. Equation \protect\hyperlink{mf:eq:B3.81}{(MF3.81)} gives \(\fitinline{E_{D-2}<\epsilon}\), simultaneously for every surviving slot subset.

Choose raw tails and finite profile/height orders backwards, then one final threshold for tails, endpoints and logarithms. Child extra masks lie on frozen columns or a frozen amplifier prime, hence stay polynomial through finite depth. Exceptional masks also include row radicals and are covered uniformly by \protect\hyperlink{mf:eq:B3.77}{(MF3.77)}. All moving support, including canceled powers, remains explicit. This proves \protect\hyperlink{mf:eq:B3.40}{(MF3.40)}.

The original \(\fitinline{3/4\le\kappa\le1}\) and refined \(\fitinline{7/10\le\kappa\le3/4}\) ranges are separate applications of this same descent. Arithmetic and exceptional estimates are \(\fitinline{\kappa}\)-independent. Its only uses are prime bound \protect\hyperlink{mf:eq:B3.47}{(MF3.47)}, requiring \(\fitinline{2\beta-1\le\kappa}\); comparison \protect\hyperlink{mf:eq:B3.49}{(MF3.49)}, where \(\fitinline{\mathfrak m'(\kappa)=2/(6\kappa-1)^2>0}\) gives minimum \(\fitinline{\mathfrak m(7/10)}\); and removal \protect\hyperlink{mf:eq:B3.72}{(MF3.72)}, whose cost/capacity is \(\fitinline{1/6}\). Bounds \(\fitinline{6\kappa-1\le5}\), \(\fitinline{\kappa\le1}\) make the mesh uniform. In the original contradiction, \(\fitinline{7/8<\beta\le1}\) permits \(\fitinline{\kappa=2\beta-1=3/4+2(\beta-7/8)}\) without an improved zero-free assumption.

For completeness, fixed nonprincipal zero-slot rows are also restored directly. For a primitive nonprincipal \(\fitinline{\vartheta}\) of the fixed modulus, orthogonality gives \(\fitinline{c_\vartheta=0}\), hence \(\fitinline{c_{\vartheta,R_*}=0}\) in \protect\hyperlink{mf:eq:B3.77}{(MF3.77)}. Each normalized nonempty \(\fitinline{X\ge1}\) plain is \(\fitinline{O(Z^{\epsilon_1}X^{-1/2}\langle t\rangle^J)}\); bounded subunit scales use volume. Prescribed valuations modulo six give \(\fitinline{O(Z^{m/6+\epsilon_1})}\) rows by the exceptional parameterization without forced-prime savings. The fixed family is finite, so multiplying and choosing small \(\fitinline{\epsilon_1}\) restores all such rows within \(\fitinline{Z^{M+\epsilon}}\), with independent heights and common masks. Principal inducing characters remain excluded. \ensuremath{\square}

\hypertarget{the-two-incumbent-ranges-and-the-physical-interfaces}{}

\hypertarget{a.7.3-the-two-incumbent-ranges-and-the-physical-interfaces}{%
\paragraph{A.7.3 The two incumbent ranges and the physical interfaces}\label{a.7.3-the-two-incumbent-ranges-and-the-physical-interfaces}}

\textbf{Proposition A.7.5} (Original bootstrap moment). \emph{\hypertarget{mf:plain-original-bootstrap}{} Under the natural family, masks, ray coefficients, individual slot bound \(\fitinline{0<z_i\le\eta_{\rm mesh}}\) and quantifier order of Theorem~\protect\hyperlink{mf:plain}{A.7.4}, suppose \(\fitinline{\beta\ge51/100}\), \(\fitinline{3/4\le\kappa\le1}\) and \(\fitinline{2\beta-1\le\kappa}\). Then the positive moment holds whenever \(\fitinline{n_1+n_2+6\kappa z\le M}\); the zero-slot assertion is unconditional. In particular the choice \(\fitinline{\kappa=2\beta-1}\) is valid in the original contradiction range \(\fitinline{7/8<\beta\le1}\).}

\emph{Proof.} Restrict Theorem \protect\hyperlink{mf:plain}{A.7.4} to \(\fitinline{[3/4,1]}\). The common proof above gives comparison margin \(\fitinline{\mathfrak m(3/4)M=M/14}\), removal ratio \(\fitinline{1/6}\) and total error \(\fitinline{E_{D-2}<\epsilon}\) by \protect\hyperlink{mf:eq:B3.81}{(MF3.81)}, with one terminal allowance. Its actual row presentation requires only \(\fitinline{2\beta-1\le\kappa}\), independently of an improved zero-free conclusion. \ensuremath{\square}

\textbf{Proposition A.7.6} (Refined incumbent moment). \emph{\hypertarget{mf:plain-refined}{} The same natural moment, with each active \(\fitinline{0<z_i\le\eta_{\rm mesh}}\), is valid for \(\fitinline{7/10\le\kappa\le3/4}\), \(\fitinline{\beta\ge51/100}\) and \(\fitinline{2\beta-1\le\kappa}\), with its mesh chosen uniformly in this interval before the slot system, profiles and target height.}

\emph{Proof.} Restrict the same proof to \(\fitinline{[7/10,3/4]}\): \(\fitinline{6\kappa-1\le7/2}\), margin \(\fitinline{\ge\mathfrak m(7/10)M>0}\), removal ratio \(\fitinline{1/6}\). Equation \protect\hyperlink{mf:eq:B3.81}{(MF3.81)} preserves the stated quantifier order. This is a separate refinement of the original range. \ensuremath{\square}

\hypertarget{literal-initial-fields-and-positivity.}{}

Literal initial fields and positivity.

For fixed source \(\fitinline{\eta}\), take scale \(\fitinline{U}\), puncture one, width \(\fitinline{1+\delta}\) once \(\fitinline{q_{\mathfrak q_\eta}\le U^\delta}\), and keep set \(\fitinline{\operatorname{initialKeep}(\eta,\mathfrak Q_0)}\). At \(\fitinline{X_1=X_2=U^m}\), Theorem \protect\hyperlink{mf:plain}{A.7.4} gives

\[\fitmath{
\sum_k\mathbf1_{\operatorname{initialKeep}(\eta,\mathfrak Q_0)}(k)
 \Re\Phi(q_k/U)\,|S_k(U^m)^2\prod_{i\in I}Q_{k,i}|^2
\ll \mathcal P_J(W_1,W_2)^{2}\mathcal D_J(\Phi)
 (1+|t_1|+|t_2|+\max_i|v_i|)^h U^{1+\delta+\epsilon}
}\]

provided \(\fitinline{2m+6\kappa\sum_{i\in I}z_i\le1}\) and \(\fitinline{0<z_i\le\eta_{\rm mesh}}\). Here \(\fitinline{\mathcal P_J}\) uses finitely many plain Euler seminorms (common slot profile fixed in the constant); \(\fitinline{\mathcal D_J}\) controls the nonnegative radial Schwartz profile. Choose \(\fitinline{J,h}\) before heights and live subset. Absorb \(\fitinline{\delta}\) and local loss into the requested loss. This physical field is \(\fitinline{|S|^4|\prod Q_i|^2}\).

Theorem \protect\hyperlink{mf:inverse}{A.7.3} admits literal relative-character slot coefficients, fixed ray Fourier coefficients and common character conjugation, all bounded and row-independent. Shared inverse primes are allocated by \protect\hyperlink{mf:eq:B3.32}{(MF3.32)}. At \(\fitinline{m=1}\) the gates are \(\fitinline{r+2z<1}\), \(\fitinline{2r+8z<3}\). For finitely many profiles/derivatives take maximum orders/thresholds and union seminorms. All energies use nonnegative measures and squared moduli, so retained subfamilies decrease them. Signed M\"{o}bius coefficients remain until their full identity is restored.

\hypertarget{every-rescaled-completed-term.}{}

Every rescaled completed term.

Theorem \protect\hyperlink{mf:marked-energy}{A.7.1} covers all original slot subsets. Priority allocation to frozen products preserves row-independent surviving coefficients, retaining exact central scalars and tuple counts outside. For physical rescaling by subset width \(\fitinline{d}\),

\[\fitmath{
M'=M-2d,\qquad z'=\ell-d,\qquad N'=1+\ell-d.
}\]

If \(\fitinline{M+\ell\le1}\), then \(\fitinline{M'+2z'\le N'}\) for every such subset. The energy exponent is therefore \(\fitinline{M'+(1+3\ell-2M+d)_+/4+\epsilon}\). After the physical Cauchy square root, its exact compensator and assigned tuple mass give the additional exponent

\[\fitmath{
-d+\frac18(1+3\ell-2M+d)_+.
}\]

Its slopes are \(\fitinline{-1,-7/8}\), so the maximum is \(\fitinline{d=0}\); this pays every rescaled term with its actual outside coefficient. Individual energies need not be \(\fitinline{\le Z^{M'}}\).

\textbf{Lemma A.7.7} (Actual finite source rows). \emph{\hypertarget{mf:finite-source-field}{} Let \(\fitinline{\mathcal R}\) be a finite row family with an injective map \(\fitinline{\iota:\mathcal R\longrightarrow\mathcal O}\). Suppose that, for each \(\fitinline{u\in\mathcal R}\), its displayed character \(\fitinline{\chi_u}\) agrees at every ideal with the natural function \(\fitinline{n\mapsto\eta(n)\chi_n(\iota(u))}\), including the fixed bad-prime zero mask. Suppose also that \(\fitinline{\iota(u)\in\operatorname{initialKeep}(\eta,\mathfrak Q_0)}\) and \(\fitinline{q_{\iota(u)}\le U}\). Choose a fixed smooth radial majorant \(\fitinline{\Phi\ge0}\) with \(\fitinline{\Phi\ge1}\) on \(\fitinline{[0,1]}\). Then the unweighted finite-row plain times marked-prime energy is at most the initial-state positive energy in Theorem~\protect\hyperlink{mf:plain}{A.7.4}. It consequently obeys its bound \(\fitinline{CU^{1+\delta+\epsilon}}\) and the stated polynomial height control whenever \(\fitinline{2m+6\kappa z\le1}\).}

\emph{Proof.} The literal character identity identifies all plain and relative-slot coefficients. Injectivity reindexes without multiplicity; the radial majorant and keep indicator dominate on the image. Enlarge only the resulting nonnegative squared-modulus sum. For profiles with at most two logarithmic derivatives, the product rule for \(\fitinline{W(y)y^{-\sigma+it}}\) gives finite derivatives of \(\fitinline{W}\) times polynomials in bounded \(\fitinline{\sigma}\) and height \(\fitinline{t}\). Combine with the moment's height polynomial. Degree and constant precede live rows, subset, ray characters and heights, proving the stated finite field and uniformity. \ensuremath{\square}

\hypertarget{ov:section}{}

\hypertarget{a.8-overlapping-prime-lists-in-the-inverse-moment}{%
\subsubsection{A.8 Overlapping prime lists in the inverse moment}\label{a.8-overlapping-prime-lists-in-the-inverse-moment}}

This coefficient-specific extension precedes affine arguments and is endpoint-independent.

\textbf{Lemma A.8.1} (Overlap adapter for the inverse moment). \emph{\hypertarget{ov:inverse-overlap}{} The estimate of Theorem~\protect\hyperlink{mf:inverse}{A.7.3} remains valid for a fixed finite number of possibly overlapping or repeated prime lists with the same bounded annular coefficients. Both admission inequalities may use the full sum of the original slot widths.}

\emph{Proof.} Partition finite prime tuples by exact equality pattern. Freeze each block of \(\fitinline{k\ge2}\) equal primes before the row norm. Its intersecting annuli have comparable \(\fitinline{D_i}\), giving central mass

\[\fitmath{
\sum_{\mathfrak p\in\bigcap_i\mathcal P_i}
       \prod_{i=1}^k D_i^{-1/2}|b_i(\mathfrak p)|
 \ll (\min_iD_i)\prod_{i=1}^k D_i^{-1/2}
 \ll D^{1-k/2}\ll1.
}\]

Constants depend only on fixed slot count and profile suprema. Keep \(\fitinline{\varphi_u(\mathfrak p)^k}\) until taking the positive norm, then contract by its modulus \(\fitinline{\le1}\) with zeros retained. The inverse polynomial is unchanged and may share that prime.

Delete frozen primes from singleton lists. For \(\fitinline{K}\) singletons, color their union independently with \(\fitinline{K}\) colors. Each distinct ordered tuple receives prescribed slot colors with probability \(\fitinline{K^{-K}}\); repeated tuples never do. Its polynomial is therefore \(\fitinline{K^K}\) times the average over products on disjoint, row-independent colored lists. Apply Theorem \protect\hyperlink{mf:inverse}{A.7.3}, then Minkowski/Jensen, costing \(\fitinline{K^{2K}}\) in energy. For \(\fitinline{K=0}\) use UM. Deleting slots decreases both admission expressions.

Triangle over bounded-mass frozen choices and finitely many equality patterns proves the claim. This coloring extension applies only to MI; MP retains its canonical-pool contract and overlap allocation. \ensuremath{\square}

\hypertarget{init:section}{}

\hypertarget{appendix-b.-the-literal-probe-supported-witnesses-and-initialization}{%
\subsection{Appendix B. The literal probe, supported witnesses and initialization}\label{appendix-b.-the-literal-probe-supported-witnesses-and-initialization}}

The same probe underlies the natural initialization and the final return. First its exact coefficients, signed Gram bound and principal signal are proved; then its nonprincipal rows supply the witnesses and weighted numerator. Finally the original complete comparison proves \(\fitinline{\beta\le7/8}\), independently of every stronger source in the main text. The Mellin transform is always that of the reciprocal signal \(\fitinline{f_\eta}\), not of the error \(\fitinline{J_\eta-f_\eta}\). Equations (188)--(223), (234)--(236) and (240) retain their original stable labels solely for local navigation.

\hypertarget{b.1-the-actual-arithmetic-coefficients-and-the-two-representations}{%
\subsubsection{B.1 The actual arithmetic coefficients and the two representations}\label{b.1-the-actual-arithmetic-coefficients-and-the-two-representations}}

Use the arithmetic chapter's primary generators, sextic symbols, angular \(\fitinline{\alpha}\), ray phases \(\fitinline{G,R}\), normalized Gauss sums \(\fitinline{\gamma_j}\) and additive measure. All residue-character powers, including zero/sixth, vanish at nonunits; rows keep six units. Ideal columns avoid fixed finite \(\fitinline{S}\), containing primes over six, primitive target conductor, ray moduli and a sufficiently large small-prime cutoff, chosen before scale/tuples/rows.

Let \(\fitinline{b_*}\) generate the product of the prime ideals in \(\fitinline{S}\). Choose a primitive residue character \(\fitinline{\xi\bmod b_*}\), quadratic at odd residue characteristic and cubic at the prime of norm four. Put

\[\fitmath{
\tau=q_{b_*}^{-1/2}g_\xi(b_*,1),\qquad
 \Xi(A)=\xi(A)\chi_A(b_*),\qquad |\tau|=1.
}\]

The finite identities proved in the arithmetic chapter give

\[\fitmath{
\gamma_2(c)^3=\mu(c)\alpha(c),\quad
 \gamma_1(c)\gamma_2(c)=\mu(c)\alpha(c)G(c),\quad
 \chi_b(a)=R(a,b)\chi_a(b)\quad ((a,b)=1).
}\]

These fix coefficient phases. Set

\[\fitmath{
V_G(r)=\frac{e^{-(\log r)^2/4}}{2\sqrt\pi},\qquad
 \widehat V_G(t)=e^{t^2},
}\]

and choose nonnegative, nonzero, smooth annular \(\fitinline{W_0,W_1}\). For \(\fitinline{A=cn^3}\), where \(\fitinline{c}\) is squarefree and \(\fitinline{c,n}\) may share primes, define

\hypertarget{dr:T}{} \hypertarget{dr:I}{}

\[\fitmath{
\begin{aligned}
 T_{m,a,\eta}(t)
 &=\sum_{c\ {\rm sf},n}
 \gamma_2(c)\overline{\alpha(A)G(A)}\eta(A)\Xi(A)^{-1}
 R(A,a)\chi_A(m)q_c^{-1/2-t}q_n^{-1-3t},\qquad\text{(187)}\\
 I_\eta(X,Y,Z)
 &=\frac1Y\sum_a
 \frac{W_1(q_a/Y)\chi_a(b_*)}
 {\tau\xi(a)\sqrt{q_{b_*}q_aX}}
 \sum_m\xi(m)q_a^{-1/2}g_{\chi_a}(a,-m)
 W_0\left(\frac{q_m}{q_{b_*}q_aX}\right)
 \frac1{2\pi i}\int_{(4)}Z^te^{t^2}T_{m,a,\eta}(t)\,dt.
 \qquad\text{(188)}
\end{aligned}
}\]

The \(\fitinline{a}\) columns include prime powers and full zero-extended imprimitive Gauss sums. Annuli make \(\fitinline{a,m\ne0}\) finite; all other sums/integrals converge absolutely. Define \(\fitinline{I_{\eta;D}}\) by inserting \(\fitinline{\mathbf1_{D\mid cn^3}}\) on the whole theta index.

Choose finitely many distinct positive \(\fitinline{\ell_i}\), \(\fitinline{P_i=Z^{\ell_i}}\), and common smooth nonnegative nonzero \(\fitinline{w}\) supported in \(\fitinline{(1,2)}\). Each slot uses all primes outside \(\fitinline{S}\); distinct lengths make supports eventually disjoint. They may be as fine as any prescribed positive mesh with prescribed total \(\fitinline{\ell=\sum_i\ell_i}\). Write \(\fitinline{w_i=w}\) to index slots, \(\fitinline{p_J=\prod_{i\in J}p_i}\), and define

\hypertarget{dr:comp}{}

\[\fitmath{
I_\eta^{\rm comp}(Z)=
 \sum_{(p_i)}\prod_iw_i(q_{p_i}/P_i)
 \sum_{J\subseteq\{1,\ldots,r_{\rm slot}\}}
 (-1)^{|J|}q_{p_J}^{-3/2}\overline{\eta(p_{J^c})}
 I_{\eta;p_{J^c}}(X/q_{p_J},Y/q_{p_J},Zq_{p_{J^c}}).
 \qquad\text{(189)}
}\]

Weights are row-independent and the coefficient is \(\fitinline{\overline\eta}\). One slot gives divisibility minus rescaled subtraction; retain all collisions with \(\fitinline{c,n}\) and later rows.

For the physical representation, expand \(\fitinline{\overline G}\) on its fixed ray group, partition \(\fitinline{a}\) into classes \(\fitinline{\sigma}\) for \(\fitinline{R(A,a)}\), set \(\fitinline{Q_0=q_{b_*}XY}\) and take \(\fitinline{\Omega=1}\) on the product annuli. Mellin inversion of \(\fitinline{W_0}\) expresses \protect\hyperlink{dr:I}{(188)} as a finite sum of

\hypertarget{dr:physical}{}

\[\fitmath{
\frac{Q_0^{-1/2}}{2\pi}\int_{\mathbb R}\widehat W_0(iv)
 \sum_{m\ne0}\Omega(q_m/Q_0)\xi(m)(q_m/Q_0)^{-iv}
 A_{m,\sigma,v}(Y)B_{m,\sigma}(Z)\,dv,
 \qquad\text{(190)}
}\]

where \(\fitinline{B}\) is a completed theta sum, with precisely the coefficients

\[\fitmath{
\mathcal C_V(N;m,\nu)=\sum_{c\ {\rm sf},n}
 \frac{\overline{\alpha(cn^3)}\gamma_2(c)\nu(cn^3)\chi_{cn^3}(m)}
 {\sqrt{q_c}\,q_n}V(q_cq_n^3/N),
}\]

and

\hypertarget{dr:A}{}

\[\fitmath{
A_{m,\sigma,v}(Y)=Y^{-3/2}\sum_{a\in\sigma}
 \frac{\chi_a(b_*)}{\tau\xi(a)}W_1(q_a/Y)
 (q_a/Y)^{-1+iv}g_{\chi_a}(a,-m).
 \qquad\text{(191)}
}\]

The identity \(\fitinline{q_m/(q_{b_*}q_aX)=(q_m/Q_0)/(q_a/Y)}\) gives \(\fitinline{Y^{-3/2}}\).

Let \(\fitinline{\widetilde W_0(|y|^2)}\) be the planar Fourier transform of \(\fitinline{W_0(|z|^2)}\) and \(\fitinline{\mathscr M(z)=\int_0^\infty\widetilde W_0(r)r^{z-1}\,dr}\), holomorphic for \(\fitinline{\Re z>0}\) with arbitrary vertical decay on compact positive substrips. Poisson gives

\hypertarget{dr:Poisson}{}

\[\fitmath{
\begin{aligned}
 I_\eta(X,Y,Z)=\frac1{(2\pi i)^3}
 \int_{(3)}\int_{(3)}\int_{(2)}
 &X^{1/2-z}Z^{s+z-1}Y^{w-1}e^{(s+z-1)^2}
 \mathscr M(z)\widehat W_1(w)\\[-2mm]
 &\quad\times\sum_u^{(6)}q_u^{-z}\overline{\xi(u)}
 \frac{\zeta_K^S(6z)L_K^S(w,\psi_u)}
 {L_K^S(s,\eta\overline{\psi_u})}H_{\eta,u}(s,w,z)
 \,dz\,dw\,ds.\qquad\text{(192)}
\end{aligned}
}\]

Its finite convolution is

\[\fitmath{
\mathcal F(a,A,H)=\sum_{\substack{d\bmod a\\H\equiv b_*Ad\bmod a}}
 \chi_a(d)g_{\xi\chi_A}
 \left(b_*A,\frac{H-b_*Ad}{a}\right).
}\]

Expanding both Gauss sums in the full transform modulo \(\fitinline{ab_*A}\) gives \(\fitinline{q_a^{1/2}\mathcal F}\). Planar dilation \(\fitinline{X/q_A}\) leaves physical factor \(\fitinline{X^{1/2}/(\sqrt{q_{b_*}}q_A)}\). The primitive factor at \(\fitinline{S}\) kills frequencies meeting \(\fitinline{S}\), including zero. Write others uniquely \(\fitinline{H=ur^6}\), with primary \(\fitinline{r}\) and the unit retained in \(\fitinline{u}\); summing \(\fitinline{r}\) gives \(\fitinline{\zeta_K^S(6z)}\) without a factor six. Mellin inversion and \(\fitinline{s=t+1-z}\) yield \protect\hyperlink{dr:Poisson}{(192)}. On its initial lines a majorant is

\[\fitmath{
C_S\sum_{u,c,n,a,r}q_u^{-2}q_c^{-5/2}q_n^{-7}q_a^{-2}q_r^{-12},
}\]

times rapidly decreasing heights, justifying absolute interchange. CRT gives \(\fitinline{g_{\xi\chi_A}(b_*A,K)=\xi(A)\chi_A(b_*)g_\xi(b_*,K)g_{\chi_A}(A,K)}\); \(\fitinline{\Xi(A)^{-1}}\) and \(\fitinline{\chi_a(b_*)/(\tau\xi(a))}\) cancel fixed-prime factors, leaving \(\fitinline{\overline{\xi(u)}}\). The displayed phase identities cancel cross-prime terms, including shared primes of \(\fitinline{a,A}\).

\hypertarget{dr:local}{}

\hypertarget{b.2-the-local-table-including-all-shared-primes}{%
\subsubsection{B.2 The local table, including all shared primes}\label{b.2-the-local-table-including-all-shared-primes}}

The complete local table for \protect\hyperlink{dr:Poisson}{(192)} and numerator averaging uses, at \(\fitinline{p\notin S}\),

\[\fitmath{
\begin{gathered}
 Q=q_p,\quad j=v_p(u)\in\{0,\ldots,5\},\quad
 \rho=\chi_p(u/p^j),\quad
 \mathfrak a=\overline{\alpha(p)}^3\eta(p)^3,\quad
 \mathfrak b=\overline{\alpha(p)}^2\eta(p)^2\overline{\chi_p(4)},
 \quad \mathfrak b^3=\mathfrak a^2,\\
 V=Q^{-6z},\quad R_p=\mathfrak a^2Q^{4-6s-6z},\quad
 W^\circ=\rho Q^{-w},\quad W=\chi_p(u)Q^{-w},\quad
 D=\eta(p)\overline{\chi_p(u)}Q^{-s}.
\end{gathered}
}\]

Thus \(\fitinline{D=W=0}\) for \(\fitinline{j>0}\). Set

\hypertarget{dr:Euler-table}{}

\[\fitmath{
\begin{aligned}
 P_p^*&=\frac1{1-R_p}
 \left\{\frac{R_p(1-Q^{-1})-\eta(p)(Q-1)Q^{-s-w}
 V^{\mathbf1_{j\le1}}}{1-V}+J_j\right\},\\
 P_p&=(1-V)^{-1}+\mathbf1_{j=0}\frac{W^\circ}{1-W^\circ}+P_p^*,
 \qquad H_p=\frac{P_p(1-V)(1-W)}{1-D},
 \quad H_{\eta,u}=\prod_{p\notin S}H_p,
 \qquad\text{(193)}
\end{aligned}
}\]

where

\[\fitmath{
\begin{array}{c|l}
 j&J_j\\ \hline
0&-\eta(p)\rho^{-1}Q^{-s}+W^\circ R_p\\
1&\eta(p)Q^{-s-w}\\
2&\mathfrak a\rho^{-3}Q^{3/2-3s}\\
3&-\eta(p)\mathfrak b\rho^{-2}Q^{2-3s-w}
                  +\mathfrak b^2\rho^{-4}Q^{2-4s}\\
4&-\eta(p)\mathfrak a\rho^{-3}Q^{5/2-4s-w}\\
5&-\mathfrak a^2Q^{3-6s}.
\end{array}
}\]

For \(\fitinline{t=e+3l}\), \(\fitinline{e\in\{0,1\}}\), the local convolution is

\[\fitmath{
C_p(t,k,j')=\sum_{\substack{d\bmod p^k\\p^{j'}\equiv p^td\bmod p^k}}
 \chi_p(d)^k g_{\chi_p^t}
 \left(p^t,\frac{p^{j'}-p^td}{p^k}\right).
}\]

Lifting a prime Gauss sum gives

\[\fitmath{
g_{\chi_p^h}(p^t,K)=
 \begin{cases}Q^{t-1}g_{\chi_p^h}(p,K/p^{t-1}),&p^{t-1}\mid K,\\
 0,&p^{t-1}\nmid K.
 \end{cases}
}\]

For \(\fitinline{t=0}\) it is one at \(\fitinline{k=0}\), \(\fitinline{\mathbf1_{j'=0}}\) at \(\fitinline{k>0}\). For \(\fitinline{t>0,k=0}\) it is the source Gauss sum; for \(\fitinline{k=1,j'\ge1}\) it is \(\fitinline{\chi_p(-1)^{-1}g_{\chi_p}(p,1)g_{\chi_p^{t-1}}(p^t,p^{j'-1})}\), while \(\fitinline{j'=0}\) has no solution. For \(\fitinline{k\ge2}\) the last lift \(\fitinline{d=d_0+p^{k-1}h}\) sums a nontrivial additive character and vanishes. Surviving \(\fitinline{t>0}\) valuations are

\[\fitmath{
\begin{array}{c|cc}
 t&k=0&k=1\\ \hline
 6r+6&j'\ge t-1&j'=t\\
 6r+1&j'=t-1&j'\ge t\\
 6r+3&j'=t-1&j'=t\\
 6r+4&j'=t-1&j'=t.
\end{array}
}\]

After CRT phase cancellation, with \(\fitinline{j'=j+6m}\) and \(\fitinline{\omega_p=\chi_p(-1)}\), the summand is

\[\fitmath{
(-1)^e\gamma_1(p)^{-e}\eta(p)^e
 (\mathfrak a\overline{\gamma_3(p)})^l
 \rho^{k-t}\omega_p^{tk-el-l(l-1)/2}C_p(t,k,j+6m)
 Q^{-(s+1/2)e-(1+3s)l-wk-6zm}.
}\]

The first two strict tails sum to \(\fitinline{R_p(1-Q^{-1})/((1-R_p)(1-V))}\) and \(\fitinline{-\eta(p)(Q-1)Q^{-s-w}V^{\mathbf1_{j\le1}}/((1-R_p)(1-V))}\). Finite Gauss identities give equality terms \(\fitinline{J_j/(1-R_p)}\); \(\fitinline{t=0}\) gives the initial terms of \(\fitinline{P_p}\). This covers every valuation.

Put \(\fitinline{c_p=\overline{\eta(p)}Q^s}\), \(\fitinline{h_p=Q^{-w}}\), and

\hypertarget{dr:G}{}

\[\fitmath{
G_p=\frac{(c_p-h_p)(1-V)(1-W)P_p^*-h_p(1-VW)}{1-D}.
 \qquad\text{(194)}
}\]

The local effect of the two alternatives in \protect\hyperlink{dr:comp}{(189)} is the identity

\[\fitmath{
\overline{\eta(p)}Q^{s+z-1}P_p^*-Q^{z-w-1}P_p
 =\frac{(1-D)Q^{z-1}G_p}{(1-V)(1-W)}.
}\]

Consequently its exact Poisson correction is

\hypertarget{dr:Hcomp}{}

\[\fitmath{
\mathcal H_{\eta,u,Z}=
 \sum_{(p_i)}\prod_iw_i(Q_i/P_i)Q_i^{z-1}G_{p_i}
 \prod_{p\notin S\cup\{p_i\}}H_p.
 \qquad\text{(195)}
}\]

For \(\fitinline{p\nmid u}\) put \(\fitinline{v=\chi_p(u)}\), \(\fitinline{E_p=P_p^*+D}\). Substituting \(\fitinline{c_pD=v^{-1}}\) gives the quotient-free identity

\hypertarget{dr:Delta}{}

\[\fitmath{
G_p+v^{-1}H_p=\frac{(1-V)(1-W)}{1-D}
 \left[(v^{-1}-h_p)\left(\frac V{1-V}-D\right)
 +(c_p+v^{-1}-h_p)E_p\right].\qquad\text{(196)}
}\]

Define \(\fitinline{\Delta_p=G_p+\overline{\psi_u(p)}H_p}\) also at row primes. Expanding \(\fitinline{G_p=-\overline{\psi_u(p)}H_p+\Delta_p}\) gives all main/error choices in \protect\hyperlink{dr:Hcomp}{(195)}, with main slot exactly \(\fitinline{-\sum_pw_i(Q/P_i)Q^{z-1}\overline{\psi_u(p)}}\), including zero character values.

The local table also proves normal convergence and \(\fitinline{H_{\eta,u}\ll q_u^\epsilon}\) on compact real boxes in

\hypertarget{dr:box}{}

\[\fitmath{
\Re s\ge51/100,\quad\Re z\ge17/50,\quad
 \Re w\ge-1/100,\quad\Re(s+w)\ge1+\epsilon_1.
 \qquad\text{(197)}
}\]

Unramified leading defects \(\fitinline{Q^{-s-w}}\) and \(\fitinline{Q^{4-6s-6z}}\) are \(\fitinline{\le Q^{-1-\epsilon_1}}\) and \(\fitinline{Q^{-11/10}}\); all others have summable margins. Row-prime defects are negative powers, giving \(\fitinline{q_u^\epsilon}\) after a fixed cutoff; \(\fitinline{1-W}\) cancels. For \(\fitinline{u=1}\) also allow \(\fitinline{\Re s\ge87/100}\), \(\fitinline{\Re w\ge3/4}\), \(\fitinline{\Re z\ge1/8}\). Substituting the six local terms in \protect\hyperlink{dr:Hcomp}{(195)} gives \(\fitinline{\ll q_u^\epsilon Z^{(17/50)\ell}}\) on \(\fitinline{(\beta+e,1/2,17/50)}\) and \(\fitinline{\ll q_u^\epsilon Z^{v\ell}}\) on \(\fitinline{(2,2,v)}\), \(\fitinline{v>2}\): selected off-row factors are bounded; the first line's exceptional \(\fitinline{O(Q^{1/2})}\) occurs only at divisor-many primes. Intervening fixed boxes have fixed polynomial height/scale bounds.

\hypertarget{dr:Gram}{}

\hypertarget{b.3-the-signed-gram-calculation-and-the-physical-bound}{%
\subsubsection{B.3 The signed Gram calculation and the physical bound}\label{b.3-the-signed-gram-calculation-and-the-physical-bound}}

\textbf{Lemma B.2.1} (Signed physical Gram estimate). \emph{\hypertarget{dr:Gram-lemma}{} For \protect\hyperlink{dr:A}{(191)}, with \(\fitinline{P_a=Y^2/Q_0\ge1}\),}

\[\fitmath{
\sum_{q_m\ll Q_0}|A_{m,\sigma,v}(Y)|^2
 \ll (1+|v|)^J\frac{Q_0}{Y}
       \left(1+P_a^{1/6}+\frac{P_a^2}{Y}\right)Z^\epsilon.
}\]

\emph{This assertion uses the displayed fixed ray coefficients.}

\emph{Proof.} Expand the signed square. Arithmetic Fourier correlation has \(\fitinline{F(u,v;0)=\mathbf1_{u=v}\phi(u)}\) and, for \(\fitinline{u=Cn_1,v=Cn_2}\), \(\fitinline{(n_1,n_2)=1}\),

\[\fitmath{
F(Cn_1,Cn_2;Ck)=\chi_{n_1}(k)\overline{\chi_{n_2}(-k)}
 R(n_1,n_2)L_C(n_1,n_2;k),\qquad |L_C|\le q_C.
}\]

Row Poisson contributes \(\fitinline{Q_0/(q_uq_v)}\), the Gauss pair \(\fitinline{q_uq_vF}\), giving prefactor \(\fitinline{Q_0/Y^3}\); zero frequency costs \(\fitinline{O(Q_0/Y)}\).

For nonzero frequencies extend the coefficient to residual pairs, setting it zero when a \(\fitinline{C}\)-prime divides both, then insert full M\"{o}bius coprimality. On divisor \(\fitinline{d}\), write \(\fitinline{n_i=dm_i}\); natural zeros force \(\fitinline{(d,Ck)=1}\). The extension cancels in the complete divisor sum. After finite ray separation the period norm is \(\fitinline{\ll q_Cq_k}\); its complete mean vanishes unless the sixth-free part of \(\fitinline{k}\) is supported on \(\fitinline{C,S}\), since every other prime gives a nonprincipal complete character sum.

On \(\fitinline{q_C\asymp C_0,q_d\asymp D_0,q_Cq_k\asymp RP_a}\), put \(\fitinline{N=Y/(C_0D_0)}\). Two-column lattice Poisson bounds zero-mean sums by

\[\fitmath{
(1+|v|)^J R^{-B}C_0N^2
           \min\{1,(RP_a/N)^A\}.
}\]

Below period use counting; above period integrate all nonzero dual frequencies by parts. There are \(\fitinline{O(D_0RP_a)}\) choices of \(\fitinline{C,d,k}\), and the remaining dyads use

\[\fitmath{
\sum_{i,j\ge0}2^{-i-j}\min\{1,(2^{i+j}/\Lambda)^A\}
 \ll_A\frac{1+\log(2+\Lambda)}\Lambda,\quad
 \Lambda=Y/(RP_a),\quad A>1.
}\]

Group by \(\fitinline{i+j=n}\) and split \(\fitinline{(n+1)2^{-n}}\) at \(\fitinline{2^n=\Lambda}\). Including \(\fitinline{Q_0/Y^3}\) gives \(\fitinline{(Q_0/Y)P_a^2/Y}\); large \(\fitinline{B}\) sums \(\fitinline{R}\).

For nonzero mean, write \(\fitinline{k=er^6}\) with sixth-free \(\fitinline{e}\) supported on \(\fitinline{C,S}\). There are divisor-many \(\fitinline{e}\) and \(\fitinline{O(Z^\epsilon(RP_a/C_0)^{1/6})}\) roots. Absolute column volume gives \(\fitinline{(Q_0/Y)P_a^{1/6}}\) with summable weights \(\fitinline{C_0^{-1/6+\epsilon}D_0^{-1+\epsilon}}\). Add zero frequency; finitely many profile derivatives give \(\fitinline{(1+|v|)^J}\). \ensuremath{\square}

We apply the completed marked energy of Lemma~\protect\hyperlink{mf:marked-energy}{A.7.1} with its following exact interface. For the displayed \(\fitinline{\mathcal C_V}\), product-form row-independent prime marks of total width \(\fitinline{z}\), completed length \(\fitinline{Z^N}\), and row range \(\fitinline{Z^M}\),

\hypertarget{dr:marked-energy-interface}{}

\[\fitmath{
M+2z\le N\quad\Longrightarrow\quad
 \sum_{0<q_m\ll Z^M}|\mathcal T(m)|^2
 \ll Z^{M+(N+2z-2M)_+/4+\epsilon}.
 \qquad\text{(198)}
}\]

Lemma \protect\hyperlink{mf:marked-energy}{A.7.1} includes both collision masks, all cusps, joint profiles and full kernel tails for these multiple slots.

Let \(\fitinline{X=Z^x,Y=Z^y,M=x+y}\) and assume

\[\fitmath{
0<\ell<x\le y,\qquad M+\ell\le1,\qquad y-2x+\ell<0.
}\]

Cauchy in \protect\hyperlink{dr:physical}{(190)}, the Gram lemma and \protect\hyperlink{dr:marked-energy-interface}{(198)} give

\hypertarget{dr:physical-bound}{}

\[\fitmath{
|I_\eta^{\rm comp}(Z)|
 \ll Z^{x/2+(y-x)/12+(1+3\ell-2M)_+/8+\epsilon}.
 \qquad\text{(199)}
}\]

For subset \(\fitinline{J}\) of width \(\fitinline{d}\), the rescaled row/mark/completed exponents are \(\fitinline{M'=M-2d}\), \(\fitinline{z'=\ell-d}\), \(\fitinline{N'=1+\ell-d}\), satisfying \(\fitinline{M'+2z'\le N'}\). Outside factor, tuple mass and energy square root add

\[\fitmath{
-d+\frac{(1+3\ell-2M+d)_+}{8}.
}\]

Its slopes \(\fitinline{-1,-7/8}\) make \(\fitinline{d=0}\) maximal. The Gram condition \(\fitinline{P_a^2/(Y/q_{p_J})\ll1}\) follows from \(\fitinline{y-2x+\ell<0}\). Integrate common Mellin profiles and sum finite rays to obtain \protect\hyperlink{dr:physical-bound}{(199)} for every subset.

\hypertarget{dr:signal}{}

\hypertarget{b.4-the-unique-principal-signal-and-its-nonvanishing-numerator}{%
\subsubsection{B.4 The unique principal signal and its nonvanishing numerator}\label{b.4-the-unique-principal-signal-and-its-nonvanishing-numerator}}

Only \(\fitinline{u=1}\) is principal in \protect\hyperlink{dr:Poisson}{(192)}: a nonunit sixth-free row has a good prime of valuation \(\fitinline{1,\ldots,5}\), hence nontrivial tame Kummer inertia. For unit \(\fitinline{u\ne1}\), CRT gives primary \(\fitinline{n\equiv-2-3\omega\pmod{36\mathcal O}}\) and \(\fitinline{n\equiv1}\) at other fixed primes. Then \(\fitinline{q_n\equiv7\pmod{36}}\) and \(\fitinline{\chi_n(u)=u^{(q_n-1)/6}=u\ne1}\). Keep all five units in the remainder.

Let \(\fitinline{\rho_S=\operatorname{Res}_{s=1}\zeta_K^S(s)>0}\). At \(\fitinline{w=1,z=1/6}\) the table simplifies, with \(\fitinline{t=Q^{-1}}\), to

\hypertarget{dr:principal-local}{}

\[\fitmath{
H_{\eta,p}(s)=\frac{1-t^2}{1-R_p}
 \left(1+\frac{t(D-R_p)}{1-D}\right),\qquad
 D=\eta(p)Q^{-s},\quad
 R_p=\overline{\alpha(p)}^6\eta(p)^6Q^{3-6s}.
 \qquad\text{(200)}
}\]

In particular, for any fixed \(\fitinline{b>2/3}\),

\[\fitmath{
|H_{\eta,p}(s)-1|\ll_b
 Q^{-2}+Q^{-1-b}+Q^{3-6b}+Q^{2-6b}\qquad(\Re s\ge b).
}\]

These powers are prime summable. Enlarge the fixed cutoff until their majorant sum is \(\fitinline{\le\log(3/2)}\); normal convergence and \(\fitinline{\prod(1+a_p)\le e^{\sum a_p}}\) give

\hypertarget{dr:H-nonzero}{}

\[\fitmath{
H_\eta(s)=\prod_{p\notin S}H_{\eta,p}(s),\qquad
 |H_\eta(s)-1|\le\tfrac12\quad (\Re s\ge b).
 \qquad\text{(201)}
}\]

Thus \(\fitinline{H_\eta}\) is holomorphic and nonzero uniformly in height; \(\fitinline{b=87/100}\) suffices for all established endpoints here.

The positive residue scalar is

\[\fitmath{
c_S=\frac{\rho_S^2}{6}\widehat W_1(1)\mathscr M(1/6)>0.
}\]

Indeed, Fourier transformation of the radial Riesz kernel, first written as a positive integral of Gaussians, gives for \(\fitinline{0<\sigma<1}\)

\[\fitmath{
\mathscr M(\sigma)=\left(\frac{2\pi}{\sqrt3}\right)^{1-2\sigma}
 \frac{\Gamma(\sigma)}{\Gamma(1-\sigma)}
 \int_0^\infty W_0(r)r^{-\sigma}\,dr>0.
}\]

Annular support/Schwartz decay justify interchange; \(\fitinline{1/6}\) in \(\fitinline{c_S}\) is the residue scaling of \(\fitinline{\zeta_K^S(6z)}\).

On this residue locus \protect\hyperlink{dr:Delta}{(196)} gives \(\fitinline{G_p/H_p=-1+O(Q^{-b})}\). Define the real nonzero normalizer

\[\fitmath{
A(Z)=(-1)^{r_{\rm slot}}Z^{-\ell/6}
       \prod_i\sum_pw_i(Q/P_i)Q^{-5/6}.
}\]

The fixed-field prime ideal theorem gives \(\fitinline{|A(Z)|\asymp(\log Z)^{-r_{\rm slot}}}\). Multiplying over fixed disjoint slots gives

\hypertarget{dr:prime-residue}{}

\[\fitmath{
\mathcal H_{\eta,1,Z}(s,1,1/6)
 =H_\eta(s)Z^{\ell/6}A(Z)(1+O(Z^{-\mu_P})),
 \quad 0<\mu_P<b\min_i\ell_i,\quad\Re s\ge b.
 \qquad\text{(202)}
}\]

Put

\hypertarget{dr:J-f}{}

\[\fitmath{
h=1+\ell-x,\quad k=\tfrac56-\tfrac x3-\tfrac\ell6,
 \qquad J_\eta(Z)=\frac{I_\eta^{\rm comp}(Z)}{c_S A(Z)},
 \quad
 f_\eta(Z)=\frac1{2\pi i}\int_{(2)}Z^{s-k}
 e^{(s-5/6)^2}\frac{H_\eta(s)}{L_K^S(s,\eta)}\,ds.
 \qquad\text{(203)}
}\]

The residue power is \(\fitinline{X^{1/3}Z^{s-5/6}Z^{\ell/6}=Z^{s-k}}\). Equation \protect\hyperlink{dr:prime-residue}{(202)} therefore identifies \(\fitinline{f_\eta}\) up to power-small error: the compensator gives \(\fitinline{Z^{\ell/6}}\) and \(\fitinline{A(Z)}\) has logarithmic size.

\hypertarget{dr:detection}{}

\hypertarget{b.5-buffered-zero-bins-and-two-witnesses-at-the-same-zero}{%
\subsubsection{B.5 Buffered zero bins and two witnesses at the same zero}\label{b.5-buffered-zero-bins-and-two-witnesses-at-the-same-zero}}

For fixed finite-ray \(\fitinline{\nu}\), let \(\fitinline{\varphi_u=\nu\psi_u^{\pm1}}\) with its natural zeros, and set

\[\fitmath{
M_u(D;W)=D^{-1/2}\sum_n\mu(n)\varphi_u(n)W(q_n/D),\qquad
 S_u(N;W)=N^{-1/2}\sum_n\varphi_u(n)W(q_n/N).
}\]

Preceding moments have common annular controls for real powers, scale derivatives and independent polynomial heights. In particular,

\hypertarget{dr:unmarked}{}

\[\fitmath{
\sum_{q_u\asymp U}^{(6)}|M_u(U^r;W)|^2
 \ll U^{\max\{1,(1+5r)/6\}+\epsilon},\qquad
 \sum_{q_u\asymp U}^{(6),\varphi_u^*\ne1}|S_u(U^m;W)|^4
 \ll U^{1+\epsilon}.
 \qquad\text{(204)}
}\]

The long inverse exponent follows from Lemma \protect\hyperlink{cf:inverse-moment}{A.4.1} by the exact amplification

\[\fitmath{
M_u(D;W)=\sum_{d\mid\operatorname{rad}(a)}
 \mu(d)\varphi_u(d)q_d^{-1/2}M_{ua^6}(D/q_d;W).
}\]

This separates squarefree column primes dividing \(\fitinline{a}\), retaining shared-prime zeros, with coefficient mass \(\fitinline{q_a^\epsilon}\). Log-scale differentiation replaces \(\fitinline{W}\) by \(\fitinline{-W/2-yW'}\). Sobolev and the inverse moment bound \(\fitinline{\sup_{0<E\le D}}\) on balls \(\fitinline{H\ge D^{1+\vartheta}}\). Average over primary \(\fitinline{q_a\le(H/U)^{1/6}}\) for \(\fitinline{H=\max(2U,D^{1+\vartheta})}\); valuations modulo six make \(\fitinline{(u,a)\mapsto ua^6}\) injective. Its \(\fitinline{\gg_S(H/U)^{1/6}}\) preimages give \(\fitinline{U^{1/6}H^{5/6}(UD)^\epsilon}\), hence \protect\hyperlink{dr:unmarked}{(204)}. Row \(\fitinline{u=1}\) gives smooth M\"{o}bius sums \(\fitinline{\ll D^{11/12+\epsilon}}\); Mellin inversion with a test nonzero at a proposed zero yields preliminary \(\fitinline{\beta\le11/12}\), before \(\fitinline{7/8}\) or marked plain moments.

Fix finite conjugation-stable \(\fitinline{\Theta}\) containing all target numerator/denominator/ray/phase twists. For nonunit sixth-free \(\fitinline{u}\), \(\fitinline{\mathcal X(u)=\{\nu\psi_u^{\pm1}:\nu\in\Theta\}}\) is nonprincipal of conductor \(\fitinline{O_S(q_u)}\) with redundant zeros. Put \(\fitinline{a_*=51/100}\), small \(\fitinline{e>0}\), \(\fitinline{I=\lceil1/e\rceil+2}\). Let \(\fitinline{B_j(u)}\) maximize \(\fitinline{a_*}\) and zero real parts of \(\fitinline{\mathcal X(u)}\) up to height \(\fitinline{3^jT}\). Some \(\fitinline{i<I}\) has \(\fitinline{B_{i+1}-B_i\le e}\); round \(\fitinline{B_i}\) down on the \(\fitinline{e}\) mesh:

\hypertarget{dr:bins}{}

\[\fitmath{
a\le B_i<a+e,\qquad B_{i+1}<a+2e.\qquad\text{(205)}
}\]

Labels \(\fitinline{(i,a)}\) are finite. Above the floor there is one actual zero \(\fitinline{\rho=\sigma+i\gamma}\), \(\fitinline{a\le\sigma<a+e}\), \(\fitinline{|\gamma|\le3^iT}\), supplying both witnesses. Put \(\fitinline{\delta=2a-1}\).

We record the analytic consequence of the buffer. The ordinary Hecke functional equation in this normalization is

\[\fitmath{
\Lambda(s,\varphi)=(3q_{\mathfrak f})^{s/2}(2\pi)^{-s}
 \Gamma(s)L_K(s,\varphi),\qquad
 \Lambda(s,\varphi)=\epsilon_\varphi\Lambda(1-s,\bar\varphi),
 \quad |\epsilon_\varphi|=1.
}\]

Euler convergence and the functional equation give polynomial conductor/height growth. Continue the Euler logarithm on a zero-free disk about \(\fitinline{2+it}\). Borel--Carath\'{e}odory bounds it by \(\fitinline{O_e(\log(2q_{\mathfrak f}(3+|t|)^2))}\) on a smaller disk; Euler convergence bounds it on radius \(\fitinline{49/100}\). Three circles with another buffer reduces the logarithm to a power \(\fitinline{<1}\). Exponentiating either sign gives

\hypertarget{dr:buffered}{}

\[\fitmath{
|L_K(s,\varphi)|+|L_K(s,\varphi)^{-1}|
 \ll_{e,\epsilon}\{2q_{\mathfrak f}(3+|t|)^2\}^{\epsilon}.
 \qquad\text{(206)}
}\]

Use radii \(\fitinline{2-a-2e,2-a-4e,2-a-6e}\); Cauchy gives logarithmic derivatives. For principal characters apply this to \(\fitinline{(s-1)\zeta_K(s)/(s+1)}\); to the right of \(\fitinline{\beta+e}\) it holds at all heights. Deleted factors cost \(\fitinline{U^\epsilon}\) on fixed positive half-planes and \(\fitinline{U^{v+\epsilon}}\) on \(\fitinline{\Re s\ge-v}\), by summing \(\fitinline{O(q_p^{-\sigma})}\) beyond a fixed cutoff and the divisor bound.

Mellin inversion and reflection now give

\hypertarget{dr:saturation-upper}{}

\[\fitmath{
|M_u(U^r)|^2\ll U^{\delta r+\epsilon_0},\qquad
 |S_u(U^m)|^2\ll U^{\delta\min(m,1-m)+\epsilon_0}.
 \qquad\text{(207)}
}\]

For prime tests apply \(\fitinline{-L'/L}\), divide by \(\fitinline{\log q_p}\) and subtract prime powers costing \(\fitinline{O(P^{1/2+\epsilon})}\) before normalization:

\hypertarget{dr:prime-upper}{}

\[\fitmath{
\left|P^{-1/2}\sum_p\overline{\psi_u(p)}w(q_p/P)
 (q_p/P)^{it}\right|
 \ll P^{\delta/2+O(e)+\epsilon}U^\epsilon(1+T)^A.
 \qquad\text{(208)}
}\]

The buffer \protect\hyperlink{dr:bins}{(205)} covers the witness neighborhood; distinct norm twists keep independent heights.

\textbf{Lemma B.5.1} (Same-zero witnesses). \emph{\hypertarget{dr:same-zero}{} Fix \(\fitinline{t_0\in[1,3/2]}\) for a nonfloor bin before selecting its witnesses. Each row has a presentation, a real power \(\fitinline{\sigma\in[a,a+e)}\), a common height and dyadic lengths \(\fitinline{U^r,U^m}\) such that}

\[\fitmath{
r\le t_0+o(1),\quad r+m\ge t_0-o(1),\quad
 0\le m\le\tfrac12+O(\epsilon_0),
}\]

\[\fitmath{
|M_u(U^r)|^2\gg U^{\delta r-O(\epsilon_0)},\qquad
 |S_u(U^m)|^2\gg U^{\delta m-O(\epsilon_0)}.
}\]

\emph{After fixing presentations and dyads, the base profiles are common.}

\emph{Proof.} Choose \(\fitinline{V_\le=1}\) on \(\fitinline{[0,1]}\), zero on \(\fitinline{[2,\infty)}\), and put \(\fitinline{C_\varphi(s)=\sum_l\mu(l)\varphi(l)V_\le(q_l/U^{t_0})q_l^{-s}}\). At the actual zero \(\fitinline{\rho}\) shift

\[\fitmath{
\frac1{2\pi i}\int_{(2)}U^{40z}\Gamma(z)
 L_K^S(\rho+z,\varphi)C_\varphi(\rho+z)\,dz
}\]

to \(\fitinline{\Re z=-1/4}\). The zero residue vanishes; polynomial growth, gamma decay and small height exponent make the new integral \(\fitinline{o(1)}\), since \(\fitinline{U^{-10}}\) dominates ideal counting. On the absolute line the unit contributes \(\fitinline{1+o(1)}\) and \(\fitinline{\mu*1=\mathbf1_1}\) cancels \(\fitinline{1<q_{lk}\le U^{t_0}}\). The remaining sum has magnitude \(\fitinline{1+o(1)}\) and product norm \(\fitinline{\gg U^{t_0}}\). Truncate smoothly at \(\fitinline{U^{41}}\) with error \(\fitinline{o(1)}\), then dyadically partition and Fourier-separate the product weight to select one pair/frequency with

\[\fitmath{
|M_u(U^r)S_u(U^m)|^2
 \gg U^{\delta(r+m)-O(\epsilon_0)}.
}\]

Both factors inherit the same character, real power and height from the convolution at \(\fitinline{\rho}\). Equation \protect\hyperlink{dr:saturation-upper}{(207)} and \(\fitinline{\delta\ge1/50}\) force \(\fitinline{m\le1/2+O(\epsilon_0)}\); division by each upper bound gives both lower bounds. The inverse cutoff \(\fitinline{V_\le(U^ry/U^{t_0})}\) has uniform annular transition derivatives. Dyads and rapidly decaying separating-frequency tails fit the loss. \ensuremath{\square}

\hypertarget{dr:count}{}

\hypertarget{b.6-actual-amplitude-selection-and-source-transport}{%
\subsubsection{B.6 Actual amplitude selection and source transport}\label{b.6-actual-amplitude-selection-and-source-transport}}

At row scale \(\fitinline{U=Z^d}\), slot lengths are \(\fitinline{z_i=\ell_i/d}\). Bin squared normalized amplitudes as \(\fitinline{|Q_{u,i}|^2\asymp U^{2z_i g_i}}\), \(\fitinline{0\le g_i\le\delta/2}\); lower-than-first and error slots have \(\fitinline{g_i=0}\). Fix the amplitude mesh within the small-power loss; put \(\fitinline{L=\ell/d}\) and

\[\fitmath{
q=L^{-1}\sum_i z_i g_i=\delta p,\qquad 0\le p\le1/2.
}\]

Order \(\fitinline{g_i}\) decreasingly. Their step-function integral over initial length \(\fitinline{z}\) is \(\fitinline{\ge qz}\); discard the last fractional slot to get a genuine subset satisfying

\hypertarget{dr:greedy}{}

\[\fitmath{
\prod_{i\in I}|Q_{u,i}|^2
 \ge U^{2qz-\delta\max_i z_i-o(1)}.
 \qquad\text{(209)}
}\]

This charges one mesh. Whole-polynomial conjugation/fixed-character multiplication changes physical \(\fitinline{\overline{\psi_u}}\) to witness orientation with row-independent coefficients.

We use Lemmas~\protect\hyperlink{mf:inverse}{A.7.3} and \protect\hyperlink{mf:plain}{A.7.4} in their precise coefficient-specific forms

\hypertarget{dr:MI}{} \hypertarget{dr:MP}{}

\[\fitmath{
\begin{aligned}
 r+2z<1,\quad 2r+8z<3
 &\ \Longrightarrow\ 
 \sum_{q_u\asymp U}|M_u(U^r)Q_{u,I}|^2\ll U^{1+\epsilon},
 \qquad\text{(210)}\\
 2m+6\kappa z<1
 &\ \Longrightarrow\ 
 \sum_{q_u\asymp U}^{\rm eligible}|S_u(U^m)|^4|Q_{u,I}|^2
 \ll U^{1+\epsilon},\qquad \kappa=2\beta-1.
 \qquad\text{(211)}
\end{aligned}
}\]

Explicitly, at \(\fitinline{z=z_0+ib_i}\) the physical main slot has modulus

\[\fitmath{
Q^-_{u,i}=P_i^{-1/2}\sum_{p\notin S}
 w(q_p/P_i)(q_p/P_i)^{-33/50+ib_i}
 \overline{\psi_u(p)}.
}\]

Extracting \(\fitinline{P_i^{z_0-1/2+ib_i}}\) from \protect\hyperlink{dr:Delta}{(196)} gives this polynomial up to unit phase. Since \(\fitinline{w}\) is real, \(\fitinline{Q^+_{u,i}=\overline{Q^-_{u,i}}}\) has the same modulus and height \(\fitinline{-b_i}\).

Apply the following physical transport to the same-zero presentations after fixing target \(\fitinline{\eta}\) and finite \(\fitinline{\Theta}\). Whole conjugation when \(\fitinline{\varepsilon=-1}\) gives both witnesses character \(\fitinline{a_j\psi_u}\), where \(\fitinline{a_j=\eta_j}\) or \(\fitinline{\overline{\eta_j}}\), preserving coefficients, profile, height and zeros.

Induce all \(\fitinline{a_j}\) to one \(\fitinline{\mathfrak q}\) containing every \(\fitinline{S}\)-prime and divisible by \(\fitinline{72\mathcal O}\) and their conductors. Use full residue-unit group \(\fitinline{H}\), \(\fitinline{Q=\mathfrak q}\), anchors \(\fitinline{\eta_0=\eta=a_j}\), and trivial \(\fitinline{\vartheta_i}\). Then \(\fitinline{Q'=Q\cap\mathfrak q_{\eta_0}=\mathfrak q\subseteq72\mathcal O}\), relative coefficient \(\fitinline{\mathbf1_{\mathfrak q}\overline{a_j}}\), and on good primes

\[\fitmath{
(a_j(p)\psi_u(p))\overline{a_j(p)}=\psi_u(p).
}\]

Nonunits retain zeros; the identity pool is the probe's full good-prime pool.

Take \(\fitinline{v_i=-b_i-t}\). Source slots are \(\fitinline{P_i^{it}Q^+_{u,i}}\) at shift \(\fitinline{33/50}\); its two plains are the transformed copies of \(\fitinline{S_u}\). Energy is exactly \(\fitinline{|S_u|^4|Q_{u,I}|^2}\) in both orientations. Growing nonunit sixth-free rows have a conductor prime outside \(\fitinline{\mathfrak q}\) and satisfy literal keep; units/small exceptions are paid below. Intersect marked HIGH and analytic meshes before slots, then maximize finite dictionary height degrees and union seminorms. A sufficiently small target-dependent height exponent works; no common degree for infinitely many targets is required.

The moment-to-count conversion is now exactly the one in \S{}7. For the plain witness use the complete source's actual clipped width \(\fitinline{\min\{z,(B-2m)/L\}}\), with its cap, witness interval and admission guards (C7). The source-domain inclusion of \S{}6 supplies every inherited face without a prior decimal endpoint. The short inverse, long inverse, plain and natural-tail cases, including all zero-width fallbacks, are proved there. The physical supply condition \(\fitinline{\ell/d_+>1/5}\) dominates both the inverse selection \(\fitinline{(1-r)/2<7/37}\) for \(\fitinline{r>23/37}\) and every actual count width \(\fitinline{z<7/37}\). Thus every selected subset comes from the actual original prime pool. The natural count and its feedback are (C1).

\hypertarget{b.7-weighted-numerator-and-conductor-allocation}{%
\subsubsection{B.7 Weighted numerator and conductor allocation}\label{b.7-weighted-numerator-and-conductor-allocation}}

For each main/error term of \protect\hyperlink{dr:Hcomp}{(195)}, define its actual numerator by full \(\fitinline{w}\) inversion:

\hypertarget{dr:N}{}

\[\fitmath{
\mathcal N_u(Y;s,z)=\frac1{2\pi i}\int_{(3)}Y^{w-1}
 \widehat W_1(w)L_K^S(w,\psi_u)
 \mathcal H^{\rm term}_{\eta,u,Z}(s,w,z)\,dw.
 \qquad\text{(216)}
}\]

Bound the \(\fitinline{w}\)-independent denominator afterwards. The following coefficient calculation admits a plain moment.

Write \(\fitinline{t=Q^{-w}}\). Away from the row let \(\fitinline{v=\chi_p(u)}\) and \(\fitinline{d_p=\eta(p)v^{-1}Q^{-s}}\). Direct expansion of \protect\hyperlink{dr:Euler-table}{(193)} gives \(\fitinline{H_p=H_{p,0}+H_{p,1}t+H_{p,2}t^2}\), where

\hypertarget{dr:Hcoeff}{}

\[\fitmath{
\begin{aligned}
 H_{p,0}&=\frac{1-R_p/Q-d_p(1-V)}{(1-R_p)(1-d_p)},\\
 H_{p,1}&=\frac{\eta(p)Q^{-s}(1-QV)-vV+vR_p/Q}
 {(1-R_p)(1-d_p)},\qquad\text{(217)}\\
 H_{p,2}&=\frac{v\eta(p)(Q-1)Q^{-s}V-v^2R_p(1-V)}
 {(1-R_p)(1-d_p)}.
\end{aligned}
}\]

On \(\fitinline{\Re s=a+16e,\Re z=z_0=17/50}\) these imply

\hypertarget{dr:Hcoeff-bound}{}

\[\fitmath{
|H_{p,0}-1|+\sum_{k=1}^2|H_{p,k}|Q^{k(a-1)}
 \ll Q^{-1-\epsilon_e}.
 \qquad\text{(218)}
}\]

The largest linear exponent is \(\fitinline{-1-16e}\); other candidates \(\fitinline{a-1-6z_0,2-4a-6z_0,4-6a-6z_0,-a-6z_0}\) are \(\fitinline{<-1}\) for \(\fitinline{a\ge51/100}\). For quotient-free error estimates write \(\fitinline{E_p=E_0+E_1t}\) with

\[\fitmath{
E_0=\frac{R_p}{1-R_p}\left(\frac{1-Q^{-1}}{1-V}-d_p\right),\qquad
 E_1=\frac{-\eta(p)(Q-1)Q^{-s}V/(1-V)+vR_p}{1-R_p}.
}\]

Insertion in \protect\hyperlink{dr:Delta}{(196)} proves

\hypertarget{dr:Delta-bound}{}

\[\fitmath{
\sum_{k=0}^3|\Delta_{p,k}|Q^{k(a-1)}\ll Q^{-a}.
 \qquad\text{(219)}
}\]

The constant is \(\fitinline{\ll Q^{-a}+Q^{-6z_0}+Q^{4-5a-6z_0}\ll Q^{-a}}\). Each additional \(\fitinline{t}\) weighs \(\fitinline{Q^{a-1}<1}\), and \(\fitinline{1-6z_0<-1}\).

At a row prime \(\fitinline{p^j\Vert u}\), \(\fitinline{1\le j\le5}\), write \(\fitinline{J_j=A_j+B_jt}\). For \(\fitinline{j\ge2}\) the exact coefficients are

\hypertarget{dr:rowcoeff}{}

\[\fitmath{
\begin{aligned}
 H_0&=\frac{1-R_p/Q+(1-V)A_j}{1-R_p},&
 H_1&=\frac{-\eta(p)(Q-1)Q^{-s}+(1-V)B_j}{1-R_p},\\
 G_0&=\frac{c_p\{R_p(1-Q^{-1})+(1-V)A_j\}}{1-R_p},&
 G_1&=\frac{-Q+R_p/Q+(1-V)(c_pB_j-A_j)}{1-R_p},
 \quad G_2=-H_1.\qquad\text{(220)}
\end{aligned}
}\]

For \(\fitinline{j=1}\) replace these by

\[\fitmath{
H_0=\frac{1-R_p/Q}{1-R_p},\quad
 H_1=\frac{\eta(p)Q^{-s}(1-QV)}{1-R_p},\quad
 G_0=\frac{c_pR_p(1-Q^{-1})}{1-R_p},\quad
 G_1=\frac{R_p/Q-QV}{1-R_p},\quad G_2=-H_1.
}\]

At row primes \(\fitinline{\Delta_p=G_p}\). The primitive conductor's good-prime exponent is one. Allocate \(\fitinline{\min(k,j-1)}\) powers of shift \(\fitinline{Q^k}\) to \(\fitinline{h_r\mid\Lambda_u=q_u/q_{\operatorname{rad}(u)}}\), the rest to \(\fitinline{h_0}\). With \(\fitinline{\delta=2a-1}\),

\hypertarget{dr:rowmass}{}

\[\fitmath{
\sum_k|H_{p,k}|Q^{-k/2}Q^{\delta(k-j+1)_+/2}\le1+CQ^{-c_e},
 \quad
 \sum_k|\Delta_{p,k}|Q^{-(k+1)/2}Q^{\delta(k-j+1)_+/2}\ll1.
 \qquad\text{(221)}
}\]

For \(\fitinline{j\ge2}\), \(\fitinline{G_1\asymp Q}\) cancels against \(\fitinline{Q^{-1}}\) and its shift lies in \(\fitinline{\Lambda_u}\). At \(\fitinline{j=2}\), quadratic error is \(\fitinline{Q^{1-a}Q^{-3/2}Q^{\delta/2}=Q^{-1}}\). Remaining nontail constant/linear exponents are \(\fitinline{-\delta}\) at \(\fitinline{j=2}\); \(\fitinline{-3\delta/2,-\delta}\) at \(\fitinline{j=3}\); \(\fitinline{-3\delta/2}\) at \(\fitinline{j=4}\); \(\fitinline{-5\delta/2}\) at \(\fitinline{j=5}\). The explicit \(\fitinline{j=1}\) terms give negative powers. At unramified error slots, removing normalizer \(\fitinline{P_i^{z_0-1/2}}\) leaves

\[\fitmath{
P_i^{-1/2}\sum_{q_p\asymp P_i}q_p^{-a}\ll P_i^{1/2-a}\le1.
}\]

Products of these inequalities and divisor bounds at the row primes prove the complete mass bound

\hypertarget{dr:mass}{}

\[\fitmath{
\sum_{h_0,h_r}|c_u(h_0,h_r)|h_0^{\delta/2}\ll U^\epsilon.
 \qquad\text{(222)}
}\]

Include all local main/error choices after normalizers, monomial shifts and fixed-prime restoration in \(\fitinline{c_u}\); keep moving coefficients out of seminorms.

Inverting monomial shift \(\fitinline{h=h_0h_r}\) gives \(\fitinline{Y^{-1/2}h^{-1/2}S_{\psi_u}(Y/h)}\). Expand fixed Euler deletions, reflect primitively and restore original zeros with bounded fixed Euler mass, only shortening scales. Central amplitude is preserved, giving

\[\fitmath{
\epsilon_u S_{\overline{\psi_u}}
 \left(\mathfrak c_u h/Y;W_1^\vee\right),\quad
 \mathfrak c_u=\frac{3q_{\mathfrak f_u}}{(2\pi)^2}
 \asymp_S q_{\operatorname{rad}(u)},\quad
 \widehat W_1^\vee(v)=\widehat W_1(1-v)\frac{\Gamma(v)}{\Gamma(1-v)}.
}\]

Since \(\fitinline{q_{\operatorname{rad}(u)}h_r\le q_u}\), dual scale is \(\fitinline{\ll_S(U/Y)h_0}\). Annular support forces \(\fitinline{h\le CY}\), keeping lengths bounded. Move the reflected Mellin line just left of zero across its simple gamma pole: the profile is bounded at zero and rapidly decreasing at infinity, with summable normalized annuli. Varying conductor changes scale only.

For \(\fitinline{m=\log_UY>1/2}\), \(\fitinline{h_0\asymp U^t}\), reflected length is \(\fitinline{\le U^{1-m+t+o(1)}}\). Take the common-profile scale supremum before the fourth moment. H\"{o}lder with count \(\fitinline{U^R}\) gives \(\fitinline{U^{1/4+3R/4}}\), selected primes give \(\fitinline{U^{-qz(t)/2}}\) after fourth root, and \protect\hyperlink{dr:mass}{(222)} gives \(\fitinline{U^{-\delta t/2}}\). Thus old capacities

\[\fitmath{
z_O(t)=\max\{0,\min(1,2m-1-2t)/(6\kappa)\},
 \qquad z_A=\frac{2\min(m,1)-1}{6\kappa},
}\]

gives

\hypertarget{dr:old-numerator}{}

\[\fitmath{
\sum_{u\text{ in class}}|\mathcal N_u|
 \ll Y^{-1/2}Z^{\ell(z_0-1/2)+\ell q}
 U^{1/4+3R/4-qz_A/2+\epsilon}(1+T)^A.
 \qquad\text{(223)}
}\]

Indeed \(\fitinline{z_A-z_O(t)\le t/(3\kappa)}\) and \(\fitinline{-\delta t/2+qt/(6\kappa)\le0}\) for \(\fitinline{q\le\delta/2}\), \(\fitinline{\kappa\ge7/10}\). Use marked moments for positive capacities and unmarked for zero. Bounded-at-zero profiles assign subunit scales length zero.

\hypertarget{dr:stage60}{}

\hypertarget{b.8-every-remaining-row-and-contour}{%
\subsubsection{B.8 Every remaining row and contour}\label{b.8-every-remaining-row-and-contour}}

For \(\fitinline{U=Z^d}\), \(\fitinline{d_0=1/200\le d\le d_+}\), bin before contours. Move \(\fitinline{z}\) to \(\fitinline{z_0=17/50}\) with \(\fitinline{s,w}\) absolute, then \(\fitinline{s}\) to \(\fitinline{\beta+20e}\) with \(\fitinline{w=3}\). Truncate while reciprocal control is global; move only retained segments to \(\fitinline{a+16e}\), buffered by \protect\hyperlink{dr:bins}{(205)}, leaving omitted segments global. The path lies in \protect\hyperlink{dr:box}{(197)}, crossing neither a nonprincipal numerator pole nor a \(\fitinline{z}\) pole. Expand \protect\hyperlink{dr:Delta}{(196)}, fix amplitudes and invert \(\fitinline{w}\) fully on its original absolute line as in \protect\hyperlink{dr:N}{(216)}.

For a physical capacity \(\fitinline{C(d)}\), the complete numerator calculation of B.7 and \S{}8 gives the factor \[\fitmath{
Y^{-1/2}Z^{\ell(z_0-1/2)+\ell\delta p}
U^{1/4+3R_*/4-\delta p C(d)/(2d)+\epsilon}(1+T)^J.
}\] Multiplying by the row weight and geometry of (192) gives precisely \(\fitinline{E_\kappa-k}\), where \(\fitinline{E_\kappa}\) is (GT9). In particular, \(\fitinline{h(z_0-1/6)=13h/75}\); the fourth-root H\"{o}lder factor is \(\fitinline{d(1/4+3R_*/4)}\). This calculation retains every local main/error choice already included in (222).

For floor \(\fitinline{a=51/100}\), trivial row count and buffered functional equation give numerator \(\fitinline{U^{a-1/2+O(e)}}\), with each slot \(\fitinline{g_i\le1/100}\) by \protect\hyperlink{dr:prime-upper}{(208)}. The unnormalized exponent is

\hypertarget{dr:floor}{}

\[\fitmath{
F_{\rm fl}(d)=
 \frac4{25}x+\left(\frac{51}{100}-\frac{33}{50}\right)
 -\frac{51}{100}y+\frac{67}{100}d-\frac3{20}\ell.
 \qquad\text{(234)}
}\]

The coefficient \(\fitinline{67/100=1+(a-1/2)-z_0}\) comprises count, reflected numerator and row weight; \(\fitinline{-3\ell/20=(z_0-1/2+1/100)\ell}\) is the prime product. Fixed cube/ray subfamilies are subsets of this positive sum with the same natural-zero bound. The positive \(\fitinline{d}\) slope reduces this full floor cost to \(\fitinline{d_+}\); the actual margin and loss budget are (GT10), (GT19)--(GT21).

For nonprincipal \(\fitinline{q_u\le Z^{d_0}}\), use global \(\fitinline{(\beta+e,1/2,z_0)}\). Functional equation/convexity give \(\fitinline{L_K^S(1/2+it,\psi_u)\ll U^{3/5+\epsilon}(1+|t|)^A}\). The quotient-free bound after \protect\hyperlink{dr:box}{(197)}, ideal counting and \(\fitinline{q_u^{-z_0}}\) give relative exponent

\hypertarget{dr:small}{}

\[\fitmath{
\frac{13h}{75}-\frac y2+\frac{63}{50}d_0+O(e)+\epsilon
.
 \qquad\text{(235)}
}\]

This includes five nonprincipal units using fixed lower \(\fitinline{U}\), with no numerator pole.

For \(\fitinline{q_u>Z^{h+\zeta}}\) take \(\fitinline{(\Re s,\Re w,\Re z)=(2,2,v)}\), \(\fitinline{v>2}\). Absolute Euler bounds and \protect\hyperlink{dr:Hcomp}{(195)} give

\hypertarget{dr:large}{}

\[\fitmath{
O\left(Z^{x/2+1+y+(h+\zeta)(1+\epsilon)-\zeta v+\epsilon}\right)
 \qquad\text{(236)}
}\]

Choose \(\fitinline{v}\) sufficiently large after fixed \(\fitinline{\zeta}\) for arbitrary power saving.

For \(\fitinline{u=1}\), move to \(\fitinline{(\beta+e,1+e,1/6+e)}\), then \(\fitinline{w=3/4}\), crossing one simple numerator pole; in its residue move \(\fitinline{z=1/8}\) across the simple \(\fitinline{1/6}\) pole of \(\fitinline{\zeta_K^S(6z)}\). The principal box and \(\fitinline{\mathscr M}\) are holomorphic otherwise. The double residue is \protect\hyperlink{dr:J-f}{(203)} by \protect\hyperlink{dr:prime-residue}{(202)}. These original contours save \(\fitinline{y/4,h/24}\); \S{}8 also proves the conservative contours used in the algebraic family. Use the residue approximation only on its locus and the quotient-free product elsewhere.

\hypertarget{dr:initialization}{}

\hypertarget{b.9-the-supremum-and-original-seven-eighths-comparison}{%
\subsubsection{B.9 The supremum and original seven-eighths comparison}\label{b.9-the-supremum-and-original-seven-eighths-comparison}}

Let \(\fitinline{\mathcal H}\) be the presented finite-order Eisenstein Hecke family with prescribed Euler factors and natural zeros. Its actual zero supremum is

\[\fitmath{
\beta=\sup\left(\{\tfrac12\}\cup
 \{\Re\rho:\chi\in\mathcal H,\ \Re\rho>0,\
 L_K(\rho,\chi)=0,\ (\chi,\rho)\ne(1,1)\}\right).
}\]

Absolute convergence of the Euler product in \(\fitinline{\Re s>1}\) gives \(\fitinline{1/2\le\beta\le1}\). The principal pole at one is excluded throughout.

\textbf{Lemma B.9.1} (Primitive targets approaching the supremum). \emph{\hypertarget{init:primitive-supremum}{} The same supremum is obtained from primitive targets. If \(\fitinline{\varepsilon>0}\) and \(\fitinline{\beta-\varepsilon\ge1/2}\), there is a primitive target with a zero \(\fitinline{\rho}\) satisfying \(\fitinline{\Re\rho>\beta-\varepsilon}\).}

\emph{Proof.} A presentation induced by \(\fitinline{\chi^*}\) has \(\fitinline{L_K(s,\chi^*)}\) times finitely many \(\fitinline{1-\chi^*(\mathfrak p)(N\mathfrak p)^{-s}}\). These are one when the value is zero and otherwise vanish only on \(\fitinline{\Re s=0}\). Primitive and presented positive-half-plane zeros coincide. Supremum approximation then supplies the asserted primitive zero; no attainment is needed. \ensuremath{\square}

The original marked-moment range is

\hypertarget{init:original-moment-range}{}

\[\fitmath{
\beta\ge\frac{51}{100},\qquad \kappa\ge\frac34,
 \qquad 2\beta-1\le\kappa.
 \qquad\text{(I.1)}
}\]

Proposition \protect\hyperlink{mf:plain-original-bootstrap}{A.7.5} proves this range before initialization, retaining full row/moving width, literal slots, fixed-inducing exclusion and profile/height control. Proposition \protect\hyperlink{mf:plain-refined}{A.7.6} supplies the later refined range separately.

\textbf{Lemma B.9.2} (Noncircular moment parameter). \emph{\hypertarget{init:moment-parameter}{} In the contradiction case \(\fitinline{\beta>7/8}\), the choice \hypertarget{init:kappa}{}}

\[\fitmath{
\Delta=\beta-\frac78,\qquad
 \kappa=\frac34+2\Delta=2\beta-1
 \qquad\text{(I.2)}
}\]

\emph{satisfies \protect\hyperlink{init:original-moment-range}{(I.1)}, with \(\fitinline{3/4<\kappa\le1}\).}

\emph{Proof.} Expansion gives the identity; \(\fitinline{\Delta>0}\) and \(\fitinline{\beta\le1}\) give \(\fitinline{3/4<\kappa\le1}\). Condition \(\fitinline{2\beta-1\le\kappa}\) holds with equality before the desired conclusion. \ensuremath{\square}

The same analytic return serves later improvements. Constants may depend on the target except those explicitly common.

\textbf{Lemma B.9.3} (A common-reserve Mellin contradiction). \emph{\hypertarget{init:mellin-return}{} Let \(\fitinline{1/2\le\theta<\beta}\), and suppose \(\fitinline{k_0}\) is real. Suppose there are common numbers \(\fitinline{\omega,\sigma_0>0}\) with \(\fitinline{\omega<\beta-\theta}\) such that each primitive target has a probe \(\fitinline{J_\eta}\) and a function \(\fitinline{H_\eta}\), holomorphic on \(\fitinline{\Re s>\theta}\), satisfying \hypertarget{init:signal}{}}

\[\fitmath{
|H_\eta(s)-1|\le\frac12,
 \qquad
 f_\eta(Z)=\frac1{2\pi i}\int_{(c)}
 Z^{s-k_0}e^{(s-5/6)^2}
 \frac{H_\eta(s)}{L_K(s,\eta)}\,ds,
 \qquad\text{(I.3)}
}\]

\emph{on an absolute line \(\fitinline{c>1}\), with the stated vertical convergence. Assume that \(\fitinline{f_\eta}\) decays faster than every power at zero and that, for large \(\fitinline{Z}\), \hypertarget{init:two-probe-bounds}{}}

\[\fitmath{
J_\eta(Z)\ll_\eta Z^{\theta-k_0+\omega},\qquad
 J_\eta(Z)-f_\eta(Z)\ll_\eta Z^{\beta-k_0-\sigma_0}.
 \qquad\text{(I.4)}
}\]

\emph{Then these hypotheses contradict the definition of \(\fitinline{\beta}\).}

\emph{Proof.} Put \(\fitinline{m=\min(\beta-\theta-\omega,\sigma_0)>0}\). The two bounds on the same probe give \(\fitinline{f_\eta(Z)\ll_\eta Z^{\beta-k_0-m}}\), so

\[\fitmath{
F_\eta(s)=\int_0^\infty f_\eta(Z)Z^{k_0-s}\frac{dZ}{Z}
}\]

is locally uniformly holomorphic for \(\fitinline{\Re s>\beta-m}\): use this bound at infinity and arbitrary decay at zero, also after logarithmic differentiation. Absolute-line Fourier--Mellin inversion identifies it with the Gaussian quotient in \protect\hyperlink{init:signal}{(I.3)}; meromorphic identity continues throughout this connected half-plane.

Since \(\fitinline{\beta-m\ge\theta+\omega>\theta}\), the numerator is holomorphic and nonzero (\(\fitinline{|H_\eta|\ge1/2}\)). Any denominator zero, of any multiplicity, would give an impossible pole. Euler deletions change no such zero by Lemma \protect\hyperlink{init:primitive-supremum}{B.9.1}; the principal pole creates no reciprocal pole.

Thus every primitive target is zero-free for \(\fitinline{\Re s>\beta-m}\), with common \(\fitinline{m}\) despite target-dependent constants. Lemma \protect\hyperlink{init:primitive-supremum}{B.9.1} at \(\fitinline{\varepsilon=m/2}\) supplies a forbidden zero, since \(\fitinline{\beta-m/2>\theta\ge1/2}\). \ensuremath{\square}

\hypertarget{init:scalar-section}{}

\hypertarget{b.9.1-the-complete-original-scalar-comparison}{%
\paragraph{B.9.1 The complete original scalar comparison}\label{b.9.1-the-complete-original-scalar-comparison}}

Let normalized amplitude \(\fitinline{p\in[0,1/2]}\) be distinct from physical length \(\fitinline{x}\), and set \(\fitinline{q=\delta p}\), \(\fitinline{\delta=2a-1}\). Before reserves, detector estimates are

\[\fitmath{
I=1-\delta[p+(1-p)r],\qquad
 P=1-\delta\left[\frac{4p}{9}
                  +\left(2-\frac{8p}{9}\right)(t-r)\right].
}\]

For \(\fitinline{1\le t\le3/2}\) they cross at

\[\fitmath{
r_c=\frac{(2-8p/9)t-5p/9}{3-17p/9},
 \quad \frac{23}{37}\le r_c\le1,
 \quad\frac13\le t-r_c\le\frac12.
}\]

Clear the positive denominator and check affine \(\fitinline{t}\) endpoints and \(\fitinline{p}\) endpoints. Use plain below the crossing, inverse above; their envelope is

\[\fitmath{
S=1-\delta+\delta
 \frac{(2-8p/9)(1-p)}{3-17p/9}\left(\frac32-t\right).
}\]

The remaining long inverse lengths have exponent \(\fitinline{L=1-\delta+(5/6-\delta)(t-1)}\).

Actual prime capacity is \(\fitinline{6\kappa=9/2+12\Delta}\). For \(\fitinline{1/3-\gamma\le m\le1/2}\), replacing \(\fitinline{9/2}\) costs

\[\fitmath{
2q(1-2m)\left(\frac29-\frac1{9/2+12\Delta}\right)
 \le\frac{\Delta}{4}+2\gamma.
}\]

Clear the denominator using \(\fitinline{q\le1/2}\), \(\fitinline{0\le\Delta\le1/8}\); discrepancy \(\fitinline{m\ge t-r-\gamma}\) costs another \(\fitinline{2\gamma}\). With capacity reserve \(\fitinline{\nu}\) and moment loss \(\fitinline{\varepsilon}\), count is at most

\hypertarget{init:count-loss}{}

\[\fitmath{
\max(S,L)+\frac{\Delta}{4}+\varepsilon+6\nu+5\gamma.
 \qquad\text{(I.5)}
}\]

Unmarked bounds cover nonpositive reserved capacities within \(\fitinline{6\nu}\); long inverse covers \(\fitinline{r\ge1}\). All witness lengths are included.

Set \(\fitinline{v=1/2-p}\) and

\[\fitmath{
D=\frac{37+34v}{18},\quad P_0=\frac{7+18v+8v^2}{9},\quad
 J=(5/6-\delta)D+\delta P_0,\quad
 t_*=1+\frac{\delta P_0}{2J}.
}\]

On \(\fitinline{0\le\delta\le5/6}\), \(\fitinline{0\le v\le1/2}\), endpoint bounds give \(\fitinline{35/54\le J\le5/2}\), \(\fitinline{1\le t_*\le3/2}\), and

\[\fitmath{
S(t_*)=L(t_*)=R_*
 =1-\delta+\frac{(5/6-\delta)\delta P_0}{2J}.
}\]

The central relative exponent is

\[\fitmath{
E_*=-\frac1{48}+\frac{2\delta}{3}
       +\frac{\delta(1/2-v)}6-\frac{13}{16}(1-R_*).
}\]

Put \(\fitinline{C=51+41v}\). Expansion gives the exact identity

\[\fitmath{
\begin{aligned}
10368CJ(-E_*)={}&(3+5v)\bigl[(4C\delta-79)^2+49\bigr]\\
 &+4v\left\{4C\delta\bigl[(1+3v)(15+32v)\delta+9-13v\bigr]
                       +265+3485v\right\}.
\end{aligned}
}\]

All terms are nonnegative (\(\fitinline{9-13v\ge5/2}\)). Since \(\fitinline{C\le17(3+5v)}\) and \(\fitinline{J\le5/2}\),

\hypertarget{init:square-margin}{}

\[\fitmath{
E_*\le-\frac{49}{440640}.
 \qquad\text{(I.6)}
}\]

This certifies the entire closed rectangle.

The exact geometric exponent supplied by the original probe is

\[\fitmath{
\begin{aligned}
 \mathscr E(a,d,R,q)=\frac3{16}
 &-\frac1{48}+\frac{2\delta}{3}+\frac q6
      -\frac{13}{16}(1-R)\\
 &+\left(d-\frac{13}{16}\right)
       \left(R+\frac\delta2-\frac{17}{50}\right).
\end{aligned}
}\]

Suppose the count is \(\fitinline{R\le R_*+\Delta/4+\mathcal L}\), \(\fitinline{d\le13/16+\zeta}\), and \(\fitinline{0\le R+\delta/2-17/50\le2}\). Then \protect\hyperlink{init:square-margin}{(I.6)} implies

\hypertarget{init:central-margin}{}

\[\fitmath{
\mathscr E-(3/16+\Delta)
 \le-\frac{49}{440640}-\frac{51\Delta}{64}
                  +\frac{13\mathcal L}{16}+2\zeta.
 \qquad\text{(I.7)}
}\]

Here \(\fitinline{51/64=1-(13/16)(1/4)}\) retains the capacity correction; slope bounds follow from \(\fitinline{1-\delta\le R\le139/96}\), \(\fitinline{0\le\delta\le1}\).

For \(\fitinline{5/6\le\delta\le1}\), unmarked inverse gives \(\fitinline{R\le1-\delta+\mathcal L}\); \(\fitinline{q\le\delta/2}\) yields

\[\fitmath{
\mathscr E-(3/16+\Delta)
 \le-\frac1{48}-\frac\delta{16}-\Delta
                  +\frac{13\mathcal L}{16}+2\zeta.
}\]

The floor \(\fitinline{a=51/100}\), \(\fitinline{\delta=1/50}\), \(\fitinline{R=1}\), \(\fitinline{q\le1/100}\) has relative constant \(\fitinline{-1/48+2/150+1/600=-7/1200}\). These cover bins outside the balanced rectangle.

The balanced loss is \(\fitinline{\mathcal L=159\varepsilon+\varepsilon_m+\mathrm{slotMesh}+7\nu\le1/32}\); high bins lose \(\fitinline{78\varepsilon+\varepsilon_m}\). At actual dyad \(\fitinline{v_{\rm row}}\) the mixed exponent is \(\fitinline{\mathscr E(a,v_{\rm row},R,q)+(d-v_{\rm row})R}\), adding at most \(\fitinline{3\mu_{\rm mesh}/2}\) when \(\fitinline{d-v_{\rm row}\le\mu_{\rm mesh}}\). Choose moment losses/common meshes within the fixed margin, then finite seminorm/height orders, then height exponent and independent analytic reserves, finally tail orders, so

\[\fitmath{
\begin{aligned}
 &\frac{13}{16}\mathcal L+2\zeta+\frac32\mu_{\rm mesh}
   +26e+(N+8)\varepsilon_{\rm arith}
   +\mathrm{heightLoss}+\frac{\mathrm{mesh}}6\\
 &\hspace{2em}+\mathrm{logCost}+\mathrm{heightCost}
   +\mathrm{momentCost}+\mathrm{saving}
 \le\frac{49}{440640}.
\end{aligned}
}\]

Every term is the corresponding full-probe cost. Fixed mesh/slot count makes amplitude counts constant; choosing height exponent after degree and tail order afterwards gives a common positive saving.

\textbf{Theorem B.9.4} (Initial half-plane). \emph{\hypertarget{init:seven-eighths}{} The Eisenstein family satisfies \(\fitinline{\beta\le7/8}\).}

\emph{Proof.} If \(\fitinline{\beta>7/8}\), use actual \protect\hyperlink{init:kappa}{(I.2)} and the probe/scalar chapters' original geometry

\[\fitmath{
h=\frac{13}{16},\quad \ell=\frac16,\quad
 x=\frac{17}{48},\quad y=\frac{23}{48},\quad
 k_0=\frac{11}{16}.
}\]

They provide common positive reserves \(\fitinline{\omega<\beta-7/8}\) and \(\fitinline{\sigma_0}\), and for every primitive target the exact signal \protect\hyperlink{init:signal}{(I.3)}, together with

\[\fitmath{
J_\eta(Z)\ll_\eta Z^{3/16+\omega},\qquad
 J_\eta(Z)-f_\eta(Z)
       \ll_\eta Z^{\beta-11/16-\sigma_0}.
}\]

These are \protect\hyperlink{init:two-probe-bounds}{(I.4)} with \(\fitinline{\theta=7/8}\), \(\fitinline{k_0=11/16}\). The probe chapter supplies zero-end decay, complete principal reconstruction, nonvanishing numerator and all row/contour errors. Lemma \protect\hyperlink{init:mellin-return}{B.9.3} contradicts \(\fitinline{\beta>7/8}\). \ensuremath{\square}

\hypertarget{b.9.2-realizing-the-initialization-with-the-literal-probe}{%
\paragraph{B.9.2 Realizing the initialization with the literal probe}\label{b.9.2-realizing-the-initialization-with-the-literal-probe}}

The constructions just given also explain how the common-probe continuation is used before the refined incumbents are available. At the original \(\fitinline{7/8}\) stage take

\[\fitmath{
(x,y,\ell,h,k)=
 (17/48,23/48,1/6,13/16,11/16).
}\]

The same multi-slot probe, table and nonvanishing numerator apply; marked energy gives \(\fitinline{x/2+(y-x)/12=3/16}\). Under \(\fitinline{\beta>7/8}\), use the original proved moment/count at \(\fitinline{\kappa=2\beta-1\ge3/4}\). Before the fourth-root refinement, move each nonprincipal \(\fitinline{w}\) to \(\fitinline{1-a-6e}\): reflection and \protect\hyperlink{dr:buffered}{(206)} give numerator \(\fitinline{U^{\delta/2+O(e)+\epsilon}}\) and factor \(\fitinline{Y^{-a+O(e)}}\). Main slots keep actual amplitudes; \protect\hyperlink{dr:Delta-bound}{(219)}, \protect\hyperlink{dr:rowmass}{(221)} bound error slots at amplitude zero. The buffered reciprocal is subpower. Counting \(\fitinline{U^R}\) rows with weight \(\fitinline{U^{-z_0}}\) gives

\[\fitmath{
x(1/2-z_0)+a+z_0-1-ay
 +d(R+\delta/2-z_0)+\ell(z_0-1/2+q)+O(e)+\epsilon.
}\]

Substituting the original geometry and subtracting its physical exponent \(\fitinline{3/16}\) gives exactly

\[\fitmath{
-\frac1{48}+\frac23\delta+\frac q6
 -\frac{13}{16}(1-R)
 +(d-13/16)(R+\delta/2-17/50).
}\]

The original balanced count leaves central reserve \(\fitinline{49/440640}\) and floor reserve \(\fitinline{7/1200}\). The same principal/small/large/floor/height decomposition therefore gives common \(\fitinline{\omega,\sigma_0>0}\), \(\fitinline{\omega<\beta-7/8}\), with

\[\fitmath{
J_\eta\ll_\eta Z^{3/16+\omega},\qquad
 J_\eta-f_\eta\ll_\eta Z^{\beta-11/16-\sigma_0}.
}\]

The next Mellin argument at \(\fitinline{\Delta<\min(\beta-7/8-\omega,\sigma_0)}\) proves initialization from these original inputs.

\hypertarget{b.10-every-presented-ordinary-character}{%
\subsubsection{B.10 Every presented ordinary character}\label{b.10-every-presented-ordinary-character}}

Let a Hecke character in modulus \(\fitinline{\mathfrak m}\) be induced by the primitive \(\fitinline{\eta^*}\) of conductor \(\fitinline{\mathfrak f\mid\mathfrak m}\). Absolute Euler products on \(\fitinline{\Re s>1}\) give

\[\fitmath{
L_{\mathfrak m}(s,\eta)=L(s,\eta^*)
 \prod_{\substack{p\mid\mathfrak m\\p\nmid\mathfrak f}}
                 (1-\eta^*(p)q_p^{-s}).
}\]

Continuation extends this identity away from the principal pole. Each deleted factor is nonzero for \(\fitinline{\Re s>0}\) since \(\fitinline{|\eta^*(p)q_p^{-s}|<1}\). Thus primitive reduction and fixed-\(\fitinline{S}\) restoration preserve zeros. Reflected numerator bounds \protect\hyperlink{dr:rowmass}{(221)}--\protect\hyperlink{dr:mass}{(222)} still retain repeated row powers/holes.

For any \(\fitinline{q\ge1}\) and complex Dirichlet \(\fitinline{\chi\bmod q}\), norm pullback \(\fitinline{\eta_\chi(\mathfrak a)=\chi(N\mathfrak a)}\) has finite order in modulus \(\fitinline{q\mathcal O}\) and is trivial on units. Let \(\fitinline{\chi_{-3}}\) be the field character and present \(\fitinline{\chi\chi_{-3}}\) at a common modulus. Euler factors give

\hypertarget{dr:norm-lift}{}

\[\fitmath{
L_K(s,\eta_\chi)=L(s,\chi)L(s,\chi\chi_{-3}).
 \qquad\text{(240)}
}\]

Split primes give two norm-\(\fitinline{p}\) factors with \(\fitinline{\chi_{-3}(p)=1}\); inert primes give a norm-\(\fitinline{p^2}\) factor equal to the product with signs \(\fitinline{\pm}\). Prime three gives one factor and \(\fitinline{\chi_{-3}(3)=0}\). At \(\fitinline{p\mid q}\) all corresponding factors are deleted. This proves \protect\hyperlink{dr:norm-lift}{(240)} absolutely, then by continuation.

For \(\fitinline{\theta<\Re s<1}\), both Dirichlet factors are holomorphic, so a zero contradicts Hecke nonvanishing. For \(\fitinline{\Re s\ge1}\) use ordinary Dirichlet nonvanishing, excepting the principal pole; treat nonprincipal \(\fitinline{s=1}\) directly even if its companion is principal. Thus every presented modulus and finite height satisfy

\[\fitmath{
\Re s>\theta,\qquad
 \neg(\chi=1\ \hbox{and}\ s=1)
 \quad\Longrightarrow\quad L(s,\chi)\ne0.
}\]

Modulus-one principal gives zeta. This transfer asserts no boundary-line or growing-modulus uniformity.

\emph{Attribution.} The probe, theta-reflection and initial source architecture originate in OpenAI's proof; the standalone same-zero, signed-Gram and continuation exposition is due to Hailey Collet. The weighted numerator is Akash Levy's refinement. The affine detector is due to DaysSky, with related count arithmetic by liqsweep/blueberryvertigo; the PR60 and PR62 width-sensitive affine refinements are Zhongpai Gao's. The marked asymmetric counts and common-radical allocations used by their moments are due to timbalud. This reconstruction preserves the supplied Apache-2.0 attribution and upstream notices; the provenance ledger distinguishes the established mathematics, its exposition, and separate formal replay or independent-review evidence.

\hypertarget{bm:section}{}

\hypertarget{appendix-c.-complete-restoration-and-source-assembly}{%
\subsection{Appendix C. Complete restoration and source assembly}\label{appendix-c.-complete-restoration-and-source-assembly}}

The proofs in this appendix discharge the specialized source obligations used in \S{}4. Their order follows the operations: finite projectors and the whole-column transfer; a reference bound; the first exact physical source; masked packing and cube-free restoration; the joint divisor budget; then the diagonal and complete generic assembly. The selected-norm foundation of \S{}3 only enters the new leaf price. It does not replace any of these complementary arguments.

Numbered equations in each C.1--C.10 module are local to that module; an external citation names the module, for example C.8(16). The older named equation families are preserved where they are already unambiguous; the correspondence table at the end of the manuscript records their local locations. No historical cycle or catalogue fetch is required to follow these proofs.

\hypertarget{c.1.-finite-projectors-and-normalized-restoring-mass}{%
\subsubsection{C.1. Finite projectors and normalized restoring mass}\label{c.1.-finite-projectors-and-normalized-restoring-mass}}

This section gives the local identities and absolute coefficient bounds used in the fourth-power source construction. It uses the Eisenstein ideal conventions of Appendix A.1--A.2. A supported ideal is nonzero and prime to the primes above 2 and 3; a supported prime has norm at least 7. All characters retain their zero values at nonunits. The results concern the actual common coefficient families; none asserts a bound for arbitrary row-dependent coefficients.

The restriction on an ordinary pair is \(\fitinline{v_{\mathfrak p}(n_1)+v_{\mathfrak p}(n_2)<5}\) at every supported prime. This is a restriction on the product, not two separate fifth-free conditions. A live marked prime is also excluded from both ordinary factors. Initially the live marks are distinct. Repeated marks will be restored by the exact equality-pattern decomposition below.

\hypertarget{local-identities-before-norms}{%
\subparagraph{1. Local identities before norms}\label{local-identities-before-norms}}

In the formal series ring in two variables put \[\fitmath{
 O=(1-x)^{-1}(1-y)^{-1},\qquad B=(1-x)(1-y),\qquad
 T_4=\sum_{i+j<5}x^iy^j.
}\] Then \[\fitmath{
 H=BT_4=1-\sum_{i+j=5}x^iy^j+xy\sum_{i+j=4}x^iy^j,
 \qquad R=H^{-1}.
}\tag{VII.2.1}\] The inverse exists formally because the constant coefficient of H is one. The nonconstant part has six terms of degree five and five of degree six. Direct multiplication gives \[\fitmath{
 HO=T_4,\qquad RT_4=O.
}\tag{VII.2.2}\] Introduce a mark variable z with \(\fitinline{z^2=0}\). The unrestricted local state is \(\fitinline{O+zO}\): a marked prime can occur in the ordinary factors. The restricted state is \(\fitinline{T_4+z}\): when the prime is marked, both ordinary valuations are zero. The two exact identities are \[\fitmath{
 [H+z(B-H)](O+zO)=T_4+z,
}\] \[\fitmath{
 [R+z(H-B)R^2](T_4+z)=O+zO.
}\tag{VII.2.3}\] These follow from BO=1 and (VII.2.2). In particular the marked coefficient in the restricted state is 1, not T\_4. Replacing it by T\_4 gives the wrong coefficient already at the monomial \(\fitinline{x^2z}\) in the inverse identity.

For finitely many distinct marked primes use separate nilpotent variables, and take the coefficient of their product after multiplying (VII.2.3) over a finite enclosing prime set. This is a finite ideal-divisor convolution at every fixed ordinary pair. Formal series cause no interchange problem: only finitely many terms can contribute to any fixed valuation vector.

Here is the reindexing that will be used analytically. Let D be a finite ordinary-pair set, k an extraction coefficient and r a residual coefficient. Complete multiplicativity, including nonunit zeros, gives \[\fitmath{
 \sum_{n\in D}\ \sum_{d_1\mid n_1,d_2\mid n_2}
 k(d)r(n/d)\chi(n_1n_2)f(n)
 =\sum_d k(d)\chi(d_1d_2)
       \sum_{a:\,da\in D}r(a)\chi(a_1a_2)f(da).
}\tag{VII.2.4}\] The outer sum is over the divisor pool of D. No character value is divided out. For a normalized row with ordinary scales X\_1,X\_2, the residual scales are exactly \(\fitinline{X_1/Nd_1,X_2/Nd_2}\), and the outside scalar is \[\fitmath{
 k(d)\chi(d_1d_2)/\sqrt{Nd_1Nd_2}.
}\tag{VII.2.5}\] If two rectangles are centered, the identical scalar and identical extraction must be used for both. Their original product volumes must be equal when a later cancellation invokes that equality.

\hypertarget{local-absolute-masses}{%
\subparagraph{2. Local absolute masses}\label{local-absolute-masses}}

For a series A write \[\fitmath{
 \|A\|_{q;s,t}=\sum_{i,j\ge0}|[x^iy^j]A|q^{-si-tj},
 \qquad q=N\mathfrak p.
}\] This positive norm is submultiplicative whenever the displayed sums converge, by the coefficient convolution and Tonelli's theorem. If \(\fitinline{s,t\ge\delta\ge0}\), (VII.2.1) gives \[\fitmath{
 \|H-1\|_{q;s,t}\le11q^{-5\delta},\quad
 \|H\|_{q;s,t}\le1+11q^{-5\delta},\quad
 \|B-H\|_{q;s,t}\le14q^{-\delta}.
}\tag{VII.2.6}\] The last inequality uses the two degree-one terms and one degree-two term of B-1, and the eleven terms of H-1.

At s=t=1/2 the quantity \(\fitinline{r=11q^{-5/2}}\) is less than 1/2 for q at least 7. The formal identity \(\fitinline{R=\sum_{k\ge0}(1-H)^k}\), whose coefficients are locally finite because H-1 starts at degree five, is therefore absolutely summable in this norm. It proves \[\fitmath{
 \|R\|\le(1-r)^{-1}\le1+2r
 \le1+400q^{-3/2},\qquad
 \|(H-B)R^2\|\le56q^{-1/2}\le10206q^{-1/2}.
}\tag{VII.2.7}\] The loose constants on the right match the imported source contracts. No analytic nonvanishing assertion for a power series is assumed to prove these estimates.

An extracted marked prime carries its own central normalization. If its norm is q, include the factor q\^{}\{-v\}. The three mass regimes needed below are consequently

\begin{longtable}[]{@{}llll@{}}
\toprule\noalign{}
Kernel & (s,t,v) & Unmarked local mass & Marked local mass \\
\midrule\noalign{}
\endhead
\bottomrule\noalign{}
\endlastfoot
Forward & (1/2,1/2,1/2) & 1+11q\^{}\{-5/2\} & at most 14/q \\
Forward, reflected & (2/3,1/3,2/3) & 1+11q\^{}\{-5/3\} & at most 14/q \\
Inverse & (1/2,1/2,1/2) & 1+400q\^{}\{-3/2\} & at most 10206/q \\
\end{longtable}

For an actual annular slot with scale P, coefficient bounded by M and prime norm q at most 2P, its central coefficient satisfies \(\fitinline{M/\sqrt P\le\sqrt2M/\sqrt q\le2M/\sqrt q}\). Thus the actual coefficient is bounded by the table times the fixed product of its slot envelopes. The character phase and natural zero stay in the exact identity and disappear only in this absolute upper bound.

\hypertarget{global-masses-and-common-hilbert-norms}{%
\subparagraph{3. Global masses and common Hilbert norms}\label{global-masses-and-common-hilbert-norms}}

The fixed-field ideal count \(\fitinline{\#\{I:NI\le X\}\ll X}\) implies \[\fitmath{
 \sum_{\mathfrak p}(N\mathfrak p)^{-\sigma}<\infty\quad(\sigma>1),
 \qquad \sum_{N\mathfrak p\le X}(N\mathfrak p)^{-1}\ll1+\log(2X).
}\tag{VII.2.8}\] The first statement follows by dyadic summation, and the second by the same argument or partial summation; replacing primes by all ideals is enough. No prime ideal theorem is needed.

For any of the table's regimes and any fixed distinct marked-prime tuple, sum the absolute coefficients over an arbitrary finite divisor cutoff. Enlarge the positive sum to the full Euler product. Its ordinary factors are bounded by \(\fitinline{\exp(C\sum_{\mathfrak p}(N\mathfrak p)^{-\sigma})}\), and each marked prime contributes at most K/q. Summing over all permitted distinct marked tuples is bounded by the unrestricted product of the harmonic sums in (VII.2.8). For a fixed number of marks and prime pools of polynomial size in Z this is \(\fitinline{O_\epsilon(Z^\epsilon)}\). Constants are chosen after the fixed slot count, envelope and polynomial cap, and before the actual row character, coefficients, divisor cutoffs and prime sets. If the slot envelope is variable, its displayed product is retained and is not hidden in Z\^{}epsilon.

The extraction scalars in (VII.2.5) may depend on the physical row u through its character and frozen-mark phases. The positive coefficient masses supply row-independent envelopes w\_d\textgreater=0 such that \(\fitinline{|c_d(u)|\le w_d}\). For a Hilbert space defined by the original nonnegative row measure, pointwise multiplication gives \[\fitmath{
 \|c_d(\cdot)F_d(\cdot)\|\le w_d\|F_d\|,\qquad
 \left\|\sum_d c_d(\cdot)F_d(\cdot)\right\|
 \le\sum_d w_d\|F_d\|.
}\tag{VII.2.9}\] This is the triangle inequality applied after retaining the actual row-dependent phases and zeros. It does not regard them as common complex scalars or divide by their possibly zero values. The coefficient family, prime failures, punctures and profiles remain inside each residual row norm. A mass estimate alone does not establish the residual energy bound.

More precisely, let Q\_f be the product of the norms of all frozen marks in an extracted term, and put \[\fitmath{
 r_d=\left(\frac{Nd_2}{Nd_1Q_f}\right)^{1/6}.
}\] Multiplying its central weight by r\_d changes the two ordinary exponents from (1/2,1/2) to (2/3,1/3), and each frozen-mark exponent from 1/2 to 2/3. Thus the first two rows of the table prove \(\fitinline{\sum_d w_d(1+r_d)\ll_\epsilon Z^\epsilon}\). If the actual residual energies are bounded by \(\fitinline{E(3+r_d^2)}\), then \(\fitinline{\sqrt{3+r_d^2}\le2(1+r_d)}\), so (VII.2.9) proves energy at most \(\fitinline{C_\epsilon Z^\epsilon E}\), after choosing half the final epsilon in the mass estimate before squaring. The whole frozen product Q\_f is needed in this identity. Dropping some frozen factors changes the required mass.

\hypertarget{common-punctures-and-collisions-of-marked-primes}{%
\subparagraph{4. Common punctures and collisions of marked primes}\label{common-punctures-and-collisions-of-marked-primes}}

At a prime that is to be excluded from both ordinary factors the local correction to T\_4 is \[\fitmath{
 J=BR=1/T_4,\qquad JT_4=1.
}\tag{VII.2.10}\] For s,t at least 1/3 its absolute mass is uniformly bounded. Indeed z=q\^{}\{-1/3\} is at most 11/20, and \[\fitmath{
 \|H-1\|_{q;s,t}\le6z^5+5z^6
 \le28183925/64000000<1/2.
}\] Hence \(\fitinline{\|R\|\le2}\), \(\fitinline{\|B\|\le(1+z)^2<3}\), and \(\fitinline{\|J\|\le6\le576}\). If all common punctures divide a nonzero ideal F with \(\fitinline{NF\le Z^{B_0}}\), their entire mass is at most \(\fitinline{576^{\omega(F)}\ll_{\epsilon,B_0} Z^\epsilon}\). For completeness, the divisor estimate here follows by choosing a norm threshold so that \(\fitinline{576\le q^\eta}\) above it; the finitely many smaller prime ideals contribute a fixed constant. Choose eta after B\_0 and epsilon. When B\_0=0, F is the unit ideal and the assertion is immediate.

The same argument covers the central and reflected ordinary weights, including \(\fitinline{(Nd_2/Nd_1)^{1/6}}\), for these common puncture kernels. The remaining live primes must be disjoint from the punctures. This disjointness is retained in the exact finite expansion, not discarded after a probability estimate.

To restore repeated marks, partition the fixed r slot indices by equality of their actual prime values. There are at most \(\fitinline{2^{r^2}}\) possible partitions. For a block J of at least two indices suppose its nonnegative coefficient envelopes satisfy \(\fitinline{v_i(p)\le B}\), \(\fitinline{\sum_pv_i(p)^2\le L B^2}\), with B,L at least one. Cauchy on two indices and the supremum bound on the others give \[\fitmath{
 \sum_p\prod_{i\in J}v_i(p)\le L B^{|J|}.
}\tag{VII.2.11}\] Freeze each such block at its actual prime. The assignment mass over all non-singleton blocks is at most \(\fitinline{(LB^r)^r}\), a convenient loose bound. Singleton marks retain their phases and are constrained to avoid the frozen primes and one another. They are precisely the distinct live family with a common puncture. If every such residual norm is at most \(\fitinline{\sqrt E}\), uniformly in the frozen choices, weighted triangle and (VII.2.11) bound the full energy by \[\fitmath{
 [2^{r^2}(LB^r)^r]^2E.
}\tag{VII.2.12}\] There is no factor equal to the number of possible prime assignments. In annular applications L and B are fixed or explicitly paid subpower envelopes. Any polynomial dependence of these quantities must be charged before using (VII.2.12).

The puncture and collision bounds are coefficient bounds. All analytic profile, height and shortened-scale conditions still belong to the residual source theorem. In particular this section has not replaced a coupled moment by a product of separate moments.

\hypertarget{c.2.-whole-column-fourth-power-transfer}{%
\subsubsection{C.2. Whole-column fourth-power transfer}\label{c.2.-whole-column-fourth-power-transfer}}

Use the supported primary-generator and Gauss conventions of Chapters 1--2. For a supported nonzero modulus a, write chi\_a for the sextic character with its natural nonunit zeros and \[\fitmath{
 \mathfrak g_a(k)=(Na)^{-1/2}\sum_{x\bmod a}\chi_a(x)e(kx/a).
}\] The unit modulus has Gauss sum one. If a is fifth-free, meaning each of its prime valuations is at most four, then for every supported ideal D with its fixed primary generator d\_D, and every frequency k, \[\fitmath{
 \mathfrak g_a(d_D^4k)=\chi_a(d_D)^2\mathfrak g_a(k).
}\tag{VII.3.1}\]

If D is coprime to a, multiplication by d\_D\^{}4 permutes the finite residue ring. Changing variables gives the multiplier chi\_a(d\_D)\^{}(-4), equal to chi\_a(d\_D)\^{}2 because its sixth power is one at units. If D shares a prime p with a, let p\^{}j exactly divide a. Then 1\textless=j\textless=4 and p\^{}j divides d\_D\^{}4k. In the exact CRT factorization, the additive character in the local p\^{}j factor is therefore trivial. Its multiplicative character is the nonprincipal j-th power of the primitive sextic residue character, lifted from the residue field. The complete local sum vanishes. Its CRT unit multipliers do not change that divisibility. Thus the left side is zero. The right side is also zero by chi\_a(d\_D)=0. This includes k=0. No character value at a nonunit has been inverted.

Let finitely many supported nonzero columns a\_i and coefficients b\_i be given, with every a\_i for which b\_i is nonzero fifth-free. Put \[\fitmath{
 P(k)=\sum_i b_i\mathfrak g_{a_i}(k).
}\] The entire coefficient vector is common to all D and k. Summing (VII.3.1), before taking absolute values, gives \[\fitmath{
 \sum_i b_i\chi_{a_i}(d_D)^2\mathfrak g_{a_i}(k)=P(d_D^4k).
}\tag{VII.3.2}\] This preserves every cross-column phase in one polynomial. It does not assert that cubic twists disappear after squaring.

Here is the exact positive-family bound. Let D range over a finite set of supported nonzero ideals with ND\textless=2W, and k over a finite set of nonzero elements with Nk\textless=K, where W,K\textgreater0. Suppose \(\fitinline{0\le w_D\le A}\) and \(\fitinline{0\le\rho(D,k)\le1}\). For every epsilon\textgreater0, \[\fitmath{
\begin{aligned}
&\frac1W\sum_Dw_D\sum_k\rho(D,k)
 \left|\sum_i b_i\chi_{a_i}(d_D)^2\mathfrak g_{a_i}(k)\right|^2\\
&\qquad\le
 \frac{C_\epsilon A(16KW^4)^\epsilon}{W}
 \sum_{r\in\{d_D^4k\}}|P(r)|^2.
\end{aligned}
}\tag{VII.3.3}\] To prove it, substitute (VII.3.2) and regroup the nonnegative sum by r=d\_D\^{}4k. Since r is nonzero, D divides the ideal (r). For each ideal D its primary generator is fixed, so k is uniquely determined. A fiber therefore has at most tau((r)) members. The elementary ideal-divisor bound gives tau((r))\textless=C\_epsilon (Nr)\^{}epsilon, and \(\fitinline{Nr=(ND)^4Nk\le16KW^4}\). Its total weight is at most A times this fiber bound. Sum over r and divide by W. If the indexing family is empty the inequality is immediate; otherwise Nr\textgreater=1 makes the displayed positive radius meaningful.

The zero frequency was excluded only from the positive pushforward (VII.3.3), where the divisor-fiber bound requires r!=0. Its exact finite identity remains (VII.3.1), and its separate source cost is retained in Appendix C.8. Actual selected-fiber weights and any outside Cauchy cardinality are also supplied there. This algebraic statement alone supplies no bound for the residual Gauss energy on the right.

The proof transcribes Sil's PR74 common-column transfer, retained at the pinned PR76 version. It uses only the previously proved finite local Gauss, CRT and ideal-divisor facts. Non-fifth-free moduli are not silently included: they require their complete prime-power defect terms.

\hypertarget{c.3.-the-sextic-sieve-and-completed-reference-estimate}{%
\subsubsection{C.3. The sextic sieve and completed reference estimate}\label{c.3.-the-sextic-sieve-and-completed-reference-estimate}}

Work over the fixed Eisenstein field \(\fitinline{K=\mathbf Q(\omega)}\), with \(\fitinline{\mathcal O=\mathbf Z[\omega]}\), the primary generators and trace Fourier normalization of Appendix A.1--A.2. Write \(\fitinline{NI}\) for the absolute norm of an ideal and \(\fitinline{Nz=|\iota(z)|^2}\) for an element.

An ideal is \textbf{supported} if it is nonzero and has no prime factor above \(\fitinline{2}\) or \(\fitinline{3}\). Additional fixed deletions, including the target conductor primes, are permitted by restriction. An ideal is sixth-power-free when every prime valuation lies in \(\fitinline{\{0,\ldots,5\}}\).

Put \[\fitmath{
\chi_I(z)=\prod_{\mathfrak p\mid I}
 \chi_{\mathfrak p}(z)^{v_{\mathfrak p}(I)}.
}\] Every local character power retains its zero at nonunits. In particular, \[\fitmath{
\chi_I(z)^6=\mathbf1_{(I,(z))=1},
}\tag{VII.4.1}\] whereas \(\fitinline{\chi_1(z)=1}\), including at \(\fitinline{z=0}\).

For finite ideal sets \(\fitinline{\mathcal R,\mathcal C}\), define \[\fitmath{
\mathcal E(\mathcal R,\mathcal C;a)
 =\sum_{R\in\mathcal R}
 \left|\sum_{I\in\mathcal C}a_I\chi_I(d_R)\right|^2,
 \qquad
 \|a\|_{\mathcal C}^2=\sum_{I\in\mathcal C}|a_I|^2.
}\] Here \(\fitinline{d_R}\) is the prescribed primary generator. The vector \(\fitinline{a}\) is common to the entire row family.

\textbf{Theorem.} For every \(\fitinline{\varepsilon>0}\), there is \(\fitinline{C_\varepsilon>0}\), depending only on \(\fitinline{\varepsilon}\) and the fixed field and conventions, such that, for \(\fitinline{M,N\ge1}\), supported sixth-power-free families satisfying \[\fitmath{
NR\le M,\qquad NI\le N
}\] obey \[\fitmath{
\mathcal E(\mathcal R,\mathcal C;a)
 \le C_\varepsilon(MN)^\varepsilon
 \left(M+N+M^{5/6}N^{1/3}+M^{1/3}N^{5/6}\right)
 \|a\|_{\mathcal C}^2.
}\tag{VII.4.2}\] The constant is chosen before the lengths, families and complex coefficient vector.

The proof below also supplies the completed reference estimate needed when neither ideal index is sixth-power-free.

\hypertarget{finite-operators-reciprocity-and-initial-control}{%
\subparagraph{1. Finite operators, reciprocity and initial control}\label{finite-operators-reciprocity-and-initial-control}}

Let \(\fitinline{\Theta(M,N)}\) be the squared operator norm of the matrix \[\fitmath{
\bigl(\chi_I(d_R)\bigr)_
 {\substack{R\ {\rm supported,\ sixth\!-\!free},\ NR\le M\\
             I\ {\rm supported,\ sixth\!-\!free},\ NI\le N}},
}\] and let \(\fitinline{\Lambda(M,N)}\) have the same columns and all supported ideal rows \(\fitinline{NR\le M}\).

These are finite matrices. Extending a vector by zero proves monotonicity in both lengths, bounds all prescribed subfamilies, and gives \[\fitmath{
\Theta(M,N)\le\Lambda(M,N).
}\tag{VII.4.3}\] If \(\fitinline{0<M<1}\), then \(\fitinline{\Lambda(M,N)=0}\).

The established reciprocity law, extended multiplicatively with its nonunit zeros, is \[\fitmath{
\chi_I(d_R)=\epsilon(\operatorname{ray}I,\operatorname{ray}R)
             \chi_R(d_I),\qquad |\epsilon|\le1.
}\tag{VII.4.4}\] The ray coordinate takes sixteen values. Split the rows into these sixteen classes. On each class, the matrix is a restriction of the transposed reverse matrix followed by a column multiplier of modulus at most one. The triangle inequality for operator norms and invariance under complex transpose therefore give \[\fitmath{
\Theta(M,N)\le256\,\Theta(N,M).
}\tag{VII.4.5}\] The same finite-family statement holds without either power-free restriction.

We first prove \[\fitmath{
\Lambda(M,N)\ll M+N^2 \qquad(M,N\ge1).
}\tag{VII.4.6}\]

Fix a nonnegative compact Schwartz majorant \(\fitinline{W_0\ge1}\) on \(\fitinline{[0,1]}\). For distinct sixth-free columns \(\fitinline{I,J}\), write \[\fitmath{
D=(I,J),\qquad I=DU,\qquad J=DV,\qquad (U,V)=1.
}\] Multiplicativity and (VII.4.1) give the exact identity \[\fitmath{
\chi_I(z)\overline{\chi_J(z)}
 =\mathbf1_{(D,(z))=1}\chi_U(z)\overline{\chi_V(z)}.
}\tag{VII.4.7}\] The coprime residual character is nonprincipal: at a prime of \(\fitinline{U}\) or \(\fitinline{V}\), its local exponent lies between one and five.

Its conductor is \(\fitinline{\operatorname{rad}(UV)}\). Primitive Poisson summation from Appendix A.4, equation (47), gives, for every \(\fitinline{A>0}\), \[\fitmath{
\left|\sum_z\chi_U(z)\overline{\chi_V(z)}W_0(Nz/A)\right|
 \ll_{W_0}\sqrt{N\operatorname{rad}U\,N\operatorname{rad}V}.
}\] Indeed, the zero Fourier coefficient vanishes, each normalized primitive Gauss factor has modulus one, and \[\fitmath{
\sum_{k\ne0}(1+tNk)^{-m}\ll_m t^{-1}\qquad(t>0,\ m>1).
}\] The latter follows directly from the fixed-lattice count \(\fitinline{O(X)}\).

Expand the mask in (VII.4.7) over squarefree divisors of \(\fitinline{D}\), and rescale the element variable in each term. This costs at most \(\fitinline{2^{\omega(D)}}\), uniformly for every resulting positive scale. Since \[\fitmath{
2^{\omega(D)}\le ND,\qquad
ND\sqrt{N\operatorname{rad}U\,N\operatorname{rad}V}
 \le\sqrt{NI\,NJ}\le N,
}\] the off-diagonal smoothed Gram entry is \(\fitinline{O_{W_0}(N)}\). The diagonal entry is \(\fitinline{O_{W_0}(M)}\), by lattice counting and \(\fitinline{|\chi_I|\le1}\).

There are \(\fitinline{O(N)}\) ideal columns. Expanding the square and using \[\fitmath{
\left(\sum_I|a_I|\right)^2
 \le\#\mathcal C\sum_I|a_I|^2
}\] proves (VII.4.6). This argument permits repeated prime factors in both columns and their gcd.

The only primitive finite-field input used here is elementary. For a nonprincipal multiplicative character \(\fitinline{\theta}\), \[\fitmath{
\left|\sum_x\theta(x)e(x)\right|^2=q.
}\] To verify it, substitute \(\fitinline{x=uy}\) in the square: the inner additive sum over \(\fitinline{y\ne0}\) is \(\fitinline{q-1}\) for \(\fitinline{u=1}\), and \(\fitinline{-1}\) otherwise; then use \(\fitinline{\sum_u\theta(u)=0}\). CRT gives the corresponding product statement, with the nonunit Fourier zeros retained.

\hypertarget{sixth-power-extraction-in-the-rows}{%
\subparagraph{2. Sixth-power extraction in the rows}\label{sixth-power-extraction-in-the-rows}}

Every supported ideal has a unique decomposition \[\fitmath{
R=R_0D^6,\qquad R_0\text{ sixth-free},
}\] with no coprimality condition between \(\fitinline{R_0,D}\). By (VII.4.1), \[\fitmath{
\chi_I(d_R)
 =\chi_I(d_{R_0})\mathbf1_{(I,D)=1}.
}\tag{VII.4.8}\] For fixed \(\fitinline{D}\), this replaces \(\fitinline{a_I}\) by the common vector \(\fitinline{a_I\mathbf1_{(I,D)=1}}\), whose mass contracts. Consequently, \[\fitmath{
\Lambda(M,N)
 \le\sum_{\substack{D\ {\rm supported}\\ND\le M^{1/6}}}
       \Theta\!\left(\frac{M}{(ND)^6},N\right).
}\tag{VII.4.9}\] This is a bound obtained from an exact partition and exact masks; it does not identify the punctured polynomial with the original polynomial.

For \(\fitinline{s>1}\), let \[\fitmath{
Z_{\rm supp}(s)=\sum_{D\ {\rm supported}}(ND)^{-s}<\infty.
}\] Convergence follows by dyadic summation from the fixed-field ideal count. We will repeatedly use \[\fitmath{
\#\{D:ND\le X\}\ll X.
}\tag{VII.4.10}\]

For \(\fitinline{\alpha\ge7/6}\), consider the exponent assertion \[\fitmath{
\Theta(M,N)\ll_\delta(MN)^\delta
 E_\alpha(M,N),\qquad
E_\alpha(M,N)=M+N^\alpha+M^{5/6}N^{1/3},
}\tag{VII.4.11}\] for every \(\fitinline{\delta>0}\). It implies the completed assertion \[\fitmath{
\Lambda(L,B)\ll_\delta(LB)^\delta
 \bigl(L^\alpha+L^{1/6}B\bigr),
 \qquad L>0,\ B\ge1.
}\tag{VII.4.12}\]

For \(\fitinline{L<1}\), this is immediate. For \(\fitinline{L\ge1}\), apply (VII.4.9), reciprocity, and (VII.4.11). Writing \(\fitinline{d=ND}\), the three resulting sums are bounded by \[\fitmath{
\begin{aligned}
\sum_{d\le L^{1/6}} B&\ll L^{1/6}B,\\
\sum_D(L/d^6)^\alpha&\le Z_{\rm supp}(6\alpha)L^\alpha,\\
\sum_D B^{5/6}(L/d^6)^{1/3}
 &\le Z_{\rm supp}(2)L^{1/3}B^{5/6}.
\end{aligned}
}\] Also \[\fitmath{
L^{1/3}B^{5/6}\le L^\alpha+L^{1/6}B:
}\] if \(\fitinline{B\le L}\), it is at most \(\fitinline{L^{7/6}\le L^\alpha}\); if \(\fitinline{B\ge L}\), it is at most \(\fitinline{L^{1/6}B}\). The loss factor in every term is at most \(\fitinline{(LB)^\delta}\). This proves (VII.4.12), including its complete length range.

The initial estimate proves (VII.4.11) with \(\fitinline{\alpha=2}\).

\hypertarget{exact-primitive-pair-poisson-formula}{%
\subparagraph{3. Exact primitive-pair Poisson formula}\label{exact-primitive-pair-poisson-formula}}

For sixth-free \(\fitinline{I}\), put \[\fitmath{
r=\operatorname{rad}I,\qquad U=I/r.
}\] The primitive character \(\fitinline{\chi_I}\) has modulus \(\fitinline{r}\): its local exponent at each prime is in \(\fitinline{\{1,\ldots,5\}}\). Let \(\fitinline{m_I}\) be the product of the prescribed prime generators, so \[\fitmath{
(m_I)=r,\qquad |m_I|=\sqrt{Nr}.
}\] Let \(\fitinline{\gamma_I}\) be its normalized trace Gauss sum; the finite-field calculation above and CRT give \(\fitinline{|\gamma_I|=1}\).

For coprime sixth-free \(\fitinline{I,J}\), primitive Poisson gives \[\fitmath{
\begin{aligned}
&\sum_z\chi_I(z)\overline{\chi_J(z)}W(Nz/M)\\
&\quad=\frac{M\Gamma(I,J)}{\sqrt{Nr\,Nt}}
 \sum_k\overline{\chi_I(k)}\chi_J(k)\,
 \mathscr FW\!\left(\frac{MNk}{Nr\,Nt}\right),
\end{aligned}
}\tag{VII.4.13}\] where \(\fitinline{t=\operatorname{rad}J}\) and \(\fitinline{\mathscr F}\) is the manuscript's radial Fourier transform, not its Mellin transform.

This is precisely equation (47) with primitive modulus \(\fitinline{m_Im_J}\) and no extra mask. Its trace character, measure, scale and nonunit zeros coincide with the existing formula. The character exponents are now the actual multiplicities \(\fitinline{e}\) and \(\fitinline{6-e}\), rather than a fixed cubic specialization.

The exact CRT multiplier is \[\fitmath{
\Gamma(I,J)
 =\overline{\chi_J(-1)}\,\gamma_I\overline{\gamma_J}\,
       \chi_I(m_J)\overline{\chi_J(m_I)}.
}\tag{VII.4.14}\] Reflection of a Gauss sum proves the inverse-character factor \(\fitinline{\gamma(\chi^{-1})=\chi^{-1}(-1)\overline{\gamma(\chi)}}\); CRT then proves (VII.4.14).

Write \(\fitinline{m_I=u_Id_r}\). Record the finite label \[\fitmath{
\ell(I)=\bigl(u_I,\ e_I,\ \operatorname{ray}r\bigr),\qquad
e_I=((NI-1)/6)\bmod6.
}\] There are \(\fitinline{6\cdot6\cdot16=576}\) labels. If \(\fitinline{h_e(u)=\iota(u^e)}\), set \[\fitmath{
\eta(\ell(I),\ell(J))
 =\overline{h_{e_J}(-1)}h_{e_I}(u_J)
   \overline{h_{e_J}(u_I)}
   \epsilon(\operatorname{ray}r,\operatorname{ray}t).
}\] Substituting \(\fitinline{m_I=u_Id_r,m_J=u_Jd_t}\) in (VII.4.14), and applying reciprocity to the squarefree radicals, gives \[\fitmath{
\Gamma(I,J)
 =\eta(\ell(I),\ell(J))\gamma_I\overline{\gamma_J}
       \chi_U(d_t)\overline{\chi_V(d_r)},
\qquad |\eta|\le1,
}\tag{VII.4.15}\] where \(\fitinline{J=tV}\). Thus the repeated parts remain coupled.

For a fixed nonzero bad-prime factor \(\fitinline{c}\), define \[\fitmath{
A_I(c)=
 \frac{\overline{a_I\gamma_I\overline{\chi_I(c)}}}{\sqrt{Nr}}.
}\] Then \[\fitmath{
|A_I(c)|^2\le\frac{|a_I|^2}{Nr}.
}\tag{VII.4.16}\] For a supported frequency ideal \(\fitinline{K}\), the summand obtained from the weighted Gram entry \(\fitinline{a_I\overline{a_J}}\) times (VII.4.13), excluding its outside factor \(\fitinline{M}\), is exactly \[\fitmath{
\eta(\ell(I),\ell(J))\,
 \overline{A_I(c)}A_J(c)\,
 \mathcal X(U,V,r,t,K)\,
 \mathscr FW\!\left(\frac{M|c|^2NK}{Nr\,Nt}\right),
}\tag{VII.4.17}\] where \[\fitmath{
\mathcal X
 =\chi_U(d_t)\overline{\chi_V(d_r)}
   \overline{\chi_{rU}(d_K)}\chi_{tV}(d_K).
}\] On fixed radical ray classes, reciprocity factors this as \[\fitmath{
\mathcal X
 =\zeta(U,V;\operatorname{ray}r,\operatorname{ray}t)\,
   \overline{\chi_U(d_K)}\chi_V(d_K)\,
   \overline{\chi_r(d_{KV})}\chi_t(d_{KU}),
\qquad |\zeta|\le1.
}\tag{VII.4.18}\] Both identities hold with their natural zeros. No division by a character value occurs. For independent left and right vectors use their corresponding separate coefficient transforms. After fixing the radical ray classes and finite labels, the factor \(\fitinline{\zeta(U,V;\operatorname{ray}r,\operatorname{ray}t) \overline{\chi_U(d_K)}\chi_V(d_K)}\) is a bounded multiplier of the common row \(\fitinline{(U,V,K)}\). Retain it through the complete inner coprime \(\fitinline{(r,t)}\)-sum, then contract it in the row norm. No separation into independent U and V factors is required or asserted; the block norms below introduce no repeated-part cardinality factor.

\hypertarget{the-coherent-block-estimate}{%
\subparagraph{4. The coherent block estimate}\label{the-coherent-block-estimate}}

We give the finite operator argument used to estimate (VII.4.17).

For nonzero column ideals \(\fitinline{R,T}\), \[\fitmath{
\mathbf1_{(R,T)=1}
 =\sum_{\substack{D\mid R\\D\mid T}}\mu(D).
}\] For arbitrary matrices \(\fitinline{A,B}\) with a common row index, expand this identity inside \[\fitmath{
\sum_x\left|
 \sum_{R,T:(R,T)=1}
   \overline{A_{xR}a_R}B_{xT}b_T\right|.
}\] The triangle inequality followed by Cauchy over \(\fitinline{(D,x)}\) bounds its square by \[\fitmath{
\left(\|A\|^2\sum_R\tau(R)|a_R|^2\right)
\left(\|B\|^2\sum_T\tau(T)|b_T|^2\right).
}\tag{VII.4.19}\] This follows because the projection onto \(\fitinline{D\mid R}\) is a column projection, and \(\fitinline{\sum_D\|a\mathbf1_{D\mid R}\|^2=\sum_R\tau(R)|a_R|^2}\). Thus ideals of norm at most \(\fitinline{X}\) cost only \(\fitinline{O_\varepsilon(X^\varepsilon)}\).

Retain the repeated parts as block indices: \[\fitmath{
a=a_{U,r},\qquad b=b_{V,t}.
}\] The left matrix has rows \(\fitinline{(V,K)}\), columns \(\fitinline{r}\), and entry \(\fitinline{\chi_r(d_{VK})}\); the right matrix has rows \(\fitinline{(U,K)}\), columns \(\fitinline{t}\). Their block extensions have common rows \(\fitinline{(U,V,K)}\). For example, \[\fitmath{
\sum_{U,V,K}\left|\sum_r A_{(V,K),r}a_{U,r}\right|^2
 \le\|A\|^2\sum_{U,r}|a_{U,r}|^2.
}\] There is no factor equal to the number of repeated-part indices.

The map \(\fitinline{(V,K)\mapsto VK}\) has at most \(\fitinline{\tau(Q)}\) representations of any nonzero ideal \(\fitinline{Q}\). Therefore, if \(\fitinline{N(VK)\le L}\) and \(\fitinline{Nr\le B}\), \[\fitmath{
\|A\|^2\ll_\varepsilon L^\varepsilon\Lambda(L,B).
}\tag{VII.4.20}\] Combining (VII.4.19)--(VII.4.20) proves the required coherent bound. It applies to the common vectors \(\fitinline{a_{U,r}}\) and \(\fitinline{b_{V,t}}\), without assigning arbitrary coefficients independently to the enlarged rows.

We also need uniform separation of the radial weight. On \(\fitinline{1\le x,y,z\le2}\), write \[\fitmath{
\mathscr FW(H z/(xy))
 =\int_{\mathbf R}b_H(v)
       e^{iv(-\log x-\log y+\log z)}\,dv,
}\] with the Fourier normalization absorbed into \(\fitinline{v}\). For every fixed \(\fitinline{m}\), \[\fitmath{
\int|b_H(v)|\,dv\ll_{W,m}(1+H)^{-m}.
}\tag{VII.4.21}\] To prove this, choose a fixed smooth cutoff equal to one on the possible interval for \(\fitinline{\log z-\log x-\log y}\), and Fourier-transform \(\fitinline{\vartheta(u)\mathscr FW(He^u)}\). Its zeroth and second derivative \(\fitinline{L^1}\) norms are \(\fitinline{O_{W,m}((1+H)^{-m})}\), by Schwartz decay. Fourier integration by parts gives (VII.4.21). Only finitely many Schwartz seminorms are used.

All separated phases have modulus one and preserve coefficient mass. Finite sums commute with the integral by its absolute \(\fitinline{L^1}\) bound. For sums of group norms, align each original group by one scalar of modulus at most one before integrating. The alignment is independent of \(\fitinline{v}\); consequently no group-cardinality loss is introduced.

\hypertarget{frequency-bands-and-the-analytic-recurrence}{%
\subparagraph{5. Frequency bands and the analytic recurrence}\label{frequency-bands-and-the-analytic-recurrence}}

Partition the columns by their 576-valued label and a dyadic radical scale: \[\fitmath{
P\le Nr\le2P,\qquad Q\le Nt\le2Q.
}\] There are \(\fitinline{O(\log N)}\) classes. The repeated parts satisfy \[\fitmath{
NU\le N/P,\qquad NV\le N/Q.
}\]

Every nonzero element frequency has a unique representation \[\fitmath{
k=c\,d_K,\qquad c=u\lambda^a2^b,\quad K\ {\rm supported}.
}\] For \(\fitinline{X/2\le Nk\le X}\), the number of possible bad-prime labels \(\fitinline{c}\) is \(\fitinline{O((1+\log X)^2)}\), hence \(\fitinline{O_\varepsilon(X^\varepsilon)}\). Writing \(\fitinline{q=|c|^2\ge1}\), its supported ideal frequency satisfies \[\fitmath{
X/(2q)\le NK\le X/q.
}\] Thus the two enlarged lengths are at most \(\fitinline{NX/Q}\) and \(\fitinline{NX/P}\). The Fourier decay scale is exactly \[\fitmath{
\frac{Mq}{PQ}\frac{X}{2q}=\frac{MX}{2PQ}.
}\] The bad-prime factor cancels from this scale and only shortens the enlarged row ranges; its character factors remain in (VII.4.16)--(VII.4.18).

For a fixed column class, (VII.4.16) contributes the mass factors \(\fitinline{P^{-1}}\), \(\fitinline{Q^{-1}}\). Define \[\fitmath{
\mathcal D_\varepsilon(P,L)
 =C_\varepsilon L^\varepsilon\Lambda(L,2P)(2P)^\varepsilon/P.
}\] Equations (VII.4.17)--(VII.4.21) consequently bound the sum of absolute dual-family values over one band by \[\fitmath{
\ll_\varepsilon
 \frac{X^\varepsilon}{(1+MX/(2PQ))^3}
 \sqrt{\mathcal D_\varepsilon(P,NX/Q)
       \mathcal D_\varepsilon(Q,NX/P)}
 \,\|a\|\,\|b\|.
}\tag{VII.4.22}\] Here \(\fitinline{a,b}\) are independent vectors. Keeping them independent is what bounds the full coprime Gram operator, including complex tests.

Assume, for some \(\fitinline{1\le\alpha\le2}\), \[\fitmath{
\Lambda(L,B)\le A(LB)^\varepsilon
                   (L^\alpha+L^{1/6}B)
\quad(L>0,\ B\ge1).
}\tag{VII.4.23}\] Put \[\fitmath{
t=\frac M{2PQ},\quad
U_0=(2N/M)^\alpha N^{\alpha-1},\quad
V_0=2(2N/M)^{1/6}N^{1/6}.
}\] Because \(\fitinline{1\le P,Q\le N}\), substitution in either density gives \[\fitmath{
\sqrt{\mathcal D_\varepsilon(P,NX/Q)
       \mathcal D_\varepsilon(Q,NX/P)}
 \ll_\varepsilon
 A(2N^2X)^{2\varepsilon}
 \bigl(U_0(tX)^\alpha+V_0(tX)^{1/6}\bigr).
}\tag{VII.4.24}\] For example, \[\fitmath{
(NX/Q)^\alpha/P
 =(2N/M)^\alpha P^{\alpha-1}(tX)^\alpha
 \le U_0(tX)^\alpha.
}\] The \(\fitinline{1/6}\) term follows in the same way. The auxiliary product length is at most \(\fitinline{2N^2X}\).

For \(\fitinline{0<\varepsilon<1/6}\), both \(\fitinline{\alpha+3\varepsilon}\) and \(\fitinline{1/6+3\varepsilon}\) lie in \(\fitinline{(0,3)}\). For \(\fitinline{0<\nu<3}\), \[\fitmath{
\sum_{j\ge0}\frac{2^{j\nu}}{(1+t2^j)^3}\ll_\nu t^{-\nu}.
}\tag{VII.4.25}\] If \(\fitinline{t\le1}\), split at \(\fitinline{t2^j=1}\) and sum the two geometric series; if \(\fitinline{t\ge1}\), the whole sum is bounded by \(\fitinline{O_\nu(t^{-3})\le O_\nu(t^{-\nu})}\).

Apply (VII.4.25) to (VII.4.22)--(VII.4.24), with \(\fitinline{X=2^j}\). A fixed pair of column classes costs \[\fitmath{
\ll_{\varepsilon,\alpha,W}
 A N^{4\varepsilon}t^{-3\varepsilon}(U_0+V_0)\|a\|\|b\|.
}\] Since \(\fitinline{P,Q\le N}\), \[\fitmath{
N^{4\varepsilon}t^{-3\varepsilon}
 \le2^{3\varepsilon}N^{10\varepsilon}M^{-3\varepsilon}.
}\] The two column-class counts cost at most another \(\fitinline{N^{2\varepsilon}}\). Thus the complete nonzero Poisson series is absolutely summable and costs \[\fitmath{
\ll A N^{12\varepsilon}M^{-3\varepsilon}
                  (U_0+V_0)\|a\|\|b\|.
}\tag{VII.4.26}\] At zero frequency, natural zeros leave only \(\fitinline{I=J=1}\), contributing at most \(\fitinline{|\mathscr FW(0)|\|a\|\|b\|}\). Restoring the outside factor \(\fitinline{M}\) in Poisson, and observing \[\fitmath{
MU_0=2^\alpha M^{1-\alpha}N^{2\alpha-1},
\qquad
MV_0=2^{7/6}M^{5/6}N^{1/3},
}\] bounds the smoothed coprime Gram operator, for every \(\fitinline{M>0,N\ge1}\), by a constant times \[\fitmath{
\mathcal T(M,N)
 =M+A N^{12\varepsilon}M^{-3\varepsilon}
       \left(M^{1-\alpha}N^{2\alpha-1}
             +M^{5/6}N^{1/3}\right).
}\tag{VII.4.27}\]

It remains to remove the coprime-pair restriction. Sort each original Gram pair by \(\fitinline{D=(I,J)}\), use (VII.4.7), expand its mask over \(\fitinline{E\mid D}\), and set \(\fitinline{z=d_Ew}\). The descendant coefficients are \[\fitmath{
a^{D,E}_{I/D}
 =\mathbf1_{D\mid I}\,a_I\chi_{I/D}(d_E),
}\] and their lengths are \[\fitmath{
M'=M/NE,\qquad N'=N/ND,\qquad 1\le NE\le ND\le N.
}\] Their total mass satisfies \[\fitmath{
\sum_D\sum_{E\mid D}\|a^{D,E}\|^2
 \le\sum_I|a_I|^2\sum_{D\mid I}\tau(D)
 \le\sum_I\tau(I)^2|a_I|^2
 \ll_\varepsilon N^\varepsilon\|a\|^2.
}\tag{VII.4.28}\] Repeated column primes and all new natural zeros are included in this inequality.

For \(\fitinline{1\le e\le d}\), \[\fitmath{
(M/e)^p(N/d)^q\le M^pN^q
}\] whenever \(\fitinline{q\ge0}\) and \(\fitinline{p+q\ge0}\). Apply this to \[\fitmath{
(p,q)=(-3\varepsilon,12\varepsilon),\
(1-\alpha,2\alpha-1),\
(5/6,1/3).
}\] Hence \(\fitinline{\mathcal T(M/e,N/d)\le\mathcal T(M,N)}\). The descendant row scale may be below one; (VII.4.27) was proved for every positive scale. Compact positive smoothing and (VII.4.28) now prove the full recurrence \[\fitmath{
\boxed{\;
\Lambda(M,N)\ll_{\varepsilon,\alpha}
 N^\varepsilon\!
 \left[M+A N^{12\varepsilon}M^{-3\varepsilon}
 \left(M^{1-\alpha}N^{2\alpha-1}
       +M^{5/6}N^{1/3}\right)\right].
\;}
}\tag{VII.4.29}\] All constants precede \(\fitinline{M,N,A}\) and the coefficient families.

\hypertarget{exponent-improvement-and-the-symmetric-majorant}{%
\subparagraph{6. Exponent improvement and the symmetric majorant}\label{exponent-improvement-and-the-symmetric-majorant}}

Suppose (VII.4.11) holds for \(\fitinline{7/6<\alpha\le2}\). Section 2 supplies (VII.4.23) for every positive loss.

Set \[\fitmath{
\alpha'=\frac{12\alpha-7}{6\alpha-1},
\qquad
a=\frac{12\alpha-8}{6\alpha-1},
\qquad T=\max(M,N^a).
}\] Then \[\fitmath{
1\le a\le\alpha',\quad
\frac56a+\frac13=\alpha',\quad
a(1-\alpha)+2\alpha-1=\alpha'.
}\] Since \(\fitinline{1-\alpha<0}\), \[\fitmath{
\begin{aligned}
T&\le M+N^{\alpha'},\\
T^{1-\alpha}N^{2\alpha-1}&\le N^{\alpha'},\\
T^{5/6}N^{1/3}
 &\le M^{5/6}N^{1/3}+N^{\alpha'}.
\end{aligned}
}\] Thus the sum of these three terms is at most \(\fitinline{3E_{\alpha'}(M,N)}\).

Given a requested loss \(\fitinline{\delta>0}\), choose \[\fitmath{
\varepsilon=\min(\delta/13,1/12).
}\] Use monotonicity \(\fitinline{\Theta(M,N)\le\Lambda(T,N)}\), apply (VII.4.29), and use \(\fitinline{T\ge1}\). The outside losses are bounded by \[\fitmath{
N^{13\varepsilon}\le(MN)^\delta.
}\] This proves (VII.4.11) with \(\fitinline{\alpha'}\).

Starting at \(\fitinline{\alpha_0=2}\), iterate this map. Direct calculation gives \[\fitmath{
\alpha_{j+1}-\frac76
 =\frac{5(\alpha_j-7/6)}{6\alpha_j-1}
 \le\frac56\left(\alpha_j-\frac76\right).
}\] Consequently \(\fitinline{7/6<\alpha_j\le2}\) and \[\fitmath{
\alpha_j-\frac76\le(5/6)^{j+1}.
}\] For any requested final \(\fitinline{\varepsilon>0}\), choose a fixed finite \(\fitinline{j}\) with \(\fitinline{\alpha_j<7/6+\varepsilon/2}\), then use (VII.4.11) with loss \(\fitinline{\varepsilon/2}\).

If \(\fitinline{N\le M}\), \[\fitmath{
N^{7/6}\le M^{5/6}N^{1/3};
}\] hence the resulting exponent shape is bounded by twice \(\fitinline{(MN)^{\varepsilon/2}}\) times the majorant in (VII.4.2). If \(\fitinline{M<N}\), first use reciprocity (VII.4.5) and interchange the lengths. The majorant is symmetric. This proves (VII.4.2), with a permissible final factor \(\fitinline{512}\) in the constant.

The iteration count and all intermediate losses are fixed after the final loss is chosen. No iteration depth or implicit constant depends on the lengths or coefficient vector.

\hypertarget{completed-rows-element-rows-and-central-columns}{%
\subparagraph{7. Completed rows, element rows and central columns}\label{completed-rows-element-rows-and-central-columns}}

Define \[\fitmath{
\mathcal R(M,N)
 =M+M^{1/6}N+M^{5/6}N^{1/3}+M^{1/3}N^{5/6}.
}\] Using (VII.4.9) with the proved symmetric bound gives \[\fitmath{
\Lambda(M,N)\ll_\varepsilon(MN)^\varepsilon\mathcal R(M,N).
}\tag{VII.4.30}\] Indeed, its four sums are bounded respectively using \(\fitinline{Z_{\rm supp}(6)}\), the count \(\fitinline{O(M^{1/6})}\), \(\fitinline{Z_{\rm supp}(5)}\), and \(\fitinline{Z_{\rm supp}(2)}\).

For literal nonzero element rows, write uniquely \(\fitinline{x=c\,d_K}\), where \(\fitinline{c=u\lambda^a2^b}\) and \(\fitinline{K}\) is supported. Twist the common coefficient by \[\fitmath{
a_I\longmapsto a_I\chi_I(c).
}\] Its mass contracts. There are \(\fitinline{O((1+\log M)^2)}\) bad-prime labels for \(\fitinline{Nx\le M}\); absorb them using half the requested loss. Apply (VII.4.30) to each supported ideal family.

At \(\fitinline{x=0}\), \(\fitinline{\chi_I(0)=\mathbf1_{I=1}}\); this row contributes at most \(\fitinline{\|a\|^2}\), already dominated by the reference majorant. Therefore \[\fitmath{
\sum_{\substack{x\in\mathcal X\\Nx\le M}}
 \left|\sum_{I\in\mathcal C}a_I\chi_I(x)\right|^2
 \ll_\varepsilon(MN)^\varepsilon\mathcal R(M,N)\|a\|^2
}\tag{VII.4.31}\] for arbitrary finite element rows and supported sixth-free columns.

Now remove the column restriction under the literal central hypothesis \[\fitmath{
|a_I|^2\le A/(NI),\qquad NI\le N.
}\tag{VII.4.32}\] Write uniquely \(\fitinline{I=JD^6}\), with \(\fitinline{J}\) sixth-free; overlap between \(\fitinline{J,D}\) is allowed. For each \(\fitinline{D}\), \[\fitmath{
F_D(x)=\sum_{I:\,\text{sixth quotient }D}a_I\chi_I(x)
 =\mathbf1_{(D,(x))=1}
     \sum_Ja_{JD^6}\chi_J(x).
}\tag{VII.4.33}\] For \(\fitinline{\delta>0}\), the reduced vector has mass at most \[\fitmath{
\sum_J|a_{JD^6}|^2
 \le\frac A{(ND)^6}\sum_{NJ\le N}(NJ)^{-1}
 \le\frac{A N^\delta Z_{\rm supp}(1+\delta)}{(ND)^6}.
}\tag{VII.4.34}\] The mask in (VII.4.33) contracts each fiber energy, so (VII.4.31) applies to that fiber.

Weighted Cauchy gives \[\fitmath{
\left|\sum_DF_D(x)\right|^2
 \le Z_{\rm supp}(2)\sum_D(ND)^2|F_D(x)|^2.
}\] After summing the fiber energies, the remaining root sum is \(\fitinline{Z_{\rm supp}(4)}\). Choose the reference loss and \(\fitinline{\delta}\) each equal to half the requested loss. Since \(\fitinline{N^\delta\le(MN)^\delta}\), we obtain \[\fitmath{
\boxed{\;
\sum_{x\in\mathcal X}
 \left|\sum_{I\in\mathcal C}a_I\chi_I(x)\right|^2
 \ll_\varepsilon(MN)^\varepsilon\mathcal R(M,N)\,A.
\;}
}\tag{VII.4.35}\] Both power-free restrictions are gone. The hypothesis replacing the column restriction is exactly (VII.4.32), not an arbitrary-coefficient assertion.

All displayed reference applications assume \(\fitinline{M,N\ge1}\). If a shortened column bound is below one, the ideal column family is empty. If a literal element row bound is below one, only \(\fitinline{x=0}\) remains, and (VII.4.32) bounds its energy by \(\fitinline{A}\). Thus subunit shortenings have an explicit terminal treatment.

\hypertarget{the-actual-common-tuple-convolution}{%
\subparagraph{8. The actual common tuple convolution}\label{the-actual-common-tuple-convolution}}

Let \(\fitinline{S}\) be any finite, row-independent set of \(\fitinline{d}\)-tuples of ideals, with arbitrary fixed collisions, overlaps and tuple restrictions. Put \[\fitmath{
I(v)=\prod_{j=1}^dv_j,\qquad
a_I=\sum_{\substack{v\in S\\I(v)=I}}\beta_v.
}\] Then, exactly, \[\fitmath{
\sum_{v\in S}\beta_v\chi_{I(v)}(x)=\sum_Ia_I\chi_I(x).
}\tag{VII.4.36}\] Every factor of a tuple with product \(\fitinline{I}\) divides \(\fitinline{I}\); hence its fiber has at most \(\fitinline{\tau(I)^d}\) members. Suppose \[\fitmath{
I(v)\text{ supported},\quad NI(v)\le N,\quad
|\beta_v|\le B/\sqrt{NI(v)}.
}\] The fixed-field divisor bound, with an arbitrarily small exponent chosen after \(\fitinline{d}\), gives \[\fitmath{
|a_I|^2\ll_{\delta,d}N^\delta B^2/(NI).
}\] Apply (VII.4.35) with half the final loss and absorb this envelope with the other half: \[\fitmath{
\sum_{x\in\mathcal X}
 \left|\sum_{v\in S}\beta_v\chi_{I(v)}(x)\right|^2
 \ll_{\varepsilon,d}(MN)^\varepsilon
          \mathcal R(M,N)B^2.
}\tag{VII.4.37}\]

For the reflected source scales assume \(\fitinline{Z\ge1, R\ge0, L>0, N\ge1, T\ge0, B\ge0}\), and \(\fitinline{Nx\le Z^R}\) for every \(\fitinline{x\in\mathcal X}\). Suppose \[\fitmath{
M=Z^R,\quad N\le Z^v,\quad N\le Z^yT,\quad
v\le4R/5,\quad y\le R/2,\quad R+v\le L.
}\] Then \[\fitmath{
\mathcal R(Z^R,N)\le Z^R(3+T^{1/3}).
}\] The terms \(\fitinline{M^{1/6}N}\) and \(\fitinline{M^{1/3}N^{5/6}}\) are at most \(\fitinline{Z^R}\); the remaining mixed term is at most \(\fitinline{Z^RT^{1/3}}\). Using loss \(\fitinline{\varepsilon/L}\) in (VII.4.37) yields \[\fitmath{
\sum_x\left|\sum_{v\in S}\beta_v\chi_{I(v)}(x)\right|^2
 \ll_{\varepsilon,L,d}Z^{R+\varepsilon}
                (3+T^{1/3})B^2.
}\tag{VII.4.38}\] Its square-root cost is bounded by \[\fitmath{
\ll Z^{(R+\varepsilon)/2}(2+T^{1/6})B.
}\]

Arbitrary fixed target twists, norm phases and angular factors may enter the common vector without changing these bounds when their modulus is at most one. Common joint profiles are handled by the already established Fourier--Minkowski adapter, provided its finite weighted \(\fitinline{L^1}\) envelope is supplied. Independent heights remain independent throughout. A row-dependent coefficient family, or a new angular reflection identity, requires its own adapter and is not asserted by (VII.4.37).

\hypertarget{c.4.-actual-selected-short-reference-estimate}{%
\subsubsection{C.4. Actual selected short-reference estimate}\label{c.4.-actual-selected-short-reference-estimate}}

This section completes the passage from Appendix C.3's common-column sieve to the selected reference used in (5.42), (5.47), and the short-root cases. It uses Appendix C.1, for the already written forward-restoration and common-puncture masses. It does not require another physical source estimate.

Write \(\fitinline{m_0=\mathrm{fixedBadMask}}\), and put \[\fitmath{
C_{\mathrm{pres}}
 =N(m_0)\,N(72),\qquad
R=\log_Z\!\bigl(C_{\mathrm{pres}}N\mathfrak q_\eta K\bigr).
}\tag{A.1}\] Here \(\fitinline{N(m_0)}\) denotes the norm of its principal ideal. The fixed presentation factor is retained in \(\fitinline{R}\).

For the selected supplier \(\fitinline{q}\), let \[\fitmath{
V_q=\mathrm{volume}(q),\qquad W_q=\mathrm{scale}(q),\qquad
d=\log_ZV_q.
}\] These are the original and padded selected volumes, respectively. In particular, \[\fitmath{
W_q=2^{|J_q|}V_q,\qquad |J_q|\le N,
}\tag{A.2}\] and every actual selected product satisfies \[\fitmath{
NF\le W_q\le2^NZ^d.
}\tag{A.3}\] We use \(\fitinline{W_q}\), not \(\fitinline{V_q}\), in the normalized grouped sum below.

For the source's displayed ordinary lengths, put \[\fitmath{
P=\log_Z(X_1X_2).
}\] Let \(\fitinline{S\ge0}\) bound its unpadded live marked product: \[\fitmath{
\prod_{i\in T}P_i\le Z^S.
}\] Assume \[\fitmath{
P+2S+4d\le R-\rho,\qquad
P+S>\frac56R-\frac73d,\qquad \rho>0.
}\tag{A.4}\] Define the physical midpoint \[\fitmath{
\ell_-=\frac32\left(P+S+\frac73d-\frac56R\right),\quad
\ell_+=\frac{P-S-R/2-d}{2},\quad
\ell=\frac{\ell_-+\ell_+}{2}.
}\tag{A.5}\] The selected short-reference case is \[\fitmath{
X_1\le Z^\ell\quad\text{or}\quad X_2\le Z^\ell.
}\tag{A.6}\]

The source \(\fitinline{s}\) retains its original prime pools, fifth-free ordinary-product restriction, distinct live marks, ordinary and marked holes, ordinary puncture \(\fitinline{R_s}\), two original Schwartz profiles \(\fitinline{W_1,W_2}\), marked profiles, frequencies and prefixes. Its marked coefficient is the literal canonical coefficient \[\fitmath{
b_i^t(I)=
 \mathrm{slotCoefficient}(\nu,R_{0},t_0+t;i,I)\,
 W_{\mathrm{slot},i}(NI/P_i).
}\tag{A.7}\] In particular the shift is \(\fitinline{t_0+t}\). Assume its envelopes satisfy \(\fitinline{M_i\le B_f}\), where \(\fitinline{B_f\ge1}\).

The explicit polynomial caps required below are \[\fitmath{
\begin{gathered}
R\le B_R,\quad 1\le V_q\le Z^{B_d},\quad W_q\le Z^{B_{\rm sel}},\\
N(F_{\rm hole}R_s)\le Z^{K_{\rm cap}},\qquad
N\mathfrak p\le Z^{K_{\rm cap}}\quad(\mathfrak p\in\mathcal S_s),\\
\prod_i2P_i\le Z^L,\qquad
bX_j\prod_i2P_i\le Z^L\quad(j=1,2),
\end{gathered}
}\tag{A.8}\] where \[\fitmath{
L=2(B_R+1),\qquad K_{\rm cap}=B_R+B_d+\rho+1.
}\tag{A.9}\] The source's finite enclosing prime set must contain every prime of every supported ordinary ideal with norm at most \(\fitinline{bX_1}\) or \(\fitinline{bX_2}\). The source hypotheses also retain \[\fitmath{
\mathfrak p\in\mathcal H_s
 \ \Longrightarrow\ 
 \mathfrak p\mid F_{\rm hole},\quad
 \mathfrak p\in R_{0,i}\quad\text{for every live }i.
}\tag{A.10}\]

For a selected \(\fitinline{F}\), let \(\fitinline{\chi_{F,z}}\) be the actual \texttt{s\allowbreak{}e\allowbreak{}l\allowbreak{}e\allowbreak{}c\allowbreak{}t\allowbreak{}e\allowbreak{}d\allowbreak{}C\allowbreak{}h\allowbreak{}a\allowbreak{}r\allowbreak{}a\allowbreak{}c\allowbreak{}t\allowbreak{}e\allowbreak{}r\allowbreak{}A\allowbreak{}t} construction. The retained row predicate may depend on \(\fitinline{F}\), but must imply \[\fitmath{
z\ne0,\qquad
\neg\mathrm{FixedInducingRow}
  \bigl(\eta,Q,m_0,d_F^2,z\bigr),
\quad d_F=\mathrm{primaryGenerator}(F).
}\tag{A.11}\]

\textbf{Theorem.} Fix the ambient finite slot type, its bound \(\fitinline{N}\), the constants in (A.8)--(A.9), \(\fitinline{a>0,b\ge0,B_f\ge1}\), and \(\fitinline{\varepsilon>0}\). There are a finite Schwartz control set \(\fitinline{\mathcal C}\), containing \(\fitinline{(0,0)}\), an integer \(\fitinline{J}\), a positive constant \(\fitinline{C}\), and a threshold for \(\fitinline{Z}\), chosen before the source, supplier, characters, profiles, heights and selected products, such that \[\fitmath{
\begin{aligned}
&\sum_{F\in\mathrm{selectedProducts}(q)}
 \frac{g_q(F)}{W_q}\,
 \mathcal E_{\mathrm{keep}_F}
 \left[
   \mathrm{heightRectangle}
       (s,\chi_{F,\bullet},t;X_1,X_2);
       \Phi,K
 \right]\\
&\qquad\le
 \mathrm{radialControl}(\Phi)\,
 C K_{\rm slot}^{2N}
 [\mathcal S_{\mathcal C}(W_1)
  \mathcal S_{\mathcal C}(W_2)]^2
 \mathcal H^J\,N\mathfrak q_\eta K\,Z^\varepsilon,
\end{aligned}
}\tag{A.12}\] where \[\fitmath{
\mathcal H=1+|t|+V+|t_1|+|t_2|,\qquad V\ge0,
}\] and \(\fitinline{\operatorname{supp}\Phi\subset(-\infty,2]}\).

The weights \(\fitinline{g_q(F)}\) are the actual grouped weights \[\fitmath{
g_q(F)=
 \sum_{\substack{y\in\mathrm{supportedSelected}(q)\\
                  \mathrm{selectedProduct}(q,y)=F}}
       |\mathrm{selectedWeight}(q,y)|^2.
}\tag{A.13}\] The supplier assumptions are \(\fitinline{|q.\nu_i(I)|\le1}\), \(\fitinline{|q.W_{\mathrm{slot},i}|\le K_{\rm slot}}\), \(\fitinline{K_{\rm slot}\ge1}\), and at most \(\fitinline{N}\) slots. No replacement coefficient sequence occurs in (A.12).

Here the radial energy uses \(\fitinline{\Re\Phi}\), exactly as in the source definition. Positivity of \(\fitinline{\Phi}\) is unnecessary for this upper bound.

\hypertarget{a-common-product-estimate-from-appendix-c.3}{%
\subparagraph{1. A common product estimate from Appendix C.3}\label{a-common-product-estimate-from-appendix-c.3}}

Suppose a finite character family \(\fitinline{\chi_z}\), on element rows \(\fitinline{Nz\le Z^R}\), satisfies \[\fitmath{
a_{\chi_z}(I)=c(I)\chi_I(z)
 \quad(I\text{ supported}),\qquad
a_{\chi_z}(I)=0\quad(I\text{ unsupported}),
}\tag{A.14}\] where \(\fitinline{c}\) is common and \(\fitinline{|c(I)|\le1}\).

A finite factor of length \(\fitinline{H}\) and amplitude \(\fitinline{B}\) means \[\fitmath{
F(\chi)=\sum_{I\in\mathcal P}a_Ia_\chi(I),\qquad
NI\le H,\qquad |a_I|\le B/\sqrt{NI}.
}\] Multiplying finitely many such factors, grouping by their complete product, and using Appendix C.3's convolution estimate gives \[\fitmath{
\sum_z\left|\prod_iF_i(\chi_z)\right|^2
 \ll Z^{R+\epsilon}
       (3+T^{1/3})\left(\prod_iB_i\right)^2
}\tag{A.15}\] provided the complete product length \(\fitinline{N_c}\) satisfies \[\fitmath{
N_c\le Z^v,\qquad N_c\le Z^yT,\qquad
R+v\le L,\quad v\le4R/5,\quad y\le R/2.
}\tag{A.16}\] Unsupported coefficients vanish before grouping; the fixed coefficient \(\fitinline{c}\) is then absorbed into the common vector.

For a smooth Hecke factor \[\fitmath{
P_\chi(W,X,u)
 =X^{-1/2}\sum_IW(NI/X)(NI/X)^{2\pi iu}a_\chi(I),
}\] support in \(\fitinline{(-\infty,b]}\) and \(\fitinline{|W|\le B}\) give length \(\fitinline{bX}\) and amplitude \(\fitinline{\sqrt b\,B}\), since \[\fitmath{
X^{-1/2}|W(NI/X)|\le\sqrt b\,B/\sqrt{NI}.
}\tag{A.17}\] Its frequency has modulus one and is retained exactly.

Thus (A.15) applies to two smooth factors and any fixed number of finite marked factors, uniformly in all their common frequencies.

\hypertarget{direct-or-reflected-selection-of-the-longer-factor}{%
\subparagraph{2. Direct or reflected selection of the longer factor}\label{direct-or-reflected-selection-of-the-longer-factor}}

Assume temporarily that every \(\fitinline{\chi_z}\) is nonprincipal and its presented modulus has norm at most a common \(\fitinline{U}\), with \[\fitmath{
1\le U\le Z^{K_{\rm cap}}.
}\] For \(\fitinline{A=Z^{\xi/2}}\), define \[\fitmath{
\mathcal B(X,U)=
 \begin{cases}
 bX,&X\le\sqrt U,\\
 4A\,U/X,&X>\sqrt U.
 \end{cases}
}\tag{A.18}\] Both bounds \[\fitmath{
\mathcal B(X,U)\le(b+4)A\sqrt U,\qquad
\mathcal B(X,U)\le(b+4)A\,U/X
}\tag{A.19}\] hold for the \textbf{same} selected branch.

The direct branch is (A.17). For completeness, the reflected branch uses the existing natural reflection of Appendix A.5, equation (71), retaining the original character's redundant zeros. Choose its primitive inducing character \(\fitinline{\psi_z}\). If \(\fitinline{\mathcal R_z}\) is the redundant natural radical, its exact dual scales have the form \[\fitmath{
Y_{D,H}
 =\frac{N\mathfrak q_{\psi_z}\,N\prod_{\mathfrak p\in D}\mathfrak p}
        {X\,NH},
\quad D\subseteq\mathcal R_z,\quad
H\text{ supported on }\mathcal R_z,
}\tag{A.20}\] with \[\fitmath{
N\mathfrak q_{\psi_z}\,N\mathcal R_z
 \le N\mathfrak q_{\chi_z}\le U.
}\] Consequently \(\fitinline{0<Y_{D,H}\le U/X}\). The coefficients are the original deletion/restoration coefficients; their absolute mass is \(\fitinline{O_\delta((N\mathfrak q_{\chi_z})^\delta)}\).

After retaining the annuli prescribed by the cutoff \(\fitinline{A}\), every nonempty retained annular scale \(\fitinline{Y'}\) obeys \[\fitmath{
1\le Y'\le4A\,U/X.
}\tag{A.21}\] If this upper bound is below one, the retained family is empty. It is not replaced by a length-one family.

Here is why the common energy (A.15) controls these row-dependent reflected scales. On an annulus, Fourier inversion in logarithmic coordinates separates the transformed profile into common norm-frequency phases. Its Fourier envelope, including finitely many derivatives, is bounded by \[\fitmath{
C\,\mathcal S_{\mathcal C}(W_{\rm long})
       (1+|t_{\rm long}|)^n
}\] times a fixed integrable decaying function. Apply Minkowski before integrating. Sobolev on the logarithmic scale interval \(\fitinline{[0,\log Y]}\), using a function and its first derivative, costs \(\fitinline{1+2\log Y}\) in squared energy. The two required tests are precisely the manuscript's \texttt{s\allowbreak{}c\allowbreak{}a\allowbreak{}l\allowbreak{}e\allowbreak{}T\allowbreak{}e\allowbreak{}s\allowbreak{}t} \(\fitinline{j=0,1}\). Gauge away the unit normalization phase before taking a logarithmic scale derivative, so no short-frequency loss is introduced. The common-frequency energies required by this argument are (A.15).

Conjugating the reflected factor inside its absolute value restores the original character orientation and conjugates its profile. This leaves the other factors' moduli unchanged. The argument therefore applies with an arbitrary fixed physical multiplier, including the product of the other finite factors.

The redundant coefficient mass and the logarithmic scale supremum cost arbitrarily small powers of \(\fitinline{Z}\). Polynomial conductor caps make this uniform over all presented characters. All profile controls are finite seminorms chosen before the profiles; the zero profile follows by homogeneity, without division by a seminorm.

Finally, the discarded natural-reflection tail is \[\fitmath{
\ll C\,\mathcal S_{\mathcal C}(W_{\rm long})
 (1+|t_{\rm long}|)^n Z^{-A_0}
}\tag{A.22}\] for any fixed \(\fitinline{A_0}\), provided \(\fitinline{N\mathfrak q_{\chi_z}\le Z^{L_{\rm tail}}X}\). In the reflected branch \(\fitinline{X>\sqrt U\ge1}\), so the coarse modulus bound supplies this with \(\fitinline{L_{\rm tail}=K_{\rm cap}}\). Taking \(\fitinline{A_0=L+1}\) pays the remaining factors' crude \(\fitinline{O(Z^L)}\) norm and the \(\fitinline{O(Z^R)}\) element-row count.

Equations (A.20)--(A.22) are the original natural reflection, redundant Euler restoration, annular truncation and profile/Sobolev argument of Appendix A.1.3 and A.5.3. The new energy supplied to that argument is (A.15); no old moment estimate is transferred.

We have proved the following selection rule: the product of a long ordinary factor, a short ordinary factor and finite live factors obeys (A.15) when its complete selected length is computed using \(\fitinline{\mathcal B}\). The resulting bound has the form \[\fitmath{
C\,\mathcal S_{\mathcal C}(W_{\rm long})^2
 (1+\mathrm{Height})^{2n}
 Z^{R+\epsilon}(3+T^{1/3})
 \left(\sqrt b\,B_{\rm short}\prod_iB_i\right)^2.
}\tag{A.23}\]

All scales in this argument are positive. A direct factor with empty ideal support is zero. A reflected range below one has no retained annuli, and only its paid tail remains.

\hypertarget{apply-the-forward-identity-to-the-actual-restricted-rectangle}{%
\subparagraph{3. Apply the forward identity to the actual restricted rectangle}\label{apply-the-forward-identity-to-the-actual-restricted-rectangle}}

First orient \(\fitinline{X_1}\) as the short scale. Put \[\fitmath{
V_s=\prod_{i\in T}2P_i.
}\]

Use Appendix C.1's exact forward identity, including its nilpotent marked variables, before taking any norm. Partition the live slot labels by the subset \(\fitinline{T_0}\) whose marked occurrence supplies a forward correction. Its actual frozen assignment is an injective map \[\fitmath{
\tau:T_0\longrightarrow\mathcal S_s,\qquad
\tau(i)\in\mathcal P_i.
}\] Write \[\fitmath{
Q_\tau=\prod_{i\in T_0}N\tau(i).
}\] For each extracted ordinary pair \(\fitinline{e=(e_1,e_2)}\), the ordinary scales are exactly \[\fitmath{
X_1/Ne_1,\qquad X_2/Ne_2.
}\tag{A.24}\] Both are positive; neither is rounded up.

The remaining marked assignments retain their mutual injectivity and avoid every frozen prime. Group their actual products into one finite factor \(\fitinline{H_{T_0,\tau}}\). This factor is formed before the physical character is supplied. Its length satisfies \[\fitmath{
\operatorname{length}(H_{T_0,\tau})
 =\prod_{i\notin T_0}2P_i
 \le V_s/Q_\tau.
}\tag{A.25}\] For any \(\fitinline{\delta_a>0}\), its amplitude is bounded by \[\fitmath{
C_{\delta_a,N}
  (\max(1,V_s))^{\delta_a}
  (\sqrt2 B_f)^N.
}\tag{A.26}\] To prove (A.26), group the remaining tuple coefficients by their whole ideal product. A product \(\fitinline{I}\) has at most \(\fitinline{\tau(I)^N}\) ordered factorizations; apply the ideal-divisor estimate after the fixed \(\fitinline{N}\). All frozen-prime holes and live collision restrictions remain in the tuple set.

The exact rectangle is consequently a finite sum \[\fitmath{
\sum_{T_0}\sum_{\tau}\sum_e
 c_{\tau,e}(\chi_z)\,
 P_{\chi_z}(W_2,X_2/Ne_2,u_2)
 P_{\chi_z}(W_1,X_1/Ne_1,u_1)
 H_{T_0,\tau}(\chi_z).
}\tag{A.27}\] The scalar is the actual forward kernel coefficient times \[\fitmath{
\frac{a_{\chi_z}(e_1e_2)}{\sqrt{Ne_1Ne_2}}
 \prod_{i\in T_0}
 \frac{b_i^t(\tau(i))a_{\chi_z}(\tau(i))}{\sqrt{P_i}}.
}\] Thus every frequency, character phase and natural zero is retained.

By (A.19), (A.24)--(A.25), the same selected child length \(\fitinline{N_c}\) satisfies \[\fitmath{
N_c\le(b+4)A\,bX_1V_s\sqrt U,
}\tag{A.28}\] and \[\fitmath{
N_c\le
 (b+4)A\,b\frac{X_1V_sU}{X_2}
 \frac{Ne_2}{Ne_1Q_\tau}.
}\tag{A.29}\] The denominator in (A.29) contains the entire frozen product \(\fitinline{Q_\tau}\).

Suppose the right sides without the final ratio are bounded respectively by \(\fitinline{Z^v,Z^y}\). Then (A.23) bounds the child energy by \[\fitmath{
E\left(3+
 \left(\frac{Ne_2}{Ne_1Q_\tau}\right)^{1/3}\right).
}\tag{A.30}\] Appendix C.1, \S{}\S{}2--3, gives a common scalar majorant \(\fitinline{w_{\tau,e}}\), independent of \(\fitinline{z}\), with \[\fitmath{
|c_{\tau,e}(\chi_z)|\le w_{\tau,e},\qquad
\sum_{\tau,e}w_{\tau,e}
 \left[1+
 \left(\frac{Ne_2}{Ne_1Q_\tau}\right)^{1/6}\right]
 \ll_\delta Z^\delta.
}\tag{A.31}\] Its central and critical regimes are exactly \[\fitmath{
(1/2,1/2;1/2),\qquad(2/3,1/3;2/3).
}\] Weighted triangle in the common row Hilbert space, followed by \(\fitinline{\sqrt{3+r^2}\le2(1+r)}\), combines (A.30)--(A.31). It does not cost the number of frozen assignments or extraction pairs.

There are finitely many \(\fitinline{T_0}\), at most \(\fitinline{2^N}\). Take the union of their finite profile controls, the sum of their height degrees, and a common constant before choosing any profiles.

For explicit loss accounting, let e\textgreater0 be the requested restricted-rectangle loss. Apply (A.23) with loss e/4. Choose the norm-mass exponent in (A.31) as e/8; its square costs e/4. Choose the tuple-divisor exponent \(\fitinline{\delta_a=e/(4L)}\); since \(\fitinline{V_s\le Z^L}\), squaring (A.26) costs e/2. These three energy losses total e. The remaining finite subset count is part of the constant. This proves the complete forward restricted bound with loss e before the separate puncture allocation in the next paragraph.

\hypertarget{the-whole-ordinary-puncture-remains-in-the-coefficients}{%
\subparagraph{4. The whole ordinary puncture remains in the coefficients}\label{the-whole-ordinary-puncture-remains-in-the-coefficients}}

The source's puncture \(\fitinline{R_s}\) is handled exactly, without changing the row character. Enlarge the common ordinary holes by the primes of \(\fitinline{R_s}\), and replace \[\fitmath{
R_{0,i}\longmapsto R_{0,i}\cup\operatorname{primeSupport}(R_s).
}\tag{A.32}\] This is a pointwise identity for the canonical marked coefficient and the ordinary coprimality restriction. The new holes divide \(\fitinline{F_{\rm hole}R_s}\), whose norm is bounded in (A.8). The new marked coefficients still have envelope \(\fitinline{B_f}\), because (A.32) only adds zeros.

At a common hole, apply Appendix C.1's exact correction \(\fitinline{J=B/H}\). Its extracted child has the same shortened scales (A.24), and its reference ratio is \(\fitinline{Ne_2/Ne_1}\). Appendix C.1, \S{}4, supplies the common positive mass for both the central and critical weights: \[\fitmath{
\sum_e w_e
 \left[1+(Ne_2/Ne_1)^{1/6}\right]\ll_\delta Z^\delta.
}\tag{A.33}\] Apply the preceding restricted estimate with half the requested loss, then spend the other half on (A.33).

This proves the actual short-rectangle estimate once \[\fitmath{
\begin{gathered}
v=\ell+(R+d_{\rm ref})/2+S+2\delta_{\rm ref},\\
y=R+d_{\rm ref}+2\ell-P+S+2\delta_{\rm ref}
\end{gathered}
}\tag{A.34}\] satisfy \(\fitinline{R+v\le L,\ v\le4R/5,\ y\le R/2}\), and \[\fitmath{
U\le Z^{R+d_{\rm ref}},\quad
(b+4)b\,2^{|T|}\le Z^{\delta_{\rm ref}},\quad
\xi/2\le\delta_{\rm ref}.
}\tag{A.35}\] Indeed, \(\fitinline{X_1X_2=Z^P}\), \(\fitinline{X_1\le Z^\ell}\), and \(\fitinline{V_s\le2^{|T|}Z^S}\) turn (A.28)--(A.29) into (A.34).

The puncture \(\fitinline{F_{\rm hole}R_s}\) has entered only (A.32)--(A.33). It has not been inserted into the presented character's conductor.

\hypertarget{construct-the-actual-character-and-its-conductor-cap}{%
\subparagraph{5. Construct the actual character and its conductor cap}\label{construct-the-actual-character-and-its-conductor-cap}}

For each supported selected product \(\fitinline{F}\), construct first a sextic row character \(\fitinline{\psi_F}\) representing \(\fitinline{I\mapsto\chi_I(d_F)}\), with presented modulus norm at most \[\fitmath{
C_{\rm cub}NF.
}\] Squaring this character preserves its presented modulus and gives \[\fitmath{
a_{\psi_F^2}(I)=\chi_I(d_F)^2
}\] on supported ideals, with the same nonunit zeros. This is the linear-in-\(\fitinline{NF}\) cubic construction. Constructing instead a new sextic character with parameter \(\fitinline{d_F^2}\) would give an unnecessarily larger bound and is not the construction used here.

Apply the established row-character construction to \(\fitinline{\eta\psi_F^2}\), the original mask \(\fitinline{m_0}\), and \(\fitinline{z\ne0}\). It produces the actual \(\fitinline{\chi_{F,z}}\) with \[\fitmath{
a_{\chi_{F,z}}(I)
 =a_\eta(I)\chi_I(d_F)^2\chi_I(m_0^6)\chi_I(z)
 \quad(I\text{ supported}),
}\tag{A.36}\] zero on unsupported ideals, and \[\fitmath{
N\mathfrak q_{\chi_{F,z}}
 \le C_{\rm cub}N\mathfrak q_\eta N(m_0)N(72)\,NF\,Nz.
}\tag{A.37}\] The fixed multiplier in (A.36) has modulus at most one and is common in \(\fitinline{z}\), so it is exactly the coefficient identification (A.14).

The construction also gives, on every element \(\fitinline{n}\), \[\fitmath{
a_{\chi_{F,z}}((n))
 =\mathrm{rowTwist}(\eta,m_0,1,d_F^2z;n),
}\tag{A.38}\] with the original unsupported-element zeros. If this character were principal, it would be induced by the primitive unit character. Since \(\fitinline{Q\le(1)}\), (A.38) would witness \(\fitinline{\mathrm{FixedInducingRow}(\eta,Q,m_0,d_F^2,z)}\). Thus (A.11) proves that the \textbf{actual} presented character is nonprincipal. No untwisted \(\fitinline{z}\)-predicate is substituted.

On the radial support \(\fitinline{Nz\le2K}\), set \[\fitmath{
U_F=\max\!\left(1,\,
 C_{\rm cub}N\mathfrak q_\eta N(m_0)N(72)\,NF\,(2K)\right).
}\tag{A.39}\] This is \texttt{a\allowbreak{}c\allowbreak{}t\allowbreak{}u\allowbreak{}a\allowbreak{}l\allowbreak{}C\allowbreak{}o\allowbreak{}n\allowbreak{}d\allowbreak{}u\allowbreak{}c\allowbreak{}t\allowbreak{}o\allowbreak{}r\allowbreak{}C\allowbreak{}a\allowbreak{}p}, and (A.37) gives the required modulus cap.

Now choose \[\fitmath{
\omega=\rho/100,\qquad
d_{\rm ref}=d+\omega,\qquad
\delta_{\rm ref}=\rho/200=\omega/2,\qquad \xi=\omega.
}\tag{A.40}\] The physical selected exponent remains \(\fitinline{d}\); only the reference call uses \(\fitinline{d_{\rm ref}}\).

By (A.1), (A.3), \[\fitmath{
C_{\rm cub}N\mathfrak q_\eta N(m_0)N(72)\,NF\,(2K)
 \le2C_{\rm cub}2^N Z^{R+d}.
}\] A single threshold, independent of all future characters and suppliers, ensures \(\fitinline{2C_{\rm cub}2^N\le Z^\omega}\). Consequently \[\fitmath{
1\le U_F\le Z^{R+d_{\rm ref}}\le Z^{K_{\rm cap}}.
}\tag{A.41}\] Also \(\fitinline{N(72)=5184}\), so \(\fitinline{C_{\rm pres}\ge2}\) and \(\fitinline{N\mathfrak q_\eta\ge1}\) give \[\fitmath{
2K\le C_{\rm pres}N\mathfrak q_\eta K=Z^R.
}\tag{A.42}\] A second common threshold ensures \(\fitinline{(b+4)b2^N\le Z^{\delta_{\rm ref}}}\). This proves every requirement in (A.35).

\hypertarget{verify-the-reference-gates-with-their-reserve}{%
\subparagraph{6. Verify the reference gates with their reserve}\label{verify-the-reference-gates-with-their-reserve}}

The interval width in (A.5) is \[\fitmath{
\ell_+-\ell_-=R-P-2S-4d\ge\rho.
}\] Therefore \(\fitinline{0<\ell<P/2}\). The midpoint can also be written \[\fitmath{
\ell=P+\frac S2+\frac{3d}2-\frac{3R}4.
}\] Substituting (A.40) into (A.34) gives \[\fitmath{
v=P+\frac32S+2d-\frac R4+\frac32\omega,
\qquad
y=P+2S+4d-\frac R2+2\omega.
}\] Using (A.4), \[\fitmath{
v\le\frac34R-\rho-\frac S2-2d+\frac32\omega
 \le\frac45R,
}\] and \[\fitmath{
y\le\frac R2-\rho+2\omega\le R/2.
}\tag{A.43}\] The same hypotheses give \(\fitinline{R>0}\); the unsimplified expressions in (A.34) show \(\fitinline{v,y\ge0}\). Finally, \[\fitmath{
R+v\le\frac95R\le L.
}\] The strengthened mixed gate also survives: \[\fitmath{
P+2S+4d_{\rm ref}
 \le R-\rho+4\omega<R.
}\tag{A.44}\] This verifies the gates of \S{}4 from the original reserve, rather than assuming them for \(\fitinline{d+\omega}\).

\hypertarget{original-profiles-independent-heights-and-orientation}{%
\subparagraph{7. Original profiles, independent heights and orientation}\label{original-profiles-independent-heights-and-orientation}}

The exact canonical identity for the row under consideration is \[\fitmath{
\begin{aligned}
\mathrm{heightRectangle}
 (s,\chi,t;X_1,X_2)
={}&\mathrm{residualPhase}(s.X_1s.X_2,s.d,1,t)\\
&\times\mathrm{rectangleRow}\!
 \left(\chi;\,
 b_i^t,\,
 \mathrm{phaseProfile}(W_1,u_1),\,
 \mathrm{phaseProfile}(W_2,u_2);
 X_1,X_2\right),
\end{aligned}
}\tag{A.45}\] where \[\fitmath{
u_j=(t_j+t)/(2\pi).
}\] The residual phase has modulus one and continues to use the original stored source scales and prefixes. They are not replaced by the displayed reference lengths.

For nonzero pool ideals, \[\fitmath{
|b_i^t(I)|=|b_i^0(I)|\le M_i\le B_f.
}\tag{A.46}\] This identity includes every canonical fixed-hole zero.

Let \(\fitinline{H_{\rm ref}=|t|+|t_1|+|t_2|}\). Then \[\fitmath{
|u_j|\le H_{\rm ref},\qquad
1+H_{\rm ref}\le\mathcal H.
}\] Choose a finite Schwartz control set that dominates the long-profile controls from \S{}\S{}2--4 and both short scale tests. Continuity of \(\fitinline{W\mapsto W}\), \(\fitinline{W\mapsto xW'}\), and the corresponding fixed scale-test combinations gives \[\fitmath{
\sup_{j=0,1;\,x}|\mathrm{scaleTest}(W_1,j,x)|
 \le D\mathcal S_{\mathcal C}(W_1).
}\] Thus the left-short estimate is homogeneous: \[\fitmath{
C[\mathcal S_{\mathcal C}(W_1)
  \mathcal S_{\mathcal C}(W_2)]^2\mathcal H^JZ^{R+\epsilon}.
}\tag{A.47}\] The controls and degree are chosen before both profiles. The zero-profile case requires no normalization.

If \(\fitinline{X_2}\) is short, exchange only the two ordinary summation variables, displayed lengths, profiles and their frequencies. The finite rectangle restrictions are symmetric in the two ordinary variables. The marked slots, original source scales, prefixes, punctures, character and residual phase are fixed. This gives exactly the same rectangle and the same estimate (A.47).

To make this uniform over all live \(\fitinline{T\subseteq\iota}\), take the finite union of their controls, the sum of their nonnegative height degrees, and \(\fitinline{1+\sum_T C_T}\). This occurs before any source or profile is supplied.

\hypertarget{radial-energy-and-the-actual-grouped-weights}{%
\subparagraph{8. Radial energy and the actual grouped weights}\label{radial-energy-and-the-actual-grouped-weights}}

For fixed \(\fitinline{F}\), retain exactly the finite row set \[\fitmath{
\mathcal Z_F=\{z:\ z\ne0,\ Nz\le2K,\ \mathrm{keep}_F(z)\}.
}\] Since \(\fitinline{\operatorname{supp}\Phi\subset(-\infty,2]}\), every other radial summand is zero. Moreover \[\fitmath{
\Re\Phi(x)\le\sup_y|\Phi(y)|
 =\mathrm{radialControl}(\Phi).
}\] Apply (A.47) on \(\fitinline{\mathcal Z_F}\), then (A.1): \[\fitmath{
\begin{aligned}
\mathcal E_{\mathrm{keep}_F}
 [\mathrm{heightRectangle};\Phi,K]
\le{}&
\mathrm{radialControl}(\Phi)\,
 C[\mathcal S_{\mathcal C}(W_1)
   \mathcal S_{\mathcal C}(W_2)]^2
 \mathcal H^J\\
&\times C_{\rm pres}N\mathfrak q_\eta K\,Z^\epsilon.
\end{aligned}
}\tag{A.48}\] Absorb the fixed \(\fitinline{C_{\rm pres}}\) into \(\fitinline{C}\). The row keep has remained literal and \(\fitinline{F}\)-dependent throughout.

For the weighted result, apply (A.48) with loss \(\fitinline{\varepsilon/2}\). Set \[\fitmath{
\delta=\frac{\varepsilon}{2(B_{\rm sel}+1)}.
}\tag{A.49}\] Every selected assignment in a fiber of product \(\fitinline{F}\) has its components among the divisors of \(\fitinline{F}\). Since there are at most \(\fitinline{N}\) components, \[\fitmath{
|\mathrm{selectedFiber}(q,F)|
 \le\tau(F)^{|J_q|}
 \ll_{\delta,N}(NF)^\delta.
}\] Each squared selected amplitude is at most \(\fitinline{K_{\rm slot}^{2N}}\). Consequently \[\fitmath{
g_q(F)\ll_{\delta,N}K_{\rm slot}^{2N}(NF)^\delta
 \ll_{\delta,N}K_{\rm slot}^{2N}Z^{B_{\rm sel}\delta},
}\tag{A.50}\] using \(\fitinline{NF\le2W_q\le2Z^{B_{\rm sel}}}\), with \(\fitinline{2^\delta}\) in the constant.

The actual ideal-ball count gives \[\fitmath{
\#\{F\in\mathrm{selectedProducts}(q)\}\le256W_q.
}\] Thus \[\fitmath{
\sum_F\frac{g_q(F)}{W_q}
 \ll_{\delta,N}K_{\rm slot}^{2N}Z^{B_{\rm sel}\delta}.
}\tag{A.51}\] All these weights are nonnegative. Multiply the common nonnegative budget in (A.48) by (A.51). The exponent fits because \[\fitmath{
\frac{\varepsilon}{2}+B_{\rm sel}\delta
 =\frac{\varepsilon}{2}
   +\frac{B_{\rm sel}\varepsilon}{2(B_{\rm sel}+1)}
 \le\varepsilon.
}\] This proves (A.12).

\hypertarget{c.5.-exact-first-source-whole-kernels-and-allocation-normalizers}{%
\subsubsection{C.5. Exact first source, whole kernels and allocation normalizers}\label{c.5.-exact-first-source-whole-kernels-and-allocation-normalizers}}

Fix \(\fitinline{U>1}\), \(\fitinline{H_0=U^M}\), a finite slot set \(\fitinline{I}\), \(\fitinline{P_i=U^{w_i}}\) and

\[\fitmath{
b_i(\mathfrak p)=\nu_i(\mathfrak p)W(q_{\mathfrak p}/P_i)
       (q_{\mathfrak p}/P_i)^{-\sigma_i+iv_i}.
}\]

Each relative \(\fitinline{\nu_i}\) retains both defining zero sets. Pools contain all support primes in the identity ray class and may overlap. Actual affine input has \(\fitinline{\sigma_i=33/50}\), \(\fitinline{\operatorname{supp}W\subset[1,2]}\). Put \(\fitinline{\chi_z(n)=\eta(n)\chi_n(m^6Az)}\) with original zeros, and

\[\fitmath{
L_z(p,X;t)=\sum_{n\ne0}\chi_z(n)q_n^{it}p(q_n/X),\qquad
 V=X_1X_2\prod_{i\in I}P_i.
}\]

The energy to reduce is

\hypertarget{sr:energy}{}

\[\fitmath{
\qquad\text{(115)}
 E_X=\frac1V\sum_{\substack{z\ne0\\\operatorname{keep}(z)}}
 \Phi(q_z/H_0)
 \left|L_z(p_1,X_1;t)L_z(p_2,X_2;t)
       \prod_i\sum_{\mathfrak p\in\mathcal P_i}
         b_i(\mathfrak p)\chi_z(\mathfrak p)q_{\mathfrak p}^{it}\right|^2 .
}\]

Use real nonnegative \(\fitinline{\Phi}\) (or \(\fitinline{\Re\Phi}\) in the complex-profile convention). Keep the actual natural predicate and the same fixed \(\fitinline{Q_{\rm int}}\) satisfying \S{}2, including \(\fitinline{Q_{\rm int}\subseteq(72)}\); later zero bins do not replace it.

Fix bounded geometric ranges: \(\fitinline{0<M_{\min}\le M\le M_{\max}}\), nonnegative slot widths with bounded total, and four positive ordinary scales \(\fitinline{X_1,X_2,Y_1,Y_2}\) with bounded clipped logarithms and \(\fitinline{X_1X_2=Y_1Y_2}\). Empty annuli contribute zero. Every nonempty ordinary annulus in the source has its central scale bounded below by \(\fitinline{4/9}\), so subunit normalization costs a fixed constant. The reference rectangles may have an empty factor; this is used in LOW below.

\hypertarget{sr:initial-scales}{}

The actual application is the root cone \(\fitinline{P+LS\le B<1}\) of Appendix C.10, including all genuine prefix shortenings. The more restrictive historical conditions \(\fitinline{.41M\le\log_U X_i\le.44M}\) are not used in the identities that follow. Define the literal difference

\hypertarget{sr:comparison}{} \[\fitmath{
D(l_1,l_2)=p_1(q_{l_1}/X_1)p_2(q_{l_2}/X_2)
          -p_1(q_{l_1}/Y_1)p_2(q_{l_2}/Y_2).
}\tag{117}\] The rectangles have the same product, whole character, masks and slots. Their two profiles remain independent. To treat original profiles \(\fitinline{f_1,f_2}\) and independent original heights \(\fitinline{t_1,t_2}\), choose the common whole-product height \(\fitinline{t=t_0}\), with \(\fitinline{|t_0|\le V_{\rm ht}}\) as in (S6) and set \[\fitmath{
 p_j(x)=f_j(x)x^{i(t_j-t)},\qquad
 L_z(f_j,X_j;t_j)=X_j^{i(t_j-t)}L_z(p_j,X_j;t).
}\tag{117a}\] The original energy is unchanged by its outside unit-modulus phase. Use these same raised profiles \(\fitinline{p_j}\) on both rectangles in (117), retaining the outside phase belonging to the original \(\fitinline{X_j}\). Equivalently, the unraised reference with heights \(\fitinline{t_j}\) is multiplied by the exact scalar \[\fitmath{
 (X_1/Y_1)^{i(t_1-t)}(X_2/Y_2)^{i(t_2-t)}.
}\tag{117b}\] Thus the principal Mellin products cancel literally: after raising, the two terms have the same common norm power and the same profile integrals, while \(\fitinline{X_1X_2=Y_1Y_2}\). Equating the two independent heights is neither necessary nor allowed. Derivatives of the raised profiles have polynomial cost in both height differences, paid by the full logarithmic kernels and the common finite profile control. Prefix extractions retain their exact unit phases. Keep the whole rectangle difference until positive child norms form.

For actual-pool tuples \(\fitinline{\mathbf l=(p_i,l_1,l_2)}\) on the union of plain annuli, put \(\fitinline{J(\mathbf l)=(\prod p_i)l_1l_2}\). With common puncture \(\fitinline{R}\) and squarefree seed \(\fitinline{s}\), define

\hypertarget{sr:beta}{}

\[\fitmath{
\qquad\text{(118)}
 \beta_{R,s}(J)=
 \sum_{\substack{\mathbf l\\J(\mathbf l)=J}}
 \left(\prod_i b_i(p_i)\right)D(l_1,l_2)
 \mathbf1_{(J,R)=1}\mathbf1_{s\mid J}.
}\]

Initially \(\fitinline{R=s=1}\). Zeros remain zero; later masks/divisors are actual extractions, not new conditions on plain quotients.

\hypertarget{sr:comparison-bound}{} \textbf{Reference inlet.} The comparison bound is supplied with its actual geometric hypotheses, after the root is classified in Appendix C.10. LOW sets one reference scale to \(\fitinline{2/9}\); support in \(\fitinline{[1/4,9/4]}\) makes that factor exactly zero, so no reference estimate is required. HIGH uses the general common-product, two-orientation estimate proved in Appendix C.4, with the reference interval and root reserves proved in Appendix C.9 and C.10(20)--(26). Those arguments retain the natural reflection masks and polynomial height/profile bounds. They do not use the old narrow numerical rectangle or an already improved endpoint.

\hypertarget{sr:first}{}

\hypertarget{the-first-exact-poisson-source}{%
\paragraph{The first exact Poisson source}\label{the-first-exact-poisson-source}}

Let \(\fitinline{E_D}\) be the energy of the whole difference in \protect\hyperlink{sr:comparison}{(117)}. Positivity gives

\hypertarget{sr:positive-comparison}{}

\[\fitmath{
\qquad\text{(119)}
 E_X\le2E_D+2E_Y.
}\]

Enlarge only \(\fitinline{E_D}\) to all rows; never transform the keep predicate.

Fix nonnegative smooth \(\fitinline{\Psi}\) supported in \(\fitinline{[-1,b_\Phi+1]}\), \(\fitinline{\Psi\ge1}\) on \(\fitinline{[0,b_\Phi]}\), so \(\fitinline{0\le\Phi\le\mathcal D(\Phi)\Psi}\) for \(\fitinline{x\ge0}\). Before live subsets/labels, expand \protect\hyperlink{sr:beta}{(118)} in the full lattice square and extract complete common support \(\fitinline{Ca,Db}\), \(\fitinline{(a,b)=(ab,CD)=1}\). Define

\[\fitmath{
c=\log_Uq_C,\ d=\log_Uq_D,\ p=\log_Uq_{\operatorname{rad}(CD)},\
 A_0=\log_UV.
}\]

Let \(\fitinline{r}\) contain nonzero net common sextic exponents, length \(\fitinline{R_1}\), primitive character \(\fitinline{\xi_r}\). Expand the remaining mask by \(\fitinline{e\mid\operatorname{rad}(CD)/r}\), length \(\fitinline{E_1}\), so \(\fitinline{R_1+E_1\le p}\). For actual allocated convolutions \(\fitinline{A_C,A_D}\), nonzero frequency gives

\hypertarget{sr:first-poisson}{}

\[\fitmath{
\qquad\text{(120)}
 \begin{aligned}
 \frac{H_0}{Vq_e\sqrt{q_r}}\sum_{h\ne0}\mathfrak g_{\xi_r}(r,h)
 \sum_{(a,b)=1}&\frac{A_C(a)\overline{A_D(b)}}{\sqrt{q_aq_b}}
 \tau_C(a)\overline{\tau_D(b)}R(a,b)\\
 &\times\mathfrak g(a,h)\overline{\mathfrak g(b,-h)}
 \widehat\Psi\!\left(\frac{H_0q_h}{q_eq_rq_aq_b}\right).
 \end{aligned}
}\]

One orientation is \(\fitinline{\tau_C(a)=\tau(a)\chi_a(er)\overline{\xi_r(a)}}\), \(\fitinline{\tau_D(b)=\tau(b)\chi_b(er)\xi_r(b)}\). Keep in \(\fitinline{A_C}\) the exact frozen CRT scalar, \(\fitinline{\mu(e)}\) and fixed character values, including zeros. Convention-independently, if \(\fitinline{f_{C,D,e}(x)}\) is the actual active common character times both residual zero-extended characters at \(\fitinline{ex}\), define

\[\fitmath{
\mathcal T_{C,D,e}(h)=\sum_{x\bmod rab}
            f_{C,D,e}(x)e(hx/(rab)).
}\]

Poisson multiplies it by \(\fitinline{H_0/(q_eq_rq_aq_b)}\); CRT defines all displayed factors simultaneously, preserving unit phases. Three Gauss sums restore \(\fitinline{\sqrt{q_rq_aq_b}}\) and central normalization gives \(\fitinline{1/V}\). The cross factor \(\fitinline{\chi_a(br)\overline{\chi_b(ar)}\overline{\xi_r(ab)}}\) gives \(\fitinline{R(a,b)}\) by reciprocity. This proves the formula with nonunit zeros; changing orientation conjugates the whole character.

On retained smooth frequency dyads define \(\fitinline{\mathcal B_q}\) as \(\fitinline{V}\) times the normalized block \protect\hyperlink{sr:first-poisson}{(120)}, removing its \(\fitinline{1/V}\) once. Include all common labels, inactive subsets and dyadic coordinates. Its normalized contribution is \(\fitinline{\mathcal B_q/V}\) and physical mass is

\hypertarget{sr:physical-mass}{}

\[\fitmath{
\qquad\text{(121)}
 \mathcal M_{\rm phys}=\sum_q|\mathcal B_q|
   =\left|\sum_q\zeta_q\mathcal B_q\right|,\qquad
 \zeta_q=\begin{cases}|\mathcal B_q|/\mathcal B_q,&\mathcal B_q\ne0,\\
 0,&\mathcal B_q=0.\end{cases}
}\]

First labels thus stay in one joint sum, without an extra count.

The powerful zero-frequency argument of \S{}\protect\hyperlink{first-poisson-from-plain-products-to-gauss-polynomials}{7.2.6}, now with full product norm \(\fitinline{V}\), gives \(\fitinline{O(H_0U^\epsilon\mathcal H_J)}\) after divisor allocations and central \(\fitinline{1/V}\), retaining every mask.

The nominal first frequency length is

\hypertarget{sr:first-scale}{}

\textbf{(122)} is \protect\hyperlink{cf:eq:D.12}{(76)} under \(\fitinline{(A,R,E,m)=(A_0,R_1,E_1,M)}\): \(\fitinline{K_0=2A_0-c-d+R_1+E_1-M}\).

Retain whole dyads through each outer box's maximum scale with enclosure \(\fitinline{\delta_1<\xi}\). Discarded genuine arguments are \(\fitinline{\ge U^{\xi/4}}\) after fixed support errors. Use \(\fitinline{|\mathfrak g(a,h)|\le\sqrt{q_a}}\) and polynomial counts, fixing raw cost before Schwartz order. For every fixed \(\fitinline{L>0}\),

\hypertarget{sr:first-remainder}{}

\[\fitmath{
\qquad\text{(123)}
 E_D\le\mathcal D(\Phi)\{\mathcal M_{\rm phys}/V+
 C\mathcal H_JH_0(U^\epsilon+U^{-L})\}.
}\]

\hypertarget{sr:initial-reduction}{} \textbf{First-source reduction.} If the chosen comparison satisfies \(\fitinline{E_Y\ll\mathcal D(\Phi)\mathcal H_J (1+|t_1|+|t_2|+|t|+\max_i|v_i|)^J U^{M+\epsilon}}\), or is identically zero, then

\hypertarget{sr:initial-inequality}{} \[\fitmath{
 E_X\le2\mathcal D(\Phi)\mathcal M_{\rm phys}/V+
 C\mathcal D(\Phi)\mathcal H_J
 (1+|t_1|+|t_2|+|t|+\max_i|v_i|)^J U^{M+\epsilon}.
}\tag{124}\] \textbf{Proof.} Combine (119) with (123) and the stated comparison bound. The displayed Poisson identity is a finite expansion followed by the exact arithmetic Poisson formula of Appendix A.6, so it imposes no numerical affine interval on an ordinary exponent. On every actual root of Appendix C.10, its full product volume is at most \(\fitinline{Z}\), and the support cap is at most \(\fitinline{(9/4)^2 2^N}\) times that volume. All extracted quotients only shorten these supports. Thus the raw counts, Gauss bounds and profile derivatives in (123) have one fixed polynomial envelope. Fix that exponent before the Schwartz tail order. The zero-frequency proof in Appendix A.7.2.6 uses the same product volume and powerful-prime condition and retains every zero mask; it has no lower interval condition on either ordinary exponent.

The full kernels (125) absorb the independent height differences at a polynomial cost. The geometric mesh precedes the finite slot dictionary; after it is fixed, take the finite maximum over its subsets and the four plain derivative pairs. Assign half the loss to the comparison and remainder and absorb the finite logarithmic partitions in the remainder. These choices give one threshold for every live shortening, mask, ordinary profile and independent height. The precise LOW/HIGH comparisons are discharged in Appendix C.10, completing the conditional inlet in the full source proof. \ensuremath{\blacksquare}

\hypertarget{sr:kernel-budget}{}

\hypertarget{whole-kernels-and-allocation-normalizers}{%
\paragraph{Whole kernels and allocation normalizers}\label{whole-kernels-and-allocation-normalizers}}

On a retained dyad the nonseparated kernel contains

\[\fitmath{
\omega_K(q_h/U^K)\widehat\Psi(H_0q_h/(q_eq_rq_aq_b))
                  (q_aq_b)^{-1/2}.
}\]

Multiply by common log-row and whole-column cutoffs, shared by both rectangles. For the resulting compact logarithmic kernel,

\hypertarget{sr:full-kernel}{}

\[\fitmath{
\qquad\text{(125)}
 \mathcal K(\mathbf x)=\int_{\mathbb R^d}\widehat{\mathcal K}(\mathbf u)
       e^{2\pi i\mathbf u\cdot\mathbf x}\,d\mathbf u,\qquad
 \int(1+\|\mathbf u\|)^J|\widehat{\mathcal K}(\mathbf u)|\,d\mathbf u<\infty.
}\]

Bounded Euler derivatives of \(\fitinline{\widehat\Psi}\) and integration by parts of order \(\fitinline{>J+d}\) give live-column-independent seminorm control. Separate the second full kernel/roots in the same order.

Extend \protect\hyperlink{sr:first-poisson}{(120)} off coprimality using separate Gauss sums, fixed bicharacter and original masks. Insert complete \(\fitinline{\sum_{s\mid(a,b)}\mu(s)}\) before factorwise estimates; artificial primes gain no common-support restrictions. Expand the ray factor and apply joint Cauchy after \protect\hyperlink{sr:full-kernel}{(125)}, giving

\hypertarget{sr:gauss-column}{}

\[\fitmath{
\qquad\text{(126)}
 \sum_h\omega_K(q_h/U^K)
 \left|U^{-(A_0-c)/2}\sum_{s\mid a}
       A_C(a)\tau'_C(a)q_a^{it}\mathfrak g(a,h)\right|^2,
}\]

and similarly on the right. Set

\[\fitmath{
B_c=(3c-5d-R_1)_+/6,\qquad B_d=(3d-5c-R_1)_+/6,
}\]

\hypertarget{sr:target-exponents}{}

\[\fitmath{
\qquad\text{(127)}
 S_C=p+A_0-c+R_1+E_1+B_c,\quad
 S_D=p+A_0-d+R_1+E_1+B_d,\quad
 O_0=M-A_0-R_1/2-E_1.
}\]

By \protect\hyperlink{mf:eq:B3.56}{(MF3.56)}, \(\fitinline{O_0+(S_C+S_D)/2\le M}\). A fixed divisor width \(\fitinline{s_0}\) leaves column allowance \(\fitinline{A_0-c+R_1+E_1+B_c-s_0}\); one common-label count \(\fitinline{p+s_0}\) gives \(\fitinline{S_C}\). Keep that count inside joint weights.

Let \(\fitinline{j}\) index actual first columns, divisors and ray characters in a common box, with chosen common ideal \(\fitinline{C_j}\). An allocation \(\fitinline{B=(B_i,B_1,B_2)}\) extracts each factor's complete \(\fitinline{C_j}\)-part, \(\fitinline{\prod_iB_iB_1B_2=C_j}\). Freeze supplying slots; survivors keep coefficients. Divide both rectangle scales of each plain by the same \(\fitinline{q_{B_k}}\) and puncture all residual factors at complete support, even if exponents cancel.

Let \(\fitinline{F_B}\) be the actual frozen coefficient with profiles, \(\fitinline{M_B=\prod_{i\,\mathrm{frozen}}M_i}\) its bound. Retain \(\fitinline{F_B=0}\) before enlargement. For slot support bound \(\fitinline{b}\) and ambient count \(\fitinline{N}\), allocation Cauchy gives

\hypertarget{sr:allocation-weight}{}

\[\fitmath{
\qquad\text{(128)}
 \lambda_{j,B}=
 \begin{cases}
 |\rho_j|^2(\max(1,b))^N
       \dfrac{\#\operatorname{Alloc}(C_j)}{q_{C_j}}M_B^2,&F_B\ne0,\\
 0,&F_B=0.
 \end{cases}
}\]

Let \(\fitinline{I_B}\) be the surviving slots. The actual child volume and exact central scalar are

\[\fitmath{
V_{j,B}=\frac{X_1X_2}{q_{B_1}q_{B_2}}\prod_{i\in I_B}P_i,\qquad
 \frac{V_{j,B}}V=\frac1{q_{C_j}}\prod_{i\notin I_B}\frac{q_{p_i}}{P_i}.
}\]

Factoring \(\fitinline{F_B}\), the original normalized sum is exactly \(\fitinline{F_B\sqrt{V_{j,B}/V}}\) times the \(\fitinline{V_{j,B}^{-1/2}}\)-normalized child. Since \(\fitinline{q_{p_i}/P_i\le b}\), its squared scalar is \(\fitinline{\le(\max(1,b))^NM_B^2/q_{C_j}}\), yielding \protect\hyperlink{sr:allocation-weight}{(128)}. Actual child volume divides its energy; the first allocation weight stays undivided.

The scalar \(\fitinline{\rho_j}\) retains the first Poisson coefficient. For original complex block scalar \(\fitinline{a_j}\), including \protect\hyperlink{sr:physical-mass}{(121)}'s phase, an asymmetric split is

\[\fitmath{
\rho_{C,j}=a_jU^{-O_0-S_C/2},\qquad \rho_{D,j}=U^{-S_D/2}.
}\]

Five geometry coordinates vary by \(\fitinline{\le\log_U2}\) per box; the positive-part allowance varies by at most \(\fitinline{3/2}\) times this width. Cross-label Cauchy costs \(\fitinline{U^{5\log_U2}}\), a constant, after freezing common coordinates---not the square of the label count.

For extracted datum \(\fitinline{s}\), let \(\fitinline{\beta_s}\) be \protect\hyperlink{sr:beta}{(118)}'s residual coefficient without \(\fitinline{F_B}\), \(\fitinline{\tau_s}\) its whole character and \(\fitinline{V_s}\) its actual volume. Define

\hypertarget{sr:normalized-gauss}{}

\[\fitmath{
\qquad\text{(129)}
 G(s;\Omega,L)=\frac1{V_s}\sum_{h\in\mathcal O}\Omega(q_h/L)
 \left|\sum_u\beta_s(u)\tau_s(u)q_u^{it}\mathfrak g(u,h)\right|^2 .
}\]

The side budget is \(\fitinline{\mathcal L=\sum_{j,B}\lambda_{j,B}G(s_{j,B};\Omega,L_j)}\), with fixed annular majorant on \(\fitinline{[1/8,2]}\), including all columns/allocations.

The column window has an independent Mellin coordinate: left/conjugate-right raw heights are \(\fitinline{t+2\pi(u-v)}\), \(\fitinline{t+2\pi(u+v)}\). Window densities have common finite weighted \(\fitinline{L^1}\) norms by integration by parts; Minkowski integrates the polynomial-height \(\fitinline{\mathcal L}\) bound through both full windows. Bounded row-profile parameters and partition scalars contract the joint norm.

\hypertarget{sr:amplifier}{}

\hypertarget{uniform-analytic-closure}{%
\paragraph{Uniform analytic closure}\label{uniform-analytic-closure}}

Choose geometric cells, mesh, strict scalar reserves and support padding before live scales, and the mesh before the dictionary. Then choose one finite profile order over subsets and the four derivative pairs, with no error charged separately to each slot. Choose seminorm and height orders backwards through both full kernels and functional equations, before live labels and heights.

Integrate retained heights against (125)'s full densities. For external cutoffs, fix the height degree before the tail order that dominates it and the raw counts. Discard genuine Poisson tails before separation, using their arguments at least \(\fitinline{U^{\xi/4}}\), \(\fitinline{|\mathfrak g|\le\sqrt{q_u}}\), and \(\fitinline{|F|\le q_uq_v}\). This gives arbitrary decay with both independent ordinary heights.

The physical mass (121) and the first allocation weights (128) enter the joint divisor and physical-column budget proved in Appendix C.8. Its diagonal, fixed-inducing and tail parts are proved in Appendix C.9; Appendix C.10(30)--(35) converts the remaining packed price and combines them with that budget. This contracts the actual first labels once, keeps the artificial divisor inside joint Cauchy, and retains all zeros, original masks, prefixes and shortened volumes. No obsolete numerical amplifier budget or additional density payment is a premise of this route.

\hypertarget{c.6.-masked-packed-twins-and-inducing-invariance}{%
\subsubsection{C.6. Masked packed twins and inducing invariance}\label{c.6.-masked-packed-twins-and-inducing-invariance}}

Fix either first-side orientation. Let \(\fitinline{C_s}\) be its first common ideal, \(\fitinline{Z_0}\) its exponent-zero prime product, \(\fitinline{V_0\ne0}\) its packed multiplier, and \(\fitinline{j_0}\) its finite dictionary index. The required first-side facts are \[\fitmath{
 Z_0\mid C_s,\qquad C_s\mid R,
}\tag{T1}\] and, for supported \(\fitinline{I}\) coprime to \(\fitinline{Z_0}\), \[\fitmath{
 h_{\mathrm{base}}(t,I)
 =h_{\mathrm{anchor}(\eta,j_0)}(t,I)\chi_I(V_0)
 \quad(t\in\mathbb R).
}\tag{T2}\]

Take the actual retained second family over \(\fitinline{\mathrm{childCharacter}(\mathrm{base},\xi)}\). Its supplying contract includes genuine Hecke characters \(\fitinline{\tau_\chi}\) and \[\fitmath{
 h_{\tau_\chi}(t,I)
 =
 h_{\mathrm{childCharacter}(\mathrm{base},\xi)}(t,I)
 \chi_I(A_2)\chi(\mathrm{primary}(I))
}\tag{T3}\] for supported \(\fitinline{I}\) satisfying \(\fitinline{(C_2,I)=1}\) or \(\fitinline{(D_2,I)=1}\), uniformly for every \(\fitinline{t}\). Here \[\fitmath{
 A_2=\mathrm{commonFrequencyGenerator}(C_2,D_2)
       \,\mathrm{nonunitFrequencyGenerator}(C_2,D_2,U)\ne0.
}\] The family also supplies the all-element row realization identifying its original nonexceptional predicate with \[\fitmath{
 z\ne0\ \land\
 \neg\mathrm{FixedInducingRow}_{\tau_\chi,Q}(m_0,1,z),
 \qquad m_0=\mathrm{fixedBadMask}.
}\tag{T4}\] This use requires the original mask \(\fitinline{m\ne0}\), with \(\fitinline{\lambda\mid m}\) and \(\fitinline{2\mid m}\), as in that realization.

Put \[\fitmath{
 j=j_0\xi\chi,\qquad \sigma=\mathrm{anchor}(\eta,j),
 \qquad V=V_0A_2\ne0.
}\] Combining (T2), (T3) and the finite-ray multiplication identity gives \[\fitmath{
 h_{\tau_\chi}(t,I)=h_\sigma(t,I)\chi_I(V)
}\tag{T5}\] when \(\fitinline{I}\) avoids \(\fitinline{Z_0}\) and either second common ideal. This remains only a masked identity.

\hypertarget{existence-of-the-twin}{%
\subparagraph{Existence of the twin}\label{existence-of-the-twin}}

Choose a genuine Hecke character \(\fitinline{\widetilde\tau}\), once for \(\fitinline{(\sigma,V)}\), with \[\fitmath{
 a_{\widetilde\tau}(I)=a_\sigma(I)\chi_I(V)
 \quad\text{for every supported }I.
}\tag{T6}\] It exists by the retained second-height construction with base \(\fitinline{\sigma}\), trivial additional ray character and multiplier \(\fitinline{V}\). Its hypotheses are precisely: \(\fitinline{\sigma}\) is a genuine Hecke character, \(\fitinline{V\ne0}\), and \(\fitinline{m_0\ne0}\) contains \(\fitinline{\lambda}\) and \(\fitinline{2}\).

To see the construction, realize the multiplicative periodic element row \[\fitmath{
 n\longmapsto
 \mathrm{rowTwist}_{\mathrm{childCharacter}(\sigma,1)}
                         (m_0,1,V;n)
}\] as a finite-modulus, unit-invariant character. On a supported ideal, the \(\fitinline{m_0^6}\) factor equals one with its original coprimality mask, and the trivial ray factor equals one. This proves (T6), and multiplication by \(\fitinline{(NI)^{it}}\) proves its height version for \textbf{all} \(\fitinline{t}\). The choice is independent of profiles and heights.

Its explicit conductor bound is \[\fitmath{
 N\mathfrak q_{\widetilde\tau}
 \le
 N\mathfrak q_{\mathrm{childCharacter}(\sigma,1)}
 N(m_0)N(72)\,N(V).
}\tag{T7}\] Thus the moving \(\fitinline{N(V)}\) is retained. This construction introduces no additional independent character-width charge after packing.

\hypertarget{why-the-original-masked-polynomial-agrees-with-the-twin}{%
\subparagraph{Why the original masked polynomial agrees with the twin}\label{why-the-original-masked-polynomial-agrees-with-the-twin}}

Let \(\fitinline{G=C_2}\) or \(\fitinline{G=D_2}\), and let \(\fitinline{c(J)}\) denote the \textbf{original} finite column coefficient, with its cube-free test, seed, ordinary and marked holes, prefixes and centered profiles intact. Define \[\fitmath{
 b(I)=1_{(G,I)=1}\,1_{D\mid I}\,c(GI).
}\tag{T8}\] A nonzero original column satisfies \[\fitmath{
 c(J)\ne0\Longrightarrow(J,R)=1.
}\tag{T9}\] Indeed, its finite tuple sum then has a nonzero summand, and that summand contains the original coprimality indicator against \(\fitinline{R}\).

By (T1), a nonzero coefficient in (T8) therefore has \[\fitmath{
 (I,Z_0)=1,\qquad(G,I)=1.
}\] Consequently (T5) and (T6) imply \[\fitmath{
 b(I)h_{\tau_\chi}(v,I)
 =b(I)h_{\widetilde\tau}(v,I)
 \quad\text{for every residual column and every }v.
}\tag{T10}\] If a mask fails, or \(\fitinline{c(GI)=0}\), both sides vanish. Thus the two complete original masked polynomials agree. No column mask has yet been removed.

\hypertarget{exact-fixed-inducing-invariance}{%
\subparagraph{Exact fixed-inducing invariance}\label{exact-fixed-inducing-invariance}}

The comparison also preserves the physical keep predicate. The needed finite-mask statement is:

\begin{quote}
If genuine characters \(\fitinline{\rho,\rho'}\) satisfy \[\fitmath{
a_{\rho'}(n)=1_{(n,w)=1}a_\rho(n),\qquad w\ne0,
}\] then existence of a primitive inducing character whose conductor divides the fixed ideal \(\fitinline{Q}\) is equivalent for \(\fitinline{\rho}\) and \(\fitinline{\rho'}\).
\end{quote}

Here is the descent argument. The forward implication from an inducing character of \(\fitinline{\rho}\) to the masked character follows directly from the zero-extension formula. Conversely, suppose the masked character is induced from a primitive character of modulus \(\fitinline{\mathfrak f}\). For a unit residue of \(\fitinline{\rho}\)'s modulus that is \(\fitinline{1}\) modulo the image of \(\fitinline{\mathfrak f}\), choose a representative \(\fitinline{a}\) with these two congruences. It is coprime to both moduli. By changing \(\fitinline{a}\) by a multiple of their product, CRT permits a representative also coprime to \(\fitinline{w}\): at primes dividing that product coprimality is automatic, and at every other prime divisor of \(\fitinline{w}\) one avoids one residue class. On this representative, the character value is \(\fitinline{1}\). Hence \(\fitinline{\rho}\) descends through the image of \(\fitinline{\mathfrak f}\). Its primitive presentation has conductor dividing \(\fitinline{\mathfrak f}\), and therefore dividing \(\fitinline{Q}\). This proves the equivalence.

For row characters this statement applies when \[\fitmath{
 m\ne0,\quad w\ne0,\quad Az\ne0,\quad
 \lambda\mid m,\quad2\mid m.
}\] The row construction supplies genuine characters for both masks, and \[\fitmath{
 \mathrm{rowTwist}(mw,1,Az;n)
 =1_{(n,w)=1}\mathrm{rowTwist}(m,1,Az;n).
}\] Thus adding or deleting the finite mask \(\fitinline{w}\) preserves the literal \texttt{F\allowbreak{}i\allowbreak{}x\allowbreak{}e\allowbreak{}d\allowbreak{}I\allowbreak{}n\allowbreak{}d\allowbreak{}u\allowbreak{}c\allowbreak{}i\allowbreak{}n\allowbreak{}g\allowbreak{}R\allowbreak{}o\allowbreak{}w} predicate, with the same \(\fitinline{Q}\).

Apply this with a nonzero generator \(\fitinline{w}\) of \(\fitinline{Z_0G}\). Away from \(\fitinline{w}\), (T5)--(T6) identify the two base coefficients. After adjoining \(\fitinline{w}\) to \(\fitinline{m_0}\), the complete row functions agree: supported unmasked terms agree, masked terms are zero, and unsupported terms are killed by \(\fitinline{m_0}\). Mask invariance then gives \[\fitmath{
 \mathrm{FixedInducingRow}_{\tau_\chi,Q}(m_0,1,z)
 \iff
 \mathrm{FixedInducingRow}_{\widetilde\tau,Q}(m_0,1,z).
}\tag{T11}\] For the twin, (T6) gives directly \[\fitmath{
 \mathrm{FixedInducingRow}_{\widetilde\tau,Q}(m_0,1,z)
 \iff
 \mathrm{FixedInducingRow}_{\sigma,Q}(m_0,1,Vz).
}\tag{T12}\] At \(\fitinline{z=0}\) both nonfixed keeps are false. Hence, for every \(\fitinline{z}\), \[\fitmath{
 \mathrm{keep}_{\tau_\chi}(z)
 =\mathrm{keep}_{\widetilde\tau}(z)
 =\mathrm{initialKeep}_{\sigma,Q}(Vz).
}\tag{T13}\] The fixed unit-generator factor in the definition of \texttt{i\allowbreak{}n\allowbreak{}i\allowbreak{}t\allowbreak{}i\allowbreak{}a\allowbreak{}l\allowbreak{}K\allowbreak{}e\allowbreak{}e\allowbreak{}p} is removed by the same exact mask invariance.

Equations (T10) and (T13) prove equality of the literal label energies for all \(\fitinline{v}\), without an inequality or an analytic loss.

\hypertarget{order-of-height-absorption-packing-and-restoration}{%
\subparagraph{Order of height absorption, packing and restoration}\label{order-of-height-absorption-packing-and-restoration}}

Use the following order: \[\fitmath{
 \boxed{
 \text{original masked }\tau_\chi\text{-row}
 \ \xrightarrow{\text{height absorption}}\
 \text{same masked row}
 \ \xrightarrow{\text{T10--T13}}\
 \text{masked twin row}
 \ \xrightarrow{y=Vz}\
 \text{masked }\sigma\text{-row}
 \ \xrightarrow{\text{restoration}}\
 \text{native }\sigma\text{-leaves}.}
}\tag{T14}\] The height absorption retains the phase \[\fitmath{
 \left(\frac{X_1X_2}{Ne_1Ne_2\,NG}\right)^{iv},
}\] shifts the slot height to \(\fitinline{t+v}\), and shifts each ordinary profile parameter to \(\fitinline{t_j+v/(2\pi)}\).

The twin satisfies the restoration packing interface with auxiliary mask \(\fitinline{1}\): its supported coefficient identity is global, and \[\fitmath{
 (G,I)=1\Longrightarrow(I,(1))=1
}\] is automatic. The temporary comparison mask \(\fitinline{Z_0G}\) is \textbf{not} used as that auxiliary mask.

Multiplication by \(\fitinline{V}\) gives an exact bijection onto \(\fitinline{V\mathcal O}\). Therefore, with the original masked polynomial still present, \[\fitmath{
 \sum_{\mathrm{keep}_{\widetilde\tau}(z)}
 |\mathcal R^{\rm orig}_{\widetilde\tau,z}|^2\Phi(Nz/K)
 =
 \sum_{\substack{V\mid y\\\mathrm{initialKeep}_{\sigma,Q}(y)}}
 |\mathcal R^{\rm orig}_{\sigma,y}|^2
 \Phi\!\left(\frac{Ny}{NVK}\right).
}\tag{T15}\] There is no Jacobian or additional power loss.

\hypertarget{part-viii-2}{}

\hypertarget{c.7.-cube-free-restoration-and-fixed-inducing-contribution}{%
\subsubsection{C.7. Cube-free restoration and fixed-inducing contribution}\label{c.7.-cube-free-restoration-and-fixed-inducing-contribution}}

Write \(\fitinline{\mathcal O=\mathbb Z[\omega]}\), and write \(\fitinline{NI}\) for the absolute norm of a nonzero integral ideal. Every character below has its original zero extension. In particular, a coefficient that vanishes because its ideal meets the presented modulus remains zero after every identity.

Fix a finite slot set \(\fitinline{\mathcal I}\), a common prime profile \(\fitinline{W}\) supported in \(\fitinline{[1,2]}\), and a compact ordinary-profile interval \(\fitinline{[a,b]}\), where \(\fitinline{a>0}\). The genuine slot coefficients have the form \[\fitmath{
 \beta_i(\mathfrak p)=
 \nu_i(\mathfrak p)\,
 W(N\mathfrak p/P_i)
 (N\mathfrak p/P_i)^{-\sigma_i+i f_i},
}\] with the original relative finite-ray character \(\fitinline{\nu_i}\), its natural zeros, its finite hole set, and its inherited norm height. The real powers \(\fitinline{\sigma_i}\) lie in a fixed compact interval. Consequently \[\fitmath{
 |\beta_i(\mathfrak p)|\leq B_f
}\tag{1}\] with one \(\fitinline{B_f\geq1}\) independent of the slot frequencies and inherited heights.

All constants in the restoration statements are chosen after the fixed field, slot number, profile \(\fitinline{W}\), real-power interval, hole-cardinality bound and fixed polynomial caps, but before the moving characters, prime pools, rows, ordinary profiles, heights, allocations and positive shortened scales.

\hypertarget{exact-local-and-labelled-identities}{%
\subparagraph{1. Exact local and labelled identities}\label{exact-local-and-labelled-identities}}

Put \[\fitmath{
 F=\frac1{(1-x)(1-y)},\qquad B=(1-x)(1-y),
}\] \[\fitmath{
 T_3=1+x+y+x^2+xy+y^2,\qquad H_3=BT_3.
}\] Multiplication gives \[\fitmath{
 H_3
 =1-(x^3+x^2y+xy^2+y^3)
   +(x^3y+x^2y^2+xy^3).
}\tag{2}\] Thus \(\fitinline{T_3}\) represents the literal condition \[\fitmath{
 v_{\mathfrak p}(I)+v_{\mathfrak p}(J)<3.
}\] Here \(\fitinline{I,J}\) are the ordinary indices of the stored source. Any original prefix ideals remain in their profiles and extraction scalars; they are not silently inserted into, or removed from, this condition.

Let \(\fitinline{\zeta^2=0}\) represent one supplying marked prime. The unrestricted marked local factor is \(\fitinline{F+\zeta F}\); the restricted factor is \(\fitinline{T_3+\zeta}\), because a supplying prime must divide neither ordinary index. Define \[\fitmath{
 K_3=H_3+\zeta(B-H_3).
}\] Then \[\fitmath{
 K_3(F+\zeta F)=T_3+\zeta,
}\tag{3}\] and \[\fitmath{
 K_3^{-1}
 =H_3^{-1}+\zeta(H_3-B)H_3^{-2},
\qquad
 K_3^{-1}(T_3+\zeta)=F+\zeta F.
}\tag{4}\] For example, the marked coefficient in (3) is \(\fitinline{H_3F+(B-H_3)F=BF=1}\). In (4), it is \[\fitmath{
 H_3^{-1}+(H_3-B)H_3^{-2}T_3=F,
}\] using \(\fitinline{H_3=BT_3}\).

These are formal-series identities: analytic division is unnecessary. For the labelled global identity, each prime either supplies no marked label or supplies exactly one label. The corresponding supplying map is injective. Square-zero notation alone does not assert that different marked labels choose different primes.

The forward marked polynomial is \[\fitmath{
 B-H_3
 =-x-y+xy+x^3+x^2y+xy^2+y^3
   -x^3y-x^2y^2-xy^3.
}\tag{5}\] The inverse marked series is exactly \[\fitmath{
 (H_3-B)H_3^{-2}.
}\tag{6}\] It multiplies \(\fitinline{T_3}\) in (4). Replacing it by a coefficient obtained from a punctured ordinary series would change the source.

Taking coefficients and then products over the finite prime enclosure proves the ideal-pair divisor identities. Reindexing by the extracted ideal pair \(\fitinline{d=(d_1,d_2)}\) gives the ordinary scales \[\fitmath{
 X_1/Nd_1,\qquad X_2/Nd_2.
}\tag{7}\] At a kernel-frozen marked prime, the remaining marked coefficients acquire the cross-assignment hole excluding that prime. In the inverse identity, this does \textbf{not} introduce an extra ordinary puncture there: the remaining restricted ordinary local factor is \(\fitinline{T_3}\).

Every coefficient, finite-ray phase, original natural zero and prefix is retained in these finite identities.

\hypertarget{central-inverse-norm}{%
\subparagraph{2. Central inverse norm}\label{central-inverse-norm}}

Write \[\fitmath{
 H_3^{-1}=\sum_{i,j\geq0}q_{ij}x^iy^j.
}\] The constant coefficient is \(\fitinline{q_{00}=1}\). Since the seven nonconstant monomials of \(\fitinline{H_3}\) have total degree at least three, \[\fitmath{
 q_{ij}=0\quad(0<i+j<3).
}\tag{8}\] There is no larger rectangular zero region.

The recurrence obtained from \(\fitinline{H_3H_3^{-1}=1}\) also gives \[\fitmath{
 |q_{ij}|\leq2^{i+j}.
}\tag{9}\] Indeed, induct on \(\fitinline{d=i+j}\). The assertion is immediate for \(\fitinline{d=0,1,2}\). For \(\fitinline{d\geq3}\), the recurrence contains at most seven preceding coefficients, each of degree at most \(\fitinline{d-3}\). Their total absolute value is at most \(\fitinline{7\,2^{d-3}\leq2^d}\).

For \(\fitinline{0\leq z\leq2/5}\), put \(\fitinline{r=2z\leq4/5}\). Equations (8)--(9) imply \[\fitmath{
\begin{aligned}
 \sum_{i,j}|q_{ij}|z^{i+j}
 &\leq1+\sum_{n\geq3}(n+1)r^n\\
 &=1+\frac{r^3(4-3r)}{(1-r)^2}\\
 &\leq1+320z^3
 \leq1+400z^3.
\end{aligned}
}\tag{10}\] The penultimate inequality follows from \[\fitmath{
 \frac{4-3r}{(1-r)^2}
 =(1-r)^{-2}+3(1-r)^{-1}\leq25+15.
}\] This proves absolute summability of the entire reciprocal, uniformly over every finite coefficient cutoff.

For the present supported primes, \[\fitmath{
 N\mathfrak p\equiv1\pmod6,\qquad N\mathfrak p\geq7.
}\] Hence \(\fitinline{z=(N\mathfrak p)^{-1/2}\leq2/5}\), and \[\fitmath{
 \sum_{i,j}|q_{ij}|(N\mathfrak p)^{-(i+j)/2}
 \leq1+400(N\mathfrak p)^{-3/2}.
}\tag{11}\] Unique ideal factorization and positivity therefore give, for every finite family of supported ideal pairs, \[\fitmath{
 \sum_d\frac{|q(d_1,d_2)|}{\sqrt{Nd_1Nd_2}}
 \leq
 \exp\!\left(400\sum_{\mathfrak p\ {\rm supported}}
                 (N\mathfrak p)^{-3/2}\right)
 =:C_{\rm inv}<\infty.
}\tag{12}\] Restrictions on the ideal pairs only reduce this positive majorant.

The inverse marked series must also be included. The ten monomials of \(\fitinline{H_3-B}\) have absolute coefficient one and positive total degree. Thus \[\fitmath{
 \|H_3-B\|_z\leq10z\leq14z.
}\] By (10), \(\fitinline{\|H_3^{-1}\|_z\leq27}\) for \(\fitinline{z\leq2/5}\). Submultiplicativity of absolute coefficient mass yields \[\fitmath{
 \|(H_3-B)H_3^{-2}\|_z
 \leq14z\cdot27^2=10206z.
}\tag{13}\] The supplying slot contributes one further central factor \(\fitinline{z}\). Its complete local inverse mass is consequently at most \[\fitmath{
 10206(N\mathfrak p)^{-1}.
}\tag{14}\] For a fixed number \(\fitinline{m}\) of frozen slots, summing all their prime choices costs an arbitrarily small power of the prime enclosure. For instance, for every \(\fitinline{\epsilon>0}\), choose \(\fitinline{\delta=\epsilon/(m+1)}\). If \(\fitinline{N\mathfrak p\leq X}\), then \[\fitmath{
 \sum_{\mathfrak p\leq X}(N\mathfrak p)^{-1}
 \leq X^\delta\sum_{\mathfrak p}
                  (N\mathfrak p)^{-1-\delta}.
}\] The \(\fitinline{m}\)-fold assignment sum is \(\fitinline{O_{\epsilon,m}(X^\epsilon)}\). Allowing optional assignments, collisions or additional forbidden primes in this \textbf{positive majorant} preserves the upper bound. The exact source continues to use its actual admissible assignments.

Combining (12)--(14) proves the complete central inverse norm, including every frozen marked correction and every finite cutoff. No reciprocal coefficient is omitted.

\hypertarget{reflected-forward-norm-and-the-prime-enclosure}{%
\subparagraph{3. Reflected forward norm and the prime enclosure}\label{reflected-forward-norm-and-the-prime-enclosure}}

For an ordinary monomial \(\fitinline{x^iy^j}\), central extraction followed by the actual reflected length ratio gives \[\fitmath{
 (N\mathfrak p)^{-(i+j)/2}
 (N\mathfrak p)^{(j-i)/6}
 =(N\mathfrak p)^{-(2i+j)/3}.
}\tag{15}\] The monomial \(\fitinline{y^3}\) therefore has weight \(\fitinline{(N\mathfrak p)^{-1}}\). The sum of the eight absolute local weights is \[\fitmath{
 1+p^{-1}+p^{-4/3}+2p^{-5/3}+2p^{-2}+p^{-7/3}
 \leq1+7p^{-1},
 \qquad p=N\mathfrak p.
}\tag{16}\]

We prove the required logarithmic prime-pool estimate directly. Fixed-field ideal counting gives \[\fitmath{
 \#\{I\ne0:NI\leq H\}\leq128H\qquad(H\geq1).
}\] A dyadic decomposition then gives, for \(\fitinline{0<t\leq1}\), \[\fitmath{
 \sum_I(NI)^{-1-t}
 \leq\frac{256}{1-2^{-t}}
 \leq\frac{768}{t}.
}\tag{17}\] The last inequality follows from \(\fitinline{1/(1-2^{-t})\leq3/t}\).

Let \(\fitinline{X\geq3}\), \(\fitinline{s=1+1/\log X}\), and let \(\fitinline{\mathcal P}\) be any finite set of nonzero prime ideals with \(\fitinline{N\mathfrak p\leq X}\). Distinct subsets of \(\fitinline{\mathcal P}\) have distinct ideal products, so \[\fitmath{
 \prod_{\mathfrak p\in\mathcal P}
       \bigl(1+(N\mathfrak p)^{-s}\bigr)
 \leq\sum_I(NI)^{-s}
 \leq768\log X.
}\] Also \[\fitmath{
 (N\mathfrak p)^{-1}\leq3(N\mathfrak p)^{-s},
 \qquad
 u/2\leq\log(1+u)\quad(0\leq u\leq1).
}\] Taking logarithms proves \[\fitmath{
 \sum_{\mathfrak p\in\mathcal P}(N\mathfrak p)^{-1}
 \leq6\log\log X+6\log768.
}\tag{18}\] Equations (16)--(18) imply the explicit source bound \[\fitmath{
 \prod_{\mathfrak p\in\mathcal P}
  \sum_{\text{eight monomials}}|\text{reflected weight}|
 \leq
 e^{42\log768}(\log X)^{42}.
}\tag{19}\]

For a marked monomial \(\fitinline{x^iy^j}\) of (5), the supplying factor and reflected ratio instead give \[\fitmath{
 p^{-(i+j+1)/2}p^{(j-i-1)/6}
 =p^{-(2i+j+2)/3}.
}\tag{20}\] Every exponent on the right is at least one. Thus the full marked local weight is at most \(\fitinline{10/p}\), and the retained source constant \(\fitinline{14/p}\) is valid.

For a fixed supplying map, multiply the ordinary product (19) by the corresponding marked weights. Summing its fixed number of labels adds at most a fixed power of the prime harmonic sum. Consequently every complete forward critical mass is bounded by \[\fitmath{
 C(1+\log X)^{42}
       (1+\log\log(3X))^{|\mathcal I|}.
}\tag{21}\] In particular, it is \(\fitinline{O_\epsilon(X^\epsilon)}\), uniformly over all admissible finite prime sets and ideal cutoffs.

The actual reflected ratio uses the \textbf{whole} frozen marked product: \[\fitmath{
 r(d,\mathbf p)=
 \left(\frac{Nd_2}
 {Nd_1\prod_{i\ {\rm frozen}}N\mathfrak p_i}\right)^{1/6}.
}\tag{22}\] Multiplying a central scalar majorant by \(\fitinline{r}\) changes its ordinary exponents to \(\fitinline{2/3,1/3}\) and every supplying exponent to \(\fitinline{2/3}\), exactly as in (15) and (20).

If the child energy is bounded by \(\fitinline{E(3+r^2)}\), then \[\fitmath{
 \sqrt{3+r^2}\leq2(1+r).
}\] Minkowski in the positive physical norm therefore bounds the parent energy by \[\fitmath{
 4\left(\sum_{\rm actual\ allocations}
            w_{\rm central}(1+r)\right)^2E.
}\tag{23}\] This sums mass-weighted child norm bounds; it does not pay the number of prime assignments.

For the application, every pool lies in \[\fitmath{
 N\mathfrak p\leq\mathrm{Cap}\,Z^{L_{\rm cap}},
}\tag{24}\] where \(\fitinline{\mathrm{Cap}}\) and \(\fitinline{L_{\rm cap}}\) are fixed before \(\fitinline{Z}\). Given a final energy loss \(\fitinline{\eta>0}\), use (21) with an exponent sufficiently smaller than \(\fitinline{\eta/[2(L_{\rm cap}+1)]}\). Squaring in (23), and absorbing the fixed power of \(\fitinline{\mathrm{Cap}}\), costs at most \(\fitinline{CZ^\eta}\).

Thus (24) is an essential hypothesis of the reflected transfer. A prime-pool-independent constant is unavailable.

\hypertarget{complete-restoration-collisions-and-natural-holes}{%
\subparagraph{4. Complete restoration, collisions and natural holes}\label{complete-restoration-collisions-and-natural-holes}}

For completeness, the positive-norm passage used above is elementary. If \[\fitmath{
 F_u=\sum_jc_j(u)G_j(u),\qquad |c_j(u)|\leq w_j,\quad w_j\geq0,
}\] and \(\fitinline{\rho(u)\geq0}\), weighted Cauchy gives \[\fitmath{
 \sum_u\rho(u)|F_u|^2
 \leq\left(\sum_jw_j\right)
       \sum_jw_j\sum_u\rho(u)|G_j(u)|^2.
}\tag{25}\] The majorants \(\fitinline{w_j}\) are independent of the physical row. All row-dependent character phases remain in \(\fitinline{c_j(u)}\).

The finite forward and inverse identities of \S{}1 supply precisely such coefficients. Their majorants are proved in \S{}\S{}2--3. Every child has a slot subset, the original ordinary profiles, the original character, and positive scales no greater than its parent scales. A shortened scale below one is retained literally; if its upper profile support contains no positive ideal norm, the child is zero.

Repeated prime slots are handled by the equality partition of the original slot map. For a block \(\fitinline{A}\) of \(\fitinline{m\geq2}\) equal slots, its complete coefficient is \[\fitmath{
 \prod_{i\in A}
 \frac{\beta_i(\mathfrak p)\chi_u(\mathfrak p)}{\sqrt{P_i}}.
}\] On its nonzero annuli, \(\fitinline{N\mathfrak p\leq2P_i}\), hence \[\fitmath{
 \left|\prod_{i\in A}
 \frac{\beta_i(\mathfrak p)\chi_u(\mathfrak p)}{\sqrt{P_i}}\right|
 \leq(2B_f)^m(N\mathfrak p)^{-m/2}
 \leq(2B_f)^m(N\mathfrak p)^{-1}.
}\tag{26}\] The original powers of the character and all coefficients remain until this positive bound is taken. Summing repeated-block primes costs a fixed power of (18). Singleton slots retain the exclusions against every frozen repeated block. Conversely, inclusion--exclusion on the finite partition lattice expresses an injective singleton polynomial in full natural products and repeated-block corrections, again with (26). The finite partition coefficients depend only on \(\fitinline{|\mathcal I|}\).

For a common ordinary puncture \(\fitinline{\mathfrak R}\), the exact formula is \[\fitmath{
 P_{\chi,\mathfrak R}(X;W)
 =\sum_{D\mid\operatorname{rad}\mathfrak R}
   \frac{\mu(D)\chi(D)}{\sqrt{ND}}\,
   P_\chi(X/ND;W).
}\tag{27}\] The character's existing zeros are unchanged. A marked hole is removed by subtracting the corresponding actual prime terms; finite holes produce slot subsets and retain their original coefficients.

If \(\fitinline{N\operatorname{rad}\mathfrak R\leq Z^{B_{\rm mask}}}\), the total coefficient mass in (27), together with these marked deletions, is \(\fitinline{O_\eta(Z^\eta)}\). To see this without a uniform-small-prime assumption, split the prime divisors at a fixed threshold. At large primes each local factor \(\fitinline{1+C p^{-u}}\), \(\fitinline{u>0}\), is at most \(\fitinline{p^\delta}\); the finitely many smaller prime factors contribute a fixed constant. The same argument bounds \(\fitinline{C^{\omega(\mathfrak R)}}\) by \(\fitinline{C_\delta(N\mathfrak R)^\delta}\).

Combining these exact decompositions with (25) proves:

\begin{quote}
\textbf{Complete cube-free restoration.} Suppose a common bound \(\fitinline{E\geq0}\) holds for every literal restricted child arising from every equality partition, repeated-block assignment, frozen supplying assignment, resulting cross-hole, slot subset, and every positive shortened ordinary pair. Then the original full natural product has physical energy at most \(\fitinline{CZ^\eta E}\). The constant precedes all those choices.
\end{quote}

This is a transfer statement with its child family displayed. It supplies no estimate for unrelated coefficients or inadmissible lengths.

\hypertarget{quantitative-reduction-of-an-actual-second-label}{%
\subparagraph{5. Quantitative reduction of an actual second label}\label{quantitative-reduction-of-an-actual-second-label}}

We now prove the bridge requested for collection.

Fix one actual second label \(\fitinline{(G,D,e)}\), where \(\fitinline{D=t_2}\) is squarefree, \(\fitinline{G\ne0}\), \(\fitinline{e_1,e_2\ne0}\), and the original seed satisfies \(\fitinline{\mathrm{seed}\mid G}\). Retain the literal common puncture \(\fitinline{R\ne0}\), ordinary-hole set \(\fitinline{H}\), slot-hole sets \(\fitinline{R_{0,i}}\), and the original centered finite-column coefficient.

The fixed parameters satisfy \[\fitmath{
a>0,\quad b\ge1,\quad
0\le L_{\rm slot}\le L_{\rm cap},\quad
L\le L_{\rm cap},\quad
\mathrm{Cap}\ge\max(2,b),\quad B_{\rm mask}\ge0.
}\] The common slot profile is smooth and supported in \(\fitinline{[1,2]}\). The slot data satisfy \[\fitmath{
P_i=Z^{w_i},\qquad 0\le w_i\le L_{\rm slot},
\qquad \mathrm{lo}\le\sigma_i\le\mathrm{hi},
\qquad |R_{0,i}|\le\mathrm{Holes}.
}\] The ordinary profiles \(\fitinline{W_1,W_2}\) are Schwartz functions supported in \(\fitinline{[a,b]}\). All four ordinary scales are positive, at most \(\fitinline{Z^L}\), and \[\fitmath{
X_1X_2=Y_1Y_2=T,\qquad
\mathcal X=T\prod_iP_i.
}\] Put \(\fitinline{\mathcal S=\operatorname{enclosure}(\mathrm{Cap}Z^{L_{\rm cap}})}\). Require \[\fitmath{
H\subseteq\mathcal S,\qquad
\mathfrak p\in H\Longrightarrow\mathfrak p\in R_{0,i}
\quad\text{for every live slot }i,
}\tag{28a}\] and \[\fitmath{
N\!\left(\prod_{\mathfrak p\in
\operatorname{primeSupport}(RGD)\,\cup\,\operatorname{holes}(H)}
\mathfrak p\right)\le Z^{B_{\rm mask}}.
}\tag{28b}\] The complete ordinary-prime enclosure is retained. Explicitly, every supported ideal in either ordinary support, including either rectangle and every subsequently shortened support, has its prime support in \(\fitinline{\mathcal S}\). This follows from the displayed caps: for example, \[\fitmath{
NI\le b\,X_j/NE\le \mathrm{Cap}Z^{L_{\rm cap}},
}\] since \(\fitinline{NE\ge1}\); each prime factor has norm at most \(\fitinline{NI}\). The same reasoning covers the \(\fitinline{Y}\)-rectangle and all further shortenings.

The canonical physical norm has \(\fitinline{K>0}\), nonnegative radial weight, and \(\fitinline{\mathrm{keep}(z)\Rightarrow z\ne0}\). In the actual application, \[\fitmath{
K=\operatorname{dyadicScale}(n_1),\qquad
\Phi=\operatorname{ballProfile},
}\] with \(\fitinline{\Phi\ge0}\) and support bounded above by \(\fitinline{2}\).

Before removing any \(\fitinline{G}\), \(\fitinline{D}\), seed or coefficient-hole restriction, absorb the independent residual height \(\fitinline{v}\) into the \textbf{original masked row}. The exact identity is \[\fitmath{
\frac{E_{G,D,e}}{\mathcal X}
=
\mathcal E_{\mathrm{keep}_\tau,\Phi,K}
\bigl[\mathcal R^{\,\mathrm{orig}}_{\tau,z}\bigr],
}\tag{35a}\] where \(\fitinline{\mathcal R^{\,\mathrm{orig}}}\) still contains all the original masks and has heights \[\fitmath{
t\longmapsto t+v,\qquad
t_j\longmapsto t_j+\frac{v}{2\pi}.
}\] Before taking its norm, the extracted scalar phase is \[\fitmath{
\left(\frac{T}{Ne_1Ne_2\,NG}\right)^{iv}.
}\] Its modulus is one; both original prefixes and \(\fitinline{G}\) remain in the identity.

At the actual second-family application, the character in (35a) is the original \(\fitinline{\tau_{\rm orig}}\). Before the next step, apply the shared packed-twin bridge (T1)--(T15), proved in section Appendix C.6. Its hypotheses include \(\fitinline{Z_0\mid C_s\mid R}\), the actual first- and second-height identities and the coefficient-sensitive implication \(\fitinline{b(I)\ne0\Rightarrow(I,Z_0)=1}\). No implication from \(\fitinline{(G,I)=1}\) alone to \(\fitinline{(I,Z_0)=1}\) is asserted. The original masked energy and keep equal those of the genuine twin \(\fitinline{\widetilde\tau}\), with coefficients \(\fitinline{a_{\widetilde\tau}(I)=a_{\sigma_A}(I)\chi_I(V)}\) on every supported ideal. In (35b)--(35d), \(\fitinline{\tau}\) denotes this twin and \(\fitinline{w_m=1}\). The temporary comparison mask \(\fitinline{Z_0G}\) is never used as its auxiliary mask.

Now apply the actual packing dictionary. Its hypotheses are \(\fitinline{V\ne0}\), \(\fitinline{w_m\ne0}\), and \[\fitmath{
a_\tau(I)=a_{\sigma_A}(I)\operatorname{idealRowHom}(V;I)
}\] for supported \(\fitinline{I}\) coprime to \(\fitinline{(w_m)}\), together with \[\fitmath{
(G,I)=1\Longrightarrow (I,(w_m))=1.
}\tag{35b}\] The dictionary transports the literal canonical keep: \[\fitmath{
\mathrm{keep}_\tau(z)
\iff\operatorname{initialKeep}_{\sigma_A,Q}(Vz).
}\tag{35c}\] On every original masked residual column, (35b) therefore identifies its natural-character coefficient at \(\fitinline{(\tau,z)}\) with that at \(\fitinline{(\sigma_A,Vz)}\). Unsupported columns vanish on both sides. Consequently \[\fitmath{
\mathcal E_{\mathrm{keep}_\tau,\Phi,K}
[\mathcal R^{\,\mathrm{orig}}_{\tau,z}]
=
\mathcal E_{\substack{V\mid y\\
\operatorname{initialKeep}_{\sigma_A,Q}(y)},\,
\Phi,\,NVK}
[\mathcal R^{\,\mathrm{orig}}_{\sigma_A,y}].
}\tag{35d}\] Multiplication by \(\fitinline{V}\) is a bijection onto \(\fitinline{V\mathcal O}\), and \(\fitinline{N(Vz)/(NVK)=Nz/K}\); there is no Jacobian factor.

The residual coefficient is still supported on \[\fitmath{
(G,I)=1,\qquad D\mid I.
}\tag{29}\] \textbf{Only after (35d)} perform the genuine-\(\fitinline{G}\) allocation, squarefree-\(\fitinline{D}\) allocation, seed identity, kernel expansion, collision completion and hole deletion. These operate on the packed \(\fitinline{\sigma_A}\)-row, with all original masks present at the start. First allocate the complete \(\fitinline{G}\)-part of every tuple coordinate. Because of \(\fitinline{(G,I)=1}\), this is the canonical prime-support decomposition, not an arbitrary allocation of a radical. Write its frozen tuple as \(\fitinline{B_j}\), so \(\fitinline{\prod_jB_j=G}\). The remaining ordinary prefixes are \[\fitmath{
 e_1B_0,\qquad e_2B_1.
}\tag{30}\] At each prime, cube-freeness splits between frozen and live ordinary products because these products are coprime. Frozen marked slots impose their actual frozen restrictions; live marked slots retain their cross-holes. This proves the coefficient identity before taking a norm.

The normalization saving is the full \(\fitinline{NG}\). Indeed, let \[\fitmath{
 \operatorname{rawReduction}(B)
   =NB_0NB_1\prod_{i\ {\rm frozen}}P_i.
}\] A nonzero frozen slot has \(\fitinline{NB_i\leq2P_i}\), while each original marked coordinate is prime. Thus \[\fitmath{
 \operatorname{rawReduction}(B)\geq2^{-|\mathcal I|}NG.
}\] If \(\fitinline{c_B}\) denotes its original frozen coefficient, \[\fitmath{
 \frac{|c_B|^2}{\operatorname{rawReduction}(B)}
 \leq\frac{2^{|\mathcal I|}B_f^{2|\mathcal I|}}{NG}.
}\tag{31}\] Together with (30), the exact full-to-live normalization gives the denominator \(\fitinline{NG\,Ne_1Ne_2}\).

The number of genuine \(\fitinline{G}\)-allocations, squared, is \(\fitinline{O_\delta((NG)^\delta)}\). For \(\fitinline{k=|\mathcal I|+2}\), it is bounded by a fixed power of the ideal divisor function. The bound \(\fitinline{\tau(G)^A\ll_{\delta,A}(NG)^\delta}\) follows prime by prime: large primes absorb \(\fitinline{(v+1)^A}\) into \(\fitinline{p^{\delta v}}\), and the finitely many small primes have bounded remaining suprema.

Next impose \(\fitinline{D\mid I}\) exactly. For each \(\fitinline{\mathfrak p\mid D}\), \[\fitmath{
 1_{\mathfrak p\mid\prod_jI_j}
 =\sum_{\varnothing\ne A\subseteq\{1,\ldots,k\}}
      (-1)^{|A|+1}\prod_{j\in A}1_{\mathfrak p\mid I_j}.
}\tag{32}\] Thus an allocation may select the same prime in several coordinates. Let \(\fitinline{D_j}\) be its selected divisors. Then \[\fitmath{
 D\mid\prod_jD_j,\qquad N\prod_jD_j\geq ND.
}\tag{33}\] All selected copies remain in the extraction denominator.

To apply (32) to the actual restricted coefficient, retain the whole \(\fitinline{D}\)-supported part of each coordinate, including exponent two where present. The shifted cube-free test is evaluated on this frozen part. The residual coordinates are coprime to \(\fitinline{D}\), so their cube-free test splits exactly. This is the bounded-label decomposition; it does not replace a shifted restriction by a fresh unshifted one. There are \(\fitinline{O_\delta((ND)^\delta)}\) squared allocation and bounded-label costs: cube-free ordinary coordinates have only the exponents \(\fitinline{0,1,2}\), and marked prime coordinates only \(\fitinline{0,1}\), giving a fixed local number of choices per prime of \(\fitinline{D}\). Equation (33) supplies \(\fitinline{1/ND}\).

After the packed-row decompositions, the resulting leaves are the full normalized natural products defined below. They have no remaining extraction prefix, seed condition, cube-free restriction, extra ordinary puncture, marked distinctness condition or finite coefficient hole. Their original natural-character zeros remain.

For a positive scale \(\fitinline{U}\), define \[\fitmath{
\mathscr P_{\chi,j}(U)
=U^{-1/2}\sum_Ia_\chi(I)\,
W_j(NI/U)(NI/U)^{\,i(2\pi t_j+v)}.
}\] For each original canonical prime pool define \[\fitmath{
\mathscr Q_{\chi,i}
=P_i^{-1/2}\sum_{\mathfrak p\in\mathcal P_i^{\rm can}}
a_\chi(\mathfrak p)\,
a_{\nu_i}(\mathfrak p)\,
\operatorname{annularWeight}
      (W,P_i,\sigma_i,t+v+f_i;\mathfrak p).
}\] These are the original relative-character and annular coefficients, with the inherited height absorbed. Expressing a height as \(\fitinline{N\mathfrak p^{\,i(t+v)}}\) instead contributes the fixed unit phase \(\fitinline{P_i^{\,i(t+v)}}\); the exact height dictionary retains this phase until the squared norm.

For \(\fitinline{A\subseteq\mathcal I}\), the packed leaf is \[\fitmath{
\mathcal L_{A,U_1,U_2}
=
\sum_{\substack{y\in V\mathcal O\\
\operatorname{initialKeep}_{\sigma_A,Q}(y)}}
\Phi\!\left(\frac{Ny}{NVK}\right)
\left|
\mathscr P_{\chi_{\sigma_A,y},1}(U_1)
\mathscr P_{\chi_{\sigma_A,y},2}(U_2)
\prod_{i\in A}\mathscr Q_{\chi_{\sigma_A,y},i}
\right|^2 .
}\tag{34a}\] Equivalently, use \(\fitinline{y=Vz}\), radial scale \(\fitinline{K}\), and keep \(\fitinline{\operatorname{initialKeep}_{\sigma_A,Q}(Vz)}\). Every leaf uses this same packed radial norm and keep predicate.

The finite generated scale family consists of positive pairs satisfying either \[\fitmath{
U_1\le X_1,\ U_2\le X_2,
\quad\text{or}\quad
U_1\le Y_1,\ U_2\le Y_2.
}\] The restoration argument proves, for every common upper bound \(\fitinline{E\ge0}\) on all these leaves, \[\fitmath{
E_{G,D,e}\le M_{G,D,e}E,
}\] where \[\fitmath{
M_{G,D,e}
=
CZ^\eta\mathcal X
\frac{(NG)^\delta(ND)^\delta}
{NG\,ND\,Ne_1Ne_2}.
}\tag{34b}\] All leaves are nonnegative and their family is finite. Hence, if that family is nonempty, a largest leaf gives \(\fitinline{E_{G,D,e}\le M_{G,D,e}\mathcal L_{A,U_1,U_2}}\). Its slot subset and scale pair may depend on the actual label, both ordinary profiles, all heights and \(\fitinline{v}\). \textbf{Only the mass constant \(\fitinline{C}\) is uniform.} If the family is empty, apply the domination with \(\fitinline{E=0}\).

Thus (34) is a finite positive combination with coefficient mass at most (34b). Subsequent collection must accommodate every generated leaf; no profile-independent maximizing leaf is asserted.

Let \(\fitinline{E_{G,D,e}}\) be the literal unnormalized one-sided residual energy. The preceding argument proves domination by a finite nonnegative combination of canonical leaves \(\fitinline{\mathcal L_\lambda}\): \[\fitmath{
 E_{G,D,e}\leq\sum_\lambda a_\lambda\mathcal L_\lambda,
 \qquad a_\lambda\geq0,
}\] \[\fitmath{
 \boxed{\;
 \sum_\lambda a_\lambda
 \leq
 C Z^\eta\mathcal X
 \frac{(NG)^\delta(ND)^\delta}
      {NG\,ND\,Ne_1Ne_2}.
 \;}
}\tag{34}\]

Finally, retain the stated extraction gain: \[\fitmath{
 gEnergy=\frac{NG\,ND\,Ne_1Ne_2\,Z^k}{\mathcal X^2}
              E_{G,D,e}.
}\] Equation (34) gives total packed-leaf mass \[\fitmath{
 \boxed{\;
 C Z^{k+\eta}\,
 \frac{(NG)^\delta(ND)^\delta}{\mathcal X}.
 \;}
}\tag{36}\] Under \(\fitinline{NG\leq Z^{2B}}\), \(\fitinline{ND\leq Z^B}\), its additional divisor loss is at most \(\fitinline{Z^{3B\delta}}\). This is the complete restoration bridge; subsequent collection of these leaves is separate.

\hypertarget{principal-cancellation-for-0leqalpha23}{%
\subparagraph{\texorpdfstring{6. Principal cancellation for \(\fitinline{0\leq\alpha<2/3}\)}{6. Principal cancellation for 0\textbackslash leq\textbackslash alpha\textless2/3}}\label{principal-cancellation-for-0leqalpha23}}

Fix the actual fixed-inducing ideal \(\fitinline{Q\ne0}\). A row satisfying the literal \texttt{F\allowbreak{}i\allowbreak{}x\allowbreak{}e\allowbreak{}d\allowbreak{}I\allowbreak{}n\allowbreak{}d\allowbreak{}u\allowbreak{}c\allowbreak{}i\allowbreak{}n\allowbreak{}g\allowbreak{}R\allowbreak{}o\allowbreak{}w} predicate admits a character induced from a primitive character whose conductor divides \(\fitinline{Q}\), together with its original natural punctures. This includes the principal primitive character. We do not substitute a condition on a separately simplified raw character.

For every positive rectangle scale \(\fitinline{U}\), use \[\fitmath{
S_{\chi,W,h}(U)
=\sum_Ia_\chi(I)\,
W(NI/U)(NI/U)^{ih}.
}\tag{37a}\] This is the same \texttt{n\allowbreak{}o\allowbreak{}r\allowbreak{}m\allowbreak{}P\allowbreak{}o\allowbreak{}w\allowbreak{}e\allowbreak{}r\allowbreak{}P\allowbreak{}r\allowbreak{}o\allowbreak{}f\allowbreak{}i\allowbreak{}l\allowbreak{}e} \(\fitinline{x\mapsto W(x)x^{ih}}\), evaluated at \textbf{that rectangle's own normalized argument} \(\fitinline{NI/U}\). In the actual height dictionary, \(\fitinline{h_j=2\pi t_j+v}\).

In particular, the first rectangle uses \(\fitinline{NI/X_j}\), the second uses \(\fitinline{NI/Y_j}\), and extracted children use their own shortened scales. The profile is never replaced by \(\fitinline{W(NI/U)(NI)^{ih}}\) without its compensating scale phase.

With this convention, the lattice main term is \[\fitmath{
c_\chi U\int_0^\infty W(x)x^{ih}\,dx.
}\] Hence the product of the two main terms depends on the scales only through \(\fitinline{U_1U_2=T}\), even when \(\fitinline{h_1,h_2}\) are independent. This is the equality used in the centered cancellation.

The fixed-ray lattice formula, followed by the exact finite puncture expansion, gives, uniformly for every \(\fitinline{X>0}\), \[\fitmath{
 S_{\chi,W,t}(X)
 =c_\chi X\,I(W,t)+
 O\!\left(Z^\epsilon
          \mathcal S(W)(1+|t|)^J\right),
}\tag{37}\] where \(\fitinline{W(x)x^{it}}\) is the profile and \(\fitinline{I(W,t)=\int W(x)x^{it}\,dx}\), with the fixed lattice measure included. The coefficient \(\fitinline{c_\chi}\) is independent of \(\fitinline{X}\).

To check uniformity, first apply Poisson to the finitely many primitive ray characters supported by \(\fitinline{Q}\). Integration by parts bounds the nonzero frequencies using finitely many profile derivatives. For \(\fitinline{X<1}\), counting and the main term give the same bounded error. If \(\fitinline{\mathfrak R}\) is the redundant natural radical, insert \[\fitmath{
 S_{\chi,W,t}(X)
 =\sum_{D\mid\mathfrak R}
       \mu(D)\psi(D)\,S_{\psi,W,t}(X/ND).
}\] Its divisor count is subpower under the presented-modulus cap. The resulting main coefficient is \(\fitinline{c_\psi\sum_{D\mid\mathfrak R}\mu(D)\psi(D)/ND}\), independent of \(\fitinline{X}\). Euler derivatives of \(\fitinline{W(x)x^{it}}\) have finite-seminorm bounds polynomial in \(\fitinline{1+|t|}\). This proves (37) with constants preceding both ordinary profiles and both independent heights.

For two rectangles with equal product \(\fitinline{T}\), their main terms in (37) agree. If all four scales are at least \(\fitinline{L\geq1}\), expanding the two products and dividing by \(\fitinline{\sqrt T}\) gives \[\fitmath{
 |\text{normalized difference}|
 \ll Z^\epsilon\mathcal S(W_1)\mathcal S(W_2)
 (1+|t_1|+|t_2|)^J\,\sqrt T\,L^{-1}.
}\] For \(\fitinline{L<1}\), absolute volume gives the same expression with \(\fitinline{L^{-1}}\) replaced by \(\fitinline{1}\). Therefore, for every \(\fitinline{0\leq\alpha\leq1}\), \[\fitmath{
 |\text{plain centered rectangle}|
 \ll Z^\epsilon\mathcal S(W_1)\mathcal S(W_2)
 (1+|t_1|+|t_2|)^J\,\sqrt T\,L^{-\alpha}.
}\tag{38}\]

Insert the cube-free restriction by its \textbf{forward} divisor identity. At primes avoided by the common holes and supplying marks, the local kernel is \(\fitinline{B}\); elsewhere it is \(\fitinline{H_3}\). If \(\fitinline{n=Nd_1Nd_2}\), each shortened scale is at least \(\fitinline{L/n}\), and its product is \(\fitinline{T/n}\). The extraction scalar and child cost combine exactly: \[\fitmath{
 n^{-1/2}\sqrt{T/n}\,(L/n)^{-\alpha}
 =n^{-(1-\alpha)}\sqrt T\,L^{-\alpha}.
}\tag{39}\] Thus the required kernel mass has exponent \[\fitmath{
 \sigma=1-\alpha.
}\] From (2), its ordinary local mass is at most \(\fitinline{1+7p^{-3\sigma}}\). It is uniformly summable precisely in the used range \[\fitmath{
 3\sigma>1,\qquad 0\leq\alpha<2/3.
}\tag{40}\] At an avoided prime, the \(\fitinline{B}\)-mass is \(\fitinline{(1+p^{-\sigma})^2\leq4}\). If these primes divide an ideal \(\fitinline{F}\) with \(\fitinline{NF\leq Z^B}\), their product \(\fitinline{4^{\omega(F)}}\) is \(\fitinline{O_\epsilon(Z^\epsilon)}\), by the fixed-support argument already proved above. Hence the complete mass in (39) is \(\fitinline{O_\epsilon(Z^\epsilon)}\).

For each fixed marked assignment, apply this bound with its augmented common-hole ideal. Its norm is bounded by the old hole norm times \(\fitinline{\prod_i2P_i}\), hence by \(\fitinline{Z^{2B}}\) under the displayed source caps. Summing the original marked coefficients absolutely costs \[\fitmath{
 \sum_{\mathfrak p\in\mathcal P_i}
       \frac{|\beta_i(\mathfrak p)|}{\sqrt{P_i}}
 \leq256M_i\sqrt{P_i},
}\tag{41}\] using ideal counting up to \(\fitinline{2P_i}\); if that range is empty, the sum is zero.

This proves the homogeneous physical principal estimate \[\fitmath{
 |\Delta\mathcal R|
 \leq
 C\,\mathcal S(W_1)\mathcal S(W_2)\,
 Z^\epsilon(1+|t_1|+|t_2|)^J
 \left(\prod_i256M_i\right)
 \sqrt{T\prod_iP_i}\,L^{-\alpha}.
}\tag{42}\] Zero seminorms are handled by homogeneity, before any normalization.

For an original extracted prefix \(\fitinline{e=(e_1,e_2)}\), put \[\fitmath{
 n_e=Ne_1Ne_2,\qquad U_e=\max(Ne_1,Ne_2).
}\] The additional central scalar and shortened volume give \[\fitmath{
 n_e^{-1/2}\sqrt{(T/n_e)\prod_iP_i}
 =\frac{\sqrt{T\prod_iP_i}}{n_e}.
}\] Combining the \(\fitinline{\alpha}\)-estimate with its \(\fitinline{\alpha=0}\) instance proves the literal all-scale form \[\fitmath{
 |\Delta\mathcal R_e|
 \leq
 C\,\mathcal S(W_1)\mathcal S(W_2)\,
 Z^\epsilon(1+|t_1|+|t_2|)^J
 \left(\prod_i256M_i\right)
 \frac{\sqrt{T\prod_iP_i}}{Ne_1Ne_2}
 \left(\max\!\left(1,\frac L{U_e}\right)\right)^{-\alpha}.
}\tag{43}\] Neither prefix is replaced by a genuine common ideal.

When the finite genuine-allocation budgets are combined, the explicit full marked factor retained by the source is \[\fitmath{
 1024^{|\mathcal I|}\left(\prod_iM_i\right)^3.
}\tag{44}\] Constants are not uniform as \(\fitinline{\alpha\uparrow2/3}\); the chosen \(\fitinline{\alpha}\) precedes the eventual threshold and height cutoff.

\hypertarget{exact-principal-cost-used-by-the-root}{%
\subparagraph{7. Exact principal cost used by the root}\label{exact-principal-cost-used-by-the-root}}

Let \(\fitinline{A=P+S}\). The separately supplied common-factor ledger gives \[\fitmath{
 F_1\geq c-3u-o,\qquad F_2\geq2b_2/3,
 \qquad c,b_2\geq0.
}\] For \(\fitinline{0\leq\alpha\leq2/3}\), (43) yields the exceptional excess \[\fitmath{
\begin{aligned}
 A-\frac56-F_1-F_2
   -\alpha\max(0,\ell-c-b_2)
 &\leq A-\frac56+3u-\alpha\ell\\
 &\quad -(1-\alpha)c-(2/3-\alpha)b_2+o\\
 &\leq A-\frac56+3u-\alpha\ell+o.
\end{aligned}
}\tag{45}\] This explains the new lower reference endpoint \[\fitmath{
 \ell_{\rm low}=\frac32(A+3u_{\max}-5/6).
}\tag{46}\]

Assume the actual selection obeys \(\fitinline{u\le u_{\max}}\), and choose the eventual threshold and small losses so that the ledger error satisfies \[\fitmath{
o\le\omega
}\] uniformly over all retained labels, profiles and admitted heights.

The root reserve satisfies \[\fitmath{
0<r_{\rm root}\le1/200.
}\] For \(\fitinline{\epsilon>0}\), set \[\fitmath{
\omega=\min(r_{\rm root}/100000,\epsilon/100000),
\qquad \alpha_P=2/3-\omega/2.
}\] Thus \(\fitinline{\omega>0}\) and \[\fitmath{
0<\alpha_P<2/3,\qquad
\alpha_P\in[2/3-\omega,2/3].
}\] On the centered branch, \[\fitmath{
\ell>\max(0,\ell_{\rm low})+20\omega,\qquad0<\ell<1,
\quad
\ell_{\rm low}=\tfrac32(A+3u_{\max}-5/6).
}\] For every admitted \(\fitinline{\alpha\in[2/3-\omega,2/3]}\), \[\fitmath{
\begin{aligned}
A-\frac56+3u-\alpha\ell+o
&\le A-\frac56+3u_{\max}-\frac23\ell
      +(2/3-\alpha)\ell+o\\
&<-\frac{40}{3}\omega+\omega+\omega\\
&=-\frac{34}{3}\omega\le-9\omega.
\end{aligned}
}\tag{47'}\] Thus the root margin includes both the replacement \(\fitinline{u\mapsto u_{\max}}\) and the uniform error \(\fitinline{o}\). On the uncentered branch, \[\fitmath{
A+3u_{\max}-5/6\le-10\omega
}\] gives the same final bound \(\fitinline{-9\omega}\) after adding \(\fitinline{o}\).

The midpoint-reference adapter remains restricted to \[\fitmath{
0\le R\le2,\qquad \rho\ge3\omega/2,
}\] together with its stated reference geometry. For \(\fitinline{\alpha=\alpha_P}\), \[\fitmath{
(2/3-\alpha_P)R=\omega R/2\le\omega\le2\rho/3.
}\] The midpoint-reference reserve and length bound are \[\fitmath{
P+S-\frac{5R}{6}+\frac{7d}{3}-\frac{2\ell}{3}\le-\rho/3,
\qquad \ell\le R/2.
}\] The displayed width condition therefore proves the required inequality \[\fitmath{
(2/3-\alpha)\ell+P+S-\frac{5R}{6}+\frac{7d}{3}
                 -\frac{2\ell}{3}\le0.
}\tag{48}\] No assertion for unrestricted \(\fitinline{R}\) or smaller \(\fitinline{\rho}\) is made. The fixed factor \(\fitinline{c_\star^{-\alpha_P}}\) remains when the reference length is \(\fitinline{c_\star Z^\ell}\).

These calculations pay the weaker cube-free principal theorem. They do not reuse the fifth-free \(\fitinline{\alpha<4/5}\) range.

\hypertarget{c.8.-joint-divisor-and-physical-packed-row-collection}{%
\subsubsection{C.8. Joint divisor and physical packed-row collection}\label{c.8.-joint-divisor-and-physical-packed-row-collection}}

The joint-divisor reduction and packed-row collection can be written without a new arithmetic assumption beyond the separately assigned cube-free restoration and priced-leaf estimates. Three normalizations are essential:

\[\fitmath{
K_{\rm child}=16K\Bigl(\prod_{j\in J}2P_j\Bigr)^4,\qquad
E_x=NG_x\,NL_x\,N d_{1,x}N d_{2,x}\frac{Z^k}{X^2}e_x,
\qquad
w_x=\frac{\rho(G_x)|p_x|}{NL_x}.
}\]

The second artificial divisor \(\fitinline{L_x}\) need not divide the packed row. Its reciprocal norm must remain in the weight.

This section proves the joint-divisor and packed-row collection. Its leaf and diagonal inputs are supplied in Appendix C.7, \S{}4 and Appendix C.10 and Appendix C.9.

\textbf{1. Data, scope, and inherited inputs.}

Work in the retained Eisenstein ring \(\fitinline{\mathcal O}\), with its fixed primary generators, completed sextic Gauss rows, reciprocity conventions, and finite ray characters. Write \(\fitinline{NI}\) for ideal norm and \(\fitinline{N(z)=N((z))}\). A supported ideal is nonzero and avoids the inherited bad primes.

The physical coefficients remain the original finite tuple coefficients: \[\fitmath{
\beta_u(I)=
 \sum_{\substack{\mathbf a\in\mathcal T\\\prod\mathbf a=I}}
 \mathbf1_{\rm tuple\ restriction}(\mathbf a)\,
 \operatorname{profileCoefficient}_{u}(\mathbf a).
}\tag{1}\] Here the tuple restriction retains the cube-free ordinary product, injective marked primes, ordinary/marked coprimality, punctures, exterior ideal, and all original seed conditions. The profile coefficient retains both ordinary reference pairs \[\fitmath{
X_1X_2=Y_1Y_2>0,
}\] both original prefixes, every marked coefficient, and independent ordinary heights. All four ordinary lengths are positive and at most \(\fitinline{Z^L}\), with fixed \(\fitinline{L}\); they may be smaller than \(\fitinline{1}\).

For the eventual grouped estimate, the marked coefficients are the literal relative-character coefficients with \[\fitmath{
P_i=Z^{w_i},\qquad 0\le w_i\le\mathfrak m,\qquad
\sigma_i=\frac{33}{50}.
}\] Holes have a fixed cardinality cap; frequencies satisfy \(\fitinline{|\nu_i|\le V_f}\). Ordinary profiles are Schwartz functions supported in \(\fitinline{[a,b]\subset[1/4,9/4]}\). The ray modulus, subgroup containing global units, fixed internal reference \(\fitinline{Q}\), slot count, enclosure exponent, and mask exponent are structural data. They precede the constants and eventual threshold.

The reused analytic inputs are the retained local Gauss/CRT table, sextic reciprocity, multiplicative primary-generator convention, fixed-inducing mask invariance, whole-kernel Fourier estimate, and primitive-sieve/finite-profile estimates. Their hypotheses remain part of every application below.

The cube-free restoration-to-leaf estimate is stated explicitly in paragraph 8. Its written proof is in Appendix C.7. The subsequent price assigned to a canonical leaf and the diagonal payment are separate dependencies.

\textbf{2. Keep the first artificial divisor inside the first Cauchy inequality.}

Let \(\fitinline{\mathcal U}\) be a finite family of squarefree ideals. The first physical expression is a sum \[\fitmath{
\sum_{u\in\mathcal U}\mu(u)\,\mathcal M_u,
}\] where \(\fitinline{\mathcal M_u}\) uses the actual residual coefficient restricted by \(\fitinline{u\mid I}\). After the inherited Fourier representation of the whole kernel, the relevant bilinear expression has joint row index \(\fitinline{(u,z)}\): \[\fitmath{
\sum_{u,z}\mu(u)\,\omega_z
          L_{u,v_L}(z)\overline{R_{u,v_R}(\pm z)},
\qquad |\omega_z|\le1.
}\] Cauchy therefore gives \[\fitmath{
\left|\sum_{u,z}\mu(u)\omega_z
 L_{u,v_L}(z)\overline{R_{u,v_R}(\pm z)}\right|
\le
\left(\sum_{u,z}|L_{u,v_L}(z)|^2\right)^{1/2}
\left(\sum_{u,z}|R_{u,v_R}(\pm z)|^2\right)^{1/2}.
}\tag{2}\] There is no factor \(\fitinline{\#\mathcal U}\) and no preliminary reciprocal-divisor completion.

For either side assume, uniformly for every real \(\fitinline{v}\), \[\fitmath{
\sum_{u\in\mathcal U}E_u(v;K)\le E(1+|v|)^{2J}.
}\tag{3}\] The compact logarithmic window has Fourier moments of every order. Fourier inversion, Minkowski, and (3) bound the windowed side by \[\fitmath{
\operatorname{windowBudget}_J(t,E)
=\sqrt E\,\operatorname{heightEnvelope}(t)^J
  \operatorname{profileMoment}_J.
}\] The right conjugate window changes its Fourier parameter to its negative. Reflection \(\fitinline{z\mapsto-z}\) is a bijection and preserves the radial weight. Thus the two sides retain their independent heights and radii.

This is the mathematical content of \texttt{F\allowbreak{}i\allowbreak{}r\allowbreak{}s\allowbreak{}t\allowbreak{}/\allowbreak{}J\allowbreak{}o\allowbreak{}i\allowbreak{}n\allowbreak{}t\allowbreak{}K\allowbreak{}e\allowbreak{}r\allowbreak{}n\allowbreak{}e\allowbreak{}l\allowbreak{}.\allowbreak{}s\allowbreak{}u\allowbreak{}m\allowbreak{}m\allowbreak{}e\allowbreak{}d\allowbreak{}\_\allowbreak{}k\allowbreak{}e\allowbreak{}r\allowbreak{}n\allowbreak{}e\allowbreak{}l\allowbreak{}\_\allowbreak{}f\allowbreak{}r\allowbreak{}o\allowbreak{}m\allowbreak{}\_\allowbreak{}r\allowbreak{}a\allowbreak{}w}. The finite ray projector is applied after (2), with its fixed coefficient mass.

\textbf{3. Exact regrouping inside the second physical block.}

Let \(\fitinline{C_2,D_2}\) be supported ideals with equal prime support. Its sector columns satisfy \[\fitmath{
(I,C_2)=(J,D_2)=(I,J)=1.
}\] For squarefree \(\fitinline{u}\), \[\fitmath{
u\mid C_2I,\quad u\mid D_2J
\quad\Longleftrightarrow\quad u\mid C_2.
}\tag{4}\] Indeed, a prime dividing \(\fitinline{u}\) but absent from the common support would have to divide both \(\fitinline{I}\) and \(\fitinline{J}\). Conversely, squarefreeness and equality of supports imply that every squarefree divisor of \(\fitinline{C_2}\) also divides \(\fitinline{D_2}\).

Suppose \(\fitinline{\mathcal U}\) contains all squarefree divisors of \(\fitinline{C_2}\), and contains only squarefree ideals. Then \[\fitmath{
\sum_{u\in\mathcal U}
 \mathbf1_{u\mid C_2I}\mathbf1_{u\mid D_2J}
=\rho(C_2),\qquad
\rho(G):=2^{\omega(G)}.
}\tag{5}\]

Put \[\fitmath{
A_2=\operatorname{commonFrequencyGenerator}(C_2,D_2)
       \operatorname{nonunitFrequencyGenerator}(C_2,D_2,U).
}\] The actual physical block is the finite sum of \[\fitmath{
\begin{aligned}
&\mathbf1_{\rm canonical\ partition}(z)\,
 \mathbf1_{(I,J)=1}\,
 \operatorname{Corr}(C_2I,D_2J;A_2z)\\
&\quad{}\times
 \beta(C_2I)h_\eta(t,I)
 \overline{\beta(D_2J)h_\eta(t,J)}
 \,\mathcal K(C_2,D_2;A_2z,I,J)
 \prod_{\ell=0}^{3}\mathcal D_{n_\ell}.
\end{aligned}
}\tag{6}\] The four dyadic weights use the actual effective scale, \(\fitinline{N(z)}\), \(\fitinline{NI}\), and \(\fitinline{NJ}\). Every factor except the two seed indicators is independent of \(\fitinline{u}\). Substituting (5) termwise proves \[\fitmath{
\boxed{\quad
\sum_{u\in\mathcal U}\operatorname{Block}[\mathbf1_{u\mid\cdot}\beta]
=\rho(C_2)\operatorname{Block}[\beta].
\quad}
}\tag{7}\] No phase, normalization, angular sector, or coefficient is replaced.

Because (7) holds for every finite row set, it commutes with any row filter. In particular it holds separately for fixed-inducing and nonfixed rows. This identity occurs after the first Hilbert-space Cauchy step; \(\fitinline{\rho(C_2)}\) is a positive multiplicity, not a surviving signed M\"{o}bius sum.

At a common prime, write the two multiplicities as \(\fitinline{l,r}\). The retained local correlation table states that an unequal pair can be nonzero only if \(\fitinline{6\mid\min(l,r)}\). The literal cube-free restriction gives \(\fitinline{l,r\in\{1,2\}}\). Therefore \(\fitinline{l\ne r}\) kills the local correlation at every frequency. CRT then kills the whole summand.

Consequently a nonzero second common contribution has \(\fitinline{C_2=D_2=G}\). This conclusion does \textbf{not} identify the two first common factors.

For equal second factors, the common frequency generator spans \(\fitinline{G}\), and the retained normalized local table gives \[\fitmath{
|\operatorname{Corr}(G,G;g\,w)|
\le \frac{NG}{N U_{\rm id}}
\quad\text{on the canonical \(\fitinline{U}\)-partition},
}\tag{8}\] where \(\fitinline{(g)=G}\) and \(\fitinline{U_{\rm id}=\prod_{P\in U}P}\). The complementary primes supply the nonunit multiplier in \(\fitinline{A_2}\). Both local unit and nonunit cases remain in the exhaustive partition.

\textbf{4. Genuine allocation, residual seeds, and fourth-power averaging.}

First extract the genuine first common ideal into the tuple components. Each tuple has a unique prime-by-prime allocation \(\fitinline{\mathcal B}\). Thus, on the actual coprime residual sector, \[\fitmath{
\beta(CI)=\sum_{\mathcal B}
 b_{\mathcal B}\,\beta_{\mathcal B}(I),
}\tag{9}\] with the frozen tuple restrictions included in \(\fitinline{b_{\mathcal B}}\). The live marked prefixes are now \(\fitinline{1}\). The artificial condition \(\fitinline{u\mid I}\) remains on the whole live product.

Cauchy over genuine allocations costs exactly \[\fitmath{
\#\mathcal B\sum_{\mathcal B}|b_{\mathcal B}|^2(\cdots).
}\tag{10}\] The elementary bound \(\fitinline{\#\mathcal B\le\tau(C)^{N+2}}\) is paid only later, using the fixed slot count and polynomial norm cap.

Fix one live allocation and selected slot set \(\fitinline{J}\). For an assignment \(\fitinline{y}\), put \[\fitmath{
\Delta_y=\prod_{j\in J}y_j,\qquad d_y=\operatorname{primary}(\Delta_y).
}\] The exact selected scalar is \[\fitmath{
s_y=
\mathbf1_{y\ {\rm injective}}\,
\operatorname{selAmp}(y)\,
\mathbf1_{(\Delta_y,R)=1}\,
h_\tau(t,\Delta_y).
}\tag{11}\] For squarefree \(\fitinline{u}\), let \(\fitinline{v(u,\Delta_y)}\) be the product of primes of \(\fitinline{u}\) absent from \(\fitinline{\Delta_y}\). Primewise, \[\fitmath{
u\mid\Delta_y I\iff v(u,\Delta_y)\mid I.
}\tag{12}\] The map \(\fitinline{u\mapsto v(u,\Delta_y)}\) has fibers of cardinality at most \(\fitinline{2^{\omega(\Delta_y)}\le2^{|J|}}\).

The retained fourth-power Gauss identity, applied to the \textbf{whole} residual vector, gives \[\fitmath{
Q_u(k)=
\sum_{y,\chi,\xi}
 s_y\,c_{\chi,\xi}\,
 \chi(d_y)\,\mathcal G_{\Delta_y}(k)\,
 Q^{\rm rest}_{v(u,\Delta_y),\,\tau\xi}(d_y^4k).
}\tag{13}\] Here \(\fitinline{c_{\chi,\xi}}\) is the fixed first reciprocity-phase table coefficient. The residual vector depends on \(\fitinline{y}\) only through its residual seed. The deleted coprimality restriction in the transferred expression is justified by the inherited Gauss zeros, rather than by arbitrary coefficient enlargement.

The scalar satisfies \[\fitmath{
|s_yc_{\chi,\xi}\chi(d_y)\mathcal G_{\Delta_y}(k)|^2
\le |c_{\chi,\xi}|^2|s_y|^2.
}\] Set \[\fitmath{
A_{\rm sel}=\sum_{y,\chi,\xi}|c_{\chi,\xi}|^2|s_y|^2.
}\tag{14}\] This is the exact Cauchy mass. If slot amplitudes have bounds \(\fitinline{M_j}\), then \[\fitmath{
A_{\rm sel}
\le c_{\rm phase}\Bigl(\prod_{j\in J}M_j\Bigr)^2
             \prod_{j\in J}256P_j.
}\tag{15}\] It has one selected-width factor.

Define the common width \[\fitmath{
W_J=\prod_{j\in J}2P_j.
}\] Then \(\fitinline{N\Delta_y\le W_J}\), and on live assignments \[\fitmath{
W_J\le2^{|J|}N\Delta_y.
}\] Since \(\fitinline{N(d_y^4)=N(\Delta_y)^4}\), the original ball \(\fitinline{N(k)<2K}\) maps inside the ball \[\fitmath{
\boxed{K_{\rm child}=16K W_J^4.}
}\tag{16}\] The majorizing ball profile equals \(\fitinline{1}\) on the transferred support. This is a fourth-power radius; replacing \(\fitinline{W_J}\) by \(\fitinline{\prod P_j}\) would omit a fixed but real factor.

For nonzero \(\fitinline{z}\), the number of pairs \(\fitinline{(y,k)}\) with \(\fitinline{d_y^4k=z}\) is at most \[\fitmath{
\#\{y:d_y^4\mid z\}\le\tau((z))^{|J|}.
}\tag{17}\] For a given \(\fitinline{y}\), multiplication by \(\fitinline{d_y^4\ne0}\) is injective. Each selected prime ideal divides \(\fitinline{(z)}\), proving the last bound.

If \(\fitinline{J\ne\varnothing}\), the zero row is killed \textbf{before} replacing coefficients by their norm bounds: an injective nonempty product of primes is not a sixth power, whereas its Gauss row at zero can be nonzero only for a sixth-power ideal. If \(\fitinline{J=\varnothing}\), there is exactly one assignment, its multiplier is \(\fitinline{1}\), and the zero row remains.

Combining (12)--(17), for any finite squarefree set \(\fitinline{\mathcal V}\) containing the residual seeds, \[\fitmath{
\boxed{
\sum_{u\in\mathcal U}E_u(K)
\le A_{\rm sel}\,2^{|J|}\,M_J
 \sum_{\chi,\xi}\sum_{v\in\mathcal V}
 E^{\rm rest}_{v,\tau\xi}(K_{\rm child}),
}
}\tag{18}\] where \(\fitinline{M_J}\) bounds (17) on the child ball, with the preceding zero-row convention. There is no remaining sum over selected assignments on the right.

Now, and only now, positivity permits completion of \(\fitinline{\mathcal V}\) to include all squarefree divisors needed by (7). This is the corrected selected-row average in \texttt{C\allowbreak{}o\allowbreak{}n\allowbreak{}c\allowbreak{}r\allowbreak{}e\allowbreak{}t\allowbreak{}e\allowbreak{}T\allowbreak{}o\allowbreak{}t\allowbreak{}a\allowbreak{}l\allowbreak{}M\allowbreak{}u\allowbreak{}l\allowbreak{}t} and \texttt{C\allowbreak{}o\allowbreak{}n\allowbreak{}c\allowbreak{}r\allowbreak{}e\allowbreak{}t\allowbreak{}e\allowbreak{}A\allowbreak{}s\allowbreak{}s\allowbreak{}e\allowbreak{}m\allowbreak{}b\allowbreak{}l\allowbreak{}y\allowbreak{}A\allowbreak{}v\allowbreak{}g}.

\textbf{5. All four Poisson contributions remain.}

Apply the retained second Poisson identity to each residual energy in (18). Define:

\begin{itemize}
\tightlist
\item
  \(\fitinline{Z_0}\): the absolute value of its actual zero-frequency term;
\item
  \(\fitinline{Z_f}\): the absolute fixed-inducing retained contribution;
\item
  \(\fitinline{Z_t}\): the absolute discarded tail;
\item
  \(\fitinline{Z_{\rm nf}}\): the sum of absolute nonfixed physical blocks after (7).
\end{itemize}

Then \[\fitmath{
\sum_{v\in\mathcal V}E^{\rm rest}_v
\le Z_0+Z_f+Z_t+Z_{\rm nf}.
}\tag{19}\] In \(\fitinline{Z_{\rm nf}}\), every block has the original coefficient with seed \(\fitinline{1}\), multiplied by \(\fitinline{\rho(C_2)}\). The seed has been removed by (7), not by estimating individual seed rows.

Thus the complete side reduction is \[\fitmath{
\sum_u E_u^{\rm original}
\le
\#\mathcal B\sum_{\mathcal B}|b_{\mathcal B}|^2
 A_{{\rm sel},\mathcal B}2^{|J_{\mathcal B}|}M_{\mathcal B}
 \sum_{\chi,\xi}(Z_0+Z_f+Z_t+Z_{\rm nf})_{\mathcal B,\xi}.
}\tag{20}\] Unit rows have norm \(\fitinline{1}\) and remain in their actual fixed/nonfixed class. Empty common supports have \(\fitinline{\rho(1)=1}\). Empty selected sets and zero-slot sources are covered by the stated conventions.

The separate diagonal supplier must pay \(\fitinline{Z_0+Z_f+Z_t}\). Nothing in the packed nonfixed argument pays these terms automatically.

\textbf{6. Exact first-side packing and \(\fitinline{q_C}\).}

Let the first common parts have multiplicities \(\fitinline{l_P,r_P\in\{1,2\}}\). Put \[\fitmath{
n_P=(l_P-r_P)\bmod6.
}\] At active primes \(\fitinline{n_P\ne0}\), the only possibilities are \[\fitmath{
\begin{array}{c|c|c|c}
(l_P,r_P)&n_P&(1+n_P)\bmod6&(1+5n_P)\bmod6\\ \hline
(2,1)&1&2&0\\
(1,2)&5&0&2 .
\end{array}
}\tag{21}\] Let \(\fitinline{\pi_P}\) be the retained primary generator, let \(\fitinline{r=\epsilon\prod_{\rm active}\pi_P}\) be the actual active conductor, and let \(\fitinline{e}\) be the first common subset generator. Define \[\fitmath{
V_C=e\epsilon\prod_{\rm active}\pi_P^{(1+n_P)\bmod6},
\qquad
V_D=e\epsilon\prod_{\rm active}\pi_P^{(1+5n_P)\bmod6}.
}\tag{22}\] The unit \(\fitinline{\epsilon}\) is retained in the character identity.

The table gives, at active primes, \[\fitmath{
k_{C,P}+r_P=1+l_P,\qquad
k_{D,P}+l_P=1+r_P,
}\] and inactive primes have \(\fitinline{l_P=r_P}\). Multiplication therefore proves \[\fitmath{
NV_C\,N c_D=Ne\,Nr\,N c_C,\qquad
NV_D\,N c_C=Ne\,Nr\,N c_D.
}\tag{23}\] For \(\fitinline{Z>1}\), putting \[\fitmath{
c=\log_ZNc_C,\quad d=\log_ZNc_D,\quad
q=\log_Z(Ne\,Nr)
}\] gives exactly \[\fitmath{
\boxed{q_C=\log_ZNV_C=q+c-d,\qquad q_D=q+d-c.}
}\tag{24}\]

The \(\fitinline{e}\)-part divides either common part. The remaining packed exponent is at most the corresponding common exponent. Hence \[\fitmath{
\boxed{(V_C)\mid c_C^2,\qquad(V_D)\mid c_D^2.}
}\tag{25}\] This proves polynomial support and divisor control for the actual packed multiplier.

For a primary column \(\fitinline{a}\) coprime to \(\fitinline{P}\), sextic reciprocity gives \[\fitmath{
(a/P)_6^{n_P}
=\operatorname{phase}(\pi_P,a)^{n_P}\chi_{(a)}(\pi_P)^{n_P}.
}\] Multiplying the existing \(\fitinline{\pi_P}\)-row symbol and reducing its exponent modulo six yields (22). On the right, conjugation replaces the exponent \(\fitinline{n_P}\) by \(\fitinline{5n_P}\); the reciprocity phase is real and remains unchanged.

If the reduced exponent is zero, the prime must remain in the puncture mask: \(\fitinline{0^6=0}\), while the reduced zeroth power equals \(\fitinline{1}\). Let \(\fitinline{Z_C,Z_D}\) be precisely these exponent-zero prime products. On supported columns avoiding them, \[\fitmath{
h_{\tau_C}(t,I)=h_{\sigma_C}(t,I)\chi_I(V_C),\qquad
h_{\tau_D}(t,I)=h_{\sigma_D}(t,I)\chi_I(V_D)
}\tag{26}\] for every real \(\fitinline{t}\). The anchors \(\fitinline{\sigma_C,\sigma_D}\) belong to the fixed finite ray dictionary obtained by multiplying the original ray twist by the phase of \(\fitinline{\prod\pi_P^{n_P}}\). They do not grow with the common ideals.

After the second ray expansion, the multiplier is \(\fitinline{V_CA_2}\), respectively \(\fitinline{V_DA_2}\), and the mask includes \(\fitinline{Z_CC_2}\), respectively \(\fitinline{Z_DD_2}\).

The independent-height absorption is also exact. If the inherited ordinary prefix is \(\fitinline{\mathbf d=(d_1,d_2)}\), then multiplication by \(\fitinline{NI^{iv}}\) changes the slot heights by \(\fitinline{v}\), the two profile parameters by \(\fitinline{v/(2\pi)}\), and leaves the scalar \[\fitmath{
\left(\frac{X_1X_2}{Nd_1Nd_2\,NG}\right)^{iv},
}\tag{27}\] of modulus \(\fitinline{1}\). Neither prefix norm nor either independent height may be omitted.

\textbf{7. Canonical keep transfers to the packed row.}

Suppose, on supported columns avoiding a finite mask \(\fitinline{w}\), \[\fitmath{
\operatorname{coeff}_{\tau}(I)
=\operatorname{coeff}_{\sigma}(I)\chi_I(V),\qquad V\ne0.
}\tag{28}\] Enlarge the fixed bad mask by \(\fitinline{w}\), using the retained fixed-inducing mask-invariance lemma. On unmasked elements, multiplicativity gives \[\fitmath{
\operatorname{coeff}_{\tau}(I)\chi_I(z)
=\operatorname{coeff}_{\sigma}(I)\chi_I(Vz).
}\] Both sides vanish on masked or unsupported elements. Their masked row characters therefore agree as functions. Consequently, for \(\fitinline{z\ne0}\), \[\fitmath{
\operatorname{FixedInducingRow}_{\tau,Q}(z)
\iff
\operatorname{FixedInducingRow}_{\sigma,Q}(Vz).
}\tag{29}\] The internal reference \(\fitinline{Q}\) is unchanged.

At \(\fitinline{z=0}\), both nonfixed keep predicates are false. Thus for every \(\fitinline{z}\), \[\fitmath{
\boxed{\operatorname{canonicalKeep}_{\tau,Q}(z)
\iff\operatorname{initialKeep}_{\sigma,Q}(Vz).}
}\tag{30}\] Natural characters agree on the same actual masked columns. This preserves their natural zeros.

For the actual second-family application, first use the shared packed-twin bridge (T1)--(T15) in section Appendix C.6. The first exponent-zero ideal \(\fitinline{Z_0}\) divides the first common factor and the retained residual puncture \(\fitinline{R}\); it need not divide the second common \(\fitinline{G}\). Original nonzero residual coefficients avoid \(\fitinline{Z_0}\). Hence the original masked energy and keep equal those of the genuine twin, whose supported coefficients are globally \(\fitinline{a_\sigma(I)\chi_I(V)}\) and whose auxiliary mask is \(\fitinline{1}\). This is the frame character used below.

The injective substitution \(\fitinline{y=Vz}\) transforms this ORIGINAL MASKED polynomial at scale \(\fitinline{K}\) into the exact sublattice energy \[\fitmath{
\sum_{\substack{V\mid y\\\operatorname{initialKeep}_{\sigma,Q}(y)}}
 |\mathcal R^{\rm orig}_{\sigma,y}|^2\,
 \Phi\!\left(\frac{Ny}{NV\,K}\right).
}\tag{31}\] Here \(\fitinline{\mathcal R^{\rm orig}_{\sigma,y}}\) retains every original common-factor, exponent-zero, tuple and puncture restriction. Only the subsequent restoration produces unrestricted canonical leaves. There is no Jacobian factor in this discrete sum.

\textbf{8. Restoration contract and extraction normalization.}

The separately required cube-free restoration statement is the following uniform contract. For a genuine common ideal \(\fitinline{G}\), second artificial squarefree divisor \(\fitinline{L}\), and inherited ordinary prefix \(\fitinline{\mathbf d}\), the literal one-sided energy \(\fitinline{e_x(v)}\) is dominated by a finite positive combination of canonical packed leaves, over every surviving slot subset and shortened ordinary pairs below either original pair, of total mass \[\fitmath{
\mathcal M_x=
C_{\rm red}Z^\eta X
 \frac{(NG)^\delta(NL)^\delta}
      {NG\,NL\,Nd_1Nd_2}.
}\tag{32}\] Here \(\fitinline{X=X_1X_2\prod_iP_i}\). The seed divides \(\fitinline{G}\); all prefixes are nonzero; the radical of the exterior ideal, \(\fitinline{G,L}\), and punctures has norm at most \(\fitinline{Z^B}\); the prime enclosure is fixed at \(\fitinline{\mathrm{Cap}\,Z^{L_{\rm cap}}}\). Heights, holes, original coefficients, and all shortenings are uniform.

The frame character is the packed twin, with auxiliary mask \(\fitinline{1}\), whose raw energy equals the original second-character energy. Apply the the restoration of Appendix C.7 only after packing the original masked polynomial. Its exact hypotheses are: nonnegative canonical radial weight and positive radial scale; keep excluding zero; nonzero prefixes and exterior ideal; seed dividing \(\fitinline{G}\); squarefree \(\fitinline{L}\); the stated radical cap; \(\fitinline{0\le w_i\le L_{\rm slot}\le L_{\rm cap}}\) and \(\fitinline{\mathrm{Cap}\ge\max(2,b)}\); complete ordinary-prime enclosure; and every ordinary-hole prime belonging to each original live slot-hole set. No hidden mask is assigned to a restored leaf. A largest positive leaf may be selected separately for each label, profiles and height \(\fitinline{v}\); only its mass constant is uniform. Empty or zero-energy families contribute zero. Appendix C.7 supplies this exact contract; a bound for unrelated unrestricted coefficients would not suffice.

Use the normalized energy \[\fitmath{
E_x(v)=NG\,NL\,Nd_1Nd_2\,Z^{k-2a}e_x(v),
\qquad X=Z^a.
}\tag{33}\] With \(\fitinline{NG\le Z^{2B}}\), \(\fitinline{NL\le Z^B}\), and \(\fitinline{\delta=\varepsilon/(48\max(1,B))}\), \[\fitmath{
NG\,NL\,Nd_1Nd_2\,\mathcal M_x
\le C_{\rm red}Z^{\eta+\varepsilon/16}X.
}\tag{34}\] The three extraction gains cancel before collection. Moving them into the label weight would destroy the harmonic fiber bound.

\textbf{9. Divisor fibers before completion.}

Fix the ray and dyadic class. Restrict to the actual finite frame image, with source, first-side, ray and dyadic data frozen. On this image a frame label is determined by \[\fitmath{
\bigl(G,(A_2),L\bigr).
}\tag{35}\] Indeed \(\fitinline{A_2}\) is exactly the multiplicative primary generator of \[\fitmath{
\prod_P P^{\min(l_P,r_P)}
 \prod_{P\in\text{nonunit partition}}P.
}\] Thus its ideal determines the element, including its normalization. The raised dyad is determined by \(\fitinline{A_2}\); the anchor is fixed by the ray class; the twin, mask \(\fitinline{1}\), exterior ideal, seed, holes, prefixes and prefactor are then prescribed frame data. This is image-restricted code injectivity, not an assertion about arbitrary frame-label records.

Even before using unequal-multiplicity vanishing, cube-freeness gives \[\fitmath{
G\mid(A_2)^2.
}\tag{36}\] At a row \(\fitinline{y\ne0}\) with \(\fitinline{A_2\mid y}\), both code ideals in (35) divide \(\fitinline{(y)^2}\). Also \(\fitinline{\rho(G)\le\tau((y)^2)}\). For fixed \(\fitinline{L}\), there are at most \(\fitinline{\tau((y)^2)^2}\) codes. Therefore, if \(\fitinline{|p_x|\le1}\), \[\fitmath{
\sum_{\substack{x\\A_{2,x}\mid y}}
 \frac{\rho(G_x)|p_x|}{NL_x}
\le
\tau((y)^2)^3
 \sum_{0<NL\le Z^B}\frac1{NL}.
}\tag{37}\] The retained ideal counting bound gives \[\fitmath{
\sum_{0<NL\le Z^B}\frac1{NL}
\le256\bigl(\lceil\log_2 Z^B\rceil+1\bigr).
}\tag{38}\] This is the required completion of the second-divisor sum. It uses reciprocal weights, not divisibility of \(\fitinline{L}\) into \(\fitinline{y}\).

For \(\fitinline{Ny\le Z^{M_{\max}+2}}\), the fixed-field divisor bound and (38) make (37) at most \(\fitinline{Z^{\varepsilon/8}}\), eventually uniformly in \(\fitinline{y}\). One can allocate \(\fitinline{\varepsilon/32}\) to the divisor power, \(\fitinline{\varepsilon/32}\) to the logarithm, and \(\fitinline{\varepsilon/16}\) to their fixed constants.

There is a separate multiplicity when original physical labels map to the same frame label. For fixed chosen common factor \(\fitinline{G}\), both original common ideals divide \(\fitinline{G^2}\), and their unit/nonunit partition has at most \(\fitinline{\rho(G)}\) choices. Hence \[\fitmath{
M_{\rm physical}\le\tau(G^2)^2\rho(G)\le\tau(G^2)^3.
}\tag{39}\] It is paid separately under the common-ideal norm cap. It is not an extra free label count.

\textbf{10. Common profiles and total grouped energy.}

For an ordinary profile \(\fitinline{W}\), let \(\fitinline{f(v)=W(e^v)}\). It is smooth and compactly supported, so \[\fitmath{
W(e^v x)=\int_{\mathbb R}\widehat f(\tau)
 e^{2\pi i\tau v}x^{2\pi i\tau}\,d\tau.
}\tag{40}\] Choose a fixed smooth two-piece partition on \(\fitinline{[1/8,9/4]}\), with scales \(\fitinline{c_0=1/4,c_1=3/2}\). The common profiles \[\fitmath{
W_{e,\tau}(u)=\psi_e(c_eu)u^{2\pi i\tau}
}\] are supported in \(\fitinline{[1/4,9/4]}\), and every required seminorm grows polynomially in \(\fitinline{|\tau|}\).

For two ordinary factors the exact scalar includes \(\fitinline{\sqrt{c_{e_1}c_{e_2}}}\), as well as the dilation phases. Its positive majorant is \[\fitmath{
M_e(\tau_1,\tau_2)
=\sqrt{c_{e_1}c_{e_2}}\,
 |\widehat f_1(\tau_1)||\widehat f_2(\tau_2)|.
}\tag{41}\] Repeated integration by parts bounds every polynomially weighted integral of (41) by finitely many seminorms of the original profiles.

Discard zero-energy children before grouping. A positive child with ordinary support in \(\fitinline{[a,b]}\) has \(\fitinline{X_i'\ge1/b}\), because a supported nonzero ideal has norm at least one. Together with \(\fitinline{X_i'\le Z^L}\), this gives the finite dyadic range counted below. Group shortened ordinary lengths into dyadic cells \(\fitinline{X_i'\in[Y_i/2,Y_i]}\). The dilation is \(\fitinline{v_i=\log(Y_i/X_i\prime)\in[0,\log2]}\). Their central-normalization ratio is \[\fitmath{
\sqrt{\frac{Y_1Y_2}{X_1'X_2'}}\le2.
}\] For a fixed anchor, surviving slot subset, and dyadic length cell, write the common-profile row as \(\fitinline{P_{e,\tau}(y)}\). Cauchy in the Fourier integral, followed by the row fiber bound \(\fitinline{D}\), yields \[\fitmath{
\sum_xw_x E_x^{\rm leaf}
\le4D\,L_0L_J\,T,
}\tag{42}\] whenever the common-profile energy is at most \(\fitinline{T\operatorname{pairWeight}_J(\tau)}\). Here \(\fitinline{L_j=\sum_e\int M_e\operatorname{pairWeight}_j}\). To verify (42), first use \[\fitmath{
\left|\int a(\tau)P(\tau)\,d\tau\right|^2
\le\left(\int|a|\right)\int|a||P|^2,
}\] then interchange nonnegative sums/integrals and apply (37) at each common row. This never sums a separate source bound over all labels.

Assume a canonical leaf at row-scale exponent \(\fitinline{H}\) has the uniform price \(\fitinline{Z^{H+p(w_0,P,s,H)+\varepsilon/8}}\), where \(\fitinline{p}\) is nonnegative, monotone in \(\fitinline{P,s}\), and Lipschitz with constant \(\fitinline{K_p}\). Put \[\fitmath{
P=\max\{\ell(X_1)+\ell(X_2),\ell(Y_1)+\ell(Y_2)\},
\quad \ell(x)=\max(0,\log_Zx),\quad s=\sum_iw_i.
}\] For a group with packed envelope within \(\fitinline{\mathfrak m}\) of \(\fitinline{H}\), the exact dual cutoff gives \[\fitmath{
H\le q_0+2a-k+2\mathfrak m.
}\] Consequently (32)--(42) give \[\fitmath{
\boxed{
\sum_xw_xE_x(v)
\le C\,\mathcal A_J\,(1+|v|)^{2J}
 Z^{a+q_0+p(w_0,P,s,H)+\varepsilon},
}
}\tag{43}\] where \[\fitmath{
\mathcal A_J=
\bigl(\|W_1\|_{\mathcal S}\|W_2\|_{\mathcal S}\bigr)^2
(1+|t|+V_f+|t_1|+|t_2|)^J.
}\] The constants and finite seminorm set precede the variable profiles, heights, lengths, labels, and \(\fitinline{Z}\).

A sufficient explicit fee allocation is:

\begin{longtable}[]{@{}lr@{}}
\toprule\noalign{}
Cost & Exponent \\
\midrule\noalign{}
\endhead
\bottomrule\noalign{}
\endlastfoot
Restoration & \(\fitinline{\varepsilon/8}\) \\
Extraction mass & \(\fitinline{\varepsilon/16}\) \\
Weighted row fibers & \(\fitinline{\varepsilon/8}\) \\
Finite anchors/subsets/dyadic length cells & \(\fitinline{\varepsilon/16}\) \\
Canonical leaf bound & \(\fitinline{\varepsilon/8}\) \\
Mesh & \(\fitinline{\varepsilon/8}\) \\
Fixed scale enlargements and price variation & \(\fitinline{\varepsilon/16}\) \\
\end{longtable}

Their sum is \(\fitinline{11\varepsilon/16}\). It suffices to choose \[\fitmath{
\mathfrak m\le\tfrac12,\qquad
(3+K_p)\mathfrak m\le\varepsilon/8,
\qquad
(1+3K_p)\log_Z3\le\varepsilon/16.
}\] The mesh choice precedes the growing family. For fixed slot count, the class count is bounded by \[\fitmath{
|\mathrm{Dict}|\,2^N
 \bigl(L\log_2Z+\log_2b+2\bigr)^2.
}\] This also shows why fixed structural size and eventuality cannot be interchanged silently.

\textbf{11. Return to the physical blocks and the six-width identity.}

In this paragraph, \(\fitinline{K}\) denotes the SECOND Poisson/Gauss scale \(\fitinline{K_{\rm child}=K_{\mathcal B}}\), with \(\fitinline{Z^k=K}\). It is not the pre-selection scale denoted by \(\fitinline{K}\) in (16).

The literal physical scalar is \[\fitmath{
\operatorname{outerScalar}
=\frac{K}{\sqrt{NC_2}\sqrt{ND_2}
               \sqrt{\operatorname{dyad}_2}\sqrt{\operatorname{dyad}_3}}.
}\] Let \(\fitinline{p_d=Nd_1Nd_2}\). The actual nonzero-column support gives \(\fitinline{c_{\rm low}X\le p_dN(C_2I)}\), and similarly on the right. Consequently the coefficient \(\fitinline{c_0}\) in the following display is \(\fitinline{p_d/c_{\rm low}}\), not a prefix-independent constant. If actual column support gives \(\fitinline{X\le c_0NC_2\operatorname{dyad}_2}\) and \(\fitinline{X\le c_0ND_2\operatorname{dyad}_3}\), then \[\fitmath{
\operatorname{outerScalar}\le c_0X\frac K{X^2}.
}\tag{44}\] The common normalizer supplies the factor bounded by either \(\fitinline{NC_2}\) or \(\fitinline{ND_2}\). In the surviving equal-second-common sector this is \(\fitinline{NG}\). Combining \(\fitinline{\operatorname{outerScalar}\le p_dK/(c_{\rm low}X)}\), that common normalizer, and the factor \(\fitinline{NG\,NL\,p_dZ^k/X^2}\) in (33) cancels the complete prefix and common factors. The resulting scalar is precisely \(\fitinline{X/c_{\rm low}}\) after the retained frequency normalization, as displayed in (45). Explicitly, with \(\fitinline{p_d=Nd_1Nd_2}\), \[\fitmath{
w_xE_x=\rho(G_x)|p_x|\,NG_x\,p_d\frac K{X^2}\,e_x;
}\] the second divisor \(\fitinline{NL_x}\) cancels against its reciprocal label weight. Hence \[\fitmath{
\rho(G_x)|p_x|\,\operatorname{outerScalar}\,
 \operatorname{normalizer}\,e_x
\le\frac X{c_{\rm low}}w_xE_x.
}\] No prefix is replaced by one.

To sum absolute block values, multiply each block by its own unit complex phase and apply the whole-kernel estimate to the joint index \(\fitinline{(\text{label},z)}\). This gives one Cauchy inequality between the two summed side energies. Both orientations use (43); no number-of-labels factor appears.

For a dyadic box \(\fitinline{[-M,M]^4}\), the resulting bound has the explicit form \[\fitmath{
Z_{\rm nf}\le
(2M+1)^4 C_{\rm ker}\frac{X}{c_{\rm low}}
 M_{\rm physical}B_0\,
 \operatorname{profileMoment}_J^2
 (1+2\pi)^{4J}(1+|t|)^{2J},
}\tag{45}\] when \(\fitinline{B_0(1+|v|)^{2J}}\) bounds both grouped frame energies. The dyadic box and (39) are separate subpower fees. Substitution into (20), together with the separately supplied diagonal bound, bounds the total first raw side.

Finally, in exact side coordinates \[\fitmath{
a=A-c-u,\qquad
k=2A-c-d+q-1-t+4u,\qquad q_C=q+c-d,
}\] direct cancellation gives \[\fitmath{
\boxed{2a-k+q_C=1+t-6u.}
}\tag{46}\] The two losses of selected width in \(\fitinline{2a}\), together with the four-power frequency enlargement, account for the coefficient \(\fitinline{6}\).

Literal physical scales require care. With \(\fitinline{u=\sum_{j\in J}w_j}\), \[\fitmath{
\log_ZK_{\rm child}
=\log_ZK+4u+(4+4|J|)\log_Z2.
}\] Frozen marked primes can also produce a volume correction \[\fitmath{
\Delta_a=\sum_{\rm frozen\ j}\log_Z(NP_j/P_j),
\qquad0\le\Delta_a\le N\log_Z2.
}\] Thus, for nominal coordinates with corrections \(\fitinline{\Delta_a,\Delta_k}\), the exact identity is \[\fitmath{
2a-k+q_C=1+t_0-6u+2\Delta_a-\Delta_k.
}\] Defining the actual defect \(\fitinline{t=t_0+2\Delta_a-\Delta_k}\) recovers (46). The raised Poisson envelope is then \[\fitmath{
H=1+(t+\mathfrak m)-6u.
}\] A permitted interval for this actual defect is a geometric input, not a consequence of cancellation. \texttt{P\allowbreak{}r\allowbreak{}o\allowbreak{}v\allowbreak{}i\allowbreak{}d\allowbreak{}e\allowbreak{}r\allowbreak{}G\allowbreak{}e\allowbreak{}o\allowbreak{}m\allowbreak{}e\allowbreak{}t\allowbreak{}r\allowbreak{}y\allowbreak{}3\allowbreak{}.\allowbreak{}b\allowbreak{}r\allowbreak{}a\allowbreak{}n\allowbreak{}c\allowbreak{}h\allowbreak{}\_\allowbreak{}w\allowbreak{}i\allowbreak{}n\allowbreak{}d\allowbreak{}o\allowbreak{}w} explicitly requires that defect window. The fixed padding must remain in its slop.

The restoration input (32) is proved in Appendix C.7; the uniform priced leaf for (43) is proved in \S{}4 and Appendix C.10; the actual defect window, support bounds and payment of \(\fitinline{Z_0+Z_f+Z_t}\) are proved in Appendix C.9. The complete assembly is Appendix C.10. No angular or arbitrary-coefficient extension is inferred. Attribution remains with Sil, ArviFoxy, Hailey Collet, OpenAI, and the inherited source authors.

\hypertarget{part-viii-4}{}

\hypertarget{c.9.-diagonal-reference-and-root-geometry}{%
\subsubsection{C.9. Diagonal, reference and root geometry}\label{c.9.-diagonal-reference-and-root-geometry}}

The stricter principal range is compatible with the diagonal and root budgets. The decisive identity is that the reference correction cancels the actual ordinary minimum \(\fitinline{L}\), leaving the root margin \[\fitmath{
\log_Z V+3u_{\max}-\frac56-\alpha\ell\le-9\omega,
\qquad \alpha=\frac23-\frac{\omega}{2}.
}\] This remains valid after every genuine common allocation and ordinary-prefix shortening.

The cube-free principal estimate is proved in Appendix C.7. The reference inlet is the general common-product, two-orientation estimate proved in Appendix C.4, \S{}\S{}1--4, equations (A.23), (A.34)--(A.35). Here it is applied with its parameters R=1+\ensuremath{\omega} and d=\ensuremath{\omega} under the explicit gates (12) below. That general estimate is distinct from the more specialized midpoint-selected theorem at the end of C.4. No additional reference theorem is imported.

\textbf{1. Uniform statement and notation.}

Fix step data with child face \[\fitmath{
P+\Lambda S\le B,\qquad s_0\le S\le c,
}\] and the literal scalar hypotheses of \S{}4.1 and Appendix C.10(13), including \(\fitinline{B<1}\), \(\fitinline{\Lambda>1}\), \(\fitinline{c<1/6}\), the monotone selection function \(\fitinline{\operatorname{sel}_\varepsilon}\), the full perturbed packed-price ledger, and the four reference inequalities below.

Use the root coordinates \[\fitmath{
P=\log_Z(X_1X_2),\qquad S=\sum_iw_i,\qquad
V=X_1X_2\prod_iP_i,
}\] so \(\fitinline{\log_ZV=P+S}\). Fix \(\fitinline{\varepsilon>0}\). The root mesh is chosen before the ambient finite slot set. After fixing that set, the ray data, slot profile, base characters and fixed inducing modulus, there are a height degree \(\fitinline{J_0}\), finite ordinary-profile control \(\fitinline{\mathcal S_0}\), and constant \(\fitinline{C_0}\). For every larger degree and control and every \(\fitinline{C_{\rm side}\ge C_0}\), one eventual threshold works for:

\begin{itemize}
\tightlist
\item
  every root shape and its surviving slot subset;
\item
  every admitted root source, including its independent ordinary heights;
\item
  every retained first block and both orientations;
\item
  every genuine allocation and selected subset;
\item
  every \(\fitinline{0\le M_c\le Z^{\varepsilon/4}}\) satisfying the selected-row multiplicity bound.
\end{itemize}

Write \[\fitmath{
\mathcal A=C_{\rm side}
 \bigl(\|W_1\|_{\mathcal S}\|W_2\|_{\mathcal S}\bigr)^2
 \operatorname{height}(t_0,V_f,t_1,t_2)^J.
}\] The diagonal contribution on a side with first common factors \(\fitinline{C,D}\) is at most \[\fitmath{
\frac12\,\mathcal A\,
 \frac{N\mathfrak f_\tau\,V^2}{NC\,ND}\,
 Z^{b(C,D;Z)+3\varepsilon}(1+|v|)^{2J}.
}\tag{1}\] Here \(\fitinline{V}\) is the original root volume, \(\fitinline{\mathfrak f_\tau}\) is the actual presented modulus of the side character, and \(\fitinline{b(C,D;Z)}\) is the retained nonnegative allowance. It is retained literally. The proof below uses only its nonnegativity.

The root source has unit original ordinary prefixes, exterior ideal \(\fitinline{1}\), common puncture product \(\fitinline{1}\), and the literal restored coefficients. Its allocated descendants generally have nonunit prefixes. Those are retained.

\textbf{2. Choose the reference before the common labels.}

Let \(\fitinline{\rho_{\rm ref}>0}\) be the \texttt{R\allowbreak{}e\allowbreak{}f\allowbreak{}e\allowbreak{}r\allowbreak{}e\allowbreak{}n\allowbreak{}c\allowbreak{}e\allowbreak{}A\allowbreak{}t} reserve. Its four inequalities are, for every admitted \(\fitinline{S}\), \[\fitmath{
\begin{aligned}
\rho_{\rm ref}&\le1-B+(\Lambda-2)S-\frac92\operatorname{sel}_\varepsilon S,\\
\rho_{\rm ref}&\le\frac1{20}+\frac32(1-B)
 +\left(\frac32\Lambda-\frac52\right)S
 -\frac92\operatorname{sel}_\varepsilon S,\\
\rho_{\rm ref}&\le1-B+(\Lambda-1)S-4\operatorname{sel}_\varepsilon S,\\
\rho_{\rm ref}&\le\frac13-2S-3\operatorname{sel}_\varepsilon S.
\end{aligned}
}\tag{2}\] Put \[\fitmath{
\rho=\frac12\min\{\rho_{\rm ref},1-B,1/100\},
\quad
\mathfrak m=\min\{\rho/1024,\varepsilon/1024,\mathfrak m_{\rm ledger}\},
}\] \[\fitmath{
\omega=\min\{\rho/100000,\varepsilon/100000\},
\qquad
u_{\max}=\operatorname{sel}_\varepsilon S+\mathfrak m.
}\tag{3}\] Choose the common frame mesh before the number of slots, with \(\fitinline{0<\mathfrak m_F\le\mathfrak m}\), also satisfying the separately proved collection and leaf tolerances. Thus both meshes are independent of the eventual number of slots. Monotonicity of selection ensures that every allocation's selected width is at most \(\fitinline{u_{\max}}\).

Define \[\fitmath{
L_-=\frac32\left(P+S+3u_{\max}-\frac56\right),\quad
U_1=\frac{P-S-1/2}{2},\quad U_2=\frac3{10}-S.
}\tag{4}\] Substitution of \(\fitinline{P\le B-\Lambda S}\) into (2), followed by \(\fitinline{\mathfrak m\le\rho/1024}\), gives \[\fitmath{
U_1-L_-\ge\rho,\qquad U_2-L_-\ge\rho,\qquad
\frac13-2S-3u_{\max}\ge\rho.
}\tag{5}\] Also \[\fitmath{
P+S\le1-\rho,\quad
5u_{\max}\le1,\quad P+S+4u_{\max}\le1.
}\tag{6}\] The last inequality follows from the third line of (2); the middle one follows from \(\fitinline{u_{\max}\le S+\mathfrak m}\), \(\fitinline{S<1/6}\), and the mesh bound.

The high branch is \[\fitmath{
P+S+3u_{\max}>\frac56-10\omega.
}\] On this branch choose \(\fitinline{\ell}\) as the midpoint of \[\fitmath{
\left(\max(0,L_-)+20\omega,\ 
      \min(U_1,U_2)-20\omega\right).
}\tag{7}\] This interval is nonempty. If \(\fitinline{L_-\ge0}\), (5) proves it. If \(\fitinline{L_-<0}\), put \(\fitinline{T=P+S+3u_{\max}-5/6>-10\omega}\). Then \[\fitmath{
U_1=\frac T2+\frac12\left(\frac13-2S-3u_{\max}\right)
>\frac{\rho}{2}-5\omega,
}\] which also leaves the required buffers. Moreover \(\fitinline{U_1<U_2}\), and \[\fitmath{
0<\ell<P/2,\qquad
\ell>L_-+20\omega,\qquad
\ell<U_i-20\omega,\qquad \ell<1.
}\tag{8}\]

Choose \[\fitmath{
\alpha=\frac23-\frac{\omega}{2}.
}\tag{9}\] Then \(\fitinline{0<\alpha<2/3}\), and (8) gives \[\fitmath{
P+S+3u_{\max}-\frac56-\alpha\ell\le-9\omega.
}\tag{10}\] Indeed the loss from replacing \(\fitinline{2/3}\) by \(\fitinline{\alpha}\) is at most \(\fitinline{\omega\ell\le\omega}\), whereas the lower-end buffer contributes more than \(\fitinline{40\omega/3}\).

On the low branch take \(\fitinline{\ell=0}\); its defining inequality gives the stronger margin \(\fitinline{-10\omega}\). Thus (10) holds on both branches.

The actual comparison rectangle uses \[\fitmath{
L_* = Z^\ell.
}\] The source minimum supplied to the diagonal estimates is bounded below by \[\fitmath{
L_{\rm ref}=\frac29 Z^\ell.
}\tag{11}\] These are different quantities, and the factor \(\fitinline{2/9}\) must remain.

\textbf{3. The same reference also meets the retained reference theorem.}

In this section only, assume the HIGH/NeedsCentering branch. The low branch uses the uncentered root and the direct low exceptional margin in (10); it invokes neither (14) nor the reference triangle below. In particular, \(\fitinline{\ell=0}\) alone does not imply the third reference gate.

Set \(\fitinline{R=1+\omega}\), \(\fitinline{d=\omega}\), and use reference tolerances \(\fitinline{\delta_{\rm ref}=\xi_{\rm ref}=\omega}\). Equations (5) and (8) imply \[\fitmath{
\begin{aligned}
P+2S+4\omega&<1+\omega,\\
\ell+\frac{1+2\omega}{2}+S+2\omega
 &\le\frac45(1+\omega),\\
1+2\omega+2\ell-P+S+2\omega
 &\le\frac12(1+\omega).
\end{aligned}
}\tag{12}\] For example, the second follows from \(\fitinline{\ell<3/10-S-20\omega}\); the third follows from \(\fitinline{2\ell<P-S-1/2-40\omega}\).

These are precisely the scalar gates of the retained homogeneous two-orientation reference estimate. Its application also requires one event, chosen before the moving rectangles, which pays \[\fitmath{
\left(\frac94+4\right)\frac94\,2^N\le Z^\omega,
\quad b_\Phi Z\le Z^{1+\omega},
}\] and \[\fitmath{
\max\{1,C_{\rm cond}(\eta)\,b_\Phi Z\}\le Z^{1+2\omega}.
}\tag{13}\] All factors on the left are fixed before \(\fitinline{Z}\). Hence (13) holds eventually.

For every rectangle of the original product, in the original enclosure, with either side at most \(\fitinline{L_*}\), the resulting kept reference energy is \[\fitmath{
\le \operatorname{radialControl}(\Phi)\,
 C\bigl(\|W_1\|_{\mathcal S}\|W_2\|_{\mathcal S}\bigr)^2
 \operatorname{height}(0,0,t,t)^J
 Z^{1+\omega+\varepsilon/100}.
}\tag{14}\] The actual predicate is \texttt{i\allowbreak{}n\allowbreak{}i\allowbreak{}t\allowbreak{}i\allowbreak{}a\allowbreak{}l\allowbreak{}K\allowbreak{}e\allowbreak{}e\allowbreak{}p} with the original internal reference. Its exclusion of fixed-inducing rows is used, rather than an unrestricted reference average.

If either original side is already at most \(\fitinline{L_*}\), (14) applies directly. Otherwise use the reference rectangle \[\fitmath{
(L_*,\,X_1X_2/L_*).
}\] Equation \(\fitinline{\ell<P/2}\), together with both original sides being at least \(\fitinline{L_*}\), keeps it inside the inherited enclosure. The exact centered-source identity and \(\fitinline{|x+y|^2\le2|x|^2+2|y|^2}\) give the original energy as at most twice the sum of the centered raw energy and (14).

The changed principal exponent also fits the inherited first-reference reserve. Put \(\fitinline{\kappa_\alpha=2/3-\alpha=\omega/2}\). For every reference width \(\fitinline{R'\le2}\) and reserve \(\fitinline{\rho'\ge\rho/2}\), \[\fitmath{
\kappa_\alpha R'\le\omega\le\frac23\rho'.
}\tag{15}\] For the retained midpoint reference, its reserve is \(\fitinline{(R'-P-2S-4d)/3\ge\rho'/3}\), while its exponent is less than \(\fitinline{R'/2}\). Thus (15) pays \(\fitinline{\kappa_\alpha\ell}\). This verifies the cube-free \texttt{A\allowbreak{}l\allowbreak{}p\allowbreak{}h\allowbreak{}a\allowbreak{}R\allowbreak{}e\allowbreak{}f\allowbreak{}e\allowbreak{}r\allowbreak{}e\allowbreak{}n\allowbreak{}c\allowbreak{}e\allowbreak{}M\allowbreak{}a\allowbreak{}r\allowbreak{}g\allowbreak{}i\allowbreak{}n}; using the old \(\fitinline{\alpha\ge2/3}\) argument would be invalid.

The analytic content of (14) is the retained two-orientation reference theorem. The new work here is its exact admission, including (12)--(15).

\textbf{4. Allocation-dependent second scales.}

Fix an actual first common pair \(\fitinline{C,D}\), a common subset \(\fitinline{E}\), and a retained first frequency dyad \(\fitinline{2^n\ge1/2}\). Let \[\fitmath{
K=2^n Z^\kappa,\qquad \kappa=\operatorname{cutoffLoss}(N,\varepsilon)>0.
}\tag{16}\] For each genuine allocation \(\fitinline{\mathcal B}\), let \(\fitinline{r_{\mathcal B}}\) be the actual residual source after selected-slot removal, before ordinary-prefix renormalization. Put \[\fitmath{
W_{\mathcal B}=\prod_{j\in J_{\mathcal B}}2P_j,\qquad
K_{\mathcal B}=16K W_{\mathcal B}^4,
}\] \[\fitmath{
Y_{\mathcal B}=\max\{1,\operatorname{normCap}(r_{\mathcal B})\},
\qquad
T_{\mathcal B}=\frac{Y_{\mathcal B}^2}{K_{\mathcal B}},
\qquad
R_{\mathcal B}=T_{\mathcal B}Z^{\mathfrak m_F/2}.
}\tag{17}\] This is the literal \texttt{T\allowbreak{}f\allowbreak{}r\allowbreak{}e\allowbreak{}q}. The tail condition \(\fitinline{Y_{\mathcal B}^2/K_{\mathcal B}\le T_{\mathcal B}}\) holds with equality.

These choices are uniform over all allocations, but their values depend on the allocation. Replacing them by an unrelated common second scale would require a new argument.

The root has \[\fitmath{
Z\le K_{\rm root}\le Z^2,\qquad
N\mathfrak f_\eta\le Z,\qquad
NC,ND\le Z^8,\qquad V\le Z.
}\tag{18}\] First common factors divide original columns, proving their norm bounds. Their common exponents are at most two, so either orientation satisfies the weaker condition \(\fitinline{ND\le(NC)^3}\) required below.

The retained first-frequency cutoff gives \[\fitmath{
K\le4\,T_{\rm first}Z^{3\kappa/2},
\quad
T_{\rm first}=\frac{Ne\,Nr\,V^2}{K_{\rm root}NC\,ND}.
}\tag{19}\] The exact character presentation is \[\fitmath{
\mathfrak f_\tau
=\mathfrak f_\eta\,(m_{\rm bad})\,(72)\,(er).
}\tag{20}\] Therefore \[\fitmath{
K\le
4Z^{3\kappa/2}\frac{N\mathfrak f_\tau\,V^2}{NC\,ND\,Z}.
}\tag{21}\]

All descendant pool caps are at most \(\fitinline{Z^8}\). The residual norm caps satisfy \[\fitmath{
\operatorname{normCap}(r_{\mathcal B})
\le \left(\frac94\right)^2 2^N V.
}\tag{22}\] Eventually \(\fitinline{Y_{\mathcal B}\le Z^8}\), \(\fitinline{1\le W_{\mathcal B}\le Z^2}\), and \(\fitinline{1/2\le K\le Z^{33}}\). Thus \[\fitmath{
T_{\mathcal B}\le Z^{16},\qquad R_{\mathcal B}\le Z^{17}.
}\tag{23}\] The diagonal theorem is invoked with norm cap exponent \(\fitinline{8}\), radial exponent \(\fitinline{17}\), and a generous polynomial scale cap \(\fitinline{40}\).

The family-conductor budget is bounded by a fixed constant times \[\fitmath{
N\mathfrak f_\tau\cdot N(r_{\mathcal B}.R)\,
 \operatorname{poolCap}(r_{\mathcal B})^2
 Z^{16}R_{\mathcal B}.
}\] Here \(\fitinline{r_{\mathcal B}.R=C}\), and the child ray character adds only the fixed \(\fitinline{(12)}\)-factor. Substitution of (18), (22), and (23), followed by payment of the fixed conductor factor, bounds this by \(\fitinline{Z^{80}}\). The common puncture product is still \(\fitinline{1}\), and the surviving marked product is at most the original marked product. This checks every cap used by the uniform principal estimate.

\textbf{5. The three per-seed bounds.}

Let \(\fitinline{v}\) now denote a residual squarefree seed, and let \(\fitinline{h\in\mathbb R}\) be the independent row height.

For zero frequency, the inherited column divisor bound gives \[\fitmath{
|\beta_v(I)|\le C_{\rm col}\operatorname{prof}(r)\,Y^{\eta/2},
}\] on nonzero coefficients, with \(\fitinline{NI\le Y}\) and \(\fitinline{v\mid I}\). The retained zero-correlation divisor estimate then gives \[\fitmath{
|Z_0(r,v)|
\le C_0\operatorname{prof}(r)^2\frac{Y^{1+\eta}}{Nv},
\qquad C_0=128C_{\rm col}^2.
}\tag{24}\] The physical zero-frequency contribution is \(\fitinline{K_{\mathcal B}|\widehat\Phi(0)|}\) times (24). Heights have modulus at most one in this estimate.

The cube-free principal input, for \(\fitinline{0\le\alpha<2/3}\), retains every prefix and supplies the common-saving factor \[\fitmath{
(\mathcal C\mathcal D)^{-1/3}
 \min\left\{
  \max(1,x/\mathcal C)^{-\alpha},
  \max(1,x/\mathcal D)^{-\alpha}
 \right\}.
}\] For \(\fitinline{\mathcal C,\mathcal D\ge1}\), this is at most \(\fitinline{x^{-\alpha}}\). To prove it, set \(\fitinline{m=\min(\mathcal C,\mathcal D)}\). Then \[\fitmath{
(\mathcal C\mathcal D)^{-1/3}
\le m^{-2/3}\le m^{-\alpha},
}\] and \(\fitinline{m\max(1,x/m)\ge x}\).

Take \(\fitinline{x=L/H}\), where \(\fitinline{H=\max(Nd_1,Nd_2)}\) is the actual inherited prefix maximum. For \(\fitinline{0<L_{\rm ref}\le L}\), \[\fitmath{
(L/H)^{-\alpha}
\le (L/H)^{-2/3}\,
 \underbrace{(L/L_{\rm ref})^{2/3}
  \max(1,L_{\rm ref})^{2/3-\alpha}}_{\mathcal R_\alpha}.
}\tag{25}\] This remains valid when the shortened length is below \(\fitinline{1}\).

The important point is that the diagonal estimates use the unrenormalized \texttt{r\allowbreak{}e\allowbreak{}s\allowbreak{}t\allowbreak{}O\allowbreak{}f} source. Its prefix is shortened, but its \(\fitinline{L}\)-field is the original root \(\fitinline{L}\). Hence \(\fitinline{\mathcal R_\alpha}\) is the same for all genuine allocations.

After the retained fixed-row count, whole-kernel integration and second common sum, the per-seed fixed-inducing contribution is bounded by \[\fitmath{
\begin{aligned}
|Z_f(r,v;h)|
\le{}&
C_f(1+|h|)^{2J_L}
 V(r)W_{\mathcal B}\,
 \operatorname{savedBudget}(r;K,W_{\mathcal B})\\
&\times \mathcal R_\alpha\,
 Z^{2B_{\rm norm}\delta}\,
 \frac{C_{\rm sum}Z^{\eta_{\rm sum}}}{Nv}.
\end{aligned}
}\tag{26}\] All finite profile and fixed-\(\fitinline{Q}\) factors are retained in \texttt{s\allowbreak{}a\allowbreak{}v\allowbreak{}e\allowbreak{}d\allowbreak{}B\allowbreak{}u\allowbreak{}d\allowbreak{}g\allowbreak{}e\allowbreak{}t}. Its exact raw geometric part is \[\fitmath{
\operatorname{shape}(V,p,H,L,K,W)
=(V/p)^{4/3}K^{5/6}W^{7/3}(L/H)^{-2/3},
}\tag{27}\] where \(\fitinline{p=Nd_1Nd_2}\). In particular \[\fitmath{
V(r)\operatorname{savedBudget}(r)
=Q_{\rm fixed}\operatorname{sourceProfile}(r)^2
 16^{5/6}a^{-4/3}\operatorname{shape}(r).
}\tag{28}\] The power \(\fitinline{W^{7/3}}\) is exactly \(\fitinline{(16KW^4)^{5/6}/W}\), after removing the fixed \(\fitinline{16^{5/6}}\).

For the tail, the retained arbitrary-decay theorem, using (17) and the polynomial caps, gives \[\fitmath{
|Z_t(r,v;h)|
\le W_{\mathcal B}\,
 C_t\operatorname{prof}(r)^2
 \|\Phi\|_{\mathcal S_t}Z^{-31}.
}\tag{29}\] Its derivative order and constant depend on the positive cutoff mesh, fixed before the family. This bound is uniform in the seed and heights; it does not claim reciprocal-seed decay.

\textbf{6. Seed sums and selected-width costs.}

Complete the residual-seed family only after the positive-energy reduction. Its members have norm at most \(\fitinline{X=Z^8}\). The retained ideal count gives \[\fitmath{
\sum_{v}\frac1{Nv}\le2304\log X,
\qquad
\#\{v\}\le128XW_{\mathcal B}.
}\tag{30}\] The latter deliberately permits the source's weaker \(\fitinline{W_{\mathcal B}}\)-factor.

Let \(\fitinline{A_{\mathcal B}}\) be the height-independent selected Cauchy mass, and let \(\fitinline{\Pi_{\mathcal B}=\prod_{j\in J_{\mathcal B}}P_j}\). The selected-row average proved in blocks 1--2 gives \[\fitmath{
A_{\mathcal B}\,2^{|J_{\mathcal B}|}M_c
 \frac{W_{\mathcal B}}{\Pi_{\mathcal B}^{4/3}}
\le C_N M_c W_{\max}^{2/3},
\quad W_{\max}=2^N Z^{u_{\max}}.
}\tag{31}\] The exponent is \[\fitmath{
u+u-\frac43u=\frac23u.
}\] There is no second assignment-count factor.

For the zero term and tail, the corresponding bounds are \[\fitmath{
\begin{aligned}
S_0&\le C_NM_c(W_{\max}+Y_r)
 W_{\max}^4\operatorname{prof}(s)^2Y_{\rm up}^{\eta},\\
S_t&\le C_NM_c\,128XW_{\max}^2\operatorname{prof}(s)^2,
\end{aligned}
}\tag{32}\] where \[\fitmath{
Y_r=\left(\frac94\right)^2 2^N V,\qquad
Y_{\rm up}=2\left(\frac94\right)^2 2^N Z.
}\tag{33}\] Thus the small exponent on the upper cap costs \(\fitinline{Z^\eta}\), not \(\fitinline{Z^{8\eta}}\).

\textbf{7. Genuine common allocations and the balanced reference.}

For an allocation \(\fitinline{\mathcal B}\) of \(\fitinline{C}\), let \(\fitinline{b_{\mathcal B}}\) be its literal frozen coefficient. The actual tuple-product identity and annular support give \[\fitmath{
\operatorname{allocationShape}_{\mathcal B}
\le \operatorname{slotCost}(s)
 (V/NC)^{4/3}K^{5/6}W^{7/3}
 \left(\frac{L}{H\,NC}\right)^{-2/3}.
}\tag{34}\] The shortened prefix maximum is bounded by \(\fitinline{H\,NC}\); the complete common ideal occurs here. The live-volume comparison supplies the fixed factor \(\fitinline{2^{4N/3}}\) in \texttt{s\allowbreak{}l\allowbreak{}o\allowbreak{}t\allowbreak{}C\allowbreak{}o\allowbreak{}s\allowbreak{}t}. It does not replace the prefix by \(\fitinline{1}\).

Define \[\fitmath{
\mathcal B_{\rm ref}
=\frac{N\mathfrak f_\tau V^2}{NC\,ND}Z^{b(C,D;Z)}.
}\tag{35}\] Using (19), (20), \(\fitinline{ND\le(NC)^3}\), and \(\fitinline{b\ge0}\), direct division of (34) by (35) yields \[\fitmath{
\frac{\operatorname{allocationShape}_{\mathcal B}}{\mathcal B_{\rm ref}}
\le
\operatorname{slotCost}(s)\,
4^{5/6}Z^{5\kappa/4}\,
V(K_{\rm root}N\mathfrak f_\eta F_{\rm pres})^{-5/6}
W^{7/3}(L/H)^{-2/3}.
}\tag{36}\] Indeed the remaining common-factor ratio is \(\fitinline{(NC)^{-1/2}(ND)^{1/6}\le1}\); the remaining modulus factor is also at most \(\fitinline{1}\).

This derives the balanced allocation reference estimate directly. Its logarithmic exponent is \[\fitmath{
E_{\rm ref}
=\log_ZV-\frac56m-\frac23\log_Z(L/H)
 +\frac73\log_ZW+\frac56x,
}\tag{37}\] where \[\fitmath{
m=\log_Z(K_{\rm root}N\mathfrak f_\eta F_{\rm pres})\ge1,
\qquad
Z^x=4Z^{3\kappa/2}.
}\] Thus \(\fitinline{Z^{5x/6}=4^{5/6}Z^{5\kappa/4}}\).

At the root \(\fitinline{H=1}\). In the following bound use \(\fitinline{W_{\mathcal B}\le W_{\max}}\) in (37); equivalently evaluate the upper reference exponent at \(\fitinline{W=W_{\max}}\). Combining (25), (31), and (37) then gives the crucial cancellation: \[\fitmath{
\begin{aligned}
\mathcal R_\alpha Z^{E_{\rm ref}}W_{\max}^{2/3}
&\le L_{\rm ref}^{-2/3}\max(1,L_{\rm ref})^{2/3-\alpha}
 2^{3N}Z^{\log_ZV-5m/6+3u_{\max}+5x/6}\\
&\le(9/2)^{2/3}8^N
 Z^{\log_ZV+3u_{\max}-5/6-\alpha\ell+5x/6}\\
&\le(9/2)^{2/3}8^N Z^{5x/6}.
\end{aligned}
}\tag{38}\] The actual root \(\fitinline{L}\) has disappeared by equality. The shortening was paid before this cancellation. Equation (10), rather than the invalid endpoint \(\fitinline{\alpha=2/3}\), proves the final inequality.

\textbf{8. Pay zero and tail against the same reference.}

By (21), \[\fitmath{
K/\mathcal B_{\rm ref}\le4Z^{3\kappa/2-1}.
}\] Equations (32)--(33) therefore reduce the zero term to fixed factors times \[\fitmath{
Z^{3\kappa/2-1}
 \left(W_{\max}^5+Y_rW_{\max}^4\right)Y_{\rm up}^{\eta}.
}\tag{39}\] Both main exponents are nonpositive by (6): \[\fitmath{
5u_{\max}-1\le0,\qquad
\log_ZV+4u_{\max}-1\le0.
}\] This pays the zero frequency even when the residual norm cap is below \(\fitinline{1}\).

For the tail, the root minimum implies \(\fitinline{V^{-1}\le Z^2}\), while \(\fitinline{NC,ND\le Z^8}\). Since the modulus norm and allowance factors are at least \(\fitinline{1}\), \[\fitmath{
\mathcal A Z^{3\varepsilon-20}
\le \mathcal A\mathcal B_{\rm ref}Z^{3\varepsilon}.
}\tag{40}\] Thus it suffices to pay the tail against the smaller budget on the left. Equations (29)--(32) supply decay \(\fitinline{Z^{-31}}\), seed count \(\fitinline{Z^8}\), and two harmless selected-width factors bounded by \(\fitinline{Z^2}\).

\textbf{9. Each contribution receives one sixth.}

Choose exactly \[\fitmath{
\delta_c=\varepsilon/80,\qquad
\delta=\min\{\varepsilon/160,1/10\},\qquad
\eta=\eta_{\rm sum}=\varepsilon/10.
}\tag{41}\] Then \[\fitmath{
\alpha+2\delta<13/15<1.
}\] The genuine allocation count satisfies \[\fitmath{
(\#\mathcal B)^2\le C_c(N,\delta_c)(NC)^{\delta_c}
\le C_cZ^{8\delta_c}.
}\] Keep \(\fitinline{M_c\le Z^{\varepsilon/4}}\). The source parameter choice gives \(\fitinline{0<\kappa\le\min(1,\varepsilon/100000)}\).

After (38), the remaining fixed-inducing exponent is bounded by \[\fitmath{
8\delta_c+\varepsilon/4+16\delta
 +\eta_{\rm sum}+2\eta+\frac54\kappa
\le\frac{60001}{80000}\varepsilon<3\varepsilon.
}\tag{42}\] The zero-frequency exponent is bounded by \[\fitmath{
8\delta_c+\varepsilon/4+\eta+\frac32\kappa
\le\frac{90003}{200000}\varepsilon<3\varepsilon.
}\tag{43}\] The tail exponent is \[\fitmath{
8\delta_c+\varepsilon/4+8+2-31
=\frac7{20}\varepsilon-21<3\varepsilon-20.
}\tag{44}\] These fractions were independently recomputed with standard-library \texttt{F\allowbreak{}r\allowbreak{}a\allowbreak{}c\allowbreak{}t\allowbreak{}i\allowbreak{}o\allowbreak{}n} arithmetic.

All remaining factors are fixed before the growing sources: the finite ray count, phase-table mass, \(\fitinline{K_{\rm slot}^N}\), \(\fitinline{256^N,1024^N}\), kernel constants, fixed-\(\fitinline{Q}\) cost, window moments, radial Fourier value, and finite profile-control conversion constants. The profile and height factors are placed in \(\fitinline{\mathcal A}\), not in an eventual threshold depending on the variable profiles or heights.

For fixed \(\fitinline{K}\) and \(\fitinline{e<c}\), \[\fitmath{
K(1+\log Z)Z^e\le\frac16 C_{\rm side}Z^c
}\] eventually. Apply this separately to (42), (43), and (44). Thus \[\fitmath{
E_0,E_f,E_t\le\frac16\mathcal A\mathcal B_{\rm ref}Z^{3\varepsilon},
}\] and their sum is the half-budget in (1).

\textbf{10. The same allocation scales supply the nonfixed geometry.}

Let \(\fitinline{V_\ell}\) be the renormalized live volume and \(\fitinline{V_s}\) the renormalized selected residual volume. On a nonzero allocation, \[\fitmath{
V\le V_\ell NC\le2^N V,\qquad V_sZ^u=V_\ell.
}\tag{45}\] The first inequality follows from the actual frozen primes lying in \(\fitinline{[P_i,2P_i]}\). Hence, with \(\fitinline{a=\log_ZV_s}\), \(\fitinline{A=\log_ZV}\), \(\fitinline{c=\log_ZNC}\), \[\fitmath{
A-c-u\le a\le A-c-u+N\log_Z2.
}\tag{46}\]

The packed identity \(\fitinline{NV_C\,ND=Ne\,Nr\,NC}\) converts (19) into \[\fitmath{
K(NC)^2 Z\le4NV_C\,V^2Z^{3\kappa/2}.
}\] Therefore \[\fitmath{
K_{\mathcal B}(NC)^2 Z
\le64\,16^N V^2NV_C\,Z^{4u+3\kappa/2}.
}\tag{47}\] Also \(\fitinline{(V_C)\mid C^2}\), so \(\fitinline{q_C\le2c}\).

Let \[\fitmath{
H_{\rm env}=2a-\log_ZK_{\mathcal B}+\mathfrak m_F+q_C.
}\] Choose the eventual event so that \[\fitmath{
\log_Z(64\,16^N)\le\kappa/2,\quad
2N\log_Z2\le\rho,\quad
\log_Z(9/2)\le\kappa.
}\tag{48}\] Then (46)--(48), \(\fitinline{K_{\mathcal B}\ge1}\), and \(\fitinline{A\le1-\rho}\) prove \[\fitmath{
-2\kappa\le H_{\rm env}-1+6u,\qquad
H_{\rm env}-1\le1+2\kappa.
}\tag{49}\] This is the actual defect window, including the fourth-power radius and fixed padding.

For ordinary lengths, division by a prefix never increases \(\fitinline{\ell_Z(X)=\log_Z\max(1,X)}\). The root lengths are at least \(\fitinline{2/9}\), giving \[\fitmath{
P_{\rm side}\le P+2\log_Z(9/2)\le P+2\kappa.
}\tag{50}\] Thus the inherited face and the actual selected/live widths satisfy the full perturbed packed-price ledger with tolerance \(\fitinline{2\kappa}\). The parameter choice ensures \[\fitmath{
2\kappa\le t_{\rm ledger},\qquad 2\kappa\le1-H_{\min}.
}\] The small branch selects nothing; the large branch selects within one frame-mesh error. Both therefore have the admitted envelope range.

If the residual norm cap is below \(\fitinline{1}\), there is no nonzero residual column. Otherwise \[\fitmath{
Y_{\mathcal B}=\operatorname{normCap}(r_{\mathcal B})
=\left(\frac94\right)^2 2^{n_{\rm rest}}V_s.
}\] Using (17), the common Poisson cutoff is met once \[\fitmath{
32\left[\left(\frac94\right)^22^N\right]^2
 Z^{\mathfrak m_F/2}\le Z^{\mathfrak m_F}.
}\tag{51}\] Because \(\fitinline{\mathfrak m_F>0}\) was fixed first, this is one eventual event for every allocation. This is why the allocation-dependent \(\fitinline{T_{\mathcal B}}\) serves both the tail and the nonfixed rows.

All second common factors and second artificial divisors have norm at most \(\fitinline{Z^8}\). Their exterior ideal is \(\fitinline{C}\), also at most \(\fitinline{Z^8}\), so \[\fitmath{
N\operatorname{rad}(C\,C_2\,L_2)\le Z^{24}.
}\tag{52}\] This proves the mask exponent \(\fitinline{B_{\rm mask}=24}\), in either orientation.

For the radius bound, (47)--(48), \(\fitinline{q_C\le2c}\) and \(\fitinline{A\le1}\) give \[\fitmath{
\log_ZK_{\mathcal B}\le1+4u+2\kappa<18,
}\] using the retained \(\fitinline{u\le1/5}\), \(\fitinline{\kappa\le1}\). This proves both the support-scale radius hypothesis and the weaker \(\fitinline{Z^{40}}\) cap used in the selected-row multiplicity. The coarse \(\fitinline{W_{\mathcal B}\le Z^2}\) alone would not prove that latter cap from \(\fitinline{K\le Z^{33}}\).

The residual columns are bounded by \(\fitinline{Z^8}\), the second twist by \(\fitinline{Z^{16}}\), and the underlying source support scales lie between \(\fitinline{Z^{-16}}\) and \(\fitinline{Z^{18}}\). Actual dyadic endpoints include fixed enlargements, at most the retained factor four. Therefore \[\fitmath{
M_{\rm box}=\lfloor19\log_2Z\rfloor+2
}\] contains every physical dyadic vector. Equations (45)--(52) prove the geometric supplier with \(\fitinline{c_b=19}\), rather than merely asserting its interface.

\textbf{11. Multiplicities, both orientations, and the full provider.}

The selected-row divisor estimate on the common polynomial radius gives \[\fitmath{
M_c=\max(C_m,1)Z^{\varepsilon/8}\le Z^{\varepsilon/4}
}\] eventually. The source uses a uniform radius cap \(\fitinline{Z^{40}}\) for this step. The empty selected set has multiplicity \(\fitinline{1}\); for nonempty selected sets the zero row vanishes before this bound is used.

The separate original-label multiplicity is bounded by \[\fitmath{
M_{\rm lab}=\left\lceil C_\delta Z^{\varepsilon/8}\right\rceil
\le Z^{\varepsilon/2}.
}\] Both constants precede the actual root shapes and labels. Genuine-allocation cardinality is paid separately.

For the right side, exchange \(\fitinline{C,D}\), transport the common subset, and use the literal right fixed-pair character. The nominal scale is invariant under this exchange, its character has the transported presentation (20), and both cube-free inequalities \(\fitinline{ND\le NC^3}\) and \(\fitinline{NC\le ND^3}\) hold. Thus the same constants prove (1) on both sides. There is no assumption that the first common factors are equal.

The priced nonfixed result of \S{}4 supplies the other half of the same balanced budget. Take the maximum height degree over the diagonal result and the finitely many surviving slot dimensions; take the union of their profile controls; choose \(\fitinline{C_{\rm side}}\) large enough for both. All choices occur before the variable root sources and heights. Adding the two half-budgets proves the complete summed raw-side contract.

The finite source closure uses this contract at internal loss \(\fitinline{\varepsilon/12}\), so its \(\fitinline{3\varepsilon}\) side loss becomes \(\fitinline{\varepsilon/4}\) in the outer units. The root-reference estimate (14), the centered triangle, and the inherited first-energy reconstruction then use the already reserved outer loss. No new assumption about the candidate endpoint enters these choices.

The analytic references used here are supplied within this manuscript:

\begin{enumerate}
\def\labelenumi{\arabic{enumi}.}
\tightlist
\item
  The cube-free centered principal estimate feeding (25)--(28), with \(\fitinline{0\le\alpha<2/3}\), all positive shortened scales, original masks, prefixes and independent heights.
\item
  The retained two-orientation homogeneous reference estimate feeding (14), under the explicitly checked gates (12)--(13).
\item
  \S{}4 and Appendix C.10 give the uniform priced-leaf statement for the nonfixed half.
\end{enumerate}

The calculations above discharge the new allocation, scale, reference-margin, geometry, multiplicity and budget obligations. The first item is Appendix C.7; the second is the general homogeneous reference theorem proved in Appendix C.4 with the gates checked above. The corrected order---pack the original masked row, then restore---is preserved for the nonfixed interface. Credit remains with Sil, ArviFoxy, Hailey Collet, OpenAI and the retained source authors.

\hypertarget{part-viii-6}{}

\hypertarget{c.10.-complete-generic-source-assembly}{%
\subsubsection{C.10. Complete generic source assembly}\label{c.10.-complete-generic-source-assembly}}

The generic native step closes with the repaired cube-free blocks. Its leaf price is exactly \[\fitmath{
\frac34\,\mathrm{deletionNeed},
}\] and its LOW root uses the fixed lower scale \(\fitinline{2/9}\). It does not obtain a principal saving from an arbitrarily large original minimum. The half-price source of \S{}4 enters the same closure through its full priced leaf and the scalar ledger proved in \S{}5.

The statement below uses the retained Gauss/reciprocity lemmas, the repaired packing--restoration block, and the repaired diagonal supplier as explicit dependencies. It supplies the complete source assembly used in \S{}\S{}4--5.

\textbf{1. Exact generic statement and uniformity.}

Write \(\fitinline{G(B,L,c)}\) for the existing ordinary \texttt{G\allowbreak{}l\allowbreak{}o\allowbreak{}b\allowbreak{}a\allowbreak{}l\allowbreak{}S\allowbreak{}o\allowbreak{}u\allowbreak{}r\allowbreak{}c\allowbreak{}e}: the literal normalized two-ordinary-factor, finitely marked natural-row energy is \(\fitinline{O(Z^{1+\epsilon})}\) whenever \[\fitmath{
\ell_Z(X_1)+\ell_Z(X_2)+L\sum_iw_i\le B,\qquad
\sum_iw_i\le c,
\quad
\ell_Z(X)=\log_Z\max(1,X).
}\tag{1}\] The marked coefficients remain \[\fitmath{
\operatorname{idealCoeff}(\chi_{\eta_0,\theta_i},I)\,
\operatorname{annularWeight}(W,Z^{w_i},33/50,v_i,I),
}\tag{2}\] with the original height coefficient, ray family, and natural zero convention. This is not an arbitrary-coefficient theorem.

For every \(\fitinline{\epsilon>0}\), the mesh is chosen before the arbitrary ambient finite slot type. After that type, ray modulus/subgroup, smooth marked profile \(\fitinline{W}\), radial support cap, fixed anchor characters, \(\fitinline{\eta_0}\), and fixed \(\fitinline{Q}\) have been fixed, choose the height degree, finite Schwartz-seminorm set, constant, and threshold. They then work simultaneously for:

\begin{itemize}
\tightlist
\item
  every slot subset, collision partition and inverse-restoration subset;
\item
  all allowed ray labels, widths, independent frequencies and heights;
\item
  all original ordinary profiles and all positive shortened ordinary lengths;
\item
  every finite set of retained rows and every allowed nonnegative radial profile.
\end{itemize}

The exact shared source contract requires \[\fitmath{
Q_{\rm int}=Q\cap\operatorname{modulus}(\eta_0),\quad
Q\subseteq\mathfrak m_{\rm ray},\quad Q_{\rm int}\ne0,\quad
Q_{\rm int}\ne\mathcal O,\quad Q_{\rm int}\subseteq(72).
}\] Require \(\fitinline{\mathfrak m_{\rm ray}\ne0}\), the ray subgroup to contain global units, and the marked profile to be smooth supported in \(\fitinline{[1,2]}\). Fix \(\fitinline{b_\Phi\ge2}\) before the constants. After constants and onset, allow every Schwartz \(\fitinline{\Phi}\) supported in \(\fitinline{( -\infty,b_\Phi]}\) with \(\fitinline{\Re\Phi\ge0}\), controlled only by the displayed \(\fitinline{\operatorname{diagControl}(\Phi)}\) factor. The live widths are nonnegative and at most the chosen mesh; the real slot power is \(\fitinline{33/50}\), and the nonnegative height bound dominates every individual slot frequency. The admitted clipped ordinary lengths satisfy the stated fixed cap.

The main step is \[\fitmath{
G(B_0,L_0,c_0)+\mathrm{StepGates}(D)
+\mathrm{PricedPackedSource}(D.\mathrm{price})
\Longrightarrow G(B,L,c),
}\tag{3}\] under \(\fitinline{51/100\le\beta}\), with envelope interval \(\fitinline{[H_{\min},H_{\max}+1]}\) and ordinary length cap \(\fitinline{4}\).

The analytic dependencies in (3) are precisely the repaired:

\begin{enumerate}
\def\labelenumi{\arabic{enumi}.}
\tightlist
\item
  joint first-divisor and physical second-column identities;
\item
  original-masked packing, followed by cube-free restoration and grouped energy;
\item
  diagonal supplier, including fixed inducing rows, zero frequency and tails.
\end{enumerate}

Their hypotheses are checked below rather than replaced by an abstract \texttt{S\allowbreak{}t\allowbreak{}e\allowbreak{}p\allowbreak{}I\allowbreak{}n\allowbreak{}p\allowbreak{}u\allowbreak{}t\allowbreak{}s} assumption.

\textbf{2. The native leaf, including its parent-domain fallback.}

Let \[\fitmath{
d(P,s,H)=
\max\!\left(0,\min\!\left(s,
\max\!\left(s-c_0H,\frac{P+L_0s-B_0H}{L_0}\right)\right)\right).
}\tag{4}\] Here \(\fitinline{P=\ell_Z(X_1)+\ell_Z(X_2)}\), \(\fitinline{s=\sum_iw_i}\), and \[\fitmath{
H=\log_Z(N(V)K).
}\] The exact packed row is evaluated at \(\fitinline{Vz}\), retained with \(\fitinline{\mathrm{initialKeep}_\eta(Q_{\rm int},Vz)}\), and radially weighted at scale \(\fitinline{K}\).

The injection \(\fitinline{z\mapsto Vz}\) and nonnegativity give \[\fitmath{
E_{\rm packed}(V,K)\le E_{\rm canonical}(N(V)K).
}\tag{5}\] There is no lattice-density compensation in this inequality. The scale changes to \(\fitinline{N(V)K=Z^H}\).

The original OpenAI bound \(\fitinline{\beta\le7/8}\) supplies \[\fitmath{
2\beta-1\le\kappa,\qquad \kappa=\frac34.
}\tag{6}\] This is proved from the original moments and literal probe in Appendix B.9; neither the endpoint nor the new child face is used.

For a deleted subset \(\fitinline{D_{\rm del}}\), factorization of the normalized slot product is literal: \[\fitmath{
R_A(z)=
\left(\prod_{i\in D_{\rm del}}\mathrm{normalizedSlot}_i(z)\right)
R_{A\setminus D_{\rm del}}(z).
}\tag{7}\] The retained natural-slot estimate, applied with zero widths outside \(\fitinline{D_{\rm del}}\), gives \[\fitmath{
\left|\prod_{i\in D_{\rm del}}\mathrm{normalizedSlot}_i(z)\right|^2
\le C(1+|t|+V_{\rm ht})^J
 Z^{\epsilon_1+\kappa\sum_{D_{\rm del}}w_i}
}\tag{8}\] on every retained row. It applies because the keep excludes the fixed inducing family, the natural conductor has the required polynomial cap, and the coefficient in (2) is unchanged. Multiplying (8) by the nonnegative radial weight and summing proves partial deletion without changing the remaining keep, ordinary profiles or lengths.

Suppose first \(\fitinline{P\le B_0H}\), with \(\fitinline{B_0,c_0>0}\). Greedily delete whole slots until their total width \(\fitinline{d'}\) first reaches (4). Each slot has width at most \(\fitinline{m}\), so \[\fitmath{
d\le d'\le d+m,\qquad s_{\rm ret}=s-d'\ge0.
}\] The definition of \(\fitinline{d}\) then gives \[\fitmath{
s_{\rm ret}\le c_0H,\qquad P+L_0s_{\rm ret}\le B_0H.
}\tag{9}\] Apply the parent source at base \(\fitinline{Z^H}\), with widths divided by \(\fitinline{H}\). Since \(\fitinline{H\ge H_{\min}>0}\), choosing the original mesh at most \(\fitinline{H_{\min}}\) times the parent mesh makes this application uniform over the entire envelope.

For leaf loss \(\fitinline{\epsilon}\), put \(\fitinline{M_{\max}=H_{\max}+1}\). Request the parent theorem with loss \(\fitinline{\epsilon/[4\max(1,M_{\max})]}\), before choosing its unit mesh and the finite arithmetic/live data. On base \(\fitinline{Z^H}\) this costs at most \(\fitinline{H\epsilon/[4\max(1,M_{\max})]\le\epsilon/4}\) in the original base. The deletion loss is \(\fitinline{\epsilon/8}\), and \(\fitinline{m\le\epsilon/8}\). Thus \[\fitmath{
E_{\rm packed}\ll
\mathcal A\,Z^{H+\kappa d+
  \epsilon/4+\epsilon/8+\kappa\epsilon/8}
\le
\mathcal A\,Z^{H+\kappa d+\epsilon}.
}\tag{10}\] Here \(\fitinline{\mathcal A}\) is the prescribed radial diagonal control, the square of the two-profile control, and a polynomial height factor.

If \(\fitinline{P>B_0H}\), then \(\fitinline{d=s}\). The parent face cannot be applied even after deleting every slot; instead the retained native canonical-energy theorem gives \[\fitmath{
E_{\rm packed}\ll\mathcal A Z^{H+\kappa s+\epsilon}.
}\tag{11}\] This fallback is essential.

For the formally allowed degenerate case \(\fitinline{B_0\le0}\) or \(\fitinline{c_0\le0}\), (4) equals \(\fitinline{s}\) on the evaluated nonnegative domain. Again (11) supplies the price. Consequently the generic native leaf is valid for every \texttt{N\allowbreak{}a\allowbreak{}t\allowbreak{}i\allowbreak{}v\allowbreak{}e\allowbreak{}G\allowbreak{}a\allowbreak{}t\allowbreak{}e\allowbreak{}s} instance, not only the positive faces in the actual chain.

This proves the native priced leaf, with the stated parent fallback and choice order.

\textbf{3. Scalar gates and actual perturbed geometry.} The direct native initiation in \S{}4.2 and the complete open-family induction in \S{}5 prove the actual price, height, parent-cap and reference inequalities. They replace the numerical native-chain table. The assembly below applies to each of those steps with its own fixed positive reserves; no numerical intermediate endpoint is a premise. Here is the entire geometric domain to which those prices are applied. If \(\fitinline{S}\) is the root marked width, \(\fitinline{w}\) its live width, \(\fitinline{f=S-w}\) its frozen width, \(\fitinline{u}\) the selected width and \(\fitinline{s}\) the retained width, the actual clipped frame defines \(\fitinline{c_g\ge f}\) and \[\fitmath{
\begin{gathered}
 P_C+(c_g-f)+LS\le B+\tau,\quad P_C\ge0,\\
 \operatorname{sel}(w)\le u\le\operatorname{sel}(w)+m,
 \quad0\le s\le w-u,\\
 -\tau\le t\le1+6u+\tau,\qquad H=1+t-6u.
\end{gathered}
}\tag{13}\] In particular \(\fitinline{P_C\le B-Lw+\tau}\), because \(\fitinline{c_g-f\ge0}\), \(\fitinline{S\ge w}\) and \(\fitinline{L>0}\). Thus the full perturbed ledger (UF18)--(UF19) applies with its \(\fitinline{W=w}\); \S{}4.2 proves the native alternative on the same domain. The height envelope and the four positive reference reserves are part of those local proofs. Below the positive width floor the native all-slot estimate pays the error, as shown in section 5 below. The symbol \(\fitinline{\mathrm{StepGates}(D)}\) in (3) abbreviates precisely these geometric, price and reference inequalities together with the fixed contract in section 1; it introduces no separate theorem or unproved code assumption.

\textbf{4. Root reference, LOW principal input, and profiles.}

For a literal restored root, use its actual shortened ordinary lengths \(\fitinline{Y_1,Y_2}\), and set \[\fitmath{
P=\log_Z(Y_1Y_2),\qquad
S=\sum_{\text{surviving singleton slots}}w_i,\qquad A=P+S.
}\tag{20}\] Thus \(\fitinline{P}\) is an actual logarithm; it need not equal the clipped ordinary width. Nevertheless shortening and subset inclusion give \[\fitmath{
P+LS\le B,\qquad s_0\le S\le c.
}\tag{21}\]

Choose a positive reference reserve \(\fitinline{r_{\rm ref}}\) for the four expressions \[\fitmath{
\begin{aligned}
R_1(S)&=1-B+(L-2)S-\tfrac92\operatorname{sel}(S),\\
R_2(S)&=\tfrac1{20}+\tfrac32(1-B)
 +(\tfrac32L-\tfrac52)S-\tfrac92\operatorname{sel}(S),\\
R_3(S)&=1-B+(L-1)S-4\operatorname{sel}(S),\\
R_4(S)&=\tfrac13-2S-3\operatorname{sel}(S).
\end{aligned}
}\tag{22}\] For native steps one may take \[\fitmath{
r_{\rm ref}=\min\{1-B,\tfrac13-(2+3a)c\}>0.
}\] The four gate inequalities above prove (22) uniformly.

Put \[\fitmath{
r_*=\tfrac12\min(r_{\rm ref},1-B,1/100),\quad
m_*\le\min(r_*/1024,\epsilon/1024,m_{\rm ledger}),
}\] \[\fitmath{
\omega=\min(r_*/10^5,\epsilon/10^5),\quad
\alpha=\tfrac23-\omega/2,\quad
u_{\max}=\operatorname{sel}(S)+m_*.
}\tag{23}\] Selection monotonicity bounds every later allocation's selected width by \(\fitinline{u_{\max}}\).

The HIGH branch is \[\fitmath{
A+3u_{\max}>5/6-10\omega.
}\] Only there choose \[\fitmath{
\ell_-=\tfrac32(A+3u_{\max}-5/6),\quad
\ell_1=(P-S-\tfrac12)/2,\quad \ell_2=\tfrac3{10}-S,
}\] \[\fitmath{
\ell=\tfrac12\left(\max(0,\ell_-)+20\omega+
                  \min(\ell_1,\ell_2)-20\omega\right).
}\tag{24}\] From (21)--(23), \[\fitmath{
0<\ell<P/2,\quad
\ell>\ell_-+20\omega,\quad
\ell<\ell_1-20\omega,\quad
\ell<\ell_2-20\omega,\quad \ell<1.
}\] Consequently \[\fitmath{
A+3u_{\max}-5/6-\alpha\ell\le-9\omega.
}\tag{25}\] The retained HIGH reference theorem is applied only here, with its actual \(\fitinline{Z^\ell}\) reference. Its numerical conditions follow from the last two upper bounds in (24).

On the LOW branch set \(\fitinline{\ell=0}\). Then directly \[\fitmath{
A+3u_{\max}-5/6\le-10\omega.
}\tag{26}\]

The literal LOW source constructor is particularly important. It sets \[\fitmath{
X_1^{\rm src}=Y_1,\quad X_2^{\rm src}=Y_2,\quad
Y_1^{\rm ref}=2/9,\quad
Y_2^{\rm ref}=\frac{Y_1Y_2}{2/9},\quad L_{\rm src}=2/9.
}\tag{27}\] The two rectangles have identical product. The reference rectangle vanishes exactly: the first ordinary profile is supported at norm at most \(\fitinline{(9/4)(2/9)=1/2}\), whereas every nonzero ideal has norm at least one. After any inherited ordinary prefix \(\fitinline{e_1}\), its first length is even smaller.

Thus the centered source in (27) is exactly the uncentered physical row. For an inherited prefix pair \(\fitinline{e}\), the principal estimate contains \[\fitmath{
\frac{\sqrt{T\prod_iP_i}}{N(e_1)N(e_2)}
\left(\frac{2/9}{\max(N(e_1),N(e_2))}\right)^{-\alpha}.
}\tag{28}\] It supplies \textbf{no original-length saving}. The prefix factor in (28) is retained in the allocation ledger, and (26) pays the root excess. This resolves the \texttt{d\allowbreak{}o\allowbreak{}C\allowbreak{}e\allowbreak{}n\allowbreak{}t\allowbreak{}e\allowbreak{}r\allowbreak{}=\allowbreak{}f\allowbreak{}a\allowbreak{}l\allowbreak{}s\allowbreak{}e} issue.

If either actual shortened length is at most \(\fitinline{2/9}\), the literal row itself is zero. Otherwise (27) is legitimate.

At this common-height inlet, the source stores its supplied profiles and common height \(\fitinline{t=t_0}\). For independently specified original profiles \(\fitinline{f_j}\) and heights \(\fitinline{t_j}\), those supplied profiles are precisely \(\fitinline{p_j(x)=f_j(x)x^{i(t_j-t_0)}}\) from C.5(117a), with the original outside phase and the exact reference correction (117b). Its dictionary applies the common norm power once; it does not raise the relative-height profiles a second time. The original independent heights are not identified. Since \(\fitinline{|t_0|\le V}\), their derivative costs are polynomial in the displayed height control of (S5). The slot profile bound is fixed from \(\fitinline{W}\) independently of all frequencies, roots and ordinary profiles. No height-dependent coefficient-mass constant or profile-dependent onset is introduced.

\textbf{5. Both kinds of leaves reach the actual physical side.}

Apply the repaired packing identity on the original masked residual row, retaining its \(\fitinline{G}\), \(\fitinline{D}\), seed and natural-zero masks. Only after height absorption, character congruence on the valid columns, and packing of the keep/radial weight does cube-free restoration produce canonical unpunctured leaves.

For each actual second frame \(\fitinline{r}\), put \[\fitmath{
\lambda=\operatorname{colLog}(r),\quad
\nu=\operatorname{twistLog}(r),\quad
P_C=\operatorname{gP}(r),\quad H=\operatorname{commonH}(r).
}\tag{29}\] Thus \(\fitinline{\operatorname{vol}(r)=Z^\lambda}\) and \(\fitinline{N(V_0)=Z^\nu}\). These are the actual frame volume and packed conductor. The previous packing block supplies \(\fitinline{V_C\mid C^2}\), the physical radius \(\fitinline{16KW^4}\), the divisor fibres, and the actual defect window. No nominal annulus identity is substituted for them.

The physical side uses clipped ordinary lengths, while \(\fitinline{P}\) in (20) is the actual logarithm of the root product. All four unrenormalized rectangle lengths in the nonzero raw construction are at least \(\fitinline{2/9}\). For each such \(\fitinline{X}\), \[\fitmath{
\ell_Z(X)\le\log_ZX+\log_Z(9/2).
}\] Division by integral prefixes cannot increase the clipped length. Hence \(\fitinline{P_C\le P+2\log_Z(9/2)}\). Choose the common onset so that \(\fitinline{2\log_Z(9/2)\le\tau}\). With \(\fitinline{f=S-w}\), define exactly \[\fitmath{
c_g=f+P+\tau-P_C,\qquad
A_{\rm aux}=\lambda+u,\quad c_{\rm aux}=0,\quad
q_{\rm aux}=\nu.
}\tag{30a}\] Then \(\fitinline{c_g\ge f}\), and \[\fitmath{
P_C+(c_g-f)+LS=P+\tau+LS\le B+\tau.
}\] This proves the tolerated root inequality (13), including negative actual \(\fitinline{P}\), with no change to physical normalization. The underlying side scalar is \(\fitinline{t_0=H-1+6u-m}\); the ledger uses \(\fitinline{t=t_0+m=H-1+6u}\). Its bounds follow from the actual defect window \(\fitinline{-\tau\le H-1+6u}\) and \(\fitinline{H-1\le1+\tau}\). The retained slot width is \(\fitinline{w-u}\), or a smaller width after restoration. These identities give the literal hypotheses of (13).

The grouped source estimate, with its prescribed profile, height and outside normalization factors retained, has exponential factor at most \(\fitinline{Z^{A_{\rm aux}-c_{\rm aux}+q_{\rm aux}+2\epsilon}}\). Consequently its exact frame conversion is \[\fitmath{
Z^{A_{\rm aux}-c_{\rm aux}+q_{\rm aux}+2\epsilon}
=\operatorname{vol}(r)N(V_0)Z^uZ^{2\epsilon}.
}\tag{30}\] In particular the full split-bound prefactor \(\fitinline{(\operatorname{vol}(r_s)/c_{\rm low})\operatorname{gAmp}}\) and the full height factor are still present. Equation (30) converts only the exponential factor; it does not delete these factors.

Below the positive floor, choose no slots. Then \(\fitinline{u=0}\), \(\fitinline{s\le\mathrm{floor}\le\epsilon/12}\), and the native all-slot price is at most \(\fitinline{3s/4\le\epsilon/16}\). This yields the same bound (30), with \(\fitinline{H\in[H_{\min},H_{\max}]}\); this branch needs no density estimate.

Ordinary rescaling within a dyadic box preserves the exact factor \(\fitinline{\sqrt{Y_1Y_2/(X_1X_2)}}\), whose square is at most four. Fourier separation uses a common finite profile-control family; the Lipschitz price bounds absorb its mesh errors. The ordinary cap four covers the actual frame cap three with this padding.

For completeness, the physical comparison after (30) uses the actual allocation volumes, not rounded scale equalities. Let \(\fitinline{\mathcal V_s}\) be the surviving frame volume before reinserting the selected slots, \(\fitinline{\mathcal V_{\rm live}}\) the live volume, and \(\fitinline{\mathcal V_{\rm root}}\) the original root volume. The allocation identities give \[\fitmath{
\mathcal V_s Z^u=\mathcal V_{\rm live},\qquad
\mathcal V_{\rm live}N(C)\le2^N\mathcal V_{\rm root}.
}\] Together with the accepted weighted frame and prefix normalizations, the converted frame factor is thus bounded by \(\fitinline{2^N\mathcal V_{\rm root}N(V_C)/N(C)\,Z^{2\epsilon}}\) in this orientation. The opposite orientation uses \(\fitinline{D,V_D}\). Let \(\fitinline{e}\) denote the actual first common-subset generator and \(\fitinline{r}\) the active conductor generator, including its retained unit. The exact packed identities are \[\fitmath{
N(V_C)N(D)=N(e)N(r)N(C),\qquad
N(V_D)N(C)=N(e)N(r)N(D).
}\tag{31}\] Their product is \(\fitinline{N(V_C)N(V_D)=(N(e)N(r))^2}\). The factor \(\fitinline{N(e)N(r)}\) varies with the allocation and remains through the joint Hilbert/divisor normalization. It is not a fixed character cost. The actual presented side modulus is \[\fitmath{
\mathfrak f_\tau=\operatorname{modulus}(\eta)\,
\mathfrak m_{\rm fixedBad}\,(72)\,(er).
}\] All factors are nonzero integral ideals. Thus \(\fitinline{N(e)N(r)\le N(\mathfrak f_\tau)}\), and (31) implies \[\fitmath{
N(V_C)N(D)\le N(\mathfrak f_\tau)N(C),\qquad
N(V_D)N(C)\le N(\mathfrak f_\tau)N(D).
}\] These inequalities connect the two physical orientations to the side budget (32), retaining all prefix, grouped-weight and reciprocal divisor factors from the packing lemma. The symbol \(\fitinline{\mathfrak f_\tau}\) always denotes this actual presented modulus; it is not replaced by a possibly smaller primitive conductor. The fixed \(\fitinline{2^N}\) enters the existing dimension-dependent constant, not an additional power of \(\fitinline{Z}\). This is the literal auxiliary-scalar adapter of \texttt{C\allowbreak{}u\allowbreak{}b\allowbreak{}e\allowbreak{}F\allowbreak{}r\allowbreak{}e\allowbreak{}e\allowbreak{}/\allowbreak{}S\allowbreak{}e\allowbreak{}c\allowbreak{}o\allowbreak{}n\allowbreak{}d\allowbreak{}/\allowbreak{}P\allowbreak{}r\allowbreak{}o\allowbreak{}v\allowbreak{}i\allowbreak{}d\allowbreak{}e\allowbreak{}r\allowbreak{}G\allowbreak{}e\allowbreak{}o\allowbreak{}m\allowbreak{}e\allowbreak{}t\allowbreak{}r\allowbreak{}y\allowbreak{}3\allowbreak{}.\allowbreak{}b\allowbreak{}r\allowbreak{}a\allowbreak{}n\allowbreak{}c\allowbreak{}h\allowbreak{}\_\allowbreak{}w\allowbreak{}i\allowbreak{}n\allowbreak{}d\allowbreak{}o\allowbreak{}w} and \texttt{P\allowbreak{}r\allowbreak{}o\allowbreak{}v\allowbreak{}i\allowbreak{}d\allowbreak{}e\allowbreak{}r\allowbreak{}L\allowbreak{}a\allowbreak{}b\allowbreak{}e\allowbreak{}l\allowbreak{}S\allowbreak{}p\allowbreak{}l\allowbreak{}i\allowbreak{}t\allowbreak{}.\allowbreak{}e\allowbreak{}x\allowbreak{}p\allowbreak{}\_\allowbreak{}l\allowbreak{}e\allowbreak{}\_\allowbreak{}o\allowbreak{}f\allowbreak{}\_\allowbreak{}g\allowbreak{}e\allowbreak{}o\allowbreak{}m\allowbreak{}e\allowbreak{}t\allowbreak{}r\allowbreak{}y}.

\textbf{6. All first-energy contributions and their constants.}

For each first-common pair \(\fitinline{C,D}\), inactive subset and retained dyad, use the actual allocation-dependent second scale \[\fitmath{
K_B=16KW_B^4,\qquad
T_B=\frac{\max(1,\operatorname{capRes}_B)^2}{K_B},
\qquad R_B=T_BZ^{m_F/2},
}\] where \(\fitinline{0<m_F\le m_*}\). A residual cap below one means there is no active residual column; that contribution is zero.

Let \(\fitinline{\mathcal V}\) denote the original source volume. The balanced budget for one orientation is \[\fitmath{
\mathcal B_\tau(C,D)
=\mathcal A\,\frac{N(\mathfrak f_\tau)\mathcal V^2}{N(C)N(D)}
 Z^{\operatorname{allowance}(C,D)}Z^{3\epsilon}.
}\tag{32}\] The repaired diagonal lemma bounds zero frequency, fixed inducing rows and tails by three shares \(\fitinline{1/6}\) of (32). Thus they consume at most one half.

For the label half, the grouped estimate (30) gives the required residual-volume factor. The actual selected multiplicity is bounded by \[\fitmath{
M_c=\max(C_m,1)Z^{\epsilon/8}\le Z^{\epsilon/4}
}\] eventually. Each first-allocation count is at most \(\fitinline{C_aZ^{\epsilon/8}}\). After squared allocation counting and the one \(\fitinline{M_c}\) factor, the label constant has exponent \[\fitmath{
2\epsilon+3\epsilon/8<3\epsilon.
}\tag{33}\] The \(\fitinline{2\epsilon}\) already includes the grouped estimate and the dyadic-box/second-label multiplicities. All remaining factors---ray cardinalities, \(\fitinline{256^N}\), \(\fitinline{4^N}\), marked-profile bounds, and the finite Fourier constants---are fixed before \(\fitinline{Z}\). Choosing the side constant at least twice their combined constant supplies the other half of (32).

Constants are chosen uniformly over residual cardinalities \(\fitinline{0,\ldots,N}\): take the maximum degree, union of seminorm sets, and sum of positive constants. The diagonal supplier is monotone under these increases. This permits both suppliers to use one amplitude and one eventual event.

Now apply joint Hilbert Cauchy to the actual first artificial-divisor/row sum. The artificial-divisor sum remains inside this inequality. The two half-budgets give the full balanced side budget, and balancing by \(\fitinline{N(C)/N(D)}\) preserves the product of the two window budgets. The ray projector costs \(\fitinline{256}\). Only the common band mass remains: \[\fitmath{
\frac{\|E_{\rm retained}\|}{\mathcal V}
\le
C_{\rm ret}\,64\,C_{\rm presentation}\,
  N(\mathfrak f_\eta)b_\Phi\,
  (\operatorname{windowMass}_J(0))^2
  \mathcal A\,Z^{1+3\epsilon+\epsilon/8}.
}\tag{34}\] No artificial-divisor mass is reintroduced.

The exact first decomposition is \[\fitmath{
E_{\rm full}=E_{\rm zero}+E_{\rm retained}+E_{\rm discarded}.
}\] The retained original zero/tail theorem bounds the other two pieces, normalized by \(\fitinline{\mathcal V}\), by \(\fitinline{C_{\rm rem}\mathcal A Z^{1+4\epsilon}}\). Combining with (34) gives \[\fitmath{
\frac{\|E_{\rm full}\|}{\mathcal V}
\le C_{\rm full}\mathcal A Z^{1+4\epsilon}.
}\tag{35}\]

The support caps used here follow from the actual row: \[\fitmath{
\mathcal V=Z^{P+S}\le Z,\qquad
\operatorname{normCap}
=(9/4)^2\,2^{\#\mathrm{slots}}\,\mathcal V\le Z^2
}\] eventually; all required weaker \(\fitinline{Z^8}\) caps follow. The radius is \(\fitinline{b_\Phi Z}\), and the fixed conductor and dimensional factors enter the threshold before moving assignments.

There are four exhaustive root-return cases:

\begin{enumerate}
\def\labelenumi{\arabic{enumi}.}
\tightlist
\item
  an actual ordinary length is at most \(\fitinline{2/9}\): exact zero;
\item
  LOW: (35) bounds the uncentered row directly through (27);
\item
  HIGH with one length below \(\fitinline{Z^\ell}\): the retained reference bound alone;
\item
  HIGH with both lengths at least \(\fitinline{Z^\ell}\): centered plus reference, using \(\fitinline{|x+y|^2\le2|x|^2+2|y|^2}\).
\end{enumerate}

The reference bound and triangle are used only in cases 3--4. Applying this construction internally at \(\fitinline{\epsilon/12}\) leaves ample room for the required restricted bound \(\fitinline{Z^{1+\epsilon}}\).

\textbf{7. Final restoration, with each row's actual surviving width.}

Finally, classify \textbf{each restored row by its own surviving singleton width} \(\fitinline{S'}\). If \(\fitinline{S'<s_0}\), use the clipped parent-admission width \(\fitinline{P_{\rm cl}=\ell_Z(Y_1)+\ell_Z(Y_2)}\), not the actual product logarithm of (20). Shortening and subset inclusion give \(\fitinline{P_{\rm cl}+LS'\le B}\), and hence \[\fitmath{
P_{\rm cl}+L_0S'
\le B+(L_0-L)s_0\le B_0,\qquad S'\le c_0.
}\tag{37}\] Each further-shortened canonical leaf has no larger clipped width, so the same parent admission holds for it. Actual \(\fitinline{P}\) remains unchanged in the physical volume and reference ledger. Forward cube-free domination produces canonical parent leaves on subsets at shortened lengths. Its \(\fitinline{Z^{\epsilon/2}}\) mass and the parent's \(\fitinline{Z^{1+\epsilon/2}}\) estimate give the restricted row bound. Profile phase shifts cost only the already uniform polynomial height factor.

If \(\fitinline{S'\ge s_0}\), use the provider just proved. Merge constants only after both families are fixed. Complete collision and inverse cube-free restoration then uses restricted loss \(\fitinline{\epsilon/2}\) and its remaining \(\fitinline{Z^{\epsilon/2}}\) mass to recover the original \(\fitinline{G(B,L,c)}\). All restoration subsets and assignments are included.

The source constructions used in this manuscript have crossing \(\fitinline{s_0=0}\). Every restored row therefore reaches the provider just proved, including width zero through its separate unmarked branch. The stated \(\fitinline{s_0>0}\) restoration remains a reusable generality of the assembly; its parent must be strictly earlier and all its displayed inequalities must hold. This completes the original global source, including every complementary sector and the stipulated order of choices.

\hypertarget{local-reference-and-foundational-input-ledger}{%
\subsection{Local reference and foundational-input ledger}\label{local-reference-and-foundational-input-ledger}}

The main statements use the labels T, S, SN, HP, UF, D, C, FB, GT, AE and K. Appendix A's locally proved theorems are numbered A.2.1--A.8.1; the probe's Gram and witness lemmas are B.2.1 and B.5.1; initialization and Mellin return are B.9.1--B.9.4. Retained arithmetic equation labels such as (67), (MF3.40) and (203) are unique named families, located below. In the source modules C.6--C.10, bare equation numbers are local to their module; references from another module state that module explicitly. Numeric ideals such as \(\fitinline{(1)}\), \(\fitinline{(12)}\) and \(\fitinline{(72)}\) are unchanged arithmetic objects, not equation labels.

\begin{longtable}[]{@{}
  >{\raggedright\arraybackslash}p{(\columnwidth - 2\tabcolsep) * \real{0.5000}}
  >{\raggedright\arraybackslash}p{(\columnwidth - 2\tabcolsep) * \real{0.5000}}@{}}
\toprule\noalign{}
\begin{minipage}[b]{\linewidth}\raggedright
Label family or interface
\end{minipage} & \begin{minipage}[b]{\linewidth}\raggedright
Local proof location
\end{minipage} \\
\midrule\noalign{}
\endhead
\bottomrule\noalign{}
\endlastfoot
(1)--(4), original sextic character and finite phases & A.1--A.2 \\
(B.*) arithmetic-reflection anchors and completed energy & A.3 \\
Inverse moment, (C.*) anchors and Theorem A.4.1 & A.4 \\
Plain moment, (67)--(83), (D.*) anchors and Theorem A.5.1 & A.5 \\
Exact finite Poisson identities and probe series & A.6 \\
MF2, MF3, theorems A.7.1--A.7.7 & A.7 \\
Literal overlapping inverse lists & A.8 \\
(188)--(203), signed Gram, table of local powers, reciprocal signal & B.1--B.4 \\
(204)--(223), same-zero witnesses and complete numerator & B.5--B.7 \\
Generic remaining rows and contours & B.8 and \S{}8 \\
I.1--I.13, original seven-eighths comparison and actual supremum & B.9 \\
Ordinary norm pullback, every presentation, line \(\fitinline{\Re s=1}\) & B.10 \\
VII.2.1--VII.2.12, complete local projectors and restoring mass & C.1 \\
VII.3.1--VII.3.3, whole fourth-power column & C.2 \\
VII.4.1--VII.4.38, sextic sieve and power-free restoration & C.3 \\
A.1--A.51 source-reference equations & C.4 \\
(115)--(129), full first Poisson kernel and allocation weights & C.5 \\
T1--T15, exact masked packed twins & C.6 \\
Cube-free projectors, both normalized masses, principal range & C.7 \\
Joint artificial divisors, weighted Hilbert norm, full physical collection & C.8 \\
Diagonal, reference, root geometry and presented side modulus & C.9 \\
Full source-step geometry, four exhaustive root returns and restoration & C.10 \\
\end{longtable}

In the retained algebraic theorem of \S{}\S{}2--10 the symbol \(\fitinline{c_*}\) denotes the then-attained endpoint (T3); Theorem H uses \(\fitinline{c_{80}}\); \S{}1 identifies the current theorem. The native source path uses \(\fitinline{c_{\rm seed}}\) for its starting cap. Other positive constants and local variables are scoped to their displayed lemmas. In particular the endpoint perturbation parameter in \S{}9 is separate from physical heights, row-conductor exponents and later analytic losses. Independent ordinary heights are implemented by common raised profiles and their exact outside phases, as in C.5(117a)--(117b); no equality of those heights is assumed.

\hypertarget{explicit-foundational-boundary}{%
\subsubsection{Explicit foundational boundary}\label{explicit-foundational-boundary}}

\begin{longtable}[]{@{}
  >{\raggedright\arraybackslash}p{(\columnwidth - 2\tabcolsep) * \real{0.5000}}
  >{\raggedright\arraybackslash}p{(\columnwidth - 2\tabcolsep) * \real{0.5000}}@{}}
\toprule\noalign{}
\begin{minipage}[b]{\linewidth}\raggedright
Input or construction
\end{minipage} & \begin{minipage}[b]{\linewidth}\raggedright
Status used in this proof
\end{minipage} \\
\midrule\noalign{}
\endhead
\bottomrule\noalign{}
\endlastfoot
Fixed-field Hecke continuation/functional equation, fixed-ray prime ideal theorem, and Dirichlet nonvanishing on \(\fitinline{\Re s=1}\) & Classical inputs with the exact versions, fixed-modulus scope and applications stated in A.1.1. \\
Quadratic and cubic large sieves and unconditional cubic-theta/cusp formulas & Precisely cited classical inputs in A.1.1; their squarefree-ideal, fixed-ray and coefficient hypotheses are checked at the uses in A.3. \\
Finite phases, completed reflection, original inverse/plain/marked moments, literal probe and seven-eighths initialization & Written proofs retained locally in A and B; no older optimized endpoint is an input. \\
Natural selected-cell norm (SN9), with arbitrary fixed positive gap and early mesh & The already adopted PR95 analytic input, stated in full in \S{}3. Its upstream proof is not newly re-proved or newly certified by this consolidation. Its source scope is unchanged. \\
Natural-row coefficient adapters, histogram, whole-root half price, cube-free source assembly and finite open-family induction & Proved locally in \S{}\S{}3--5 and C, from the explicit preceding inputs. \\
Count faces/windows and algebraic finite return data & Exact unchanged data with the mathematical interpretation and whole-domain verification method in \S{}\S{}6--10. Historical arithmetic receipts remain bound to those data, not to rewritten prose. \\
\end{longtable}

These are the complete external-input boundaries of the attained route. Optional catalogue sources in the research guide supply no missing lemma here. The original OpenAI framework, Collet's reflection/probe/return framework, Sil's fourth-power construction, ArviFoxy's marked-source work, and Bebel's natural selected norm and subsequent source extensions retain their supplied attribution. Exact source notices and assistance credits remain in the source records; the old-to-new map identifies their recovery without requiring historical narrative for this proof.

\hypertarget{appendix-d}{}

\hypertarget{appendix-d.-marked-anchor-induction-literal-sources-and-proof-spine}{%
\subsection{Appendix D. Marked-anchor induction: literal sources and proof spine}\label{appendix-d.-marked-anchor-induction-literal-sources-and-proof-spine}}

Sil's accepted PR137, \emph{Anchor induction for marked inverse moments}, at head 9d81a4c0\allowbreak{}b0e8338e\allowbreak{}f458f3ba\allowbreak{}5139657a\allowbreak{}e5404fde supplies simultaneous actual inverse and paired ordinary marked moments on rectangles extending beyond \(\fitinline{P+2S=1}\). These are additional suppliers, not new members of the old sloping family \(\fitinline{G(B,L,c_s)}\). This local integration claims no new endpoint. The accepted submission's ordinary Dirichlet endpoint is \(\fitinline{2183367/2500000}\); its physical return is another lane's responsibility.

\hypertarget{evidence-and-foundational-boundary}{%
\subsubsection{1. Evidence and foundational boundary}\label{evidence-and-foundational-boundary}}

SOURCE\_USE.json binds exact source paths, declarations, SHA-256 and Git-blob hashes. ARITHMETIC\_RECEIPT.json independently reconstructs the rational recursion, fourteen endpoint gates, earlier-line minorants and return-line minorants. This lane inspected source text and proof bodies; it did not build Lean, execute contributor code, audit axioms or reproduce the accepted verification pipeline. The controller owns acceptance receipts and licence preservation.

The argument uses the inherited Eisenstein framework of Cycle068 Appendix A: ideal factorization, sextic reciprocity with natural zero extensions, cubic Gauss identities, masked lattice Poisson summation, radial Fourier decay, the functional equation with finite profile controls, supported same-zero detectors, and the completed inverse/ordinary energy dictionaries. The ordinary half also uses the completed common-divisor, collision, second-divisor and diagonal assembly in its fork of the density-split proof. Those are explicit admitted analytic inputs at their inherited evidence scope. This note explains the new induction and its exact analytic contracts; it is not an independent new proof of these foundations.

The catalogue was searched before development: directory, mathematical records and Lean-module records, including marked-anchor, marked inverse, density-split and global-source terms. Related entries include local:G304 and local:G90. The owner's fetch receipt, read back here, confirms clean HEAD = origin/main = 44824138\allowbreak{}c975e345\allowbreak{}52d35e31\allowbreak{}4cda60b3\allowbreak{}2ad346b5 at 2026-11-15T20:12:40.684646+00:00. The owner's discovery readback reports no exact PR129/137/141 head records. Both receipts are in .local/catalogue-manager-20261018/coordination-local069-20261115/. Catalogue status supplies discovery and provenance, not an analytic premise.

\hypertarget{quantities-and-coefficient-classes}{%
\subsubsection{2. Quantities and coefficient classes}\label{quantities-and-coefficient-classes}}

Work in \(\fitinline{\mathcal O=\mathbf Z[\omega]}\), with \(\fitinline{\omega=e^{2\pi i/3}}\), ideal norm \(\fitinline{NI}\), and element norm \(\fitinline{|u|^2}\) in the fixed embedding. Count each ideal once through its fixed primary generator and all six units in element-row sums. Every presented character retains primitive-conductor, redundant-presentation and bad-prime zeros. Even a zeroth or sixth residue-character power remains zero at a nonunit.

For the actual presented row character \(\fitinline{\chi_u}\), use \[\fitmath{
 M_u(X;\Phi,\theta)=X^{-1/2}\sum_{I\ne0}
       \mu(I)\chi_u(I)\Phi(NI/X)(NI/X)^{i\theta},
}\] \[\fitmath{
 A_u(X;p,t)=X^{-1/2}\sum_{I\ne0}\chi_u(I)(NI)^{it}p(NI/X).
}\tag{A1}\] The inverse polynomial retains its exact M\"{o}bius sign. Its \(\fitinline{\chi_u}\) is precisely the character of the fixed RowData, including the deleted base and its presentation masks; no replacement inverse coefficient vector is asserted. A whole-product phase \(\fitinline{X^{i\theta}}\) does not affect a row norm.

A direct normalized prime slot is \[\fitmath{
 Q_i(u)=P_i^{-1/2}\sum_{\mathfrak p\in\mathcal P_i}
                 b_i(\mathfrak p)\chi_{\mathfrak p}(u),
 \qquad P_i=Z^{w_i}. 
}\tag{A2}\] The residue factor is the natural ideal-row homomorphism, with its zeros.

The \textbf{new inverse} source requires a fixed-ray canonical slot class. Before the scale fix a nonzero ray modulus \(\fitinline{\mathfrak M}\), subgroups \(\fitinline{H_{\rm twist}\le H_{\rm pool}}\) containing global units, one smooth \(\fitinline{W}\) supported in \(\fitinline{[1,2]}\), and real part \(\fitinline{33/50}\). Each pool is the whole original identity-class prime pool for \(\fitinline{H_{\rm pool}}\) at scale \(\fitinline{P_i}\), intersected with \(\fitinline{N\mathfrak p/P_i\in[1,2]}\). Its coefficient is \[\fitmath{
 b_i(I)=\nu_{\vartheta,\eta}(I)\eta(I)
       W(NI/P_i)(NI/P_i)^{-33/50+iv_i}. 
}\tag{A3}\] One common relative ray character \(\fitinline{\vartheta}\) is used for all slots; the heights \(\fitinline{v_i}\) are independent. The modulus and subgroups belong to the fixed class, not to live scale data. Allowing a new modulus after \(\fitinline{Z}\) could disguise arbitrary prime subsets and is not this theorem. Whole slots may be omitted, but arbitrary subsets of primes in one slot are not admitted. Pools may overlap.

The inverse inlet additionally requires \(\fitinline{|b_i(\mathfrak p)|\le1}\), prime-annulus support and \(\fitinline{N(\operatorname{modulus}(\mathrm{deletedBase}))<P_i}\). A fixed bump may be normalized by its fixed supremum and its finite-slot homogeneity restored in the late constant. Every finite-character zero in (A3) stays.

The \textbf{ordinary} source is literally the original UnitFaceSource: original coefficient-one ordinary sums, two independent profiles/heights, common whole-product height, anchors, punctures, masks, natural/presentation zeros and initialKeep. Its original per-slot relative twists need not equal one common twist. Do not impose the inverse's narrower class on this separate ordinary theorem; conversely, do not transfer the native inverse's arbitrary bounded coefficient class to the new inverse rectangle.

For the ordinary keep, put \(\fitinline{Q_0=Q\cap\mathfrak q_{\eta_0}}\), with \(\fitinline{Q_0\ne0,\mathcal O}\) and \(\fitinline{Q_0\subseteq(72)}\). Exclude precisely rows whose actual natural presentation is induced by a primitive character of conductor dividing \(\fitinline{Q_0}\). The inverse source has no such keep: all nonzero ball rows, including fixed-inducing and principal rows, are inside its conclusion.

\hypertarget{reusable-theorem-and-order-of-choices}{%
\subsubsection{3. Reusable theorem and order of choices}\label{reusable-theorem-and-order-of-choices}}

For \(\fitinline{d>0}\), define the anchor rectangle \[\fitmath{
 \mathcal R(R,z;d)=\{(P,S):0\le P\le R-d,\ 
                     0\le S\le(1-R)/2+z-d\}. 
}\tag{A4}\] Write \(\fitinline{\mathrm{Anchor}(R,z)}\) when both following complete contracts hold for every \(\fitinline{d>0}\).

\textbf{Inverse.} For every \(\fitinline{\epsilon>0}\), choose a mesh \(\fitinline{m>0}\) before the slot bound \(\fitinline{K}\), canonical class, smooth profile \(\fitinline{\Phi}\) supported in some \(\fitinline{[a,b]\subset(0,\infty)}\), and row data. After these are fixed choose an integer \(\fitinline{J}\), \(\fitinline{C>0}\), and \(\fitinline{Z_0>1}\). They work simultaneously for every \(\fitinline{Z\ge Z_0}\), every at-most-\(\fitinline{K}\) slot vector, every \(\fitinline{h_v\ge0}\), admissible coefficients (A3) with \(\fitinline{|v_i|\le h_v}\), widths \(\fitinline{0\le w_i\le m}\), finite nonzero row set \(\fitinline{\mathcal U}\) with \(\fitinline{|u|^2\le Z}\), real \(\fitinline{\theta}\), and \(\fitinline{(r,\sum_iw_i)\in\mathcal R(R,z;d)}\): \[\fitmath{
 \sum_{u\in\mathcal U}|M_u(Z^r;\Phi,\theta)\prod_iQ_i(u)|^2
       \le C Z^{1+\epsilon}(1+|\theta|+h_v)^{2J}. 
}\tag{A5}\] The profile and row data precede \(\fitinline{J,C,Z_0}\); the live heights do not.

\textbf{Ordinary.} The full positive original energy \[\fitmath{
 \sum_{u:\mathrm{initialKeep}}|A_u(X_1;p_1,t_1)A_u(X_2;p_2,t_2)|^2
       \prod_i|Q_i(u)|^2\Re\Psi(|u|^2/Z) 
}\tag{A6}\] is bounded by \[\fitmath{
 C\mathcal D(\Psi)
 [\mathcal S_{\mathcal F}(p_1)\mathcal S_{\mathcal F}(p_2)]^2
 (1+|t_1|+|t_2|+V)^J Z^{1+\epsilon}
}\] whenever \(\fitinline{(\ell_Z(X_1)+\ell_Z(X_2),\sum_iw_i)}\) is in (A4), where \(\fitinline{\ell_Z(X)=\log_Z\max(1,X)}\). The \(\fitinline{p_j}\) have the original annular supports inside \(\fitinline{[1/4,9/4]}\); \(\fitinline{\Psi}\) is Schwartz, supported above by a fixed cap and \(\fitinline{\Re\Psi\ge0}\). The mesh precedes the ambient finite slot type. The fixed arithmetic data, marked bump and radial cap precede the finite seminorm set \(\fitinline{\mathcal F}\), height degree, constant and onset, which precede live profiles, heights, lengths, widths, radial tests and retained rows. Finite row restrictions follow by positivity.

Both domains are downward closed: shorten the inverse exponent, either ordinary length, or omit whole slots. Inverse (A5) is stated for \(\fitinline{X=Z^r\ge1}\), not all subunit \(\fitinline{X>0}\). Ordinary (A6) already includes all positive shortened lengths; bounded nonempty scales retain the unit, and empty annuli vanish. On a fixed radius interval \(\fitinline{0<M_{\min}\le M\le M_{\max}}\), put \(\fitinline{U=Z^M}\), take mesh at most \(\fitinline{M_{\min}}\) times the unit mesh and source loss at most \(\fitinline{\epsilon/M_{\max}}\). This gives \(\fitinline{Z^{M+\epsilon}}\), uniformly over that interval.

For each fixed \(\fitinline{K}\), late controls may depend on \(\fitinline{K}\). No constant, degree or onset uniform in \(\fitinline{K=K(Z)\to\infty}\) is claimed. A small early mesh does not remove this restriction. Inverse profile-fixed uniformity and ordinary finite-seminorm uniformity are different.

\textbf{Theorem A (PR137).} Under the standing case \(\fitinline{\beta\ge51/100}\), set \[\fitmath{
 I=[1467/2000,311/400],\qquad
 g(r)=552429/10^7-(7227847/10^8)r. 
}\tag{A7}\] Then \(\fitinline{\mathrm{Anchor}(r,g(r))}\) holds for every \(\fitinline{r\in I}\). Every source call has a positive decrement \(\fitinline{d}\). Constants need not remain uniform as \(\fitinline{d\downarrow0}\). The same affine line is supplied natively on \(\fitinline{[311/400,311/400+1/1000]}\), for the finite-grid consumer transfer.

\hypertarget{the-one-stage-invariant}{%
\subsubsection{4. The one-stage invariant}\label{the-one-stage-invariant}}

Suppose \(\fitinline{\mathrm{Anchor}(R,z)}\) is proved. Put \(\fitinline{S_R=(1-R)/2+z}\), choose \(\fitinline{r\le R}\), and define \[\fitmath{
 H=r/R,\quad u=(1-H)/6,\quad
 Y(r)=Hz+(1/k-2)u,\quad S(r)=(1-r)/2+Y(r). 
}\tag{A8}\] The decisive identity is \[\fitmath{
 s:=S(r)-u=H S_R+u/k. 
}\tag{A9}\] After selecting width \(\fitinline{u}\), erasing the excess over the rescaled old anchor costs exactly \(\fitinline{k(s-HS_R)=u}\). This pays LOW rows. HIGH rows cost at most the same \(\fitinline{u}\), using earlier anchors. These costs occur on disjoint positive row sets and combine by their maximum. The stage proves \(\fitinline{\mathrm{Anchor}(r,Y(r)-D)}\), with \(\fitinline{D=1/25000000}\).

\hypertarget{high-rows-and-the-floored-detector-price}{%
\subsubsection{5. HIGH rows and the floored detector price}\label{high-rows-and-the-floored-detector-price}}

Let \(\fitinline{\nu\in[k,3/4]}\) be a buffered zero-bin exponent and \(\fitinline{e\in[0,\nu]}\) its lower slot-mean amplitude. Pay the actual ordinary pair, or inverse factor, pointwise at exponent \(\fitinline{\nu P}\), retaining the slot product inside the row sum; its amplitude price is \(\fitinline{eS}\). Replacing this by the cruder \(\fitinline{\nu S}\) loses the gain.

For a line \(\fitinline{l(w)=a_0-a_1w}\) below earlier proved gains, put \[\fitmath{
 \sigma(w)=\max(0,(1-w)/2+l(w)-10^{-7}),\quad
 w\in[2/3,79/100]. 
}\tag{A10}\] The same-zero detector and largest-amplitude whole-slot selection give a short count exponent \[\fitmath{
 C_{\rm short}=M-\nu Mw-e\min(S,M\sigma(w)). 
}\tag{A11}\] To see this, select whole slots of total width at most the cap. Their common lower product is \(\fitinline{U^{ez-O(\mathrm{mesh})}}\), the inverse or paired-ordinary witness gives \(\fitinline{U^{\nu w-O(\epsilon)}}\), and their marked moment is \(\fitinline{O(U^{1+\epsilon})}\). Divide the summed energy by the two lower bounds. Derivatives through order two and real shifts in \(\fitinline{[0,1]}\) stay inside the detector energy inequality. Both witnesses arise from the same actual zero.

The long-inverse count, at \(\fitinline{t=3w/2+10^{-6}}\), has exponent \[\fitmath{
 C_{\rm long}=M/6+(5/6-\nu)M(3w/2+10^{-6}). 
}\tag{A12}\] The supported-witness theorem permits an unbounded height index \(\fitinline{i}\) under the actual budget \[\fitmath{
 (3+(3i+2)Z^\tau)^4\le(Z^d)^\eta;
}\] this does not permit an unbounded number of prime slots.

For tiny sixth-free cores of norm at most \(\fitinline{Z^{2/25}}\), the count and pointwise price give excess \[\fitmath{
 E_{\rm tiny}=M/6+1/15+(3/4)P+S-M. 
}\tag{A13}\] For the other cores use \[\fitmath{
 E(\nu,e,w)=\nu P+eS+\max(C_{\rm short},C_{\rm long})-M. 
}\tag{A14}\] Every row is restored as \(\fitinline{u_0b^6}\), including its extra Euler mask and every cofactor; \(\fitinline{u_0,b}\) need not be coprime. The sixth-power cofactor sum restores radius \(\fitinline{M}\) with arbitrarily small divisor loss. It explains the \(\fitinline{M/6}\) and \(\fitinline{1/15}\) terms and cannot be omitted.

The HIGH price is the maximum of zero, (A13), and \(\fitinline{\sup_{\nu,e}\inf_w E}\). At stage \(\fitinline{n}\) the infimum is constrained by \[\fitmath{
 w\ge R_n+10^{-6}. 
}\tag{A15}\] The floor raises the price. Below the top \(\fitinline{311/400}\), every short witness must use an earlier stage \(\fitinline{j<n}\); above it the line is nonpositive and native sources suffice. An unfloored optimization does not prove this induction.

Here is the scalar price proof at \(\fitinline{(P,S,M)=(r,s,H)}\) of (A9). Write \(\fitinline{D_1=1/2-a_1\ge1/10}\) and \[\fitmath{
 G_\nu=\frac{r+s-H/2-Ha_0-u/\nu}{D_1},\qquad
 w_\nu=G_\nu/H+10^{-6}. 
}\tag{A16}\] Since \(\fitinline{l(R)\le z}\), \[\fitmath{
 D_1(G_\nu-r)=H(z-a_0+a_1R)+u(1/k-1/\nu)\ge0. 
}\tag{A17}\] Thus \(\fitinline{w_\nu\ge R+10^{-6}}\). It increases with \(\fitinline{\nu}\); its endpoint gates put it in \(\fitinline{[2/3,79/100]}\). Before paddings the short excess over \(\fitinline{H+u}\) is \[\fitmath{
 \nu(r-G_\nu)+e\{s-(H-G_\nu)/2-Ha_0+a_1G_\nu\}-u\le0. 
}\tag{A18}\] The braces equal \(\fitinline{G_\nu-r+u/\nu\ge0}\); replacing \(\fitinline{e\le\nu}\) by \(\fitinline{\nu}\) makes the expression zero. The width and capacity paddings are paid because \(\fitinline{a_1\le2/5}\) and \[\fitmath{
 (1/2+a_1)10^{-6}+10^{-7}\le10^{-6}.
}\] The positive part in (A10) only helps. Hence the short excess is at most \(\fitinline{u}\).

For the long branch put \[\fitmath{
 \mathcal G(\nu)=H+u-H/6-(5/6-\nu)(3G_\nu/2)-\nu(r+s).
}\] At \(\fitinline{\nu=k}\), \(\fitinline{G_k=Hw_b}\), \(\fitinline{\mathcal G(k)=HF_0}\), where \[\fitmath{
 w_b=(R/2+z-a_0)/D_1,\quad
 F_0=5/6-(5/4-3k/2)w_b-k(1+R)/2-kz. 
}\tag{A19}\] Its exact difference is \[\fitmath{
 \mathcal G(\nu)-\mathcal G(k)=(\nu-k)
 \{\tfrac32G_\nu-\tfrac32(5/6-k)u/(k\nu D_1)-(r+s)\}.
}\] It is nonnegative if \(\fitinline{r/2-s-\tfrac32(5/6-k)u/(k^2D_1)\ge0}\). The certificate supplies this with reserve \(\fitinline{1/10}\) and \(\fitinline{HF_0>3\cdot10^{-5}}\). Long-cutoff paddings cost at most \(\fitinline{1/400000}\), leaving excess at most \(\fitinline{u-11/400000}\). The tiny-core gate gives excess at most \(\fitinline{u-3/100}\). Thus the \textbf{floored} HIGH price is at most \(\fitinline{u}\). Monotonicity in free/marked widths and the cofactor radius comparison extend the ledger to the required shortened children.

\hypertarget{actual-inverse-selection-poisson-and-restoration}{%
\subsubsection{6. Actual inverse selection, Poisson and restoration}\label{actual-inverse-selection-poisson-and-restoration}}

The inverse step re-runs the actual initialization with injective prime marks. Overlapping canonical pools cannot first be coloured into arbitrary prime subsets and then be treated as canonical sources. Instead expand injective tuples; collision restoration afterwards recovers the full product, with a fixed finite combinatorial cost and small divisor loss. The initialized dictionary keeps the actual M\"{o}bius coefficients and natural zeros.

For a live-mark state with column width \(\fitinline{N}\) and label width \(\fitinline{V}\), put \(\fitinline{F=N+V}\), \(\fitinline{K=2F-1}\). The native canonical energy supplies the Gauss-row ball of exponent \(\fitinline{K}\), enlarged by a small slack. If live width is at least \(\fitinline{u}\), greedily select whole slots of width \(\fitinline{u'\in[u-O(\mathrm{mesh}),u]}\), one distinct prime per selected slot. For their prime product \(\fitinline{D}\), use the exact Gauss identity \[\fitmath{
 \gamma_2(U'D)=\gamma_2(U')\gamma_2(D)\chi_{U'}(D)^4. 
}\tag{A20}\] The injective tuple split moves \(\fitinline{D}\) into the row \(\fitinline{kD^4}\). The residual squarefree column avoids those selected primes; its other primes remain free even if they occur in the selected slots' pools. The zero \(\fitinline{\chi_{U'}(D)=0}\) at a common prime stays. The selected coefficient is \[\fitmath{
 \gamma_2(D)\Psi(D)\mathbf1_{(D,m)=1}
                  \chi_D(f)^4\chi_D(k)\,\mathrm{mark}_{\rm selected}(D).
}\] Cauchy and compact logarithmic separation leave one common residual row at column width \(\fitinline{N'=N-u'}\), Gauss radius \(\fitinline{K'=K+4u'}\). A second masked Poisson summation on that physical row yields a child radius \[\fitmath{
 H_c=2(N'+V)-K'+\text{slack}=1-6u'+\text{slack}. 
}\tag{A21}\]

For labels \(\fitinline{Nf\le Z^V}\) and nonnegative weights \(\fitinline{D_f\le D_{\max}}\), a bound \(\fitinline{B}\) for all physical inverse children gives \[\fitmath{
\begin{aligned}
 E_{\rm selected}\le{}&CZ^{K'-N'-V+\epsilon}D_{\max}
                        (1+|s_0|)^{2J}B\\
 &+CZ^{K'+\epsilon}D_{\max}(Z^{-V}\#\{f\})
       [Z^{-N'}\sum_U|\mathrm{mark}_{\rm inj}(U)|^2
                           |W(NU/Z^{N'})|^2].
\end{aligned} 
}\tag{A22}\] The second line is the already paid positive diagonal. The first line's child is the physical M\"{o}bius column \[\fitmath{
 \sum_A\mu(A)\psi(A)\mathbf1_{(A,m)=1}\chi_A(y)^{\pm1}
       \mathrm{mark}_{\rm residual}(A)\Phi(NA/X_c)/\sqrt{NA},
}\] summed over all nonzero \(\fitinline{Ny\le Z^{H_c}}\). It includes both orientations, every residual projector puncture, every shortened \(\fitinline{0<X_c\le X_{\max}}\), and every member and height of the fixed finite profile family. The ray \(\fitinline{\psi}\) belongs to a fixed finite ray-twist family.

For completeness, (A22) follows by expanding the squared canonical row, separating its common and coprime squarefree parts, applying the admitted masked Poisson identity to the row variable, and grouping the resulting Gauss and quadratic cross phases by the finite ray dictionary. The zero frequency is the displayed diagonal. On the child ball the grouped row is precisely the displayed M\"{o}bius column, so its positive sum is bounded by \(\fitinline{B}\), retaining all puncture terms and label weights. Outside the ball the radial Fourier argument is at least \(\fitinline{Z^\kappa/b^2}\). Choose decay order after the fixed height degree; lattice counting with a fixed polynomial envelope pays the complement. Divisor bounds and finite profile separation account for \(\fitinline{Z^\epsilon}\). The primitive Poisson and Gauss identities are the foundational inputs in section 1, rather than consequences of a theorem name or successful source import.

At zero slack the bookkeeping identities are \[\fitmath{
 K'-N'-V=K-F+5u,\qquad H=1-6u,\qquad
 (K-F+5u)+(H+u)=F. 
}\tag{A23}\] A child of size \(\fitinline{Z^{H+u+\epsilon}}\), rather than \(\fitinline{Z^H}\), suffices. This is precisely the excess paid by the LOW identity (A9) and HIGH ledger.

Restoring that child to the actual inverse is substantive. For a finite prime pool \(\fitinline{S}\), write \[\fitmath{
 C_S(X)=\sum_{A\subseteq S}\prod_{p\in A}c_p\,
             \Phi(\prod_{p\in A}Np/X),\quad
 c_p=-\psi(p)\chi_p(y)/\sqrt{Np}.
}\] For \(\fitinline{p\notin S}\), \[\fitmath{
 C_{S\cup\{p\}}(X)=C_S(X)+c_pC_S(X/Np). 
}\tag{A24}\] Iterating solves the avoided-prime column by a finite reciprocal Euler series of unrestricted columns. The remainder vanishes when \(\fitinline{X/(Np)^J<1/b}\), since \(\fitinline{\Phi}\) is supported in \(\fitinline{[a,b]\subset(0,\infty)}\) and all ideal norms are at least one. Required primes are restored by the corresponding difference in (A24). For overlapping pools group injective tuples by exact image before the recursion; do not factor that tuple sum as independent slot sums.

The unrestricted column is the actual inverse with profile \(\fitinline{\Phi(x)x^{-1/2}}\), since \(\fitinline{X^{-1/2}(NI/X)^{-1/2}=(NI)^{-1/2}}\). Mellin separation makes the profile independent of frozen tuple scales. Use a bump equal to one on \(\fitinline{[a',2^{2K+1}b']}\) and Schwartz Fourier weights for the finite child profiles and their conjugates. Their weighted integrals absorb auxiliary heights, leaving a polynomial in the original row and slot heights. Choose the source mesh first; restoration uses at most \(\fitinline{\min(m_{\rm child}H_{\rm lo},1)}\), independent of later classes and child profiles. Small column scales receive a bounded finite-support estimate, not an assertion that every such scale is empty.

The Euler recursion retains its damping: \(\fitinline{|c_p|\le(Np)^{-1/2}}\), so its \(\fitinline{j}\)-th term carries \(\fitinline{(Np)^{-j/2}}\). It is not an undamped count of prime powers. Finite required-prime restorations use this factor; the potentially large puncture support is treated by the following norm-power bound.

The puncture/mask restoration costs satisfy, for every fixed \(\fitinline{K}\) and \(\fitinline{\eta>0}\), \[\fitmath{
 \mathrm{cost}\le C_{K,\eta}(Nt\,|m|^2)^\eta. 
}\tag{A25}\] Indeed the product norm of distinct primes in \(\fitinline{t\cup\{p:p\mid m\}}\) is at most \(\fitinline{Nt|m|^2}\). A fixed exponential in their number is absorbed into an arbitrary norm power; the additional fixed-order polynomial combinatorics is absorbed likewise. There is no hypothesis that the full small-prime primorial has polynomial norm. Negative orientations use conjugated profiles with the same masks.

For parent caps \(\fitinline{P_{\max},S_{\max}}\) and slack \(\fitinline{\kappa>0}\), the literal normalized child is \[\fitmath{
 \left(\frac{P_{\max}+2\kappa}{1-6u-\kappa},
       \frac{S_{\max}-u+3\kappa}{1-6u-\kappa}\right), 
}\tag{A26}\] with excess \(\fitinline{u/(1-6u-\kappa)}\). Its gates are \[\fitmath{
 P_{\max}+S_{\max}+4u+10\kappa<1,\qquad6u+2\kappa<1. 
}\tag{A27}\] States with insufficient live width use the native canonical-energy estimate; the other states use (A20)--(A26). Their positive union closes the initialization oracle.

At a stage put \(\fitinline{P_{\max}=r-d}\), \(\fitinline{S_{\max}=S(r)-D-d}\), \(\fitinline{\kappa=d/8}\). The old anchor with decrement \(\fitinline{d/2}\) covers the child's free cap. For the marked cap its LOW erasure ledger is \[\fitmath{
 k(S_{\rm child}-S_{\rm old})H_c
 \le u-kD-kd/2+k(3+S_R+d/2)\kappa\le u. 
}\tag{A28}\] The stage bounds make the last inequality valid; HIGH children have the same \(\fitinline{H+u}\) budget from section 5. Positive reserves absorb all mesh, profile and divisor losses. Principal inverse rows use their principal-character bound; they and fixed-inducing rows were never silently deleted. All roots and units remain in the all-row moment. This proves the inverse half of the new anchor.

\hypertarget{ordinary-rectangle-and-reflected-reference-lane}{%
\subsubsection{7. Ordinary rectangle and reflected reference lane}\label{ordinary-rectangle-and-reflected-reference-lane}}

The target is \(\fitinline{G(r-d,0,S(r)-D-d)}\), a rectangle with marked slope zero. Substituting \(\fitinline{L=0}\) into the old centered-reference proof is insufficient: its reference interval is empty at \(\fitinline{P+2S>1}\). PR137 forks the assembly and supplies a new reference lane.

The inherited assembly still performs the actual common-divisor factorization, complete collision and cube/ray restoration, first and second divisors, zero frequency, fixed-inducing diagonal and tails. The following calculation pays every packed first-common side; the stage-point identity (A9) alone would not suffice.

It is enough to treat \(\fitinline{0<d\le1/1000}\); larger endpoint decrements follow by rectangle restriction. Distinguish the stage target \(\fitinline{u_0=(1-r/R)/6}\), \(\fitinline{H_0=r/R=1-6u_0}\), from a side's actual selected width \(\fitinline{u_{\rm s}}\). Let \(\fitinline{\Sigma}\) be its original marked width, \(\fitinline{f}\) its frozen marked width, \(\fitinline{c_g}\) its extracted common-norm width, \(\fitinline{P_C}\) its ordinary width, and \(\fitinline{s_{\rm s}}\) its retained marked width. The exact side geometry gives \[\fitmath{
\begin{gathered}
 \Sigma\le S(r)-D-d,\quad 0\le f\le c_g,\quad
 0\le P_C,\quad P_C+(c_g-f)\le r-d+\tau,\\
 w=\Sigma-f,\quad
 \min(u_0,w)\le u_{\rm s}\le\min(u_0,w)+\mu,\quad
 0\le s_{\rm s}\le w-u_{\rm s},\quad t\ge-\tau .
\end{gathered}
}\] Here \(\fitinline{\mu}\) is the slot mesh and \(\fitinline{\tau}\) the frequency tolerance. Choose both at most \[\fitmath{
 m=\min\{d/64,\epsilon/[256(L_{\rm high}+1)]\},
}\] where \(\fitinline{L_{\rm high}\ge0}\) is the HIGH price's Lipschitz constant. Set \(\fitinline{\xi=\tau+6\mu}\). Then \(\fitinline{0\le\xi\le7d/64}\), and the packed row exponent satisfies \[\fitmath{
 H_{\rm s}=1+t-6u_{\rm s}\ge H_0-\xi.
}\] The side's upper frequency bound \(\fitinline{t\le1+6u_{\rm s}+\tau}\) also gives \(\fitinline{H_{\rm s}\le2+\tau\le3}\). The stage bound \(\fitinline{u_0\le1/10000}\) gives \(\fitinline{H_0-\xi>1/2}\), so the packed radius lies in the assembly's required interval.

Soften the old rectangle by \[\fitmath{
 d_1=\min(D,d)/4,\quad B_1=R-d_1,\quad
 c_1=S_R-d_1,\quad \ell=d/4>0.
}\] The source \(\fitinline{G(B_1,\ell,c_1)}\) is a restriction of the old rectangle \(\fitinline{G(B_1,0,c_1)}\): adding the nonnegative term \(\fitinline{\ell s}\) only shrinks its admission domain. The stage bounds give \(\fitinline{B_1\ge0}\) and \(\fitinline{0\le c_1\le1}\). Its LOW price is \(\fitinline{k}\) times the clipped deletion need: erase whole slots until both \(\fitinline{s\le c_1H}\) and \(\fitinline{P+\ell s\le B_1H}\), with the native zero-slot bound as the delete-everything fallback.

First suppose \(\fitinline{w\ge u_0}\). Then \(\fitinline{u_{\rm s}\ge u_0}\) and \[\fitmath{
 P_C\le r-d+\tau,\qquad
 s_{\rm s}\le\bar s:=S(r)-D-d-u_0.
}\] The deletion need increases with \(\fitinline{P,s}\), decreases with the radius for \(\fitinline{B_1,c_1\ge0}\), and is bounded above by the positive part of its two defects. Therefore \[\fitmath{
\begin{aligned}
 \mathrm{need}(P_C,s_{\rm s},H_{\rm s})
 \le\max\left(0,\max\left\{
  \bar s-c_1(H_0-\xi),\
  \frac{r-d+\tau+\ell\bar s-B_1(H_0-\xi)}{\ell}
 \right\}\right).
\end{aligned}
}\] For the cap defect, identity (A9) gives the exact equality \[\fitmath{
 k[\bar s-c_1(H_0-\xi)]
 =u_0+k(H_0d_1-D-d+c_1\xi)\le u_0.
}\] Indeed \(\fitinline{H_0\le1}\), \(\fitinline{d_1\le d/4}\), \(\fitinline{c_1\le1}\), and \(\fitinline{\xi\le7d/64}\), so the bracket is at most \(\fitinline{-D-41d/64<0}\). For the budget defect, \(\fitinline{RH_0=r}\) gives \[\fitmath{
\begin{aligned}
 r-d+\tau+\ell\bar s-B_1(H_0-\xi)
 &=-d+\tau+\ell\bar s+d_1H_0+R\xi-d_1\xi\\
 &\le-d+d/64+d/24+d/4+7d/64\\
 &=-7d/12<0.
\end{aligned}
}\] Here \(\fitinline{\bar s\le S(r)<137/1000<1/6}\), hence \(\fitinline{\ell\bar s\le d/24}\); also \(\fitinline{R<1}\) and \(\fitinline{d_1\xi\ge0}\). Division by \(\fitinline{\ell>0}\) preserves negativity. The second defect therefore contributes nothing, and the first one's positive part satisfies \[\fitmath{
 k\,\mathrm{need}(P_C,s_{\rm s},H_{\rm s})\le u_0\le u_{\rm s},
}\] including the case where that first defect is nonpositive. This is the complete packed-side LOW payment.

For the HIGH part, \(\fitinline{s_{\rm s}\le S(r)-u_0}\). Monotonicity in \(\fitinline{P,s}\), antitonicity in the radius and the same Lipschitz bounds compare the side with the stage point: \[\fitmath{
\begin{aligned}
 \mathrm{high}(w,P_C,s_{\rm s},H_{\rm s})
 &\le\mathrm{high}(w,r,S(r)-u_0,H_0)
                  +L_{\rm high}(\tau+\xi)\\
 &\le u_0+\epsilon/32\le u_{\rm s}+\epsilon/32.
\end{aligned}
}\] The error bound follows from \(\fitinline{\tau+\xi\le8m}\) and \(\fitinline{L_{\rm high}m\le\epsilon/256}\). Thus the maximum of the LOW and HIGH prices is at most \(\fitinline{u_{\rm s}+\epsilon/16}\), as required.

If \(\fitinline{w<u_0}\), the greedy bounds instead give \(\fitinline{u_{\rm s}\ge w}\), hence \(\fitinline{s_{\rm s}=0}\). The clipped deletion need is zero, and the HIGH price at \(\fitinline{(r,0,H_0)}\) is zero; the same comparison costs at most \(\fitinline{\epsilon/32}\). At the analytic leaf this is the erase-all-slots/native-zero-slot payment, not an application of the old marked source outside its ordinary cap. These two cases exhaust all sides above the positive live-width floor \(\fitinline{\epsilon/12}\). Below that floor, use no selection and the native all-slot price: its coefficient \(\fitinline{3/4}\) makes the cost at most \(\fitinline{(3/4)(\epsilon/12)=\epsilon/16}\), including zero marked width.

This is the calculation in the pinned Price/Gates.lean stage\_ledger, with its \(\fitinline{\beta=d}\), tl \(\fitinline{=\tau}\), mesh \(\fitinline{=\mu}\), \(\fitinline{x=\xi}\), and \(\fitinline{u_0}\) distinguished from the actual side selection. It preserves the source's softened face, clipped fallback and full frequency tolerance.

The new reference parameters are \[\fitmath{
 l_{\rm ref}=7/50,\quad \alpha=3/5,\quad P_{\rm ref}=2/3.
}\] Use a separately proved parent \(\fitinline{G(99/100,49/20,7/50)}\). Forward \(\fitinline{K_3}\) packet expansion dominates the complete reference rectangle by native leaves at shortened lengths, with one factor at most \(\fitinline{Z^{l_{\rm ref}}}\). If the long factor has exponent at most \(\fitinline{1/2}\), the parent pays it directly. Otherwise naturally reflect it; the dual exponent is at most \(\fitinline{1/2+\delta+\xi}\). For stage cap \(\fitinline{c<137/1000}\), \[\fitmath{
 \frac7{50}+\frac12+\frac{49}{20}c+\frac1{100}<\frac{99}{100},
 \quad c\le\frac7{50},\quad\frac7{50}+c\le1. 
}\tag{A29}\] Choose reflection paddings below this reserve. Retain natural conductor zeros and the fixed-inducing predicate. Choose finite profile controls before Fourier heights, estimate original tails with a fixed envelope, and preserve the height factor \(\fitinline{1+2|t|+V}\). This pays every native leaf; positivity and the \(\fitinline{K_3}\) kernel bound restore the whole reference rectangle.

The corrected reference contract chooses its mesh before ambient data, but degree, controls and constant \textbf{after} the fixed radial cap, characters and internal ideal \(\fitinline{Q_0}\). An older stronger ordering put these constants before that data and omitted slot-frequency dependence; this route does not prove that stronger statement. The actual consumer fixes the data early and so needs only the corrected contract.

Finally the fixed preceding parent \(\fitinline{G(1-10^{-12},2+10^{-11},3/20)}\) pays inheritance below crossing \(\fitinline{(B_0-(r-d))/L_0}\); above it use the new reference and priced leaf. The completed common-divisor/second-divisor/diagonal assembly from section 1 remains an explicit dependency. The finite stage gates verify its exceptional, root-reference and second-diagonal reserves. Consequently the output is the full ordinary source (A6) on (A4), including all shortened lengths and complementary sectors, rather than a bound only for favorable packed rows.

\hypertarget{finite-chain-and-acyclicity}{%
\subsubsection{8. Finite chain and acyclicity}\label{finite-chain-and-acyclicity}}

Start from \(\fitinline{(R_0,z_0,a_{0,0},a_{1,0})=(311/400,0,0,0)}\). The inverse decremented base rectangle satisfies \(\fitinline{P+2S<1}\), \(\fitinline{2P+8S<3}\); the native theorem supplies it even for arbitrary bounded slot coefficients. The ordinary base is covered by reserve-parametric sources approaching \(\fitinline{(B,L,c)=(1,2,3/20)}\). Cycle068's complete family already supplies these parents; the accepted PR137 proof uses its inherited PR103/PR129 sources.

For \(\fitinline{n=0,\ldots,175}\), put \(\fitinline{R_{n+1}=R_n-1/4000}\). Compute \(\fitinline{w_b}\) by (A19), then \[\fitmath{
 T_n=\frac{5w_b/4-5/6}{3w_b/2-(1+R_n)/2-z_n},\quad
 k_n=10^{-6}\lceil10^6(T_n+1/5000)\rceil. 
}\tag{A30}\] Use (A8) with \(\fitinline{k=k_n}\), and set \[\fitmath{
 z_{n+1}=10^{-10}\lfloor10^{10}(Y_n(R_{n+1})-D)\rfloor. 
}\tag{A31}\] The next line is the affine function through \(\fitinline{Y_n(r)-D-2\cdot10^{-7}}\) at the two stage endpoints. The exact certificate checks it is a strict \(\fitinline{10^{-7}}\) minorant on every earlier stage, is negative at \(\fitinline{311/400}\), has slope in \(\fitinline{(0,2/5)}\), and is below the next entry gain. The initial line is zero.

For audit, the fourteen strictly positive endpoint quantities are the following, with \(\fitinline{s=S-u}\), \(\fitinline{H=r/R}\), \(\fitinline{u=(1-H)/6}\), \(\fitinline{w_{\rm lo}=w_k}\), \(\fitinline{w_{\rm hi}=w_{3/4}}\):

\begin{longtable}[]{@{}
  >{\raggedright\arraybackslash}p{(\columnwidth - 2\tabcolsep) * \real{0.5000}}
  >{\raggedright\arraybackslash}p{(\columnwidth - 2\tabcolsep) * \real{0.5000}}@{}}
\toprule\noalign{}
\begin{minipage}[b]{\linewidth}\raggedright
Gate
\end{minipage} & \begin{minipage}[b]{\linewidth}\raggedright
Quantity required positive
\end{minipage} \\
\midrule\noalign{}
\endhead
\bottomrule\noalign{}
\endlastfoot
1 & \(\fitinline{HF_0-3/100000}\) \\
2 & \(\fitinline{137/1000-S}\) \\
3 & \(\fitinline{14/100-S/H}\) \\
4 & \(\fitinline{9/10-r-S-4u}\) \\
5 & \(\fitinline{91/100-r-S-4u}\) \\
6 & \(\fitinline{5/6+21/250-r-S-3u-1/100}\) \\
7 & \(\fitinline{3H-2r-8s-1/5}\) \\
8 & \(\fitinline{r-1/2-Y+u-1/5}\) \\
9 & \(\fitinline{r/2-s-\tfrac32(5/6-k)u/(k^2(1/2-a_1))-1/10}\) \\
10 & \(\fitinline{H+u-H/6-1/15-(3/4)r-s-3/100}\) \\
11 & \(\fitinline{w_{\rm lo}-2/3}\) \\
12 & \(\fitinline{79/100-w_{\rm hi}}\) \\
13 & \(\fitinline{1-(3/4)w_{\rm hi}-(3/4)\max(0,(1-w_{\rm hi})/2+a_0-a_1w_{\rm hi})-1/6}\) \\
14 & \(\fitinline{79/100-r-(2+10^{-11})u}\) \\
\end{longtable}

The certificate also checks \(\fitinline{R\le w_b}\), \(\fitinline{45457/125000\le k<1/2}\), \(\fitinline{k>9/25}\), and nonnegative next gain. Gates transfer through each interval by affine dependence, multiplying by positive \(\fitinline{H}\), and treating the two affine branches of gate 13 separately. The own line and return line are affine, so endpoint minorants imply full-interval minorants. The independent rational reconstruction checks these endpoint facts and all earlier-line inequalities; it does not certify the analytic foundations.

Strong induction now proves every whole stage. Its entry anchor follows from the preceding stage and downward rounding (A31). Short HIGH witnesses below the top use strictly earlier stages by (A15); above the top the nonpositive line permits native sources. Sections 6 and 7 prove both new halves. Thus no stage assumes its own moment and no limiting source is invoked. The 176 intervals cover \(\fitinline{I}\), since \(\fitinline{311/400-176/4000=1467/2000}\). The return line (A7) is below each \(\fitinline{Y_n-D}\), with the stated positive pad. Rectangle restriction gives Theorem A.

This is the fixed-decrement PR137 construction. Removing its fixed \(\fitinline{D}\), or replacing the fixed \(\fitinline{10^{-6}}\) cutoff padding by a variable debit, requires the separate PR141 analytic extension; it is not a consequence of sending \(\fitinline{D}\) to zero in this finite certificate.

\hypertarget{consumer-transfer-strict-gain-and-retained-suppliers}{%
\subsubsection{9. Consumer transfer, strict gain and retained suppliers}\label{consumer-transfer-strict-gain-and-retained-suppliers}}

Compact logarithmic Fourier separation converts inverse reference profiles to detector inverse profiles with both orientations, all log-derivatives \(\fitinline{n\le2}\), shifts \(\fitinline{s\in[0,1]}\), and heights \(\fitinline{|t|\le h}\), with \(\fitinline{C(1+h)^J}\). Its degree may depend on the fixed row data. The estimate \[\fitmath{
 (1+|\theta|+h_v)^{2J}\le(1+h_v)^{2J}(1+|\theta|)^{2J}
}\] lets Schwartz Fourier integrals absorb auxiliary heights. The constant does not acquire arbitrary dependence on a live height. The physical slot is the conjugate direct slot at the physical real part \(\fitinline{17/50}\): conjugating its annular shift gives real exponent \(\fitinline{-33/50}\), with the exact ray-character and zero masks. No substitute coefficient vector occurs.

For a detector fiber at scale \(\fitinline{U}\), the new inverse inlet is \(\fitinline{r\in I}\) and selected marked width at most \(\fitinline{(1-r)/2+g(r)-\nu_0}\). The ordinary inlet is \textbf{doubled} per-factor width \(\fitinline{2m\in I}\), cap evaluated at \(\fitinline{2m}\), and derivative tests \(\fitinline{j+k\le2}\). All original row subsets, profiles, natural zeros, real shifts and height tests remain.

To make constants uniform over live \(\fitinline{r}\), use a finite grid just to the right of each requested width, with a smaller common decrement; take the minimum mesh and maximum/union of degrees and controls over that grid. The native extension above \(\fitinline{311/400}\) covers the top cells. The consumer's fixed positive slack pays the capacity loss. This gives one mesh before weighted-source data and polynomial-height constants before scale. Choosing unrelated constants separately for every real \(\fitinline{r}\) would not prove this inlet.

There is a strict source enlargement. At \(\fitinline{r=3/4}\), \[\fitmath{
 g(r)=413619/400000000>0.
}\] With \(\fitinline{d=1/100000}\), its admitted corner is \[\fitmath{
 P=74999/100000,\quad S=50409619/400000000,\quad
 P+2S-1=407619/200000000>0. 
}\tag{A32}\] It is outside every Cycle068 complete source \(\fitinline{G(B,L,c_s)}\), \(\fitinline{B<1}\), \(\fitinline{L>\max(1,2B)}\): old admission would imply \(\fitinline{P<B(1-2S)<1-2S}\), contradicting (A32). It also violates the first native inverse gate, while the second has reserve \(\fitinline{24591381/50000000>0}\). This is a genuine supplier improvement outside the union of the old complete family, regardless of whether a particular endpoint certificate uses that corner.

The positive affine gain holds only for \(\fitinline{r<5524290/7227847}\); the line is negative near the top, where native inputs already suffice. The immediate new consumer is the selected-mark capacity in the same-zero detector, not an arbitrary larger sloping global source without a domain-cover proof. The theorem does not pay unrelated residual inverse bilinear targets, admit arbitrary prime subsets or row-dependent coefficients, remove moving presentation zeros, allow a live growing slot number, or make boundary constants uniform.

The anchor theorem changes source admission, not the physical return's prime-shift damping. Any return using it must retain the original shift weights, root restoration, intrinsic floor, tails and conductor sectors; their payment is not implied by (A32).

Keep the older arbitrary-coefficient native inverse theorem, inverse-specific stronger special moments, coherent inverse products, zero-width complements and useful source geometries. The new source does not dominate all of them. This section is ready for source/anchor integration; the count and physical-return lane must verify its own guards before a new endpoint can be reported.

Attribution: Sil's PR137 supplies this induction. The inherited OpenAI, Hailey Collet, ArviFoxy, Joseph Bebel, Joe Richey and Jacob Sussman contributions and notices remain as recorded in the submission metadata and licence files. The repository's Mathematical Writing for AI Agents guidance was used with its existing attribution preserved in that file.

\hypertarget{appendix-e}{}

\hypertarget{appendix-e.-continuum-marked-anchors-the-averaged-floor-and-the-rational-physical-return}{%
\subsection{Appendix E. Continuum marked anchors, the averaged floor, and the rational physical return}\label{appendix-e.-continuum-marked-anchors-the-averaged-floor-and-the-rational-physical-return}}

This appendix preserves the complete PR141 rational stage. Its local statement about the exact contact marks that stage's scope; Appendix G proves the subsequent exact-contact theorem using the same full foundations.

This is mathematical Appendix E of the current ordinary proof. It uses the retained \S{}\S{}2--10 and embedded Appendices A--C. Those sections and appendices remain the full arithmetic, theta, selected-norm, source-assembly, witness, Poisson, complementary-conductor and Mellin foundations. The coupled anchor theorem reconstructed in the parallel anchor note is the additional analytic inlet isolated below. The full working proof since OpenAI remains available, and a properly derived alternative to any earlier step is permitted. Nothing here changes the separately paused canonical controller.

Sil contributed the central marked-anchor method in PR137. Joseph Bebel's accepted PR141 extends it to the continuum supplier and supplies the averaged floor and exact return. OpenAI, Hailey Collet, ArviFoxy, Joe Richey, Jacob Sussman and the other inherited contributors retain their credit and notices. The exact intake revision is 03a1e5ad\allowbreak{}0cab7779\allowbreak{}69313090\allowbreak{}ee78b35d\allowbreak{}232ef355. SOURCE\_USE.json binds the files, declarations and hashes. Mathematical writing follows the retained Mathematical Writing for AI Agents, including its Sepia/CWR attribution. This is a mathematical reconstruction, not a new Lean build or axiom audit.

\hypertarget{admitted-theorem-and-exact-analytic-boundary}{%
\subsubsection{1. Admitted theorem and exact analytic boundary}\label{admitted-theorem-and-exact-analytic-boundary}}

Write \(\fitinline{t=N/D}\), where \[\fitmath{
N=9908119238306612133369374975000622988619121777878229799765263366946606045502455625911569003364955967645745531862430068266267717,
}\] \[\fitmath{
D=11345012023616779600791582323373494582287250283647959965785871011074347975117040113347744503437234376916434871293982689424818700.
}\] The accepted theorem is: for every positive integer \(\fitinline{q}\), every ordinary complex Dirichlet character \(\fitinline{\chi}\) modulo \(\fitinline{q}\), and every \(\fitinline{s\in\mathbf C}\), \[\fitmath{
\Re s>t,\qquad \neg(\chi\text{ principal and }s=1)
\quad\Longrightarrow\quad L(s,\chi)\ne0.                 
}\tag{1}\] Thus \(\fitinline{c=1-t=.126654143893281423198815203934927643\ldots}\), an increase over the retained \(\fitinline{c_*=(75+\sqrt{1081})/852}\). The assertion is at all heights, with each modulus fixed in its analytic estimates. It asserts neither the boundary line nor constants uniform in a growing modulus.

Use \(\fitinline{K=\mathbf Q(\sqrt{-3})}\), natural presented Hecke characters, all their Euler zeros, and the zero supremum \(\fitinline{\beta}\) of the retained proof. In particular \[\fitmath{
\beta\le7/8.                                             
}\tag{2}\] The contrary case \(\fitinline{\beta>t>.51}\) supplies the reserve sources already constructed from the original inputs. The additional analytic theorem needed below is the \textbf{coupled anchor extension}, with inverse and ordinary contracts both proved from the same earlier paired anchors. It includes the full original inverse row sum, fixed-inducing inverse rows, the ordinary reference lane after its physical reflection, all HIGH classes, common masks and polynomial height control. A result only about selected subrows or disjoint pools is insufficient. The parallel anchor note reconstructs this inlet. Reading a formal declaration or checking the scalar calculations below does not prove its analytic content.

A normalized canonical mark is \[\fitmath{
Q_i(u)=P_i^{-1/2}\sum_{\mathfrak p\in\mathcal P_i}
 b_i(\mathfrak p)c_u(\mathfrak p),\qquad P_i=Z^{\ell_i}.    
}\tag{3}\] The pool is the whole identity-ray prime pool with \(\fitinline{P_i\le N\mathfrak p\le2P_i}\). Its ray modulus, pool subgroup, relative twist group and smooth marked profile are fixed before the scale. A twist subgroup may be smaller than the pool subgroup. The inverse canonical class uses one common relative ray character across its slots; the ordinary initial source allows its original independent relative characters in the different slots. Coefficients are the actual relative ray character times the row base character and common annular profile at real part \(\fitinline{33/50}\), with independent imaginary heights bounded by \(\fitinline{h_v}\), and \(\fitinline{|b_i|\le1}\). Natural zeros remain. Whole slots may be selected; whole pools may not be replaced by arbitrary scale-dependent prime subsets. Inverse pools need not be disjoint.

Let \(\fitinline{\mathcal A(R,z)}\) mean that, for every decrement \(\fitinline{d>0}\), both literal row-sized estimates hold on \[\fitmath{
0\le a\le R-d,\qquad
0\le S=\sum_i\ell_i\le(1-R)/2+z-d.                         
}\tag{4}\] The inverse estimate is \[\fitmath{
\sum_{0\ne u,\,Nu\le Z}
 \left|M_u(Z^a;\Phi,\vartheta)\prod_iQ_i(u)\right|^2
 \le C Z^{1+\epsilon}(1+|\vartheta|+h_v)^{2j},              
}\tag{5}\] where \(\fitinline{M_u}\) is the actual inverse Hecke polynomial with the original child-log profile. Every finite subset of the ball is allowed. The ordinary estimate is the original positive initial energy \[\fitmath{
\sum_{u:\mathrm{initialKeep}(u)}
 |A_{1,u}(Z^{a_1})A_{2,u}(Z^{a_2})|^2
 \prod_i|Q_i(u)|^2\Re\Phi(Nu/Z)
 \le C\mathcal D(\Phi)[\mathcal S_{\mathcal F}(p_1)
 \mathcal S_{\mathcal F}(p_2)]^2
 Z^{1+\epsilon}(1+|t_1|+|t_2|+V)^J,
}\tag{6}\] where \(\fitinline{a_1,a_2\ge0}\), \(\fitinline{a_1+a_2\le R-d}\), and the ordinary profiles \(\fitinline{p_1,p_2}\) and their heights \(\fitinline{t_1,t_2}\) are independent. Here \(\fitinline{V}\) bounds the marked and common whole-product heights, and \(\fitinline{J}\) is the finite ordinary height degree. The finite seminorm controls are those of Appendix D (A6). Its fixed-inducing predicate is the original ordinary predicate. Equations (5)--(6) do not assert a moment for \(\fitinline{M_uA_u}\).

For the inverse estimate (5), the order of choices is \[\fitmath{
\forall d>0\ \forall\epsilon>0\ \exists\mu>0\
\forall K\ \forall\text{fixed admissible class/profile/row data}\
\exists j,C,Z_0\ \forall Z\ge Z_0.                         
}\tag{7}\] Here its radial profile is fixed before the height degree and constant. For the ordinary estimate (6), after the early mesh, fixed slot type, marked profile and fixed arithmetic/ray data, one chooses a finite set of seminorm controls, a height degree and a constant \textbf{before} the live two ordinary profiles and row profile. Their support and smoothness are the original ones, and their seminorm factors are paid explicitly. The bound is therefore uniform in the reflected and Fourier profiles used later; pointwise constants chosen anew for each such profile would not suffice. There are at most \(\fitinline{K}\) marks and each \(\fitinline{0\le\ell_i\le\mu}\). Both early meshes precede \(\fitinline{K}\), profiles and arithmetic data. Live permitted heights are paid polynomially. Constants may depend on fixed \(\fitinline{K}\) and outer data; this is not uniform for a number of prime factors growing with \(\fitinline{Z}\). Downward shortening in both coordinates is part of the contract.

\hypertarget{finite-construction-of-the-continuum-supplier}{%
\subsubsection{2. Finite construction of the continuum supplier}\label{finite-construction-of-the-continuum-supplier}}

For \(\fitinline{r>2/3}\), put \[\fitmath{
q(r)=[r^4(r-2/3)]^{1/5},\quad P(r)=11r/16-7/24,
}\] \[\fitmath{
A(r)=\frac4{5r}+\frac1{5(r-2/3)},\quad
b(r)=\frac{3r-7/3}{r(15r-10)},\quad F(r,z)=A(r)z+b(r).      
}\tag{8}\] Fifth roots are positive real roots. Differentiation gives \[\fitmath{
q'=Aq,\quad P'=F(r,P(r)),\quad
q''=-\frac{16q}{225r^2(r-2/3)^2}<0.                         
}\tag{9}\] Fix \(\fitinline{0<\lambda<1}\), \(\fitinline{3/4\le T<7/9}\), and define \[\fitmath{
Z_T(r)=P(r)-\frac{P(T)}{q(T)}q(r),\quad
W(r)=\lambda Z_T(r),\quad v(r)=W'(r).                      
}\tag{10}\] Then \(\fitinline{W(T)=0}\), \(\fitinline{W}\) is convex, and \[\fitmath{
v(r)-F(r,W(r))=(\lambda-1)b(r)>0
\quad(3/4\le r\le T).                                    
}\tag{11}\] Both strict choices matter: \(\fitinline{b<0}\) only below \(\fitinline{7/9}\), and shrinking the curve creates its positive differential reserve. Compactness supplies \(\fitinline{0<\gamma\le1/100}\) for which this reserve is at least \(\fitinline{3\gamma}\) everywhere. Uniform continuity of \(\fitinline{v}\) supplies \(\fitinline{\eta>0}\) such that \[\fitmath{
W(r)\le W(R)+(v(R)-\gamma)(r-R)
\quad(3/4\le r\le R\le T,\ R-r\le\eta).                    
}\tag{12}\] Indeed, on such an interval the derivative of the difference between \(\fitinline{W}\) and the affine expression is nonnegative, and the difference is zero at \(\fitinline{R}\).

The exact scalar output of the coupled analytic step is \[\fitmath{
H=r/R,\quad u=(1-H)/6,\quad
Y(r)=Hz+(1/k-2)u,\quad S(r)=(1-r)/2+Y(r),\quad s=S-u.       
}\tag{13}\] At the earlier anchor \(\fitinline{z=W(R)}\), \(\fitinline{v=v(R)}\), choose \[\fitmath{
k=\{2+6(z-Rv+R\gamma)\}^{-1}.                             
}\tag{14}\] The affine output has \(\fitinline{Y(R)=z}\) and derivative \(\fitinline{v-\gamma}\), so (12) is exactly \(\fitinline{W(r)\le Y(r)}\). Choose one fixed positive tolerance \[\fitmath{
0<q_{\rm tol}\le\min(10^{-12},\gamma/1000),\quad
w_{\rm lo}=R+q_{\rm tol},\quad
a_0=z-Rv-q_{\rm tol}/5,\quad a_1=-v.                       
}\tag{15}\] The elementary strict-curve bounds are \(\fitinline{0\le z\le11/10000}\), \(\fitinline{-4/45\le v\le0}\), and \(\fitinline{45457/125000\le k<1/2}\). Reduce \(\fitinline{\eta\le1/10000}\). They give the retained stage gates \[\fitmath{
S<.137,\quad r+S+4u<.9,\quad
r+S+3u+.01<5/6+21/250,                                   
}\tag{16}\] as well as \(\fitinline{u\le10^{-4}}\), \(\fitinline{(1-R)/2+z\le.137}\), \(\fitinline{R\in[7/10,7/9]}\), and \(\fitinline{r\in[1467/2000,R]}\). The tiny HIGH price is at most \(\fitinline{u-3/100}\). Each non-tiny HIGH class has an allowed cutoff \(\fitinline{w\ge R+q_{\rm tol}}\), total source width at most six, and its recentered detector price at most \(\fitinline{u}\). These are the scalar class inequalities of FamilyGates.curve\_stage\_gates; their analytic consumer is PairedStage.paired\_anchor\_of\_earlier. Convexity gives \[\fitmath{
a_0-a_1x<W(x)\quad(R<x\le T),\qquad a_0-a_1T\le0.           
}\tag{17}\] Thus HIGH uses only previously proved paired anchors. Variable positive tolerance and rectangle decrements pay their errors; no fixed gain loss is subtracted at each stage. A fixed-decrement stage theorem would not suffice. Part A of FINITE\_SOURCE\_APPENDIX.md derives the full class bound: the unpadded reserve is at least \(\fitinline{\gamma/3}\), the actual pad costs at most \(\fitinline{q_{\rm tol}}\), and \(\fitinline{HF_0>30q_{\rm tol}}\). The short pad debit cancels against its witness-shift gain, while the long pad costs at most \(\fitinline{5q_{\rm tol}/2}\). Its exact recentering and the zero fixed-stage-loss analytic adapter are given there.

Choose an integer \(\fitinline{m>(T-3/4)/\eta}\) and divide \(\fitinline{[3/4,T]}\) into \(\fitinline{m}\) intervals. The native inverse rectangle and reserve ordinary source give \(\fitinline{\mathcal A(T,0)}\). Assume the anchors on the already treated interval \(\fitinline{[R,T]}\) have been proved. Equations (14)--(17) and the coupled analytic step give \(\fitinline{\mathcal A(r,Y(r))}\) throughout the next interval, and (12) restricts this to \(\fitinline{\mathcal A(r,W(r))}\). Finite induction reaches \(\fitinline{3/4}\). No stage uses an unproved later anchor or an unproved point in its current interval.

The ideal gain is \[\fitmath{
I(r)=P(r)-\frac5{16}[7r^4(r-2/3)]^{1/5},
\qquad3/4\le r<7/9.                                     
}\tag{18}\] For fixed \(\fitinline{r,d>0}\) with \(\fitinline{r<7/9}\), continuity at \(\fitinline{(\lambda,T)=(1,7/9)}\) supplies one strict finite curve with \(\fitinline{W(r)>I(r)-d/2}\). Use its anchor with decrement \(\fitinline{d/2}\), then restrict its rectangle to (4) with gain \(\fitinline{I(r)}\) and decrement \(\fitinline{d}\). Therefore \[\fitmath{
\mathcal A(r,I(r))\quad(3/4\le r<7/9).                     
}\tag{19}\] This attains an open-rectangle contract. Each requested \(\fitinline{r,d}\) uses its own finite curve; no analytic constant passes through a limit. The identity \[\fitmath{
[P(r)-I(r)]^5=(5/16)^5\,7r^4(r-2/3)                       
}\tag{20}\] allows exact certification of affine minorants.

\hypertarget{tangent-and-admitted-same-zero-count}{%
\subsubsection{3. Tangent and admitted same-zero count}\label{tangent-and-admitted-same-zero-count}}

Let \(\fitinline{\alpha}\) be the root in \(\fitinline{(13/15,7/8)}\) of \[\fitmath{
f(x)=3888\,19^5(15x-13)^4(1-x)^5
       -7(7-6x)^4(9x-7)^5.                               
}\tag{21}\] Rational sign substitution gives \[\fitmath{
\frac{109168232013339822100148}{125000000000000000000000}
<\alpha<
\frac{174669171221343715360237}{200000000000000000000000}. 
}\tag{22}\] For uniqueness, logarithmically differentiate the ratio of the two positive terms of (21): \[\fitmath{
\frac{60}{15x-13}-\frac5{1-x}+\frac{24}{7-6x}-\frac{45}{9x-7}>0.
}\] On this interval the first term is at least \(\fitinline{480}\), the second has absolute value at most \(\fitinline{40}\), the fourth at most \(\fitinline{225/4}\), and the third is positive. The intermediate-value theorem and strict monotonicity prove the assertion.

Set \[\fitmath{
w=\frac{12\alpha-14}{9(12\alpha-11)},\quad
g=\frac{w+12-(45/2)w^2}{45w-14},\quad
\delta_*=\frac{45w-14}{54w+6},\quad R_*=5/9+\delta_*/6.     
}\tag{23}\] Then \(\fitinline{.7523<w<.7524}\), \(\fitinline{.4258<\delta_*<.4259}\). Substitution into (21) gives \[\fitmath{
19^5(3w-2)^4(9w+2)^5-567w^4(45w-14)^5=0.
}\] Combining this with (20), retaining the positive fifth-root branch, gives \(\fitinline{g=I(w)}\). Let \(\fitinline{a_0=I(w)-I'(w)w}\), \(\fitinline{a_1=-I'(w)}\), and \(\fitinline{L(r)=a_0-a_1r}\). Convexity proves \(\fitinline{L\le I}\) on \(\fitinline{[3/4,7/9]}\), and exact rational interval substitution gives \[\fitmath{
.05805911<a_0<.05805913,\qquad
.07599618<a_1<.0759962.                                  
}\tag{24}\] Below \(\fitinline{3/4}\), 113 legal finite paired stages (8079 through 8191) supply a lower patch on \(\fitinline{[1467/2000,3/4]}\). At both endpoints of each stage their affine output lies strictly above \[\fitmath{
\frac{725738991}{12500000000}+10^{-7}
 -\frac{7599618487}{100000000000}r.                       
}\tag{25}\] LOWER\_CHAIN\_CERTIFICATE.json stores the complete required 8192-stage six-rational data, beginning at the native top anchor. Part B of FINITE\_SOURCE\_APPENDIX.md specifies every stage formula, all fourteen endpoint gates, the adjacent nesting conditions, finite recursion and strong-induction proof that each stage uses only earlier paired anchors. Its exact linear arithmetic receipt includes all 226 new padded-line endpoint comparisons. Affineness gives the comparison throughout each stage. Equation (24) puts the exact tangent below (25). This uses the actual paired stage supplier, not scalar inequalities alone. On \(\fitinline{[311/400,311/400+1/1000]}\), \(\fitinline{L(r)\le-9/10000}\), so native inverse and reserve ordinary rectangles give the small upper overlap required by the detector fiber's open endpoint. Thus its full paired fiber contract is supplied.

For density \(\fitinline{d\in[0,3/4]}\), amplitude \(\fitinline{p\in[0,1/2]}\), put \[\fitmath{
k(p)=1-p(1+2a_1),\quad U(p)=1-p(1+6a_0-4a_1),\quad V(d)=5/6-d,
}\] \[\fitmath{
J=3V+2dk,\quad N=(1-d)J+VdU,\quad T=J+dU.                  
}\tag{26}\] On the improved window \[\fitmath{
\mathcal W=[189/512,15/32]\times[15/32,1/2],                
}\tag{27}\] the count is \(\fitinline{N/J}\), its long cutoff is \(\fitinline{T/J}\), and its short crossing is \(\fitinline{r_c=2T/(3J)}\). With \(\fitinline{J>0}\), direct algebra gives \[\fitmath{
1-d\{p(1+2a_0)+k(p)r_c\}=N/J
 =1/6+(3/2)V(d)r_c.                                     
}\tag{28}\] These are witnesses at the same zero. Whole-slot selection retains the original coefficient class and loses only its already charged final mesh slot; its cap is \(\fitinline{(1-r)/2+L(r)}\) with a positive rectangle decrement.

Admission also requires every continuation. Define \(\fitinline{B=1-10^{-12}}\), \(\fitinline{\Lambda=2+10^{-11}}\), \(\fitinline{c_0=3/20}\), \(\fitinline{r_0=1467/2000}\), \(\fitinline{r_1=311/400}\), \(\fitinline{s_-=27/1000}\), and \(\fitinline{m_H=(1-(9/2)s_-)/2}\). The ten positive guards are \[\fitmath{
\begin{aligned}
&k,\ J,\ T-J,\ (12499/10000)J-T,\ 2T-3(r_0+10^{-6})J,\\
&N-[1-d(r_1+(1-r_1)p)]J,\\
&N-[1-d(r_1+2(B-r_1)p/\Lambda)]J,\\
&N-[1-d(r_1+2(c_0-10^{-7})p)]J,\\
&N-[1-2m_Hd-2s_-dp]J,\\
&U[-2(1+2a_1)d]+(1+6a_0-4a_1)J.                          
\end{aligned}
}\tag{29}\] The last makes \(\fitinline{N/J}\) antitone in \(\fitinline{p}\). Every guard is separately affine in \(\fitinline{a_0,a_1}\). At all four corners of (24), every tensor-Bernstein coefficient on (27) is positive; the independent checker recomputes all forty guards. The common-witness theorem, (28)--(29), and the strict cap \(\fitinline{c_0-10^{-7}<c_0}\) therefore admit the branch. Native branches remain available on \[\fitmath{
[15/64,9/16]\times[9/32,1/2],\qquad
[15/64,1/2]\times[3/16,1/2].                              
}\tag{30}\] Let \(\fitinline{\mathcal E_\kappa(d,p)}\) be the minimum of the old count and precisely these available native and improved branches, with their full admission conditions. A branch contributes nothing outside its window. This is the conductor-independent antitone envelope consumed by the physical return.

\hypertarget{the-averaged-physical-floor}{%
\subsubsection{4. The averaged physical floor}\label{the-averaged-physical-floor}}

The old pointwise floor premise fails at the new geometry. The replacement must be proved for the actual physical numerator, not obtained by deleting a geometry field. The original numerator reconstruction and unmarked physical fourth moment give it as follows.

Fix the original arithmetic source data, finite main-slot number \(\fitinline{n}\), total width \(\fitinline{\ell=\sum_j\ell_j}\), bad-prime set, profiles and one actual reflected coefficient packet. Its nonprincipal rows satisfy \(\fitinline{U^{1/100}\le Nu\le U}\), with every natural mask retained. Write its positive coefficient scale as \(\fitinline{N_u}\) and its actual combined conductor as \(\fitinline{C_u\le C_{\rm fix}Nu}\). Put \(\fitinline{L_u=C_u/(C_{\rm fix}Nu)\le1}\). For fixed \(\fitinline{\xi>0}\), partition \(\fitinline{\max(1,L_uN_u)}\) into boxes \[\fitmath{
U^{j\xi}\le\max(1,L_uN_u)\le U^{(j+1)\xi}.
}\] For \(\fitinline{Y=U^v}\), \(\fitinline{0<v\le99}\), the active condition \(\fitinline{N_u\le bY}\) implies \[\fitmath{
C_u\le C_{\rm fix}U^m(Y/N_u),\quad
m=\max(0,1-v+(j+1)\xi)\le1+2\xi.                          
}\tag{31}\] Choose positive annular slack with \(\fitinline{2\xi+\mathrm{slack}\le1/2}\). On inactive rows use a positive safe scale to invoke the full fourth moment, then discard those rows by positivity. This handles the actual row-dependent scales; they are not replaced by one common dyadic coefficient polynomial. The original certified fourth moment, with \textbf{empty selected product}, gives on each active box \[\fitmath{
\sum_u|A_u(Y/N_u)|^4\le C U^{1+\tau+\epsilon}.             
}\tag{32}\] Here \(\fitinline{\tau>0}\) is the fixed source slack. The moment parameter is any \(\fitinline{\kappa_m\in[.7,.75]}\) with \(\fitinline{2\beta-1\le\kappa_m}\), supplied by (2). Constants precede the scale and are uniform over the allowed finite row sets, active boxes and coefficient scales, for the fixed outer source data.

The ball contains at most \(\fitinline{128U}\) actual rows. H\"{o}lder, (32), and the original exact coefficient-packet mass bound give the empty-selection packet estimate \[\fitmath{
\sum_u|\mathcal P_u(Y)P_u|
 \ll Y^{-1/2}\prod_{i\in J}Y_i^{-4/25}
             U^{1+(38/25)\tau},\qquad |P_u|\le1.           
}\tag{33}\] Before harmless absorption its row exponent is \(\fitinline{1+\epsilon_m/4+\xi/100+\epsilon_{\rm mass}}\), and it has the finite number of dilation boxes as a prefactor. Thus cardinality is used with power \(\fitinline{3/4}\), fourth moment with power \(\fitinline{1/4}\). Take \(\fitinline{\epsilon_m=2\tau}\), \(\fitinline{\xi=\epsilon_{\rm mass}=\tau}\) to obtain (33). This is the empty-selection specialization of the retained literal numerator packet theorem. It uses no density gain or improved affine source.

For a physical error mask \(\fitinline{J}\), the complementary main product has the deterministic positive cap \[\fitmath{
B_Z=Z^{(1/100+\tau)\ell}.                                 
}\tag{34}\] Apply (33) with that product divided by \(\fitinline{B_Z}\), then restore \(\fitinline{B_Z}\) outside the positive sum. One never divides by a row amplitude. All zeros are retained. The exact front contributes \(\fitinline{Z^{-(4/25)\sum_{i\notin J}\ell_i}}\) and the packet contributes \(\fitinline{Z^{-(4/25)\sum_{i\in J}\ell_i}}\); including (34), every mask has \[\fitmath{
Z^{(-3/20+\tau)\ell}.                                    
}\tag{35}\] Sum all masks using the exact \(\fitinline{2\pi}\) numerator-reconstruction factor. At contour real parts \(\fitinline{\Re z=17/50}\), \(\fitinline{\Re x=51/100+16e}\), the retained frequency/row scalar bound contributes \(\fitinline{U^{-17/50+2\epsilon}}\), its row-annulus slack, and the polynomial modulus/height factors. The original remaining contour integrations then give \[\fitmath{
\begin{aligned}
|\mathrm{FullW}_{\rm floor}|\ll{}&
X^{4/25}Z^{-3/20+16e}
(N\mathfrak q_\eta)^{2\epsilon}(3+H)^{2\epsilon}Y^{-1/2}\\
&\cdot Z^{(17/25)\tau+(-3/20+\tau)\ell}
U^{33/50+2\epsilon+(38/25)\tau}.                           
\end{aligned}
}\tag{36}\] The inclusive floor event says its detector maximum at the next padded height is less than \(\fitinline{51/100+2e}\). It supplies (34) along the entire contour strip. The current height is at most \(\fitinline{(3i+2)Z^\upsilon}\), while the event is at \(\fitinline{3(i+1)Z^\upsilon}\). The original support, coprimality and calibration assumptions remain. No selected large-value mark or count saving is used.

Set \(\fitinline{X=Z^x,Y=Z^y,U=Z^D}\) and divide by the original nonzero normalizer. After the small analytic losses the exponent is \[\fitmath{
E_f(D)=4x/25-3/20-y/2-3\ell/20+33D/50
       =\mathrm{signalShift}+F_f(D),                     
}\tag{37}\] \[\fitmath{
F_f(D)=51/100-49h/150+51\ell/100+33D/50,\qquad h=1+\ell-x.
}\] The equality uses \(\fitinline{x+y=1-\ell}\) and includes the full signal normalizer. The chosen-source loss is \[\fitmath{
16e+17\tau/25+\tau\ell+D(2\epsilon+38\tau/25)+\tau+\tau/8
<10^{-5}.                                               
}\tag{38}\] The retained coarse box \[\fitmath{
t>.873,\quad .793<h<.816,\quad .1672<\ell<.175,\quad
D\le h+10^{-10}
}\] gives \[\fitmath{
t-F_f(D)-10^{-5}-10^{-12}>10^{-3}.                        
}\tag{39}\] Indeed the coefficient of \(\fitinline{h}\), after substituting \(\fitinline{D=h+10^{-10}}\), is \(\fitinline{1/3}\), so upper bounds for \(\fitinline{h,\ell}\) suffice. Each normalized floor cell has exponent at most \(\fitinline{\mathrm{signal}(t)-1/2000}\). The dyadic row sum costs at most an additional \(\fitinline{Z^{1/4000}}\); the finite height partition only changes its fixed constant. The full Fourier-tail error from cube to FullW is \(\fitinline{O((N\mathfrak q_\eta)^2Z^{-1})}\). Thus \[\fitmath{
\left|\frac{\mathrm{floorCube}}{\mathrm{normalizer}}\right|
\ll (N\mathfrak q_\eta)^{2+2\epsilon}
             Z^{\mathrm{signal}(t)-1/4000}.               
}\tag{40}\] The exact supported row split subtracts this actual floor from the full normalized transport, leaving only the true nonfloor numerator.

The source is chosen first. There is then a finite scalar degree \(\fitinline{J\ge0}\). Next choose the positive height exponent \(\fitinline{\upsilon}\) satisfying its height-admissibility conditions. For each finite height-grid bound \(\fitinline{n}\) there is \(\fitinline{C>0}\), and for every fixed \(\fitinline{\eta}\) the estimate holds for all sufficiently large \(\fitinline{Z}\), uniformly in all admissible buffered selections and supported finite row subsets. Choosing height before the degree, or discarding an unsupported row to obtain (40), would not establish this result.

\hypertarget{exact-return-including-the-zero-margin-contact}{%
\subsubsection{5. Exact return including the zero-margin contact}\label{exact-return-including-the-zero-margin-contact}}

Put \(\fitinline{\kappa=2\alpha-1}\), \(\fitinline{\ell=11/3-4\alpha}\). The contact identities are \[\fitmath{
1-\delta_*\{(1+w)/2+I(w)\}=R_*
 =1/6+(3/2)(5/6-\delta_*)w.                              
}\tag{41}\] Write \(\fitinline{\rho(d)=N(d,1/2)/J(d,1/2)}\), and choose \[\fitmath{
h=\frac{1/2+\ell}{1/8-(3/4)\rho'(\delta_*)},\quad
y=h-2\ell,\quad x=1+\ell-h.                              
}\tag{42}\] The denominator is positive. Exact interval propagation gives \[\fitmath{
.815324503<h<.815324505.                                  
}\tag{43}\] For conductor exponent \(\fitinline{D}\), capacity \(\fitinline{C}\), density \(\fitinline{d}\), amplitude \(\fitinline{p}\), the complete physical row is \[\fitmath{
\mathcal R(D,C,d,p)=1/2+13h/75-y/2-\ell/2-9D/100
 +d/2+dp(\ell-C/2)+(3D/4)\mathcal E_\kappa(d,p).            
}\tag{44}\] Its two capacities are \[\fitmath{
C_o(D)=\frac{\min(D,2y-D)}{6\kappa},\qquad
C_n(D)=\min(3D/40,(2y-D)/2).                              
}\tag{45}\] The four edges are \(\fitinline{D=1/200,19/25,19/25,h}\), with capacities respectively \[\fitmath{
\frac1{200\cdot6\kappa},\quad
\frac{2y-19/25}{6\kappa},\quad
\frac3{40}\frac{19}{25},\quad\frac{h-4\ell}{2}.             
}\tag{46}\] RETURN\_CERTIFICATE.json contains all 57 rectangular cells covering \(\fitinline{[0,3/4]\times[0,1/2]}\) at those four edges, with each admitted branch and the five rational coefficient-lower polynomials used by 37 noncritical cells. For branch denominator \(\fitinline{J_b>0}\), the margin is \[\fitmath{
J_b(d,p)\{\alpha-\mathrm{base}(D,C,d,p)\}
                      -(3D/4)N_b(d,p).                  
}\tag{47}\] The old branch has \[\fitmath{
D_o=3-p-2p/(3\kappa),\quad C_o'=2(1-p)(1-p/(3\kappa)),
\quad J_o=VD_o+dC_o',\quad N_o=(1-d)J_o+dC_o'V/2.          
}\tag{48}\] The native branch is the original affine-source count on (30); the marked branch is (26) on (27). Denominators are positive before clearing them. Exact rational parameter boxes enclose every coefficient under interval addition, multiplication and positive-denominator division. Coefficientwise lower bounds preserve a polynomial lower bound since \(\fitinline{d,p\ge0}\). Positive tensor-Bernstein coefficients prove each noncritical cell including its boundary. Our checker verifies the 37 saved rational-polynomial inequalities and full cover/admission. Parameter containment and saved coefficient-lower identities are the exact source argument; these finite checks do not establish the analytic moments.

The twenty remaining cells lie in (27), at \(\fitinline{D=h}\), using the marked branch. Their limiting margin is nonnegative, not uniformly positive. Let \[\fitmath{
B_h(d,p)=1/2+\ell/2-5h/12+d/2+dp(2\ell-h/4),\quad
M_h=J(\alpha-B_h)-(3h/4)N.
}\] Put \(\fitinline{a=1+2a_1}\), \(\fitinline{b=1+6a_0-4a_1}\), and \[\fitmath{
q_{10}=2a(\alpha-1/2-\ell/2)+5\ell-h(2a/3+5/8+5b/8),
}\] \[\fitmath{
q_{20}=-a-(1+a)(2\ell-h/4)+(3h/4)(2a+b),\qquad
q_{21}=-2a(2\ell-h/4).
}\] Equations (41)--(42) make the maximum-amplitude quadratic and its density derivative vanish at \(\fitinline{\delta_*}\). Put \(\fitinline{j'=-2-2a_1}\) and \[\fitmath{
K_h=-j'(1/2+\ell-h/8)+(3h/4)(j'+U(1/2))>1/5.
}\] Its exact quadratic expansion gives \[\fitmath{
M_h(d,p)=K_h(d-\delta_*)^2
 +(1/2-p)d\{q_{10}+d(q_{20}+q_{21}p)\}.                   
}\tag{49}\] Exact bounds \(\fitinline{q_{10}>-.143}\), \(\fitinline{q_{20}>.585}\), \(\fitinline{q_{21}>-.331}\) make the braces at least \(\fitinline{1/100}\) throughout (27). Hence the whole critical rectangle has nonnegative margin, with equality at \(\fitinline{(\delta_*,1/2)}\). This pays all twenty cells without a false strict certificate. The positive denominators and \(\fitinline{\mathcal E_\kappa\le N_b/J_b}\) then imply \(\fitinline{\mathcal R\le\alpha}\) at each edge.

For fixed \(\fitinline{d,p}\), the row is antitone in \(\fitinline{C}\). Substituting either affine branch of (45) makes it affine in \(\fitinline{D}\). The minimum capacity makes the row the maximum of those affine expressions, hence convex in \(\fitinline{D}\). Endpoint bounds therefore fill the whole conductor intervals, retaining every branch switch: \[\fitmath{
\mathcal R(D,C_o(D),d,p)\le\alpha\quad(1/200\le D\le19/25),
}\] \[\fitmath{
\mathcal R(D,C_n(D),d,p)\le\alpha\quad(19/25\le D\le h).     
}\tag{50}\]

\hypertarget{fixed-positive-costs-and-the-rational-theorem}{%
\subsubsection{6. Fixed positive costs and the rational theorem}\label{fixed-positive-costs-and-the-rational-theorem}}

The equality case of (49) has no budget for source cost, row fee or positive conductor extension. The accepted result keeps all three. Choose \[\fitmath{
s_0=\frac1{100\,10^{256}}=10^{-258}>0.                    
}\tag{51}\] The genuine physical source is supplied at \[\fitmath{
B=1-s_0,\quad\Lambda=2+10s_0,\quad
c_{\rm source}=3/20-1000s_0/1143.
}\] It restricts to the half-width physical face with cap \(\fitinline{3/40}\), with the original coefficients, zeros, profiles and height controls. For every positive loss it supplies an early mesh before arithmetic data; constants need not be uniform as \(\fitinline{s_0}\) shrinks. Taking the finite minimum of count and physical meshes gives their common early mesh.

Set \(\fitinline{\ell_t=11/3-4t}\), \(\fitinline{y_t=h-2\ell_t}\), \(\fitinline{x_t=1+\ell_t-h}\), \(\fitinline{K_t=2t-1}\). The actual upper capacity is \[\fitmath{
C_{s,t}(D)=\min\left(3D/40,
                \frac{2y_t-(1+s)D}{2+10s}\right).         
}\tag{52}\] For all \(\fitinline{d\in[0,3/4]}\), \(\fitinline{p\in[0,1/2]}\), the exact changes are \[\fitmath{
C_{s,\alpha}(D)\ge C_n(D)-5s,\qquad
\text{row debit}\le15s/16;                               
}\tag{53}\] \[\fitmath{
\mathrm{row}_t\big|_{\rm fixed\ capacity\ and\ count}
 =\mathcal R-(2+4dp)(t-\alpha);                           
}\tag{54}\] \[\fitmath{
\mathcal E_K-\mathcal E_\kappa\le5(K-\kappa)/216,\qquad
\text{row count debit}\le(K-\kappa)/50\quad(D\le1);          
}\tag{55}\] \[\fitmath{
\mathrm{row}(D)-\mathrm{row}(h)\le2(D-h)\quad(D\ge h).       
}\tag{56}\] For (53), compare the second fractions with positive denominators and retain the unchanged \(\fitinline{3D/40}\) minimum branch. The cost is \(\fitinline{dp/2}\) times lost capacity, and \(\fitinline{dp\le3/8}\). Equation (54) is an identity with the same capacity and count on both sides. For the actual capacity, \(\fitinline{C_{s,t}(D)\ge C_{s,\alpha}(D)}\) because \(\fitinline{t\ge\alpha}\) increases \(\fitinline{y_t}\); the row is antitone in capacity, so this change only improves that bound. Equation (55) is the original quantitative old-count bound; native and marked branches stay fixed, so taking their common minimum preserves it. Old-count monotonicity cannot replace it. For (56), capacity falls by at most \(\fitinline{(D-h)/2}\), \(\fitinline{\mathcal E_K\le2}\), and substitution in (44) bounds the increase by \(\fitinline{2(D-h)}\). Consequently upper rows on the extended interval satisfy \[\fitmath{
\mathrm{row}_t(D)\le
\alpha+15s/16-(49/25)(t-\alpha)+2\Delta
\quad(19/25\le D\le h+\Delta).                            
}\tag{57}\]

The old interval must also move its denominator: \[\fitmath{
C_{o,t}(D)=\frac{\min(D,2y_t-D)}{6K_t}
           \ge C_o(D)-(K_t-\kappa).
}\] Its capacity debit is at most \(\fitinline{(3/16)(K_t-\kappa)}\). Combining with (54)--(55) gives \[\fitmath{
\mathrm{row}_t(D)\le\alpha-\tfrac32(t-\alpha)
\quad(1/200\le D\le19/25).                               
}\tag{58}\] The rational certificate proves the stronger separation \[\fitmath{
\alpha+25600s_0<t.                                      
}\tag{59}\] Indeed, put \(\fitinline{t_b=t-256/10^{256}}\). Exact arithmetic gives \(\fitinline{f(t_b)>0}\), whereas the lower Farey neighbor \(\fitinline{t_-}\) has \(\fitinline{f(t_-)<0}\); both belong to \(\fitinline{(13/15,7/8)}\). Uniqueness proves \(\fitinline{\alpha<t_b}\), and \(\fitinline{25600s_0=256/10^{256}}\). The lower neighbor and determinant are in check\_exact.py and its receipt. The rational field contains exactly 256 characters. That is a declaration-format limit, not an analytic optimum.

Use \(\fitinline{s=s_0}\), conductor extension \(\fitinline{\Delta=s_0}\), and row fee \(\fitinline{s_0}\). Equations (57)--(59) imply strictly \(\fitinline{\mathrm{row}_t<t-s_0}\) on both complete intervals. All original geometric obligations hold: \(\fitinline{1/6<\ell_t<1/5}\), \(\fitinline{\ell_t<x_t<y_t<1/2}\), \(\fitinline{x_t+y_t=1-\ell_t}\), the signal identity and switch \(\fitinline{19/25}\), and actual prime supply \(\fitinline{\ell_t/(h+s_0)>1/5}\). A nonpositive shortened plain length uses the native cone instead of a positive-length source.

Suppose now \(\fitinline{\beta>t}\), put \(\fitinline{\mu=\beta-t>0}\), and write \(\fitinline{\kappa_\beta=2\beta-1=K_t+2\mu}\). The count sources, (19), its lower patch and the positive-cost physical source give one chosen source at a common early mesh. The natural cap uses (2), \(\fitinline{\kappa_\beta\le3/4}\). The fixed branches, source line and geometry remain those at target \(\fitinline{t}\); the old count and natural inverse capacity must be evaluated at the actual parameter \(\fitinline{\kappa_\beta}\).

The retained feedback calculation gives the quantitative changes \[\fitmath{
E_{\rm new}(\kappa_\beta)-E_{\rm new}(K_t)
\le(\kappa_\beta-K_t)/50,
}\] \[\fitmath{
E_{\rm old}(\kappa_\beta)-E_{\rm old}(K_t)
\le(\kappa_\beta-K_t)/8.                                 
}\tag{60}\] For the new row the capacity stays fixed and the admitted-envelope bound (55) suffices. For the old row include its capacity change: writing \(\fitinline{m=\min(D,2y_t-D)\le1/2}\), the extra term is exactly \[\fitmath{
\frac{dp\,m}{12}\left(\frac1{K_t}-\frac1{\kappa_\beta}\right).
}\] It is bounded using \(\fitinline{dp\le3/8}\), \(\fitinline{K_t,\kappa_\beta\ge.7}\) and \[\fitmath{
1/K_t-1/\kappa_\beta
=(\kappa_\beta-K_t)/(K_t\kappa_\beta).
}\] Together with the count debit, this is at most the old allowance in (60). The nominal rows are at most \(\fitinline{t-s_0}\); subtracting \(\fitinline{\beta=t+\mu}\) gives \[\fitmath{
E_{\rm new}(\kappa_\beta)-\beta\le-s_0-24\mu/25,\qquad
E_{\rm old}(\kappa_\beta)-\beta\le-s_0-3\mu/4.              
}\tag{61}\] Thus the strict rows hold at the actual supremum, not just at the target count parameter. The genuine finite nonfloor numerator consumes (61), while (40) removes the actual floor. After their finite height degrees are fixed, the original joint-height theorem chooses one admissible positive height scale. The exact physical probe then differs from its nonzero reciprocal signal by a fixed power saving. The retained common-probe/Mellin argument improves the actual zero supremum strictly, contradicting its definition. Hence \(\fitinline{\beta\le t}\), and the original all-moduli norm pullback proves (1).

This passage proves no nonvanishing at the exact algebraic root \(\fitinline{\alpha}\). The zero-margin scalar return at \(\fitinline{\alpha}\) and the complete analytic theorem at fixed rational \(\fitinline{t}\) have different scope. Another endpoint requires its own actual parameterization and positive-loss order.

\hypertarget{verification-scope-and-integration-boundary}{%
\subsubsection{7. Verification scope and integration boundary}\label{verification-scope-and-integration-boundary}}

The independent arithmetic receipt verifies 53 named groups: 13 source-line intervals, forty count guards, 37 strict return boxes, all 57 cell admissions and complete cover, rational root/cost separation, floor margin and the outside-old-cone consumer in MISSED\_ADVANCE.md. It runs only our own standard-library arithmetic and reads literal contributor data.

This lane has run neither Lean nor contributor programs, providers or a library build. It checked the finite chain's scalar data by one exact linear arithmetic pass, without independently recompiling its analytic source proof or complete closure. Exact source statements, attribution and evidence scope are retained. The manuscript's integration is complete only with the coupled-anchor/reference-lane argument in the parallel anchor note and the retained original numerator, fourth-moment, reconstruction, height and Mellin foundations. A polynomial certificate or catalogue approval cannot repair a missing physical coefficient adapter.

\hypertarget{appendix-f}{}

\hypertarget{appendix-f.-two-finite-source-obligations-used-by-the-continuum-return}{%
\subsection{Appendix F. Two finite-source obligations used by the continuum return}\label{appendix-f.-two-finite-source-obligations-used-by-the-continuum-return}}

This appendix supplies the variable-tolerance HIGH calculation and the exact finite lower-patch construction referenced in RESEARCH.md. Sil's original paired source theorem and Joseph Bebel's zero fixed-stage-loss adaptation remain the analytic step; the calculations here establish its scalar hypotheses and the acyclic order of its source inputs. Original OpenAI, Collet, ArviFoxy and other inherited foundations and notices remain.

\hypertarget{a.-the-variable-tolerance-high-price}{%
\subsubsection{A. The variable-tolerance HIGH price}\label{a.-the-variable-tolerance-high-price}}

Fix a strict curve, a right endpoint \(\fitinline{R}\), and \(\fitinline{r\le R}\) as in RESEARCH.md \S{}2. Use \[\fitmath{
z=W(R),\quad v=W'(R),\quad H=r/R,\quad u=(1-H)/6,\quad
s=H((1-R)/2+z)+u/k,
}\] \[\fitmath{
k=\{2+6(z-Rv+R\gamma)\}^{-1},\quad
a_0=z-Rv-q/5,\quad a_1=-v,\quad
0<q\le\min(10^{-12},\gamma/1000).
}\] The strict-curve bounds give \[\fitmath{
3/4\le R\le7/9,\quad .9\le H\le1,\quad 0\le z\le.0011,
\quad .363656\le k<.5,\quad0\le a_1\le.1,\quad u\le10^{-4}.
}\tag{A1}\] Here \(\fitinline{s}\) is the child mark width after removing the bare \(\fitinline{u}\) term; it is not the total \(\fitinline{S=s+u}\) of the ordinary source.

Write \(\fitinline{B=1/2-a_1=1/2+v>0}\), and define \[\fitmath{
w_b=(R/2+z-a_0)/B,\qquad
F_0=5/6-(5/4-3k/2)w_b-k(1+R)/2-kz.
}\] The unpadded reserve has the exact identity \[\fitmath{
B_0:=5/6-5R/4+k(R-1/2-z)
 =\frac{kR(15R-10)}2\{v-\gamma-F(R,z)\}.                  
}\tag{A2}\] The cutoff relation defining \(\fitinline{k}\) proves the equality simply by expansion. The differential reserve is at least \(\fitinline{3\gamma}\). Since \(\fitinline{kR(15R-10)\ge1/3}\), (A2) gives \(\fitinline{B_0\ge\gamma/3}\). The actual padded line satisfies \[\fitmath{
w_b=R+\frac{q}{5(1/2+v)},\qquad
F_0=B_0-\frac{(5/4-3k/2)q}{5(1/2+v)}\ge B_0-q.
}\tag{A3}\] The inequality follows from \(\fitinline{v\ge-4/45}\), \(\fitinline{k\ge0}\), \(\fitinline{q\ge0}\). Thus \(\fitinline{F_0\ge\gamma/4}\), and \[\fitmath{
HF_0\ge9\gamma/40>30q.                                   
}\tag{A4}\] This is the actual positive reserve after the source-pad debit.

For each non-tiny HIGH class \(\fitinline{k\le\nu\le3/4,\ 0\le e\le\nu}\), choose \[\fitmath{
g_\nu=\frac{r+s-H/2-Ha_0-u/\nu}{B},\qquad
w_\nu=g_\nu/H+q,\qquad
\sigma_q(w)=\max(0,(1-w)/2+a_0-a_1w-q/10).                
}\tag{A5}\] The identity \[\fitmath{
B(g_\nu-r)=H(z-a_0+a_1R)+u(1/k-1/\nu)\ge0
}\tag{A6}\] gives \(\fitinline{w_\nu\ge R+q}\). Also \(\fitinline{g_\nu}\) increases with \(\fitinline{\nu}\). Its formula and (A1) give \(\fitinline{g_{3/4}/H\le R+.01}\); hence \(\fitinline{2/3<w_\nu\le.79}\). Positivity and the earlier-anchor inequality \(\fitinline{a_0-a_1R\le z}\) imply \[\fitmath{
H\sigma_q(w_\nu)\le H((1-R)/2+z)\le s.                   
}\tag{A7}\] In particular this is actual mark supply, not a free selected-width parameter. Since \(\fitinline{a_0\le1}\), \(\fitinline{w_\nu\ge0}\) and \(\fitinline{a_1\ge0}\), one may use the conservative \(\fitinline{\sigma_q\le3/2}\), so \(\fitinline{w_\nu+\sigma_q<6}\).

The two detector counts at this cutoff are \[\fitmath{
C_{\rm short}=H-\nu Hw_\nu-e\min(s,H\sigma_q(w_\nu)),
}\] \[\fitmath{
C_{\rm long}=H/6+(5/6-\nu)H(3w_\nu/2+q).
}\tag{A8}\] Their price is \(\fitinline{\nu r+es+\max(C_{\rm short},C_{\rm long})-H}\). Both counts, including both positive pads, must be paid.

For the short count let \[\fitmath{
X_0=(H-g_\nu)/2+Ha_0-a_1g_\nu,\qquad
c_q=(1/2+a_1)q+q/10.
}\] Then \(\fitinline{0\le c_q\le q}\), and the raw selected expression has \[\fitmath{
H\{(1-w_\nu)/2+a_0-a_1w_\nu-q/10\}=X_0-Hc_q.
}\] It is bounded above by both \(\fitinline{s}\) and \(\fitinline{H\sigma_q}\). Moreover (A5) gives \[\fitmath{
s-X_0=g_\nu-r+u/\nu\ge0,\qquad
\nu(r-g_\nu)+e(s-X_0)\le u.
}\] Consequently the short price is at most \[\fitmath{
\nu(r-g_\nu)+e(s-X_0)-\nu Hq+eHc_q\le u,
}\tag{A9}\] using \(\fitinline{e\le\nu}\), \(\fitinline{c_q\le q}\). The lost cap is explicitly paid by the positive witness shift; it has not disappeared.

For the long count define \[\fitmath{
G(\nu)=H+u-H/6-(5/6-\nu)(3g_\nu/2)-\nu(r+s).
}\] Direct expansion gives \[\fitmath{
G(k)=HF_0,
}\] \[\fitmath{
G(\nu)-G(k)=(\nu-k)
\left\{\frac32g_\nu-r-s
 -\frac32(5/6-k)\frac{u}{k\nu B}\right\}.                 
}\tag{A10}\] The braces are nonnegative: \(\fitinline{g_\nu\ge r}\), \(\fitinline{\nu\ge k}\), and the derivative gate is \[\fitmath{
r/2-s-\frac32(5/6-k)\frac{u}{k^2B}\ge0.                  
}\tag{A11}\] For completeness, (A1) bounds its left side below by \[\fitmath{
H(R-1/2-z)-u/k-30u
\ \ge\ 27/125-3u-30u>0.
}\] Thus \(\fitinline{G(\nu)\ge HF_0>30q}\). Substituting \(\fitinline{Hw_\nu=g_\nu+Hq}\) in the long count shows that its extra padding cost is \[\fitmath{
(5/6-\nu)H(5q/2)\le5q/2.
}\] The long price is therefore at most \[\fitmath{
u-G(\nu)+5q/2<u-55q/2<u.                                
}\tag{A12}\] Finally, the tiny class has reserve at least \(\fitinline{3/100}\): its gate is \[\fitmath{
H+u-H/6-1/15-(3/4)r-s>3/100.
}\] Its left side equals \(\fitinline{H(1/3-R/4-z)+u-u/k-1/15}\), which exceeds \(\fitinline{3/100}\) by (A1). Equations (A9)--(A12) prove the class price bound, including tiny cores.

The recentering used by the physical analytic consumer is exact. Put \[\fitmath{
\widetilde a_0=a_0+10^{-13}-q/10.
}\] Then \[\fitmath{
\max(0,(1-w)/2+\widetilde a_0-a_1w-10^{-13})
=\sigma_q(w).
}\tag{A13}\] The generic consumer is instantiated with the same actual witness shift \(\fitinline{q}\), so its short and long counts are exactly (A8), and its detector price is exactly the one bounded in (A9)--(A12). It changes no coefficient, mask, profile, row or selected subset. The complete earlier-anchor analytic supplier also uses the same positive source pad \(\fitinline{q/10}\); a scalar identity by itself would not supply that physical theorem. PR141's ToleranceSources and ToleranceConsumer perform this parameterized transport. The fixed PR137 pad is therefore not being silently used at a new value.

The remaining zero-fixed-loss adapter retains the requested rectangle decrement. For every sufficiently small \(\fitinline{d>0}\), take the child-radius shortening \(\fitinline{\varkappa=d/8}\) and set the fixed gain decrement to zero. With \(\fitinline{S_R=(1-R)/2+z}\), the original LOW ledger becomes \[\fitmath{
u-kd/2+k(3+S_R+d/2)d/8<u,
}\] because \(\fitinline{S_R\le.137}\) and \(\fitinline{d}\) is small. Thus the requested positive rectangle decrement, rather than a permanent per-stage gain loss, pays LOW. The ordinary predecessor has its own positive decrement \(\fitinline{d/4}\): restrict the already proved ordinary anchor to the source \[\fitmath{
G(R-d/4,\ d/4,\ S_R-d/4).
}\] This follows from its rectangle with decrement \(\fitinline{d/4}\); the extra positive slope only narrows that rectangle because its mark width is bounded by \(\fitinline{S_R-d/4}\). It does not set the old expression \(\fitinline{\min(D_{\rm fixed},d)/4}\) to zero when \(\fitinline{D_{\rm fixed}=0}\). For \(\fitinline{c_1=S_R-d/4}\), choose the positive tolerance and mesh with \(\fitinline{\xi=\mathrm{tol}+6\,\mathrm{mesh}\le7d/64}\). The exact packed-side LOW cap identity is \[\fitmath{
u+k\{Hd/4-d+c_1\xi\}\le u.
}\] Indeed \(\fitinline{H\le1}\), \(\fitinline{0\le c_1\le1}\), so its braces are at most \(\fitinline{d/4-d+7d/64<0}\). The other LOW deletion-budget defect is strictly negative at the same predecessor face. The ordinary fork's original reference gates and positive reserve therefore remain available. For HIGH the actual child free width decreases by \(\fitinline{3d/4}\), and its row radius decreases by \(\fitinline{d/8}\). On each selected class, the already proved \(\fitinline{w+\sigma_q(w)\le6}\) and \(\fitinline{0\le e\le\nu}\) pay that radius change in the literal inverse and ordinary fiber moment. These are restrictions of the original child source, not new source estimates at an excluded endpoint. The reference lane is unchanged and keeps its reflected ordinary coefficient class. The full positive mesh is chosen after these fixed decrements and before slot types/profiles/arithmetic data. A proof for small \(\fitinline{d}\) suffices: larger decrements shrink the rectangle. This supplies the actual zero-fixed-stage-loss paired output used by the finite curve.

\hypertarget{b.-exact-finite-chain-supplying-the-lower-patch}{%
\subsubsection{B. Exact finite chain supplying the lower patch}\label{b.-exact-finite-chain-supplying-the-lower-patch}}

LOWER\_CHAIN\_CERTIFICATE.json contains all 8192 exact ordered rational tuples \[\fitmath{
(R_i,l_i,z_i,k_i,a_i,b_i),\qquad 0\le i<8192.              
}\tag{B1}\] It is mathematical data, not executable code. Its bytes are bound to the exact pinned CertificateData source in SOURCE\_USE.json. Index \(\fitinline{i}\) is the source's block \(\fitinline{i\mathbin{//}64}\), position \(\fitinline{i\bmod64}\). These tuples, with the formulas below, completely specify the construction; one need not trust or rerun the program that originally selected them.

The base is \[\fitmath{
R_0=311/400,\quad z_0=a_0=b_0=0,\quad
l_0=248799/320000,\quad k_0=249624624679/500000000000.
}\tag{B2}\] Its anchor is the native inverse plus reserve ordinary anchor. The last left endpoint is \(\fitinline{l_{8191}=1467/2000}\). Fixed positive quantities are \[\fitmath{
p_{\rm src}=10^{-13},\qquad q_{\rm wit}=10^{-12},\qquad
\rho_{\rm high}=3\cdot10^{-11}.                           
}\tag{B3}\] The old comparison line has \[\fitmath{
A=2902544862471/(5\cdot10^{13}),\quad
B=1899637920999/(25\cdot10^{12}),\quad L_*(r)=A-Br.
}\] For stage \(\fitinline{i}\) and any \(\fitinline{r\in[l_i,R_i]}\), define \[\fitmath{
H_i=r/R_i,\quad u_i=(1-H_i)/6,\quad
Y_i(r)=H_i z_i+(1/k_i-2)u_i,\quad
S_i=(1-r)/2+Y_i,\quad s_i=S_i-u_i,
}\] \[\fitmath{
B_i=1/2-b_i,\quad w_{b,i}=(R_i/2+z_i-a_i)/B_i,
}\] \[\fitmath{
F_i=5/6-(5/4-3k_i/2)w_{b,i}-k_i(1+R_i)/2-k_i z_i,
}\] \[\fitmath{
g_i(\nu)=\{r+s_i-H_i/2-H_i a_i-u_i/\nu\}/B_i,\quad
w_i(\nu)=g_i(\nu)/H_i+q_{\rm wit},\quad
L_i(r)=a_i-b_i r.
}\tag{B4}\] The data certificate requires the following basic inequalities at every stage: \[\fitmath{
.7\le R_i\le311/400,\quad1467/2000\le l_i<R_i,\quad
45457/125000\le k_i<1/2,\quad z_i\ge0,\quad0\le b_i<2/5,
}\] \[\fitmath{
u_i(l_i)\le10^{-4},\quad R_i\le w_{b,i},\quad
L_i(R_i)\le z_i,
}\] \[\fitmath{
(a_i=b_i=0)\ \text{or}\ (b_i>0,\ L_i(311/400)<0),\qquad
L_*(r)+p_{\rm src}<Y_i(r)\quad(r=l_i,R_i).                 
}\tag{B5}\] It also requires all fourteen following expressions to be positive at both \(\fitinline{r=l_i,R_i}\); put \(\fitinline{w_+=w_i(3/4)}\) and \(\fitinline{\sigma_+=(1-w_+)/2+L_i(w_+)}\): \[\fitmath{
\begin{aligned}
&H_iF_i-\rho_{\rm high},\quad .137-S_i,\quad .14-S_i/H_i,\\
&.9-r-S_i-4u_i,\quad .91-r-S_i-4u_i,\\
&5/6+21/250-r-S_i-3u_i-.01,\\
&3H_i-2r-8s_i-.2,\quad r-.5-Y_i+u_i-.2,\\
&r/2-s_i-\tfrac32(5/6-k_i)u_i/(k_i^2B_i)-.1,\\
&H_i+u_i-H_i/6-1/15-(3/4)r-s_i-.03,\\
&w_i(k_i)-2/3,\quad .79-w_+,\\
&1-(3/4)w_+-(3/4)\max(0,\sigma_+)-1/6,\\
&.79-r-(2+10^{-11})u_i.                                  
\end{aligned}
}\tag{B6}\] These are exact rational checks, without floating-point rounding. Any rounding used by the original data selector has already been frozen into the displayed rational tuple. No unrecorded rounding algorithm is needed.

The adjacent conditions, for \(\fitinline{i<8191}\), are \[\fitmath{
R_{i+1}=l_i,\quad z_{i+1}\le Y_i(l_i),\quad b_i\le b_{i+1},
\quad L_{i+1}(R_i)\le L_i(R_i),
}\] \[\fitmath{
L_{i+1}(l_i)+p_{\rm src}<Y_i(l_i),\quad
L_{i+1}(R_i)+p_{\rm src}<Y_i(R_i).                         
}\tag{B7}\] The independent linear exact pass has checked all (B5), all 229376 endpoint inequalities in (B6), and all 8191 adjacent contracts. Its receipt records positive minima, including the minimum strict adjacent-line margin \(\fitinline{10^{-13}}\). This arithmetic is evidence about these data, not an analytic theorem.

Here is why finite endpoint certification suffices. The ordinary gates, \(\fitinline{H_iF_i}\), tiny gate and derivative gate are affine in \(\fitinline{r}\). For a cutoff inequality, multiply by the positive \(\fitinline{H_i}\) to make it affine before interpolating. The terms \(\fitinline{S_i/H_i}\) and the other fractional-affine gates are handled in the same way. For the max gate, check each affine alternative after clearing its positive denominator. Thus the positive endpoint gates imply the required real-width stage gates on the full closed interval. In particular (B6) implies the HIGH algebra of Part A at the fixed parameters (B3), and the ordinary/reference source gates. The source and witness pads and HIGH reserve remain positive.

The crucial earlier-anchor condition does not require a quadratic number of certificates. If \(\fitinline{b_i\le b_{i+1}}\) and \(\fitinline{L_{i+1}(R_i)\le L_i(R_i)}\), then \[\fitmath{
L_{i+1}(x)\le L_i(x)\quad(x\ge R_i).                      
}\tag{B8}\] On the immediately preceding interval, the last two inequalities of (B7) and affineness give \(\fitinline{L_{i+1}(x)+p_{\rm src}<Y_i(x)}\). Iterating (B8) therefore gives \[\fitmath{
L_n(x)+p_{\rm src}<Y_j(x)
\quad(j<n,\ l_j\le x\le R_j).                            
}\tag{B9}\]

Prove actual paired stage outputs by strong finite induction. Stage zero enters at the native base (B2). At stage \(\fitinline{n>0}\), the previous output at \(\fitinline{l_{n-1}=R_n}\), the second inequality of (B7), and downward restriction give its entry anchor \(\fitinline{\mathcal A(R_n,z_n)}\). If \(\fitinline{x>R_n}\), the already covered earlier intervals select some \(\fitinline{j<n}\) containing \(\fitinline{x}\). Its proved anchor \(\fitinline{\mathcal A(x,Y_j(x))}\) dominates \(\fitinline{L_n(x)}\) by (B9). The top-line condition in (B5) handles the upper zero line. The complete paired analytic step can now be applied on \(\fitinline{[l_n,R_n]}\), yielding \[\fitmath{
\mathcal A(r,Y_n(r))\qquad(l_n\le r\le R_n).               
}\tag{B10}\] Every analytic source input has index strictly less than \(\fitinline{n}\). In particular, none is justified by its own later output.

Finally \(\fitinline{R_{8079}>3/4}\), and the consecutive intervals 8079 through 8191 cover \(\fitinline{[1467/2000,3/4]}\). Let \[\fitmath{
G_0=725738991/12500000000,\qquad
G_1=7599618487/100000000000.
}\] The new comparisons are precisely \[\fitmath{
G_0+10^{-7}-G_1r<Y_i(r)
\quad(i=8079,\ldots,8191,\ r=l_i,R_i).                    
}\tag{B11}\] All 226 rational inequalities were checked independently; the minimum margin is \(\fitinline{280216855113/32000000000000000>0}\). Affineness and (B10) give the paired source on that entire lower line. RESEARCH.md's tangent coefficient enclosure puts the exact tangent below this padded line, completing its lower fiber supplier.

\hypertarget{appendix-g}{}

\hypertarget{appendix-g.-the-exact-contact-from-finite-positive-cost-returns}{%
\subsection{Appendix G. The exact contact from finite positive-cost returns}\label{appendix-g.-the-exact-contact-from-finite-positive-cost-returns}}

This appendix completes the ordinary theorem at the exact contact. It uses the full preceding proof, particularly Appendix E's literal paired sources, averaged floor, exact count and zero-margin scalar return, and Appendices A--D and F's analytic and finite-source foundations. All original coefficients, natural zero sets and character conventions remain those already defined. This is Joseph Bebel's PR143 extension of his PR141 continuum return and Sil's PR137 marked-anchor method, reconstructed from the pinned primary source at commit \texttt{c\allowbreak{}7\allowbreak{}4\allowbreak{}d\allowbreak{}8\allowbreak{}5\allowbreak{}5\allowbreak{}c\allowbreak{}3\allowbreak{}b\allowbreak{}d\allowbreak{}f\allowbreak{}c\allowbreak{}e\allowbreak{}e\allowbreak{}e\allowbreak{}e\allowbreak{}0\allowbreak{}5\allowbreak{}9\allowbreak{}5\allowbreak{}0\allowbreak{}1\allowbreak{}a\allowbreak{}e\allowbreak{}a\allowbreak{}7\allowbreak{}6\allowbreak{}c\allowbreak{}8\allowbreak{}b\allowbreak{}1\allowbreak{}6\allowbreak{}2\allowbreak{}7\allowbreak{}b\allowbreak{}9\allowbreak{}a\allowbreak{}b\allowbreak{}d}. OpenAI, Hailey Collet, ArviFoxy, Joe Richey, Jacob Sussman and inherited contributors retain their credit and original notices. No new mean-value estimate is introduced here.

\hypertarget{g.1.-the-exact-root-and-the-theorem-being-strengthened}{%
\subsubsection{G.1. The exact root and the theorem being strengthened}\label{g.1.-the-exact-root-and-the-theorem-being-strengthened}}

Put \[\fitmath{
f(t)=3888\,19^5(15t-13)^4(1-t)^5-7(7-6t)^4(9t-7)^5.
}\tag{G1}\] At \(\fitinline{13/15}\) it is negative, and at \(\fitinline{7/8}\) it is positive. On this interval all four factors inside the logarithms below are positive. Define \[\fitmath{
H(t)=\log(3888\,19^5)+4\log(15t-13)+5\log(1-t)
 -\log7-4\log(7-6t)-5\log(9t-7).
}\] Then \[\fitmath{
H'(t)=\frac{60}{15t-13}-\frac5{1-t}
 +\frac{24}{7-6t}-\frac{45}{9t-7}
 >480-40-225/4>0.                                       
}\tag{G2}\] Indeed \(\fitinline{15t-13<1/8}\), \(\fitinline{1-t>1/8}\), \(\fitinline{9t-7>4/5}\), and the omitted third term is positive. Thus \(\fitinline{H}\) is strictly increasing. A root of \(\fitinline{f}\) has \(\fitinline{H=0}\), so the intermediate value theorem gives a unique root \(\fitinline{\theta\in(13/15,7/8)}\). This identifies Appendix E's contact \(\fitinline{\alpha}\) with \(\fitinline{\theta}\).

In the tracker's ascending-coefficient convention the primitive polynomial with positive leading coefficient is \(\fitinline{-f}\), with coefficients \[\fitmath{
\begin{aligned}
(&-274959111914881,\ 2643837682856883,\ -11291220876239226,\\
 &28111661257098222,\ -44964545308854309,\ 47916748658170215,\\
 &-34020380857037616,\ 15517896164111592,\ -4126408710053184,\\
 &487371101862528).
\end{aligned}                                             
}\tag{G3}\] Expansion proves this identity. Existence and uniqueness make the infimum of its roots in \(\fitinline{(13/15,7/8)}\) equal to \(\fitinline{\theta}\); the infimum presentation creates no additional analytic assumption.

The conclusion to prove is \(\fitinline{\beta\le\theta}\), where \(\fitinline{\beta}\) is the same Hecke zero supremum as in \S{}2. The ordinary norm pullback will then give \(\fitinline{L(s,\chi)\ne0}\) for \(\fitinline{\Re s>\theta}\), for every fixed positive presented modulus, at every height, except the actual principal pole. It gives no assertion on \(\fitinline{\Re s=\theta}\).

Let \(\fitinline{R_{141}=N/D}\) be Appendix E's exact rational. The independent integer/Fraction certificate \texttt{E\allowbreak{}X\allowbreak{}A\allowbreak{}C\allowbreak{}T\allowbreak{}\_\allowbreak{}C\allowbreak{}O\allowbreak{}N\allowbreak{}T\allowbreak{}A\allowbreak{}C\allowbreak{}T\allowbreak{}\_\allowbreak{}R\allowbreak{}E\allowbreak{}C\allowbreak{}E\allowbreak{}I\allowbreak{}P\allowbreak{}T\allowbreak{}.\allowbreak{}j\allowbreak{}s\allowbreak{}o\allowbreak{}n} encloses \(\fitinline{\theta}\) between exact rationals of width less than \(\fitinline{10^{-330}}\), proves \(\fitinline{\theta<R_{141}}\), and gives \[\fitmath{
\begin{aligned}
1-\theta&=.12665414\allowbreak{}38932814\allowbreak{}23198815\allowbreak{}20393492\allowbreak{}76432295\allowbreak{}58060156\allowbreak{}0188104\ldots,\\
R_{141}-\theta&=3.27684902\allowbreak{}51484275\allowbreak{}31001825\allowbreak{}73402592\allowbreak{}10629791\allowbreak{}20476632\allowbreak{}0\ldots
                 \times10^{-254}>0.
\end{aligned}                                             
}\tag{G4}\] The exact gain is the algebraic number \(\fitinline{R_{141}-\theta}\), not a displayed decimal. PR143 uses exactly the same rational ranking witness as PR141.

\hypertarget{g.2.-a-separate-positive-schedule-for-each-strict-target}{%
\subsubsection{G.2. A separate positive schedule for each strict target}\label{g.2.-a-separate-positive-schedule-for-each-strict-target}}

Fix \[\fitmath{
\theta<t\le T:=87334587/10^8,\qquad g=t-\theta>0.
}\] The retained root isolation gives \(\fitinline{\theta<T}\). Keep the exact tangent, stationary cutoff \(\fitinline{h}\), admitted branch list and count windows from Appendix E. Move the geometric quantities with the target: \[\fitmath{
\ell_t=11/3-4t,\quad y_t=h-2\ell_t,\quad
x_t=1+\ell_t-h,\quad K_t=2t-1.                           
}\tag{G5}\] Appendix E's interval bounds imply throughout this corridor \[\fitmath{
.17328<\ell_t<.17329,\quad
.815324503<h<.815324505,\quad .458<y_t<.473.
}\] Consequently the original literal geometry holds: \(\fitinline{1/6<\ell_t<1/5}\), \(\fitinline{\ell_t<x_t<y_t<1/2}\), \(\fitinline{x_t+y_t=1-\ell_t}\), the signal identity, the dilation and tuple margins, and \(\fitinline{y_t<19/25<2y_t}\). These checks use the averaged floor, rather than the obsolete pointwise floor premise.

The fixed contour reconstruction's row-extension guard is \(\fitinline{F=10^{-258}>0}\). Choose, for this one target, \[\fitmath{
b=\min(F,g/100),\qquad s=b/2,\qquad \Delta=b,\qquad a=b.
}\tag{G6}\] Here \(\fitinline{s}\) is the positive source cost, \(\fitinline{\Delta}\) the positive conductor extension, and \(\fitinline{a}\) the positive row fee. All are at most \(\fitinline{10^{-9}}\), and \(\fitinline{\Delta/2\le F}\). The prime-supply inequality \(\fitinline{\ell_t/(h+\Delta)>1/5}\) follows from the displayed coarse bounds. The symbol \(\fitinline{a}\) in this appendix denotes a fee, not a row amplitude; the latter remains \(\fitinline{p\in[0,1/2]}\).

The actual capacities are \[\fitmath{
C_{o,t}(D)=\frac{\min(D,2y_t-D)}{6K_t},\qquad
C_{s,t}(D)=\min\left(\frac{3D}{40},
                  \frac{2y_t-(1+s)D}{2+10s}\right).       
}\tag{G7}\] For every \(\fitinline{d\in[0,3/4]}\), \(\fitinline{p\in[0,1/2]}\), Appendix E's exact capacity, target-shift, count and extension calculations give \[\fitmath{
\begin{aligned}
\mathcal R_t(D,C_{o,t}(D),d,p)
 &\le\theta-\tfrac32g
       &&(1/200\le D\le19/25),\\
\mathcal R_t(D,C_{s,t}(D),d,p)
 &\le\theta+15s/16-(49/25)g+2\Delta
       &&(19/25\le D\le h+\Delta).
\end{aligned}                                             
}\tag{G8}\] These are not new numerical certificates. They follow from the zero-margin 57-cell scalar theorem at \(\fitinline{\theta}\), retaining the moving denominator \(\fitinline{6K_t}\) in the old capacity, the capacity debit \(\fitinline{5s}\), the exact row shift \(\fitinline{-(2+4dp)g}\), and the quantitative count debit. Convexity in \(\fitinline{D}\) still fills both complete conductor intervals. Substituting (G6) proves both bounds strictly less than \(\fitinline{t-a}\). In particular equality at the stationary scalar contact has acquired a positive budget at this target.

\hypertarget{g.3.-the-positive-costs-reach-the-actual-physical-coefficients}{%
\subsubsection{G.3. The positive costs reach the actual physical coefficients}\label{g.3.-the-positive-costs-reach-the-actual-physical-coefficients}}

It remains to justify an analytic consumer for (G8) with the costs (G6). Throughout this construction assume \(\fitinline{\beta>t>\theta>.873}\); in particular \(\fitinline{\beta\ge.51}\), as required by the physical reserve source. The fixed numbers in Appendix E cannot simply be replaced in a final inequality. The following construction exposes where the smaller budgets enter the already proved physical argument.

\textbf{The genuine source.} Appendix D's reserve source and Appendix E's restriction give, for every fixed \(\fitinline{0<s\le10^{-9}}\), the actual uniform initial energies on \[\fitmath{
(B,\Lambda,c)=
(1-s,\ 2+10s,\ 3/20-1000s/1143).
}\tag{G9}\] Restriction to the physical half-width face uses the smaller cap \(\fitinline{3/40}\). The inverse polynomial retains its M\"{o}bius coefficients; each ordinary factor and each whole canonical prime slot retains its original ray character, common smooth profile, zero extension and independent imaginary height. A reflected ordinary factor is restored to the same row-feature orientation before applying the source. Fixed-inducing inverse rows and the ordinary reference lane are treated by the paired anchor argument. There is no permission to replace a whole prime pool by arbitrary prime subsets or add an arbitrary unsigned rough cofactor. For each positive source loss an early mesh exists before the modulus, profiles and finite slot data. Neither that mesh nor its constants is asserted uniform as \(\fitinline{s\downarrow0}\).

The exact tangent fiber moments, native count windows and their lower patch are unchanged. Restriction of a valid source to a smaller mesh preserves it. Taking a finite minimum therefore gives one early mesh for every count window, the exact tangent fiber moments and (G9). The physical inlet applies (G9) to the original conjugated profile, the fixed source-principal character, and the original fixed ideal; all natural zeros and derivative seminorms remain in that contract. This supplies the reflected initial energy for the actual physical numerator, not merely a scalar source profile.

\textbf{Data, padding and loss order.} Suppose \(\fitinline{\beta>t}\), and put \(\fitinline{\mu=\beta-t>0}\). The natural cap is still \(\fitinline{\kappa_\beta=2\beta-1\le3/4}\), from the original initialization. The original baseline constructor accepts every positive loss budget and every positive extra-mesh function. Request it with the minimum of its old transport budget, \(\fitinline{\mu/10000}\), \(\fitinline{a/4}\), and \(\fitinline{\Delta}\). Its contour spacing can then be decreased after the finite slot partition is fixed. In the original baseline notation \(\fitinline{\varepsilon_D,t_D>0}\), this gives actual finite data with \[\fitmath{
159\varepsilon_D+609t_D\le
        \min(\mu/10000,a/4,\Delta),\quad
t_D\le\Delta/4,\quad \Delta/2+2t_D\le\Delta.             
}\tag{G10}\] The slot lengths are at most one two-hundredth of the common mesh; the finite positive minimum slot length exists after the partition is fixed. The earlier construction of phase, contour, source and floor budgets is retained by also including its positive old bounds in the minimum. Thus taking a smaller target budget causes no circular choice of slot number, arithmetic data or height. Arithmetic source data are then chosen with the same original coefficients and count parameters.

Use actual conductor padding \(\fitinline{\Delta/2}\). Its original fixed upper guard follows from \(\fitinline{\Delta/2\le F}\); (G10) gives the full padded cutoff \(\fitinline{h+\Delta/2+2t_D\le h+\Delta}\). The two complete intervals in (G8) therefore cover every selected physical row. Small conductors, fixed inducing sectors and principal terms keep the complementary estimates already established in the preceding proof.

\textbf{Feedback at the actual supremum.} The source, tangent and geometric choices remain those at \(\fitinline{t}\). The old count and its inverse capacity are evaluated at \(\fitinline{\kappa_\beta}\), rather than frozen at \(\fitinline{K_t}\). Appendix E's quantitative feedback charges at most \(\fitinline{(\kappa_\beta-K_t)/8}\) in the old interval and \(\fitinline{(\kappa_\beta-K_t)/50}\) in the new interval. Since \(\fitinline{\kappa_\beta-K_t=2\mu}\), (G8) yields \[\fitmath{
\begin{aligned}
\mathrm{row}_{\rm old}(\kappa_\beta)-\beta
 &\le-a-3\mu/4,\\
\mathrm{row}_{\rm new}(\kappa_\beta)-\beta
 &\le-a-24\mu/25.
\end{aligned}                                             
}\tag{G11}\] The old allowance includes exactly the lost denominator capacity \(\fitinline{dp\min(D,2y_t-D)(1/K_t-1/\kappa_\beta)/12}\). It is not a consequence of qualitative monotonicity alone.

\textbf{Finite numerator and averaged floor.} The source-energy substitution, finite Fourier labels, whole-slot selections, amplitude bins, masks and contour reconstruction are the original finite operations. Their positive source-face cap now gives exactly the second capacity in (G7). For a nonpositive shortened plain length the original native cone pays the term; it is not passed to a positive-length source theorem. Finite summation retains the common profiles, derivative losses and independent heights. Fourier truncation is removed using the original summable majorants. The full supported nonfloor numerator and the empty-selection floor consequently refer to the same actual physical coefficients.

The scalar packet identity is unchanged. To see explicitly that no fixed fee remains hidden in it, write \(\fitinline{E_{\rm num}}\) for its numerator exponent, \(\fitinline{E_{\rm ker}}\) for its kernel exponent and \(\fitinline{L}\) for its charged source loss. The original identity expresses their sum by the physical row plus the fixed signal shift. Choose the positive contour, phase, source and height losses within \[\fitmath{
16e+L+\mathrm{height}\le\mu/4.                           
}\tag{G12}\] This is possible by (G10) and the original freely reducible positive parameters. A count loss at most \(\fitinline{a/4}\) costs \(\fitinline{3D\,\mathrm{loss}/4\le a/2}\), since \(\fitinline{0\le D\le1}\). Substitution of (G11) into the unchanged identity then gives \[\fitmath{
E_{\rm num}+E_{\rm ker}+L
 <\beta+\mathrm{signalShift}-a/2-\mu/2-\mathrm{height}.
}\tag{G13}\] This is the positive saving consumed by the actual finite numerator. The averaged floor keeps its separate fixed coarse margin from Appendix E, whose literal price is \(\fitinline{51/100-49h/150+51\ell_t/100+33D/50}\). The preceding interval bounds and its \(\fitinline{10^{-5}}\) transport-error budget imply the retained saving \(\fitinline{1/4000}\). The full physical row split subtracts this actual floor, including its error masks. No unsupported row is discarded in applying (G13).

The finite numerator, source moments and transport each first give a finite polynomial height degree. Add those degrees and then apply the original height-choice lemma. One positive height exponent absorbs them into the already fixed positive power budgets. This produces the same complete common probe, with a positive saving and the original nonzero reciprocal numerator, at this target. The common-probe/Mellin contradiction of Appendix B.9 improves the actual zero supremum strictly whenever \(\fitinline{\beta>t}\), an impossibility. Hence \[\fitmath{
\beta\le t\qquad(\theta<t\le T).                         
}\tag{G14}\] Constants and onset scales can depend on this one finite construction, the fixed character and its modulus. This argument takes no limit in any source inequality, sum, integral, height or implicit constant.

\hypertarget{g.4.-the-order-step-and-ordinary-all-moduli-conclusion}{%
\subsubsection{G.4. The order step and ordinary all-moduli conclusion}\label{g.4.-the-order-step-and-ordinary-all-moduli-conclusion}}

If \(\fitinline{\beta>\theta}\), choose \[\fitmath{
t=\min\{(\theta+\beta)/2,(\theta+T)/2\}.
}\] Then \(\fitinline{\theta<t<\beta}\) and \(\fitinline{t\le T}\), contradicting (G14). Therefore \(\fitinline{\beta\le\theta}\). This order argument uses only one finite strict-target construction under the contrary hypothesis; it does not assert a source at zero cost.

For any ordinary character \(\fitinline{\chi}\) presented modulo a positive integer \(\fitinline{q}\), the original norm pullback factors the corresponding Eisenstein Hecke \(\fitinline{L}\)-function into the ordinary factor and its quadratic-twist companion. When \(\fitinline{\theta<\Re s<1}\), both factors are holomorphic, so an ordinary zero would contradict \(\fitinline{\beta\le\theta}\). For \(\fitinline{\Re s\ge1}\) the original classical nonvanishing and pole treatment apply, including the nonprincipal value at \(\fitinline{s=1}\). Euler deletion factors of imprimitive presentations have no zero in the positive half-plane. Thus every presentation, every finite height and the actual principal-pole exclusion in Theorem T are restored. This proves the strict half-plane \(\fitinline{\Re s>\theta}\) for every positive modulus.

\hypertarget{g.5.-source-and-verification-boundaries}{%
\subsubsection{G.5. Source and verification boundaries}\label{g.5.-source-and-verification-boundaries}}

The exact root and uniqueness are scalar facts. The analytic extension is the variable-cost construction in G.2--G.3, which replaces the fixed positive-cost consumer of Appendix E while preserving its proved rational theorem. The source counterpart is PR143's \texttt{P\allowbreak{}o\allowbreak{}s\allowbreak{}i\allowbreak{}t\allowbreak{}i\allowbreak{}v\allowbreak{}e\allowbreak{}T\allowbreak{}r\allowbreak{}a\allowbreak{}n\allowbreak{}s\allowbreak{}p\allowbreak{}o\allowbreak{}r\allowbreak{}t\allowbreak{}D\allowbreak{}a\allowbreak{}t\allowbreak{}a}, \texttt{P\allowbreak{}o\allowbreak{}s\allowbreak{}i\allowbreak{}t\allowbreak{}i\allowbreak{}v\allowbreak{}e\allowbreak{}C\allowbreak{}h\allowbreak{}o\allowbreak{}s\allowbreak{}e\allowbreak{}n\allowbreak{}D\allowbreak{}a\allowbreak{}t\allowbreak{}a}, \texttt{P\allowbreak{}o\allowbreak{}s\allowbreak{}i\allowbreak{}t\allowbreak{}i\allowbreak{}v\allowbreak{}e\allowbreak{}A\allowbreak{}f\allowbreak{}f\allowbreak{}i\allowbreak{}n\allowbreak{}e\allowbreak{}I\allowbreak{}n\allowbreak{}p\allowbreak{}u\allowbreak{}t}, \texttt{P\allowbreak{}o\allowbreak{}s\allowbreak{}i\allowbreak{}t\allowbreak{}i\allowbreak{}v\allowbreak{}e\allowbreak{}S\allowbreak{}c\allowbreak{}h\allowbreak{}e\allowbreak{}d\allowbreak{}u\allowbreak{}l\allowbreak{}e}, \texttt{P\allowbreak{}o\allowbreak{}s\allowbreak{}i\allowbreak{}t\allowbreak{}i\allowbreak{}v\allowbreak{}e\allowbreak{}B\allowbreak{}u\allowbreak{}f\allowbreak{}f\allowbreak{}e\allowbreak{}r\allowbreak{}e\allowbreak{}d\allowbreak{}D\allowbreak{}a\allowbreak{}t\allowbreak{}a}, the literal physical/floor/finite-numerator chain, \texttt{P\allowbreak{}o\allowbreak{}s\allowbreak{}i\allowbreak{}t\allowbreak{}i\allowbreak{}v\allowbreak{}e\allowbreak{}J\allowbreak{}o\allowbreak{}i\allowbreak{}n\allowbreak{}t\allowbreak{}H\allowbreak{}e\allowbreak{}i\allowbreak{}g\allowbreak{}h\allowbreak{}t}, \texttt{P\allowbreak{}o\allowbreak{}s\allowbreak{}i\allowbreak{}t\allowbreak{}i\allowbreak{}v\allowbreak{}e\allowbreak{}S\allowbreak{}o\allowbreak{}u\allowbreak{}r\allowbreak{}c\allowbreak{}e\allowbreak{}C\allowbreak{}l\allowbreak{}o\allowbreak{}s\allowbreak{}u\allowbreak{}r\allowbreak{}e}, and \texttt{E\allowbreak{}x\allowbreak{}a\allowbreak{}c\allowbreak{}t\allowbreak{}E\allowbreak{}n\allowbreak{}d\allowbreak{}p\allowbreak{}o\allowbreak{}i\allowbreak{}n\allowbreak{}t}. \texttt{P\allowbreak{}o\allowbreak{}l\allowbreak{}y\allowbreak{}n\allowbreak{}o\allowbreak{}m\allowbreak{}i\allowbreak{}a\allowbreak{}l\allowbreak{}C\allowbreak{}o\allowbreak{}n\allowbreak{}t\allowbreak{}a\allowbreak{}c\allowbreak{}t} and \texttt{C\allowbreak{}a\allowbreak{}n\allowbreak{}d\allowbreak{}i\allowbreak{}d\allowbreak{}a\allowbreak{}t\allowbreak{}e\allowbreak{}.\allowbreak{}M\allowbreak{}a\allowbreak{}i\allowbreak{}n} identify the canonical exact-root statement. Their original paths, aliases and byte hashes are in the source manifest. PR142 supplies the root grammar and isolation obligation; its parser is not a proof of (G14).

The local calculator checks exact polynomial expansion, rational isolation, strict gain, cost coefficients and feedback identities. These checks establish their stated scalar scope. This intake executed no contributor program, Lean build, kernel checker or provider model. Public GitHub check records report successful PR143 verification and admission for the pinned head; those are public historical evidence, not a locally repeated check. The original full proof and properly proved alternatives to its earlier steps remain available. No other research owner's pause is changed.

\hypertarget{appendix-h}{}

\hypertarget{appendix-h.-the-return-at-1746691720000000}{%
\section{Appendix H. The return at 17466917/20000000}\label{appendix-h.-the-return-at-1746691720000000}}

The retained foundations of Appendices A--G imply the improved bound \(\fitinline{\beta\le17466917/20000000}\). We prove the new signed fixed-row estimate and two-sided raw source bound, return them to the original reciprocal-containing physical integral, and recover the ordinary all-height conclusion. The finite calculations below certify only their stated rational inequalities; no new end-to-end formal verification is asserted.

This appendix preserves the full proof since the original OpenAI argument and permits a properly derived replacement of any earlier step. The original ordinary and inverse moments, arithmetic conventions, Poisson identities, literal probe, restored finite kernels, source axioms and transport results proved or explicitly adopted in A--G remain in force. In particular SN9 is the adopted natural-feature large-sieve input, not a theorem proved here. A catalogue entry, a Lean declaration, a rational check and a draft assertion do not enlarge that foundational boundary.

\hypertarget{appendix-h-1}{}

\hypertarget{h.1.-the-target-and-the-literal-source}{%
\subsection{H.1. The target and the literal source}\label{h.1.-the-target-and-the-literal-source}}

Put \[\fitmath{
 t=17466917/20000000=.87334585,\qquad
 \ell=11/3-4t.
}\tag{H1}\] Let \(\fitinline{\alpha\in(13/15,7/8)}\) be the exact root from Appendix G: \[\fitmath{
 3888\,19^5(15\alpha-13)^4(1-\alpha)^5
 -7(7-6\alpha)^4(9\alpha-7)^5=0.
}\tag{H2}\] The established conclusion used here is \(\fitinline{\beta\le\alpha}\), where \(\fitinline{\beta}\) is the actual supremum of real parts of primitive target zeros. It is not assumed that \(\fitinline{\beta\le t}\). The retained exact contact data give \[\fitmath{
 0<\alpha-t<611/10^{11},\qquad
 .815324503<h<.815324505.
}\tag{H3}\] Use the same stationary \(\fitinline{h}\) as in G, and set \[\fitmath{
 y_t=h-2\ell,\quad x_t=1+\ell-h,\quad
 k=5/6-x_t/3-\ell/6.
}\tag{H4}\] The natural pricing parameter is \(\fitinline{\kappa_*=2\alpha-1}\), which dominates \(\fitinline{2\beta-1}\). The actual baseline source retains its equality \(\fitinline{D.\mathrm{momentKappa}=2\beta-1}\); H.7 explains the monotone price/capacity transfer, so no source-data equality is changed.

All arithmetic is over \(\fitinline{K=\mathbb Q(\sqrt{-3})}\), \(\fitinline{\mathcal O_K=\mathbb Z[\omega]}\). Ideal columns use their primary representatives once; element rows retain all six units. Every presented character and every residue power, including powers zero, six and negative, is zero at a nonunit of its presented modulus. Primitive induction does not erase these holes. The finite bad set, ray characters, projectors and twists are fixed before taking asymptotics. Each row sum has its actual common punctures. A changing support restriction is charged to column width, not hidden in a profile seminorm.

The original prime slots are complete identity-ray annuli. In a fixed normalization they are \[\fitmath{
 Q_i(z)=P_i^{-1/2}\sum_{\mathfrak p\in\mathcal P_i}
 \nu_i(\mathfrak p)c_z(\mathfrak p)
 W_i(N\mathfrak p/P_i)(N\mathfrak p/P_i)^{-33/50+iv_i},
 \qquad
 c_z(I)=\eta(I)\chi_I(\mathfrak m_{\rm bad}^{6}z).
}\tag{H5}\] The coefficient \(\fitinline{\nu_i}\), the whole-prime supplying map and every shared-prime contribution are retained. A selected subset of the primes in a pool cannot replace the pool in the canonical source estimate. Ordinary factors retain their own two rays and independent annular heights; an external common norm twist does not identify those heights.

Here is the relevant quantifier boundary from the earlier appendices. Native marked inverse and ordinary-pair estimates have an early source mesh, then a later finite slot type, profiles and fixed ray dictionary. They give the corresponding \(\fitinline{U^{1+\epsilon}}\) positive energies on actual rows, with finite seminorms and a polynomial in all independent heights. The source is available with arbitrarily small strict width decrements. For SN9, fixed \(\fitinline{A>0}\), \(\fitinline{P\ge3A+1}\), a fixed gap \(\fitinline{0<g<1/2}\) and \(\fitinline{\epsilon>0}\) precede the early mesh. Later finite descriptors have common coefficient vectors of bounded squared mass, supports \(\fitinline{NI\le2Z^{w_i}}\), \(\fitinline{0\le w_i\le\omega}\), total width bounded below, and common lower responses with \(\fitinline{v\sum w_i\le\sum e_i}\). For \(\fitinline{1/2+g\le v\le3/4}\), arbitrary selected natural rows of norm at most \(\fitinline{Y\le Z^A}\), and arbitrary ordinary source columns \(\fitinline{NI\le X\le Z^P}\), the adopted conclusion is \[\fitmath{
 \mathcal N(\mathcal F,\mathcal S)
 \ll Z^\epsilon\{X+Y^{1-v}X^v\}.
}\tag{H6}\] The feature in H6 is the original zero-extended sextic feature. Adjoint, restriction and full conjugation preserve it. A merged vector is obtained by actual ideal convolution; fixed-factor divisor Cauchy costs \(\fitinline{U^\epsilon}\). These facts do not permit division by a character value or removal of its zeros.

Fix the baseline source data \(\fitinline{D,F_{\rm src}}\) and transport data \(\fitinline{P}\) \textbf{before} choosing the target \(\fitinline{\eta}\). Here \(\fitinline{D}\) denotes baseline data, not an ideal divisor or a conductor exponent. The source data object \(\fitinline{F_{\rm src}}\) is distinct from the numerical padding \(\fitinline{F}\) below. Use the fixed pad \[\fitmath{
 F=\Delta=10^{-258},\qquad s=10^{-260},\qquad
 0<D.t\le F/4.
}\tag{H7}\] For a contrary \(\fitinline{\mu=\beta-t>0}\), request baseline loss below \(\fitinline{\min(\mu/10000,10^{-9}/4,F)}\). The literal donor cube has pad \(\fitinline{F/2=5\cdot10^{-259}}\); H7 gives \(\fitinline{F/2+2D.t\le F}\). A smaller variable pad \(\fitinline{10^{-260}}\) would not cover this cube.

\hypertarget{appendix-h-2}{}

\hypertarget{h.2.-signed-centering-on-a-fixed-inducing-row}{%
\subsection{H.2. Signed centering on a fixed inducing row}\label{h.2.-signed-centering-on-a-fixed-inducing-row}}

\textbf{Lemma H.2.} Consider the actual cube-free two-ordinary source on a fixed inducing row, with fixed primitive inducing characters \(\fitinline{\psi_1,\psi_2}\). They may coincide or both be principal. Fix the live prime assignment first and delete its supplying primes in both ordinary variables. Keep all other original holes, seeds, congruences and fixed bad places. Compare ordinary scales \(\fitinline{(X_1,X_2)}\) and \(\fitinline{(Y_1,Y_2)}\) satisfying \(\fitinline{X_1X_2=Y_1Y_2=T}\), using the same two scaled profiles and their same independent normalized heights. If \(\fitinline{L=\min(X_1,X_2,Y_1,Y_2)}\), their normalized signed difference is \[\fitmath{
 \ll_\epsilon U^\epsilon\mathcal S(W_1)\mathcal S(W_2)
 (1+|h_1|+|h_2|)^J\sqrt T\max(1,L)^{-3/5}.
}\tag{H8}\] The implied constants may depend on the fixed character dictionary and source caps. A polynomial norm cap on the combined exceptional radical is required.

\textbf{Proof.} The exact good-prime restricted generating function, after extracting the two ordinary Euler factors, has numerator \[\fitmath{
 H_3(x,y)=1-(x^3+x^2y+xy^2+y^3)
             +(x^3y+x^2y^2+xy^3),
}\tag{H9}\] where \(\fitinline{x=\psi_1(\mathfrak p)(N\mathfrak p)^{-s_1}}\) and \(\fitinline{y=\psi_2(\mathfrak p)(N\mathfrak p)^{-s_2}}\). Thus the complete series is \[\fitmath{
 L(s_1,\psi_1)L(s_2,\psi_2)\,
 \prod_{\mathfrak p\ {\rm good}}H_3(x,y)
 \prod_{\mathfrak p\ {\rm exceptional}} E_{\mathfrak p}(s_1,s_2).
}\tag{H10}\] Remove exceptional Euler factors from the displayed \(\fitinline{L}\)'s first, and let \(\fitinline{E_{\mathfrak p}}\) be the actual polynomial of permitted residual valuations. This prescription covers deletion in just one coordinate and never divides by \(\fitinline{H_3}\) or a character value. On \(\fitinline{\Re s_i\ge2/5}\), the good product converges absolutely, since its first nonconstant terms have degree three. The exceptional product costs \(\fitinline{C^{\omega(\mathfrak f)}\ll_\epsilon U^\epsilon}\) under the stated polynomial radical cap. The ordinary fixed-character \(\fitinline{L}\)-functions have polynomial vertical growth there, apart from their possible principal poles at one. No zero-free reciprocal estimate is needed for this lemma.

Write \(\fitinline{\Phi_i(u)=W_i(u)u^{ih_i}}\). On the first arm the double Mellin weight is \(\fitinline{\widehat\Phi_1(s_1)\widehat\Phi_2(s_2)X_1^{s_1}X_2^{s_2}}\); on the second it is exactly the same pair of transforms with \(\fitinline{Y_i}\). If the incoming height is written as \(\fitinline{(NI)^{it_i}}\), take its outside phase \(\fitinline{X_i^{it_i}}\) \textbf{before} installing a reference arm. With actual prefixes \(\fitinline{d_1,d_2}\), genuine common ideal \(\fitinline{G}\), and common residual height \(\fitinline{v}\), the remaining common phase is \[\fitmath{
       (T/(Nd_1Nd_2NG))^{iv}.
}\tag{H11}\] A profile prefix introduces no ray character unless that ideal was literally extracted from the character's argument. This is necessary for independent \(\fitinline{t_1,t_2}\).

Shift both Mellin lines to \(\fitinline{2/5}\). If both characters are principal, the double residue equals \[\fitmath{
 T\widehat\Phi_1(1)\widehat\Phi_2(1)
 \operatorname{Res}_1 L(\cdot,\psi_1)
 \operatorname{Res}_1 L(\cdot,\psi_2)\mathcal H(1,1),
}\] where \(\fitinline{\mathcal H}\) includes this row's actual punctures. It cancels exactly between the equal-product arms. Different puncture densities in the two coordinates cause no problem: both arms use the same two densities. The single residues and shifted double integral on either arm are bounded by \[\fitmath{
 X_1^{2/5}X_2+X_1X_2^{2/5}+T^{2/5},
}\tag{H12}\] times the height, profile and subpower factors in H8. This follows directly from rapid Mellin decay against the polynomial vertical bounds; choose the finite number of profile derivatives after those degrees are fixed. Divide by \(\fitinline{\sqrt T}\) and use \(\fitinline{X_i,Y_i\ge L\ge1}\) to obtain H8. For \(\fitinline{L<1}\), nonempty fixed annuli have scales bounded below by the reciprocal upper support endpoint, so absolute ideal counting gives \(\fitinline{O(\sqrt T)}\); an empty coordinate contributes zero. Finally sum the frozen marked assignments with their actual coefficient masses \(\fitinline{\ll\prod_i M_i\sqrt{P_i}}\). The deletion and all zeros were preserved before summation. This proves the lemma. \(\fitinline{\square}\)

The stronger exponent \(\fitinline{.7}\) in the direct research note is unnecessary here. The proof above gives exactly the exponent needed below and works for the two independent physical heights and the shortened reference profiles. It does not assert a bound for either positive arm of a doubly principal pair separately.

\hypertarget{appendix-h-3}{}

\hypertarget{h.3.-both-divisor-prices-survive-centering}{%
\subsection{H.3. Both divisor prices survive centering}\label{h.3.-both-divisor-prices-survive-centering}}

The divisor in this subsection is the \textbf{second} artificial M\"{o}bius divisor, the one in H26. The first artificial divisor \(\fitinline{u_{\rm art}}\) in H25 is unweighted. The genuine common code factors \(\fitinline{C_2,D_2}\) below are separate ideals; the font \(\fitinline{\mathfrak D}\) distinguishes the second artificial divisor from them.

Use \(\fitinline{\mathfrak D}\) for the artificial squarefree divisor, \(\fitinline{G}\) for the genuine second common ideal, and \(\fitinline{d_1,d_2}\) for incoming ordinary prefixes. Put \(\fitinline{p_d=Nd_1Nd_2}\). A nonzero seed must divide \(\fitinline{G}\). If \(\fitinline{\mathfrak D}\) meets \(\fitinline{G}\), the residual conditions \(\fitinline{(I,G)=1}\), \(\fitinline{\mathfrak D\mid I}\) make the label zero. Otherwise expand the whole \(\fitinline{G}\)-part into \(\fitinline{\prod B_j=G}\) and the whole \(\fitinline{\mathfrak D}\)-primary part into \(\fitinline{C_j}\), with the original signed nonempty-subset allocation at each prime of \(\fitinline{\mathfrak D}\). Every surviving allocation has \(\fitinline{\mathfrak D\mid\prod C_j}\). Cubefreeness bounds all local valuations. The finite allocation count and its square are \(\fitinline{U^\epsilon}\). The signs, ordinary ray factors, frozen prime weights and height phases are common on the two arms.

Write \(\fitinline{\mathcal X=T\prod_iP_i}\), let \(\fitinline{L_f}\) be the original minimum ordinary scale, and put \[\fitmath{
 q=p_d\prod_jN(B_jC_j)\ge q_0=p_dNGN\mathfrak D.
}\tag{H13}\] The shortened minimum is at least \(\fitinline{L_f/q}\). The coefficient normalization in H8 and the extracted prefixes give amplitude \(\fitinline{\ll\sqrt{\mathcal X}\,q^{-1}\max(1,L_f/q)^{-3/5}}\), with the common height, profile and subpower factors suppressed. For \(\fitinline{q\ge q_0\ge1}\), \[\fitmath{
 q^{-2}\max(1,L_f/q)^{-6/5}
 \le \frac1{q_0\max(1,L_f)}.
}\tag{H14}\] Indeed if \(\fitinline{L_f\le1}\), use \(\fitinline{q^{-2}\le q_0^{-1}}\). If \(\fitinline{1\le q\le L_f}\), the left side is \(\fitinline{q^{-4/5}L_f^{-6/5}\le q^{-1}L_f^{-1}}\); if \(\fitinline{q\ge L_f}\), it is \(\fitinline{q^{-2}\le q^{-1}L_f^{-1}}\). This proves H14 in every case. The only analytic requirement was twice the cancellation exponent at least one.

Consequently the \textbf{unnormalized}, actual per-label row energy has the bound \[\fitmath{
 e_x[y]\ll H U^\epsilon
 \frac{\mathcal X^2}{NGN\mathfrak D\,p_d\max(1,L_f)}
 \mathbf1_{(V_x)\mid y}\,
 \Phi\!\left(\frac{Ny}{NV_xK_x}\right).
}\tag{H15}\] Here \(\fitinline{y=V_xz}\) is the original injective packed row, with no Jacobian and no replacement by a shorter radial ball; \(\fitinline{H}\) denotes the common polynomial height/profile factor. The actual second common ideals are \(\fitinline{C_2=D_2=G}\), whose code ideals divide \(\fitinline{(y)^2}\). They are not the first common divisors. After multiplying by the native weight \(\fitinline{NGN\mathfrak D\,p_dK/\mathcal X^2}\), the remaining multiplicity at fixed \(\fitinline{y}\) is bounded by \[\fitmath{
 \tau((y)^2)^3\sum_{N\mathfrak D\le U^B}\frac1{N\mathfrak D}
 \ll_\epsilon U^\epsilon.
}\tag{H16}\] There is no requirement that \(\fitinline{\mathfrak D\mid y}\). The displayed harmonic ideal sum, not an extra inverse divisor inserted in the raw side, pays its freedom. The \(\fitinline{\rho(G)=2^{\omega(G)}}\) allocation and finite code multiplicities are absorbed in the divisor factor.

Fixed inducing cores have \(\fitinline{O(U^{M/6+\epsilon})}\) sixth-power lifts in a row ball of exponent \(\fitinline{M}\). Thus H15--H16 give the fixed part price \[\fitmath{
 U^{k_0+M/6-\log_U\max(1,L_f)+\epsilon}
}\tag{H17}\] in an actual frame with \(\fitinline{k_0=\log_U K}\), the selected second-row radius exponent. Here \(\fitinline{A_0}\) is the original root-volume exponent, \(\fitinline{c=\log_U NC}\) is the first common ideal exponent in this orientation, \(\fitinline{a=\log_U\mathcal X}\), and \(\fitinline{q_C=\log_U NV_0}\). For the retained frame relations \(\fitinline{a=A_0-c-u'+o(1)}\), \(\fitinline{M=2a-k_0+q_C+o(1)}\), \(\fitinline{M\ge1-6u'-o(1)}\), \(\fitinline{A_0<.903}\), \(\fitinline{u'\le.0006+o(1)}\), and \(\fitinline{\log_U\max(1,L_f)\ge(.139-c)_+-O_N(1/\log U)}\), its excess over \(\fitinline{a+q_C+u'}\) is at most \[\fitmath{
 A_0+3u'-5/6-c-(.139-c)_+<-.0675.
}\tag{H18}\] If \(\fitinline{A_0+3u'\le5/6}\), the volume bound already pays without the minimum saving. Otherwise retain both long coordinates for centering and the short reference for the regular part. A coordinate of scale at most \(\fitinline{2/9}\) with upper profile support \(\fitinline{9/4}\) is exactly empty; bounded scales use the volume estimate. Thus H18 does not silently assume that every shortened scale grows.

\hypertarget{appendix-h-4}{}

\hypertarget{h.4.-the-regular-terminal-input-and-its-finite-inequalities}{%
\subsection{H.4. The regular terminal input and its finite inequalities}\label{h.4.-the-regular-terminal-input-and-its-finite-inequalities}}

We record the specific regular estimate consumed in H.5. For the actual two ordinary rays put \[\fitmath{
 \begin{aligned}
 .42505771&\le n\le.42506657,\\
 1.12796/3-.00004&\le y\le1.12903/3+.00001,\\
 P(V)&=V/2-n,\\
 S(V)&=(1/2+g_0)V-(1+2g_1)y,\\
 W_*&=(P(1)+S(1))/2+.000042,
 \end{aligned}
}\tag{H19}\] where \(\fitinline{.05805911<g_0<.05805913}\), \(\fitinline{.07599618<g_1<.07599620}\). Let \(\fitinline{H_0=.9964}\), \(\fitinline{u_0=.0006}\). The regular selected-row terminal bound on the actual common ball \(\fitinline{2U^{M+o(1)}}\), for live width below \(\fitinline{W_*-u_0}\), is \[\fitmath{
 U^{M+u_0-7\cdot10^{-6}+(5/6)(H_0-M)_++\epsilon}.
}\tag{H20}\] Here ``regular'' excludes the finite inducing classes for both initial ordinary keeps. It does not mean that a merely distinct presentation is nonprincipal.

The proof uses the two separate positive ordinary fourth moments, not a fourth moment of a mixture of characters. On a sixth-free core of relative exponent \(\fitinline{V}\), take the clean maximal prime squared-response exponent \(\fitinline{e}\). With \(\fitinline{L}\) the live width, the two fourth moments give excess \[\fitmath{
 C_0(V)=e\{L-(P(V)+S(V))/2\}.
}\tag{H21}\] Use H21 for \(\fitinline{e\le.4403}\). Repeated use of a maximal slot is allowed by the native source. Rounding a quota loses at most one early mesh per source call; zero repetitions use the unmarked face. For \(\fitinline{.4403\le e\le.51}\), a common buffered zero band for the finite union of the two ray families and prime projectors has \(\fitinline{e\le\nu\le3/4}\) and \(\fitinline{|A_2|^2\ll U^{\nu y+o(1)}}\). The floor cannot produce this amplitude. The \textbf{same attaining zero} supplies the inverse and ordinary witnesses. At \(\fitinline{w=.7575}\), the two alternatives have deficits \[\fitmath{
 D_s=\nu w+e(1-w)/2,\qquad
 D_l=5/6-(5/6-\nu)3w/2.
}\tag{H22}\] For an inverse witness of length at least one use its amplified second moment; below one the relation \(\fitinline{r+m\ge3w/2-o(1)}\) forces \(\fitinline{r}\) or \(\fitinline{2m}\) at least \(\fitinline{w-o(1)}\), and the native width \(\fitinline{(1-w')/2}\) or its ordinary/unmarked fallback gives \(\fitinline{D_s}\). With both log tests included, Cauchy gives excess \[\fitmath{
 C_i(V)=\nu y+e\{L-P(V)/2\}-VD_i/2,\qquad i=s,l.
}\tag{H23}\] If the native inverse requires disjoint prime lists, Newton identities give \(\fitinline{j!e_j(z_p)=(\sum z_p)^j+O_j(1+|\sum z_p|^{j-2})}\), since normalized higher power sums are bounded. Independent random coloring and Jensen then apply the arbitrary-bounded-coefficient inverse source on disjoint lists; the positive row event is fixed before coloring. This is not a change of a canonical prime pool. For \(\fitinline{e\ge.51}\), H6 with \(\fitinline{v=.505}\), fixed gap \(\fitinline{.005}\), and the actual merged coefficient vector pays the bin. Its squared coefficient mass is \(\fitinline{U^\epsilon}\), including shared primes; the relevant column ceiling is \(\fitinline{.900593101004}\), giving a strict saving greater than \(\fitinline{.04010}\).

At \(\fitinline{V=H_0}\), the lower bounds for \(\fitinline{u_0-C_0,u_0-C_s,u_0-C_l}\) are respectively \[\fitmath{
 .000007133217,\quad .000011681338,\quad .000010172505.
}\tag{H24}\] For H21 use the endpoint \(\fitinline{e=.4403}\). For H23 the four vertices \(\fitinline{(\nu,e)=(.4403,.4403),(3/4,.4403),(3/4,.51),(.51,.51)}\) suffice, with the corners of the independent \(\fitinline{n,y,g_0,g_1}\) boxes. Each expression is affine in the first two variables and separately affine in the other four. Their \(\fitinline{V}\)-slopes have magnitude less than \(\fitinline{5/6}\) and are nonpositive, giving H20 on \(\fitinline{[.98,1]}\). For \(\fitinline{V\ge1}\), the native quotas \(\fitinline{1/2-n,1/2-y}\) cover \(\fitinline{L}\). For \(\fitinline{.875\le V<.98}\), the quotas \(\fitinline{V/2-n,V/2-y}\) leave at most \(\fitinline{3(1-V)/8}\), below the allowed \(\fitinline{5(H_0-V)/6}\). For \(\fitinline{.08\le V<.875}\), the two nonprincipal unmarked fourths and the global \(\fitinline{7/8}\) prime bound cost at most \(\fitinline{3L/4}\), which fits at the upper endpoint and by slopes. Below \(\fitinline{.08}\), count lifts directly; their exponent is at most \(\fitinline{M/6+(5/6).08+(3/4)(n+y)+L<M/6+(5/6)H_0}\). Writing the actual row as \(\fitinline{b^6c}\) gives \(\fitinline{O(U^{b_0+\epsilon})}\) cofactors and \(\fitinline{V=M-6b_0}\); this is paid by the same slope bound. No coprimality of \(\fitinline{b,c}\) is imposed, and the factor \(\fitinline{\chi_I(b)^6}\) retains its zeros. These arguments prove H20 from the retained native moments and H6, subject to the finite corners H24.

The actual short reference has short ordinary width at most \(\fitinline{.139}\) and marks at most \(\fitinline{.101}\). Its squared merged column has width \(\fitinline{2(.139+.101)=.480}\); since \(\fitinline{M\ge H_0>2(.480)}\), smooth Poisson gives arbitrary decay for nonprincipal column pairs. Open redundant masks by divisors before Poisson. Principal pairs have matching good sixth-free parts and normalized mass \(\fitinline{\ll U^\epsilon\sum_c(Nc)^{-1}(\sum_j(Nj)^{-3})^2\ll U^{2\epsilon}}\). This proves its short marked fourth, paired by Cauchy with the longer ordinary's unmarked fourth on regular rows. The actual longer reference is below \(\fitinline{.663}\). An arbitrary polynomially long reference is used only in the regular unmarked fourth. Fixed rows are instead paid by H15--H18.

We spell out fixed-profile re-anchoring, including the quantifiers. Let \(\fitinline{[a,b]\subset(0,\infty)}\) be the fixed annulus containing an incoming test, and choose a fixed smooth bump \(\fitinline{V}\) equal to one there. Extend \(\fitinline{F(e^x)e^{\sigma x}}\) smoothly by zero in the log variable. Fourier inversion gives \[\fitmath{
 F(u)u^{ih}=V(u)u^{-\sigma}
       \int_{\mathbb R}\widehat F_\sigma(\xi)u^{i(h+\xi)}\,d\xi,
 \qquad
 \int(1+|\xi|)^J|\widehat F_\sigma(\xi)|\,d\xi
 \ll_{J,\sigma,a,b}\|F\|_{C^{J+2}}.
}\tag{H24a}\] The weighted bound follows by integration by parts \(\fitinline{J+2}\) times and the fixed log-support length. Apply H24a separately in the two ordinary variables, with their own heights, then Minkowski to the positive row norm; the cost is the product of the two finite controls. The first log test is obtained by replacing \(\fitinline{F}\) by \(\fitinline{F\log u}\), with the same support and finite controls. Slot profiles, when separated, get their own independent Fourier variables and polynomial weights. The original slot polynomial and its lower response are kept intact in selecting a row class; no individual Fourier summand is presumed large. An incoming raw phase \(\fitinline{X_i^{ih_i}}\) is factored out before making a reference, as in H11. Native canonical-slot dictionaries use these normalized relative heights, with the actual whole-pool coefficients; reversing orientation conjugates the entire polynomial, profile and height together.

First use the source with all analytic heights at most \(\fitinline{U^\zeta}\), after choosing \(\fitinline{\zeta>0}\) below the fixed buffered allowance and the power reserves divided by their finite degrees. The source's early mesh was already fixed; it does not depend on the later number of slots or on \(\fitinline{\zeta}\). For an outer height above \(\fitinline{U^\zeta}\), absolute ideal counting has some fixed cost \(\fitinline{U^B}\), absorbed in \(\fitinline{(1+|h|)^{B/\zeta}}\). For large auxiliary Fourier variables, apply this same bound under H24a and take sufficiently many derivatives, making the tail any required fixed power. Thus no bound is imposed on a common auxiliary test height, and all two-height/slot-height integrals have finite weighted mass. A finite ray dictionary takes the maximum of finitely many constants and degrees. This proves the exact independent-height/profile extension used in H20 and H27, without uniformity over later varying conductors.

The three exact corner minima in H24 are respectively \(\fitinline{35666085549/5000000000000000}\), \(\fitinline{2920334609/250000000000000}\), and \(\fitinline{762937913731/75000000000000000}\). They follow by substituting the two endpoints of each interval in H19 in H21 and the four stated vertices in H23; subtract each from \(\fitinline{u_0}\) and take the minimum. \texttt{C\allowbreak{}H\allowbreak{}E\allowbreak{}C\allowbreak{}K\allowbreak{}\_\allowbreak{}F\allowbreak{}I\allowbreak{}N\allowbreak{}I\allowbreak{}T\allowbreak{}E\allowbreak{}.\allowbreak{}p\allowbreak{}y} performs these rational substitutions, together with the interval and outer-range guards used above. The earlier exact evidence is \texttt{r\allowbreak{}e\allowbreak{}s\allowbreak{}e\allowbreak{}a\allowbreak{}r\allowbreak{}c\allowbreak{}h\allowbreak{}/\allowbreak{}l\allowbreak{}o\allowbreak{}c\allowbreak{}a\allowbreak{}l\allowbreak{}\_\allowbreak{}c\allowbreak{}y\allowbreak{}c\allowbreak{}l\allowbreak{}e\allowbreak{}0\allowbreak{}7\allowbreak{}9\allowbreak{}/\allowbreak{}d\allowbreak{}i\allowbreak{}r\allowbreak{}e\allowbreak{}c\allowbreak{}t\allowbreak{}/\allowbreak{}A\allowbreak{}R\allowbreak{}I\allowbreak{}T\allowbreak{}H\allowbreak{}M\allowbreak{}E\allowbreak{}T\allowbreak{}I\allowbreak{}C\allowbreak{}.\allowbreak{}j\allowbreak{}s\allowbreak{}o\allowbreak{}n}, key \texttt{.\allowbreak{}0\allowbreak{}0\allowbreak{}0\allowbreak{}0\allowbreak{}4\allowbreak{}2}; the present proof specifies all inputs and the finite operation, and does not assume a decimal table or a contributor program as a theorem.

\hypertarget{appendix-h-5}{}

\hypertarget{h.5.-the-two-raw-sides-and-the-original-marked-root}{%
\subsection{H.5. The two raw sides and the original marked root}\label{h.5.-the-two-raw-sides-and-the-original-marked-root}}

Fix a first common block \(\fitinline{C,D_1}\), its common frequency partition and either of its actual ray orientations \(\fitinline{\omega}\). Let \(\fitinline{V}\) be the original ordinary product times all its marked scales. Write \(\fitinline{u_{\rm art}}\) for the \textbf{first} squarefree M\"{o}bius divisor; \(\fitinline{\mathfrak D}\) is the \textbf{second} one in H13--H16. Denote by \(\fitinline{E^\omega_{\tau,u_{\rm art}}(h)}\) the unnormalized positive first Gauss energy at raw height \(\fitinline{h}\): in (126) use the actual decorated coefficient, the restriction \(\fitinline{u_{\rm art}\mid I}\), the presented side character, and that side's radial window, before division by its column volume. With the amplitude \(\fitinline{\mathcal A}\) and nonnegative allowance of C.10(32), we claim, on each side, \[\fitmath{
 \sum_{u_{\rm art}} E^\omega_{\tau,u_{\rm art}}(h)
 \ll \mathcal A\,\frac{N\mathfrak f_\tau\,V^2}{NC\,ND_1}
 U^{\operatorname{allowance}(C,D_1)+O(\epsilon)}
 (1+|h|)^{2J}.
}\tag{H25}\] The left side has no reciprocal-divisor weight. The opposite side uses its own column pool, inverse ray and reflected row. The bound is for every \(\fitinline{h}\) \textbf{after} summing \(\fitinline{u_{\rm art}}\).

\textbf{Proof of H25.} At the first stage Fourier-invert the entire dyadic kernel (125), including its row and two column windows. At each Fourier point the finite expression is a sum of \(\fitinline{\mu(u_{\rm art})a(z)L_{u_{\rm art},h_L}(z)\overline{R_{u_{\rm art},h_R}(-z)}}\), over \(\fitinline{(u_{\rm art},z)}\); \(\fitinline{|a|\le1}\) after extraction of the displayed physical scalar. Joint Cauchy gives the square roots of the two sums of squared columns. The right Fourier window is conjugated: in the common-height notation its raw heights are \(\fitinline{h+2\pi(s-v)}\) and \(\fitinline{h+2\pi(s+v)}\), respectively. This neither identifies the initial ordinary heights nor moves a supremum inside the divisor sum.

Consider one of these positive sides. Allocate the whole \(\fitinline{C}\)-part of each tuple uniquely as \(\fitinline{\prod B_i=C}\) (use \(\fitinline{D_1}\) on the other side). This gives the identity \(\fitinline{\beta(CI)=\sum_{\mathcal B}F_{\mathcal B}\beta_{\mathcal B}(I)}\) on the actual coprime pool. Each extracted ordinary factor contributes its own \(\fitinline{\omega_i(B_i)}\); a zero stays zero. Shorten both rectangles in coordinate \(\fitinline{i}\) by the \emph{same} \(\fitinline{NB_i}\). The normalized child scalar is exactly \(\fitinline{F_{\mathcal B}\sqrt{V_{\rm live}/V}}\), where \(\fitinline{V_{\rm live}/V=(NC)^{-1}\prod_{\rm frozen}Np_j/P_j}\). Equivalently use the unnormalized child, since \(\fitinline{W_i(Nd_iNI/X_i)=W_i(NI/(X_i/Nd_i))}\), also for the normalized height factor. Do not charge a prefix in both conventions. Allocation Cauchy costs the actual finite multiplicity \(\fitinline{\le\tau(C)^{N+2}}\).

For selected whole slots \(\fitinline{J_s}\), put \(\fitinline{\Pi=\prod_{j\in J_s}P_j=U^{u'}}\). On an injective assignment with ideal product \(\fitinline{\Delta}\) and primary generator \(\fitinline{d_\Delta}\), the fourth-power Gauss identity C.8(13) acts on the \textbf{whole} residual coefficient. Its scalar contains the actual selection amplitude, reciprocity table and Gauss factor. It takes \(\fitinline{z}\) to \(\fitinline{d_\Delta^4 z}\), and \(\fitinline{u_{\rm art}\mid\Delta I}\) to \(\fitinline{v\mid I}\), \(\fitinline{v=u_{\rm art}/(u_{\rm art},\operatorname{rad}\Delta)}\). Gauss zeros enforce the removed coprimalities. There are at most \(\fitinline{2^{|J_s|}}\) preimages of \(\fitinline{v}\); at a nonzero image \(\fitinline{z'}\) the assignment multiplicity is \(\fitinline{\le\tau((z'))^{|J_s|}}\). A nonempty selection kills zero by its actual Gauss factor; the empty selection has one assignment and leaves zero. Cauchy over assignments costs \(\fitinline{\mathcal A_{\rm sel}\ll_N\Pi\prod M_j^2}\), with no further unrestricted sum over assignments. Only now enlarge the residual seed set to all squarefree seeds under its norm cap, in the \textbf{positive} energy, at the actual radius \(\fitinline{K=K_B=16K_{\rm first}(\prod_{j\in J_s}2P_j)^4}\).

Apply second Poisson to this completed sum. Separate frequency zero and the discarded tail, but keep the fixed and nonfixed parts \textbf{together} on each finite retained nonzero set. To see precisely what is being bounded, a retained summand before separation of the columns is \[\fitmath{
 {\bf1}_{\rm part}(z){\bf1}_{(I,I')=1}
 {\rm Corr}(C_2I,D_2I';A_2z)\,
 \beta(C_2I)h_\tau(h,I)\,
 \overline{\beta(D_2I')h_\tau(h,I')}\,
 \mathcal K(A_2z,I,I')\prod_{i=0}^3\mathcal D_i ,
}\] times \(\fitinline{{\bf1}_{v\mid C_2I}{\bf1}_{v\mid D_2I'}}\). The four dyadic weights are on the effective scale, the row, and the two column norms. The two common supports agree; the residuals are coprime to their own common parts and to one another. Consequently these two indicators hold exactly when \(\fitinline{v\mid\operatorname{rad}C_2}\). Their completed sum is \(\fitinline{\rho(C_2)=2^{\omega(C_2)}}\) times the \textbf{same signed summand}, before absolute values. At a common prime the local correlation for unequal multiplicities requires \(\fitinline{6\mid\min(l,r)}\), impossible for \(\fitinline{l,r\in\{1,2\}}\). Thus the surviving \emph{second} factors are \(\fitinline{C_2=D_2=G}\). Open \(\fitinline{{\bf1}_{(I,I')=1}}\) with its full second M\"{o}bius divisor \(\fitinline{\mathfrak D}\); its sign goes into the block scalar, not the coefficients' absolute values.

Here is the arbitrary-row step. In one dyadic four-box expand the finite ray table. Its residual columns are sums over the actual ideals, with coefficients \(\fitinline{{\bf1}_{(G,I)=1,\ \mathfrak D\mid I}\beta(GI)h_{\tau}(h,I)}\) times the corresponding ray and row character; use the opposite ray on the right. Retain the canonical row partition and scalar from the correlation. Let \(\fitinline{x}\) denote a resulting label, \(\fitinline{S_x}\) any subset of its finite nonzero rows. Invert the \emph{whole} compact logarithmic kernel as in (125). Each Fourier term has row multiplier of modulus at most one and norm twists on these two columns; a window on the right has the negative Fourier parameter. More explicitly, after retaining the known nonnegative weight \(\fitinline{b_x}\), the absolute block sum is bounded by integrating terms of the form \[\fitmath{
 \left|\sum_{x,z}{\bf1}_{S_x}(z)b_x a_x(z)
       F_{x,L}(z;h_L)\overline{F_{x,R}(-z;h_R)}\right|
 \le\prod_{\iota=L,R}
 \left(\sum_{x,z\in S_x}b_x
     |F_{x,\iota}(\pm z;h_\iota)|^2\right)^{1/2}.
}\] Here \(\fitinline{|a_x(z)|\le1}\) includes the unit phase of each block used to sum absolute values; insert \(\fitinline{\sqrt{b_x}}\) in both vectors. The notation on the right means \(\fitinline{x}\) ranges over labels and \(\fitinline{z}\) ranges over \(\fitinline{S_x}\); its signs are \(\fitinline{+,-}\). On the shared logarithmic cutoffs the normalized kernels have uniform derivative bounds (125). Integration by parts gives the common pointwise envelope \(\fitinline{D(\xi)=C(1+\|\xi\|)^{-m}}\) for their Fourier densities, with \(\fitinline{m>J+\dim\mathcal K}\) chosen to absorb the required combined degrees. A label-dependent density divided by \(\fitinline{D}\) is a contraction in \(\fitinline{a_x}\); integrate against \(\fitinline{D}\). The accompanying column-window densities satisfy the same weighted estimate. Their norm phases produce the indicated raw heights. Thus a uniform polynomial bound for either \textbf{summed positive side} suffices. This proof applies to arbitrary \(\fitinline{S_x}\) because its indicator multiplies a row, independently of both column indices. Taking the full retained nonzero set and then majorizing by the actual positive radial windows costs no number of labels. The finite ray projector (mass \(\fitinline{256}\)), its ray pairs and the four-box count remain the ones in C.5 and C.8.

On each active decorated column the exponent-zero primes of the first reduction belong to the original residual puncture; they need not divide \(\fitinline{G}\). Off that puncture the exact reciprocity identity (C.8(26)) replaces the moving presentation by its fixed anchor times \(\fitinline{\chi_I(V_0)}\), where \(\fitinline{V_0=V_C}\) or \(\fitinline{V_D}\); unsupported terms vanish. Hence multiplication by \(\fitinline{V_x=V_0A_2}\) injects the appropriate signed row \(\fitinline{z}\) into the common row \(\fitinline{y=V_xz}\). The positive measure becomes exactly \(\fitinline{{\bf1}_{(V_x)\mid y}\Phi(Ny/(NV_xK_x))}\), with \(\fitinline{K_x}\) its actual radial scale. There is no Jacobian. Within each first orientation the grouped envelope uses the raised \textbf{same row ball} for both ordinary factors, with reflection where required; the opposite first orientation uses its own \(\fitinline{V_0}\) and scale. The norm phase of a residual column contributes \(\fitinline{(T/(Nd_1Nd_2\,NG))^{iv}}\) on \emph{both} equal-product arms, besides the shifted heights. All equality statements are on the supported ideals; none divides by a character value.

Only now partition the positive energy into the finite dictionary of fixed inducing cores and its complement, where both ordinary initial-keep predicates hold. Use \(\fitinline{p_d=Nd_1Nd_2}\), \(\fitinline{\mathcal X=U^a}\) for the residual product, \(\fitinline{K=U^{k_0}}\), \(\fitinline{q_C=\log_U NV_0}\) in the chosen orientation. On an active column the tuple annuli give \(\fitinline{c_{\rm low}\mathcal X\le p_d N(GI)}\) and its right counterpart. Thus the outer scalar \(\fitinline{K/(\sqrt{NC_2ND_2}\sqrt{\operatorname{dyad}_2\operatorname{dyad}_3})}\) is at most \(\fitinline{p_dK/(c_{\rm low}\mathcal X)}\); the common normalizer is at most \(\fitinline{NG}\). For the \textbf{unnormalized} energy \(\fitinline{e_x[y]}\) on either side this gives \[\fitmath{
 \rho(G)|p_x|\,\operatorname{outerScalar}\,
                 \operatorname{normalizer}\,e_x[y]
 \ \le\ \frac{\mathcal X}{c_{\rm low}}\,
 \frac{\rho(G)|p_x|}{N\mathfrak D}
 \left\{\frac{NGN\mathfrak D\,p_dK}{\mathcal X^2}\,e_x[y]\right\},
 \qquad |p_x|\le1.
}\tag{H26}\] The reciprocal in the label weight cancels the displayed \(\fitinline{N\mathfrak D}\) in braces. It is not an additional weight in H25. Original physical labels over one frame have multiplicity \(\fitinline{\le\tau(G^2)^2\rho(G)}\), separately from the weighted fiber H16.

On the fixed part H15 bounds the braces by \(\fitinline{HU^\epsilon K/\max(1,L_f)}\) times the displayed sublattice measure. Sum labels at each \(\fitinline{y}\) using H16 and then the \(\fitinline{O(U^{M/6+\epsilon})}\) fixed lifts. This is H17; H18 bounds it by \(\fitinline{HU^{a+q_C+u'+O(\epsilon)}}\), even if a shortened coordinate is bounded. We use centering whenever both original coordinates are long. On a volume-only LOW branch with \(\fitinline{A_0+3u'\le5/6}\) the same calculation needs no minimum saving.

On the regular part separate the two arms by the squared triangle inequality \textbf{after packing}. Take the full \(\fitinline{G}\)- and \(\fitinline{\mathfrak D}\)-primary allocations in H13, including their signs before Cauchy. Ordinary shortening contributes an inverse norm in the squared normalization; each forced prime costs \(\fitinline{M_j^2/P_j\ll M_j^2/Np_j}\). Thus \(\fitinline{\prod B_j=G,\ \mathfrak D\mid\prod C_j}\) give the full reciprocal \(\fitinline{1/(NGN\mathfrak D\,p_d)}\). Apply the forward cube-free identity C.7 to the remaining two variables, with the substitutions \(\fitinline{x_i\mapsto\omega_i(\mathfrak p)x_i}\) also at zero values. In its collision expansion the row multiplier is the product of the two \emph{actual} ordinary polynomials. Remove their punctures by the two M\"{o}bius identities and each slot hole by subtracting the excluded terms. The polynomial radical cap, fixed slot count and absolute central kernel mass in C.7 bound the resulting positive comparison mass by \[\fitmath{
 U^\epsilon\mathcal X\,(NGN\mathfrak D)^\delta/(NGN\mathfrak D\,p_d)
}\] for the raw energy, with arbitrarily small fixed \(\fitinline{\delta>0}\). Its leaves have a surviving subproduct of the original slots and both ordinary lengths shortened below \textbf{one} original or reference pair. They retain the same sublattice. Group by this choice, the two dyadic lengths, slot subset and finite anchor. For each length dilation Fourier-invert \(\fitinline{W_i(e^s)}\); the dilations contribute unit phases, the central ratios are bounded, and the weighted Fourier masses are finite. Cauchy in these integrals permits H16 at each common \(\fitinline{y}\) before applying H20 to the common-profile sum. This pays no free count of labels. The small \(\fitinline{(NGN\mathfrak D)^\delta}\) is absorbed under the polynomial cap. For a filled selection \(\fitinline{u'=u_0+O(\mathrm{mesh})}\), \(\fitinline{M\ge1-6u'-O(\mathrm{mesh})\ge H_0-O(\mathrm{mesh})}\); H20 saves \(\fitinline{7\cdot10^{-6}}\) before the mesh fees. Leave a strict width decrement before selecting, so the live-width hypothesis survives the one-slot rounding. An exhausted selection has no residual marks and uses the unmarked fourths; if the tiny-width floor \(\fitinline{s}\) is used instead, absolute prime bounds cost at most \(\fitinline{3s/4+O(\epsilon)}\), within its preassigned fee. The short reference uses H.4's short-reference estimate. Multiplication by \(\fitinline{K/\mathcal X}\) in braces and \(\fitinline{M=2a-k_0+q_C+O(\mathrm{mesh})}\) now give \(\fitinline{HU^{a+q_C+u'+O(\epsilon)}}\) also on the regular part. The fixed and regular estimates are applied to disjoint positive rows on each side, not substituted for the two sides of Cauchy.

Restore the factor \(\fitinline{\mathcal X/c_{\rm low}}\) of H26 and then \(\fitinline{\mathcal A_{\rm sel}}\). The resulting unnormalized scale on the first side is \(\fitinline{\mathcal X^2NV_C\Pi^2}\); \(\fitinline{\mathcal X\Pi=V_{\rm live}\le2^NV/(NC)}\). The first phase table C.8(21)--(23) gives \(\fitinline{NV_C\,ND_1=Ne\,Nr\,NC}\). On the opposite side the same argument has \(\fitinline{V_{\rm live}\le2^NV/(ND_1)}\) and \(\fitinline{NV_D\,NC=Ne\,Nr\,ND_1}\). Both raw bounds are therefore \(\fitinline{\ll V^2Ne\,Nr/(NC\,ND_1)}\) times their stated fees. The actual presented modulus contains \(\fitinline{er}\), so \(\fitinline{Ne\,Nr\le N\mathfrak f_\tau}\); no primitive-conductor replacement was used.

The remaining second terms precede this nonzero-row collection. With \(\fitinline{Y=\max(1,\operatorname{capRes}_{\mathcal B})}\), the second zero term has its actual coefficient \(\fitinline{K_B|\widehat\Phi(0)|}\) times \(\fitinline{\ll\mathcal S^2Y^{1+\epsilon}/Nv}\), by the zero-correlation divisor bound on the actual seeded columns \(\fitinline{v\mid I,\ NI\le Y}\) in C.9(24). Its harmonic seed sum costs \(\fitinline{O(\log U)}\). The zero-mode comparisons \(\fitinline{5u'<1,\ A_0+4u'<1}\) hold with room since \(\fitinline{u'\le .0006+O(\mathrm{mesh})}\) and \(\fitinline{A_0<.903}\). The discarded tail starts at \(\fitinline{R_B=(Y^2/K_B)U^{m_F/2}}\). Writing \(\fitinline{W_B=\prod_{j\in J_s}2P_j}\), C.9(29)--(30) give \(\fitinline{O(W_B\mathcal S^2\|\Phi\|_{\mathcal S_t}U^{-31})}\) per seed and count \(\fitinline{O(U^8W_B)}\). This pays the \textbf{count} and selected width factors, without a reciprocal-seed tail estimate. An empty residual pool contributes zero. One active decorated coefficient requires a legal supported tuple; the original norm cap and absolute divisor bound used here survive because all added ray factors have modulus at most one. These are the zero and tail terms of C.10(32); the fixed nonzero part is already paid above. Four dyadic indices, allocations and divisor multiplicities cost \(\fitinline{U^{O(\epsilon)}}\). Take one finite seminorm set and height degree over both sides and all surviving subsets; C.5(125) integrates the independent height polynomials on the full windows. Restoring the native nonnegative allowance proves H25. \(\fitinline{\square}\)

\textbf{Recovery of the uncentered, kept root.} At the first stage insert H25 on the joint \(\fitinline{(u_{\rm art},z)}\) space above. Keep the first frequency scalar and ray projector; balancing by \(\fitinline{NC/ND_1}\) preserves the product of the two window budgets. The actual common-band sum C.10(34) bounds the retained energy divided by \(\fitinline{V}\) by \(\fitinline{C\,N\mathfrak f_\eta b_\Phi\,\operatorname{windowMass}_J(0)^2\mathcal A\,U^{1+O(\epsilon)}}\). Here \(\fitinline{N\mathfrak f_\eta}\) is the original fixed presentation factor in that band sum, not \(\fitinline{N\mathfrak f_\tau}\) declared fixed. The first zero and discarded terms obey C.10(35) at the same normalization: the legal decorated tuples have the same product cap and coefficient majorant and \(\fitinline{V<U^{.903}}\). The conclusion controls the \textbf{whole centered square} on all element rows, including the original two signed rectangles with the same profiles, for every admissible shortening.

Fix now one original retained physical sixth-free class \(\fitinline{\mathcal C}\), including its error mask, target and contour point. A positive fixed ball cutoff \(\fitinline{\Phi}\) majorizes its shell. At growing norm both original initial keeps hold: a sixth-free ideal supported on any fixed finite set has bounded norm, and a prime outside that set ramifies in the natural character. Fixed ray twists and the fixed punctures cannot remove this prime. For two independent ordinary profiles and heights let \(\fitinline{B_{z,i}(X_i)}\) have its actual \(\fitinline{X_i^{-1/2}}\) normalization, and \(\fitinline{D_z(I)=\prod_{j\in I}Q_j(z)}\) for a subproduct of whole original slots. We bound \[\fitmath{
 \sum_{z\in\mathcal C}\Phi\,|B_{z,1}(X_1)B_{z,2}(X_2)D_z(I)|^2. \qquad\text{(A)}
}\] If \(\fitinline{\min X_i\le U^{.139}}\), apply Cauchy directly: \[\fitmath{
 \left(\sum_{\mathcal C}\Phi |B_{\rm long}|^4\right)^{1/2}
 \left(\sum_{\text{all }z}\Phi |B_{\rm short}D_z(I)|^4\right)^{1/2}. \qquad\text{(B)}
}\] The first factor uses its own ordinary unmarked fourth on its nonprincipal keep. For the second the squared merged column has width \(\fitinline{b\le2(.139+.101)=.480}\) and coefficients \(\fitinline{\ll U^\epsilon/\sqrt{NI}}\) by divisor counting. In a nonprincipal column pair the reduced conductor times an opened redundant-mask divisor has norm \(\fitinline{\ll U^{2b+O(\mathrm{mesh})}}\); smooth Poisson at radius \(\fitinline{\asymp U}\) has zero constant term and arbitrary decay. Principal pairs share their good sixth-free part \(\fitinline{c}\) and have mass \(\fitinline{\ll U^\epsilon\sum_{Nc\ll U^b}(Nc)^{-1}(\sum_j(Nj)^{-3})^2}\). Hence (B) is \(\fitinline{\ll U^{1+O(\epsilon)}}\) with the stated profiles/heights. A coordinate of scale at most \(\fitinline{2/9}\) is empty.

Otherwise open the inverse cube-free identity C.7 on the full product in (A). The local bidegree \(\fitinline{(a,b)}\) has factor \(\fitinline{\omega_1(\mathfrak p)^a\omega_2(\mathfrak p)^b}\) under substitution in the formal power series, preserving zero-extended ray coefficients of all positive degrees. Both the inverse and forward identities keep their square-zero mark variables, collision terms and full hole subtractions. Their absolute central masses cost \(\fitinline{U^\epsilon}\) under the stated caps (Minkowski before squaring); a surviving coordinate or mark list only shortens. A moving row phase or zero indicator is extracted \textbf{exactly} before its modulus is bounded. On a short restricted child use the forward identity and (B). For a child with both coordinates long, take on \textbf{both sides of the first kernel} the same reference \(\fitinline{(Y_1,Y_2)=(U^{.139},X_1X_2/U^{.139})}\). The longer reference is \(\fitinline{\le U^{.663+O(\mathrm{mesh})}}\) here. Put any incoming \(\fitinline{(NI_i)^{it_i}}\) into the normalized profile first, with its outside \(\fitinline{X_i^{it_i}}\) unchanged on the reference. The signed residual phase stated above is then the same on both arms. For the actually restricted row \(\fitinline{R_X}\) and its reference \(\fitinline{R_Y}\), \[\fitmath{
 \sum_{\mathcal C}\Phi |R_X|^2
 \le 2\sum_{\mathcal C}\Phi |R_Y|^2+
      2\sum_{\text{all }z}\Phi |R_X-R_Y|^2. \qquad\text{(C)}
}\] The forward identity and (B) pay the first term; the full centered square just proved pays the second. Restore the inverse comparison over all its children. This proves (A) on the original keep without a separate all-row estimate for either long positive arm.

On that class the reflected numerator is exactly the dyad \(\fitinline{P_z(n,v)}\): its ordinary weight is \(\fitinline{X^{-1/2}X^{-iv}W(NI/X)(NI/X)^{-1/2-iv}}\). The denominator ordinary dyad \(\fitinline{S_z(y,t')}\) has coefficient \(\fitinline{f(I)=\eta_{\rm ray}(I)\overline{\chi_I(z)}}\) of H34, with its original Euler deletions. If orientations differ, conjugate this \textbf{entire} factor, including ray, profile and height. The relative anchor in every \(\fitinline{Q_j}\) is the original one in H5; moving a common height into the relative phase leaves its actual outside unit \(\fitinline{P_j^{ih}}\). The two ordinary heights remain independent. The ordinary-witness length range \(\fitinline{m-4\lambda<y<m+\lambda}\) of H.6 and the reflected-numerator window in the actual physical ticket H37 lie within H19's separate length caps; every componentwise shortening in (A) is permitted. For a row-dependent scale first group \(\fitinline{X_z\in[X_0/(9/8),X_0]}\) (similarly for the other coordinate). Its original \(\fitinline{[1/2,2]}\) profile is supported on \(\fitinline{NI/X_0\in[4/9,2]}\). Insert a common bump equal to one there, supported in \(\fitinline{(1/4,9/4)}\); log-Fourier inversion of the original dilation leaves a unit Fourier phase and central scalar of modulus \(\fitinline{\le\sqrt{9/8}}\). Weighted Fourier masses are finite original seminorms. Apply the original absolute packet bound B.7(222) on each physical row, including its genuine scale/dilation damping, before taking the common norm. There are only \(\fitinline{O(\log U)}\) scale bins per coordinate. The argument also applies to the first normalized log tests, since \(\fitinline{\partial_{t'}}\) inserts \(\fitinline{i\log(NI/Y)}\) in the denominator and likewise in the normalized numerator test. Thus, with \(\fitinline{T}\) bounding the other physical frequencies, for any original selected subproduct below \(\fitinline{W_*}\) with a strict decrement, \[\fitmath{
 \sum_{z\in\mathcal C}|P_z(n,v)S_z(y,t')\prod_{j\in I}Q_j(z)|^2
 \ll U^{1+\epsilon}(1+|v|+|t'|+T)^J .
}\tag{H27}\] Finite profile seminorms are understood. To remove marks, on the same frozen amplitude class sort its distinct whole slots in decreasing order of lower logarithmic gain per unit width. Its total original width is \(\fitinline{>1/5>W_*}\) and its average lower gain is \(\fitinline{q=dp}\), up to the bin error. A prefix reaching \(\fitinline{W_*}\) minus the strict decrement has at least this average and loses at most one slot to rounding; if positive slots run out, their total gain already exceeds \(\fitinline{qW_*}\) up to the bin error. Divide H27 by the \textbf{squared} common lower response on this class. The resulting comparison, including its first log tests, costs \(\fitinline{U^{1-2qW_*+O(\mathrm{mesh}+\mathrm{decrement}+\epsilon)}(1+|v|+|t'|+T)^J}\). No repeated slot used to prove H20 is used in this physical lower product. Fix the source losses, decrement and early mesh before the later finite slot dictionary and target \(\fitinline{\Gamma}\); then choose uniform seminorm and height degrees. Choose the one small physical height power \(\fitinline{\tau_\Gamma>0}\) after these degrees and before the arbitrary positive buffered height index. H27 is at each fixed independent height; H.6--H.8 use only the one ordinary witness actually supplied on each retained row. \(\fitinline{\square}\)

\hypertarget{appendix-h-6}{}

\hypertarget{h.6.-full-inverse-witnesses-and-the-donor-exclusions}{%
\subsection{H.6. Full inverse witnesses and the donor exclusions}\label{h.6.-full-inverse-witnesses-and-the-donor-exclusions}}

We give the needed exclusions at the ordinary source contract; they are not imported as a theorem from the public draft. Fix an actual density/amplitude bin \[\fitmath{
 .425\le d\le.4267,\qquad .49995\le p\le1/2,
 \qquad \lambda=10^{-5},\quad \eta_0=10^{-7}.
}\tag{H28}\] The retained subpatch later has \(\fitinline{p>.499996}\). Set \[\fitmath{
 \begin{aligned}
 a&=p(1+2g_0),& k_1&=1-p(1+2g_1),\\
 V_0&=5/6-d,& U_0&=1-p(1+6g_0-4g_1),\\
 J_0&=3V_0+2dk_1,& t_0&=1+dU_0/J_0,\\
 b&=2t_0/3,&m&=t_0/3.
 \end{aligned}
}\tag{H29}\] The exact crossing identities are \[\fitmath{
 R=1-d(a+k_1b)=1/6+(5/6-d)t_0.
}\tag{H30}\] The tangent source on doubled ordinary length \(\fitinline{2y}\) or inverse length \(\fitinline{x}\), in the anchor interval \(\fitinline{[1467/2000,311/400]}\), gives counts \(\fitinline{U^{1-d(a+2k_1y)+\epsilon}}\) and \(\fitinline{U^{1-d(a+k_1x)+\epsilon}}\), respectively, after dividing by the available saturated response. Therefore inverse lengths \(\fitinline{x\ge b+\lambda}\) and ordinary lengths \(\fitinline{y\ge m+\lambda}\) save at least \(\fitinline{dk_1\lambda}\) and \(\fitinline{2dk_1\lambda}\), until the anchor endpoint.

Above the anchor, the inverse count up to \(\fitinline{.90}\) has coefficient \(\fitinline{x+p(1-x)}\), from native width \(\fitinline{(1-x)/2}\) with a strict decrement. Its gates are \(\fitinline{x+2w<1}\), \(\fitinline{2x+8w<3}\); the latter is \(\fitinline{4-2x<3}\), valid here. A small interior decrement gives the first strict inequality. The unmarked inverse gives coefficient \(\fitinline{x}\) from \(\fitinline{.90}\) to 1. From 1 to \(\fitinline{t_0-\lambda}\), the amplified inverse gives count exponent \(\fitinline{1/6+(5/6-d)x}\), below H30 by at least \(\fitinline{(5/6-d)\lambda}\).

For ordinary \(\fitinline{y}\) above half the anchor endpoint, use \[\fitmath{
 B=1-10^{-12},\quad \Lambda=2+10^{-11},\quad
 c_* =3/20-10^{-7},\quad s_-=.027,
 \quad y_*=(1-(9/2)s_-)/2=.43925.
}\] Up to \(\fitinline{y_*}\), the ordinary count-face width is \(\fitinline{\min(c_*,(B-2y)/\Lambda)}\), with savings coefficient \(\fitinline{2y+2p\min(c_*,(B-2y)/\Lambda)}\). From \(\fitinline{y_*}\) to \(\fitinline{.4398}\), use native width \(\fitinline{(1-2y)/(9/2)}\), whose lower bound is greater than \(\fitinline{.02}\). Thereafter the unmarked fourth gives coefficient \(\fitinline{2y}\). These coefficients are increasing in \(\fitinline{y}\) on each piece. The endpoints that must exceed the crossing coefficient \(\fitinline{C=a+k_1b}\) are \[\fitmath{
 \begin{gathered}
 r_1+p(1-r_1),\quad .90,\quad
 r_1+2pc_*,\quad r_1+2p(B-r_1)/\Lambda,\\
 2y_*+2ps_-,\quad .8796,
 \qquad r_1=311/400.
 \end{gathered}
}\tag{H31}\] Their associated count margins are multiplied by \(\fitinline{d}\).

Here is a deliberately conservative rational verification. Use independent enclosing endpoints \(\fitinline{g_{0,-},g_{0,+},g_{1,-},g_{1,+}}\) from H19 and the full rectangle H28. Bound \(\fitinline{t_0}\) by substituting the lower numerator/upper denominator and conversely in H29, then put \(\fitinline{C_+=p_+(1+2g_{0,+})+k_+b_+}\). For H31 subtract \(\fitinline{C_+}\), use \(\fitinline{p_-}\) in positive terms and multiply by \(\fitinline{d_-}\). The resulting nine raw margins (including the three moving endpoints) all exceed \(\fitinline{10\eta_0}\); the smallest is \[\fitmath{
 d_-k_-\lambda=1.80201615\cdot10^{-6}>10^{-6}.
}\tag{H32}\] The anchor contains \(\fitinline{[b_-,b_++2\lambda]}\). The affine and native slopes are positive because \(\fitinline{2-4p_+/\Lambda>0}\) and \(\fitinline{2-4p_+/(9/2)>0}\). \texttt{C\allowbreak{}H\allowbreak{}E\allowbreak{}C\allowbreak{}K\allowbreak{}\_\allowbreak{}F\allowbreak{}I\allowbreak{}N\allowbreak{}I\allowbreak{}T\allowbreak{}E\allowbreak{}.\allowbreak{}p\allowbreak{}y} records the exact fractions and substitutions for H31--H32 in \texttt{F\allowbreak{}I\allowbreak{}N\allowbreak{}I\allowbreak{}T\allowbreak{}E\allowbreak{}.\allowbreak{}j\allowbreak{}s\allowbreak{}o\allowbreak{}n}; its code is an independent arithmetic calculation and executes no donor program. This is a finite proof of the displayed inequalities, not a proof of the input moments.

The saturated squared-response debit is at most \(\fitinline{3\eta_0}\) in a second moment and \(\fitinline{6\eta_0}\) in a fourth. Choose the remaining strict source decrements, amplitude bin widths, dyadic, profile and height losses below the unused part of H32. The fixed-profile native sources apply: the inverse anchor annulus is separated from its truncation boundary; Mellin separation reduces the ordinary profiles to the original fixed annular tests; independent norm powers are handled by the two log tests and Sobolev. A union over the finite dictionary, logarithmic dyads and Fourier truncations costs \(\fitinline{U^\epsilon}\). Hence the exceptional rows admitting a saturated inverse in \(\fitinline{[b+\lambda,t_0-\lambda]}\), or a saturated ordinary of length at least \(\fitinline{m+\lambda}\), are \[\fitmath{
 O(U^{R-\eta_0}).
}\tag{H33}\] The exceptional set is a set of original rows fixed before the positive physical marked fourth estimate. Applying that estimate on this smaller set saves \(\fitinline{3\eta_0/4}\) in physical mass. It is not a new full-class mixed count.

For completeness the required factorization is an identity of the physical denominator, independently of which twist supplied the original detector zero. Let \(\fitinline{f=\psi_{\rm den}}\), \(\fitinline{h=\mu\psi_{\rm den}}\), with \(\fitinline{h*f=\mathbf1}\), all with their actual Euler deletions. Take \(\fitinline{h_A}\) to be the smooth cutoff at \(\fitinline{A=U^{t_0-2\lambda}}\). The coefficients of \(\fitinline{f*h_A-\mathbf1}\) vanish below \(\fitinline{A}\), and \[\fitmath{
 h=h_A-(f*h_A-\mathbf1)*h.
}\tag{H34}\] On the long annulus \(\fitinline{h_A}\) vanishes. Insert a smooth cutoff equal to zero through \(\fitinline{A/2}\) and one from \(\fitinline{A}\), split the three factors dyadically, and Mellin-separate total norm and first-two-factor norm. The two Fourier variables have bounded weighted \(\fitinline{L^1}\) norms. Their heights are the actual shared inverse/ordinary shifts \(\fitinline{t+\xi_1+\xi_2}\), and the short inverse height \(\fitinline{t+\xi_1}\). The original numerator keeps its independent height \(\fitinline{v}\). A large full inverse therefore supplies at least one saturated triple with \[\fitmath{
 x\le t_0-2\lambda+o(1),\quad x+y\ge t_0-2\lambda-o(1),
 \quad x+y+z=r+o(1).
}\] Here is the saturation argument and its precise loss. On the retained nonprincipal physical denominator, the common buffered zero-free band of B.5--B.6 bounds each normalized ordinary or inverse dyad of length \(\fitinline{U^a}\) by \(\fitinline{U^{da/2+\epsilon}}\), uniformly on the permitted heights and fixed profile family. Indeed shift its Mellin integral for \(\fitinline{L(s,\psi_{\rm den})}\) or \(\fitinline{1/L(s,\psi_{\rm den})}\) to \(\fitinline{\Re s=(1+d)/2+\epsilon}\); reciprocal control and the buffer give polynomial height growth, and the rapid transform supplies the integral bound. For the buffered estimate B.5(206), write \(\fitinline{a_{\rm band}=(1+d)/2}\). Choose its buffer \(\fitinline{e>0}\) first sufficiently small, and then choose the displacement \(\fitinline{\epsilon\ge6e}\), so that \(\fitinline{a_{\rm band}+\epsilon\ge a_{\rm band}+6e}\) is in the permitted strip. All descendant lengths are bounded by \(\fitinline{2.26+o(1)}\); choose these preliminary parameters so that their bounded-length shift costs fit inside the fixed witness tolerance below. The principal class and heights outside this buffer are among the separately paid complements, not silently included in this contour shift. The uniformly smooth factors \(\fitinline{V(NI/A)}\) and \(\fitinline{G(NI/A)}\) have bounded log derivatives on each dyad (they are constant outside their transitions); H24a therefore re-anchors them with the original fixed tests and independent common shifts.

There are \(\fitinline{O((\log U)^3)}\) triples. Their Fourier masses and the absolute separation tails are bounded as just proved; truncate each new shift at \(\fitinline{T_1/16}\), keeping the original height at most \(\fitinline{T_1/4}\), then choose the finite tail order. Starting even from the weaker full-inverse response \(\fitinline{U^{dr/2-\eta_0}}\), the triangle inequality supplies some triple and frequencies with product at least \(\fitinline{U^{d(x+y+z)/2-\eta_0-o(1)}}\). Divide by the two companion upper bounds \(\fitinline{U^{da/2+\epsilon}}\). Choosing all preliminary losses below the spare \(\fitinline{\eta_0/2}\) proves, for each of its three factors, a lower magnitude \(\fitinline{U^{da/2-3\eta_0/2}}\). Thus each normalized \textbf{squared} response has debit \(\fitinline{3\eta_0}\). This proves an existential simultaneous witness; it neither restricts the original full convolution to a selected region nor assigns largeness to an arbitrary Fourier term.

Remove H33. The remaining triple has \(\fitinline{b-4\lambda<x<b+\lambda}\), \(\fitinline{m-4\lambda<y<m+\lambda}\), with its actual ordinary witness satisfying \[\fitmath{
 |S_z(y,t_z)|\ge U^{dy/2-3\eta_0/2}.
}\tag{H35}\] This says that a witness exists for each retained full inverse. It does not bound a whole convolution by restricting it to the existence of one witness.

There is also no hidden two-inverse exclusion in this argument. If the second detector is used, choose cutoff \(\fitinline{2t_0-5\lambda<2.26}\). Its original \(\fitinline{U^{20}}\) truncation still works, since \(\fitinline{-3+2.26(5/4-.51)<-1}\). It gives a long inverse length \(\fitinline{r}\) and outer ordinary \(\fitinline{y_2}\) with \(\fitinline{r+y_2\ge2t_0-5\lambda-o(1)}\). Outside H33, \(\fitinline{y_2<m+\lambda}\), so the inverse annulus is strictly between \(\fitinline{2A}\) and \(\fitinline{A^2}\). The exact identity \[\fitmath{
 h-2h_A+h_A*f*h_A=(h-h_A)*f*(h-h_A)
}\tag{H36}\] has right side zero below \(\fitinline{A^2}\). On that annulus it gives \(\fitinline{h=-h_A*f*h_A}\). Apply the same common-frequency separation. Each inverse descendant is below \(\fitinline{t_0-\lambda}\); H33 removes any descendant above \(\fitinline{b+\lambda}\), or inner ordinary above \(\fitinline{m+\lambda}\). Subtracting the other three upper bounds from the total gives each lower bound \(\fitinline{b-9\lambda}\) or \(\fitinline{m-9\lambda}\), including support errors for large \(\fitinline{U}\). These exclusions need only the single-factor ordinary/inverse contracts just proved. They do \textbf{not} prove the separate two-inverse mixed moment proposed in the draft. That proposed moment is not used in H27 or below.

\hypertarget{appendix-h-7}{}

\hypertarget{h.7.-the-complete-finite-ledger-and-its-physical-consumer}{%
\subsection{H.7. The complete finite ledger and its physical consumer}\label{h.7.-the-complete-finite-ledger-and-its-physical-consumer}}

We return the marked two-ordinary estimate H27 to the physical integral, using the original sources and count admissions of A--G.

Write \(\fitinline{D_c=\log_ZU}\), \(\fitinline{q=dp}\), \(\fitinline{\varepsilon_p=1/2-p}\), \(\fitinline{b=2\cdot10^{-6}}\), and \(\fitinline{\varepsilon_0=4\cdot10^{-6}}\). Keep the stationary \(\fitinline{h}\), \(\fitinline{\ell_t=11/3-4t}\), \(\fitinline{y_t=h-2\ell_t}\), and \(\fitinline{0<\xi=\alpha-t<\xi_0=611/10^{11}}\). The natural pricing parameter \(\fitinline{\kappa_*=2\alpha-1}\) is an upper parameter, while the actual source data retain \(\fitinline{D.\mathrm{momentKappa}=2\beta-1}\). Under \(\fitinline{t<\beta\le\alpha}\), FB2--FB6 imply \(\fitinline{R_O(2\beta-1)\le R_O(\kappa_*)}\), and the actual natural capacity is at least the capacity priced at \(\fitinline{\kappa_*}\). Enhanced branches are independent of \(\fitinline{\kappa}\). Thus all prices below apply to the actual source without altering its data equality.

The complete returned exponent is \[\fitmath{
 E(D_c,C;d,p)=\tfrac12+13h/75-y_t/2-\ell_t/2-9D_c/100
       +d/2+dp(\ell_t-C/2)+(3D_c/4)R.                  
}\tag{H37a}\] Here \(\fitinline{R}\) is the conductor-independent minimum of precisely the admitted original branches. In the critical comparisons H30, H37a on E(27), and H42, one may instead put \(\fitinline{R=R_b=1-d[p(1+2g_0)+2k_1m]}\), the particular admissible count upper bound at the crossing in H29--H30. The actual count is bounded by this branch whether or not it minimizes the envelope, and E(49) prices the physical return with this very branch. No equality between the minimum and the tangent branch is required. On \(\fitinline{[1/200,19/25]}\) use \(\fitinline{C_O=\min(D_c,2y_t-D_c)/(6\kappa_*)}\); on \(\fitinline{[19/25,h]}\) the native capacity is \(\fitinline{C_N=\min(3D_c/40,(2y_t-D_c)/2)}\).

\hypertarget{exact-finite-data-and-the-original-57-boxes}{%
\subsubsection{Exact finite data and the original 57 boxes}\label{exact-finite-data-and-the-original-57-boxes}}

The attached finite objects are named in the packed proof as follows:

\begin{itemize}
\tightlist
\item
  \texttt{r\allowbreak{}e\allowbreak{}s\allowbreak{}e\allowbreak{}a\allowbreak{}r\allowbreak{}c\allowbreak{}h\allowbreak{}/\allowbreak{}l\allowbreak{}o\allowbreak{}c\allowbreak{}a\allowbreak{}l\allowbreak{}\_\allowbreak{}c\allowbreak{}y\allowbreak{}c\allowbreak{}l\allowbreak{}e\allowbreak{}0\allowbreak{}6\allowbreak{}9\allowbreak{}/\allowbreak{}c\allowbreak{}o\allowbreak{}n\allowbreak{}t\allowbreak{}i\allowbreak{}n\allowbreak{}u\allowbreak{}u\allowbreak{}m\allowbreak{}/\allowbreak{}R\allowbreak{}E\allowbreak{}T\allowbreak{}U\allowbreak{}R\allowbreak{}N\allowbreak{}\_\allowbreak{}C\allowbreak{}E\allowbreak{}R\allowbreak{}T\allowbreak{}I\allowbreak{}F\allowbreak{}I\allowbreak{}C\allowbreak{}A\allowbreak{}T\allowbreak{}E\allowbreak{}.\allowbreak{}j\allowbreak{}s\allowbreak{}o\allowbreak{}n}: all 57 original closed rectangles, branch assignments and exact \texttt{l\allowbreak{}o\allowbreak{}w\allowbreak{}e\allowbreak{}r\allowbreak{}\_\allowbreak{}p\allowbreak{}o\allowbreak{}l\allowbreak{}y\allowbreak{}n\allowbreak{}o\allowbreak{}m\allowbreak{}i\allowbreak{}a\allowbreak{}l\allowbreak{}s}; SHA256 \texttt{8\allowbreak{}c\allowbreak{}3\allowbreak{}b\allowbreak{}0\allowbreak{}e\allowbreak{}d\allowbreak{}f\allowbreak{}6\allowbreak{}c\allowbreak{}2\allowbreak{}8\allowbreak{}4\allowbreak{}c\allowbreak{}1\allowbreak{}3\allowbreak{}4\allowbreak{}e\allowbreak{}5\allowbreak{}c\allowbreak{}a\allowbreak{}e\allowbreak{}0\allowbreak{}4\allowbreak{}0\allowbreak{}5\allowbreak{}c\allowbreak{}9\allowbreak{}8\allowbreak{}2\allowbreak{}3\allowbreak{}e\allowbreak{}8\allowbreak{}7\allowbreak{}8\allowbreak{}a\allowbreak{}9\allowbreak{}f\allowbreak{}c\allowbreak{}5\allowbreak{}a\allowbreak{}8\allowbreak{}4\allowbreak{}7\allowbreak{}0\allowbreak{}4\allowbreak{}c\allowbreak{}0\allowbreak{}0\allowbreak{}8\allowbreak{}4\allowbreak{}4\allowbreak{}0\allowbreak{}7\allowbreak{}1\allowbreak{}3\allowbreak{}5\allowbreak{}d\allowbreak{}c\allowbreak{}6\allowbreak{}5\allowbreak{}7\allowbreak{}d\allowbreak{}5}.
\item
  \texttt{r\allowbreak{}e\allowbreak{}s\allowbreak{}e\allowbreak{}a\allowbreak{}r\allowbreak{}c\allowbreak{}h\allowbreak{}/\allowbreak{}l\allowbreak{}o\allowbreak{}c\allowbreak{}a\allowbreak{}l\allowbreak{}\_\allowbreak{}c\allowbreak{}y\allowbreak{}c\allowbreak{}l\allowbreak{}e\allowbreak{}0\allowbreak{}7\allowbreak{}9\allowbreak{}/\allowbreak{}n\allowbreak{}e\allowbreak{}a\allowbreak{}r\allowbreak{}\_\allowbreak{}r\allowbreak{}e\allowbreak{}c\allowbreak{}o\allowbreak{}r\allowbreak{}d\allowbreak{}/\allowbreak{}R\allowbreak{}E\allowbreak{}T\allowbreak{}U\allowbreak{}R\allowbreak{}N\allowbreak{}\_\allowbreak{}F\allowbreak{}I\allowbreak{}N\allowbreak{}A\allowbreak{}L\allowbreak{}.\allowbreak{}j\allowbreak{}s\allowbreak{}o\allowbreak{}n}: the 37 strict lower coefficients, parameter enclosures, critical constants and four reserves; SHA256 \texttt{5\allowbreak{}1\allowbreak{}0\allowbreak{}d\allowbreak{}c\allowbreak{}7\allowbreak{}a\allowbreak{}6\allowbreak{}3\allowbreak{}b\allowbreak{}2\allowbreak{}0\allowbreak{}1\allowbreak{}8\allowbreak{}b\allowbreak{}c\allowbreak{}4\allowbreak{}5\allowbreak{}4\allowbreak{}0\allowbreak{}6\allowbreak{}f\allowbreak{}f\allowbreak{}b\allowbreak{}1\allowbreak{}f\allowbreak{}f\allowbreak{}8\allowbreak{}8\allowbreak{}a\allowbreak{}3\allowbreak{}c\allowbreak{}8\allowbreak{}b\allowbreak{}a\allowbreak{}5\allowbreak{}2\allowbreak{}3\allowbreak{}a\allowbreak{}f\allowbreak{}1\allowbreak{}e\allowbreak{}7\allowbreak{}6\allowbreak{}1\allowbreak{}7\allowbreak{}f\allowbreak{}2\allowbreak{}5\allowbreak{}6\allowbreak{}8\allowbreak{}0\allowbreak{}b\allowbreak{}2\allowbreak{}9\allowbreak{}9\allowbreak{}5\allowbreak{}0\allowbreak{}f\allowbreak{}d\allowbreak{}e\allowbreak{}1\allowbreak{}b\allowbreak{}2}.
\item
  \texttt{r\allowbreak{}e\allowbreak{}s\allowbreak{}e\allowbreak{}a\allowbreak{}r\allowbreak{}c\allowbreak{}h\allowbreak{}/\allowbreak{}l\allowbreak{}o\allowbreak{}c\allowbreak{}a\allowbreak{}l\allowbreak{}\_\allowbreak{}c\allowbreak{}y\allowbreak{}c\allowbreak{}l\allowbreak{}e\allowbreak{}0\allowbreak{}8\allowbreak{}0\allowbreak{}/\allowbreak{}n\allowbreak{}e\allowbreak{}a\allowbreak{}r\allowbreak{}\_\allowbreak{}r\allowbreak{}e\allowbreak{}c\allowbreak{}o\allowbreak{}r\allowbreak{}d\allowbreak{}/\allowbreak{}R\allowbreak{}E\allowbreak{}T\allowbreak{}U\allowbreak{}R\allowbreak{}N\allowbreak{}\_\allowbreak{}F\allowbreak{}I\allowbreak{}N\allowbreak{}A\allowbreak{}L\allowbreak{}.\allowbreak{}j\allowbreak{}s\allowbreak{}o\allowbreak{}n}: final fixed-pad repricing of every sector, six additional reserves, source-domain guards, averaged floor and principal/small margins; SHA256 \texttt{4\allowbreak{}c\allowbreak{}5\allowbreak{}8\allowbreak{}e\allowbreak{}7\allowbreak{}d\allowbreak{}a\allowbreak{}6\allowbreak{}1\allowbreak{}7\allowbreak{}4\allowbreak{}3\allowbreak{}1\allowbreak{}5\allowbreak{}7\allowbreak{}0\allowbreak{}4\allowbreak{}0\allowbreak{}a\allowbreak{}1\allowbreak{}2\allowbreak{}c\allowbreak{}0\allowbreak{}5\allowbreak{}3\allowbreak{}d\allowbreak{}8\allowbreak{}c\allowbreak{}1\allowbreak{}a\allowbreak{}1\allowbreak{}f\allowbreak{}2\allowbreak{}c\allowbreak{}5\allowbreak{}5\allowbreak{}4\allowbreak{}4\allowbreak{}f\allowbreak{}3\allowbreak{}9\allowbreak{}0\allowbreak{}1\allowbreak{}d\allowbreak{}7\allowbreak{}2\allowbreak{}1\allowbreak{}e\allowbreak{}2\allowbreak{}1\allowbreak{}3\allowbreak{}a\allowbreak{}d\allowbreak{}c\allowbreak{}d\allowbreak{}a\allowbreak{}0\allowbreak{}4\allowbreak{}3\allowbreak{}f\allowbreak{}0\allowbreak{}6\allowbreak{}4}.
\end{itemize}

The exact-contact input is \texttt{r\allowbreak{}e\allowbreak{}s\allowbreak{}e\allowbreak{}a\allowbreak{}r\allowbreak{}c\allowbreak{}h\allowbreak{}/\allowbreak{}l\allowbreak{}o\allowbreak{}c\allowbreak{}a\allowbreak{}l\allowbreak{}\_\allowbreak{}c\allowbreak{}y\allowbreak{}c\allowbreak{}l\allowbreak{}e\allowbreak{}0\allowbreak{}7\allowbreak{}7\allowbreak{}/\allowbreak{}i\allowbreak{}n\allowbreak{}t\allowbreak{}a\allowbreak{}k\allowbreak{}e\allowbreak{}/\allowbreak{}E\allowbreak{}X\allowbreak{}A\allowbreak{}C\allowbreak{}T\allowbreak{}\_\allowbreak{}C\allowbreak{}O\allowbreak{}N\allowbreak{}T\allowbreak{}A\allowbreak{}C\allowbreak{}T\allowbreak{}\_\allowbreak{}R\allowbreak{}E\allowbreak{}C\allowbreak{}E\allowbreak{}I\allowbreak{}P\allowbreak{}T\allowbreak{}.\allowbreak{}j\allowbreak{}s\allowbreak{}o\allowbreak{}n}, SHA256 \texttt{b\allowbreak{}d\allowbreak{}b\allowbreak{}2\allowbreak{}b\allowbreak{}a\allowbreak{}b\allowbreak{}4\allowbreak{}b\allowbreak{}4\allowbreak{}a\allowbreak{}3\allowbreak{}9\allowbreak{}7\allowbreak{}c\allowbreak{}b\allowbreak{}a\allowbreak{}f\allowbreak{}8\allowbreak{}0\allowbreak{}9\allowbreak{}b\allowbreak{}e\allowbreak{}8\allowbreak{}a\allowbreak{}0\allowbreak{}e\allowbreak{}9\allowbreak{}5\allowbreak{}2\allowbreak{}0\allowbreak{}0\allowbreak{}7\allowbreak{}7\allowbreak{}d\allowbreak{}5\allowbreak{}4\allowbreak{}d\allowbreak{}2\allowbreak{}9\allowbreak{}e\allowbreak{}b\allowbreak{}6\allowbreak{}2\allowbreak{}e\allowbreak{}1\allowbreak{}5\allowbreak{}f\allowbreak{}e\allowbreak{}1\allowbreak{}1\allowbreak{}b\allowbreak{}2\allowbreak{}4\allowbreak{}c\allowbreak{}4\allowbreak{}4\allowbreak{}0\allowbreak{}c\allowbreak{}7\allowbreak{}c\allowbreak{}e\allowbreak{}c\allowbreak{}0}. The portable079 checker uses only these exact data and the rational polynomial helper \texttt{r\allowbreak{}e\allowbreak{}s\allowbreak{}e\allowbreak{}a\allowbreak{}r\allowbreak{}c\allowbreak{}h\allowbreak{}/\allowbreak{}c\allowbreak{}y\allowbreak{}c\allowbreak{}l\allowbreak{}e\allowbreak{}0\allowbreak{}6\allowbreak{}6\allowbreak{}/\allowbreak{}r\allowbreak{}e\allowbreak{}t\allowbreak{}u\allowbreak{}r\allowbreak{}n\allowbreak{}\_\allowbreak{}c\allowbreak{}o\allowbreak{}r\allowbreak{}e\allowbreak{}.\allowbreak{}p\allowbreak{}y}, SHA256 \texttt{3\allowbreak{}4\allowbreak{}8\allowbreak{}1\allowbreak{}8\allowbreak{}b\allowbreak{}0\allowbreak{}2\allowbreak{}7\allowbreak{}a\allowbreak{}3\allowbreak{}b\allowbreak{}8\allowbreak{}c\allowbreak{}4\allowbreak{}a\allowbreak{}3\allowbreak{}e\allowbreak{}1\allowbreak{}9\allowbreak{}5\allowbreak{}4\allowbreak{}d\allowbreak{}a\allowbreak{}d\allowbreak{}7\allowbreak{}7\allowbreak{}a\allowbreak{}a\allowbreak{}2\allowbreak{}d\allowbreak{}4\allowbreak{}a\allowbreak{}9\allowbreak{}9\allowbreak{}b\allowbreak{}f\allowbreak{}4\allowbreak{}8\allowbreak{}2\allowbreak{}e\allowbreak{}6\allowbreak{}6\allowbreak{}a\allowbreak{}a\allowbreak{}d\allowbreak{}e\allowbreak{}1\allowbreak{}8\allowbreak{}b\allowbreak{}b\allowbreak{}4\allowbreak{}6\allowbreak{}b\allowbreak{}b\allowbreak{}c\allowbreak{}9\allowbreak{}3\allowbreak{}6\allowbreak{}0\allowbreak{}8\allowbreak{}8\allowbreak{}1\allowbreak{}3}. The080 checker accepts the packed079 receipt through \texttt{-\allowbreak{}\/\allowbreak{}-\allowbreak{}p\allowbreak{}r\allowbreak{}e\allowbreak{}v\allowbreak{}i\allowbreak{}o\allowbreak{}u\allowbreak{}s\allowbreak{}-\allowbreak{}s\allowbreak{}o\allowbreak{}u\allowbreak{}r\allowbreak{}c\allowbreak{}e}. These names identify mathematical input data; their finite checking does not replace the analytic contracts below.

Edges 0,1,2,3 mean, respectively, old \(\fitinline{D_c=1/200}\), old \(\fitinline{D_c=19/25}\), native \(\fitinline{D_c=19/25}\), and native \(\fitinline{D_c=h}\). They contain 1,2,2,52 boxes. Their closed union at each edge is \(\fitinline{[0,3/4]\times[0,1/2]}\), with every nonold branch used only inside its original full admission window. There are 17 old, 4 nativeLow, 16 native and 20 improved boxes. The 20 improved boxes are all on edge3, with IDs \[\fitmath{
27,29,35,38,39,40,41,42,43,44,45,46,47,48,49,50,51,52,53,54.
}\] Their rectangles lie in E(27). The remaining 37 IDs use the saved strict polynomials. Closed endpoint overlaps are harmless; no uncovered open cell or admission beyond those windows is introduced.

Appendix E proves that each saved polynomial is a lower bound for the corresponding cleared return margin throughout the retained parameter enclosure. Its branch denominator satisfies \(\fitinline{0<J\le5/2}\). On the box \(\fitinline{[a,b]\times[c,e]}\), substitute \(\fitinline{d=a+(b-a)x,p=c+(e-c)y}\). If the resulting polynomial is \(\fitinline{\sum a_{ij}x^iy^j}\), of bidegree at most \(\fitinline{(m,n)}\), its exact tensor-Bernstein coefficient is \[\fitmath{
 B_{kl}=\sum_{i\le k,j\le l}a_{ij}
        {\binom{k}{i}\binom{l}{j}\over\binom{m}{i}\binom{n}{j}}.
}\] The tensor basis is nonnegative and sums to one, so the polynomial is at least \(\fitinline{\min B_{kl}}\). This is the implication used for each \texttt{s\allowbreak{}t\allowbreak{}r\allowbreak{}i\allowbreak{}c\allowbreak{}t\allowbreak{}\_\allowbreak{}c\allowbreak{}e\allowbreak{}l\allowbreak{}l\allowbreak{}s\allowbreak{}.\allowbreak{}c\allowbreak{}l\allowbreak{}e\allowbreak{}a\allowbreak{}r\allowbreak{}e\allowbreak{}d\allowbreak{}\_\allowbreak{}m\allowbreak{}a\allowbreak{}r\allowbreak{}g\allowbreak{}i\allowbreak{}n}; the coefficients and boxes remain attached, rather than being replaced by a success flag. Dividing by \(\fitinline{J\le5/2}\) and charging \(\fitinline{6\xi_0}\) gives \[\fitmath{
 r_{37}={637501751112735390608378416271343917\over
 589824000002949120000000000000000000000000}>10^{-6}.     
}\tag{H37b}\]

For the 20 critical boxes, E(49) gives the exact identity \[\fitmath{
 M_h(d,p)=K_h(d-\delta_*)^2+\varepsilon_p dQ(d,p),
 \quad K_h>1/5,\quad .4258<\delta_*<.4259.
}\] On E(27), \(\fitinline{Q\ge.01,J\le2.5}\). On \(\fitinline{d\in[.425,.4267]}\), exact interval propagation of the original \(\fitinline{q_{10},q_{20},q_{21}}\) gives \[\fitmath{
 Q\ge {45947938059979804034282897191\over
 1171875000000000000000000000000}>.039,
 \quad J\le1294053644991/800000000000<1.63.
}\] Thus \(\fitinline{\varepsilon_p dQ/J\ge A\varepsilon_p}\), \(\fitinline{A=663/65200}\). Outside that density patch the distance to \(\fitinline{\delta_*}\) is at least \(\fitinline{.0008}\). After the target debit, the two unrestricted critical complements have reserves \[\fitmath{
 r_d={(1/5)(.0008)^2\over5/2}-6\xi_0
       ={727\over50000000000},\qquad
 r_p=A\varepsilon_0-6\xi_0
       ={32721\over8150000000000}>4.0148\cdot10^{-9}.    
}\tag{H37c}\] The second case is \(\fitinline{\varepsilon_p\ge\varepsilon_0}\).

\hypertarget{the-whole-radical-sector-and-its-affine-endpoints}{%
\subsubsection{The whole radical sector and its affine endpoints}\label{the-whole-radical-sector-and-its-affine-endpoints}}

For \(\fitinline{N\operatorname{rad}(u)\le U^{1-b}}\), the original masked-mean identity bounds the ACTUAL error-slot width by \(\fitinline{2\varepsilon_p\ell_t/D_c}\). For an injective error tuple let \(\fitinline{d_T}\) be the product of norms of its primes with row valuation at least two. The complete mass bound is \[\fitmath{
 \sum_e|A_e|h_e^{-1/2}\max(1,h_e/d_T)^{d/2}\le2B_T,
 \qquad \sum_{T,\mathrm{masks}}B_T\ll U^\epsilon.
}\] It retains every ramified coefficient, main/error choice and shift. Since \(\fitinline{C_u\ll N\operatorname{rad}(u)}\), \(\fitinline{C_ud_T\ll U^{1-b+2\varepsilon_p\ell_t/D_c+o(1)}}\). The genuine source \(\fitinline{G(1-s,2+10s,3/20-1000s/1143)}\), \(\fitinline{s=10^{-260}}\), therefore supplies \[\fitmath{
 C_b(D_c,p)=\min\left\{3D_c/40,
 {2y_t-(1+s-2b)D_c-4\varepsilon_p\ell_t\over2+10s}\right\}.
}\tag{H39}\] On a shift bin of exponent \(\fitinline{v}\), the mass pays \(\fitinline{dv/2}\), while capacity clipping costs at most \(\fitinline{qv/(2+10s)}\). Their difference is positive for \(\fitinline{q\le d/2}\). This proves the full shifted consumer, including all shortenings, native nonpositive-length fallbacks and empty/zero-scale cases.

Use (H39) only on \(\fitinline{[.425,.4267]\times[.5-\varepsilon_0,.5]}\). The exact guard \(\fitinline{b(19/25)-2\varepsilon_0\ell_t>0}\) makes its lower conductor edge no worse than the native source. At the top, the native affine term is below its cap by more than \(\fitinline{5.357\cdot10^{-5}}\), whereas \(\fitinline{bh<1.631\cdot10^{-6}}\). Hence the top capacity gain at the \(\fitinline{\alpha}\) geometry is exactly \(\fitinline{bh-2\varepsilon_p\ell_\alpha}\). Write \(\fitinline{h_- =815324503/10^9}\) and \(\fitinline{\ell_{\alpha,+}=4061325989995133424889/ 23437500000000000000000}\). The remaining reserve is the minimum of an AFFINE expression in \(\fitinline{\varepsilon_p}\): \[\fitmath{
 r_b=\min_{a\in\{0,\varepsilon_0\}}
 \left[Aa+{.425(.5-\varepsilon_0)\over2}
                  (bh_- -2a\ell_{\alpha,+})\right]-7\xi_0
 ={227980394593663235296917992419\over
 9550781250000000000000000000000000000}>2.38\cdot10^{-8}. 
}\tag{H39a}\] The \(\fitinline{6\xi_0}\) and \(\fitinline{7\xi_0}\) charges in (H37b)--(H39a) retain the live capacity denominators and target change. The extra one in (H39a) pays the \(\fitinline{-4\varepsilon_p\ell_t}\) dependence. For fixed \(\fitinline{d,p}\), (H37a) is the maximum of affine functions of \(\fitinline{D_c}\), because each capacity is a minimum of affine functions and its coefficient is \(\fitinline{-dp/2\le0}\). The endpoint bounds fill both complete original conductor intervals. Thus these four cases cover every original box, not merely a contact or top edge.

\hypertarget{fixed-padding-and-the-final-physical-complements}{%
\subsubsection{Fixed padding and the final physical complements}\label{fixed-padding-and-the-final-physical-complements}}

Put \(\fitinline{F=\Delta=10^{-258}}\), and choose actual finite baseline data with \(\fitinline{0<D.t\le F/4}\), using loss budget \(\fitinline{\min((\beta-t)/10000,10^{-9}/4,F)}\) before live heights. Then the literal cube pad satisfies \(\fitinline{F/2+2D.t\le F}\). This is GEOMETRIC coverage. The distinct exponent charge is \(\fitinline{15s/16+2F}\). It is charged to every margin, not bounded by \(\fitinline{F}\). Indeed natural/native/radical capacities are 1-Lipschitz in \(\fitinline{D_c}\), and \(\fitinline{0\le R\le1}\), so (H37a) changes by at most \[\fitmath{
 (9/100+3/4+3/16)|\Delta D_c|=(411/400)|\Delta D_c|
                              <2|\Delta D_c|.
}\] The080 receipt verifies each of (H37b)--(H39a), AFTER subtracting \(\fitinline{15s/16+2F}\), exceeds \(\fitinline{3\cdot10^{-9}}\), along with \(\fitinline{h+F<\min(1,2y_t)}\), \(\fitinline{y_t<.76<h}\), and \(\fitinline{\ell_t/(h+F)>.2}\). All 57 sectors therefore reach the fixed-pad cutoff. The earlier variable-pad \(\fitinline{10^{-260}}\) extension is not used for this cube.

On the remaining critical strip let \[\fitmath{
 B_t=(3+4(.4267))\xi_0=7189637/250000000000000,
 \quad D_- =.815324503-10^{-7}.
}\tag{H38}\] At \(\fitinline{D_c=h}\), \(\fitinline{E_t-t\le B_t}\). Both native capacity branches have return slope at least \(\fitinline{a_D=531159/1600000}\); use \(\fitinline{R_b\ge1-d}\) in (H37a). The complete conductor complement \(\fitinline{D_c\le h-10^{-7}}\) is therefore paid by \(\fitinline{a_D10^{-7}-B_t}\).

For the remaining terms, the six BARE reserves below already subtract \(\fitinline{B_t}\), but precede the separate source/extension and reconstruction charges. Their exact fractions are in the080 receipt.

\begin{longtable}[]{@{}
  >{\raggedright\arraybackslash}p{(\columnwidth - 4\tabcolsep) * \real{0.3000}}
  >{\raggedright\arraybackslash}p{(\columnwidth - 4\tabcolsep) * \real{0.3000}}
  >{\raggedleft\arraybackslash}p{(\columnwidth - 4\tabcolsep) * \real{0.4000}}@{}}
\toprule\noalign{}
\begin{minipage}[b]{\linewidth}\raggedright
Removed original contribution
\end{minipage} & \begin{minipage}[b]{\linewidth}\raggedright
Reserve formula
\end{minipage} & \begin{minipage}[b]{\linewidth}\raggedleft
Bare value
\end{minipage} \\
\midrule\noalign{}
\endhead
\bottomrule\noalign{}
\endlastfoot
\(\fitinline{D_c\le h-10^{-7}}\) & \(\fitinline{a_D10^{-7}-B_t}\) & \(\fitinline{4.4388895\cdot10^{-9}}\) \\
\(\fitinline{n\le n_0-.5\cdot10^{-6}}\) & \(\fitinline{.425(.499996)(.49\cdot10^{-6})D_-/2-B_t}\) & \(\fitinline{1.36889\allowbreak{}39148\allowbreak{}57365\cdot10^{-8}}\) \\
dilation \(\fitinline{v_{\rm dil}>.4\cdot10^{-6}}\) & \(\fitinline{(.425/4)(.4\cdot10^{-6})D_- -B_t}\) & \(\fitinline{5.89273\allowbreak{}91275\cdot10^{-9}}\) \\
\(\fitinline{r<1/D_c-4\cdot10^{-7}}\) & \(\fitinline{.10(4\cdot10^{-7})D_--B_t}\) & \(\fitinline{3.85442812\cdot10^{-9}}\) \\
\(\fitinline{r>1/D_c+1.5\cdot10^{-7}}\) & \(\fitinline{.28(1.5\cdot10^{-7})D_--B_t}\) & \(\fitinline{5.485076926\cdot10^{-9}}\) \\
low inverse response & \(\fitinline{(5\cdot10^{-8})D_--B_t}\) & \(\fitinline{1.200767215\cdot10^{-8}}\) \\
\end{longtable}

Here \(\fitinline{n_0=1-y_t/D_c}\). The short-numerator quota is \(\fitinline{z_0+.49\cdot10^{-6}}\), where \(\fitinline{z_0=(1-s-2n_0)/(2+10s)}\); it stays below \(\fitinline{3/40}\) and satisfies \(\fitinline{2(.5\cdot10^{-6})>(2+10s)(.49\cdot10^{-6})}\). For the actual coefficient dilation \(\fitinline{v_{\rm dil}=\log_U\max(1,L_iK)}\), mass and clipping save \(\fitinline{dv_{\rm dil}/4}\). Retained terms have reflected scale \(\fitinline{\ll U^{n_0+.4\cdot10^{-6}}}\); annular reconstruction with fixed extra slack \(\fitinline{.01\cdot10^{-6}}\) gives \(\fitinline{n\le n_0+.41\cdot10^{-6}}\). No coefficient is divided out.

For reciprocal tails first open the ACTUAL denominator on \(\fitinline{\Re s=3}\), at fixed \(\fitinline{z}\), mask and amplitude class, after full \(\fitinline{w}\) inversion. The class is independent of \(\fitinline{s,w}\). The quotient-free local expansion is holomorphic on the move. Each inverse dyad includes \(\fitinline{X^{1/2-s}M_u(r;s)}\). Move it by \(\fitinline{\Delta_s=-.10}\) or \(\fitinline{.28}\). Variable damping \(\fitinline{\rho=d/2+\Delta_s}\) pays all shifts since \(\fitinline{\rho-q/2\ge.005825}\). Restoring its exterior power gives \(\fitinline{U^{\Delta_s(1/D_c-r)+o(1)}}\), proving the two length-tail prices. At \(\fitinline{\Delta_s=0}\), inverse response below \(\fitinline{U^{dr/2-5\cdot10^{-8}}}\) gives the final price. This amplitude restriction is made after contour motion. Every table entry minus \(\fitinline{15s/16+2F}\) is \(\fitinline{>3\cdot10^{-9}}\). Choose total remaining reconstruction loss below \(\fitinline{10^{-9}}\); a common final reserve \(\fitinline{10^{-9}}\) remains.

Use E(31)--(40)'s AVERAGED floor, whose exponent is \(\fitinline{F_f(D_c)=.51-49h/150+.51\ell_t+.66D_c}\). At \(\fitinline{D_c=h+F}\), its margin after \(\fitinline{10^{-5}+10^{-12}}\) is \(\fitinline{>10^{-3}}\), giving the supported floor cube's saving \(\fitinline{1/4000}\). The080 receipt separately gives small-row margin \(\fitinline{3570096049/50000000000}\) and principal contour margins \(\fitinline{64117721941/3000000000000}\), \(\fitinline{65324503/600000000000}\). Keep B.8's five nonprincipal units, GT23's large-row estimate with fixed positive \(\fitinline{F}\), and the complete principal residues, quotient-free other contours, finite prime correction and tails. The older pointwise floor is not used at this geometry.

\hypertarget{the-original-witness-and-the-mark-divided-return}{%
\subsubsection{The original witness and the mark-divided return}\label{the-original-witness-and-the-mark-divided-return}}

The remaining original ticket has precisely \[\fitmath{
 \begin{gathered}
 h-10^{-7}<D_c\le h+F,\quad .425\le d\le.4267,
 \quad p>.499996,\quad N\operatorname{rad}(u)>U^{.999998},\\
 n_0-.5\cdot10^{-6}<n\le n_0+.41\cdot10^{-6},
 \quad n_0=1-y_t/D_c,\\
 -4\cdot10^{-7}\le r-1/D_c\le1.5\cdot10^{-7},
 \quad |M_{\rm den}|\ge U^{dr/2-5\cdot10^{-8}}.
 \end{gathered}
}\tag{H37}\]

The only remaining original rows have \(\fitinline{N\operatorname{rad}(u)>U^{.999998}}\), \(\fitinline{d\in[.425,.4267]}\), \(\fitinline{p>.499996}\), and the retained conductor and dyad ranges above. H.6's actual inverse identity gives, outside its small exceptional set, SOME denominator ordinary witness with \(\fitinline{m-4\lambda<y<m+\lambda}\), \(\fitinline{\lambda=10^{-5}}\), and \(\fitinline{|S_u(y,t_u)|\ge U^{dy/2-1.5\eta_0}}\), \(\fitinline{\eta_0=10^{-7}}\). This does not truncate the entire inverse convolution. High/middle disallowed witnesses have the original single-factor count \(\fitinline{O(U^{R-\eta_0})}\): tangent/fallback margins exceed \(\fitinline{10\eta_0}\), witness debits are at most \(\fitinline{6\eta_0}\) or \(\fitinline{3\eta_0}\), and the finite loss fits. Original fourth/H\"{o}lder then saves \(\fitinline{3\eta_0/4}\) in their physical mass, which exceeds \(\fitinline{B_t}\) after conversion to \(\fitinline{Z}\).

For H27 choose whole original slots of quota \(\fitinline{W_*=(W_P+W_S)/2+42\cdot10^{-6}<.101}\) in decreasing lower gain. The TOTAL original width is \(\fitinline{L=\ell_t/D_c>1/5}\), and its full gain is \(\fitinline{qL}\). A fractional prefix has gain at least \(\fitinline{qW_*}\). Rounding costs at most one slot's charged mesh. If positive-gain slots are exhausted, take them all: their unchanged gain is \(\fitinline{qL\ge qW_*}\). No lower bound on total positive-gain width is asserted. Dividing the POSITIVE H27 estimate by this squared common response therefore gives on the SAME original class \[\fitmath{
 \sum|P_u(n,v)S_u(y,t')|^2
 \ll U^{1-2qW_*+\epsilon}(1+|v|+|t'|+T)^J.             
}\tag{H42a}\] Apply unit-interval unweighted Sobolev to (H42a) and its first log test over the full allowed \(\fitinline{t_u}\) interval enlarged by one. The derivative inserts \(\fitinline{i\log(NI/U^y)}\); the finite polynomial height cost is paid after its degree is fixed. Keep numerator height \(\fitinline{v}\) independent and integrate its original rapid weight. Cauchy with the count \(\fitinline{U^R}\) gives \[\fitmath{
 E_{\rm wit}=(R+1-dy+3\eta_0-2qW_*)/2.                
}\tag{H42b}\] The full \(\fitinline{+3\eta_0}\) remains. With \(\fitinline{R=1-dC}\), \(\fitinline{C=p(1+2g_0)+2k_1m}\), \(\fitinline{k_1=1-p(1+2g_1)}\), as in H29, the saving beyond \(\fitinline{3\eta_0/4}\) over \(\fitinline{(3R+1)/4-qz_0/2}\) is at least \[\fitmath{
 { .425[4(.499996)(42\cdot10^{-6})-8k_+10^{-5}
                         -.41\cdot10^{-6}]-9\eta_0\over4}
 ={63153680778247\over12500000000000000000},
 \quad k_+=1-.499996(1+2(.07599618)).
}\tag{H42}\] This is \(\fitinline{>5.0522\cdot10^{-6}}\) in the \(\fitinline{U}\) exponent. Multiply the FULL saving by \(\fitinline{D_-}\) and subtract \(\fitinline{B_t}\): the remaining \(\fitinline{Z}\) saving is \(\fitinline{51895621855589310661541/12500000000000000000000000000 >4.1516\cdot10^{-6}}\). The attached \texttt{r\allowbreak{}e\allowbreak{}s\allowbreak{}e\allowbreak{}a\allowbreak{}r\allowbreak{}c\allowbreak{}h\allowbreak{}/\allowbreak{}l\allowbreak{}o\allowbreak{}c\allowbreak{}a\allowbreak{}l\allowbreak{}\_\allowbreak{}c\allowbreak{}y\allowbreak{}c\allowbreak{}l\allowbreak{}e\allowbreak{}0\allowbreak{}8\allowbreak{}0\allowbreak{}/\allowbreak{}n\allowbreak{}e\allowbreak{}a\allowbreak{}r\allowbreak{}\_\allowbreak{}r\allowbreak{}e\allowbreak{}c\allowbreak{}o\allowbreak{}r\allowbreak{}d\allowbreak{}/\allowbreak{}J\allowbreak{}O\allowbreak{}I\allowbreak{}N\allowbreak{}.\allowbreak{}j\allowbreak{}s\allowbreak{}o\allowbreak{}n} records these exact fractions. Source decrements, packet damping and rowwise scale/profile/Fourier losses fit this reserve. This is the actual one-witness consumer of H27, with both independent heights and the original physical reciprocal still accounted for. H.8--H.9 then supply the stated source-first height quantifiers and original common-probe/Mellin return; no additional scalar affirmation of the endpoint is a premise here.

The finite inputs preserve Bebel's count/contact work, Sil's paired source, Levy's numerator and Collet's physical return, with the full OpenAI foundation and the inherited source notices and assistance disclosures. Only this integration ledger is newly written.

\hypertarget{appendix-h-8}{}

\hypertarget{h.8.-the-actual-finite-integral-and-its-quantifiers}{%
\subsection{H.8. The actual finite integral and its quantifiers}\label{h.8.-the-actual-finite-integral-and-its-quantifiers}}

The finite cube in this argument is the literal truncation, original buffered segments, row masks and grids of the compensated probe in B.1--B.4. Before truncation its exact triple Poisson representation has scalar kernel \[\fitmath{
 X^{1/2-z}Z^{s+z-1}Y^{w-1}e^{(s+z-1)^2}M(z)\widehat W_1(w)
}\] times \[\fitmath{
 \sum_u^{(6)}(Nu)^{-z}\bar\xi(u)
 \frac{\zeta_K^{\mathcal S}(6z)L_K^{\mathcal S}(w,\psi_u)}
 {L_K^{\mathcal S}(s,\eta\bar\psi_u)}
 H^{\rm comp}_{\eta,u,Z}(s,w,z),
}\tag{H43}\] integrated with \(\fitinline{(2\pi i)^{-3}}\) on the original absolute lines \(\fitinline{(3,3,2)}\). The selected-prime part of \(\fitinline{H^{\rm comp}}\) is \[\fitmath{
 \sum_{(\mathfrak p_i)}\prod_iw_i(N\mathfrak p_i/P_i)
 (N\mathfrak p_i)^{z-1}G_{\mathfrak p_i}
 \prod_{\mathfrak p\notin\mathcal S,\,\mathfrak p\ {\rm not\ selected}}H_{\mathfrak p},
 \qquad G_{\mathfrak p}=-\bar\psi_u(\mathfrak p)H_{\mathfrak p}+\Delta_{\mathfrak p}.
}\tag{H44}\] This is quotient-free, including primes meeting the row and every main/error mask. The Euler product converges normally on the original region \(\fitinline{\Re s\ge.51}\), \(\fitinline{\Re z\ge.34}\), \(\fitinline{\Re w\ge-.01}\), \(\fitinline{\Re(s+w)>1+\epsilon}\). The unique principal unit is \(\fitinline{u=1}\); the other five are nonprincipal. At a good row prime the row valuation runs from one through five; its primitive conductor exponent is one. Formula H43, not an unweighted mixed moment, is what the bounds above must estimate.

The raw source estimate H25, the full finite ledger of H.7 and all the complementary sectors give the following literal source-first statement: \[\fitmath{
 \begin{gathered}
 \forall\eta\ \exists\tau_\eta>0\ \forall n>0\ \exists C_n>0\quad
 \text{eventually in }Z:\quad\forall\mathrm{idx,grid},\\
 \mathrm{ActualBufferedSelection}\Longrightarrow
 \left|\frac{\mathrm{cube}_{D,F_{\rm src},\eta}(Z,\tau_\eta,\mathrm{idx,grid})}
 {\mathrm{normalizer}_{D,F_{\rm src}}(Z)}\right|
 \le C_n Z^{\beta-k-P.\sigma}.
\end{gathered}
}\tag{H45}\] Here is the selection in H45 explicitly. Write \(\fitinline{e=D.e}\), \(\fitinline{T=Z^{\tau_\eta}}\), \(\fitinline{i_u=\mathrm{idx}(u)}\), \(\fitinline{j_u=\mathrm{grid}(u)}\), and \(\fitinline{a_u=.51+e j_u}\). For each free row the indices are natural, \(\fitinline{1\le i_u\le n}\) and \(\fitinline{j_u\le\lfloor .49/e\rfloor}\); its cube has height \(\fitinline{(3i_u+1)T}\), with \(\fitinline{(3i_u+1)T+T/2\le(3i_u+2)T}\). We require \(\fitinline{.51\le a_u\le7/8}\). The maximum for the original source detector family (the actual target and the full ray-cube family, as in B.5) lies in \(\fitinline{[a_u,a_u+e)}\) at height \(\fitinline{3i_uT}\) and is \(\fitinline{<a_u+2e}\) at \(\fitinline{3(i_u+1)T}\). If \(\fitinline{a_u>.51}\), some member has an actual zero \(\fitinline{\rho}\) other than the principal pole, with \(\fitinline{a_u\le\Re\rho<a_u+e}\) and \(\fitinline{|\Im\rho|\le3i_uT}\). In the original row band with cutoffs \(\fitinline{Z^{1/100}, Z^{h+F/2}}\), each row with nonzero calibration residue has every indicated ray twist nonprincipal, of presented modulus norm at most the original conductor constant times \(\fitinline{N M\,N(u)}\) for the fixed ray modulus \(\fitinline{M}\); its ray-family zero maximum at \(\fitinline{3(i_u+1)T}\) is \(\fitinline{<a_u+2e}\). These are conditions on the \emph{same} maps and rows provided by transport. The estimate is uniform over those maps.

Here \(\fitinline{P.\sigma>0}\) is fixed below the common reserve and existing transport budgets. Source data \(\fitinline{D,F_{\rm src},P}\) were chosen in H.1 before \(\fitinline{\eta}\), with numerical padding \(\fitinline{F}\) as in H7. The symbol \(\fitinline{k}\) is the physical shift defined in H4; \(\fitinline{k_1}\) in H29 is the separate count coefficient. Fix meshes and count/source/physical reserves before the finite original slots. After those slots fix the character dictionary, internal profile, Sobolev and moment orders. Only then choose \(\fitinline{\tau_\eta}\) beneath their finitely many polynomial degree caps, before the grid bound \(\fitinline{n}\). Constants and eventual thresholds may depend on the fixed target and \(\fitinline{n}\); there is no uniformity in growing moduli. Independent retained segments are truncated while reciprocal control is global. Horizontal joins have Gaussian decay in \(\fitinline{\Im(s+z)}\), and the independent \(\fitinline{w}\)-weight supplies its own decay. External decay orders follow the finite polynomial envelopes. These choices give the quantifiers in H45, rather than a bound for one preselected grid.

\hypertarget{appendix-h-9}{}

\hypertarget{h.9.-mellin-return-and-ordinary-characters}{%
\subsection{H.9. Mellin return and ordinary characters}\label{h.9.-mellin-return-and-ordinary-characters}}

The principal double residue in H43 has nonzero normalizer \[\fitmath{
 c_{\mathcal S}\mathcal A(Z),\qquad
 \mathcal A(Z)=(-1)^NZ^{-\ell/6}
 \prod_i\sum_{\mathfrak p}w_i(N\mathfrak p/P_i)(N\mathfrak p)^{-5/6},
}\tag{H46}\] of logarithmic size in the fixed prime annuli. Write \(\fitinline{J_\eta}\) for the original compensated probe divided by this full normalizer. The exact signal is \[\fitmath{
 f_\eta(Z)=\frac1{2\pi i}\int_{(2)}Z^{s-k}
 e^{(s-5/6)^2}\frac{H_\eta(s)}{L_K^{\mathcal S}(s,\eta)}\,ds.
}\tag{H47}\] The numerator is the original Euler numerator from B.4, holomorphic and satisfying \(\fitinline{|H_\eta-1|\le1/2}\) on \(\fitinline{\Re s\ge.87}\) after the originally fixed bad cutoff. It is not a newly centered polynomial or the transform of an error term.

H45, with exactly this probe and normalizer, is the input of the existing \texttt{h\allowbreak{}i\allowbreak{}g\allowbreak{}h\allowbreak{}\_\allowbreak{}o\allowbreak{}f\allowbreak{}\_\allowbreak{}c\allowbreak{}u\allowbreak{}b\allowbreak{}e} transport argument. It gives \(\fitinline{J_\eta-f_\eta=O_\eta(Z^{\beta-k-\sigma})}\) for a common \(\fitinline{\sigma>0}\). The retained signed physical Gram estimate gives \(\fitinline{J_\eta=O_\eta(Z^{t-k+\omega})}\) for every \(\fitinline{\omega>0}\). Under \(\fitinline{\mu=\beta-t>0}\), choose \(\fitinline{\omega<\mu/4}\) and \[\fitmath{
 0<m<\min(\mu-\omega,\sigma,\beta-.87).
}\] Then \(\fitinline{f_\eta(Z)=O_\eta(Z^{\beta-k-m})}\) at infinity and has the original arbitrary-power decay at zero. Its Mellin transform \(\fitinline{\int_0^\infty f_\eta(Z)Z^{k-s}\,dZ/Z}\) is holomorphic for \(\fitinline{\Re s>\beta-m}\), by those two bounds, locally uniformly on compact subsets. On the absolute line it equals the Gaussian times \(\fitinline{H_\eta/L_\eta}\) by Mellin inversion; uniqueness gives that continuation throughout the half-plane. But the definition of the supremum supplies a primitive target zero with real part above \(\fitinline{\beta-m/2}\). The numerator is nonzero there, so the quotient cannot be holomorphic. No target attaining \(\fitinline{\beta}\) was assumed. This contradiction proves the bound.

Consequently \(\fitinline{\beta\le t}\). For an ordinary presented Dirichlet character use the original norm pullback \(\fitinline{L_K(s,\eta_\chi)=L(s,\chi)L(s,\chi\chi_{-3})}\), restore the fixed deleted Euler factors, which have no zeros for \(\fitinline{\Re s>0}\), and use the usual nonvanishing on \(\fitinline{\Re s\ge1}\). The possible principal pole at one is stated separately. This preserves all heights, every fixed presented modulus and the strict boundary \(\fitinline{\Re s>t}\); it supplies no growing-modulus uniformity.

\hypertarget{appendix-h-10}{}

\hypertarget{h.10.-attribution-and-exact-inputs}{%
\subsection{H.10. Attribution and exact inputs}\label{h.10.-attribution-and-exact-inputs}}

The original OpenAI proof, the full retained A--G foundations, Sil's native paired source, Joseph Bebel's PR141/PR143 exact contact, Akash Levy's numerator and Hailey Collet's physical return retain their source credit. The donor reductions are adapted from the public \textbf{draft} PR43 at immutable head \texttt{c\allowbreak{}8\allowbreak{}3\allowbreak{}8\allowbreak{}a\allowbreak{}7\allowbreak{}5\allowbreak{}0\allowbreak{}0\allowbreak{}7\allowbreak{}9\allowbreak{}f\allowbreak{}7\allowbreak{}f\allowbreak{}8\allowbreak{}6\allowbreak{}a\allowbreak{}5\allowbreak{}7\allowbreak{}3\allowbreak{}d\allowbreak{}f\allowbreak{}9\allowbreak{}c\allowbreak{}9\allowbreak{}f\allowbreak{}7\allowbreak{}2\allowbreak{}9\allowbreak{}b\allowbreak{}8\allowbreak{}b\allowbreak{}c\allowbreak{}9\allowbreak{}2\allowbreak{}d\allowbreak{}1\allowbreak{}2\allowbreak{}5\allowbreak{}e}, repository \texttt{S\allowbreak{}o\allowbreak{}v\allowbreak{}e\allowbreak{}r\allowbreak{}e\allowbreak{}i\allowbreak{}g\allowbreak{}n\allowbreak{}S\allowbreak{}t\allowbreak{}e\allowbreak{}a\allowbreak{}k\allowbreak{}/\allowbreak{}q\allowbreak{}u\allowbreak{}a\allowbreak{}s\allowbreak{}i\allowbreak{}-\allowbreak{}r\allowbreak{}h\allowbreak{}-\allowbreak{}h\allowbreak{}i\allowbreak{}l\allowbreak{}l\allowbreak{}c\allowbreak{}l\allowbreak{}i\allowbreak{}m\allowbreak{}b}, principally:

\begin{itemize}
\tightlist
\item
  \href{https://github.com/SovereignSteak/quasi-rh-hillclimb/blob/c838a750079f7f86a573df9c9f729b8bc92d125e/auto-research/continuum-local-return/physical-denominator.md}{physical-denominator.md}, the actual denominator and retained physical source;
\item
  \href{https://github.com/SovereignSteak/quasi-rh-hillclimb/blob/c838a750079f7f86a573df9c9f729b8bc92d125e/auto-research/continuum-local-return/short-physical-inverse.md}{short-physical-inverse.md}, exact inverse identity, common frequencies and witness reduction;
\item
  \href{https://github.com/SovereignSteak/quasi-rh-hillclimb/blob/c838a750079f7f86a573df9c9f729b8bc92d125e/auto-research/continuum-local-return/four-factors.md}{four-factors.md}, ``Single-factor counts'' and ``Second detector''; H.6 supplies the required derivation here, without assuming the proposed two-inverse estimate;
\item
  \href{https://github.com/SovereignSteak/quasi-rh-hillclimb/blob/c838a750079f7f86a573df9c9f729b8bc92d125e/auto-research/continuum-local-return/marked-ray-transfer.md}{marked-ray-transfer.md}, exact paired rays, restoring factors and finite source geometry;
\item
  \href{https://github.com/SovereignSteak/quasi-rh-hillclimb/blob/c838a750079f7f86a573df9c9f729b8bc92d125e/auto-research/continuum-local-return/marked-window.md}{marked-window.md}, original ordinary source faces;
\item
  the linked tracker source at \texttt{c\allowbreak{}7\allowbreak{}4\allowbreak{}d\allowbreak{}8\allowbreak{}5\allowbreak{}5\allowbreak{}c\allowbreak{}3\allowbreak{}b\allowbreak{}d\allowbreak{}f\allowbreak{}c\allowbreak{}e\allowbreak{}e\allowbreak{}e\allowbreak{}e\allowbreak{}0\allowbreak{}5\allowbreak{}9\allowbreak{}5\allowbreak{}0\allowbreak{}1\allowbreak{}a\allowbreak{}e\allowbreak{}a\allowbreak{}7\allowbreak{}6\allowbreak{}c\allowbreak{}8\allowbreak{}b\allowbreak{}1\allowbreak{}6\allowbreak{}2\allowbreak{}7\allowbreak{}b\allowbreak{}9\allowbreak{}a\allowbreak{}b\allowbreak{}d}, \texttt{A\allowbreak{}n\allowbreak{}c\allowbreak{}h\allowbreak{}o\allowbreak{}r\allowbreak{}Z\allowbreak{}e\allowbreak{}r\allowbreak{}o\allowbreak{}L\allowbreak{}o\allowbreak{}s\allowbreak{}s\allowbreak{}/\allowbreak{}A\allowbreak{}f\allowbreak{}f\allowbreak{}i\allowbreak{}n\allowbreak{}e\allowbreak{}C\allowbreak{}o\allowbreak{}u\allowbreak{}n\allowbreak{}t\allowbreak{}F\allowbreak{}i\allowbreak{}b\allowbreak{}e\allowbreak{}r\allowbreak{}.\allowbreak{}l\allowbreak{}e\allowbreak{}a\allowbreak{}n} and \texttt{A\allowbreak{}n\allowbreak{}c\allowbreak{}h\allowbreak{}o\allowbreak{}r\allowbreak{}C\allowbreak{}o\allowbreak{}n\allowbreak{}t\allowbreak{}i\allowbreak{}n\allowbreak{}u\allowbreak{}u\allowbreak{}m\allowbreak{}/\allowbreak{}E\allowbreak{}x\allowbreak{}a\allowbreak{}c\allowbreak{}t\allowbreak{}C\allowbreak{}o\allowbreak{}u\allowbreak{}n\allowbreak{}t\allowbreak{}W\allowbreak{}i\allowbreak{}n\allowbreak{}d\allowbreak{}o\allowbreak{}w\allowbreak{}.\allowbreak{}l\allowbreak{}e\allowbreak{}a\allowbreak{}n}, for the stated source contracts and finite crossing/endpoint identities. Their declarations are not a local end-to-end formal verification of this appendix.
\end{itemize}

Those donor notes disclose OpenAI Codex assistance. The retained Quasi-RH Hillclimb notice credits Copyright (c) 2026 Rohan Arun and contributors, Apache-2.0, and preserves its historical MIT notice and external authors' notices. This appendix adapts those source reductions with the complete argument and coefficient contracts given above. Exposition also uses \emph{Mathematical Writing for AI Agents}, an independent mathematical adaptation of selected principles from \href{https://github.com/Nanako0129/sepia}{Sepia} (Nanako Tsai, MIT) and \href{https://github.com/haileystorm/creative-writing-rubrics}{Creative Writing Rubrics} (Hailey Collet, Apache-2.0); those source terms and the original notices remain retained.

The two additional arithmetic support files are \texttt{r\allowbreak{}e\allowbreak{}s\allowbreak{}e\allowbreak{}a\allowbreak{}r\allowbreak{}c\allowbreak{}h\allowbreak{}/\allowbreak{}l\allowbreak{}o\allowbreak{}c\allowbreak{}a\allowbreak{}l\allowbreak{}\_\allowbreak{}c\allowbreak{}y\allowbreak{}c\allowbreak{}l\allowbreak{}e\allowbreak{}0\allowbreak{}8\allowbreak{}0\allowbreak{}/\allowbreak{}m\allowbreak{}a\allowbreak{}n\allowbreak{}u\allowbreak{}s\allowbreak{}c\allowbreak{}r\allowbreak{}i\allowbreak{}p\allowbreak{}t\allowbreak{}\_\allowbreak{}c\allowbreak{}a\allowbreak{}n\allowbreak{}d\allowbreak{}i\allowbreak{}d\allowbreak{}a\allowbreak{}t\allowbreak{}e\allowbreak{}/\allowbreak{}C\allowbreak{}H\allowbreak{}E\allowbreak{}C\allowbreak{}K\allowbreak{}\_\allowbreak{}F\allowbreak{}I\allowbreak{}N\allowbreak{}I\allowbreak{}T\allowbreak{}E\allowbreak{}.\allowbreak{}p\allowbreak{}y} and \texttt{r\allowbreak{}e\allowbreak{}s\allowbreak{}e\allowbreak{}a\allowbreak{}r\allowbreak{}c\allowbreak{}h\allowbreak{}/\allowbreak{}l\allowbreak{}o\allowbreak{}c\allowbreak{}a\allowbreak{}l\allowbreak{}\_\allowbreak{}c\allowbreak{}y\allowbreak{}c\allowbreak{}l\allowbreak{}e\allowbreak{}0\allowbreak{}8\allowbreak{}0\allowbreak{}/\allowbreak{}m\allowbreak{}a\allowbreak{}n\allowbreak{}u\allowbreak{}s\allowbreak{}c\allowbreak{}r\allowbreak{}i\allowbreak{}p\allowbreak{}t\allowbreak{}\_\allowbreak{}c\allowbreak{}a\allowbreak{}n\allowbreak{}d\allowbreak{}i\allowbreak{}d\allowbreak{}a\allowbreak{}t\allowbreak{}e\allowbreak{}/\allowbreak{}F\allowbreak{}I\allowbreak{}N\allowbreak{}I\allowbreak{}T\allowbreak{}E\allowbreak{}.\allowbreak{}j\allowbreak{}s\allowbreak{}o\allowbreak{}n}. H.7 identifies the existing exact count-box and return data by their packed addresses and hashes. The new program uses rational arithmetic and standard-library operations for the displayed endpoint/corner substitutions, independent of donor programs; its output does not replace the analytic proofs.

The author sections used for this assembly are the Local079 regular two-ray calculation and Local080 signed source calculation, the latter's \texttt{H\allowbreak{}2\allowbreak{}5\allowbreak{}\_\allowbreak{}P\allowbreak{}R\allowbreak{}O\allowbreak{}O\allowbreak{}F\allowbreak{}.\allowbreak{}m\allowbreak{}d} at SHA256 \texttt{8\allowbreak{}4\allowbreak{}4\allowbreak{}7\allowbreak{}3\allowbreak{}2\allowbreak{}4\allowbreak{}3\allowbreak{}3\allowbreak{}f\allowbreak{}9\allowbreak{}a\allowbreak{}8\allowbreak{}e\allowbreak{}1\allowbreak{}2\allowbreak{}c\allowbreak{}9\allowbreak{}2\allowbreak{}3\allowbreak{}d\allowbreak{}1\allowbreak{}d\allowbreak{}8\allowbreak{}e\allowbreak{}9\allowbreak{}d\allowbreak{}0\allowbreak{}2\allowbreak{}e\allowbreak{}8\allowbreak{}a\allowbreak{}d\allowbreak{}3\allowbreak{}b\allowbreak{}3\allowbreak{}f\allowbreak{}9\allowbreak{}6\allowbreak{}e\allowbreak{}f\allowbreak{}e\allowbreak{}f\allowbreak{}e\allowbreak{}8\allowbreak{}1\allowbreak{}9\allowbreak{}3\allowbreak{}6\allowbreak{}9\allowbreak{}8\allowbreak{}6\allowbreak{}9\allowbreak{}8\allowbreak{}d\allowbreak{}d\allowbreak{}1\allowbreak{}f\allowbreak{}8\allowbreak{}2\allowbreak{}6\allowbreak{}8\allowbreak{}7\allowbreak{}c}, and the Local080 near finite ledger \texttt{H\allowbreak{}7\allowbreak{}\_\allowbreak{}L\allowbreak{}E\allowbreak{}D\allowbreak{}G\allowbreak{}E\allowbreak{}R\allowbreak{}.\allowbreak{}m\allowbreak{}d} at SHA256 \texttt{9\allowbreak{}5\allowbreak{}d\allowbreak{}d\allowbreak{}9\allowbreak{}0\allowbreak{}c\allowbreak{}9\allowbreak{}0\allowbreak{}a\allowbreak{}7\allowbreak{}e\allowbreak{}f\allowbreak{}d\allowbreak{}7\allowbreak{}7\allowbreak{}7\allowbreak{}c\allowbreak{}c\allowbreak{}6\allowbreak{}8\allowbreak{}d\allowbreak{}5\allowbreak{}7\allowbreak{}8\allowbreak{}6\allowbreak{}9\allowbreak{}7\allowbreak{}9\allowbreak{}a\allowbreak{}0\allowbreak{}1\allowbreak{}5\allowbreak{}b\allowbreak{}a\allowbreak{}2\allowbreak{}2\allowbreak{}f\allowbreak{}2\allowbreak{}6\allowbreak{}2\allowbreak{}8\allowbreak{}e\allowbreak{}6\allowbreak{}d\allowbreak{}5\allowbreak{}4\allowbreak{}a\allowbreak{}c\allowbreak{}0\allowbreak{}7\allowbreak{}d\allowbreak{}2\allowbreak{}9\allowbreak{}5\allowbreak{}d\allowbreak{}0\allowbreak{}4\allowbreak{}7\allowbreak{}e\allowbreak{}9\allowbreak{}4\allowbreak{}a\allowbreak{}d}. The calculations and proofs required from them are included here; their private research status is not invoked as a mathematical premise.

\hypertarget{appendix-i}{}

\hypertarget{appendix-i.-the-whole-convolution-return-at-174669200000}{%
\section{Appendix I. The whole-convolution return at 174669/200000}\label{appendix-i.-the-whole-convolution-return-at-174669200000}}

\hypertarget{appendix-i-1}{}

\hypertarget{i.1.-the-improvement-and-the-retained-foundations}{%
\subsection{I.1. The improvement and the retained foundations}\label{i.1.-the-improvement-and-the-retained-foundations}}

Put \[\fitmath{
 \theta_1=\frac{174669}{200000}=.873345,\qquad
 c_1=1-\theta_1=\frac{25331}{200000}=.126655.
}\tag{I1}\] We prove the ordinary all-height conclusion at this value on the retained foundations of A--H. Appendix H is the previous result, at \(\fitinline{\theta_H=.87334585}\). The exact lesser coefficient \(\fitinline{1266543/10^7=.12665430}\) follows by inclusion of half-planes. The case \(\fitinline{c=.127}\) remains open.

Let \(\fitinline{\alpha,h,g_0,g_1}\) be the exact contact data of E and G, with their rational enclosures retained in H: \[\fitmath{
 \begin{gathered}
 .815324503<h<.815324505,\quad
 .05805911<g_0<.05805913,\quad
 .07599618<g_1<.07599620,\\
 \alpha<
 \frac{174669171221343715360237}{200000000000000000000000}
 <.873345856107.
 \end{gathered}
}\tag{I2}\] Set \(\fitinline{\ell_1=11/3-4\theta_1}\), \(\fitinline{y_1=h-2\ell_1}\), \(\fitinline{x_1=1+\ell_1-h}\), and \[\fitmath{
 k_{\mathrm{phys}}=5/6-x_1/3-\ell_1/6.
}\tag{I3}\] The actual target-zero supremum satisfies \(\fitinline{\beta\le\theta_H<\alpha}\). In a contrary argument \(\fitinline{\beta>\theta_1}\), the source data retain their required equality \(\fitinline{D_{\mathrm{src}}.\mathrm{momentKappa}=2\beta-1}\). Only the natural count and capacity are conservatively priced at \(\fitinline{\kappa_*=2\alpha-1}\), using the monotonicity proved in FB2--FB6. No new equality for the source data is imposed.

All coefficient conventions of H.1 remain literal: ideal columns, six element units, zero extension at every presented nonunit, original complete canonical prime pools, their common relative ray where required by the inverse source, independent ordinary heights, main/error masks, actual deleted primes, seeds and fixed bad set. In particular the complete original \(\fitinline{Q_j}\)'s, not arbitrary prime subsets, provide the common mark responses. The new consumer below does not modify a M\"{o}bius coefficient or replace the physical denominator by the presentation that supplied a detector zero. SN9 remains the same adopted analytic input with an arbitrary fixed positive gap. No new formal verification or discharge of that foundational input is asserted.

The new ingredients are a larger two-ray marked moment, obtained from the already proved curved supplier E(19), and a weighted use of the entire denominator convolution. Sections I.4--I.6 verify that these estimates reach every original physical sector.

\hypertarget{appendix-i-2}{}

\hypertarget{i.2.-moments-before-markov-and-their-physical-weight}{%
\subsection{I.2. Moments before Markov and their physical weight}\label{i.2.-moments-before-markov-and-their-physical-weight}}

Fix one original physical row class \(\fitinline{\mathcal C}\), one main/error allocation, the actual permitted contour point and denominator presentation \(\fitinline{\psi_{\mathrm{den}}}\), and write \(\fitinline{U=Z^{D_c}}\), \(\fitinline{q=dp}\). Let \(\fitinline{R}\) be any admitted upper count exponent on this class. It need not be the minimum of all available counts. A separated denominator factor of length \(\fitinline{U^a}\) is normalized as \[\fitmath{
 f_u=U^{-da/2}|F_u|.
}\] For an inverse factor take \(\fitinline{j=2}\); for an ordinary factor take \(\fitinline{j=4}\). The original marked moment, divided by its genuine squared common mark response, has the form \[\fitmath{
 \sum_{u\in\mathcal C}f_u^j
 \ll U^{R-g+\epsilon}\mathcal H.
}\tag{I4}\] Here \(\fitinline{\mathcal H}\) is the actual fixed-profile seminorm and polynomial height factor. This is a moment statement for every magnitude threshold, before Markov.

We check its source scope. Appendix D(A3)--(A6), E24--E29, B.5(204) and B.6(209)--(211) apply to the same presented physical denominator, its deleted Euler factors and its common actual relative ray. For total marked width \(\fitinline{w}\), the normalized inverse exponent is \[\fitmath{
 1-dx-2qw,
}\] with \(\fitinline{\max(1,(1+5x)/6)}\) in place of \(\fitinline{1}\) on the amplified unmarked branch. The normalized ordinary fourth exponent is \[\fitmath{
 1-2dy-2qw.
}\tag{I5}\] The canonical inverse inlet retains its coefficient bound, deleted-base norm gate, whole identity-ray pools, common relative ray and independent relative heights. Here is its actual B.6 matching. After a permissible whole conjugation, write the common fixed base ray as \(\fitinline{a}\); take every first/base ray to be \(\fitinline{a}\), with common relative multiplier \(\fitinline{\mathbf1_{\mathfrak q}\overline a}\). On every actual physical good-prime slot, \[\fitmath{
 [a(\mathfrak p)\psi_u(\mathfrak p)]\overline{a(\mathfrak p)}
       =\psi_u(\mathfrak p).
}\] No character is divided at a nonunit: the fixed \(\fitinline{\mathfrak q}\)-mask and the original deleted primes remain. If a physical slot has norm height \(\fitinline{b_i}\) and the inverse has height \(\fitinline{t}\), choose its relative height \(\fitinline{v_i=-b_i-t}\), retaining the actual outside unit. The matching therefore gives the whole original \(\fitinline{Q_i}\), including its profile and coefficient. The deleted base is fixed before asymptotics and its \textbf{whole presented modulus norm} is below every live \(\fitinline{P_i}\) for sufficiently large \(\fitinline{U}\), since the finite positive slot widths have already been fixed. The empty slot list satisfies that condition vacuously. A finite normalization of a fixed bump is restored by homogeneity. The ordinary source retains its own two-profile uniformity and original initial keep. These are the global source contracts; a conditional improved-inverse field for one detector family, label and profile would not supply I4 for an arbitrary physical taper.

For later use put \[\fitmath{
 \begin{aligned}
 k_1&=1-p(1+2g_1),& V_d&=5/6-d,\\
 U_d&=1-p(1+6g_0-4g_1),&
 J_d&=3V_d+2dk_1,\\
 t_0&=1+dU_d/J_d,& b&=2t_0/3,\quad m=t_0/3.
 \end{aligned}
}\tag{I6}\] The admitted tangent count is \[\fitmath{
 R_b=1-d\{p(1+2g_0)+k_1b\}
     =1/6+V_dt_0.
}\tag{I7}\] The inverse and ordinary tangent margins in I4 are respectively \(\fitinline{g=dk_1(x-b)}\) and \(\fitinline{g=2dk_1(y-m)}\). The amplified inverse margin is \(\fitinline{g=V_d(t_0-x)}\). The six other faces are exactly H31 at the current fixed bin, with their actual admissions checked in I.6.

Suppose a physical numerator term has its genuine fourth-power price \(\fitinline{U^P}\). H\"{o}lder with exponents \(\fitinline{4,j,(3/4-1/j)^{-1}}\) gives \[\fitmath{
 \sum_{\mathcal C}|P_u|f_u
 \ll U^{P/4+3R/4-g/j+\epsilon}\mathcal H'.
}\tag{I8}\] Indeed their three reciprocals sum to one, and the cardinality factor is \(\fitinline{U^{R(3/4-1/j)}}\). Thus the saving is \(\fitinline{g/j}\), not the saving of a saturated-row count.

The same statement holds for the actual coefficient packet without assuming a fourth moment at an undilated quota. First freeze its genuine numerator-dilation band, masks and row class, and retain the original fourth at that dilation and its damping. Equivalently partition \(\fitinline{f_u\asymp U^{-a}}\). Markov applied to I4 gives \(\fitinline{\#\{f_u\asymp U^{-a}\}\ll U^{R-g+ja+\epsilon}\mathcal H}\). On that subset the original packet fourth and its absolute packet mass comparison in B.7/H.7 remain valid. Relative to the unweighted physical price, their saving is \[\fitmath{
 a+\frac34(g-ja)_+\ \ge\ g/j,\qquad j=2,4.
}\tag{I9}\] For \(\fitinline{a\le g/j}\) the left side is affine and attains its minimum \(\fitinline{g/j}\) at the endpoint; for \(\fitinline{a\ge g/j}\) it is at least \(\fitinline{a}\). Buffered pointwise bounds give \(\fitinline{f_u\ll U^\epsilon}\), so a negative level of size \(\fitinline{O(\epsilon)}\) is harmless. Use a sufficiently small fixed exponent mesh; levels below \(\fitinline{U^{-K}}\), \(\fitinline{K>g+1}\), are paid by the original unweighted comparison. Apply the packet's dilation damping on each class before summing its coefficients. This proves I8 for the genuine packet and all its original main/error allocations. A mark used in I4 and a numerator mark used in its own fourth are separate positive comparisons; neither is inserted into the physical integrand a second time.

Now take \(\fitinline{f=\psi_{\mathrm{den}}}\), \(\fitinline{h=\mu\psi_{\mathrm{den}}}\) with their actual nonunit zeros, so \(\fitinline{f*h=\mathbf1}\). Let \(\fitinline{h_A(I)=h(I)V(NI/A)}\), where \(\fitinline{V=1}\) through \(\fitinline{1}\), vanishes from \(\fitinline{2}\), and \[\fitmath{
 A=U^{t_0-\gamma}.
}\] The exact identity \[\fitmath{
 h=h_A-(f*h_A-\mathbf1)*h
}\tag{I10}\] retains the M\"{o}bius sign. Its first-two-product coefficient vanishes up to \(\fitinline{A}\); insert the same smooth product cutoff as in H34. On a long annulus with \(\fitinline{r>t_0-\gamma}\) by a fixed margin the \(\fitinline{h_A}\) term vanishes. The original three annular partitions, including units, and two log-Fourier separations give \textbf{all} triples \[\fitmath{
 \begin{gathered}
 M_u(x;t+\xi_1+\xi_2)\,
 S_u(y;t+\xi_1+\xi_2)\,
 M_u(z;t+\xi_1),\\
 x\le t_0-\gamma+o(1),\quad
 x+y\ge t_0-\gamma-o(1),\quad x+y+z=r+o(1).
 \end{gathered}
}\tag{I11}\] The numerator height stays independent. The signs and filters remain attached until the exact identity is applied.

On the original central line, the exterior scalar of this inverse dyad is \(\fitinline{U^{-dr/2+O(\epsilon)}}\). Bounding two companions by the buffered B.5--B.6 bounds \(\fitinline{U^{d\,\mathrm{length}/2+\epsilon}}\) leaves precisely the third factor normalized as in I4. This retains the size of the full inverse at its physical weight.

Partition all triples in order. If \(\fitinline{x\ge b+\lambda_x}\), use I8 for that inverse. For the remaining triples with \(\fitinline{y\ge m+\lambda_y}\), use I8 for that ordinary. All other triples satisfy \[\fitmath{
 m-\gamma-\lambda_x-o(1)\le y<m+\lambda_y.
}\tag{I12}\] Suppose the actual two-ordinary marked source at width \(\fitinline{W}\) has positive energy \(\fitinline{U^{1+\epsilon}}\), with both independent heights and first log tests. Mark division on the same class gives the energy \(\fitinline{U^{1-2qW+\epsilon}}\). In the last case of I12, remove the two inverse companions using the normalization just proved, and apply Cauchy with the original count \(\fitinline{U^R}\). The physical exponent is \[\fitmath{
 E_{\mathrm{mid}}=\frac{R+1-dy-2qW}{2}.
}\tag{I13}\] There is no saturation debit: every triple has been paid, including the first two cases. No existential witness was used to restrict the full convolution.

For completeness the varying taper on a dyad is covered by the actual global inlet. Write it on one fixed compact logarithmic bump, Fourier-invert its ratio \(\fitinline{U^x/A}\), and integrate against a common weighted \(\fitinline{L^1}\) envelope as in H24a. On transition dyads the ratio is bounded; off the transition the taper is constant, so all fixed log derivatives are uniformly bounded. Near \(\fitinline{x=t_0-\gamma}\) use the native unmarked/amplified inverse source, whose coefficient remains \(\fitinline{\mu\psi_{\mathrm{den}}}\). The improved inverse anchor lies a fixed power below the taper and uses its fixed annular profile. Truncate each separation shift at \(\fitinline{T_1/16}\) inside the original \(\fitinline{T_1/4}\) buffer, then choose the decay order against the finite polynomial envelope. For the buffered reciprocal estimate first choose \(\fitinline{e>0}\) sufficiently small and then \(\fitinline{\epsilon\ge6e}\); the bounded descendant lengths make its line-shift cost fit the fixed positive reserve. Unit-interval Sobolev includes the first normalized log test on every varying-height factor. These arguments prove I4--I13 with the original independent heights, not with a new profile-specific constant chosen at each Fourier point.

\hypertarget{appendix-i-3}{}

\hypertarget{i.3.-a-curved-two-ray-moment-with-100-micro-units-of-surplus}{%
\subsection{I.3. A curved two-ray moment with 100 micro-units of surplus}\label{i.3.-a-curved-two-ray-moment-with-100-micro-units-of-surplus}}

For \[\fitmath{
 .42495\le n\le.42522,\qquad .37445\le y\le.378,
}\tag{I14}\] put \[\fitmath{
 P(V)=V/2-n,\quad
 S(V)=(1/2+g_0)V-(1+2g_1)y,\quad
 W=\frac{P(1)+S(1)}2+10^{-4}.
}\tag{I15}\] We prove the H27 conclusion at every width strictly below \(\fitinline{W}\), including componentwise shortenings, the actual two ordinary rays and independent heights, and both first log tests. We use only the curved source already proved in E(19).

Write \[\fitmath{
 \begin{aligned}
 I(r)&=11r/16-7/24-\frac5{16}[7r^4(r-2/3)]^{1/5},\\
 s(r)&=(1-r)/2+I(r),\qquad 3/4\le r<7/9.
 \end{aligned}
}\tag{I16}\] E(19) supplies both actual inverse and ordinary open rectangles with this gain. To use a requested free width \(\fitinline{r}\), take a nearby larger anchor and a positive interior decrement. On a compact interval below \(\fitinline{7/9}\), finitely many such rectangles cover the widths with an arbitrarily small prescribed quota loss. Fix this finite bank before the early mesh and later slot type; no analytic constant passes through a limiting curve.

Set \[\fitmath{
 u_0=3/2000,\quad H_0=1-6u_0=991/1000,\quad
 \rho=1/500000.
}\tag{I17}\] On a regular sixth-free core of exponent \(\fitinline{V}\), let \(\fitinline{e}\) be the squared-response exponent of the clean maximizing prime slot. For \(\fitinline{.98\le V\le1}\), the shorter ordinary's available absolute quota is \[\fitmath{
 B(V)=V\,s\!\left(\max\{3/4,2y/V\}\right).
}\tag{I18}\] The factor two is required for its fourth moment. At the lower corner \(\fitinline{2y/V=.7489}\), use downward inclusion at \(\fitinline{3/4}\), not a curve value below its domain. The two formulas agree at the crossing. The upper ratio is \(\fitinline{2(.378)/.98<7/9}\).

Repeat the same clean maximizing canonical slot, separately re-anchored to each tested ordinary, up to its respective decremented quota. On common good primes the character identity is \(\fitinline{a\,\overline{(a\bar\gamma)}=\gamma}\); retain all presented zeros and fixed prime deletions. H24a supplies the normalized relative-frequency adjustment and keeps the outside phase. Each fourth has its own initial keep, profile and height. Repetition loses at most one early slot width; zero repetitions use the unmarked face. These repetitions are source calls only, not the physical lower-response product. Cauchy on the same positive class gives excess \[\fitmath{
 C_0(V)=e\{L-(P(V)+B(V))/2\},\qquad L=W-u_0.
}\tag{I19}\]

For the middle response range use the common buffered finite-family zero band from H.4, with \(\fitinline{e\le\nu\le3/4}\). Its one attaining zero supplies both supported witnesses. At a fixed cutoff \(\fitinline{3w/2}\), \(\fitinline{.75\le w\le.775}\), a short inverse or doubled ordinary witness has width \(\fitinline{w'\ge w-o(1)}\). E(19) gives its deficit \(\fitinline{\nu w'+e s(w')}\) up to \(\fitinline{w'=.7775}\). For the inverse this is the same fixed-profile M\"{o}bius source; the cutoff equals one on this annulus, since \(\fitinline{.7775<1.125}\). The original common relative ray, complete pool, \(\fitinline{33/50}\) profile, independent relative frequency and outside unit are unchanged. Fixed compact Fourier transfer covers its real powers and log tests. Hold the external mark fixed during witness-height Sobolev.

The radical in I16 is concave, so \(\fitinline{I}\) is convex and decreasing, with \[\fitmath{
 I'(3/4)>-.09,\qquad 0<I(r)\le I(3/4)<.0011,\qquad I'(7/9)=0.
}\tag{I20}\] Hence \(\fitinline{\nu+e s'(r)>0}\). Above \(\fitinline{.7775}\), native width \(\fitinline{(1-w')/2}\), with its strict decrement and unmarked fallback near one, also pays the desired deficit: \(\fitinline{\nu-e/2\ge e/2}\), and \(\fitinline{(.7775-.775)/2=.00125>\max I}\). The disjoint-list inverse convention is handled by H.4's Newton/color expansion. If the doubled ordinary width is at least one, its unmarked fourth pays \(\fitinline{\nu}\), which dominates the same deficit. The long inverse uses its amplified second moment. Thus the two alternative count deficits and energy excesses are \[\fitmath{
 \begin{aligned}
 D_s&=\nu w+e s(w),&
 D_l&=5/6-(5/6-\nu)3w/2,\\
 C_i(V)&=\nu y+e[L-P(V)/2]-VD_i/2,\qquad i=s,l.
 \end{aligned}
}\tag{I21}\] This Cauchy uses the longer ordinary's marked fourth and the shorter ordinary's buffered pointwise bound. Both witnesses belong to the same zero. A finite bank of cutoffs is chosen before rows are evaluated.

Here is the exact finite proof on the central radius \(\fitinline{V=H_0}\). It also specifies the mathematical meaning of the data in research/\allowbreak{}local\_cycle081/\allowbreak{}direct/\allowbreak{}CHECK.py and CERTIFICATE.json. Use the worst upper-width endpoints \(\fitinline{g=.05805913}\), \(\fitinline{h_1=.07599618}\). Split the \(\fitinline{y}\)-interval in I14 into intervals \(\fitinline{[a,b]}\) of width \(\fitinline{1/40000}\). Let \(\fitinline{r_*=(a+b)/H_0}\), and enclose the radical \[\fitmath{
 c(r)=\frac5{16}[7r^4(r-2/3)]^{1/5}
}\] by the least multiple \(\fitinline{\widehat c(r)}\) of \(\fitinline{10^{-13}}\) above it. This enclosure is proved by the two adjacent integer fifth-power comparisons. Concavity and \(\fitinline{c'=A(r)c}\), \(\fitinline{A(r)=4/(5r)+1/[5(r-2/3)]}\), give the rational minorant \[\fitmath{
 I_-(r)=11r/16-7/24-\widehat c(r_*)[1+A(r_*)(r-r_*)].
}\tag{I22}\] The bracket is positive on each interval. This same enclosure at \(\fitinline{3/4}\), the explicit derivative formula and \(\fitinline{I'(7/9)=0}\) prove I20.

The variable \(\fitinline{n}\) cancels exactly from I19 and I21 at \(\fitinline{H_0}\). Put \[\fitmath{
 \begin{aligned}
 B_*&=\max_{z=a,b}\{(1-H_0)/2+g/2-h_1z
           -H_0I_-(2z/H_0)/2+10^{-4}-u_0\},\\
 k_*&=10^{-6}
 \left\lfloor10^6(u_0-\rho-10^{-8})/B_*\right\rfloor .
 \end{aligned}
}\tag{I23}\] Since the substituted expression is affine, \(\fitinline{e\le k_*}\) gives \(\fitinline{C_0(H_0)\le u_0-\rho}\). For \(\fitinline{k_*\le e\le.5002}\), \(\fitinline{e\le\nu\le3/4}\), cover the closed quadrilateral by its two triangles with vertices \[\fitmath{
 A=(k_*,k_*),\quad B=(.5002,.5002),\quad
 C=(3/4,.5002),\quad D=(3/4,k_*).
}\] The first coordinate is \(\fitinline{\nu}\). The expression \(\fitinline{\nu y+e[L-P(H_0)/2]}\) increases in \(\fitinline{y}\), since \(\fitinline{\nu\ge e}\) and \(\fitinline{h_1<1/2}\). At \(\fitinline{y=b}\) define \[\fitmath{
 A_b=(1-H_0)/4+[1/2+g-(1+2h_1)b]/2+10^{-4}-u_0,\qquad
 T(v,e)=2(vb+eA_b-u_0+\rho)/H_0.
}\tag{I24}\] On a triangle choose \(\fitinline{w}\) by rounding down to a multiple of \(\fitinline{10^{-7}}\) the minimum of \(\fitinline{.775}\) and \[\fitmath{
 \frac{5/6-T(v,e)}{(3/2)(5/6-v)}
}\] over its vertices. At all vertices check \[\fitmath{
 vw+e\{(1-w)/2+11w/16-7/24-\widehat c(w)\}\ge T(v,e),\quad
 5/6-(5/6-v)3w/2\ge T(v,e).
}\tag{I25}\] For fixed \(\fitinline{w}\) both inequalities are affine in \(\fitinline{(v,e)}\), so these vertex checks prove the entire closed triangle. If a check fails, bisect its longest edge into two closed children and repeat; the children cover their parent including the diagonal. This specifies the whole finite cover without an interpolation assumption.

The resulting exact data are 142 length intervals, 522 accepted triangles, and at most six bisections on a path. Every \(\fitinline{w}\) lies in \(\fitinline{[.75,.775]}\), and \(\fitinline{.447073\le k_*\le.458949}\). The low, short and long minima of \(\fitinline{u_0-C_i(H_0)}\) are respectively \[\fitmath{
 \frac{67427650524626278446191}{33545494700000000000000000000},\qquad
 \frac{678242247812771}{300000000000000000000},\qquad
 \frac{375003939239}{187500000000000000},
}\tag{I26}\] all exceeding \(\fitinline{\rho=2\cdot10^{-6}}\). The finite program records these rational substitutions; the proof of I22 and the affine implications I23--I25 explain why they certify the full domains.

For \(\fitinline{e\ge.5002}\), use H6 with \(\fitinline{v_{\mathrm{SN}}=.5001}\), fixed gap \(\fitinline{10^{-4}}\), on the actual merged two-ordinary/live-prime coefficient and one clean maximizing descriptor. The merged squared mass is \(\fitinline{U^\epsilon}\) by divisor-fibre Cauchy. The coefficient width ceiling below gives a saving exceeding \(\fitinline{.038}\) for \(\fitinline{V\ge.98}\). Fix this positive gap and the early mesh before the finite slot list; after that list is fixed take its bin tolerance below the \(\fitinline{.0001}\) separation, including division by the smallest positive slot width. No slope-\(\fitinline{1/2}\) SN9 claim is used.

The outer-radius proof of H.4 transfers with these verified constants. On \(\fitinline{[.98,1]}\), I19 and I21 decrease in \(\fitinline{V}\), with absolute slope less than \(\fitinline{5/6}\). On the curved piece \(\fitinline{B'=1/2+I-rI'<.573}\); on the clipped piece \(\fitinline{B'=s(3/4)}\). Both pieces agree at their crossing. For \(\fitinline{V\ge1}\) the native absolute quotas have mean greater than \(\fitinline{L}\). For \(\fitinline{.875\le V<.98}\) the positive native quotas leave at most \(\fitinline{3(1-V)/8<5(H_0-V)/6}\). Below \(\fitinline{.875}\) the two regular unmarked fourths and the original \(\fitinline{7/8}\) prime bound cost at most \(\fitinline{3L/4}\); below \(\fitinline{.08}\) the lift/core cost is bounded as in H.4. The corresponding strict guards, in the same order, include \[\fitmath{
 \frac{827204601}{10^{12}},\quad \frac1{600},\quad
 \frac{264744841409}{12000000000000},\quad
 \frac{171836613803}{3000000000000}.
}\tag{I27}\] Writing the actual element as \(\fitinline{b^6c}\), with no coprimality assumption, restores every cofactor and both punctures by the same \(\fitinline{5/6}\) slope comparison. Hence the regular energy on its actual common row ball is \[\fitmath{
 \ll \mathcal H U^{M+u_0-\rho+(5/6)(H_0-M)_++\epsilon}.
}\tag{I28}\]

It remains to check the literal H.5 restoration rather than infer a raw estimate from one leaf. The same certificate gives \[\fitmath{
 W\le\frac{100972795399}{10^{12}}<.101,\qquad
 A_0\le\frac{2818790653}{3125000000}<.903.
}\tag{I29}\] The original selected radius satisfies \(\fitinline{M\ge1-6u'-O(\mathrm{mesh})}\), where a filled quota has \(\fitinline{u'=u_0+O(\mathrm{mesh})}\). Thus I28 pays it with a positive reserve. Exhaustion uses the two actual unmarked fourths. The signed fixed-row proof H8--H18 is unchanged: \(\fitinline{2(3/5)\ge1}\) still pays the full second-divisor extraction, and its coarse spare \(\fitinline{5/6+.139-.903-3u_0>.0648}\) is positive. The first artificial divisor stays unweighted; the second label's \(\fitinline{1/N\mathfrak D}\) cancels its matching H26 normalizer. Both first orientations retain their own reflected ray, radial ball and sublattice, and both budgets have denominator \(\fitinline{NC\,ND_1}\) by the two distinct identities in H.5.

The short reference has its original \(\fitinline{.139}\) coordinate and companion \(\fitinline{n+y-.139\le.66422<.665}\). The earlier \(\fitinline{.663}\) was a box ceiling, not a restriction of the arbitrary bounded-length unmarked fourth or fixed signed H8. Its short marked fourth only needs \(\fitinline{4(.139+W)<H_0}\), whose strict slack is \(\fitinline{7777204601/250000000000}\). The two zero gates \(\fitinline{5u'<1}\), \(\fitinline{A_0+4u'<1}\), the genuine zero-frequency coefficient, harmonic seed sum, and the counted \(\fitinline{U^{-31}}\) tail retain their H.5 prices. The first zero/tail still has \(\fitinline{A_0<.903}\). Thus H25--H27, inverse restoration and the kept short/all-row centered comparison H.5(A)--(C) prove \[\fitmath{
 \sum_{u\in\mathcal C}
 |P_u(n,v)S_u(y,t')\prod_{j\in I}Q_j(u)|^2
 \ll U^{1+\epsilon}(1+|v|+|t'|+T)^J
}\tag{I30}\] on the original kept class throughout I14, with actual independent heights, profiles and normalized first log tests.

\hypertarget{appendix-i-4}{}

\hypertarget{i.4.-the-complete-exterior-and-the-additional-admitted-wing}{%
\subsection{I.4. The complete exterior and the additional admitted wing}\label{i.4.-the-complete-exterior-and-the-additional-admitted-wing}}

Use the original 57 closed count boxes, capacities and cleared lower polynomials of E/G/H. Their exact data remain at research/\allowbreak{}local\_cycle069/\allowbreak{}continuum/\allowbreak{}RETURN\_CERTIFICATE.json. The natural intervals are \(\fitinline{D_c\in[1/200,19/25]}\), with \(\fitinline{C_O=\min(D_c,2y_1-D_c)/(6\kappa_*)}\), and \(\fitinline{D_c\in[19/25,h]}\), with \(\fitinline{C_N=\min(3D_c/40,(2y_1-D_c)/2)}\). The returned exponent is the same algebraic expression as H37a: \[\fitmath{
 E(D_c,C;d,p)=\tfrac12+13h/75-y_1/2-\ell_1/2-9D_c/100
 +d/2+dp(\ell_1-C/2)+(3D_c/4)R.
}\tag{I31}\] The target changes the capacities and geometric powers; it does not change a count identity.

For a saved lower polynomial on a box, substitute its affine unit-square coordinates and expand in the tensor Bernstein basis as in H.7. Its coefficients are exact rational finite sums; nonnegativity of the basis bounds the polynomial by their minimum. The branch denominator is \(\fitinline{0<J\le5/2}\), so multiplying that lower bound by \(\fitinline{2/5}\) gives a margin. At a native edge the target debit is at most \(\fitinline{(3+8d_+p_+)(\alpha-\theta_1)}\); at an old edge it is at most \(\fitinline{\{3+(4+4/(3\cdot.746))d_+p_+\}(\alpha-\theta_1)}\). These bounds follow by keeping the actual capacity denominator in I31, not by substituting a nominal final capacity.

One previously noncritical box needs a better supplier. Box 28 is the closed wing \[\fitmath{
 \mathcal W=[93/256,189/512]\times[63/128,1/2]
}\tag{I32}\] on the top native edge. Apply the same tangent paired source E24--E29 there. We give the full admission check. Set \(\fitinline{d,p}\) in this wing, \[\fitmath{
 V=5/6-d,\quad k_1=1-p(1+2g_1),\quad
 U=1-p(1+6g_0-4g_1),\quad
 J=3V+2dk_1,\quad T=J+dU,\quad N=(1-d)J+VdU.
}\] With \(\fitinline{r_0=1467/2000}\), \(\fitinline{r_1=311/400}\), \(\fitinline{B_0=1-10^{-12}}\), \(\fitinline{\Lambda=2+10^{-11}}\), \(\fitinline{c_*=3/20-10^{-7}}\), \(\fitinline{s_-=.027}\), \(\fitinline{m_H=(1-(9/2)s_-)/2}\), the ten strict guards are \[\fitmath{
 \begin{gathered}
 k_1,\ J,\ T-J,\ 1.2499J-T,\ 2T-3(r_0+10^{-6})J,\\
 N-\{1-d[r_1+(1-r_1)p]\}J,\\
 N-\{1-d[r_1+2p(B_0-r_1)/\Lambda]\}J,\\
 N-\{1-d(r_1+2pc_*)\}J,\quad
 N-(1-2m_Hd-2s_-dp)J,\\
 (1+6g_0-4g_1)J-2(1+2g_1)dU.
 \end{gathered}
}\tag{I33}\] Each is positive on the entire closed wing, at all four \(\fitinline{g_0,g_1}\) enclosure corners. The exact lower bounds are the four attached guard records in research/\allowbreak{}local\_cycle081/\allowbreak{}near\_record/\allowbreak{}EXTERIOR\_FINAL.json. They use the same rational lower-polynomial/Bernstein operation just described. The expressions are separately affine in the enclosed source constants after the stated substitution; their corner bounds therefore cover those constants as well. These are the positivity, anchor placement, inverse and ordinary fallback, and crossing-monotonicity gates of E29, with the real fixed positive source decrements. Thus the new count branch is admitted on the whole wing, rather than inferred from a numerical improvement.

On the 20 original improved boxes the identity \[\fitmath{
 M_h(d,p)=K_h(d-\delta_*)^2+(1/2-p)dQ(d,p),
 \quad K_h>1/5,\quad .4258<\delta_*<.4259
}\tag{I34}\] is unchanged. Interval propagation of the actual \(\fitinline{q_{10},q_{20},q_{21}}\) gives \(\fitinline{Q>0}\) also on I32. Since the wing lies below \(\fitinline{\delta_*}\), its first term alone gives \[\fitmath{
 \frac{(1/5)(.4258-189/512)^2}{5/2}
 -(3+8(189/512)/2)(\alpha-\theta_1)>0.
}\tag{I35}\] All 36 remaining noncritical boxes retain their saved branches. Hence the cover consists of 36 saved noncritical boxes, the newly admitted wing, and the 20 original improved boxes. The original edge counts \(\fitinline{1,2,2,52}\) and all closed rectangles are unchanged. Their union is checked by sorting their rational coordinate endpoints, testing each resulting elementary cell in an assigned closed rectangle, and retaining the closed boundary edges; no missing cell is supplied by continuity outside an admission window.

The critical square identity gives two unrestricted complements. Outside \[\fitmath{
 d_-=.4185\le d\le d_+=.4333,\qquad p>p_-=.49954,
}\tag{I36}\] use respectively the square-distance lower bound in I34 and \(\fitinline{(1/2-p)dQ/J}\). For every \(\fitinline{d\in[d_-,d_+]}\), throughout the entire original critical amplitude window \(\fitinline{15/32\le p\le1/2}\), exact interval propagation gives \[\fitmath{
 Q\ge\frac{266993408049562748296338051173}
 {7324218750000000000000000000000}>.03645,\qquad
 J\le\frac{1314511282809}{800000000000}<1.644.
}\tag{I37}\] The full exterior certificate retains the sharper actual values. Subtracting the complete target debit gives density reserve greater than \(\fitinline{8.96\cdot10^{-8}}\) and amplitude reserve greater than \(\fitinline{2.18\cdot10^{-7}}\). Unlike the narrower previous return, no radical filter is needed: every original row, including every good valuation one through five and every error-slot tuple, remains in the middle ticket. All original masks and coefficient masses are retained.

For fixed \(\fitinline{d,p}\), I31 is the maximum of affine functions of \(\fitinline{D_c}\): its capacities are minima of affine functions and have nonpositive coefficient. Thus the endpoint comparisons fill both original conductor intervals. This proves the full exterior on all 57 original boxes.

\hypertarget{appendix-i-5}{}

\hypertarget{i.5.-geometry-contour-complements-and-the-actual-middle-ticket}{%
\subsection{I.5. Geometry, contour complements and the actual middle ticket}\label{i.5.-geometry-contour-complements-and-the-actual-middle-ticket}}

We construct the generic geometric instance directly at the new target. In exact arithmetic \(\fitinline{\ell_1=25993/150000}\), so \(\fitinline{.17328<\ell_1<.17329}\). Together with I2, the definitions \(\fitinline{x_1=1+\ell_1-h}\), \(\fitinline{y_1=h-2\ell_1}\) satisfy every field of Geometry.Parameters: \[\fitmath{
 \begin{gathered}
 0<\ell_1,\quad 1/6<\ell_1<1/5,\quad
 \ell_1<x_1<y_1<1/2,\quad x_1+y_1=1-\ell_1,\\
 0<2x_1-y_1-\ell_1,\qquad
 1+3\ell_1-2(x_1+y_1)<0,\\
 \theta_1-(4+y_1-x_1)/6=x_1/2+(y_1-x_1)/12.
 \end{gathered}
}\tag{I38a}\] For the nontrivial inequalities, \(\fitinline{x_1-\ell_1=1-h>0}\), \(\fitinline{y_1-x_1=2h-1-3\ell_1>.110779006}\), \(\fitinline{y_1<.468764505<1/2}\), and \(\fitinline{2x_1-y_1-\ell_1=2+3\ell_1-3h>.073866485}\). The tuple margin is \(\fitinline{5\ell_1-1=-4007/30000<0}\). The last equality is equivalent to \(\fitinline{\theta_1=2/3+(x_1+y_1)/4=11/12-\ell_1/4}\), the exact defining identity. It also identifies \(\fitinline{k_{\mathrm{phys}}=(4+y_1-x_1)/6}\).

The five BoundsCore fields follow separately: \[\fitmath{
 \begin{gathered}
 .873<\theta_1<7/8,\quad .793<h<.816,\quad
 .458<y_1<.473,\quad .1672<\ell_1<.175,\\
 .07<y_1/2-h(17/50-1/6)-1/50.
 \end{gathered}
}\tag{I38b}\] Indeed \(\fitinline{.468744503<y_1<.468764505}\), and the last expression is \(\fitinline{49h/150-\ell_1-1/50>.073049}\). Thus the exact dilation, tuple, signal and coarse hypotheses are established independently at \(\fitinline{\theta_1}\). The earlier targetGeometry/targetCoarse constructor, whose corridor starts at the previous algebraic endpoint, is not used.

Fix the numerical pad \(\fitinline{F_{\mathrm{pad}}=\Delta=10^{-258}}\) and the positive source parameter \(\fitinline{s=10^{-260}}\). The baseline data object \(\fitinline{F_{\mathrm{src}}}\) is different from this number. In a contrary argument \(\fitinline{\beta>\theta_1}\), request the actual finite baseline data \(\fitinline{D_{\mathrm{src}},F_{\mathrm{src}}}\) and transport data \(\fitinline{P_{\mathrm{src}}}\) with loss below \[\fitmath{
 \min\{(\beta-\theta_1)/10000,\ 10^{-9}/4,\ F_{\mathrm{pad}}\},
 \qquad 0<D_{\mathrm{src}}.t\le F_{\mathrm{pad}}/4,
}\tag{I38}\] before choosing the target character and live heights. Then the literal buffered pad obeys \(\fitinline{F_{\mathrm{pad}}/2+2D_{\mathrm{src}}.t\le F_{\mathrm{pad}}}\). The separate exponent debit is \(\fitinline{15s/16+2F_{\mathrm{pad}}>F_{\mathrm{pad}}}\), charged to every sector. Capacities are 1-Lipschitz in \(\fitinline{D_c}\), \(\fitinline{0\le R\le1}\), and I31 has conductor variation at most \(\fitinline{(411/400)|\Delta D_c|<2|\Delta D_c|}\). The full exterior therefore extends to \(\fitinline{D_c\le h+F_{\mathrm{pad}}}\).

For the remaining critical patch the top-edge debit is at most \[\fitmath{
 B_*=(3+4d_+)(\alpha-\theta_1)
 < \frac{10130314131}{2500000000000000}
 =4.0521256524\cdot10^{-6}=:B.
}\tag{I39}\] We use this slightly larger rational \(\fitinline{B}\) for the whole-convolution join. The exterior certificate uses the sharper exact contact bound from I2. Put \(\fitinline{D_-=.815324503-15\cdot10^{-6}=.815309503}\). The conductor return slope is at least \[\fitmath{
 a_D=\tfrac34(1-d_+)-9/100-3d_+/160.
}\] Consequently the sector \(\fitinline{D_c\le h-15\cdot10^{-6}}\) saves \(\fitinline{a_D(15\cdot10^{-6})-B_*}\).

The other original complements use the following unchanged mechanisms, with their new widths:

\begin{longtable}[]{@{}
  >{\raggedright\arraybackslash}p{(\columnwidth - 4\tabcolsep) * \real{0.2727}}
  >{\raggedleft\arraybackslash}p{(\columnwidth - 4\tabcolsep) * \real{0.3636}}
  >{\raggedleft\arraybackslash}p{(\columnwidth - 4\tabcolsep) * \real{0.3636}}@{}}
\toprule\noalign{}
\begin{minipage}[b]{\linewidth}\raggedright
Complement
\end{minipage} & \begin{minipage}[b]{\linewidth}\raggedleft
Width or response
\end{minipage} & \begin{minipage}[b]{\linewidth}\raggedleft
Bare reserve after \(\fitinline{B_*}\)
\end{minipage} \\
\midrule\noalign{}
\endhead
\bottomrule\noalign{}
\endlastfoot
low conductor & \(\fitinline{15\cdot10^{-6}}\) & \(\fitinline{>8.51\cdot10^{-7}}\) \\
short numerator & \(\fitinline{n\le n_0-50\cdot10^{-6}}\) & \(\fitinline{>1.23\cdot10^{-7}}\) \\
large coefficient dilation & \(\fitinline{v_{\mathrm{dil}}>50\cdot10^{-6}}\) & \(\fitinline{>2.12\cdot10^{-7}}\) \\
short inverse length & \(\fitinline{r<1/D_c-55\cdot10^{-6}}\) & \(\fitinline{>4.32\cdot10^{-7}}\) \\
long inverse length & \(\fitinline{r>1/D_c+20\cdot10^{-6}}\) & \(\fitinline{>5.13\cdot10^{-7}}\) \\
\end{longtable}

Here \(\fitinline{n_0=1-y_1/D_c}\). For the short numerator use quota \[\fitmath{
 z_0+49\cdot10^{-6},\qquad
 z_0=\frac{1-s-2n_0}{2+10s};
}\] it is below \(\fitinline{3/40}\), and \(\fitinline{2(50\cdot10^{-6})>(2+10s)(49\cdot10^{-6})}\). The scalar for a live coefficient is not divided out. Its genuine dilation \(\fitinline{v_{\mathrm{dil}}=\log_U\max(1,L_iK)}\), the absolute packet bound and capacity clipping save \(\fitinline{dv_{\mathrm{dil}}/4}\). Retaining the complementary dilation and allowing annular slack \(\fitinline{0.1\cdot10^{-6}}\) gives \(\fitinline{n\le n_0+50.1\cdot10^{-6}}\).

For inverse-length tails, open the actual denominator after full \(\fitinline{w}\)-inversion on its original common class. The quotient-free Euler expansion is holomorphic on the two contour moves \(\fitinline{\Delta_s=-.10,.28}\). The exterior inverse power restores \(\fitinline{U^{\Delta_s(1/D_c-r)+o(1)}}\). The variable damping \(\fitinline{\rho_d=d/2+\Delta_s}\) pays all coefficient shifts, since \[\fitmath{
 \rho_d-q/2\ge .4185/2-.10-.4333/4
 =37/40000>0.
}\tag{I40}\] Both moved lines remain in the original contour region. Their two length margins give the table's last two prices. This partition is made after the legal contour motion and with the original Euler deletions. No lower amplitude bound for the full inverse is required by I10--I13; the optionally small inverse-response sector need not be removed.

All these bare reserves and the 57-box reserves, after subtracting \(\fitinline{15s/16+2F_{\mathrm{pad}}}\), exceed \(\fitinline{3\cdot10^{-9}}\). Choose the sum of the remaining source, mesh and reconstruction losses below \(\fitinline{10^{-9}}\). The domain guards remain \[\fitmath{
 h+F_{\mathrm{pad}}<\min(1,2y_1),\quad y_1<.76<h,\quad
 \ell_1/(h+F_{\mathrm{pad}})>1/5.
}\tag{I41}\] The averaged floor of E(31)--(40), including its original masks and finite grids, still has exponent \(\fitinline{.51-49h/150+.51\ell_1+.66D_c}\). At the upper cutoff its margin after \(\fitinline{10^{-5}+10^{-12}}\) exceeds \(\fitinline{.001}\), so its supported saving \(\fitinline{1/4000}\) is retained. The small-row margin and both principal contour margins remain strictly positive. The five nonprincipal units, fixed-pad large-row estimate, actual principal residues, quotient-free remaining contours, finite prime corrections, compensators, zero frequencies and Fourier tails retain their B/H proofs with these checked bounds. No principal pole is paid by a row-count saving.

Thus the only remaining ticket is \[\fitmath{
 \begin{gathered}
 h-15\cdot10^{-6}<D_c\le h+F_{\mathrm{pad}},\quad
 d\in[.4185,.4333],\quad p>.49954,\\
 n_0-50\cdot10^{-6}<n\le n_0+50.1\cdot10^{-6},\\
 -55\cdot10^{-6}\le r-1/D_c\le20\cdot10^{-6}.
 \end{gathered}
}\tag{I42}\] There is no radical predicate or inverse-amplitude saturation restriction on I42.

\hypertarget{appendix-i-6}{}

\hypertarget{i.6.-the-entire-convolution-on-the-retained-ticket}{%
\subsection{I.6. The entire convolution on the retained ticket}\label{i.6.-the-entire-convolution-on-the-retained-ticket}}

Use I10--I13 with \[\fitmath{
 \gamma=26\cdot10^{-6},\qquad
 \lambda_x=\lambda_y=60\cdot10^{-6},\qquad
 \Delta W=10^{-4}.
}\tag{I43}\] The exact rational enclosure in research/\allowbreak{}local\_cycle081/\allowbreak{}recombination/\allowbreak{}FINITE.json, case stretch, gives \[\fitmath{
 \begin{gathered}
 1.12401727168200<t_0<1.13327756075298,\\
 .42501352643752<n<.42512420492227,\\
 .37458642389400<y<.37781918691767
 \quad\hbox{in the middle case I12}.
 \end{gathered}
}\tag{I44}\] These displayed outer decimals are deliberately enlarged rational enclosures. The exact fractions in the attached data decide the comparisons. The \(\fitinline{n}\)-bound even replaces the positive pad by the larger \(\fitinline{10^{-12}}\), and the independent endpoint enclosure for \(\fitinline{t_0}\) is wider than needed. Hence I14 covers the whole middle rectangle, not just an optimally correlated pair of lengths. The lower retained inverse length exceeds \(\fitinline{1.22645051267163}\); it is above \(\fitinline{t_0-\gamma}\) by more than \(\fitinline{.09319}\), so the long-annulus use of I10 is valid.

We check every outside source admission. The cut \(\fitinline{b+\lambda_x}\), and doubled cut \(\fitinline{2(m+\lambda_y)}\), remain inside the genuine improved anchor interval \(\fitinline{[1467/2000,311/400]}\). At the two tangent edges, the raw moment margins are \(\fitinline{dk_1\lambda_x}\) and \(\fitinline{2dk_1\lambda_y}\); the amplified inverse edge is \(\fitinline{V_d\gamma}\). For the upper inverse use native width \(\fitinline{(1-x)/2}\) through \(\fitinline{.90}\), then the unmarked inverse through \(\fitinline{1}\), and the amplified inverse with unmarked width \(\fitinline{w=0}\) on \(\fitinline{1\le x\le t_0-\gamma}\). For the ordinary use the same capped/affine quota \(\fitinline{\min(c_*,(B_0-2y)/\Lambda)}\), then the native width \(\fitinline{(1-2y)/(9/2)}\) through \(\fitinline{.4398}\), and its unmarked fourth thereafter. All widths retain their original strict decrements, positive slopes and fixed profile contracts. They are precisely the six H31 faces and the three moving edges; no new inverse product estimate is inserted.

For clarity, one exact calculation of all these margins uses the independent enclosing endpoints of I6. Let \[\fitmath{
 \begin{gathered}
 V_-=5/6-d_+,\quad
 k_-=1-\tfrac12(1+2g_{1,+}),\quad
 k_+=1-p_-(1+2g_{1,-}),\\
 C_+=\tfrac12(1+2g_{0,+})+k_+(2t_{0,+}/3).
 \end{gathered}
}\] The inverse native and unmarked raw gaps are at least \[\fitmath{
 d_-\{r_1+p_-(1-r_1)-C_+\},\qquad d_-(.90-C_+).
}\] The four ordinary raw gaps are at least \[\fitmath{
 \begin{gathered}
 d_-\{r_1+2p_-c_*-C_+\},\quad
 d_-\{r_1+2p_-(B_0-r_1)/\Lambda-C_+\},\\
 d_-\{2m_H+2p_-s_--C_+\},\quad
 d_-(.8796-C_+).
 \end{gathered}
}\tag{I45}\] Apply order \(\fitinline{j=2}\) to the first two and \(\fitinline{j=4}\) to the four ordinary gaps. The latter also have the stated increasing width coefficients, so checking their left endpoints proves the full intervals. The smallest raw gap is greater than \(\fitinline{.00033389}\). After division by the correct moment orders every fixed-face \(\fitinline{g/j}\) is greater than \(\fitinline{.00008347}\). Near the long taper only the native unmarked/amplified inverse is needed; the improved profile is identically untapered on its annulus. Thus the global moment proof of I4, rather than a conditional detector field, applies at every outside length.

The three moving \(\fitinline{U}\)-savings in I8 are bounded below by \[\fitmath{
 \frac{d_-k_-\lambda_x}{2},\quad
 \frac{d_-k_-\lambda_y}{2},\quad
 \frac{V_-\gamma}{2};
}\] their minimum is \[\fitmath{
 s_{\mathrm{out}}=\frac{156013}{30000000000}.
}\tag{I46}\] Converting to \(\fitinline{Z}\) with \(\fitinline{D_c\ge D_-}\) and subtracting the entire debit \(\fitinline{B}\) of I39 leaves \[\fitmath{
 D_-s_{\mathrm{out}}-B
 =\frac{5635111919539}{30000000000000000000}
 >1.87837\cdot10^{-7}.
}\tag{I47}\] The fixed faces have a much larger margin. This pays the first two cases of the whole convolution at their actual amplitudes.

In the middle case apply I30 and I13. At the crossing I7, compare with the inherited exponent \(\fitinline{E_0=(3R_b+1)/4-qz_0/2}\). At the limiting quota \(\fitinline{z_0=1/2-n_0}\), substitution of I15 gives exactly \[\fitmath{
 E_0-E_{\mathrm{mid}}
 =d\{p\Delta W+k_1(y-m)/2-p(n-n_0)/2\}.
}\tag{I48}\] The genuine strict \(\fitinline{z_0=(1-s-2n_0)/(2+10s)\le1/2-n_0}\) only increases this saving. With \(\fitinline{y-m\ge-\gamma-\lambda_x-o(1)}\), \(\fitinline{n-n_0\le50.1\cdot10^{-6}}\), and \(\fitinline{p/2\le1/4}\), it is at least \[\fitmath{
 \begin{aligned}
 s_{\mathrm{mid}}
 &=d_-\{p_-\Delta W-k_+(\gamma+\lambda_x)/2
                      -(50.1\cdot10^{-6})/4\}\\
 &=\frac{20060849112683463}{2500000000000000000000}.
 \end{aligned}
}\tag{I49}\] Consequently \[\fitmath{
 D_-s_{\mathrm{mid}}-B
 =\frac{6225486788819945214848889}
 {2500000000000000000000000000000}
 >2.49019\cdot10^{-6}.
}\tag{I50}\] There is no \(\fitinline{+3\eta}\) saturation loss in I48--I50. The first two cases were bounded with I8, not removed by an existential witness argument. The surplus required by I49 is less than \(\fitinline{.000085391<.000100}\), within the proved source I15.

Greedy mark division used here is on the same original class: its \textbf{total} original slot width is \(\fitinline{\ell_1/D_c>1/5}\), and its full lower gain is \(\fitinline{q\ell_1/D_c}\), up to the charged bin error. Sort distinct whole slots by lower gain per width. A fractional prefix of quota \(\fitinline{W}\) has at least the average gain \(\fitinline{qW}\); rounding loses one early mesh. If positive slots are exhausted, their unchanged total gain already exceeds \(\fitinline{qW}\). No assertion that their positive width exceeds \(\fitinline{1/5}\) is needed. Dividing the positive I30 by the squared common response leaves the mark-divided comparison \(\fitinline{U^{1-2qW+\epsilon}}\) on that same class. Repetitions used in I19/I21 never enter this lower product.

Finally sum all three cases of I11, not only a chosen triple. Their number is \(\fitinline{O((\log U)^3)}\), their two separation densities have a common weighted integrable envelope, and the fixed positive reserves I47/I50 pay the small support, source and bin errors. Numerator-dilation classes use I9 with their actual packet damping before their absolute coefficient masses are summed. Unit-interval Sobolev, including normalized first log derivatives, pays row-dependent heights on the original enlarged intervals; the numerator frequency is independent and keeps its original rapid weight. Thus this is a bound for the entire reciprocal-containing retained ticket I42.

\hypertarget{appendix-i-7}{}

\hypertarget{i.7.-source-first-physical-return-and-ordinary-characters}{%
\subsection{I.7. Source-first physical return and ordinary characters}\label{i.7.-source-first-physical-return-and-ordinary-characters}}

Choose the exact finite banks, strict decrements and source/physical meshes before the later original finite slot type and arithmetic dictionary. After fixing that type, choose the finite profile, Fourier and Sobolev orders and common polynomial degrees. Only then choose the one positive target height exponent \(\fitinline{\tau_\eta}\), before every positive buffered height/grid bound \(\fitinline{n}\). The high external heights are paid by H24a's polynomial absorption and the independent rapid tails. The early source mesh depends on the fixed positive SN9 gap, curve bank and loss budgets; it does not depend on a growing slot count. All constants may depend on the fixed presented target modulus, not on live heights or the finite grids.

The complete exterior I.4, contour and independently established generic geometry I.5, and entire retained convolution I.6 now give the literal H45 statement at the new geometry: \[\fitmath{
 \begin{gathered}
 \forall\eta\ \exists\tau_\eta>0\ \forall n>0\ \exists C_n>0
 \quad\text{eventually in }Z:\quad \forall\mathrm{idx,grid},\\
 \mathrm{ActualBufferedSelection}\ \Longrightarrow\
 \left|\frac{\mathrm{cube}_{D_{\mathrm{src}},F_{\mathrm{src}},\eta}
 (Z,\tau_\eta,\mathrm{idx,grid})}
 {\mathrm{normalizer}_{D_{\mathrm{src}},F_{\mathrm{src}}}(Z)}\right|
 \le C_nZ^{\beta-k_{\mathrm{phys}}-P_{\mathrm{src}}.\sigma}.
 \end{gathered}
}\tag{I51}\] Choose a common \(\fitinline{P_{\mathrm{src}}.\sigma>0}\) below \(\fitinline{10^{-9}}\) and the original transport reserves. The fixed data in I38 precede \(\fitinline{\eta}\). The predicate is the complete one in H.8: actual source slots and calibration, the same buffered maps, finite ray dictionary, presented nonprincipal rows, source modulus cap, zero band and attaining zero at its original allowed height, with its original upper-family zero guard. None of those maps or guards is replaced by the existence of an unrelated detector witness.

The cube is the literal finite truncation of H43--H44, including both independent retained segments, every main/error coefficient and the actual reciprocal of the physical denominator. The only substitution is the new \(\fitinline{x_1,y_1,\ell_1}\) geometry. The quotient-free Euler numerator, finite prime corrections and double principal residue are unchanged. In particular its nonzero normalizer is exactly \[\fitmath{
 c_{\mathcal S}(-1)^NZ^{-\ell_1/6}
 \prod_i\sum_{\mathfrak p}w_i(N\mathfrak p/P_i)
                         (N\mathfrak p)^{-5/6}.
}\tag{I52}\] It has the same logarithmic size and sign as in H46. No new normalized probe or cancelled principal pole is substituted.

Apply the generic AsymmetricPhysicalReturn.high\_of\_cube common-probe transport under the Geometry.Parameters and BoundsCore just proved, with I51 and I52. The same \(\fitinline{D_{\mathrm{src}},F_{\mathrm{src}},P_{\mathrm{src}}}\), actual cutoff and normalizer are used in the cube, normalized source probe and signal. The same signed Gram computation at the new geometry gives \(\fitinline{J_\eta=O_\eta(Z^{\theta_1-k_{\mathrm{phys}}+\omega})}\) for every \(\fitinline{\omega>0}\), while \(\fitinline{J_\eta-f_\eta=O_\eta(Z^{\beta-k_{\mathrm{phys}}-\sigma})}\). Here \(\fitinline{f_\eta}\) is exactly H47 with \(\fitinline{k_{\mathrm{phys}}}\), the same nonvanishing \(\fitinline{H_\eta(s)}\), and the original reciprocal \(\fitinline{1/L_\eta(s)}\). To see the last contradiction without an attained supremum, suppose \(\fitinline{\mu=\beta-\theta_1>0}\), choose \(\fitinline{\omega<\mu/4}\), and then \[\fitmath{
 0<m<\min\{\mu-\omega,\sigma,\beta-.87\}.
}\] The Mellin transform of \(\fitinline{f_\eta}\), using its arbitrary decay at zero and its \(\fitinline{Z^{\beta-k_{\mathrm{phys}}-m}}\) bound at infinity, is holomorphic for \(\fitinline{\Re s>\beta-m}\). On the original absolute line it equals the fixed Gaussian times \(\fitinline{H_\eta/L_\eta}\). Since \(\fitinline{|H_\eta-1|\le1/2}\) on \(\fitinline{\Re s\ge.87}\), the quotient cannot be holomorphic at a primitive target zero above \(\fitinline{\beta-m/2}\), which exists by the definition of \(\fitinline{\beta}\). This proves \(\fitinline{\beta\le\theta_1}\).

The original norm pullback, fixed Euler restoration and nonvanishing on \(\fitinline{\Re s\ge1}\) therefore give the stated ordinary conclusion for every fixed presented modulus and every finite height, in the strict half-plane \(\fitinline{\Re s>\theta_1}\), with the actual principal pole at one excepted. The same conclusion at \(\fitinline{\theta=.87334570}\), or \(\fitinline{c=.12665430}\), follows by inclusion. There is no growing-modulus uniformity or assertion of nonvanishing on the boundary.

\hypertarget{appendix-i-8}{}

\hypertarget{i.8.-exact-data-and-source-attribution}{%
\subsection{I.8. Exact data and source attribution}\label{i.8.-exact-data-and-source-attribution}}

The complete preceding manuscript is the retained 080 source at SHA256 cafc83fb\allowbreak{}58539ac2\allowbreak{}8a27c4d7\allowbreak{}f4446f20\allowbreak{}0a272afa\allowbreak{}dfddd9a7\allowbreak{}6038a0cf\allowbreak{}022ee4e6. Appendix H is used as a previous proved source, with its full original foundations and notices. The new source-to-consumer arguments are supplied by the whole-convolution comparison of controller081, the curved two-ray source of direct081, the complete exterior of near081 and the adaptive physical comparison of recombination081. Their necessary mathematics is included above.

The exact finite objects are:

\begin{itemize}
\tightlist
\item
  research/\allowbreak{}local\_cycle081/\allowbreak{}direct/\allowbreak{}CHECK.py and CERTIFICATE.json: the rational radical enclosures, 142 length cells, 522 closed affine triangles, central minima and all outer-radius/restoration guards in I.3.
\item
  research/\allowbreak{}local\_cycle081/\allowbreak{}near\_record/\allowbreak{}CHECK.py and EXTERIOR\_FINAL.json: the actual 57 boxes, the additional paired wing 28, original lower-polynomial substitutions, target-specific exterior and physical complements, and source-domain guards in I.4--I.5.
\item
  research/\allowbreak{}local\_cycle081/\allowbreak{}recombination/\allowbreak{}CHECK.py and FINITE.json, case stretch: all three moving and six fixed weighted faces, the actual source rectangle, cutoff separation and exact margins I47/I50.
\item
  research/\allowbreak{}local\_cycle069/\allowbreak{}continuum/\allowbreak{}RETURN\_CERTIFICATE.json; research/\allowbreak{}local\_cycle077/\allowbreak{}intake/\allowbreak{}EXACT\_CONTACT\_RECEIPT.json; research/\allowbreak{}local\_cycle079/\allowbreak{}near\_record/\allowbreak{}RETURN\_FINAL.json; and research/\allowbreak{}cycle066/\allowbreak{}return\_core.py: the already retained count boxes, exact-contact enclosure, saved polynomial data and rational polynomial operations, with the immutable hashes identified in H.7 and the new exterior receipt.
\item
  The literal generic geometry fields are the \href{https://github.com/GTimG/QuasiRiemannTracker/blob/c74d855c3bdfceeee059501aea76c8b1627b9abd/submissions/2026-10-21-hailey-collet-1252270103/src/Submissions/hailey_collet_1252270103/Record1252270103Return/Geometry.lean\#L10-L19}{base Geometry.Parameters}, with the five coarse fields retained at research/\allowbreak{}local\_cycle069/\allowbreak{}intake/\allowbreak{}pr141/\allowbreak{}submissions/\allowbreak{}2026-11-13-continuum-anchor/\allowbreak{}src/\allowbreak{}Candidate/\allowbreak{}Local/\allowbreak{}Qrh/\allowbreak{}TwoSlot/\allowbreak{}AnchorZeroLoss/\allowbreak{}BoundsCore.lean. The generic physical consumer is research/\allowbreak{}local\_cycle078/\allowbreak{}near\_record/\allowbreak{}primary/\allowbreak{}lean/\allowbreak{}AsymmetricPhysicalReturn.lean, Geometry section; the direct inequalities I38a--I38b, not a previous-target constructor, provide its hypotheses here.
\end{itemize}

These finite programs establish only the stated exact arithmetic implications. The global source contracts, true moment orders, Fourier/taper transfer, divisor assembly and physical return are the analytic arguments given here and in the explicitly retained A--H proofs.

Credits and original terms remain with OpenAI and the inherited authors, including Sil's paired canonical source, Joseph Bebel's exact contact and count work, Zhongpai Gao, Liu, Akash Levy's numerator and Hailey Collet's reflection/probe/return framework. H.10 preserves the pinned SovereignSteak draft source head c838a750\allowbreak{}079f7f86\allowbreak{}a573df9c\allowbreak{}9f729b8b\allowbreak{}c92d125e, its Rohan Arun/contributor and historical notices, and its Codex assistance disclosure; none of its proposed endpoint assertions is an extra premise. The Sepia/CWR mathematical-writing source attribution and terms remain as in H.10. Catalogue and Lean source contracts were used for discovery and exact scope; no new Lean theorem or end-to-end formal verification is claimed.

\hypertarget{appendix-j}{}

\hypertarget{appendix-j.-a-cellwise-physical-return-with-simultaneous-moments}{%
\section{Appendix J. A cellwise physical return with simultaneous moments}\label{appendix-j.-a-cellwise-physical-return-with-simultaneous-moments}}

The argument below retains the complete proof through Appendix I at SHA-256 \texttt{c\allowbreak{}0\allowbreak{}0\allowbreak{}2\allowbreak{}5\allowbreak{}3\allowbreak{}c\allowbreak{}4\allowbreak{}f\allowbreak{}7\allowbreak{}7\allowbreak{}1\allowbreak{}3\allowbreak{}e\allowbreak{}d\allowbreak{}c\allowbreak{}1\allowbreak{}9\allowbreak{}7\allowbreak{}b\allowbreak{}3\allowbreak{}1\allowbreak{}b\allowbreak{}4\allowbreak{}b\allowbreak{}f\allowbreak{}2\allowbreak{}8\allowbreak{}e\allowbreak{}4\allowbreak{}a\allowbreak{}a\allowbreak{}1\allowbreak{}d\allowbreak{}0\allowbreak{}8\allowbreak{}d\allowbreak{}4\allowbreak{}d\allowbreak{}5\allowbreak{}a\allowbreak{}7\allowbreak{}5\allowbreak{}5\allowbreak{}f\allowbreak{}8\allowbreak{}6\allowbreak{}b\allowbreak{}d\allowbreak{}b\allowbreak{}8\allowbreak{}d\allowbreak{}d\allowbreak{}a\allowbreak{}0\allowbreak{}a\allowbreak{}0\allowbreak{}7\allowbreak{}8\allowbreak{}a\allowbreak{}9\allowbreak{}9\allowbreak{}1\allowbreak{}8}, including its stated A--H and SN9 foundations. Its complete ticket is \[\fitmath{
 c_J=\frac{126657}{1000000}=.126657,\qquad
 t_J=1-c_J=.873343,\qquad h=h_*.                              
}\tag{J1}\] Here \(\fitinline{h_*}\) is the exact contact conductor of E and G, not a newly optimized scalar. On the stated foundations the proof gives the ordinary fixed-modulus, all-height conclusion strictly to the right of \(\fitinline{t_J}\), with an actual principal pole excepted. No new Lean verification, growing-modulus uniformity or boundary nonvanishing is asserted.

\hypertarget{j.1.-the-physical-comparison-as-a-parametric-lemma}{%
\subsection{J.1. The physical comparison as a parametric lemma}\label{j.1.-the-physical-comparison-as-a-parametric-lemma}}

Use the original coefficients, main/error masks, ideal columns, six units, nonunit zero extension, fixed Euler deletions, whole canonical prime pools and independent ordinary heights of H--I. Fix a target \(\fitinline{t}\), an admissible conductor \(\fitinline{h}\), and \[\fitmath{
 \ell=11/3-4t,\qquad y_t=h-2\ell,\qquad x_t=1+\ell-h.
}\tag{J2}\] The finite cell lemma below applies to any such geometry once the literal geometric and physical guards are established. It is not an extrapolation of a saturated-row count.

\textbf{Cell lemma.} Let \(\fitinline{\mathcal B=[a,b]\times[f,e]}\) be a closed density/amplitude cell, with \(\fitinline{0<a\le b<5/6}\), \(\fitinline{0<f\le e\le1/2}\). Suppose the paired count E27--E29 is admitted on this cell with its full ordinary and inverse continuations. For its actual parameters \(\fitinline{d,p}\), put \[\fitmath{
 \begin{aligned}
 q&=dp,& k&=1-p(1+2g_1),& V&=5/6-d,\\
 J&=3V+2dk,& U_d&=1-p(1+6g_0-4g_1),\\
 t_0&=1+dU_d/J,& m&=t_0/3,& b_0&=2m,\\
 R&=1-d[p(1+2g_0)+kb_0]=1/6+Vt_0.
 \end{aligned}                                                   
}\tag{J3}\] The exponent \(\fitinline{R}\) is a proved upper count. It need not equal the minimum of the available counts. Suppose the old upper physical price on this cell, including its previously unused arithmetic reserve, has excess at most \(\fitinline{B}\) over the new target. Negative \(\fitinline{B}\) is allowed. Choose positive conductor, dilation and inverse-length cutoffs \(\fitinline{\omega,v_*,v_-,v_+}\), and convolution cutoffs \(\fitinline{\gamma,\lambda_x,\lambda_y}\). Set \(\fitinline{D_-=h_- -\omega}\), \(\fitinline{k_-=1-e(1+2g_{1,+})}\), \(\fitinline{k_+=1-f(1+2g_{1,-})}\), and \(\fitinline{V_-=5/6-b}\). Assume the following strict prices exceed \(\fitinline{B}\), with a common positive reserve: \[\fitmath{
 \begin{gathered}
 a_D\omega,\quad D_-av_*/4,\quad .10D_-v_-,\quad .27D_-v_+,\\
 D_-ak_-\lambda_x/2,\quad D_-ak_-\lambda_y/2,\quad
 D_-V_-\gamma/2,\qquad
 a_D=3(1-b)/4-9/100-3b/160.                                   
 \end{gathered}
}\tag{J4}\] The six fixed continuation prices are also required at their genuine change points, divided by moment order two for an inverse and four for an ordinary factor. Require the full tangent \(\fitinline{x}\)-interval \(\fitinline{[b_0-\gamma-\lambda_y,b_0+\lambda_x]}\), with its small support collars, to lie inside the admitted inverse anchor; require the long reciprocal annulus to lie above \(\fitinline{t_0-\gamma}\).

For each retained coefficient and scale, put \(\fitinline{n_0=1-y_t/D}\) and \(\fitinline{n_{\rm ref}=\max(n,n_0)}\). Suppose the actual two-ordinary source I30 is proved, with componentwise shortening, both independent heights, common profiles and normalized first log tests, at every required pair \(\fitinline{(n_{\rm ref},y)}\), and with width \[\fitmath{
 W=\tfrac12[1+g_0-n_{\rm ref}-(1+2g_1)y]+\Delta.
}\tag{J5}\] It suffices that, throughout the cell, the further price \[\fitmath{
                  D_-a(f\Delta-k_+\gamma/2)>B                 
}\tag{J6}\] has a common positive reserve, after the directed support and annular costs. All source admissions, contour gates, original total mark supply and fixed-pad charges must be paid as well. Under these hypotheses the entire original reciprocal-containing physical contribution on this cell has a strict saving at the new target.

\textbf{Proof.} First move contours on the original pre-amplitude class. The conductor loss follows from the paired upper price: its two native-capacity slopes are at least \(\fitinline{a_D}\), since this particular \(\fitinline{R\ge1-d}\). The true coefficient dilation \(\fitinline{v=\log_U\max(1,L_iK)}\) has the original large-dilation saving \(\fitinline{dv/4}\). The full \(\fitinline{w}\)-inversion and the quotient-free expansion give inverse-length weights \(\fitinline{U^{\Delta_s(1/D-r)}}\) on the two legal shifts \(\fitinline{\Delta_s=-.10,.27}\). These give the first four prices J4. The discarded regions and the retained region are disjoint.

On the remaining long inverse use the exact identity I10, with cutoff \(\fitinline{A=U^{t_0-\gamma}}\), and retain every separated triple: \[\fitmath{
 x\le t_0-\gamma+o(1),\qquad x+y\ge t_0-\gamma-o(1),\qquad
 x+y+z=r+o(1).                                                
}\tag{J7}\] If \(\fitinline{x\ge b_0+\lambda_x}\), retain the actual inverse factor and use its second moment. On the remaining triples with \(\fitinline{y\ge m+\lambda_y}\), retain the actual ordinary factor and use its fourth moment. I8--I9 give gain \(\fitinline{g/j}\) for a normalized factor whose moment exponent is \(\fitinline{R-g}\), including the original coefficient packet at its actual dilation. Thus the two tangent gains are \(\fitinline{dk\lambda_x/2}\), \(\fitinline{dk\lambda_y/2}\). The amplified inverse gain is \(\fitinline{V\gamma/2}\); between the anchor and amplified branch use the six actual continuations. This proves the remaining exclusions in J4. No inverse-amplitude threshold is imposed.

In the middle case normalize \(\fitinline{f_x=U^{-dx/2}|M_x|}\). Its independently marked second moment and the original pointwise bound plus count give \[\fitmath{
 \sum_{\mathcal C} f_x^2\ll
 U^{R-dk(x-b_0)_++\epsilon}\mathcal H.
}\] Bound only the \(\fitinline{z}\)-inverse pointwise. The original exterior weight and the relation in J7 leave \(\fitinline{U^{-dy/2}f_x|P(n)S(y)|}\). Divide the genuine lower marks in each of the two positive source inequalities separately, and then apply Cauchy on the same original rows. An overlap of their chosen prefixes introduces no new physical mark. The resulting exponent is \[\fitmath{
 E_{\rm mid}=\tfrac12\{1-dy-2qW+R-dk(x-b_0)_+\}.             
}\tag{J8}\] The original physical price is \(\fitinline{(3R+1)/4-qz_0/2}\), with \(\fitinline{z_0=(1-s-2n_0)/(2+10s)\le1/2-n_0}\). Substitution of J3 and J5 gives saving at least \[\fitmath{
 d\{p\Delta+\tfrac k2[y-m+(x-b_0)_+]
                       -\tfrac p2(n_{\rm ref}-n_0)\}.         
}\tag{J9}\] In J7, \(\fitinline{y-m+(x-b_0)_+\ge-\gamma-o(1)}\). This removes the former \(\fitinline{\lambda_x}\) debit while retaining the actual inverse amplitude.

The last term of J9 is paid by the coefficient that enlarged the numerator. At the central contour the fully summed original physical label mass satisfies, separately in every row, \[\fitmath{
 \sum_e |A_i(e)|N_e^{-1/2}\max(1,L_iN_e)^{d/2}
       \le C U^\epsilon\prod_jY_{p,j}^{-4/25}.                 
}\tag{J9a}\] Here \(\fitinline{e}\) includes the actual coefficient and error-tuple indices, their exterior smooth Gauss weights and multiplicities; \(\fitinline{L_i}\) contains the true conductor and repeated-prime deficit. This is B.7(221)--(222), with the explicit rowwise contraction in \texttt{P\allowbreak{}o\allowbreak{}s\allowbreak{}i\allowbreak{}t\allowbreak{}i\allowbreak{}v\allowbreak{}e\allowbreak{}L\allowbreak{}a\allowbreak{}b\allowbreak{}e\allowbreak{}l\allowbreak{}P\allowbreak{}a\allowbreak{}c\allowbreak{}k\allowbreak{}e\allowbreak{}t\allowbreak{}.\allowbreak{}l\allowbreak{}e\allowbreak{}a\allowbreak{}n} 89--96. The preceding bare \texttt{e\allowbreak{}r\allowbreak{}r\allowbreak{}o\allowbreak{}r\allowbreak{}\_\allowbreak{}p\allowbreak{}a\allowbreak{}c\allowbreak{}k\allowbreak{}e\allowbreak{}t\allowbreak{}\_\allowbreak{}d\allowbreak{}a\allowbreak{}m\allowbreak{}p\allowbreak{}e\allowbreak{}d\allowbreak{}\_\allowbreak{}m\allowbreak{}a\allowbreak{}s\allowbreak{}s} still has an \texttt{e\allowbreak{}r\allowbreak{}r\allowbreak{}o\allowbreak{}r\allowbreak{}P\allowbreak{}r\allowbreak{}o\allowbreak{}d\allowbreak{}u\allowbreak{}c\allowbreak{}t\allowbreak{}S\allowbreak{}i\allowbreak{}z\allowbreak{}e} on its right side. Only its \(\fitinline{579^{\omega}}\) prefactor and fixed constant are uniformly subpower; that fact alone does not prove J9a. The full physical error-label sum proves J9a while retaining its displayed slot factor. In one \(\fitinline{v}\)-bin the unweighted mass is therefore bounded by the right side of J9a times \(\fitinline{U^{-dv/2}}\). The actual reflected support implies \[\fitmath{
       n\le n_0+v+\epsilon_{\rm ann}+\epsilon_{\rm bin}.        
}\tag{J10}\] The two dilation terms in J9 and the mass bound leave at least \(\fitinline{(d-q)v/2}\), apart from those directed costs. This is nonnegative. The damping used on these retained bins is not added a second time to the discarded large-dilation bins. For \(\fitinline{n<n_0}\), use the proved downward mixed source anchored at \(\fitinline{n_0}\) on the actual shorter factor; its null/unit and signed equal-product reference cases are included in I30.

Here the packet argument is rowwise before the coefficient sum. Fixed-ratio \(\fitinline{9/8}\) reanchoring and H24a log-Fourier inversion separate the resulting scale profiles with a common weighted integrable envelope and both independent frequencies. Alternatively sum each row's positive absolute coefficient mass before its tested common-profile supremum. The scale/dilation supremum is controlled by this common compact log-Fourier envelope, equivalently by the actual scale derivative \(\fitinline{xW'(x)}\) and its fixed finite seminorms. The separate height supremum is controlled by unit-interval Sobolev and the normalized first norm-log test. Neither step puts a row-dependent arithmetic coefficient into a fixed-ray operator theorem. Equations J9--J10, with the fully summed J9a, prove J6 and the lemma. \(\fitinline{\square}\)

\hypertarget{j.2.-paying-the-old-boxes-at-the-changed-geometry}{%
\subsection{J.2. Paying the old boxes at the changed geometry}\label{j.2.-paying-the-old-boxes-at-the-changed-geometry}}

Retain the original four conductor edges and all 57 closed density/amplitude rectangles. Write \(\fitinline{\alpha}\) for the exact contact and \(\fitinline{\xi=\alpha_+-t}\). The actual source data still satisfy \(\fitinline{D_{\rm src}.\mathrm{momentKappa}=2\beta-1}\). In the contrary argument \(\fitinline{t<\beta\le t_I<\alpha}\), FB2--FB6 permit conservative natural-count and capacity pricing at \(\fitinline{2\alpha-1}\); they do not change that equality.

For an admitted paired box \(\fitinline{[a,b]\times[f,e]}\), E49 gives \[\fitmath{
 M=K(d-\delta_*)^2+(1/2-p)dQ,\qquad
 Q=q_{10}+d(q_{20}+q_{21}p).
}\] The retained exact-contact isolation yields \[\fitmath{
 \begin{gathered}
 .425830047444893\le\delta_*\le.425830047444894,\quad
 K\ge.205830466908273,\\
 (q_{10},q_{20},q_{21})\ge
 (-.140952330598569,.588338843157365,-.328860096935511)
 \quad\hbox{componentwise}.
 \end{gathered}                                                
}\tag{J11}\] On every actual use, \(\fitinline{Q_-=q_{10,-}+a(q_{20,-}+q_{21,-}e)\ge0}\) and \(\fitinline{q_{20,-}+q_{21,-}/2>0}\). Consequently the old unconsumed margin is at least \[\fitmath{
 S_{\mathcal B}=
 \frac{K_-\operatorname{dist}([a,b],[\delta_{*,-},\delta_{*,+}])^2
             +(1/2-e)aQ_-}
      {5/2-a[1+2f(1+2g_{1,-})]} .                             
}\tag{J12}\] This positive margin is retained in central cells as well as outside them. The new target costs at most \(\fitinline{(3+8be)\xi}\) on the affine top capacity. Increasing the conductor by \(\fitinline{\eta_h=h-h_*}\) adds at most \(\fitinline{\eta_h A_{\mathcal B}}\), where \[\fitmath{
 A_{\mathcal B}=-5/12-af/4+3R_+/4.                            
}\tag{J13}\] Here an upper bound \(\fitinline{R_+}\) for the particular paired count follows directly from E28, using the monotonic corner bounds for \(\fitinline{t_0}\) in J3. Thus \(\fitinline{B=(3+8be)\xi+\eta_h A_{\mathcal B}-S_{\mathcal B}}\) is a valid cellwise upper excess. No saving from changing \(\fitinline{h}\) at the contact is presumed: at \(\fitinline{d=\delta_*,p=1/2}\), \(\fitinline{R_*=5/9+\delta_*/6}\) makes J13 vanish exactly.

For a saved old/native lower polynomial \(\fitinline{L}\), retain \(\fitinline{2L/5}\) and the I31 perturbation coefficient: \(\fitinline{3+8be}\) on the two upper edges and \(\fitinline{3+(4+4/(3\cdot.746))be}\) on the lower edges. Its top-conductor slope can use the globally available old E48 count, even if the saved margin used a native count, because the full conductor-independent pointwise minimum is used throughout the conductor interval. An explicit upper bound is \[\fitmath{
 R_{\rm old,+}=\min\!\left(1,1-a+
 \frac{b}{2[D_0/C_0+a/(5/6-a)]}\right),\quad
 D_0=3-e-\frac{2e}{3\kappa_-},\quad
 C_0=2(1-f)\left(1-\frac f{3\kappa_+}\right),                 
}\tag{J14}\] where \(\fitinline{\kappa_\pm=2\alpha_\pm-1}\). Positive denominators follow from the retained isolation. Substitute J14 in J13. The three lower conductor edges improve under \(\fitinline{\eta_h\ge0}\): their base slope is \(\fitinline{-49/150}\), their capacities increase or remain constant, and their count is independent of \(\fitinline{h}\). At fixed \(\fitinline{d,p}\) the full price is convex in \(\fitinline{D}\); using the same admitted envelope at both ends fills each original conductor interval.

Use the paired alternative only on a proved full upward admission window. At J1 these are precisely E27 and the additional I32 wing; no new window is needed. I33 checks all ten literal E29 guards at each of the four \(\fitinline{(g_0,g_1)}\) corners, including the same-zero crossing, monotonicity and all native/affine/capped/unmarked continuations. Cell37, although originally saved with a native count, lies wholly in E27 and can use J12. Every other saved box retains its actual lower polynomial until a genuinely admitted paired alternative improves it.

The rational construction splits only unresolved upper-edge boxes, with density step at most \(\fitinline{1/2048}\) and amplitude step at most \(\fitinline{1/4096}\). It either proves a strict exterior margin or invokes J.1 on that whole closed subbox. For a central subbox, let \(\fitinline{B}\) be rounded outward and set \(\fitinline{\mathcal B_0=\max(0,B)+10^{-8}}\). Choose each cutoff by rounding upward to multiples of \(\fitinline{10^{-9}}\) the equality case of J4 with right side \(\fitinline{\mathcal B_0}\). Set \[\fitmath{
 \Delta=\left\lceil
 \frac{\mathcal B_0/(D_-a)+k_+(\gamma+3\epsilon_{\rm ann})/2
                         +e\epsilon_{\rm ann}/2}{f}
 \right\rceil_{10^{-9}},\qquad \epsilon_{\rm ann}=10^{-8}.     
}\tag{J15}\] This pays the directed collars in J7--J10 and all J4 exclusions. The six fixed continuations are checked at their genuine change points and orders. This rule reconstructs the finite construction; it is not a new analytic hypothesis.

At J1 there are 178 central subboxes. The 57 original boxes produce 35 saved-polynomial leaves, 14189 E49 exterior leaves and these 178 full-convolution leaves. Their respective least bare reserves are \[\fitmath{
 2.815889856\cdot10^{-6},\qquad
 6.5272227\cdot10^{-8},\qquad
 1.0000013\cdot10^{-8}.                                      
}\tag{J16}\] The complete central demand is contained in \[\fitmath{
 \begin{gathered}
 D\ge.815284035,\quad .41552734375\le d\le.43603515625,\quad
 p\ge.49853515625,\\
 .425065190112956\le n_{\rm ref}\le.425248701009802,\\
 .374528263703056\le y\le.377835826988264,\qquad
 \Delta\le.000111913.                                       
 \end{gathered}
}\tag{J17}\] These are outward rational enclosures, not one artificial rectangular cell used to reprice all subboxes. Split the actual length coordinate at \(\fitinline{y=151/400=.3775}\) before assigning widths or choosing fixed-ratio Fourier anchors:

\begin{longtable}[]{@{}
  >{\raggedright\arraybackslash}p{(\columnwidth - 6\tabcolsep) * \real{0.2308}}
  >{\raggedright\arraybackslash}p{(\columnwidth - 6\tabcolsep) * \real{0.2308}}
  >{\raggedleft\arraybackslash}p{(\columnwidth - 6\tabcolsep) * \real{0.3077}}
  >{\raggedright\arraybackslash}p{(\columnwidth - 6\tabcolsep) * \real{0.2308}}@{}}
\toprule\noalign{}
\begin{minipage}[b]{\linewidth}\raggedright
Closed part
\end{minipage} & \begin{minipage}[b]{\linewidth}\raggedright
Required \(\fitinline{n_{\rm ref}}\)
\end{minipage} & \begin{minipage}[b]{\linewidth}\raggedleft
Required surplus
\end{minipage} & \begin{minipage}[b]{\linewidth}\raggedright
Proved supplier
\end{minipage} \\
\midrule\noalign{}
\endhead
\bottomrule\noalign{}
\endlastfoot
\(\fitinline{.37452\allowbreak{}82637\allowbreak{}03056\le y\le.3775}\) & \(\fitinline{[.42506\allowbreak{}51901\allowbreak{}12956,.42524\allowbreak{}87010\allowbreak{}09802]}\) & \(\fitinline{\le111.913\cdot10^{-6}}\) & The lower source proved below, with \(\fitinline{112\cdot10^{-6}}\) on \(\fitinline{[.42506,.42525]\times[.37452,.3775]}\). \\
\(\fitinline{.3775\le y\le.37783\allowbreak{}58269\allowbreak{}88264}\) & \(\fitinline{[.42507\allowbreak{}95129\allowbreak{}86348,.42516\allowbreak{}83270\allowbreak{}09802]}\) & \(\fitinline{\le54.483\cdot10^{-6}}\) & I30, with \(\fitinline{100\cdot10^{-6}}\) on I14. \\
\end{longtable}

There are 175 lower intersections and nine upper intersections, including all six straddlers. Anchor a seam scale at the seam itself and choose its profile envelope within the assigned closed piece; a fixed open smooth-support gap is unnecessary. One can equivalently use the weaker I30 width on each entire straddling physical subbox: its numerator interval is within I14 and its required surplus is at most \(\fitinline{54.483\cdot10^{-6}}\). Thus no term is lost at the seam. The freely small width decrement for the lower source is chosen below the physical spare \(\fitinline{87\cdot10^{-9}}\), and all resulting exponent costs below the common reserve in J16. Its terminal energy reserve \(\fitinline{10^{-6}}\) is a different quantity and is not treated as an allowable width decrement of that size.

The largest required total width is below \(\fitinline{.100766}\), and its required root volume below \(\fitinline{.901796}\). The actual suppliers below and in I.3 have width below \(\fitinline{.101}\) and volume below \(\fitinline{.903}\). The whole original slot width remains greater than \(\fitinline{1/5}\). The common lower response and the same whole-slot greedy argument in I.6 therefore give exactly the needed squared gain \(\fitinline{2qW}\), including the case where positive slots are exhausted. No radical filter or inverse saturation restriction remains.

The central extrema also give \(\fitinline{\gamma\le.000081041}\), \(\fitinline{\lambda_x=\lambda_y\le.000182737}\), \(\fitinline{v_*\le.000154962}\), and \(\fitinline{1.226340569673141\le r\le1.226627485354473}\). These bounds come after the individual cell prices. Every tangent cutoff is interior to \(\fitinline{[.7335,.7775]}\), every retained long annulus is above its cutoff, and every fixed continuation has physical reserve greater than \(\fitinline{.0000817464}\).

The new lower supplier is a retuning of I.3 using only the already proved curved anchor E(19). Set \(\fitinline{u=29/20000=.00145}\), \(\fitinline{H=1-6u=.9913}\), \(\fitinline{\rho=10^{-6}}\), and retain the actual mean in J5 with \(\fitinline{\Delta=.000112}\). For a regular core radius exponent \(\fitinline{V}\), use \[\fitmath{
 P(V)=V/2-n,\qquad B(V)=V s(\max(3/4,2y/V)),\quad
 s(r)=(1-r)/2+I(r),                                         
}\tag{J17a}\] where \(\fitinline{I}\) is E(19), explicitly given in I16. The low and two same-zero middle energy excesses are the unchanged I19/I21 expressions \[\fitmath{
 \begin{aligned}
 C_0(V)&=e\{W-u-(P(V)+B(V))/2\},\\
 C_i(V)&=\nu y+e[W-u-P(V)/2]-VD_i/2,\\
 D_s&=\nu w+e s(w),\qquad
 D_l=5/6-(5/6-\nu)3w/2.
 \end{aligned}                                                
}\tag{J17b}\] At \(\fitinline{V=H}\), \(\fitinline{n}\) cancels. Partition \(\fitinline{[.37452,.3775]}\) into 120 intervals of width at most \(\fitinline{1/40000}\). On each, enclose the positive fifth root in E(19) between adjacent integer fifth powers at denominator \(\fitinline{10^{13}}\); the radical's supporting tangent minorizes \(\fitinline{I}\) throughout that interval. Choose the low threshold rounded down at \(\fitinline{10^{-6}}\), and cover \(\fitinline{k\le e\le.5002,\ e\le\nu\le3/4}\) by longest-edge bisection into 584 closed triangles. At each vertex test both affine deficits with one rational \(\fitinline{w}\), rounded down at \(\fitinline{10^{-7}}\), on that triangle. These substitutions imply \[\fitmath{
 \min(u-C_0(H),u-C_s(H),u-C_l(H))>10^{-6};                    
}\tag{J17c}\] the exact three minima are, in order, \[\fitmath{
 \frac{125389711143134171766499889}{124140142005000000000000000000000},
 \quad\frac{131207512060881139}{120000000000000000000000},
 \quad\frac{15000043262857}{15000000000000000000}.
}\] The finite proof is affine on each closed triangle, not a sample at detector points. The anchor ratios stay below \(\fitinline{7/9}\); the short witness has \(\fitinline{.75\le w\le.775}\), strictly below .7775 and the actual detector cutoff. For \(\fitinline{e\ge.5002}\) the original SN9 norm is used at slope .5001 and fixed positive gap .0001, on the actual merged vector with its shared-prime terms. All decrements and this gap precede the early mesh.

The outer-radius and restoration arguments of I.3 apply with these verified constants. On \(\fitinline{[.98,1]}\) the energy excesses decrease with \(\fitinline{V}\), with absolute slope below \(\fitinline{5/6}\); the native ranges \(\fitinline{[.875,.98]}\) and \(\fitinline{[1,\infty)}\), the unmarked fourths below .875 and small cores below .08 have strictly positive exact guards in the certificate. The actual maxima are \[\fitmath{
 W\le\frac{126111844583}{1250000000000}<.101,\qquad
 A_0\le\frac{18036560141}{20000000000}<.903.                  
}\tag{J17d}\] Thus the regular energy is at most \(\fitinline{\mathcal H U^{M+u-\rho+(5/6)(H-M)_++\epsilon}}\). Restore every sixth-power factor, including common primes, with this same \(\fitinline{5/6}\) slope. Both raw sides use the literal selected radius \(\fitinline{16K_{\rm first}(\prod2P_j)^4}\), giving \(\fitinline{M\ge1-6u'-o(1)}\), and the unweighted first artificial divisor. The second reciprocal weight cancels its matching H26 divisor exactly; each orientation retains its original reflected ray and the common denominator \(\fitinline{NC\,ND_1}\). The signed equal-product fixed-core treatment retains both principal arms, independent profiles and the \(\fitinline{3/5}\) cancellation, with \(\fitinline{2(3/5)\ge1}\) paying full extraction. Its fixed-row guard exceeds \(\fitinline{.06615}\). The actual reference has short coordinate .139, long coordinate at most .66375, and its short fourth-support guard is positive. All zero and tail gates \(\fitinline{5u'<1}\), \(\fitinline{A_0+4u'<1}\), and the first zero/tail cap are checked. I.3's inverse cube-free restoration and H.5(A)--(C) therefore give the original uncentered kept root, including componentwise shortenings and both first-log tests. This is a new nonrecursive source, with no use of a later paired-source idea.

\hypertarget{j.3.-generic-geometry-tails-and-the-full-return}{%
\subsection{J.3. Generic geometry, tails and the full return}\label{j.3.-generic-geometry-tails-and-the-full-return}}

The original narrow numerical interval for \(\fitinline{\ell}\) in I.5 is not reused. A direct field proof works uniformly for \[\fitmath{
 .87334\le t\le.873345,\qquad
 .815324503\le h\le.815994505.                               
}\tag{J18}\] Indeed J2 gives \(\fitinline{25993/150000\le\ell\le6499/37500}\), and hence \[\fitmath{
 \begin{gathered}
 1/6<\ell<1/5,\quad x_t-\ell=1-h>.18400549,\\
 y_t-x_t=2h-1-3\ell>.11072900,\quad
 1/2-y_t>.03057882,\\
 2x_t-y_t-\ell=2+3\ell-3h>.07187648,\quad
 1-5\ell\ge1001/7500>0,\\
 x_t+y_t=1-\ell,\qquad
 t-(4+y_t-x_t)/6=x_t/2+(y_t-x_t)/12.                         
 \end{gathered}
}\tag{J19}\] These are every field of the generic \texttt{G\allowbreak{}e\allowbreak{}o\allowbreak{}m\allowbreak{}e\allowbreak{}t\allowbreak{}r\allowbreak{}y\allowbreak{}.\allowbreak{}P\allowbreak{}a\allowbreak{}r\allowbreak{}a\allowbreak{}m\allowbreak{}e\allowbreak{}t\allowbreak{}e\allowbreak{}r\allowbreak{}s}. The five \texttt{B\allowbreak{}o\allowbreak{}u\allowbreak{}n\allowbreak{}d\allowbreak{}s\allowbreak{}C\allowbreak{}o\allowbreak{}r\allowbreak{}e} fields follow from \(\fitinline{.873<t<7/8}\), \(\fitinline{.793<h<.816}\), \(\fitinline{.458<y_t<.473}\), \(\fitinline{.1672<\ell<.175}\), and \[\fitmath{
 y_t/2-h(17/50-1/6)-1/50
       =49h/150-\ell-1/50>.07303267>.07.                     
}\tag{J20}\] Thus no previous-target geometry constructor is used. At J1 the additional physical source guards are directly checked: \(\fitinline{h+F<\min(1,2y_t)}\), \(\fitinline{y_t<.76<h}\), \(\fitinline{\ell/(h+F)>1/5}\), and \(\fitinline{h/2-2\ell<3h/40}\) with outward endpoints. The last inequality justifies the top affine capacity used in J12--J13.

Keep the literal numerical pad \(\fitinline{F=\Delta_{\rm pad}=10^{-258}}\) and source cost \(\fitinline{s=10^{-260}}\). In a contrary argument \(\fitinline{\beta>t}\), choose the actual baseline data with \(\fitinline{0<D_{\rm src}.t\le F/4}\) and losses below \(\fitinline{\min((\beta-t)/10000,10^{-9}/4,F)}\), before the target character and live heights. Then \(\fitinline{F/2+2D_{\rm src}.t\le F}\). Charge \(\fitinline{15s/16+2F}\) separately to every J16 reserve. The same 1-Lipschitz capacities, \(\fitinline{0\le R\le1}\), and absolute conductor slope below \(\fitinline{411/400<2}\) extend every sector to \(\fitinline{D\le h+F}\). These charges leave more than \(\fitinline{3\cdot10^{-9}}\); choose the sum of the remaining finite analytic losses below \(\fitinline{10^{-9}}\).

The contours in J4 remain legal on every central cell. With \(\fitinline{\rho=d/2+\Delta_s}\), the actual two-line mass condition on the lower line is \[\fitmath{
 \rho-q/2=d(1-p)/2-.10\ge159/40960>0.                       
}\tag{J21}\] The same \(\fitinline{d,p}\) occur on both sides. The upper contour has \(\fitinline{(1+d)/2+.27<1}\), the lower has \(\fitinline{(1+d)/2-.10>.51}\), and \(\fitinline{\Re w=1/2,\Re z=17/50}\) remain fixed. All quotient-free tuple corrections, repeated-prime factors and original Euler zeros are retained. These are contour statements before amplitude subdivision. The original coefficient prefactor is subpower uniformly in the row; the chosen tail orders absorb its polynomial seminorm and height costs.

The averaged floor E31--E40 has price \(\fitinline{.51-49h/150+.51\ell+.66D}\). At J1 its margin after subtracting \(\fitinline{10^{-5}+10^{-12}}\) is greater than \(\fitinline{.00317788}\). The small-row margin \(\fitinline{49h/150-5\ell/4+13/600}\) is greater than \(\fitinline{.07138767}\); the two principal margins \(\fitinline{.049h-\ell/10-1/800}\) and \(\fitinline{h/600-1/800}\) exceed \(\fitinline{.02137143}\) and \(\fitinline{.0001088741}\). Thus the averaged floor, fixed-pad large rows, five nonprincipal units, actual principal residues, finite prime corrections and compensators keep their original B/H estimates. No principal pole is charged to a row count. The signed Gram identity (198)--(199), fixed numerator nonvanishing on \(\fitinline{\Re s\ge.87}\), zero frequencies and rapid Fourier tails are unchanged.

For the retained convolution sum every triple, not only a large witness. There are polynomially many logarithmic dyads; the two separation integrals have the original common weighted integrable envelope. First fix the positive source decrements, cutoff and contour margins, then the early mesh and slot geometry, then the actual finite dictionary and seminorm/Sobolev degrees. Only afterward choose one positive height exponent, before every buffered height/grid parameter. Unit-interval Sobolev includes each normalized first log test; ordinary heights remain independent. Truncate each new frequency shift at \(\fitinline{T_1/16}\) inside the original \(\fitinline{T_1/4}\) buffer, and choose the rapid-decay orders after the finite polynomial envelopes. This proves a uniform bound for the actual rowwise packet, profiles, all retained heights and both original physical segments.

Consequently the complete cube bound I51 holds at J2, with some fixed \(\fitinline{\sigma>0}\) below \(\fitinline{10^{-9}}\) and the original transport reserves, in its literal order \(\fitinline{\forall\eta\ \exists\tau_\eta>0\ \forall n>0\ \exists C_n}\), eventually in \(\fitinline{Z}\), for every original \texttt{A\allowbreak{}c\allowbreak{}t\allowbreak{}u\allowbreak{}a\allowbreak{}l\allowbreak{}B\allowbreak{}u\allowbreak{}f\allowbreak{}f\allowbreak{}e\allowbreak{}r\allowbreak{}e\allowbreak{}d\allowbreak{}S\allowbreak{}e\allowbreak{}l\allowbreak{}e\allowbreak{}c\allowbreak{}t\allowbreak{}i\allowbreak{}o\allowbreak{}n}. The predicate retains the actual slots, calibration, finite ray dictionary, modulus cap, zero-height and upper-family guards. Use the same baseline/transport data, cutoff and nonzero normalizer I52, substituting only the new \(\fitinline{\ell,x_t,y_t}\). The generic \texttt{A\allowbreak{}s\allowbreak{}y\allowbreak{}m\allowbreak{}m\allowbreak{}e\allowbreak{}t\allowbreak{}r\allowbreak{}i\allowbreak{}c\allowbreak{}P\allowbreak{}h\allowbreak{}y\allowbreak{}s\allowbreak{}i\allowbreak{}c\allowbreak{}a\allowbreak{}l\allowbreak{}R\allowbreak{}e\allowbreak{}t\allowbreak{}u\allowbreak{}r\allowbreak{}n\allowbreak{}.\allowbreak{}h\allowbreak{}i\allowbreak{}g\allowbreak{}h\allowbreak{}\_\allowbreak{}o\allowbreak{}f\allowbreak{}\_\allowbreak{}c\allowbreak{}u\allowbreak{}b\allowbreak{}e}, under J19--J20, gives the same normalized physical probe and signal. Its Gram price is \(\fitinline{t-k_{\rm phys}}\), where \(\fitinline{k_{\rm phys}=(4+y_t-x_t)/6}\).

The final Mellin contradiction is I.7 with these parameters: for \(\fitinline{\beta>t}\), choose \(\fitinline{\omega<(\beta-t)/4}\) and \(\fitinline{0<m<\min(\beta-t-\omega,\sigma,\beta-.87)}\). The exact signal has Mellin transform equal on its absolute line to the original Gaussian times \(\fitinline{H_\eta/L_\eta}\), and its improved growth makes that transform holomorphic on \(\fitinline{\Re s>\beta-m}\). Since \(\fitinline{|H_\eta-1|\le1/2}\) on \(\fitinline{\Re s\ge.87}\), a primitive target zero above \(\fitinline{\beta-m/2}\), supplied by the definition of the supremum, contradicts holomorphy. An attained supremum is unnecessary. Original norm pullback and restoration of fixed Euler factors give the ordinary conclusion for every fixed presented modulus and every finite height in the strict half-plane \(\fitinline{\Re s>t_J}\), with the actual principal pole excepted.

\hypertarget{j.4.-exact-data-and-provenance}{%
\subsection{J.4. Exact data and provenance}\label{j.4.-exact-data-and-provenance}}

The finite object for J1 is \texttt{r\allowbreak{}e\allowbreak{}s\allowbreak{}e\allowbreak{}a\allowbreak{}r\allowbreak{}c\allowbreak{}h\allowbreak{}/\allowbreak{}l\allowbreak{}o\allowbreak{}c\allowbreak{}a\allowbreak{}l\allowbreak{}\_\allowbreak{}c\allowbreak{}y\allowbreak{}c\allowbreak{}l\allowbreak{}e\allowbreak{}0\allowbreak{}8\allowbreak{}2\allowbreak{}/\allowbreak{}n\allowbreak{}e\allowbreak{}a\allowbreak{}r\allowbreak{}\_\allowbreak{}r\allowbreak{}e\allowbreak{}c\allowbreak{}o\allowbreak{}r\allowbreak{}d\allowbreak{}/\allowbreak{}F\allowbreak{}I\allowbreak{}N\allowbreak{}A\allowbreak{}L\allowbreak{}\_\allowbreak{}1\allowbreak{}2\allowbreak{}6\allowbreak{}6\allowbreak{}5\allowbreak{}7\allowbreak{}.\allowbreak{}j\allowbreak{}s\allowbreak{}o\allowbreak{}n}, SHA-256 \texttt{c\allowbreak{}7\allowbreak{}9\allowbreak{}6\allowbreak{}d\allowbreak{}6\allowbreak{}9\allowbreak{}1\allowbreak{}0\allowbreak{}c\allowbreak{}0\allowbreak{}1\allowbreak{}6\allowbreak{}b\allowbreak{}d\allowbreak{}3\allowbreak{}a\allowbreak{}0\allowbreak{}e\allowbreak{}3\allowbreak{}2\allowbreak{}1\allowbreak{}9\allowbreak{}7\allowbreak{}3\allowbreak{}5\allowbreak{}4\allowbreak{}f\allowbreak{}1\allowbreak{}6\allowbreak{}6\allowbreak{}3\allowbreak{}9\allowbreak{}e\allowbreak{}5\allowbreak{}5\allowbreak{}4\allowbreak{}6\allowbreak{}6\allowbreak{}0\allowbreak{}b\allowbreak{}7\allowbreak{}3\allowbreak{}e\allowbreak{}0\allowbreak{}3\allowbreak{}e\allowbreak{}e\allowbreak{}1\allowbreak{}0\allowbreak{}6\allowbreak{}e\allowbreak{}7\allowbreak{}d\allowbreak{}0\allowbreak{}7\allowbreak{}2\allowbreak{}e\allowbreak{}0\allowbreak{}0\allowbreak{}5\allowbreak{}d\allowbreak{}3\allowbreak{}1} (15133 bytes). Its generators, kept in the same directory, are \texttt{C\allowbreak{}H\allowbreak{}E\allowbreak{}C\allowbreak{}K\allowbreak{}.\allowbreak{}p\allowbreak{}y} (\texttt{5\allowbreak{}3\allowbreak{}4\allowbreak{}2\allowbreak{}a\allowbreak{}5\allowbreak{}6\allowbreak{}b\allowbreak{}0\allowbreak{}f\allowbreak{}3\allowbreak{}1\allowbreak{}2\allowbreak{}2\allowbreak{}a\allowbreak{}7\allowbreak{}3\allowbreak{}b\allowbreak{}c\allowbreak{}f\allowbreak{}a\allowbreak{}5\allowbreak{}6\allowbreak{}d\allowbreak{}2\allowbreak{}a\allowbreak{}4\allowbreak{}e\allowbreak{}3\allowbreak{}2\allowbreak{}1\allowbreak{}e\allowbreak{}a\allowbreak{}1\allowbreak{}8\allowbreak{}0\allowbreak{}9\allowbreak{}f\allowbreak{}9\allowbreak{}b\allowbreak{}f\allowbreak{}1\allowbreak{}5\allowbreak{}9\allowbreak{}7\allowbreak{}9\allowbreak{}a\allowbreak{}5\allowbreak{}0\allowbreak{}6\allowbreak{}d\allowbreak{}f\allowbreak{}5\allowbreak{}8\allowbreak{}6\allowbreak{}1\allowbreak{}f\allowbreak{}e\allowbreak{}0\allowbreak{}6\allowbreak{}6\allowbreak{}3\allowbreak{}f\allowbreak{}c}) and \texttt{C\allowbreak{}H\allowbreak{}E\allowbreak{}C\allowbreak{}K\allowbreak{}\_\allowbreak{}F\allowbreak{}I\allowbreak{}N\allowbreak{}A\allowbreak{}L\allowbreak{}.\allowbreak{}p\allowbreak{}y} (\texttt{9\allowbreak{}0\allowbreak{}9\allowbreak{}1\allowbreak{}f\allowbreak{}4\allowbreak{}9\allowbreak{}b\allowbreak{}6\allowbreak{}b\allowbreak{}9\allowbreak{}a\allowbreak{}c\allowbreak{}e\allowbreak{}d\allowbreak{}8\allowbreak{}1\allowbreak{}6\allowbreak{}0\allowbreak{}4\allowbreak{}6\allowbreak{}0\allowbreak{}3\allowbreak{}d\allowbreak{}4\allowbreak{}3\allowbreak{}3\allowbreak{}3\allowbreak{}5\allowbreak{}6\allowbreak{}8\allowbreak{}4\allowbreak{}9\allowbreak{}5\allowbreak{}d\allowbreak{}f\allowbreak{}5\allowbreak{}c\allowbreak{}1\allowbreak{}a\allowbreak{}2\allowbreak{}d\allowbreak{}c\allowbreak{}c\allowbreak{}e\allowbreak{}8\allowbreak{}4\allowbreak{}1\allowbreak{}8\allowbreak{}f\allowbreak{}3\allowbreak{}1\allowbreak{}6\allowbreak{}e\allowbreak{}a\allowbreak{}6\allowbreak{}a\allowbreak{}2\allowbreak{}a\allowbreak{}6\allowbreak{}8\allowbreak{}e\allowbreak{}7\allowbreak{}f}). Run \texttt{python\ -B\ CHECK\_FINAL.py\ -\/-source-root\ \textless{}unpacked-proof-root\textgreater{}}, where the root contains \texttt{r\allowbreak{}e\allowbreak{}s\allowbreak{}e\allowbreak{}a\allowbreak{}r\allowbreak{}c\allowbreak{}h\allowbreak{}/}; the latter fixes target .126657 and conductor shift zero. It reconstructs the complete 57-box routes and 178 central cells, all cutoffs, six fixed continuations, contour/floor/principal margins and the directed source-coordinate extrema, then intersects each actual length interval with the two closed suppliers. Both source matches are verified before serialization. The canonical rational leaf-price hash is \texttt{a\allowbreak{}1\allowbreak{}9\allowbreak{}8\allowbreak{}a\allowbreak{}8\allowbreak{}3\allowbreak{}8\allowbreak{}b\allowbreak{}6\allowbreak{}7\allowbreak{}b\allowbreak{}6\allowbreak{}5\allowbreak{}0\allowbreak{}b\allowbreak{}d\allowbreak{}c\allowbreak{}7\allowbreak{}8\allowbreak{}4\allowbreak{}2\allowbreak{}4\allowbreak{}c\allowbreak{}c\allowbreak{}f\allowbreak{}c\allowbreak{}0\allowbreak{}c\allowbreak{}9\allowbreak{}2\allowbreak{}1\allowbreak{}e\allowbreak{}e\allowbreak{}2\allowbreak{}6\allowbreak{}9\allowbreak{}0\allowbreak{}2\allowbreak{}8\allowbreak{}4\allowbreak{}9\allowbreak{}b\allowbreak{}b\allowbreak{}6\allowbreak{}8\allowbreak{}a\allowbreak{}5\allowbreak{}4\allowbreak{}5\allowbreak{}6\allowbreak{}e\allowbreak{}9\allowbreak{}6\allowbreak{}9\allowbreak{}3\allowbreak{}4\allowbreak{}0\allowbreak{}e\allowbreak{}b\allowbreak{}f\allowbreak{}a\allowbreak{}b\allowbreak{}a}. The program checks all four original edge tilings and reconstructs the leaves deterministically; the compact object need not repeat their coordinates. The retained local and portable-root replays produce the same object bytes.

Its immutable inputs are the original \texttt{r\allowbreak{}e\allowbreak{}s\allowbreak{}e\allowbreak{}a\allowbreak{}r\allowbreak{}c\allowbreak{}h\allowbreak{}/\allowbreak{}c\allowbreak{}y\allowbreak{}c\allowbreak{}l\allowbreak{}e\allowbreak{}0\allowbreak{}6\allowbreak{}6\allowbreak{}/\allowbreak{}r\allowbreak{}e\allowbreak{}t\allowbreak{}u\allowbreak{}r\allowbreak{}n\allowbreak{}\_\allowbreak{}c\allowbreak{}o\allowbreak{}r\allowbreak{}e\allowbreak{}.\allowbreak{}p\allowbreak{}y} (\texttt{3\allowbreak{}4\allowbreak{}8\allowbreak{}1\allowbreak{}8\allowbreak{}b\allowbreak{}0\allowbreak{}2\allowbreak{}7\allowbreak{}a\allowbreak{}3\allowbreak{}b\allowbreak{}8\allowbreak{}c\allowbreak{}4\allowbreak{}a\allowbreak{}3\allowbreak{}e\allowbreak{}1\allowbreak{}9\allowbreak{}5\allowbreak{}4\allowbreak{}d\allowbreak{}a\allowbreak{}d\allowbreak{}7\allowbreak{}7\allowbreak{}a\allowbreak{}a\allowbreak{}2\allowbreak{}d\allowbreak{}4\allowbreak{}a\allowbreak{}9\allowbreak{}9\allowbreak{}b\allowbreak{}f\allowbreak{}4\allowbreak{}8\allowbreak{}2\allowbreak{}e\allowbreak{}6\allowbreak{}6\allowbreak{}a\allowbreak{}a\allowbreak{}d\allowbreak{}e\allowbreak{}1\allowbreak{}8\allowbreak{}b\allowbreak{}b\allowbreak{}4\allowbreak{}6\allowbreak{}b\allowbreak{}b\allowbreak{}c\allowbreak{}9\allowbreak{}3\allowbreak{}6\allowbreak{}0\allowbreak{}8\allowbreak{}8\allowbreak{}1\allowbreak{}3}), \texttt{r\allowbreak{}e\allowbreak{}s\allowbreak{}e\allowbreak{}a\allowbreak{}r\allowbreak{}c\allowbreak{}h\allowbreak{}/\allowbreak{}l\allowbreak{}o\allowbreak{}c\allowbreak{}a\allowbreak{}l\allowbreak{}\_\allowbreak{}c\allowbreak{}y\allowbreak{}c\allowbreak{}l\allowbreak{}e\allowbreak{}0\allowbreak{}6\allowbreak{}9\allowbreak{}/\allowbreak{}c\allowbreak{}o\allowbreak{}n\allowbreak{}t\allowbreak{}i\allowbreak{}n\allowbreak{}u\allowbreak{}u\allowbreak{}m\allowbreak{}/\allowbreak{}R\allowbreak{}E\allowbreak{}T\allowbreak{}U\allowbreak{}R\allowbreak{}N\allowbreak{}\_\allowbreak{}C\allowbreak{}E\allowbreak{}R\allowbreak{}T\allowbreak{}I\allowbreak{}F\allowbreak{}I\allowbreak{}C\allowbreak{}A\allowbreak{}T\allowbreak{}E\allowbreak{}.\allowbreak{}j\allowbreak{}s\allowbreak{}o\allowbreak{}n} (\texttt{8\allowbreak{}c\allowbreak{}3\allowbreak{}b\allowbreak{}0\allowbreak{}e\allowbreak{}d\allowbreak{}f\allowbreak{}6\allowbreak{}c\allowbreak{}2\allowbreak{}8\allowbreak{}4\allowbreak{}c\allowbreak{}1\allowbreak{}3\allowbreak{}4\allowbreak{}e\allowbreak{}5\allowbreak{}c\allowbreak{}a\allowbreak{}e\allowbreak{}0\allowbreak{}4\allowbreak{}0\allowbreak{}5\allowbreak{}c\allowbreak{}9\allowbreak{}8\allowbreak{}2\allowbreak{}3\allowbreak{}e\allowbreak{}8\allowbreak{}7\allowbreak{}8\allowbreak{}a\allowbreak{}9\allowbreak{}f\allowbreak{}c\allowbreak{}5\allowbreak{}a\allowbreak{}8\allowbreak{}4\allowbreak{}7\allowbreak{}0\allowbreak{}4\allowbreak{}c\allowbreak{}0\allowbreak{}0\allowbreak{}8\allowbreak{}4\allowbreak{}4\allowbreak{}0\allowbreak{}7\allowbreak{}1\allowbreak{}3\allowbreak{}5\allowbreak{}d\allowbreak{}c\allowbreak{}6\allowbreak{}5\allowbreak{}7\allowbreak{}d\allowbreak{}5}) and \texttt{r\allowbreak{}e\allowbreak{}s\allowbreak{}e\allowbreak{}a\allowbreak{}r\allowbreak{}c\allowbreak{}h\allowbreak{}/\allowbreak{}l\allowbreak{}o\allowbreak{}c\allowbreak{}a\allowbreak{}l\allowbreak{}\_\allowbreak{}c\allowbreak{}y\allowbreak{}c\allowbreak{}l\allowbreak{}e\allowbreak{}0\allowbreak{}7\allowbreak{}7\allowbreak{}/\allowbreak{}i\allowbreak{}n\allowbreak{}t\allowbreak{}a\allowbreak{}k\allowbreak{}e\allowbreak{}/\allowbreak{}E\allowbreak{}X\allowbreak{}A\allowbreak{}C\allowbreak{}T\allowbreak{}\_\allowbreak{}C\allowbreak{}O\allowbreak{}N\allowbreak{}T\allowbreak{}A\allowbreak{}C\allowbreak{}T\allowbreak{}\_\allowbreak{}R\allowbreak{}E\allowbreak{}C\allowbreak{}E\allowbreak{}I\allowbreak{}P\allowbreak{}T\allowbreak{}.\allowbreak{}j\allowbreak{}s\allowbreak{}o\allowbreak{}n} (\texttt{b\allowbreak{}d\allowbreak{}b\allowbreak{}2\allowbreak{}b\allowbreak{}a\allowbreak{}b\allowbreak{}4\allowbreak{}b\allowbreak{}4\allowbreak{}a\allowbreak{}3\allowbreak{}9\allowbreak{}7\allowbreak{}c\allowbreak{}b\allowbreak{}a\allowbreak{}f\allowbreak{}8\allowbreak{}0\allowbreak{}9\allowbreak{}b\allowbreak{}e\allowbreak{}8\allowbreak{}a\allowbreak{}0\allowbreak{}e\allowbreak{}9\allowbreak{}5\allowbreak{}2\allowbreak{}0\allowbreak{}0\allowbreak{}7\allowbreak{}7\allowbreak{}d\allowbreak{}5\allowbreak{}4\allowbreak{}d\allowbreak{}2\allowbreak{}9\allowbreak{}e\allowbreak{}b\allowbreak{}6\allowbreak{}2\allowbreak{}e\allowbreak{}1\allowbreak{}5\allowbreak{}f\allowbreak{}e\allowbreak{}1\allowbreak{}1\allowbreak{}b\allowbreak{}2\allowbreak{}4\allowbreak{}c\allowbreak{}4\allowbreak{}4\allowbreak{}0\allowbreak{}c\allowbreak{}7\allowbreak{}c\allowbreak{}e\allowbreak{}c\allowbreak{}0}). The inherited upper source is I.3 and its packed \texttt{r\allowbreak{}e\allowbreak{}s\allowbreak{}e\allowbreak{}a\allowbreak{}r\allowbreak{}c\allowbreak{}h\allowbreak{}/\allowbreak{}l\allowbreak{}o\allowbreak{}c\allowbreak{}a\allowbreak{}l\allowbreak{}\_\allowbreak{}c\allowbreak{}y\allowbreak{}c\allowbreak{}l\allowbreak{}e\allowbreak{}0\allowbreak{}8\allowbreak{}1\allowbreak{}/\allowbreak{}d\allowbreak{}i\allowbreak{}r\allowbreak{}e\allowbreak{}c\allowbreak{}t\allowbreak{}/\allowbreak{}\{\allowbreak{}C\allowbreak{}H\allowbreak{}E\allowbreak{}C\allowbreak{}K\allowbreak{}.\allowbreak{}p\allowbreak{}y\allowbreak{},\allowbreak{}C\allowbreak{}E\allowbreak{}R\allowbreak{}T\allowbreak{}I\allowbreak{}F\allowbreak{}I\allowbreak{}C\allowbreak{}A\allowbreak{}T\allowbreak{}E\allowbreak{}.\allowbreak{}j\allowbreak{}s\allowbreak{}o\allowbreak{}n\allowbreak{}\}}; the certificate SHA-256 is \texttt{4\allowbreak{}4\allowbreak{}d\allowbreak{}d\allowbreak{}9\allowbreak{}2\allowbreak{}4\allowbreak{}9\allowbreak{}c\allowbreak{}b\allowbreak{}4\allowbreak{}2\allowbreak{}0\allowbreak{}a\allowbreak{}6\allowbreak{}6\allowbreak{}9\allowbreak{}8\allowbreak{}5\allowbreak{}e\allowbreak{}e\allowbreak{}4\allowbreak{}9\allowbreak{}4\allowbreak{}2\allowbreak{}5\allowbreak{}a\allowbreak{}4\allowbreak{}c\allowbreak{}0\allowbreak{}b\allowbreak{}f\allowbreak{}9\allowbreak{}b\allowbreak{}f\allowbreak{}e\allowbreak{}6\allowbreak{}3\allowbreak{}6\allowbreak{}8\allowbreak{}e\allowbreak{}6\allowbreak{}a\allowbreak{}9\allowbreak{}0\allowbreak{}3\allowbreak{}d\allowbreak{}a\allowbreak{}2\allowbreak{}b\allowbreak{}b\allowbreak{}2\allowbreak{}2\allowbreak{}6\allowbreak{}4\allowbreak{}f\allowbreak{}2\allowbreak{}0\allowbreak{}0\allowbreak{}2\allowbreak{}c\allowbreak{}6\allowbreak{}b\allowbreak{}7}.

The new lower source is \texttt{r\allowbreak{}e\allowbreak{}s\allowbreak{}e\allowbreak{}a\allowbreak{}r\allowbreak{}c\allowbreak{}h\allowbreak{}/\allowbreak{}l\allowbreak{}o\allowbreak{}c\allowbreak{}a\allowbreak{}l\allowbreak{}\_\allowbreak{}c\allowbreak{}y\allowbreak{}c\allowbreak{}l\allowbreak{}e\allowbreak{}0\allowbreak{}8\allowbreak{}2\allowbreak{}/\allowbreak{}d\allowbreak{}i\allowbreak{}r\allowbreak{}e\allowbreak{}c\allowbreak{}t\allowbreak{}/\allowbreak{}R\allowbreak{}E\allowbreak{}S\allowbreak{}E\allowbreak{}A\allowbreak{}R\allowbreak{}C\allowbreak{}H\allowbreak{}.\allowbreak{}m\allowbreak{}d} (\texttt{2\allowbreak{}6\allowbreak{}5\allowbreak{}1\allowbreak{}4\allowbreak{}b\allowbreak{}c\allowbreak{}2\allowbreak{}5\allowbreak{}d\allowbreak{}f\allowbreak{}9\allowbreak{}b\allowbreak{}9\allowbreak{}e\allowbreak{}2\allowbreak{}2\allowbreak{}9\allowbreak{}6\allowbreak{}9\allowbreak{}f\allowbreak{}2\allowbreak{}b\allowbreak{}9\allowbreak{}1\allowbreak{}d\allowbreak{}8\allowbreak{}b\allowbreak{}b\allowbreak{}b\allowbreak{}7\allowbreak{}5\allowbreak{}e\allowbreak{}3\allowbreak{}c\allowbreak{}c\allowbreak{}d\allowbreak{}a\allowbreak{}d\allowbreak{}f\allowbreak{}e\allowbreak{}3\allowbreak{}0\allowbreak{}7\allowbreak{}5\allowbreak{}d\allowbreak{}d\allowbreak{}5\allowbreak{}b\allowbreak{}a\allowbreak{}b\allowbreak{}5\allowbreak{}9\allowbreak{}b\allowbreak{}8\allowbreak{}6\allowbreak{}b\allowbreak{}c\allowbreak{}7\allowbreak{}0\allowbreak{}9\allowbreak{}4\allowbreak{}2\allowbreak{}f}), its \texttt{C\allowbreak{}H\allowbreak{}E\allowbreak{}C\allowbreak{}K\allowbreak{}.\allowbreak{}p\allowbreak{}y} (\texttt{3\allowbreak{}e\allowbreak{}0\allowbreak{}7\allowbreak{}b\allowbreak{}2\allowbreak{}b\allowbreak{}6\allowbreak{}3\allowbreak{}9\allowbreak{}7\allowbreak{}9\allowbreak{}a\allowbreak{}c\allowbreak{}3\allowbreak{}2\allowbreak{}f\allowbreak{}0\allowbreak{}4\allowbreak{}2\allowbreak{}6\allowbreak{}6\allowbreak{}d\allowbreak{}0\allowbreak{}1\allowbreak{}d\allowbreak{}2\allowbreak{}2\allowbreak{}5\allowbreak{}8\allowbreak{}0\allowbreak{}d\allowbreak{}c\allowbreak{}e\allowbreak{}5\allowbreak{}d\allowbreak{}c\allowbreak{}e\allowbreak{}9\allowbreak{}b\allowbreak{}b\allowbreak{}4\allowbreak{}7\allowbreak{}9\allowbreak{}3\allowbreak{}f\allowbreak{}3\allowbreak{}c\allowbreak{}5\allowbreak{}5\allowbreak{}e\allowbreak{}4\allowbreak{}9\allowbreak{}4\allowbreak{}1\allowbreak{}c\allowbreak{}b\allowbreak{}5\allowbreak{}5\allowbreak{}d\allowbreak{}b\allowbreak{}7\allowbreak{}0\allowbreak{}4}) and \texttt{C\allowbreak{}E\allowbreak{}R\allowbreak{}T\allowbreak{}I\allowbreak{}F\allowbreak{}I\allowbreak{}C\allowbreak{}A\allowbreak{}T\allowbreak{}E\allowbreak{}.\allowbreak{}j\allowbreak{}s\allowbreak{}o\allowbreak{}n} (\texttt{8\allowbreak{}3\allowbreak{}d\allowbreak{}b\allowbreak{}6\allowbreak{}b\allowbreak{}0\allowbreak{}2\allowbreak{}8\allowbreak{}0\allowbreak{}e\allowbreak{}6\allowbreak{}4\allowbreak{}b\allowbreak{}5\allowbreak{}c\allowbreak{}8\allowbreak{}4\allowbreak{}4\allowbreak{}b\allowbreak{}a\allowbreak{}8\allowbreak{}6\allowbreak{}f\allowbreak{}7\allowbreak{}3\allowbreak{}c\allowbreak{}b\allowbreak{}6\allowbreak{}b\allowbreak{}d\allowbreak{}0\allowbreak{}9\allowbreak{}c\allowbreak{}5\allowbreak{}e\allowbreak{}3\allowbreak{}5\allowbreak{}6\allowbreak{}9\allowbreak{}2\allowbreak{}c\allowbreak{}3\allowbreak{}2\allowbreak{}a\allowbreak{}c\allowbreak{}c\allowbreak{}f\allowbreak{}f\allowbreak{}5\allowbreak{}9\allowbreak{}3\allowbreak{}0\allowbreak{}0\allowbreak{}7\allowbreak{}b\allowbreak{}5\allowbreak{}f\allowbreak{}f\allowbreak{}2\allowbreak{}f\allowbreak{}3\allowbreak{}9\allowbreak{}8}). The certificate also records a separately proved 100-micro-unit upper widening, 14 length cells and 102 triangles; J1 uses the genuine inherited I30 upper source instead. Neither a .126660 numerical request, the earlier .12665668 fallback nor a future recursive source is a premise.

The new analytic comparisons are controller082 D1--D3 and near082 (2)--(9), packed respectively at \texttt{r\allowbreak{}e\allowbreak{}s\allowbreak{}e\allowbreak{}a\allowbreak{}r\allowbreak{}c\allowbreak{}h\allowbreak{}/\allowbreak{}l\allowbreak{}o\allowbreak{}c\allowbreak{}a\allowbreak{}l\allowbreak{}\_\allowbreak{}c\allowbreak{}y\allowbreak{}c\allowbreak{}l\allowbreak{}e\allowbreak{}0\allowbreak{}8\allowbreak{}2\allowbreak{}/\allowbreak{}c\allowbreak{}o\allowbreak{}n\allowbreak{}t\allowbreak{}r\allowbreak{}o\allowbreak{}l\allowbreak{}l\allowbreak{}e\allowbreak{}r\allowbreak{}/\allowbreak{}R\allowbreak{}E\allowbreak{}S\allowbreak{}E\allowbreak{}A\allowbreak{}R\allowbreak{}C\allowbreak{}H\allowbreak{}.\allowbreak{}m\allowbreak{}d} and \texttt{r\allowbreak{}e\allowbreak{}s\allowbreak{}e\allowbreak{}a\allowbreak{}r\allowbreak{}c\allowbreak{}h\allowbreak{}/\allowbreak{}l\allowbreak{}o\allowbreak{}c\allowbreak{}a\allowbreak{}l\allowbreak{}\_\allowbreak{}c\allowbreak{}y\allowbreak{}c\allowbreak{}l\allowbreak{}e\allowbreak{}0\allowbreak{}8\allowbreak{}2\allowbreak{}/\allowbreak{}n\allowbreak{}e\allowbreak{}a\allowbreak{}r\allowbreak{}\_\allowbreak{}r\allowbreak{}e\allowbreak{}c\allowbreak{}o\allowbreak{}r\allowbreak{}d\allowbreak{}/\allowbreak{}R\allowbreak{}E\allowbreak{}S\allowbreak{}E\allowbreak{}A\allowbreak{}R\allowbreak{}C\allowbreak{}H\allowbreak{}.\allowbreak{}m\allowbreak{}d}, with its final matching statement \texttt{F\allowbreak{}I\allowbreak{}N\allowbreak{}A\allowbreak{}L\allowbreak{}.\allowbreak{}m\allowbreak{}d} and version bindings \texttt{S\allowbreak{}O\allowbreak{}U\allowbreak{}R\allowbreak{}C\allowbreak{}E\allowbreak{}\_\allowbreak{}U\allowbreak{}S\allowbreak{}E\allowbreak{}\_\allowbreak{}F\allowbreak{}I\allowbreak{}N\allowbreak{}A\allowbreak{}L\allowbreak{}.\allowbreak{}j\allowbreak{}s\allowbreak{}o\allowbreak{}n}. The latter supplies simultaneous Cauchy and the complete adaptive exterior, rather than the incomplete squarefree-column improvement recorded by recombination082. Catalogue G137 is an explicitly conditional witness reduction and is not a premise here. The retained \texttt{A\allowbreak{}f\allowbreak{}f\allowbreak{}i\allowbreak{}n\allowbreak{}e\allowbreak{}C\allowbreak{}o\allowbreak{}u\allowbreak{}n\allowbreak{}t\allowbreak{}F\allowbreak{}i\allowbreak{}b\allowbreak{}e\allowbreak{}r}, \texttt{P\allowbreak{}o\allowbreak{}s\allowbreak{}i\allowbreak{}t\allowbreak{}i\allowbreak{}v\allowbreak{}e\allowbreak{}S\allowbreak{}o\allowbreak{}u\allowbreak{}r\allowbreak{}c\allowbreak{}e\allowbreak{}M\allowbreak{}o\allowbreak{}m\allowbreak{}e\allowbreak{}n\allowbreak{}t\allowbreak{}s}, \texttt{B\allowbreak{}o\allowbreak{}u\allowbreak{}n\allowbreak{}d\allowbreak{}s\allowbreak{}C\allowbreak{}o\allowbreak{}r\allowbreak{}e}, \texttt{N\allowbreak{}u\allowbreak{}m\allowbreak{}e\allowbreak{}r\allowbreak{}a\allowbreak{}t\allowbreak{}o\allowbreak{}r\allowbreak{}M\allowbreak{}a\allowbreak{}s\allowbreak{}s\allowbreak{}D\allowbreak{}a\allowbreak{}m\allowbreak{}p\allowbreak{}i\allowbreak{}n\allowbreak{}g}, \texttt{P\allowbreak{}o\allowbreak{}s\allowbreak{}i\allowbreak{}t\allowbreak{}i\allowbreak{}v\allowbreak{}e\allowbreak{}L\allowbreak{}a\allowbreak{}b\allowbreak{}e\allowbreak{}l\allowbreak{}P\allowbreak{}a\allowbreak{}c\allowbreak{}k\allowbreak{}e\allowbreak{}t} and generic physical-return Lean sources identify the exact source/height/coefficient contracts; their declarations do not constitute end-to-end formal verification of J1. In particular the fully summed contraction J9a is retained from \texttt{r\allowbreak{}e\allowbreak{}s\allowbreak{}e\allowbreak{}a\allowbreak{}r\allowbreak{}c\allowbreak{}h\allowbreak{}/\allowbreak{}l\allowbreak{}o\allowbreak{}c\allowbreak{}a\allowbreak{}l\allowbreak{}\_\allowbreak{}c\allowbreak{}y\allowbreak{}c\allowbreak{}l\allowbreak{}e\allowbreak{}0\allowbreak{}7\allowbreak{}7\allowbreak{}/\allowbreak{}i\allowbreak{}n\allowbreak{}t\allowbreak{}a\allowbreak{}k\allowbreak{}e\allowbreak{}/\allowbreak{}p\allowbreak{}r\allowbreak{}1\allowbreak{}4\allowbreak{}3\allowbreak{}/\allowbreak{}s\allowbreak{}u\allowbreak{}b\allowbreak{}m\allowbreak{}i\allowbreak{}s\allowbreak{}s\allowbreak{}i\allowbreak{}o\allowbreak{}n\allowbreak{}s\allowbreak{}/\allowbreak{}2\allowbreak{}0\allowbreak{}2\allowbreak{}6\allowbreak{}-\allowbreak{}1\allowbreak{}1\allowbreak{}-\allowbreak{}1\allowbreak{}3\allowbreak{}-\allowbreak{}c\allowbreak{}o\allowbreak{}n\allowbreak{}t\allowbreak{}i\allowbreak{}n\allowbreak{}u\allowbreak{}u\allowbreak{}m\allowbreak{}-\allowbreak{}e\allowbreak{}x\allowbreak{}a\allowbreak{}c\allowbreak{}t\allowbreak{}-\allowbreak{}c\allowbreak{}o\allowbreak{}n\allowbreak{}t\allowbreak{}a\allowbreak{}c\allowbreak{}t\allowbreak{}/\allowbreak{}s\allowbreak{}r\allowbreak{}c\allowbreak{}/\allowbreak{}C\allowbreak{}a\allowbreak{}n\allowbreak{}d\allowbreak{}i\allowbreak{}d\allowbreak{}a\allowbreak{}t\allowbreak{}e\allowbreak{}/\allowbreak{}L\allowbreak{}o\allowbreak{}c\allowbreak{}a\allowbreak{}l\allowbreak{}/\allowbreak{}Q\allowbreak{}r\allowbreak{}h\allowbreak{}/\allowbreak{}T\allowbreak{}w\allowbreak{}o\allowbreak{}S\allowbreak{}l\allowbreak{}o\allowbreak{}t\allowbreak{}/\allowbreak{}A\allowbreak{}n\allowbreak{}c\allowbreak{}h\allowbreak{}o\allowbreak{}r\allowbreak{}C\allowbreak{}o\allowbreak{}n\allowbreak{}t\allowbreak{}i\allowbreak{}n\allowbreak{}u\allowbreak{}u\allowbreak{}m\allowbreak{}/\allowbreak{}P\allowbreak{}o\allowbreak{}s\allowbreak{}i\allowbreak{}t\allowbreak{}i\allowbreak{}v\allowbreak{}e\allowbreak{}L\allowbreak{}a\allowbreak{}b\allowbreak{}e\allowbreak{}l\allowbreak{}P\allowbreak{}a\allowbreak{}c\allowbreak{}k\allowbreak{}e\allowbreak{}t\allowbreak{}.\allowbreak{}l\allowbreak{}e\allowbreak{}a\allowbreak{}n}, not from a subpower assertion about its unsummed error factor.

The full proof since OpenAI is retained, and properly derived alternatives remain permitted. Credits and terms of A--I remain, including Sil's canonical marked sources, Joseph Bebel's exact contact and count work, Akash Levy's actual numerator, and Hailey Collet's signed reflection/probe/return framework. The physical reductions retain SovereignSteak's draft PR43 at head \texttt{c\allowbreak{}8\allowbreak{}3\allowbreak{}8\allowbreak{}a\allowbreak{}7\allowbreak{}5\allowbreak{}0\allowbreak{}0\allowbreak{}7\allowbreak{}9\allowbreak{}f\allowbreak{}7\allowbreak{}f\allowbreak{}8\allowbreak{}6\allowbreak{}a\allowbreak{}5\allowbreak{}7\allowbreak{}3\allowbreak{}d\allowbreak{}f\allowbreak{}9\allowbreak{}c\allowbreak{}9\allowbreak{}f\allowbreak{}7\allowbreak{}2\allowbreak{}9\allowbreak{}b\allowbreak{}8\allowbreak{}b\allowbreak{}c\allowbreak{}9\allowbreak{}2\allowbreak{}d\allowbreak{}1\allowbreak{}2\allowbreak{}5\allowbreak{}e}, its contributor notices and Codex assistance disclosure; no conditional donor endpoint affirmation is assumed. The Sepia/CWR attribution and terms of \emph{Mathematical Writing for AI Agents} remain part of the retained packet.

\hypertarget{appendix-k}{}

\hypertarget{appendix-k.-actual-zero-witnesses-and-one-earlier-pair-source}{%
\section{Appendix K. Actual zero witnesses and one earlier-pair source}\label{appendix-k.-actual-zero-witnesses-and-one-earlier-pair-source}}

On the retained foundations of A--J, including the stated natural selected-cell input SN9, we prove the ordinary fixed-modulus, all-height conclusion at \[\fitmath{
 c_K=\frac{63329}{500000}=.126658,\qquad
 t_K=1-c_K=.873342,\qquad h=h_*.                              
}\tag{K1}\] Here \(\fitinline{h_*}\), \(\fitinline{\alpha}\), \(\fitinline{g_0}\) and \(\fitinline{g_1}\) are the exact contact data already retained in E, G and I. The conclusion is in the strict half-plane \(\fitinline{\Re s>t_K}\), with an actual principal pole excepted. No boundary nonvanishing, growing-modulus uniformity or new end-to-end Lean verification is asserted.

There are two new comparisons. A zero from the original finite family supplies a genuinely saturated inverse and ordinary pair on each nonfloor row. It is used to bound that row's numerator, without replacing the physical denominator. Separately, a previously proved marked-pair source can be used once on certain regular sixth-free cores. This is an input from a strictly earlier source, before lower-mark division. The complete source below covers 40 closed length bands; neither a conjectural fixed point nor the prospective .126660 calculation is used.

\hypertarget{k.1.-the-actual-reciprocal-and-its-three-contour-variables}{%
\subsection{K.1. The actual reciprocal and its three contour variables}\label{k.1.-the-actual-reciprocal-and-its-three-contour-variables}}

Keep the original physical cube H43, its complete main/error expansion, all units, Euler deletions, nonunit zeros and coefficient weights. Its reciprocal is \[\fitmath{
                    1/L_K^{\mathcal S}(s,\eta\overline{\psi_u}).
}\tag{K2}\] Its argument is \(\fitinline{s}\). The separate condition \(\fitinline{\Re(s+w)>1+\epsilon}\) is the normal-convergence guard for the quotient-free Euler product \(\fitinline{H^{\rm comp}_{\eta,u,Z}(s,w,z)}\); it is not a shifted argument of K2.

Follow B.8 in its original order. Move \(\fitinline{z}\) to \(\fitinline{17/50}\) while \(\fitinline{s,w}\) are absolute. Use the global line \(\fitinline{\Re s=\beta+20e}\), \(\fitinline{\Re w=3}\), before truncation, where \(\fitinline{\beta}\) is the actual primitive-family supremum and \(\fitinline{e>0}\) is the small buffered loss. Only the retained \(\fitinline{s}\) segment is moved to \(\fitinline{\Re s=a+16e}\); B(206) bounds the actual reciprocal to the right of \(\fitinline{a+6e}\). The independent imaginary variables are \(\fitinline{t_s,t_w,t_z}\). Keep the Gaussian in \(\fitinline{t_s+t_z}\), the independent Schwartz weights and the horizontal joins on their original buffered segments. The omitted \(\fitinline{s}\) segments remain on the global line. The local reciprocal estimate is uniform before the later detector and numerator Fourier frequencies are chosen.

In H45's finite-grid convention the detected zero has height at most \(\fitinline{3i_uT}\), the upper family zero guard extends to \(\fitinline{3(i_u+1)T}\), and the retained height is at most \(\fitinline{(3i_u+1)T}\). An additional local shift of at most \(\fitinline{T/2}\) stays below \(\fitinline{(3i_u+2)T}\). B(205)'s alternative pigeonhole convention instead uses \(\fitinline{3^jT}\) and \(\fitinline{3^{j+1}T}\), with \(\fitinline{3^jT+T\le3^{j+1}T}\). Both provide the needed buffer; their indices are not identified. Thus K2 is bounded for the actual physical character in the fixed family, with its imprimitive holes, even when a different family member realizes the detector zero. The floor, bounded rows, principal residues, other five units and fixed-pad high tails retain their separate B/H proofs.

Full \(\fitinline{w}\)-inversion gives the reflected numerator and its genuine coefficient packet. Before improvement, B.8 bounds this term by \[\fitmath{
 Y^{-1/2}Z^{\ell(z_0-1/2)+\ell q}
 U^{1/4+3R/4-qC/(2D)+\epsilon}(1+T)^J,\qquad U=Z^D.           
}\tag{K3}\] Here \(\fitinline{d=2a-1}\), \(\fitinline{q=dp}\), and \(\fitinline{R}\) is an admitted upper count on the actual class. The external factors in H43 are unchanged; in particular its base price includes the \(\fitinline{d/2}\) from the retained \(\fitinline{s}\)-line. This is a bound for the original reciprocal at \(\fitinline{s}\), not for \(\fitinline{1/L(s+w)}\).

The exact scalar certificate conservatively keeps J's two reciprocal-length tails as well. After full \(\fitinline{w}\)-inversion return to the absolute \(\fitinline{s}\)-line, open the actual denominator and retain its annuli. The shifts \(\fitinline{-.10,+.27}\) pay the omitted lengths as in J4. On a retained annulus, annular Mellin inversion and the same local buffer give \(\fitinline{|M_{\rm den}(U^r)|\ll U^{dr/2+\epsilon}}\), whose power is cancelled by its original exterior scalar. The row comparison below then applies to the numerator. This second implementation does not restrict the physical denominator by a detector witness.

\hypertarget{k.2.-a-genuine-family-zero-pays-a-stronger-row-count-price}{%
\subsection{K.2. A genuine family zero pays a stronger row-count price}\label{k.2.-a-genuine-family-zero-pays-a-stronger-row-count-price}}

On a fixed nonfloor density/amplitude cell retain J3's definitions \(\fitinline{k,V,J,t_0,m,b_0,R}\). In particular \(\fitinline{R=1/6+Vt_0=1-d[p(1+2g_0)+kb_0]}\) and \(\fitinline{b_0=2m}\). Choose \(\fitinline{T_*=t_0-\gamma}\) before selecting witnesses. Compactness permits a sufficiently fine fixed subdivision in \(\fitinline{d,p}\) and a representative cutoff in each bin, with an arbitrarily small paid error. The certificate checks \(\fitinline{1<T_*<3/2}\) with a strict margin.

Lemma B.5.1 supplies to every such row a presentation \(\fitinline{\varphi=\nu\psi_u^{\pm1}}\) from the original finite dictionary, one actual zero with real part in \(\fitinline{[a,a+e)}\), a common height, and two dyads satisfying \[\fitmath{
 \begin{gathered}
 x\le t_0-\gamma+o(1),\qquad x+y\ge t_0-\gamma-o(1),\\
 |M_\varphi(U^x)|^2\gg U^{dx-O(\epsilon_0)},\qquad
 |S_\varphi(U^y)|^2\gg U^{dy-O(\epsilon_0)}.                 
 \end{gathered}
}\tag{K4}\] For clarity, the proof uses the original \(\fitinline{\mu*1}\) identity at that zero with a cutoff \(\fitinline{U^{T_*}}\). Its unit contribution and cancellation below the cutoff give a large product after dyadic/Fourier separation. B(207) supplies both upper bounds, hence both lower bounds in K4, at the same zero. Fixing presentation and dyads gives common base profiles. Real powers, outside phases, scale derivatives and norm heights remain in their common profile envelopes. On the tangent inverse annulus the cutoff is identically one; its transition belongs to the native long branch. These are proved detector witnesses, not saturation hypotheses about K2.

Partition rows in order. First take those whose selected inverse has \(\fitinline{x\ge b_0+\lambda_x}\), including its fixed continuations and amplified branch; then those whose ordinary has \(\fitinline{y\ge m+\lambda_y}\). The original marked moments and K4 give raw row-count savings \(\fitinline{dk\lambda_x,2dk\lambda_y,V\gamma}\). The genuine numerator fourth on a subset costs the \(\fitinline{3/4}\) power of that subset's count. Consequently the three physical savings in \(\fitinline{U}\) are \[\fitmath{
          3dk\lambda_x/4,\qquad 3dk\lambda_y/2,
                       \qquad 3V\gamma/4.                    
}\tag{K5}\] The six fixed H31 gaps also earn \(\fitinline{3/4}\) of their raw count saving. This is a distinct use of a genuinely saturated detector. J's arbitrary physical factors still have their original moment-order prices \(\fitinline{g/2}\) or \(\fitinline{g/4}\); K5 does not reprice them.

The remaining rows have \(\fitinline{b_0-\gamma-\lambda_y-o(1)\le x<b_0+\lambda_x}\) and \(\fitinline{m-\gamma-\lambda_x-o(1)\le y<m+\lambda_y}\), wholly inside the actual anchors. Their inverse lower bound and second moment give count at most \(\fitinline{U^{R-dk(x-b_0)_++O(\epsilon_0)}}\). Their ordinary lower bound gives the rowwise comparison \[\fitmath{
 |P_{\rm true}(n)|\ll U^{-dy/2+O(\epsilon_0)}
                  |P_{\rm true}(n)S_\varphi(y)|.               
}\tag{K6}\] This applies to the entire numerator term on that row; its detector is not substituted into K2. A proved mixed source for these two fixed presentations, followed by its genuine lower-mark division and Cauchy on the same rows, has exponent \[\fitmath{
 \tfrac12\{1-dy-2qW+R-dk(x-b_0)_+\},\qquad
 W=\tfrac12[1+g_0-n_{\rm ref}-(1+2g_1)y]+\Delta.             
}\tag{K7}\] Here \(\fitinline{n_0=1-y_t/D}\), \(\fitinline{n_{\rm ref}=\max(n,n_0)}\), and actual shorter factors use the proved componentwise shortening. The count in K7 is not an extra squared moment of the selected inverse.

Subtract K7 from \(\fitinline{(3R+1)/4-qz_0/2}\), with \(\fitinline{z_0\le1/2-n_0}\), exactly as in J9. Use \(\fitinline{y-m+(x-b_0)_+\ge-\gamma-o(1)}\). The fully summed rowwise physical packet mass J9a pays \(\fitinline{dv/2}\), where \(\fitinline{v=\log_U\max(1,L_iK)}\), and the actual support has \(\fitinline{n_{\rm ref}-n_0\le v+\mathrm{ann}}\). Thus the middle saving is again at least \[\fitmath{
                           d(p\Delta-k\gamma/2),              
}\tag{K8}\] with directed annular costs, but now with the smaller \(\fitinline{\gamma}\) paid by K5. The complete physical label sum, its true external error weights and its \(\fitinline{\prod Y_p^{-4/25}}\) factor are retained. No \texttt{e\allowbreak{}r\allowbreak{}r\allowbreak{}o\allowbreak{}r\allowbreak{}P\allowbreak{}r\allowbreak{}o\allowbreak{}d\allowbreak{}u\allowbreak{}c\allowbreak{}t\allowbreak{}S\allowbreak{}i\allowbreak{}z\allowbreak{}e} is declared subpower and no dilation mass is charged twice.

Row-dependent choices of presentation, dyad and height are handled before the fixed-height source call. Freeze the finite dictionary, partition scales by bounded ratios, keep all outside phases, and use the common compact log-Fourier envelope, or its actual \(\fitinline{xW'(x)}\) derivative, for scale/dilation suprema. Unit-interval Sobolev and the normalized first log tests handle the two separate height variables. Their mixed derivative is the product of the two allowed profiles. Choose the resulting finite seminorm and height degrees before the one target-dependent height exponent. This preserves the physical numerator height independently of the detector height and the three contour variables above.

\hypertarget{k.3.-the-earlier-pair-is-an-actual-source-before-mark-division}{%
\subsection{K.3. The earlier pair is an actual source before mark division}\label{k.3.-the-earlier-pair-is-an-actual-source-before-mark-division}}

We need the pre-division source H.5(A) for two separate fixed finite rays multiplying the same moving sextic feature. This version follows within its proof. Insert \(\fitinline{\omega_1(l_1)\omega_2(l_2)}\) in the actual ordinary divisor columns before either Poisson identity. Their frequencies depend on the common product, while the separate multiplicative factors remain on their columns. Every allocation retains \(\fitinline{\omega_i(d_i)\omega_i(l_i/d_i)}\), including its zeros. In C.7's two Euler identities substitute \(\fitinline{X\mapsto\omega_1(\mathfrak p)X}\), \(\fitinline{Y\mapsto\omega_2(\mathfrak p)Y}\) before coefficient extraction. Its bidegree \(\fitinline{(i,j)}\) coefficient gains that exact scalar, so both signed identities, prime-local correlations and H25--H26 divisor budgets remain valid.

On a fixed inducing core use H8 for these two actual inducing characters and the same independent profiles. Equal-product centering cancels their product means, including a zero mean or two principal means. The regular fourths retain their own keeps; each single anchor reproduces the original prime ray by its own relative multiplier on the common good support, with its outside unit and complete annulus. The fixed family includes both presentations. H6 is applied to the merged true coefficient vector. The original short-reference, all-row centered-square, zero and tail estimates therefore prove the same pre-division marked product estimate for this pair. The fixed modulus is chosen before asymptotics, so the whole deleted-modulus norm is below every positive live prime scale eventually. Whole-factor conjugation retains all zeros. There is no division at a nonunit or source theorem for an arbitrary moving arithmetic multiplier.

This proves the required version both for \(\fitinline{P_{\rm true},S_\varphi}\) in K6 and for the regular grouped core. In the latter write \(\fitinline{z=b^6c}\) and hold \(\fitinline{b}\), the outer label, rays, heights and Fourier variables fixed. The actual grouped map \(\fitinline{z=V_0A_2z'}\) gives, on supported ideals, two fixed multipliers times the same power-one feature \(\fitinline{\chi_I(c)}\), with the genuine \(\fitinline{\chi_I(b)^6}\) zero mask. It does not give a new square or cube feature. The common-ball label/sublattice count has already been paid by H16. Restrict the positive source to its actual regular subrows; the short reference gets smaller and the centered square keeps its full second term.

The earlier sources extend to the following first-coordinate ceilings: \[\fitmath{
 \begin{array}{c|c|c}
 n&y&\beta_{\rm old}\\ \hline
 [.42495,.4271]&[.37445,.378]&.000100\\{}
 [.42506,.4271]&[.37452,.3775]&.000112.
 \end{array}                                                 
}\tag{K9}\] Indeed in I19/I21 and J17b, \(\fitinline{L-P(V)/2}\) and the entire low bracket are independent of \(\fitinline{n}\); their terminal certificates are unchanged. The remaining ancillary expressions are affine at the two coordinate endpoints. Their exact check gives maximum volumes \(\fitinline{2821728153/3125000000}\) and \(\fitinline{18055060141/20000000000}\), both below .903, and widths below .101. It checks \(\fitinline{H>4(.139+W)}\), \(\fitinline{5/6+.139>A_0+3u}\), \(\fitinline{5u<1}\), \(\fitinline{A_0+4u<1}\), the native ranges, tiny cores and fixed SN9 gap. The longer reference can reach .6661; its original nonprincipal unmarked fourth applies at every fixed bounded-length exponent. It is the other, .139 coordinate whose marked fourth needs the short gate. Thus K9 is a proved enlargement of earlier sources, including their independent heights, first log tests and componentwise shortening.

Remove the \(\fitinline{b}\)-holes of the two ordinary factors by their exact ideal M\"{o}bius expansions. Each \(\fitinline{l_i\mid\operatorname{rad}(b)}\) shortens only its own coordinate and contributes modulus at most \(\fitinline{(Nl_i)^{-1/2}}\); sum these positive scalars after the identity. Their mass is \(\fitinline{U^\epsilon}\) under the genuine polynomial cap on \(\fitinline{Nb}\). Keep common-prime zeros; no \(\fitinline{(b,c)=1}\) restriction is made. Similarly write each punctured prime pool as its clean complete annulus minus the actual deleted terms and expand in the positive norm. At a fixed positive slot width the latter sum is \(\fitinline{O_j(P_j^{-1/2}(1+\log U/\log P_j))}\). Fixed holes obey the same bound. Zero-width pools retain their bounded terms, and the finitely many sufficiently small slots have total arbitrarily small cost after the slot family is fixed. No prime deletion mask is assumed smooth and no signed arm is silently cleaned.

Let the actual core-radius exponent be \(\fitinline{V}\), and put \(\fitinline{Z_0=U^V}\). At an admitted point the earlier source permits width \[\fitmath{
 w(n/V,y/V)=[1+g_0-n/V-(1+2g_1)y/V]/2+\beta_{\rm old}.
}\] From the genuine residual slots choose a whole prefix of width at most \(\fitinline{Vw}\), with its strict decrement. The prefix can be chosen from the frozen descriptors before estimating the row sum. On a clean maximal squared-response bin \(\fitinline{e}\), the complementary product costs \(\fitinline{U^{e\,\mathrm{width}+\epsilon}}\). Thus the positive regular-core energy has excess \[\fitmath{
 C_{\rm pair}(V,e)=e[T(V)]_+,\qquad
 T(V)=\beta-u+A(1-V)-V\beta_{\rm old},\quad A=(1+g_0)/2.       
}\tag{K10}\] Here the new live width is \(\fitinline{L=W_{\rm base}(n,y)+\beta-u}\). The positive part is essential: exhaustion gives zero excess, not a negative energy price. Slots missing from the live product have no residual charge.

Source rescaling is literal: \(\fitinline{Z_0^{w_j/V}=U^{w_j}=P_j}\). Retain the actual normalized heights and all fixed profile constants. On inner calls \(\fitinline{V>.98}\); choose the new early maximum slot width with \(\fitinline{w_{\max,\rm new}/V_{\min}<w_{\max,\rm old}}\) and below every prescribed old decrement/reserve before the later slot type. Choose the earlier source's finite polynomial degrees at the actual dictionary first, then a new common height exponent below their rescaled admissible exponents. This is a finite dependency order with one strictly earlier source call, not self-use.

\hypertarget{k.4.-every-core-radius-including-the-clamp}{%
\subsection{K.4. Every core radius, including the clamp}\label{k.4.-every-core-radius-including-the-clamp}}

The 40-band source uses \(\fitinline{n\in[.42506,.42532]}\). If a whole closed band fits one of K9's sources and that source's surplus strictly exceeds the required physical width, use it. Otherwise set \[\fitmath{
 \beta=10^{-6}(\lceil10^6\Delta_{\rm required}\rceil+1).
}\tag{K10a}\] The five extreme bands with \(\fitinline{\beta\le.0001}\) use the nonrecursive calculation at \(\fitinline{u=.0015}\). Twenty bands use K10 once with \(\fitinline{\beta_{\rm old}=.0001}\), taking \[\fitmath{
 u=\min\!\left(.0006,\frac1{20000}
       \left\lfloor20000\left[\frac{1-y_+/.378}{6}-.00001\right]\right\rfloor\right),
 \qquad H=1-6u.                                               
}\tag{K10b}\] Their \(\fitinline{u}\) ranges from .00035 to .0006. These are finite band choices before the live rows or response bins. Their entry radius is \[\fitmath{
 V_{\rm ent}=\left\lceil\max\{.42532/.4271,y_+/.378\}\right\rceil_{10^{-6}}
                 +10^{-6},\qquad .98<V_{\rm ent}<H<1.         
}\tag{K11}\] This admits every normalized upper coordinate on \(\fitinline{[V_{\rm ent},1]}\). Below entry the pair is not presumed available.

For each of these bands the exact low and same-zero curved alternatives J17b are checked both at \(\fitinline{H}\) and at \(\fitinline{V_{\rm ent}}\). At \(\fitinline{H}\), the low coefficient is the smaller of the curved bracket and \(\fitinline{[T(H)]_+}\); the two middle alternatives remain the genuine nonrecursive ones. The budget is \(\fitinline{u-\rho}\), with \(\fitinline{\rho=10^{-7}}\). At entry use the curved alternatives alone and budget \(\fitinline{u-\rho+(5/6)(H-V_{\rm ent})}\). As in I22--I25, integer fifth-power bounds at denominator \(\fitinline{10^{13}}\), radical supporting tangents, downward rounding of the low threshold at \(\fitinline{10^{-6}}\), and a rational cutoff rounded down at \(\fitinline{10^{-7}}\) establish each closed affine triangle. The changed main and entry tests accept 56 closed triangles without further bisection; every residual in these polygons exceeds \(\fitinline{10^{-9}}\). This proves all response and common-zero parameters, rather than selected detector values.

This proves whole intervals, as follows. The curved energy excesses are decreasing in \(\fitinline{V}\) on \(\fitinline{[.98,1]}\), with absolute slope less than \(\fitinline{5/6}\). For the pair low branch, \(\fitinline{e\le.5002}\) and \(\fitinline{C'_{\rm pair}\ge-.5002(A+\beta_{\rm old})>-5/6}\) wherever its positive part is active; its derivative is zero elsewhere. Therefore the target allowance \(\fitinline{u-\rho+(5/6)(H-V)_+}\) grows at least as fast as the chosen bound when \(\fitinline{V}\) decreases. The \(\fitinline{H}\) check covers \(\fitinline{[V_{\rm ent},H]}\), and both bounds decrease above \(\fitinline{H}\). The independent curved check at entry covers \(\fitinline{[.98,V_{\rm ent}]}\). Its threshold and triangle choice may differ from the \(\fitinline{H}\) choice; the two closed covers meet at entry. This is a full radius argument, not a check of the slice \(\fitinline{V=H}\).

For \(\fitinline{V\ge1}\), actual coordinates can fall below the lower edge of K9. Use downward restriction at \(\fitinline{n^\dagger=\max(n/V,.42495)}\), \(\fitinline{y^\dagger=\max(y/V,.37445)}\). The actual uncovered width is now \[\fitmath{
 T_{\rm cl}(V)=T(V)+\tfrac12(V\cdot.42495-n)_+
                    +\tfrac{1+2g_1}{2}(V\cdot.37445-y)_+ .    
}\tag{K12}\] At \(\fitinline{V=1}\), \(\fitinline{T_{\rm cl}\le\beta-u-.0001<0}\). On every affine piece its derivative is at most \[\fitmath{
 -\frac{1+.05805911}{2}-.0001+\frac{.42495}{2}
              +\frac{1+2(.07599620)}2\,.37445
       =-\frac{10097277791}{10^{11}}<0.                       
}\tag{K13}\] The negative term uses the lower bound for \(\fitinline{g_0}\); the positive term uses the upper bound for \(\fitinline{g_1}\). Thus \(\fitinline{[T_{\rm cl}(V)]_+=0}\) for every \(\fitinline{V\ge1}\): the old pair takes all residual slots. A clamped comparison with an uncharged larger radius would not suffice; K12 contains its true width cost.

The other 20 bands use six held100, nine held112 and five new nonrecursive curved calls. The five latter use \(\fitinline{u=.0015}\), \(\fitinline{\rho=10^{-7}}\), and the same closed polygon construction. Their native \(\fitinline{V\ge1}\) guard is checked directly. On every new band the peripheral native ranges down to .98, the two unmarked fourths down to .875, tiny cores below .08, and the SN9 range \(\fitinline{e\ge.5002}\) at slope .5001 with fixed positive gap .0001 retain positive exact guards. Widths are below .101 and volumes below .903. The source energy on every radius is consequently \[\fitmath{
 \ll\mathcal H U^{V+u-\rho+(5/6)(H-V)_++\epsilon}.             
}\tag{K14}\] All earlier held bands already have the stronger I/J source statement.

Restore \(\fitinline{z=b^6c}\), with \(\fitinline{V=M-6\log_U Nb}\), using exactly this \(\fitinline{5/6}\) slope and all common primes. At the actual selected radius \(\fitinline{16K_{\rm first}(\prod2P_j)^4}\), \(\fitinline{M\ge1-6u'-o(1)}\); a filled selection has \(\fitinline{u'=u+o(1)}\). K14 pays the grouped target; exhaustion uses its actual unmarked face. The first artificial divisor stays unweighted in joint Cauchy. The two orientations retain their own reflected rays and sublattices, the H26 second-label reciprocal cancellation and their same final denominator \(\fitinline{NC\,ND_1}\). Fixed cores use the equal-product signed centering for both actual inducing characters, including both principal arms; the \(\fitinline{3/5}\) cancellation pays the whole second-divisor extraction. The .139 short reference and its actual longer partner use their original marked/unmarked fourths. All zero, seed, principal and rapid-tail terms retain their exact scalar before restoration. The source is thus the full kept uncentered marked root, with the original independent heights, two first-log tests and componentwise shortenings, before lower-mark division. It is not just a regular terminal leaf.

\hypertarget{k.5.-all-original-boxes-and-the-source-first-physical-return}{%
\subsection{K.5. All original boxes and the source-first physical return}\label{k.5.-all-original-boxes-and-the-source-first-physical-return}}

Use J11--J15's exact E49 margin, saved polynomials and conductor convexity at the new target, retaining the actual unused reserve in each box. The source data still have \(\fitinline{\mathrm{momentKappa}=2\beta-1}\); only conservative count/capacity pricing uses \(\fitinline{2\alpha-1}\), by FB2--FB6 and the previously attained endpoint. No equality between count alternatives is imposed.

One additional paired admission is needed, on original cell34: \[\fitmath{
                3/8\le d\le27/64,\qquad7/16\le p\le1/2.      
}\tag{K15}\] It uses the existing E29 theorem, not a postulated new source. Put \(\fitinline{T=J+dU_d}\), \(\fitinline{N=(1-d)J+VdU_d}\), \(\fitinline{r_0=.7335}\), \(\fitinline{r_1=.7775}\), \(\fitinline{B_0=1-10^{-12}}\), \(\fitinline{L_0=2+10^{-11}}\), \(\fitinline{c_0=.15-10^{-7}}\), \(\fitinline{s_m=.027}\), \(\fitinline{y_m=(1-4.5s_m)/2}\). On K15 the ten cleared guards are \[\fitmath{
 \begin{gathered}
 k,\ J,\ T-J,\ 1.2499J-T,\ 2T-3(r_0+10^{-6})J,\\
 N-[1-d(r_1+(1-r_1)p)]J,\\
 N-[1-d(r_1+2(B_0-r_1)p/L_0)]J,\\
 N-[1-d(r_1+2c_0p)]J,\quad N-[1-2y_md-2s_mdp]J,\\
 -2(1+2g_1)dU_d+(1+6g_0-4g_1)J .
 \end{gathered}                                              
}\tag{K16}\] Tensor Bernstein lower bounds on the full closed rectangle, at each of the four \(\fitinline{(g_0,g_1)}\) corners, are positive; the least of these forty bounds is \(\fitinline{69989628971/10^{13}>.0069989}\). Clearing uses positive denominators. These are the same-zero crossing, all six actual continuations and antitone amplitude conditions of E29, so the pointwise count is admitted throughout K15. The earlier E27 window and I32 wing remain in force. No other new admission is required for K1.

The 57 original boxes now give 34 saved leaves, 14105 E49 exterior leaves and 263 central witness cells. The respective least bare physical reserves are \[\fitmath{
 3.14287383\cdot10^{-7},\qquad1.3943222\cdot10^{-8},
                         \qquad1.0000049\cdot10^{-8}.          
}\tag{K17}\] For each central cell the excess \(\fitinline{B}\) and the four conductor/dilation/length cutoffs are exactly J's directed construction with budget \(\fitinline{B_+=\max(0,B)+10^{-8}}\). Replace only its moving cutoffs by \[\fitmath{
 \gamma=\left\lceil\frac{4B_+}{3D_-V_-}\right\rceil_{10^{-9}},\quad
 \lambda_x=\left\lceil\frac{4B_+}{3D_-ak_-}\right\rceil_{10^{-9}},\quad
 \lambda_y=\left\lceil\frac{2B_+}{3D_-ak_-}\right\rceil_{10^{-9}}.
}\tag{K18}\] Use K5 and \(\fitinline{3/4}\) of all six raw fixed gaps, retaining their original length and mark admissions. J15 then supplies the width needed in K8. Annular slack is \(\fitinline{10^{-8}}\); all remaining support, source and bin costs are chosen with total below \(\fitinline{10^{-9}}\).

The central union is enclosed by \[\fitmath{
 \begin{gathered}
 D\ge.815269865,\quad .41357421875\le d\le.43798828125,\quad
 p\ge.498046875,\\
 .425065010017674\le n_{\rm ref}\le.425312776053928,\\
 .374243484648087\le y\le.378142238298387,\quad
 \Delta_{\rm needed}\le.000135586.                            
 \end{gathered}
}\tag{K19}\] Do not charge this maximum width to every length. Intersect each central cell with consecutive closed horizontal bands of width at most \(\fitinline{10^{-4}}\), and take the maximum required width only over cells meeting that band. This gives the 40-band profile. The source just proved supplies its entire numerator interval within \(\fitinline{[.42506,.42532]}\) and a strictly larger width on every band. The least unused width is \(\fitinline{251/250000000=1.004\cdot10^{-6}}\). Insert each exact band edge in the actual scale partition before fixing its source and Fourier anchors; all seam terms and straddling cells are retained. Source width decrements and finite profile costs are chosen below this physical spare and the smaller positive terminal residuals. The terminal energy reserve \(\fitinline{10^{-7}}\) is a separate quantity.

The ordinary/detector anchor and taper guards, all six fixed continuations, and \(\fitinline{1<T_*<3/2}\) are checked on every cell. The least fixed physical continuation reserve exceeds .00017675. The original total slot width is greater than \(\fitinline{1/5}\); greedy whole-slot selection and its positive-exhaustion alternative give the squared common lower response \(\fitinline{2qW}\) on the actual physical class. No radical predicate, squarefree-row restriction or physical-inverse saturation is introduced.

For completeness every changed geometric and peripheral guard is checked at \(\fitinline{t_K}\). Define \(\fitinline{\ell=11/3-4t_K}\), \(\fitinline{y_t=h_*-2\ell}\), \(\fitinline{x_t=1+\ell-h_*}\). These lie in J18, so the direct field proof J19--J20 supplies \texttt{G\allowbreak{}e\allowbreak{}o\allowbreak{}m\allowbreak{}e\allowbreak{}t\allowbreak{}r\allowbreak{}y\allowbreak{}.\allowbreak{}P\allowbreak{}a\allowbreak{}r\allowbreak{}a\allowbreak{}m\allowbreak{}e\allowbreak{}t\allowbreak{}e\allowbreak{}r\allowbreak{}s} and all five \texttt{B\allowbreak{}o\allowbreak{}u\allowbreak{}n\allowbreak{}d\allowbreak{}s\allowbreak{}C\allowbreak{}o\allowbreak{}r\allowbreak{}e} fields, including the exact signal boundary; no previous-target constructor is invoked. The capacity and source guards \(\fitinline{h+F<\min(1,2y_t)}\), \(\fitinline{y_t<.76<h}\), \(\fitinline{\ell/(h+F)>1/5}\), \(\fitinline{h/2-2\ell<3h/40}\) hold directly. The lower contour satisfies \(\fitinline{d(1-p)/2-.10\ge.0033935546875>0}\), and \(\fitinline{(1+d)/2+.27<1}\), \(\fitinline{(1+d)/2-.10>.51}\). The actual tuple damping, quotient-free expansions and Euler deletions thus remain legal.

The averaged floor margin, after \(\fitinline{10^{-5}+10^{-12}}\), exceeds .00317484; the small-row margin exceeds .07138267; the principal margins exceed .02137103 and .0001088741. The other five units, fixed-pad large rows, zero frequencies, finite-prime corrections, compensators and independent Fourier tails retain their B/H proofs. Keep \(\fitinline{F=10^{-258}}\), \(\fitinline{s=10^{-260}}\), and choose \(\fitinline{0<D_{\rm src}.t\le F/4}\) before the live character/heights. Then \(\fitinline{F/2+2D_{\rm src}.t\le F}\). Charge \(\fitinline{15s/16+2F}\) separately to K17, using the same 1-Lipschitz capacities and conductor variation below \(\fitinline{411/400<2}\). All sectors extend to \(\fitinline{D\le h+F}\) with more than \(\fitinline{3\cdot10^{-9}}\) left. No principal pole is paid by a count gain.

Fix first the earlier source banks and their strict decrements, then the new bank, actual contour/cutoff gaps and early width mesh. After the finite slot type and ray dictionary are fixed, choose the common finite profile, Fourier and Sobolev orders. Only then choose the one positive physical height exponent, before every later buffered grid bound. The independent rapid tails and H24a's polynomial absorption retain this order. Rowwise positive masses are summed with their actual exterior factors, not converted into moving common coefficients.

These estimates give the literal normalized cube inequality I51 at the new geometry for every original \texttt{A\allowbreak{}c\allowbreak{}t\allowbreak{}u\allowbreak{}a\allowbreak{}l\allowbreak{}B\allowbreak{}u\allowbreak{}f\allowbreak{}f\allowbreak{}e\allowbreak{}r\allowbreak{}e\allowbreak{}d\allowbreak{}S\allowbreak{}e\allowbreak{}l\allowbreak{}e\allowbreak{}c\allowbreak{}t\allowbreak{}i\allowbreak{}o\allowbreak{}n}, in its original order \(\fitinline{\forall\eta\ \exists\tau_\eta>0\ \forall n>0\ \exists C_n}\), eventually in \(\fitinline{Z}\). Choose a common \(\fitinline{\sigma>0}\) below \(\fitinline{10^{-9}}\) and the transport reserves, with baseline losses below \(\fitinline{\min((\beta-t_K)/10000,10^{-9}/4,F)}\). Keep the same actual data, cutoff, complete H43--H44 cube and nonzero normalizer I52. The generic \texttt{h\allowbreak{}i\allowbreak{}g\allowbreak{}h\allowbreak{}\_\allowbreak{}o\allowbreak{}f\allowbreak{}\_\allowbreak{}c\allowbreak{}u\allowbreak{}b\allowbreak{}e} under the directly proved geometry gives the same physical probe and signal, with Gram exponent \(\fitinline{t_K-k_{\rm phys}}\). For a contrary \(\fitinline{\beta>t_K}\), choose \(\fitinline{\omega<(\beta-t_K)/4}\) and \(\fitinline{0<m<\min(\beta-t_K-\omega,\sigma,\beta-.87)}\). The exact H47 signal then has a holomorphic Mellin transform on \(\fitinline{\Re s>\beta-m}\), while on the absolute line it equals the original Gaussian times \(\fitinline{H_\eta/L_\eta}\). Since \(\fitinline{|H_\eta-1|\le1/2}\) on \(\fitinline{\Re s\ge.87}\), a target zero above \(\fitinline{\beta-m/2}\) contradicts holomorphy. The supremum need not be attained. Original norm pullback and fixed Euler restoration give K1 for every ordinary character, including imprimitive presentations.

\hypertarget{k.6.-finite-reconstruction-and-source-scope}{%
\subsection{K.6. Finite reconstruction and source scope}\label{k.6.-finite-reconstruction-and-source-scope}}

The physical input is \texttt{r\allowbreak{}e\allowbreak{}s\allowbreak{}e\allowbreak{}a\allowbreak{}r\allowbreak{}c\allowbreak{}h\allowbreak{}/\allowbreak{}l\allowbreak{}o\allowbreak{}c\allowbreak{}a\allowbreak{}l\allowbreak{}\_\allowbreak{}c\allowbreak{}y\allowbreak{}c\allowbreak{}l\allowbreak{}e\allowbreak{}0\allowbreak{}8\allowbreak{}3\allowbreak{}/\allowbreak{}n\allowbreak{}e\allowbreak{}a\allowbreak{}r\allowbreak{}\_\allowbreak{}r\allowbreak{}e\allowbreak{}c\allowbreak{}o\allowbreak{}r\allowbreak{}d\allowbreak{}/\allowbreak{}T\allowbreak{}I\allowbreak{}C\allowbreak{}K\allowbreak{}E\allowbreak{}T\allowbreak{}\_\allowbreak{}1\allowbreak{}2\allowbreak{}6\allowbreak{}6\allowbreak{}5\allowbreak{}8\allowbreak{}.\allowbreak{}j\allowbreak{}s\allowbreak{}o\allowbreak{}n}, SHA-256 \texttt{f\allowbreak{}7\allowbreak{}2\allowbreak{}2\allowbreak{}a\allowbreak{}4\allowbreak{}a\allowbreak{}b\allowbreak{}d\allowbreak{}d\allowbreak{}e\allowbreak{}8\allowbreak{}b\allowbreak{}9\allowbreak{}2\allowbreak{}b\allowbreak{}1\allowbreak{}9\allowbreak{}6\allowbreak{}0\allowbreak{}9\allowbreak{}0\allowbreak{}7\allowbreak{}8\allowbreak{}d\allowbreak{}7\allowbreak{}8\allowbreak{}a\allowbreak{}8\allowbreak{}9\allowbreak{}d\allowbreak{}9\allowbreak{}c\allowbreak{}8\allowbreak{}b\allowbreak{}1\allowbreak{}d\allowbreak{}5\allowbreak{}7\allowbreak{}1\allowbreak{}6\allowbreak{}0\allowbreak{}0\allowbreak{}f\allowbreak{}6\allowbreak{}4\allowbreak{}c\allowbreak{}f\allowbreak{}2\allowbreak{}a\allowbreak{}5\allowbreak{}f\allowbreak{}2\allowbreak{}0\allowbreak{}0\allowbreak{}b\allowbreak{}e\allowbreak{}b\allowbreak{}1\allowbreak{}6\allowbreak{}6\allowbreak{}5\allowbreak{}a\allowbreak{}6}, generated by \texttt{C\allowbreak{}H\allowbreak{}E\allowbreak{}C\allowbreak{}K\allowbreak{}.\allowbreak{}p\allowbreak{}y} (\texttt{1\allowbreak{}6\allowbreak{}e\allowbreak{}2\allowbreak{}8\allowbreak{}c\allowbreak{}4\allowbreak{}c\allowbreak{}0\allowbreak{}6\allowbreak{}f\allowbreak{}2\allowbreak{}c\allowbreak{}2\allowbreak{}5\allowbreak{}4\allowbreak{}6\allowbreak{}7\allowbreak{}6\allowbreak{}b\allowbreak{}d\allowbreak{}2\allowbreak{}c\allowbreak{}2\allowbreak{}0\allowbreak{}5\allowbreak{}f\allowbreak{}0\allowbreak{}6\allowbreak{}b\allowbreak{}1\allowbreak{}8\allowbreak{}2\allowbreak{}d\allowbreak{}4\allowbreak{}6\allowbreak{}8\allowbreak{}c\allowbreak{}c\allowbreak{}1\allowbreak{}2\allowbreak{}7\allowbreak{}b\allowbreak{}3\allowbreak{}d\allowbreak{}e\allowbreak{}4\allowbreak{}f\allowbreak{}1\allowbreak{}2\allowbreak{}8\allowbreak{}b\allowbreak{}2\allowbreak{}a\allowbreak{}1\allowbreak{}2\allowbreak{}2\allowbreak{}9\allowbreak{}3\allowbreak{}4\allowbreak{}b\allowbreak{}f\allowbreak{}d\allowbreak{}8}). Its compact source profile is \texttt{W\allowbreak{}I\allowbreak{}T\allowbreak{}N\allowbreak{}E\allowbreak{}S\allowbreak{}S\allowbreak{}\_\allowbreak{}W\allowbreak{}I\allowbreak{}N\allowbreak{}D\allowbreak{}O\allowbreak{}W\allowbreak{}S\allowbreak{}\_\allowbreak{}1\allowbreak{}2\allowbreak{}6\allowbreak{}6\allowbreak{}5\allowbreak{}8\allowbreak{}.\allowbreak{}j\allowbreak{}s\allowbreak{}o\allowbreak{}n} (\texttt{e\allowbreak{}b\allowbreak{}d\allowbreak{}1\allowbreak{}8\allowbreak{}2\allowbreak{}5\allowbreak{}1\allowbreak{}9\allowbreak{}6\allowbreak{}d\allowbreak{}a\allowbreak{}3\allowbreak{}d\allowbreak{}1\allowbreak{}7\allowbreak{}1\allowbreak{}f\allowbreak{}6\allowbreak{}0\allowbreak{}9\allowbreak{}9\allowbreak{}b\allowbreak{}6\allowbreak{}f\allowbreak{}3\allowbreak{}a\allowbreak{}2\allowbreak{}b\allowbreak{}4\allowbreak{}d\allowbreak{}4\allowbreak{}e\allowbreak{}b\allowbreak{}f\allowbreak{}c\allowbreak{}d\allowbreak{}b\allowbreak{}2\allowbreak{}0\allowbreak{}5\allowbreak{}6\allowbreak{}b\allowbreak{}7\allowbreak{}0\allowbreak{}6\allowbreak{}0\allowbreak{}3\allowbreak{}c\allowbreak{}8\allowbreak{}e\allowbreak{}f\allowbreak{}4\allowbreak{}f\allowbreak{}d\allowbreak{}c\allowbreak{}3\allowbreak{}3\allowbreak{}9\allowbreak{}9\allowbreak{}4\allowbreak{}1\allowbreak{}3\allowbreak{}e}). The physical checker reconstructs all 57 closed original edge tilings, all forty new admission guards, the E49/saved routes and every central price. The profile takes the stated closed intersections of its reconstructed cells, with no projection dropping a requested numerator or length.

\texttt{r\allowbreak{}e\allowbreak{}s\allowbreak{}e\allowbreak{}a\allowbreak{}r\allowbreak{}c\allowbreak{}h\allowbreak{}/\allowbreak{}l\allowbreak{}o\allowbreak{}c\allowbreak{}a\allowbreak{}l\allowbreak{}\_\allowbreak{}c\allowbreak{}y\allowbreak{}c\allowbreak{}l\allowbreak{}e\allowbreak{}0\allowbreak{}8\allowbreak{}3\allowbreak{}/\allowbreak{}n\allowbreak{}e\allowbreak{}a\allowbreak{}r\allowbreak{}\_\allowbreak{}r\allowbreak{}e\allowbreak{}c\allowbreak{}o\allowbreak{}r\allowbreak{}d\allowbreak{}/\allowbreak{}B\allowbreak{}I\allowbreak{}N\allowbreak{}D\allowbreak{}\_\allowbreak{}1\allowbreak{}2\allowbreak{}6\allowbreak{}6\allowbreak{}5\allowbreak{}8\allowbreak{}.\allowbreak{}p\allowbreak{}y} (\texttt{2\allowbreak{}c\allowbreak{}5\allowbreak{}7\allowbreak{}b\allowbreak{}2\allowbreak{}0\allowbreak{}a\allowbreak{}a\allowbreak{}6\allowbreak{}7\allowbreak{}d\allowbreak{}9\allowbreak{}6\allowbreak{}5\allowbreak{}8\allowbreak{}9\allowbreak{}c\allowbreak{}6\allowbreak{}8\allowbreak{}c\allowbreak{}7\allowbreak{}0\allowbreak{}e\allowbreak{}9\allowbreak{}7\allowbreak{}a\allowbreak{}6\allowbreak{}f\allowbreak{}e\allowbreak{}c\allowbreak{}f\allowbreak{}6\allowbreak{}6\allowbreak{}b\allowbreak{}0\allowbreak{}e\allowbreak{}2\allowbreak{}5\allowbreak{}c\allowbreak{}f\allowbreak{}8\allowbreak{}9\allowbreak{}1\allowbreak{}a\allowbreak{}6\allowbreak{}5\allowbreak{}3\allowbreak{}3\allowbreak{}e\allowbreak{}8\allowbreak{}a\allowbreak{}7\allowbreak{}1\allowbreak{}e\allowbreak{}d\allowbreak{}1\allowbreak{}d\allowbreak{}9\allowbreak{}1\allowbreak{}f\allowbreak{}3\allowbreak{}e\allowbreak{}9}) reproduces that scalar object and checks the 40 source entries against it, including each endpoint, required width and numerator enclosure. Its final receipt \texttt{F\allowbreak{}I\allowbreak{}N\allowbreak{}A\allowbreak{}L\allowbreak{}\_\allowbreak{}M\allowbreak{}A\allowbreak{}T\allowbreak{}C\allowbreak{}H\allowbreak{}\_\allowbreak{}1\allowbreak{}2\allowbreak{}6\allowbreak{}6\allowbreak{}5\allowbreak{}8\allowbreak{}.\allowbreak{}j\allowbreak{}s\allowbreak{}o\allowbreak{}n} (\texttt{1\allowbreak{}c\allowbreak{}2\allowbreak{}8\allowbreak{}2\allowbreak{}8\allowbreak{}b\allowbreak{}3\allowbreak{}6\allowbreak{}4\allowbreak{}b\allowbreak{}3\allowbreak{}8\allowbreak{}5\allowbreak{}9\allowbreak{}c\allowbreak{}4\allowbreak{}6\allowbreak{}7\allowbreak{}6\allowbreak{}0\allowbreak{}3\allowbreak{}0\allowbreak{}0\allowbreak{}7\allowbreak{}b\allowbreak{}8\allowbreak{}8\allowbreak{}3\allowbreak{}0\allowbreak{}3\allowbreak{}4\allowbreak{}1\allowbreak{}4\allowbreak{}2\allowbreak{}5\allowbreak{}e\allowbreak{}a\allowbreak{}4\allowbreak{}e\allowbreak{}2\allowbreak{}e\allowbreak{}3\allowbreak{}7\allowbreak{}b\allowbreak{}d\allowbreak{}a\allowbreak{}8\allowbreak{}4\allowbreak{}6\allowbreak{}4\allowbreak{}9\allowbreak{}5\allowbreak{}a\allowbreak{}5\allowbreak{}f\allowbreak{}7\allowbreak{}3\allowbreak{}4\allowbreak{}7\allowbreak{}b\allowbreak{}1\allowbreak{}d\allowbreak{}b}, 1237 bytes) records the exact source kinds and minimum spare width, bound to the released direct certificate below. It supports \texttt{-\/-source-root\ \textless{}unpacked-proof-root\textgreater{}} containing \texttt{r\allowbreak{}e\allowbreak{}s\allowbreak{}e\allowbreak{}a\allowbreak{}r\allowbreak{}c\allowbreak{}h\allowbreak{}/}; its immutable core/cover/contact inputs are the same three pinned files listed in J.4. Retained local and actual portable-root reproductions agree. The earlier \texttt{M\allowbreak{}A\allowbreak{}T\allowbreak{}C\allowbreak{}H\allowbreak{}\_\allowbreak{}1\allowbreak{}2\allowbreak{}6\allowbreak{}6\allowbreak{}5\allowbreak{}8\allowbreak{}.\allowbreak{}j\allowbreak{}s\allowbreak{}o\allowbreak{}n} is retained historical evidence and is not the final source binding. Numerical reconstruction verifies arithmetic and source matching; the analytic applicability, all-radius slope proof and physical return are the arguments above.

The analytic physical source is \texttt{r\allowbreak{}e\allowbreak{}s\allowbreak{}e\allowbreak{}a\allowbreak{}r\allowbreak{}c\allowbreak{}h\allowbreak{}/\allowbreak{}l\allowbreak{}o\allowbreak{}c\allowbreak{}a\allowbreak{}l\allowbreak{}\_\allowbreak{}c\allowbreak{}y\allowbreak{}c\allowbreak{}l\allowbreak{}e\allowbreak{}0\allowbreak{}8\allowbreak{}3\allowbreak{}/\allowbreak{}n\allowbreak{}e\allowbreak{}a\allowbreak{}r\allowbreak{}\_\allowbreak{}r\allowbreak{}e\allowbreak{}c\allowbreak{}o\allowbreak{}r\allowbreak{}d\allowbreak{}/\allowbreak{}R\allowbreak{}E\allowbreak{}S\allowbreak{}E\allowbreak{}A\allowbreak{}R\allowbreak{}C\allowbreak{}H\allowbreak{}.\allowbreak{}m\allowbreak{}d} (\texttt{9\allowbreak{}d\allowbreak{}2\allowbreak{}5\allowbreak{}a\allowbreak{}3\allowbreak{}e\allowbreak{}a\allowbreak{}5\allowbreak{}9\allowbreak{}8\allowbreak{}2\allowbreak{}f\allowbreak{}8\allowbreak{}9\allowbreak{}1\allowbreak{}5\allowbreak{}d\allowbreak{}b\allowbreak{}1\allowbreak{}c\allowbreak{}1\allowbreak{}b\allowbreak{}1\allowbreak{}0\allowbreak{}8\allowbreak{}a\allowbreak{}e\allowbreak{}b\allowbreak{}a\allowbreak{}2\allowbreak{}5\allowbreak{}6\allowbreak{}1\allowbreak{}6\allowbreak{}1\allowbreak{}4\allowbreak{}6\allowbreak{}e\allowbreak{}b\allowbreak{}8\allowbreak{}9\allowbreak{}a\allowbreak{}8\allowbreak{}f\allowbreak{}1\allowbreak{}c\allowbreak{}6\allowbreak{}8\allowbreak{}f\allowbreak{}4\allowbreak{}d\allowbreak{}7\allowbreak{}f\allowbreak{}2\allowbreak{}2\allowbreak{}2\allowbreak{}5\allowbreak{}a\allowbreak{}7\allowbreak{}0\allowbreak{}2\allowbreak{}9\allowbreak{}d}). The pre-division fixed-ray and inner-source transfer is controller083 K1--K3 in \texttt{r\allowbreak{}e\allowbreak{}s\allowbreak{}e\allowbreak{}a\allowbreak{}r\allowbreak{}c\allowbreak{}h\allowbreak{}/\allowbreak{}l\allowbreak{}o\allowbreak{}c\allowbreak{}a\allowbreak{}l\allowbreak{}\_\allowbreak{}c\allowbreak{}y\allowbreak{}c\allowbreak{}l\allowbreak{}e\allowbreak{}0\allowbreak{}8\allowbreak{}3\allowbreak{}/\allowbreak{}c\allowbreak{}o\allowbreak{}n\allowbreak{}t\allowbreak{}r\allowbreak{}o\allowbreak{}l\allowbreak{}l\allowbreak{}e\allowbreak{}r\allowbreak{}/\allowbreak{}R\allowbreak{}E\allowbreak{}S\allowbreak{}E\allowbreak{}A\allowbreak{}R\allowbreak{}C\allowbreak{}H\allowbreak{}.\allowbreak{}m\allowbreak{}d} (\texttt{c\allowbreak{}e\allowbreak{}7\allowbreak{}2\allowbreak{}d\allowbreak{}6\allowbreak{}5\allowbreak{}7\allowbreak{}5\allowbreak{}9\allowbreak{}7\allowbreak{}9\allowbreak{}3\allowbreak{}0\allowbreak{}d\allowbreak{}2\allowbreak{}d\allowbreak{}7\allowbreak{}3\allowbreak{}a\allowbreak{}9\allowbreak{}a\allowbreak{}8\allowbreak{}5\allowbreak{}6\allowbreak{}e\allowbreak{}3\allowbreak{}4\allowbreak{}4\allowbreak{}9\allowbreak{}3\allowbreak{}c\allowbreak{}d\allowbreak{}f\allowbreak{}d\allowbreak{}5\allowbreak{}5\allowbreak{}e\allowbreak{}1\allowbreak{}f\allowbreak{}c\allowbreak{}c\allowbreak{}f\allowbreak{}b\allowbreak{}e\allowbreak{}7\allowbreak{}d\allowbreak{}4\allowbreak{}f\allowbreak{}7\allowbreak{}9\allowbreak{}6\allowbreak{}8\allowbreak{}f\allowbreak{}8\allowbreak{}8\allowbreak{}a\allowbreak{}c\allowbreak{}8\allowbreak{}5\allowbreak{}1\allowbreak{}b\allowbreak{}4\allowbreak{}a}), using only those comparisons and its earlier100/112 enlargement. Its ancillary endpoint verification is \texttt{C\allowbreak{}O\allowbreak{}R\allowbreak{}E\allowbreak{}\_\allowbreak{}B\allowbreak{}A\allowbreak{}S\allowbreak{}E\allowbreak{}\_\allowbreak{}C\allowbreak{}H\allowbreak{}E\allowbreak{}C\allowbreak{}K\allowbreak{}.\allowbreak{}p\allowbreak{}y} (\texttt{3\allowbreak{}1\allowbreak{}4\allowbreak{}b\allowbreak{}3\allowbreak{}9\allowbreak{}2\allowbreak{}9\allowbreak{}f\allowbreak{}c\allowbreak{}a\allowbreak{}9\allowbreak{}4\allowbreak{}e\allowbreak{}8\allowbreak{}6\allowbreak{}d\allowbreak{}2\allowbreak{}0\allowbreak{}d\allowbreak{}3\allowbreak{}2\allowbreak{}0\allowbreak{}c\allowbreak{}8\allowbreak{}1\allowbreak{}f\allowbreak{}e\allowbreak{}d\allowbreak{}f\allowbreak{}6\allowbreak{}8\allowbreak{}d\allowbreak{}5\allowbreak{}f\allowbreak{}1\allowbreak{}a\allowbreak{}6\allowbreak{}2\allowbreak{}9\allowbreak{}3\allowbreak{}9\allowbreak{}d\allowbreak{}5\allowbreak{}c\allowbreak{}0\allowbreak{}8\allowbreak{}7\allowbreak{}7\allowbreak{}5\allowbreak{}7\allowbreak{}d\allowbreak{}0\allowbreak{}8\allowbreak{}9\allowbreak{}6\allowbreak{}6\allowbreak{}b\allowbreak{}4\allowbreak{}7\allowbreak{}b\allowbreak{}6\allowbreak{}8\allowbreak{}a}) and \texttt{C\allowbreak{}O\allowbreak{}R\allowbreak{}E\allowbreak{}\_\allowbreak{}B\allowbreak{}A\allowbreak{}S\allowbreak{}E\allowbreak{}.\allowbreak{}j\allowbreak{}s\allowbreak{}o\allowbreak{}n} (\texttt{3\allowbreak{}a\allowbreak{}a\allowbreak{}8\allowbreak{}8\allowbreak{}1\allowbreak{}6\allowbreak{}b\allowbreak{}9\allowbreak{}8\allowbreak{}2\allowbreak{}2\allowbreak{}f\allowbreak{}a\allowbreak{}d\allowbreak{}f\allowbreak{}f\allowbreak{}2\allowbreak{}2\allowbreak{}5\allowbreak{}8\allowbreak{}9\allowbreak{}1\allowbreak{}e\allowbreak{}0\allowbreak{}2\allowbreak{}6\allowbreak{}f\allowbreak{}5\allowbreak{}f\allowbreak{}9\allowbreak{}2\allowbreak{}2\allowbreak{}8\allowbreak{}e\allowbreak{}6\allowbreak{}a\allowbreak{}4\allowbreak{}9\allowbreak{}c\allowbreak{}c\allowbreak{}b\allowbreak{}0\allowbreak{}0\allowbreak{}9\allowbreak{}9\allowbreak{}e\allowbreak{}8\allowbreak{}1\allowbreak{}2\allowbreak{}c\allowbreak{}6\allowbreak{}a\allowbreak{}7\allowbreak{}4\allowbreak{}e\allowbreak{}e\allowbreak{}d\allowbreak{}7\allowbreak{}7\allowbreak{}f\allowbreak{}9\allowbreak{}7\allowbreak{}5}), in that same directory.

The issued source is \texttt{r\allowbreak{}e\allowbreak{}s\allowbreak{}e\allowbreak{}a\allowbreak{}r\allowbreak{}c\allowbreak{}h\allowbreak{}/\allowbreak{}l\allowbreak{}o\allowbreak{}c\allowbreak{}a\allowbreak{}l\allowbreak{}\_\allowbreak{}c\allowbreak{}y\allowbreak{}c\allowbreak{}l\allowbreak{}e\allowbreak{}0\allowbreak{}8\allowbreak{}3\allowbreak{}/\allowbreak{}d\allowbreak{}i\allowbreak{}r\allowbreak{}e\allowbreak{}c\allowbreak{}t\allowbreak{}/\allowbreak{}R\allowbreak{}E\allowbreak{}S\allowbreak{}E\allowbreak{}A\allowbreak{}R\allowbreak{}C\allowbreak{}H\allowbreak{}.\allowbreak{}m\allowbreak{}d} (\texttt{f\allowbreak{}b\allowbreak{}e\allowbreak{}a\allowbreak{}6\allowbreak{}e\allowbreak{}7\allowbreak{}c\allowbreak{}a\allowbreak{}7\allowbreak{}d\allowbreak{}a\allowbreak{}9\allowbreak{}7\allowbreak{}e\allowbreak{}6\allowbreak{}1\allowbreak{}4\allowbreak{}0\allowbreak{}9\allowbreak{}4\allowbreak{}6\allowbreak{}d\allowbreak{}c\allowbreak{}d\allowbreak{}b\allowbreak{}6\allowbreak{}6\allowbreak{}a\allowbreak{}d\allowbreak{}9\allowbreak{}8\allowbreak{}f\allowbreak{}6\allowbreak{}8\allowbreak{}6\allowbreak{}1\allowbreak{}9\allowbreak{}f\allowbreak{}9\allowbreak{}0\allowbreak{}a\allowbreak{}9\allowbreak{}1\allowbreak{}5\allowbreak{}e\allowbreak{}5\allowbreak{}e\allowbreak{}a\allowbreak{}0\allowbreak{}5\allowbreak{}2\allowbreak{}3\allowbreak{}d\allowbreak{}4\allowbreak{}2\allowbreak{}f\allowbreak{}0\allowbreak{}4\allowbreak{}c\allowbreak{}e\allowbreak{}1\allowbreak{}4\allowbreak{}f}), \texttt{C\allowbreak{}H\allowbreak{}E\allowbreak{}C\allowbreak{}K\allowbreak{}.\allowbreak{}p\allowbreak{}y} (\texttt{3\allowbreak{}d\allowbreak{}0\allowbreak{}2\allowbreak{}b\allowbreak{}8\allowbreak{}2\allowbreak{}3\allowbreak{}6\allowbreak{}0\allowbreak{}c\allowbreak{}0\allowbreak{}8\allowbreak{}3\allowbreak{}3\allowbreak{}d\allowbreak{}f\allowbreak{}9\allowbreak{}b\allowbreak{}9\allowbreak{}b\allowbreak{}2\allowbreak{}8\allowbreak{}0\allowbreak{}e\allowbreak{}3\allowbreak{}4\allowbreak{}f\allowbreak{}0\allowbreak{}4\allowbreak{}2\allowbreak{}3\allowbreak{}1\allowbreak{}b\allowbreak{}4\allowbreak{}f\allowbreak{}e\allowbreak{}9\allowbreak{}1\allowbreak{}6\allowbreak{}5\allowbreak{}2\allowbreak{}8\allowbreak{}c\allowbreak{}d\allowbreak{}c\allowbreak{}5\allowbreak{}0\allowbreak{}9\allowbreak{}c\allowbreak{}0\allowbreak{}4\allowbreak{}6\allowbreak{}5\allowbreak{}4\allowbreak{}f\allowbreak{}7\allowbreak{}3\allowbreak{}8\allowbreak{}0\allowbreak{}0\allowbreak{}d\allowbreak{}f\allowbreak{}c}) and \texttt{C\allowbreak{}E\allowbreak{}R\allowbreak{}T\allowbreak{}I\allowbreak{}F\allowbreak{}I\allowbreak{}C\allowbreak{}A\allowbreak{}T\allowbreak{}E\allowbreak{}.\allowbreak{}j\allowbreak{}s\allowbreak{}o\allowbreak{}n} (\texttt{8\allowbreak{}e\allowbreak{}a\allowbreak{}1\allowbreak{}7\allowbreak{}b\allowbreak{}2\allowbreak{}5\allowbreak{}8\allowbreak{}3\allowbreak{}b\allowbreak{}e\allowbreak{}3\allowbreak{}a\allowbreak{}5\allowbreak{}1\allowbreak{}1\allowbreak{}5\allowbreak{}1\allowbreak{}e\allowbreak{}d\allowbreak{}0\allowbreak{}3\allowbreak{}e\allowbreak{}7\allowbreak{}b\allowbreak{}b\allowbreak{}4\allowbreak{}6\allowbreak{}6\allowbreak{}f\allowbreak{}6\allowbreak{}3\allowbreak{}2\allowbreak{}7\allowbreak{}b\allowbreak{}c\allowbreak{}9\allowbreak{}4\allowbreak{}6\allowbreak{}8\allowbreak{}5\allowbreak{}3\allowbreak{}9\allowbreak{}3\allowbreak{}3\allowbreak{}7\allowbreak{}c\allowbreak{}c\allowbreak{}1\allowbreak{}b\allowbreak{}a\allowbreak{}1\allowbreak{}c\allowbreak{}5\allowbreak{}c\allowbreak{}2\allowbreak{}e\allowbreak{}e\allowbreak{}f\allowbreak{}9\allowbreak{}4\allowbreak{}d\allowbreak{}d}). Replay the direct checker with \texttt{-\allowbreak{}\/\allowbreak{}-\allowbreak{}p\allowbreak{}r\allowbreak{}o\allowbreak{}f\allowbreak{}i\allowbreak{}l\allowbreak{}e} naming the packed witness-window input and \texttt{-\allowbreak{}\/\allowbreak{}-\allowbreak{}o\allowbreak{}u\allowbreak{}t\allowbreak{}p\allowbreak{}u\allowbreak{}t} naming a fresh file. It checks that exact input hash and reconstructs every band, threshold, closed triangle, entry test and ancillary gate. Its actual mathematical derivation is K.3--K.4, not its status label. The original I100 and J112 certificates remain the earlier scientific inputs; the immediate Appendix J has SHA-256 \texttt{a\allowbreak{}e\allowbreak{}e\allowbreak{}8\allowbreak{}8\allowbreak{}a\allowbreak{}0\allowbreak{}6\allowbreak{}4\allowbreak{}1\allowbreak{}a\allowbreak{}b\allowbreak{}9\allowbreak{}9\allowbreak{}0\allowbreak{}9\allowbreak{}b\allowbreak{}e\allowbreak{}7\allowbreak{}c\allowbreak{}7\allowbreak{}1\allowbreak{}2\allowbreak{}8\allowbreak{}6\allowbreak{}4\allowbreak{}5\allowbreak{}b\allowbreak{}7\allowbreak{}7\allowbreak{}1\allowbreak{}3\allowbreak{}a\allowbreak{}d\allowbreak{}f\allowbreak{}8\allowbreak{}5\allowbreak{}4\allowbreak{}e\allowbreak{}5\allowbreak{}e\allowbreak{}c\allowbreak{}f\allowbreak{}c\allowbreak{}5\allowbreak{}2\allowbreak{}7\allowbreak{}1\allowbreak{}e\allowbreak{}2\allowbreak{}f\allowbreak{}7\allowbreak{}5\allowbreak{}7\allowbreak{}c\allowbreak{}3\allowbreak{}a\allowbreak{}6\allowbreak{}1\allowbreak{}0\allowbreak{}3\allowbreak{}9\allowbreak{}e\allowbreak{}b}, on the complete I input \texttt{c\allowbreak{}0\allowbreak{}0\allowbreak{}2\allowbreak{}5\allowbreak{}3\allowbreak{}c\allowbreak{}4\allowbreak{}f\allowbreak{}7\allowbreak{}7\allowbreak{}1\allowbreak{}3\allowbreak{}e\allowbreak{}d\allowbreak{}c\allowbreak{}1\allowbreak{}9\allowbreak{}7\allowbreak{}b\allowbreak{}3\allowbreak{}1\allowbreak{}b\allowbreak{}4\allowbreak{}b\allowbreak{}f\allowbreak{}2\allowbreak{}8\allowbreak{}e\allowbreak{}4\allowbreak{}a\allowbreak{}a\allowbreak{}1\allowbreak{}d\allowbreak{}0\allowbreak{}8\allowbreak{}d\allowbreak{}4\allowbreak{}d\allowbreak{}5\allowbreak{}a\allowbreak{}7\allowbreak{}5\allowbreak{}5\allowbreak{}f\allowbreak{}8\allowbreak{}6\allowbreak{}b\allowbreak{}d\allowbreak{}b\allowbreak{}8\allowbreak{}d\allowbreak{}d\allowbreak{}a\allowbreak{}0\allowbreak{}a\allowbreak{}0\allowbreak{}7\allowbreak{}8\allowbreak{}a\allowbreak{}9\allowbreak{}9\allowbreak{}1\allowbreak{}8}. The separate wider 106-micro-unit bank and larger multi-cutoff requests are outside this chain.

The retained exact grouped interface is \texttt{r\allowbreak{}e\allowbreak{}s\allowbreak{}e\allowbreak{}a\allowbreak{}r\allowbreak{}c\allowbreak{}h\allowbreak{}/\allowbreak{}l\allowbreak{}o\allowbreak{}c\allowbreak{}a\allowbreak{}l\allowbreak{}\_\allowbreak{}c\allowbreak{}y\allowbreak{}c\allowbreak{}l\allowbreak{}e\allowbreak{}0\allowbreak{}7\allowbreak{}8\allowbreak{}/\allowbreak{}n\allowbreak{}e\allowbreak{}a\allowbreak{}r\allowbreak{}\_\allowbreak{}r\allowbreak{}e\allowbreak{}c\allowbreak{}o\allowbreak{}r\allowbreak{}d\allowbreak{}/\allowbreak{}p\allowbreak{}r\allowbreak{}i\allowbreak{}m\allowbreak{}a\allowbreak{}r\allowbreak{}y\allowbreak{}/\allowbreak{}l\allowbreak{}e\allowbreak{}a\allowbreak{}n\allowbreak{}/\allowbreak{}A\allowbreak{}s\allowbreak{}y\allowbreak{}m\allowbreak{}m\allowbreak{}e\allowbreak{}t\allowbreak{}r\allowbreak{}i\allowbreak{}c\allowbreak{}R\allowbreak{}a\allowbreak{}y\allowbreak{}G\allowbreak{}r\allowbreak{}o\allowbreak{}u\allowbreak{}p\allowbreak{}e\allowbreak{}d\allowbreak{}S\allowbreak{}i\allowbreak{}d\allowbreak{}e\allowbreak{}.\allowbreak{}l\allowbreak{}e\allowbreak{}a\allowbreak{}n} (\texttt{2\allowbreak{}a\allowbreak{}4\allowbreak{}6\allowbreak{}a\allowbreak{}d\allowbreak{}8\allowbreak{}a\allowbreak{}f\allowbreak{}0\allowbreak{}b\allowbreak{}f\allowbreak{}e\allowbreak{}f\allowbreak{}b\allowbreak{}2\allowbreak{}9\allowbreak{}2\allowbreak{}6\allowbreak{}5\allowbreak{}b\allowbreak{}b\allowbreak{}a\allowbreak{}e\allowbreak{}f\allowbreak{}5\allowbreak{}6\allowbreak{}f\allowbreak{}5\allowbreak{}a\allowbreak{}2\allowbreak{}9\allowbreak{}3\allowbreak{}3\allowbreak{}6\allowbreak{}6\allowbreak{}c\allowbreak{}a\allowbreak{}f\allowbreak{}b\allowbreak{}5\allowbreak{}8\allowbreak{}7\allowbreak{}a\allowbreak{}c\allowbreak{}f\allowbreak{}1\allowbreak{}8\allowbreak{}2\allowbreak{}4\allowbreak{}1\allowbreak{}9\allowbreak{}3\allowbreak{}c\allowbreak{}c\allowbreak{}5\allowbreak{}6\allowbreak{}5\allowbreak{}c\allowbreak{}a\allowbreak{}c\allowbreak{}b\allowbreak{}2\allowbreak{}e}), at its pinned PR43 source. Its \texttt{r\allowbreak{}a\allowbreak{}y\allowbreak{}C\allowbreak{}h\allowbreak{}i\allowbreak{}l\allowbreak{}d} and \texttt{S\allowbreak{}h\allowbreak{}a\allowbreak{}p\allowbreak{}e\allowbreak{}L\allowbreak{}e\allowbreak{}a\allowbreak{}f} keep two separate fixed multipliers, one moving feature, all subproduct/profile variables and the actual common-ball radius. Its labelwise reduction and leaf-energy assumptions are supplied by the proofs above; the interface alone does not discharge them.

The complete proof since OpenAI is retained, including properly derived alternatives. Credits and original terms remain with OpenAI, Sil's canonical marked sources, Joseph Bebel's contact/count work, Akash Levy's numerator, Hailey Collet's signed physical return and the other inherited authors. The SovereignSteak draft is used at its retained PR43 head \texttt{c\allowbreak{}8\allowbreak{}3\allowbreak{}8\allowbreak{}a\allowbreak{}7\allowbreak{}5\allowbreak{}0\allowbreak{}0\allowbreak{}7\allowbreak{}9\allowbreak{}f\allowbreak{}7\allowbreak{}f\allowbreak{}8\allowbreak{}6\allowbreak{}a\allowbreak{}5\allowbreak{}7\allowbreak{}3\allowbreak{}d\allowbreak{}f\allowbreak{}9\allowbreak{}c\allowbreak{}9\allowbreak{}f\allowbreak{}7\allowbreak{}2\allowbreak{}9\allowbreak{}b\allowbreak{}8\allowbreak{}b\allowbreak{}c\allowbreak{}9\allowbreak{}2\allowbreak{}d\allowbreak{}1\allowbreak{}2\allowbreak{}5\allowbreak{}e}, with its contributor notices and assistance disclosure; no proposed endpoint affirmation is a premise. Catalogue mathematical and Lean interfaces, including the actual coefficient and finite-family contracts, were consulted for scope. G137's conditional narrow simultaneous-witness entry is not promoted to an unconditional source. No finite receipt, declaration or import is an end-to-end formal proof of K1. The Sepia/CWR mathematical-writing attribution and terms remain as retained in H--J.

\hypertarget{appendix-l}{}

\hypertarget{appendix-l.-a-finite-quadratic-source-and-a-bank-of-actual-zero-witnesses}{%
\section{Appendix L. A finite quadratic source and a bank of actual zero witnesses}\label{appendix-l.-a-finite-quadratic-source-and-a-bank-of-actual-zero-witnesses}}

Retain the complete proof through K, including the adopted natural selected-cell input SN9, the original fixed presented modulus, all six units, imprimitive Euler holes and the actual physical denominator. We prove the ordinary all-height conclusion at \[\fitmath{
 c_L=\frac{6333207}{50000000}=.12666414,\qquad
 t_L=1-c_L=.87333586,\qquad h=h_*+\frac{27}{40000}<.815999505<.816.
}\tag{L0}\] The conclusion is strictly to the right of \(\fitinline{t_L}\), with an actual principal pole excepted. We prove the positive marked source and its complete physical comparison below. Both retain the independent ordinary heights, normalized first log tests, every main/error coefficient and all original prime pools. There is no growing-modulus uniformity, boundary nonvanishing or new end-to-end Lean verification assertion.

\hypertarget{l.1.-the-quadratic-source-and-its-finite-induction}{%
\subsection{L.1. The quadratic source and its finite induction}\label{l.1.-the-quadratic-source-and-its-finite-induction}}

Let \(\fitinline{g_0,g_1}\) be the exact contact constants of E--G, with \(\fitinline{.05805911<g_0<.05805913}\), \(\fitinline{.07599618<g_1<.07599620}\). Put \[\fitmath{
 \begin{gathered}
 A=(1+g_0)/2,\qquad B=(1+2g_1)/2,\qquad l=.38145,\\
 a=.00017,\quad c_2=11/2,\quad
 \beta(y)=a+c_2(y-l)^2,\quad W(n,y)=A-n/2-By+\beta(y),\\
 u=1/50000,\qquad H=1-6u=24997/25000,\qquad \rho=10^{-9}.
 \end{gathered}                                                
}\tag{L1}\] For every upper pair \[\fitmath{
             .373\le y\le.387,\qquad1.11\le n/y\le1.16,       
}\tag{L2}\] every strictly decremented width \(\fitinline{W(n,y)}\) is admitted by the actual two-ordinary marked energy \textbf{before lower-mark division}, on its original kept rows. The factors have two fixed finite rays multiplying the same moving sextic feature, as proved in K.3. Actual componentwise shorter factors use this upper pair and its width; their ratios need not satisfy L2. This is a full uncentered kept root after the two raw-side restorations, not an all-row estimate for either separate long signed arm.

Here is a finite proof. At and above \(\fitinline{y_*=.387}\), through \(\fitinline{y_*/H}\), use only the regular-core positive energy. The two native fourth quotas \(\fitinline{1/2-n}\), \(\fitinline{1/2-y}\) are positive and below their native cap on that strip. Their mean is \(\fitinline{(1-n-y)/2}\). Each fourth can amplify the actual attaining clean slot at its own relative ray up to its strictly decremented quota; H.4's upper comparison bounds the remaining genuine mark product. Equivalently allocate whole sufficiently small slots between the two repeated-label fourth comparisons. This is a positive regular-core statement with its actual keep, not a new bound on fixed inducing rows. The mean exceeds \(\fitinline{W(n,s)}\) by \[\fitmath{
                   g_1s-g_0/2-\beta(s).
}\] This is concave on the starting strip; its smaller endpoint is at least \(\fitinline{4154291/10^{11}>0}\), using the rational worst contact endpoints. Empty/tiny slots, punctures and rounding fit a strict decrement, as in K.3.

Set \(\fitinline{b_j=.387H^j}\). Prove the source successively on \(\fitinline{[\max(.373,b_{j+1}),b_j]}\). Exact rational powers give \(\fitinline{b_{307}>.373\ge b_{308}}\), so there are 308 stages. If \(\fitinline{y}\) is in the current interval, \(\fitinline{y/H}\) is in the preceding higher interval or the native starting strip, and \(\fitinline{(n/H)/(y/H)=n/y}\). At a shared endpoint choose its earlier proof. No point of a new interval is assumed in its own proof.

Fix one stage and an actual regular core \(\fitinline{z=b^6c}\), with radius \(\fitinline{U^V}\), before taking its positive norm. Keep the unit, both fixed rays, both outer heights and every hole. K.3's exact ordinary M\"{o}bius expansions and positive prime-hole expansions leave actual shortened coordinates and subpower scalar masses, without imposing \(\fitinline{(b,c)=1}\). Let \(\fitinline{e}\) be the clean maximizing \textbf{squared} prime-response exponent and \(\fitinline{L=W(n,y)-u+o(1)}\) the live width. Whenever an earlier upper pair \(\fitinline{(n',s)}\) dominates the actual normalized coordinates, its pre-division source gives the excess over \(\fitinline{U^V}\) \[\fitmath{
                          e[L-VW(n',s)]_+.                    
}\tag{L3}\] Choose a whole residual subproduct by width, with its strict decrement and one early-slot rounding fee; bound only the complement by its upper response on this same clean bin. There is no lower-mark hypothesis in L3. Scale identities are literal: \(\fitinline{(U^V)^{w_j/V}=U^{w_j}}\), with the actual prime profiles and frequencies.

At \(\fitinline{V=H}\) the earlier pair is \(\fitinline{(n/H,y/H)}\). The bracket in L3 is exactly \(\fitinline{uG(y)}\), apart from the small fees, where \[\fitmath{
 G(y)=2+3g_0-6\{c_2(y^2/H-l^2)-a\}.                          
}\tag{L4}\] Indeed \(\fitinline{H\beta(y/H)-\beta(y)=6u[c_2(y^2/H-l^2)-a]}\); the two linear length terms cancel first. On a length interval \(\fitinline{[a_1,b_1]}\) of width \(\fitinline{1/80000}\), take \[\fitmath{
 k=10^{-7}\left\lfloor
 \frac{10^7(1-2\rho/u)}
 {2+3(.05805913)-6\{c_2(a_1^2/H-l^2)-a\}}
 \right\rfloor .                                             
}\tag{L5}\] The positive denominator bounds \(\fitinline{G}\) above. Thus \(\fitinline{e\le k}\) costs at most \(\fitinline{u-2\rho}\). The finite range of these thresholds is \(\fitinline{[.4192409,.4915489]}\).

This low-response proof covers every radius. For \(\fitinline{V\ge H}\), keep the earlier upper pair fixed at \(\fitinline{(n/H,y/H)}\); its positive absolute quota \(\fitinline{VW(n/H,y/H)}\) increases, while the actual inputs shorten relative to that cap. Use the genuine componentwise-shortening statement at each actual input scale, with its own original normalization and H24a compact log-Fourier envelope. A constant upper cap is not one smooth bump evaluated at a largest fictitious scale. Bounded and null factors retain their original cases. Thus L3 does not increase. For \(\fitinline{V<H}\) while \(\fitinline{s=y/V\le.387}\), use the true earlier cap \(\fitinline{(n/V,y/V)}\). Its absolute quota derivative is \[\fitmath{
             A+\beta(s)-s\beta'(s)=A+a+c_2(l^2-s^2).           
}\tag{L6}\] This is greater than .50, and its upper bound times \(\fitinline{e\le.5002}\) is strictly below \(\fitinline{5/6}\), with more than .55 spare. The positive part preserves that Lipschitz bound. When \(\fitinline{s}\) reaches .387, switch to the native two-fourth mean at that endpoint, which is larger by the strictly positive starting margin. Below that seam use its actual mean, with any negative individual quota replaced by the unmarked face. On \(\fitinline{[.98,H]}\) both quotas are positive, their absolute mean has derivative \(\fitinline{1/2}\), and both remain below their native caps. Therefore the low excess is at most \(\fitinline{u-2\rho+5(H-V)_+/6}\). The quadratic profile is never extrapolated above its starting edge.

\hypertarget{l.2.-the-common-zero-part-all-remaining-radii-and-the-true-root}{%
\subsection{L.2. The common-zero part, all remaining radii and the true root}\label{l.2.-the-common-zero-part-all-remaining-radii-and-the-true-root}}

For \(\fitinline{k\le e\le.5002}\), use I.3's actual finite-family nonfloor zero band \(\fitinline{e\le\nu\le3/4}\). The threshold \(\fitinline{k>.419}\) excludes the floor. At its one attaining zero, a fixed cutoff \(\fitinline{3w/2}\), \(\fitinline{.75\le w\le.775}\), gives the two supported deficits \[\fitmath{
 D_s=\nu w+e[(1-w)/2+I(w)],\qquad
 D_l=5/6-(3/2)(5/6-\nu)w,                                   
}\tag{L7}\] where \(\fitinline{I}\) is E(19), explicitly I16. Its actual M\"{o}bius inverse profile, native/ordinary continuations and .7775 switch are exactly the original source contracts. We use the curve for these witness widths; no curve at \(\fitinline{2y/V}\) outside its domain is assumed. With the genuine longer ordinary fourth, the two energy excesses are \[\fitmath{
 C_i(V)=\nu y+e\{W(n,y)-u-(V/2-n)/2\}-VD_i/2,
                  \qquad i=s,l.                             
}\tag{L8}\] On \(\fitinline{[.98,1]}\) both decrease with \(\fitinline{V}\), with absolute slope below \(\fitinline{5/6}\). For \(\fitinline{V\ge1}\), use instead the genuinely available absolute long quota \(\fitinline{V(1/2-n)}\), by downward restriction at the original upper length \(\fitinline{n}\). Its test and normalization stay at the actual shortened scale, using H24a; only the source's upper label is held fixed. It agrees at one and increases. Keep the same \(\fitinline{w}\) and the actual relative zero argument. The short ordinary bound remains \(\fitinline{\nu y}\), since its absolute length is \(\fitinline{U^y}\); it does not acquire a factor \(\fitinline{V}\). Both expressions continue to decrease. This supplies the large-radius join with a legitimate cap.

Here is the complete finite calculation at \(\fitinline{H}\). The first coordinate cancels. With worst contact endpoints put \[\fitmath{
 P_b=(1-H)/4+[1/2+.05805913-(1+2(.07599618))b_1]/2
                            +\beta(b_1)-u.
}\] The expression \(\fitinline{\nu y+eP(y)}\) increases in \(\fitinline{y}\): its derivative is positive because \(\fitinline{\nu\ge e}\) and \(\fitinline{1-(1+2g_1)/2+\min\beta'>.33}\). Hence it is at most \(\fitinline{\nu b_1+eP_b}\). Define \(\fitinline{D(\nu,e)=2(\nu b_1+eP_b-u+\rho)/H}\). Cover \(\fitinline{k\le e\le.5002,\ e\le\nu\le.75}\) by the two closed triangles with vertices \(\fitinline{(k,k),(.5002,.5002),(.75,k)}\) and \(\fitinline{(.5002,.5002),(.75,.5002),(.75,k)}\), in coordinates \(\fitinline{(\nu,e)}\). On each take the floor to \(\fitinline{10^{-9}}\) of the minimum of .775 and \[\fitmath{
                   \frac{5/6-D(\nu,e)}{(3/2)(5/6-\nu)}
}\] over its three vertices. The resulting \(\fitinline{w}\) is in \(\fitinline{[.75,.775]}\). Bound \(\fitinline{(5/16)[7w^4(w-2/3)]^{1/5}}\) upward at denominator \(\fitinline{10^{13}}\) by integer fifth-power tests; subtract it from \(\fitinline{11w/16-7/24}\) to minorize \(\fitinline{I(w)}\). With this fixed cutoff both deficits and \(\fitinline{D}\) are affine on the entire triangle, so the vertex tests prove it.

There are 1120 length intervals and 2240 triangles, all passing without subdivision. The minimum low, short and long values of \(\fitinline{u-C}\) are respectively \[\fitmath{
 \frac{2499710571340103}{1249850000000000000000000},\qquad
 \frac{19827188214753383021}{750000000000000000000000},\qquad
 \frac{750200839421}{750000000000000000000},                    
}\tag{L9}\] the first greater than \(\fitinline{2\rho}\), the latter two greater than \(\fitinline{\rho}\). The smallest terminal room above \(\fitinline{\rho}\) is greater than \(\fitinline{2.67\cdot10^{-13}}\). Choose the total new local exponent fees below \(\fitinline{10^{-14}}\); \(\fitinline{\rho}\) itself is not that spare room. In the finite induction choose the inner strict decrements and meshes in dependency order, making their accumulated costs smaller than each outer positive remainder.

For \(\fitinline{e\ge.5002}\) on \(\fitinline{V\ge.98}\), use H6/SN9 at slope .5001 with fixed positive gap .0001, on the actual signed merged column and its true squared coefficient mass. The root volume \(\fitinline{A_0=n+y+W<.919}\) gives \(\fitinline{.5001(.98-A_0)>.030}\). On \(\fitinline{[.875,.98]}\) the native fourths and squared prime bound \(\fitinline{3/4}\) cost at most \[\fitmath{
 \tfrac34[(1-V)/2+.05805913/2-.07599618(.373)
                           +\max\beta-u]_+.
}\] Clipping a negative ordinary quota at zero increases the available allowance, so this remains an upper cost. At .98 the reserve against \(\fitinline{5(H-V)/6}\) exceeds .008 and its slope is at most \(\fitinline{3/8}\). On \(\fitinline{[.08,.875]}\) use both unmarked nonprincipal fourths; their cost is at most \(\fitinline{3(W_{\max}-u)/4}\), leaving more than .023 at .875. Below .08 the actual small-core/lift count leaves more than .031. Thus all radii have regular-core energy at most \[\fitmath{
                  \mathcal H U^{V+u-\rho+(5/6)(H-V)_++\epsilon}.
}\tag{L10}\] For \(\fitinline{z=b^6c}\), count all cofactors by \(\fitinline{U^{(M-V)/6+\epsilon}}\). The earned \(\fitinline{5/6}\) slope gives the required whole-ball price with \(\fitinline{M}\) in place of \(\fitinline{V}\). No common-prime or unit restriction is added.

The full H.5 assembly is required at each stage. Its actual C.10 support test is \(\fitinline{\mathcal V<U^{.919}<U}\) and \(\fitinline{(9/4)^22^N\mathcal V<U^2}\) eventually for the fixed slot type; the earlier convenient .903 cap is not used. The .139 short reference has \(\fitinline{4(.139+W_{\max})<H-.012}\); its other partner is below .71 and uses the arbitrary bounded-length unmarked nonprincipal fourth. Both raw orientations keep H25's unweighted first divisor, H26's compensated second divisor, the exact \(\fitinline{V_C,V_D}\) norm identities, both independent phases and the actual numerator/denominator orientation. The signed fixed-row estimate at the common four-coordinate minimum retains both principal arms and has \(\fitinline{5/6+.139-A_0-3u>.054}\). The original zeros, seed sums, counted rapid tails and full label sum remain, with \(\fitinline{1-5u>.9998}\) and \(\fitinline{1-A_0-4u>.081}\). Reconstruction H.5(A)--(C) returns the actual uncentered kept root with an open width. This finishes the source proof, not only its regular terminal leaf.

For uniformity treat L2 as a family of upper caps. On a compact earlier rectangle cover the needed \(\fitinline{(s,n/s)}\) pairs by a finite mesh rounded upward, retaining its closed upper edges and a sufficiently small continuous width change. This dominates both actual lengths, including every shortening. Split true radius and length bins at source seams. Flatten the 308 earlier stages into a finite list; at an outer call choose its early maximum slot width below the minimum inner width times its positive lower core-radius ratio and below every prescribed loss. Finite ray dictionaries and their dual products can be fixed to depth 308 before external height bounds. After the late finite slot family choose the common profile/Fourier/Sobolev orders, then the one positive height exponent, then every later buffered grid bound. Distribute the final \(\fitinline{U^\epsilon}\) costs over the finite calls. There is no growing-depth or growing-modulus uniformity assumption.

\hypertarget{l.3.-a-finite-maximum-of-genuine-saturated-deficits}{%
\subsection{L.3. A finite maximum of genuine saturated deficits}\label{l.3.-a-finite-maximum-of-genuine-saturated-deficits}}

Use K.1--K.2's original reciprocal at \(\fitinline{s}\), its three independent contour heights and its actual finite-family buffer. A detector does not replace the physical denominator. Fix a nonfloor \(\fitinline{(d,p)}\) cell, its paired crossing \(\fitinline{R,b_0=2m,t_0=b_0+m}\), \(\fitinline{k=1-p(1+2g_1)}\), \(\fitinline{V_d=5/6-d}\), and actual residual \(\fitinline{Z}\)-debit \(\fitinline{B_*}\) after its unused E49 and conductor credits. Put \(\fitinline{B_+=\max(0,B_*)+10^{-8}}\). Before selecting any witness choose a finite bank \[\fitmath{
 T_j=t_0-\gamma_j,\qquad
 \gamma_j=j\gamma_{\max}/N,\quad0\le j\le N,
 \qquad\gamma_{\max}\ge4B_+/(3D_-V_-),                         
}\tag{L11}\] with the directed source/support collars. All cutoffs are strictly in B.5.1's domain \(\fitinline{[1,3/2]}\). A fine predetermined \(\fitinline{(d,p)}\) subdivision permits fixed representative cutoffs. At each cutoff use the same actual attaining zero of that row's finite family. B.5.1 gives genuine saturated inverse/ordinary lengths \(\fitinline{x_j,y_j}\), with \(\fitinline{x_j\le T_j+o(1)}\), \(\fitinline{x_j+y_j\ge T_j-o(1)}\).

Take \(\fitinline{\mathrm{ann}=10^{-8}}\), and round \(\fitinline{\gamma_{\max}}\) upward to \(\fitinline{10^{-9}}\). Likewise choose \(\fitinline{\lambda_x\ge4B_+/(3D_-a_dk_-)}\) and \(\fitinline{\lambda_y\ge2B_+/(3D_-a_dk_-)}\), with that same upward rounding. After the paid outside faces, the length interval at cutoff \(\fitinline{j}\) is enclosed by \(\fitinline{[m_- -\gamma_j-\lambda_x-2\mathrm{ann},m_++\lambda_y+\mathrm{ann}]}\), with directed endpoint rounding. These are actual source-length enclosures, not independent freely chosen factors.

Let \(\fitinline{G}\) be the largest nonnegative \textbf{raw count deficit} furnished by any of these observed saturated witnesses and any actually admitted tangent, native, affine, cap or unmarked branch. Freeze the finite witness/length descriptors: their maximum certified deficit gives the minimum of their count upper bounds on the same designated rows. Equivalently use the qualifying event and a finite descriptor union. No counts are multiplied and no joint moment of witnesses is presumed. If \(\fitinline{G\ge4B_+/(3D_-)}\), K3's numerator fourth pays the target by the genuine \(\fitinline{3G/4}\) gain, up to the preallocated witness/support losses. Any short inverse or ordinary in a fixed outside face is paid in this way. If all inverse witnesses are amplified, the last cutoff supplies deficit \(\fitinline{V_d\gamma_{\max}-o(1)}\), giving \(\fitinline{B_+-o(1)>B_*}\).

Otherwise take the first index \(\fitinline{j}\) whose inverse is short, \(\fitinline{x_j<1}\), after the paid outside faces have been removed. Its two lengths lie in the central anchors. The current inverse gives \(\fitinline{dk(x_j-b_0)_+}\); if \(\fitinline{j>0}\), the preceding amplified inverse gives \(\fitinline{V_d\gamma_{j-1}}\), with \(\fitinline{\gamma_{-1}=0}\) for \(\fitinline{j=0}\). The same current ordinary also gives \(\fitinline{2dk(y_j-m)_+}\). Hence on these same rows \[\fitmath{
 G\ge\max\{dk(x_j-b_0)_+,\ V_d\gamma_{j-1},\ 2dk(y_j-m)_+\}.
}\tag{L12}\] The bound for \(\fitinline{G}\) in L12 is up to the same preallocated witness/support losses, paid by the explicit collars in L16. It is obtained from the whole predetermined bank; the first-short index merely identifies usable faces of it. This is distinct from selecting a detector by binning \(\fitinline{G}\) itself.

Use the genuinely saturated ordinary at this index in the rowwise comparison K6. The source in L1--L10 applies to the actual two fixed rays before division by the original common marks. After that division and Cauchy, the extra \(\fitinline{U}\)-saving over the tangent mean is at least \[\fitmath{
 dp\Delta+\frac12\left[dk(y_j-m)+
   \max\{dk(x_j-b_0)_+,V_d\gamma_{j-1},2dk(y_j-m)_+\}\right].   
}\tag{L13}\] Using the support bound gives the convenient lower price \[\fitmath{
 dp\Delta+\frac d2\max\left\{-k\gamma_j,
     k(y_j-m)+\max\{(V_d/d)\gamma_{j-1},2k(y_j-m)_+\}\right\},
}\tag{L14}\] up to the directed finite errors. The original full positive coefficient mass J9a pays \(\fitinline{q(n_{\rm ref}-n_0)/2}\) just as in K8. No new analytic loss is charged for every tested cutoff, and no physical term is discarded because a different detector was favorable.

There is one genuine lower-coordinate adjustment. If the actual ordinary length \(\fitinline{y<.373}\), use \(\fitinline{y_{\rm use}=.373}\); otherwise use \(\fitinline{y_{\rm use}=y}\). The upper source pair is \(\fitinline{(n_{\rm ref},y_{\rm use})}\), and must satisfy L2. Its effective extra width relative to the \emph{actual} tangent at \(\fitinline{y}\) is \[\fitmath{
 \Delta_{\rm eff}(y)=\beta(y_{\rm use})
               -\frac{1+2g_1}{2}(y_{\rm use}-y).              
}\tag{L15}\] The saturated lower in K6 still uses the actual \(\fitinline{y}\). Thus raising the source coordinate incurs the displayed true tangent debit; it is not free clamping. Actual smaller numerator coordinates use componentwise downward restriction at \(\fitinline{n_{\rm ref}=\max(n,n_0)}\), retaining its width and all null/unit cases.

On a cell \(\fitinline{[a_d,b_d]\times[f,e_p]}\), enclose \(\fitinline{m}\) by \(\fitinline{m_-\le m\le m_+}\), \(\fitinline{k_-\le k\le k_+}\), and \(\fitinline{V_d\ge V_-}\). A sufficient normalized price on every actual length interval is \[\fitmath{
 f\Delta_{\rm eff}(y)+\tfrac12\max\left\{
  -k_+(\gamma_j+3\mathrm{ann}),\quad
  k_{\rm sign}(y-m_+-\mathrm{ann})+
  \max\{(V_-/b_d)(\gamma_{j-1}-3\mathrm{ann})_+,
                    2k_-(y-m_+-\mathrm{ann})_+\}\right\}
  -e_p\mathrm{ann}/2-B_+/(D_-a_d)>0,                         
}\tag{L16}\] provided \(\fitinline{\Delta_{\rm eff}\ge0}\); take \(\fitinline{k_{\rm sign}=k_+}\) for a negative argument and \(\fitinline{k_-}\) for a positive one. If a genuine application has negative effective width, price its \(\fitinline{p}\)-factor with the appropriate upper endpoint instead. The final ticket must check these signs, every ratio \(\fitinline{n_{\rm ref}/y_{\rm use}}\), every actual shortening and every contour/source admission.

This is a finite exact minimization, not sampling. Write \(\fitinline{\delta=y-m_+-\mathrm{ann}}\) and \(\fitinline{g_m=(V_-/b_d)(\gamma_{j-1}-3\mathrm{ann})_+}\). Partition at \(\fitinline{.373}\), \(\fitinline{m_++\mathrm{ann}}\), the crossover of the constant and sloping faces, \(\fitinline{m_++\mathrm{ann}+g_m/(2k_-)}\), and their sign switches. Each remaining piece is affine or a convex quadratic with coefficient \(\fitinline{fc_2}\). Check both endpoints and each interior derivative-zero vertex. The latter are \(\fitinline{l}\), \(\fitinline{l-k_-/(4fc_2)}\), \(\fitinline{l-k_+/(4fc_2)}\), or \(\fitinline{l-3k_-/(4fc_2)}\), according to the active face. This is the full candidate list.

A further exact reduction pays most cells without enumerating every cutoff. Suppose its whole length interval \(\fitinline{[y_-,y_+]}\) has \(\fitinline{y_-\ge.373}\), \(\fitinline{y_+<l}\), and \[\fitmath{
 \begin{gathered}
 2fc_2(l-y_-)<k_-/2,\qquad
 2fc_2(l-y_-)\left(1+\frac{V_-}{b_dk_+}\right)<k_+/2,\\
 \lambda_x\ge\frac{V_-}{b_dk_+}(\gamma_{\max}-3\mathrm{ann})_+.
 \end{gathered}                                               
}\tag{L16a}\] At cutoff \(\fitinline{j}\), the minimum lies at \[\fitmath{
 y_c=m_++\mathrm{ann}-\gamma_j-3\mathrm{ann}
          -\frac{V_-}{b_dk_+}(\gamma_{j-1}-3\mathrm{ann})_+ .   
}\tag{L16b}\] It is within the actual length interval by the last inequality and \(\fitinline{m_-\le m_+}\). Below this crossover the constant face is active and the decreasing quadratic decreases the price; above it the sloping face's derivative dominates the quadratic by the first inequality. At and below \(\fitinline{y_c<m_++\mathrm{ann}}\), the ordinary positive face vanishes; above it that face can only help. As \(\fitinline{j}\) increases, the crossover moves left with speed at most \(\fitinline{1+V_-/(b_dk_+)}\), while the negative constant face loses \(\fitinline{k_+/2}\) per unit \(\fitinline{\gamma}\). The second inequality therefore makes the minimum price decrease with \(\fitinline{j}\). It suffices to check \(\fitinline{j=N}\) at L16b. Cells not satisfying L16a use the complete finite candidate enumeration, including every cutoff. The actual receipt records which proof is used on each cell.

\hypertarget{l.4.-changed-exterior-and-direct-geometric-contracts}{%
\subsection{L.4. Changed exterior and direct geometric contracts}\label{l.4.-changed-exterior-and-direct-geometric-contracts}}

The following changes use the original positive-cost source and counts. Choose physical cap \(\fitinline{\sigma=.10}\) and declared cap \(\fitinline{\bar\sigma=.11}\), independently of the enhanced count's cap. With source \(\fitinline{G(1-s,2+10s,3/20-1000s/1143)}\), the actual GT4--GT7 inequalities are \[\fitmath{
 .10<.11<c_{\rm src},\quad \ell/(h+F)>.2,\quad
 B_{\rm src}/2+(9/2-L_{\rm src}/2)\bar\sigma<1,\quad
 2y_t-(2-B_{\rm src})(h+F)>0.                                
}\tag{L17}\] The same GT13--GT15 then supply real shortenings, shifted profiles, zero-width and nonpositive-length fallbacks and the full rowwise coefficient damping. The old physical capacity stays available, while the larger cap may be used on the same affine numerator branch. In particular it is the chosen .10 source, rather than an assumed .075 ceiling, that pays any \(\fitinline{z_0>.075}\). Its upper-conductor slope is at least \(\fitinline{3(1-b_d)/4-9/100-b_d\sigma/4}\).

For a strictly positive saved polynomial lower bound \(\fitinline{L_i}\), divide by its actual positive denominator. On a subbox \(\fitinline{[a_d,b_d]\times[f,e_p]}\), write \[\fitmath{
 J(d,p)=(5/6-d)D(p)+dC(p).
}\] For a native/nativeLow branch use \(\fitinline{D(p)=3-p-4p/L_0}\), \(\fitinline{C(p)=(1-p)(2-4p/L_0)}\), \(\fitinline{L_0=200000000001/10^{11}}\); for old E48 use \(\fitinline{D(p)=3-p-2p/(3\kappa)}\), \(\fitinline{C(p)=2(1-p)(1-p/(3\kappa))}\), conservatively at the upper \(\fitinline{\kappa}\) endpoint. Both coefficients decrease with \(\fitinline{p}\), and \(\fitinline{D-C>0}\): respectively \(\fitinline{1+p-4p^2/L_0}\) and \(\fitinline{1+p-2p^2/(3\kappa)}\). Hence \[\fitmath{
 J(d,p)\le(5/6-a_d)D(f)+a_dC(f),                             
}\tag{L18}\] giving a valid improved lower bound \(\fitinline{L_i/J_+}\). It is used only for \(\fitinline{L_i>0}\), on its actual edge and admitted branch. Retain J12's genuine E49 reserve on paired cells, all ten E29 guards at four contact corners on every newly used upward window, and conductor convexity with the same conductor-independent count envelope. None of these operations changes the actual equality \(\fitinline{D_{\rm src}.\mathrm{momentKappa}=2\beta-1}\).

The fixed native ordinary continuation extends through \(\fitinline{y=.45}\). On \(\fitinline{y_*=(1-(9/2).027)/2=.43925\le y\le.45}\) take \(\fitinline{z=(1-2y)/(9/2)}\) with its independently affordable strict decrement. Then \(\fitinline{2y+6\kappa z<1}\) for \(\fitinline{\kappa\le3/4}\), and \(\fitinline{z>.02}\) even at .45. The coefficient \(\fitinline{2y+2pz}\) increases from its old initial value. Beyond .45 the unmarked coefficient is at least .90. This is the native B(211) branch, not an extended enhanced tangent or a deleted row sector. All six fixed continuation prices still use genuine saturated count savings and the original numerator's \(\fitinline{3/4}\) exponent.

For the final target put \(\fitinline{\ell=11/3-4t}\), \(\fitinline{y_t=h-2\ell}\), \(\fitinline{x_t=1+\ell-h}\). The earlier J18 corridor is not assumed. A direct field proof is available whenever \[\fitmath{
             .8733\le t\le.873345,\qquad .815324503\le h<.816.
}\tag{L19}\] Indeed \(\fitinline{25993/150000\le\ell\le1301/7500<.175}\), so \(\fitinline{1/6<\ell<1/5}\), \(\fitinline{x_t-\ell=1-h>.184}\), \(\fitinline{y_t-x_t=2h-1-3\ell>.110249}\), \(\fitinline{1/2-y_t>.0305733}\), \(\fitinline{2x_t-y_t-\ell=2+3\ell-3h>.07186}\), and \(\fitinline{5\ell-1<0}\). Also \(\fitinline{x_t+y_t=1-\ell}\) and \(\fitinline{t-(4+y_t-x_t)/6=x_t/2+(y_t-x_t)/12}\). These are every \texttt{G\allowbreak{}e\allowbreak{}o\allowbreak{}m\allowbreak{}e\allowbreak{}t\allowbreak{}r\allowbreak{}y\allowbreak{}.\allowbreak{}P\allowbreak{}a\allowbreak{}r\allowbreak{}a\allowbreak{}m\allowbreak{}e\allowbreak{}t\allowbreak{}e\allowbreak{}r\allowbreak{}s} field. The five \texttt{B\allowbreak{}o\allowbreak{}u\allowbreak{}n\allowbreak{}d\allowbreak{}s\allowbreak{}C\allowbreak{}o\allowbreak{}r\allowbreak{}e} fields follow from \(\fitinline{.873<t<7/8}\), \(\fitinline{.793<h<.816}\), \(\fitinline{.458<y_t<.473}\), \(\fitinline{.1672<\ell<.175}\), and \[\fitmath{
 y_t/2-h(17/50-1/6)-1/50=49h/150-\ell-1/50>.07287267>.07.
}\tag{L20}\] All further cap, conductor, floor, contour and fixed-pad guards must be checked at the actual numerical \(\fitinline{(t,h)}\). The literal \(\fitinline{h<.816}\) interface is retained; a larger conductor would require replacing every authentic wrapper consumer.

\hypertarget{l.5.-the-complete-ticket-at-l0-and-the-original-return}{%
\subsection{L.5. The complete ticket at L0 and the original return}\label{l.5.-the-complete-ticket-at-l0-and-the-original-return}}

Apply these arguments at the exact L0 geometry, with the original physical source \(\fitinline{s=10^{-260}}\), numerical pad \(\fitinline{F=10^{-258}}\), physical cap .10 and declared cap .11. There are still exactly 57 original boxes and four closed conductor-edge tilings. Use J12's E49 reserve, J13's actual \(\fitinline{h-h_*}\) perturbation, the old globally admitted count for the saved-edge slope, and the actual positive denominator L18. The actual moment parameter remains \(\fitinline{2\beta-1}\); under a contrary \(\fitinline{t_L<\beta}\) the earlier attained bound gives \(\fitinline{\beta\le t_K<\alpha}\), so FB2--FB6 justify only the conservative pricing parameter \(\fitinline{2\alpha-1}\).

Besides the original E27 window and I32 wing, the complete recorded upward admissions are \[\fitmath{
 \begin{array}{c|c}
 d\text{ interval}&p\text{ interval}\\ \hline
 [93/256,3/8]&[15/32,1/2]\\{}
 [741/2048,93/256]&[127/256,1/2]\\{}
 [741/2048,93/256]&[255/512,1/2]\\{}
 [93/256,189/512]&[31/64,1/2]\\{}
 [3/8,27/64]&[7/16,1/2].
 \end{array}                                                 
}\tag{L21}\] For each whole closed rectangle, all ten cleared K16 guards at each of the four contact corners have positive tensor-Bernstein lower bounds. The least recorded guard across these windows is \(\fitinline{310983969823/(5\cdot10^{14})>0}\). These are applications of E29 with its actual same-zero crossing, native/affine/capped/unmarked and amplitude-monotonicity gates, not invented paid windows. Conductor convexity fills each original interval using the same admitted pointwise envelope at its two ends. All original saved or paired complementary bins are covered.

The coarse physical construction has 1078 central cells. Subdivide only a cell whose quadratic price has insufficient stated reserve into four equal \(\fitinline{(d,p)}\) children; recheck each child's paired admission, E49 debit, conductor collar, dilation/length prices, fixed continuations and source support independently. There are 75 such four-way subdivisions, 1378 inspected nodes and 1303 final central cells, with no refined child removed as an exterior shortcut. On 1292 final cells L16a--L16b prove the full length-and-cutoff minimum; the remaining eleven check every cutoff and all endpoints, kinks and quadratic vertices.

The bank has \(\fitinline{N=256}\) spacings and \textbf{257 predetermined cutoffs} on each final cell. The worst normalized price L16 is at least \[\fitmath{
 \frac{10168881}{10^{15}}=1.0168881\cdot10^{-8},                
}\tag{L22}\] after the separately included \(\fitinline{10^{-8}}\) budget in \(\fitinline{B_+}\). Every cell has \(\fitinline{Dd>.3}\); hence the additional \(\fitinline{Z}\)-saving is at least \(\fitinline{381333/125000000000000>3.05066\cdot10^{-9}}\), on top of that budget. The exact worst occurs on original cell52, a depth-three child, at \(\fitinline{j=256}\), \(\fitinline{y\approx.3762291458784}\). That decimal identifies the exact rational minimizer recorded in the certificate; it is not the value used to decide the sign.

The actual source inputs over all final cells satisfy \[\fitmath{
 \begin{gathered}
 .4247101972508802<n_{\rm ref}\le.425366994805076,\\
 .37280873983177\le y\le.378956998253669,\qquad
 y_{\rm use}=\max(.373,y),\\
 1.1209910584<n_{\rm ref}/y_{\rm use}<1.1392644660.
 \end{gathered}                                               
}\tag{L23}\] Thus every upper pair is strictly inside L2's ratio interval and in its closed length range. L15's true debit is included on every shortened lower-length case; its effective width is nonnegative wherever priced with the lower \(\fitinline{p}\)-endpoint. The source covers every actual componentwise shortening. There is no radical, squarefree-row, inverse-amplitude or shared-height restriction. The original total slot width is greater than \(\fitinline{1/5}\); L1's global width maximum is below .108, so the same whole-slot greedy lower response, including positive exhaustion, supplies the required squared mark gain. Its prefix remains a subproduct of the unchanged whole original prime pools.

The complete fixed and moving complementary prices, including every refined child, have respective bare minima \[\fitmath{
 \begin{array}{c|c}
 \text{saved original branch}&133950211/10^{15}\\
 \text{E49 exterior}&1021723/(5\cdot10^{13})\\
 \text{conductor/dilation/length/fixed and moving gates}&10000011/10^{15}.
 \end{array}                                                 
}\tag{L24}\] These are cellwise checks; no single box formed from independent global worst endpoints is used to reprice the union. Coarse enclosing parameters are \(\fitinline{D\ge.815854982}\), \(\fitinline{.404296875\le d\le.44287109375}\), \(\fitinline{p\ge.4951171875}\). The lower contour has \(\fitinline{d(1-p)/2-.10\ge.00107421875>0}\), while \(\fitinline{(1+d)/2-.10>.51}\) and \(\fitinline{(1+d)/2+.27<1}\). The same absolute contour motion, full \(\fitinline{w}\)-inversion and quotient-free tuple mass as K.1 therefore pay both inverse-length tails on their original classes, before amplitude subdivision. Large dilation uses its original clipped native capacity and \(\fitinline{dv/4}\) saving; the retained packet uses its fully summed J9a mass once inside the mixed comparison. The .10 source also handles all short-numerator, null and bounded-scale cases at their authentic quota.

The actual \(\fitinline{t_L,h}\) satisfy L19--L20 and L17, \(\fitinline{h+F<\min(1,2y_t)}\), \(\fitinline{y_t<.76<h}\), and \(\fitinline{\ell/(h+F)>.2}\). Thus the generic geometry is constructed directly below J's old corridor, with every \texttt{G\allowbreak{}e\allowbreak{}o\allowbreak{}m\allowbreak{}e\allowbreak{}t\allowbreak{}r\allowbreak{}y\allowbreak{}.\allowbreak{}P\allowbreak{}a\allowbreak{}r\allowbreak{}a\allowbreak{}m\allowbreak{}e\allowbreak{}t\allowbreak{}e\allowbreak{}r\allowbreak{}s} and \texttt{B\allowbreak{}o\allowbreak{}u\allowbreak{}n\allowbreak{}d\allowbreak{}s\allowbreak{}C\allowbreak{}o\allowbreak{}r\allowbreak{}e} field and the exact signal boundary. The averaged-floor margin after \(\fitinline{10^{-5}+10^{-12}}\) exceeds .00293117; the small-row margin exceeds .07157247; the two principal margins exceed .02140165 and .0001099991. The floor, bounded rows, five nonprincipal units, actual principal residues, zero frequencies, finite prime corrections, compensators and rapid Fourier tails retain their original B/H arguments with these checked bounds. No pole is paid by a row-count gain.

Before live target characters and heights, choose baseline data with \(\fitinline{0<D_{\rm src}.t\le F/4}\) and losses below \(\fitinline{\min((\beta-t_L)/10000,10^{-9}/4,F)}\). Then the literal cube pad satisfies \(\fitinline{F/2+2D_{\rm src}.t\le F}\). Charge the separate price \(\fitinline{15s/16+2F}\) on every L24 sector. Capacities are still 1-Lipschitz and the original conductor variation is below \(\fitinline{411/400<2}\), extending the full partition to \(\fitinline{D\le h+F}\). These charges leave more than \(\fitinline{3\cdot10^{-9}}\); choose all other finite physical losses with total below \(\fitinline{10^{-9}}\). The much smaller inner source losses in L9 are also respected. Neither the larger cap nor the finite bank changes the order of choices.

Fix the complete finite source dependency list, its strict decrements, the physical cutoff banks and the early mesh before the later actual slot family and ray dictionary. After that family is fixed choose common compact Fourier envelopes, seminorms and Sobolev orders. Scale/dilation suprema use the actual log-profile derivative/envelope, height suprema the separate normalized first log tests. Only then choose one positive target-dependent height exponent, before every later positive buffered grid bound. All three physical contour frequencies and the independent ordinary/detector frequencies keep their original weights. Large tails are chosen against these finite polynomial envelopes. These estimates concern the entire physical reciprocal and numerator, not a favorable term selected by a witness.

They give the literal I51 cube statement at L0, for every original \texttt{A\allowbreak{}c\allowbreak{}t\allowbreak{}u\allowbreak{}a\allowbreak{}l\allowbreak{}B\allowbreak{}u\allowbreak{}f\allowbreak{}f\allowbreak{}e\allowbreak{}r\allowbreak{}e\allowbreak{}d\allowbreak{}S\allowbreak{}e\allowbreak{}l\allowbreak{}e\allowbreak{}c\allowbreak{}t\allowbreak{}i\allowbreak{}o\allowbreak{}n}, in the order \(\fitinline{\forall\eta\ \exists\tau_\eta>0\ \forall n>0\ \exists C_n}\), eventually in \(\fitinline{Z}\). Choose \(\fitinline{0<\sigma<10^{-9}}\) below the retained transport reserves. Keep the same data, cutoff and nonzero normalizer I52, with only the new \(\fitinline{\ell,x_t,y_t}\). The generic \texttt{h\allowbreak{}i\allowbreak{}g\allowbreak{}h\allowbreak{}\_\allowbreak{}o\allowbreak{}f\allowbreak{}\_\allowbreak{}c\allowbreak{}u\allowbreak{}b\allowbreak{}e} under the directly proved geometry gives the original physical probe and signal, with Gram exponent \(\fitinline{t_L-k_{\rm phys}}\). For a contrary \(\fitinline{\beta>t_L}\), choose \(\fitinline{\omega<(\beta-t_L)/4}\) and \(\fitinline{0<m<\min(\beta-t_L-\omega,\sigma,\beta-.87)}\). The H47 signal then has a Mellin transform holomorphic on \(\fitinline{\Re s>\beta-m}\), but on the absolute line this is the original Gaussian times \(\fitinline{H_\eta/L_\eta}\). The numerator satisfies \(\fitinline{|H_\eta-1|\le1/2}\) on \(\fitinline{\Re s\ge.87}\). A target zero above \(\fitinline{\beta-m/2}\), supplied by the definition of the supremum, contradicts holomorphy. An attained supremum is unnecessary. Norm pullback and fixed Euler restoration prove L0 for every ordinary character, including imprimitive presentations, at every finite height.

\hypertarget{l.6.-exact-finite-objects-and-provenance}{%
\subsection{L.6. Exact finite objects and provenance}\label{l.6.-exact-finite-objects-and-provenance}}

The source derivation is \texttt{r\allowbreak{}e\allowbreak{}s\allowbreak{}e\allowbreak{}a\allowbreak{}r\allowbreak{}c\allowbreak{}h\allowbreak{}/\allowbreak{}l\allowbreak{}o\allowbreak{}c\allowbreak{}a\allowbreak{}l\allowbreak{}\_\allowbreak{}c\allowbreak{}y\allowbreak{}c\allowbreak{}l\allowbreak{}e\allowbreak{}0\allowbreak{}8\allowbreak{}4\allowbreak{}/\allowbreak{}c\allowbreak{}o\allowbreak{}n\allowbreak{}t\allowbreak{}r\allowbreak{}o\allowbreak{}l\allowbreak{}l\allowbreak{}e\allowbreak{}r\allowbreak{}/\allowbreak{}Q\allowbreak{}U\allowbreak{}A\allowbreak{}D\allowbreak{}R\allowbreak{}A\allowbreak{}T\allowbreak{}I\allowbreak{}C\allowbreak{}\_\allowbreak{}P\allowbreak{}R\allowbreak{}O\allowbreak{}F\allowbreak{}I\allowbreak{}L\allowbreak{}E\allowbreak{}\_\allowbreak{}P\allowbreak{}R\allowbreak{}O\allowbreak{}O\allowbreak{}F\allowbreak{}.\allowbreak{}m\allowbreak{}d} (\texttt{d\allowbreak{}a\allowbreak{}d\allowbreak{}0\allowbreak{}6\allowbreak{}c\allowbreak{}3\allowbreak{}7\allowbreak{}1\allowbreak{}0\allowbreak{}f\allowbreak{}7\allowbreak{}5\allowbreak{}f\allowbreak{}d\allowbreak{}0\allowbreak{}9\allowbreak{}6\allowbreak{}c\allowbreak{}3\allowbreak{}f\allowbreak{}1\allowbreak{}f\allowbreak{}1\allowbreak{}7\allowbreak{}e\allowbreak{}e\allowbreak{}2\allowbreak{}9\allowbreak{}4\allowbreak{}c\allowbreak{}0\allowbreak{}7\allowbreak{}7\allowbreak{}1\allowbreak{}d\allowbreak{}8\allowbreak{}a\allowbreak{}d\allowbreak{}d\allowbreak{}7\allowbreak{}e\allowbreak{}e\allowbreak{}0\allowbreak{}c\allowbreak{}3\allowbreak{}8\allowbreak{}b\allowbreak{}2\allowbreak{}9\allowbreak{}e\allowbreak{}d\allowbreak{}5\allowbreak{}8\allowbreak{}4\allowbreak{}e\allowbreak{}9\allowbreak{}3\allowbreak{}a\allowbreak{}3\allowbreak{}5\allowbreak{}9\allowbreak{}d\allowbreak{}2}, 14861 bytes), with the rational construction \texttt{P\allowbreak{}R\allowbreak{}O\allowbreak{}F\allowbreak{}I\allowbreak{}L\allowbreak{}E\allowbreak{}\_\allowbreak{}C\allowbreak{}H\allowbreak{}E\allowbreak{}C\allowbreak{}K\allowbreak{}.\allowbreak{}p\allowbreak{}y} (\texttt{f\allowbreak{}a\allowbreak{}a\allowbreak{}0\allowbreak{}e\allowbreak{}d\allowbreak{}4\allowbreak{}1\allowbreak{}5\allowbreak{}1\allowbreak{}7\allowbreak{}5\allowbreak{}a\allowbreak{}d\allowbreak{}e\allowbreak{}8\allowbreak{}5\allowbreak{}c\allowbreak{}9\allowbreak{}1\allowbreak{}4\allowbreak{}2\allowbreak{}0\allowbreak{}1\allowbreak{}2\allowbreak{}7\allowbreak{}c\allowbreak{}d\allowbreak{}2\allowbreak{}c\allowbreak{}a\allowbreak{}c\allowbreak{}0\allowbreak{}4\allowbreak{}6\allowbreak{}a\allowbreak{}7\allowbreak{}2\allowbreak{}6\allowbreak{}9\allowbreak{}0\allowbreak{}5\allowbreak{}4\allowbreak{}e\allowbreak{}4\allowbreak{}2\allowbreak{}f\allowbreak{}1\allowbreak{}1\allowbreak{}2\allowbreak{}9\allowbreak{}7\allowbreak{}6\allowbreak{}5\allowbreak{}3\allowbreak{}c\allowbreak{}2\allowbreak{}c\allowbreak{}8\allowbreak{}9\allowbreak{}d\allowbreak{}1\allowbreak{}4\allowbreak{}5}) and compact output \texttt{P\allowbreak{}R\allowbreak{}O\allowbreak{}F\allowbreak{}I\allowbreak{}L\allowbreak{}E\allowbreak{}\_\allowbreak{}C\allowbreak{}E\allowbreak{}R\allowbreak{}T\allowbreak{}I\allowbreak{}F\allowbreak{}I\allowbreak{}C\allowbreak{}A\allowbreak{}T\allowbreak{}E\allowbreak{}.\allowbreak{}j\allowbreak{}s\allowbreak{}o\allowbreak{}n} (\texttt{6\allowbreak{}e\allowbreak{}e\allowbreak{}b\allowbreak{}e\allowbreak{}0\allowbreak{}2\allowbreak{}7\allowbreak{}c\allowbreak{}e\allowbreak{}1\allowbreak{}2\allowbreak{}a\allowbreak{}c\allowbreak{}b\allowbreak{}8\allowbreak{}7\allowbreak{}3\allowbreak{}f\allowbreak{}7\allowbreak{}9\allowbreak{}8\allowbreak{}6\allowbreak{}c\allowbreak{}4\allowbreak{}5\allowbreak{}f\allowbreak{}c\allowbreak{}1\allowbreak{}c\allowbreak{}2\allowbreak{}1\allowbreak{}0\allowbreak{}8\allowbreak{}8\allowbreak{}7\allowbreak{}2\allowbreak{}7\allowbreak{}4\allowbreak{}1\allowbreak{}7\allowbreak{}2\allowbreak{}6\allowbreak{}4\allowbreak{}6\allowbreak{}7\allowbreak{}9\allowbreak{}f\allowbreak{}4\allowbreak{}3\allowbreak{}4\allowbreak{}2\allowbreak{}0\allowbreak{}e\allowbreak{}2\allowbreak{}2\allowbreak{}e\allowbreak{}c\allowbreak{}e\allowbreak{}3\allowbreak{}7\allowbreak{}3\allowbreak{}0\allowbreak{}f}). The finite object records L1's exact formula, 308 induction stages, 1120 length intervals, 2240 affine triangles, all three minima L9, threshold range and every peripheral/raw/reference/fixed/zero margin. The old I source checker supplies only its pinned integer fifth-root function. Those integer comparisons establish the numerical inequalities used above; the finite induction, source applicability and all-height return are the mathematical arguments L.1--L.5.

The actual physical output is \texttt{r\allowbreak{}e\allowbreak{}s\allowbreak{}e\allowbreak{}a\allowbreak{}r\allowbreak{}c\allowbreak{}h\allowbreak{}/\allowbreak{}l\allowbreak{}o\allowbreak{}c\allowbreak{}a\allowbreak{}l\allowbreak{}\_\allowbreak{}c\allowbreak{}y\allowbreak{}c\allowbreak{}l\allowbreak{}e\allowbreak{}0\allowbreak{}8\allowbreak{}4\allowbreak{}/\allowbreak{}n\allowbreak{}e\allowbreak{}a\allowbreak{}r\allowbreak{}\_\allowbreak{}r\allowbreak{}e\allowbreak{}c\allowbreak{}o\allowbreak{}r\allowbreak{}d\allowbreak{}/\allowbreak{}M\allowbreak{}A\allowbreak{}X\allowbreak{}Q\allowbreak{}\_\allowbreak{}1\allowbreak{}2\allowbreak{}6\allowbreak{}6\allowbreak{}6\allowbreak{}4\allowbreak{}1\allowbreak{}4\allowbreak{}.\allowbreak{}j\allowbreak{}s\allowbreak{}o\allowbreak{}n}, SHA-256 \texttt{b\allowbreak{}d\allowbreak{}6\allowbreak{}d\allowbreak{}e\allowbreak{}0\allowbreak{}c\allowbreak{}6\allowbreak{}4\allowbreak{}6\allowbreak{}f\allowbreak{}8\allowbreak{}f\allowbreak{}0\allowbreak{}4\allowbreak{}9\allowbreak{}2\allowbreak{}0\allowbreak{}d\allowbreak{}c\allowbreak{}3\allowbreak{}8\allowbreak{}7\allowbreak{}8\allowbreak{}a\allowbreak{}9\allowbreak{}e\allowbreak{}3\allowbreak{}f\allowbreak{}8\allowbreak{}6\allowbreak{}1\allowbreak{}b\allowbreak{}9\allowbreak{}e\allowbreak{}4\allowbreak{}d\allowbreak{}d\allowbreak{}2\allowbreak{}2\allowbreak{}3\allowbreak{}e\allowbreak{}7\allowbreak{}5\allowbreak{}9\allowbreak{}4\allowbreak{}d\allowbreak{}2\allowbreak{}0\allowbreak{}4\allowbreak{}7\allowbreak{}5\allowbreak{}8\allowbreak{}4\allowbreak{}7\allowbreak{}6\allowbreak{}c\allowbreak{}4\allowbreak{}d\allowbreak{}a\allowbreak{}2\allowbreak{}9\allowbreak{}a\allowbreak{}a}. Its two required generators are \texttt{F\allowbreak{}U\allowbreak{}L\allowbreak{}L\allowbreak{}\_\allowbreak{}P\allowbreak{}R\allowbreak{}O\allowbreak{}F\allowbreak{}I\allowbreak{}L\allowbreak{}E\allowbreak{}.\allowbreak{}p\allowbreak{}y} (\texttt{3\allowbreak{}f\allowbreak{}4\allowbreak{}4\allowbreak{}7\allowbreak{}d\allowbreak{}6\allowbreak{}f\allowbreak{}9\allowbreak{}f\allowbreak{}0\allowbreak{}b\allowbreak{}b\allowbreak{}6\allowbreak{}9\allowbreak{}5\allowbreak{}7\allowbreak{}2\allowbreak{}6\allowbreak{}8\allowbreak{}8\allowbreak{}1\allowbreak{}2\allowbreak{}4\allowbreak{}d\allowbreak{}e\allowbreak{}0\allowbreak{}d\allowbreak{}5\allowbreak{}0\allowbreak{}9\allowbreak{}1\allowbreak{}2\allowbreak{}8\allowbreak{}e\allowbreak{}2\allowbreak{}a\allowbreak{}5\allowbreak{}6\allowbreak{}a\allowbreak{}5\allowbreak{}5\allowbreak{}f\allowbreak{}6\allowbreak{}e\allowbreak{}1\allowbreak{}7\allowbreak{}2\allowbreak{}e\allowbreak{}d\allowbreak{}8\allowbreak{}5\allowbreak{}6\allowbreak{}7\allowbreak{}c\allowbreak{}5\allowbreak{}2\allowbreak{}7\allowbreak{}0\allowbreak{}d\allowbreak{}f\allowbreak{}6\allowbreak{}4\allowbreak{}c}) and \texttt{S\allowbreak{}H\allowbreak{}I\allowbreak{}F\allowbreak{}T\allowbreak{}\_\allowbreak{}E\allowbreak{}N\allowbreak{}G\allowbreak{}I\allowbreak{}N\allowbreak{}E\allowbreak{}.\allowbreak{}p\allowbreak{}y} (\texttt{b\allowbreak{}6\allowbreak{}2\allowbreak{}7\allowbreak{}5\allowbreak{}a\allowbreak{}7\allowbreak{}c\allowbreak{}3\allowbreak{}8\allowbreak{}7\allowbreak{}d\allowbreak{}b\allowbreak{}3\allowbreak{}d\allowbreak{}f\allowbreak{}9\allowbreak{}b\allowbreak{}6\allowbreak{}b\allowbreak{}6\allowbreak{}7\allowbreak{}e\allowbreak{}a\allowbreak{}c\allowbreak{}2\allowbreak{}a\allowbreak{}1\allowbreak{}1\allowbreak{}4\allowbreak{}0\allowbreak{}f\allowbreak{}1\allowbreak{}2\allowbreak{}0\allowbreak{}d\allowbreak{}3\allowbreak{}e\allowbreak{}b\allowbreak{}d\allowbreak{}7\allowbreak{}7\allowbreak{}4\allowbreak{}e\allowbreak{}2\allowbreak{}2\allowbreak{}5\allowbreak{}f\allowbreak{}4\allowbreak{}9\allowbreak{}3\allowbreak{}5\allowbreak{}1\allowbreak{}4\allowbreak{}4\allowbreak{}6\allowbreak{}c\allowbreak{}a\allowbreak{}7\allowbreak{}3\allowbreak{}8\allowbreak{}0\allowbreak{}2\allowbreak{}5}) in that same directory. Run \texttt{python\ -B\ FULL\_PROFILE.py\ -\/-target\ .12666414\ -\/-h-add\ .000675\ -\/-steps\ 256\ -\/-depth\ 5\ -\/-source-root\ \textless{}unpacked-proof-root\textgreater{}\ -\/-output\ \textless{}fresh-file\textgreater{}}. The unpacked root contains \texttt{r\allowbreak{}e\allowbreak{}s\allowbreak{}e\allowbreak{}a\allowbreak{}r\allowbreak{}c\allowbreak{}h\allowbreak{}/}. The three immutable inputs remain the J.4 core/cover/contact files at their exact hashes. The outer engine reconstructs all original tilings, saved positive polynomials and true denominators, five admission windows, actual E49/h prices and every original physical complement. The matcher independently rechecks each refined child through that engine, proves the applicable L16a shortcut or enumerates the full candidate set, verifies every source shape and price, and records the exact worst rational value. Its coverage digest is \texttt{8\allowbreak{}b\allowbreak{}e\allowbreak{}6\allowbreak{}a\allowbreak{}f\allowbreak{}4\allowbreak{}5\allowbreak{}f\allowbreak{}a\allowbreak{}a\allowbreak{}2\allowbreak{}f\allowbreak{}3\allowbreak{}e\allowbreak{}8\allowbreak{}b\allowbreak{}2\allowbreak{}e\allowbreak{}4\allowbreak{}1\allowbreak{}0\allowbreak{}a\allowbreak{}0\allowbreak{}3\allowbreak{}e\allowbreak{}e\allowbreak{}f\allowbreak{}6\allowbreak{}e\allowbreak{}3\allowbreak{}a\allowbreak{}d\allowbreak{}a\allowbreak{}5\allowbreak{}a\allowbreak{}c\allowbreak{}2\allowbreak{}a\allowbreak{}f\allowbreak{}3\allowbreak{}3\allowbreak{}0\allowbreak{}0\allowbreak{}d\allowbreak{}5\allowbreak{}f\allowbreak{}4\allowbreak{}4\allowbreak{}5\allowbreak{}2\allowbreak{}0\allowbreak{}5\allowbreak{}a\allowbreak{}1\allowbreak{}8\allowbreak{}c\allowbreak{}c\allowbreak{}d\allowbreak{}7\allowbreak{}4\allowbreak{}f\allowbreak{}9\allowbreak{}9}.

The analytic physical derivation is also retained at \texttt{r\allowbreak{}e\allowbreak{}s\allowbreak{}e\allowbreak{}a\allowbreak{}r\allowbreak{}c\allowbreak{}h\allowbreak{}/\allowbreak{}l\allowbreak{}o\allowbreak{}c\allowbreak{}a\allowbreak{}l\allowbreak{}\_\allowbreak{}c\allowbreak{}y\allowbreak{}c\allowbreak{}l\allowbreak{}e\allowbreak{}0\allowbreak{}8\allowbreak{}4\allowbreak{}/\allowbreak{}n\allowbreak{}e\allowbreak{}a\allowbreak{}r\allowbreak{}\_\allowbreak{}r\allowbreak{}e\allowbreak{}c\allowbreak{}o\allowbreak{}r\allowbreak{}d\allowbreak{}/\allowbreak{}F\allowbreak{}I\allowbreak{}N\allowbreak{}A\allowbreak{}L\allowbreak{}.\allowbreak{}m\allowbreak{}d}, with the cap/count extensions in its \texttt{R\allowbreak{}E\allowbreak{}S\allowbreak{}E\allowbreak{}A\allowbreak{}R\allowbreak{}C\allowbreak{}H\allowbreak{}.\allowbreak{}m\allowbreak{}d}. The maximum is formed from actual saturated descriptors in the whole fixed bank; first-short selection supplies the three known faces in L12. This is not the separate \texttt{G\allowbreak{}A\allowbreak{}I\allowbreak{}N\allowbreak{}\_\allowbreak{}B\allowbreak{}A\allowbreak{}N\allowbreak{}K\allowbreak{}.\allowbreak{}p\allowbreak{}y} detector-selection rule. The earlier two-step 48-band .126660 source, the cubic full-column/low-radical comparison and any broader \(\fitinline{h}\) experiment are alternatives, not premises of this proof. There is no .12666415 conclusion here.

The complete proof since OpenAI and properly derived alternatives remain retained. Original credits and terms include Sil's paired canonical source, Joseph Bebel's exact contact and continuum count, Akash Levy's genuine numerator, Hailey Collet's signed reflection/probe/return and the other inherited authors. The actual physical reductions preserve SovereignSteak's draft PR43 head \texttt{c\allowbreak{}8\allowbreak{}3\allowbreak{}8\allowbreak{}a\allowbreak{}7\allowbreak{}5\allowbreak{}0\allowbreak{}0\allowbreak{}7\allowbreak{}9\allowbreak{}f\allowbreak{}7\allowbreak{}f\allowbreak{}8\allowbreak{}6\allowbreak{}a\allowbreak{}5\allowbreak{}7\allowbreak{}3\allowbreak{}d\allowbreak{}f\allowbreak{}9\allowbreak{}c\allowbreak{}9\allowbreak{}f\allowbreak{}7\allowbreak{}2\allowbreak{}9\allowbreak{}b\allowbreak{}8\allowbreak{}b\allowbreak{}c\allowbreak{}9\allowbreak{}2\allowbreak{}d\allowbreak{}1\allowbreak{}2\allowbreak{}5\allowbreak{}e}, its contributor notices and Codex-assistance disclosure; no proposed donor endpoint affirmation is assumed. Mathematical/Lean catalogue contracts were consulted for discovery and exact scope, not used instead of the source proof above. No declaration, finite receipt or model review constitutes a new end-to-end formal verification. The retained \emph{Mathematical Writing for AI Agents} attribution and Sepia/CWR terms remain.

\hypertarget{appendix-m}{}

\hypertarget{appendix-m.-same-zero-ordinary-substitution-and-the-current-native-return}{%
\section{Appendix M. Same-zero ordinary substitution and the current native return}\label{appendix-m.-same-zero-ordinary-substitution-and-the-current-native-return}}

Retain the complete OpenAI/Collet proof through L, including its explicitly retained natural selected-cell input H6/SN9, original proper keep, six units, fixed inducing collection, imprimitive Euler holes and actual physical reciprocal. We prove the strict ordinary conclusion at \[\fitmath{
 c_M=\frac{63333}{500000}=.126666,\qquad
 t_M=1-c_M=\frac{436667}{500000}=.873334,\qquad
 h=h_*+\frac3{5000},\quad .815924503\le h\le.815924505<.816.
}\tag{M0}\] For each fixed positive presented modulus this means every ordinary character and every finite height in the open half-plane \(\fitinline{\Re s>t_M}\), with only the actual principal pole excepted. The new analytic source is proved below, separately from its finite physical comparison. The retained input status is unchanged; there is no boundary, growing-modulus uniformity or new end-to-end formal verification assertion.

\hypertarget{m.1.-a-new-undivided-two-ordinary-source}{%
\subsection{M.1. A new undivided two-ordinary source}\label{m.1.-a-new-undivided-two-ordinary-source}}

Put \[\fitmath{
 \begin{gathered}
 f(y)=\frac{114}{25}(.3888-y)^2+
                \frac{101}{2}(.3888-y)^3-10^{-7},\\
 W_{85}(n,y)=\frac{1-n-y}{2}+f(y),\qquad
 u=\frac1{50000},\quad H=1-6u=\frac{24997}{25000},\quad\rho=10^{-9}.
 \end{gathered}                                                
}\tag{M1}\] Every strictly smaller width than \(\fitinline{W_{85}}\) is admitted by the original positive two-ordinary marked energy \textbf{before lower-mark division}, for upper shape labels \[\fitmath{
              .3735\le y\le.3888,\qquad1.115\le n/y\le1.155.
}\tag{M2}\] This includes every componentwise shortening under these labels, even when the actual shorter ratio leaves M2. The two fixed rays, their independent heights and normalized first log tests, original whole prime slots and all coefficient normalizers remain. This is a positive regular-core source with its subsequent complete raw restoration; it is not a bound for a new positive long signed arm.

The previously proved source L.1 is \[\fitmath{
 \begin{gathered}
 W_o(t,s)=A-t/2-Bs+\beta_o(s),\qquad
 A=(1+g_0)/2,\quad B=(1+2g_1)/2,\\
 \beta_o(s)=.00017+5.5(s-.38145)^2,
 \qquad .373\le s\le.387,\quad1.11\le t/s\le1.16,\\
 .05805911<g_0<.05805913,\qquad .07599618<g_1<.07599620.
 \end{gathered}                                                
}\tag{M3}\] Use its full componentwise-shortening statement. In the numerical and derivative comparisons below we may replace \(\fitinline{W_o}\) by the conservative lower envelope \(\fitinline{W_o^-=A_- -t/2-B_+s+\beta_o(s)}\), where \(\fitinline{A_-=(1+.05805911)/2}\) and \(\fitinline{B_+=(1+2(.07599620))/2}\). Thus an occurrence of \(\fitinline{B_+}\) in a clamp derivative differentiates this envelope, rather than asserting that \(\fitinline{B_+}\) is a lower bound for the true \(\fitinline{B}\).

Fix the actual cofactor in \(\fitinline{z=b^6c}\), retaining the unit and the full sixth-free support of \(\fitinline{c}\), and write its clean radius as \(\fitinline{Z=U^V}\). On a nonfloor regular bin let \(\fitinline{e\le\nu\le3/4}\) be the clean maximal squared prime-response exponent and the common buffered squared zero exponent. At one early cutoff \(\fitinline{T\in[1,3/2]}\), the actual B.5.1 convolution gives simultaneously, at one attaining zero, presentation and common height, \[\fitmath{
 \begin{gathered}
 r\le T+o(1),\quad r+m\ge T-o(1),\quad0\le m\le1/2+o(1),\\
 |M_r|^2\gg U^{V\nu r-o(1)},\qquad
 |S_m|^2\gg U^{V\nu m-o(1)}.
 \end{gathered}                                                
}\tag{M4}\] The original shorter ordinary factor has \(\fitinline{|A_{\rm short}|^2\ll U^{\nu y+o(1)}}\). Consequently on this same positive class, \[\fitmath{
 |A_{\rm long}A_{\rm short}Q|^2
 \ll U^{\nu(y-Vm)+o(1)}|A_{\rm long}S_mQ|^2.                 
}\tag{M5}\] The witness lower bound makes the division legitimate. The outside exponent in M5 has no extra factor \(\fitinline{1/2}\). If needed conjugate the whole \(\fitinline{S_m}\), including its test and height, to match the longer ordinary's moving sextic orientation. Its fixed twist lies in the early finite conjugation-stable dictionary. The original incoming height and the witness height are independent, as allowed by M3. The annular real powers and separated dyadic tests have B.5.1/H24a's original Fourier and seminorm controls.

This application has a fixed-dictionary requirement. Before slots or heights choose a finite bad set containing the conductors of all original, attaining and conjugate presentations. At a good prime of a regular sixth-free core the row valuation is one through five and the natural primitive conductor exponent is one; a fixed bad-set-supported twist cannot cancel it. Thus the auxiliary attaining rays remain nonprincipal on the proper regular keep. Cores without such a good prime belong to the original fixed-inducing collection. H8--H18 treats that collection on the \textbf{original two rectangles}, with their two original inducing characters, independent profile integrals and complete bad Euler factors, before M5. It is not dropped or assigned the means of the substituted product.

The auxiliary detector in M4 is on the clean natural \(\fitinline{c}\), with only its fixed deleted base and natural \(\fitinline{c}\)-zeros. No growing \(\fitinline{\operatorname{rad}(b)}\) conductor is inserted into an E(19) inverse inlet: its fixed deleted-base modulus remains below every live positive prime scale. The original variable \(\fitinline{b}\)-masks stay on \(\fitinline{A_{\rm long},A_{\rm short}}\); their deleted Euler factors cost \(\fitinline{U^\epsilon}\) in the buffered upper bound. M5 compares positive magnitudes and needs no agreement between the redundant holes of \(\fitinline{A_{\rm short}}\) and \(\fitinline{S_m}\). When M3 is applied to \(\fitinline{A_{\rm long}S_m}\), K.3 opens the original ordinary mask by its exact ideal M\"{o}bius identity and handles the original prime punctures at their true scales. Only then are the scalars bounded inside the positive norm. Each surviving mark is a genuine original whole-slot subproduct; a missing slot has no residual charge.

Canonical source matching is also literal. After whole conjugation use the attaining base ray \(\fitinline{a}\) with relative multiplier \(\fitinline{\mathbf1_q\overline a}\); at every good slot the product of \(\fitinline{a(p)\psi_c(p)}\) and \(\fitinline{\overline{a(p)}}\) is \(\fitinline{\psi_c(p)}\). Keep the outside unit, fixed reference coefficient, actual annular profile and adjusted relative frequency. B.6 thus supplies this source call; it is not an arbitrary row-dependent coefficient estimate.

For live width \(\fitinline{L=W_{85}(n,y)-u+O(\mathrm{mesh})}\), take a genuine subproduct at the earlier admitted absolute quota \(\fitinline{VW_o(t,s)}\), with \(\fitinline{t\ge n/V,s\ge m}\), and pay the complement by its actual upper response. The new excess is \[\fitmath{
 C_{\rm sub}(V;m,t,s)
       =\nu(y-Vm)+e[L-VW_o(t,s)]_+.                           
}\tag{M6}\] The first term is charged even if positive. A ceiling \(\fitinline{s>m}\) changes the source quota, not the true witness exponent \(\fitinline{V\nu m}\). The native two-fourth source at mean quota \(\fitinline{(1-t-s)/2}\) is a separate available supplier on the same regular keep. Its two individual quotas are positive and below their native caps in the uses below.

\hypertarget{m.2.-the-exhaustive-witness-split-and-finite-induction}{%
\subsection{M.2. The exhaustive witness split and finite induction}\label{m.2.-the-exhaustive-witness-split-and-finite-induction}}

Use the genuine E(19) quota \[\fitmath{
 q(w)=\frac{1-w}{2}+\frac{11w}{16}-\frac7{24}
                  -\frac5{16}\{7w^4(w-2/3)\}^{1/5}
                 \quad(3/4\le w<7/9).                       
}\tag{M7}\] Beyond \(\fitinline{7/9}\) use the native continuation \(\fitinline{(1-w)/2}\); never extrapolate the radical improvement. The ordinary native fourth comes from the actual positive family \(\fitinline{G(1-s,2+10s,3/20-1000s/1143)}\), with positive \(\fitinline{s}\) fixed early relative to strict decrements. The inverse retains its own two native gates and fixed annular test. B(211)'s smaller \(\fitinline{6\kappa}\) face is not substituted for this ordinary fourth. A dying quota uses its unmarked face and allocated collar.

On each closed upper \(\fitinline{y}\)-cell choose rational \(\fitinline{r_0,m_0,w_p}\) and set \(\fitinline{T=r_0+m_0}\). The split at this \textbf{one cutoff} is exhaustive:

\begin{itemize}
\tightlist
\item
  If \(\fitinline{r\ge1}\), B.5(204) gives \(\fitinline{D_l=5/6-(5/6-\nu)T}\).
\item
  If \(\fitinline{r_0\le r<1}\), the marked inverse gives \(\fitinline{D_r=\nu r_0+e q(r_0)}\). Its continuation is monotone through the curved/native/unmarked regimes, and \(\fitinline{T>1.12}\) separates the curved inverse call from its taper.
\item
  If \(\fitinline{r<r_0}\), the same saturated pair gives \(\fitinline{m\ge m_0-o(1)}\). Use M6 until the ordinary count alone pays.
\end{itemize}

Writing \(\fitinline{P(V)=V/2-n}\), the first two costs are \[\fitmath{
 C_{r,l}(V)=\nu y+e[L-P(V)/2]-VD_{r,l}/2.                     
}\tag{M8}\] Choose \(\fitinline{w_p}\) so the analogous ordinary count, with deficit \(\fitinline{\nu w_p+e q(w_p)}\), pays the whole response quadrilateral. Then \(\fitinline{m_1=\max(m_0,w_p/2)}\). Only \(\fitinline{[m_0,m_1]}\) requires M6; \textbf{all \(\fitinline{m\ge m_1}\)} use this paid ordinary escape. The tested ranges are \[\fitmath{
 \begin{gathered}
 .4167821\le k\le.4999827,\quad .75\le r_0\le.7776934,\\
 .3738319\le m_0\le.3888467,\quad .375\le m_1\le.3888467.
 \end{gathered}                                                
}\tag{M9}\] Both changed moment thresholds \(\fitinline{r_0,w_p}\) are below \(\fitinline{7/9}\); actual longer witnesses still require the native continuation.

The old supplier is used only at genuinely admitted caps \[\fitmath{
 s=\max\{m,.373,n/(1.16V)\},\qquad t=\max\{n/V,1.11s\},
                    \qquad s\le.387.                        
}\tag{M10}\] At \(\fitinline{V=H}\), the finite cells first verify \(\fitinline{1.11\le(n/H)/\max(m_0,.373)\le1.16}\), over their full original cap ratio. Both ceiling penalties in M10 are retained. For the conservative envelope in M3, an unclamped active M6 has derivative \(\fitinline{V[e(B_+-\beta_o')-\nu]<0}\); a fixed larger second ceiling gives \(\fitinline{-V\nu}\). On the first-coordinate lower clamp, \(\fitinline{Vm>y}\) because the original ratio is at least 1.115. The worst response there is \(\fitinline{\nu=e}\), with active derivative \(\fitinline{eV[B_++1.11/2-1-\beta_o'(m)]>0}\). Where its positive part vanishes M6 decreases. Thus its largest value is bounded by the two actual endpoint costs \(\fitinline{m_0,m_1}\), with the proper ceiling penalties, even when \(\fitinline{m_0}\) is outside the clamped region. There are 929 cells using this supplier and 295 using the native mean, whose cost decreases in \(\fitinline{m}\). Nothing presumes an old source extending to \(\fitinline{m=1/2}\).

For the new source set \(\fitinline{b_j=.3888H^j}\). There are 335 descending bands down to .3735, clipping the last. On the starting strip \(\fitinline{[.3888,.3888/H]}\), \(\fitinline{f<0}\), so the two native fourths already prove M1. In a later band, the low-response branch at \(\fitinline{V=H}\) uses the \textbf{preceding} new source at \(\fitinline{(n/H,y/H)}\), with exact leftover \[\fitmath{
 T_{\rm low}(y)=\frac{1-H}{2}+f(y)-Hf(y/H)-u.                 
}\tag{M11}\] Since \(\fitinline{f''>0}\), it decreases in \(\fitinline{y}\). For a closed cell \(\fitinline{[a,b]}\) of width \(\fitinline{1/80000}\), put \[\fitmath{
 k=10^{-7}\left\lfloor10^7\frac{u-3\rho}{T_{\rm low}(a)}
                                               \right\rfloor.
}\tag{M12}\] All \(\fitinline{0\le e\le k}\), including no-zero floor and empty/tiny bins, cost at most \(\fitinline{u-3\rho}\). The high branch uses only L's old source or native fourths. It never assumes the new source in its current band.

For \(\fitinline{k\le e\le.5002,e\le\nu\le3/4}\), M8 is affine in \(\fitinline{(\nu,e)}\), and M6 is the maximum of two affine functions. The four vertices of this quadrilateral therefore suffice. The count and initial unclamped costs increase in \(\fitinline{y}\), since \(\fitinline{\nu/e-1/2+f'>0}\); the final clamped cost is convex, so its two \(\fitinline{y}\)-endpoints suffice. The historical cutoff certificate fixes all 1224 cells and their \(\fitinline{10^{-7}}\)-grid choices. The independent exact rational replay identified in M.8 verifies those fixed choices with upward fifth-power rounding and rational zero in every positive part. It reevaluates the low and count branches and all 39,168 substitution prices: both \(\fitinline{y}\)-endpoints, both upper-ratio endpoints, all four response vertices and both \(\fitinline{m}\)-endpoints, with both ceiling penalties. The resulting exact reserves after the designated \(\fitinline{\rho}\), for low, inverse, amplified-long, substitution and final substitution endpoint, exceed \[\fitmath{
 2.0000\cdot10^{-9},\quad3.0040\cdot10^{-9},\quad
 2.9295\cdot10^{-5},\quad3.0113\cdot10^{-9},\quad
 3.2974\cdot10^{-9}.                                         
}\tag{M13}\] These are exact finite exponent inequalities from that replay, distinct from the later physical budget. The last endpoint's exact minimum is \(\fitinline{6594927754987/2000000000000000000000=3.2974638774935\cdot10^{-9}}\); the historical certificate's affected decimal spellings are not used as rational proof values.

Every radius is covered. Below \(\fitinline{H}\), the low branch's absolute capacity has derivative \(\fitinline{1/2+f(s)-sf'(s)}\), \(\fitinline{s=y/V}\); it is positive and, multiplied by .5002, is below \(\fitinline{5/6}\). At \(\fitinline{s=.3888}\), switch to the larger native mean, whose absolute slope is \(\fitinline{1/2}\). On \(\fitinline{[.98,H]}\), the three old-cap derivatives in M10 are \[\fitmath{
 A_- -B_+m+\beta_o(m),\quad
 A_--(B_++1.11/2)m+\beta_o(m),\quad
 A_-+\beta_o(s)-s\beta_o'(s).
}\] They lie in \(\fitinline{(0,.61)}\). When the second ceiling reaches .387, the native mean is larger: the seam margin \(\fitinline{g_1(.387)-g_0/2-\beta_o(.387)>4154291/10^{11}}\). Including the exterior \(\fitinline{\nu(y-Vm)}\), every substitution piece has slope magnitude below \(\fitinline{.75(.39)+.5002(.61)<.6<5/6}\); positive parts and seams preserve the upper bound. Counts M8 have decreasing slopes of magnitude below \(\fitinline{5/6}\). Thus all extend against \[\fitmath{
                   G(V)=u-\rho+\frac56(H-V)_+.                
}\tag{M14}\] For \(\fitinline{V\ge H}\), freeze the \textbf{entire} earlier admitted M10 pair at \(\fitinline{H}\), including its first-coordinate clamp: \(\fitinline{s_H=\max\{m,.373,n/(1.16H)\}}\), \(\fitinline{t_H=\max\{n/H,1.11s_H\}}\). This pair dominates the actual normalized lengths. The absolute quota increases while \(\fitinline{\nu(y-Vm)}\) decreases. Use the native supplier where that was the admitted pair. Freeze the low pair at \(\fitinline{(n/H,y/H)}\) likewise. Above \(\fitinline{V=1}\), continue the longer ordinary at its genuine fixed relative native quota \(\fitinline{1/2-n}\), of absolute width \(\fitinline{V(1/2-n)}\); the tested short exponent remains \(\fitinline{\nu y}\), not \(\fitinline{\nu Vy}\). The ordinary escape has this same continuation.

Response \(\fitinline{e>.5002}\) uses H6/SN9 at .5001, with fixed .0001 gap and volume reserve exceeding .0305. On \(\fitinline{[.875,.98]}\) the native/unmarked cost \(\fitinline{\frac34[(1-V)/2+\max f-u]_+}\) has .98 reserve over .00816 and slope at most \(\fitinline{3/8}\). Negative formal quotas are replaced by the legitimate bounded-length nonprincipal unmarked fourth. Below .875 the unmarked cost \(\fitinline{3(W_{\max}-u)/4}\) has reserve over .0243; below .08 the original counted tiny-core bound has over .0319. Restoring \(\fitinline{b^6c}\) uses \(\fitinline{(H-V)_+-(H-M)_+\le M-V=6\log_U Nb}\); the \(\fitinline{5/6}\) slope pays the full cofactor count, without imposing \(\fitinline{(b,c)=1}\).

\hypertarget{m.3.-raw-restoration-signed-rows-and-source-heights}{%
\subsection{M.3. Raw restoration, signed rows and source heights}\label{m.3.-raw-restoration-signed-rows-and-source-heights}}

Keep H25/H26's two orientations and coefficients: the first artificial divisor has no reciprocal weight; the second M\"{o}bius divisor retains its matching reciprocal normalizer; the separate norm identities yield \(\fitinline{NC\,ND_1}\). Keep full sublattice multiplicity, seed/harmonic masses, outside unit, masks and projectors. M5 occurs only inside the positive regular norm after the original fixed-row split. H8's signed centered comparison still has the common four-coordinate minimum, both independent integrals and both original principal inducing means.

The actual root bounds are \[\fitmath{
 W_{\max}=\frac{212543940077}{2000000000000},\qquad
 A_{0,\max}=\frac{9189319}{10000000}<.919.
}\tag{M15}\] The original gates \(\fitinline{H>4(.139+W_{\max})}\) and \(\fitinline{5/6+.139>A_0+3u}\) have reserves over .01879 and .05334. The equal-product companion is below .71 and uses the proved arbitrary bounded-length nonprincipal unmarked fourth. The second-zero gates are \(\fitinline{5u<1,A_0+4u<1}\). For the first zero and tails, C.10 requires literally root volume at most \(\fitinline{U}\), with its fixed norm-cap multiple at most \(\fitinline{U^2}\); M15 supplies these, without an obsolete .903 convenience cap. Counted tails retain \(\fitinline{W_BU^{-31}}\) per seed against \(\fitinline{O(U^8W_B)}\) possibilities. H.5(A)--(C) consequently restore the full original uncentered product. Only the later physical operation divides lower marks.

Fix the finite 335-stage dependency list, cutoff/length banks and strict decrements before the live slots. Each early maximum slot width is below the earlier mesh times the positive lower core radius and below all decrements. The combined inner rounding, bin, radial and height losses can be below \(\fitinline{10^{-12}}\), within M13. The finite ray dictionary is fixed next; genuine fixed-inducing rows retain their original signed treatment. Shortened tests remain anchored at actual scales. Compact log-Fourier envelopes on bounded ratio groups, with original dyadic/polylogarithmic masses, give shortening uniformity; a compressed test is not a single bump at a fictitious upper scale. Separate normalized first log tests control height derivatives. Fix all envelope degrees and Sobolev orders before the one positive physical height exponent, itself preceding every later positive buffered-grid index. This proves M1--M2 at the full original height and coefficient scope.

\hypertarget{m.4.-price-the-current-native-count-on-its-original-admission}{%
\subsection{M.4. Price the current native count on its original admission}\label{m.4.-price-the-current-native-count-on-its-original-admission}}

The physical improvement also needs a coupled comparison of the actual current native count. PR141's \texttt{E\allowbreak{}x\allowbreak{}a\allowbreak{}c\allowbreak{}t\allowbreak{}R\allowbreak{}e\allowbreak{}t\allowbreak{}u\allowbreak{}r\allowbreak{}n\allowbreak{}P\allowbreak{}o\allowbreak{}l\allowbreak{}y\allowbreak{}n\allowbreak{}o\allowbreak{}m\allowbreak{}i\allowbreak{}a\allowbreak{}l\allowbreak{}.\allowbreak{}l\allowbreak{}e\allowbreak{}a\allowbreak{}n} identifies \texttt{.\allowbreak{}n\allowbreak{}a\allowbreak{}t\allowbreak{}i\allowbreak{}v\allowbreak{}e} and \texttt{.\allowbreak{}n\allowbreak{}a\allowbreak{}t\allowbreak{}i\allowbreak{}v\allowbreak{}e\allowbreak{}L\allowbreak{}o\allowbreak{}w} with the lift of \texttt{(\allowbreak{}t\allowbreak{}o\allowbreak{}Q\allowbreak{}\ \allowbreak{}.\allowbreak{}C\allowbreak{})\allowbreak{}.\allowbreak{}N\allowbreak{}q\allowbreak{}/\allowbreak{}J\allowbreak{}q}; \texttt{C\allowbreak{}o\allowbreak{}u\allowbreak{}n\allowbreak{}t\allowbreak{}S\allowbreak{}c\allowbreak{}h\allowbreak{}e\allowbreak{}m\allowbreak{}e\allowbreak{}.\allowbreak{}l\allowbreak{}e\allowbreak{}a\allowbreak{}n} supplies their original windows \[\fitmath{
 [15/64,9/16]\times[9/32,1/2],\qquad
 [15/64,1/2]\times[3/16,1/2].                                 
}\tag{M16}\] Every assigned original native rectangle is contained in one of them. The PR141 windows use its .15-cap fork. Although the uncapped \(\fitinline{J,N}\) are independent of that cap, their admission is not. We therefore prove all the real guards directly at the \textbf{available} positive source \(\fitinline{B=1-s,L=2+10s,c_G=3/20-1000s/1143}\), \(\fitinline{s=10^{-260}}\), using PR129 \texttt{G\allowbreak{}u\allowbreak{}a\allowbreak{}r\allowbreak{}d\allowbreak{}s\allowbreak{}.\allowbreak{}l\allowbreak{}e\allowbreak{}a\allowbreak{}n}'s universal real formulas, not its target-typed \texttt{D\allowbreak{}a\allowbreak{}t\allowbreak{}a\allowbreak{}.\allowbreak{}l\allowbreak{}e\allowbreak{}a\allowbreak{}n} transfer or smaller .1342894 cap. No old target geometry is imported. These pinned declaration reads identify the existing analytic conditions; no new Lean execution is asserted.

For \(\fitinline{v=5/6-d}\), write \[\fitmath{
 \begin{gathered}
 D_0=3-p-4p/L,\quad A_0'=2-4p/L,\quad C_0=1-2pB/L,\\
 J=vD_0+d(1-p)A_0',\qquad
 N=(1-d)J+vd(1-p)C_0,\qquad R_{\rm nat}=N/J.
 \end{gathered}                                                
}\tag{M17}\] Here is the actual admission proof. Set \(\fitinline{U_0=(1-p)C_0}\), \(\fitinline{Z_0=1-C_0}\), \(\fitinline{T=J+dU_0}\), \(\fitinline{D_J=D_0J}\), \(\fitinline{R_N=A_0'T+(Z_0-p)J}\), and \(\fitinline{M_N=D_0T-R_N}\). For \(\fitinline{\eta_0=10^{-7},s_0=.027,m_0^{\rm nat}=(1-(9/2)s_0)/2=.43925}\), put \[\fitmath{
 \begin{aligned}
 m_-^{\rm nat}&=\max\{.3+\eta_0,[B-L(c_G-\eta_0)]/2+\eta_0\},\\
 m_+^{\rm nat}&=\min\{.44-\eta_0,m_0^{\rm nat}-\eta_0,
                         [B-L(s_0+\eta_0)]/2-\eta_0\}.
 \end{aligned}
}\] The fifteen cleared real guards, with derivatives in \(\fitinline{p}\), are \[\fitmath{
 \begin{gathered}
 D_0,\ A_0',\ C_0,\ J,\ T-J,\ (12499/10000)J-T,\\
 R_N-(23/37)D_J,\ D_J-R_N,\ M_N-m_-^{\rm nat}D_J,
       \ m_+^{\rm nat}D_J-M_N,\\
 N-(1-2dm_0^{\rm nat}-2dps_0)J,\quad
 U_0\partial_pJ-(\partial_pU_0)J,\\
 .44-m_0^{\rm nat}-\eta_0,\quad
 B-2m_0^{\rm nat}-L(s_0+\eta_0)-2\eta_0,\quad1.
 \end{gathered}
}\] Exact tensor-Bernstein bounds on the two complete windows M16 give all fifteen strictly positive, with common minima at least \(\fitinline{585388183593/10^{15}}\) and \(\fitinline{7499/10^7}\), respectively. These are the literal full-source guards, including both length ceilings, the native crossing and amplitude derivative. \texttt{N\allowbreak{}A\allowbreak{}T\allowbreak{}I\allowbreak{}V\allowbreak{}E\allowbreak{}\_\allowbreak{}A\allowbreak{}D\allowbreak{}M\allowbreak{}I\allowbreak{}S\allowbreak{}S\allowbreak{}I\allowbreak{}O\allowbreak{}N\allowbreak{}.\allowbreak{}p\allowbreak{}y} and its exact receipt supply every individual bound. The optional .1342894 restriction actually fails the ninth guard on both windows; no inference from it is used. The positive source \(\fitinline{G(B,L,c_G)}\) and these real guard inequalities admit precisely the count M17 on the existing support.

At the actual geometry put \(\fitinline{\ell=11/3-4t_M}\), \(\fitinline{y_t=h-2\ell}\), \(\fitinline{x_t=1+\ell-h}\). The old and physical capacities are \[\fitmath{
 C_o(D)=\frac{\min(D,2y_t-D)}{6\kappa_*},\qquad
 C_p(D)=\min\{.10D,[2y_t-(2-B)D]/L\},\quad\kappa_*=2\alpha_+-1.
}\] Only the conservative FB2--FB6 pricing uses \(\fitinline{\kappa_*}\); the actual moment parameter is still \(\fitinline{2\beta-1}\). The returned exponent is exactly \[\fitmath{
 E=\frac12+\frac{13h}{75}-\frac{y_t}{2}-\frac\ell2-.09D
       +d/2+dp(\ell-C/2)+\frac{3D}{4}R_{\rm nat}.              
}\tag{M18}\] Clear the current margin \(\fitinline{M=(t_M-E_{\rm base})J-3DN/4}\), at both directed \(\fitinline{h}\) endpoints and the four original conductor edges \(\fitinline{D=.005,.76,.76,h}\). The count is independent of \(\fitinline{h}\); the minimum of affine capacities enters E with a nonpositive coefficient. Hence E is convex in \(\fitinline{h}\), and these endpoints suffice. Original conductor convexity fills both closed intervals with the same admitted count envelope.

On each original admitted density/amplitude box tensor-Bernstein bounds give \(\fitinline{M\ge M_-}\) and \(\fitinline{0<J_-\le J\le J_+}\). Divide by \(\fitinline{J_+}\) for \(\fitinline{M_-\ge0}\), by \(\fitinline{J_-}\) for \(\fitinline{M_-<0}\). This prices the entire current box. It is not a separate saved-contact shift plus an independent worst slope. In the final cover it pays 24 pieces at margin at least \(\fitinline{1554468211/10^{15}}\). Its use is confined to the original native/nativeLow assignments M16; the other saved and paired routes retain their own payers.

\hypertarget{m.5.-the-actual-finite-bank-and-continuous-length-comparison}{%
\subsection{M.5. The actual finite bank and continuous length comparison}\label{m.5.-the-actual-finite-bank-and-continuous-length-comparison}}

Use L.3's finite maximum of \textbf{observed} raw count deficits, with \(\fitinline{N_b=1024}\) spacings and 1025 cutoffs prescribed before observing a row. At every \(\fitinline{T_j=t_0-\gamma_j}\), \(\fitinline{\gamma_j=j\gamma/N_b}\), B.5.1 gives the simultaneous actual inverse and ordinary witness. On its finite descriptor class let \(\fitinline{G}\) be the maximum of the certified nonnegative count deficits. An attained qualifying descriptor bounds that class by \(\fitinline{U^{R-G+o(1)}}\); a finite union preserves the exponent. No two independent counts or raw moments are multiplied.

If \(\fitinline{3G/4}\) pays the local excess, use the original numerator fourth. This includes all six original fixed native/unmarked/affine continuations, too-large short witnesses, and the last amplified witness if every inverse is long. Otherwise select the first short inverse, index \(\fitinline{j}\). Its current inverse face, the preceding amplified inverse and the current ordinary face are lower bounds for the \textbf{same} \(\fitinline{G}\). The auxiliary inverse is not the physical reciprocal. Same-zero saturation gives \(\fitinline{|P|\ll U^{-dy/2+o(1)}|PS_j|}\); simultaneous Cauchy on the original row subset supplies half of this count saving. The pre-division source in M1 is applied on that same physical subset with its independent original heights.

Here is the complete finite price. On \(\fitinline{[a,b]\times[f_p,e_p]}\) put \[\fitmath{
 \begin{gathered}
 B_+=\max(0,B_{\rm loc})+10^{-8},\quad v_-=5/6-b,
 \quad\epsilon_a=10^{-8},\\
 k_-=1-e_p(1+2g_{1,+}),\quad k_+=1-f_p(1+2g_{1,-}),\quad
 m_\pm=t_{0,\pm}/3,\quad M_0=m_++\epsilon_a.
 \end{gathered}
}\] Here \(\fitinline{B_{\rm loc}}\) is the directed top excess after the genuine paired E49 reserve and current geometry charge; \(\fitinline{D_-}\) is the actual directed retained-conductor lower bound. Round upwards to \(\fitinline{10^{-9}}\) \[\fitmath{
 \gamma\ge\frac{4B_+}{3D_-v_-},\qquad
 \lambda_x\ge\frac{4B_+}{3D_-ak_-},\qquad
 \lambda_y\ge\frac{2B_+}{3D_-ak_-}.                           
}\tag{M19}\] At cutoff \(\fitinline{j}\), the actual ordinary length lies in the outward interval \(\fitinline{[m_- -\gamma_j-\lambda_x-2\epsilon_a,m_++\lambda_y+\epsilon_a]}\). Define \(\fitinline{\gamma_{-1}=0}\) and \[\fitmath{
 \begin{gathered}
 F_j=-k_+(\gamma_j+3\epsilon_a),\qquad
 H_j=(v_-/b)[\gamma_{j-1}-3\epsilon_a]_+,\\
 X_j(y)=\max\{F_j,\ k_z(y-M_0)+\max(H_j,2k_-[y-M_0]_+)\},
 \end{gathered}                                                
}\tag{M20}\] where \(\fitinline{k_z=k_+}\) below \(\fitinline{M_0}\) and \(\fitinline{k_-}\) above it. For the true upper source pair use \(\fitinline{n_{\rm ref}=\max(n,n_0)}\), \(\fitinline{n_0=1-y_t/D}\), and \(\fitinline{\widehat y=\max(y,.3735)}\). Its conservative additional width above the actual tangent is \[\fitmath{
 \Delta_{85}(y)=g_{1,-}\widehat y-g_{0,+}/2+f(\widehat y)
                    -\frac{1+2g_{1,+}}2(\widehat y-y).         
}\tag{M21}\] In particular the lower clamp has its true debit. Saturation and \(\fitinline{U^{-dy/2}}\) still use the actual \(\fitinline{y}\). The already fully summed positive packet mass J9a compensates the numerator cap change before profile/height suprema; there is no invented shorter numerator. The whole condition is \[\fitmath{
 f_p\Delta_{85}(y)+\frac{X_j(y)}2
                     -e_p\epsilon_a/2-\frac{B_+}{D_-a}>0.     
}\tag{M22}\] The checker establishes \(\fitinline{\Delta_{85}\ge0}\) before replacing actual \(\fitinline{p}\) by \(\fitinline{f_p}\).

The source cubic is convex. On 1244 final cells the strict L16 derivative inequalities, with \(\fitinline{-\Delta_{85}'}\) in place of the quadratic slope, show that the minimum is the floor/sloping crossover and the worst cutoff is \(\fitinline{j=1024}\). Their source interval is below the zero of \(\fitinline{\Delta_{85}'}\), and its increasing derivative bounds the adverse slope at the left endpoint. The other 135 cells check all 1025 cutoffs, every endpoint, kink and lower clamp, and every stationary point of the four affine-plus-cubic pieces. Their increasing derivatives have at most one zero, enclosed by 70 exact rational bisections. Evaluating both bracket endpoints and subtracting ten times their separation gives a rigorous lower value, since the absolute derivative is below ten. Below the clamp the price is piecewise affine. These are continuous interval and response bounds, not sampled lengths or amplitudes.

\hypertarget{m.6.-whole-cover-and-all-physical-complementary-cases}{%
\subsection{M.6. Whole cover and all physical complementary cases}\label{m.6.-whole-cover-and-all-physical-complementary-cases}}

The exact cover reconstructs all 57 original rectangles and their four closed conductor-edge tilings: 18 old saved pieces, 24 current-native pieces, and 19393 paired exterior pieces. Retain the original E27 and I32 wing. In addition to L21 the only new upward window is \(\fitinline{[2961/8192,741/2048]\times[1023/2048,1/2]}\); all six recorded upward windows have all ten E29 tensor-Bernstein guards at all four actual \(\fitinline{g_0,g_1}\) corners, with common lower bound exceeding \(\fitinline{113992769667/10^{15}}\). Thus each used count has its original genuine admission.

There are 1361 coarse central cells, six four-child refinements and 1379 final central cells. Every child rechecks its full source, contour, conductor, dilation and fixed-continuation contract. None is removed by an unrecorded restriction. There are no unpaid cells. The minimum of M22, above the distinct \(\fitinline{10^{-8}}\) budget in \(\fitinline{B_+}\), is at least \[\fitmath{
 \frac{24447707}{10^{15}}=2.4447707\cdot10^{-8}.
}\tag{M23}\] Every final cell has \(\fitinline{D_-a>.3}\); its additional original \(\fitinline{Z}\)-exponent reserve is at least \(\fitinline{916789/125000000000000>7.3343\cdot10^{-9}}\). The exact worst is original cell52 at the last cutoff, actual \(\fitinline{y}\) approximately .3759332815586077; the sign uses its recorded exact rational value.

Across the complete cover, \[\fitmath{
 \begin{gathered}
 .42474914701958<n_{\rm ref}<.42552334362808,\\
 .37258992644642<y<.37928196596324,\qquad
 1.1201828365683<n_{\rm ref}/\widehat y<1.1387982988704.
 \end{gathered}                                                
}\tag{M24}\] All raised pairs are strictly inside M2's ratio range. Actual shortenings are included. The unchanged prime pools have total slot width greater than \(\fitinline{1/5}\); M15 gives positive exhaustion and a genuine greedy whole-slot prefix. No lower mark is transferred to a different character, and no squarefree/radical/amplitude row filter is imposed.

The saved, current-native, E49 exterior and other fixed/moving route margins are respectively at least \[\fitmath{
 \frac{5923437039}{25000000000000},\quad
 \frac{1554468211}{10^{15}},\quad
 \frac{28053869}{5\cdot10^{14}},\quad
 \frac{2000001}{2\cdot10^{14}}.                               
}\tag{M25}\] The fixed-continuation reserve exceeds .00238985. The original six continuations retain their native/affine/unmarked faces, proper caps and ordinary escape. The conductor is at least .815753978, density is in \(\fitinline{[.40283203125,.44482421875]}\), and \(\fitinline{p\ge.494384765625}\). The lower contour damping is at least \(\fitinline{29/40960}\); also \(\fitinline{(1+d)/2-.10>.51}\) and \(\fitinline{(1+d)/2+.27<1}\). The original absolute contour motion, complete \(\fitinline{w}\)-inversion and full error-tuple sum pay both inverse-length tails before amplitude subdivision. Large numerator dilation uses its authentic \(\fitinline{dv/4}\) saving.

In particular preserve the actual rowwise packet inequality \[\fitmath{
 \sum_{\mathbf e}|A_i(\mathbf e)|N_{\mathbf e}^{-1/2}
      \max(1,L_iN_{\mathbf e})^{d/2}
 \le C U^\epsilon\prod_{p\in\mathcal E}Y_p^{-4/25}.             
}\tag{M26}\] This is B.7(221--222)/\texttt{P\allowbreak{}o\allowbreak{}s\allowbreak{}i\allowbreak{}t\allowbreak{}i\allowbreak{}v\allowbreak{}e\allowbreak{}L\allowbreak{}a\allowbreak{}b\allowbreak{}e\allowbreak{}l\allowbreak{}P\allowbreak{}a\allowbreak{}c\allowbreak{}k\allowbreak{}e\allowbreak{}t\allowbreak{}.\allowbreak{}l\allowbreak{}e\allowbreak{}a\allowbreak{}n} 89--96 with the actual external error weights and full tuple sum; a subpower factor alone does not bound \texttt{e\allowbreak{}r\allowbreak{}r\allowbreak{}o\allowbreak{}r\allowbreak{}P\allowbreak{}r\allowbreak{}o\allowbreak{}d\allowbreak{}u\allowbreak{}c\allowbreak{}t\allowbreak{}S\allowbreak{}i\allowbreak{}z\allowbreak{}e}. Apply M26 once in the original positive mixed norm, keeping the exterior \(\fitinline{U^{-dv/2}}\). Retain the two raw denominator identities, \(\fitinline{NC\,ND_1}\), the unweighted first artificial divisor and matching second reciprocal normalizer, all seeds, projectors, outside unit and full lattice multiplicity.

The physical cap is .10, declared .11, under the separate GT4--GT7 positive family, not the old .075 cap. The actual M0 geometry directly satisfies every generic \texttt{G\allowbreak{}e\allowbreak{}o\allowbreak{}m\allowbreak{}e\allowbreak{}t\allowbreak{}r\allowbreak{}y\allowbreak{}.\allowbreak{}P\allowbreak{}a\allowbreak{}r\allowbreak{}a\allowbreak{}m\allowbreak{}e\allowbreak{}t\allowbreak{}e\allowbreak{}r\allowbreak{}s} and \texttt{B\allowbreak{}o\allowbreak{}u\allowbreak{}n\allowbreak{}d\allowbreak{}s\allowbreak{}C\allowbreak{}o\allowbreak{}r\allowbreak{}e} field: \(\fitinline{1/6<\ell<1/5}\), \(\fitinline{\ell<x_t<y_t<1/2}\), \(\fitinline{2x_t-y_t-\ell>0}\), \(\fitinline{5\ell<1}\), \(\fitinline{x_t+y_t=1-\ell}\), and the exact signal boundary. In particular \(\fitinline{49h/150-\ell-1/50>.07}\). It also satisfies \(\fitinline{h+F<\min(1,2y_t)}\), \(\fitinline{y_t<.76<h}\), \(\fitinline{\ell/(h+F)>.2}\), \(\fitinline{2y_t-(2-B)(h+F)>0}\), and \(\fitinline{h/2-2\ell<.10h}\). The literal \(\fitinline{h<.816}\) consumer is retained. J's old narrower target corridor is not assumed.

The floor margin after \(\fitinline{10^{-5}+10^{-12}}\) exceeds .0029505249; the small-row margin exceeds .07153867; principal \(\fitinline{w,z}\) margins exceed .02139723 and .0001098741. Short numerators, bounded scales, null terms, the five nonprincipal units, genuine principal residues, zero frequencies, finite Euler factors and compensators retain their distinct original B/H payers. No pole is charged to a count gain.

Choose positive source cost \(\fitinline{s=10^{-260}}\), pad \(\fitinline{F=10^{-258}}\), baseline tolerance \(\fitinline{D_{\rm src}.t\le F/4}\), and \(\fitinline{\Delta=F}\). The literal cube pad is then \(\fitinline{F/2+2D_{\rm src}.t\le F}\). Charge \(\fitinline{15s/16+2F}\) separately. The capacities are 1-Lipschitz and conductor variation is below two; this extends the whole partition to \(\fitinline{D\le h+F}\). M23--M25 pay these charges with over \(\fitinline{3\cdot10^{-9}}\) left. Allocate other outer finite losses below \(\fitinline{10^{-10}}\), and inner source decrements also below M13's strict margins. The direct terminal reserve is not an uncharged width decrement.

Keep the physical reciprocal at B.8's retained \(\fitinline{s=a+16e+it}\); \(\fitinline{s+w}\) is its separate Euler-product guard. During truncation retain the original global \(\fitinline{\beta+20e}\) line and move only the retained segment. Three physical contour frequencies and two ordinary/witness heights keep their independent weights. Source choices and the finite bank precede slots; common compact log-Fourier envelopes and seminorm orders precede the single positive physical height exponent and all subsequent buffered-grid indices. Scale/dilation suprema use their actual profile envelope or \(\fitinline{xW'(x)}\); height suprema use separate normalized first log tests. Tails are chosen against these fixed polynomial envelopes. A favorable detector term never replaces the whole physical inverse.

\hypertarget{m.7.-original-cube-probe-and-mellin-return}{%
\subsection{M.7. Original cube, probe and Mellin return}\label{m.7.-original-cube-probe-and-mellin-return}}

The preceding complete estimate supplies I51 at M0 for every original \texttt{A\allowbreak{}c\allowbreak{}t\allowbreak{}u\allowbreak{}a\allowbreak{}l\allowbreak{}B\allowbreak{}u\allowbreak{}f\allowbreak{}f\allowbreak{}e\allowbreak{}r\allowbreak{}e\allowbreak{}d\allowbreak{}S\allowbreak{}e\allowbreak{}l\allowbreak{}e\allowbreak{}c\allowbreak{}t\allowbreak{}i\allowbreak{}o\allowbreak{}n}, in the original order \[\fitmath{
 \forall\eta\ \exists\tau_\eta>0\ \forall n>0\ \exists C_n,
                         \qquad Z\ \hbox{eventually large}.   
}\tag{M27}\] Choose the common source reserve \(\fitinline{0<\sigma=P_{\rm src}.\sigma<10^{-9}}\), below all applicable original I51 transport reserves. The literal reciprocal-containing cube, divided by its original nonzero normalizer, is bounded by \(\fitinline{C_n Z^{\beta-k_{\rm phys}-\sigma}}\). Baseline data and the finite source choices are fixed before the target \(\fitinline{\eta}\). Keep I52's original nonzero normalizer and cutoff, now with \(\fitinline{\ell,x_t,y_t}\) from M0. The directly constructed generic geometry lets \texttt{h\allowbreak{}i\allowbreak{}g\allowbreak{}h\allowbreak{}\_\allowbreak{}o\allowbreak{}f\allowbreak{}\_\allowbreak{}c\allowbreak{}u\allowbreak{}b\allowbreak{}e} give the actual physical probe and signal Gram exponent \(\fitinline{t_M-k_{\rm phys}}\).

For a contrary \(\fitinline{\beta>t_M}\), choose \(\fitinline{\omega<(\beta-t_M)/4}\) and \(\fitinline{0<m<\min(\beta-t_M-\omega,\sigma,\beta-.87)}\), using that same \(\fitinline{\sigma=P_{\rm src}.\sigma}\). H47 gives a Mellin transform holomorphic on \(\fitinline{\Re s>\beta-m}\), whereas on the absolute line it is the original Gaussian times \(\fitinline{H_\eta/L_\eta}\). Its original numerator satisfies \(\fitinline{|H_\eta-1|\le1/2}\) on \(\fitinline{\Re s\ge.87}\), so no denominator zero there can be cancelled. The definition of the true zero supremum supplies a zero above \(\fitinline{\beta-m/2}\), contradicting holomorphy; no attained supremum is assumed. Original norm pullback and fixed Euler restoration prove M0 for all ordinary and imprimitive characters of every fixed positive presented modulus and all finite heights. This completes the H45--H47 return, with the actual principal pole excluded.

\hypertarget{m.8.-exact-evidence-and-attribution}{%
\subsection{M.8. Exact evidence and attribution}\label{m.8.-exact-evidence-and-attribution}}

The complete predecessor is the actual closed Local084 packet member \texttt{C\allowbreak{}L\allowbreak{}O\allowbreak{}S\allowbreak{}E\allowbreak{}D\allowbreak{}\_\allowbreak{}P\allowbreak{}R\allowbreak{}O\allowbreak{}O\allowbreak{}F\allowbreak{}.\allowbreak{}m\allowbreak{}d}, SHA-256 \texttt{3\allowbreak{}c\allowbreak{}e\allowbreak{}3\allowbreak{}1\allowbreak{}7\allowbreak{}6\allowbreak{}9\allowbreak{}2\allowbreak{}8\allowbreak{}9\allowbreak{}4\allowbreak{}2\allowbreak{}0\allowbreak{}b\allowbreak{}4\allowbreak{}b\allowbreak{}0\allowbreak{}8\allowbreak{}7\allowbreak{}a\allowbreak{}9\allowbreak{}1\allowbreak{}9\allowbreak{}e\allowbreak{}3\allowbreak{}d\allowbreak{}2\allowbreak{}f\allowbreak{}6\allowbreak{}e\allowbreak{}1\allowbreak{}7\allowbreak{}5\allowbreak{}b\allowbreak{}2\allowbreak{}1\allowbreak{}c\allowbreak{}9\allowbreak{}0\allowbreak{}4\allowbreak{}7\allowbreak{}2\allowbreak{}6\allowbreak{}f\allowbreak{}7\allowbreak{}e\allowbreak{}3\allowbreak{}d\allowbreak{}2\allowbreak{}2\allowbreak{}5\allowbreak{}4\allowbreak{}9\allowbreak{}0\allowbreak{}6\allowbreak{}9\allowbreak{}1\allowbreak{}4\allowbreak{}f\allowbreak{}2\allowbreak{}9\allowbreak{}0\allowbreak{}7}. The successor delta reconstructs those exact bytes as \texttt{h\allowbreak{}i\allowbreak{}s\allowbreak{}t\allowbreak{}o\allowbreak{}r\allowbreak{}y\allowbreak{}/\allowbreak{}p\allowbreak{}r\allowbreak{}o\allowbreak{}o\allowbreak{}f\allowbreak{}\_\allowbreak{}v\allowbreak{}e\allowbreak{}r\allowbreak{}s\allowbreak{}i\allowbreak{}o\allowbreak{}n\allowbreak{}s\allowbreak{}/\allowbreak{}v\allowbreak{}i\allowbreak{}r\allowbreak{}t\allowbreak{}u\allowbreak{}a\allowbreak{}l\allowbreak{}/\allowbreak{}3\allowbreak{}c\allowbreak{}e\allowbreak{}3\allowbreak{}1\allowbreak{}7\allowbreak{}6\allowbreak{}9\allowbreak{}2\allowbreak{}8\allowbreak{}9\allowbreak{}4\allowbreak{}2\allowbreak{}0\allowbreak{}b\allowbreak{}4\allowbreak{}b\allowbreak{}0\allowbreak{}8\allowbreak{}7\allowbreak{}a\allowbreak{}9\allowbreak{}1\allowbreak{}9\allowbreak{}e\allowbreak{}3\allowbreak{}d\allowbreak{}2\allowbreak{}f\allowbreak{}6\allowbreak{}e\allowbreak{}1\allowbreak{}7\allowbreak{}5\allowbreak{}b\allowbreak{}2\allowbreak{}1\allowbreak{}c\allowbreak{}9\allowbreak{}0\allowbreak{}4\allowbreak{}7\allowbreak{}2\allowbreak{}6\allowbreak{}f\allowbreak{}7\allowbreak{}e\allowbreak{}3\allowbreak{}d\allowbreak{}2\allowbreak{}2\allowbreak{}5\allowbreak{}4\allowbreak{}9\allowbreak{}0\allowbreak{}6\allowbreak{}9\allowbreak{}1\allowbreak{}4\allowbreak{}f\allowbreak{}2\allowbreak{}9\allowbreak{}0\allowbreak{}7\allowbreak{}.\allowbreak{}m\allowbreak{}d}. Append M to that complete proof; retain L's earlier rational result and proof as historical evidence. The controller-local \texttt{c\allowbreak{}y\allowbreak{}c\allowbreak{}l\allowbreak{}e\allowbreak{}0\allowbreak{}8\allowbreak{}4\allowbreak{}/\allowbreak{}i\allowbreak{}n\allowbreak{}t\allowbreak{}e\allowbreak{}g\allowbreak{}r\allowbreak{}a\allowbreak{}t\allowbreak{}i\allowbreak{}o\allowbreak{}n\allowbreak{}/\allowbreak{}p\allowbreak{}r\allowbreak{}o\allowbreak{}o\allowbreak{}f\allowbreak{}-\allowbreak{}3\allowbreak{}c\allowbreak{}e\allowbreak{}3\allowbreak{}1\allowbreak{}7\allowbreak{}6\allowbreak{}9\allowbreak{}2\allowbreak{}8\allowbreak{}9\allowbreak{}4\allowbreak{}.\allowbreak{}m\allowbreak{}d} is a working copy, not an asserted archive path.

The analytic source is \texttt{r\allowbreak{}e\allowbreak{}s\allowbreak{}e\allowbreak{}a\allowbreak{}r\allowbreak{}c\allowbreak{}h\allowbreak{}/\allowbreak{}l\allowbreak{}o\allowbreak{}c\allowbreak{}a\allowbreak{}l\allowbreak{}\_\allowbreak{}c\allowbreak{}y\allowbreak{}c\allowbreak{}l\allowbreak{}e\allowbreak{}0\allowbreak{}8\allowbreak{}5\allowbreak{}/\allowbreak{}d\allowbreak{}i\allowbreak{}r\allowbreak{}e\allowbreak{}c\allowbreak{}t\allowbreak{}/\allowbreak{}R\allowbreak{}E\allowbreak{}S\allowbreak{}E\allowbreak{}A\allowbreak{}R\allowbreak{}C\allowbreak{}H\allowbreak{}.\allowbreak{}m\allowbreak{}d}, SHA-256 \texttt{b\allowbreak{}8\allowbreak{}e\allowbreak{}2\allowbreak{}2\allowbreak{}2\allowbreak{}f\allowbreak{}6\allowbreak{}b\allowbreak{}1\allowbreak{}a\allowbreak{}f\allowbreak{}2\allowbreak{}9\allowbreak{}c\allowbreak{}8\allowbreak{}d\allowbreak{}0\allowbreak{}6\allowbreak{}e\allowbreak{}5\allowbreak{}1\allowbreak{}d\allowbreak{}e\allowbreak{}a\allowbreak{}8\allowbreak{}0\allowbreak{}8\allowbreak{}5\allowbreak{}6\allowbreak{}0\allowbreak{}a\allowbreak{}d\allowbreak{}4\allowbreak{}1\allowbreak{}c\allowbreak{}8\allowbreak{}c\allowbreak{}a\allowbreak{}e\allowbreak{}4\allowbreak{}0\allowbreak{}1\allowbreak{}1\allowbreak{}6\allowbreak{}d\allowbreak{}a\allowbreak{}e\allowbreak{}b\allowbreak{}7\allowbreak{}a\allowbreak{}f\allowbreak{}a\allowbreak{}0\allowbreak{}7\allowbreak{}5\allowbreak{}c\allowbreak{}2\allowbreak{}0\allowbreak{}d\allowbreak{}a\allowbreak{}5\allowbreak{}3\allowbreak{}7}, with historical candidate checker \texttt{6\allowbreak{}0\allowbreak{}3\allowbreak{}2\allowbreak{}8\allowbreak{}e\allowbreak{}0\allowbreak{}2\allowbreak{}d\allowbreak{}9\allowbreak{}d\allowbreak{}b\allowbreak{}3\allowbreak{}9\allowbreak{}e\allowbreak{}d\allowbreak{}f\allowbreak{}3\allowbreak{}5\allowbreak{}c\allowbreak{}2\allowbreak{}3\allowbreak{}d\allowbreak{}4\allowbreak{}8\allowbreak{}3\allowbreak{}d\allowbreak{}a\allowbreak{}f\allowbreak{}b\allowbreak{}8\allowbreak{}d\allowbreak{}d\allowbreak{}c\allowbreak{}0\allowbreak{}e\allowbreak{}5\allowbreak{}3\allowbreak{}4\allowbreak{}e\allowbreak{}b\allowbreak{}8\allowbreak{}f\allowbreak{}7\allowbreak{}f\allowbreak{}a\allowbreak{}3\allowbreak{}f\allowbreak{}d\allowbreak{}d\allowbreak{}4\allowbreak{}9\allowbreak{}1\allowbreak{}6\allowbreak{}7\allowbreak{}c\allowbreak{}7\allowbreak{}e\allowbreak{}3\allowbreak{}9\allowbreak{}8\allowbreak{}5\allowbreak{}5\allowbreak{}d} and \texttt{C\allowbreak{}E\allowbreak{}R\allowbreak{}T\allowbreak{}I\allowbreak{}F\allowbreak{}I\allowbreak{}C\allowbreak{}A\allowbreak{}T\allowbreak{}E\allowbreak{}.\allowbreak{}j\allowbreak{}s\allowbreak{}o\allowbreak{}n} \texttt{c\allowbreak{}2\allowbreak{}b\allowbreak{}6\allowbreak{}4\allowbreak{}a\allowbreak{}3\allowbreak{}2\allowbreak{}c\allowbreak{}7\allowbreak{}f\allowbreak{}6\allowbreak{}d\allowbreak{}7\allowbreak{}5\allowbreak{}3\allowbreak{}2\allowbreak{}1\allowbreak{}6\allowbreak{}c\allowbreak{}8\allowbreak{}7\allowbreak{}d\allowbreak{}d\allowbreak{}8\allowbreak{}5\allowbreak{}0\allowbreak{}8\allowbreak{}8\allowbreak{}2\allowbreak{}a\allowbreak{}e\allowbreak{}a\allowbreak{}d\allowbreak{}1\allowbreak{}8\allowbreak{}6\allowbreak{}f\allowbreak{}6\allowbreak{}5\allowbreak{}6\allowbreak{}b\allowbreak{}c\allowbreak{}b\allowbreak{}5\allowbreak{}2\allowbreak{}a\allowbreak{}3\allowbreak{}4\allowbreak{}e\allowbreak{}6\allowbreak{}6\allowbreak{}f\allowbreak{}0\allowbreak{}4\allowbreak{}5\allowbreak{}8\allowbreak{}6\allowbreak{}f\allowbreak{}f\allowbreak{}b\allowbreak{}9\allowbreak{}1\allowbreak{}5}. These retain M9--M15's 335 stages, 1224 cells and fixed cutoff choices. In the generator's final old-source clamp, \texttt{m\allowbreak{}a\allowbreak{}x\allowbreak{}(\allowbreak{}0\allowbreak{},\allowbreak{}H\allowbreak{}*\allowbreak{}O\allowbreak{}L\allowbreak{}D\allowbreak{}R\allowbreak{}L\allowbreak{}O\allowbreak{}*\allowbreak{}s\allowbreak{}-\allowbreak{}n\allowbreak{})\allowbreak{}/\allowbreak{}2} can produce floating zero. All 929 old-source upper-endpoint reserve strings are therefore retained as historical output, rather than described as wholly exact rational inequalities.

The exact supplement is \texttt{r\allowbreak{}e\allowbreak{}s\allowbreak{}e\allowbreak{}a\allowbreak{}r\allowbreak{}c\allowbreak{}h\allowbreak{}/\allowbreak{}l\allowbreak{}o\allowbreak{}c\allowbreak{}a\allowbreak{}l\allowbreak{}\_\allowbreak{}c\allowbreak{}y\allowbreak{}c\allowbreak{}l\allowbreak{}e\allowbreak{}0\allowbreak{}8\allowbreak{}5\allowbreak{}/\allowbreak{}r\allowbreak{}e\allowbreak{}t\allowbreak{}u\allowbreak{}r\allowbreak{}n\allowbreak{}\_\allowbreak{}r\allowbreak{}e\allowbreak{}v\allowbreak{}i\allowbreak{}e\allowbreak{}w\allowbreak{}/\allowbreak{}V\allowbreak{}E\allowbreak{}R\allowbreak{}I\allowbreak{}F\allowbreak{}Y\allowbreak{}\_\allowbreak{}D\allowbreak{}I\allowbreak{}R\allowbreak{}E\allowbreak{}C\allowbreak{}T\allowbreak{}\_\allowbreak{}E\allowbreak{}N\allowbreak{}D\allowbreak{}P\allowbreak{}O\allowbreak{}I\allowbreak{}N\allowbreak{}T\allowbreak{}S\allowbreak{}.\allowbreak{}p\allowbreak{}y}, SHA-256 \texttt{e\allowbreak{}d\allowbreak{}e\allowbreak{}e\allowbreak{}0\allowbreak{}d\allowbreak{}e\allowbreak{}3\allowbreak{}6\allowbreak{}4\allowbreak{}f\allowbreak{}8\allowbreak{}f\allowbreak{}9\allowbreak{}3\allowbreak{}6\allowbreak{}9\allowbreak{}d\allowbreak{}8\allowbreak{}0\allowbreak{}4\allowbreak{}1\allowbreak{}a\allowbreak{}3\allowbreak{}1\allowbreak{}8\allowbreak{}7\allowbreak{}0\allowbreak{}1\allowbreak{}2\allowbreak{}b\allowbreak{}c\allowbreak{}2\allowbreak{}3\allowbreak{}1\allowbreak{}0\allowbreak{}2\allowbreak{}5\allowbreak{}6\allowbreak{}7\allowbreak{}5\allowbreak{}0\allowbreak{}1\allowbreak{}d\allowbreak{}e\allowbreak{}4\allowbreak{}d\allowbreak{}e\allowbreak{}e\allowbreak{}5\allowbreak{}1\allowbreak{}6\allowbreak{}9\allowbreak{}2\allowbreak{}9\allowbreak{}a\allowbreak{}1\allowbreak{}2\allowbreak{}b\allowbreak{}b\allowbreak{}4\allowbreak{}9\allowbreak{}0\allowbreak{}7}, and \texttt{D\allowbreak{}I\allowbreak{}R\allowbreak{}E\allowbreak{}C\allowbreak{}T\allowbreak{}\_\allowbreak{}E\allowbreak{}X\allowbreak{}A\allowbreak{}C\allowbreak{}T\allowbreak{}\_\allowbreak{}E\allowbreak{}N\allowbreak{}D\allowbreak{}P\allowbreak{}O\allowbreak{}I\allowbreak{}N\allowbreak{}T\allowbreak{}S\allowbreak{}.\allowbreak{}j\allowbreak{}s\allowbreak{}o\allowbreak{}n}, SHA-256 \texttt{b\allowbreak{}7\allowbreak{}8\allowbreak{}0\allowbreak{}f\allowbreak{}9\allowbreak{}b\allowbreak{}6\allowbreak{}d\allowbreak{}7\allowbreak{}8\allowbreak{}c\allowbreak{}1\allowbreak{}2\allowbreak{}e\allowbreak{}1\allowbreak{}a\allowbreak{}f\allowbreak{}7\allowbreak{}6\allowbreak{}b\allowbreak{}2\allowbreak{}6\allowbreak{}e\allowbreak{}5\allowbreak{}0\allowbreak{}2\allowbreak{}c\allowbreak{}a\allowbreak{}1\allowbreak{}2\allowbreak{}9\allowbreak{}0\allowbreak{}2\allowbreak{}0\allowbreak{}d\allowbreak{}a\allowbreak{}c\allowbreak{}e\allowbreak{}b\allowbreak{}c\allowbreak{}4\allowbreak{}e\allowbreak{}c\allowbreak{}4\allowbreak{}a\allowbreak{}e\allowbreak{}d\allowbreak{}a\allowbreak{}a\allowbreak{}6\allowbreak{}e\allowbreak{}7\allowbreak{}e\allowbreak{}3\allowbreak{}8\allowbreak{}b\allowbreak{}c\allowbreak{}d\allowbreak{}3\allowbreak{}d\allowbreak{}e\allowbreak{}5\allowbreak{}3}. It reads the held certificate only for cutoff hints, without importing or running its generator, and recomputes the relevant M6/M8/M11 prices entirely as rational numbers, including the full M10 caps and \(\fitinline{H(1.11)s-n}\) penalty. It checks every vertex described above on all 929 old-source and 295 native cells. The largest discrepancy from an old reported upper reserve is \(\fitinline{26188241/(2\cdot10^{26})}\); every M13 threshold remains strict. The exact worst upper endpoint is cell905 (zero based), with the rational minimum given after M13. This supplement supplies the finite inequalities; M.1--M.3 supply the continuous-length monotonicity, finite-stage induction, original raw restoration and all-height analytic return. Its earlier supplier remains the separately proved L source, not the concurrent 085 controller cubic.

The complete physical derivation is \texttt{r\allowbreak{}e\allowbreak{}s\allowbreak{}e\allowbreak{}a\allowbreak{}r\allowbreak{}c\allowbreak{}h\allowbreak{}/\allowbreak{}l\allowbreak{}o\allowbreak{}c\allowbreak{}a\allowbreak{}l\allowbreak{}\_\allowbreak{}c\allowbreak{}y\allowbreak{}c\allowbreak{}l\allowbreak{}e\allowbreak{}0\allowbreak{}8\allowbreak{}5\allowbreak{}/\allowbreak{}n\allowbreak{}e\allowbreak{}a\allowbreak{}r\allowbreak{}\_\allowbreak{}r\allowbreak{}e\allowbreak{}c\allowbreak{}o\allowbreak{}r\allowbreak{}d\allowbreak{}/\allowbreak{}R\allowbreak{}E\allowbreak{}S\allowbreak{}E\allowbreak{}A\allowbreak{}R\allowbreak{}C\allowbreak{}H\allowbreak{}.\allowbreak{}m\allowbreak{}d}, SHA-256 \texttt{c\allowbreak{}2\allowbreak{}b\allowbreak{}8\allowbreak{}6\allowbreak{}d\allowbreak{}8\allowbreak{}3\allowbreak{}8\allowbreak{}8\allowbreak{}8\allowbreak{}0\allowbreak{}8\allowbreak{}d\allowbreak{}b\allowbreak{}9\allowbreak{}f\allowbreak{}3\allowbreak{}1\allowbreak{}1\allowbreak{}1\allowbreak{}e\allowbreak{}0\allowbreak{}d\allowbreak{}e\allowbreak{}5\allowbreak{}9\allowbreak{}c\allowbreak{}3\allowbreak{}0\allowbreak{}b\allowbreak{}4\allowbreak{}a\allowbreak{}7\allowbreak{}b\allowbreak{}8\allowbreak{}1\allowbreak{}0\allowbreak{}c\allowbreak{}3\allowbreak{}9\allowbreak{}5\allowbreak{}6\allowbreak{}3\allowbreak{}d\allowbreak{}a\allowbreak{}c\allowbreak{}2\allowbreak{}1\allowbreak{}8\allowbreak{}3\allowbreak{}f\allowbreak{}f\allowbreak{}a\allowbreak{}8\allowbreak{}e\allowbreak{}9\allowbreak{}2\allowbreak{}0\allowbreak{}0\allowbreak{}2\allowbreak{}3\allowbreak{}3\allowbreak{}8}. Its terminal receipt \texttt{F\allowbreak{}I\allowbreak{}N\allowbreak{}A\allowbreak{}L\allowbreak{}\_\allowbreak{}1\allowbreak{}2\allowbreak{}6\allowbreak{}6\allowbreak{}6\allowbreak{}6\allowbreak{}.\allowbreak{}j\allowbreak{}s\allowbreak{}o\allowbreak{}n} is \texttt{d\allowbreak{}0\allowbreak{}3\allowbreak{}6\allowbreak{}0\allowbreak{}9\allowbreak{}d\allowbreak{}5\allowbreak{}f\allowbreak{}6\allowbreak{}2\allowbreak{}e\allowbreak{}2\allowbreak{}2\allowbreak{}9\allowbreak{}f\allowbreak{}e\allowbreak{}a\allowbreak{}c\allowbreak{}d\allowbreak{}8\allowbreak{}0\allowbreak{}3\allowbreak{}0\allowbreak{}9\allowbreak{}a\allowbreak{}6\allowbreak{}0\allowbreak{}7\allowbreak{}8\allowbreak{}5\allowbreak{}b\allowbreak{}2\allowbreak{}d\allowbreak{}8\allowbreak{}5\allowbreak{}d\allowbreak{}6\allowbreak{}f\allowbreak{}1\allowbreak{}e\allowbreak{}a\allowbreak{}2\allowbreak{}1\allowbreak{}4\allowbreak{}1\allowbreak{}2\allowbreak{}8\allowbreak{}0\allowbreak{}e\allowbreak{}c\allowbreak{}9\allowbreak{}4\allowbreak{}a\allowbreak{}d\allowbreak{}2\allowbreak{}d\allowbreak{}f\allowbreak{}3\allowbreak{}c\allowbreak{}6\allowbreak{}e\allowbreak{}e\allowbreak{}f}, with coverage digest \texttt{c\allowbreak{}0\allowbreak{}4\allowbreak{}1\allowbreak{}1\allowbreak{}c\allowbreak{}d\allowbreak{}1\allowbreak{}2\allowbreak{}f\allowbreak{}c\allowbreak{}b\allowbreak{}9\allowbreak{}c\allowbreak{}9\allowbreak{}f\allowbreak{}4\allowbreak{}3\allowbreak{}c\allowbreak{}0\allowbreak{}2\allowbreak{}7\allowbreak{}8\allowbreak{}4\allowbreak{}d\allowbreak{}b\allowbreak{}5\allowbreak{}8\allowbreak{}a\allowbreak{}4\allowbreak{}4\allowbreak{}c\allowbreak{}3\allowbreak{}0\allowbreak{}c\allowbreak{}5\allowbreak{}6\allowbreak{}d\allowbreak{}e\allowbreak{}e\allowbreak{}2\allowbreak{}d\allowbreak{}a\allowbreak{}a\allowbreak{}5\allowbreak{}1\allowbreak{}c\allowbreak{}2\allowbreak{}d\allowbreak{}5\allowbreak{}2\allowbreak{}d\allowbreak{}3\allowbreak{}c\allowbreak{}2\allowbreak{}1\allowbreak{}9\allowbreak{}8\allowbreak{}0\allowbreak{}9\allowbreak{}a\allowbreak{}c\allowbreak{}5\allowbreak{}0}. The eight-file portable front is that note, receipt, \texttt{S\allowbreak{}O\allowbreak{}U\allowbreak{}R\allowbreak{}C\allowbreak{}E\allowbreak{}\_\allowbreak{}U\allowbreak{}S\allowbreak{}E\allowbreak{}\_\allowbreak{}F\allowbreak{}I\allowbreak{}N\allowbreak{}A\allowbreak{}L\allowbreak{}.\allowbreak{}j\allowbreak{}s\allowbreak{}o\allowbreak{}n}, \texttt{F\allowbreak{}I\allowbreak{}N\allowbreak{}A\allowbreak{}L\allowbreak{}\_\allowbreak{}C\allowbreak{}H\allowbreak{}E\allowbreak{}C\allowbreak{}K\allowbreak{}.\allowbreak{}p\allowbreak{}y}, \texttt{C\allowbreak{}O\allowbreak{}U\allowbreak{}P\allowbreak{}L\allowbreak{}E\allowbreak{}D\allowbreak{}\_\allowbreak{}E\allowbreak{}N\allowbreak{}G\allowbreak{}I\allowbreak{}N\allowbreak{}E\allowbreak{}.\allowbreak{}p\allowbreak{}y}, \texttt{D\allowbreak{}I\allowbreak{}R\allowbreak{}E\allowbreak{}C\allowbreak{}T\allowbreak{}\_\allowbreak{}S\allowbreak{}O\allowbreak{}U\allowbreak{}R\allowbreak{}C\allowbreak{}E\allowbreak{}\_\allowbreak{}C\allowbreak{}H\allowbreak{}E\allowbreak{}C\allowbreak{}K\allowbreak{}.\allowbreak{}p\allowbreak{}y}, \texttt{N\allowbreak{}A\allowbreak{}T\allowbreak{}I\allowbreak{}V\allowbreak{}E\allowbreak{}\_\allowbreak{}A\allowbreak{}D\allowbreak{}M\allowbreak{}I\allowbreak{}S\allowbreak{}S\allowbreak{}I\allowbreak{}O\allowbreak{}N\allowbreak{}.\allowbreak{}p\allowbreak{}y}, and \texttt{N\allowbreak{}A\allowbreak{}T\allowbreak{}I\allowbreak{}V\allowbreak{}E\allowbreak{}\_\allowbreak{}A\allowbreak{}D\allowbreak{}M\allowbreak{}I\allowbreak{}S\allowbreak{}S\allowbreak{}I\allowbreak{}O\allowbreak{}N\allowbreak{}.\allowbreak{}j\allowbreak{}s\allowbreak{}o\allowbreak{}n}. The matcher reconstructs every physical cell and source minimum; \texttt{FINAL\_CHECK.py\ -\/-source-root\ \textless{}unpacked-proof-root\textgreater{}\ -\/-recompute\ -\/-output\ \textless{}fresh-json\textgreater{}} checks original input identities, native support, all fifteen real full-source guards, geometry and byte-identical full replay. The native-guard receipt is \texttt{6\allowbreak{}c\allowbreak{}8\allowbreak{}1\allowbreak{}8\allowbreak{}7\allowbreak{}a\allowbreak{}c\allowbreak{}1\allowbreak{}7\allowbreak{}4\allowbreak{}5\allowbreak{}8\allowbreak{}7\allowbreak{}c\allowbreak{}6\allowbreak{}3\allowbreak{}1\allowbreak{}3\allowbreak{}8\allowbreak{}8\allowbreak{}f\allowbreak{}f\allowbreak{}a\allowbreak{}b\allowbreak{}5\allowbreak{}0\allowbreak{}f\allowbreak{}5\allowbreak{}6\allowbreak{}a\allowbreak{}b\allowbreak{}1\allowbreak{}a\allowbreak{}7\allowbreak{}6\allowbreak{}e\allowbreak{}9\allowbreak{}3\allowbreak{}0\allowbreak{}5\allowbreak{}f\allowbreak{}0\allowbreak{}9\allowbreak{}a\allowbreak{}3\allowbreak{}c\allowbreak{}7\allowbreak{}6\allowbreak{}f\allowbreak{}6\allowbreak{}8\allowbreak{}9\allowbreak{}a\allowbreak{}0\allowbreak{}e\allowbreak{}4\allowbreak{}d\allowbreak{}8\allowbreak{}7\allowbreak{}e\allowbreak{}0\allowbreak{}b\allowbreak{}e}. The unpacked root contains the exact \texttt{r\allowbreak{}e\allowbreak{}s\allowbreak{}e\allowbreak{}a\allowbreak{}r\allowbreak{}c\allowbreak{}h\allowbreak{}/} inputs. The physical receipt's retained \texttt{P\allowbreak{}A\allowbreak{}S\allowbreak{}S\allowbreak{}\_\allowbreak{}W\allowbreak{}H\allowbreak{}O\allowbreak{}L\allowbreak{}E\allowbreak{}\_\allowbreak{}R\allowbreak{}E\allowbreak{}T\allowbreak{}U\allowbreak{}R\allowbreak{}N\allowbreak{}\_\allowbreak{}C\allowbreak{}O\allowbreak{}N\allowbreak{}D\allowbreak{}I\allowbreak{}T\allowbreak{}I\allowbreak{}O\allowbreak{}N\allowbreak{}A\allowbreak{}L\allowbreak{}\_\allowbreak{}O\allowbreak{}N\allowbreak{}\_\allowbreak{}P\allowbreak{}R\allowbreak{}O\allowbreak{}F\allowbreak{}I\allowbreak{}L\allowbreak{}E} label is arithmetic scope; M.1--M.3 and the bound source discharge that analytic premise. Neither a finite receipt nor declaration read proves the analytic chain by itself.

The larger .1266665 run leaves 139 computed leaves unpaid; it supplies no extra digit. The .126670 width request is also unpaid by this source. The separate .127 and .12725 campaigns remain open. The independent 085 controller cubic and the 084 arithmetic column source are useful alternatives, not hidden premises here.

Preserve the complete working proof since OpenAI and properly derived alternatives, original OpenAI, Collet, Sil, Bebel, Levy, Gao, Liu and other contributor credits, SovereignSteak's retained draft PR43 head \texttt{c\allowbreak{}8\allowbreak{}3\allowbreak{}8\allowbreak{}a\allowbreak{}7\allowbreak{}5\allowbreak{}0\allowbreak{}0\allowbreak{}7\allowbreak{}9\allowbreak{}f\allowbreak{}7\allowbreak{}f\allowbreak{}8\allowbreak{}6\allowbreak{}a\allowbreak{}5\allowbreak{}7\allowbreak{}3\allowbreak{}d\allowbreak{}f\allowbreak{}9\allowbreak{}c\allowbreak{}9\allowbreak{}f\allowbreak{}7\allowbreak{}2\allowbreak{}9\allowbreak{}b\allowbreak{}8\allowbreak{}b\allowbreak{}c\allowbreak{}9\allowbreak{}2\allowbreak{}d\allowbreak{}1\allowbreak{}2\allowbreak{}5\allowbreak{}e}, its terms/notices and Codex-assistance disclosure. Preserve \emph{Mathematical Writing for AI Agents} attribution and Sepia/CWR terms. Mathematical and Lean catalogue entries guide discovery and source scope; neither their status nor declaration names establish applicability. The source-use receipt binds the exact original evidence. No public action or new formal verification is asserted.

\hypertarget{appendix-n}{}

\hypertarget{appendix-n.-reusing-the-completed-ordinary-source-with-a-clipped-lower-branch}{%
\section{Appendix N. Reusing the completed ordinary source with a clipped lower branch}\label{appendix-n.-reusing-the-completed-ordinary-source-with-a-clipped-lower-branch}}

Retain the complete proof through M, with its explicitly retained natural selected-cell input H6/SN9, original fixed presented modulus, proper keep, all six units, imprimitive Euler holes and actual physical reciprocal. The composition below is at \[\fitmath{
 c_N=\frac{2533329}{20000000}=.12666645,\qquad
 t_N=\frac{17466671}{20000000}=.87333355,\qquad
 h=h_*+\frac{63}{100000},\quad
 .815954503\le h\le.815954505<.816.                            
}\tag{N0}\] Its conclusion has the same strict open-half-plane, fixed positive presented-modulus and all-height scope as M, with only the actual principal pole excepted. The new source and full physical match are separate arguments. No boundary, growing-modulus uniformity, new upstream analytic admission or end-to-end formal verification is asserted.

\hypertarget{n.1.-the-completed-source-may-be-used-in-a-later-source-proof}{%
\subsection{N.1. The completed source may be used in a later source proof}\label{n.1.-the-completed-source-may-be-used-in-a-later-source-proof}}

Write \(\fitinline{Y=.3888,a=.3735,y_0=.3700}\), and put \[\fitmath{
 \begin{gathered}
 f_M(s)=4.56(Y-s)^2+50.5(Y-s)^3-10^{-7},\\
 F(y)=
 \begin{cases}
 4.59(Y-y)^2+50.5(Y-y)^3+.00031(Y-y)-10^{-7},&a\le y\le Y/H,\\
 F(a)+.155(a-y),&y_0\le y<a,
 \end{cases}\\
 W_N(n,y)=\frac{1-n-y}{2}+F(y),\qquad
 u=\frac1{50000},\quad H=1-6u=\frac{24997}{25000},\quad\rho=10^{-9}.
 \end{gathered}                                                
}\tag{N1}\] Every strictly smaller absolute width than \(\fitinline{W_N}\) is admitted by the original positive two-ordinary marked energy \textbf{before lower-mark division}, for upper shape labels \[\fitmath{
              .3700\le y\le.3888,\qquad1.118\le n/y\le1.15.  
}\tag{N2}\] All actual componentwise shortenings are included under those labels. The algebraic cubic is defined through Y/H for the starting strip and low-budget expression; the admitted upper labels in N2 still end at Y. The two ordinary coefficients and fixed rays retain their independent real heights, annular powers, normalized first log tests, original whole prime slots, natural zeros and true coefficient normalizers.

This uses completed M as a mathematical input: its source is \[\fitmath{
 W_M(t,s)=(1-t-s)/2+f_M(s),\qquad
 .3735\le s\le.3888,\quad1.115\le t/s\le1.155,                 
}\tag{N3}\] with full componentwise shortening and K.3/H.5 raw restoration. The old quadratic supplier L retains its own earlier role in M; the new high branch does not assume the source N1 currently being proved. Its low branch will use N1 only on strictly earlier, higher length bands. These are two ordered uses of completed sources, not a mixed inverse--ordinary moment assumption. On the common ratio range, the additional width at y=.376 is exactly \(\fitinline{.03(.0128)^2+.00031(.0128)=.0000088832}\).

The tangent and curved moment sources are the actual internal E(25), E(19) and I.3 inputs. E(25)'s 113 proved paired stages supply the genuine lower tangent below 3/4; it is not obtained by extending the curve there. K.3 preserves both original signed inducing rectangles before positive substitution, L proves the earlier supplier, and M proves N3 with its exact endpoint supplement. The full earlier source contracts and their original notices remain.

\hypertarget{n.2.-same-zero-substitution-with-both-genuine-ceilings}{%
\subsection{N.2. Same-zero substitution with both genuine ceilings}\label{n.2.-same-zero-substitution-with-both-genuine-ceilings}}

Fix the actual cofactor in \(\fitinline{z=b^6c}\), keeping the unit and sixth-free core c, without imposing \(\fitinline{(b,c)=1}\). Let its radius be \(\fitinline{U^V}\), and let \(\fitinline{0\le e\le\nu\le3/4}\) be the clean squared prime-response and common zero-band exponents. B.5.1 at one early cutoff T gives a simultaneous inverse and ordinary at the same attaining zero, presentation and allowed profile: \[\fitmath{
 r\le T+o(1),\quad r+m\ge T-o(1),\quad m\le1/2+o(1),\qquad
 |S_m|^2\gg U^{V\nu m-o(1)}.                                
}\tag{N4}\] The original shorter ordinary has squared upper bound \(\fitinline{U^{\nu y+o(1)}}\). M5 therefore gives the exact positive comparison \[\fitmath{
 |A_{\rm long}A_{\rm short}Q|^2
 \ll U^{\nu(y-Vm)+o(1)}|A_{\rm long}S_mQ|^2.                  
}\tag{N5}\] There is no additional half factor in this exponent. The auxiliary detector is on the clean natural core with only its fixed deleted base. The growing cofactor's redundant holes remain on the original ordinary factors and are opened by K.3's exact ordinary/prime identities at their true scales. They are not inserted into an E(19) M\"{o}bius inverse inlet. The substitute need not have the shorter factor's redundant holes. Whole conjugation, if needed, includes the witness test and height, and matches the original moving orientation through M's common relative prime multiplier. Fixed dictionary conductors are chosen early; regular nonprincipal cores and the original signed fixed-inducing collection keep M.1's separate treatment.

For a closed upper-length cell \(\fitinline{[b_0,b]}\), the tested ordinary endpoints satisfy \(\fitinline{m\ge b/H}\). At \(\fitinline{V=H}\) take the completed M source at \[\fitmath{
               s=\max(m,a),\qquad t=\max(n/H,1.115s).         
}\tag{N6}\] Since \(\fitinline{n/y\le1.15<1.155}\) and \(\fitinline{Hm\ge b\ge y}\), its upper-ratio condition holds. Use it only for \(\fitinline{s\le Y}\). Both ceilings enter its actual capacity \(\fitinline{HW_M(t,s)}\); raising s changes the capacity, not the witness lower response \(\fitinline{V\nu m}\).

Set \(\fitinline{L=W_N(n,y)-u}\), apart from allocated mesh loss. The exact old-source cost \(\fitinline{\nu(y-Hm)+e[L-HW_M(t,s)]_+}\) is largest in the first-coordinate ceiling penalty at \(\fitinline{n=1.118y}\). It is therefore bounded by \[\fitmath{
 \begin{aligned}
 C_o(y,m;\nu,e)={}&-\nu(Hm-y)\\
 &+e\left[\frac{1-H}{2}+\frac{Hs-y}{2}+F(y)-u-Hf_M(s)
                    +\frac{(1.115Hs-1.118y)_+}{2}\right]_+.
 \end{aligned}                                                  
}\tag{N7}\] The native two-fourth alternative is the same expression with \(\fitinline{s=m}\) and the \(\fitinline{f_M}\) and first-clamp terms removed. Its individual quotas are genuinely positive and below their native caps in the uses here. The ordinary native fourth comes from the full positive family \(\fitinline{G(1-s_0,2+10s_0,3/20-1000s_0/1143)}\), with positive \(\fitinline{s_0}\) fixed early relative to strict decrements; the inverse has its own two native inequalities. A smaller cap attached to a specialized old target is not substituted for that source. Near a vanishing quota keep the true unmarked face and its collar.

\hypertarget{n.3.-exhaustive-cutoff-tests-the-lower-kink-and-finite-induction}{%
\subsection{N.3. Exhaustive cutoff tests, the lower kink, and finite induction}\label{n.3.-exhaustive-cutoff-tests-the-lower-kink-and-finite-induction}}

For \(\fitinline{.734\le w\le1}\), use \(\fitinline{q(w)=(1-w)/2+J(w)}\), with \[\fitmath{
 J(w)=
 \begin{cases}
 .05805911-.07599620w,&.734\le w<3/4,\\
 11w/16-7/24-(5/16)\{7w^4(w-2/3)\}^{1/5},&3/4\le w<7/9,\\
 0,&7/9\le w\le1.
 \end{cases}                                                  
}\tag{N8}\] The first is a lower quota supplied by E(25)'s actual paired rectangles, the second by E(19). The count \(\fitinline{\nu w+eq(w)}\) increases on each piece when \(\fitinline{\nu\ge e}\); at 3/4 the lower tangent lies strictly below the curve, and the curved gain tends to zero at 7/9, where the native continuation applies. The largest tested inverse or ordinary threshold is .7776934\textless7/9. I20's strictly increasing deficit therefore leaves a uniform native collar before that endpoint; fix the collar and a finite bank of strictly decremented E(19) rectangles first, with all losses inside the finite reserves below. No limiting source constant, curve value below 3/4, or curve extrapolation beyond 7/9 is used.

For each closed length cell choose rational \(\fitinline{k,r_0,m_0,w_p}\), put \(\fitinline{T=r_0+m_0}\), and \(\fitinline{m_1=\max(m_0,w_p/2)}\). The test enforces \(\fitinline{r_0\ge2b/H,m_0\ge b/H}\). If \(\fitinline{r\ge1}\), use the original amplified inverse deficit \(\fitinline{D_l=5/6-(5/6-\nu)T}\). If \(\fitinline{r_0\le r<1}\), use \(\fitinline{D_r=\nu r_0+eq(r_0)}\). Otherwise the same pair N4 has \(\fitinline{m\ge m_0-o(1)}\): use N7 or its fixed native alternative on \(\fitinline{[m_0,m_1]}\), and the ordinary fourth deficit \(\fitinline{D_p=\nu w_p+eq(w_p)}\) on \textbf{every \(\fitinline{m\ge m_1}\)}. All witnesses belong to this one T. The count cost at H is \[\fitmath{
 C(D)=\nu y+e[1/2-H/4-y/2+F(y)-u]-HD/2.                       
}\tag{N9}\] It increases in y because \(\fitinline{1/2+\inf F'>0}\). The exact minimum of the 1504 selected cutoffs is \(\fitinline{T_{\min}=111799/100000}\), and \(\fitinline{T_{\min}-1.117=99/100000>0}\). Hence every actual short inverse r\textless1 is a fixed distance greater than .117 below its taper, and the curved calls at r at most 7/9 are farther below it. In the exact certificate 107 cells have T\textless1.12, so the old numerical sentence T\textgreater1.12 is not used. The tangent is required in 279 cells; those tests are genuine E(25) calls.

For low response the earlier new source at \(\fitinline{(n/H,y/H)}\) gives the cost \[\fitmath{
 k\,B(y),\qquad B(y)=(1-H)/2+F(y)-HF(y/H)-u.                  
}\tag{N10}\] The formula for F is continuous but has a \textbf{downward derivative jump} at a; global convexity is not asserted. Partition into 1504 closed cells of width \(\fitinline{1/80000}\), with a as an endpoint. On a low-budget cell include both endpoints and \(\fitinline{Ha}\) when interior. With both arguments cubic, \(\fitinline{B'=F'(y)-F'(y/H)\le0}\); with both affine it is constant. On the crossing piece \(\fitinline{Ha<y<a}\), the right cubic derivative on \(\fitinline{[a,a/H]}\) is below -.17, so \(\fitinline{B'=-.155-F'(y/H)>0}\). These points therefore capture the whole low-budget maximum. The affine plateau is at least \(\fitinline{B(Y)>4.0004\cdot10^{-5}}\), so the whole bracket is positive. Take k as the downward \(\fitinline{10^{-7}}\) rounding of \(\fitinline{(u-3\rho)/\max B}\).

For the high branch, the remaining response quadrilateral has vertices \(\fitinline{(\nu,e)=(k,k),(.5002,.5002),(3/4,.5002),(3/4,k)}\). N9 is affine and N7 is convex in these variables, so those vertices suffice. On each individual y-cell F is convex; N7's extra first-cap debit is a positive part of an affine function. Thus both y-endpoints suffice even below a.

Both ordinary endpoints \(\fitinline{m_0,m_1}\) suffice as well. Below a the caps are fixed and the cost decreases in m. With neither lower cap active its active derivative is \(\fitinline{H[-\nu+e(1/2-f_M'(m))]<0}\). A first-cap-active segment with \(\fitinline{s=m}\) has \(\fitinline{Hm>y}\), so its worst response is \(\fitinline{\nu=e}\), where the active derivative is \(\fitinline{eH[(1.115-1)/2-f_M'(m)]\ge0}\). The inactive expression decreases. The resulting maximum of the decreasing inactive and increasing active pieces is bounded by the two endpoints; entering this segment from a decreasing segment introduces no new maximum. The native alternative decreases in m. This explicitly pays both cap penalties; it does not assume first-cap nonactivation at \(\fitinline{m_0}\).

The exact rational construction fixes all cutoffs, tests these vertices and endpoints, and rounds the radical by adjacent integer fifth powers. Every zero in its positive parts is rational. It uses completed M on 1412 cells and the native pair on 92. The low, inverse, ordinary, long and substitution reserves after the designated \(\fitinline{\rho}\) each exceed \(\fitinline{1.9\cdot10^{-9}}\), with long reserve exceeding \(\fitinline{4.25\cdot10^{-6}}\). These finite inequalities support the analytic proof; they are not source admission by themselves.

Set \(\fitinline{b_j=YH^j}\) and induct downward on its 413 bands to y\_0, clipping the last. On \(\fitinline{[Y,Y/H]}\), F is negative by convexity and its two negative endpoints, so native fourths prove the starting source. In every subsequent band y/H has an earlier proof. High calls invoke only completed M or native fourths. No stage assumes its own conclusion.

\hypertarget{n.4.-every-radius-original-signs-coefficients-and-heights}{%
\subsection{N.4. Every radius, original signs, coefficients and heights}\label{n.4.-every-radius-original-signs-coefficients-and-heights}}

For \(\fitinline{V\ge H}\), freeze the \textbf{entire} earlier M pair tested at H, including its first clamp, so it dominates the actual shortened pair. Its positive absolute quota increases with V, while \(\fitinline{\nu(y-Vm)}\) decreases. Freeze the low pair \(\fitinline{(n/H,y/H)}\) likewise. Above V=1, M's genuine longer-ordinary native continuation is used; the tested short upper exponent is \(\fitinline{\nu y}\), not \(\fitinline{\nu Vy}\). The ordinary escape has the same continuation.

For \(\fitinline{.98\le V\le H}\), use true moving caps in a high M call: \[\fitmath{
 s=\max\{m,a,n/(1.155V)\},\qquad t=\max\{n/V,1.115s\}.         
}\tag{N11}\] The possible derivatives of \(\fitinline{VW_M(t,s)}\) are \[\fitmath{
 (1-s)/2+f_M(s),\quad[1-(1+1.115)s]/2+f_M(s),\quad
                         1/2+f_M(s)-sf_M'(s).
}\] They lie in (0,.61) on the genuine admitted domain. At s=Y switch to the native pair, whose quota is larger since \(\fitinline{f_M(Y)=-10^{-7}}\) and its caps can be shorter. Including the outside exponent gives slope magnitude at most \(\fitinline{(3/4)(.39)+(.5002)(.61)<5/6}\). Count slopes obey the same bound. The low absolute capacity has derivative \(\fitinline{1/2+F(s)-sF'(s)}\), using its one-sided values at the kink; it is continuous, increasing, and its slope times .5002 is below 5/6. At s=Y its native replacement is larger. Positive parts and these source seams preserve the bound. Consequently every H comparison extends against \[\fitmath{
                       u-\rho+(5/6)(H-V)_+.                 
}\tag{N12}\] This is uniform in each actual witness m; it is not continuation of one frozen numerical candidate instead of the witness.

For e\textgreater.5002 retain H6/SN9 at .5001 with its fixed response gap. On \(\fitinline{[.875,.98]}\), the true native/unmarked cost \(\fitinline{(3/4)[(1-V)/2+\max F-u]_+}\) has the checked positive .98 reserve and slope at most 3/8. Below .875 use the two genuine unmarked fourths and true prime-power upper bound; below .08 use the original counted-core estimate. Negative formal mark quotas are replaced by the legitimate unmarked face. The exact global maxima are \[\fitmath{
 W_{\max}=\frac{219944971477}{2000000000000},\qquad
 A_{\max}=\frac{9179599}{10000000}<1.                          
}\tag{N13}\] The certificate verifies separately \(\fitinline{H>4(.139+W_{\max})}\), the .71 short-reference gate, \(\fitinline{5/6+.139>A_{\max}+3u}\), the first-zero literal norm-cap gates and \(\fitinline{5u<1,A_{\max}+4u<1}\). Thus M.3's actual short reference, signed fixed comparison, first zero/tail, second zero/off-zero and counted remainders apply at these verified arguments. There is no new long fixed arm, improved short reference or generic cancellation premise.

Keep both H25/H26 signed arms on their individual original rectangles and inducing characters, their four-coordinate minimum, both independent profile integrals, full bad Euler factors and means. Positive witness substitution changes neither arm. Keep the unweighted first artificial divisor, the second M\"{o}bius divisor with its reciprocal normalizer, both raw identities yielding \(\fitinline{NC\,ND_1}\), full sublattice multiplicity, all seeds/harmonic masses, units, projectors and original signs. In \(\fitinline{z=b^6c}\), the relative exponent is \(\fitinline{6\log_U Nb}\); N12's 5/6 slope pays the full cofactor count without new coprimality. K.3/H.5 restore the original uncentered positive source before any physical lower-mark division.

Fix the finite source dependency list, collars, cutoff/response banks and meshes before live slots; finite strict losses can be below \(\fitinline{10^{-12}}\) at each call, with earlier collars made smaller in reverse dependency order. Keep actual shortened profiles at actual scales, separated by the common compact log-Fourier envelopes of H24a/M.3; do not treat compressed profiles as one bump at an upper ceiling. Fix their seminorm/Sobolev orders and separate normalized log tests before the later one positive physical height exponent and every later buffered-grid index. All real powers, fixed moduli, complete Euler masks and independent heights persist. This proves N1--N2 at M's original source and restoration scope.

\hypertarget{n.5.-the-physical-comparison-at-n0}{%
\subsection{N.5. The physical comparison at N0}\label{n.5.-the-physical-comparison-at-n0}}

Use the actual M.4 current-native rational count and its full-source fifteen-guard proof on only the two original native windows. At the new t,h clear the actual current margin \(\fitinline{(t-E_{\rm base})J-3DN/4}\) at both retained h endpoints and all four original conductor edges; keep the sign-aware division by J. Original conductor/h convexity fills the closed intervals. No old saved-contact shift is substituted for this coupled current price, and the rejected smaller-cap PR129 admission is not used.

The remaining central classes use the actual finite bank of 1025 saturated B.5.1 cutoffs and the maximum of observed raw count deficits G, exactly as M.5. Sufficient G is paid by the original numerator fourth with saving 3G/4. Otherwise choose the first short inverse; its current inverse face, the preceding amplified face and current ordinary face are three lower bounds for that same G. Do not multiply independent counts or raw moments. The genuine ordinary saturation and simultaneous Cauchy use the same physical row subset, while N1 bounds the full uncentered mixed source before mark division. The detector never replaces the physical reciprocal.

For a directed box \(\fitinline{[a_d,b_d]\times[f_p,e_p]}\), retain M19--M22's actual conductor lower bound, annular slack, rounded gamma/lambda choices and faces \(\fitinline{X_j(y)}\). Replace only the surplus by \[\fitmath{
 \begin{gathered}
 \widehat y=\max(y,.3700),\qquad
 \beta_N(v)=g_{1,-}v-g_{0,+}/2+F(v),\\
 \Delta_N(y)=\beta_N(\widehat y)
                 -(1+2g_{1,+})(\widehat y-y)/2,\\
 f_p\Delta_N(y)+X_j(y)/2-e_p\epsilon_a/2-B_+/(D_-a_d)>0.
 \end{gathered}                                                
}\tag{N14}\] The general lower source-clamp debit is retained in N14. Every final ordinary interval lies above .3700, so \(\fitinline{\widehat y=y}\) on this ticket: the physical floor is not used. The larger earlier-M caps \textbf{inside the source proof} remain charged in N7. Source ceilings do not change the actual saturation exponent or exterior \(\fitinline{U^{-dy/2}}\). As in J9a/M26 the fully summed original positive packet mass pays the numerator upper-cap change before profiles or heights are maximized. Positivity holds on every actual length, before p is replaced by its lower endpoint. The pointwise lower function \(\fitinline{b_*(y)=.07599618y-.05805913/2+F(y)}\) has derivative \(\fitinline{.07599618-.155<0}\) on the affine branch; on the cubic branch through the outward upper endpoint .3793594 its derivative is increasing but still negative there. Continuity at the seam therefore gives the global lower bound \[\fitmath{
 \Delta_N(y)\ge
 \frac{63675147649377127}{250000000000000000000}
                         >.0002547                           
}\tag{N14a}\] using the conservative contact endpoints \(\fitinline{g_{1,-}\ge.07599618}\), \(\fitinline{g_{0,+}\le.05805913}\). Thus positivity is established on the whole interval, rather than only at a tested crossover.

The derivative of \(\fitinline{\beta_N}\) is increasing on the upper cubic branch and constant on the lower affine branch, with a downward jump at .3735. The whole-interval adverse-slope bound is the larger of the left-end value and the cubic right derivative at the seam when it lies inside. The strict M/L floor/sloping-crossover derivative inequalities with this bound prove the continuous y minimum and worst cutoff j=1024 on every final cell. This handles the kink without asserting global convexity. There are 81 final windows crossing the seam, 78 wholly below it and 1683 wholly above, so the affine branch is actually used. The general matcher also retains every cutoff, lower clamp, seam, kink and admissible cubic stationary root when that sufficient shortcut fails; it does not continue the upper cubic below .3735.

The exact physical cover has all 57 original rectangles and their four closed conductor-edge tilings, 1440 coarse central cells and 134 full four-way refinements. All 1842 final central cells pass the derivative reduction; no refined child is dropped as an exterior shortcut. Each child rechecks its full count admission, conductor, contour, dilation, source-support and fixed-continuation contract. There is no unpaid cell. The normalized minimum in N14, above the distinct \(\fitinline{10^{-8}}\) target budget already in \(\fitinline{B_+}\), is at least \[\fitmath{
 \frac{5356037}{500000000000000}=1.0712074\cdot10^{-8}.          
}\tag{N15}\] Every cell has \(\fitinline{D_-a_d>.3}\), and the additional original Z-exponent reserve is at least \(\fitinline{1606811/500000000000000>3.2136\cdot10^{-9}}\). All source pairs have ratio strictly between 1.11993 and 1.14059, inside N2; the actual ordinary range is enclosed by {[}.3725154,.3793594{]}, and the upper numerator range by {[}.4247333,.4255376{]}. Every actual shortening is included. The unchanged whole prime pools retain positive greedy exhaustion and true subproducts.

The saved, current-native, E49 exterior and other fixed/moving gate minima, including all refined children, are respectively at least \[\fitmath{
 \frac{7124685929}{31250000000000},\quad
 \frac{79046583}{10^{15}},\quad
 \frac{9570937}{10^{15}},\quad
 \frac{10000017}{10^{15}}.                                   
}\tag{N16}\] There are nine additional upward paired windows: M's six, together with \(\fitinline{[741/2048,93/256]\times[253/512,1/2]}\), \(\fitinline{[2961/8192,741/2048]\times[511/1024,1/2]}\), and \(\fitinline{[23685/65536,2961/8192]\times[8191/16384,1/2]}\). Every one retains all ten authentic E29 guards at all four contact corners, with common directed lower bound \(\fitinline{50495873419/10^{15}>0}\). The original noncentral routes comprise 18 saved, 31 current-native and 19317 E49 exterior pieces. The actual M.4 native windows remain the only native support. No own-band alternative or .12725 partial enters this endpoint composition.

\hypertarget{n.6.-original-complementary-cases-and-all-height-return}{%
\subsection{N.6. Original complementary cases and all-height return}\label{n.6.-original-complementary-cases-and-all-height-return}}

Retain the physical cap .10, declared .11, with the separate GT4--GT7 positive family; actual moment parameter remains \(\fitinline{2\beta-1}\), with the conservative FB2--FB6 contact upper parameter used only in pricing. The original six fixed native/affine/unmarked continuations, ordinary escape, conductor collar, bounded/null/short-numerator/floor cases, units, actual principal residues, zero frequencies, finite Euler corrections and compensators keep their distinct M.6 payers. The actual retained conductor is at least .815777455, density lies in {[}.40234375,.4453125{]}, and p is at least .494140625. Lower-contour damping is at least .0005859375 and the fixed-continuation reserve exceeds .002374704. No pole is paid by a count gain.

Keep M26's entire error-tuple sum with its \(\fitinline{N_{\mathbf e}^{-1/2}}\) and damping weights inside the original positive mixed norm once. A subpower factor alone does not bound errorProductSize. Keep both raw denominator identities, full seeds and multiplicities, genuine dilation saving \(\fitinline{dv/4}\), and the actual absolute contour motion and complete w-inversion before amplitude subdivision. The physical reciprocal stays at B.8's retained \(\fitinline{a+16e+it}\), with its separate s+w Euler-product guard; truncation retains the global \(\fitinline{\beta+20e}\) line and moves only the retained segment. Three physical contour frequencies and independent ordinary/witness heights keep their separate weights.

At the actual N0 parameters construct the generic geometry directly: \(\fitinline{\ell=11/3-4t_N}\), \(\fitinline{x_t=1+\ell-h}\), \(\fitinline{y_t=h-2\ell}\). It lies in L19's proved corridor, has every Geometry.Parameters and BoundsCore field and the exact signal boundary, and keeps the literal h\textless.816 wrapper. The actual GT capacity, floor/principal, contour/damping, conductor and pad guards are checked at N0, not borrowed from M's numerical reserves. The floor reserve after \(\fitinline{10^{-5}+10^{-12}}\) exceeds .0029391569, small rows exceed .07154622, and the two principal reserves exceed .02139852 and .0001099241. Choose \(\fitinline{s_0=10^{-260},F_{\rm pad}=10^{-258}}\), \(\fitinline{D_{\rm src}.t\le F_{\rm pad}/4}\), \(\fitinline{\Delta=F_{\rm pad}}\); the literal cube pad is then \(\fitinline{F_{\rm pad}/2+2D_{\rm src}.t\le F_{\rm pad}}\). Capacities are 1-Lipschitz and conductor variation below two; charge \(\fitinline{15s_0/16+2F_{\rm pad}}\) separately, extend to \(\fitinline{D\le h+F_{\rm pad}}\), and allocate further outer finite losses below \(\fitinline{10^{-10}}\), within N15--N16. Source losses separately fit N.3's strict bounds.

Fix the complete finite source/cutoff dependency list and the common profile envelopes before the one positive physical height exponent, which precedes every later positive buffered-grid bound. Scale/dilation suprema use the actual common Fourier envelope or \(\fitinline{xW'(x)}\), height suprema the separate normalized first log tests; choose tails against these fixed polynomial degrees. This estimates the entire original reciprocal and numerator, not an advantageous witness term.

The literal I51 cube statement therefore holds at N0 for every original ActualBufferedSelection, in the same order \(\fitinline{\forall\eta\ \exists\tau_\eta>0\ \forall n>0\ \exists C_n}\), eventually in Z. Choose one common \(\fitinline{0<\sigma=P_{\rm src}.\sigma<10^{-9}}\) below all original transport reserves and bound the cube divided by I52's original nonzero normalizer by \(\fitinline{C_n Z^{\beta-k_{\rm phys}-\sigma}}\). Generic high\_of\_cube under the directly proved geometry gives the actual physical probe and Gram exponent \(\fitinline{t_N-k_{\rm phys}}\).

For \(\fitinline{\beta>t_N}\), choose \(\fitinline{\omega<(\beta-t_N)/4}\) and \(\fitinline{0<m<\min(\beta-t_N-\omega,\sigma,\beta-.87)}\). H47 gives holomorphy on \(\fitinline{\Re s>\beta-m}\), while the absolute-line Mellin expression is the original Gaussian times \(\fitinline{H_\eta/L_\eta}\), with \(\fitinline{|H_\eta-1|\le1/2}\) on \(\fitinline{\Re s\ge.87}\). A zero above \(\fitinline{\beta-m/2}\), supplied by the definition of the true supremum, cannot cancel, contradicting holomorphy. No attained supremum is assumed. Original norm pullback and fixed Euler restoration give N0 at every finite height for every ordinary and imprimitive character of each fixed positive presented modulus, excluding only its actual principal pole.

\hypertarget{n.7.-retained-evidence-and-source-scope}{%
\subsection{N.7. Retained evidence and source scope}\label{n.7.-retained-evidence-and-source-scope}}

The predecessor is the installed complete M paper, canonical active member \texttt{C\allowbreak{}L\allowbreak{}O\allowbreak{}S\allowbreak{}E\allowbreak{}D\allowbreak{}\_\allowbreak{}P\allowbreak{}R\allowbreak{}O\allowbreak{}O\allowbreak{}F\allowbreak{}.\allowbreak{}m\allowbreak{}d}, 932849 bytes, SHA-256 \texttt{5\allowbreak{}d\allowbreak{}5\allowbreak{}2\allowbreak{}c\allowbreak{}a\allowbreak{}4\allowbreak{}a\allowbreak{}7\allowbreak{}d\allowbreak{}e\allowbreak{}2\allowbreak{}a\allowbreak{}e\allowbreak{}c\allowbreak{}1\allowbreak{}f\allowbreak{}9\allowbreak{}a\allowbreak{}2\allowbreak{}3\allowbreak{}d\allowbreak{}0\allowbreak{}f\allowbreak{}6\allowbreak{}3\allowbreak{}7\allowbreak{}e\allowbreak{}2\allowbreak{}4\allowbreak{}7\allowbreak{}0\allowbreak{}0\allowbreak{}e\allowbreak{}2\allowbreak{}6\allowbreak{}4\allowbreak{}7\allowbreak{}7\allowbreak{}3\allowbreak{}5\allowbreak{}9\allowbreak{}1\allowbreak{}8\allowbreak{}6\allowbreak{}5\allowbreak{}8\allowbreak{}e\allowbreak{}0\allowbreak{}7\allowbreak{}1\allowbreak{}8\allowbreak{}9\allowbreak{}9\allowbreak{}2\allowbreak{}3\allowbreak{}d\allowbreak{}0\allowbreak{}6\allowbreak{}d\allowbreak{}f\allowbreak{}d\allowbreak{}8\allowbreak{}d}. Its prospective successor virtual address is \texttt{h\allowbreak{}i\allowbreak{}s\allowbreak{}t\allowbreak{}o\allowbreak{}r\allowbreak{}y\allowbreak{}/\allowbreak{}p\allowbreak{}r\allowbreak{}o\allowbreak{}o\allowbreak{}f\allowbreak{}\_\allowbreak{}v\allowbreak{}e\allowbreak{}r\allowbreak{}s\allowbreak{}i\allowbreak{}o\allowbreak{}n\allowbreak{}s\allowbreak{}/\allowbreak{}v\allowbreak{}i\allowbreak{}r\allowbreak{}t\allowbreak{}u\allowbreak{}a\allowbreak{}l\allowbreak{}/\allowbreak{}5\allowbreak{}d\allowbreak{}5\allowbreak{}2\allowbreak{}c\allowbreak{}a\allowbreak{}4\allowbreak{}a\allowbreak{}7\allowbreak{}d\allowbreak{}e\allowbreak{}2\allowbreak{}a\allowbreak{}e\allowbreak{}c\allowbreak{}1\allowbreak{}f\allowbreak{}9\allowbreak{}a\allowbreak{}2\allowbreak{}3\allowbreak{}d\allowbreak{}0\allowbreak{}f\allowbreak{}6\allowbreak{}3\allowbreak{}7\allowbreak{}e\allowbreak{}2\allowbreak{}4\allowbreak{}7\allowbreak{}0\allowbreak{}0\allowbreak{}e\allowbreak{}2\allowbreak{}6\allowbreak{}4\allowbreak{}7\allowbreak{}7\allowbreak{}3\allowbreak{}5\allowbreak{}9\allowbreak{}1\allowbreak{}8\allowbreak{}6\allowbreak{}5\allowbreak{}8\allowbreak{}e\allowbreak{}0\allowbreak{}7\allowbreak{}1\allowbreak{}8\allowbreak{}9\allowbreak{}9\allowbreak{}2\allowbreak{}3\allowbreak{}d\allowbreak{}0\allowbreak{}6\allowbreak{}d\allowbreak{}f\allowbreak{}d\allowbreak{}8\allowbreak{}d\allowbreak{}.\allowbreak{}m\allowbreak{}d}; that alias is registered only by the owning packager after actual predecessor closure and the approved successor proof splice. No local-only predecessor path is represented as an archive member. Retain the full I,K,L,M proofs, including M's original fixed-cutoff exact rational supplement, as historical and logical source evidence.

The new source is \texttt{r\allowbreak{}e\allowbreak{}s\allowbreak{}e\allowbreak{}a\allowbreak{}r\allowbreak{}c\allowbreak{}h\allowbreak{}/\allowbreak{}l\allowbreak{}o\allowbreak{}c\allowbreak{}a\allowbreak{}l\allowbreak{}\_\allowbreak{}c\allowbreak{}y\allowbreak{}c\allowbreak{}l\allowbreak{}e\allowbreak{}0\allowbreak{}8\allowbreak{}6\allowbreak{}/\allowbreak{}c\allowbreak{}o\allowbreak{}n\allowbreak{}t\allowbreak{}r\allowbreak{}o\allowbreak{}l\allowbreak{}l\allowbreak{}e\allowbreak{}r\allowbreak{}/\allowbreak{}T\allowbreak{}W\allowbreak{}O\allowbreak{}\_\allowbreak{}E\allowbreak{}P\allowbreak{}O\allowbreak{}C\allowbreak{}H\allowbreak{}\_\allowbreak{}S\allowbreak{}O\allowbreak{}U\allowbreak{}R\allowbreak{}C\allowbreak{}E\allowbreak{}.\allowbreak{}m\allowbreak{}d}, SHA-256 \texttt{8\allowbreak{}0\allowbreak{}9\allowbreak{}0\allowbreak{}3\allowbreak{}b\allowbreak{}2\allowbreak{}1\allowbreak{}2\allowbreak{}d\allowbreak{}e\allowbreak{}0\allowbreak{}9\allowbreak{}3\allowbreak{}9\allowbreak{}1\allowbreak{}0\allowbreak{}f\allowbreak{}4\allowbreak{}0\allowbreak{}8\allowbreak{}8\allowbreak{}c\allowbreak{}0\allowbreak{}3\allowbreak{}2\allowbreak{}6\allowbreak{}b\allowbreak{}f\allowbreak{}6\allowbreak{}b\allowbreak{}b\allowbreak{}e\allowbreak{}7\allowbreak{}8\allowbreak{}a\allowbreak{}b\allowbreak{}1\allowbreak{}e\allowbreak{}6\allowbreak{}9\allowbreak{}d\allowbreak{}c\allowbreak{}b\allowbreak{}9\allowbreak{}2\allowbreak{}d\allowbreak{}4\allowbreak{}9\allowbreak{}2\allowbreak{}0\allowbreak{}d\allowbreak{}e\allowbreak{}5\allowbreak{}2\allowbreak{}6\allowbreak{}d\allowbreak{}e\allowbreak{}5\allowbreak{}a\allowbreak{}e\allowbreak{}2\allowbreak{}8\allowbreak{}f}, with \texttt{T\allowbreak{}W\allowbreak{}O\allowbreak{}\_\allowbreak{}E\allowbreak{}P\allowbreak{}O\allowbreak{}C\allowbreak{}H\allowbreak{}\_\allowbreak{}C\allowbreak{}H\allowbreak{}E\allowbreak{}C\allowbreak{}K\allowbreak{}.\allowbreak{}p\allowbreak{}y} \texttt{2\allowbreak{}3\allowbreak{}3\allowbreak{}4\allowbreak{}d\allowbreak{}4\allowbreak{}d\allowbreak{}f\allowbreak{}6\allowbreak{}1\allowbreak{}5\allowbreak{}a\allowbreak{}a\allowbreak{}c\allowbreak{}9\allowbreak{}e\allowbreak{}4\allowbreak{}4\allowbreak{}7\allowbreak{}f\allowbreak{}8\allowbreak{}f\allowbreak{}d\allowbreak{}4\allowbreak{}7\allowbreak{}a\allowbreak{}3\allowbreak{}3\allowbreak{}4\allowbreak{}4\allowbreak{}a\allowbreak{}7\allowbreak{}9\allowbreak{}d\allowbreak{}f\allowbreak{}b\allowbreak{}c\allowbreak{}1\allowbreak{}f\allowbreak{}f\allowbreak{}c\allowbreak{}e\allowbreak{}4\allowbreak{}3\allowbreak{}d\allowbreak{}8\allowbreak{}3\allowbreak{}3\allowbreak{}e\allowbreak{}c\allowbreak{}b\allowbreak{}7\allowbreak{}0\allowbreak{}9\allowbreak{}9\allowbreak{}e\allowbreak{}e\allowbreak{}2\allowbreak{}8\allowbreak{}c\allowbreak{}1\allowbreak{}4\allowbreak{}8\allowbreak{}7} and \texttt{T\allowbreak{}W\allowbreak{}O\allowbreak{}\_\allowbreak{}E\allowbreak{}P\allowbreak{}O\allowbreak{}C\allowbreak{}H\allowbreak{}\_\allowbreak{}C\allowbreak{}E\allowbreak{}R\allowbreak{}T\allowbreak{}I\allowbreak{}F\allowbreak{}I\allowbreak{}C\allowbreak{}A\allowbreak{}T\allowbreak{}E\allowbreak{}.\allowbreak{}j\allowbreak{}s\allowbreak{}o\allowbreak{}n} \texttt{0\allowbreak{}6\allowbreak{}6\allowbreak{}3\allowbreak{}f\allowbreak{}8\allowbreak{}1\allowbreak{}8\allowbreak{}0\allowbreak{}b\allowbreak{}4\allowbreak{}e\allowbreak{}f\allowbreak{}8\allowbreak{}c\allowbreak{}d\allowbreak{}4\allowbreak{}b\allowbreak{}8\allowbreak{}2\allowbreak{}9\allowbreak{}1\allowbreak{}3\allowbreak{}2\allowbreak{}a\allowbreak{}c\allowbreak{}6\allowbreak{}4\allowbreak{}0\allowbreak{}b\allowbreak{}1\allowbreak{}6\allowbreak{}b\allowbreak{}7\allowbreak{}6\allowbreak{}2\allowbreak{}b\allowbreak{}a\allowbreak{}9\allowbreak{}8\allowbreak{}a\allowbreak{}9\allowbreak{}a\allowbreak{}8\allowbreak{}6\allowbreak{}4\allowbreak{}a\allowbreak{}a\allowbreak{}2\allowbreak{}2\allowbreak{}b\allowbreak{}c\allowbreak{}7\allowbreak{}5\allowbreak{}b\allowbreak{}7\allowbreak{}7\allowbreak{}a\allowbreak{}b\allowbreak{}6\allowbreak{}b\allowbreak{}b\allowbreak{}a\allowbreak{}6}: 1504 closed cells, 413 stages, all selected rational cutoffs and outward radical bounds. The independent fixed-cutoff verifier \texttt{r\allowbreak{}e\allowbreak{}s\allowbreak{}e\allowbreak{}a\allowbreak{}r\allowbreak{}c\allowbreak{}h\allowbreak{}/\allowbreak{}l\allowbreak{}o\allowbreak{}c\allowbreak{}a\allowbreak{}l\allowbreak{}\_\allowbreak{}c\allowbreak{}y\allowbreak{}c\allowbreak{}l\allowbreak{}e\allowbreak{}0\allowbreak{}8\allowbreak{}6\allowbreak{}/\allowbreak{}r\allowbreak{}e\allowbreak{}t\allowbreak{}u\allowbreak{}r\allowbreak{}n\allowbreak{}\_\allowbreak{}r\allowbreak{}e\allowbreak{}v\allowbreak{}i\allowbreak{}e\allowbreak{}w\allowbreak{}/\allowbreak{}V\allowbreak{}E\allowbreak{}R\allowbreak{}I\allowbreak{}F\allowbreak{}Y\allowbreak{}\_\allowbreak{}T\allowbreak{}W\allowbreak{}O\allowbreak{}\_\allowbreak{}E\allowbreak{}P\allowbreak{}O\allowbreak{}C\allowbreak{}H\allowbreak{}.\allowbreak{}p\allowbreak{}y}, SHA-256 \texttt{e\allowbreak{}a\allowbreak{}3\allowbreak{}d\allowbreak{}e\allowbreak{}8\allowbreak{}d\allowbreak{}6\allowbreak{}c\allowbreak{}1\allowbreak{}5\allowbreak{}a\allowbreak{}6\allowbreak{}2\allowbreak{}8\allowbreak{}6\allowbreak{}b\allowbreak{}4\allowbreak{}7\allowbreak{}9\allowbreak{}5\allowbreak{}5\allowbreak{}6\allowbreak{}a\allowbreak{}5\allowbreak{}6\allowbreak{}0\allowbreak{}0\allowbreak{}1\allowbreak{}e\allowbreak{}d\allowbreak{}4\allowbreak{}4\allowbreak{}0\allowbreak{}7\allowbreak{}b\allowbreak{}f\allowbreak{}9\allowbreak{}b\allowbreak{}7\allowbreak{}3\allowbreak{}8\allowbreak{}f\allowbreak{}c\allowbreak{}8\allowbreak{}f\allowbreak{}2\allowbreak{}7\allowbreak{}1\allowbreak{}9\allowbreak{}d\allowbreak{}7\allowbreak{}c\allowbreak{}c\allowbreak{}6\allowbreak{}d\allowbreak{}7\allowbreak{}1\allowbreak{}e\allowbreak{}7\allowbreak{}6\allowbreak{}4\allowbreak{}1\allowbreak{}f}, reads the certificate without importing its generator. Its \texttt{T\allowbreak{}W\allowbreak{}O\allowbreak{}\_\allowbreak{}E\allowbreak{}P\allowbreak{}O\allowbreak{}C\allowbreak{}H\allowbreak{}\_\allowbreak{}R\allowbreak{}E\allowbreak{}A\allowbreak{}D\allowbreak{}B\allowbreak{}A\allowbreak{}C\allowbreak{}K\allowbreak{}.\allowbreak{}j\allowbreak{}s\allowbreak{}o\allowbreak{}n}, \texttt{6\allowbreak{}7\allowbreak{}d\allowbreak{}7\allowbreak{}b\allowbreak{}8\allowbreak{}a\allowbreak{}2\allowbreak{}e\allowbreak{}4\allowbreak{}a\allowbreak{}d\allowbreak{}4\allowbreak{}b\allowbreak{}d\allowbreak{}8\allowbreak{}7\allowbreak{}5\allowbreak{}0\allowbreak{}b\allowbreak{}3\allowbreak{}6\allowbreak{}2\allowbreak{}5\allowbreak{}3\allowbreak{}3\allowbreak{}9\allowbreak{}d\allowbreak{}1\allowbreak{}e\allowbreak{}e\allowbreak{}2\allowbreak{}a\allowbreak{}4\allowbreak{}2\allowbreak{}c\allowbreak{}0\allowbreak{}c\allowbreak{}8\allowbreak{}7\allowbreak{}8\allowbreak{}7\allowbreak{}6\allowbreak{}8\allowbreak{}6\allowbreak{}c\allowbreak{}1\allowbreak{}8\allowbreak{}1\allowbreak{}4\allowbreak{}9\allowbreak{}7\allowbreak{}f\allowbreak{}a\allowbreak{}2\allowbreak{}0\allowbreak{}8\allowbreak{}9\allowbreak{}f\allowbreak{}3\allowbreak{}9\allowbreak{}3\allowbreak{}c\allowbreak{}8}, reproduces the exact five minima and all 48128 substitution prices across both y and ratio endpoints, four response vertices and both ordinary endpoints, with the genuine M caps. It separately checks the one-sided source seam, taper and raw guards. These finite checks do not replace the source and restoration arguments above.

The physical derivation is \texttt{r\allowbreak{}e\allowbreak{}s\allowbreak{}e\allowbreak{}a\allowbreak{}r\allowbreak{}c\allowbreak{}h\allowbreak{}/\allowbreak{}l\allowbreak{}o\allowbreak{}c\allowbreak{}a\allowbreak{}l\allowbreak{}\_\allowbreak{}c\allowbreak{}y\allowbreak{}c\allowbreak{}l\allowbreak{}e\allowbreak{}0\allowbreak{}8\allowbreak{}6\allowbreak{}/\allowbreak{}n\allowbreak{}e\allowbreak{}a\allowbreak{}r\allowbreak{}\_\allowbreak{}r\allowbreak{}e\allowbreak{}c\allowbreak{}o\allowbreak{}r\allowbreak{}d\allowbreak{}/\allowbreak{}R\allowbreak{}E\allowbreak{}S\allowbreak{}E\allowbreak{}A\allowbreak{}R\allowbreak{}C\allowbreak{}H\allowbreak{}.\allowbreak{}m\allowbreak{}d}, SHA-256 \texttt{0\allowbreak{}a\allowbreak{}6\allowbreak{}4\allowbreak{}8\allowbreak{}f\allowbreak{}e\allowbreak{}c\allowbreak{}5\allowbreak{}5\allowbreak{}6\allowbreak{}c\allowbreak{}1\allowbreak{}8\allowbreak{}1\allowbreak{}5\allowbreak{}c\allowbreak{}5\allowbreak{}8\allowbreak{}0\allowbreak{}c\allowbreak{}f\allowbreak{}d\allowbreak{}4\allowbreak{}3\allowbreak{}0\allowbreak{}f\allowbreak{}b\allowbreak{}8\allowbreak{}f\allowbreak{}c\allowbreak{}f\allowbreak{}2\allowbreak{}7\allowbreak{}d\allowbreak{}d\allowbreak{}3\allowbreak{}c\allowbreak{}8\allowbreak{}c\allowbreak{}8\allowbreak{}3\allowbreak{}b\allowbreak{}b\allowbreak{}c\allowbreak{}5\allowbreak{}a\allowbreak{}8\allowbreak{}0\allowbreak{}2\allowbreak{}8\allowbreak{}f\allowbreak{}e\allowbreak{}4\allowbreak{}b\allowbreak{}9\allowbreak{}9\allowbreak{}6\allowbreak{}4\allowbreak{}9\allowbreak{}1\allowbreak{}6\allowbreak{}5\allowbreak{}3}. Its \texttt{R\allowbreak{}E\allowbreak{}T\allowbreak{}U\allowbreak{}R\allowbreak{}N\allowbreak{}\_\allowbreak{}F\allowbreak{}I\allowbreak{}N\allowbreak{}A\allowbreak{}L\allowbreak{}.\allowbreak{}j\allowbreak{}s\allowbreak{}o\allowbreak{}n} has SHA-256 \texttt{f\allowbreak{}3\allowbreak{}5\allowbreak{}d\allowbreak{}e\allowbreak{}7\allowbreak{}4\allowbreak{}3\allowbreak{}5\allowbreak{}1\allowbreak{}d\allowbreak{}2\allowbreak{}a\allowbreak{}a\allowbreak{}d\allowbreak{}b\allowbreak{}1\allowbreak{}3\allowbreak{}4\allowbreak{}b\allowbreak{}3\allowbreak{}b\allowbreak{}c\allowbreak{}9\allowbreak{}1\allowbreak{}3\allowbreak{}1\allowbreak{}a\allowbreak{}4\allowbreak{}9\allowbreak{}4\allowbreak{}c\allowbreak{}0\allowbreak{}7\allowbreak{}9\allowbreak{}1\allowbreak{}6\allowbreak{}1\allowbreak{}6\allowbreak{}9\allowbreak{}3\allowbreak{}d\allowbreak{}7\allowbreak{}5\allowbreak{}5\allowbreak{}d\allowbreak{}0\allowbreak{}8\allowbreak{}f\allowbreak{}1\allowbreak{}3\allowbreak{}0\allowbreak{}4\allowbreak{}6\allowbreak{}c\allowbreak{}7\allowbreak{}1\allowbreak{}f\allowbreak{}b\allowbreak{}f\allowbreak{}d\allowbreak{}f\allowbreak{}8\allowbreak{}1}, with coverage digest \texttt{5\allowbreak{}9\allowbreak{}6\allowbreak{}8\allowbreak{}8\allowbreak{}5\allowbreak{}b\allowbreak{}4\allowbreak{}0\allowbreak{}9\allowbreak{}6\allowbreak{}1\allowbreak{}1\allowbreak{}5\allowbreak{}3\allowbreak{}b\allowbreak{}6\allowbreak{}2\allowbreak{}e\allowbreak{}3\allowbreak{}6\allowbreak{}c\allowbreak{}7\allowbreak{}f\allowbreak{}2\allowbreak{}a\allowbreak{}7\allowbreak{}9\allowbreak{}4\allowbreak{}f\allowbreak{}3\allowbreak{}4\allowbreak{}4\allowbreak{}8\allowbreak{}d\allowbreak{}4\allowbreak{}f\allowbreak{}7\allowbreak{}4\allowbreak{}4\allowbreak{}1\allowbreak{}5\allowbreak{}d\allowbreak{}0\allowbreak{}0\allowbreak{}c\allowbreak{}b\allowbreak{}6\allowbreak{}9\allowbreak{}b\allowbreak{}b\allowbreak{}2\allowbreak{}b\allowbreak{}7\allowbreak{}1\allowbreak{}7\allowbreak{}d\allowbreak{}b\allowbreak{}9\allowbreak{}c\allowbreak{}b\allowbreak{}d\allowbreak{}9\allowbreak{}0}. The portable front is that note, receipt, \texttt{S\allowbreak{}O\allowbreak{}U\allowbreak{}R\allowbreak{}C\allowbreak{}E\allowbreak{}\_\allowbreak{}U\allowbreak{}S\allowbreak{}E\allowbreak{}\_\allowbreak{}F\allowbreak{}I\allowbreak{}N\allowbreak{}A\allowbreak{}L\allowbreak{}.\allowbreak{}j\allowbreak{}s\allowbreak{}o\allowbreak{}n}, \texttt{B\allowbreak{}A\allowbreak{}S\allowbreak{}E\allowbreak{}\_\allowbreak{}E\allowbreak{}N\allowbreak{}G\allowbreak{}I\allowbreak{}N\allowbreak{}E\allowbreak{}.\allowbreak{}p\allowbreak{}y}, \texttt{P\allowbreak{}I\allowbreak{}E\allowbreak{}C\allowbreak{}E\allowbreak{}W\allowbreak{}I\allowbreak{}S\allowbreak{}E\allowbreak{}\_\allowbreak{}P\allowbreak{}R\allowbreak{}O\allowbreak{}F\allowbreak{}I\allowbreak{}L\allowbreak{}E\allowbreak{}\_\allowbreak{}C\allowbreak{}H\allowbreak{}E\allowbreak{}C\allowbreak{}K\allowbreak{}.\allowbreak{}p\allowbreak{}y}, and \texttt{F\allowbreak{}I\allowbreak{}N\allowbreak{}A\allowbreak{}L\allowbreak{}\_\allowbreak{}C\allowbreak{}H\allowbreak{}E\allowbreak{}C\allowbreak{}K\allowbreak{}.\allowbreak{}p\allowbreak{}y}; the latter checks the bound input identities and with \texttt{-\/-source-root\ \textless{}unpacked-proof-root\textgreater{}\ -\/-recompute\ -\/-output\ \textless{}fresh-json\textgreater{}} reconstructs the complete physical calculation byte-identically. It retains the inherited actual full-source native-guard receipt rather than a smaller-cap transfer.

The independent \texttt{r\allowbreak{}e\allowbreak{}s\allowbreak{}e\allowbreak{}a\allowbreak{}r\allowbreak{}c\allowbreak{}h\allowbreak{}/\allowbreak{}l\allowbreak{}o\allowbreak{}c\allowbreak{}a\allowbreak{}l\allowbreak{}\_\allowbreak{}c\allowbreak{}y\allowbreak{}c\allowbreak{}l\allowbreak{}e\allowbreak{}0\allowbreak{}8\allowbreak{}6\allowbreak{}/\allowbreak{}r\allowbreak{}e\allowbreak{}t\allowbreak{}u\allowbreak{}r\allowbreak{}n\allowbreak{}\_\allowbreak{}r\allowbreak{}e\allowbreak{}v\allowbreak{}i\allowbreak{}e\allowbreak{}w\allowbreak{}/\allowbreak{}V\allowbreak{}E\allowbreak{}R\allowbreak{}I\allowbreak{}F\allowbreak{}Y\allowbreak{}\_\allowbreak{}P\allowbreak{}H\allowbreak{}Y\allowbreak{}S\allowbreak{}I\allowbreak{}C\allowbreak{}A\allowbreak{}L\allowbreak{}.\allowbreak{}p\allowbreak{}y}, SHA-256 \texttt{e\allowbreak{}f\allowbreak{}5\allowbreak{}8\allowbreak{}2\allowbreak{}1\allowbreak{}4\allowbreak{}2\allowbreak{}e\allowbreak{}5\allowbreak{}4\allowbreak{}2\allowbreak{}0\allowbreak{}2\allowbreak{}c\allowbreak{}e\allowbreak{}c\allowbreak{}c\allowbreak{}c\allowbreak{}e\allowbreak{}4\allowbreak{}1\allowbreak{}2\allowbreak{}f\allowbreak{}5\allowbreak{}9\allowbreak{}7\allowbreak{}9\allowbreak{}9\allowbreak{}8\allowbreak{}5\allowbreak{}6\allowbreak{}5\allowbreak{}2\allowbreak{}7\allowbreak{}a\allowbreak{}5\allowbreak{}a\allowbreak{}3\allowbreak{}2\allowbreak{}7\allowbreak{}f\allowbreak{}f\allowbreak{}4\allowbreak{}f\allowbreak{}1\allowbreak{}a\allowbreak{}4\allowbreak{}d\allowbreak{}6\allowbreak{}1\allowbreak{}7\allowbreak{}0\allowbreak{}f\allowbreak{}b\allowbreak{}9\allowbreak{}7\allowbreak{}0\allowbreak{}c\allowbreak{}6\allowbreak{}e\allowbreak{}6\allowbreak{}f\allowbreak{}a}, reuses its pinned earlier independent edge algorithm and checks every cutoff at its actual continuous minimum, with both one-sided seam slopes. \texttt{P\allowbreak{}H\allowbreak{}Y\allowbreak{}S\allowbreak{}I\allowbreak{}C\allowbreak{}A\allowbreak{}L\allowbreak{}\_\allowbreak{}A\allowbreak{}L\allowbreak{}L\allowbreak{}\_\allowbreak{}C\allowbreak{}U\allowbreak{}T\allowbreak{}O\allowbreak{}F\allowbreak{}F\allowbreak{}S\allowbreak{}.\allowbreak{}j\allowbreak{}s\allowbreak{}o\allowbreak{}n}, SHA-256 \texttt{f\allowbreak{}a\allowbreak{}8\allowbreak{}4\allowbreak{}3\allowbreak{}7\allowbreak{}8\allowbreak{}b\allowbreak{}6\allowbreak{}c\allowbreak{}2\allowbreak{}7\allowbreak{}7\allowbreak{}c\allowbreak{}7\allowbreak{}c\allowbreak{}9\allowbreak{}6\allowbreak{}a\allowbreak{}6\allowbreak{}7\allowbreak{}1\allowbreak{}d\allowbreak{}a\allowbreak{}9\allowbreak{}4\allowbreak{}c\allowbreak{}6\allowbreak{}7\allowbreak{}8\allowbreak{}a\allowbreak{}9\allowbreak{}9\allowbreak{}7\allowbreak{}6\allowbreak{}1\allowbreak{}9\allowbreak{}9\allowbreak{}4\allowbreak{}e\allowbreak{}1\allowbreak{}b\allowbreak{}7\allowbreak{}f\allowbreak{}b\allowbreak{}a\allowbreak{}3\allowbreak{}d\allowbreak{}5\allowbreak{}9\allowbreak{}7\allowbreak{}d\allowbreak{}2\allowbreak{}6\allowbreak{}f\allowbreak{}3\allowbreak{}a\allowbreak{}0\allowbreak{}7\allowbreak{}0\allowbreak{}8\allowbreak{}d\allowbreak{}e\allowbreak{}2}, records 2025400 node/cutoff pairs and 1888050 terminal-cell/cutoff pairs with the same directed reserves. Exact dependencies and logical manuscript span addresses are in the companion source-use receipt. The finite receipt retains \texttt{P\allowbreak{}A\allowbreak{}S\allowbreak{}S\allowbreak{}\_\allowbreak{}W\allowbreak{}H\allowbreak{}O\allowbreak{}L\allowbreak{}E\allowbreak{}\_\allowbreak{}R\allowbreak{}E\allowbreak{}T\allowbreak{}U\allowbreak{}R\allowbreak{}N\allowbreak{}\_\allowbreak{}C\allowbreak{}O\allowbreak{}N\allowbreak{}D\allowbreak{}I\allowbreak{}T\allowbreak{}I\allowbreak{}O\allowbreak{}N\allowbreak{}A\allowbreak{}L\allowbreak{}\_\allowbreak{}O\allowbreak{}N\allowbreak{}\_\allowbreak{}P\allowbreak{}R\allowbreak{}O\allowbreak{}F\allowbreak{}I\allowbreak{}L\allowbreak{}E}: N.1--N.4 discharge its analytic premise and N.5--N.6 perform its full return. Neither arithmetic nor a source name is source admission or a new Lean run.

The separate own-band component, two-prime .12725 inlet and selected/unselected common-column partials are useful research, not premises here. The full .12725 and .127 campaigns remain open. Preserve the full working proof since OpenAI and permission for properly derived alternatives, the original OpenAI/Collet/Sil/Bebel/Levy/Gao/Liu and other source credits, contributor licenses and reported assistance, retained PR43 source terms and Codex disclosure, and Mathematical Writing for AI Agents / Sepia-CWR attribution. Catalogue and Lean interfaces remain discovery and provenance, not replacements for source applicability or end-to-end formal verification.

\hypertarget{appendix-o}{}

\hypertarget{appendix-o.-a-lower-paired-anchor-and-the-return-above-the-old-conductor-corridor}{%
\section{Appendix O. A lower paired anchor and the return above the old conductor corridor}\label{appendix-o.-a-lower-paired-anchor-and-the-return-above-the-old-conductor-corridor}}

Retain the complete proof through N, including the explicitly retained natural selected-cell input H6/SN9, every original source status, fixed presented modulus, initial keep, all six units and finite Euler holes. The result of this appendix is \[\fitmath{
 c_O=\frac{126667}{10^6}=.126667,\qquad
 t_O=\frac{873333}{10^6}=.873333.                              
}\tag{O0}\] For every fixed positive presented modulus, every ordinary character, including an imprimitive presentation, and every finite height, the ordinary L-function has no zero in the strict half-plane \(\fitinline{\Re s>t_O}\), with its actual principal pole excepted. The gain over N is exactly \(\fitinline{11/20000000}\) in c.~There is no assertion on the boundary, growing-modulus uniformity or new end-to-end formal verification.

Two changes pay this return. A genuine ordinary reference source with slope \(\fitinline{23/10}\) permits 68 further paired stages down to .72. A fixed interior cutoff from those sources pays several original conductor edges without an optimized-crossing admission. The remaining central rows use the completed N source. The higher conductor is treated by its actual primitive inequalities; the former \(\fitinline{h<.816}\) wrapper is not used. None of the open .127 or .12725 campaigns, signed Gauss hypotheses or rough-column partials is a premise.

It suffices to exclude a contrary actual Hecke-family zero supremum \(\fitinline{\beta>t_O}\). N already proves \(\fitinline{\beta\le t_N=.87333355}\) before ordinary pullback. Thus throughout the source construction \(\fitinline{t_O<\beta\le t_N<\alpha}\), and in particular \(\fitinline{\beta>.873>.51}\), as required by the physical and native-source inputs. Only pricing uses the larger contact upper parameter; the actual moment data remain those of \(\fitinline{\beta}\).

\hypertarget{o.1.-a-genuine-reference-source-with-positive-erasure}{%
\subsection{O.1. A genuine reference source with positive erasure}\label{o.1.-a-genuine-reference-source-with-positive-erasure}}

Let \(\fitinline{G(B,L,c)}\) denote the original global \textbf{ordinary} positive source, with its full natural coefficient class, common radial keep, fixed rays, ordinary and prime-slot masks, independent heights and finite homogeneous profile controls. Put \[\fitmath{
 s_0=10^{-260},\quad B_0=1-s_0,\quad L_0=2+10s_0,\quad
 c_0=\frac3{20}-\frac{1000s_0}{1143},\qquad
 L_R=\frac{23}{10},\quad c_R=.149,\quad S_R=.146.             
}\tag{O1}\] The full native source already used in M--N is \(\fitinline{G(B_0,L_0,c_0)}\). Restrict its original rectangles to the smaller face \(\fitinline{G(B_0,L_R,c_R)}\); \(\fitinline{L_R>L_0}\) and \(\fitinline{c_R<c_0}\) give this restriction directly. This statement concerns the ordinary source, not the smaller cap of an unrelated old count admission. Inverse calls retain their fixed deleted base and natural conductor; the growing redundant holes of original ordinary factors are restored by K.3/H.5 at their true scales.

Keep three separate collars: \[\fitmath{
 \begin{array}{c|c}
 \text{native cap to declared reference cap}&c_0-c_R=.001-1000s_0/1143\\
 \text{declared reference cap to child cap}&c_R-S_R=.003\\
 \text{operational reflection/erasure reserve}&\rho_R=.001.
 \end{array}                                                
}\tag{O2}\] Thus \(\fitinline{S+\rho_R\le.147<c_R}\) for every child \(\fitinline{0\le S\le S_R}\). The separate .01 below is a geometric slack, not an available cap collar or a substitute for the actual \texttt{R\allowbreak{}e\allowbreak{}f\allowbreak{}G\allowbreak{}a\allowbreak{}t\allowbreak{}e\allowbreak{}s\allowbreak{}.\allowbreak{}c\allowbreak{}a\allowbreak{}p} inequality.

Freeze the common radial keep and every original mark and error label before Hilbert Cauchy. With \(\fitinline{\kappa_*=3/4}\), the norm tilt is \[\fitmath{
             \tau_R=\frac{\kappa_*}{2L_R}=\frac{15}{92}.
}\tag{O3}\] If extracting an ordinary divisor of logarithmic width \(\fitinline{t_2}\) and freezing marked width f creates the positive capacity defect \(\fitinline{(t_2-L_Rf)_+}\), erasing live width pays energy at most \(\fitinline{\kappa_*(t_2-L_Rf)_+/L_R}\), apart from the whole-slot rounding collar. Its Hilbert-norm cost must retain the positive part: \[\fitmath{
 e^{\tau_R(t_2-L_Rf)_+}
       \le 1+e^{\tau_R(t_2-L_Rf)}.                            
}\tag{O4}\] The two resulting forward kernels are both kept. In the tilted kernel a prime of ordinary degree k and mark multiplicity m has exponent, after prime summation, \[\fitmath{
 E_R(k,m)=1-\frac{k+m}{2}+\tau_R(k-L_Rm).
}\] The actual possibilities and their bounds are \[\fitmath{
 \begin{array}{c|c}
 m=0,\ k\ge3&E_R(k,0)\le-1/92\\
 m=1,\ k\ge1&E_R(k,1)\le-39/184\\
 m\ge2,\ k\ge0&E_R(k,m)\le-3/4.
 \end{array}                                                
}\tag{O5}\] Hence this kernel has the positive prime-power gap \(\fitinline{1/92}\). In the untilted kernel the exponent is \(\fitinline{1-(k+m)/2}\). The cases \(\fitinline{(k,m)=(1,1),(0,2)}\) are marginal and are not assigned a negative Euler gap. Their primes belong to the finitely labeled marked pools of the original expansion. Their dyadic harmonic sums, with the fixed divisor multiplicities, cost at most fixed powers of logarithms and therefore every prescribed positive loss. All other untilted cases have a strict negative exponent. The slot/label number is fixed before the scale; no unbounded assignment sum is replaced by one fixed tuple. Original error coefficients, multiplicities and log tests stay inside these majorants.

Where an affordable full-slot erasure exists, use the original selection, rounding erased width upward on its early mesh. Choose that mesh and its finite profile tolerances below the positive \(\fitinline{\rho_R}\) allowance. The erasure lemma applies only with its actual hypotheses \(\fitinline{S-f-e\ge0}\) and \(\fitinline{t_2\le L_R(f+e)}\); then the shortened reflected rectangle lies in the parent face, including \(\fitinline{S-f-e\le c_R}\). In particular this does not supply an erasure when \(\fitinline{t_2>L_RS}\).

Keep the original short reference \(\fitinline{l_R=.14}\), root threshold \(\fitinline{P_R=2/3}\), and ordinary StepGates/geometric reserve .01. Direct or naturally reflected leaves fit because \[\fitmath{
 l_R+\frac12+L_RS_R+.01=.9858<B_0,
\qquad
 1-P_R+2l_R+L_RS_R+.01
       =\frac13+.28+.3358+.01<B_0.                            
}\tag{O6}\] Also \(\fitinline{l_R+S_R<1,\ 2l_R<P_R,\ .6<P_R}\). These inequalities give an unconditional payment when erasure is unavailable, including \(\fitinline{t_2>L_RS}\). The actual forward \(\fitinline{K_3}\) expansion already produces native leaves at their \textbf{shortened} lengths, with one factor at most \(\fitinline{Z^{l_R}}\). If the resulting long exponent is at most 1/2, apply the parent directly. Otherwise reflect that actual shortened factor; its dual exponent is at most \(\fitinline{1/2+\delta+\xi}\). The half gate O6 pays this leaf with every surviving mark, since their total is at most S\_R. This is the retained A29/\texttt{R\allowbreak{}e\allowbreak{}f\allowbreak{}e\allowbreak{}r\allowbreak{}e\allowbreak{}n\allowbreak{}c\allowbreak{}e\allowbreak{}.\allowbreak{}L\allowbreak{}e\allowbreak{}a\allowbreak{}f\allowbreak{}C\allowbreak{}o\allowbreak{}r\allowbreak{}e} argument, with no affordability premise. The actual reflection annuli and errors are chosen inside its strict geometric gaps, while operational cap/erasure losses, where used, have O2's separate reserve. Apply the parent source before restoring each profile. The common fixed radial majorant bounds the full positive original row sum; no prime quotient or new keep is introduced.

The reflected and unreflected ordinary heights remain independent. Differences between a common height and either independent height are absorbed separately in its compact norm-power profile and normalized first logarithmic test. The early mesh is chosen before ambient data. After the permitted fixed radial cap, presented characters and internal ideal \(\fitinline{Q_0}\) are given, choose the finite degrees, controls and constants, taking one maximum over the finitely many suppliers. These precede the live frequencies, physical height exponent and grid indices; they are not claimed independent of that fixed data. The parent source has polynomial height factors at those controls. Thus the usual reference energy and ordinary zero-stage-loss input are proved for O1, with every prescribed positive output decrement. This redoes the changed reference estimate; it does not instantiate the old numerical Lean reference theorem at new constants.

\hypertarget{o.2.-sixty-eight-finite-paired-stages}{%
\subsection{O.2. Sixty-eight finite paired stages}\label{o.2.-sixty-eight-finite-paired-stages}}

Use the actual last tuple of the retained 8192-stage PR141 lower chain, evaluate its output at \(\fitinline{1467/2000=.7335}\), and keep its cutoff and earlier affine line. The paired-anchor contract is the original one: for each positive rectangle decrement, it supplies the inverse moment and the positive two-ordinary energy on their individual original rectangles, not a mixed inverse--ordinary moment.

For an entry anchor \(\fitinline{(R,z)}\), cutoff k and earlier line \(\fitinline{a-bw}\), define, for \(\fitinline{r\le R}\), \[\fitmath{
 \begin{gathered}
 H=r/R,\quad u=(1-H)/6,\quad
 Y=Hz+(1/k-2)u,\quad S=(1-r)/2+Y,\quad s'=S-u,\quad D=1/2-b,\\
 w_b=(R/2+z-a)/D,\qquad
 F_{\rm high}=\frac56-(\frac54-\frac32k)w_b-\frac{k(1+R)}2-kz,\\
 w(\nu)=\frac{r+s'-H/2-Ha-u/\nu}{DH}+q,\qquad
 q=10^{-12},\quad p_0=10^{-13}.
 \end{gathered}                                             
}\tag{O7}\] For \(\fitinline{k\le\nu\le3/4,\ 0\le e\le\nu}\), the original short detector price is at most u. If \(\fitinline{g_\nu=H(w(\nu)-q)}\), its long-price difference is exactly \[\fitmath{
 G(k)=HF_{\rm high},\quad
 G(\nu)-G(k)=(\nu-k)
 \left(\frac32g_\nu-r-s'-\frac{3(5/6-k)u}{2k\nu D}\right).
}\tag{O8}\] The derivative gate below makes the bracket nonnegative. The true positive witness width is \(\fitinline{\max(0,(1-w)/2+a-bw-p_0)}\), not an independently assigned mark. In the short branch the actual selected-cap debit is \(\fitinline{(1/2+b)q+p_0\le q}\), since \(\fitinline{b<.4}\) and \(\fitinline{p_0=q/10}\); the positive witness shift and \(\fitinline{e\le\nu}\) pay it. The long q-padding costs at most \(\fitinline{5q/2}\). The \(\fitinline{3\cdot10^{-11}}\) high reserve pays it, and the original tiny class keeps its .03 reserve.

For both endpoints of each stage the exact construction proves the following fourteen quantities positive: \[\fitmath{
 \begin{aligned}
 &HF_{\rm high}-3\cdot10^{-11},\quad .146-S,\quad .149-S/H,\quad
 .9-r-S-4u,\quad .91-r-S-4u,\\
 &\frac56+.6(.14)-r-S-3u-.01,\quad
 3H-2r-8s'-.2,\quad r-.5-Y+u-.2,\\
 &r/2-s'-\frac{3(5/6-k)u}{2k^2D}-.1,\quad
 H+u-H/6-1/15-(3/4)r-s'-.03,\\
 &w(k)-2/3,\quad .79-w(3/4),\\
 &1-(3/4)w(3/4)-(3/4)\max(0,(1-w(3/4))/2+a-bw(3/4))-1/6,\\
 &.79-r-L_0u.
 \end{aligned}                                               
}\tag{O9}\] Clear the positive H,D,k denominators and interpolate each affine alternative of the maximum to obtain the continuous gates. O6 pays both direct and reflected ordinary leaves. The inverse child keeps the original restored inverse source and masks. The former lower endpoint .7335 is replaced at its actual consumers by \(\fitinline{r\ge.72>0}\) and \[\fitmath{
                  (B_0-.72+.001)/L_0<c_0.                   
}\tag{O10}\] The inverse child's radius bound follows from its positive radius denominator and \(\fitinline{r\le R\le311/400}\); it does not require the old .7335 endpoint.

The other former child-cap uses have direct replacements. For a requested small reconstruction decrement \(\fitinline{0<\delta_{\rm rec}\le.001}\), put \(\fitinline{\varkappa=\delta_{\rm rec}/8}\). Then \[\fitmath{
 6u+\varkappa<1,\quad(2+R)\varkappa<\delta_{\rm rec}/2,\quad
 [3+S(R)+\delta_{\rm rec}/2]\varkappa<\delta_{\rm rec}/2.
}\] Both R and \(\fitinline{S(R)\ge(1-.7775)/2}\) exceed \(\fitinline{\delta_{\rm rec}/2}\). These are the actual inverse-child primitive's size conditions. Ordinary LOW deletion has \(\fitinline{S-u-\delta_{\rm rec}<1/6}\) and entry cap below one, since S\textless.146. Larger requested decrements follow by restricting a smaller positive rectangle. The ordinary .01 reserve is separately paid by \(\fitinline{5/6-P_R-S-3u\ge1/6-.146-.0003>.01}\).

Partition {[}.72,.7335{]} into 68 equal intervals, in decreasing order. Each next entry gain is exactly the preceding stage's output at its left endpoint. The exact rational receipt proves O9 and \[\fitmath{
 \begin{gathered}
 u\le10^{-4},\quad R\le w_b,\quad
 Y(r)>a-br+p_0,\\
 Y(r)>.05805913-.07599618r+10^{-7}.
 \end{gathered}                                             
}\tag{O11}\] The actual HIGH calls use \(\fitinline{w\ge w(k)>R}\), since q\textgreater0 and \(\fitinline{R\le w_b}\), and they have the strict p0 output collar. Thus they lie in a \textbf{finite} cover of already completed earlier-anchor rectangles. Fix the cover, positive decrements and early meshes, then, after the permitted fixed radial cap, characters and \(\fitinline{Q_0}\), take the maximum of its finitely many degrees, controls and constants before any live heights or grid indices. This does not assume uniform constants over a continuum of anchors or place those constants before their allowed fixed data. Above the earlier table's genuine endpoint \(\fitinline{311/400}\), its negative affine line is covered by the actual native inverse/ordinary continuation. No tangent line is asserted beyond its source.

Finite induction therefore proves both paired sources on {[}.72,.7335{]}. The final gain is exactly \(\fitinline{11148401676729938999/2084336624804687500000}\). Together with E's genuine tangent supplier this gives the tangent paired fiber on \(\fitinline{[.72,311/400]}\), with all positive rectangle decrements, original measures, fixed rays, natural zeros, masks, independent heights and profile/log tests. The smallest high reserve after its fixed debit exceeds .00024255, and the child-cap gap exceeds .00065134. A count rectangle still requires all forty authentic E29 guards; merely lowering a crossing cutoff admits no rectangle.

For this full closed fiber use a second finite right-anchor grid on the requested child widths, inside O11's strict line surplus and the consumer's retained width collars. Together with the finite earlier-HIGH cover it supplies one minimum mesh and finitely many profile controls. No continuum of independently chosen constants is passed to a physical supremum.

\hypertarget{o.3.-the-actual-higher-conductor-geometry-and-finite-padding}{%
\subsection{O.3. The actual higher-conductor geometry and finite padding}\label{o.3.-the-actual-higher-conductor-geometry-and-finite-padding}}

Choose \[\fitmath{
 h=h_*+\frac1{500},\quad
 .817324503\le h\le.817324505,\quad
 \ell=\frac{130001}{750000}=\frac{11}3-4t_O,\quad
 x=1+\ell-h,\quad y=h-2\ell.                              
}\tag{O12}\] Here x,y are the outer geometric exponents, distinct from the detector lengths below. Keep \(\fitinline{F_{\rm pad}=10^{-258}}\), source cost s0, actual physical cap \(\fitinline{17/200=.085}\), and separately declared cap .11. Native count/source caps remain their genuine O1 values. The old \texttt{B\allowbreak{}o\allowbreak{}u\allowbreak{}n\allowbreak{}d\allowbreak{}s\allowbreak{}C\allowbreak{}o\allowbreak{}r\allowbreak{}e\allowbreak{}.\allowbreak{}h\allowbreak{}\_\allowbreak{}i\allowbreak{}n\allowbreak{}t\allowbreak{}e\allowbreak{}r\allowbreak{}v\allowbreak{}a\allowbreak{}l\allowbreak{}.\allowbreak{}2} is false at O12. We construct and use the primitive geometry directly.

The directed geometry receipt proves the 33 geometric inequalities and two actual padding comparisons throughout the retained h interval. Their substantive replacements are \[\fitmath{
 \begin{gathered}
 .873<t_O<.875,\quad .1672<\ell<.175,\quad .458<y<.473,\\
 (49/150)h-\ell-.02>.07,\qquad
 h>1/2,\quad h+F_{\rm pad}<.865<7/8<1,\\
 h+F_{\rm pad}<2y,\quad \ell/(h+F_{\rm pad})>1/5,\quad
 y/(1/200)<99,\quad 5(h+F_{\rm pad})/288<1/50,\\
 \ell<x<y<1/2,\quad2x-y-\ell>0,\quad5\ell<1,\quad
 t_O=11/12-\ell/4.
 \end{gathered}                                             
}\tag{O13}\] The full source census binds the 24 direct h-interval consumer files in the PR141/PR143 interfaces. In finite geometry/summation, source actual/ledger, full-w errors, canonical/common/cube/nonfloor transport and positive scheduling, replace each use of the old upper endpoint by its actual consequence in O13. Positivity uses h\textgreater0; source calls use \(\fitinline{h+F_{\rm pad}<2y}\); rowwise prime and floor moments use \(\fitinline{D\le7/8}\); the shift base and active schedule use \(\fitinline{h+F_{\rm pad}<1}\); feedback uses \(\fitinline{5d_+/288<1/50}\). None requires an instance of the false coarse class or the stationary \texttt{E\allowbreak{}x\allowbreak{}a\allowbreak{}c\allowbreak{}t\allowbreak{}E\allowbreak{}n\allowbreak{}d\allowbreak{}p\allowbreak{}o\allowbreak{}i\allowbreak{}n\allowbreak{}t} adapter.

The actual pinned \texttt{T\allowbreak{}r\allowbreak{}a\allowbreak{}n\allowbreak{}s\allowbreak{}p\allowbreak{}o\allowbreak{}r\allowbreak{}t\allowbreak{}P\allowbreak{}a\allowbreak{}r\allowbreak{}a\allowbreak{}m\allowbreak{}e\allowbreak{}t\allowbreak{}e\allowbreak{}r\allowbreak{}s\allowbreak{}.\allowbreak{}l\allowbreak{}e\allowbreak{}a\allowbreak{}n} defines row padding \(\fitinline{F_{\rm pad}/2}\), and \texttt{T\allowbreak{}r\allowbreak{}a\allowbreak{}n\allowbreak{}s\allowbreak{}p\allowbreak{}o\allowbreak{}r\allowbreak{}t\allowbreak{}G\allowbreak{}e\allowbreak{}o\allowbreak{}m\allowbreak{}e\allowbreak{}t\allowbreak{}r\allowbreak{}y\allowbreak{}.\allowbreak{}l\allowbreak{}e\allowbreak{}a\allowbreak{}n} defines row extension \(\fitinline{F_{\rm pad}}\). Retain those definitions and the stricter literal primitive guard \(\fitinline{\zeta\le F_{\rm pad}}\): \[\fitmath{
 p_{\rm row}=F_{\rm pad}/2,\quad
 0<p_{\rm row}\le F_{\rm pad}=\text{rowExtension},\quad
 D_{\rm src}.t\le F_{\rm pad}/4.
}\] Then \(\fitinline{p_{\rm row}+2D_{\rm src}.t\le F_{\rm pad}}\), and \[\fitmath{
              h+p_{\rm row}<7/8-.01.                       
}\tag{O14}\] Its signal is above .1 by O13. Use the per-character \texttt{n\allowbreak{}o\allowbreak{}r\allowbreak{}m\allowbreak{}a\allowbreak{}l\allowbreak{}i\allowbreak{}z\allowbreak{}e\allowbreak{}d\allowbreak{}P\allowbreak{}r\allowbreak{}o\allowbreak{}b\allowbreak{}e\allowbreak{}\_\allowbreak{}h\allowbreak{}i\allowbreak{}g\allowbreak{}h\allowbreak{}\_\allowbreak{}o\allowbreak{}f\allowbreak{}\_\allowbreak{}f\allowbreak{}i\allowbreak{}n\allowbreak{}i\allowbreak{}t\allowbreak{}e\allowbreak{}\_\allowbreak{}n\allowbreak{}o\allowbreak{}n\allowbreak{}f\allowbreak{}l\allowbreak{}o\allowbreak{}o\allowbreak{}r\allowbreak{}\_\allowbreak{}n\allowbreak{}u\allowbreak{}m\allowbreak{}e\allowbreak{}r\allowbreak{}a\allowbreak{}t\allowbreak{}o\allowbreak{}r\allowbreak{}\_\allowbreak{}a\allowbreak{}t} and \texttt{s\allowbreak{}o\allowbreak{}u\allowbreak{}r\allowbreak{}c\allowbreak{}e\allowbreak{}P\allowbreak{}r\allowbreak{}o\allowbreak{}b\allowbreak{}e\allowbreak{}\_\allowbreak{}h\allowbreak{}i\allowbreak{}g\allowbreak{}h\allowbreak{}\_\allowbreak{}o\allowbreak{}f\allowbreak{}\_\allowbreak{}f\allowbreak{}i\allowbreak{}n\allowbreak{}i\allowbreak{}t\allowbreak{}e\allowbreak{}\_\allowbreak{}n\allowbreak{}u\allowbreak{}m\allowbreak{}e\allowbreak{}r\allowbreak{}a\allowbreak{}t\allowbreak{}o\allowbreak{}r\allowbreak{}\_\allowbreak{}a\allowbreak{}t} primitive proofs, with this padding and the same common sigma, replacing the old h-ceiling deduction by O14. Obsolete header comments with larger padding constants are not current definitions. The scale threshold may depend on the fixed positive padding, as allowed. This preserves the original order \(\fitinline{\forall\eta\ \exists\tau_\eta>0\ \forall n>0}\) followed by an eventual bound uniform in the allowed indices/grid; neither the padding nor tau is chosen after n or a live grid label. The literal cube pad is bounded by the actual F-based budget.

The averaged floor is priced directly, with its complete original numerator and keep: \[\fitmath{
 F_h(D)=.51-49h/150+.51\ell+.66D,\qquad
 t_O-F_h(D)-10^{-5}-10^{-12}>10^{-3}\quad(D\le h+F_{\rm pad}).
}\tag{O15}\] Substitute the upper D endpoint to verify this affine bound. O13 supplies its required \(\fitinline{D\le7/8, D\le1, 0<y/D\le99}\) for \(\fitinline{D\ge1/200}\). This replaces the numerical floor-margin lemmas rather than applying them outside their corridor.

For every \(\fitinline{0<e\le.001}\), the two separate principal remainders retain \[\fitmath{
 y/20-e(1+h)>1/4000,\qquad h/600-e>1/4000.                   
}\tag{O16}\] The original costs \(\fitinline{(49/1000)h-\ell/10-1/800}\), \(\fitinline{h/600-1/800}\), and small-row cost \(\fitinline{(49/150)h-5\ell/4+13/600}\) are positive at both directed h endpoints. These are their actual contour powers, not a symmetric-height approximation or a count gain used to pay a pole.

\hypertarget{o.4.-the-fixed-interior-detector-with-every-length-case}{%
\subsection{O.4. The fixed interior detector, with every length case}\label{o.4.-the-fixed-interior-detector-with-every-length-case}}

Use the original B.5.1 same-zero decomposition at the fixed early cutoff \[\fitmath{
 r_*=.720002,\quad T_*=3r_*/2=1.080003,\quad m_*=r_*/2.
}\tag{O17}\] For each original high buffered class and its actual attaining zero, let v be the inverse witness length and m its companion ordinary length. The original construction gives \(\fitinline{v\le T_*+o(1)}\), \(\fitinline{v+m\ge T_*-o(1)}\), with the original taper, both witnessing tests and independent heights. This is a fixed genuine cutoff, not an optimized E29 crossing. Its separation from every short inverse is \(\fitinline{T_*-1=.080003>.08}\), so the taper remains separated throughout the following cases.

Write d for the actual zero-band parameter and p for the actual relative mark response, and retain directed contact endpoints \(\fitinline{g_{0,-},g_{1,+}}\). The two available count exponents, with the allocated \(\fitinline{10^{-11}}\) count/mark allowance, are \[\fitmath{
 \begin{aligned}
 S_*&=1-d\{p(1+2g_{0,-})+[1-p(1+2g_{1,+})]r_*\},\\
 A_*&=1/6+(5/6-d)T_*.
 \end{aligned}                                                
}\tag{O18}\] These bound the complete alternatives, as follows.

If \(\fitinline{v\ge1}\), C9f's original amplified-long inverse source and B.5(204) give \[\fitmath{
 (1+5v)/6-dv\le1/6+(5/6-d)T_*+o(1)=A_*+o(1).
}\] Here \(\fitinline{5/6-d>0}\), and the cutoff and actual inverse are the same ones used in the other cases. If \(\fitinline{r_*\le v<1}\), use O.2's genuine paired inverse while its anchor is in the proved range, followed by the true native inverse continuation through 1. The coefficient \(\fitinline{1-p(1+2g_1)}\) is positive, so the paired short-count exponent is largest at r*. Above \(\fitinline{311/400}\) the tangent correction is negative; the actual native face dominates the needed bound there, without continuing a positive tangent beyond its source.

If \(\fitinline{v<r_*}\), the same original witness has \(\fitinline{m\ge T_*-v-o(1)\ge m_*-o(1)}\). Use the paired ordinary fourth on its genuine range, then the actual native ordinary upper continuation. C10's original continuation includes the unmarked ordinary fourth for \(\fitinline{m\ge1/2}\) and the bounded-witness collar. Its monotonic count comparison is again largest at \(\fitinline{2m=r_*}\), giving S*. Thus every original length case is covered by \(\fitinline{\max(S_*,A_*)+10^{-11}}\); no short-inverse-only invocation stands in for the ordinary alternatives.

More explicitly, set \(\fitinline{c_*=p(1+2g_{0,-})+[1-p(1+2g_{1,+})]r_*}\), and restrict the genuine native family to \(\fitinline{\bar B=1-10^{-12},\ \bar L=2+10^{-11},\ \bar c=.15-10^{-7}}\). Its ordinary quota is \(\fitinline{\min(\bar c,(\bar B-2m)/\bar L)}\). The natural continuation uses \(\fitinline{s_-=.027}\), \(\fitinline{m_H=[1-(9/2)s_-]/2=.43925}\), at the actual incumbent \(\fitinline{\kappa\in[.7,.75]}\). The endpoint comparisons paying all those arms are \[\fitmath{
 \begin{aligned}
 c_*&\le311/400+(1-311/400)p,\\
 c_*&\le311/400+2p\bar c,\qquad
 c_*\le311/400+2p(\bar B-311/400)/\bar L,\\
 c_*&\le2m_H+2ps_-,\qquad c_*<.9<1.
 \end{aligned}
}\] Each is affine in \(\fitinline{0\le p\le1/2}\), so its two endpoints suffice; the smallest strict continuation gap exceeds .0254. Their brackets increase along the corresponding length pieces. The actual unmarked ordinary fiber already pays \(\fitinline{m\ge.45}\), with bracket \(\fitinline{2m\ge.9>c_*}\), including the entire \(\fitinline{2m=1}\) collar and larger lengths. The upper inverse quota is tapered to zero near one and its unmarked fiber pays that collar. No selected width at a vanishing quota, and no fourth-moment rectangle beyond its true native budget, is assumed.

Take actual finite collars inside the \(\fitinline{2\cdot10^{-6}}\) interior anchor gap and the strict O11 output surplus. They pay the B.5.1 annular o(1), rectangle decrements, whole slots, both log tests and fixed finite earlier-anchor cover, with losses smaller than the \(\fitinline{10^{-11}}\) allowance. Keep the original measures, coefficient normalizers, natural/presentation zeros and masks. The attaining inverse is on the clean regular core and fixed deleted base; original growing ordinary holes are opened at their true scales, and are not installed in an inverse inlet. The original signed fixed-inducing collection and its two individual inducing rays are paid before the positive comparison. The detector is not the physical reciprocal.

To price this count on an original conductor edge D, define its actual capacity \[\fitmath{
 C_D=\begin{cases}
 \min(D,2y-D)/(6\kappa_+),&\text{on either unmarked edge},\\
 \min(.085D,[2y-(2-B_0)D]/L_0),&\text{on either native edge},
 \end{cases}
}\] where \(\fitinline{\kappa_+}\) is the conservative contact upper parameter used only for pricing. The exterior exponent is \[\fitmath{
 E_{\rm out}=\frac12+\frac{13}{75}h-\frac y2-\frac\ell2-.09D
                  +\frac d2+dp(\ell-C_D/2).                  
}\tag{O19}\] The exact box proof clears both \(\fitinline{t_O-E_{\rm out}-(3/4)D(S_*+10^{-11})}\) and its A* counterpart at both h endpoints and the original conductor edges. Conductor convexity fills the intervals. The minimum fixed-interior reserve is at least \(\fitinline{78796110673/10^{15}}\). A bound on an optimized paired crossing is neither needed nor asserted for this route.

\hypertarget{o.5.-the-complete-physical-cover-and-the-inherited-mixed-source}{%
\subsection{O.5. The complete physical cover and the inherited mixed source}\label{o.5.-the-complete-physical-cover-and-the-inherited-mixed-source}}

On an original native branch price the actual coupled numerator and denominator: clear \(\fitinline{(t_O-E_{\rm out})J-3DN/4}\) at both h endpoints, using its authentic full native admission and sign-aware division by J. The rejected smaller-cap PR129 transfer is not used. Saved routes use their true conductor-independent upper count before their conductor slope is applied. A newly used paired rectangle requires all ten authentic E29 guards at all four contact corners, with O.2's lower endpoint .72. The original four closed conductor tilings still contain all 57 rectangles.

The unresolved central pieces use precisely N's completed positive two-ordinary source \[\fitmath{
 W_N(n,z)=\frac{1-n-z}{2}+F_N(z),\quad
 .3700\le z\le.3888,\quad1.118\le n/z\le1.15,                
}\tag{O20}\] with all actual componentwise shortenings and every strictly smaller quota. With \(\fitinline{Y=.3888,a=.3735}\), its function is \[\fitmath{
 F_N(z)=\begin{cases}
 4.59(Y-z)^2+50.5(Y-z)^3+.00031(Y-z)-10^{-7},&a\le z\le Y,\\
 F_N(a)+.155(a-z),&.3700\le z<a.
 \end{cases}                                                
}\tag{O21}\] Its upper cubic's starting-strip extension, completed-M caps and both clamps stay in N's source proof. O20 is used before lower-mark division; it is not a mixed inverse--ordinary moment or an improved cutoff theorem. The actual matched shapes lie inside O20, including every shorter component. In particular the lower linear branch is used where required, with no cubic extrapolation below a.

Fix the 4097-cutoff bank before actual responses. On a fixed original row subclass let G be the largest genuine raw count deficit from the current inverse, preceding amplified inverse and current ordinary witnesses. It is one cardinality bound, obtained as the minimum of three count upper bounds. The sufficient-G sector saves \(\fitinline{3G/4}\) in the original numerator-fourth argument. On the complement use the first short witness, its actual inverse amplitude and the full N source before actual complete prime-slot mark division. This is the controller's complete simultaneous consumer; no product of unrelated row counts, restriction of an entire convolution, or identification of independent marks is made.

The normalized central comparison is N14 with F\_N unchanged and the newly computed debit, gamma/lambda lengths and source ceilings. The continuous y minimum uses the maximum adverse \textbf{one-sided} slope. At the seam a the derivative jumps downward; no global convexity is assumed. All 2147 final central cells satisfy this derivative reduction, including their actual source support, contour, dilation and continuation gates. All 4097 cutoffs are thereby paid on the whole continuous interval; the receipt's zero explicit full sweeps is the outcome of the proved reduction, not a missing cutoff.

The exact cover starts with 1904 coarse central cells and makes 81 complete four-way subdivisions, giving 2147 final cells and 2228 visited nodes. No refined child is discarded as an exterior shortcut. There are no unpaid cells. Above the distinct \(\fitinline{10^{-8}}\) normalized target budget already charged in the debit, its common directed normalized surplus is \[\fitmath{
 \frac{5214079}{125000000000000}
        =4.1712632\cdot10^{-8}.                              
}\tag{O22}\] Since \(\fitinline{D_-d_->.3}\), the additional original Z-exponent reserve is at least \(\fitinline{12513789/10^{15}>1.25\cdot10^{-8}}\). The fixed/saved/current-native/E49/central gate minima, including every refined child, are respectively at least \[\fitmath{
 \frac{78796110673}{10^{15}},\quad
 \frac{5413472629}{250000000000000},\quad
 \frac{1692483179}{125000000000000},\quad
 \frac{9357613}{500000000000000},\quad
 \frac{5000011}{500000000000000}.                            
}\tag{O23}\] The final source ratios are between 1.1195 and 1.1415, safely inside O20; exact rational ranges are retained in the receipt. All newly used upward windows retain their forty real guards. The same whole good-prime pools, greedy selection/exhaustion and initial keep are used; there is no radial filter on the physical rows.

The .085 physical cap is distinct from the declared .11 and native c0. Its top affine capacity is unchanged because \(\fitinline{h/2-2\ell<.085h}\); all other edges are actually repriced with .085. The GT source exists since \[\fitmath{
 .085<.11<c_0,\qquad
 B_0/2+(9/2-L_0/2)(.11)<1,\qquad \ell/(h+F_{\rm pad})>1/5.
}\tag{O24}\] In every retained central cell the actual damping is \(\fitinline{d(1-p)/2-.085>0}\). The original fixed inverse/ordinary continuations, r- and n-contours, amplitude and coefficient-dilation tails, conductor collar, short/null/bounded numerators, ordinary escape, floor, units and principal corrections keep their separate payers. Actual moment data use \(\fitinline{2\beta-1}\); the contact upper parameter is a conservative pricing bound as in FB2--FB6, not an equality imposed on the data.

\hypertarget{o.6.-original-signs-all-heights-and-the-final-mellin-contradiction}{%
\subsection{O.6. Original signs, all heights and the final Mellin contradiction}\label{o.6.-original-signs-all-heights-and-the-final-mellin-contradiction}}

Keep both raw denominator identities, the unweighted first artificial divisor and second M\"{o}bius divisor with its reciprocal normalizer, full sublattice multiplicities and seeds, the unit in \(\fitinline{z=b^6c}\), and the original signed H25/H26 arms on their own rectangles and inducing characters. K.3/H.5 restore the uncentered positive source before mark division. The fully summed error packet stays inside the actual mixed norm with \[\fitmath{
 \sum_{\mathbf e}|A_i(\mathbf e)|N_{\mathbf e}^{-1/2}
       \max(1,L_iN_{\mathbf e})^{d/2}
    \ll U^\epsilon\prod_{p\ \mathrm{main}}Y_p^{-4/25}.         
}\tag{O25}\] A subpower label factor does not bound \texttt{e\allowbreak{}r\allowbreak{}r\allowbreak{}o\allowbreak{}r\allowbreak{}P\allowbreak{}r\allowbreak{}o\allowbreak{}d\allowbreak{}u\allowbreak{}c\allowbreak{}t\allowbreak{}S\allowbreak{}i\allowbreak{}z\allowbreak{}e} separately. Retain the genuine coefficient-dilation saving, every local main/error weight, primitivity and complete finite Euler factors. The absolute contour motion and full w-inversion precede amplitude subdivision.

The physical reciprocal remains on B.8's actual buffered s-line, with its separate s+w guard; omitted segments stay on the global \(\fitinline{\beta+20e}\) line. The actual three contour frequencies and independent ordinary/witness heights are distinct. The normalized first log test and unit-interval Sobolev control each independent height supremum; common compact log-Fourier envelopes or actual scale derivatives control dilation/profile suprema. Choose their finite polynomial degrees, source dependency lists, collars, earlier-anchor cover, response/cutoff banks and meshes before the one positive physical height exponent and every later buffered-grid bound. The fixed construction has positive strict room, so these finite losses, source cost and the separate \(\fitinline{15s_0/16+2F_{\rm pad}}\) charge fit O22--O24. No shared profile is compressed to an upper-scale bump.

The actual finite schedule and primitive inequalities O13--O16 now give the literal I51 cube estimate for every original \texttt{A\allowbreak{}c\allowbreak{}t\allowbreak{}u\allowbreak{}a\allowbreak{}l\allowbreak{}B\allowbreak{}u\allowbreak{}f\allowbreak{}f\allowbreak{}e\allowbreak{}r\allowbreak{}e\allowbreak{}d\allowbreak{}S\allowbreak{}e\allowbreak{}l\allowbreak{}e\allowbreak{}c\allowbreak{}t\allowbreak{}i\allowbreak{}o\allowbreak{}n}, in its original quantifier order. Use O14's F-based padding, not a hard-coded old wrapper. Choose one common \(\fitinline{0<\sigma=P_{\rm src}.\sigma<10^{-9}}\) below all transport reserves. The cube divided by I52's original nonzero normalizer is bounded by \(\fitinline{C_nZ^{\beta-k_{\rm phys}-\sigma}}\). Reproduce the generic high-of-cube conversion through its parameterized finite-nonfloor numerator primitive and the directly proved geometry. It gives the original physical probe and Gram exponent \(\fitinline{t_O-k_{\rm phys}}\); it does not instantiate a record containing the false old h ceiling.

In this contrary case choose \(\fitinline{\omega<(\beta-t_O)/4}\) and \[\fitmath{
       0<m<\min(\beta-t_O-\omega,\sigma,\beta-.87).
}\] H47 then gives a Mellin transform holomorphic on \(\fitinline{\Re s>\beta-m}\). On the absolute line it is the original Gaussian times \(\fitinline{H_\eta(s)/L_\eta(s)}\), with the unchanged \(\fitinline{|H_\eta-1|\le1/2}\) on \(\fitinline{\Re s\ge.87}\). The definition of the actual Hecke-family zero supremum supplies a zero above \(\fitinline{\beta-m/2}\), without assuming the supremum is attained. Its numerator cannot vanish, contradicting holomorphy. The ordinary norm pullback and restoration of every original finite Euler factor prove O0 for all ordinary and imprimitive characters of each fixed positive presented modulus and all finite heights, excluding only the actual principal pole.

\hypertarget{o.7.-exact-sources-and-review-boundary}{%
\subsection{O.7. Exact sources and review boundary}\label{o.7.-exact-sources-and-review-boundary}}

The complete predecessor through N has 961917 bytes, SHA-256 \texttt{5\allowbreak{}2\allowbreak{}4\allowbreak{}a\allowbreak{}e\allowbreak{}c\allowbreak{}5\allowbreak{}d\allowbreak{}6\allowbreak{}4\allowbreak{}6\allowbreak{}8\allowbreak{}d\allowbreak{}9\allowbreak{}5\allowbreak{}c\allowbreak{}0\allowbreak{}3\allowbreak{}f\allowbreak{}f\allowbreak{}9\allowbreak{}5\allowbreak{}2\allowbreak{}f\allowbreak{}9\allowbreak{}b\allowbreak{}1\allowbreak{}1\allowbreak{}8\allowbreak{}3\allowbreak{}0\allowbreak{}6\allowbreak{}f\allowbreak{}5\allowbreak{}7\allowbreak{}1\allowbreak{}e\allowbreak{}9\allowbreak{}2\allowbreak{}a\allowbreak{}4\allowbreak{}b\allowbreak{}f\allowbreak{}3\allowbreak{}4\allowbreak{}b\allowbreak{}0\allowbreak{}4\allowbreak{}e\allowbreak{}f\allowbreak{}b\allowbreak{}6\allowbreak{}f\allowbreak{}3\allowbreak{}9\allowbreak{}a\allowbreak{}0\allowbreak{}6\allowbreak{}c\allowbreak{}c\allowbreak{}6\allowbreak{}4\allowbreak{}7\allowbreak{}4}, and is identified by the original \texttt{c\allowbreak{}y\allowbreak{}c\allowbreak{}l\allowbreak{}e\allowbreak{}0\allowbreak{}8\allowbreak{}6\allowbreak{}/\allowbreak{}i\allowbreak{}n\allowbreak{}t\allowbreak{}e\allowbreak{}g\allowbreak{}r\allowbreak{}a\allowbreak{}t\allowbreak{}i\allowbreak{}o\allowbreak{}n\allowbreak{}/\allowbreak{}P\allowbreak{}R\allowbreak{}O\allowbreak{}O\allowbreak{}F\allowbreak{}\_\allowbreak{}H\allowbreak{}A\allowbreak{}N\allowbreak{}D\allowbreak{}O\allowbreak{}F\allowbreak{}F\allowbreak{}.\allowbreak{}j\allowbreak{}s\allowbreak{}o\allowbreak{}n}. In a successor packet its canonical retained proof is resolved by that exact identity; the owning packager alone registers any \texttt{h\allowbreak{}i\allowbreak{}s\allowbreak{}t\allowbreak{}o\allowbreak{}r\allowbreak{}y\allowbreak{}/\allowbreak{}p\allowbreak{}r\allowbreak{}o\allowbreak{}o\allowbreak{}f\allowbreak{}\_\allowbreak{}v\allowbreak{}e\allowbreak{}r\allowbreak{}s\allowbreak{}i\allowbreak{}o\allowbreak{}n\allowbreak{}s\allowbreak{}/\allowbreak{}v\allowbreak{}i\allowbreak{}r\allowbreak{}t\allowbreak{}u\allowbreak{}a\allowbreak{}l\allowbreak{}/\allowbreak{}5\allowbreak{}2\allowbreak{}4\allowbreak{}a\allowbreak{}e\allowbreak{}c\allowbreak{}5\allowbreak{}d\allowbreak{}6\allowbreak{}4\allowbreak{}6\allowbreak{}8\allowbreak{}d\allowbreak{}9\allowbreak{}5\allowbreak{}c\allowbreak{}0\allowbreak{}3\allowbreak{}f\allowbreak{}f\allowbreak{}9\allowbreak{}5\allowbreak{}2\allowbreak{}f\allowbreak{}9\allowbreak{}b\allowbreak{}1\allowbreak{}1\allowbreak{}8\allowbreak{}3\allowbreak{}0\allowbreak{}6\allowbreak{}f\allowbreak{}5\allowbreak{}7\allowbreak{}1\allowbreak{}e\allowbreak{}9\allowbreak{}2\allowbreak{}a\allowbreak{}4\allowbreak{}b\allowbreak{}f\allowbreak{}3\allowbreak{}4\allowbreak{}b\allowbreak{}0\allowbreak{}4\allowbreak{}e\allowbreak{}f\allowbreak{}b\allowbreak{}6\allowbreak{}f\allowbreak{}3\allowbreak{}9\allowbreak{}a\allowbreak{}0\allowbreak{}6\allowbreak{}c\allowbreak{}c\allowbreak{}6\allowbreak{}4\allowbreak{}7\allowbreak{}4\allowbreak{}.\allowbreak{}m\allowbreak{}d} alias. No root-local integration filename is asserted to be an archive member. Retain the complete A--N proof and all immutable prior source and exact-check evidence.

The new source derivations are \texttt{r\allowbreak{}e\allowbreak{}s\allowbreak{}e\allowbreak{}a\allowbreak{}r\allowbreak{}c\allowbreak{}h\allowbreak{}/\allowbreak{}l\allowbreak{}o\allowbreak{}c\allowbreak{}a\allowbreak{}l\allowbreak{}\_\allowbreak{}c\allowbreak{}y\allowbreak{}c\allowbreak{}l\allowbreak{}e\allowbreak{}0\allowbreak{}8\allowbreak{}7\allowbreak{}/\allowbreak{}n\allowbreak{}e\allowbreak{}a\allowbreak{}r\allowbreak{}\_\allowbreak{}r\allowbreak{}e\allowbreak{}c\allowbreak{}o\allowbreak{}r\allowbreak{}d\allowbreak{}/\allowbreak{}P\allowbreak{}A\allowbreak{}I\allowbreak{}R\allowbreak{}\_\allowbreak{}R\allowbreak{}E\allowbreak{}F\allowbreak{}E\allowbreak{}R\allowbreak{}E\allowbreak{}N\allowbreak{}C\allowbreak{}E\allowbreak{}.\allowbreak{}m\allowbreak{}d}, SHA-256 \texttt{a\allowbreak{}d\allowbreak{}d\allowbreak{}b\allowbreak{}d\allowbreak{}9\allowbreak{}c\allowbreak{}0\allowbreak{}7\allowbreak{}d\allowbreak{}0\allowbreak{}8\allowbreak{}9\allowbreak{}0\allowbreak{}9\allowbreak{}e\allowbreak{}3\allowbreak{}f\allowbreak{}3\allowbreak{}2\allowbreak{}3\allowbreak{}c\allowbreak{}2\allowbreak{}5\allowbreak{}7\allowbreak{}a\allowbreak{}c\allowbreak{}b\allowbreak{}5\allowbreak{}8\allowbreak{}6\allowbreak{}7\allowbreak{}9\allowbreak{}9\allowbreak{}8\allowbreak{}9\allowbreak{}8\allowbreak{}e\allowbreak{}2\allowbreak{}2\allowbreak{}f\allowbreak{}c\allowbreak{}e\allowbreak{}a\allowbreak{}f\allowbreak{}1\allowbreak{}4\allowbreak{}8\allowbreak{}6\allowbreak{}e\allowbreak{}0\allowbreak{}7\allowbreak{}f\allowbreak{}9\allowbreak{}a\allowbreak{}f\allowbreak{}4\allowbreak{}3\allowbreak{}3\allowbreak{}b\allowbreak{}5\allowbreak{}a\allowbreak{}8\allowbreak{}f}; \texttt{H\allowbreak{}\_\allowbreak{}G\allowbreak{}E\allowbreak{}O\allowbreak{}M\allowbreak{}E\allowbreak{}T\allowbreak{}R\allowbreak{}Y\allowbreak{}.\allowbreak{}m\allowbreak{}d}, \texttt{4\allowbreak{}8\allowbreak{}c\allowbreak{}3\allowbreak{}6\allowbreak{}7\allowbreak{}a\allowbreak{}5\allowbreak{}4\allowbreak{}a\allowbreak{}8\allowbreak{}0\allowbreak{}e\allowbreak{}3\allowbreak{}6\allowbreak{}4\allowbreak{}2\allowbreak{}e\allowbreak{}9\allowbreak{}8\allowbreak{}a\allowbreak{}1\allowbreak{}d\allowbreak{}e\allowbreak{}3\allowbreak{}8\allowbreak{}a\allowbreak{}f\allowbreak{}d\allowbreak{}5\allowbreak{}5\allowbreak{}a\allowbreak{}5\allowbreak{}c\allowbreak{}4\allowbreak{}b\allowbreak{}d\allowbreak{}4\allowbreak{}9\allowbreak{}d\allowbreak{}2\allowbreak{}c\allowbreak{}f\allowbreak{}b\allowbreak{}3\allowbreak{}3\allowbreak{}9\allowbreak{}d\allowbreak{}c\allowbreak{}d\allowbreak{}f\allowbreak{}b\allowbreak{}5\allowbreak{}5\allowbreak{}5\allowbreak{}1\allowbreak{}7\allowbreak{}a\allowbreak{}e\allowbreak{}f\allowbreak{}5\allowbreak{}8\allowbreak{}1\allowbreak{}a}; and \texttt{F\allowbreak{}I\allowbreak{}X\allowbreak{}E\allowbreak{}D\allowbreak{}\_\allowbreak{}D\allowbreak{}E\allowbreak{}T\allowbreak{}E\allowbreak{}C\allowbreak{}T\allowbreak{}O\allowbreak{}R\allowbreak{}.\allowbreak{}m\allowbreak{}d}, \texttt{5\allowbreak{}5\allowbreak{}0\allowbreak{}a\allowbreak{}a\allowbreak{}8\allowbreak{}e\allowbreak{}b\allowbreak{}3\allowbreak{}4\allowbreak{}d\allowbreak{}7\allowbreak{}d\allowbreak{}9\allowbreak{}c\allowbreak{}0\allowbreak{}d\allowbreak{}8\allowbreak{}f\allowbreak{}c\allowbreak{}6\allowbreak{}7\allowbreak{}9\allowbreak{}1\allowbreak{}a\allowbreak{}1\allowbreak{}c\allowbreak{}c\allowbreak{}f\allowbreak{}a\allowbreak{}d\allowbreak{}c\allowbreak{}9\allowbreak{}1\allowbreak{}1\allowbreak{}b\allowbreak{}d\allowbreak{}3\allowbreak{}1\allowbreak{}8\allowbreak{}d\allowbreak{}8\allowbreak{}2\allowbreak{}1\allowbreak{}6\allowbreak{}d\allowbreak{}6\allowbreak{}5\allowbreak{}2\allowbreak{}c\allowbreak{}1\allowbreak{}7\allowbreak{}7\allowbreak{}c\allowbreak{}5\allowbreak{}9\allowbreak{}4\allowbreak{}a\allowbreak{}8\allowbreak{}c\allowbreak{}8\allowbreak{}2\allowbreak{}0\allowbreak{}1}. Their paired receipt is \texttt{P\allowbreak{}A\allowbreak{}I\allowbreak{}R\allowbreak{}E\allowbreak{}D\allowbreak{}\_\allowbreak{}E\allowbreak{}X\allowbreak{}T\allowbreak{}E\allowbreak{}N\allowbreak{}S\allowbreak{}I\allowbreak{}O\allowbreak{}N\allowbreak{}\_\allowbreak{}F\allowbreak{}I\allowbreak{}N\allowbreak{}A\allowbreak{}L\allowbreak{}.\allowbreak{}j\allowbreak{}s\allowbreak{}o\allowbreak{}n}, \texttt{f\allowbreak{}0\allowbreak{}a\allowbreak{}5\allowbreak{}e\allowbreak{}4\allowbreak{}e\allowbreak{}c\allowbreak{}1\allowbreak{}b\allowbreak{}c\allowbreak{}e\allowbreak{}6\allowbreak{}1\allowbreak{}b\allowbreak{}d\allowbreak{}c\allowbreak{}7\allowbreak{}7\allowbreak{}9\allowbreak{}0\allowbreak{}1\allowbreak{}d\allowbreak{}a\allowbreak{}6\allowbreak{}f\allowbreak{}e\allowbreak{}c\allowbreak{}f\allowbreak{}c\allowbreak{}e\allowbreak{}4\allowbreak{}a\allowbreak{}f\allowbreak{}0\allowbreak{}e\allowbreak{}7\allowbreak{}5\allowbreak{}b\allowbreak{}b\allowbreak{}4\allowbreak{}1\allowbreak{}3\allowbreak{}6\allowbreak{}0\allowbreak{}d\allowbreak{}e\allowbreak{}7\allowbreak{}7\allowbreak{}a\allowbreak{}4\allowbreak{}f\allowbreak{}6\allowbreak{}9\allowbreak{}1\allowbreak{}d\allowbreak{}9\allowbreak{}2\allowbreak{}f\allowbreak{}d\allowbreak{}0\allowbreak{}b\allowbreak{}4\allowbreak{}5}; the primitive geometry receipt is \texttt{G\allowbreak{}E\allowbreak{}O\allowbreak{}M\allowbreak{}E\allowbreak{}T\allowbreak{}R\allowbreak{}Y\allowbreak{}\_\allowbreak{}F\allowbreak{}I\allowbreak{}N\allowbreak{}A\allowbreak{}L\allowbreak{}.\allowbreak{}j\allowbreak{}s\allowbreak{}o\allowbreak{}n}, \texttt{9\allowbreak{}f\allowbreak{}3\allowbreak{}9\allowbreak{}6\allowbreak{}f\allowbreak{}6\allowbreak{}e\allowbreak{}9\allowbreak{}6\allowbreak{}9\allowbreak{}a\allowbreak{}a\allowbreak{}a\allowbreak{}9\allowbreak{}1\allowbreak{}1\allowbreak{}a\allowbreak{}c\allowbreak{}7\allowbreak{}7\allowbreak{}4\allowbreak{}e\allowbreak{}3\allowbreak{}a\allowbreak{}3\allowbreak{}d\allowbreak{}0\allowbreak{}7\allowbreak{}d\allowbreak{}b\allowbreak{}d\allowbreak{}b\allowbreak{}2\allowbreak{}5\allowbreak{}5\allowbreak{}8\allowbreak{}2\allowbreak{}2\allowbreak{}f\allowbreak{}9\allowbreak{}b\allowbreak{}8\allowbreak{}5\allowbreak{}0\allowbreak{}e\allowbreak{}a\allowbreak{}c\allowbreak{}a\allowbreak{}a\allowbreak{}7\allowbreak{}7\allowbreak{}a\allowbreak{}7\allowbreak{}4\allowbreak{}4\allowbreak{}1\allowbreak{}8\allowbreak{}5\allowbreak{}8\allowbreak{}7\allowbreak{}6\allowbreak{}b\allowbreak{}b}; the full fixed-detector receipt is \texttt{F\allowbreak{}I\allowbreak{}X\allowbreak{}E\allowbreak{}D\allowbreak{}\_\allowbreak{}D\allowbreak{}E\allowbreak{}T\allowbreak{}E\allowbreak{}C\allowbreak{}T\allowbreak{}O\allowbreak{}R\allowbreak{}\_\allowbreak{}F\allowbreak{}I\allowbreak{}N\allowbreak{}A\allowbreak{}L\allowbreak{}.\allowbreak{}j\allowbreak{}s\allowbreak{}o\allowbreak{}n}, \texttt{c\allowbreak{}f\allowbreak{}5\allowbreak{}e\allowbreak{}6\allowbreak{}a\allowbreak{}4\allowbreak{}8\allowbreak{}6\allowbreak{}c\allowbreak{}f\allowbreak{}e\allowbreak{}9\allowbreak{}b\allowbreak{}c\allowbreak{}e\allowbreak{}3\allowbreak{}2\allowbreak{}4\allowbreak{}9\allowbreak{}1\allowbreak{}2\allowbreak{}9\allowbreak{}4\allowbreak{}7\allowbreak{}8\allowbreak{}e\allowbreak{}b\allowbreak{}d\allowbreak{}a\allowbreak{}b\allowbreak{}f\allowbreak{}6\allowbreak{}7\allowbreak{}c\allowbreak{}b\allowbreak{}0\allowbreak{}0\allowbreak{}d\allowbreak{}6\allowbreak{}6\allowbreak{}9\allowbreak{}d\allowbreak{}0\allowbreak{}3\allowbreak{}8\allowbreak{}1\allowbreak{}c\allowbreak{}4\allowbreak{}2\allowbreak{}9\allowbreak{}a\allowbreak{}4\allowbreak{}b\allowbreak{}c\allowbreak{}7\allowbreak{}c\allowbreak{}8\allowbreak{}9\allowbreak{}2\allowbreak{}7\allowbreak{}4\allowbreak{}b\allowbreak{}c}. The actual 4097-cutoff physical receipt is \texttt{R\allowbreak{}E\allowbreak{}T\allowbreak{}U\allowbreak{}R\allowbreak{}N\allowbreak{}\_\allowbreak{}F\allowbreak{}I\allowbreak{}N\allowbreak{}A\allowbreak{}L\allowbreak{}.\allowbreak{}j\allowbreak{}s\allowbreak{}o\allowbreak{}n}, \texttt{0\allowbreak{}6\allowbreak{}8\allowbreak{}7\allowbreak{}8\allowbreak{}9\allowbreak{}3\allowbreak{}0\allowbreak{}8\allowbreak{}2\allowbreak{}2\allowbreak{}a\allowbreak{}3\allowbreak{}e\allowbreak{}9\allowbreak{}4\allowbreak{}4\allowbreak{}d\allowbreak{}9\allowbreak{}c\allowbreak{}6\allowbreak{}7\allowbreak{}f\allowbreak{}f\allowbreak{}1\allowbreak{}a\allowbreak{}b\allowbreak{}1\allowbreak{}a\allowbreak{}c\allowbreak{}2\allowbreak{}4\allowbreak{}c\allowbreak{}6\allowbreak{}a\allowbreak{}2\allowbreak{}e\allowbreak{}3\allowbreak{}3\allowbreak{}4\allowbreak{}c\allowbreak{}a\allowbreak{}6\allowbreak{}2\allowbreak{}6\allowbreak{}a\allowbreak{}2\allowbreak{}2\allowbreak{}8\allowbreak{}d\allowbreak{}7\allowbreak{}5\allowbreak{}4\allowbreak{}b\allowbreak{}d\allowbreak{}d\allowbreak{}d\allowbreak{}5\allowbreak{}c\allowbreak{}7\allowbreak{}5\allowbreak{}a\allowbreak{}d\allowbreak{}4}, with coverage digest \texttt{f\allowbreak{}d\allowbreak{}f\allowbreak{}a\allowbreak{}9\allowbreak{}b\allowbreak{}0\allowbreak{}0\allowbreak{}3\allowbreak{}2\allowbreak{}2\allowbreak{}a\allowbreak{}f\allowbreak{}3\allowbreak{}a\allowbreak{}2\allowbreak{}d\allowbreak{}b\allowbreak{}2\allowbreak{}4\allowbreak{}c\allowbreak{}6\allowbreak{}e\allowbreak{}f\allowbreak{}b\allowbreak{}7\allowbreak{}9\allowbreak{}d\allowbreak{}d\allowbreak{}6\allowbreak{}f\allowbreak{}2\allowbreak{}f\allowbreak{}0\allowbreak{}3\allowbreak{}a\allowbreak{}7\allowbreak{}6\allowbreak{}a\allowbreak{}6\allowbreak{}8\allowbreak{}9\allowbreak{}d\allowbreak{}a\allowbreak{}4\allowbreak{}5\allowbreak{}a\allowbreak{}f\allowbreak{}d\allowbreak{}4\allowbreak{}1\allowbreak{}d\allowbreak{}e\allowbreak{}3\allowbreak{}a\allowbreak{}4\allowbreak{}0\allowbreak{}d\allowbreak{}3\allowbreak{}d\allowbreak{}2\allowbreak{}c\allowbreak{}2\allowbreak{}3}. The earlier \texttt{E\allowbreak{}X\allowbreak{}T\allowbreak{}E\allowbreak{}N\allowbreak{}D\allowbreak{}E\allowbreak{}D\allowbreak{}8\allowbreak{}5\allowbreak{}\_\allowbreak{}1\allowbreak{}2\allowbreak{}6\allowbreak{}6\allowbreak{}6\allowbreak{}7\allowbreak{}\_\allowbreak{}H\allowbreak{}2\allowbreak{}0\allowbreak{}0\allowbreak{}0\allowbreak{}\_\allowbreak{}F\allowbreak{}I\allowbreak{}N\allowbreak{}E\allowbreak{}.\allowbreak{}j\allowbreak{}s\allowbreak{}o\allowbreak{}n} is retained as original evidence; the final receipt binds the repaired reference metadata at the same scientific rows. Source-use maps record the exact note/program identities and source-relative replay paths. Finite arithmetic proves its stated inequalities and coverage, not source admission, height integration or a Lean run; O.1--O.6 supply the corresponding written implications.

The existing N two-epoch proof and its independent exact source/physical supplements remain unchanged. New paired stages use the literal retained PR141 entry chain and original source interfaces with their notices; no old Lean theorem is reported to accept the changed reference face or conductor. The catalogue's paired-anchor, ordinary-reference and actual averaged-floor components are provenance/discovery contracts, not a replacement for the reproof above.

The frozen near-source manifest is \texttt{r\allowbreak{}e\allowbreak{}s\allowbreak{}e\allowbreak{}a\allowbreak{}r\allowbreak{}c\allowbreak{}h\allowbreak{}/\allowbreak{}l\allowbreak{}o\allowbreak{}c\allowbreak{}a\allowbreak{}l\allowbreak{}\_\allowbreak{}c\allowbreak{}y\allowbreak{}c\allowbreak{}l\allowbreak{}e\allowbreak{}0\allowbreak{}8\allowbreak{}7\allowbreak{}/\allowbreak{}n\allowbreak{}e\allowbreak{}a\allowbreak{}r\allowbreak{}\_\allowbreak{}r\allowbreak{}e\allowbreak{}c\allowbreak{}o\allowbreak{}r\allowbreak{}d\allowbreak{}/\allowbreak{}S\allowbreak{}O\allowbreak{}U\allowbreak{}R\allowbreak{}C\allowbreak{}E\allowbreak{}\_\allowbreak{}U\allowbreak{}S\allowbreak{}E\allowbreak{}\_\allowbreak{}F\allowbreak{}I\allowbreak{}N\allowbreak{}A\allowbreak{}L\allowbreak{}.\allowbreak{}j\allowbreak{}s\allowbreak{}o\allowbreak{}n}, SHA-256 \texttt{2\allowbreak{}3\allowbreak{}8\allowbreak{}8\allowbreak{}7\allowbreak{}1\allowbreak{}a\allowbreak{}6\allowbreak{}b\allowbreak{}c\allowbreak{}4\allowbreak{}3\allowbreak{}e\allowbreak{}5\allowbreak{}5\allowbreak{}6\allowbreak{}2\allowbreak{}a\allowbreak{}6\allowbreak{}b\allowbreak{}a\allowbreak{}6\allowbreak{}0\allowbreak{}4\allowbreak{}d\allowbreak{}f\allowbreak{}5\allowbreak{}5\allowbreak{}8\allowbreak{}0\allowbreak{}6\allowbreak{}8\allowbreak{}c\allowbreak{}5\allowbreak{}e\allowbreak{}a\allowbreak{}3\allowbreak{}5\allowbreak{}b\allowbreak{}c\allowbreak{}3\allowbreak{}f\allowbreak{}9\allowbreak{}5\allowbreak{}1\allowbreak{}a\allowbreak{}7\allowbreak{}0\allowbreak{}a\allowbreak{}d\allowbreak{}a\allowbreak{}f\allowbreak{}1\allowbreak{}1\allowbreak{}4\allowbreak{}d\allowbreak{}1\allowbreak{}d\allowbreak{}b\allowbreak{}b\allowbreak{}b\allowbreak{}b\allowbreak{}c\allowbreak{}3}; its handoff is \texttt{P\allowbreak{}R\allowbreak{}O\allowbreak{}O\allowbreak{}F\allowbreak{}\_\allowbreak{}H\allowbreak{}A\allowbreak{}N\allowbreak{}D\allowbreak{}O\allowbreak{}F\allowbreak{}F\allowbreak{}.\allowbreak{}m\allowbreak{}d}, \texttt{9\allowbreak{}1\allowbreak{}9\allowbreak{}4\allowbreak{}1\allowbreak{}0\allowbreak{}0\allowbreak{}3\allowbreak{}7\allowbreak{}a\allowbreak{}8\allowbreak{}d\allowbreak{}b\allowbreak{}e\allowbreak{}5\allowbreak{}4\allowbreak{}2\allowbreak{}9\allowbreak{}d\allowbreak{}0\allowbreak{}4\allowbreak{}f\allowbreak{}0\allowbreak{}b\allowbreak{}b\allowbreak{}e\allowbreak{}2\allowbreak{}e\allowbreak{}2\allowbreak{}7\allowbreak{}4\allowbreak{}c\allowbreak{}9\allowbreak{}6\allowbreak{}a\allowbreak{}e\allowbreak{}2\allowbreak{}d\allowbreak{}3\allowbreak{}9\allowbreak{}b\allowbreak{}3\allowbreak{}5\allowbreak{}3\allowbreak{}3\allowbreak{}c\allowbreak{}9\allowbreak{}2\allowbreak{}5\allowbreak{}4\allowbreak{}3\allowbreak{}a\allowbreak{}e\allowbreak{}5\allowbreak{}c\allowbreak{}d\allowbreak{}1\allowbreak{}3\allowbreak{}1\allowbreak{}d\allowbreak{}d\allowbreak{}8\allowbreak{}0\allowbreak{}0}. They bind fifteen owned proof/program/receipt/handoff files and 53 retained primary inputs, including the actual padding definitions and per-character high-source primitives. The curated front and \texttt{C\allowbreak{}H\allowbreak{}E\allowbreak{}C\allowbreak{}K\allowbreak{}\_\allowbreak{}F\allowbreak{}I\allowbreak{}N\allowbreak{}A\allowbreak{}L\allowbreak{}.\allowbreak{}p\allowbreak{}y} specify the four-receipt portable reconstruction; preparatory scouts, superseded proof notes and repeated replay outputs remain historical evidence outside the active payload.

Preserve the full proof since OpenAI and permission for properly derived replacements, all original OpenAI/Collet/Sil/Bebel/Levy/Gao/Liu and other contributor credits, licenses and reported assistance, retained PR43 source terms and Codex disclosure, and Mathematical Writing for AI Agents / Sepia-CWR attribution. This appendix is the single amendment span for the parent's joined proof. Separate open high-ceiling research does not acquire endpoint or formal-verification status by this composition.

\hypertarget{appendix-p}{}

\hypertarget{appendix-p.-ordered-ordinary-caps-and-the-complete-physical-return}{%
\section{Appendix P. Ordered ordinary caps and the complete physical return}\label{appendix-p.-ordered-ordinary-caps-and-the-complete-physical-return}}

Retain the complete proof through O, including the adopted natural selected-cell input H6/SN9, the proper initial keep, all six units, every fixed Euler puncture and the actual physical reciprocal. The strict ordinary result of this appendix is at \[\fitmath{
 c_P=\frac{25333441}{200000000}=.126667205,\qquad
 t_P=\frac{174666559}{200000000}=.873332795,\qquad
 h=h_*+.00225953.                                             
}\tag{P0}\] For each fixed positive presented modulus and every ordinary character, including an imprimitive presentation, there is no zero in \(\fitinline{\Re s>t_P}\) at any finite height, with the actual principal pole excepted. The gain over O is \(\fitinline{41/200000000}\) in \(\fitinline{c}\). There is no assertion on the boundary or of uniformity as the modulus grows. The inherited standing input has its previous scope.

Two changes pay the return. In the positive two-ordinary source, order the upper length labels and use the earlier short label even when the ordinary witness and the retained ordinary exchange roles. In the physical estimate, clear the current paired count at \(\fitinline{D=h}\), then bound the remaining debit in the actual conductor units. O's paired sources, fixed detector, and primitive transport remain available with their original scopes. O also gives \(\fitinline{\beta\le t_O=.873333}\) before ordinary pullback; thus, in the contrary case, \(\fitinline{t_P<\beta\le t_O}\). The actual moment parameter remains \(\fitinline{2\beta-1}\); the contact upper parameter used in the arithmetic is a conservative price. The open higher-target arguments are independent of this proof.

\hypertarget{p.1.-the-ordered-two-ordinary-source}{%
\subsection{P.1. The ordered two-ordinary source}\label{p.1.-the-ordered-two-ordinary-source}}

Set \(\fitinline{Y_0=.369}\), \(\fitinline{Y=.3888}\), \(\fitinline{a=.3735}\), and \(\fitinline{\lambda=1.15}\). With \(\fitinline{y\le n\le\lambda y}\) and \(\fitinline{Y_0\le y\le Y}\), the pre-division positive two-ordinary marked energy of N.1 admits every nonnegative width strictly less than \[\fitmath{
 \begin{aligned}
 W(n,y)&=(1-n-y)/2+F(y),\\
 F(y)&=\begin{cases}
 4.618(Y-y)^2+50.5(Y-y)^3+.00031(Y-y)-10^{-7},&y\ge a,\\
 F(a)+.174(a-y)+6.7(a-y)^2,&y<a.
 \end{cases}
 \end{aligned}                                                  
}\tag{P1}\] It admits all componentwise shortenings of the two labelled caps. In H.5's original positive sum, the two factors are the actual norm-normalized ordinary sums on their respective supported ideals; the prime factor is the surviving subproduct of the original whole slots. The proper initial keep, holes and actual phases remain, with independent ordinary heights, annular profiles and normalized first log tests. This is a bound on the original marked energy before division by the complete slot amplitude. Ordering the labels does not order the actual lengths or heights. Compared with N.1, the source gains \(\fitinline{.028(Y-y)^2}\) above \(\fitinline{a}\), and \(\fitinline{.028(Y-a)^2+.019(a-y)+6.7(a-y)^2}\) below it, on the common domain. The new ratio interval includes \(\fitinline{[1,1.118)}\).

Use the positive regular sixth-free core of K.3 and N.2 at radius \(\fitinline{U^V}\), after fixing its actual cofactor. Let \(\fitinline{0\le e\le\nu\le3/4}\) be the exponents of the clean maximal squared slot response and the common zero band. Put \(\fitinline{u=3\cdot10^{-6}}\), \(\fitinline{H=1-6u=.999982}\), \(\fitinline{\rho=2\cdot10^{-11}}\) and \(\fitinline{L=W(n,y)-u}\), up to the early mesh. The regular input needed by H.5 is normalized excess at most \[\fitmath{
                    u-\rho+\tfrac56(H-V)_+.                  
}\tag{P2}\] For \(\fitinline{Y\le y\le Y/H}\) the two native fourths at their mean width give the starting strip: the continued cubic is convex and negative at both endpoints. The genuine native fourths are the reserve family \(\fitinline{G(1-s_0,2+10s_0,3/20-1000s_0/1143)}\), with \(\fitinline{s_0>0}\) sufficiently small for each prescribed decrement, plus the true unmarked cases when needed. Induct over the closed descending intervals with endpoints \(\fitinline{YH^j}\), clipping the last at \(\fitinline{Y_0}\). At \(\fitinline{V=H}\) the low response uses the earlier pair \(\fitinline{(n/H,y/H)}\). Its excess is at most \[\fitmath{
                   e[2u+F(y)-HF(y/H)]_+.                    
}\tag{P3}\] Indeed \(\fitinline{L-HW(n/H,y/H)}\) equals the expression in brackets: the two absolute length charges cancel. Keep an inward-rounded whole-slot subset of the actual live product on its original ray. The remaining absolute widths cost at most \(\fitinline{e}\) times their sum on the positive frozen subclass. No lower estimate is placed on a Fourier summand.

For high response fix a rational cutoff \(\fitinline{T=r_0+m_0}\) in the finite early length and response bank. B.5.1 supplies an inverse and an ordinary witness at one actual zero, with the same presentation, \(\fitinline{r\le T+o(1)}\), \(\fitinline{r+m\ge T-o(1)}\), \(\fitinline{m\le1/2+o(1)}\), and \(\fitinline{|S_m|^2\gg U^{V\nu m-o(1)}}\). The original ordinary under label \(\fitinline{y}\) has the squared bound \(\fitinline{U^{\nu y+o(1)}}\). Hence on that frozen subclass, with the original other factor \(\fitinline{A_n}\) and clean mark product \(\fitinline{Q}\), \[\fitmath{
 |A_nA_yQ|^2\ll U^{\nu(y-Vm)+o(1)} |A_nS_mQ|^2.               
}\tag{P4}\] The witness has only its actual fixed deleted base. Variable cofactor holes stay on the original ordinary factor and are opened by the exact K.3 identities at their shortened scales, never inserted into an inverse inlet. Re-anchor the witness and the retained ordinary in the fixed dictionary, with each anchor's own relative multiplier; if orientation requires conjugation, conjugate the whole witness including its height, mask and test. Invert each actual compact log profile with its own outside norm phase. Now permute the two tuples consisting of the upper cap, actual ordinary factor and mask, fixed ray and relative multiplier, annular profile, outside norm phase and height. On this subclass the product is identical after the permutation; \(\fitinline{Q}\) and the keep are unaltered. A good prime in the primitive regular core does not disappear under a fixed finite ray. Fixed-inducing rows never use P4: they use the original signed comparison on the two original rectangles in P.3.

At \(\fitinline{H}\), if \(\fitinline{m\ge m_0\ge y/H}\) and \(\fitinline{m\le Y}\), the ordered caps \(\fitinline{n/H,m}\) have ratio at most \(\fitinline{\lambda}\): \(\fitinline{(n/H)/m\le\lambda}\) and \(\fitinline{m/(n/H)\le Y/Y_0<\lambda}\). Their shorter cap \(\fitinline{s=\min(n/H,m)}\) is at least \(\fitinline{y/H}\), hence already proved or in the starting strip. Since \(\fitinline{F}\) decreases, \(\fitinline{F(s)\ge F(m)}\). With native mean quota for \(\fitinline{m>Y}\), the cost in P4 is bounded by \[\fitmath{
 C_o(H;m)=-\nu(Hm-y)+e[2u+(Hm-y)/2+F(y)-HJ(m)]_+,
 \qquad J(m)=\begin{cases}F(m),&m\le Y,\\0,&m>Y.\end{cases}   
}\tag{P5}\] The native branch uses the true near-unit unmarked collar. No improved width is used above \(\fitinline{Y}\). The active derivative in \(\fitinline{m}\) is \(\fitinline{H[-\nu+e(1/2-F'(m))]\le0}\) below \(\fitinline{Y}\), and the inactive derivative is \(\fitinline{-H\nu}\). Both one-sided slopes use \(\fitinline{-.25<F'<0}\). The change from \(\fitinline{F(Y)<0}\) to zero at the native continuation is favorable. Thus the checked cost at \(\fitinline{m_0}\) pays every larger actual ordinary witness.

If \(\fitinline{r\ge1}\), B(204)'s genuine amplified count gives deficit \(\fitinline{D_l=A-(A-\nu)T}\), with \(\fitinline{A=5/6}\). If \(\fitinline{r_0\le r<1}\), the deficit is at least \(\fitinline{D_i=\nu r_0+e q(r_0)}\), with \(\fitinline{q(w)=(1-w)/2+I(w)}\) as in N8: E(25)'s lower tangent for \(\fitinline{.734\le w<.75}\), E(19)'s actual curve up to \(\fitinline{7/9}\) and the genuine native and unmarked continuation thereafter. The deficit \(\fitinline{\nu w+e q(w)}\) increases on each arm, with the favorable tangent-to-curve transition at \(\fitinline{.75}\); this uses I.3's derivative and the tangential slope on \(\fitinline{\nu=e}\). The checked cutoffs lie in \(\fitinline{[1.1126804197,1.1664209957]}\), so the short inverse remains in its untapered domain. The count on the same rows, the retained longer-ordinary native fourth and the buffered shorter pointwise estimate give the H23 price \[\fitmath{
             C_D(V)=\nu y+e[L-(V/2-n)/2]-VD/2                 
}\tag{P6}\] for the indicated deficit \(\fitinline{D}\). This credit uses H23's repetitions of the same attaining mark from the genuine native source; P3 and P5 instead keep the positive part in case the original marked subproduct is exhausted. The last case \(\fitinline{r<r_0}\) forces \(\fitinline{m\ge m_0-o(1)}\) from the same pair. For a collar \(\fitinline{m\ge m_0-\gamma_0}\), use the dominating cap \(\fitinline{\bar m=\max(m,m_0)}\) in the ordered pair while retaining the actual \(\fitinline{-V\nu m}\) in P4. The bounds on the width slopes below and the factor \(\fitinline{\nu V}\) charge \(\fitinline{O(\gamma_0)}\) to the preassigned strict collars. This also works across \(\fitinline{Y}\) by the favorable native transition. These three cases exhaust the actual witness pair.

\hypertarget{p.2.-directed-checks-radial-continuation-and-uniformity}{%
\subsection{P.2. Directed checks, radial continuation and uniformity}\label{p.2.-directed-checks-radial-continuation-and-uniformity}}

On each of 3960 closed length cells \(\fitinline{[b_0,b]}\) of width \(\fitinline{.000005}\), the low bracket in P3 is maximal at an endpoint or \(\fitinline{Ha}\) if it crosses the seam: on a common convex piece it decreases, while the crossing derivative is positive since \(\fitinline{F'_-(Ha)>F'_+(a/H)}\). Round \(\fitinline{(u-3\rho)}\) divided by this maximum down to a rational threshold \(\fitinline{k\le.5002}\). The remaining response range, whose lower end exceeds \(\fitinline{.38}\), is covered by eight closed bins with a rational cutoff per bin. Their tests have \(\fitinline{r_0\ge2b/H}\), \(\fitinline{m_0\ge b/H}\), \(\fitinline{HT/2>b}\). Thus P5 and P6 are largest on \(\fitinline{\nu=e}\) for fixed \(\fitinline{e}\) and are nondecreasing in \(\fitinline{y}\); on an active branch the \(\fitinline{y}\)-derivative is \(\fitinline{\nu-e/2+eF'(y)\ge0}\). Test \(\fitinline{y=b}\), \(\fitinline{\nu=e}\), and both response endpoints (the resulting functions of \(\fitinline{e}\) are affine). The rational fifth-root enclosures compare adjacent integer fifth powers. The 31,680 tests give reserves above \(\fitinline{\rho}\) respectively greater than \(\fitinline{3.99\cdot10^{-11},6\cdot10^{-11},5.56\cdot10^{-6},6\cdot10^{-11}}\) for low, short inverse, long inverse and ordinary substitution. These strict reserves pay the finite profile collars and inward-rounded losses.

For \(\fitinline{V\ge H}\) in P3 and P5 freeze the whole earlier pair at the checked \(\fitinline{H}\) labels, which dominate the true shorter components. Its positive absolute width grows with \(\fitinline{V}\), and the factor \(\fitinline{\nu(y-Vm)}\) improves while the original upper stays \(\fitinline{\nu y}\). In P6 the genuine native cap for \(\fitinline{V\ge1}\) is \(\fitinline{V(1/2-n)}\), agreeing at one; the count prices also decrease. For \(\fitinline{.98\le V\le H}\), the low case uses \(\fitinline{(n/V,y/V)}\) until \(\fitinline{y/V=Y}\) and then the native pair. Below that seam its absolute width has derivative \(\fitinline{1/2+F(s)-sF'(s)}\) at \(\fitinline{s=y/V}\), between zero and \(\fitinline{.61}\), including the one-sided derivatives at \(\fitinline{a}\).

For the substitute below \(\fitinline{H}\), use \(\fitinline{n/V,m'=\max(\bar m,n/(\lambda V))}\) while \(\fitinline{m'\le Y}\). The shorter reordered cap is at least \(\fitinline{y/H}\). The reverse ratio is at most \(\fitinline{Y/Y_0}\), so this is still an earlier source. When \(\fitinline{m'=\bar m}\), the derivative of its absolute capacity is \(\fitinline{(1-\bar m)/2+F(\bar m)}\) if the witness cap is shorter and \(\fitinline{(1-\bar m)/2+F(n/V)-(n/V)F'(n/V)}\) otherwise. When \(\fitinline{m'=n/(\lambda V)}\), it is \(\fitinline{1/2+F(m')-m'F'(m')}\). These derivatives are positive and below \(\fitinline{.61}\); the caps exchange order continuously. At \(\fitinline{m'=Y}\) the native quota uses the actual \(\fitinline{m\le m'}\), with \(\fitinline{F(Y)<0}\), so the change is favorable in the decreasing-radius direction. For \(\fitinline{m>Y}\) the pair is native throughout. The saturation comparison still uses the actual \(\fitinline{m}\). The total adverse derivative has magnitude at most \(\fitinline{(3/4)/2+.5002(.61)<5/6}\), with the same bound for either derivative across the positive part. P6 has smaller adverse slopes. This continues all three checked prices to P2 below \(\fitinline{H}\).

For \(\fitinline{e>.5002}\) at \(\fitinline{V\ge.98}\), use precisely H6/SN9 at slope \(\fitinline{.5001}\) and its fixed positive margin, on the original merged ordinary and marked coefficient; the squared coefficient mass is paid by the original divisor-fiber bound. For \(\fitinline{.875\le V<.98}\) the two native or genuinely unmarked fourths cost at most \(\fitinline{(3/4)[(1-V)/2+F(Y_0)-u]_+}\), strictly below P2 by the checked edge at \(\fitinline{.98}\) and the smaller slope \(\fitinline{3/8}\). For \(\fitinline{.08\le V<.875}\) use the two unmarked fourths and the genuine prime bound, costing at most \(\fitinline{(3/4)(W_{\max}-u)}\). Below \(\fitinline{.08}\) use the original core count with pointwise bound \(\fitinline{(3/4)(n+y)+W_{\max}-u}\). All edge guards are strict. In \(\fitinline{z=b^6c}\), keep the actual unit and cofactor and prime holes. The \(\fitinline{5/6}\) slope pays the full \(\fitinline{U^{(M-V)/6+o(1)}}\) cofactor multiplicity, without requiring \(\fitinline{(b,c)=1}\). This gives the grouped-ball bound \(\fitinline{U^{M+u-\rho+(5/6)(H-M)_++\epsilon}}\) times the original controls.

There are 2904 induction stages. On preceding closed stages choose finite cap and ratio covers, rounding up within their supported closed edges and taking smaller positive width decrements than the strict reserves. The low and swapped high branches both use short labels at least \(\fitinline{y/H}\), including an endpoint already covered by the preceding stage. Fix these finite covers, geometric losses and the common early mesh before the ambient arithmetic data; after fixing the finite slot family, radial cap, presented characters, internal inducing ideal and ray dictionary, take a finite union of profile, Fourier and log-test controls and maxima of their polynomial degrees. These choices precede live heights, arrays and later buffers. Use separate compact profiles with their actual outside phases and common weighted Fourier envelopes, including first-log tests; bounded or empty lengths retain their original cases. No height is obtained by differentiating a discrete cap choice.

\hypertarget{p.3.-the-short-reference-original-signs-and-raw-sides}{%
\subsection{P.3. The short reference, original signs and raw sides}\label{p.3.-the-short-reference-original-signs-and-raw-sides}}

In the actual equal-product reference of H.2--H.5, replace the short width by \(\fitinline{b_{\rm ref}=.115}\). For a regular child with a coordinate at most \(\fitinline{b_{\rm ref}}\), use the marked short fourth at its actual scale and the other factor's nonprincipal unmarked fourth. If both are longer, use the equal-product reference scales \(\fitinline{(U^{b_{\rm ref}},X_1X_2/U^{b_{\rm ref}})}\), with the two original profiles and pitches after extracting the genuine outside norm phases. The short marked column and its squared coefficient have combined nonprincipal conductor and opened-mask width at most \(\fitinline{4(b_{\rm ref}+W_{\max})+o(1)<H}\), with spare greater than \(\fitinline{.00724}\). The zero-frequency-free smooth Poisson estimate and the unchanged sixth-free diagonal divisor bound pay that fourth. The reference companion has exponent at most \(\fitinline{(1+\lambda)Y-b_{\rm ref}=.72092}\). It uses only the true bounded-length nonprincipal \textbf{unmarked} fourth; the short marked restriction does not apply to it.

On the fixed-inducing rows apply H8 to the original signed rectangles, with the same supported ideals, mark assignment, profiles and independent heights before any substitution. The equal product cancels the common double principal residue where it occurs; H8 bounds each remaining single residue and shifted integral. It gives the \(\fitinline{3/5}\) minimum-length saving with the two genuinely punctured factors. Its square pays the full second divisor extraction of H13--H16, since \(\fitinline{2(3/5)>1}\). After pulling the first divisor of width \(\fitinline{c_1}\) the four-scale minimum contributes \(\fitinline{(b_{\rm ref}-c_1)_+}\), with the original bounded-scale loss. Since \(\fitinline{c_1+(b_{\rm ref}-c_1)_+\ge b_{\rm ref}}\), H18's fixed-row excess is at most \[\fitmath{
                    A_0+3u-5/6-b_{\rm ref}<-.03036,            
}\tag{P7}\] apart from the freely small rounding fee. The source check has \(\fitinline{A_0\le .9179599}\), \(\fitinline{5u<1}\) and \(\fitinline{A_0+4u<1}\). Thus the actual root volume \(\fitinline{\mathcal V\le U^{.9179599}<U}\) and \(\fitinline{(9/4)^2 2^{N_{\rm slots}}\mathcal V<U^2}\) eventually for the fixed slot type, as required by C.10. The fixed small-scale volume case stays separate from the regular short-reference fourth.

For both original raw orientations retain the unweighted first artificial M\"{o}bius divisor and the full joint (row, divisor) Cauchy in H.5. The second M\"{o}bius divisor \(\fitinline{\mathfrak D}\) keeps its sign in the block scalar; its reciprocal label weight cancels the \(\fitinline{N\mathfrak D}\) factor in H26. Both complete seed sums, the actual off-zero supports and counted tails stay in place. Packing uses the original injection into the same grouped ball, with no Jacobian and with \(\fitinline{M=2a_{\rm vol}-k_0+q_C+O(\mathrm{mesh})}\). On a filled selection \(\fitinline{M\ge H-O(\mathrm{mesh})}\), so the P2 bound and P7 pay the disjoint regular and fixed rows in each positive side. At exhaustion use the actual unmarked fourths, or the original tiny-width prime estimate. Allocation restores the original \(\fitinline{NC\,ND_1}\) denominator using the separate left and right H25 identities. H.5's zero and discarded pieces use their own norm and count estimates, and both raw side estimates are inserted before the first signed sum is bounded. This restores the whole uncentered kept marked energy, establishing P1 with the original coefficient normalizers and every smaller admissible quota.

\hypertarget{p.4.-geometry-and-the-current-paired-top-edge}{%
\subsection{P.4. Geometry and the current paired top edge}\label{p.4.-geometry-and-the-current-paired-top-edge}}

Set \(\fitinline{\ell=11/3-4t_P}\), \(\fitinline{x=1+\ell-h}\) and \(\fitinline{y_t=h-2\ell}\). Distinguish \(\fitinline{y_t}\) from the ordinary length labels. The directed conductor endpoints are \(\fitinline{h_- = .817584033}\), \(\fitinline{h_+= .817584035}\). Keep the physical cap \(\fitinline{.085}\) declared below \(\fitinline{.11}\), the genuine native cap, source cost \(\fitinline{s_0=10^{-260}}\) and \(\fitinline{F_{\rm pad}=10^{-258}}\). At these endpoints the checks give the actual primitive inequalities of O.3: \[\fitmath{
 \begin{gathered}
 h>1/2,\quad h+F_{\rm pad}<.865<7/8<1,\quad h+F_{\rm pad}<2y_t,\quad
 \ell/(h+F_{\rm pad})>1/5,\quad y_t/(1/200)<99,\quad 5(h+F_{\rm pad})/288<1/50,\\
 \ell<x<y_t<1/2,\quad 2x-y_t-\ell>0,\quad 5\ell<1,\quad t_P=11/12-\ell/4.
 \end{gathered}                                                
}\tag{P8}\] The signed signal is above \(\fitinline{.1}\). The averaged floor price \(\fitinline{.51-49h/150+.51\ell+.66D}\), at \(\fitinline{D\le h+F_{\rm pad}}\), leaves more than \(\fitinline{10^{-3}}\) after \(\fitinline{10^{-5}+10^{-12}}\). The small-row saving and both principal comparisons are separately strict: for \(\fitinline{0<e_0\le.001}\), \(\fitinline{y_t/20-e_0(1+h)>1/4000}\) and \(\fitinline{h/600-e_0>1/4000}\), as well as the original principal and transport margins in O.3. Thus the literal consumers can be reproduced without the false \(\fitinline{h<.816}\) field. In pinned PR141 the actual \texttt{r\allowbreak{}o\allowbreak{}w\allowbreak{}P\allowbreak{}a\allowbreak{}d\allowbreak{}d\allowbreak{}i\allowbreak{}n\allowbreak{}g} is \(\fitinline{F_{\rm pad}/2}\), \texttt{r\allowbreak{}o\allowbreak{}w\allowbreak{}E\allowbreak{}x\allowbreak{}t\allowbreak{}e\allowbreak{}n\allowbreak{}s\allowbreak{}i\allowbreak{}o\allowbreak{}n} is \(\fitinline{F_{\rm pad}}\), and the high-source cap is \(\fitinline{0<\zeta\le F_{\rm pad}}\). Take \(\fitinline{0<D_{\rm src}.t\le F_{\rm pad}/4}\); the supported rows then have \(\fitinline{D\le h+F_{\rm pad}}\) by \(\fitinline{p_{\rm row}+2D_{\rm src}.t\le F_{\rm pad}}\). The loose historical \(\fitinline{10^{-10}}\) guard does not enter this support argument. Charge \(\fitinline{15s_0/16+2F_{\rm pad}}\) separately, including the genuine extension above \(\fitinline{h}\).

On an admitted paired cell put \[\fitmath{
 v=5/6-d,\quad k=1-p(1+2g_1),\quad U_0=1-p(1+6g_0-4g_1),\quad
 J=3v+2dk,\quad N=J/6+v(J+dU_0),\quad t_0=1+dU_0/J.            
}\tag{P9}\] The true paired count is bounded by \(\fitinline{R=N/J=1/6+v t_0}\), with short crossing \(\fitinline{2t_0/3}\). O's extension supplies the actual paired lower rectangles where needed. The full E29 admission remains, with ten guards at each of its four enclosing contact corners. The count need not minimize the other available counts. At the upper conductor range its current price is \[\fitmath{
 \begin{aligned}
 C(D)&=\min(.085D,[2y_t-(1+s_0)D]/(2+10s_0)),\\
 E(D)&=E_b(D)+3DR/4,\\
 E_b(D)&=1/2+13h/75-y_t/2-\ell/2-.09D+d/2+dp(\ell-C(D)/2).
 \end{aligned}                                                  
}\tag{P10}\] At \(\fitinline{D=h}\), clear the actual margin \(\fitinline{(t_P-E_b(h))J-3hN/4}\). On the admission \(\fitinline{J,U_0>0}\); the count decreases with \(\fitinline{g_0}\) and increases with \(\fitinline{g_1}\), so \(\fitinline{(g_{0,-},g_{1,+})}\) gives the upper count in that clearing. The narrowed enclosures use the isolated E(22) root \(\fitinline{\alpha}\), \(\fitinline{w=(12\alpha-14)/(9(12\alpha-11))}\), and \[\fitmath{
 I(w)=\frac{w+12-(45/2)w^2}{45w-14},\qquad c_w=11w/16-7/24-I(w).
}\] Logarithmic differentiation of the E(19) radical gives \(\fitinline{g_1=(4/(5w)+1/(5(w-2/3)))c_w-11/16}\) and \(\fitinline{g_0=I(w)+wg_1}\). The root signs and outward substitutions are rationally checked. These enclosures lie within the original admitted ones.

For fixed \(\fitinline{d,p}\), the top price \(\fitinline{E(h)}\) is convex in \(\fitinline{h}\), since \(\fitinline{C(h)}\) is a minimum of affine capacities with a nonpositive coefficient in the price. Evaluate at the two directed endpoints. On each closed \(\fitinline{(d,p)}\) box, tensor Bernstein gives a lower bound \(\fitinline{M_-}\) for the cleared margin and \(\fitinline{0<J_-\le J\le J_+}\). Divide by \(\fitinline{J_+}\) when \(\fitinline{M_-\ge0}\) and by \(\fitinline{J_-}\) otherwise. The larger of the resulting lower margin and the old directed E49 margin is still a lower bound on this same admitted count price. Let \(\fitinline{B}\) be the smaller corresponding upper excess. Directly cleared rows with strict room are paid by the original marked numerator fourth.

\hypertarget{p.5.-the-debit-and-the-finite-same-zero-bank}{%
\subsection{P.5. The debit and the finite same-zero bank}\label{p.5.-the-debit-and-the-finite-same-zero-bank}}

Reprice O's verified raw-ticket construction at P0, including its lower conductor bound \(\fitinline{D_-}\), true numerator dilation and actual amplitude, count, auxiliary contour and reciprocal support. For a final box \(\fitinline{[a_d,b_d]\times[f_p,e_p]}\), the slope of P10 is at least \[\fitmath{
 S=\tfrac34(1-b_d)-.09-b_d(.085)/4,
}\] because \(\fitinline{R\ge1-d}\) and the two capacity slopes are at most \(\fitinline{.085}\). Put \(\fitinline{\epsilon_a=10^{-10}}\), \(\fitinline{\epsilon_b=10^{-9}}\), and \[\fitmath{
             q_0=(\max(0,B)+\epsilon_b)/h_- .                 
}\tag{P11}\] The directed check gives \(\fitinline{S>q_0}\). For \(\fitinline{D_-\le D\le h}\) the excess with the inward fee is at most \(\fitinline{B+\epsilon_b-S(h-D)\le Dq_0}\); its positive part also obeys this bound. Thus \(\fitinline{q_0}\) is a debit in the actual units \(\fitinline{U=Z^D}\). The top correction through \(\fitinline{h+F_{\rm pad}}\) is paid by the separate exponent allowance in P.4. At the lower edges of the ticket retain the repriced exterior and raw support estimates. The new inward losses, rounded displacements and thresholds are no larger than those obtained by O's rule at the same parameters, since \(\fitinline{B\le B_{\rm old}}\) and \(\fitinline{h_-\ge D_-}\). Each refined child has its own raw-ticket checks.

Take \(\fitinline{k_-\le k\le k_+}\), \(\fitinline{m_-\le t_0/3\le m_+}\), \(\fitinline{v_-=5/6-b_d}\). Choose \(\fitinline{\gamma,\lambda_x,\lambda_y}\) by rounding upward to \(\fitinline{10^{-13}}\) the respective requirements \[\fitmath{
 \gamma\ge\frac{4q_0}{3v_-},\qquad
 \lambda_x\ge\frac{4q_0}{3a_d k_-},\qquad
 \lambda_y\ge\frac{2q_0}{3a_d k_-}.                            
}\tag{P12}\] There are 16,385 offsets \(\fitinline{\gamma_j=j\gamma/16384}\) for the cutoffs \(\fitinline{t_0-\gamma_j}\). Use the original finite scale bank: for the fixed length and amplitude cells put a sufficiently fine rational grid on the allowed \(\fitinline{t_0}\), and fix the corresponding cutoff lists in advance. Exhaust the closed endpoints and pay the grid, response and annular errors within the allocated collars \(\fitinline{\epsilon_a,\epsilon_b}\), by taking the early meshes smaller than the finitely many strict reserves. A row uses its fixed coarse class to choose a list, before its witness responses are evaluated. The endpoint and short-inverse taper guards, including the stricter check that the physical reciprocal's retained lower length exceeds every bank cutoff, hold for each ticket. No physical reciprocal is evaluated at the auxiliary inverse witness.

For each row fix the actual attaining zero of its original finite family and its presentation. At every cutoff B.5.1 supplies the simultaneous genuine inverse and ordinary at that zero. Freeze the descriptors on the fixed row subclass. The maximum of the genuine available count deficits is one upper bound on that common subclass: select the strongest of the available cardinality bounds (or sum over the fixed finite choice of cases). This includes the preceding amplified inverse and the current inverse and ordinary at the relevant cutoffs. In the sector \(\fitinline{3G/4\ge q_0}\) the original marked numerator fourth pays, since it takes three quarters of the count deficit \(\fitinline{G}\). If all inverse witnesses are long, the last cutoff enforces that sector. Otherwise use the first short inverse and the three original L12 deficit faces on the remaining rows. The argument uses no product of separately obtained counts.

For the remaining rows use the actual ordinary witness at length \(\fitinline{z}\) in K6. Its saturated lower bound supplies the rowwise comparison with the product of the true numerator and that ordinary; the reciprocal and its height are untouched. Use P1 at \(\fitinline{n_{\rm ref}=\max(n,1-y_t/D)}\) and \(\fitinline{z}\), with all componentwise shortenings and separate ray, profile and height controls. The checked length and ratio domains lie within P1. Only after applying the positive source divide by the true complete original whole-slot amplitude. The extra width over the tangent mean is bounded below by \[\fitmath{
               \beta_W(z)=g_{1,-}z-g_{0,+}/2+F(z)>0.           
}\tag{P13}\] The original J9a/K8 coefficient sum pays the possible dilation \(\fitinline{n_{\rm ref}-n_0\le v_{\rm dil}+\epsilon_a}\), where \(\fitinline{v_{\rm dil}}\) is its true packet dilation. It uses the full error-tuple mass and reciprocal weights; neither a single long error product nor a larger comparison source can replace that sum.

\hypertarget{p.6.-every-mixed-piece-and-its-complement}{%
\subsection{P.6. Every mixed piece and its complement}\label{p.6.-every-mixed-piece-and-its-complement}}

Write \(\fitinline{M_0=m_++\epsilon_a}\), \(\fitinline{F_j=-k_+(\gamma_j+3\epsilon_a)}\), and \(\fitinline{H_j=(v_-/b_d)[\gamma_{j-1}-3\epsilon_a]_+}\), with \(\fitinline{\gamma_{-1}=0}\). For \(\fitinline{\Delta=z-M_0}\), the remaining directed comparison in the actual conductor units is L16 with the new debit: \[\fitmath{
 \begin{aligned}
 f_p\beta_W(z)&+\tfrac12\max\{F_j,\ k_{\rm sign}\Delta+\max(H_j,2k_-\Delta_+)\}\\
             &-e_p\epsilon_a/2-q_0/a_d,\qquad
 k_{\rm sign}=\begin{cases}k_+,&\Delta<0,\\ k_-,&\Delta\ge0.\end{cases}
 \end{aligned}                                                  
}\tag{P14}\] On the tested length interval the width decreases, with \(\fitinline{f_p\max(-\beta_W')<k_-/2}\), checked on each side of the seam. Thus P14 falls on the constant face and rises after the first increasing face takes over. The least value is at \(\fitinline{z_j=M_0+(F_j-H_j)/k_+}\), within the tested support by the checked length inequalities. Here \(\fitinline{\Delta\le0}\), so the positive ordinary face is zero; on the larger lengths it can only help.

If the adverse derivative bound times the maximum speed of this crossover is less than \(\fitinline{k_+/2}\), the price at the crossover decreases along the bank and the last offset suffices. Otherwise minimize over the closed continuous interval of offsets: the crossover speeds are \(\fitinline{-1}\) and \(\fitinline{-1-v_-/(b_dk_+)}\), with switch at \(\fitinline{\gamma/16384+3\epsilon_a}\). Split also where \(\fitinline{z_j}\) crosses \(\fitinline{a}\). On each piece the objective is affine plus a convex quadratic or cubic. Check both endpoints and all interior stationary points with rational bisection, subtracting ten times the root enclosure's length because the absolute derivative is less than ten. This gives a directed bound for every actual cutoff and every intermediate length, including the seam and bank endpoints.

The four original closed conductor tilings contain 57 rectangles. Each edge is covered by its original saved route, the fixed detector, the actual current native count, the full admitted paired price, or a remaining mixed ticket. O's fixed detector keeps both branches of its maximum, along with its full native and unmarked continuations; the actual native routes keep their own \(\fitinline{J,N}\), caps, and directed denominator bounds. The original order of roots and children is preserved. The return checks 2044 central parents and 557 full four-way refinements, leaving 3669 final mixed pieces and 46 additional exterior pieces paid by the old or current cleared paired price. Each child is independently checked against its raw ticket. The identity \(\fitinline{3669+46=2044+3(557)}\) exhausts the partition; no child is unpaid. The minimum normalized mixed remainder is at least \(\fitinline{104029/200000000000000=5.20145\cdot10^{-10}}\), and its directed physical \(\fitinline{Z}\)-exponent remainder is at least \(\fitinline{90523/500000000000000=1.81046\cdot10^{-10}}\), with the inward fees already charged. The true retained reciprocal length clears all cutoffs by more than \(\fitinline{.0867}\). The source length interval is contained in \(\fitinline{[.37144,.37880]}\), and the ratio interval in \(\fitinline{[1.11939,1.14162]}\), inside P1.

Each exterior contribution retains its own positive saving. The \(\fitinline{.085}\) cap fits the actual affine capacity and the declared \(\fitinline{.11}\) native source guard; the central coefficient damping remains \(\fitinline{d(1-p)/2-.085>0}\). Keep the actual independent reciprocal contour and its \(\fitinline{s+w}\) guard, both amplitude and reciprocal-length tails, whole-slot costs, unmarked tails, short, bounded and null numerators, and the numerator ordinary escape. The small and averaged floor rows, principal residues and prime compensators are on their own contours and keep the bounds of P.4. In the nonprincipal pieces the absolute contour motion and full \(\fitinline{w}\)-inversion precede subdivision. With the true signs, weights, and scale of the packets, the bound remains \[\fitmath{
 \sum_{\mathbf e}|A_i(\mathbf e)|N_{\mathbf e}^{-1/2}
        \max(1,L_iN_{\mathbf e})^{d/2}
 \ll U^\epsilon\prod_{p\ \mathrm{main}}Y_p^{-4/25}.             
}\tag{P15}\] Apply it before summing the separation and dyadic classes. The original fixed-ray lower response and both signed fixed rays precede complete amplitude division. The original six continuations, count admissibility and conductor collar are paid by the verified ticket gates and P.4. In particular a gain checked on the central mixed price does not pay a discarded tail or any principal term.

\hypertarget{p.7.-heights-and-the-strict-zero-free-region}{%
\subsection{P.7. Heights and the strict zero-free region}\label{p.7.-heights-and-the-strict-zero-free-region}}

Fix the common finite source dependencies, strict decrements, loss budgets, cutoff and response lists, and early meshes. After the actual fixed slot and ray dictionary, radial cap, characters and internal ideals, take the finite prescribed seminorms, normalized first log tests, Sobolev orders and a common polynomial degree. Use the H24a compact log-Fourier envelopes at the actual independent dilations and unit-interval Sobolev on the actual independent height intervals. The reciprocal remains on the original B.8 buffered line; its two omitted segments use the original global \(\fitinline{\beta+20e_0}\) line. The profiles, complete main/error terms, prime corrections and signed fixed residues are kept at their actual scales and their own heights. Positive scalar and source margins allow the early discretization errors and remaining freely small exponents to be chosen within the stated reserves. In particular charge the separate \(\fitinline{15s_0/16+2F_{\rm pad}}\), and take the remaining outer exponent losses and the common reserve \(\fitinline{\sigma>0}\) together below \(\fitinline{10^{-11}}\) and within the original transport constraints. For the inner source choose the meshes and decrements within the strict P.2 margins, independently of the outward loss budget. No additional \(\fitinline{10^{-9}}\) fee is charged against the final physical remainder; that inward fee is already in P11.

For each fixed presented character \(\fitinline{\eta}\), after the finite common degrees choose one admissible \(\fitinline{\tau_\eta>0}\). For all positive truncation orders \(\fitinline{n}\) there is a constant \(\fitinline{C_n}\) such that, eventually in \(\fitinline{Z}\), the literal cube divided by the original nonzero I52 normalizer is bounded by \(\fitinline{C_n Z^{\beta-k_{\rm phys}-\sigma}}\) for all original \texttt{A\allowbreak{}c\allowbreak{}t\allowbreak{}u\allowbreak{}a\allowbreak{}l\allowbreak{}B\allowbreak{}u\allowbreak{}f\allowbreak{}f\allowbreak{}e\allowbreak{}r\allowbreak{}e\allowbreak{}d\allowbreak{}S\allowbreak{}e\allowbreak{}l\allowbreak{}e\allowbreak{}c\allowbreak{}t\allowbreak{}i\allowbreak{}o\allowbreak{}n} indices and grids. The reserve \(\fitinline{\sigma}\) is common and \(\fitinline{\tau_\eta}\) is the same for every order, index and grid of this character. This is I51. Apply the one-character \texttt{n\allowbreak{}o\allowbreak{}r\allowbreak{}m\allowbreak{}a\allowbreak{}l\allowbreak{}i\allowbreak{}z\allowbreak{}e\allowbreak{}d\allowbreak{}P\allowbreak{}r\allowbreak{}o\allowbreak{}b\allowbreak{}e\allowbreak{}\_\allowbreak{}h\allowbreak{}i\allowbreak{}g\allowbreak{}h\allowbreak{}\_\allowbreak{}o\allowbreak{}f\allowbreak{}\_\allowbreak{}f\allowbreak{}i\allowbreak{}n\allowbreak{}i\allowbreak{}t\allowbreak{}e\allowbreak{}\_\allowbreak{}n\allowbreak{}o\allowbreak{}n\allowbreak{}f\allowbreak{}l\allowbreak{}o\allowbreak{}o\allowbreak{}r\allowbreak{}\_\allowbreak{}n\allowbreak{}u\allowbreak{}m\allowbreak{}e\allowbreak{}r\allowbreak{}a\allowbreak{}t\allowbreak{}o\allowbreak{}r\allowbreak{}\_\allowbreak{}a\allowbreak{}t} to reproduce the generic high-of-cube transport with P.4's geometry and support. It gives the original fixed probe and signed Gram exponent \(\fitinline{t_P-k_{\rm phys}}\), using the actual moment parameter \(\fitinline{2\beta-1}\).

If \(\fitinline{\beta>t_P}\), choose \(\fitinline{\omega< (\beta-t_P)/4}\) and \(\fitinline{0<m<\min(\beta-t_P-\omega,\sigma,\beta-.87)}\). H47 gives the Mellin transform holomorphic on \(\fitinline{\Re s>\beta-m}\), agreeing on the absolute line with the original Gaussian times \(\fitinline{H_\eta(s)/L_\eta(s)}\); the numerator still has \(\fitinline{|H_\eta-1|\le1/2}\) for \(\fitinline{\Re s\ge.87}\). The supremum supplies an actual zero above \(\fitinline{\beta-m/2}\) without needing to be attained. The numerator does not vanish there, contradicting holomorphy. The original norm pullback and fixed Euler restoration give P0 at every finite height for each fixed positive ordinary presentation, including imprimitive ones, with its actual principal pole excepted.

\hypertarget{p.8.-contact-test-and-exact-sources}{%
\subsection{P.8. Contact test and exact sources}\label{p.8.-contact-test-and-exact-sources}}

The following limits only the sufficient tangent-bank comparison. At \(\fitinline{p=1/2}\), with \(\fitinline{d_*=(t-1/2-\ell/2)/(1/2+\ell)}\), the ideal zero-cost affine price in P10 has \(\fitinline{E(h)-t=A_0+hA_1}\) with \(\fitinline{A_0=0}\). Thus \(\fitinline{B/h=A_1}\), \(\fitinline{\gamma=4A_1/(3v)}\), and \(\fitinline{z=t_0/3-\gamma(1+v/(d_*k))}\) are independent of \(\fitinline{h}\) on the scanned corridor in the ideal comparison. Its continuous zero-fee price is \(\fitinline{\tfrac12\beta_W(z)-k\gamma/2-A_1/d_*}\). The isolated exact bracket gives a positive sign at \(\fitinline{.126667206}\) and a negative sign at \(\fitinline{.126667207}\). The former does not pay the finite charged return; at \(\fitinline{.12668}\) this one comparison needs at least another \(\fitinline{.0003838822}\) of source width. Neither sign asserts an obstruction to a different proof or an endpoint beyond P0.

The full predecessor through O has 996201 bytes, SHA-256 \texttt{7\allowbreak{}f\allowbreak{}a\allowbreak{}8\allowbreak{}d\allowbreak{}0\allowbreak{}f\allowbreak{}a\allowbreak{}5\allowbreak{}4\allowbreak{}f\allowbreak{}a\allowbreak{}3\allowbreak{}6\allowbreak{}1\allowbreak{}7\allowbreak{}a\allowbreak{}9\allowbreak{}8\allowbreak{}a\allowbreak{}3\allowbreak{}6\allowbreak{}1\allowbreak{}5\allowbreak{}e\allowbreak{}0\allowbreak{}b\allowbreak{}6\allowbreak{}1\allowbreak{}a\allowbreak{}c\allowbreak{}9\allowbreak{}5\allowbreak{}8\allowbreak{}e\allowbreak{}9\allowbreak{}5\allowbreak{}d\allowbreak{}b\allowbreak{}4\allowbreak{}d\allowbreak{}0\allowbreak{}c\allowbreak{}3\allowbreak{}a\allowbreak{}d\allowbreak{}5\allowbreak{}8\allowbreak{}3\allowbreak{}c\allowbreak{}6\allowbreak{}7\allowbreak{}9\allowbreak{}4\allowbreak{}b\allowbreak{}6\allowbreak{}5\allowbreak{}e\allowbreak{}5\allowbreak{}7\allowbreak{}e\allowbreak{}f\allowbreak{}f\allowbreak{}4}. Its scientific handoff, complete N predecessor, frozen O source interfaces and their exact evidence remain inherited. The Cycle 088 source map \texttt{r\allowbreak{}e\allowbreak{}s\allowbreak{}e\allowbreak{}a\allowbreak{}r\allowbreak{}c\allowbreak{}h\allowbreak{}/\allowbreak{}l\allowbreak{}o\allowbreak{}c\allowbreak{}a\allowbreak{}l\allowbreak{}\_\allowbreak{}c\allowbreak{}y\allowbreak{}c\allowbreak{}l\allowbreak{}e\allowbreak{}0\allowbreak{}8\allowbreak{}8\allowbreak{}/\allowbreak{}n\allowbreak{}e\allowbreak{}a\allowbreak{}r\allowbreak{}\_\allowbreak{}r\allowbreak{}e\allowbreak{}c\allowbreak{}o\allowbreak{}r\allowbreak{}d\allowbreak{}/\allowbreak{}S\allowbreak{}O\allowbreak{}U\allowbreak{}R\allowbreak{}C\allowbreak{}E\allowbreak{}\_\allowbreak{}B\allowbreak{}I\allowbreak{}N\allowbreak{}D\allowbreak{}I\allowbreak{}N\allowbreak{}G\allowbreak{}S\allowbreak{}.\allowbreak{}j\allowbreak{}s\allowbreak{}o\allowbreak{}n} has SHA-256 \texttt{f\allowbreak{}5\allowbreak{}8\allowbreak{}d\allowbreak{}5\allowbreak{}4\allowbreak{}f\allowbreak{}8\allowbreak{}0\allowbreak{}9\allowbreak{}3\allowbreak{}6\allowbreak{}5\allowbreak{}f\allowbreak{}f\allowbreak{}a\allowbreak{}5\allowbreak{}f\allowbreak{}9\allowbreak{}c\allowbreak{}e\allowbreak{}8\allowbreak{}2\allowbreak{}e\allowbreak{}a\allowbreak{}0\allowbreak{}3\allowbreak{}0\allowbreak{}a\allowbreak{}c\allowbreak{}d\allowbreak{}3\allowbreak{}a\allowbreak{}f\allowbreak{}8\allowbreak{}2\allowbreak{}1\allowbreak{}3\allowbreak{}5\allowbreak{}e\allowbreak{}2\allowbreak{}0\allowbreak{}8\allowbreak{}2\allowbreak{}4\allowbreak{}e\allowbreak{}b\allowbreak{}7\allowbreak{}8\allowbreak{}d\allowbreak{}2\allowbreak{}5\allowbreak{}9\allowbreak{}1\allowbreak{}8\allowbreak{}6\allowbreak{}b\allowbreak{}a\allowbreak{}f\allowbreak{}7\allowbreak{}a\allowbreak{}2\allowbreak{}7\allowbreak{}2}. \texttt{S\allowbreak{}Y\allowbreak{}M\allowbreak{}M\allowbreak{}E\allowbreak{}T\allowbreak{}R\allowbreak{}I\allowbreak{}C\allowbreak{}\_\allowbreak{}S\allowbreak{}O\allowbreak{}U\allowbreak{}R\allowbreak{}C\allowbreak{}E\allowbreak{}.\allowbreak{}m\allowbreak{}d} (\texttt{b\allowbreak{}a\allowbreak{}d\allowbreak{}e\allowbreak{}7\allowbreak{}d\allowbreak{}9\allowbreak{}d\allowbreak{}4\allowbreak{}7\allowbreak{}6\allowbreak{}3\allowbreak{}e\allowbreak{}3\allowbreak{}a\allowbreak{}3\allowbreak{}0\allowbreak{}7\allowbreak{}e\allowbreak{}2\allowbreak{}c\allowbreak{}9\allowbreak{}f\allowbreak{}8\allowbreak{}b\allowbreak{}5\allowbreak{}1\allowbreak{}2\allowbreak{}f\allowbreak{}9\allowbreak{}1\allowbreak{}4\allowbreak{}5\allowbreak{}5\allowbreak{}9\allowbreak{}a\allowbreak{}7\allowbreak{}7\allowbreak{}b\allowbreak{}6\allowbreak{}b\allowbreak{}f\allowbreak{}1\allowbreak{}1\allowbreak{}f\allowbreak{}b\allowbreak{}8\allowbreak{}d\allowbreak{}8\allowbreak{}e\allowbreak{}a\allowbreak{}5\allowbreak{}1\allowbreak{}c\allowbreak{}4\allowbreak{}3\allowbreak{}f\allowbreak{}5\allowbreak{}b\allowbreak{}b\allowbreak{}1\allowbreak{}2\allowbreak{}2\allowbreak{}4}) gives the source derivation; \texttt{S\allowbreak{}Y\allowbreak{}M\allowbreak{}M\allowbreak{}E\allowbreak{}T\allowbreak{}R\allowbreak{}I\allowbreak{}C\allowbreak{}\_\allowbreak{}S\allowbreak{}O\allowbreak{}U\allowbreak{}R\allowbreak{}C\allowbreak{}E\allowbreak{}.\allowbreak{}j\allowbreak{}s\allowbreak{}o\allowbreak{}n} (\texttt{f\allowbreak{}3\allowbreak{}c\allowbreak{}9\allowbreak{}9\allowbreak{}a\allowbreak{}9\allowbreak{}5\allowbreak{}7\allowbreak{}6\allowbreak{}5\allowbreak{}5\allowbreak{}f\allowbreak{}7\allowbreak{}3\allowbreak{}c\allowbreak{}7\allowbreak{}6\allowbreak{}d\allowbreak{}d\allowbreak{}5\allowbreak{}0\allowbreak{}4\allowbreak{}7\allowbreak{}2\allowbreak{}8\allowbreak{}8\allowbreak{}7\allowbreak{}c\allowbreak{}8\allowbreak{}2\allowbreak{}1\allowbreak{}f\allowbreak{}5\allowbreak{}2\allowbreak{}9\allowbreak{}c\allowbreak{}f\allowbreak{}2\allowbreak{}b\allowbreak{}e\allowbreak{}5\allowbreak{}8\allowbreak{}1\allowbreak{}2\allowbreak{}5\allowbreak{}e\allowbreak{}9\allowbreak{}7\allowbreak{}6\allowbreak{}a\allowbreak{}d\allowbreak{}a\allowbreak{}6\allowbreak{}5\allowbreak{}d\allowbreak{}3\allowbreak{}9\allowbreak{}0\allowbreak{}6\allowbreak{}2\allowbreak{}0\allowbreak{}4\allowbreak{}7}) gives its exact scalar checks. The new physical derivation \texttt{R\allowbreak{}E\allowbreak{}T\allowbreak{}U\allowbreak{}R\allowbreak{}N\allowbreak{}\_\allowbreak{}P\allowbreak{}R\allowbreak{}O\allowbreak{}O\allowbreak{}F\allowbreak{}.\allowbreak{}m\allowbreak{}d} is \texttt{d\allowbreak{}e\allowbreak{}a\allowbreak{}9\allowbreak{}7\allowbreak{}6\allowbreak{}2\allowbreak{}a\allowbreak{}f\allowbreak{}4\allowbreak{}1\allowbreak{}5\allowbreak{}b\allowbreak{}3\allowbreak{}2\allowbreak{}f\allowbreak{}a\allowbreak{}f\allowbreak{}3\allowbreak{}4\allowbreak{}b\allowbreak{}8\allowbreak{}c\allowbreak{}5\allowbreak{}3\allowbreak{}1\allowbreak{}f\allowbreak{}9\allowbreak{}9\allowbreak{}b\allowbreak{}c\allowbreak{}e\allowbreak{}6\allowbreak{}a\allowbreak{}4\allowbreak{}3\allowbreak{}7\allowbreak{}e\allowbreak{}b\allowbreak{}3\allowbreak{}4\allowbreak{}c\allowbreak{}5\allowbreak{}7\allowbreak{}d\allowbreak{}3\allowbreak{}7\allowbreak{}2\allowbreak{}8\allowbreak{}6\allowbreak{}8\allowbreak{}9\allowbreak{}b\allowbreak{}9\allowbreak{}f\allowbreak{}8\allowbreak{}9\allowbreak{}a\allowbreak{}6\allowbreak{}f\allowbreak{}f\allowbreak{}b\allowbreak{}a\allowbreak{}9}. Its receipt \texttt{R\allowbreak{}E\allowbreak{}T\allowbreak{}U\allowbreak{}R\allowbreak{}N\allowbreak{}\_\allowbreak{}1\allowbreak{}2\allowbreak{}6\allowbreak{}6\allowbreak{}6\allowbreak{}7\allowbreak{}2\allowbreak{}0\allowbreak{}5\allowbreak{}.\allowbreak{}j\allowbreak{}s\allowbreak{}o\allowbreak{}n} is \texttt{f\allowbreak{}d\allowbreak{}5\allowbreak{}7\allowbreak{}3\allowbreak{}c\allowbreak{}e\allowbreak{}c\allowbreak{}6\allowbreak{}5\allowbreak{}1\allowbreak{}e\allowbreak{}c\allowbreak{}0\allowbreak{}b\allowbreak{}3\allowbreak{}d\allowbreak{}f\allowbreak{}9\allowbreak{}5\allowbreak{}8\allowbreak{}5\allowbreak{}d\allowbreak{}9\allowbreak{}e\allowbreak{}9\allowbreak{}4\allowbreak{}0\allowbreak{}c\allowbreak{}3\allowbreak{}2\allowbreak{}1\allowbreak{}c\allowbreak{}b\allowbreak{}4\allowbreak{}0\allowbreak{}b\allowbreak{}c\allowbreak{}7\allowbreak{}d\allowbreak{}1\allowbreak{}c\allowbreak{}8\allowbreak{}1\allowbreak{}6\allowbreak{}7\allowbreak{}4\allowbreak{}9\allowbreak{}c\allowbreak{}0\allowbreak{}a\allowbreak{}f\allowbreak{}d\allowbreak{}4\allowbreak{}c\allowbreak{}0\allowbreak{}3\allowbreak{}5\allowbreak{}5\allowbreak{}e\allowbreak{}b\allowbreak{}e\allowbreak{}2\allowbreak{}e} with coverage digest \texttt{7\allowbreak{}7\allowbreak{}0\allowbreak{}f\allowbreak{}2\allowbreak{}c\allowbreak{}f\allowbreak{}9\allowbreak{}d\allowbreak{}0\allowbreak{}5\allowbreak{}2\allowbreak{}3\allowbreak{}d\allowbreak{}2\allowbreak{}f\allowbreak{}5\allowbreak{}6\allowbreak{}6\allowbreak{}b\allowbreak{}5\allowbreak{}f\allowbreak{}9\allowbreak{}8\allowbreak{}8\allowbreak{}1\allowbreak{}4\allowbreak{}a\allowbreak{}d\allowbreak{}1\allowbreak{}2\allowbreak{}c\allowbreak{}5\allowbreak{}1\allowbreak{}6\allowbreak{}3\allowbreak{}5\allowbreak{}b\allowbreak{}a\allowbreak{}9\allowbreak{}7\allowbreak{}6\allowbreak{}6\allowbreak{}8\allowbreak{}a\allowbreak{}3\allowbreak{}0\allowbreak{}9\allowbreak{}c\allowbreak{}b\allowbreak{}b\allowbreak{}c\allowbreak{}0\allowbreak{}0\allowbreak{}7\allowbreak{}5\allowbreak{}f\allowbreak{}1\allowbreak{}d\allowbreak{}f\allowbreak{}d\allowbreak{}a\allowbreak{}8\allowbreak{}5}. \texttt{C\allowbreak{}O\allowbreak{}N\allowbreak{}T\allowbreak{}A\allowbreak{}C\allowbreak{}T\allowbreak{}\_\allowbreak{}B\allowbreak{}R\allowbreak{}A\allowbreak{}C\allowbreak{}K\allowbreak{}E\allowbreak{}T\allowbreak{}.\allowbreak{}j\allowbreak{}s\allowbreak{}o\allowbreak{}n} (\texttt{d\allowbreak{}c\allowbreak{}7\allowbreak{}0\allowbreak{}f\allowbreak{}1\allowbreak{}1\allowbreak{}d\allowbreak{}5\allowbreak{}c\allowbreak{}a\allowbreak{}0\allowbreak{}2\allowbreak{}2\allowbreak{}2\allowbreak{}3\allowbreak{}c\allowbreak{}4\allowbreak{}5\allowbreak{}0\allowbreak{}b\allowbreak{}c\allowbreak{}c\allowbreak{}2\allowbreak{}5\allowbreak{}b\allowbreak{}0\allowbreak{}c\allowbreak{}b\allowbreak{}c\allowbreak{}8\allowbreak{}4\allowbreak{}c\allowbreak{}b\allowbreak{}8\allowbreak{}f\allowbreak{}4\allowbreak{}7\allowbreak{}5\allowbreak{}7\allowbreak{}1\allowbreak{}1\allowbreak{}4\allowbreak{}7\allowbreak{}e\allowbreak{}e\allowbreak{}3\allowbreak{}4\allowbreak{}9\allowbreak{}4\allowbreak{}a\allowbreak{}1\allowbreak{}7\allowbreak{}b\allowbreak{}b\allowbreak{}0\allowbreak{}a\allowbreak{}b\allowbreak{}d\allowbreak{}7\allowbreak{}e\allowbreak{}2\allowbreak{}4\allowbreak{}5}) is the separate ideal comparison. The source map and \texttt{H\allowbreak{}A\allowbreak{}N\allowbreak{}D\allowbreak{}O\allowbreak{}F\allowbreak{}F\allowbreak{}.\allowbreak{}m\allowbreak{}d} bind the checker identities, original input paths and replay parameters.

The finite checks establish the stated rational inequalities and coverage. The written source, fixed and regular joins and height implications above use the complete original OpenAI/Collet proof and its retained Sil, Bebel, and other original contributions, notices and attributions. The catalogue references remain discovery and provenance. This appendix supplies no new end-to-end Lean verification or public release.

\hypertarget{appendix-q}{}

\hypertarget{appendix-q.-the-refined-full-return-at-c.126667206}{%
\section{\texorpdfstring{Appendix Q. The refined full return at \(\fitinline{c=.126667206}\)}{Appendix Q. The refined full return at c=.126667206}}\label{appendix-q.-the-refined-full-return-at-c.126667206}}

Retain the complete A--P proof, including the adopted H6/SN9 input and P's ordered two-ordinary source with its \(\fitinline{.115}\) short marked reference. Its source, raw signs and coefficient normalizers are unchanged. Set \[\fitmath{
 c_Q=\frac{63333603}{500000000}=.126667206,\qquad
 t_Q=\frac{436666397}{500000000}=.873332794,\qquad
 h=h_*+.0022595743.                                            
}\tag{Q0}\] For every fixed positive presented modulus and every ordinary character, including an imprimitive presentation, the zero-free region is \(\fitinline{\Re s>t_Q}\) at every finite height, except for the actual principal pole. The gain over P is \(\fitinline{10^{-9}}\) in \(\fitinline{c}\). The statement is strict and has no growing-modulus uniformity.

\hypertarget{q.1.-directed-top-clearing-and-the-normalized-debit}{%
\subsection{Q.1. Directed top clearing and the normalized debit}\label{q.1.-directed-top-clearing-and-the-normalized-debit}}

Use \(\fitinline{\ell=11/3-4t_Q,\ x=1+\ell-h,\ y_t=h-2\ell}\). The directed \(\fitinline{h}\) interval is \(\fitinline{[.8175840773,.8175840793]}\). Write \(\fitinline{s_0=10^{-260}}\) and \(\fitinline{F_{\rm pad}=10^{-258}}\), and keep the actual physical cap \(\fitinline{.085}\) declared below \(\fitinline{.11}\). The original geometry P8 is strict on these endpoints, as are the small, floor, Gram and two separate principal comparisons. For the parameterized PR141 transport the actual row pad is \(\fitinline{F_{\rm pad}/2}\), extension is \(\fitinline{F_{\rm pad}}\), and the high-source cap is \(\fitinline{0<\zeta\le F_{\rm pad}}\). Choose \(\fitinline{0<D_{\rm src}.t\le F_{\rm pad}/4}\); then \(\fitinline{p_{\rm row}+2D_{\rm src}.t\le F_{\rm pad}}\) gives the true support through \(\fitinline{h+F_{\rm pad}}\). Charge \(\fitinline{15s_0/16+2F_{\rm pad}}\) separately. This uses neither the obsolete \(\fitinline{h<.816}\) wrapper nor the loose historical \(\fitinline{10^{-10}}\) padding guard.

On an admitted paired cell use P9: \(\fitinline{v=5/6-d}\), \(\fitinline{k=1-p(1+2g_1)}\), \(\fitinline{U_0=1-p(1+6g_0-4g_1)}\), \(\fitinline{J=3v+2dk>0}\), \(\fitinline{N=J/6+v(J+dU_0)}\), and \(\fitinline{t_0=1+dU_0/J}\). Thus the count bound is \(\fitinline{R=N/J=1/6+vt_0\ge1-d}\). At the upper conductor edge the repriced P10 is \[\fitmath{
 \begin{aligned}
 C(D)&=\min\bigl(.085D,[2y_t-(1+s_0)D]/(2+10s_0)\bigr),\\
 E(D)&=E_b(D)+3DR/4,\qquad
 E_b(D)=\tfrac12+\tfrac{13h}{75}-\tfrac{y_t}{2}-\tfrac\ell2-.09D
               +\tfrac d2+dp(\ell-C(D)/2).
 \end{aligned}                                                  
}\tag{Q1}\] At \(\fitinline{D=h}\) clear \(\fitinline{(t_Q-E_b(h))J-3hN/4}\) at both directed \(\fitinline{h}\) endpoints, using the same rational contact enclosures as P.4. The paired count decreases with \(\fitinline{g_0}\) and increases with \(\fitinline{g_1}\), so use \(\fitinline{g_{0,-},g_{1,+}}\) here. The price at \(\fitinline{D=h}\) is convex in \(\fitinline{h}\) because the minimum of the affine capacities enters with a nonpositive coefficient. On each closed \(\fitinline{d,p}\) box a tensor-Bernstein lower bound \(\fitinline{M_-}\) for the cleared numerator is divided by \(\fitinline{J_+}\) if nonnegative, by \(\fitinline{J_-}\) otherwise, with \(\fitinline{0<J_-\le J\le J_+}\). Take the larger lower margin of this clearing and the original directed E49 bound on the same admitted count. Let \(\fitinline{B}\) be the resulting upper excess.

Use \(\fitinline{\epsilon_a=10^{-13}}\) for the witness collar and \(\fitinline{\epsilon_b=10^{-12}}\) for the allocated inward debit. On a box \(\fitinline{[a_d,b_d]\times[f_p,e_p]}\), put \[\fitmath{
 S=\tfrac34(1-b_d)-.09-b_d(.085)/4,\qquad
 q_0=(\max(0,B)+\epsilon_b)/h_-.
}\] The slope of Q1 at fixed \(\fitinline{h}\) is at least \(\fitinline{S}\). The raw ticket gives \(\fitinline{S>0}\), so when \(\fitinline{B<-\epsilon_b}\) the original marked numerator fourth pays the exterior throughout this interval at the same count. On the remaining mixed boxes the certificate gives \(\fitinline{S>q_0}\). Hence, for the retained \(\fitinline{D_-\le D\le h}\), \[\fitmath{
 [E(D)-t_Q+\epsilon_b]_+\le [B+\epsilon_b-S(h-D)]_+\le Dq_0.
}\tag{Q2}\] This is the debit for \(\fitinline{U=Z^D}\). The old raw \(\fitinline{D_-}\) band and all its support, tail and contour conditions are repriced with their \textbf{original} fees and are kept when the witness debit shrinks. The separate padding debit pays the extension above \(\fitinline{h}\).

\hypertarget{q.2.-all-offsets-and-mixed-lengths}{%
\subsection{Q.2. All offsets and mixed lengths}\label{q.2.-all-offsets-and-mixed-lengths}}

Round \(\fitinline{\gamma,\lambda_x,\lambda_y}\) up to \(\fitinline{10^{-13}}\) from the P12 requirements \(\fitinline{4q_0/(3v_-),4q_0/(3a_dk_-),2q_0/(3a_dk_-)}\), where \(\fitinline{v_-=5/6-b_d}\) and \(\fitinline{k_-\le k\le k_+}\). Use the \(\fitinline{65{,}537}\) cutoffs \(\fitinline{t_0-j\gamma/65536}\), including both endpoints of the \(\fitinline{65{,}536}\) subintervals. Fix the finite lists and meshes before seeing any response, with the \(\fitinline{t_0}\) and response rounding and annular losses inside their strict allocations. The real reciprocal's retained lower length is above \textbf{every} cutoff. The narrowed cutoffs still have the genuine short-inverse taper gap and remain within \(\fitinline{(1,3/2)}\).

Use B.5.1 at the same actual attaining zero and presentation throughout the chosen list. The strongest genuine count deficit on a fixed subclass gives one cardinality bound there. The original marked numerator fourth pays whenever three quarters of that deficit covers \(\fitinline{q_0}\), including the case when all inverses are long by the last cutoff. Otherwise the first short inverse gives the preceding long, current inverse, and current ordinary deficit faces of L12. Use K6 and P1 on the actual saturated ordinary length \(\fitinline{z}\) and \(\fitinline{n_{\rm ref}=\max(n,1-y_t/D)}\), retaining the genuine shortenings. The source ratios lie in \(\fitinline{[1.11939,1.14162]}\), and the tested \(\fitinline{z}\)'s lie in \(\fitinline{[.37144,.37880]}\). The gain is \(\fitinline{\beta_W(z)=g_{1,-}z-g_{0,+}/2+F(z)}\) as in P13. Divide the true whole-slot amplitude only after the positive P1 application. The physical reciprocal is unchanged, and the complete P15 coefficient sum still pays the numerator dilation.

For P14 with the new \(\fitinline{\epsilon_a,q_0}\), the price decreases on the constant face and increases after the first rising face because \(\fitinline{f_p\max(-\beta_W')<k_-/2}\) on the tested interval, using both sides of the source seam. At offset \(\fitinline{\xi=j\gamma/65536}\) the crossing is at \[\fitmath{
 z(\xi)=m_++\epsilon_a-\xi-3\epsilon_a
       -\frac{v_-}{b_dk_+}\,[\xi-\gamma/65536-3\epsilon_a]_+ ,
}\] where \(\fitinline{m_+}\) is the upper bound for \(\fitinline{t_0/3}\). There the positive ordinary face vanishes. The remaining price is \(\fitinline{f_p\beta_W(z(\xi))-k_+(\xi+3\epsilon_a)/2-e_p\epsilon_a/2-q_0/a_d}\). When the adverse width derivative times the maximum crossover speed is below \(\fitinline{k_+/2}\), the last offset minimizes it. Otherwise split the full closed offset interval at the speed switch and the source seam and check endpoints and stationary points on each piece. The objective there is affine plus a convex quadratic or cubic. On these boxes \(\fitinline{b_d>.39}\), \(\fitinline{k_+>.4}\), \(\fitinline{v_-<5/6}\) and P1 gives \(\fitinline{|\beta_W'|<1}\) on each side of the seam. The crossover speed is at most seven, so ten bounds the absolute price derivative. Use 42 steps of rational bisection for each interior stationary point and subtract ten times the root enclosure length from the lesser endpoint price of that enclosure. This certifies every offset and every admissible intermediate length.

\hypertarget{q.3.-physical-cover-and-margins}{%
\subsection{Q.3. Physical cover and margins}\label{q.3.-physical-cover-and-margins}}

All four original conductor tilings are covered on their 57 original closed cells. The old saved, fixed detector, native and E49 routes retain their actual caps, both fixed-count branches where present, and their existing contour and coefficient obligations. The four used extensions keep all ten E29 guards at each of their four enclosing contact corners; the recorded least directed guard is \(\fitinline{69989628971/10000000000000>0}\). Every visited mixed child has a fresh original raw ticket. The traversal has 2045 central roots, 12289 visited nodes and 2561 complete four-way refinements; 9681 mixed terminals and 47 added exterior pieces exhaust it, since \(\fitinline{9681+47=2045+3\cdot2561}\). There are 50 continuous stationary-node checks and no unpaid terminal. There were 0 exterior children discarded before a visit.

The following are directed rational lower bounds, with the witness costs already included in the first two:

\begin{longtable}[]{@{}
  >{\raggedright\arraybackslash}p{(\columnwidth - 2\tabcolsep) * \real{0.4286}}
  >{\raggedleft\arraybackslash}p{(\columnwidth - 2\tabcolsep) * \real{0.5714}}@{}}
\toprule\noalign{}
\begin{minipage}[b]{\linewidth}\raggedright
Quantity
\end{minipage} & \begin{minipage}[b]{\linewidth}\raggedleft
Lower bound
\end{minipage} \\
\midrule\noalign{}
\endhead
\bottomrule\noalign{}
\endlastfoot
Normalized mixed price & \(\fitinline{500541/10000\allowbreak{}00000\allowbreak{}00000\allowbreak{}0}\) \\
Mixed physical \(\fitinline{Z}\)-exponent reserve & \(\fitinline{174177/10000\allowbreak{}00000\allowbreak{}00000\allowbreak{}0}\) \\
Current paired exterior surplus & \(\fitinline{19714351/10000\allowbreak{}00000\allowbreak{}00000\allowbreak{}0}\) \\
\(\fitinline{S-q_0}\) & \(\fitinline{63837\allowbreak{}55322\allowbreak{}7731/20000\allowbreak{}00000\allowbreak{}00000}\) \\
Retained reciprocal length minus every cutoff & \(\fitinline{17348\allowbreak{}94922\allowbreak{}9623/20000\allowbreak{}00000\allowbreak{}00000}\) \\
Original paid-route reserves before separate padding & \(\fitinline{1245083/20000\allowbreak{}00000\allowbreak{}00000}\) \\
\end{longtable}

The mixed physical bound is the minimum of the terminal normalized lower value multiplied by its own \(\fitinline{D_-a_d}\), with downward rounding after taking the minimum. The least normalized terminal is in original cell 46 at depth 11. In particular the mixed reserve exceeds the separate pad debit. Even using only the exterior acceptance threshold \(\fitinline{10^{-12}}\), take the remaining outer exponent losses and common reserve \(\fitinline{\sigma>0}\) with sum \(\fitinline{<10^{-14}}\), within the original transport and other analytic allowances. The inner source meshes fit the independent P.2 reserves. There is no second inward fee on the mixed margin. The original raw margins still pay their own shifts, floor and principal contours, and all amplitude, reciprocal-length and dilation tails. Keep both original signed sides and their full error-tuple weights, including P15, until their original points of insertion.

\hypertarget{q.4.-the-original-height-and-mellin-implication}{%
\subsection{Q.4. The original height and Mellin implication}\label{q.4.-the-original-height-and-mellin-implication}}

Keep the literal B.8 reciprocal on the retained buffered line with the separate \(\fitinline{s+w}\) guard and the omitted segments on the global \(\fitinline{\beta+20e_0}\) line. The physical contour frequencies, witness height and reflected ordinary height remain independent. Fix the source parameters, fit the early meshes to the strict reserves, and form the finite cutoff and response banks before the live rows and buffers. After the actual fixed slot and ray dictionary, radial cap, characters and internal ideals, take the finite profile and Fourier controls, actual first log tests, Sobolev orders and common height degree. As in P.7, choose one \(\fitinline{\tau_\eta>0}\) for each fixed presented character \(\fitinline{\eta}\) \textbf{before} all live height intervals and grid indices. For every positive truncation order \(\fitinline{n}\) there is a constant \(\fitinline{C_{\eta,n}}\) such that, eventually in \(\fitinline{Z}\), the literal cube divided by the original nonzero I52 normalizer is bounded by \(\fitinline{C_{\eta,n}Z^{\beta-k_{\rm phys}-\sigma}}\) for all original buffered selections and grids of that character. The same \(\fitinline{\tau_\eta}\) works for every order and grid; separate H24a log-Fourier and unit-interval Sobolev bounds handle the actual independent dilation and height suprema. All moment data still use \(\fitinline{2\beta-1}\). Apply the one-character high-of-cube primitive \(\fitinline{\texttt{n\allowbreak{}o\allowbreak{}r\allowbreak{}m\allowbreak{}a\allowbreak{}l\allowbreak{}i\allowbreak{}z\allowbreak{}e\allowbreak{}d\allowbreak{}P\allowbreak{}r\allowbreak{}o\allowbreak{}b\allowbreak{}e\allowbreak{}\_\allowbreak{}h\allowbreak{}i\allowbreak{}g\allowbreak{}h\allowbreak{}\_\allowbreak{}o\allowbreak{}f\allowbreak{}\_\allowbreak{}f\allowbreak{}i\allowbreak{}n\allowbreak{}i\allowbreak{}t\allowbreak{}e\allowbreak{}\_\allowbreak{}n\allowbreak{}o\allowbreak{}n\allowbreak{}f\allowbreak{}l\allowbreak{}o\allowbreak{}o\allowbreak{}r\allowbreak{}\_\allowbreak{}n\allowbreak{}u\allowbreak{}m\allowbreak{}e\allowbreak{}r\allowbreak{}a\allowbreak{}t\allowbreak{}o\allowbreak{}r\allowbreak{}\_\allowbreak{}a\allowbreak{}t}}\) with Q.1's geometry. It bounds the difference between the original signed normalized probe \(\fitinline{J_\eta}\) and its H47 reciprocal signal \(\fitinline{f_\eta}\): \[\fitmath{
 |J_\eta(Z)-f_\eta(Z)|\ll_\eta Z^{\beta-k_{\rm phys}-\sigma},
 \qquad k_{\rm phys}=5/6-x/3-\ell/6 .
}\] Independently, the original positive marked-energy bound B(198) and the signed Gram bound B(199) estimate \(\fitinline{J_\eta}\) itself. Here \(\fitinline{x+y_t=1-\ell}\), with \(\fitinline{\ell<x<y_t}\) and \(\fitinline{y_t-2x+\ell<0}\) by Q.1. The additional term in B(199) vanishes, since \(\fitinline{1+3\ell-2(x+y_t)=5\ell-1<0}\). The exact exponent identity is \[\fitmath{
 k_{\rm phys}+x/2+(y_t-x)/12
   =2/3+(x+y_t)/4=11/12-\ell/4=t_Q .
}\] Dividing by the original nonzero logarithmic normalizer gives \(\fitinline{|J_\eta(Z)|\ll_{\eta,\omega} Z^{t_Q-k_{\rm phys}+\omega}}\) for every \(\fitinline{\omega>0}\), absorbing the logarithms and choosing the Gram loss smaller than \(\fitinline{\omega}\).

By P the actual zero supremum in the contrary case satisfies \(\fitinline{0<\beta-t_Q\le10^{-9}}\); no uniform positive lower bound on this gap is needed. Choose \(\fitinline{0<\omega<(\beta-t_Q)/4}\) and \(\fitinline{0<m<\min(\beta-t_Q-\omega,\sigma,\beta-.87)}\). The two bounds on the same probe give \(\fitinline{f_\eta(Z)\ll_\eta Z^{\beta-k_{\rm phys}-m}}\) at infinity. Together with the original arbitrary-power decay at zero, this makes the H47 signal's Mellin transform \(\fitinline{\int_0^\infty f_\eta(Z)Z^{k_{\rm phys}-s}\,dZ/Z}\) holomorphic on \(\fitinline{\Re s>\beta-m}\). On the absolute line it equals the original Gaussian times \(\fitinline{H_\eta(s)/L_K^{\mathcal S}(s,\eta)}\). An actual zero above \(\fitinline{\beta-m/2}\) exists by the supremum, and its numerator is nonzero by \(\fitinline{|H_\eta-1|\le1/2}\) on \(\fitinline{\Re s\ge.87}\). This contradicts holomorphy. The original norm pullback and Euler-factor restoration give Q0.

The exact Cycle 089 handoff (\texttt{0\allowbreak{}8\allowbreak{}f\allowbreak{}0\allowbreak{}6\allowbreak{}3\allowbreak{}8\allowbreak{}9\allowbreak{}5\allowbreak{}b\allowbreak{}e\allowbreak{}4\allowbreak{}5\allowbreak{}8\allowbreak{}c\allowbreak{}1\allowbreak{}5\allowbreak{}7\allowbreak{}b\allowbreak{}b\allowbreak{}8\allowbreak{}9\allowbreak{}6\allowbreak{}2\allowbreak{}d\allowbreak{}4\allowbreak{}4\allowbreak{}4\allowbreak{}a\allowbreak{}1\allowbreak{}e\allowbreak{}0\allowbreak{}1\allowbreak{}a\allowbreak{}b\allowbreak{}c\allowbreak{}e\allowbreak{}2\allowbreak{}6\allowbreak{}0\allowbreak{}8\allowbreak{}1\allowbreak{}9\allowbreak{}7\allowbreak{}a\allowbreak{}7\allowbreak{}d\allowbreak{}e\allowbreak{}a\allowbreak{}4\allowbreak{}e\allowbreak{}e\allowbreak{}3\allowbreak{}8\allowbreak{}8\allowbreak{}3\allowbreak{}a\allowbreak{}7\allowbreak{}c\allowbreak{}1\allowbreak{}c\allowbreak{}4\allowbreak{}b\allowbreak{}7} as SHA-256) retains without change the 1,030,405-byte full proof through P (SHA-256 \texttt{6\allowbreak{}7\allowbreak{}4\allowbreak{}4\allowbreak{}d\allowbreak{}f\allowbreak{}4\allowbreak{}e\allowbreak{}a\allowbreak{}9\allowbreak{}6\allowbreak{}0\allowbreak{}a\allowbreak{}d\allowbreak{}0\allowbreak{}d\allowbreak{}4\allowbreak{}4\allowbreak{}9\allowbreak{}3\allowbreak{}6\allowbreak{}e\allowbreak{}e\allowbreak{}4\allowbreak{}d\allowbreak{}f\allowbreak{}b\allowbreak{}5\allowbreak{}5\allowbreak{}8\allowbreak{}4\allowbreak{}b\allowbreak{}c\allowbreak{}0\allowbreak{}7\allowbreak{}0\allowbreak{}6\allowbreak{}c\allowbreak{}a\allowbreak{}2\allowbreak{}f\allowbreak{}d\allowbreak{}0\allowbreak{}7\allowbreak{}7\allowbreak{}8\allowbreak{}5\allowbreak{}5\allowbreak{}d\allowbreak{}4\allowbreak{}1\allowbreak{}9\allowbreak{}7\allowbreak{}e\allowbreak{}f\allowbreak{}e\allowbreak{}6\allowbreak{}2\allowbreak{}5\allowbreak{}d\allowbreak{}2\allowbreak{}9\allowbreak{}8\allowbreak{}d}). The active refined receipt is \texttt{R\allowbreak{}E\allowbreak{}F\allowbreak{}I\allowbreak{}N\allowbreak{}E\allowbreak{}D\allowbreak{}\_\allowbreak{}1\allowbreak{}2\allowbreak{}6\allowbreak{}6\allowbreak{}6\allowbreak{}7\allowbreak{}2\allowbreak{}0\allowbreak{}6\allowbreak{}.\allowbreak{}j\allowbreak{}s\allowbreak{}o\allowbreak{}n} (Cycle 090 controller, 31632 bytes, SHA-256 \texttt{4\allowbreak{}f\allowbreak{}0\allowbreak{}0\allowbreak{}c\allowbreak{}2\allowbreak{}1\allowbreak{}0\allowbreak{}d\allowbreak{}4\allowbreak{}3\allowbreak{}0\allowbreak{}a\allowbreak{}6\allowbreak{}e\allowbreak{}8\allowbreak{}a\allowbreak{}5\allowbreak{}4\allowbreak{}3\allowbreak{}b\allowbreak{}9\allowbreak{}3\allowbreak{}7\allowbreak{}5\allowbreak{}8\allowbreak{}7\allowbreak{}2\allowbreak{}5\allowbreak{}e\allowbreak{}8\allowbreak{}3\allowbreak{}d\allowbreak{}9\allowbreak{}b\allowbreak{}6\allowbreak{}a\allowbreak{}5\allowbreak{}0\allowbreak{}9\allowbreak{}e\allowbreak{}4\allowbreak{}1\allowbreak{}c\allowbreak{}3\allowbreak{}d\allowbreak{}1\allowbreak{}7\allowbreak{}8\allowbreak{}4\allowbreak{}a\allowbreak{}5\allowbreak{}e\allowbreak{}0\allowbreak{}f\allowbreak{}8\allowbreak{}b\allowbreak{}d\allowbreak{}b\allowbreak{}2\allowbreak{}a\allowbreak{}d\allowbreak{}6\allowbreak{}6}), with coverage digest \texttt{c\allowbreak{}c\allowbreak{}3\allowbreak{}3\allowbreak{}0\allowbreak{}0\allowbreak{}b\allowbreak{}e\allowbreak{}4\allowbreak{}7\allowbreak{}e\allowbreak{}d\allowbreak{}d\allowbreak{}8\allowbreak{}3\allowbreak{}1\allowbreak{}b\allowbreak{}a\allowbreak{}7\allowbreak{}b\allowbreak{}5\allowbreak{}b\allowbreak{}e\allowbreak{}8\allowbreak{}9\allowbreak{}9\allowbreak{}7\allowbreak{}1\allowbreak{}0\allowbreak{}a\allowbreak{}8\allowbreak{}a\allowbreak{}7\allowbreak{}4\allowbreak{}7\allowbreak{}8\allowbreak{}d\allowbreak{}0\allowbreak{}c\allowbreak{}b\allowbreak{}5\allowbreak{}8\allowbreak{}d\allowbreak{}7\allowbreak{}3\allowbreak{}5\allowbreak{}8\allowbreak{}4\allowbreak{}9\allowbreak{}6\allowbreak{}5\allowbreak{}3\allowbreak{}c\allowbreak{}c\allowbreak{}3\allowbreak{}9\allowbreak{}e\allowbreak{}0\allowbreak{}7\allowbreak{}0\allowbreak{}9\allowbreak{}7\allowbreak{}c\allowbreak{}c}. Its wrapper \texttt{R\allowbreak{}E\allowbreak{}F\allowbreak{}I\allowbreak{}N\allowbreak{}E\allowbreak{}\_\allowbreak{}C\allowbreak{}O\allowbreak{}N\allowbreak{}T\allowbreak{}A\allowbreak{}C\allowbreak{}T\allowbreak{}.\allowbreak{}p\allowbreak{}y} is \texttt{f\allowbreak{}f\allowbreak{}f\allowbreak{}5\allowbreak{}6\allowbreak{}f\allowbreak{}e\allowbreak{}5\allowbreak{}0\allowbreak{}5\allowbreak{}c\allowbreak{}a\allowbreak{}c\allowbreak{}9\allowbreak{}3\allowbreak{}b\allowbreak{}1\allowbreak{}d\allowbreak{}a\allowbreak{}3\allowbreak{}9\allowbreak{}7\allowbreak{}1\allowbreak{}f\allowbreak{}a\allowbreak{}6\allowbreak{}6\allowbreak{}9\allowbreak{}d\allowbreak{}9\allowbreak{}7\allowbreak{}d\allowbreak{}6\allowbreak{}a\allowbreak{}a\allowbreak{}9\allowbreak{}a\allowbreak{}3\allowbreak{}c\allowbreak{}3\allowbreak{}4\allowbreak{}0\allowbreak{}d\allowbreak{}b\allowbreak{}4\allowbreak{}6\allowbreak{}b\allowbreak{}c\allowbreak{}a\allowbreak{}7\allowbreak{}6\allowbreak{}9\allowbreak{}4\allowbreak{}7\allowbreak{}f\allowbreak{}1\allowbreak{}d\allowbreak{}8\allowbreak{}7\allowbreak{}f\allowbreak{}d\allowbreak{}f\allowbreak{}2\allowbreak{}5}, its upstream \texttt{C\allowbreak{}H\allowbreak{}E\allowbreak{}C\allowbreak{}K\allowbreak{}\_\allowbreak{}F\allowbreak{}U\allowbreak{}L\allowbreak{}L\allowbreak{}.\allowbreak{}p\allowbreak{}y} \texttt{f\allowbreak{}7\allowbreak{}b\allowbreak{}f\allowbreak{}7\allowbreak{}1\allowbreak{}b\allowbreak{}a\allowbreak{}5\allowbreak{}f\allowbreak{}a\allowbreak{}b\allowbreak{}7\allowbreak{}d\allowbreak{}2\allowbreak{}e\allowbreak{}4\allowbreak{}4\allowbreak{}a\allowbreak{}2\allowbreak{}0\allowbreak{}d\allowbreak{}a\allowbreak{}b\allowbreak{}7\allowbreak{}c\allowbreak{}f\allowbreak{}1\allowbreak{}2\allowbreak{}8\allowbreak{}e\allowbreak{}1\allowbreak{}9\allowbreak{}a\allowbreak{}4\allowbreak{}8\allowbreak{}4\allowbreak{}2\allowbreak{}d\allowbreak{}4\allowbreak{}4\allowbreak{}e\allowbreak{}d\allowbreak{}b\allowbreak{}8\allowbreak{}e\allowbreak{}a\allowbreak{}9\allowbreak{}1\allowbreak{}6\allowbreak{}b\allowbreak{}e\allowbreak{}1\allowbreak{}4\allowbreak{}c\allowbreak{}8\allowbreak{}5\allowbreak{}b\allowbreak{}1\allowbreak{}0\allowbreak{}7\allowbreak{}6\allowbreak{}d\allowbreak{}4}, the expanded program \texttt{b\allowbreak{}5\allowbreak{}a\allowbreak{}e\allowbreak{}6\allowbreak{}a\allowbreak{}e\allowbreak{}0\allowbreak{}a\allowbreak{}7\allowbreak{}e\allowbreak{}9\allowbreak{}e\allowbreak{}b\allowbreak{}3\allowbreak{}4\allowbreak{}2\allowbreak{}4\allowbreak{}d\allowbreak{}1\allowbreak{}6\allowbreak{}2\allowbreak{}f\allowbreak{}c\allowbreak{}0\allowbreak{}f\allowbreak{}3\allowbreak{}e\allowbreak{}6\allowbreak{}0\allowbreak{}6\allowbreak{}2\allowbreak{}e\allowbreak{}7\allowbreak{}e\allowbreak{}e\allowbreak{}3\allowbreak{}a\allowbreak{}1\allowbreak{}3\allowbreak{}2\allowbreak{}9\allowbreak{}c\allowbreak{}3\allowbreak{}9\allowbreak{}0\allowbreak{}a\allowbreak{}9\allowbreak{}b\allowbreak{}c\allowbreak{}7\allowbreak{}4\allowbreak{}c\allowbreak{}1\allowbreak{}3\allowbreak{}8\allowbreak{}f\allowbreak{}6\allowbreak{}d\allowbreak{}2\allowbreak{}c\allowbreak{}a\allowbreak{}f\allowbreak{}b}, and the original ticket engine \texttt{a\allowbreak{}3\allowbreak{}0\allowbreak{}8\allowbreak{}a\allowbreak{}e\allowbreak{}9\allowbreak{}7\allowbreak{}e\allowbreak{}e\allowbreak{}8\allowbreak{}a\allowbreak{}5\allowbreak{}a\allowbreak{}0\allowbreak{}f\allowbreak{}7\allowbreak{}7\allowbreak{}c\allowbreak{}0\allowbreak{}6\allowbreak{}7\allowbreak{}b\allowbreak{}d\allowbreak{}6\allowbreak{}f\allowbreak{}2\allowbreak{}d\allowbreak{}4\allowbreak{}9\allowbreak{}4\allowbreak{}f\allowbreak{}7\allowbreak{}c\allowbreak{}c\allowbreak{}f\allowbreak{}4\allowbreak{}0\allowbreak{}a\allowbreak{}0\allowbreak{}7\allowbreak{}1\allowbreak{}b\allowbreak{}c\allowbreak{}8\allowbreak{}0\allowbreak{}2\allowbreak{}9\allowbreak{}f\allowbreak{}0\allowbreak{}8\allowbreak{}a\allowbreak{}6\allowbreak{}1\allowbreak{}6\allowbreak{}5\allowbreak{}0\allowbreak{}2\allowbreak{}4\allowbreak{}0\allowbreak{}6\allowbreak{}1\allowbreak{}6\allowbreak{}9}. The wrapper changes only the two witness allowances and the identity fields in the program, and runs the parameters in Q0 and Q.2; the original tickets keep their distinct fees. The analytic sources and written transports remain those bound by P, including its source map \texttt{f\allowbreak{}5\allowbreak{}8\allowbreak{}d\allowbreak{}5\allowbreak{}4\allowbreak{}f\allowbreak{}8\allowbreak{}0\allowbreak{}9\allowbreak{}3\allowbreak{}6\allowbreak{}5\allowbreak{}f\allowbreak{}f\allowbreak{}a\allowbreak{}5\allowbreak{}f\allowbreak{}9\allowbreak{}c\allowbreak{}e\allowbreak{}8\allowbreak{}2\allowbreak{}e\allowbreak{}a\allowbreak{}0\allowbreak{}3\allowbreak{}0\allowbreak{}a\allowbreak{}c\allowbreak{}d\allowbreak{}3\allowbreak{}a\allowbreak{}f\allowbreak{}8\allowbreak{}2\allowbreak{}1\allowbreak{}3\allowbreak{}5\allowbreak{}e\allowbreak{}2\allowbreak{}0\allowbreak{}8\allowbreak{}2\allowbreak{}4\allowbreak{}e\allowbreak{}b\allowbreak{}7\allowbreak{}8\allowbreak{}d\allowbreak{}2\allowbreak{}5\allowbreak{}9\allowbreak{}1\allowbreak{}8\allowbreak{}6\allowbreak{}b\allowbreak{}a\allowbreak{}f\allowbreak{}7\allowbreak{}a\allowbreak{}2\allowbreak{}7\allowbreak{}2}, notices and attributions. The finite certificate proves the repriced cover and rational inequalities; the analytic implication is Q.1--Q.4 with the retained written sources. The unrefined \(\fitinline{.126667206}\) receipt is unpaid, and the old ideal contact test is still separate. The \(\fitinline{.12668}\) and \(\fitinline{.12725}\) proposals are not consequences of Q.

\end{document}
